\IfFileExists{stacks-project-book.cls}{%
\documentclass{stacks-project-book}
}{%
\documentclass{amsbook}
}

\usepackage{verbatim}
\newenvironment{reference}{\comment}{\endcomment}
\newenvironment{slogan}{\comment}{\endcomment}
\newenvironment{history}{\comment}{\endcomment}

\usepackage[all]{xy}

\xyoption{2cell}
\UseAllTwocells

\usepackage{multicol}

\usepackage{lmodern}
\usepackage[T1]{fontenc}
\usepackage[french]{babel}
\AtBeginDocument{\shorthandoff{?}}


\usepackage{hyperref}


\theoremstyle{plain}
\newtheorem{theorem}[subsection]{Théorème}
\newtheorem{proposition}[subsection]{Proposition}
\newtheorem{lemma}[subsection]{Lemme}

\theoremstyle{definition}
\newtheorem{definition}[subsection]{Définition}
\newtheorem{example}[subsection]{Exemple}
\newtheorem{exercise}[subsection]{Exercice}
\newtheorem{situation}[subsection]{Situation}

\theoremstyle{remark}
\newtheorem{remark}[subsection]{Remarque}
\newtheorem{remarks}[subsection]{Remarques}

\renewcommand{\proofname}{Démonstration}

\numberwithin{equation}{subsection}

\def\lim{\mathop{\mathrm{lim}}\nolimits}
\def\colim{\mathop{\mathrm{colim}}\nolimits}
\def\Spec{\mathop{\mathrm{Spec}}}
\def\Hom{\mathop{\mathrm{Hom}}\nolimits}
\def\Ext{\mathop{\mathrm{Ext}}\nolimits}
\def\SheafHom{\mathop{\mathcal{H}\!\mathit{om}}\nolimits}
\def\SheafExt{\mathop{\mathcal{E}\!\mathit{xt}}\nolimits}
\def\Sch{\mathit{Sch}}
\def\Mor{\mathop{\mathrm{Mor}}\nolimits}
\def\Ob{\mathop{\mathrm{Ob}}\nolimits}
\def\Sh{\mathop{\mathit{Sh}}\nolimits}
\def\NL{\mathop{N\!L}\nolimits}
\def\CH{\mathop{\mathrm{CH}}\nolimits}
\def\proetale{{pro\text{-}\acute{e}tale}}
\def\etale{{\acute{e}tale}}
\def\QCoh{\mathit{QCoh}}
\def\Ker{\mathop{\mathrm{Ker}}}
\def\Im{\mathop{\mathrm{Im}}}
\def\Coker{\mathop{\mathrm{Coker}}}
\def\Coim{\mathop{\mathrm{Coim}}}

\DeclareMathSymbol{\boxtimes}{\mathbin}{AMSa}{"02}

\def\QCohstack{\mathcal{QC}\!\mathit{oh}}
\def\Cohstack{\mathcal{C}\!\mathit{oh}}
\def\Spacesstack{\mathcal{S}\!\mathit{paces}}
\def\Quotfunctor{\mathrm{Quot}}
\def\Hilbfunctor{\mathrm{Hilb}}
\def\Curvesstack{\mathcal{C}\!\mathit{urves}}
\def\Polarizedstack{\mathcal{P}\!\mathit{olarized}}
\def\Complexesstack{\mathcal{C}\!\mathit{omplexes}}
\def\Pic{\mathop{\mathrm{Pic}}\nolimits}
\def\Picardstack{\mathcal{P}\!\mathit{ic}}
\def\Picardfunctor{\mathrm{Pic}}
\def\Deformationcategory{\mathcal{D}\!\mathit{ef}}
\begin{document}
\title{Le projet Stacks --- édition française, partie 9}
\maketitle
\tableofcontents

% OK, start here.
%

\chapter{Cohomologie des espaces algébriques}


\phantomsection
\label{spaces-cohomology-section-phantom}





\section{Introduction}
\label{spaces-cohomology-section-introduction}

\noindent
Dans ce chapitre, nous traitons de la cohomologie des espaces algébriques.
Bien que nous démontrions certains résultats sur la cohomologie des faisceaux abéliens,
nous nous intéressons principalement à la cohomologie des faisceaux quasi-cohérents, c'est-à-dire
que nous démontrons les analogues des résultats du chapitre ``Cohomologie des schémas''.
Certains résultats de ce chapitre se trouvent dans \cite{Kn}.

\medskip\noindent
Un ingrédient important qui manque dans ce chapitre est le
{\it principe d'induction}, c'est-à-dire l'analogue, pour les espaces algébriques quasi-compacts
et quasi-séparés, du
chapitre Cohomologie des schémas, lemme \ref{coherent-lemma-induction-principle}.
Ce principe est formulé précisément et démontré en détail dans
Catégories dérivées des espaces, section \ref{spaces-perfect-section-induction}.
Au lieu du principe d'induction, nous utilisons dans ce chapitre le
complexe de {\v C}ech alterné, voir
section \ref{spaces-cohomology-section-alternating-cech}.
Il est conçu pour démontrer des énoncés d'annulation tels que
la proposition \ref{spaces-cohomology-proposition-vanishing},
mais, dans certains cas, le principe d'induction est un outil plus puissant
et peut-être plus « standard ». Nous encourageons le lecteur
à consulter le principe d'induction
après avoir lu une partie du contenu de cette section.



\section{Conventions}
\label{spaces-cohomology-section-conventions}

\noindent
L'hypothèse générale est que tous les schémas sont contenus dans
un grand site fppf $\Sch_{fppf}$. Tous les anneaux $A$ considérés
ont la propriété que $\Spec(A)$ est (isomorphe à) un
objet de ce grand site.

\medskip\noindent
Soit $S$ un schéma et soit $X$ un espace algébrique sur $S$.
Dans ce chapitre et le suivant, nous écrirons $X \times_S X$
pour le produit de $X$ par lui-même (dans la catégorie des espaces algébriques
sur $S$), plutôt que $X \times X$.








\section{Images directes supérieures}
\label{spaces-cohomology-section-higher-direct-image}

\noindent
Soit $S$ un schéma. Soit $X$ un espace algébrique représentable sur $S$.
Soit $\mathcal{F}$ un module quasi-cohérent sur $X$ (voir
Propriétés des espaces, section \ref{spaces-properties-section-quasi-coherent}).
D'après
Descente, proposition \ref{descent-proposition-same-cohomology-quasi-coherent},
les groupes de cohomologie $H^i(X, \mathcal{F})$ coïncident avec le groupe de
cohomologie usuel calculé dans la topologie de Zariski du
module quasi-cohérent sur le schéma représentant $X$.

\medskip\noindent
Plus généralement, soit $f : X \to Y$ un morphisme quasi-compact et quasi-séparé
d'espaces algébriques représentables $X$ et $Y$. Soit
$\mathcal{F}$ un module quasi-cohérent sur $X$. D'après
Descente, lemme \ref{descent-lemma-higher-direct-images-small-etale},
le faisceau $R^if_*\mathcal{F}$ coïncide avec l'image
directe supérieure usuelle calculée pour la topologie de Zariski
du module quasi-cohérent sur le schéma représentant $X$
qui se déduit vers le schéma représentant $Y$.

\medskip\noindent
Plus généralement encore, supposons que $f : X \to Y$ soit un
morphisme représentable, quasi-compact et
quasi-séparé d'espaces algébriques sur $S$. Soit $V$ un schéma
et soit $V \to Y$ un morphisme étale surjectif. Posons $U = V \times_Y X$
et soit $f' : U \to V$ le changement de base de $f$. Pour tout
Pour tout $\mathcal{O}_X$-module quasi-cohérent $\mathcal{F}$, nous avons
\begin{equation}
\label{spaces-cohomology-equation-representable-higher-direct-image}
R^if'_*(\mathcal{F}|_U) = (R^if_*\mathcal{F})|_V,
\end{equation}
voir
Propriétés des espaces,
lemme \ref{spaces-properties-lemma-pushforward-etale-base-change-modules}.
Comme $f' : U \to V$ est un morphisme quasi-compact et quasi-séparé
de schémas, la remarque du paragraphe précédent permet de
calculer $R^if'_*(\mathcal{F}|_U)$ en considérant $\mathcal{F}|_U$ comme un
faisceau quasi-cohérent sur le schéma $U$, et $f'$ comme un morphisme de schémas.
Nous utiliserons souvent ce fait sans autre mention.

\medskip\noindent
Nous démontrons ensuite que les images directes supérieures de faisceaux quasi-cohérents sont
quasi-cohérentes pour tout morphisme quasi-compact et quasi-séparé d'
espaces algébriques. Dans la démonstration, nous utilisons une astuce ; une démonstration « meilleure » utiliserait
un complexe de {\v C}ech relatif, comme cela est expliqué dans
Faisceaux sur les champs, sections \ref{stacks-sheaves-section-cech} et
\ref{stacks-sheaves-section-sheaf-cech-complex} ff.

\begin{lemma}
\label{spaces-cohomology-lemma-higher-direct-image}
Soit $S$ un schéma. Soit $f : X \to Y$ un morphisme d'espaces algébriques
sur $S$. Si $f$ est quasi-compact et quasi-séparé, alors $R^if_*$
transforme les $\mathcal{O}_X$-modules quasi-cohérents en
$\mathcal{O}_Y$-modules quasi-cohérents.
\end{lemma}

\begin{proof}
Soit $V \to Y$ un morphisme étale où $V$ est un schéma affine. Posons
$U = V \times_Y X$ et notons $f' : U \to V$ le morphisme induit.
Soit $\mathcal{F}$ un $\mathcal{O}_X$-module quasi-cohérent. D'après
Propriétés des espaces, lemme
\ref{spaces-properties-lemma-pushforward-etale-base-change-modules}
nous avons
$R^if'_*(\mathcal{F}|_U) = (R^if_*\mathcal{F})|_V$.
Puisque la propriété d'être un module quasi-cohérent est locale pour la
topologie étale sur $Y$ (voir
Propriétés des espaces, lemme
\ref{spaces-properties-lemma-characterize-quasi-coherent})
nous pouvons remplacer $Y$ par $V$, c'est-à-dire supposer que $Y$ est un schéma affine.

\medskip\noindent
Supposons $Y$ affine. Comme $f$ est quasi-compact, $X$ est
quasi-compact. Nous pouvons donc choisir un schéma affine $U$ et un
morphisme étale surjectif $g : U \to X$, voir
Propriétés des espaces,
lemme \ref{spaces-properties-lemma-quasi-compact-affine-cover}.
Figure
$$
\xymatrix{
U \ar[r]_g \ar[rd]_{f \circ g} & X \ar[d]^f \\
& Y
}
$$
Le morphisme $g : U \to X$ est représentable, séparé
et quasi-compact puisque $X$ est quasi-séparé. Le lemme
s'applique donc à $g$ (par la discussion précédente).
Il s'applique également à $f \circ g : U \to Y$ (car il s'agit d'un morphisme
de schémas affines).

\medskip\noindent
Dans la situation décrite au paragraphe précédent, nous montrerons par
récurrence sur $n$ que $IH_n$ : pour tout faisceau quasi-cohérent $\mathcal{F}$
sur $X$, les faisceaux $R^if\mathcal{F}$
sont quasi-cohérents pour $i \leq n$.
Le cas $n = 0$ résulte de
Morphismes des espaces, lemme \ref{spaces-morphisms-lemma-pushforward}.
Supposons $IH_n$. Dans le reste de la démonstration, nous montrons que $IH_{n + 1}$ est valable.

\medskip\noindent
Soit $\mathcal{H}$ un $\mathcal{O}_U$-module quasi-cohérent.
Considérons la suite spectrale de Leray
$$
E_2^{p, q} = R^pf_* R^qg_* \mathcal{H}
\Rightarrow
R^{p + q}(f \circ g)_*\mathcal{H}
$$
Cohomologie sur les sites, lemme \ref{sites-cohomology-lemma-relative-Leray}.
Comme $R^qg_*\mathcal{H}$ est quasi-cohérent par $IH_n$, les faisceaux
$R^pf_*R^qg_*\mathcal{H}$ sont quasi-cohérents pour $p \leq n$.
Les faisceaux $R^{p + q}(f \circ g)_*\mathcal{H}$ sont tous
quasi-cohérents (en fait nuls pour $p + q > 0$, mais cela ne nous est pas nécessaire).
En degrés $\leq n + 1$, le seul module dont nous n'avons pas encore
établi la quasi-cohérence est $E_2^{n + 1, 0} = R^{n + 1}f_*g_*\mathcal{H}$.
De plus, les différentielles
$d_r^{n + 1, 0} : E_r^{n + 1, 0} \to E_r^{n + 1 + r, 1 - r}$
sont nulles puisque le terme cible est nul. Comme $\QCoh(\mathcal{O}_X)$
est une sous-catégorie de Serre faible de $\textit{Mod}(\mathcal{O}_X)$
(Propriétés des espaces, lemme
\ref{spaces-properties-lemma-properties-quasi-coherent})
il s'ensuit que $R^{n + 1}f_*g_*\mathcal{H}$
est quasi-cohérent (les détails sont omis).

\medskip\noindent
Soit $\mathcal{F}$ un $\mathcal{O}_X$-module quasi-cohérent.
Posons $\mathcal{H} = g^*\mathcal{F}$. Le morphisme d'adjonction
$\mathcal{F} \to g_*g^*\mathcal{F} = g_*\mathcal{H}$ est injectif
car $U \to X$ est étale surjectif. Considérons la suite exacte
$$
0 \to \mathcal{F} \to g_*\mathcal{H} \to \mathcal{G} \to 0
$$
où $\mathcal{G}$ est le conoyau de la première flèche et, en particulier,
quasi-cohérent. En appliquant la suite exacte longue de cohomologie, on obtient
$$
R^nf_*g_*\mathcal{H} \to
R^nf_*\mathcal{G} \to
R^{n + 1}f_*\mathcal{F} \to
R^{n + 1}f_*g_*\mathcal{H} \to
R^{n + 1}f_*\mathcal{G}
$$
Le conoyau de la première flèche est quasi-cohérent et
nous avons vu ci-dessus que $R^{n + 1}f_*g_*\mathcal{H}$ est quasi-cohérent.
Ainsi, $R^{n + 1}f_*\mathcal{F}$ possède une filtration à $2$ étapes,
dont le premier quotient est quasi-cohérent et le second un sous-module d'un
faisceau quasi-cohérent. Comme $\mathcal{F}$ est un module quasi-cohérent arbitraire
$\mathcal{O}_X$-module, le même résultat vaut pour $\mathcal{G}$.
Nous pouvons donc choisir une suite exacte
$0 \to \mathcal{A} \to R^{n + 1}f_*\mathcal{G} \to \mathcal{B}$
où $\mathcal{A}$ et $\mathcal{B}$ sont des $\mathcal{O}_Y$-modules quasi-cohérents.
Alors le noyau $\mathcal{K}$ de
$R^{n + 1}f_*g_*\mathcal{H} \to R^{n + 1}f_*\mathcal{G}
\to \mathcal{B}$ est quasi-cohérent ; on obtient donc un morphisme
$\mathcal{K} \to \mathcal{A}$ dont le noyau $\mathcal{K}'$ est
également quasi-cohérent. Ainsi $R^{n + 1}f_*\mathcal{F}$ s'insère dans une suite
exacte
$$
R^nf_*g_*\mathcal{H} \to
R^nf_*\mathcal{G} \to
R^{n + 1}f_*\mathcal{F} \to \mathcal{K}' \to 0
$$
où tous les modules sont quasi-cohérents, sauf éventuellement $R^{n + 1}f_*\mathcal{F}$.
Nous concluons que $R^{n + 1}f_*\mathcal{F}$ est quasi-cohérent, c'est-à-dire que
$IH_{n + 1}$ est vérifiée, comme souhaité.
\end{proof}

\begin{lemma}
\label{spaces-cohomology-lemma-quasi-coherence-higher-direct-images-application}
Soit $S$ un schéma. Soit $f : X \to Y$ un morphisme quasi-séparé et quasi-compact
d'espaces algébriques sur $S$. Pour tout
$\mathcal{O}_X$-module quasi-cohérent $\mathcal{F}$ et tout objet affine $V$ de
$Y_\etale$, on a
$$
H^q(V \times_Y X, \mathcal{F}) = H^0(V, R^qf_*\mathcal{F})
$$
pour tout $q \in \mathbf{Z}$.
\end{lemma}

\begin{proof}
Comme la formation de $Rf_*$ commute à la localisation étale,
(Propriétés des espaces, lemme
\ref{spaces-properties-lemma-pushforward-etale-base-change-modules})
nous pouvons remplacer $Y$ par $V$ et supposer que $Y = V$ est affine.
Considérons la suite spectrale de Leray
$E_2^{p, q} = H^p(Y, R^qf_*\mathcal{F})$
convergeant vers $H^{p + q}(X, \mathcal{F})$, voir
Cohomologie sur les sites, lemme \ref{sites-cohomology-lemma-Leray}.
Par le lemme \ref{spaces-cohomology-lemma-higher-direct-image},
les faisceaux $R^qf_*\mathcal{F}$ sont quasi-cohérents. D'après
Cohomologie des schémas, lemme
\ref{coherent-lemma-quasi-coherent-affine-cohomology-zero}
on a $E_2^{p, q} = 0$ lorsque $p > 0$.
La suite spectrale dégénère donc en $E_2$, ce qui conclut.
\end{proof}




\section{Morphismes finis}
\label{spaces-cohomology-section-finite-morphisms}

\noindent
Voici quelques résultats valables pour tous les faisceaux abéliens
(en particulier pour les modules quasi-cohérents).
Nous {\bf avertissons} le lecteur que ces lemmes ne sont pas valables pour
les morphismes finis de schémas munis de la topologie de Zariski.

\begin{lemma}
\label{spaces-cohomology-lemma-finite-higher-direct-image-zero}
Soit $S$ un schéma. Soit $f : X \to Y$ un morphisme intégral (par exemple fini)
d'espaces algébriques. Alors
$f_* : \textit{Ab}(X_\etale) \to \textit{Ab}(Y_\etale)$
est un foncteur exact et $R^pf_* = 0$ pour $p > 0$.
\end{lemma}

\begin{proof}
D'après Propriétés des espaces, lemme
\ref{spaces-properties-lemma-pushforward-etale-base-change}
nous pouvons calculer les images directes supérieures sur un recouvrement étale de $Y$.
Nous pouvons donc supposer que $Y$ est un schéma.
Cela implique que $X$ est un schéma (Morphismes des espaces, lemme
\ref{spaces-morphisms-lemma-integral-local}).
Dans ce cas, nous pouvons appliquer
Cohomologie étale, lemme \ref{etale-cohomology-lemma-what-integral}.
Dans le cas fini, le lecteur pourra consulter la
proposition moins technique de Cohomologie étale,
\ref{etale-cohomology-proposition-finite-higher-direct-image-zero}.
\end{proof}

\begin{lemma}
\label{spaces-cohomology-lemma-stalk-push-finite}
Soit $S$ un schéma. Soit $f : X \to Y$ un morphisme fini d'espaces algébriques
sur $S$. Soit $\overline{y}$ un point géométrique de $Y$ dont les
relèvements notés $\overline{x}_1, \ldots, \overline{x}_n$ dans $X$. Alors
$$
(f_*\mathcal{F})_{\overline{y}} =
\prod\nolimits_{i = 1, \ldots, n}
\mathcal{F}_{\overline{x}_i}
$$
pour tout faisceau $\mathcal{F}$ sur $X_\etale$.
\end{lemma}

\begin{proof}
Choisissons un voisinage étale $(V, \overline{v})$ de $\overline{y}$.
Alors la fibre $(f_*\mathcal{F})_{\overline{y}}$
est la fibre de $f_*\mathcal{F}|_V$ en $\overline{v}$.
D'après Propriétés des espaces,
lemme \ref{spaces-properties-lemma-pushforward-etale-base-change}
nous pouvons remplacer $Y$ par $V$ et $X$ par $X \times_Y V$.
Alors $Z \to X$ est un morphisme fini de schémas et le résultat est
la proposition de Cohomologie étale
\ref{etale-cohomology-proposition-finite-higher-direct-image-zero}.
\end{proof}

\begin{lemma}
\label{spaces-cohomology-lemma-finite-rings}
Soit $S$ un schéma. Soit $\pi : X \to Y$ un morphisme fini d'espaces
algébriques sur $S$. Soit $\mathcal{A}$ un faisceau d'anneaux sur $X_\etale$.
Soit $\mathcal{B}$ un faisceau d'anneaux sur $Y_\etale$.
Soit $\varphi : \mathcal{B} \to \pi_*\mathcal{A}$
un homomorphisme de faisceaux d'anneaux, de sorte que l'on obtient un
morphisme de topos annelés
$$
f = (\pi, \varphi) :
(\Sh(X_\etale), \mathcal{A})
\longrightarrow
(\Sh(Y_\etale), \mathcal{B}).
$$
Pour un faisceau de $\mathcal{A}$-modules $\mathcal{F}$ et un
faisceau de $\mathcal{B}$-modules $\mathcal{G}$, l'application canonique
$$
\mathcal{G} \otimes_\mathcal{B} f_*\mathcal{F}
\longrightarrow
f_*(f^*\mathcal{G} \otimes_\mathcal{A} \mathcal{F}).
$$
est un isomorphisme.
\end{lemma}

\begin{proof}
L'application est adjointe à l'application
$$
f^*\mathcal{G} \otimes_\mathcal{A}
f^* f_*\mathcal{F} =
f^*(\mathcal{G} \otimes_\mathcal{B} f_*\mathcal{F})
\longrightarrow
f^*\mathcal{G} \otimes_\mathcal{A} \mathcal{F}
$$
provenant de $\text{id} : f^*\mathcal{G} \to f^*\mathcal{G}$
et du morphisme d'adjonction $f^* f_*\mathcal{F} \to \mathcal{F}$.
Pour voir que cette application est un isomorphisme, nous pouvons vérifier sur les fibres
(Propriétés des espaces, théorème
\ref{spaces-properties-theorem-exactness-stalks}).
Soit $\overline{y}$ un point géométrique de $Y$ et
soient $\overline{x}_1, \ldots, \overline{x}_n$ les points géométriques
de $X$ situés au-dessus de $\overline{y}$.
En calculant l'effet de nos applications sur les fibres, nous voyons qu'il faut
démontrer
$$
\mathcal{G}_{\overline{y}}
\otimes_{\mathcal{B}_{\overline{y}}}
\left(
\bigoplus\nolimits_{i = 1, \ldots, n} \mathcal{F}_{\overline{x}_i}
\right) =
\bigoplus\nolimits_{i = 1, \ldots, n}
(\mathcal{G}_{\overline{y}}
\otimes_{\mathcal{B}_{\overline{x}}}
\mathcal{A}_{\overline{x}_i}) \otimes_{\mathcal{A}_{\overline{x}_i}}
\mathcal{F}_{\overline{x}_i}
$$
ce qui est vrai. Nous avons utilisé ici que
le produit tensoriel commute au passage aux fibres, ainsi que le
comportement des fibres par image inverse,
Propriétés des espaces, lemme \ref{spaces-properties-lemma-stalk-pullback}, et
le comportement des fibres sous poussée le long d'une immersion fermée, voir
lemme \ref{spaces-cohomology-lemma-stalk-push-finite}.
\end{proof}

\noindent
Nous terminons cette section par une formule de projection très générale
pour les morphismes finis.

\begin{lemma}
\label{spaces-cohomology-lemma-projection-formula-finite}
Avec $S$, $X$, $Y$, $\pi$, $\mathcal{A}$, $\mathcal{B}$, $\varphi$ et $f$
comme dans le lemme \ref{spaces-cohomology-lemma-finite-rings}, on a
$$
K \otimes_\mathcal{B}^\mathbf{L} Rf_*M =
Rf_*(Lf^*K \otimes_\mathcal{A}^\mathbf{L} M)
$$
dans $D(\mathcal{B})$ pour tout $K \in D(\mathcal{B})$ et
$M \in D(\mathcal{A})$.
\end{lemma}

\begin{proof}
Comme $f_*$ est exact (lemme \ref{spaces-cohomology-lemma-finite-higher-direct-image-zero}),
le foncteur $Rf_*$ se calcule en appliquant $f_*$ à un complexe représentant quelconque.
Choisissons un complexe $\mathcal{K}^\bullet$ de $\mathcal{B}$-modules
représentant $K$, K-plat à termes plats, voir
Cohomologie sur les sites, lemme \ref{sites-cohomology-lemma-K-flat-resolution}.
Alors $f^*\mathcal{K}^\bullet$ est K-plat à termes plats, voir
Cohomologie sur les sites, lemme \ref{sites-cohomology-lemma-pullback-K-flat}.
Choisissons un complexe quelconque $\mathcal{M}^\bullet$ de $\mathcal{A}$-modules
représentant $M$. Alors
il faut montrer
$$
\text{Tot}(\mathcal{K}^\bullet \otimes_\mathcal{B} f_*\mathcal{M}^\bullet)
=
f_*\text{Tot}(f^*\mathcal{K}^\bullet \otimes_\mathcal{A} \mathcal{M}^\bullet)
$$
car, par nos choix, ces complexes représentent respectivement les membres de droite et de gauche
de la formule du lemme.
Comme $f_*$ commute aux sommes directes
(par exemple d'après la description des fibres dans
le lemme \ref{spaces-cohomology-lemma-stalk-push-finite}),
il suffit de vérifier les égalités
$$
\mathcal{K}^n \otimes_\mathcal{B} f_*\mathcal{M}^m
=
f_*(f^*\mathcal{K}^n \otimes_\mathcal{A} \mathcal{M}^m)
$$
qui résultent du lemme \ref{spaces-cohomology-lemma-finite-rings}.
\end{proof}









\section{Colimites et cohomologie}
\label{spaces-cohomology-section-colimits}

\noindent
Le lemme suivant s'applique en particulier aux diagrammes de
faisceaux quasi-cohérents.

\begin{lemma}
\label{spaces-cohomology-lemma-colimits}
Soit $S$ un schéma. Soit $X$ un espace algébrique sur $S$.
Si $X$ est quasi-compact et quasi-séparé, alors
$$
\colim_i H^p(X, \mathcal{F}_i)
\longrightarrow
H^p(X, \colim_i \mathcal{F}_i)
$$
est un isomorphisme
pour tout diagramme filtrant de faisceaux abéliens sur $X_\etale$.
\end{lemma}

\begin{proof}
C'est une conséquence du lemme
Cohomologie sur les sites,
\ref{sites-cohomology-lemma-colim-works-over-collection}.
À savoir, soit $\mathcal{B} \subset \Ob(X_{spaces, \etale})$
l'ensemble des espaces quasi-compacts et quasi-séparés étales au-dessus de $X$.
Si $U \in \mathcal{B}$, comme $U$ est quasi-compact,
la famille des recouvrements finis $\{U_i \to U\}$ avec $U_i \in \mathcal{B}$
est cofinale dans l'ensemble des recouvrements de $U$ dans $X_{spaces, \etale}$. Par
le lemme de Morphismes des espaces
\ref{spaces-morphisms-lemma-quasi-compact-quasi-separated-permanence}
l'ensemble $\mathcal{B}$ vérifie toutes les hypothèses du
lemme de Cohomologie sur les sites
\ref{sites-cohomology-lemma-colim-works-over-collection}.
Comme $X \in \mathcal{B}$, le résultat s'ensuit.
\end{proof}

\begin{lemma}
\label{spaces-cohomology-lemma-colimit-cohomology}
\begin{slogan}
Les images directes supérieures des morphismes qcqs commutent aux colimites filtrantes
de faisceaux.
\end{slogan}
Soit $S$ un schéma. Soit $f : X \to Y$ un morphisme quasi-compact et quasi-séparé
d'espaces algébriques sur $S$. Soit $\mathcal{F} = \colim \mathcal{F}_i$
une colimite filtrante de faisceaux abéliens sur $X_\etale$.
Alors, pour tout $p \geq 0$, on a
$$
R^pf_*\mathcal{F} = \colim R^pf_*\mathcal{F}_i.
$$
\end{lemma}

\begin{proof}
Nous utiliserons le fait que le morphisme de topos $f_{small} : X_{small} \to Y_{small}$
provient du morphisme de sites
$f_{spaces, \etale} : X_{spaces, \etale} \to Y_{spaces, \etale}$
associé au foncteur continu $V \longmapsto X \times_Y V$,
voir Propriétés des espaces, lemme
\ref{spaces-properties-lemma-functoriality-etale-site}.
Nous appliquerons le lemme
\ref{sites-cohomology-lemma-higher-direct-image-colimit}
à ce morphisme de sites. Comme tout objet de $Y_{spaces, \etale}$
admet un recouvrement par des objets affines, il suffit de montrer que
pour $V$ affine et étale au-dessus de $Y$, on a
$H^p(X \times_Y V, \mathcal{F}) = \colim H^p(X \times_Y V, \mathcal{F}_i)$.
Comme $V$ est affine, l'espace algébrique $X \times_Y V$ est
quasi-compact et quasi-séparé.
Nous pouvons donc appliquer le lemme \ref{spaces-cohomology-lemma-colimits}
pour conclure.
\end{proof}

\noindent
Le lemme suivant montre que les modules de présentation finie se comportent
comme attendu sur les espaces algébriques quasi-compacts et quasi-séparés.

\begin{lemma}
\label{spaces-cohomology-lemma-finite-presentation-quasi-compact-colimit}
Soit $S$ un schéma. Soit $X$ un espace algébrique quasi-compact et quasi-séparé
sur $S$. Soit $I$ un ensemble dirigé et
soit $(\mathcal{F}_i, \varphi_{ii'})$ un système indexé par $I$
de $\mathcal{O}_X$-modules. Soit $\mathcal{G}$ un
$\mathcal{O}_X$-module de présentation finie. Alors
$$
\colim_i \Hom_X(\mathcal{G}, \mathcal{F}_i)
=
\Hom_X(\mathcal{G}, \colim_i \mathcal{F}_i).
$$
En particulier, $\Hom_X(\mathcal{G}, -)$ commute aux colimites filtrantes
dans $\QCoh(\mathcal{O}_X)$.
\end{lemma}

\begin{proof}
L'égalité affichée est un cas particulier du lemme de
Modules sur les sites, lemme \ref{sites-modules-lemma-finite-presentation-quasi-compact-colimit}.
Pour l'appliquer, il faut vérifier les hypothèses du
lemme de Sites \ref{sites-lemma-directed-colimits-global-sections}, partie (4),
pour le site $X_\etale$.
Pour cela, nous vérifierons les hypothèses
(2)(a), (2)(b) et (2)(c) de
la remarque de Sites \ref{sites-remark-stronger-conditions}.
À savoir, soit $\mathcal{B} \subset \Ob(X_\etale)$ l'ensemble des objets affines.
Alors
\begin{enumerate}
\item Puisque $X$ est quasi-compact, il existe un $U \in \mathcal{B}$
tel que $U \to X$ soit surjectif
(Propriétés des espaces, lemme
\ref{spaces-properties-lemma-quasi-compact-affine-cover}),
d'où $h_U^\# \to *$ est surjectif.
\item Pour $U \in \mathcal{B}$, tout recouvrement étale $\{U_i \to U\}_{i \in I}$
de $U$ admet un raffinement par un recouvrement étale fini
$\{U_j \to U\}_{j = 1, \ldots, m}$ avec $U_j \in \mathcal{B}$
(Topologies, lemme \ref{topologies-lemma-etale-affine}).
\item Pour $U, U' \in \Ob(X_\etale)$, on a
$h_U^\# \times h_{U'}^\# = h_{U \times_X U'}^\#$.
Si $U, U' \in \mathcal{B}$, alors $U \times_X U'$ est quasi-compact
car $X$ est quasi-séparé ; voir par exemple le lemme de Morphismes des espaces
\ref{spaces-morphisms-lemma-quasi-compact-quasi-separated-permanence}
Par conséquent, il existe un morphisme étale surjectif
$U'' \to U \times_X U'$ avec $U'' \in \mathcal{B}$
Propriétés des espaces, lemme
\ref{spaces-properties-lemma-quasi-compact-affine-cover}).
Autrement dit, il existe des morphismes $U'' \to U$ et $U'' \to U'$
tels que l'application $h_{U''}^\# \to h_U^\# \times h_{u'}^\#$ soit
surjective.
\end{enumerate}
Pour l'énoncé final, observons que le foncteur d'inclusion
$\QCoh(\mathcal{O}_X) \to \textit{Mod}(\mathcal{O}_X)$
commute aux colimites et que les modules de présentation finie
sont quasi-cohérents. Voir
Propriétés des Espaces, lemme
\ref{spaces-properties-lemma-properties-quasi-coherent}.
\end{proof}





\section{Le complexe de {\v C}ech alterné}
\label{spaces-cohomology-section-alternating-cech}

\noindent
Soit $S$ un schéma. Soit $f : U \to X$ un morphisme étale d'espaces
algébriques sur $S$. Le foncteur
$$
j : U_{spaces, \etale} \longrightarrow X_{spaces, \etale},\quad
V/U \longmapsto V/X
$$
induit une équivalence entre $U_{spaces, \etale}$ et la catégorie localisée
$X_{spaces, \etale}/U$, voir
Propriétés des espaces, section \ref{spaces-properties-section-localize}.
Il existe donc des foncteurs
$$
f_! : \textit{Ab}(U_\etale) \longrightarrow
\textit{Ab}(X_\etale),\quad
f_! : \textit{Mod}(\mathcal{O}_U) \longrightarrow \textit{Mod}(\mathcal{O}_X),
$$
qui sont adjoints à gauche des foncteurs
$$
f^{-1} : \textit{Ab}(X_\etale) \longrightarrow
\textit{Ab}(U_\etale),\quad
f^* : \textit{Mod}(\mathcal{O}_X) \longrightarrow \textit{Mod}(\mathcal{O}_U)
$$
voir
Modules sur les sites, section \ref{sites-modules-section-localize}.
Avertissement : a priori, ce foncteur n'a
rien à voir avec la cohomologie à supports compacts !
Dans la référence ci-dessus, nous avons appelé ce foncteur « extension par zéro ».
Notons que les deux versions de $f_!$ coïncident puisque $f^* = f^{-1}$ pour
les faisceaux de $\mathcal{O}_X$-modules.

\medskip\noindent
Comme nous utiliserons cette construction ci-dessous, rappelons quelques-unes de ses
propriétés. Soit $\mathcal{G}$ un faisceau abélien sur $U_\etale$ ;
le faisceau $f_!$ est le faisceautisé du préfaisceau
$$
V/X \longmapsto
f_!\mathcal{G}(V) =
\bigoplus\nolimits_{\varphi \in \Mor_X(V, U)}
\mathcal{G}(V \xrightarrow{\varphi} U),
$$
voir
Modules sur les sites, lemme \ref{sites-modules-lemma-extension-by-zero}.
De plus, si $\mathcal{G}$ est un $\mathcal{O}_U$-module, alors $f_!\mathcal{G}$
est le faisceautisé du même préfaisceau de groupes abéliens,
muni de façon évidente d'une structure de $\mathcal{O}_X$-module
(voir loc. cit.). Soit $\overline{x} : \Spec(k) \to X$
un point géométrique. Il existe alors une identification canonique
$$
(f_!\mathcal{G})_{\overline{x}} =
\bigoplus\nolimits_{\overline{u}} \mathcal{G}_{\overline{u}}
$$
où la somme porte sur tous les $\overline{u} : \Spec(k) \to U$ tels que
$f \circ \overline{u} = \overline{x}$, voir
Modules sur les sites, lemme \ref{sites-modules-lemma-stalk-j-shriek}
et
Propriétés des espaces, lemme
\ref{spaces-properties-lemma-points-small-etale-site}.
Dans la suite, nous étudions le faisceau $f_!\underline{\mathbf{Z}}$.
Ici, $\underline{\mathbf{Z}}$ désigne le faisceau constant sur
$X_\etale$ ou $U_\etale$.

\begin{lemma}
\label{spaces-cohomology-lemma-product-is-tensor-product}
Soit $S$ un schéma. Soient $f_i : U_i \to X$ des morphismes étales
d'espaces algébriques sur $S$. Alors il existe des isomorphismes
$$
f_{1, !}\underline{\mathbf{Z}} \otimes_{\mathbf{Z}}
f_{2, !}\underline{\mathbf{Z}}
\longrightarrow
f_{12, !}\underline{\mathbf{Z}}
$$
où $f_{12} : U_1 \times_X U_2 \to X$ est le morphisme structural
et
$$
(f_1 \amalg f_2)_! \underline{\mathbf{Z}}
\longrightarrow
f_{1, !}\underline{\mathbf{Z}} \oplus
f_{2, !}\underline{\mathbf{Z}}
$$
\end{lemma}

\begin{proof}
Une fois l'application définie, ce sera un isomorphisme d'après la description
des fibres ci-dessus. Pour définir l'application, il suffit de travailler au niveau
des préfaisceaux. Nous devons donc définir une application
$$
\left(\bigoplus\nolimits_{\varphi_1 \in \Mor_X(V, U_1)} \mathbf{Z}\right)
\otimes_{\mathbf{Z}}
\left(\bigoplus\nolimits_{\varphi_2 \in \Mor_X(V, U_2)} \mathbf{Z}\right)
\longrightarrow
\bigoplus\nolimits_{\varphi \in \Mor_X(V, U_1 \times_X U_2)}
\mathbf{Z}
$$
Nous envoyons l'élément $1_{\varphi_1} \otimes 1_{\varphi_2}$ sur l'élément
$1_{\varphi_1 \times \varphi_2}$, avec les notations évidentes. Nous omettons la preuve
de la seconde égalité.
\end{proof}

\noindent
Une autre propriété importante est l'application trace
$$
\text{Tr}_f : f_!\underline{\mathbf{Z}} \longrightarrow \underline{\mathbf{Z}}.
$$
L'application trace est adjointe à
l'application $\mathbf{Z} \to f^{-1}\underline{\mathbf{Z}}$ (qui est un isomorphisme).
Si $\overline{x}$ est comme ci-dessus, l'application $\text{Tr}_f$ sur la fibre en $\overline{x}$
est l'application
$$
(\text{Tr}_f)_{\overline{x}} :
(f_!\underline{\mathbf{Z}})_{\overline{x}} =
\bigoplus\nolimits_{\overline{u}} \mathbf{Z}
\longrightarrow
\mathbf{Z} = \underline{\mathbf{Z}}_{\overline{x}}
$$
qui fait la somme des entiers donnés. Cela résulte de l'adjonction à l'application
$1 : \mathbf{Z} \to f^{-1}\underline{\mathbf{Z}}$. En particulier, si
$f$ est étale et surjectif, alors $\text{Tr}_f$ est surjectif.

\medskip\noindent
Supposons que $f : U \to X$ soit un morphisme étale surjectif
d'espaces algébriques. Considérons le {\it complexe de Koszul}
associé à l'application trace définie ci-dessus
$$
\ldots \to \wedge^3f_!\underline{\mathbf{Z}} \to
\wedge^2f_!\underline{\mathbf{Z}} \to f_!\underline{\mathbf{Z}} \to
\underline{\mathbf{Z}} \to 0
$$
Ici, les puissances extérieures sont prises sur le faisceau d'anneaux $\underline{\mathbf{Z}}$.
Les applications sont définies par la règle
$$
e_1 \wedge \ldots \wedge e_n \longmapsto
\sum\nolimits_{i = 1, \ldots, n} (-1)^{i + 1}
\text{Tr}_f(e_i)
e_1 \wedge \ldots \wedge \widehat{e_i} \wedge \ldots \wedge e_n
$$
où $e_1, \ldots, e_n$ sont des sections locales de $f_!\underline{\mathbf{Z}}$.
Soit $\overline{x}$ un point géométrique de $X$ et posons
$M_{\overline{x}} = (f_!\underline{\mathbf{Z}})_{\overline{x}} =
\bigoplus_{\overline{u}} \mathbf{Z}$. La fibre du complexe ci-dessus en
$\overline{x}$ est le complexe
$$
\ldots \to \wedge^3 M_{\overline{x}} \to \wedge^2 M_{\overline{x}}
\to M_{\overline{x}} \to \mathbf{Z} \to 0
$$
qui est exact puisque $M_{\overline{x}} \to \mathbf{Z}$ est surjectif, voir
Compléments d'algèbre, lemme \ref{more-algebra-lemma-homotopy-koszul-abstract}.
Ainsi, si $K^\bullet = K^\bullet(f)$ désigne le complexe défini par
$K^i = \wedge^{i + 1}f_!\underline{\mathbf{Z}}$, on obtient un
quasi-isomorphisme
\begin{equation}
\label{spaces-cohomology-equation-quasi-isomorphism}
K^\bullet \longrightarrow \underline{\mathbf{Z}}[0]
\end{equation}
Nous utilisons le complexe $K^\bullet$ pour définir ce que nous appelons
le complexe de {\v C}ech alterné associé à $f : U \to X$.

\begin{definition}
\label{spaces-cohomology-definition-alternating-cech-complex}
Soit $S$ un schéma. Soit $f : U \to X$ un morphisme étale surjectif
d'espaces algébriques sur $S$. Soit $\mathcal{F}$ un objet de
$\textit{Ab}(X_\etale)$. Le
{\it complexe de {\v C}ech alterné}\footnote{Cette notation peut être non standard}
$\check{\mathcal{C}}^\bullet_{alt}(f, \mathcal{F})$
associé à $\mathcal{F}$ et $f$ est le complexe
$$
\Hom(K^0, \mathcal{F}) \to \Hom(K^1, \mathcal{F}) \to
\Hom(K^2, \mathcal{F}) \to \ldots
$$
dont les groupes de morphismes sont calculés dans $\textit{Ab}(X_\etale)$.
\end{definition}

\noindent
Le lecteur peut vérifier que si $U = \coprod U_i$ et si $f|_{U_i} : U_i \to X$
est l'immersion ouverte d'un sous-espace, alors
$\check{\mathcal{C}}_{alt}^\bullet(f, \mathcal{F})$ coïncide avec le complexe
introduit dans
Cohomologie, section \ref{cohomology-section-alternating-cech}
pour le recouvrement de Zariski $X = \bigcup U_i$ et la restriction
de $\mathcal{F}$ au site de Zariski de $X$. Mais surtout,
il importe cependant de relier la cohomologie du
{\v C}ech alterné à la cohomologie.

\begin{lemma}
\label{spaces-cohomology-lemma-alternating-cech-to-cohomology}
Soit $S$ un schéma. Soit $f : U \to X$ un morphisme étale surjectif
d'espaces algébriques sur $S$. Soit $\mathcal{F}$ un objet de
$\textit{Ab}(X_\etale)$. Il existe une application canonique
$$
\check{\mathcal{C}}^\bullet_{alt}(f, \mathcal{F})
\longrightarrow
R\Gamma(X, \mathcal{F})
$$
dans $D(\textit{Ab})$. De plus, il existe une suite spectrale dont la page $E_1$ est
$$
E_1^{p, q} =
\Ext_{\textit{Ab}(X_\etale)}^q(K^p, \mathcal{F})
$$
convergeant vers $H^{p + q}(X, \mathcal{F})$, où
$K^p = \wedge^{p + 1}f_!\underline{\mathbf{Z}}$.
\end{lemma}

\begin{proof}
Rappelons le quasi-isomorphisme
$K^\bullet \to \underline{\mathbf{Z}}[0]$, voir
(\ref{spaces-cohomology-equation-quasi-isomorphism}).
Choisissons une résolution injective $\mathcal{F} \to \mathcal{I}^\bullet$
dans $\textit{Ab}(X_\etale)$. Considérons le bicomplexe
$\Hom(K^\bullet, \mathcal{I}^\bullet)$ dont les termes sont
$\Hom(K^p, \mathcal{I}^q)$. La différentielle
$d_1^{p, q} : A^{p, q} \to A^{p + 1, q}$
est induite par la différentielle $K^{p + 1} \to K^p$
et la différentielle $d_2^{p, q} : A^{p, q} \to A^{p, q + 1}$ est induite par
la différentielle
$\mathcal{I}^q \to \mathcal{I}^{q + 1}$.
Notons $\text{Tot}(\Hom(K^\bullet, \mathcal{I}^\bullet))$
le complexe total associé, voir
Homologie, section \ref{homology-section-double-complexes}.
Nous utiliserons les deux suites spectrales
$({}'E_r, {}'d_r)$ et $({}''E_r, {}''d_r)$
associées à ce bicomplexe, voir
Homologie, section \ref{homology-section-double-complex}.

\medskip\noindent
Comme $K^\bullet$ est une résolution de $\underline{\mathbf{Z}}$,
on voit que les complexes
$$
\Hom(K^\bullet, \mathcal{I}^q) :
\Hom(K^0, \mathcal{I}^q) \to
\Hom(K^1, \mathcal{I}^q) \to
\Hom(K^2, \mathcal{I}^q) \to \ldots
$$
sont acycliques en degrés strictement positifs et ont pour $H^0$
$\Gamma(X, \mathcal{I}^q)$. Par conséquent, d'après
Homologie, lemme \ref{homology-lemma-double-complex-gives-resolution}
l'application naturelle
$$
\mathcal{I}^\bullet(X) \longrightarrow
\text{Tot}(\Hom(K^\bullet, \mathcal{I}^\bullet))
$$
est un quasi-isomorphisme de complexes de groupes abéliens.
En particulier, on conclut que
$H^n(\text{Tot}(\Hom(K^\bullet, \mathcal{I}^\bullet))) = H^n(X, \mathcal{F})$.

\medskip\noindent
L'application $\check{\mathcal{C}}^\bullet_{alt}(f, \mathcal{F}) \to
R\Gamma(X, \mathcal{F})$ du lemme est la composée de
$\check{\mathcal{C}}^\bullet_{alt}(f, \mathcal{F}) \to
\text{Tot}(\Hom(K^\bullet, \mathcal{I}^\bullet))$
avec l'inverse du quasi-isomorphisme affiché.

\medskip\noindent
Considérons enfin la suite spectrale $({}'E_r, {}'d_r)$.
Nous avons
$$
E_1^{p, q} = q\text{-ième cohomologie de }
\Hom(K^p, \mathcal{I}^0) \to
\Hom(K^p, \mathcal{I}^1) \to
\Hom(K^p, \mathcal{I}^2) \to \ldots
$$
Ceci démontre le lemme.
\end{proof}

\noindent
Il résulte du lemme qu'il est important de comprendre
les groupes d'extensions $\Ext_{\textit{Ab}(X_\etale)}(K^p, \mathcal{F})$,
c'est-à-dire les foncteurs dérivés à droite de
$\mathcal{F} \mapsto \Hom(K^p, \mathcal{F})$.

\begin{lemma}
\label{spaces-cohomology-lemma-compute}
Soit $S$ un schéma. Soit $f : U \to X$ un morphisme étale, surjectif et séparé
d'espaces algébriques sur $S$. Pour $p \geq 0$, posons
$$
W_p = U \times_X \ldots \times_X U \setminus \text{toutes les diagonales}
$$
où le produit fibré comporte $p + 1$ facteurs.
Le groupe $S_{p + 1}$ agit librement sur $W_p$ au-dessus de $X$ et
$$
\Hom(K^p, \mathcal{F}) = S_{p + 1}\text{-anti-invariants de }
\mathcal{F}(W_p)
$$
fonctoriellement en $\mathcal{F}$, où
$K^p = \wedge^{p + 1}f_!\underline{\mathbf{Z}}$.
\end{lemma}

\begin{proof}
Comme $U \to X$ est séparé, la diagonale $U \to U \times_X U$ est une
immersion fermée. Comme $U \to X$ est étale, cette diagonale
$U \to U \times_X U$ est une immersion ouverte, voir
Morphismes des espaces, lemmes
\ref{spaces-morphisms-lemma-etale-unramified} et
\ref{spaces-morphisms-lemma-diagonal-unramified-morphism}.
Ainsi $W_p$ est un sous-espace à la fois ouvert et fermé de
$U^{p + 1} = U \times_X \ldots \times_X U$. L'action de $S_{p + 1}$
sur $W_p$ est libre, car nous avons supprimé les points fixes de l'action.
D'après
lemme \ref{spaces-cohomology-lemma-product-is-tensor-product}
on obtient
$$
(f_!\underline{\mathbf{Z}})^{\otimes p + 1} =
f^{p + 1}_!\underline{\mathbf{Z}} = (W_p \to X)_!\underline{\mathbf{Z}}
\oplus Rest
$$
où $f^{p + 1} : U^{p + 1} \to X$ est le morphisme structural.
En examinant les fibres au-dessus d'un point géométrique $\overline{x}$ de $X$,
on voit que
$$
\left(
\bigoplus\nolimits_{\overline{u} \mapsto \overline{x}} \mathbf{Z}
\right)^{\otimes p + 1}
\longrightarrow
(W_p \to X)_!\underline{\mathbf{Z}}_{\overline{x}}
$$
est le quotient dont le noyau est engendré par tous les tenseurs
$1_{\overline{u}_0} \otimes \ldots \otimes 1_{\overline{u}_p}$
où $\overline{u}_i = \overline{u}_j$ pour certains $i \not = j$.
Ainsi, l'application quotient
$$
(f_!\underline{\mathbf{Z}})^{\otimes p + 1}
\longrightarrow
\wedge^{p + 1}f_!\underline{\mathbf{Z}}
$$
se factorise par $(W_p \to X)_!\underline{\mathbf{Z}}$ ; on obtient donc
$$
(f_!\underline{\mathbf{Z}})^{\otimes p + 1}
\longrightarrow
(W_p \to X)_!\underline{\mathbf{Z}}
\longrightarrow
\wedge^{p + 1}f_!\underline{\mathbf{Z}}
$$
Cela montre déjà que $\Hom(K^p, \mathcal{F})$ est (fonctoriellement) un
sous-groupe de
$$
\Hom((W_p \to X)_!\underline{\mathbf{Z}}, \mathcal{F}) = \mathcal{F}(W_p)
$$
Pour l'identifier aux éléments anti-invariants sous $S_{p + 1}$, il faut montrer que
la surjection $(W_p \to X)_!\underline{\mathbf{Z}}
\to \wedge^{p + 1}f_!\underline{\mathbf{Z}}$ est le quotient anti-invariant maximal
$S_{p + 1}$-quotient anti-invariant. Autrement dit, il faut montrer que
$\wedge^{p + 1}f_!\underline{\mathbf{Z}}$ est le quotient de
$(W_p \to X)_!\underline{\mathbf{Z}}$ par le sous-faisceau engendré par
les sections locales $s - \text{signature}(\sigma)\sigma(s)$, où $s$ est
une section locale de $(W_p \to X)_!\underline{\mathbf{Z}}$.
Cela se vérifie sur les fibres, où c'est clair.
\end{proof}

\begin{lemma}
\label{spaces-cohomology-lemma-twist}
Soit $S$ un schéma. Soit $W$ un espace algébrique sur $S$.
Soit $G$ un groupe fini agissant librement sur $W$.
Soit $U = W/G$, voir
Propriétés des espaces, lemme \ref{spaces-properties-lemma-quotient}.
Soit $\chi : G \to \{+1, -1\}$ un caractère.
Alors il existe un faisceau de $\mathbf{Z}$-modules localement libre de rang 1
$\underline{\mathbf{Z}}(\chi)$ sur $U_\etale$ tel que, pour tout
faisceau abélien $\mathcal{F}$ sur $U_\etale$, on ait
$$
H^0(W, \mathcal{F}|_W)^\chi =
H^0(U, \mathcal{F} \otimes_{\mathbf{Z}} \underline{\mathbf{Z}}(\chi))
$$
\end{lemma}

\begin{proof}
Le morphisme quotient $q : W \to U$ est un $G$-torseur ; autrement dit, il existe
un morphisme étale surjectif $U' \to U$ tel que
$W \times_U U' = \coprod_{g \in G} U'$ comme espaces munis de l'action de $G$ au-dessus de $U'$.
(On peut notamment prendre $U' = W$.) Ainsi, $q_*\underline{\mathbf{Z}}$ est un
module de $\mathbf{Z}$ localement libre de rang fini muni d'une action de $G$.
Pour tout point géométrique $\overline{u}$ de $U$, on obtient des $G$-équivariants
isomorphismes
$$
(q_*\underline{\mathbf{Z}})_{\overline{u}}
= \bigoplus\nolimits_{\overline{w} \mapsto \overline{u}} \mathbf{Z}
= \bigoplus\nolimits_{g \in G} \mathbf{Z} = \mathbf{Z}[G]
$$
où le second signe $=$ utilise un point géométrique
$\overline{w}_0$ au-dessus de $\overline{u}$ et
envoie le facteur correspondant à $g \in G$ sur le facteur
correspondant à $g(\overline{w}_0)$. Nous avons
$$
H^0(W, \mathcal{F}|_W) =
H^0(U, \mathcal{F} \otimes_\mathbf{Z} q_*\underline{\mathbf{Z}})
$$
car
$q_*\mathcal{F}|_W = \mathcal{F} \otimes_\mathbf{Z} q_*\underline{\mathbf{Z}}$
comme on peut le vérifier après restriction à $U'$.
$$
\underline{\mathbf{Z}}(\chi) =
(q_*\underline{\mathbf{Z}})^\chi \subset
q_*\underline{\mathbf{Z}}
$$
qui est le sous-faisceau des sections se transformant selon $\chi$.
Pour tout point géométrique $\overline{u}$ de $U$, on a
$$
\underline{\mathbf{Z}}(\chi)_{\overline{u}} =
\mathbf{Z} \cdot \sum\nolimits_g \chi(g) g
\subset
\mathbf{Z}[G] = (q_*\underline{\mathbf{Z}})_{\overline{u}}
$$
Il s'ensuit que $\underline{\mathbf{Z}}(\chi)$ est localement libre de
rang 1 (plus précisément, cela doit être vérifié après restriction à $U'$).
Notons que, pour tout $\mathbf{Z}$-module $M$, les $\chi$-semi-invariants
de $M[G]$ sont les éléments de la forme $m \cdot \sum\nolimits_g \chi(g) g$.
Ainsi, pour tout faisceau abélien $\mathcal{F}$ sur $U$, on a
$$
\left(\mathcal{F} \otimes_\mathbf{Z} q_*\underline{\mathbf{Z}}\right)^\chi
=
\mathcal{F} \otimes_\mathbf{Z} \underline{\mathbf{Z}}(\chi)
$$
car l'égalité se vérifie sur toutes les fibres. Le résultat du lemme s'obtient en
prenant les sections globales.
\end{proof}

\noindent
Nous pouvons maintenant tout rassembler et obtenir le résultat suivant
particulièrement agréable.

\begin{lemma}
\label{spaces-cohomology-lemma-alternating-spectral-sequence}
Soit $S$ un schéma. Soit $f : U \to X$ un morphisme d'espaces algébriques
sur $S$, surjectif, \'etale et séparé. Pour $p \geq 0$, posons
$$
W_p = U \times_X \ldots \times_X U \setminus \text{toutes les diagonales}
$$
(avec $p + 1$ facteurs), comme dans le lemme \ref{spaces-cohomology-lemma-compute}.
Soit $\chi_p : S_{p + 1} \to \{+1, -1\}$ le caractère signe.
Soient $U_p = W_p/S_{p + 1}$ et $\underline{\mathbf{Z}}(\chi_p)$ les objets définis comme dans le
lemme \ref{spaces-cohomology-lemma-twist}.
Alors la suite spectrale du
lemme \ref{spaces-cohomology-lemma-alternating-cech-to-cohomology}
a pour page $E_1$
$$
E_1^{p, q} =
H^q(U_p, \mathcal{F}|_{U_p} \otimes_\mathbf{Z} \underline{\mathbf{Z}}(\chi_p))
$$
et aboutit à $H^{p + q}(X, \mathcal{F})$.
\end{lemma}

\begin{proof}
Remarquons que, puisque l'action de $S_{p + 1}$ sur $W_p$ est au-dessus de $X$, on
obtient bien un morphisme $U_p \to X$. Puisque $W_p \to X$ est \'etale et que
$W_p \to U_p$ est \'etale surjectif, il s'ensuit
que $U_p \to X$ est lui aussi \'etale ; voir
Morphismes d'espaces, lemme \ref{spaces-morphisms-lemma-etale-local}.
Par conséquent, un objet injectif de
$\textit{Ab}(X_\etale)$ devient par restriction un objet injectif de
$\textit{Ab}(U_{p, \etale})$ ; voir
Cohomologie sur les sites, lemme \ref{sites-cohomology-lemma-cohomology-of-open}.
De plus, le foncteur
$\mathcal{G} \mapsto
\mathcal{G} \otimes_\mathbf{Z} \underline{\mathbf{Z}}(\chi_p))$
est une auto-équivalence de $\textit{Ab}(U_p)$, donc transforme les objets injectifs
en objets injectifs et est exact (car
$\underline{\mathbf{Z}}(\chi_p)$ est un
$\underline{\mathbf{Z}}$-module inversible). Ainsi, étant donnée une résolution injective
$\mathcal{F} \to \mathcal{I}^\bullet$ dans $\textit{Ab}(X_\etale)$,
le complexe
$$
\Gamma(U_p,
\mathcal{I}^0|_{U_p} \otimes_\mathbf{Z} \underline{\mathbf{Z}}(\chi_p))
\to
\Gamma(U_p,
\mathcal{I}^1|_{U_p} \otimes_\mathbf{Z} \underline{\mathbf{Z}}(\chi_p))
\to
\Gamma(U_p,
\mathcal{I}^2|_{U_p} \otimes_\mathbf{Z} \underline{\mathbf{Z}}(\chi_p))
\to \ldots
$$
calcule
$H^*(U_p,
\mathcal{F}|_{U_p} \otimes_\mathbf{Z} \underline{\mathbf{Z}}(\chi_p))$.
D'autre part, d'après
le lemme \ref{spaces-cohomology-lemma-twist},
il est égal au complexe des éléments anti-invariants sous $S_{p + 1}$ de
$$
\Gamma(W_p, \mathcal{I}^0) \to
\Gamma(W_p, \mathcal{I}^1) \to
\Gamma(W_p, \mathcal{I}^2) \to \ldots
$$
qui, d'après
le lemme \ref{spaces-cohomology-lemma-compute},
est égal au complexe
$$
\Hom(K^p, \mathcal{I}^0) \to
\Hom(K^p, \mathcal{I}^1) \to
\Hom(K^p, \mathcal{I}^2) \to \ldots
$$
qui calcule
$\Ext^*_{\textit{Ab}(X_\etale)}(K^p, \mathcal{F})$.
En combinant ce qui précède, on conclut.
\end{proof}





\section{Annulation en degrés supérieurs pour les faisceaux quasi-cohérents}
\label{spaces-cohomology-section-higher-vanishing}

\noindent
Dans cette section, nous montrons que, pour tout espace algébrique quasi-compact et
quasi-séparé $X$, il existe un entier
$n = n(X)$ tel que la cohomologie de tout faisceau quasi-cohérent
sur $X$ est nulle en degrés strictement supérieurs à $n$.

\begin{lemma}
\label{spaces-cohomology-lemma-quasi-coherent-twist}
Soient $S$, $W$, $G$, $U$ et $\chi$ comme dans le
lemme \ref{spaces-cohomology-lemma-twist}.
Si $\mathcal{F}$ est un $\mathcal{O}_U$-module quasi-cohérent,
alors $\mathcal{F} \otimes_{\mathbf{Z}} \underline{\mathbf{Z}}(\chi)$ l'est aussi.
\end{lemma}

\begin{proof}
La structure de $\mathcal{O}_U$-module est claire. Pour vérifier que
$\mathcal{F} \otimes_{\mathbf{Z}} \underline{\mathbf{Z}}(\chi)$
est quasi-cohérent, il suffit de le vérifier localement pour la topologie \'etale.
Comme $\underline{\mathbf{Z}}(\chi)$
est un $\underline{\mathbf{Z}}$-module localement libre de rang fini, on en déduit le lemme.
\end{proof}

\noindent
La proposition suivante est intéressante même lorsque $X$ est un schéma.
Elle constitue la généralisation naturelle du
lemme \ref{coherent-lemma-vanishing-nr-affines} du chapitre Cohomologie des schémas.
Avant de l'énoncer, remarquons que, pour tout morphisme \'etale
$f : U \to X$ d'un schéma affine vers un espace algébrique
quasi-séparé $X$, les fibres de $f$ sont universellement bornées ; en particulier,
il existe un entier $d$ tel que toutes les fibres de $|U| \to |X|$
soient de cardinal au plus $d$ ; c'est l'implication
$(\eta) \Rightarrow (\delta)$ du lemme
\ref{decent-spaces-lemma-bounded-fibres} du chapitre Espaces décents.

\begin{proposition}
\label{spaces-cohomology-proposition-vanishing}
Soit $S$ un schéma. Soit $X$ un espace algébrique sur $S$.
Supposons $X$ quasi-compact et séparé.
Soit $U$ un schéma affine, et soit
$f : U \to X$ un morphisme \'etale surjectif.
Soit $d$ un majorant du cardinal des fibres de
$|U| \to |X|$. Alors, pour tout $\mathcal{O}_X$-module quasi-cohérent $\mathcal{F}$,
on a $H^q(X, \mathcal{F}) = 0$ pour tout $q \geq d$.
\end{proposition}

\begin{proof}
Nous utiliserons la suite spectrale du
lemme \ref{spaces-cohomology-lemma-alternating-spectral-sequence}.
Le lemme s'applique, car $f$ est séparé puisque $U$ est séparé ; voir
Morphismes d'espaces, lemme
\ref{spaces-morphisms-lemma-compose-after-separated}.
Puisque $X$ est séparé, le schéma $U \times_X \ldots \times_X U$ est un sous-schéma fermé
de
$U \times_{\Spec(\mathbf{Z})} \ldots \times_{\Spec(\mathbf{Z})} U$
et est donc affine. Ainsi, $W_p$ est affine. Par conséquent, $U_p = W_p/S_{p + 1}$ est un
schéma affine d'après
Groupoïdes, proposition \ref{groupoids-proposition-finite-flat-equivalence}.
La discussion de la
section \ref{spaces-cohomology-section-higher-direct-image}
montre que la cohomologie des faisceaux quasi-cohérents sur $W_p$ (considéré comme espace
algébrique) coïncide avec la cohomologie du faisceau quasi-cohérent correspondant
sur le schéma affine sous-jacent ; elle est donc nulle en degrés positifs d'après
Cohomologie des schémas, lemme
\ref{coherent-lemma-quasi-coherent-affine-cohomology-zero}.
D'après
le lemme \ref{spaces-cohomology-lemma-quasi-coherent-twist},
les faisceaux
$\mathcal{F}|_{U_p} \otimes_\mathbf{Z} \underline{\mathbf{Z}}(\chi_p)$
sont quasi-cohérents. Ainsi,
$H^q(W_p,
\mathcal{F}|_{U_p} \otimes_\mathbf{Z} \underline{\mathbf{Z}}(\chi_p))$
est nul pour $q > 0$. Par définition de l'entier $d$, on voit que
$W_p = \emptyset$ pour $p \geq d$. Par conséquent,
$H^0(W_p,
\mathcal{F}|_{U_p} \otimes_\mathbf{Z} \underline{\mathbf{Z}}(\chi_p))$
est également nul pour $p \geq d$.
Cela démontre la proposition.
\end{proof}

\noindent
Dans le lemme suivant, nous établissons qu'un espace algébrique quasi-compact et
quasi-séparé est de dimension cohomologique finie
pour les modules quasi-cohérents. Nous ne précisons la borne que parce que
nous l'utiliserons plus loin pour démontrer un résultat analogue pour les images directes
supérieures.

\begin{lemma}
\label{spaces-cohomology-lemma-vanishing-quasi-separated}
Soit $S$ un schéma. Soit $X$ un espace algébrique sur $S$.
Supposons $X$ quasi-compact et quasi-séparé.
On peut alors choisir
\begin{enumerate}
\item un schéma affine $U$,
\item un morphisme \'etale surjectif $f : U \to X$,
\item un entier $d$ majorant les degrés des fibres de $U \to X$,
\item pour tout $p = 0, 1, \ldots, d$, un morphisme \'etale surjectif
$V_p \to U_p$, où $V_p$ est un schéma affine et $U_p$ est défini comme dans le
lemme \ref{spaces-cohomology-lemma-alternating-spectral-sequence}, et
\item un entier $d_p$ majorant le degré des fibres de $V_p \to U_p$.
\end{enumerate}
De plus, chaque fois que l'on dispose de (1) -- (5), pour tout
$\mathcal{O}_X$-module quasi-cohérent $\mathcal{F}$, on a $H^q(X, \mathcal{F}) = 0$ pour tout
$q \geq \max(d_p + p)$.
\end{lemma}

\begin{proof}
Puisque $X$ est quasi-compact, on peut trouver un morphisme \'etale surjectif
$U \to X$ avec $U$ affine ; voir
Propriétés d'espaces, lemme
\ref{spaces-properties-lemma-quasi-compact-affine-cover}.
D'après
Espaces décents, lemme \ref{decent-spaces-lemma-bounded-fibres},
les fibres de $f$ sont universellement bornées ; on peut donc trouver $d$.
On a $U_p = W_p/S_{p + 1}$ et $W_p \subset U \times_X \ldots \times_X U$
est ouvert et fermé. Puisque $X$ est quasi-séparé, les schémas
$W_p$ sont quasi-compacts, donc $U_p$ est quasi-compact.
Puisque $U$ est séparé, les schémas $W_p$ sont séparés ; ainsi,
$U_p$ est séparé d'après (la version absolue du)
Espaces, lemme \ref{spaces-lemma-quotient-finite-separated}.
D'après
Propriétés d'espaces, lemme
\ref{spaces-properties-lemma-quasi-compact-affine-cover},
on peut trouver les morphismes $V_p \to W_p$.
D'après
Espaces décents, lemme \ref{decent-spaces-lemma-bounded-fibres},
on peut trouver les entiers $d_p$.

\medskip\noindent
On utilise alors la suite spectrale
$$
E_1^{p, q} =
H^q(U_p, \mathcal{F}|_{U_p} \otimes_\mathbf{Z} \underline{\mathbf{Z}}(\chi_p))
\Rightarrow
H^{p + q}(X, \mathcal{F})
$$
voir
le lemme \ref{spaces-cohomology-lemma-alternating-spectral-sequence}.
Par définition de l'entier $d$, on voit que $U_p = 0$ pour $p \geq d$.
D'après la proposition \ref{spaces-cohomology-proposition-vanishing}
et
le lemme \ref{spaces-cohomology-lemma-quasi-coherent-twist},
on voit que
$H^q(U_p,
\mathcal{F}|_{U_p} \otimes_\mathbf{Z} \underline{\mathbf{Z}}(\chi_p))$
est nul pour $q \geq d_p$ lorsque $p = 0, \ldots, d$.
On en déduit le lemme.
\end{proof}









\section{Annulation des images directes supérieures}
\label{spaces-cohomology-section-vanishing-higher-direct-images}

\noindent
On applique les résultats de la
section \ref{spaces-cohomology-section-higher-vanishing}
afin d'obtenir l'annulation des images directes supérieures des faisceaux quasi-cohérents
pour les morphismes quasi-compacts et quasi-séparés. Ce résultat est utile, car
il permet, dans certaines situations, de raisonner par récurrence descendante sur le degré
cohomologique.

\begin{lemma}
\label{spaces-cohomology-lemma-vanishing-higher-direct-images}
Soit $S$ un schéma. Soit $f : X \to Y$ un
morphisme d'espaces algébriques sur $S$.
Supposons que
\begin{enumerate}
\item $f$ soit quasi-compact et quasi-séparé, et
\item $Y$ soit quasi-compact.
\end{enumerate}
Alors il existe un entier $n(X \to Y)$ tel que,
pour tout espace algébrique $Y'$, tout morphisme $Y' \to Y$
et tout faisceau quasi-cohérent $\mathcal{F}'$ sur $X' = Y' \times_Y X$,
les images directes supérieures $R^if'_*\mathcal{F}'$ soient nulles pour tout
$i \geq n(X \to Y)$.
\end{lemma}

\begin{proof}
Soit $V \to Y$ un morphisme \'etale surjectif, où $V$ est un schéma
affine ; voir
Propriétés d'espaces, lemme
\ref{spaces-properties-lemma-quasi-compact-affine-cover}.
Supposons le résultat prouvé pour le changement de base $f_V : V \times_Y X \to V$.
Le résultat vaut alors pour $f$ avec $n(X \to Y) = n(X_V \to V)$.
En effet, si $Y' \to Y$ et $\mathcal{F}'$ sont comme dans le lemme, alors
$R^if'_*\mathcal{F}'|_{V \times_Y Y'}$ est égal à
$R^if'_{V, *}\mathcal{F}'|_{X'_V}$, où
$f'_V : X'_V = V \times_Y Y' \times_Y X \to V \times_Y Y' = Y'_V$ ; voir
Propriétés d'espaces,
lemme \ref{spaces-properties-lemma-pushforward-etale-base-change-modules}.
On peut donc supposer que $Y$ est un schéma affine.

\medskip\noindent
De plus, pour démontrer l'annulation pour tout $Y' \to Y$ et toute
$\mathcal{F}'$, il suffit de le faire lorsque $Y'$ est un schéma affine.
Dans ce cas, $R^if'_*\mathcal{F}'$ est quasi-cohérent d'après
le lemme \ref{spaces-cohomology-lemma-higher-direct-image}.
Il suffit donc de démontrer que $H^i(X', \mathcal{F}') = 0$, car
$H^i(X', \mathcal{F}') = H^0(Y', R^if'_*\mathcal{F}')$ d'après
Cohomologie sur les sites, lemme \ref{sites-cohomology-lemma-apply-Leray}
et le fait que la cohomologie des faisceaux quasi-cohérents
sur les espaces algébriques affines est nulle en degrés supérieurs
(proposition \ref{spaces-cohomology-proposition-vanishing}).

\medskip\noindent
Choisissons $U \to X$, $d$, $V_p \to U_p$ et $d_p$ comme dans
le lemme \ref{spaces-cohomology-lemma-vanishing-quasi-separated}.
Pour tout schéma affine $Y'$ et tout morphisme $Y' \to Y$, notons
$X' = Y' \times_Y X$, $U' = Y' \times_Y U$, $V'_p = Y' \times_Y V_p$.
Alors $U' \to X'$, $d' = d$, $V'_p \to U'_p$ et $d'_p = d$
forment un ensemble de choix comme dans
le lemme \ref{spaces-cohomology-lemma-vanishing-quasi-separated}
pour l'espace algébrique $X'$ (détails omis).
On voit donc que $H^i(X', \mathcal{F}') = 0$ pour $i \geq \max(p + d_p)$,
et l'on gagne.
\end{proof}

\begin{lemma}
\label{spaces-cohomology-lemma-affine-vanishing-higher-direct-images}
Soit $S$ un schéma. Soit $f : X \to Y$ un morphisme affine
d'espaces algébriques sur $S$. Alors
$R^if_*\mathcal{F} = 0$ pour $i > 0$ et tout
$\mathcal{O}_X$-module quasi-cohérent $\mathcal{F}$.
\end{lemma}

\begin{proof}
Rappelons qu'un morphisme affine d'espaces algébriques est représentable.
L'assertion résulte donc de (\ref{spaces-cohomology-equation-representable-higher-direct-image}) et du
lemme \ref{coherent-lemma-relative-affine-vanishing} du chapitre Cohomologie des schémas.
\end{proof}

\begin{lemma}
\label{spaces-cohomology-lemma-relative-affine-cohomology}
Soit $S$ un schéma. Soit $f : X \to Y$ un morphisme affine
d'espaces algébriques sur $S$.
Soit $\mathcal{F}$ un $\mathcal{O}_X$-module quasi-cohérent.
Alors $H^i(X, \mathcal{F}) = H^i(Y, f_*\mathcal{F})$ pour tout $i \geq 0$.
\end{lemma}

\begin{proof}
L'assertion résulte du lemme \ref{spaces-cohomology-lemma-affine-vanishing-higher-direct-images}
et de la suite spectrale de Leray. Voir
Cohomologie sur les sites, lemme \ref{sites-cohomology-lemma-apply-Leray}.
\end{proof}






\section{Cohomologie à support dans un sous-espace fermé}
\label{spaces-cohomology-section-cohomology-support}

\noindent
Cette section est l'analogue des sections du chapitre Cohomologie
\ref{cohomology-section-cohomology-support} et
\ref{cohomology-section-cohomology-support-bis}
et de la section du chapitre Cohomologie \'etale
\ref{etale-cohomology-section-cohomology-support}
pour les faisceaux abéliens sur les espaces algébriques.

\medskip\noindent
Soit $S$ un schéma.
Soit $X$ un espace algébrique sur $S$ et soit $Z \subset X$ un
sous-espace fermé. Soit $\mathcal{F}$ un faisceau abélien sur $X_\etale$. On pose
$$
\Gamma_Z(X, \mathcal{F}) =
\{s \in \mathcal{F}(X) \mid \text{Supp}(s) \subset Z\}
$$
l'ensemble des sections à support dans $Z$
(Propriétés d'espaces, définition \ref{spaces-properties-definition-support}).
C'est un foncteur exact à gauche qui n'est pas exact en général.
On obtient donc un foncteur dérivé
$$
R\Gamma_Z(X, -) : D(X_\etale) \longrightarrow D(\textit{Ab})
$$
et les groupes de cohomologie à support dans $Z$ définis par
$H^q_Z(X, \mathcal{F}) = R^q\Gamma_Z(X, \mathcal{F})$.

\medskip\noindent
Soit $\mathcal{I}$ un faisceau abélien injectif sur $X_\etale$. Soit
$U \subset X$ le sous-espace ouvert complémentaire de $Z$.
Alors l'application de restriction $\mathcal{I}(X) \to \mathcal{I}(U)$ est surjective
(Cohomologie sur les sites, lemme
\ref{sites-cohomology-lemma-restriction-along-monomorphism-surjective})
et a pour noyau $\Gamma_Z(X, \mathcal{I})$. Il s'ensuit immédiatement que
pour $K \in D(X_\etale)$ il existe un triangle distingué
$$
R\Gamma_Z(X, K) \to R\Gamma(X, K) \to R\Gamma(U, K) \to R\Gamma_Z(X, K)[1]
$$
dans $D(\textit{Ab})$. On obtient ainsi une suite exacte longue
de cohomologie
$$
\ldots \to H^i_Z(X, K) \to H^i(X, K) \to H^i(U, K) \to
H^{i + 1}_Z(X, K) \to \ldots
$$
pour tout $K$ dans $D(X_\etale)$.

\medskip\noindent
Pour un faisceau abélien $\mathcal{F}$ sur $X_\etale$, on peut considérer
le {\it sous-faisceau des sections à support dans $Z$}, noté
$\mathcal{H}_Z(\mathcal{F})$, défini par la formule
$$
\mathcal{H}_Z(\mathcal{F})(U) =
\{s \in \mathcal{F}(U) \mid \text{Supp}(s) \subset U \times_X Z\}
$$
Nous utilisons ici le support d'une section défini dans
Propriétés d'espaces, définition \ref{spaces-properties-definition-support}.
D'après l'équivalence de Morphismes d'espaces, lemme
\ref{spaces-morphisms-lemma-closed-immersion-push-pull}
on peut considérer $\mathcal{H}_Z(\mathcal{F})$ comme un faisceau abélien sur
$Z_\etale$. On obtient ainsi un foncteur
$$
\textit{Ab}(X_\etale) \longrightarrow \textit{Ab}(Z_\etale),\quad
\mathcal{F} \longmapsto \mathcal{H}_Z(\mathcal{F})
$$
qui est exact à gauche, mais n'est pas exact en général.

\begin{lemma}
\label{spaces-cohomology-lemma-sections-with-support-acyclic}
Soit $S$ un schéma.
Soit $i : Z \to X$ une immersion fermée d'espaces algébriques sur $S$.
Soit $\mathcal{I}$ un faisceau abélien injectif sur $X_\etale$.
Alors $\mathcal{H}_Z(\mathcal{I})$ est un faisceau abélien injectif
sur $Z_\etale$.
\end{lemma}

\begin{proof}
Remarquons que, pour tout faisceau abélien $\mathcal{G}$ sur $Z_\etale$,
on a
$$
\Hom_Z(\mathcal{G}, \mathcal{H}_Z(\mathcal{F})) =
\Hom_X(i_*\mathcal{G}, \mathcal{F})
$$
car toute section de $i_*\mathcal{G}$ est à support dans $Z$.
Comme $i_*$ est exact (lemme \ref{spaces-cohomology-lemma-finite-higher-direct-image-zero})
et que $\mathcal{I}$ est injectif sur $X_\etale$, on conclut que
$\mathcal{H}_Z(\mathcal{I})$ est injectif sur $Z_\etale$.
\end{proof}

\noindent
Notons
$$
R\mathcal{H}_Z : D(X_\etale) \longrightarrow D(Z_\etale)
$$
le foncteur dérivé. On pose
$\mathcal{H}^q_Z(\mathcal{F}) = R^q\mathcal{H}_Z(\mathcal{F})$, de sorte que
$\mathcal{H}^0_Z(\mathcal{F}) = \mathcal{H}_Z(\mathcal{F})$.
D'après le lemme précédent, on a une suite spectrale de Grothendieck
$$
E_2^{p, q} = H^p(Z, \mathcal{H}^q_Z(\mathcal{F}))
\Rightarrow H^{p + q}_Z(X, \mathcal{F})
$$

\begin{lemma}
\label{spaces-cohomology-lemma-cohomology-with-support-sheaf-on-support}
Soit $S$ un schéma. Soit $i : Z \to X$ une immersion fermée
d'espaces algébriques sur $S$. Soit $\mathcal{G}$ un faisceau abélien injectif
sur $Z_\etale$. Alors $\mathcal{H}^p_Z(i_*\mathcal{G}) = 0$ pour $p > 0$.
\end{lemma}

\begin{proof}
Cela résulte du fait que le foncteur $i_*$ est exact
(lemme \ref{spaces-cohomology-lemma-finite-higher-direct-image-zero}) et transforme
les faisceaux abéliens injectifs en faisceaux abéliens injectifs
(Cohomologie sur les sites, lemme
\ref{sites-cohomology-lemma-pushforward-injective-flat}).
\end{proof}

\begin{lemma}
\label{spaces-cohomology-lemma-etale-localization-sheaf-with-support}
Soit $S$ un schéma. Soit $f : X \to Y$ un morphisme \'etale
d'espaces algébriques sur $S$. Soit $Z \subset Y$ un sous-espace fermé
tel que $f^{-1}(Z) \to Z$ soit un isomorphisme d'espaces algébriques.
Soit $\mathcal{F}$ un faisceau abélien sur $X$. Alors
$$
\mathcal{H}^q_Z(\mathcal{F}) = \mathcal{H}^q_{f^{-1}(Z)}(f^{-1}\mathcal{F})
$$
en tant que faisceaux abéliens sur $Z = f^{-1}(Z)$ et on
a $H^q_Z(Y, \mathcal{F}) = H^q_{f^{-1}(Z)}(X, f^{-1}\mathcal{F})$.
\end{lemma}

\begin{proof}
Comme $f$ est \'etale, l'image inverse d'une résolution injective de $\mathcal{F}$
est une résolution injective de $f^{-1}\mathcal{F}$.
Il suffit donc de vérifier l'égalité pour $\mathcal{H}_Z(-)$,
ce qui résulte des définitions. La démonstration pour la cohomologie à
support est la même. Certains détails sont omis.
\end{proof}

\noindent
Soit $S$ un schéma et soit $X$ un espace algébrique sur $S$.
Soit $T \subset |X|$ un sous-ensemble fermé.
On note $D_T(X_\etale)$ la
sous-catégorie triangulée strictement pleine et saturée de $D(X_\etale)$
formée des objets dont les faisceaux de cohomologie sont à support dans $T$.

\begin{lemma}
\label{spaces-cohomology-lemma-complexes-with-support-on-closed}
Soit $S$ un schéma.
Soit $i : Z \to X$ une immersion fermée d'espaces algébriques sur $S$.
L'application $Ri_* = i_* : D(Z_\etale) \to D(X_\etale)$
induit une équivalence $D(Z_\etale) \to D_{|Z|}(X_\etale)$ de quasi-inverse
$$
i^{-1}|_{D_Z(X_\etale)} = R\mathcal{H}_Z|_{D_{|Z|}(X_\etale)}
$$
\end{lemma}

\begin{proof}
Rappelons que $i^{-1}$ et $i_*$ forment une paire de foncteurs adjoints
exacts telle que $i^{-1}i_*$ est isomorphe au foncteur identité
sur les faisceaux abéliens. Voir
Propriétés d'espaces, lemme
\ref{spaces-properties-lemma-stalk-pullback} et
Morphismes d'espaces, lemme
\ref{spaces-morphisms-lemma-closed-immersion-push-pull}.
Ainsi $i_* : D(Z_\etale) \to D_Z(X_\etale)$ est pleinement fidèle et
$i^{-1}$ fournit
un inverse à gauche. D'autre part, supposons que $K$ soit un objet de
$D_Z(X_\etale)$ et considérons le morphisme d'adjonction
$K \to i_*i^{-1}K$.
Par exactitude de $i_*$ et de $i^{-1}$,
celui-ci induit les morphismes d'adjonction
$H^n(K) \to i_*i^{-1}H^n(K)$ sur les faisceaux de cohomologie.
Comme ces faisceaux de cohomologie
sont à support dans $Z$, ces morphismes d'adjonction sont des isomorphismes
et on conclut que $D(Z_\etale) \to D_Z(X_\etale)$ est une équivalence.

\medskip\noindent
Pour terminer la démonstration, il faut montrer que $R\mathcal{H}_Z(K) = i^{-1}K$
si $K$ est un objet de $D_Z(X_\etale)$. Pour cela, on peut utiliser l'égalité
$K = i_*i^{-1}K$
que nous venons de démontrer. On peut alors
choisir un représentant K-injectif $\mathcal{I}^\bullet$ de
$i^{-1}K$.
Puisque $i_*$ est l'adjoint à droite du foncteur exact
$i^{-1}$, le
complexe $i_*\mathcal{I}^\bullet$ est K-injectif
(Catégories dérivées, lemme \ref{derived-lemma-adjoint-preserve-K-injectives}).
On voit que $R\mathcal{H}_Z(K)$ est calculé par
$\mathcal{H}_Z(i_*\mathcal{I}^\bullet) = \mathcal{I}^\bullet$,
comme voulu.
\end{proof}






\section{Annulation au-delà de la dimension}
\label{spaces-cohomology-section-vanishing-above-dimension}

\noindent
Soit $S$ un schéma. Soit $X$ un espace algébrique quasi-compact et quasi-séparé
sur $S$. Dans ce cas, $|X|$ est un espace spectral, voir
Propriétés d'espaces, lemme
\ref{spaces-properties-lemma-quasi-compact-quasi-separated-spectral}.
De plus, la dimension de $X$ (telle qu'elle est définie dans
Propriétés d'espaces, définition \ref{spaces-properties-definition-dimension})
est égale à la dimension de Krull de $|X|$, voir
Espaces décents, lemme \ref{decent-spaces-lemma-dimension-decent-space}.
Nous allons montrer que, pour les faisceaux quasi-cohérents sur $X$, la
cohomologie s'annule au-delà de la dimension. Ce résultat est déjà intéressant pour les
espaces algébriques quasi-séparés de type fini sur un corps.

\begin{lemma}
\label{spaces-cohomology-lemma-vanishing-above-dimension}
Soit $S$ un schéma. Soit $X$ un espace algébrique quasi-compact et quasi-séparé
sur $S$. Supposons que $\dim(X) \leq d$ pour un entier $d$.
Soit $\mathcal{F}$ un faisceau quasi-cohérent $\mathcal{F}$ sur $X$.
\begin{enumerate}
\item $H^q(X, \mathcal{F}) = 0$ pour $q > d$,
\item $H^d(X, \mathcal{F}) \to H^d(U, \mathcal{F})$ est surjective
pour tout ouvert quasi-compact $U \subset X$,
\item $H^q_Z(X, \mathcal{F}) = 0$ pour $q > d$ pour tout sous-espace fermé
$Z \subset X$ dont le complémentaire est quasi-compact.
\end{enumerate}
\end{lemma}

\begin{proof}
D'après Propriétés d'espaces, lemme
\ref{spaces-properties-lemma-dimension-decent-invariant-under-etale}
tout espace algébrique $Y$ \'etale sur $X$ est de dimension $\leq d$.
Si $Y$ est quasi-séparé, la dimension de $Y$ est égale à la
dimension de Krull de $|Y|$ d'après 
Espaces décents, lemme \ref{decent-spaces-lemma-dimension-decent-space}.
De plus, si $Y$ est un schéma, alors la cohomologie \'etale de $\mathcal{F}$
sur $Y$, resp.\ la cohomologie \'etale de $\mathcal{F}$ à support dans un
sous-schéma fermé coïncide avec la cohomologie usuelle de $\mathcal{F}$,
resp.\ avec la cohomologie usuelle à support dans le sous-schéma fermé.
Voir
Descente, proposition \ref{descent-proposition-same-cohomology-quasi-coherent}
et
Cohomologie \'etale, lemme
\ref{etale-cohomology-lemma-cohomology-with-support-quasi-coherent}.
Nous utiliserons ces faits sans autre mention.

\medskip\noindent
D'après Espaces décents, lemme
\ref{decent-spaces-lemma-filter-quasi-compact-quasi-separated}
il existe un entier $n$ et des sous-espaces ouverts
$$
\emptyset = U_{n + 1} \subset
U_n \subset U_{n - 1} \subset \ldots \subset U_1 = X
$$
ayant la propriété suivante : si l'on pose $T_p = U_p \setminus U_{p + 1}$
(muni de la structure réduite induite), il existe un schéma séparé
quasi-compact $V_p$ et un morphisme \'etale surjectif $f_p : V_p \to U_p$
tel que $f_p^{-1}(T_p) \to T_p$ soit un isomorphisme.

\medskip\noindent
Comme $U_n = V_n$ est un schéma, nos remarques initiales montrent que la cohomologie de
$\mathcal{F}$ sur $U_n$ s'annule en degrés $> d$ d'après
Cohomologie, proposition
\ref{cohomology-proposition-cohomological-dimension-spectral}.
Supposons avoir montré, par récurrence, que
$H^q(U_{p + 1}, \mathcal{F}|_{U_{p + 1}}) = 0$ pour $q > d$.
Il suffit de montrer que $H_{T_p}^q(U_p, \mathcal{F})$ est nul pour
$q > d$ afin d'en déduire l'annulation de la cohomologie
de $\mathcal{F}$ sur $U_p$ en degrés $> d$.
Cependant, on a
$$
H^q_{T_p}(U_p, \mathcal{F}) = H^q_{f_p^{-1}(T_p)}(V_p, \mathcal{F})
$$
d'après le lemme \ref{spaces-cohomology-lemma-etale-localization-sheaf-with-support}
et, puisque $V_p$ est un schéma, l'annulation voulue résulte de
Cohomologie, proposition
\ref{cohomology-proposition-cohomological-dimension-spectral}.
On conclut ainsi que (1) est vraie.

\medskip\noindent
Pour démontrer (2), soit $U \subset X$ un sous-espace ouvert quasi-compact.
Considérons le sous-espace ouvert $U' = U \cup U_n$. Posons $Z = U' \setminus U$.
Alors $g : U_n \to U'$ est un morphisme \'etale tel que
$g^{-1}(Z) \to Z$ soit un isomorphisme. Par conséquent, d'après le
lemme \ref{spaces-cohomology-lemma-etale-localization-sheaf-with-support},
on a $H^q_Z(U', \mathcal{F}) = H^q_Z(U_n, \mathcal{F})$,
qui s'annule en degré $> d$ puisque $U_n$ est un schéma
et que l'on peut appliquer
Cohomologie, proposition
\ref{cohomology-proposition-cohomological-dimension-spectral}.
On conclut que $H^d(U', \mathcal{F}) \to H^d(U, \mathcal{F})$
est surjective. Supposons, par récurrence, avoir ramené
notre problème au cas où $U$ contient $U_{p + 1}$.
On pose alors $U' = U \cup U_p$, on pose $Z = U' \setminus U$, et
on raisonne au moyen du morphisme $f_p : V_p \to U'$ qui est \'etale
et tel que $f_p^{-1}(Z) \to Z$ soit un isomorphisme.
Autrement dit, on voit de nouveau que
$$
H^q_Z(U', \mathcal{F}) = H^q_{f_p^{-1}(Z)}(V_p, \mathcal{F})
$$
et l'on voit de nouveau que ce groupe s'annule en degrés $> d$.
On conclut que $H^d(U', \mathcal{F}) \to H^d(U, \mathcal{F})$
est surjective. On atteint finalement le stade où $U_1 = X \subset U$,
ce qui achève la démonstration.

\medskip\noindent
Un argument formel montre que (2) implique (3).
\end{proof}








\section{Cohomologie et changement de base, I}
\label{spaces-cohomology-section-cohomology-and-base-change}

\noindent
Soit $S$ un schéma.
Soit $f : X \to Y$ un morphisme d'espaces algébriques au-dessus de $S$.
Soit $\mathcal{F}$ un faisceau quasi-cohérent sur $X$.
Supposons en outre que $g : Y' \to Y$ soit un morphisme d'espaces algébriques au-dessus de
$S$. Notons $X' = X_{Y'} = Y' \times_Y X$ le changement de base de $X$ et notons
$f' : X' \to Y'$ le changement de base de $f$.
Notons aussi $g' : X' \to X$ la projection,
et posons $\mathcal{F}' = (g')^*\mathcal{F}$.
La situation est représentée par le diagramme suivant :
\begin{equation}
\label{spaces-cohomology-equation-base-change-diagram}
\vcenter{
\xymatrix{
\mathcal{F}' = (g')^*\mathcal{F} &
X' \ar[r]_{g'} \ar[d]_{f'} &
X \ar[d]^f &
\mathcal{F} \\
Rf'_*\mathcal{F}' &
Y' \ar[r]^g &
Y &
Rf_*\mathcal{F}
}
}
\end{equation}
Voici le cas le plus simple de la propriété de changement de base que nous avons en vue.

\begin{lemma}
\label{spaces-cohomology-lemma-affine-base-change}
Soit $S$ un schéma. Soit $f : X \to Y$ un morphisme affine d'espaces
algébriques au-dessus de $S$. Soit $\mathcal{F}$ un $\mathcal{O}_X$-module quasi-cohérent.
Dans ce cas, $f_*\mathcal{F} \cong Rf_*\mathcal{F}$ est un faisceau quasi-cohérent,
et, pour tout diagramme (\ref{spaces-cohomology-equation-base-change-diagram}),
on a
$$
g^*f_*\mathcal{F} = f'_*(g')^*\mathcal{F}.
$$
\end{lemma}

\begin{proof}
D'après la discussion autour de
(\ref{spaces-cohomology-equation-representable-higher-direct-image}),
on se ramène au cas d'un morphisme affine de schémas, qui
est traité dans Cohomologie des schémas, lemme
\ref{coherent-lemma-affine-base-change}.
\end{proof}

\begin{lemma}[Changement de base plat]
\label{spaces-cohomology-lemma-flat-base-change-cohomology}
Soit $S$ un schéma. Considérons un diagramme cartésien d'espaces algébriques
$$
\xymatrix{
X' \ar[d]_{f'} \ar[r]_{g'} & X \ar[d]^f \\
Y' \ar[r]^g & Y
}
$$
au-dessus de $S$.
Soit $\mathcal{F}$ un $\mathcal{O}_X$-module quasi-cohérent
dont l'image inverse est $\mathcal{F}' = (g')^*\mathcal{F}$.
Supposons que $g$ soit plat et que $f$ soit quasi-compact et quasi-séparé.
Pour tout $i \geq 0$,
\begin{enumerate}
\item la flèche de changement de base de
Cohomologie sur les sites, lemme
\ref{sites-cohomology-lemma-base-change-map-flat-case}
est un isomorphisme
$$
g^*R^if_*\mathcal{F} \longrightarrow R^if'_*\mathcal{F}',
$$
\item si $Y = \Spec(A)$ et $Y' = \Spec(B)$, alors
$H^i(X, \mathcal{F}) \otimes_A B = H^i(X', \mathcal{F}')$.
\end{enumerate}
\end{lemma}

\begin{proof}
Le morphisme $g'$ est plat d'après
Morphismes d'espaces, lemme \ref{spaces-morphisms-lemma-base-change-flat}.
Notons que la platitude de $g$ et de $g'$ équivaut à la platitude
des morphismes de petits sites étales annelés, voir
Morphismes d'espaces, lemme \ref{spaces-morphisms-lemma-flat-morphism-sites}.
On peut donc appliquer
Cohomologie sur les sites, lemme
\ref{sites-cohomology-lemma-base-change-map-flat-case}
pour obtenir une flèche de changement de base
$$
g^*R^pf_*\mathcal{F} \longrightarrow R^pf'_*\mathcal{F}'
$$
Pour montrer que cette flèche est un isomorphisme, on peut travailler localement pour la
topologie étale sur $Y'$. On peut donc supposer que $Y$ et $Y'$ sont des schémas
affines. Posons $Y = \Spec(A)$ et $Y' = \Spec(B)$.
Dans ce cas, il s'agit en fait de montrer que la flèche
$$
H^p(X, \mathcal{F}) \otimes_A B \longrightarrow H^p(X_B, \mathcal{F}_B)
$$
est un isomorphisme, où $X_B = \Spec(B) \times_{\Spec(A)} X$ et
$\mathcal{F}_B$ est l'image inverse de $\mathcal{F}$ sur $X_B$.
Autrement dit, il suffit de démontrer (2).

\medskip\noindent
Fixons un homomorphisme plat d'anneaux $A \to B$ et soit $X$ un espace algébrique
quasi-compact et quasi-séparé sur $A$. Notons que $g' : X_B \to X$
est affine en tant que changement de base de $\Spec(B) \to \Spec(A)$. Par conséquent,
les images directes supérieures $R^i(g')_*\mathcal{F}_B$ sont nulles d'après le
lemme \ref{spaces-cohomology-lemma-affine-vanishing-higher-direct-images}.
Ainsi, $H^p(X_B, \mathcal{F}_B) = H^p(X, g'_*\mathcal{F}_B)$, voir
Cohomologie sur les sites, lemme \ref{sites-cohomology-lemma-apply-Leray}.
De plus, on a
$$
g'_*\mathcal{F}_B = \mathcal{F} \otimes_{\underline{A}} \underline{B}
$$
où $\underline{A}$ et $\underline{B}$ désignent les faisceaux constants d'anneaux
de valeurs respectives $A$ et $B$. En effet, il est clair qu'il existe une flèche
du membre de droite vers celui de gauche. Pour tout schéma affine $U$ étale sur $X$, on a
\begin{align*}
g'_*\mathcal{F}_B(U) & = \mathcal{F}_B(\Spec(B) \times_{\Spec(A)} U) \\
& =
\Gamma(\Spec(B) \times_{\Spec(A)} U,
(\Spec(B) \times_{\Spec(A)} U \to U)^*\mathcal{F}|_U) \\
& =
B \otimes_A \mathcal{F}(U)
\end{align*}
donc cette flèche est un isomorphisme. Écrivons $B = \colim M_i$ comme limite
inductive filtrante de $A$-modules libres de type fini $M_i$ en utilisant le théorème de Lazard, voir
Algèbre, théorème \ref{algebra-theorem-lazard}.
On en déduit que
\begin{align*}
H^p(X, g'_*\mathcal{F}_B) &
= H^p(X, \mathcal{F} \otimes_{\underline{A}} \underline{B}) \\
& = H^p(X, \colim_i \mathcal{F} \otimes_{\underline{A}} \underline{M_i}) \\
& = \colim_i H^p(X, \mathcal{F} \otimes_{\underline{A}} \underline{M_i}) \\
& = \colim_i H^p(X, \mathcal{F}) \otimes_A M_i \\
& = H^p(X, \mathcal{F}) \otimes_A \colim_i M_i \\
& = H^p(X, \mathcal{F}) \otimes_A B
\end{align*}
La première égalité vient de ce que
$g'_*\mathcal{F}_B = \mathcal{F} \otimes_{\underline{A}} \underline{B}$,
comme on l'a vu ci-dessus.
La deuxième vient de ce que $\otimes$ commute aux limites inductives.
La troisième égalité vient de ce que la cohomologie sur $X$ commute aux
limites inductives (voir le
lemme \ref{spaces-cohomology-lemma-colimits}).
La quatrième égalité vient de ce que $M_i$ est libre de type fini (c'est-à-dire de ce que la cohomologie
commute aux sommes directes finies).
La cinquième vient de ce que $\otimes$ commute aux limites inductives.
La sixième résulte du choix de notre système.
\end{proof}

\begin{lemma}
\label{spaces-cohomology-lemma-etale-pull-push-split}
Soit $f : X \to Y$ un morphisme quasi-compact, séparé et \'etale
d'espaces algébriques. Alors, pour tout $\mathcal{O}_X$-module quasi-cohérent
$\mathcal{F}$, le morphisme $f^*f_*\mathcal{F} \to \mathcal{F}$ est scindé.
\end{lemma}

\begin{proof}
Considérons le diagramme cartésien
$$
\xymatrix{
X \times_Y X \ar[r]_-p \ar[d]_q & X \ar[d]^f \\
X \ar[r]^f & Y
}
$$
D'après le Lemme \ref{spaces-cohomology-lemma-affine-base-change}, on a
$f^*f_*\mathcal{F} = q_*p^*\mathcal{F}$. Le morphisme
$\Delta : X \to X \times_Y X$ est ouvert (car c'est un morphisme entre
espaces algébriques \'etales sur $Y$) et fermé puisque $f$ est séparé.
On voit donc que $p^*\mathcal{F}$ est somme directe de
$\Delta_*\mathcal{F}$ et d'un module quasi-cohérent à support
dans le complémentaire (ouvert et fermé) de $\Delta(X)$.
En suivant les morphismes, le lecteur vérifie que
$$
\mathcal{F} = q_*\Delta_*\mathcal{F} \to
q_*p^*\mathcal{F} = f^*f_*\mathcal{F} \to
\mathcal{F}
$$
est le morphisme identité. Nous omettons les détails.
\end{proof}



\section{Modules cohérents sur les espaces algébriques localement noethériens}
\label{spaces-cohomology-section-coherent}

\noindent
Cette section est l'analogue de
Cohomologie des schémas, section \ref{coherent-section-coherent-sheaves}.
Dans Modules sur les sites, Définition \ref{sites-modules-definition-site-local},
nous avons défini les modules cohérents sur tout topos annelé. Nous employons cette notion
pour définir les modules cohérents sur les espaces algébriques localement noethériens.
Bien qu'il soit possible de travailler plus généralement avec des modules cohérents,
nous nous abstiendrons de le faire.

\begin{definition}
\label{spaces-cohomology-definition-coherent}
Soit $S$ un schéma. Soit $X$ un espace algébrique localement noethérien sur $S$.
Un module quasi-cohérent $\mathcal{F}$ sur $X$ est dit {\it cohérent}
si $\mathcal{F}$ est un $\mathcal{O}_X$-module cohérent sur le site
$X_\etale$ au sens de
Modules sur les sites, Définition \ref{sites-modules-definition-site-local}.
\end{definition}

\noindent
Cette définition est compatible avec la notion déjà existante de
module cohérent sur un schéma localement noethérien; voir l'assertion (5) de
Propriétés des espaces, section
\ref{spaces-properties-section-properties-modules} (ou, plus directement,
Descente, Lemme \ref{descent-lemma-equivalence-quasi-coherent-properties}).
Désormais, si $X$ est un schéma localement noethérien sur $S$,
nous ne distinguerons donc pas un module cohérent sur $X$ considéré
comme schéma d'un module cohérent sur $X$ considéré comme espace algébrique;
ceci est compatible avec les identifications correspondantes des catégories
de modules quasi-cohérents décrites dans
Propriétés des espaces, section \ref{spaces-properties-section-quasi-coherent}.

\medskip\noindent
Cela étant, le lemme suivant donne une caractérisation concrète
des modules cohérents sur les espaces algébriques localement
noethériens.

\begin{lemma}
\label{spaces-cohomology-lemma-coherent-Noetherian}
Soit $S$ un schéma.
Soit $X$ un espace algébrique localement noethérien sur $S$.
Soit $\mathcal{F}$ un $\mathcal{O}_X$-module.
Les conditions suivantes sont équivalentes :
\begin{enumerate}
\item $\mathcal{F}$ est cohérent,
\item $\mathcal{F}$ est un $\mathcal{O}_X$-module quasi-cohérent de type fini,
\item $\mathcal{F}$ est un $\mathcal{O}_X$-module de présentation finie,
\item pour tout morphisme \'etale $\varphi : U \to X$, où $U$ est un schéma,
l'image inverse $\varphi^*\mathcal{F}$ est un module cohérent sur $U$, et
\item il existe un morphisme \'etale surjectif $\varphi : U \to X$,
où $U$ est un schéma, tel que l'image inverse $\varphi^*\mathcal{F}$ soit
un module cohérent sur $U$.
\end{enumerate}
En particulier, $\mathcal{O}_X$ est cohérent, tout
$\mathcal{O}_X$-module inversible est cohérent et, plus généralement, tout
$\mathcal{O}_X$-module localement libre de type fini est cohérent.
\end{lemma}

\begin{proof}
Notons que, si $X$ est un espace algébrique localement noethérien et si
$U \to X$ est un morphisme \'etale, alors $U$ est localement noethérien; voir
Propriétés des espaces, section \ref{spaces-properties-section-types-properties}.
Le lemme résulte alors des points (1) -- (5) établis dans
Propriétés des espaces, section \ref{spaces-properties-section-properties-modules},
et du résultat correspondant pour les modules cohérents sur les schémas
localement noethériens; voir
Cohomologie des schémas, Lemme \ref{coherent-lemma-coherent-Noetherian}.
\end{proof}

\begin{lemma}
\label{spaces-cohomology-lemma-coherent-abelian-Noetherian}
Soit $S$ un schéma. Soit $X$ un espace algébrique localement noethérien sur $S$.
La catégorie des $\mathcal{O}_X$-modules cohérents est abélienne. Plus précisément,
le noyau et le conoyau d'un morphisme de $\mathcal{O}_X$-modules cohérents sont
cohérents. Toute extension de faisceaux cohérents est cohérente.
\end{lemma}

\begin{proof}
Choisissons un schéma $U$ et un morphisme \'etale surjectif $f : U \to X$.
L'image inverse $f^*$ est un foncteur exact, car elle est un foncteur de restriction; voir
Propriétés des espaces, Équation
(\ref{spaces-properties-equation-restrict-modules}).
D'après le
Lemme \ref{spaces-cohomology-lemma-coherent-Noetherian}, pour vérifier qu'un
$\mathcal{O}_X$-module $\mathcal{F}$ est
cohérent, il suffit de vérifier que $f^*\mathcal{F}$ est cohérent. Le
lemme résulte donc du cas des schémas, qui est
Cohomologie des schémas, Lemme \ref{coherent-lemma-coherent-abelian-Noetherian}.
\end{proof}

\noindent
Les modules cohérents forment une sous-catégorie de Serre de la
catégorie des $\mathcal{O}_X$-modules quasi-cohérents. Cela n'est pas vrai
pour les modules sur un topos annelé quelconque.

\begin{lemma}
\label{spaces-cohomology-lemma-coherent-Noetherian-quasi-coherent-sub-quotient}
Soit $S$ un schéma.
Soit $X$ un espace algébrique localement noethérien sur $S$.
Soit $\mathcal{F}$ un $\mathcal{O}_X$-module cohérent.
Tout sous-module quasi-cohérent de $\mathcal{F}$ est cohérent.
Tout module quotient quasi-cohérent de $\mathcal{F}$ est cohérent.
\end{lemma}

\begin{proof}
Choisissons un schéma $U$ et un morphisme \'etale surjectif $f : U \to X$.
L'image inverse $f^*$ est un foncteur exact, car elle est un foncteur de restriction; voir
Propriétés des espaces, Équation
(\ref{spaces-properties-equation-restrict-modules}).
D'après le
Lemme \ref{spaces-cohomology-lemma-coherent-Noetherian}, pour vérifier qu'un
$\mathcal{O}_X$-module $\mathcal{G}$ est
cohérent, il suffit de vérifier que $f^*\mathcal{H}$ est cohérent. Le
lemme résulte donc du cas des schémas, qui est
Cohomologie des schémas, Lemme
\ref{coherent-lemma-coherent-Noetherian-quasi-coherent-sub-quotient}.
\end{proof}

\begin{lemma}
\label{spaces-cohomology-lemma-tensor-hom-coherent}
Soit $S$ un schéma.
Soit $X$ un espace algébrique localement noethérien sur $S$.
Soient $\mathcal{F}$, $\mathcal{G}$ des $\mathcal{O}_X$-modules cohérents.
Les $\mathcal{O}_X$-modules $\mathcal{F} \otimes_{\mathcal{O}_X} \mathcal{G}$
et $\SheafHom_{\mathcal{O}_X}(\mathcal{F}, \mathcal{G})$ sont
cohérents.
\end{lemma}

\begin{proof}
Par le Lemme \ref{spaces-cohomology-lemma-coherent-Noetherian}, cela résulte du résultat
pour les schémas; voir
Cohomologie des schémas, Lemme \ref{coherent-lemma-tensor-hom-coherent}.
\end{proof}

\begin{lemma}
\label{spaces-cohomology-lemma-local-isomorphism}
Soit $S$ un schéma. Soit $X$ un espace algébrique localement noethérien sur $S$.
Soient $\mathcal{F}$, $\mathcal{G}$ des $\mathcal{O}_X$-modules cohérents.
Soit $\varphi : \mathcal{G} \to \mathcal{F}$ un homomorphisme
de $\mathcal{O}_X$-modules. Soit $\overline{x}$ un point géométrique de $X$
au-dessus de $x \in |X|$.
\begin{enumerate}
\item Si $\mathcal{F}_{\overline{x}} = 0$, il existe un voisinage
ouvert $X' \subset X$ de $x$ tel que $\mathcal{F}|_{X'} = 0$.
\item Si $\varphi_{\overline{x}} : \mathcal{G}_{\overline{x}} \to
\mathcal{F}_{\overline{x}}$ est injectif, il existe un voisinage
ouvert $X' \subset X$ de $x$ tel que $\varphi|_{X'}$ soit injectif.
\item Si $\varphi_{\overline{x}} : \mathcal{G}_{\overline{x}} \to
\mathcal{F}_{\overline{x}}$ est surjectif, il existe un voisinage
ouvert $X' \subset X$ de $x$ tel que $\varphi|_{X'}$ soit surjectif.
\item Si $\varphi_{\overline{x}} : \mathcal{G}_{\overline{x}} \to
\mathcal{F}_{\overline{x}}$ est bijectif, il existe un voisinage
ouvert $X' \subset X$ de $x$ tel que $\varphi|_{X'}$ soit un isomorphisme.
\end{enumerate}
\end{lemma}

\begin{proof}
Soit $\varphi : U \to X$ un morphisme \'etale, où $U$ est un schéma, et
soit $u \in U$ un point d'image $x$. D'après
Propriétés des espaces, Lemmes
\ref{spaces-properties-lemma-stalk-quasi-coherent} et
\ref{spaces-properties-lemma-describe-etale-local-ring}
ainsi que
Compléments d'algèbre, Lemme \ref{more-algebra-lemma-dumb-properties-henselization},
on voit que $\varphi_{\overline{x}}$ est injectif, surjectif ou bijectif
si et seulement si $\varphi_u : \varphi^*\mathcal{F}_u \to \varphi^*\mathcal{G}_u$
possède la propriété correspondante. On peut donc appliquer la version de
ce lemme pour les schémas pour voir que (quitte à restreindre $U$) le morphisme
$\varphi^*\mathcal{F} \to \varphi^*\mathcal{G}$ est injectif, surjectif
ou un isomorphisme. Soit $X' \subset X$ le sous-espace ouvert correspondant
à $|\varphi|(|U|) \subset |X|$; voir
Propriétés des espaces, Lemme \ref{spaces-properties-lemma-open-subspaces}.
Puisque $\{U \to X'\}$ est un recouvrement pour la topologie \'etale, on conclut
que $\varphi|_{X'}$ est injectif, surjectif ou un isomorphisme, ce qui donne le résultat voulu.
Enfin, remarquons que (1) résulte de (2) en considérant le morphisme
$\mathcal{F} \to 0$.
\end{proof}

\begin{lemma}
\label{spaces-cohomology-lemma-coherent-support-closed}
Soit $S$ un schéma. Soit $X$ un espace algébrique localement noethérien sur $S$.
Soit $\mathcal{F}$ un $\mathcal{O}_X$-module cohérent. Soit $i : Z \to X$
le support schématique de $\mathcal{F}$, et soit $\mathcal{G}$
le $\mathcal{O}_Z$-module quasi-cohérent tel que
$i_*\mathcal{G} = \mathcal{F}$; voir
Morphismes d'espaces, Définition
\ref{spaces-morphisms-definition-scheme-theoretic-support}.
Alors $\mathcal{G}$ est un $\mathcal{O}_Z$-module cohérent.
\end{lemma}

\begin{proof}
L'énoncé du lemme a un sens, car un module cohérent est en
particulier de type fini. De plus, comme $Z \to X$ est une immersion fermée,
ce morphisme est localement de type fini, et $Z$ est donc localement noethérien; voir
Morphismes d'espaces, Lemmes
\ref{spaces-morphisms-lemma-immersion-locally-finite-type} et
\ref{spaces-morphisms-lemma-locally-finite-type-locally-noetherian}.
Enfin, puisque $\mathcal{G}$ est de type fini, c'est un
$\mathcal{O}_Z$-module cohérent d'après le
Lemme \ref{spaces-cohomology-lemma-coherent-Noetherian}.
\end{proof}

\begin{lemma}
\label{spaces-cohomology-lemma-i-star-equivalence}
Soit $S$ un schéma. Soit $i : Z \to X$ une immersion fermée d'espaces
algébriques localement noethériens sur $S$.
Soit $\mathcal{I} \subset \mathcal{O}_X$ le faisceau quasi-cohérent d'idéaux
définissant $Z$. Le foncteur $i_*$ induit une équivalence entre la
catégorie des $\mathcal{O}_X$-modules cohérents annulés par $\mathcal{I}$
et la catégorie des $\mathcal{O}_Z$-modules cohérents.
\end{lemma}

\begin{proof}
Le foncteur est pleinement fidèle d'après
Morphismes d'espaces, Lemme \ref{spaces-morphisms-lemma-i-star-equivalence}.
Soit $\mathcal{F}$ un $\mathcal{O}_X$-module cohérent
annulé par $\mathcal{I}$. D'après
Morphismes d'espaces, Lemme \ref{spaces-morphisms-lemma-i-star-equivalence},
on peut écrire $\mathcal{F} = i_*\mathcal{G}$ pour un faisceau quasi-cohérent
$\mathcal{G}$ sur $Z$. Pour vérifier que $\mathcal{G}$ est cohérent,
on peut travailler localement pour la topologie \'etale (Lemme \ref{spaces-cohomology-lemma-coherent-Noetherian}).
En choisissant un recouvrement \'etale par un schéma, on conclut que
$\mathcal{G}$ est cohérent par le cas des schémas
(Cohomologie des schémas, Lemme \ref{coherent-lemma-i-star-equivalence}).
Ainsi le foncteur est pleinement fidèle, ce qui termine la démonstration.
\end{proof}

\begin{lemma}
\label{spaces-cohomology-lemma-finite-pushforward-coherent}
Soit $S$ un schéma. Soit $f : X \to Y$ un morphisme fini d'espaces
algébriques sur $S$, avec $Y$ localement noethérien. Soit $\mathcal{F}$ un
$\mathcal{O}_X$-module cohérent. Supposons $f$ fini et $Y$ localement
noethérien. Alors $R^pf_*\mathcal{F} = 0$ pour $p > 0$, et
$f_*\mathcal{F}$ est cohérent.
\end{lemma}

\begin{proof}
Choisissons un schéma $V$ et un morphisme \'etale surjectif $V \to Y$.
Alors $V \times_Y X \to V$ est un morphisme fini de schémas localement
noethériens. D'après (\ref{spaces-cohomology-equation-representable-higher-direct-image}), on se ramène
au cas des schémas, qui est
Cohomologie des schémas, Lemme \ref{coherent-lemma-finite-pushforward-coherent}.
\end{proof}






\section{Faisceaux cohérents sur les espaces noethériens}
\label{spaces-cohomology-section-coherent-quasi-compact}

\noindent
Dans cette section, nous donnons quelques propriétés des faisceaux cohérents sur
les espaces algébriques noethériens.

\begin{lemma}
\label{spaces-cohomology-lemma-acc-coherent}
Soit $S$ un schéma. Soit $X$ un espace algébrique noethérien sur $S$.
Soit $\mathcal{F}$ un $\mathcal{O}_X$-module cohérent.
Les sous-modules quasi-cohérents de $\mathcal{F}$ vérifient la condition
de chaîne ascendante. Autrement dit, pour toute suite
$$
\mathcal{F}_1 \subset \mathcal{F}_2 \subset \ldots \subset \mathcal{F}
$$
de sous-modules quasi-cohérents, on a
$\mathcal{F}_n = \mathcal{F}_{n + 1} = \ldots $ pour un certain $n \geq 0$.
\end{lemma}

\begin{proof}
Choisissons un schéma affine $U$ et un morphisme \'etale surjectif $U \to X$
(voir Propriétés des espaces, Lemme
\ref{spaces-properties-lemma-quasi-compact-affine-cover}).
Alors $U$ est un schéma noethérien (d'après
Morphismes d'espaces, Lemme
\ref{spaces-morphisms-lemma-locally-finite-type-locally-noetherian}).
Si $\mathcal{F}_n|_U = \mathcal{F}_{n + 1}|_U = \ldots$,
alors $\mathcal{F}_n = \mathcal{F}_{n + 1} = \ldots$.
Le résultat découle donc du cas des schémas; voir
Cohomologie des schémas, Lemme \ref{coherent-lemma-acc-coherent}.
\end{proof}

\begin{lemma}
\label{spaces-cohomology-lemma-power-ideal-kills-sheaf}
Soit $S$ un schéma. Soit $X$ un espace algébrique noethérien sur $S$.
Soit $\mathcal{F}$ un faisceau cohérent sur $X$. Soit
$\mathcal{I} \subset \mathcal{O}_X$ un faisceau quasi-cohérent d'idéaux
correspondant à un sous-espace fermé $Z \subset X$. Il existe alors un
$n \geq 0$ tel que $\mathcal{I}^n\mathcal{F} = 0$ si et seulement si
$\text{Supp}(\mathcal{F}) \subset Z$ ensemblistement.
\end{lemma}

\begin{proof}
Choisissons un schéma affine $U$ et un morphisme \'etale surjectif $U \to X$
(voir Propriétés des espaces, Lemme
\ref{spaces-properties-lemma-quasi-compact-affine-cover}).
Alors $U$ est un schéma noethérien (d'après
Morphismes d'espaces, Lemme
\ref{spaces-morphisms-lemma-locally-finite-type-locally-noetherian}).
Notons que $\mathcal{I}^n\mathcal{F}|_U = 0$ si et seulement si
$\mathcal{I}^n\mathcal{F} = 0$, et il en va de même pour la condition sur
le support. Le résultat découle donc du cas des schémas; voir
Cohomologie des schémas, Lemme \ref{coherent-lemma-power-ideal-kills-sheaf}.
\end{proof}

\begin{lemma}[Artin-Rees]
\label{spaces-cohomology-lemma-Artin-Rees}
Soit $S$ un schéma. Soit $X$ un espace algébrique noethérien sur $S$.
Soit $\mathcal{F}$ un faisceau cohérent sur $X$. Soit
$\mathcal{G} \subset \mathcal{F}$ un sous-faisceau quasi-cohérent.
Soit $\mathcal{I} \subset \mathcal{O}_X$ un faisceau quasi-cohérent
d'idéaux. Il existe alors un $c \geq 0$ tel que, pour tout $n \geq c$, on
ait
$$
\mathcal{I}^{n - c}(\mathcal{I}^c\mathcal{F} \cap \mathcal{G})
=
\mathcal{I}^n\mathcal{F} \cap \mathcal{G}
$$
\end{lemma}

\begin{proof}
Choisissons un schéma affine $U$ et un morphisme \'etale surjectif $U \to X$
(voir Propriétés des espaces, Lemme
\ref{spaces-properties-lemma-quasi-compact-affine-cover}).
Alors $U$ est un schéma noethérien (d'après
Morphismes d'espaces, Lemme
\ref{spaces-morphisms-lemma-locally-finite-type-locally-noetherian}).
L'égalité du lemme est vérifiée si et seulement si elle l'est après
restriction à $U$. Le résultat découle donc du cas des schémas; voir
Cohomologie des schémas, Lemme \ref{coherent-lemma-Artin-Rees}.
\end{proof}

\begin{lemma}
\label{spaces-cohomology-lemma-homs-over-open}
Soit $S$ un schéma. Soit $X$ un espace algébrique noethérien sur $S$.
Soit $\mathcal{F}$ un $\mathcal{O}_X$-module quasi-cohérent.
Soit $\mathcal{G}$ un $\mathcal{O}_X$-module cohérent.
Soit $\mathcal{I} \subset \mathcal{O}_X$ un faisceau quasi-cohérent
d'idéaux. Notons $Z \subset X$ le sous-espace fermé correspondant et
posons $U = X \setminus Z$. Il existe un isomorphisme canonique
$$
\colim_n \Hom_{\mathcal{O}_X}(\mathcal{I}^n\mathcal{G}, \mathcal{F})
\longrightarrow
\Hom_{\mathcal{O}_U}(\mathcal{G}|_U, \mathcal{F}|_U).
$$
En particulier, on a un isomorphisme
$$
\colim_n \Hom_{\mathcal{O}_X}(\mathcal{I}^n, \mathcal{F})
\longrightarrow
\Gamma(U, \mathcal{F}).
$$
\end{lemma}

\begin{proof}
Soit $W$ un schéma affine et soit $W \to X$ un morphisme \'etale
surjectif (voir Propriétés des espaces, Lemme
\ref{spaces-properties-lemma-quasi-compact-affine-cover}).
Posons $R = W \times_X W$. Alors $W$ et $R$ sont des schémas noethériens; voir
Morphismes d'espaces, Lemme
\ref{spaces-morphisms-lemma-locally-finite-type-locally-noetherian}.
Le résultat vaut donc pour les restrictions de $\mathcal{F}$, $\mathcal{G}$,
ainsi que de $\mathcal{I}$, $U$, $Z$, à $W$ et à $R$, d'après
Cohomologie des schémas, Lemme \ref{coherent-lemma-homs-over-open}.
Il en résulte formellement que le résultat vaut sur $X$.
\end{proof}






\section{Dévissage des faisceaux cohérents}
\label{spaces-cohomology-section-devissage}

\noindent
Cette section est l'analogue de
Cohomologie des schémas, section \ref{coherent-section-devissage}.

\begin{lemma}
\label{spaces-cohomology-lemma-prepare-filter-support}
Soit $S$ un schéma. Soit $X$ un espace algébrique noethérien sur $S$.
Soit $\mathcal{F}$ un faisceau cohérent sur $X$. Supposons que
$\text{Supp}(\mathcal{F}) = Z \cup Z'$, où $Z$, $Z'$ sont fermés.
Il existe alors une suite exacte courte de faisceaux cohérents
$$
0 \to \mathcal{G}' \to \mathcal{F} \to \mathcal{G} \to 0
$$
avec $\text{Supp}(\mathcal{G}') \subset Z'$ et
$\text{Supp}(\mathcal{G}) \subset Z$.
\end{lemma}

\begin{proof}
Soit $\mathcal{I} \subset \mathcal{O}_X$ le faisceau d'idéaux
définissant sur $Z$ la structure réduite induite de sous-espace fermé; voir
Propriétés des espaces, Lemme
\ref{spaces-properties-lemma-reduced-closed-subspace}.
Considérons les sous-faisceaux
$\mathcal{G}'_n = \mathcal{I}^n\mathcal{F}$ et les
quotients $\mathcal{G}_n = \mathcal{F}/\mathcal{I}^n\mathcal{F}$.
Pour tout $n$, on a une suite exacte courte
$$
0 \to \mathcal{G}'_n \to \mathcal{F} \to \mathcal{G}_n \to 0
$$
Pour tout point géométrique $\overline{x}$ de $Z' \setminus Z$, on a
$\mathcal{I}_{\overline{x}} = \mathcal{O}_{X, \overline{x}}$
et donc $\mathcal{G}_{n, \overline{x}} = 0$. Ainsi, on voit que
$\text{Supp}(\mathcal{G}_n) \subset Z$. Notons que $X \setminus Z'$
est un espace algébrique noethérien. D'après le
Lemme \ref{spaces-cohomology-lemma-power-ideal-kills-sheaf},
il existe donc un $n$ tel que $\mathcal{G}'_n|_{X \setminus Z'} =
\mathcal{I}^n\mathcal{F}|_{X \setminus Z'} = 0$.
Pour un tel $n$, on voit que $\text{Supp}(\mathcal{G}'_n) \subset Z'$.
Il suffit donc de poser $\mathcal{G}' = \mathcal{G}'_n$ et $\mathcal{G} = \mathcal{G}_n$
pour conclure.
\end{proof}

\noindent
Dans la suite, nous emploierons librement le support schématique des
modules de type fini défini dans Morphismes d'espaces, Définition
\ref{spaces-morphisms-definition-scheme-theoretic-support}.

\begin{lemma}
\label{spaces-cohomology-lemma-prepare-filter-irreducible}
Soit $S$ un schéma. Soit $X$ un espace algébrique noethérien sur $S$.
Soit $\mathcal{F}$ un faisceau cohérent sur $X$. Supposons que le support
schématique de $\mathcal{F}$ soit un sous-espace réduit $Z \subset X$ tel que
$|Z|$ soit irréductible. Il existe alors un entier $r > 0$, un
faisceau non nul d'idéaux $\mathcal{I} \subset \mathcal{O}_Z$ et un
morphisme injectif de faisceaux cohérents
$$
i_*\left(\mathcal{I}^{\oplus r}\right) \to \mathcal{F}
$$
dont le conoyau est à support dans un sous-espace fermé strict de $Z$.
\end{lemma}

\begin{proof}
Par hypothèse, il existe un $\mathcal{O}_Z$-module cohérent
$\mathcal{G}$ de support $Z$ tel que $\mathcal{F} \cong i_*\mathcal{G}$; voir
Lemme \ref{spaces-cohomology-lemma-coherent-support-closed}. Il suffit donc de démontrer le
lemme dans le cas $Z = X$ et $i = \text{id}$.

\medskip\noindent
Par Propriétés des espaces, Proposition
\ref{spaces-properties-proposition-locally-quasi-separated-open-dense-scheme}
il existe un sous-espace ouvert dense $U \subset X$ qui est un schéma. Notons que
$U$ est un schéma intègre noethérien. Quitte à restreindre $U$, on peut supposer
que $\mathcal{F}|_U \cong \mathcal{O}_U^{\oplus r}$ (par exemple d'après
Cohomologie des schémas, Lemme \ref{coherent-lemma-prepare-filter-irreducible},
ou par un argument direct d'algèbre). Soit $\mathcal{I} \subset \mathcal{O}_X$
un faisceau quasi-cohérent d'idéaux dont le sous-espace fermé associé
est le complémentaire de $U$ dans $X$ (voir, par exemple,
Propriétés des espaces, section \ref{spaces-properties-section-reduced}). 
D'après le Lemme \ref{spaces-cohomology-lemma-homs-over-open}, il existe un $n \geq 0$ et un
morphisme $\mathcal{I}^n(\mathcal{O}_X^{\oplus r}) \to \mathcal{F}$
qui redonne notre isomorphisme sur $U$. Puisque
$\mathcal{I}^n(\mathcal{O}_X^{\oplus r}) = (\mathcal{I}^n)^{\oplus r}$,
on obtient un morphisme comme dans le lemme. Il est injectif : en effet, si $\sigma$ est
une section non nulle de $\mathcal{I}^{\oplus r}$ sur un schéma $W$ \'etale
sur $X$, alors, puisque $X$, et donc $W$, est réduit, le support de $\sigma$
contient un ouvert non vide de $W$. Or le noyau de
$(\mathcal{I}^n)^{\oplus r} \to \mathcal{F}$ est nul
sur un ouvert dense; $\sigma$ ne peut donc pas être une section du noyau.
\end{proof}

\begin{lemma}
\label{spaces-cohomology-lemma-coherent-filter}
Soit $S$ un schéma. Soit $X$ un espace algébrique noethérien sur $S$.
Soit $\mathcal{F}$ un faisceau cohérent sur $X$. Il existe une filtration
$$
0 = \mathcal{F}_0 \subset \mathcal{F}_1 \subset
\ldots \subset \mathcal{F}_m = \mathcal{F}
$$
par des sous-faisceaux cohérents telle que, pour chaque $j = 1, \ldots, m$,
il existe un sous-espace fermé réduit $Z_j \subset X$ tel que $|Z_j|$
soit irréductible, et un faisceau d'idéaux $\mathcal{I}_j \subset \mathcal{O}_{Z_j}$
tels que
$$
\mathcal{F}_j/\mathcal{F}_{j - 1}
\cong (Z_j \to X)_* \mathcal{I}_j
$$
\end{lemma}

\begin{proof}
Considérons l'ensemble
$$
\mathcal{T} =
\left\{
\begin{matrix}
T \subset |X|
\text{ fermé tel qu'il existe un faisceau cohérent }
\mathcal{F} \\
\text{ avec }
\text{Supp}(\mathcal{F}) = T
\text{ pour lequel le lemme est faux}
\end{matrix}
\right\}
$$
Nous voulons montrer que $\mathcal{T}$ est vide. Sinon,
puisque $|X|$ est noethérien (Propriétés des espaces, Lemme
\ref{spaces-properties-lemma-Noetherian-topology})
on peut choisir un élément minimal $T \in \mathcal{T}$. Cela signifie qu'il
existe un faisceau cohérent $\mathcal{F}$ sur $X$, de support $T$,
pour lequel le lemme n'est pas vrai. Il est clair que $T \not = \emptyset$, car
le seul faisceau de support vide est le faisceau nul, pour lequel le
lemme est vrai (avec $m = 0$).

\medskip\noindent
Si $T$ n'est pas irréductible, on peut écrire $T = Z_1 \cup Z_2$,
où $Z_1, Z_2$ sont fermés et strictement contenus dans $T$.
On peut alors appliquer le Lemme \ref{spaces-cohomology-lemma-prepare-filter-support}
pour obtenir une suite exacte courte de faisceaux cohérents
$$
0 \to
\mathcal{G}_1 \to
\mathcal{F} \to
\mathcal{G}_2 \to 0
$$
avec $\text{Supp}(\mathcal{G}_i) \subset Z_i$. Par minimalité de
$T$, chacun des $\mathcal{G}_i$ admet une filtration comme dans l'énoncé
du lemme. En considérant la filtration induite sur $\mathcal{F}$,
on aboutit à une contradiction. On en conclut
que $T$ est irréductible.

\medskip\noindent
Supposons $T$ irréductible. Soit $\mathcal{J}$ le faisceau d'idéaux
définissant sur $T$ la structure réduite induite de sous-espace fermé;
voir Propriétés des espaces, Lemme
\ref{spaces-properties-lemma-reduced-closed-subspace}.
D'après le Lemme \ref{spaces-cohomology-lemma-power-ideal-kills-sheaf}, il existe
un $n \geq 0$ tel que $\mathcal{J}^n\mathcal{F} = 0$. On obtient donc
une filtration
$$
0 = \mathcal{I}^n\mathcal{F} \subset \mathcal{I}^{n - 1}\mathcal{F}
\subset \ldots \subset \mathcal{I}\mathcal{F} \subset \mathcal{F}
$$
dont chacun des sous-quotients successifs est annulé par $\mathcal{J}$.
Si chacun de ces sous-quotients admet une filtration comme dans l'énoncé
du lemme, il en va de même pour $\mathcal{F}$. Autrement dit, on peut
supposer que $\mathcal{J}$ annule $\mathcal{F}$.

\medskip\noindent
Supposons $T$ irréductible et $\mathcal{J}\mathcal{F} = 0$, où
$\mathcal{J}$ est comme ci-dessus. Alors le support schématique de
$\mathcal{F}$ est $T$; voir
Morphismes d'espaces, Lemme \ref{spaces-morphisms-lemma-i-star-equivalence}.
On peut donc appliquer le Lemme \ref{spaces-cohomology-lemma-prepare-filter-irreducible}.
On obtient une suite exacte courte
$$
0 \to
i_*(\mathcal{I}^{\oplus r}) \to
\mathcal{F} \to
\mathcal{Q} \to 0
$$
où le support de $\mathcal{Q}$ est une partie fermée stricte de $T$.
On voit donc que $\mathcal{Q}$ admet une filtration du type voulu
par minimalité de $T$. Mais il en va alors manifestement de même pour $\mathcal{F}$,
d'où la contradiction finale.
\end{proof}

\begin{lemma}
\label{spaces-cohomology-lemma-property-initial}
Soit $S$ un schéma. Soit $X$ un espace algébrique noethérien sur $S$.
Soit $\mathcal{P}$ une propriété des faisceaux cohérents sur $X$. Supposons que
\begin{enumerate}
\item pour toute suite exacte courte de faisceaux cohérents
$$
0 \to \mathcal{F}_1 \to \mathcal{F} \to \mathcal{F}_2 \to 0
$$
si les $\mathcal{F}_i$, $i = 1, 2$, vérifient la propriété $\mathcal{P}$,
il en va de même pour $\mathcal{F}$;
\item pour tout sous-espace fermé réduit $Z \subset X$ tel que $|Z|$ soit irréductible
et tout faisceau quasi-cohérent d'idéaux $\mathcal{I} \subset \mathcal{O}_Z$,
la propriété $\mathcal{P}$ est vérifiée par $i_*\mathcal{I}$.
\end{enumerate}
Alors la propriété $\mathcal{P}$ est vérifiée par tout faisceau cohérent sur $X$.
\end{lemma}

\begin{proof}
Notons d'abord que, si $\mathcal{F}$ est un faisceau cohérent muni d'une filtration
$$
0 = \mathcal{F}_0 \subset \mathcal{F}_1 \subset
\ldots \subset \mathcal{F}_m = \mathcal{F}
$$
par des sous-faisceaux cohérents telle que chaque $\mathcal{F}_i/\mathcal{F}_{i - 1}$
vérifie la propriété $\mathcal{P}$, il en va de même pour $\mathcal{F}$.
Cela résulte de la propriété (1) de $\mathcal{P}$.
D'autre part, d'après le Lemme \ref{spaces-cohomology-lemma-coherent-filter},
on peut munir tout $\mathcal{F}$ d'une filtration
dont les sous-quotients successifs sont comme dans (2).
Le lemme en résulte.
\end{proof}

\noindent
Voici une variante plus utile du lemme précédent.

\begin{lemma}
\label{spaces-cohomology-lemma-property-higher-rank-cohomological}
Soit $S$ un schéma. Soit $X$ un espace algébrique noethérien sur $S$.
Soit $\mathcal{P}$ une propriété des faisceaux cohérents sur $X$. Supposons que
\begin{enumerate}
\item pour toute suite exacte courte de faisceaux cohérents
$$
0 \to \mathcal{F}_1 \to \mathcal{F} \to \mathcal{F}_2 \to 0
$$
si les $\mathcal{F}_i$, $i = 1, 2$, vérifient la propriété $\mathcal{P}$,
il en va de même pour $\mathcal{F}$;
\item si $\mathcal{P}$ est vérifiée par $\mathcal{F}^{\oplus r}$ pour
un certain $r \geq 1$, elle est vérifiée par $\mathcal{F}$;
\item pour tout sous-espace fermé réduit $i : Z \to X$ tel que
$|Z|$ soit irréductible, il existe un faisceau cohérent $\mathcal{G}$ sur $Z$
tel que
\begin{enumerate}
\item $\text{Supp}(\mathcal{G}) = Z$,
\item pour tout faisceau quasi-cohérent non nul d'idéaux
$\mathcal{I} \subset \mathcal{O}_Z$, il existe un sous-faisceau quasi-cohérent
$\mathcal{G}' \subset \mathcal{I}\mathcal{G}$ tel que
$\text{Supp}(\mathcal{G}/\mathcal{G}')$ soit une partie fermée stricte de $|Z|$
et que la propriété $\mathcal{P}$ soit vérifiée par $i_*\mathcal{G}'$.
\end{enumerate}
\end{enumerate}
Alors la propriété $\mathcal{P}$ est vérifiée par tout faisceau cohérent sur $X$.
\end{lemma}

\begin{proof}
Considérons l'ensemble
$$
\mathcal{T} =
\left\{
\begin{matrix}
T \subset |X|
\text{ fermé non vide tel qu'il existe un faisceau cohérent } \\
\mathcal{F}
\text{ avec }
\text{Supp}(\mathcal{F}) = T
\text{ pour lequel le lemme est faux}
\end{matrix}
\right\}
$$
Nous voulons montrer que $\mathcal{T}$ est vide. Sinon,
puisque $|X|$ est noethérien (Propriétés des espaces, Lemme
\ref{spaces-properties-lemma-Noetherian-topology})
on peut choisir un élément minimal $T \in \mathcal{T}$. Cela signifie qu'il
existe un faisceau cohérent $\mathcal{F}$ sur $X$, de support $T$,
pour lequel le lemme n'est pas vrai.

\medskip\noindent
Si $T$ n'est pas irréductible, on peut écrire $T = Z_1 \cup Z_2$,
où $Z_1, Z_2$ sont fermés et strictement contenus dans $T$.
On peut alors appliquer le Lemme \ref{spaces-cohomology-lemma-prepare-filter-support}
pour obtenir une suite exacte courte de faisceaux cohérents
$$
0 \to
\mathcal{G}_1 \to
\mathcal{F} \to
\mathcal{G}_2 \to 0
$$
avec $\text{Supp}(\mathcal{G}_i) \subset Z_i$. Par minimalité de
$T$, chacun des $\mathcal{G}_i$ vérifie $\mathcal{P}$. Ainsi $\mathcal{F}$
vérifie la propriété $\mathcal{P}$ d'après (1), ce qui est une contradiction.

\medskip\noindent
Supposons $T$ irréductible. Soit $\mathcal{J}$ le faisceau d'idéaux
définissant sur $T$ la structure réduite induite de sous-espace fermé;
voir Propriétés des espaces, Lemme
\ref{spaces-properties-lemma-reduced-closed-subspace}.
D'après le Lemme \ref{spaces-cohomology-lemma-power-ideal-kills-sheaf}, il existe
un $n \geq 0$ tel que $\mathcal{J}^n\mathcal{F} = 0$. On obtient donc
une filtration
$$
0 = \mathcal{J}^n\mathcal{F} \subset \mathcal{J}^{n - 1}\mathcal{F}
\subset \ldots \subset \mathcal{J}\mathcal{F} \subset \mathcal{F}
$$
dont chacun des sous-quotients successifs est annulé par $\mathcal{J}$.
Si chacun de ces sous-quotients admet une filtration comme dans l'énoncé
du lemme, il en va de même pour $\mathcal{F}$ d'après (1). Autrement dit, on peut
supposer que $\mathcal{J}$ annule $\mathcal{F}$.

\medskip\noindent
Supposons $T$ irréductible et $\mathcal{J}\mathcal{F} = 0$, où
$\mathcal{J}$ est comme ci-dessus. Notons $i : Z \to X$ le sous-espace fermé
correspondant à $\mathcal{J}$. Alors $\mathcal{F} = i_*\mathcal{H}$
pour un certain $\mathcal{O}_Z$-module cohérent $\mathcal{H}$; voir
Morphismes d'espaces, Lemme \ref{spaces-morphisms-lemma-i-star-equivalence},
et Lemme \ref{spaces-cohomology-lemma-coherent-support-closed}.
Soit $\mathcal{G}$ le faisceau cohérent sur $Z$ vérifiant
(3)(a) et (3)(b). Appliquons le Lemme \ref{spaces-cohomology-lemma-prepare-filter-irreducible}
pour obtenir des morphismes injectifs
$$
\mathcal{I}_1^{\oplus r_1} \to \mathcal{H}
\quad\text{et}\quad
\mathcal{I}_2^{\oplus r_2} \to \mathcal{G}
$$
dont les conoyaux sont à support dans des parties fermées strictes de $Z$. On trouve donc
un ouvert non vide $V \subset Z$ tel que
$$
\mathcal{H}^{\oplus r_2}_V \cong \mathcal{G}^{\oplus r_1}_V
$$
Soit $\mathcal{I} \subset \mathcal{O}_Z$ un faisceau quasi-cohérent d'idéaux
définissant $Z \setminus V$; d'après
le Lemme \ref{spaces-cohomology-lemma-homs-over-open}, on obtient
un morphisme
$$
\mathcal{I}^n\mathcal{G}^{\oplus r_1} \longrightarrow \mathcal{H}^{\oplus r_2}
$$
qui est un isomorphisme sur $V$. Le noyau est à support dans $Z \setminus V$,
donc est annulé par une puissance de $\mathcal{I}$; voir
Lemme \ref{spaces-cohomology-lemma-power-ideal-kills-sheaf}. Ainsi, quitte à augmenter
$n$, on peut supposer que le morphisme affiché est injectif; voir
Lemme \ref{spaces-cohomology-lemma-Artin-Rees}. En appliquant (3)(b), on trouve
$\mathcal{G}' \subset \mathcal{I}^n\mathcal{G}$ tel que
$$
(i_*\mathcal{G}')^{\oplus r_1} \longrightarrow
i_*\mathcal{H}^{\oplus r_2} = \mathcal{F}^{\oplus r_2}
$$
soit injectif, avec un conoyau à support dans une partie fermée stricte de $Z$,
et que la propriété $\mathcal{P}$ soit vérifiée par $i_*\mathcal{G}'$.
D'après (1), la propriété $\mathcal{P}$ est vérifiée par $(i_*\mathcal{G}')^{\oplus r_1}$.
D'après (1) et la minimalité de $T = |Z|$, la propriété $\mathcal{P}$ est vérifiée par
$\mathcal{F}^{\oplus r_2}$. Enfin, d'après (2), la propriété $\mathcal{P}$
est vérifiée par $\mathcal{F}$, ce qui est la contradiction cherchée.
\end{proof}

\begin{lemma}
\label{spaces-cohomology-lemma-property-higher-rank-cohomological-variant}
Soit $S$ un schéma. Soit $X$ un espace algébrique noethérien sur $S$.
Soit $\mathcal{P}$ une propriété des faisceaux cohérents sur $X$. Supposons que
\begin{enumerate}
\item pour toute suite exacte courte de faisceaux cohérents sur $X$,
si deux des trois termes vérifient la propriété $\mathcal{P}$, le troisième la vérifie aussi;
\item si $\mathcal{P}$ est vérifiée par $\mathcal{F}^{\oplus r}$ pour
un certain $r \geq 1$, elle est vérifiée par $\mathcal{F}$;
\item pour tout sous-espace fermé réduit $i : Z \to X$ tel que
$|Z|$ soit irréductible, il existe un faisceau cohérent $\mathcal{G}$ sur $X$
dont le support schématique est $Z$ et tel que $\mathcal{P}$ soit vérifiée par
$\mathcal{G}$.
\end{enumerate}
Alors la propriété $\mathcal{P}$ est vérifiée par tout faisceau cohérent sur $X$.
\end{lemma}

\begin{proof}
Nous allons montrer que les conditions (1) et (2) du
Lemme \ref{spaces-cohomology-lemma-property-initial} sont vérifiées. C'est clair pour la condition (1).
Pour établir (2), posons
$$
\mathcal{T} =
\left\{
\begin{matrix}
i : Z \to X \text{ sous-espace fermé réduit tel que }|Z|\text{ soit irréductible et}\\
\text{ que }i_*\mathcal{I}\text{ ne vérifie pas }\mathcal{P}
\text{ pour un certain faisceau quasi-cohérent }\mathcal{I} \subset \mathcal{O}_Z
\end{matrix}
\right\}
$$
Si $\mathcal{T}$ n'est pas vide, comme $X$ est noethérien, on peut
trouver un $i : Z \to X$ minimal dans $\mathcal{T}$. Nous allons montrer
que cela conduit à une contradiction.

\medskip\noindent
Soit $\mathcal{G}$ le faisceau de support schématique $Z$ dont
l'existence est postulée dans l'hypothèse (3). Soit
$\varphi : i_*\mathcal{I}^{\oplus r} \to \mathcal{G}$ comme dans le
Lemme \ref{spaces-cohomology-lemma-prepare-filter-irreducible}. Soit
$$
0 = \mathcal{F}_0 \subset \mathcal{F}_1 \subset
\ldots \subset \mathcal{F}_m = \Coker(\varphi)
$$
une filtration comme dans le Lemme \ref{spaces-cohomology-lemma-coherent-filter}. Par minimalité
de $Z$ et d'après l'hypothèse (1), on voit que $\Coker(\varphi)$ vérifie
la propriété $\mathcal{P}$. Puisque $\varphi$ est injectif, l'hypothèse (1)
montre encore que $i_*\mathcal{I}^{\oplus r}$ vérifie la propriété
$\mathcal{P}$. L'hypothèse (2) montre alors que $i_*\mathcal{I}$
vérifie la propriété $\mathcal{P}$.

\medskip\noindent
Enfin, si $\mathcal{J} \subset \mathcal{O}_Z$ est un second faisceau quasi-cohérent
d'idéaux, posons $\mathcal{K} = \mathcal{I} \cap \mathcal{J}$
et considérons les suites exactes courtes
$$
0 \to \mathcal{K} \to \mathcal{I} \to \mathcal{I}/\mathcal{K} \to 0
\quad
\text{et}
\quad
0 \to \mathcal{K} \to \mathcal{J} \to \mathcal{J}/\mathcal{K} \to 0
$$
En raisonnant comme ci-dessus et en utilisant la minimalité de $Z$, on voit que
$i_*\mathcal{I}/\mathcal{K}$ et $i_*\mathcal{J}/\mathcal{K}$
vérifient $\mathcal{P}$. L'hypothèse (1) montre donc que
$i_*\mathcal{K}$, puis $i_*\mathcal{J}$, vérifient $\mathcal{P}$.
Autrement dit, $Z$ n'est pas un élément de $\mathcal{T}$, ce qui est
la contradiction cherchée.
\end{proof}






\section{Limites de modules cohérents}
\label{spaces-cohomology-section-limits}

\noindent
Une limite inductive de modules cohérents (sur un espace algébrique
localement noethérien) n'est en général pas cohérente. Elle est toutefois quasi-cohérente,
car toute limite inductive de modules quasi-cohérents sur un espace algébrique est
quasi-cohérente, voir Propriétés des espaces, Lemme
\ref{spaces-properties-lemma-properties-quasi-coherent}.
Réciproquement, si l'espace algébrique est noethérien, tout module quasi-cohérent
est une limite inductive filtrante de modules cohérents.

\begin{lemma}
\label{spaces-cohomology-lemma-directed-colimit-coherent}
Soit $S$ un schéma. Soit $X$ un espace algébrique noethérien sur $S$.
Tout $\mathcal{O}_X$-module quasi-cohérent est la limite inductive filtrante
de ses sous-modules cohérents.
\end{lemma}

\begin{proof}
Soit $\mathcal{F}$ un $\mathcal{O}_X$-module quasi-cohérent.
Si $\mathcal{G}, \mathcal{H} \subset \mathcal{F}$ sont des
sous-$\mathcal{O}_X$-modules cohérents, alors l'image de
$\mathcal{G} \oplus \mathcal{H} \to \mathcal{F}$ est encore un
sous-$\mathcal{O}_X$-module cohérent qui les contient tous deux
(voir les Lemmes \ref{spaces-cohomology-lemma-coherent-abelian-Noetherian} et
\ref{spaces-cohomology-lemma-coherent-Noetherian-quasi-coherent-sub-quotient}).
On voit ainsi que le système est filtrant.
Il suffit donc de montrer que $\mathcal{F}$ peut s'écrire comme
une limite inductive filtrante de modules cohérents : en prenant alors les
images de ces modules dans $\mathcal{F}$, on conclut qu'il y en a
suffisamment.

\medskip\noindent
Soit $U$ un schéma affine et $U \to X$ un morphisme \'etale surjectif.
Posons $R = U \times_X U$, de sorte que $X = U/R$ comme d'habitude. D'après
Propriétés des espaces, Proposition
\ref{spaces-properties-proposition-quasi-coherent},
on a $\QCoh(\mathcal{O}_X) = \QCoh(U, R, s, t, c)$.
On se ramène donc à établir l'assertion correspondante pour
$\QCoh(U, R, s, t, c)$. Le résultat découle alors du
Lemme plus général \ref{groupoids-lemma-colimit-coherent} du chapitre Groupoïdes.
\end{proof}

\begin{lemma}
\label{spaces-cohomology-lemma-direct-colimit-finite-presentation}
Soit $S$ un schéma. Soit $f : X \to Y$ un morphisme affine d'espaces
algébriques sur $S$, avec $Y$ noethérien. Alors tout
$\mathcal{O}_X$-module quasi-cohérent est une limite inductive filtrante de
$\mathcal{O}_X$-modules de présentation finie.
\end{lemma}

\begin{proof}
Soit $\mathcal{F}$ un $\mathcal{O}_X$-module quasi-cohérent.
Écrivons $f_*\mathcal{F} = \colim \mathcal{H}_i$, où $\mathcal{H}_i$ est
un $\mathcal{O}_Y$-module cohérent, voir le
Lemme \ref{spaces-cohomology-lemma-directed-colimit-coherent}.
D'après le Lemme \ref{spaces-cohomology-lemma-coherent-Noetherian}, les modules $\mathcal{H}_i$
sont des $\mathcal{O}_Y$-modules de présentation finie. Par conséquent,
$f^*\mathcal{H}_i$ est un $\mathcal{O}_X$-module de présentation finie, voir
Propriétés des espaces, section
\ref{spaces-properties-section-properties-modules}.
Nous affirmons que le morphisme
$$
\colim f^*\mathcal{H}_i = f^*f_*\mathcal{F} \to \mathcal{F}
$$
est surjectif puisque $f$ est supposé affine. En effet, choisissons
un schéma $V$ et un morphisme \'etale surjectif $V \to Y$. Posons
$U = X \times_Y V$. Alors $U$ est un schéma, $f' : U \to V$ est affine, et
$U \to X$ est \'etale surjectif. D'après
Propriétés des espaces, Lemme
\ref{spaces-properties-lemma-pushforward-etale-base-change-modules},
on a $f'_*(\mathcal{F}|_U) = f_*\mathcal{F}|_V$, et de même
pour les images réciproques. Ainsi, la restriction de $f^*f_*\mathcal{F} \to \mathcal{F}$
à $U$ est le morphisme
$$
f^*f_*\mathcal{F}|_U = (f')^*(f_*\mathcal{F})|_V) =
(f')^*f'_*(\mathcal{F}|_U) \to \mathcal{F}|_U
$$
qui est surjectif puisque $f'$ est un morphisme affine de schémas.
L'affirmation est donc démontrée.

\medskip\noindent
On conclut que tout module quasi-cohérent sur $X$ est quotient d'une
limite inductive filtrante de modules de présentation finie. En particulier,
$\mathcal{F}$ est le conoyau d'un morphisme
$$
\colim_{j \in J} \mathcal{G}_j \longrightarrow \colim_{i \in I} \mathcal{H}_i
$$
où les $\mathcal{G}_j$ et les $\mathcal{H}_i$ sont de présentation finie. Notons
que, pour tout $j \in I$, il existe $i \in I$ et un morphisme
$\alpha : \mathcal{G}_j \to \mathcal{H}_i$ tels que le diagramme
$$
\xymatrix{
\mathcal{G}_j \ar[r]_\alpha \ar[d] & \mathcal{H}_i \ar[d] \\
\colim_{j \in J} \mathcal{G}_j \ar[r] &
\colim_{i \in I} \mathcal{H}_i
}
$$
soit commutatif, voir le
Lemme \ref{spaces-cohomology-lemma-finite-presentation-quasi-compact-colimit}.
Dans cette situation, $\Coker(\alpha)$ est un
$\mathcal{O}_X$-module de présentation finie muni d'un morphisme
$\Coker(\alpha) \to \mathcal{F}$. Considérons l'ensemble $K$ des
triplets $(i, j, \alpha)$ comme ci-dessus. On pose
$(i, j, \alpha) \leq (i', j', \alpha')$ si et seulement si
$i \leq i'$, $j \leq j'$, et si le diagramme
$$
\xymatrix{
\mathcal{G}_j \ar[r]_\alpha \ar[d] & \mathcal{H}_i \ar[d] \\
\mathcal{G}_{j'} \ar[r]^{\alpha'} &
\mathcal{H}_{i'}
}
$$
est commutatif. Il résulte de ce qui précède que $K$ est
un ensemble ordonné filtrant,
$$
\mathcal{F} = \colim_{(i, j, \alpha) \in K} \Coker(\alpha),
$$
ce qui conclut la démonstration.
\end{proof}






\section{Annulation de la cohomologie}
\label{spaces-cohomology-section-vanishing}

\noindent
Dans cette section, nous montrons qu'un espace algébrique quasi-compact et quasi-séparé
est affine si sa cohomologie supérieure s'annule
pour tout faisceau quasi-cohérent. Nous procédons par une suite de lemmes
qui deviendront tous superflus une fois démontrée la
Proposition \ref{spaces-cohomology-proposition-vanishing-affine}.

\begin{situation}
\label{spaces-cohomology-situation-vanishing}
Ici, $S$ est un schéma et $X$ est un espace algébrique quasi-compact et quasi-séparé
sur $S$ ayant la propriété suivante : pour tout
$\mathcal{O}_X$-module quasi-cohérent $\mathcal{F}$, on a
$H^1(X, \mathcal{F}) = 0$. Posons $A = \Gamma(X, \mathcal{O}_X)$.
\end{situation}

\noindent
Nous voudrions montrer que le morphisme canonique
$$
p : X \longrightarrow \Spec(A)
$$
(voir Propriétés des espaces, Lemme
\ref{spaces-properties-lemma-morphism-to-affine-scheme}) est un isomorphisme.
Si $M$ est un $A$-module, nous notons $M \otimes_A \mathcal{O}_X$
le module quasi-cohérent $p^*\tilde M$.

\begin{lemma}
\label{spaces-cohomology-lemma-vanishing-compute}
Dans la Situation \ref{spaces-cohomology-situation-vanishing}, pour tout $A$-module $M$, on a
$p_*(M \otimes_A \mathcal{O}_X) = \widetilde{M}$ et
$\Gamma(X, M \otimes_A \mathcal{O}_X) = M$.
\end{lemma}

\begin{proof}
L'égalité $p_*(M \otimes_A \mathcal{O}_X) = \widetilde{M}$ résulte
de l'égalité $\Gamma(X, M \otimes_A \mathcal{O}_X) = M$, car
$p_*(M \otimes_A \mathcal{O}_X)$ est un module quasi-cohérent sur
$\Spec(A)$ d'après Morphismes d'espaces, Lemme
\ref{spaces-morphisms-lemma-pushforward}.
Observons que $\Gamma(X, \bigoplus_{i \in I} \mathcal{O}_X) =
\bigoplus_{i \in I} A$ d'après le Lemme \ref{spaces-cohomology-lemma-colimits}. Le
lemme est donc vrai pour les modules libres. Choisissons une suite exacte courte
$F_1 \to F_0 \to M$, où $F_0, F_1$ sont des $A$-modules libres. Puisque
$H^1(X, -)$ est nul, le foncteur des sections globales est exact à droite.
De plus, l'image réciproque $p^*$ est elle aussi exacte à droite. On voit donc
que
$$
\Gamma(X, F_1 \otimes_A \mathcal{O}_X) \to
\Gamma(X, F_0 \otimes_A \mathcal{O}_X) \to
\Gamma(X, M \otimes_A \mathcal{O}_X) \to 0
$$
est exacte. Le résultat en découle.
\end{proof}

\noindent
Le lemme suivant montre que la Situation \ref{spaces-cohomology-situation-vanishing}
est préservée par le changement de base de $X \to \Spec(A)$ par $\Spec(A') \to \Spec(A)$.

\begin{lemma}
\label{spaces-cohomology-lemma-vanishing-base-change}
Dans la Situation \ref{spaces-cohomology-situation-vanishing}.
\begin{enumerate}
\item Pour tout morphisme affine $X' \to X$ d'espaces algébriques, on a
$H^1(X', \mathcal{F}') = 0$ pour tout
$\mathcal{O}_{X'}$-module quasi-cohérent $\mathcal{F}'$.
\item Pour toute $A$-algèbre $A'$, en posant $X' = X \times_{\Spec(A)} \Spec(A')$,
le morphisme $X' \to X$ est affine et $\Gamma(X', \mathcal{O}_{X'}) = A'$.
\end{enumerate}
\end{lemma}

\begin{proof}
L'assertion (1) résulte du Lemme \ref{spaces-cohomology-lemma-affine-vanishing-higher-direct-images}
et de la suite spectrale de Leray (Cohomologie sur les sites, Lemme
\ref{sites-cohomology-lemma-Leray}). Soit $A \to A'$ comme en (2).
Alors $X' \to X$ est affine, car les morphismes affines sont stables par
changement de base (Morphismes d'espaces, Lemme
\ref{spaces-morphisms-lemma-base-change-affine}) et parce que
tout morphisme de schémas affines est affine. L'égalité
$\Gamma(X', \mathcal{O}_{X'}) = A'$ résulte de
$(X' \to X)_*\mathcal{O}_{X'} = A' \otimes_A \mathcal{O}_X$
d'après le Lemme \ref{spaces-cohomology-lemma-affine-base-change}, puis de
$$
\Gamma(X', \mathcal{O}_{X'}) =
\Gamma(X, (X' \to X)_*\mathcal{O}_{X'}) =
\Gamma(X, A' \otimes_A \mathcal{O}_X) = A'
$$
d'après le Lemme \ref{spaces-cohomology-lemma-vanishing-compute}.
\end{proof}

\begin{lemma}
\label{spaces-cohomology-lemma-vanishing-separate-closed}
Dans la Situation \ref{spaces-cohomology-situation-vanishing}, soient $Z_0, Z_1 \subset |X|$
deux parties fermées disjointes. Il existe alors $a \in A$ tel que
$Z_0 \subset V(a)$ et $Z_1 \subset V(a - 1)$.
\end{lemma}

\begin{proof}
On peut, et l'on va, munir $Z_0$, $Z_1$ de la structure d'espace réduit induite
(Propriétés des espaces, Définition
\ref{spaces-properties-definition-reduced-induced-space}); notons
$i_0 : Z_0 \to X$ et $i_1 : Z_1 \to X$ les immersions fermées correspondantes.
Comme $Z_0 \cap Z_1 = \emptyset$, le morphisme canonique de
$\mathcal{O}_X$-modules quasi-cohérents
$$
\mathcal{O}_X
\longrightarrow
i_{0, *}\mathcal{O}_{Z_0} \oplus i_{1, *}\mathcal{O}_{Z_1}
$$
est surjectif (il suffit de regarder les germes aux points géométriques). Puisque $H^1(X, -)$
est nul sur le noyau de ce morphisme, le morphisme induit sur les sections globales est
surjectif. On peut donc trouver $a \in A$ dont l'image est la section globale
$(0, 1)$ du membre de droite.
\end{proof}

\begin{lemma}
\label{spaces-cohomology-lemma-vanishing-injective}
Dans la Situation \ref{spaces-cohomology-situation-vanishing}, le morphisme $p : X \to \Spec(A)$ est
universellement injectif.
\end{lemma}

\begin{proof}
Soit $A \to k$ un homomorphisme d'anneaux, où $k$ est un corps. Il suffit de
montrer que $\Spec(k) \times_{\Spec(A)} X$ a au plus un point (voir
Morphismes d'espaces, Lemme
\ref{spaces-morphisms-lemma-universally-injective-local}).
En utilisant le Lemme \ref{spaces-cohomology-lemma-vanishing-base-change}, on peut supposer que $A$
est un corps, et il faut montrer que $|X|$ a au plus un point.

\medskip\noindent
Considérons $X$ comme un espace algébrique sur $\Spec(k)$ et utilisons
la notation $X(K)$ pour désigner les $K$-points de $X$
pour toute extension $K/k$, voir
Morphismes d'espaces, section \ref{spaces-morphisms-section-points-fields}.
Si $K/k$ est une extension algébriquement close
de degré de transcendance assez grand, alors $X(K) \to |X|$
est surjective, voir Morphismes d'espaces, Lemme
\ref{spaces-morphisms-lemma-large-enough}. Après avoir remplacé $k$
par $K$, il suffit donc de prouver que $X(k)$ est réduit à un seul élément
(dans le cas $A = k)$.

\medskip\noindent
Soient $x, x' \in X(k)$. D'après Espaces décents, Lemme
\ref{decent-spaces-lemma-algebraic-residue-field-extension-closed-point},
$x$ et $x'$ sont des points fermés de $|X|$. Ainsi $x$ et $x'$
auraient des images distinctes dans $\Spec(k)$ si $x \not = x'$, d'après le
Lemme \ref{spaces-cohomology-lemma-vanishing-separate-closed}. On conclut que
$x = x'$, comme voulu.
\end{proof}

\begin{lemma}
\label{spaces-cohomology-lemma-vanishing-separated}
Dans la Situation \ref{spaces-cohomology-situation-vanishing}, le morphisme $p : X \to \Spec(A)$ est
séparé.
\end{lemma}

\begin{proof}
D'après Espaces décents, Lemme
\ref{decent-spaces-lemma-there-is-a-scheme-integral-over},
il existe un schéma $Y$ et un morphisme entier surjectif
$Y \to X$. Comme un morphisme entier est affine, on peut appliquer le
Lemme \ref{spaces-cohomology-lemma-vanishing-base-change}
pour voir que $H^1(Y, \mathcal{G}) = 0$ pour tout
$\mathcal{O}_Y$-module quasi-cohérent $\mathcal{G}$.
Puisque $Y \to X$ est quasi-compact et que $X$ est quasi-compact,
$Y$ est quasi-compact.
Comme $Y$ est un schéma, on peut appliquer
Cohomologie des schémas, Lemme
\ref{coherent-lemma-quasi-compact-h1-zero-covering}
pour voir que $Y$ est affine. Ainsi $Y$ est séparé.
Notons qu'un morphisme entier est affine et universellement fermé, voir
Morphismes d'espaces, Lemme
\ref{spaces-morphisms-lemma-integral-universally-closed}.
D'après Morphismes d'espaces, Lemme
\ref{spaces-morphisms-lemma-image-universally-closed-separated},
$X$ est un espace algébrique séparé.
\end{proof}

\begin{proposition}
\label{spaces-cohomology-proposition-vanishing-affine}
\begin{slogan}
Le critère d'affinité de Serre pour les espaces algébriques.
\end{slogan}
Un espace algébrique quasi-compact et quasi-séparé est affine
si et seulement si tous les groupes de cohomologie supérieure des faisceaux quasi-cohérents
s'annulent. Plus précisément, tout espace algébrique comme dans la
Situation \ref{spaces-cohomology-situation-vanishing} est un schéma affine.
\end{proposition}

\begin{proof}
Choisissons un schéma affine $U = \Spec(B)$ et un morphisme \'etale surjectif
$\varphi : U \to X$. Posons $R = U \times_X U$. Comme $p$ est séparé
(Lemme \ref{spaces-cohomology-lemma-vanishing-separated}), $R$ est un
sous-schéma fermé de $U \times_{\Spec(A)} U = \Spec(B \otimes_A B)$.
Ainsi $R = \Spec(C)$ est lui aussi affine, et l'homomorphisme d'anneaux
$$
B \otimes_A B \longrightarrow C
$$
est surjectif. Notons comme d'habitude $s, t : B \to C$ les deux homomorphismes. Choisissons
$g_1, \ldots, g_m \in B$ tels que $s(g_1), \ldots, s(g_m)$ engendrent $C$
sur $t : B \to C$ (ce qui est possible puisque $t : B \to C$ est de présentation
finie et que l'homomorphisme affiché est surjectif). Alors $g_1, \ldots, g_m$
définissent des sections globales de $\varphi_*\mathcal{O}_U$, et le morphisme
$$
\mathcal{O}_X[z_1, \ldots, z_n] \longrightarrow \varphi_*\mathcal{O}_U,
\quad
z_j \longmapsto g_j
$$
est surjectif : on peut le vérifier après restriction à $U$.
En effet, $\varphi^*\varphi_*\mathcal{O}_U = t_*\mathcal{O}_R$
(d'après le Lemme \ref{spaces-cohomology-lemma-flat-base-change-cohomology});
on obtient donc exactement la condition que les $s(g_i)$ engendrent $C$
sur $t : B \to C$. L'annulation de $H^1$ du noyau montre alors que
$$
\Gamma(X, \mathcal{O}_X[x_1, \ldots, x_n]) =
A[x_1, \ldots, x_n] \longrightarrow
\Gamma(X, \varphi_*\mathcal{O}_U) = \Gamma(U, \mathcal{O}_U) = B
$$
est surjectif. On conclut que $B$ est une $A$-algèbre de type fini.
Ainsi $X \to \Spec(A)$ est de type fini et séparé.
D'après le Lemme \ref{spaces-cohomology-lemma-vanishing-injective}
et le
Lemme \ref{spaces-morphisms-lemma-locally-quasi-finite} de Morphismes d'espaces,
il est aussi localement quasi-fini. Par conséquent, $X \to \Spec(A)$ est représentable d'après
Morphismes d'espaces, Lemme
\ref{spaces-morphisms-lemma-locally-quasi-finite-separated-representable},
et $X$ est un schéma. Enfin, $X$ est affine, donc égal à $\Spec(A)$,
par application de Cohomologie des schémas, Lemme
\ref{coherent-lemma-quasi-compact-h1-zero-covering}.
\end{proof}

\begin{lemma}
\label{spaces-cohomology-lemma-Noetherian-h1-zero}
Soit $S$ un schéma. Soit $X$ un espace algébrique noethérien sur $S$.
Supposons que, pour tout $\mathcal{O}_X$-module cohérent
$\mathcal{F}$, on ait $H^1(X, \mathcal{F}) = 0$.
Alors $X$ est un schéma affine.
\end{lemma}

\begin{proof}
L'hypothèse entraîne que $H^1(X, \mathcal{F}) = 0$ pour tout
$\mathcal{O}_X$-module quasi-cohérent $\mathcal{F}$, d'après les
Lemmes \ref{spaces-cohomology-lemma-directed-colimit-coherent} et \ref{spaces-cohomology-lemma-colimits}.
Alors $X$ est affine d'après la
Proposition \ref{spaces-cohomology-proposition-vanishing-affine}.
\end{proof}

\begin{lemma}
\label{spaces-cohomology-lemma-Noetherian-h1-zero-invertible}
Soit $S$ un schéma. Soit $X$ un espace algébrique noethérien sur $S$.
Soit $\mathcal{L}$ un $\mathcal{O}_X$-module inversible.
Supposons que, pour tout $\mathcal{O}_X$-module cohérent
$\mathcal{F}$, il existe $n \geq 1$ tel que
$H^1(X, \mathcal{F} \otimes_{\mathcal{O}_X} \mathcal{L}^{\otimes n}) = 0$.
Alors $X$ est un schéma et $\mathcal{L}$ est ample sur $X$.
\end{lemma}

\begin{proof}
Soit $s \in H^0(X, \mathcal{L}^{\otimes d})$ une section globale.
Soit $U \subset X$ le sous-espace ouvert sur lequel $s$ engendre
$\mathcal{L}^{\otimes d}$. En particulier, on a
$\mathcal{L}^{\otimes d}|_U \cong \mathcal{O}_U$.
Nous affirmons que $U$ est affine.

\medskip\noindent
Démonstration de l'affirmation. Montrons que $H^1(U, \mathcal{F}) = 0$
pour tout $\mathcal{O}_U$-module quasi-cohérent $\mathcal{F}$.
Cela démontrera l'affirmation d'après la Proposition \ref{spaces-cohomology-proposition-vanishing-affine}.
Notons $j : U \to X$ le morphisme d'inclusion.
Comme le morphisme $j$ est \'etale-localement affine
(d'après Morphismes, Lemme \ref{morphisms-lemma-affine-s-open}),
on voit que $j$ est affine (Morphismes d'espaces, Lemme
\ref{spaces-morphisms-lemma-affine-local}).
On a donc
$$
H^1(U, \mathcal{F}) = H^1(X, j_*\mathcal{F})
$$
d'après le Lemme \ref{spaces-cohomology-lemma-affine-vanishing-higher-direct-images}
(et Cohomologie sur les sites, Lemme \ref{sites-cohomology-lemma-apply-Leray}).
Écrivons $j_*\mathcal{F} = \colim \mathcal{F}_i$ comme limite inductive filtrante
de $\mathcal{O}_X$-modules cohérents, voir le
Lemme \ref{spaces-cohomology-lemma-directed-colimit-coherent}. Alors
$$
H^1(X, j_*\mathcal{F}) = \colim H^1(X, \mathcal{F}_i)
$$
d'après le Lemme \ref{spaces-cohomology-lemma-colimits}.
Il suffit donc de montrer que $H^1(X, \mathcal{F}_i)$ s'envoie
sur zéro dans $H^1(U, j^*\mathcal{F}_i)$. Par hypothèse, il existe
$n \geq 1$ tel que
$$
H^1(X,
\mathcal{F}_i \otimes_{\mathcal{O}_X}
(\mathcal{O}_X \oplus \mathcal{L} \oplus \ldots
\oplus \mathcal{L}^{\otimes d - 1})
\otimes_{\mathcal{O}_X} \mathcal{L}^{\otimes n}) = 0
$$
Il existe donc $a \geq 0$ tel que
$H^1(X, \mathcal{F}_i \otimes_{\mathcal{O}_X} \mathcal{L}^{\otimes ad}) = 0$.
D'autre part, le morphisme
$$
s^a : \mathcal{F}_i \longrightarrow
\mathcal{F}_i \otimes_{\mathcal{O}_X} \mathcal{L}^{\otimes ad}
$$
est un isomorphisme après restriction à $U$. En considérant le
diagramme commutatif
$$
\xymatrix{
H^1(X, \mathcal{F}_i) \ar[r] \ar[d]_{s^a} & H^1(U, j^*\mathcal{F}_i)
\ar[d]^{\cong} \\
H^1(X, \mathcal{F}_i \otimes_{\mathcal{O}_X} \mathcal{L}^{\otimes ad}) \ar[r] &
H^1(U,
j^*(\mathcal{F}_i \otimes_{\mathcal{O}_X} \mathcal{L}^{\otimes ad}))
}
$$
on conclut que le morphisme $H^1(X, \mathcal{F}_i) \to H^1(U, j^*\mathcal{F}_i)$
est nul, et l'affirmation est démontrée.

\medskip\noindent
Soit $x \in |X|$ un point fermé. D'après Espaces décents, Lemme
\ref{decent-spaces-lemma-decent-space-closed-point},
on peut représenter $x$ par une immersion fermée $i : \Spec(k) \to X$
(cela utilise aussi le fait qu'un espace algébrique quasi-séparé est
décent, voir Espaces décents, section
\ref{decent-spaces-section-reasonable-decent}).
Ainsi $\mathcal{O}_X \to i_*\mathcal{O}_{\Spec(k)}$ est surjectif.
Soit $\mathcal{I} \subset \mathcal{O}_X$ son noyau, et choisissons
$d \geq 1$ tel que
$H^1(X, \mathcal{I} \otimes_{\mathcal{O}_X} \mathcal{L}^{\otimes d}) = 0$.
Alors
$$
H^0(X, \mathcal{L}^{\otimes d}) \to
H^0(X,
i_*\mathcal{O}_{\Spec(k)} \otimes_{\mathcal{O}_X} \mathcal{L}^{\otimes d}) =
H^0(\Spec(k), i^*\mathcal{L}^{\otimes d}) \cong k
$$
est surjectif d'après la suite exacte longue de cohomologie. Il existe donc
$s \in H^0(X, \mathcal{L}^{\otimes d})$
tel que $x \in U$, où $U$ est le sous-espace ouvert associé à $s$
comme ci-dessus. D'après notre affirmation, $x$ appartient donc au lieu schématique
(voir Propriétés des espaces, Lemme \ref{spaces-properties-lemma-subscheme})
de $X$.

\medskip\noindent
Pour conclure que $X$ est un schéma, il suffit de montrer que
toute partie ouverte de $|X|$ qui contient tous les points fermés
est égale à $|X|$. Cela résulte du fait que $|X|$
est un espace topologique noethérien, voir
Propriétés des espaces, Lemme \ref{spaces-properties-lemma-Noetherian-sober}.
Enfin, lorsque $X$ est un schéma, on peut appliquer
Cohomologie des schémas, Lemme
\ref{coherent-lemma-quasi-compact-h1-zero-invertible}
pour conclure que $\mathcal{L}$ est ample.
\end{proof}









\section{Morphismes finis et espaces affines}
\label{spaces-cohomology-section-finite-affine}

\noindent
Cette section est l'analogue de la
section \ref{coherent-section-finite-affine} du chapitre Cohomologie des schémas.

\begin{lemma}
\label{spaces-cohomology-lemma-finite-morphism-Noetherian}
Soit $S$ un schéma. Soit $f : Y \to X$ un morphisme d'espaces algébriques
sur $S$. Supposons $f$ fini et surjectif, et $X$ localement noethérien.
Soit $i : Z \to X$ une immersion fermée. Notons $i' : Z' \to Y$
l'image réciproque de $Z$
(Morphismes d'espaces, section \ref{spaces-morphisms-section-closed-immersions})
et $f' : Z' \to Z$
le morphisme induit. Alors $\mathcal{G} = f'_*\mathcal{O}_{Z'}$
est un $\mathcal{O}_Z$-module cohérent dont le support est $Z$.
\end{lemma}

\begin{proof}
Observons que $f'$ est le changement de base de $f$; il est donc fini
et surjectif d'après Morphismes d'espaces, Lemmes
\ref{spaces-morphisms-lemma-base-change-surjective} et
\ref{spaces-morphisms-lemma-base-change-integral}.
Notons que $Y$, $Z$ et $Z'$ sont localement noethériens d'après
Morphismes d'espaces, Lemme
\ref{spaces-morphisms-lemma-locally-finite-type-locally-noetherian}
(et parce que les immersions fermées et les morphismes finis
sont de type fini). D'après le Lemme \ref{spaces-cohomology-lemma-finite-pushforward-coherent},
$\mathcal{G}$ est un $\mathcal{O}_Z$-module cohérent.
Le support de $\mathcal{G}$ est fermé dans $|Z|$, voir
Morphismes d'espaces, Lemme
\ref{spaces-morphisms-lemma-support-finite-type}.
Si le support de $\mathcal{G}$ n'était pas égal à $|Z|$,
on pourrait donc, après avoir remplacé $X$ par un sous-espace ouvert, supposer
$\mathcal{G} = 0$ mais $Z \not = \emptyset$.
On aurait alors $f'_*\mathcal{O}_{Z'} = 0$. En particulier, la section
$1 \in \Gamma(Z', \mathcal{O}_{Z'}) = \Gamma(Z, f'_*\mathcal{O}_{Z'})$
serait nulle, ce qui entraînerait que $Z' = \emptyset$ est l'espace
algébrique vide. C'est impossible puisque $Z' \to Z$ est surjectif.
\end{proof}

\begin{lemma}
\label{spaces-cohomology-lemma-affine-morphism-projection-ideal}
Soit $S$ un schéma. Soit $f : Y \to X$ un morphisme d'espaces algébriques
sur $S$. Soit $\mathcal{F}$ un faisceau quasi-cohérent sur $Y$.
Soit $\mathcal{I}$ un faisceau quasi-cohérent d'idéaux sur $X$.
Si $f$ est affine, alors
$\mathcal{I}f_*\mathcal{F} = f_*(f^{-1}\mathcal{I}\mathcal{F})$
(avec les notations expliquées dans la démonstration).
\end{lemma}

\begin{proof}
Les notations signifient ce qui suit. Puisque $f^{-1}$ est un foncteur exact,
$f^{-1}\mathcal{I}$ est un faisceau
d'idéaux de $f^{-1}\mathcal{O}_X$. Par le morphisme
$f^\sharp : f^{-1}\mathcal{O}_X \to \mathcal{O}_Y$
sur $Y_\etale$, il opère sur $\mathcal{F}$.
Alors $f^{-1}\mathcal{I}\mathcal{F}$ est le sous-faisceau
engendré par les sommes de sections locales de la forme $as$, où $a$
est une section locale de $f^{-1}\mathcal{I}$ et $s$ une section locale
de $\mathcal{F}$. C'est un sous-$\mathcal{O}_Y$-module quasi-cohérent
de $\mathcal{F}$, car c'est aussi l'image d'un morphisme naturel
$f^*\mathcal{I} \otimes_{\mathcal{O}_Y} \mathcal{F} \to \mathcal{F}$.

\medskip\noindent
Cela étant, la démonstration est immédiate. En effet, la question est
\'etale-locale sur $X$, et l'on peut donc supposer que $X$ est un schéma affine.
Dans ce cas, le résultat découle du résultat correspondant
pour les schémas, voir Cohomologie des schémas, Lemme
\ref{coherent-lemma-affine-morphism-projection-ideal}.
\end{proof}

\begin{lemma}
\label{spaces-cohomology-lemma-image-affine-finite-morphism-affine-Noetherian}
Soit $S$ un schéma. Soit $f : Y \to X$ un morphisme d'espaces
algébriques sur $S$. Supposons que
\begin{enumerate}
\item $f$ soit fini,
\item $f$ soit surjectif,
\item $Y$ soit affine, et
\item $X$ soit noethérien.
\end{enumerate}
Alors $X$ est affine.
\end{lemma}

\begin{proof}
Nous allons montrer que, sous les hypothèses du lemme, pour tout
$\mathcal{O}_X$-module cohérent $\mathcal{F}$, on a $H^1(X, \mathcal{F}) = 0$.
Cela entraîne que $H^1(X, \mathcal{F}) = 0$ pour tout
$\mathcal{O}_X$-module quasi-cohérent $\mathcal{F}$ d'après les
Lemmes \ref{spaces-cohomology-lemma-directed-colimit-coherent} et \ref{spaces-cohomology-lemma-colimits}.
Il s'ensuit alors que $X$ est affine d'après la
Proposition \ref{spaces-cohomology-proposition-vanishing-affine}.

\medskip\noindent
Soit $\mathcal{P}$ la propriété des faisceaux cohérents
$\mathcal{F}$ sur $X$ définie par
$$
\mathcal{P}(\mathcal{F}) \Leftrightarrow H^1(X, \mathcal{F}) = 0.
$$
Nous allons appliquer le Lemme \ref{spaces-cohomology-lemma-property-higher-rank-cohomological}.
Il faut donc vérifier les propriétés (1), (2) et (3) de ce lemme pour $\mathcal{P}$.
La propriété (1) résulte de la suite exacte longue de cohomologie associée
à une suite exacte courte de faisceaux. La propriété (2) résulte du fait que
$H^1(X, -)$ est un foncteur additif. Pour vérifier (3), soit $i : Z \to X$
un sous-espace fermé réduit tel que $|Z|$ soit irréductible. Soient $i' : Z' \to Y$
et $f' : Z' \to Z$ comme dans le Lemme \ref{spaces-cohomology-lemma-finite-morphism-Noetherian},
et posons $\mathcal{G} = f'_*\mathcal{O}_{Z'}$. Nous affirmons que
$\mathcal{G}$ vérifie les propriétés (3)(a) et (3)(b) du
Lemme \ref{spaces-cohomology-lemma-property-higher-rank-cohomological},
ce qui achèvera la démonstration. La propriété (3)(a) a été établie dans le
Lemme \ref{spaces-cohomology-lemma-finite-morphism-Noetherian}. Pour vérifier (3)(b), soit
$\mathcal{I}$ un faisceau quasi-cohérent non nul d'idéaux sur $Z$.
Notons $\mathcal{I}' \subset \mathcal{O}_{Z'}$ l'idéal quasi-cohérent
$(f')^{-1}\mathcal{I} \mathcal{O}_{Z'}$, c'est-à-dire l'
image de $(f')^*\mathcal{I} \to \mathcal{O}_{Z'}$.
D'après le Lemme \ref{spaces-cohomology-lemma-affine-morphism-projection-ideal}, on a
$f_*\mathcal{I}' = \mathcal{I} \mathcal{G}$.
Nous affirmons que la valeur commune
$\mathcal{G}' = \mathcal{I} \mathcal{G} = f'_*\mathcal{I}'$
satisfait la condition énoncée en (3)(b).
Tout d'abord, il est clair que le support de $\mathcal{G}/\mathcal{G}'$
est contenu dans le support de $\mathcal{O}_Z/\mathcal{I}$,
qui est un sous-espace propre de $|Z|$, puisque $\mathcal{I}$ est un
faisceau d'idéaux non nul sur l'espace algébrique réduit et irréductible $Z$.
Le morphisme $f'$ est affine, donc $R^1f'_*\mathcal{I}' = 0$ d'après le
Lemme \ref{spaces-cohomology-lemma-affine-vanishing-higher-direct-images}.
Comme $Z'$ est affine (en tant que sous-schéma fermé d'un schéma affine),
on a $H^1(Z', \mathcal{I}') = 0$. La suite spectrale de Leray
(sous la forme du
Lemme \ref{sites-cohomology-lemma-apply-Leray} de Cohomologie sur les sites)
entraîne donc que $H^1(Z, f'_*\mathcal{I}') = 0$.
Comme $i : Z \to X$ est affine, on conclut que
$R^1i_*f'_*\mathcal{I}' = 0$, puis que $H^1(X, i_*f'_*\mathcal{I}') = 0$
par une nouvelle application de Leray. Autrement dit, $H^1(X, i_*\mathcal{G}') = 0$,
comme voulu.
\end{proof}








\section{Une version faible du lemme de Chow}
\label{spaces-cohomology-section-weak-chow}

\noindent
Dans cette section, nous démontrons rapidement le lemme suivant afin
de faciliter la démonstration des résultats fondamentaux sur la cohomologie des
modules cohérents sur les espaces algébriques propres.

\begin{lemma}
\label{spaces-cohomology-lemma-weak-chow}
Soit $A$ un anneau. Soit $X$ un espace algébrique sur $\Spec(A)$
dont le morphisme structural $X \to \Spec(A)$ est séparé et de type fini.
Il existe alors un morphisme propre surjectif $X' \to X$,
où $X'$ est un schéma H-quasi-projectif sur $\Spec(A)$.
\end{lemma}

\begin{proof}
Soit $W$ un schéma affine et soit $f : W \to X$ un morphisme
\'etale surjectif. Il existe un entier $d$ tel que toutes les fibres géométriques
de f aient $\leq d$ points (car $X$ est un espace algébrique séparé,
donc raisonnable, voir
Espaces décents, Lemme \ref{decent-spaces-lemma-bounded-fibres}).
En choisissant $d$ minimal, on obtient un ouvert non vide $U \subset X$ tel que
$f^{-1}(U) \to U$ soit fini \'etale de degré $d$, voir
Espaces décents, Lemme
\ref{decent-spaces-lemma-quasi-compact-reasonable-stratification}.
Soit
$$
V \subset W \times_X W \times_X \ldots \times_X W
$$
(avec $d$ facteurs dans le produit fibré) le complémentaire de toutes les diagonales.
Comme $W \to X$ est séparé, la diagonale $W \to W \times_X W$ est une
immersion fermée. Comme $W \to X$ est \'etale, la diagonale
$W \to W \times_X W$ est une immersion ouverte, voir
Morphismes d'espaces, Lemmes
\ref{spaces-morphisms-lemma-etale-unramified} et
\ref{spaces-morphisms-lemma-diagonal-unramified-morphism}.
Les diagonales sont donc des sous-schémas ouverts et fermés
du schéma quasi-compact $W \times_X \ldots \times_X W$.
On conclut en particulier que $V$ est un schéma quasi-compact.
Choisissons une immersion ouverte $W \subset Y$, où $Y$ est H-projectif sur
$A$ (c'est possible puisque $W$ est affine et de type fini sur $A$;
on peut par exemple utiliser
Morphismes, Lemmes
\ref{morphisms-lemma-quasi-affine-finite-type-over-S} et
\ref{morphisms-lemma-H-quasi-projective-open-H-projective}).
Soit
$$
Z \subset Y \times_A Y \times_A \ldots \times_A Y
$$
l'image schématique du composé
$V \to W \times_X \ldots \times_X W \to Y \times_A \ldots \times_A Y$.
Observons que ce morphisme est quasi-compact, puisque $V$ est quasi-compact
et que $Y \times_A \ldots \times_A Y$ est séparé.
Notons que $V \to Z$ est une immersion ouverte, car
$V \to Y \times_A \ldots \times_A Y$ est une immersion, voir
Morphismes, Lemme \ref{morphisms-lemma-quasi-compact-immersion}.
Les morphismes de projection donnent $d$ morphismes $g_i : Z \to Y$.
Ces morphismes $g_i$ sont projectifs puisque $Y$ est projectif sur $A$, voir
les résultats de Morphismes, section \ref{morphisms-section-projective}.
Posons
$$
X' = \bigcup g_i^{-1}(W) \subset Z
$$
Il existe un morphisme $X' \to X$ dont la restriction à $g_i^{-1}(W)$ est
le composé $g_i^{-1}(W) \to W \to X$.
En effet, ces morphismes coïncident sur $V$, donc sur
$g_i^{-1}(W) \cap g_j^{-1}(W)$ d'après
Morphismes d'espaces, Lemme \ref{spaces-morphisms-lemma-equality-of-morphisms}.
Affirmation : le morphisme $X' \to X$ est propre.

\medskip\noindent
Si l'affirmation est vraie, le lemme résulte d'une récurrence sur $d$.
En effet, par construction, $X'$ est H-quasi-projectif sur $\Spec(A)$.
L'image de $X' \to X$ contient l'ouvert $U$, puisque $V$ se projette surjectivement sur $U$.
Notons $T$ l'espace algébrique muni de la structure réduite induite sur $X \setminus U$.
Alors $T \times_X W$ est un sous-schéma fermé de $W$, donc est affine.
De plus, le morphisme $T \times_X W \to T$ est \'etale et toute fibre géométrique
a $< d$ points. Par hypothèse de récurrence, il existe un morphisme propre
surjectif $T' \to T$, où $T'$ est un schéma H-quasi-projectif
sur $\Spec(A)$. Comme $T$ est un sous-espace fermé de $X$,
$T' \to X$ est un morphisme propre. Le lemme résulte donc du
morphisme propre surjectif $X' \amalg T' \to X$.

\medskip\noindent
Démonstration de l'affirmation. Par construction, le morphisme $X' \to X$ est séparé
et de type fini. Nous allons vérifier les conditions (1) -- (4) de
Morphismes d'espaces, Lemme
\ref{spaces-morphisms-lemma-refined-valuative-criterion-universally-closed},
pour les morphismes $V \to X'$ et $X' \to X$.
Les conditions (1) et (2) ont été établies ci-dessus.
La condition (3) est satisfaite parce que $X' \to X$ est séparé (son
espace algébrique source étant séparé). Il suffit donc de vérifier
l'existence d'un relèvement dans $X'$ pour les diagrammes
$$
\xymatrix{
\Spec(K) \ar[r] \ar[d] & V \ar[d] \\
\Spec(R) \ar[r] & X
}
$$
où $R$ est un anneau de valuation de corps des fractions $K$.
Notons que le morphisme horizontal supérieur est donné par $d$ points à valeurs dans
$K$, deux à deux distincts, $w_1, \ldots, w_d$, de $W$. En fait, ils
sont exactement les antécédents du point $x \in X(K)$
provenant du diagramme. Comme $W \to X$ est surjectif,
on peut, après avoir éventuellement remplacé $R$ par une extension d'anneaux de valuation,
relever le morphisme $\Spec(R) \to X$ en un morphisme $w : \Spec(R) \to W$, voir
Morphismes d'espaces, Lemme
\ref{spaces-morphisms-lemma-lift-valuation-ring-through-flat-morphism}.
Puisque $w_1, \ldots, w_d$ forment l'ensemble de tous les antécédents de
$x$, la restriction $w|_{\Spec(K)}$ est égale à l'un d'eux, disons $w_i$.
On obtient ainsi un diagramme commutatif
$$
\xymatrix{
\Spec(K) \ar[r] \ar[d] &  Z \ar[d]_{g_i}\\
\Spec(R) \ar[r]^w & Y
}
$$
D'après le critère valuatif de propreté appliqué au morphisme projectif
$g_i$, on peut relever $w$ en $z : \Spec(R) \to Z$, voir
Morphismes, Lemme \ref{morphisms-lemma-locally-projective-proper}
et
Schémas, Proposition \ref{schemes-proposition-characterize-universally-closed}.
L'image de $z$ est contenue dans $g_i^{-1}(W) \subset X'$, ce qui achève la démonstration.
\end{proof}







\section{Critère valuatif noethérien}
\label{spaces-cohomology-section-Noetherian-valuative-criterion}

\noindent
Nous démontrons une version du critère valuatif de propreté
faisant intervenir des anneaux de valuation discrète. On trouvera des versions plus précises
(et donc plus techniques) dans
Limites d'espaces, section
\ref{spaces-limits-section-Noetherian-valuative-criterion}.

\begin{lemma}
\label{spaces-cohomology-lemma-check-separated-dvr}
Soit $S$ un schéma. Soit $f : X \to Y$ un morphisme d'espaces algébriques
sur $S$. Supposons que
\begin{enumerate}
\item $Y$ soit localement noethérien,
\item $f$ soit localement de type fini et quasi-séparé,
\item pour tout diagramme commutatif
$$
\xymatrix{
\Spec(K) \ar[r] \ar[d] & X \ar[d] \\
\Spec(A) \ar[r] \ar@{-->}[ru] & Y
}
$$
où $A$ est un anneau de valuation discrète et $K$ son corps des fractions,
il existe au plus une flèche pointillée rendant le diagramme commutatif.
\end{enumerate}
Alors $f$ est séparé.
\end{lemma}

\begin{proof}
Il faut montrer que la diagonale $\Delta : X \to X \times_Y X$ est une immersion
fermée. On sait déjà que $\Delta$ est représentable, séparée,
un monomorphisme et localement de type fini, voir Morphismes d'espaces, Lemme
\ref{spaces-morphisms-lemma-properties-diagonal}.
Choisissons un schéma affine $U$ et un morphisme \'etale
$U \to X \times_Y X$. Posons $V = X \times_{\Delta, X \times_Y X} U$.
Il suffit de montrer que $V \to U$ est une immersion fermée
(Morphismes d'espaces, Lemme
\ref{spaces-morphisms-lemma-closed-immersion-local}).
Comme $X \times_Y X$ est localement de type fini sur $Y$,
$U$ est noethérien (utiliser Morphismes d'espaces, Lemmes
\ref{spaces-morphisms-lemma-composition-finite-type},
\ref{spaces-morphisms-lemma-base-change-finite-type} et
\ref{spaces-morphisms-lemma-locally-finite-type-locally-noetherian}).
Notons que $V$ est un schéma, puisque $\Delta$ est représentable.
En outre, $V$ est quasi-compact puisque $f$ est quasi-séparé.
Ainsi $V \to U$ est de type fini. Considérons un diagramme commutatif
$$
\xymatrix{
\Spec(K) \ar[r] \ar[d] & V \ar[d] \\
\Spec(A) \ar[r] \ar@{-->}[ru] & U
}
$$
de morphismes de schémas,
où $A$ est un anneau de valuation discrète de corps des fractions $K$.
On peut interpréter le composé $\Spec(A) \to U \to X \times_Y X$
comme une paire de morphismes $a, b : \Spec(A) \to X$ qui coïncident comme morphismes
dans $Y$ et dont les restrictions à $\Spec(K)$ sont égales. L'hypothèse
(3) garantit donc que $a = b$, et l'on obtient la flèche pointillée du diagramme.
D'après Limites, Lemme \ref{limits-lemma-Noetherian-dvr-valuative-proper},
on conclut que $V \to U$ est propre. Autrement dit, $\Delta$ est propre.
Comme $\Delta$ est un monomorphisme, on en déduit que $\Delta$ est une
immersion fermée (Morphismes \'etales, Lemme
\ref{etale-lemma-characterize-closed-immersion}), comme voulu.
\end{proof}

\begin{lemma}
\label{spaces-cohomology-lemma-check-proper-dvr}
Soit $S$ un schéma. Soit $f : X \to Y$ un morphisme d'espaces algébriques
sur $S$. Supposons que
\begin{enumerate}
\item $Y$ soit localement noethérien,
\item $f$ soit de type fini et quasi-séparé,
\item pour tout diagramme commutatif
$$
\xymatrix{
\Spec(K) \ar[r] \ar[d] & X \ar[d] \\
\Spec(A) \ar[r] \ar@{-->}[ru] & Y
}
$$
où $A$ est un anneau de valuation discrète et $K$ son corps des fractions,
il existe une unique flèche pointillée rendant le diagramme commutatif.
\end{enumerate}
Alors $f$ est propre.
\end{lemma}

\begin{proof}
Il suffit de prouver que $f$ est universellement fermé, car $f$ est séparé d'après le
Lemme \ref{spaces-cohomology-lemma-check-separated-dvr}.
Pour cela, on peut travailler \'etale-localement sur $Y$
(Morphismes d'espaces, Lemme
\ref{spaces-morphisms-lemma-universally-closed-local}).
On peut donc supposer que $Y = \Spec(A)$ est un schéma affine noethérien.
Choisissons $X' \to X$ comme dans la version faible du lemme de Chow
(Lemme \ref{spaces-cohomology-lemma-weak-chow}). Nous affirmons que $X' \to \Spec(A)$
est universellement fermé. L'affirmation entraîne le lemme d'après
Morphismes d'espaces, Lemme \ref{spaces-morphisms-lemma-image-proper-is-proper}.
Pour la démontrer, d'après
Limites, Lemme \ref{limits-lemma-check-universally-closed-Noetherian},
il suffit de prouver que, dans tout diagramme commutatif à flèches pleines
$$
\xymatrix{
\Spec(K) \ar[r] \ar[d] & X' \ar[r] & X \ar[d] \\
\Spec(A) \ar[rr] \ar@{-->}[ru]^a \ar@{-->}[rru]_b & & Y
}
$$
où $A$ est un anneau de valuation discrète de corps des fractions $K$, on peut trouver la
flèche pointillée $a$. Par hypothèse, on peut trouver la flèche pointillée $b$.
Alors le morphisme $X' \times_{X, b} \Spec(A) \to \Spec(A)$
est un morphisme propre de schémas et, d'après le critère valuatif
pour les morphismes de schémas, le morphisme $b$ admet le relèvement voulu $a$.
\end{proof}

\begin{remark}[Variante pour les anneaux de valuation discrète complets]
\label{spaces-cohomology-remark-variant}
Dans les Lemmes \ref{spaces-cohomology-lemma-check-separated-dvr} et \ref{spaces-cohomology-lemma-check-proper-dvr},
il suffit de considérer les anneaux de valuation discrète complets.
Plus précisément, dans le Lemme \ref{spaces-cohomology-lemma-check-separated-dvr}, on peut remplacer
la condition (3) par la condition suivante : pour tout diagramme commutatif
$$
\xymatrix{
\Spec(K) \ar[r] \ar[d] & X \ar[d] \\
\Spec(A) \ar[r] \ar@{-->}[ru] & Y
}
$$
où $A$ est un anneau de valuation discrète complet de corps des fractions $K$,
il existe au plus une flèche pointillée rendant le diagramme commutatif. En effet, pour
tout diagramme comme dans le Lemme \ref{spaces-cohomology-lemma-check-separated-dvr} (3),
le complété $A^\wedge$ est un anneau de valuation discrète
(Compléments d'algèbre, Lemme \ref{more-algebra-lemma-completion-dvr}),
et l'unicité de la flèche $\Spec(A^\wedge) \to X$
entraîne celle de la flèche $\Spec(A) \to X$,
par exemple d'après Propriétés des espaces, Proposition
\ref{spaces-properties-proposition-sheaf-fpqc}.
De même, dans le Lemme \ref{spaces-cohomology-lemma-check-proper-dvr},
on peut remplacer la condition (3) par la condition suivante :
pour tout diagramme commutatif
$$
\xymatrix{
\Spec(K) \ar[r] \ar[d] & X \ar[d] \\
\Spec(A) \ar[r] & Y
}
$$
où $A$ est un anneau de valuation discrète complet de corps des fractions $K$,
il existe une extension $A \subset A'$ d'anneaux de valuation discrète complets
induisant une extension des corps des fractions $K \subset K'$ telle qu'il existe une
unique flèche $\Spec(A') \to X$ rendant le diagramme
$$
\xymatrix{
\Spec(K') \ar[r] \ar[d] & \Spec(K) \ar[r] & X \ar[d] \\
\Spec(A') \ar[r] \ar[rru] & \Spec(A) \ar[r] & Y
}
$$
commutatif. En effet, pour tout diagramme comme dans le Lemme \ref{spaces-cohomology-lemma-check-proper-dvr}
partie (3), l'existence d'un diagramme commutatif
$$
\xymatrix{
\Spec(L) \ar[r] \ar[d] & \Spec(K) \ar[r] & X \ar[d] \\
\Spec(B) \ar[r] \ar[rru] & \Spec(A) \ar[r] & Y
}
$$
pour {\it toute} extension $A \subset B$ d'anneaux de valuation discrète
entraîne l'existence d'une flèche $\Spec(A) \to X$ s'insérant dans
le diagramme. Cela a été démontré dans
Morphismes d'espaces, Lemme \ref{spaces-morphisms-lemma-push-down-solution}.
En fait, ces considérations montrent qu'il suffit de chercher
les flèches pointillées dans les diagrammes pour toute classe d'anneaux de valuation discrète
telle que, pour tout anneau de valuation discrète, il en existe une extension
qui appartienne à la classe. On pourrait, par exemple, prendre les anneaux de valuation
discrète complets à corps résiduel algébriquement clos.
\end{remark}







\section{Images directes supérieures des faisceaux cohérents}
\label{spaces-cohomology-section-proper-pushforward}

\noindent
Dans cette section, nous démontrons le fait fondamental que les images
directes supérieures d'un faisceau cohérent par un morphisme propre
sont cohérentes. Commençons par un lemme auxiliaire.

\begin{lemma}
\label{spaces-cohomology-lemma-kill-by-twisting}
Soit $S$ un schéma. Considérons un diagramme commutatif
$$
\xymatrix{
X \ar[r]_i \ar[rd]_f & \mathbf{P}^n_Y \ar[d] \\
& Y
}
$$
d'espaces algébriques au-dessus de $S$. Supposons que $i$ soit une immersion fermée
et que $Y$ soit noethérien. Posons $\mathcal{L} = i^*\mathcal{O}_{\mathbf{P}^n_Y}(1)$.
Soit $\mathcal{F}$ un module cohérent sur $X$.
Il existe alors un entier $d_0$ tel que, pour tout $d \geq d_0$, on ait
$R^pf_*(\mathcal{F} \otimes_{\mathcal{O}_X} \mathcal{L}^{\otimes d}) = 0$
pour tout $p > 0$.
\end{lemma}

\begin{proof}
La nullité de $R^pf_*(\mathcal{F} \otimes \mathcal{L}^{\otimes d})$
se vérifie localement pour la topologie \'etale sur $Y$, voir
l'équation (\ref{spaces-cohomology-equation-representable-higher-direct-image}).
On peut donc supposer que $Y$ est le spectre d'un anneau noethérien. Dans ce cas,
$X$ est un schéma et le résultat découle de
Cohomologie des schémas, lemme \ref{coherent-lemma-kill-by-twisting}.
\end{proof}

\begin{lemma}
\label{spaces-cohomology-lemma-proper-pushforward-coherent}
Soit $S$ un schéma. Soit $f : X \to Y$ un morphisme propre
d'espaces algébriques au-dessus de $S$, avec $Y$ localement noethérien.
Soit $\mathcal{F}$ un $\mathcal{O}_X$-module cohérent.
Alors $R^if_*\mathcal{F}$ est un $\mathcal{O}_Y$-module cohérent
pour tout $i \geq 0$.
\end{lemma}

\begin{proof}
Remarquons d'abord que $X$ est un espace algébrique localement noethérien
d'après Morphismes d'espaces, lemme
\ref{spaces-morphisms-lemma-locally-finite-type-locally-noetherian}.
L'énoncé du lemme a donc un sens. De plus, la formation de
$R^if_*\mathcal{F}$ commute à la localisation \'etale sur $Y$
(Propriétés des espaces, lemme
\ref{spaces-properties-lemma-pushforward-etale-base-change-modules})
et la cohérence de $R^if_*\mathcal{F}$ se vérifie
localement pour la topologie \'etale sur $Y$ (lemme \ref{spaces-cohomology-lemma-coherent-Noetherian}).
On peut donc supposer que $Y = \Spec(A)$ est un schéma affine noethérien.

\medskip\noindent
Supposons que $Y = \Spec(A)$ soit un schéma affine. Notons que $f$ est localement
de présentation finie
(Morphismes d'espaces, lemme
\ref{spaces-morphisms-lemma-noetherian-finite-type-finite-presentation}).
Il est donc de présentation finie, et par suite $X$ est noethérien
(Morphismes d'espaces, lemme
\ref{spaces-morphisms-lemma-finite-presentation-noetherian}).
Le lemme \ref{spaces-cohomology-lemma-property-higher-rank-cohomological-variant}
s'applique donc à la catégorie des modules cohérents sur $X$.
Pour un faisceau cohérent $\mathcal{F}$ sur $X$, disons que $\mathcal{P}$ est vérifiée
si et seulement si $R^if_*\mathcal{F}$ est un module cohérent sur $\Spec(A)$.
Nous allons montrer que les conditions (1), (2) et (3) du
lemme \ref{spaces-cohomology-lemma-property-higher-rank-cohomological-variant} sont vérifiées
pour cette propriété, ce qui achèvera la démonstration du lemme.

\medskip\noindent
Vérification de la condition (1). Soit
$$
0 \to \mathcal{F}_1 \to \mathcal{F}_2 \to \mathcal{F}_3 \to 0
$$
une suite exacte courte de faisceaux cohérents sur $X$.
Considérons la suite exacte longue des images directes supérieures
$$
R^{p - 1}f_*\mathcal{F}_3 \to
R^pf_*\mathcal{F}_1 \to
R^pf_*\mathcal{F}_2 \to
R^pf_*\mathcal{F}_3 \to
R^{p + 1}f_*\mathcal{F}_1
$$
Il est alors clair que, si deux des trois faisceaux $\mathcal{F}_i$
vérifient la propriété $\mathcal{P}$, les images directes supérieures du
troisième sont encadrées, dans ce complexe exact, par deux faisceaux
cohérents. Ces images directes supérieures sont donc elles aussi cohérentes d'après les
lemmes \ref{spaces-cohomology-lemma-coherent-abelian-Noetherian} et
\ref{spaces-cohomology-lemma-coherent-Noetherian-quasi-coherent-sub-quotient}.
Ainsi, la propriété $\mathcal{P}$ est également vérifiée pour le troisième.

\medskip\noindent
Vérification de la condition (2). Elle résulte immédiatement de ce que
$R^if_*(\mathcal{F}_1 \oplus \mathcal{F}_2) =
R^if_*\mathcal{F}_1 \oplus R^if_*\mathcal{F}_2$ et de ce qu'un facteur direct
d'un module cohérent est cohérent (voir les lemmes cités ci-dessus).

\medskip\noindent
Vérification de la condition (3). Soit $i : Z \to X$ une immersion fermée,
où $Z$ est réduit et $|Z|$ irréductible. Posons $g = f \circ i : Z \to \Spec(A)$.
Soit $\mathcal{G}$ un module cohérent sur $Z$ dont le support schématique
est égal à $Z$, tel que $R^pg_*\mathcal{G}$ soit cohérent pour tout $p$.
Alors $\mathcal{F} = i_*\mathcal{G}$ est un module cohérent sur
$X$, de support schématique $Z$, tel que
$R^pf_*\mathcal{F} = R^pg_*\mathcal{G}$. Pour le voir, utilisons
la suite spectrale de Leray
(Cohomologie sur les sites, lemme \ref{sites-cohomology-lemma-relative-Leray})
et le fait que $R^qi_*\mathcal{G} = 0$ pour $q > 0$ d'après le
lemme \ref{spaces-cohomology-lemma-affine-vanishing-higher-direct-images},
ainsi que le fait qu'une immersion fermée est affine
(Morphismes d'espaces, lemme
\ref{spaces-morphisms-lemma-closed-immersion-affine}).
On se ramène donc à trouver un faisceau cohérent $\mathcal{G}$ sur $Z$
de support égal à $Z$, tel que $R^pg_*\mathcal{G}$ soit cohérent
pour tout $p$.

\medskip\noindent
Appliquons le lemme \ref{spaces-cohomology-lemma-weak-chow} au morphisme $Z \to \Spec(A)$.
On obtient ainsi un diagramme
$$
\xymatrix{
Z \ar[rd]_g & Z' \ar[d]^-{g'} \ar[l]^\pi \ar[r]_i & \mathbf{P}^n_A \ar[dl] \\
& \Spec(A) &
}
$$
où $\pi : Z' \to Z$ est propre et surjectif, et $i$ une immersion.
Puisque $Z \to \Spec(A)$ est propre, on conclut que $g'$ est propre
(Morphismes d'espaces, lemme \ref{spaces-morphisms-lemma-composition-proper}).
Par suite, $i$ est une immersion fermée
(Morphismes d'espaces, lemmes
\ref{spaces-morphisms-lemma-universally-closed-permanence} et
\ref{spaces-morphisms-lemma-immersion-when-closed}).
Il s'ensuit que le morphisme
$i' = (i, \pi) : \mathbf{P}^n_A \times_{\Spec(A)} Z' = \mathbf{P}^n_Z$ est
une immersion fermée
(Morphismes d'espaces, lemme \ref{spaces-morphisms-lemma-semi-diagonal}).
Posons
$$
\mathcal{L} =
i^*\mathcal{O}_{\mathbf{P}^n_A}(1) =
(i')^*\mathcal{O}_{\mathbf{P}^n_Z}(1)
$$
On peut appliquer le lemme \ref{spaces-cohomology-lemma-kill-by-twisting}
à $\mathcal{L}$ et $\pi$, ainsi qu'à $\mathcal{L}$ et $g'$.
Ainsi, pour tout $d \gg 0$, on a
$R^p\pi_*\mathcal{L}^{\otimes d} = 0$ pour tout $p > 0$ et
$R^p(g')_*\mathcal{L}^{\otimes d} = 0$ pour tout $p > 0$.
Posons $\mathcal{G} = \pi_*\mathcal{L}^{\otimes d}$.
D'après la suite spectrale de Leray
(Cohomologie sur les sites, lemme \ref{sites-cohomology-lemma-relative-Leray}),
on a
$$
E_2^{p, q} = R^pg_* R^q\pi_*\mathcal{L}^{\otimes d}
\Rightarrow
R^{p + q}(g')_*\mathcal{L}^{\otimes d}
$$
et, par le choix de $d$, les seuls termes non nuls de $E_2^{p, q}$ sont
ceux pour lesquels $q = 0$, tandis que les seuls termes non nuls de
$R^{p + q}(g')_*\mathcal{L}^{\otimes d}$ sont ceux pour lesquels $p = q = 0$.
Il en résulte que $R^pg_*\mathcal{G} = 0$ pour $p > 0$ et
que $g_*\mathcal{G} = (g')_*\mathcal{L}^{\otimes d}$.
En appliquant
Cohomologie des schémas, lemme
\ref{coherent-lemma-locally-projective-pushforward}
on voit que $g_*\mathcal{G} = (g')_*\mathcal{L}^{\otimes d}$ est
cohérent.

\medskip\noindent
Il reste à vérifier que le support de $\mathcal{G}$ est $Z$.
Cela résulte du fait que $\mathcal{L}^{\otimes d}$ possède
beaucoup de sections globales. Donnons les détails.
Notons que $\mathcal{L}^{\otimes d}$ est engendré par ses sections globales pour tout $d \geq 0$,
car il en est de même de $\mathcal{O}_{\mathbf{P}^n}(d)$.
Choisissons un point $z \in Z'$ s'envoyant sur le point générique $\xi$ de $Z$ ;
c'est possible puisque $\pi$ est surjectif.
(Observons que $Z$ possède bien un point générique, car $|Z|$ est irréductible
et $Z$ est noethérien, donc quasi-séparé ; ainsi $|Z|$ est un espace topologique
sobre d'après
Propriétés des espaces, lemme
\ref{spaces-properties-lemma-quasi-separated-sober}.)
Choisissons $s \in \Gamma(Z', \mathcal{L}^{\otimes d})$ qui ne s'annule pas en $z$.
Puisque $\Gamma(Z, \mathcal{G}) = \Gamma(Z', \mathcal{L}^{\otimes d})$,
on peut considérer $s$ comme une section globale de $\mathcal{G}$.
Choisissons un point géométrique $\overline{z}$ de $Z'$ au-dessus de $z$
et notons $\overline{\xi} = g' \circ \overline{z}$
le point géométrique correspondant de $Z$. L'application d'adjonction
$$
(g')^*\mathcal{G} =
(g')^*g'_*\mathcal{L}^{\otimes d} \longrightarrow \mathcal{L}^{\otimes d}
$$
induit une application entre les fibres
$\mathcal{G}_{\overline{\xi}} \to \mathcal{L}_{\overline{z}}$, voir
Propriétés des espaces, lemme
\ref{spaces-properties-lemma-stalk-pullback-quasi-coherent}.
De plus, l'application d'adjonction envoie l'image réciproque de $s$ (vu comme section
de $\mathcal{G}$) sur $s$ (vu comme section de $\mathcal{L}^{\otimes d}$).
Ainsi, l'image de $s$ dans l'espace vectoriel source de la flèche
$$
\mathcal{G}_{\overline{\xi}} \otimes \kappa(\overline{\xi})
\longrightarrow
\mathcal{L}^{\otimes d}_{\overline{z}} \otimes \kappa(\overline{z})
$$
n'est pas nulle, puisque, par le choix de $s$, son image dans la cible de la flèche
est non nulle. Ainsi, $\xi$ appartient au support de $\mathcal{G}$
(Morphismes d'espaces, lemme \ref{spaces-morphisms-lemma-support-finite-type}).
Puisque $|Z|$ est irréductible et $Z$ réduit, on conclut que
le support schématique de $\mathcal{G}$ est $Z$ tout entier, comme voulu.
\end{proof}

\begin{lemma}
\label{spaces-cohomology-lemma-proper-over-affine-cohomology-finite}
Soit $A$ un anneau noethérien.
Soit $f : X \to \Spec(A)$ un morphisme propre d'espaces algébriques.
Soit $\mathcal{F}$ un $\mathcal{O}_X$-module cohérent.
Alors $H^i(X, \mathcal{F})$ est un $A$-module de type fini pour tout $i \geq 0$.
\end{lemma}

\begin{proof}
Il s'agit simplement du cas affine du lemme \ref{spaces-cohomology-lemma-proper-pushforward-coherent}.
En effet, le lemme \ref{spaces-cohomology-lemma-higher-direct-image} montre que
$R^if_*\mathcal{F}$ est un faisceau quasi-cohérent. C'est donc le faisceau
quasi-cohérent associé au $A$-module
$\Gamma(\Spec(A), R^if_*\mathcal{F}) = H^i(X, \mathcal{F})$.
L'égalité résulte de
Cohomologie sur les sites, lemme \ref{sites-cohomology-lemma-apply-Leray},
et de l'annulation des groupes de cohomologie supérieurs des modules quasi-cohérents
sur les schémas affines (Cohomologie des schémas, lemme
\ref{coherent-lemma-quasi-coherent-affine-cohomology-zero}).
D'après le lemme \ref{spaces-cohomology-lemma-coherent-Noetherian}, le faisceau $R^if_*\mathcal{F}$ est
cohérent si et seulement si $H^i(X, \mathcal{F})$
est un $A$-module de type fini. Le
lemme \ref{spaces-cohomology-lemma-proper-pushforward-coherent} donne donc la conclusion.
\end{proof}

\begin{lemma}
\label{spaces-cohomology-lemma-graded-finiteness}
Soit $A$ un anneau noethérien.
Soit $B$ une $A$-algèbre graduée de type fini.
Soit $f : X \to \Spec(A)$ un morphisme propre d'espaces algébriques.
Posons $\mathcal{B} = f^*\widetilde B$.
Soit $\mathcal{F}$ un $\mathcal{B}$-module gradué
quasi-cohérent de type fini.
Pour tout $p \geq 0$, le $B$-module gradué $H^p(X, \mathcal{F})$
est un $B$-module de type fini.
\end{lemma}

\begin{proof}
Pour le démontrer, considérons le diagramme de produit fibré
$$
\xymatrix{
X' = \Spec(B) \times_{\Spec(A)} X
\ar[r]_-\pi \ar[d]_{f'} &
X \ar[d]^f \\
\Spec(B) \ar[r] &
\Spec(A)
}
$$
Notons que $f'$ est un morphisme propre, voir
Morphismes d'espaces, lemme \ref{spaces-morphisms-lemma-base-change-proper}.
De plus, $B$ est une $A$-algèbre de type fini, et est donc
noethérienne (Algèbre, lemme \ref{algebra-lemma-Noetherian-permanence}).
Il en résulte que $X'$ est un espace algébrique noethérien
(Morphismes d'espaces, lemme
\ref{spaces-morphisms-lemma-finite-presentation-noetherian}).
Notons que $X'$ est le spectre relatif de la
$\mathcal{O}_X$-algèbre quasi-cohérente $\mathcal{B}$ d'après
Morphismes d'espaces, lemme
\ref{spaces-morphisms-lemma-affine-equivalence-algebras}.
Puisque $\mathcal{F}$ est un $\mathcal{B}$-module quasi-cohérent,
il existe un unique
$\mathcal{O}_{X'}$-module quasi-cohérent $\mathcal{F}'$ tel que
$\pi_*\mathcal{F}' = \mathcal{F}$, voir
Morphismes d'espaces, lemme
\ref{spaces-morphisms-lemma-affine-equivalence-modules}.
Puisque $\mathcal{F}$ est de type fini comme $\mathcal{B}$-module, on
conclut que $\mathcal{F}'$ est un
$\mathcal{O}_{X'}$-module de type fini (détails omis). Autrement dit,
$\mathcal{F}'$ est un $\mathcal{O}_{X'}$-module cohérent
(lemme \ref{spaces-cohomology-lemma-coherent-Noetherian}).
Puisque le morphisme $\pi : X' \to X$ est affine, on a
$$
H^p(X, \mathcal{F}) = H^p(X', \mathcal{F}')
$$
d'après le
lemme \ref{spaces-cohomology-lemma-affine-vanishing-higher-direct-images}
et
Cohomologie sur les sites, lemme \ref{sites-cohomology-lemma-apply-Leray}.
Le lemme résulte donc du
lemme \ref{spaces-cohomology-lemma-proper-over-affine-cohomology-finite}.
\end{proof}









\section{Faisceaux inversibles amples et cohomologie}
\label{spaces-cohomology-section-ample-cohomology}

\noindent
Voici un critère d'amplitude pour les espaces algébriques propres sur une base affine,
exprimé par l'annulation de la cohomologie après tensorisation.

\begin{lemma}
\label{spaces-cohomology-lemma-vanshing-gives-ample}
Soit $R$ un anneau noethérien. Soit $X$ un espace algébrique propre
sur $R$. Soit $\mathcal{L}$ un $\mathcal{O}_X$-module inversible.
Les conditions suivantes sont équivalentes :
\begin{enumerate}
\item $X$ est un schéma et $\mathcal{L}$ est ample sur $X$ ;
\item pour tout $\mathcal{O}_X$-module cohérent $\mathcal{F}$, il existe
un $n_0 \geq 0$ tel que
$H^p(X, \mathcal{F} \otimes \mathcal{L}^{\otimes n}) = 0$ pour tout $n \geq n_0$
et tout $p > 0$ ; et
\item pour tout $\mathcal{O}_X$-module cohérent $\mathcal{F}$, il existe
un $n \geq 1$ tel que
$H^1(X, \mathcal{F} \otimes \mathcal{L}^{\otimes n}) = 0$.
\end{enumerate}
\end{lemma}

\begin{proof}
L'implication (1) $\Rightarrow$ (2) résulte de
Cohomologie des schémas, lemme \ref{coherent-lemma-vanshing-gives-ample}.
L'implication (2) $\Rightarrow$ (3) est triviale.
L'implication (3) $\Rightarrow$ (1) est le
lemme \ref{spaces-cohomology-lemma-Noetherian-h1-zero-invertible}.
\end{proof}

\begin{lemma}
\label{spaces-cohomology-lemma-surjective-finite-morphism-ample}
Soit $R$ un anneau noethérien. Soit $f : Y \to X$ un morphisme entre espaces
algébriques propres sur $R$. Soit $\mathcal{L}$ un
$\mathcal{O}_X$-module inversible. Supposons $f$ fini et surjectif.
Les conditions suivantes sont équivalentes :
\begin{enumerate}
\item $X$ est un schéma et $\mathcal{L}$ est ample ; et
\item $Y$ est un schéma et $f^*\mathcal{L}$ est ample.
\end{enumerate}
\end{lemma}

\begin{proof}
Supposons (1). Alors $Y$ est un schéma, car un morphisme fini est représentable
(par des schémas), voir Morphismes d'espaces, lemme
\ref{spaces-morphisms-lemma-integral-local}. Par suite, (2)
résulte de Cohomologie des schémas, lemme
\ref{coherent-lemma-surjective-finite-morphism-ample}.

\medskip\noindent
Supposons (2). Soit $P$ la propriété suivante des
$\mathcal{O}_X$-modules cohérents $\mathcal{F}$ : il existe $n_0$
tel que $H^p(X, \mathcal{F} \otimes \mathcal{L}^{\otimes n}) = 0$
pour tout $n \geq n_0$ et tout $p > 0$. Nous allons montrer que $P$ est vérifiée
par tout $\mathcal{O}_X$-module cohérent $\mathcal{F}$, ce qui entraîne que
$\mathcal{L}$ est ample d'après le lemme \ref{spaces-cohomology-lemma-vanshing-gives-ample}.
Nous allons appliquer le lemme \ref{spaces-cohomology-lemma-property-higher-rank-cohomological}.
Il faut donc vérifier les conditions (1), (2) et (3) de ce lemme pour $P$.
La propriété (1) résulte de la suite exacte longue de cohomologie associée
à une suite exacte courte de faisceaux et du fait que la tensorisation par
un faisceau inversible est un foncteur exact. La propriété (2) résulte de ce que
$H^p(X, -)$ est un foncteur additif.

\medskip\noindent
Pour vérifier (3), soit $i : Z \to X$ un sous-espace fermé réduit tel que $|Z|$
soit irréductible. Soient $i' : Z' \to Y$ et $f' : Z' \to Z$ comme dans le
lemme \ref{spaces-cohomology-lemma-finite-morphism-Noetherian}, et posons
$\mathcal{G} = f'_*\mathcal{O}_{Z'}$. Nous affirmons que
$\mathcal{G}$ vérifie les propriétés (3)(a) et (3)(b) du
lemme \ref{spaces-cohomology-lemma-property-higher-rank-cohomological},
ce qui achèvera la démonstration. La propriété (3)(a) a été établie au
lemme \ref{spaces-cohomology-lemma-finite-morphism-Noetherian}. Pour voir (3)(b), soit
$\mathcal{I}$ un faisceau quasi-cohérent d'idéaux non nul sur $Z$.
Notons $\mathcal{I}' \subset \mathcal{O}_{Z'}$ l'idéal quasi-cohérent
$(f')^{-1}\mathcal{I} \mathcal{O}_{Z'}$, c'est-à-dire
l'image de $(f')^*\mathcal{I} \to \mathcal{O}_{Z'}$.
D'après le lemme \ref{spaces-cohomology-lemma-affine-morphism-projection-ideal}, on a
$f_*\mathcal{I}' = \mathcal{I} \mathcal{G}$. Nous affirmons que la valeur commune
$\mathcal{G}' = \mathcal{I} \mathcal{G} = f'_*\mathcal{I}'$
vérifie la condition énoncée en (3)(b).
D'abord, il est clair que le support de $\mathcal{G}/\mathcal{G}'$
est contenu dans celui de $\mathcal{O}_Z/\mathcal{I}$,
qui est un sous-espace propre de $|Z|$, puisque $\mathcal{I}$ est un
faisceau d'idéaux non nul sur l'espace algébrique réduit et irréductible $Z$.
Rappelons que $f'_*$, $i_*$ et $i'_*$ transforment les modules cohérents en
modules cohérents, voir les lemmes \ref{spaces-cohomology-lemma-finite-pushforward-coherent} et
\ref{spaces-cohomology-lemma-i-star-equivalence}.
Puisque $Y$ est un schéma et que $\mathcal{L}$ est ample,
le lemme \ref{spaces-cohomology-lemma-vanshing-gives-ample} montre
qu'il existe $n_0$ tel que
$$
H^p(Y, i'_*\mathcal{I}' \otimes_{\mathcal{O}_Y} f^*\mathcal{L}^{\otimes n}) = 0
$$
pour $n \geq n_0$ et $p > 0$. On obtient alors
\begin{align*}
H^p(X, i_*\mathcal{G}' \otimes_{\mathcal{O}_X} \mathcal{L}^{\otimes n})
& =
H^p(Z, \mathcal{G'} \otimes_{\mathcal{O}_Z} i^*\mathcal{L}^{\otimes n}) \\
& =
H^p(Z,
f'_*\mathcal{I}' \otimes_{\mathcal{O}_Z} i^*\mathcal{L}^{\otimes n})) \\
& =
H^p(Z, f'_*(\mathcal{I}' \otimes_{\mathcal{O}_{Z'}}
(f')^*i^*\mathcal{L}^{\otimes n})) \\
& =
H^p(Z, f'_*(\mathcal{I}' \otimes_{\mathcal{O}_{Z'}}
(i')^*f^*\mathcal{L}^{\otimes n})) \\
& =
H^p(Z',
\mathcal{I}' \otimes_{\mathcal{O}_{Z'}} (i')^*f^*\mathcal{L}^{\otimes n})) \\
& =
H^p(Y, i'_*\mathcal{I}' \otimes_{\mathcal{O}_Y} f^*\mathcal{L}^{\otimes n}) = 0
\end{align*}
Nous avons utilisé ici la formule de projection et la suite spectrale de Leray
(voir Cohomologie sur les sites, sections
\ref{sites-cohomology-section-projection-formula} et
\ref{sites-cohomology-section-leray})
ainsi que le lemme \ref{spaces-cohomology-lemma-finite-higher-direct-image-zero}.
Cela vérifie la propriété (3)(b) du
lemme \ref{spaces-cohomology-lemma-property-higher-rank-cohomological}, comme voulu.
\end{proof}










\section{Le théorème des fonctions formelles}
\label{spaces-cohomology-section-theorem-formal-functions}

\noindent
Cette section est l'analogue de
Cohomologie des schémas, section \ref{coherent-section-theorem-formal-functions}.
Nous invitons le lecteur à lire d'abord cette section.

\begin{situation}
\label{spaces-cohomology-situation-formal-functions}
Ici, $A$ est un anneau noethérien et $I \subset A$ un idéal.
De plus, $f : X \to \Spec(A)$ est un morphisme propre d'espaces algébriques
et $\mathcal{F}$ un faisceau cohérent sur $X$.
\end{situation}

\noindent
Dans cette situation, on note $I^n\mathcal{F}$ le sous-module quasi-cohérent
de $\mathcal{F}$ engendré, comme $\mathcal{O}_X$-module,
par les produits de sections locales de $\mathcal{F}$ et d'éléments de $I^n$.
Autrement dit, c'est l'image de l'application
$f^*\widetilde{I} \otimes_{\mathcal{O}_X} \mathcal{F} \to \mathcal{F}$.

\begin{lemma}
\label{spaces-cohomology-lemma-cohomology-powers-ideal-times-F}
Dans la situation \ref{spaces-cohomology-situation-formal-functions}.
Posons $B = \bigoplus_{n \geq 0} I^n$.
Alors, pour tout $p \geq 0$, le $B$-module gradué
$\bigoplus_{n \geq 0} H^p(X, I^n\mathcal{F})$ est
un $B$-module de type fini.
\end{lemma}

\begin{proof}
Posons $\mathcal{B} = \bigoplus I^n\mathcal{O}_X = f^*\widetilde{B}$.
Alors $\bigoplus I^n\mathcal{F}$ est un $\mathcal{B}$-module
gradué de type fini. Le résultat découle donc
du lemme \ref{spaces-cohomology-lemma-graded-finiteness}.
\end{proof}

\begin{lemma}
\label{spaces-cohomology-lemma-cohomology-powers-ideal-application}
Dans la situation \ref{spaces-cohomology-situation-formal-functions}.
Pour tout $p \geq 0$, il existe un entier $c \geq 0$ tel que
\begin{enumerate}
\item l'application de multiplication
$I^{n - c} \otimes H^p(X, I^c\mathcal{F}) \to H^p(X, I^n\mathcal{F})$
soit surjective pour tout $n \geq c$ ; et
\item l'image de $H^p(X, I^{n + m}\mathcal{F}) \to H^p(X, I^n\mathcal{F})$
soit contenue dans le sous-module $I^{m - c} H^p(X, I^n\mathcal{F})$
pour tout $n \geq 0$ et tout $m \geq c$.
\end{enumerate}
\end{lemma}

\begin{proof}
D'après le lemme \ref{spaces-cohomology-lemma-cohomology-powers-ideal-times-F},
on peut trouver $d_1, \ldots, d_t \geq 0$ et
$x_i \in H^p(X, I^{d_i}\mathcal{F})$ tels que
$\bigoplus_{n \geq 0} H^p(X, I^n\mathcal{F})$ soit engendré
par $x_1, \ldots, x_t$ sur $B = \bigoplus_{n \geq 0} I^n$.
Posons $c = \max\{d_i\}$. Il est clair que (1) est vérifiée.
Pour (2), posons $b = \max(0, n - c)$.
Considérons le diagramme commutatif de $A$-modules
$$
\xymatrix{
I^{n + m - c - b} \otimes I^b \otimes
H^p(X, I^c\mathcal{F}) \ar[r] \ar[d] &
I^{n + m - c} \otimes H^p(X, I^c\mathcal{F}) \ar[r] &
H^p(X, I^{n + m}\mathcal{F}) \ar[d] \\
I^{n + m - c - b} \otimes H^p(X, I^n\mathcal{F}) \ar[rr] & &
H^p(X, I^n\mathcal{F})
}
$$
D'après la partie (1) du lemme, le composé des flèches horizontales
est surjectif si $n + m \geq c$. D'autre part, il est clair
que $n + m - c - b \geq m - c$. D'où la partie (2).
\end{proof}

\begin{lemma}
\label{spaces-cohomology-lemma-ML-cohomology-powers-ideal}
Dans la situation \ref{spaces-cohomology-situation-formal-functions}.
Fixons $p \geq 0$.
\begin{enumerate}
\item Il existe $c_1 \geq 0$ tel que, pour tout $n \geq c_1$,
on ait
$$
\Ker(
H^p(X, \mathcal{F}) \to H^p(X, \mathcal{F}/I^n\mathcal{F})
)
\subset
I^{n - c_1}H^p(X, \mathcal{F}).
$$
\item Le système projectif
$$
\left(H^p(X, \mathcal{F}/I^n\mathcal{F})\right)_{n \in \mathbf{N}}
$$
vérifie la condition de Mittag-Leffler (voir
Homologie, définition \ref{homology-definition-Mittag-Leffler}).
\item En fait, quels que soient $p$ et $n$, il existe $c_2(n) \geq n$
tel que
$$
\Im(H^p(X, \mathcal{F}/I^k\mathcal{F})
\to H^p(X, \mathcal{F}/I^n\mathcal{F}))
=
\Im(H^p(X, \mathcal{F})
\to H^p(X, \mathcal{F}/I^n\mathcal{F}))
$$
pour tout $k \geq c_2(n)$.
\end{enumerate}
\end{lemma}

\begin{proof}
Posons $c_1 = \max\{c_p, c_{p + 1}\}$, où $c_p, c_{p +1}$ sont les entiers
obtenus dans le lemme \ref{spaces-cohomology-lemma-cohomology-powers-ideal-application} pour
$H^p$ et $H^{p + 1}$. Nous utiliserons cette constante dans les démonstrations de
(1), (2) et (3).

\medskip\noindent
Démontrons la partie (1). Considérons la suite exacte courte
$$
0 \to I^n\mathcal{F} \to \mathcal{F} \to \mathcal{F}/I^n\mathcal{F} \to 0
$$
La suite exacte longue de cohomologie donne
$$
\Ker(
H^p(X, \mathcal{F}) \to H^p(X, \mathcal{F}/I^n\mathcal{F})
)
=
\Im(
H^p(X, I^n\mathcal{F}) \to H^p(X, \mathcal{F})
)
$$
D'après le choix de $c_1$, ce sous-module est donc contenu dans
$I^{n - c_1}H^p(X, \mathcal{F})$ pour $n \geq c_1$.

\medskip\noindent
Notons que la partie (3) entraîne la partie (2) par définition de la condition
de Mittag-Leffler.

\medskip\noindent
Démontrons la partie (3).
Fixons $n$ pour le reste de la démonstration.
Considérons le diagramme commutatif
$$
\xymatrix{
0 \ar[r] &
I^n\mathcal{F} \ar[r] &
\mathcal{F} \ar[r] &
\mathcal{F}/I^n\mathcal{F} \ar[r] &
0 \\
0 \ar[r] &
I^{n + m}\mathcal{F} \ar[r] \ar[u] &
\mathcal{F} \ar[r] \ar[u] &
\mathcal{F}/I^{n + m}\mathcal{F} \ar[r] \ar[u] &
0
}
$$
Il induit le diagramme commutatif suivant
$$
\xymatrix{
H^p(X, I^n\mathcal{F}) \ar[r] &
H^p(X, \mathcal{F}) \ar[r] &
H^p(X, \mathcal{F}/I^n\mathcal{F}) \ar[r]_\delta &
H^{p + 1}(X, I^n\mathcal{F}) \\
H^p(X, I^{n + m}\mathcal{F}) \ar[r] \ar[u] &
H^p(X, \mathcal{F}) \ar[r] \ar[u]^1 &
H^p(X, \mathcal{F}/I^{n + m}\mathcal{F}) \ar[r] \ar[u] &
H^{p + 1}(X, I^{n + m}\mathcal{F}) \ar[u]^a
}
$$
Si $m \geq c_1$, on voit que l'image de $a$ est
contenue dans $I^{m - c_1} H^{p + 1}(X, I^n\mathcal{F})$.
D'après le lemme d'Artin--Rees (voir Algèbre, lemme \ref{algebra-lemma-map-AR}),
il existe un entier $c_3(n)$ tel que
$$
I^N H^{p + 1}(X, I^n\mathcal{F}) \cap \Im(\delta)
\subset
\delta\left(I^{N - c_3(n)}H^p(X, \mathcal{F}/I^n\mathcal{F})\right)
$$
pour tout $N \geq c_3(n)$. Comme $H^p(X, \mathcal{F}/I^n\mathcal{F})$
est annulé par $I^n$, on voit que, si $m \geq c_3(n) + c_1 + n$,
alors
$$
\Im(H^p(X, \mathcal{F}/I^{n + m}\mathcal{F})
\to H^p(X, \mathcal{F}/I^n\mathcal{F}))
=
\Im(H^p(X, \mathcal{F})
\to H^p(X, \mathcal{F}/I^n\mathcal{F}))
$$
Autrement dit, la partie (3) est vérifiée avec $c_2(n) = c_3(n) + c_1 + n$.
\end{proof}

\begin{theorem}[Théorème des fonctions formelles]
\label{spaces-cohomology-theorem-formal-functions}
Dans la situation \ref{spaces-cohomology-situation-formal-functions}. Fixons $p \geq 0$.
Le système d'applications
$$
H^p(X, \mathcal{F})/I^nH^p(X, \mathcal{F})
\longrightarrow
H^p(X, \mathcal{F}/I^n\mathcal{F})
$$
définit un isomorphisme de limites
$$
H^p(X, \mathcal{F})^\wedge
\longrightarrow
\lim_n H^p(X, \mathcal{F}/I^n\mathcal{F})
$$
où le membre de gauche est le complété du $A$-module
$H^p(X, \mathcal{F})$ relativement à l'idéal $I$, voir
Algèbre, section \ref{algebra-section-completion}.
De plus, c'est un homéomorphisme pour les topologies de limite projective.
\end{theorem}

\begin{proof}
Cela résulte en fait immédiatement du
lemme \ref{spaces-cohomology-lemma-ML-cohomology-powers-ideal}. Donnons les détails.
Posons $M = H^p(X, \mathcal{F})$ et $M_n = H^p(X, \mathcal{F}/I^n\mathcal{F})$.
Notons $N_n = \Im(M \to M_n)$.
D'après la description de la limite dans Homologie, section
\ref{homology-section-inverse-systems}, on a
$$
\lim_n M_n
=
\{(x_n) \in \prod M_n \mid \varphi_i(x_n) = x_{n - 1}, \ n = 2, 3, \ldots\}
$$
Choisissons un élément $x = (x_n) \in \lim_n M_n$.
D'après la partie (3) du lemme \ref{spaces-cohomology-lemma-ML-cohomology-powers-ideal},
on a $x_n \in N_n$ pour tout $n$, puisque, par
définition, $x_n$ est l'image d'un certain $x_{n + m} \in M_{n + m}$ pour
tout $m$. D'après la partie (1) du lemme \ref{spaces-cohomology-lemma-ML-cohomology-powers-ideal},
il existe une factorisation
$$
M \to N_n \to M/I^{n - c_1}M
$$
de l'application de réduction. Notons $y_n \in M/I^{n - c_1}M$ l'image de $x_n$
pour $n \geq c_1$. Puisque, pour $n' \geq n$, le composé
$M \to M_{n'} \to M_n$ est l'application donnée $M \to M_n$, on voit que
$y_{n'}$ s'envoie sur $y_n$ par l'application canonique
$M/I^{n' - c_1}M \to M/I^{n - c_1}M$. Ainsi, $y = (y_{n + c_1})$
définit un élément de $\lim_n M/I^nM$.
Nous omettons de vérifier que $y$ s'envoie sur $x$ par
l'application
$$
M^\wedge = \lim_n M/I^nM \longrightarrow \lim_n M_n
$$
du lemme. Nous omettons également la vérification concernant les topologies.
\end{proof}

\begin{lemma}
\label{spaces-cohomology-lemma-spell-out-theorem-formal-functions}
Soit $A$ un anneau. Soit $I \subset A$ un idéal. Supposons $A$
noethérien et complet pour la topologie $I$-adique.
Soit $f : X \to \Spec(A)$ un morphisme propre d'espaces algébriques.
Soit $\mathcal{F}$ un faisceau cohérent sur $X$. Alors
$$
H^p(X, \mathcal{F}) = \lim_n H^p(X, \mathcal{F}/I^n\mathcal{F})
$$
pour tout $p \geq 0$.
\end{lemma}

\begin{proof}
C'est une reformulation du théorème des fonctions formelles
(théorème \ref{spaces-cohomology-theorem-formal-functions}) dans le
cas d'un anneau de base noethérien complet. En effet, dans ce cas, le
$A$-module $H^p(X, \mathcal{F})$ est de type fini
(lemme \ref{spaces-cohomology-lemma-proper-over-affine-cohomology-finite}), donc
complet pour la topologie $I$-adique (Algèbre, lemme \ref{algebra-lemma-completion-tensor}),
et l'on voit qu'il est inutile de compléter le membre de gauche.
\end{proof}

\begin{lemma}
\label{spaces-cohomology-lemma-formal-functions-stalk}
Soit $S$ un schéma. Soit $f : X \to Y$ un morphisme d'espaces algébriques
au-dessus de $S$, et soit $\mathcal{F}$ un faisceau quasi-cohérent sur $X$. Supposons que
\begin{enumerate}
\item $Y$ soit localement noethérien,
\item $f$ soit propre, et
\item $\mathcal{F}$ soit cohérent.
\end{enumerate}
Soit $\overline{y}$ un point géométrique de $Y$.
Considérons les « voisinages infinitésimaux »
$$
\xymatrix{
X_n =
\Spec(\mathcal{O}_{Y, \overline{y}}/\mathfrak m_{\overline{y}}^n) \times_Y X
\ar[r]_-{i_n} \ar[d]_{f_n} &
X \ar[d]^f \\
\Spec(\mathcal{O}_{Y, \overline{y}}/\mathfrak m_{\overline{y}}^n)
\ar[r]^-{c_n} & Y
}
$$
de la fibre $X_1 = X_{\overline{y}}$, et posons
$\mathcal{F}_n = i_n^*\mathcal{F}$. On a alors
$$
\left(R^pf_*\mathcal{F}\right)_{\overline{y}}^\wedge
\cong
\lim_n H^p(X_n, \mathcal{F}_n)
$$
comme $\mathcal{O}_{Y, \overline{y}}^\wedge$-modules.
\end{lemma}

\begin{proof}
Il s'agit simplement d'une reformulation d'un cas particulier du théorème
des fonctions formelles, théorème \ref{spaces-cohomology-theorem-formal-functions}.
Donnons les détails. Notons que $\mathcal{O}_{Y, \overline{y}}$
est un anneau local noethérien, voir
Propriétés des espaces, lemme
\ref{spaces-properties-lemma-Noetherian-local-ring-Noetherian}.
Considérons le morphisme canonique
$c : \Spec(\mathcal{O}_{Y, \overline{y}}) \to Y$.
C'est un morphisme plat, car il identifie les anneaux locaux.
Notons $f' : X' \to \Spec(\mathcal{O}_{Y, \overline{y}})$
le changement de base de $f$ à cet anneau local. On a
$c^*R^pf_*\mathcal{F} = R^pf'_*\mathcal{F}'$ d'après le
lemme \ref{spaces-cohomology-lemma-flat-base-change-cohomology}.
De plus, on dispose d'identifications canoniques $X_n = X'_n$
pour tout $n \geq 1$.

\medskip\noindent
On peut donc supposer que $Y = \Spec(A)$ est le spectre d'un
anneau local noethérien strictement hensélien $A$, d'idéal maximal
$\mathfrak m$, et que $\overline{y} \to Y$ est égal à
$\Spec(A/\mathfrak m) \to Y$. Il en résulte que
$$
\left(R^pf_*\mathcal{F}\right)_{\overline{y}} =
\Gamma(Y, R^pf_*\mathcal{F}) =
H^p(X, \mathcal{F})
$$
car $(Y, \overline{y})$ est un objet initial de la catégorie
des voisinages \'etales de $\overline{y}$.
Les morphismes $c_n$ sont tous des immersions fermées.
Leurs changements de base $i_n$ sont donc eux aussi des immersions fermées.
Notons que $i_{n, *}\mathcal{F}_n = i_{n, *}i_n^*\mathcal{F}
= \mathcal{F}/\mathfrak m^n\mathcal{F}$. La suite spectrale de Leray
pour $i_n$ et le lemme \ref{spaces-cohomology-lemma-finite-pushforward-coherent} donnent
$$
H^p(X_n, \mathcal{F}_n) =
H^p(X, i_{n, *}\mathcal{F}) =
H^p(X, \mathcal{F}/\mathfrak m^n\mathcal{F})
$$
On peut donc bien appliquer le théorème des fonctions formelles pour calculer
la limite figurant dans l'énoncé du lemme, ce qui conclut.
\end{proof}

\noindent
Voici un lemme que nous généraliserons ensuite aux fibres de
dimension $ > 0$, dans le lemme suivant.

\begin{lemma}
\label{spaces-cohomology-lemma-higher-direct-images-zero-finite-fibre}
Soit $S$ un schéma. Soit $f : X \to Y$ un morphisme d'espaces algébriques
au-dessus de $S$. Soit $\overline{y}$ un point géométrique de $Y$.
Supposons que
\begin{enumerate}
\item $Y$ soit localement noethérien,
\item $f$ soit propre, et
\item l'espace topologique sous-jacent à $X_{\overline{y}}$ soit discret.
\end{enumerate}
Alors, pour tout faisceau cohérent $\mathcal{F}$ sur $X$, on a
$(R^pf_*\mathcal{F})_{\overline{y}} = 0$ pour tout $p > 0$.
\end{lemma}

\begin{proof}
Soit $\kappa(\overline{y})$ le corps résiduel de l'anneau
local $\mathcal{O}_{Y, \overline{y}}$. Comme dans le
lemme \ref{spaces-cohomology-lemma-formal-functions-stalk},
posons $X_{\overline{y}} = X_1 = \Spec(\kappa(\overline{y})) \times_Y X$.
D'après Morphismes d'espaces, lemme
\ref{spaces-morphisms-lemma-quasi-finite-at-point}
le morphisme $f : X \to Y$ est quasi-fini en chacun
des points de la fibre de $X \to Y$ au-dessus de $\overline{y}$.
Il en résulte que $X_{\overline{y}} \to \overline{y}$ est séparé et
quasi-fini. Par suite, $X_{\overline{y}}$ est un schéma
d'après Morphismes d'espaces, proposition
\ref{spaces-morphisms-proposition-locally-quasi-finite-separated-over-scheme}.
Comme il est quasi-compact, son espace topologique sous-jacent est un espace
discret fini. C'est donc un schéma affine d'après
Schémas, lemme \ref{schemes-lemma-scheme-finite-discrete-affine}.
Le lemme \ref{spaces-cohomology-lemma-image-affine-finite-morphism-affine-Noetherian}
montre alors que les espaces algébriques $X_n$ sont eux aussi des schémas affines.
De plus, l'espace topologique sous-jacent à chaque $X_n$ est le même
que celui de $X_1$. Il s'ensuit que $H^p(X_n, \mathcal{F}_n) = 0$
pour tout $p > 0$. Ainsi,
$(R^pf_*\mathcal{F})_{\overline{y}}^\wedge = 0$
d'après le lemme \ref{spaces-cohomology-lemma-formal-functions-stalk}.
Notons que $R^pf_*\mathcal{F}$ est cohérent d'après le
lemme \ref{spaces-cohomology-lemma-proper-pushforward-coherent} ;
ainsi $R^pf_*\mathcal{F}_{\overline{y}}$ est un
$\mathcal{O}_{Y, \overline{y}}$-module de type fini.
D'après Algèbre, lemme \ref{algebra-lemma-completion-tensor},
cela implique que $(R^pf_*\mathcal{F})_{\overline{y}} = 0$.
\end{proof}

\begin{lemma}
\label{spaces-cohomology-lemma-higher-direct-images-zero-above-dimension-fibre}
\begin{slogan}
Pour les morphismes propres, les fibres des images directes supérieures sont nulles aux degrés
strictement supérieurs à la dimension de la fibre.
\end{slogan}
Soit $S$ un schéma.
Soit $f : X \to Y$ un morphisme d'espaces algébriques au-dessus de $S$.
Soit $\overline{y}$ un point géométrique de $Y$.
Supposons que
\begin{enumerate}
\item $Y$ soit localement noethérien,
\item $f$ soit propre, et
\item $\dim(X_{\overline{y}}) = d$.
\end{enumerate}
Alors, pour tout faisceau cohérent $\mathcal{F}$ sur $X$, on a
$(R^pf_*\mathcal{F})_{\overline{y}} = 0$ pour tout $p > d$.
\end{lemma}

\begin{proof}
Soit $\kappa(\overline{y})$ le corps résiduel de l'anneau
local $\mathcal{O}_{Y, \overline{y}}$. Comme dans le
lemme \ref{spaces-cohomology-lemma-formal-functions-stalk},
posons $X_{\overline{y}} = X_1 = \Spec(\kappa(\overline{y})) \times_Y X$.
De plus, l'espace topologique sous-jacent à chaque voisinage
infinitésimal $X_n$ est le même que celui de $X_{\overline{y}}$.
Ainsi, $H^p(X_n, \mathcal{F}_n) = 0$ pour tout $p > d$ d'après le
lemme \ref{spaces-cohomology-lemma-vanishing-above-dimension}.
Ainsi, on voit que $(R^pf_*\mathcal{F})_{\overline{y}}^\wedge = 0$
d'après le lemme \ref{spaces-cohomology-lemma-formal-functions-stalk} pour $p > d$.
Notons que $R^pf_*\mathcal{F}$ est cohérent d'après le
lemme \ref{spaces-cohomology-lemma-proper-pushforward-coherent} ;
ainsi $R^pf_*\mathcal{F}_{\overline{y}}$ est un
$\mathcal{O}_{Y, \overline{y}}$-module de type fini.
D'après Algèbre, lemme \ref{algebra-lemma-completion-tensor},
cela implique que $(R^pf_*\mathcal{F})_{\overline{y}} = 0$.
\end{proof}






\section{Applications du théorème des fonctions formelles}
\label{spaces-cohomology-section-applications-formal-functions}

\noindent
Nous compléterons cette section selon les besoins.

\begin{lemma}
\label{spaces-cohomology-lemma-characterize-finite}
(Pour une version plus générale, voir
Compléments sur les morphismes d'espaces, lemme
\ref{spaces-more-morphisms-lemma-characterize-finite}).
Soit $S$ un schéma.
Soit $f : X \to Y$ un morphisme d'espaces algébriques au-dessus de $S$.
Supposons $Y$ localement noethérien.
Les conditions suivantes sont équivalentes :
\begin{enumerate}
\item $f$ est fini ; et
\item $f$ est propre et $|X_k|$ est un espace discret
pour tout morphisme $\Spec(k) \to Y$, où $k$ est un corps.
\end{enumerate}
\end{lemma}

\begin{proof}
Un morphisme fini est propre d'après
Morphismes d'espaces, lemme \ref{spaces-morphisms-lemma-finite-proper}.
Un morphisme fini est quasi-fini d'après
Morphismes d'espaces, lemme \ref{spaces-morphisms-lemma-finite-quasi-finite}.
Un morphisme quasi-fini a des fibres discrètes $X_k$, voir
Morphismes d'espaces, lemme \ref{spaces-morphisms-lemma-locally-quasi-finite}.
Ainsi, un morphisme fini est propre et a des fibres discrètes $X_k$.

\medskip\noindent
Supposons $f$ propre, à fibres discrètes $X_k$.
Nous voulons montrer que $f$ est fini.
Il suffit en fait de montrer que $f$ est affine.
En effet, si $f$ est affine, il s'ensuit que
$f$ est entier d'après
Morphismes d'espaces, lemme
\ref{spaces-morphisms-lemma-integral-universally-closed}
puis, d'après
Morphismes d'espaces, lemme \ref{spaces-morphisms-lemma-finite-integral},
que $f$ est fini.

\medskip\noindent
Pour montrer que $f$ est affine, on peut supposer $Y$ affine ; notre
but est de montrer que $X$ est également affine.
Puisque $f$ est propre, $X$ est séparé et quasi-compact.
Nous allons montrer que, pour tout
$\mathcal{O}_X$-module cohérent $\mathcal{F}$, on a $H^1(X, \mathcal{F}) = 0$.
Cela implique que $H^1(X, \mathcal{F}) = 0$ pour tout
$\mathcal{O}_X$-module quasi-cohérent $\mathcal{F}$ d'après les
lemmes \ref{spaces-cohomology-lemma-directed-colimit-coherent} et \ref{spaces-cohomology-lemma-colimits}.
On en déduit que $X$ est affine d'après la
proposition \ref{spaces-cohomology-proposition-vanishing-affine}. D'après le
lemme \ref{spaces-cohomology-lemma-higher-direct-images-zero-finite-fibre},
on conclut que les fibres de $R^1f_*\mathcal{F}$
sont nulles en tout point géométrique de $Y$.
Autrement dit, $R^1f_*\mathcal{F} = 0$. La suite spectrale de Leray
pour $f$ donne donc
$H^1(X , \mathcal{F}) = H^1(Y, f_*\mathcal{F})$.
Puisque $Y$ est affine et $f_*\mathcal{F}$ quasi-cohérent
(Morphismes d'espaces, lemme \ref{spaces-morphisms-lemma-pushforward}),
on conclut que $H^1(Y, f_*\mathcal{F}) = 0$
d'après
Cohomologie des schémas, lemme
\ref{coherent-lemma-quasi-coherent-affine-cohomology-zero}.
Ainsi, $H^1(X, \mathcal{F}) = 0$, comme voulu.
\end{proof}

\noindent
On en déduit le résultat utile suivant.

\begin{lemma}
\label{spaces-cohomology-lemma-proper-finite-fibre-finite-in-neighbourhood}
(Pour une version plus générale, voir
Compléments sur les morphismes d'espaces, lemme
\ref{spaces-more-morphisms-lemma-proper-finite-fibre-finite-in-neighbourhood}).
Soit $S$ un schéma. Soit $f : X \to Y$ un morphisme d'espaces algébriques
au-dessus de $S$. Soit $\overline{y}$ un point géométrique de $Y$.
Supposons que
\begin{enumerate}
\item $Y$ soit localement noethérien,
\item $f$ soit propre, et
\item $|X_{\overline{y}}|$ soit fini.
\end{enumerate}
Alors il existe un voisinage ouvert $V \subset Y$ de $\overline{y}$
tel que $f|_{f^{-1}(V)} : f^{-1}(V) \to V$ soit fini.
\end{lemma}

\begin{proof}
Le morphisme $f$ est quasi-fini en tous les points géométriques de $X$
au-dessus de $\overline{y}$ d'après
Morphismes d'espaces, lemme \ref{spaces-morphisms-lemma-quasi-finite-at-point}.
D'après Morphismes d'espaces, lemme
\ref{spaces-morphisms-lemma-locally-finite-type-quasi-finite-part}, l'ensemble
des points où $f$ est quasi-fini est un sous-espace ouvert $U \subset X$.
Posons $Z = X \setminus U$. Alors $\overline{y} \not \in f(Z)$. Puisque $f$
est propre, l'ensemble $f(Z) \subset Y$ est fermé. Choisissons un voisinage ouvert
$V \subset Y$ de $\overline{y}$ tel que $Z \cap V = \emptyset$. Alors
$f^{-1}(V) \to V$ est localement quasi-fini et propre.
Ainsi, $f^{-1}(V) \to V$ a des fibres discrètes $X_k$
(Morphismes d'espaces, lemme
\ref{spaces-morphisms-lemma-locally-quasi-finite})
qui sont quasi-compactes, donc finies.
Par conséquent, $f^{-1}(V) \to V$
est fini d'après le lemme \ref{spaces-cohomology-lemma-characterize-finite}.
\end{proof}








% OK, start here.
%

\chapter{Limites des espaces algébriques}



\phantomsection
\label{spaces-limits-section-phantom}


\section{Introduction}
\label{spaces-limits-section-introduction}

\noindent
Dans ce chapitre, nous rassemblons des résultats relatifs aux limites des espaces algébriques.
Un premier thème est la caractérisation des espaces algébriques $F$ localement
de présentation finie sur la base $S$ comme foncteurs compatibles aux limites.
Nous poursuivons par l'étude des limites de systèmes projectifs indexés par des
ensembles ordonnés filtrants (Catégories, définition \ref{categories-definition-directed-set})
dont les morphismes de transition sont affines. Nous discutons l'approximation
noethérienne absolue des espaces algébriques quasi-compacts et quasi-séparés,
suivant \cite{CLO}. Une autre approche est due à David Rydh (voir
\cite{rydh_approx}); ses résultats couvrent également l'approximation noethérienne
absolue de certains champs algébriques.


\section{Conventions}
\label{spaces-limits-section-conventions}

\noindent
L'hypothèse générale est que tous les schémas sont contenus dans
un grand site fppf $\Sch_{fppf}$. Tous les anneaux $A$ considérés
ont la propriété que $\Spec(A)$ est (isomorphe à) un
objet de ce grand site.

\medskip\noindent
Soit $S$ un schéma et soit $X$ un espace algébrique sur $S$.
Dans ce chapitre et dans le suivant, nous écrirons $X \times_S X$
pour le produit de $X$ par lui-même (dans la catégorie des espaces algébriques
sur $S$), au lieu de $X \times X$.











\section{Morphismes de présentation finie}
\label{spaces-limits-section-finite-presentation}

\noindent
Dans cette section, nous généralisons
Limites, proposition
\ref{limits-proposition-characterize-locally-finite-presentation}
aux morphismes d'espaces algébriques.
La définition suivante est motivée par
la proposition que nous venons de citer.

\begin{definition}
\label{spaces-limits-definition-locally-finite-presentation}
Soit $S$ un schéma.
\begin{enumerate}
\item Un foncteur $F : (\Sch/S)_{fppf}^{opp} \to \textit{Ens}$
est dit {\it compatible aux limites} ou {\it localement de présentation finie} si
pour tout schéma affine $T$ sur $S$ qui est une limite $T = \lim T_i$
d'un système projectif filtrant de schémas affines $T_i$ sur $S$, on a
$$
F(T) = \colim F(T_i).
$$
Nous disons parfois que $F$ est {\it localement de présentation finie sur $S$}.
\item Soient $F, G : (\Sch/S)_{fppf}^{opp} \to \textit{Ens}$.
Une transformation de foncteurs $a : F \to G$
est {\it compatible aux limites} ou {\it localement de présentation finie}
si, pour tout schéma $T$ sur $S$ et tout $y \in G(T)$, le foncteur
$$
F_y : (\Sch/T)_{fppf}^{opp} \longrightarrow \textit{Ens}, \quad
T'/T \longmapsto \{x \in F(T') \mid a(x) = y|_{T'}\}
$$
est localement de présentation finie sur $T$\footnote{La caractérisation (2) du
lemme \ref{spaces-limits-lemma-characterize-relative-limit-preserving}
est peut-être plus facile à lire.}. Nous disons parfois que
$F$ est {\it compatible aux limites relativement à $G$}.
\end{enumerate}
\end{definition}

\noindent
Le foncteur $F_y$ est en un certain sens la fibre de
$a : F \to G$ au-dessus de $y$, sauf qu'il s'agit d'un préfaisceau sur le grand site fppf
de $T$. Une formule pour ce foncteur est la suivante :
\begin{equation}
\label{spaces-limits-equation-fibre-map-functors}
F_y =
F|_{(\Sch/T)_{fppf}}
{\times}_{G|_{(\Sch/T)_{fppf}}}
*
\end{equation}
Ici, $*$ est l'objet final de la catégorie des (pré)faisceaux
sur $(\Sch/T)_{fppf}$ (voir
Sites, exemple \ref{sites-example-singleton-sheaf})
et le morphisme $* \to G|_{(\Sch/T)_{fppf}}$ est donné par $y$.
Notons que si $j : (\Sch/T)_{fppf} \to (\Sch/S)_{fppf}$
est le foncteur de localisation, alors la formule ci-dessus devient
$F_y = j^{-1}F \times_{j^{-1}G} *$ et $j_!F_y$ est simplement le produit fibré
$F \times_{G, y} T$. (Voir
Sites, section \ref{sites-section-localize},
pour des informations sur la localisation, et en particulier
Sites, remarque \ref{sites-remark-localize-presheaves}
pour des informations sur $j_!$ pour les préfaisceaux.)

\medskip\noindent
À ce stade, nous disposons temporairement de deux définitions de ce que signifie,
pour un morphisme $X \to Y$ d'espaces algébriques sur $S$, être localement de présentation
finie. L'une est donnée par
Morphismes d'espaces,
définition \ref{spaces-morphisms-definition-locally-finite-presentation}
et l'autre utilise le fait que $X \to Y$ est une transformation de foncteurs, de sorte que
la définition \ref{spaces-limits-definition-locally-finite-presentation}
s'applique (nous emploierons autant que possible la terminologie « compatible aux limites »
pour cette notion). Nous montrerons dans
la proposition \ref{spaces-limits-proposition-characterize-locally-finite-presentation}
que ces deux définitions coïncident.

\begin{lemma}
\label{spaces-limits-lemma-characterize-relative-limit-preserving}
Soit $S$ un schéma. Soit $a : F \to G$ une transformation de foncteurs
$(\Sch/S)_{fppf}^{opp} \to \textit{Ens}$.
Les conditions suivantes sont équivalentes :
\begin{enumerate}
\item $a : F \to G$ est compatible aux limites ;
\item pour tout schéma affine $T$ sur $S$ qui est une
limite $T = \lim T_i$ d'un système projectif filtrant de
schémas affines $T_i$ sur $S$, le diagramme d'ensembles
$$
\xymatrix{
\colim_i F(T_i) \ar[r] \ar[d]_a & F(T) \ar[d]^a \\
\colim_i G(T_i) \ar[r] & G(T)
}
$$
est cartésien.
\end{enumerate}
\end{lemma}

\begin{proof}
Supposons (1). Considérons $T = \lim_{i \in I} T_i$ comme dans (2). Soit
$(y, x_T)$ un élément du produit fibré
$\colim_i G(T_i) \times_{G(T)} F(T)$.
Alors $y$ provient d'un $y_i \in G(T_i)$ pour un certain $i$.
Considérons le foncteur $F_{y_i}$ sur $(\Sch/T_i)_{fppf}$ comme dans
la définition \ref{spaces-limits-definition-locally-finite-presentation}.
On voit que $x_T \in F_{y_i}(T)$. De plus, $T = \lim_{i' \geq i} T_{i'}$
est un système projectif filtrant de schémas affines sur $T_i$. Ainsi (1) implique
que $x_T$ est l'image d'un unique élément $x$ de
$\colim_{i' \geq i} F_{y_i}(T_{i'})$. Ainsi $x$ est l'unique
élément de $\colim F(T_i)$ ayant pour image le couple $(y, x_T)$.
Ceci prouve (2).

\medskip\noindent
Supposons (2). Soit $T$ un schéma et $y_T \in G(T)$. Nous devons montrer que
$F_{y_T}$ est compatible aux limites. Soit $T' = \lim_{i \in I} T'_i$ un
schéma affine sur $T$, limite d'un système projectif filtrant de schémas affines $T'_i$
sur $T$. Soit $x_{T'} \in F_{y_T}$. Choisissons $i \in I$, ce qui est possible car
$I$ est un ensemble ordonné filtrant. Notons $y_i \in F(T'_i)$ l'image
de $y_{T'}$. Alors $(y_i, x_{T'})$ est un
élément du produit fibré
$\colim_i G(T'_i) \times_{G(T')} F(T')$.
Par (2), il existe donc un unique élément $x$ de $\colim_i F(T'_i)$
ayant pour image $(y_i, x_{T'})$. Il est clair que $x$ définit un élément
de $\colim_i F_y(T'_i)$ ayant pour image $x_{T'}$, ce qui conclut.
\end{proof}

\begin{lemma}
\label{spaces-limits-lemma-composition-locally-finite-presentation}
Soit $S$ un schéma contenu dans $\Sch_{fppf}$.
Soient $F, G, H : (\Sch/S)_{fppf}^{opp} \to \textit{Ens}$.
Soient $a : F \to G$, $b : G \to H$ des transformations de foncteurs.
Si $a$ et $b$ sont compatibles aux limites, alors
$$
b \circ a : F \longrightarrow H
$$
est compatible aux limites.
\end{lemma}

\begin{proof}
Soit $T = \lim_{i \in I} T_i$ comme dans la caractérisation (2) du
lemme \ref{spaces-limits-lemma-characterize-relative-limit-preserving}.
Considérons le diagramme d'ensembles
$$
\xymatrix{
\colim_i F(T_i) \ar[r] \ar[d]_a & F(T) \ar[d]^a \\
\colim_i G(T_i) \ar[r] \ar[d]_b & G(T) \ar[d]^b \\
\colim_i H(T_i) \ar[r] & H(T)
}
$$
Par hypothèse, les deux carrés sont cartésiens. Le
rectangle extérieur est donc lui aussi cartésien, ce qui prouve le lemme.
\end{proof}

\begin{lemma}
\label{spaces-limits-lemma-locally-finite-presentation-permanence}
Soit $S$ un schéma contenu dans $\Sch_{fppf}$.
Soient $F, G, H : (\Sch/S)_{fppf}^{opp} \to \textit{Ens}$.
Soient $a : F \to G$, $b : G \to H$ des transformations de foncteurs.
Si $b \circ a$ et $b$ sont compatibles aux limites, alors $a$
est compatible aux limites.
\end{lemma}

\begin{proof}
Soit $T = \lim_{i \in I} T_i$ comme dans la caractérisation (2) du
lemme \ref{spaces-limits-lemma-characterize-relative-limit-preserving}.
Considérons le diagramme d'ensembles
$$
\xymatrix{
\colim_i F(T_i) \ar[r] \ar[d]_a & F(T) \ar[d]^a \\
\colim_i G(T_i) \ar[r] \ar[d]_b & G(T) \ar[d]^b \\
\colim_i H(T_i) \ar[r] & H(T)
}
$$
Par hypothèse, le carré inférieur et le rectangle extérieur
sont cartésiens. Le carré supérieur
est donc lui aussi cartésien, ce qui prouve le lemme.
\end{proof}

\begin{lemma}
\label{spaces-limits-lemma-base-change-locally-finite-presentation}
Soit $S$ un schéma contenu dans $\Sch_{fppf}$.
Soient $F, G, H : (\Sch/S)_{fppf}^{opp} \to \textit{Ens}$.
Soient $a : F \to G$, $b : H \to G$ des transformations de foncteurs.
Considérons le diagramme cartésien
$$
\xymatrix{
H \times_{b, G, a} F \ar[r]_-{b'} \ar[d]_{a'} & F \ar[d]^a \\
H \ar[r]^b & G
}
$$
Si $a$ est compatible aux limites, alors le changement de base $a'$ l'est aussi.
\end{lemma}

\begin{proof}
Preuve omise. Indication : c'est formel.
\end{proof}

\begin{lemma}
\label{spaces-limits-lemma-fibre-product-locally-finite-presentation}
Soit $S$ un schéma contenu dans $\Sch_{fppf}$.
Soient $E, F, G, H : (\Sch/S)_{fppf}^{opp} \to \textit{Ens}$.
Soient $a : F \to G$, $b : H \to G$ et $c : G \to E$
des transformations de foncteurs. Si $c$, $c \circ a$ et $c \circ b$
sont compatibles aux limites, alors $F \times_G H \to E$ l'est aussi.
\end{lemma}

\begin{proof}
Soit $T = \lim_{i \in I} T_i$ comme dans la caractérisation (2) du
lemme \ref{spaces-limits-lemma-characterize-relative-limit-preserving}.
Nous avons alors
$$
\colim (F \times_G H)(T_i) =
\colim F(T_i) \times_{\colim G(T_i)} \colim H(T_i)
$$
car les limites inductives filtrantes commutent aux produits finis. Notre but est donc de
montrer que
$$
\xymatrix{
\colim F(T_i) \times_{\colim G(T_i)} \colim H(T_i) \ar[r] \ar[d] &
F(T) \times_{G(T)} H(T) \ar[d] \\
\colim_i E(T_i) \ar[r] & E(T)
}
$$
est cartésien. Cela découle de l'observation que,
pour des applications d'ensembles $E' \to E$, $F \to G$, $H \to G$ et $G \to E$,
nous avons
$$
E' \times_E (F \times_G H) =
(E' \times_E F) \times_{(E' \times_E G)} (E' \times_E H)
$$
Quelques détails sont omis.
\end{proof}

\begin{lemma}
\label{spaces-limits-lemma-sheafify-finite-presentation}
Soit $S$ un schéma contenu dans $\Sch_{fppf}$.
Soit $F : (\Sch/S)_{fppf}^{opp} \to \textit{Ens}$ un foncteur.
Si $F$ est compatible aux limites, alors le faisceau associé $F^\#$ l'est aussi.
\end{lemma}

\begin{proof}
Supposons que $F$ soit compatible aux limites.
Il suffit de montrer que $F^+$ est compatible aux limites, puisque
$F^\# = (F^+)^+$, voir
Sites, théorème \ref{sites-theorem-plus}.
Soit $T$ un schéma affine sur $S$, et écrivons $T = \lim T_i$
comme limite d'un système projectif filtrant de schémas affines sur $S$.
Rappelons que $F^+(T)$ est la limite inductive des $\check H^0(\mathcal{V}, F)$
où cette limite inductive porte sur tous les recouvrements de $T$ dans $(\Sch/S)_{fppf}$.
Tout recouvrement fppf d'un schéma affine admet un raffinement qui est un recouvrement
fppf standard, voir
Topologies, lemme \ref{topologies-lemma-fppf-affine}.
Nous pouvons donc écrire
$$
F^+(T)
=
\colim_{\mathcal{V}\text{ recouvrement standard de }T}
\check H^0(\mathcal{V}, F).
$$
Tout recouvrement $\mathcal{V} = \{T_k \to T\}_{k = 1, \ldots, n}$
intervenant dans cette limite inductive peut s'écrire
$\mathcal{V}_i \times_{T_i} T$ pour un certain $i$ et un recouvrement fppf standard
$\mathcal{V}_i = \{T_{i, k} \to T_i\}_{k = 1, \ldots, n}$ de $T_i$.
Notons $\mathcal{V}_{i'} = \{T_{i', k} \to T_{i'}\}_{k = 1, \ldots, n}$
le recouvrement obtenu par changement de base pour $i' \geq i$. On a alors
\begin{align*}
\colim_{i' \geq i} \check H^0(\mathcal{V}_i, F)
& =
\colim_{i' \geq i}
\text{Noyau du couple}
\left(
\xymatrix{
\prod F(T_{i', k})
\ar@<1ex>[r] \ar@<-1ex>[r] &
\prod F(T_{i', k} \times_{T_{i'}} T_{i', l})
}
\right)
\\
& =
\text{Noyau du couple}
\left(
\xymatrix{
\colim_{i' \geq i}
\prod F(T_{i', k})
\ar@<1ex>[r] \ar@<-1ex>[r] &
\colim_{k' \geq k}
\prod F(T_{i', k} \times_{T_{i'}} T_{i', l})
}
\right)
\\
& =
\text{Noyau du couple}
\left(
\xymatrix{
\prod F(T_k)
\ar@<1ex>[r] \ar@<-1ex>[r] &
\prod F(T_k \times_T T_l)
}
\right)
\\
& =
\check H^0(\mathcal{V}, F)
\end{align*}
La deuxième égalité résulte de l'exactitude des limites inductives filtrantes.
La troisième résulte de ce que $F$ est compatible aux limites et de ce que
$\lim_{i' \geq i} T_{i', k} = T_k$ et
$\lim_{i' \geq i} T_{i', k} \times_{T_{i'}} T_{i', l} = T_k \times_T T_l$
d'après Limites, lemme \ref{limits-lemma-scheme-over-limit}.
En utilisant ceci simultanément pour tous les recouvrements, nous obtenons
\begin{align*}
F^+(T)
& =
\colim_{\mathcal{V}\text{ recouvrement standard de }T} \check H^0(\mathcal{V}, F) \\
& =
\colim_{i \in I}
\colim_{\mathcal{V}_i\text{ recouvrement standard de }T_i}
\check H^0(T \times_{T_i}\mathcal{V}_i, F) \\
& =
\colim_{i \in I} F^+(T_i)
\end{align*}
On peut permuter l'ordre des limites inductives en vertu de
Catégories, lemme \ref{categories-lemma-colimits-commute}.
\end{proof}

\begin{lemma}
\label{spaces-limits-lemma-sheaf-finite-presentation}
Soit $S$ un schéma.
Soit $F : (\Sch/S)_{fppf}^{opp} \to \textit{Ens}$ un foncteur.
Supposons que
\begin{enumerate}
\item $F$ soit un faisceau, et
\item il existe un recouvrement fppf $\{U_j \to S\}_{j \in J}$ tel que
$F|_{(\Sch/U_j)_{fppf}}$ soit compatible aux limites.
\end{enumerate}
Alors $F$ est compatible aux limites.
\end{lemma}

\begin{proof}
Soit $T$ un schéma affine sur $S$.
Soit $I$ un ensemble ordonné filtrant et soit
$T_i$ un système projectif de schémas affines sur $S$ tel que
$T = \lim T_i$. Il s'agit de montrer que l'application canonique
$\colim F(T_i) \to F(T)$ est bijective.

\medskip\noindent
Choisissons un élément $0 \in I$ et un recouvrement fppf standard
$\{V_{0, k} \to T_{0}\}_{k = 1, \ldots, m}$ qui raffine
le recouvrement $\{U_j \times_S T_0 \to T_0\}$ obtenu par changement de base à partir du recouvrement donné de $S$.
Pour tout $i \geq 0$, posons $V_{i, k} = T_i \times_{T_0} V_{0, k}$, et
posons $V_k = T \times_{T_0} V_{0, k}$. On a
$V_k = \lim_{i \geq 0} V_{i, k}$, voir
Limites, lemme \ref{limits-lemma-scheme-over-limit}.

\medskip\noindent
Supposons que $x, x' \in \colim F(T_i)$ aient la même image dans
$F(T)$. Représentons $x, x'$ par des éléments $x_i, x'_i \in F(T_i)$
pour un certain $i \in I$ (on peut choisir le même $i$ pour les deux puisque $I$ est filtrant).
L'hypothèse (2) et l'égalité des images de $x_i, x'_i$ dans
$F(T)$ montrent que
$$
x_i|_{V_{i', k}} = x'_i|_{V_{i', k}}
$$
pour un certain $i' \in I$ suffisamment grand. On peut choisir le même $i'$ pour tout
$k$, car $k \in \{1, \ldots, m\}$ parcourt un ensemble fini.
Puisque $\{V_{i', k} \to T_{i'}\}$
est un recouvrement fppf et que $F$ est un faisceau, il vient
$x_i|_{T_{i'}} = x'_i|_{T_{i'}}$. L'application
$\colim F(T_i) \to F(T)$ est injective.

\medskip\noindent
La surjectivité se démontre de manière analogue.
Soit $x \in F(T)$. Pour chaque $k$, l'hypothèse (2) permet de
choisir un indice $i$ tel que $x|_{V_k}$ provienne d'un
élément $x_{i, k} \in F(V_{i, k})$. Comme précédemment, on peut choisir un
même $i$ qui convienne pour tous les $k$. D'après l'injectivité
établie ci-dessus, on a
$$
x_{i, k}|_{V_{i', k} \times_{T_{i'}} V_{i', l}}
=
x_{i, l}|_{V_{i', k} \times_{T_{i'}} V_{i', l}}
$$
pour un $i'$ suffisamment grand. La condition de faisceau pour $F$ implique alors que
les éléments $x_{i, k}|_{V_{i', k}}$ se recollent en un élément $x_{i'} \in F(T_{i'})$,
ce qui conclut.
\end{proof}

\begin{lemma}
\label{spaces-limits-lemma-sheafify-finite-presentation-map}
Soit $S$ un schéma contenu dans $\Sch_{fppf}$.
Soient $F, G : (\Sch/S)_{fppf}^{opp} \to \textit{Ens}$ des foncteurs.
Si la transformation $a : F \to G$ est compatible aux limites,
il en est de même de la transformation induite entre les faisceaux associés
$F^\# \to G^\#$.
\end{lemma}

\begin{proof}
Soit $T$ un schéma et soit $y \in G^\#(T)$.
Il s'agit de montrer que le foncteur
$F^\#_y : (\Sch/T)_{fppf}^{opp} \to \textit{Ens}$
construit à partir de $F^\# \to G^\#$ et de $y$ comme dans
la définition \ref{spaces-limits-definition-locally-finite-presentation}
est compatible aux limites.
D'après l'équation (\ref{spaces-limits-equation-fibre-map-functors}),
$F^\#_y$ est un faisceau. Choisissons un recouvrement fppf
$\{V_j \to T\}_{j \in J}$ tel que $y|_{V_j}$ provienne
d'un élément $y_j \in F(V_j)$.
La restriction de $F^\#$ à $(\Sch/V_j)_{fppf}$
n'est autre que $F^\#_{y_j}$. Si nous montrons que $F^\#_{y_j}$
est compatible aux limites, alors
le lemme \ref{spaces-limits-lemma-sheaf-finite-presentation}
garantit que $F^\#_y$ est compatible aux limites,
ce qui conclut. On est donc ramené au cas $y \in G(T)$.

\medskip\noindent
Soit $y \in G(T)$. Dans ce cas, on a $F^\#_y = (F_y)^\#$.
Cette égalité résulte de
l'équation (\ref{spaces-limits-equation-fibre-map-functors}).
Le résultat découle donc du
lemme \ref{spaces-limits-lemma-sheafify-finite-presentation}.
\end{proof}

\begin{proposition}
\label{spaces-limits-proposition-characterize-locally-finite-presentation}
Soit $S$ un schéma. Soit $f : X \to Y$ un morphisme d'espaces algébriques
sur $S$. Les conditions suivantes sont équivalentes :
\begin{enumerate}
\item Le morphisme $f$ est localement de présentation finie en tant que
morphisme d'espaces algébriques, voir
Morphismes d'espaces,
définition \ref{spaces-morphisms-definition-locally-finite-presentation}.
\item Le morphisme $f : X \to Y$ est compatible aux limites en tant que
transformation de foncteurs, voir
définition \ref{spaces-limits-definition-locally-finite-presentation}.
\end{enumerate}
\end{proposition}

\begin{proof}
Supposons (1). Soit $T$ un schéma et soit $y \in Y(T)$. Nous devons montrer que
$T \times_Y X$ est compatible aux limites sur $T$ au sens de
la définition \ref{spaces-limits-definition-locally-finite-presentation}.
Nous sommes donc réduits à démontrer que si $X$ est un espace algébrique qui
est localement de présentation finie sur $S$ en tant qu'espace algébrique, alors il
est compatible aux limites en tant que foncteur
$X : (\Sch/S)_{fppf}^{opp} \to \textit{Ens}$.
Pour le voir, choisissons une présentation $X = U/R$, voir
Espaces, définition \ref{spaces-definition-presentation}.
Il découle de
Morphismes d'espaces,
définition \ref{spaces-morphisms-definition-locally-finite-presentation}
que $U$ et $R$ sont des schémas localement de présentation finie
sur $S$. Donc, par
Limites, proposition
\ref{limits-proposition-characterize-locally-finite-presentation},
nous avons
$$
U(T) = \colim U(T_i), \quad
R(T) = \colim R(T_i)
$$
chaque fois que $T = \lim_i T_i$ dans $(\Sch/S)_{fppf}$. Il s'ensuit
que le préfaisceau
$$
(\Sch/S)_{fppf}^{opp} \longrightarrow \textit{Ens}, \quad
W \longmapsto U(W)/R(W)
$$
est compatible aux limites. Donc, par le
lemme \ref{spaces-limits-lemma-sheafify-finite-presentation},
le faisceau associé $X = U/R$ est lui aussi compatible aux limites.

\medskip\noindent
Supposons (2). Choisissons un schéma $V$ et un morphisme \'etale surjectif
$V \to Y$. Ensuite, choisissons un schéma $U$ et un morphisme \'etale surjectif
$U \to V \times_Y X$. Par le
lemme \ref{spaces-limits-lemma-base-change-locally-finite-presentation}, la transformation de foncteurs
$V \times_Y X \to V$ est compatible aux limites. Par
Morphismes d'espaces,
lemme \ref{spaces-morphisms-lemma-etale-locally-finite-presentation},
le morphisme d'espaces algébriques $U \to V \times_Y X$ est localement
de présentation finie, donc est compatible aux limites en tant que
transformation de foncteurs par la première partie de la preuve. Par
lemme \ref{spaces-limits-lemma-composition-locally-finite-presentation},
la composition $U \to V \times_Y X \to V$ est compatible aux limites
en tant que transformation de foncteurs. Donc
le morphisme de schémas $U \to V$ est localement de présentation finie par
Limites, proposition
\ref{limits-proposition-characterize-locally-finite-presentation}
(à ceci près d'une remarque ensembliste, voir le dernier paragraphe de la preuve).
Cela signifie, par définition, que la condition (1) est satisfaite.

\medskip\noindent
Remarque ensembliste. Soit $U \to V$ un morphisme de
$(\Sch/S)_{fppf}$. Dans l'énoncé de
Limites, proposition
\ref{limits-proposition-characterize-locally-finite-presentation}
le fait que $U \to V$ soit localement de présentation finie
se caractérise ainsi : pour {\it tout} système projectif filtrant $(T_i, f_{ii'})$ de schémas affines
sur $V$, on a $U(T) = \colim V(T_i)$ ; mais, dans le contexte actuel,
nous ne pouvons considérer que des schémas affines $T_i$ sur $V$ qui sont (isomorphes à)
un objet de $(\Sch/S)_{fppf}$. Nous devons donc nous assurer qu'il y a suffisamment
de schémas affines dans $(\Sch/S)_{fppf}$ pour que la preuve fonctionne.
En examinant la preuve de (2) $\Rightarrow$ (1) de
Limites, proposition
\ref{limits-proposition-characterize-locally-finite-presentation}
nous voyons que la question se réduit au cas où $U$ et $V$ sont affines.
Écrivons $U = \Spec(A)$ et $V = \Spec(B)$. Par construction
de $(\Sch/S)_{fppf}$, le spectre de tout anneau de cardinal
$\leq |B|$ est isomorphe à un objet de $(\Sch/S)_{fppf}$.
Il suffit donc d'observer que, dans la démonstration du sens nécessaire de
Algèbre, lemme \ref{algebra-lemma-characterize-finite-presentation},
seules des $A$-algèbres de cardinal $\leq |B|$ sont utilisées.
\end{proof}

\begin{remark}
\label{spaces-limits-remark-limit-preserving}
Voici un cas particulier important de
proposition \ref{spaces-limits-proposition-characterize-locally-finite-presentation}.
Soit $S$ un schéma. Soit $X$ un espace algébrique sur $S$.
Alors $X$ est localement de présentation finie sur $S$ si et seulement
si $X$, en tant que foncteur $(\Sch/S)^{opp} \to \textit{Ens}$,
est compatible aux limites. Comparer avec
Limites, remarque \ref{limits-remark-limit-preserving}.
En fait, nous verrons dans le lemme \ref{spaces-limits-lemma-surjection-is-enough}
ci-dessous qu'il suffit que l'application
$$
\colim X(T_i) \longrightarrow X(T)
$$
soit surjective lorsque $T = \lim T_i$ est la limite d'un système projectif filtrant de
schémas affines sur $S$.
\end{remark}

\begin{lemma}
\label{spaces-limits-lemma-surjection-is-enough}
Soit $S$ un schéma.
Soit $f : X \to Y$ un morphisme d'espaces algébriques sur $S$.
Si, pour toute limite $T = \lim_{i \in I} T_i$ d'un système projectif filtrant
de schémas affines sur $S$, l'application
$$
\colim X(T_i) \longrightarrow X(T) \times_{Y(T)} \colim Y(T_i)
$$
est surjective, alors $f$ est localement de présentation finie.
Autrement dit, dans la partie (2) de la
proposition \ref{spaces-limits-proposition-characterize-locally-finite-presentation},
il suffit de vérifier la surjectivité dans le critère du
lemme \ref{spaces-limits-lemma-characterize-relative-limit-preserving}.
\end{lemma}

\begin{proof}
Choisissons un schéma $V$ et un morphisme \'etale surjectif $g : V \to Y$.
Ensuite, choisissons un schéma $U$ et un morphisme \'etale surjectif
$h : U \to V \times_Y X$. Il suffit de montrer, pour $T = \lim T_i$
comme dans le lemme, que l'application
$$
\colim U(T_i) \longrightarrow U(T) \times_{V(T)} \colim V(T_i)
$$
est surjective, car alors $U \to V$ sera localement de présentation
finie par Limites, lemme \ref{limits-lemma-surjection-is-enough}
(à ceci près d'une remarque ensembliste exactement comme dans la preuve de
la proposition \ref{spaces-limits-proposition-characterize-locally-finite-presentation}).
Ainsi, prenons $a : T \to U$ et $b_i : T_i \to V$ qui déterminent
le même morphisme $T \to V$. Voici le diagramme
$$
\xymatrix{
T \ar[d]_a \ar[rr]_{p_i} & & T_i \ar[d]^{b_i} \ar@{..>}[ld] \\
U \ar[r]^-h & X \times_Y V \ar[d] \ar[r] & V \ar[d]^g \\
& X \ar[r]^f & Y
}
$$
D'après l'hypothèse du lemme, quitte à augmenter $i$, il existe
un morphisme $c_i : T_i \to X$ tel que
$h \circ a = (b_i, c_i) \circ p_i : T_i \to V \times_Y X$
et tel que $f \circ c_i = g \circ b_i$.
Puisque $h$ est un morphisme \'etale d'espaces algébriques
(et donc localement de présentation finie), nous avons la surjectivité de
$$
\colim U(T_i) \longrightarrow U(T) \times_{(X \times_Y V)(T)}
\colim (X \times_Y V)(T_i)
$$
par la proposition \ref{spaces-limits-proposition-characterize-locally-finite-presentation}.
Par conséquent, quitte à augmenter encore $i$, il existe un
morphisme $a_i : T_i \to U$ tel que $a = a_i \circ p_i$ et
$b_i = (U \to V) \circ a_i$.
\end{proof}















\section{Limites d'espaces algébriques}
\label{spaces-limits-section-limits}

\noindent
Le lemme suivant explique comment nous concevons les limites des espaces
algébriques dans ce chapitre. Nous utiliserons (sans autre mention) que le
changement de base d'un morphisme affine d'espaces algébriques est affine (voir
Morphismes d'espaces, lemme \ref{spaces-morphisms-lemma-base-change-affine}).

\begin{lemma}
\label{spaces-limits-lemma-directed-inverse-system-has-limit}
Soit $S$ un schéma. Soit $I$ un ensemble ordonné filtrant.
Soit $(X_i, f_{ii'})$ un système projectif indexé par $I$
dans la catégorie des espaces algébriques sur $S$.
Si les morphismes $f_{ii'} : X_i \to X_{i'}$ sont affines, alors la
limite $X = \lim_i X_i$ (en tant que faisceau fppf) est un espace algébrique.
De plus,
\begin{enumerate}
\item chacun des morphismes $f_i : X \to X_i$ est affine,
\item pour tout $i \in I$ et tout morphisme d'espaces algébriques
$T \to X_i$, nous avons
$$
X \times_{X_i} T = \lim_{i' \geq i} X_{i'} \times_{X_i} T.
$$
Cette égalité étant entendue dans la catégorie des espaces algébriques sur $S$.
\end{enumerate}
\end{lemma}

\begin{proof}
La partie (2) est une conséquence formelle de l'existence de la
limite $X = \lim X_i$ en tant qu'espace algébrique sur $S$.
Choisissons un élément $0 \in I$ (c'est possible puisque tout ensemble filtrant est non vide).
Choisissons un schéma $U_0$ et un morphisme
\'etale surjectif $U_0 \to X_0$. Posons $R_0 = U_0 \times_{X_0} U_0$
de sorte que $X_0 = U_0/R_0$. Pour $i \geq 0$, posons
$U_i = X_i \times_{X_0} U_0$ et
$R_i = X_i \times_{X_0} R_0 = U_i \times_{X_i} U_i$.
D'après Limites, lemme \ref{limits-lemma-directed-inverse-system-has-limit}
on voit que $U = \lim_{i \geq 0} U_i$ et $R = \lim_{i \geq 0} R_i$
sont des schémas. De plus, les deux morphismes $s, t : R \to U$ sont les
changements de base des deux projections $R_0 \to U_0$ par le morphisme
$U \to U_0$, en particulier ils sont \'etales. Le morphisme $R \to U \times_S U$
définit une relation d'équivalence, car une limite projective de relations d'équivalence
indexée par un ensemble filtrant est une relation d'équivalence. Par conséquent, le morphisme
$R \to U \times_S U$ est une relation d'équivalence \'etale. Nous affirmons que
le morphisme naturel
\begin{equation}
\label{spaces-limits-equation-isomorphism-sheaves}
U/R \longrightarrow \lim X_i
\end{equation}
est un isomorphisme de faisceaux fppf sur la catégorie des schémas sur $S$.
Cette assertion implique que $X = \lim X_i$ est un espace algébrique
d'après Espaces, théorème \ref{spaces-theorem-presentation}.

\medskip\noindent
Soit $Z$ un schéma et soit $a : Z \to \lim X_i$ un morphisme.
Alors $a = (a_i)$, où $a_i : Z \to X_i$. Posons $W_0 = Z \times_{a_0, X_0} U_0$.
Remarquons que $W_0 = Z \times_{a_i, X_i} U_i$ pour tout $i \geq 0$, par notre
choix de $U_i \to X_i$ ci-dessus. Nous obtenons donc un morphisme
$W_0 \to \lim_{i \geq 0} U_i = U$. Comme $W_0 \to Z$ est surjectif
et \'etale, nous concluons que (\ref{spaces-limits-equation-isomorphism-sheaves})
est un morphisme surjectif de faisceaux. Enfin, soit
$Z$ un schéma et soient $a, b : Z \to U/R$ deux morphismes
dont les images par (\ref{spaces-limits-equation-isomorphism-sheaves}) sont égales.
Il s'agit de montrer que $a = b$.
Après avoir remplacé $Z$ par les objets d'un recouvrement fppf,
on peut supposer qu'il existe des morphismes $a', b' : Z \to U$ qui
donnent $a$ et $b$. Le fait que les images de $a, b$ par
(\ref{spaces-limits-equation-isomorphism-sheaves}) soient égales signifie que
pour chaque $i \geq 0$ les compositions $a_i', b_i' : Z \to U \to U_i$
sont égales comme morphismes à valeurs dans $U_i/R_i = X_i$. Par conséquent,
$(a_i', b_i') : Z \to U_i \times_S U_i$ se factorise par
$R_i$ au moyen d'un morphisme $c_i : Z \to R_i$. Comme
$R = \lim_{i \geq 0} R_i$, on voit que $c = \lim c_i : Z \to R$
est un morphisme qui montre que $a, b$ sont égaux comme morphismes
de $Z$ vers $U/R$.

\medskip\noindent
La partie (1) découle du fait que nous avons vu ci-dessus
$U_i \times_{X_i} X = U$ et que $U \to U_i$ est affine par
construction.
\end{proof}

\begin{lemma}
\label{spaces-limits-lemma-space-over-limit}
Soit $S$ un schéma. Soit $I$ un ensemble ordonné filtrant.
Soit $(X_i, f_{ii'})$ un système projectif indexé par $I$ d'espaces algébriques
sur $S$ à morphismes de transition affines.
Soit $X = \lim_i X_i$. Soit $0 \in I$. Supposons que $T \to X_0$ soit un
morphisme d'espaces algébriques. Alors
$$
T \times_{X_0} X = \lim_{i \geq 0} T \times_{X_0} X_i
$$
en tant qu'espaces algébriques sur $S$.
\end{lemma}

\begin{proof}
La limite $X$ est un espace algébrique d'après
le lemme \ref{spaces-limits-lemma-directed-inverse-system-has-limit}.
L'égalité est formelle, voir
Catégories, lemme \ref{categories-lemma-colimits-commute}.
\end{proof}

\begin{lemma}
\label{spaces-limits-lemma-directed-inverse-system-closed-immersions}
Soit $S$ un schéma. Soit $I$ un ensemble ordonné filtrant.
Soit $(X_i, f_{i'i}) \to (Y_i, g_{i'i})$ un morphisme
de systèmes projectifs indexés par $I$ d'espaces algébriques sur $S$.
Supposons que
\begin{enumerate}
\item les morphismes $f_{i'i} : X_{i'} \to X_i$ soient affines,
\item les morphismes $g_{i'i} : Y_{i'} \to Y_i$ soient affines,
\item les morphismes $X_i \to Y_i$ soient des immersions fermées.
\end{enumerate}
Alors $\lim X_i \to \lim Y_i$ est une immersion fermée.
\end{lemma}

\begin{proof}
Observons que $\lim X_i$ et $\lim Y_i$ existent d'après
le lemme \ref{spaces-limits-lemma-directed-inverse-system-has-limit}.
Choisissons $0 \in I$ et un schéma affine $V_0$ ainsi qu'un morphisme \'etale
$V_0 \to Y_0$. Alors les morphismes
$V_i = Y_i \times_{Y_0} V_0 \to U_i = X_i \times_{Y_0} V_0$
sont des immersions fermées de schémas affines.
Par conséquent, le morphisme $V = Y \times_{Y_0} V_0 \to U = X \times_{Y_0} V_0$
est une immersion fermée, car $V = \lim V_i$, $U = \lim U_i$
et parce qu'une limite projective d'immersions fermées de schémas affines est une
immersion fermée : une limite inductive filtrante d'homomorphismes d'anneaux surjectifs
est surjective. Comme les morphismes \'etales $V \to Y$ forment un
recouvrement \'etale de $Y$ lorsque nous faisons varier notre choix de $V_0 \to Y_0$
le lemme s'ensuit.
\end{proof}

\begin{lemma}
\label{spaces-limits-lemma-directed-inverse-system-reduced}
Soit $S$ un schéma. Soit $I$ un ensemble ordonné filtrant.
Soit $(X_i, f_{i'i})$ un système projectif indexé par $I$
d'espaces algébriques sur $S$. Si $X_i$ est réduit
pour tout $i$, alors $X$ est réduit.
\end{lemma}

\begin{proof}
Observons que $\lim X_i$ existe d'après
le lemme \ref{spaces-limits-lemma-directed-inverse-system-has-limit}.
Choisissons $0 \in I$ et un schéma affine $V_0$ ainsi qu'un morphisme \'etale
$U_0 \to X_0$. Alors les schémas affines
$U_i = X_i \times_{X_0} U_0$ sont réduits.
Par conséquent, $U = X \times_{X_0} U_0$
est un schéma affine réduit, comme limite projective de schémas affines réduits :
une limite inductive filtrante d'anneaux réduits est réduite.
Comme les morphismes \'etales $U \to X$ forment un
recouvrement \'etale de $X$ lorsque nous faisons varier notre choix de $U_0 \to X_0$
le lemme s'ensuit.
\end{proof}

\begin{lemma}
\label{spaces-limits-lemma-better-characterize-relative-limit-preserving}
Soit $S$ un schéma. Soit $X \to Y$ un morphisme d'espaces algébriques
sur $S$. Les conditions équivalentes (1) et (2) de la
proposition \ref{spaces-limits-proposition-characterize-locally-finite-presentation}
sont également équivalentes à
\begin{enumerate}
\item[(3)] pour toute limite $T = \lim T_i$ d'un système projectif filtrant
d'espaces algébriques quasi-compacts et quasi-séparés $T_i$ sur $S$ à morphismes de transition affines, le diagramme
d'ensembles
$$
\xymatrix{
\colim_i \Mor(T_i, X) \ar[r] \ar[d] & \Mor(T, X) \ar[d] \\
\colim_i \Mor(T_i, Y) \ar[r] & \Mor(T, Y)
}
$$
est un diagramme cartésien.
\end{enumerate}
\end{lemma}

\begin{proof}
Il est clair que (3) implique (2). Nous supposerons (2) et démontrerons (3).
La démonstration est assez formelle et nous laissons au lecteur le soin
d'en trouver une.

\medskip\noindent
Commençons par démontrer que (3) est valable
lorsque $T_i$ est en outre supposé séparé pour tout $i$.
Choisissons $i \in I$ et un morphisme \'etale surjectif $U_i \to T_i$
où $U_i$ est affine. En utilisant le lemme \ref{spaces-limits-lemma-space-over-limit},
on voit qu'en posant $U = U_i \times_{T_i} T$ et
$U_{i'} = U_i \times_{T_i} T_{i'}$, on a $U = \lim_{i' \geq i} U_{i'}$.
Bien sûr, $U$ et $U_{i'}$ sont affines (voir
le lemme \ref{spaces-limits-lemma-directed-inverse-system-has-limit}).
Comme $T_i$ est séparé, le produit fibré $V_i = U_i \times_{T_i} U_i$
est également un schéma affine et nous obtenons des schémas affines
$V = V_i \times_{T_i} T$ et
$V_{i'} = V_i \times_{T_i} T_{i'}$ tels que $V = \lim_{i' \geq i} V_{i'}$.
Observons que $U \to T$ et $U_i \to T_i$ sont surjectifs et \'etales
et que $V = U \times_T U$ et $V_{i'} = U_{i'} \times_{T_{i'}} U_{i'}$.
Remarquons que $\Mor(T, X)$ est le noyau du couple formé par les deux applications
$\Mor(U, X) \to \Mor(V, X)$ ; cela résulte par exemple du fait que
$X$, comme faisceau sur $(\Sch/S)_{fppf}$, est le conoyau du couple formé par
les deux applications $h_V \to h_u$. De même,
$\Mor(T_{i'}, X)$ est le noyau du couple formé par les
deux applications $\Mor(U_{i'}, X) \to \Mor(V_{i'}, X)$.
Et bien sûr, la même chose vaut en remplaçant $X$ par $Y$.
La condition (2) dit que les diagrammes de (3) sont cartésiens
dans le cas de $U = \lim U_i$ et $V = \lim V_i$.
Il s'ensuit formellement que la même chose vaut pour $T = \lim T_i$.

\medskip\noindent
Dans le cas général, choisissons un schéma affine $U$, un $i \in I$,
et un morphisme \'etale surjectif $U \to T_i$. En répétant
l'argument du paragraphe précédent, on conclut encore :
les schémas $V_{i'}$, $V$ ne sont plus affines, mais ils sont
toujours quasi-compacts et
séparés, et le résultat du paragraphe précédent s'applique.
\end{proof}




\section{Descente des propriétés}
\label{spaces-limits-section-descent}

\noindent
Cette section est l'analogue de Limites, section \ref{limits-section-descent}.

\begin{lemma}
\label{spaces-limits-lemma-inverse-limit-sets}
Soit $S$ un schéma. Soit $X = \lim_{i \in I} X_i$ la limite d'un système projectif filtrant
d'espaces algébriques sur $S$ à morphismes de transition affines
(lemme \ref{spaces-limits-lemma-directed-inverse-system-has-limit}). Si chaque $X_i$
est décent (par exemple quasi-séparé ou localement séparé),
alors $|X| = \lim_i |X_i|$ en tant qu'ensembles.
\end{lemma}

\begin{proof}
Il existe une application canonique $|X| \to \lim |X_i|$. Choisissons $0 \in I$.
Si $W_0 \subset X_0$ est un sous-espace ouvert, alors
$f_0^{-1}W_0 = \lim_{i \geq 0} f_{i0}^{-1}W_0$, voir
le lemme \ref{spaces-limits-lemma-directed-inverse-system-has-limit}.
Ainsi, le cas général se ramène à celui des systèmes projectifs pour lesquels $X_0$
est quasi-compact. On peut donc supposer
désormais que $X_0$ est quasi-compact.

\medskip\noindent
Choisissons un schéma affine $U_0$ et un morphisme \'etale surjectif $U_0 \to X_0$.
Posons $U_i = X_i \times_{X_0} U_0$ et $U = X \times_{X_0} U_0$.
Posons $R_i = U_i \times_{X_i} U_i$ et $R = U \times_X U$.
Rappelons que $U = \lim U_i$ et $R = \lim R_i$, voir la démonstration du
lemme \ref{spaces-limits-lemma-directed-inverse-system-has-limit}.
Rappelons que $|X| = |U|/|R|$ et $|X_i| = |U_i|/|R_i|$. D'après
Limites, lemme \ref{limits-lemma-topology-limit}, nous avons
$|U| = \lim |U_i|$ et $|R| = \lim |R_i|$.

\medskip\noindent
Surjectivité de $|X| \to \lim |X_i|$. Soit $(x_i) \in \lim |X_i|$. Notons
$S_i \subset |U_i|$ l'image réciproque de $x_i$. Il s'agit d'un ensemble fini non vide
par définition des espaces décents
(Espaces décents, définition \ref{decent-spaces-definition-very-reasonable}).
Donc $\lim S_i$ est non vide, voir
Catégories, lemme \ref{categories-lemma-nonempty-limit}.
Soit $(u_i) \in \lim S_i \subset \lim |U_i|$. D'après ce qui précède, cela détermine
un point $u \in |U|$ ayant pour image un point $x \in |X|$ dont l'image est l'élément donné
$(x_i)$ de $\lim |X_i|$.

\medskip\noindent
Injectivité de $|X| \to \lim |X_i|$. Supposons que $x, x' \in |X|$
aient la même image dans $\lim |X_i|$. Choisissons des relèvements $u, u' \in |U|$
et notons $u_i, u'_i \in |U_i|$ leurs images.
Pour chaque $i$, soit $T_i \subset |R_i|$ l'ensemble des points ayant pour image
sur $(u_i, u'_i) \in |U_i| \times |U_i|$. Il s'agit d'un ensemble fini
par définition des espaces décents
(Espaces décents, définition \ref{decent-spaces-definition-very-reasonable}).
De plus, $T_i$ est non vide puisque nous avons supposé que $x$ et $x'$ avaient
le même point de $X_i$. Donc $\lim T_i$ est non vide, voir
Catégories, lemme \ref{categories-lemma-nonempty-limit}.
Comme précédemment, soit $r \in |R| = \lim |R_i|$ un point correspondant à un
élément de $\lim T_i$. Alors $r$ a pour image $(u, u')$ dans $|U| \times |U|$
par construction et nous voyons que $x = x'$ dans $|X|$, comme voulu.

\medskip\noindent
Pour justifier l'assertion entre parenthèses, un espace algébrique quasi-séparé est décent, voir
Espaces décents, section \ref{decent-spaces-section-reasonable-decent}
(l'observation clé à cet égard est le lemme de Propriétés des espaces,
\ref{spaces-properties-lemma-finite-fibres-presentation}).
Un espace algébrique localement séparé est décent d'après
Espaces décents, lemme \ref{decent-spaces-lemma-locally-separated-decent}.
\end{proof}

\begin{lemma}
\label{spaces-limits-lemma-topology-limit}
Avec les mêmes notations et hypothèses que dans le lemme \ref{spaces-limits-lemma-inverse-limit-sets},
nous avons $|X| = \lim_i |X_i|$ en tant qu'espaces topologiques.
\end{lemma}

\begin{proof}
Nous utiliserons le critère de
Topologie, lemme \ref{topology-lemma-characterize-limit}.
Nous avons vu que $|X| = \lim_i |X_i|$ en tant qu'ensembles dans
le lemme \ref{spaces-limits-lemma-inverse-limit-sets}.
Les $f_i : X \to X_i$ sont des morphismes d'espaces algébriques
et induisent donc des applications continues $|X| \to |X_i|$.
Ainsi $f_i^{-1}(U_i)$ est ouvert pour tout
ouvert $U_i \subset |X_i|$. Enfin,
soit $x \in |X|$ et soit $x \in V \subset |X|$ un voisinage ouvert.
Nous devons trouver un $i$ et un voisinage ouvert
$W_i \subset |X_i|$ de l'image de $x$
tel que $f_i^{-1}(W_i) \subset V$.
Choisissons $0 \in I$. Choisissons un schéma $U_0$ et un morphisme
\'etale surjectif $U_0 \to X_0$. Posons $U = X \times_{X_0} U_0$
et $U_i = X_i \times_{X_0} U_0$ pour $i \geq 0$.
Alors $U = \lim_{i \geq 0} U_i$ dans la catégorie des schémas d'après
le lemme \ref{spaces-limits-lemma-directed-inverse-system-has-limit}.
Choisissons $u \in U$ s'envoyant sur $x$. D'après le résultat pour les schémas
(Limites, lemme \ref{limits-lemma-inverse-limit-top})
nous pouvons trouver un $i \geq 0$ et un voisinage ouvert
$E_i \subset U_i$ de l'image de $u$
dont l'image réciproque dans $U$ est contenue dans
l'image réciproque de $V$ dans $U$. Nous pouvons alors poser $W_i \subset |X_i|$
égal à l'image de $E_i$. Cela fonctionne car $|U_i| \to |X_i|$ est ouvert.
\end{proof}

\begin{lemma}
\label{spaces-limits-lemma-limit-nonempty}
Soit $S$ un schéma. Soit $X = \lim_{i \in I} X_i$ la limite d'un système
projectif filtrant d'espaces algébriques sur $S$ à morphismes de transition affines
(lemme \ref{spaces-limits-lemma-directed-inverse-system-has-limit}). Si chaque $X_i$
est quasi-compact et non vide, alors $|X|$ est non vide.
\end{lemma}

\begin{proof}
Choisissons $0 \in I$.
Choisissons un schéma affine $U_0$ et un morphisme \'etale surjectif $U_0 \to X_0$.
Posons $U_i = X_i \times_{X_0} U_0$ et $U = X \times_{X_0} U_0$.
Alors chaque $U_i$ est un schéma affine non vide. Donc $U = \lim U_i$
est non vide (Limites, lemme \ref{limits-lemma-limit-nonempty}) et ainsi
$X$ est non vide.
\end{proof}

\begin{lemma}
\label{spaces-limits-lemma-inverse-limit-irreducibles}
Soit $S$ un schéma. Soit $X = \lim_{i \in I} X_i$ la limite d'un système
projectif filtrant d'espaces algébriques sur $S$ à morphismes de transition affines
(lemme \ref{spaces-limits-lemma-directed-inverse-system-has-limit}).
Soit $x \in |X|$ d'images $x_i \in |X_i|$. Si chaque $X_i$ est décent,
alors $\overline{\{x\}} = \lim_i \overline{\{x_i\}}$ en tant qu'ensembles
et comme espaces algébriques lorsqu'ils sont munis de leur structure induite réduite.
\end{lemma}

\begin{proof}
Posons $Z = \overline{\{x\}} \subset |X|$ et
$Z_i = \overline{\{x_i\}} \subset |X_i|$.
Comme l'application $|X| \to |X_i|$ est continue, l'image de $Z$ est contenue dans $Z_i$
pour tout $i$. On obtient donc une application injective $Z \to \lim Z_i$
car $|X| = \lim |X_i|$ en tant qu'ensembles (lemme \ref{spaces-limits-lemma-inverse-limit-sets}).
Supposons que $x' \in |X|$ n'appartienne pas à $Z$.
Il existe alors un ouvert $U \subset |X|$ tel que $x' \in U$
et $x \not \in U$. Comme
$|X| = \lim |X_i|$ en tant qu'espaces topologiques (lemme \ref{spaces-limits-lemma-topology-limit}),
on peut écrire $U = \bigcup_{j \in J} f_j^{-1}(U_j)$
pour une partie $J \subset I$ et des ouverts $U_j \subset |X_j|$, voir
Topologie, lemme \ref{topology-lemma-describe-limits}.
Il existe donc $j \in J$ tel que $f_j(x') \in U_j$
et $f_j(x) \not \in U_j$. Autrement dit, $f_j(x') \not \in Z_j$.
Ainsi $Z = \lim Z_i$ en tant qu'ensembles.

\medskip\noindent
Munissons ensuite $Z$ et $Z_i$ de leur structure induite réduite, voir
Propriétés des espaces, définition
\ref{spaces-properties-definition-reduced-induced-space}.
Les morphismes de transition $X_{i'} \to X_i$ induisent des morphismes affines
$Z_{i'} \to Z_i$, et les projections $X \to X_i$
induisent des morphismes compatibles $Z \to Z_i$.
On obtient donc des morphismes $Z \to \lim Z_i \to X$ d'espaces algébriques.
D'après le lemme \ref{spaces-limits-lemma-directed-inverse-system-closed-immersions},
le morphisme $\lim Z_i \to X$ est une
immersion fermée. D'après le lemme \ref{spaces-limits-lemma-directed-inverse-system-reduced},
l'espace algébrique $\lim Z_i$ est réduit.
D'après ce qui précède, $Z \to \lim Z_i$ est bijectif sur les points.
Par unicité de la structure induite réduite de sous-espace fermé,
ce morphisme est un isomorphisme d'espaces algébriques.
\end{proof}

\begin{situation}
\label{spaces-limits-situation-descent}
Soit $S$ un schéma. Soit $X = \lim_{i \in I} X_i$ la limite d'un système projectif
filtrant d'espaces algébriques sur $S$ à morphismes de transition affines
(lemme \ref{spaces-limits-lemma-directed-inverse-system-has-limit}).
On suppose que $X_i$ est quasi-compact et quasi-séparé pour tout $i \in I$.
On choisit aussi un élément $0 \in I$.
\end{situation}

\begin{lemma}
\label{spaces-limits-lemma-descend-section}
Notations et hypothèses comme dans la situation \ref{spaces-limits-situation-descent}.
Supposons que $\mathcal{F}_0$ soit un faisceau quasi-cohérent sur $X_0$.
Posons $\mathcal{F}_i = f_{0i}^*\mathcal{F}_0$ pour $i \geq 0$ et posons
$\mathcal{F} = f_0^*\mathcal{F}_0$. Alors
$$
\Gamma(X, \mathcal{F}) = \colim_{i \geq 0} \Gamma(X_i, \mathcal{F}_i)
$$
\end{lemma}

\begin{proof}
Choisissons un morphisme étale surjectif $U_0 \to X_0$, où $U_0$ est un schéma
affine (Propriétés des espaces, lemme
\ref{spaces-properties-lemma-quasi-compact-affine-cover}).
Posons $U_i = X_i \times_{X_0} U_0$.
Posons $R_0 = U_0 \times_{X_0} U_0$ et $R_i = R_0 \times_{X_0} X_i$.
Dans la démonstration du lemme \ref{spaces-limits-lemma-directed-inverse-system-has-limit}, nous avons
vu qu'il existe une présentation $X = U/R$ avec
$U = \lim U_i$ et $R = \lim R_i$.
Notons que $U_i$ et $U$ sont affines, et que $R_i$ et $R$ sont
quasi-compacts et séparés (puisque $X_i$ est quasi-séparé). Par conséquent,
Limites, lemme \ref{limits-lemma-descend-section}
donne
$$
\mathcal{F}(U) = \colim \mathcal{F}_i(U_i)
\quad\text{et}\quad
\mathcal{F}(R) = \colim \mathcal{F}_i(R_i).
$$
Le lemme en résulte car
$\Gamma(X, \mathcal{F}) = \Ker(\mathcal{F}(U) \to \mathcal{F}(R))$
et, de même,
$\Gamma(X_i, \mathcal{F}_i) =
\Ker(\mathcal{F}_i(U_i) \to \mathcal{F}_i(R_i))$
\end{proof}

\begin{lemma}
\label{spaces-limits-lemma-descend-opens}
Notations et hypothèses comme dans la situation \ref{spaces-limits-situation-descent}.
Pour tout sous-espace ouvert quasi-compact $U \subset X$, il existe un $i$
et un ouvert quasi-compact $U_i \subset X_i$ dont l'image inverse dans $X$ est $U$.
\end{lemma}

\begin{proof}
Cela résulte formellement de la construction des limites dans le
lemme \ref{spaces-limits-lemma-directed-inverse-system-has-limit}
et du résultat correspondant pour les schémas :
Limites, lemme \ref{limits-lemma-descend-opens}.
\end{proof}

\noindent
Le lemme suivant sera supplanté par le lemme plus fort
\ref{spaces-limits-lemma-descend-isomorphism}.

\begin{lemma}
\label{spaces-limits-lemma-descend-equality}
Notations et hypothèses comme dans la situation \ref{spaces-limits-situation-descent}.
Soit $f_0 : Y_0 \to Z_0$ un morphisme d'espaces algébriques sur $X_0$.
Supposons (a) que $Y_0 \to X_0$ et $Z_0 \to X_0$ soient représentables, (b) que
$Y_0$ et $Z_0$ soient quasi-compacts et quasi-séparés, (c) que
$f_0$ soit localement de présentation finie, et
(d) que $Y_0 \times_{X_0} X \to Z_0 \times_{X_0} X$ soit un isomorphisme.
Il existe alors un $i \geq 0$ tel que
$Y_0 \times_{X_0} X_i \to Z_0 \times_{X_0} X_i$ soit un isomorphisme.
\end{lemma}

\begin{proof}
Choisissons un schéma affine $U_0$ et un morphisme étale surjectif $U_0 \to X_0$.
Posons $U_i = U_0 \times_{X_0} X_i$ et $U = U_0 \times_{X_0} X$.
Appliquons Limites, lemme \ref{limits-lemma-descend-isomorphism} :
on obtient que $Y_0 \times_{X_0} U_i \to Z_0 \times_{X_0} U_i$
est un isomorphisme de schémas pour un certain $i \geq 0$ (détails omis).
Comme $U_i \to X_i$ est étale surjectif, il s'ensuit que
$Y_0 \times_{X_0} X_i \to Z_0 \times_{X_0} X_i$ est un isomorphisme
(détails omis).
\end{proof}

\begin{lemma}
\label{spaces-limits-lemma-descend-separated}
Notations et hypothèses comme dans la situation \ref{spaces-limits-situation-descent}.
Si $X$ est séparé, alors $X_i$ est séparé pour un certain $i \in I$.
\end{lemma}

\begin{proof}
Choisissons un schéma affine $U_0$ et un morphisme étale surjectif $U_0 \to X_0$.
Pour $i \geq 0$, posons $U_i = U_0 \times_{X_0} X_i$, et posons
$U = U_0 \times_{X_0} X$. Notons que $U_i$ et $U$ sont des schémas affines
munis de morphismes étales surjectifs $U_i \to X_i$
et $U \to X$. Posons $R_i = U_i \times_{X_i} U_i$ et $R = U \times_X U$,
de projections $s_i, t_i : R_i \to U_i$ et $s, t : R \to U$.
Notons que $R_i$ et $R$ sont des schémas quasi-compacts et séparés (puisque les
espaces algébriques $X_i$ et $X$ sont quasi-séparés). Les morphismes
$s_i : R_i \to U_i$ et $s : R \to U$ sont de type fini.
Par définition, $X_i$ est séparé si et seulement si
$(t_i, s_i) : R_i \to U_i \times U_i$
est une immersion fermée; comme $X$ est séparé par hypothèse,
le morphisme $(t, s) : R \to U \times U$ est une immersion fermée. Comme
$R \to U$ est de type fini, il existe un
$i$ tel que le morphisme $R \to U_i \times U$ soit une immersion fermée
(Limites, lemme \ref{limits-lemma-finite-type-eventually-closed}).
Fixons un tel $i \in I$. Appliquons Limites, lemme
\ref{limits-lemma-descend-closed-immersion-finite-presentation}
au système de morphismes $R_{i'} \to U_i \times U_{i'}$ pour $i' \geq i$
(ce qui est permis puisque
$R_{i'} = R_i \times_{U_i \times U_i} U_i \times U_{i'}$).
On voit ainsi que $R_{i'} \to U_i \times U_{i'}$ est une immersion fermée
pour $i'$ assez grand. Il s'ensuit immédiatement
que $R_{i'} \to U_{i'} \times U_{i'}$ est une immersion fermée,
ce qui achève la démonstration du lemme.
\end{proof}

\begin{lemma}
\label{spaces-limits-lemma-limit-is-affine}
Notations et hypothèses comme dans la situation \ref{spaces-limits-situation-descent}.
Si $X$ est affine, il existe un $i$ tel que $X_i$ soit affine.
\end{lemma}

\begin{proof}
Choisissons $0 \in I$, puis un schéma affine $U_0$ et un morphisme
étale surjectif $U_0 \to X_0$. Posons $U = U_0 \times_{X_0} X$
et $U_i = U_0 \times_{X_0} X_i$ pour $i \geq 0$. Comme les morphismes de transition
sont affines, les espaces algébriques $U_i$ et $U$ sont affines.
Ainsi, $U \to X$ est un morphisme étale de schémas affines. On peut donc
écrire $X = \Spec(A)$, $U = \Spec(B)$ et
$$
B = A[x_1, \ldots, x_n]/(g_1, \ldots, g_n)
$$
où $\Delta = \det(\partial g_\lambda/\partial x_\mu)$ est inversible
dans $B$; voir Algèbre, lemme \ref{algebra-lemma-etale-standard-smooth}.
Posons $A_i = \mathcal{O}_{X_i}(X_i)$. On a $A = \colim A_i$ d'après le
lemme \ref{spaces-limits-lemma-descend-section}. Quitte à augmenter $0$, on peut supposer
donnés $g_{1, i}, \ldots, g_{n, i} \in A_i[x_1, \ldots, x_n]$ dont les images sont
$g_1, \ldots, g_n$. Posons
$$
B_i = A_i[x_1, \ldots, x_n]/(g_{1, i}, \ldots, g_{n, i})
$$
pour tout $i \geq 0$. Quitte à augmenter $0$ si nécessaire, on peut supposer que
$\Delta_i = \det(\partial g_{\lambda, i}/\partial x_\mu)$ est inversible
dans $B_i$ pour tout $i \geq 0$. Ainsi, $A_i \to B_i$ est un homomorphisme étale.
Quitte à augmenter $0$, on peut aussi supposer que
$\Spec(B_i) \to \Spec(A_i)$ est surjectif; voir
Limites, lemme \ref{limits-lemma-descend-surjective}. Quitte à augmenter
encore $0$, choisissons des éléments
$h_{1, i}, \ldots, h_{n, i} \in \mathcal{O}_{U_i}(U_i)$ dont les images sont les
classes de $x_1, \ldots, x_n$ dans $B = \mathcal{O}_U(U)$ et tels que
$g_{\lambda, i}(h_{\nu, i}) = 0$ dans $\mathcal{O}_{U_i}(U_i)$. Ainsi,
on obtient un diagramme commutatif
\begin{equation}
\label{spaces-limits-equation-to-show-cartesian}
\vcenter{
\xymatrix{
X_i \ar[d] & U_i \ar[l] \ar[d] \\
\Spec(A_i) & \Spec(B_i) \ar[l]
}
}
\end{equation}
Par construction, $B_i = B_0 \otimes_{A_0} A_i$ et
$B = B_0 \otimes_{A_0} A$. Considérons le morphisme
$$
f_0 : U_0 \longrightarrow X_0 \times_{\Spec(A_0)} \Spec(B_0)
$$
C'est un morphisme d'espaces algébriques quasi-compacts et quasi-séparés,
représentable, séparé et étale sur $X_0$. Par nos choix, le changement de base de $f_0$
à $X$ est un isomorphisme. Le lemme
\ref{spaces-limits-lemma-descend-equality}
assure donc qu'il existe un $i$ tel que le changement de base de $f_0$
à $X_i$ soit un isomorphisme; autrement dit, le diagramme
(\ref{spaces-limits-equation-to-show-cartesian}) est cartésien. Par conséquent,
Descente, lemme \ref{descent-lemma-descent-data-sheaves},
appliqué au recouvrement fppf $\{\Spec(B_i) \to \Spec(A_i)\}$
et combiné avec Descente, lemme \ref{descent-lemma-affine},
montre que $X_i \to \Spec(A_i)$ est représentable par un schéma
affine sur $\Spec(A_i)$, comme voulu. (Il s'ensuit bien sûr également
que $X_i = \Spec(A_i)$, mais nous n'en avons pas besoin.)
\end{proof}

\begin{lemma}
\label{spaces-limits-lemma-limit-is-scheme}
Notations et hypothèses comme dans la situation \ref{spaces-limits-situation-descent}.
Si $X$ est un schéma, il existe un $i$ tel que $X_i$ soit un schéma.
\end{lemma}

\begin{proof}
Choisissons un recouvrement ouvert affine fini $X = \bigcup W_j$.
D'après le lemme \ref{spaces-limits-lemma-descend-opens},
il existe un $i \in I$ et des sous-espaces ouverts $W_{j, i} \subset X_i$
dont le changement de base à $X$ est $W_j \to X$. D'après le
lemme \ref{spaces-limits-lemma-limit-is-affine}, on peut supposer que
chaque $W_{j, i}$ est un schéma affine. Il en résulte que $X_i$
est un schéma (voir, par exemple,
Propriétés des espaces, section \ref{spaces-properties-section-schematic}).
\end{proof}

\begin{lemma}
\label{spaces-limits-lemma-finite-type-eventually-closed}
Soit $S$ un schéma. Soit $B$ un espace algébrique sur $S$.
Soit $X = \lim X_i$ la limite d'un système projectif filtrant
d'espaces algébriques sur $B$ à morphismes de transition affines.
Soit $Y \to X$ un morphisme d'espaces algébriques sur $B$.
\begin{enumerate}
\item Si $Y \to X$ est une immersion fermée, si $X_i$ est quasi-compact et si
$Y \to B$ est localement de type fini, alors $Y \to X_i$ est une immersion
fermée pour $i$ assez grand.
\item Si $Y \to X$ est une immersion, si $X_i$ est quasi-séparé, si $Y \to B$
est localement de type fini et si $Y$ est quasi-compact, alors $Y \to X_i$ est
une immersion pour $i$ assez grand.
\item Si $Y \to X$ est un isomorphisme, si $X_i$ est quasi-compact,
si $X_i \to B$ est localement de type fini, si les morphismes de transition
$X_{i'} \to X_i$ sont des immersions fermées et si $Y \to B$ est localement
de présentation finie, alors $Y \to X_i$ est un isomorphisme pour $i$
assez grand.
\item Si $Y \to X$ est un monomorphisme, si $X_i$ est quasi-séparé,
si $Y \to B$ est localement de type fini et si $Y$ est quasi-compact, alors
$Y \to X_i$ est un monomorphisme pour $i$ assez grand.
\end{enumerate}
\end{lemma}

\begin{proof}
Démonstration de (1). Choisissons $0 \in I$. Comme $X_0$ est quasi-compact,
on peut choisir un schéma affine $W$ et un morphisme étale $W \to B$ tels que
l'image de $|X_0| \to |B|$ soit contenue dans celle de $|W| \to |B|$.
Choisissons un schéma affine $U_0$ et un morphisme étale
$U_0 \to X_0 \times_B W$ tels que $U_0 \to X_0$ soit surjectif.
(Cela est possible par le choix de $W$ et la quasi-compacité de $X_0$;
les détails sont omis.)
Notons $V \to Y$, resp.\ $U \to X$, resp.\ $U_i \to X_i$, les changements
de base de $U_0 \to X_0$ (pour $i \geq 0$). Il suffit de prouver que
$V \to U_i$ est une immersion fermée pour $i$ assez grand. On est donc ramené
à démontrer le résultat pour $V \to U = \lim U_i$ sur $W$. Le résultat découle
alors du cas des schémas, traité dans Limites, lemme
\ref{limits-lemma-finite-type-eventually-closed}.

\medskip\noindent
Démonstration de (2). Choisissons $0 \in I$. Choisissons un sous-espace ouvert
quasi-compact $X'_0 \subset X_0$ tel que $Y \to X_0$ se factorise par $X'_0$.
Après avoir remplacé $X_i$ par l'image inverse de $X'_0$ pour $i \geq 0$,
on peut supposer que tous les $X_i'$ sont quasi-compacts et quasi-séparés.
Soit $U \subset X$ un ouvert quasi-compact tel que $Y \to X$ se factorise
par une immersion fermée $Y \to U$ (un tel $U$ existe puisque $Y$ est
quasi-compact). D'après le lemme \ref{spaces-limits-lemma-descend-opens},
on peut supposer que $U = \lim U_i$, où $U_i \subset X_i$ est un ouvert
quasi-compact. D'après (1), $Y \to U_i$ est une immersion fermée pour un
certain $i$. Ainsi, (2) est établi.

\medskip\noindent
Démonstration de (3). Choisissons $0 \in I$, puis un schéma affine $U_0$
et un morphisme étale surjectif $U_0 \to X_0$.
Posons $U_i = X_i \times_{X_0} U_0$ et
$U = X \times_{X_0} U_0 = Y \times_{X_0} U_0$. Alors $U = \lim U_i$ est une
limite de schémas affines, les morphismes de transition du système sont des
immersions fermées, et $U \to U_0$ est de présentation finie (car
$U \to B$ est localement de présentation finie, $U_0 \to B$ est localement
de type fini, et d'après
Morphismes d'espaces, lemme
\ref{spaces-morphisms-lemma-finite-presentation-permanence}).
On est ainsi ramené au fait algébrique suivant : si $A = \lim A_i$
est une limite inductive filtrante de $R$-algèbres à morphismes de transition
surjectifs et si $A$ est de présentation finie sur $A_0$, alors $A = A_i$
pour un certain $i$. En effet, écrivons $A = A_0/(f_1, \ldots, f_n)$.
Choisissons $i$ tel que $f_1, \ldots, f_n$ aient une image nulle par le morphisme surjectif $A_0 \to A_i$.

\medskip\noindent
Démonstration de (4). Posons $Z_i = Y \times_{X_i} Y$.
Comme les morphismes de transition $X_{i'} \to X_i$ sont affines, donc
séparés, les morphismes de transition $Z_{i'} \to Z_i$ sont des immersions
fermées; voir Morphismes d'espaces, lemme
\ref{spaces-morphisms-lemma-fibre-product-after-map}.
On a $\lim Z_i = Y \times_X Y = Y$, puisque $Y \to X$ est un monomorphisme.
Choisissons $0 \in I$. Comme $Y \to X_0$ est localement de type fini
(Morphismes d'espaces, lemme
\ref{spaces-morphisms-lemma-permanence-finite-type})
le morphisme $Y \to Z_0$ est localement de présentation finie
(Morphismes d'espaces, lemme
\ref{spaces-morphisms-lemma-diagonal-morphism-finite-type}).
Les morphismes $Z_i \to Z_0$ sont localement de type fini
(ce sont des immersions fermées).
Enfin, $Z_i = Y \times_{X_i} Y$ est quasi-compact puisque
$X_i$ est quasi-séparé et $Y$ est quasi-compact.
La partie (3) s'applique donc à $Y = \lim_{i \geq 0} Z_i$ sur $Z_0$,
et l'on conclut que $Y = Z_i$ pour un certain $i$. Cela démontre (4) et le lemme.
\end{proof}

\begin{lemma}
\label{spaces-limits-lemma-eventually-separated}
Soit $S$ un schéma. Soit $Y$ un espace algébrique sur $S$.
Soit $X = \lim X_i$ la limite d'un système projectif filtrant d'espaces
algébriques sur $Y$ à morphismes de transition affines. Supposons que
\begin{enumerate}
\item $Y$ soit quasi-séparé,
\item $X_i$ soit quasi-compact et quasi-séparé,
\item le morphisme $X \to Y$ soit séparé.
\end{enumerate}
Alors $X_i \to Y$ est séparé pour tout $i$ assez grand.
\end{lemma}

\begin{proof}
Soit $0 \in I$. Choisissons un schéma affine $W$ et un morphisme étale
$W \to Y$ tels que l'image de $|W| \to |Y|$ contienne celle de
$|X_0| \to |Y|$. Cela est possible puisque $X_0$ est quasi-compact.
Il suffit de vérifier que $W \times_Y X_i \to W$ est séparé
pour un certain $i \geq 0$, car la diagonale de $W \times_Y X_i$
sur $W$ est le changement de base de $X_i \to X_i \times_Y X_i$ par
le morphisme étale surjectif $(X_i \times_Y X_i) \times_Y W \to
X_i \times_Y X_i$. Comme $Y$ est quasi-séparé, les espaces algébriques
$W \times_Y X_i$ sont quasi-compacts (et aussi quasi-séparés).
On peut donc effectuer le changement de base à $W$ et supposer que $Y$ est un schéma affine.
Lorsque $Y$ est un schéma affine, il faut montrer que $X_i$ est un espace
algébrique séparé pour $i$ assez grand, sachant que $X$ est un espace
algébrique séparé. Ce cas résulte donc du
lemme \ref{spaces-limits-lemma-descend-separated}.
\end{proof}

\begin{lemma}
\label{spaces-limits-lemma-eventually-affine}
Soit $S$ un schéma. Soit $Y$ un espace algébrique sur $S$.
Soit $X = \lim X_i$ la limite d'un système projectif filtrant d'espaces
algébriques sur $Y$ à morphismes de transition affines. Supposons que
\begin{enumerate}
\item $Y$ soit quasi-compact et quasi-séparé,
\item $X_i$ soit quasi-compact et quasi-séparé,
\item $X \to Y$ soit affine.
\end{enumerate}
Alors $X_i \to Y$ est affine pour $i$ assez grand.
\end{lemma}

\begin{proof}
Choisissons un schéma affine $W$ et un morphisme étale surjectif $W \to Y$.
Alors $X \times_Y W$ est affine, et il suffit de vérifier que
$X_i \times_Y W$ est affine pour un certain $i$ (Morphismes d'espaces,
lemme \ref{spaces-morphisms-lemma-affine-local}).
Cela résulte du lemme \ref{spaces-limits-lemma-limit-is-affine}.
\end{proof}

\begin{lemma}
\label{spaces-limits-lemma-eventually-finite}
Soit $S$ un schéma. Soit $Y$ un espace algébrique sur $S$.
Soit $X = \lim X_i$ la limite d'un système projectif filtrant d'espaces
algébriques sur $Y$ à morphismes de transition affines. Supposons que
\begin{enumerate}
\item $Y$ soit quasi-compact et quasi-séparé,
\item $X_i$ soit quasi-compact et quasi-séparé,
\item les morphismes de transition $X_{i'} \to X_i$ soient finis,
\item $X_i \to Y$ soit localement de type fini,
\item $X \to Y$ soit intégral.
\end{enumerate}
Alors $X_i \to Y$ est fini pour $i$ assez grand.
\end{lemma}

\begin{proof}
Choisissons un schéma affine $W$ et un morphisme étale surjectif $W \to Y$.
Alors $X \times_Y W$ est fini sur $W$, et il suffit de vérifier que
$X_i \times_Y W$ est fini sur $W$ pour un certain $i$ (Morphismes d'espaces,
lemme \ref{spaces-morphisms-lemma-integral-local}). D'après le
lemme \ref{spaces-limits-lemma-limit-is-scheme}, on est ramené au cas des schémas.
Dans ce cas, le résultat découle de
Limites, lemme \ref{limits-lemma-eventually-finite}.
\end{proof}

\begin{lemma}
\label{spaces-limits-lemma-eventually-closed-immersion}
Soit $S$ un schéma. Soit $Y$ un espace algébrique sur $S$.
Soit $X = \lim X_i$ la limite d'un système projectif filtrant d'espaces
algébriques sur $Y$ à morphismes de transition affines. Supposons que
\begin{enumerate}
\item $Y$ soit quasi-compact et quasi-séparé,
\item $X_i$ soit quasi-compact et quasi-séparé,
\item les morphismes de transition $X_{i'} \to X_i$ soient des immersions fermées,
\item $X_i \to Y$ soit localement de type fini,
\item $X \to Y$ soit une immersion fermée.
\end{enumerate}
Alors $X_i \to Y$ est une immersion fermée pour $i$ assez grand.
\end{lemma}

\begin{proof}
Choisissons un schéma affine $W$ et un morphisme étale surjectif $W \to Y$.
Alors $X \times_Y W$ est un sous-espace fermé de $W$, et il suffit de vérifier
que $X_i \times_Y W$ est un sous-espace fermé de $W$ pour un certain $i$
(Morphismes d'espaces,
lemme \ref{spaces-morphisms-lemma-closed-immersion-local}). D'après le
lemme \ref{spaces-limits-lemma-limit-is-scheme}, on est ramené au cas des schémas.
Dans ce cas, le résultat découle de
Limites, lemme \ref{limits-lemma-eventually-closed-immersion}.
\end{proof}

















\section{Descente de propriétés des morphismes}
\label{spaces-limits-section-descent-of-properties}

\noindent
Cette section est l'analogue de la section \ref{spaces-limits-section-descent}
pour les propriétés des morphismes. Nous nous placerons dans la situation suivante.

\begin{situation}
\label{spaces-limits-situation-descent-property}
Soit $S$ un schéma. Soit $B = \lim B_i$ la limite d'un système projectif filtrant
d'espaces algébriques sur $S$ à morphismes de transition affines
(lemme \ref{spaces-limits-lemma-directed-inverse-system-has-limit}).
Soit $0 \in I$ et soit $f_0 : X_0 \to Y_0$ un morphisme d'espaces algébriques
sur $B_0$. Supposons $B_0$, $X_0$, $Y_0$ quasi-compacts et quasi-séparés.
Soit $f_i : X_i \to Y_i$ le changement de base de $f_0$ à $B_i$ et
soit $f : X \to Y$ le changement de base de $f_0$ à $B$.
\end{situation}

\begin{lemma}
\label{spaces-limits-lemma-descend-etale}
Notations et hypothèses comme dans la
situation \ref{spaces-limits-situation-descent-property}. Si
\begin{enumerate}
\item $f$ est étale,
\item $f_0$ est localement de présentation finie,
\end{enumerate}
alors $f_i$ est étale pour un certain $i \geq 0$.
\end{lemma}

\begin{proof}
Choisissons un schéma affine $V_0$ et un morphisme étale surjectif
$V_0 \to Y_0$. Choisissons un schéma affine $U_0$ et un morphisme étale
surjectif $U_0 \to V_0 \times_{Y_0} X_0$. On a le diagramme
$$
\xymatrix{
U_0 \ar[d] \ar[r] & V_0 \ar[d] \\
X_0 \ar[r] & Y_0
}
$$
Les flèches verticales sont étales et surjectives par construction.
Après changement de base de ce diagramme à $B_i$ ou à $B$, on obtient
$$
\vcenter{
\xymatrix{
U_i \ar[d] \ar[r] & V_i \ar[d] \\
X_i \ar[r] & Y_i
}
}
\quad\text{et}\quad
\vcenter{
\xymatrix{
U \ar[d] \ar[r] & V \ar[d] \\
X \ar[r] & Y
}
}
$$
Remarquons que $U_i, V_i, U, V$ sont des schémas affines,
que les morphismes verticaux sont étales et surjectifs, et que la limite des
morphismes $U_i \to V_i$ est $U \to V$. Rappelons que $X_i \to Y_i$ est étale
si et seulement si $U_i \to V_i$ est
étale et, de même, que $X \to Y$ est étale si et seulement si
$U \to V$ est étale
(Morphismes des espaces, lemme \ref{spaces-morphisms-lemma-etale-local}).
Comme $f_0$ est localement de présentation
finie, le morphisme $U_0 \to V_0$ l'est aussi. Le lemme résulte donc
de Limites, lemme \ref{limits-lemma-descend-etale}.
\end{proof}

\begin{lemma}
\label{spaces-limits-lemma-descend-smooth}
Notations et hypothèses comme dans la
situation \ref{spaces-limits-situation-descent-property}. Si
\begin{enumerate}
\item $f$ est lisse,
\item $f_0$ est localement de présentation finie,
\end{enumerate}
alors $f_i$ est lisse pour un certain $i \geq 0$.
\end{lemma}

\begin{proof}
Choisissons un schéma affine $V_0$ et un morphisme étale surjectif
$V_0 \to Y_0$. Choisissons un schéma affine $U_0$ et un morphisme étale
surjectif $U_0 \to V_0 \times_{Y_0} X_0$. On a le diagramme
$$
\xymatrix{
U_0 \ar[d] \ar[r] & V_0 \ar[d] \\
X_0 \ar[r] & Y_0
}
$$
Les flèches verticales sont étales et surjectives par construction.
Après changement de base de ce diagramme à $B_i$ ou à $B$, on obtient
$$
\vcenter{
\xymatrix{
U_i \ar[d] \ar[r] & V_i \ar[d] \\
X_i \ar[r] & Y_i
}
}
\quad\text{et}\quad
\vcenter{
\xymatrix{
U \ar[d] \ar[r] & V \ar[d] \\
X \ar[r] & Y
}
}
$$
Remarquons que $U_i, V_i, U, V$ sont des schémas affines,
que les morphismes verticaux sont étales et surjectifs, et que la limite des
morphismes $U_i \to V_i$ est $U \to V$. Rappelons que $X_i \to Y_i$ est lisse
si et seulement si $U_i \to V_i$ est lisse et, de même,
$X \to Y$ est lisse si et seulement si $U \to V$ est lisse
(Morphismes des espaces, définition \ref{spaces-morphisms-definition-smooth}).
Comme $f_0$ est localement de présentation
finie, le morphisme $U_0 \to V_0$ l'est aussi. Le lemme résulte donc
de Limites, lemme \ref{limits-lemma-descend-smooth}.
\end{proof}

\begin{lemma}
\label{spaces-limits-lemma-descend-surjective}
Notations et hypothèses comme dans la
situation \ref{spaces-limits-situation-descent-property}. Si
\begin{enumerate}
\item $f$ est surjectif,
\item $f_0$ est localement de présentation finie,
\end{enumerate}
alors $f_i$ est surjectif pour un certain $i \geq 0$.
\end{lemma}

\begin{proof}
Choisissons un schéma affine $V_0$ et un morphisme étale surjectif
$V_0 \to Y_0$. Choisissons un schéma affine $U_0$ et un morphisme étale
surjectif $U_0 \to V_0 \times_{Y_0} X_0$. On a le diagramme
$$
\xymatrix{
U_0 \ar[d] \ar[r] & V_0 \ar[d] \\
X_0 \ar[r] & Y_0
}
$$
Les flèches verticales sont étales et surjectives par construction.
Après changement de base de ce diagramme à $B_i$ ou à $B$, on obtient
$$
\vcenter{
\xymatrix{
U_i \ar[d] \ar[r] & V_i \ar[d] \\
X_i \ar[r] & Y_i
}
}
\quad\text{et}\quad
\vcenter{
\xymatrix{
U \ar[d] \ar[r] & V \ar[d] \\
X \ar[r] & Y
}
}
$$
Remarquons que $U_i, V_i, U, V$ sont des schémas affines, que les morphismes verticaux sont
étales et surjectifs, que la limite des morphismes $U_i \to V_i$ est $U \to V$,
et que les morphismes $U_i \to X_i \times_{Y_i} V_i$ et
$U \to X \times_Y V$ sont surjectifs (car ce sont les changements de base de
$U_0 \to X_0 \times_{Y_0} V_0$). En particulier, on voit que
$X_i \to Y_i$ est surjectif si et seulement si $U_i \to V_i$ est surjectif
et, de même, que $X \to Y$ est surjectif si et seulement si $U \to V$ est surjectif.
Comme $f_0$ est localement de présentation
finie, le morphisme $U_0 \to V_0$ l'est aussi. Le lemme résulte donc
du cas des schémas (Limites, lemme \ref{limits-lemma-descend-surjective}).
\end{proof}

\begin{lemma}
\label{spaces-limits-lemma-descend-universally-injective}
Notations et hypothèses comme dans la situation \ref{spaces-limits-situation-descent-property}. Si
\begin{enumerate}
\item $f$ est universellement injectif,
\item $f_0$ est localement de type fini,
\end{enumerate}
alors $f_i$ est universellement injectif pour un certain $i \geq 0$.
\end{lemma}

\begin{proof}
Rappelons qu'un morphisme $X \to Y$ est universellement injectif si et
seulement si la diagonale $X \to X \times_Y X$ est surjective
(Morphismes des espaces, définition
\ref{spaces-morphisms-definition-universally-injective} et
lemme \ref{spaces-morphisms-lemma-universally-injective}).
Remarquons que $X_0 \to X_0 \times_{Y_0} X_0$ est localement de
présentation finie (Morphismes des espaces, lemme
\ref{spaces-morphisms-lemma-diagonal-morphism-finite-type}).
Le lemme résulte donc du lemme \ref{spaces-limits-lemma-descend-surjective}
appliqué au morphisme $X_0 \to X_0 \times_{Y_0} X_0$.
\end{proof}

\begin{lemma}
\label{spaces-limits-lemma-descend-affine}
Notations et hypothèses comme dans la situation \ref{spaces-limits-situation-descent-property}. Si
$f$ est affine, alors $f_i$ est affine pour un certain $i \geq 0$.
\end{lemma}

\begin{proof}
Choisissons un schéma affine $V_0$ et un morphisme étale surjectif $V_0 \to Y_0$.
Posons $V_i = V_0 \times_{Y_0} Y_i$ et $V = V_0 \times_{Y_0} Y$.
Comme $f$ est affine, on voit que $V \times_Y X = \lim V_i \times_{Y_i} X_i$
est affine. D'après le lemme \ref{spaces-limits-lemma-limit-is-affine},
$V_i \times_{Y_i} X_i$ est affine pour un certain $i \geq 0$. Pour cet $i$, le morphisme
$f_i$ est affine
(Morphismes des espaces, lemme \ref{spaces-morphisms-lemma-affine-local}).
\end{proof}

\begin{lemma}
\label{spaces-limits-lemma-descend-finite}
Notations et hypothèses comme dans la situation \ref{spaces-limits-situation-descent-property}. Si
\begin{enumerate}
\item $f$ est fini,
\item $f_0$ est localement de type fini,
\end{enumerate}
alors $f_i$ est fini pour un certain $i \geq 0$.
\end{lemma}

\begin{proof}
Choisissons un schéma affine $V_0$ et un morphisme étale surjectif $V_0 \to Y_0$.
Posons $V_i = V_0 \times_{Y_0} Y_i$ et $V = V_0 \times_{Y_0} Y$.
Comme $f$ est fini, on voit que $V \times_Y X = \lim V_i \times_{Y_i} X_i$
est un schéma fini sur $V$. D'après le lemme \ref{spaces-limits-lemma-limit-is-affine},
$V_i \times_{Y_i} X_i$ est affine pour un certain $i \geq 0$. Quitte à augmenter $i$ si
nécessaire, on trouve que $V_i \times_{Y_i} X_i \to V_i$ est fini d'après
Limites, lemme \ref{limits-lemma-descend-finite-finite-presentation}.
Pour cet $i$, le morphisme $f_i$ est fini
(Morphismes des espaces, lemme \ref{spaces-morphisms-lemma-integral-local}).
\end{proof}

\begin{lemma}
\label{spaces-limits-lemma-descend-closed-immersion}
Notations et hypothèses comme dans la situation \ref{spaces-limits-situation-descent-property}. Si
\begin{enumerate}
\item $f$ est une immersion fermée,
\item $f_0$ est localement de type fini,
\end{enumerate}
alors $f_i$ est une immersion fermée pour un certain $i \geq 0$.
\end{lemma}

\begin{proof}
Choisissons un schéma affine $V_0$ et un morphisme étale surjectif $V_0 \to Y_0$.
Posons $V_i = V_0 \times_{Y_0} Y_i$ et $V = V_0 \times_{Y_0} Y$.
Comme $f$ est une immersion fermée, on voit que
$V \times_Y X = \lim V_i \times_{Y_i} X_i$
est un sous-schéma fermé du schéma affine $V$. D'après le
lemme \ref{spaces-limits-lemma-limit-is-affine},
$V_i \times_{Y_i} X_i$ est affine pour un certain $i \geq 0$. Quitte à augmenter $i$ si
nécessaire, on trouve que $V_i \times_{Y_i} X_i \to V_i$ est une immersion fermée d'après
Limites, lemme \ref{limits-lemma-descend-closed-immersion-finite-presentation}.
Pour cet $i$, le morphisme $f_i$ est une immersion fermée
(Morphismes des espaces, lemme \ref{spaces-morphisms-lemma-integral-local}).
\end{proof}

\begin{lemma}
\label{spaces-limits-lemma-descend-separated-morphism}
Notations et hypothèses comme dans la situation \ref{spaces-limits-situation-descent-property}.
Si $f$ est séparé, alors $f_i$ est séparé pour un certain $i \geq 0$.
\end{lemma}

\begin{proof}
Appliquons le lemme \ref{spaces-limits-lemma-descend-closed-immersion}
au morphisme diagonal $\Delta_{X_0/Y_0} : X_0 \to X_0 \times_{Y_0} X_0$.
(Les morphismes diagonaux sont localement de type fini
et le produit fibré $X_0 \times_{Y_0} X_0$ est quasi-compact et
quasi-séparé. Certains détails sont omis.)
\end{proof}

\begin{lemma}
\label{spaces-limits-lemma-descend-isomorphism}
Notations et hypothèses comme dans la situation \ref{spaces-limits-situation-descent-property}. Si
\begin{enumerate}
\item $f$ est un isomorphisme,
\item $f_0$ est localement de présentation finie,
\end{enumerate}
alors $f_i$ est un isomorphisme pour un certain $i \geq 0$.
\end{lemma}

\begin{proof}
Être un isomorphisme équivaut à être étale, universellement injectif
et surjectif, voir
Morphismes des espaces, lemme
\ref{spaces-morphisms-lemma-etale-universally-injective-open}.
Le lemme résulte donc des
lemmes \ref{spaces-limits-lemma-descend-etale},
\ref{spaces-limits-lemma-descend-surjective} et
\ref{spaces-limits-lemma-descend-universally-injective}.
\end{proof}

\begin{lemma}
\label{spaces-limits-lemma-descend-monomorphism}
Notations et hypothèses comme dans la situation \ref{spaces-limits-situation-descent-property}. Si
\begin{enumerate}
\item $f$ est un monomorphisme,
\item $f_0$ est localement de type fini,
\end{enumerate}
alors $f_i$ est un monomorphisme pour un certain $i \geq 0$.
\end{lemma}

\begin{proof}
Rappelons qu'un morphisme est un monomorphisme si et seulement si sa diagonale est
un isomorphisme. Le morphisme $X_0 \to X_0 \times_{Y_0} X_0$ est localement de
présentation finie d'après
Morphismes des espaces, lemme
\ref{spaces-morphisms-lemma-diagonal-morphism-finite-type}.
Comme $X_0 \times_{Y_0} X_0$ est quasi-compact et quasi-séparé,
on déduit du
lemme \ref{spaces-limits-lemma-descend-isomorphism}
que $\Delta_i : X_i \to X_i \times_{Y_i} X_i$ est un isomorphisme pour
un certain $i \geq 0$. Pour cet $i$, le morphisme $f_i$ est un monomorphisme.
\end{proof}

\begin{lemma}
\label{spaces-limits-lemma-descend-flat}
Notations et hypothèses comme dans la situation \ref{spaces-limits-situation-descent-property}.
Soit $\mathcal{F}_0$ un $\mathcal{O}_{X_0}$-module quasi-cohérent;
notons $\mathcal{F}_i$ son image inverse sur $X_i$ et $\mathcal{F}$
son image inverse sur $X$. Si
\begin{enumerate}
\item $\mathcal{F}$ est plat sur $Y$,
\item $\mathcal{F}_0$ est de présentation finie, et
\item $f_0$ est localement de présentation finie,
\end{enumerate}
alors $\mathcal{F}_i$ est plat sur $Y_i$ pour un certain $i \geq 0$.
En particulier, si $f_0$ est localement de présentation finie et si
$f$ est plat, alors $f_i$ est plat pour un certain $i \geq 0$.
\end{lemma}

\begin{proof}
Choisissons un schéma affine $V_0$ et un morphisme étale surjectif
$V_0 \to Y_0$. Choisissons un schéma affine $U_0$ et un morphisme étale
surjectif $U_0 \to V_0 \times_{Y_0} X_0$. On a le diagramme
$$
\xymatrix{
U_0 \ar[d] \ar[r] & V_0 \ar[d] \\
X_0 \ar[r] & Y_0
}
$$
Les flèches verticales sont étales et surjectives par construction.
Après changement de base de ce diagramme à $B_i$ ou à $B$, on obtient
$$
\vcenter{
\xymatrix{
U_i \ar[d] \ar[r] & V_i \ar[d] \\
X_i \ar[r] & Y_i
}
}
\quad\text{et}\quad
\vcenter{
\xymatrix{
U \ar[d] \ar[r] & V \ar[d] \\
X \ar[r] & Y
}
}
$$
Remarquons que $U_i, V_i, U, V$ sont des schémas affines, que les morphismes verticaux sont
étales et surjectifs, et que la limite des morphismes $U_i \to V_i$ est
$U \to V$. Rappelons que $\mathcal{F}_i$ est plat sur $Y_i$ si et seulement si
$\mathcal{F}_i|_{U_i}$ est plat sur $V_i$ et, de même, que $\mathcal{F}$ est plat
sur $Y$ si et seulement si $\mathcal{F}|_U$ est plat sur $V$
(Morphismes des espaces, définition \ref{spaces-morphisms-definition-flat}).
Comme $f_0$ est localement de présentation finie, le morphisme
$U_0 \to V_0$ l'est aussi. Le lemme résulte donc
de Limites, lemme \ref{limits-lemma-descend-module-flat-finite-presentation}.
\end{proof}

\begin{lemma}
\label{spaces-limits-lemma-eventually-proper}
Hypothèses et notations comme dans la situation \ref{spaces-limits-situation-descent-property}.
Si
\begin{enumerate}
\item $f$ est propre, et
\item $f_0$ est localement de type fini,
\end{enumerate}
alors il existe un $i$ tel que $f_i$ soit propre.
\end{lemma}

\begin{proof}
Choisissons un schéma affine $V_0$ et un morphisme étale surjectif $V_0 \to Y_0$.
Posons $V_i = Y_i \times_{Y_0} V_0$ et $V = Y \times_{Y_0} V_0$.
Il suffit de montrer que le changement de base de $f_i$ à $V_i$ est
propre, voir Morphismes des espaces, lemme
\ref{spaces-morphisms-lemma-proper-local}.
On peut donc supposer $Y_0$ affine.

\medskip\noindent
D'après le lemme \ref{spaces-limits-lemma-descend-separated-morphism}, on voit que
$f_i$ est séparé pour un certain $i \geq 0$. En remplaçant
$0$ par $i$, on peut supposer que $f_0$ est séparé.
Remarquons que $f_0$ est quasi-compact. Ainsi, $f_0$ est séparé et
de type fini. D'après
Cohomologie des espaces, lemme \ref{spaces-cohomology-lemma-weak-chow},
on peut choisir un diagramme
$$
\xymatrix{
X_0 \ar[rd] & X_0' \ar[d] \ar[l]^\pi \ar[r] & \mathbf{P}^n_{Y_0} \ar[dl] \\
& Y_0 &
}
$$
où $X_0' \to \mathbf{P}^n_{Y_0}$ est une immersion, et
$\pi : X_0' \to X_0$ est propre et surjectif. Introduisons
$X' = X_0' \times_{Y_0} Y$ et $X_i' = X_0' \times_{Y_0} Y_i$.
D'après Morphismes des espaces, lemmes
\ref{spaces-morphisms-lemma-composition-proper} et
\ref{spaces-morphisms-lemma-base-change-proper}
on voit que $X' \to Y$ est propre. Ainsi, $X' \to \mathbf{P}^n_Y$ est
une immersion fermée (Morphismes des espaces, lemme
\ref{spaces-morphisms-lemma-universally-closed-permanence}). D'après
Morphismes des espaces, lemme \ref{spaces-morphisms-lemma-image-proper-is-proper},
il suffit de montrer que $X'_i \to Y_i$ est propre pour un certain $i$.
D'après le lemme \ref{spaces-limits-lemma-descend-closed-immersion},
on trouve que $X'_i \to \mathbf{P}^n_{Y_i}$ est
une immersion fermée pour $i$ assez grand. Alors $X'_i \to Y_i$
est propre, ce qui conclut la démonstration.
\end{proof}

\begin{lemma}
\label{spaces-limits-lemma-eventually-relative-dimension}
Hypothèses et notations comme dans la situation \ref{spaces-limits-situation-descent-property}.
Soit $d \geq 0$. Si
\begin{enumerate}
\item $f$ est de dimension relative $\leq d$
(Morphismes des espaces, définition
\ref{spaces-morphisms-definition-relative-dimension}), et
\item $f_0$ est localement de type fini,
\end{enumerate}
alors il existe un $i$ tel que $f_i$ soit de dimension relative $\leq d$.
\end{lemma}

\begin{proof}
Choisissons un schéma affine $V_0$ et un morphisme étale surjectif
$V_0 \to Y_0$. Choisissons un schéma affine $U_0$ et un morphisme étale
surjectif $U_0 \to V_0 \times_{Y_0} X_0$. On a le diagramme
$$
\xymatrix{
U_0 \ar[d] \ar[r] & V_0 \ar[d] \\
X_0 \ar[r] & Y_0
}
$$
Les flèches verticales sont étales et surjectives par construction.
Après changement de base de ce diagramme à $B_i$ ou à $B$, on obtient
$$
\vcenter{
\xymatrix{
U_i \ar[d] \ar[r] & V_i \ar[d] \\
X_i \ar[r] & Y_i
}
}
\quad\text{et}\quad
\vcenter{
\xymatrix{
U \ar[d] \ar[r] & V \ar[d] \\
X \ar[r] & Y
}
}
$$
Remarquons que $U_i, V_i, U, V$ sont des schémas affines,
que les morphismes verticaux sont étales et surjectifs, et que la limite des
morphismes $U_i \to V_i$ est $U \to V$.
Dans cette situation, $X_i \to Y_i$ est de dimension relative $\leq d$
si et seulement si $U_i \to V_i$ est de dimension relative $\leq d$
(au sens de Morphismes, définition
\ref{morphisms-definition-relative-dimension-d}).
Pour voir l'équivalence, utilisons le fait que la définition pour les morphismes
d'espaces algébriques fait intervenir Morphismes des espaces, définition
\ref{spaces-morphisms-definition-dimension-fibre}
qui utilise la localisation étale. Il en va de même pour $X \to Y$ et $U \to V$.
Comme $f_0$ est localement de type fini, le morphisme $U_0 \to V_0$ l'est aussi.
Le lemme résulte donc du résultat plus général
Limites, lemme \ref{limits-lemma-limit-dimension}.
\end{proof}














\section{Descente des objets relatifs}
\label{spaces-limits-section-descending-relative}

\noindent
Le lemme suivant est représentatif du type de résultats de cette section.

\begin{lemma}
\label{spaces-limits-lemma-descend-finite-presentation}
Soit $S$ un schéma. Soit $I$ un ensemble préordonné filtrant.
Soit $(X_i, f_{ii'})$ un système projectif indexé par $I$ d'espaces algébriques
sur $S$. Supposons que
\begin{enumerate}
\item les morphismes $f_{ii'} : X_i \to X_{i'}$ soient affines,
\item les espaces $X_i$ soient quasi-compacts et quasi-séparés.
\end{enumerate}
Soit $X = \lim_i X_i$. Alors la catégorie des espaces algébriques
de présentation finie sur $X$ est la limite inductive sur $I$ des
catégories des espaces algébriques de présentation finie sur $X_i$.
\end{lemma}

\begin{proof}
Fixons $0 \in I$. Choisissons un morphisme étale surjectif $U_0 \to X_0$, où
$U_0$ est un schéma affine (Propriétés des espaces, lemme
\ref{spaces-properties-lemma-quasi-compact-affine-cover}).
Posons $U_i = X_i \times_{X_0} U_0$. Posons $R_0 = U_0 \times_{X_0} U_0$ et
$R_i = R_0 \times_{X_0} X_i$. Notons $s_i, t_i : R_i \to U_i$ et
$s, t : R \to U$ les deux projections. Dans la démonstration du
lemme \ref{spaces-limits-lemma-directed-inverse-system-has-limit}, on a
vu qu'il existe une présentation $X = U/R$ avec
$U = \lim U_i$ et $R = \lim R_i$. Remarquons que $U_i$ et $U$ sont affines et
que $R_i$ et $R$ sont quasi-compacts et séparés (puisque $X_i$ est
quasi-séparé). Soit $Y$ un espace algébrique sur $S$ et soit
$Y \to X$ un morphisme de présentation finie. Posons $V = U \times_X Y$.
C'est un espace algébrique de présentation finie sur $U$.
Choisissons un schéma affine $W$ et un morphisme étale surjectif $W \to V$.
Alors $W \to Y$ est également étale et surjectif. Posons $R' = W \times_Y W$,
de sorte que $Y = W/R'$ (voir Espaces, section \ref{spaces-section-presentations}).
Remarquons que $W$ est un schéma de présentation finie sur $U$ et que $R'$
est un schéma de présentation finie sur $R$ (détails omis).
D'après Limites, lemme \ref{limits-lemma-descend-finite-presentation},
on peut trouver un indice $i$ et un morphisme de schémas $W_i \to U_i$ de
présentation finie dont le changement de base à $U$ donne $W \to U$. De même,
on peut trouver, quitte à augmenter $i$, un schéma $R'_i$ de présentation
finie sur $R_i$ dont le changement de base à $R$ est $R'$.
Les morphismes de projection $s', t' : R' \to W$ sont des morphismes au-dessus des
morphismes de projection $s, t : R \to U$. On peut donc considérer $s'$,
resp.\ $t'$, comme un morphisme entre schémas de présentation finie sur
$U$ (le morphisme structural $R' \to U$ étant donné par $R' \to R$ suivi
de $s$, resp.\ $t$). On peut donc appliquer
de nouveau le lemme \ref{limits-lemma-descend-finite-presentation} de Limites,
pour voir que, quitte à augmenter $i$, il existe des
morphismes $s'_i, t'_i : R'_i \to W_i$ qui, après changement de base à $U$,
sont $S', t'$. D'après Limites, lemmes \ref{limits-lemma-descend-etale} et
\ref{limits-lemma-descend-monomorphism},
on peut supposer que $s'_i, t'_i$ sont étales et que
$j'_i : R'_i \to W_i \times_{X_i} W_i$ est un monomorphisme (ici, on
considère $j'_i$ comme un morphisme de schémas de présentation finie sur $U_i$ via
l'une des projections -- peu importe laquelle). En posant
$Y_i = W_i/R'_i$ (voir Espaces, théorème \ref{spaces-theorem-presentation}),
on obtient un espace algébrique de présentation finie
sur $X_i$ dont le changement de base à $X$ est isomorphe à $Y$.

\medskip\noindent
Cela montre que tout espace algébrique de présentation finie sur $X$ provient
d'un espace algébrique de présentation finie sur un certain $X_i$, c'est-à-dire
que le foncteur du lemme est essentiellement surjectif. Pour
montrer qu'il est pleinement fidèle, considérons un indice $0 \in I$ et deux
espaces algébriques $Y_0, Z_0$ de présentation finie sur $X_0$.
Posons $Y_i = X_i \times_{X_0} Y_0$, $Y = X \times_{X_0} Y_0$,
$Z_i = X_i \times_{X_0} Z_0$ et $Z = X \times_{X_0} Z_0$. Soit
$\alpha : Y \to Z$ un morphisme d'espaces algébriques sur $X$.
Choisissons un morphisme étale surjectif $V_0 \to Y_0$, où $V_0$ est
un schéma affine. Posons $V_i = V_0 \times_{Y_0} Y_i$ et
$V = V_0 \times_{Y_0} Y$, qui sont des schémas affines munis de
morphismes étales surjectifs vers $Y_i$ et $Y$. Le composé
$V \to Y \to Z \to Z_0$ provient d'un morphisme (essentiellement unique)
$V_i \to Z_0$ pour un certain $i \geq 0$ d'après la
proposition \ref{spaces-limits-proposition-characterize-locally-finite-presentation}
(appliquée à $Z_0 \to X_0$, qui est de présentation finie par hypothèse).
Quitte à augmenter $i$, les deux composés
$$
V_i \times_{Y_i} V_i \to V_i \to Z_0
$$
sont égaux, puisque c'est vrai à la limite. On obtient donc un morphisme
(essentiellement unique) $Y_i \to Z_0$. Comme c'est un morphisme sur $X_0$,
il induit un morphisme vers $Z_i = Z_0 \times_{X_0} X_i$, comme souhaité.
\end{proof}

\begin{lemma}
\label{spaces-limits-lemma-descend-modules-finite-presentation}
Avec les notations et hypothèses du
lemme \ref{spaces-limits-lemma-descend-finite-presentation}.
La catégorie des $\mathcal{O}_X$-modules de présentation finie est la
limite inductive sur $I$ des catégories des $\mathcal{O}_{X_i}$-modules de
présentation finie.
\end{lemma}

\begin{proof}
Choisissons $0 \in I$. Choisissons un schéma affine $U_0$ et un morphisme
étale surjectif $U_0 \to X_0$. Posons $U_i = X_i \times_{X_0} U_0$.
Posons $R_0 = U_0 \times_{X_0} U_0$ et $R_i = R_0 \times_{X_0} X_i$.
Notons $s_i, t_i : R_i \to U_i$ et $s, t : R \to U$ les deux
projections. Dans la démonstration du
lemme \ref{spaces-limits-lemma-directed-inverse-system-has-limit}, on a
vu qu'il existe une présentation $X = U/R$ avec
$U = \lim U_i$ et $R = \lim R_i$. Remarquons que $U_i$ et $U$ sont affines et
que $R_i$ et $R$ sont quasi-compacts et séparés (puisque $X_i$ est
quasi-séparé). En outre, on a également
$R \times_{s, U, t} R = \colim R_i \times_{s_i, U_i, t_i} R_i$.
On sait donc que
$\QCoh(\mathcal{O}_U) = \colim \QCoh(\mathcal{O}_{U_i})$,
$\QCoh(\mathcal{O}_R) = \colim \QCoh(\mathcal{O}_{R_i})$,
et
$\QCoh(\mathcal{O}_{R \times_{s, U, t} R}) =
\colim \QCoh(\mathcal{O}_{R_i \times_{s_i, U_i, t_i} R_i})$
d'après Limites, lemme \ref{limits-lemma-descend-modules-finite-presentation}.
On a $\QCoh(\mathcal{O}_X) = \QCoh(U, R, s, t, c)$ et
$\QCoh(\mathcal{O}_{X_i}) = \QCoh(U_i, R_i, s_i, t_i, c_i)$,
voir Propriétés des espaces, proposition
\ref{spaces-properties-proposition-quasi-coherent}.
Le résultat en découle donc formellement.
\end{proof}

\begin{lemma}
\label{spaces-limits-lemma-descend-invertible-modules}
Avec les notations et hypothèses du
lemme \ref{spaces-limits-lemma-descend-finite-presentation}. Alors
\begin{enumerate}
\item tout $\mathcal{O}_X$-module localement libre de type fini est l'image inverse
d'un $\mathcal{O}_{X_i}$-module localement libre de type fini pour un certain $i$,
\item tout $\mathcal{O}_X$-module inversible est l'image inverse d'un
$\mathcal{O}_{X_i}$-module inversible pour un certain $i$.
\end{enumerate}
\end{lemma}

\begin{proof}
Démonstration de (2).
Soit $\mathcal{L}$ un $\mathcal{O}_X$-module inversible. Comme les
modules inversibles sont de présentation finie, on peut trouver un $i$
et des modules $\mathcal{L}_i$ et $\mathcal{N}_i$ de présentation finie
sur $X_i$ tels que $f_i^*\mathcal{L}_i \cong \mathcal{L}$ et
$f_i^*\mathcal{N}_i \cong \mathcal{L}^{\otimes -1}$, voir le
lemme \ref{spaces-limits-lemma-descend-modules-finite-presentation}.
Comme l'image inverse commute au produit tensoriel, on voit que
$f_i^*(\mathcal{L}_i \otimes_{\mathcal{O}_{X_i}} \mathcal{N}_i)$
est isomorphe à $\mathcal{O}_X$. Comme le produit tensoriel de
modules de présentation finie est de présentation finie, le même
lemme montre que
$f_{i'i}^*\mathcal{L}_i
\otimes_{\mathcal{O}_{X_{i'}}} f_{i'i}^*\mathcal{N}_i$
est isomorphe à $\mathcal{O}_{X_{i'}}$ pour un certain $i' \geq i$.
Il s'ensuit que $f_{i'i}^*\mathcal{L}_i$ est inversible
(Modules sur les sites, lemme \ref{sites-modules-lemma-invertible})
et la démonstration est terminée.

\medskip\noindent
Démonstration de (1). Omise. Indication : raisonner comme dans la démonstration de (2)
en utilisant qu'un module (sur un site localement annelé) est localement libre
de type fini si et seulement s'il admet un dual, voir
Modules sur les sites, section \ref{sites-modules-section-duals}.
On peut aussi raisonner comme dans la démonstration pour les schémas, voir
Limites, lemme \ref{limits-lemma-descend-invertible-modules}.
\end{proof}














\section{Approximation noethérienne absolue}
\label{spaces-limits-section-approximation}

\noindent
Le résultat suivant est \cite[théorème 1.2.2]{CLO}.
Un ingrédient essentiel de la démonstration est
Espaces décents, lemme
\ref{decent-spaces-lemma-filter-quasi-compact-quasi-separated}.

\begin{proposition}
\label{spaces-limits-proposition-approximate}
\begin{reference}
Notre démonstration suit de près celle donnée dans \cite[théorème 1.2.2]{CLO}.
\end{reference}
Soit $X$ un espace algébrique quasi-compact et quasi-séparé sur
$\Spec(\mathbf{Z})$. Il existe un ensemble préordonné filtrant $I$
et un système projectif d'espaces algébriques $(X_i, f_{ii'})$ indexé par $I$
tels que
\begin{enumerate}
\item les morphismes de transition $f_{ii'}$ soient affines,
\item chaque $X_i$ soit quasi-séparé et de type fini sur
$\mathbf{Z}$, et
\item $X = \lim X_i$.
\end{enumerate}
\end{proposition}

\begin{proof}
Appliquons Espaces décents, lemme
\ref{decent-spaces-lemma-filter-quasi-compact-quasi-separated}
pour obtenir des sous-espaces ouverts $U_p \subset X$, des schémas $V_p$ et des morphismes
$f_p : V_p \to U_p$ possédant les propriétés énoncées. Remarquons que
$f_n : V_n \to U_n$ est un morphisme étale d'espaces algébriques
dont la restriction à l'image inverse de $T_n = (V_n)_{red}$ est un
isomorphisme. Ainsi $f_n$ est un isomorphisme, par exemple d'après
Morphismes des espaces, lemme
\ref{spaces-morphisms-lemma-etale-universally-injective-open}.
En particulier, $U_n$ est un schéma quasi-compact et séparé.
On peut donc écrire $U_n = \lim U_{n, i}$ comme limite projective filtrante
de schémas de type fini sur $\mathbf{Z}$ dont les morphismes de transition sont
affines, voir Limites, proposition \ref{limits-proposition-approximate}.
En procédant par récurrence descendante sur $p$, on voit donc qu'on a ramené
la question au problème posé dans le paragraphe suivant.

\medskip\noindent
On dispose ici de $U \subset X$, $U = \lim U_i$, $Z \subset X$ et
$f : V \to X$ vérifiant les propriétés suivantes
\begin{enumerate}
\item $X$ est un espace algébrique quasi-compact et quasi-séparé,
\item $V$ est un schéma quasi-compact et séparé,
\item $U \subset X$ est un sous-espace ouvert quasi-compact,
\item $(U_i, g_{ii'})$ est un système projectif filtrant d'espaces algébriques
quasi-séparés
de type fini sur $\mathbf{Z}$, à morphismes de transition affines,
dont la limite est $U$,
\item $Z \subset X$ est un sous-espace fermé tel que $|X| = |U| \amalg |Z|$,
\item $f : V \to X$ est un morphisme étale surjectif tel que
$f^{-1}(Z) \to Z$ soit un isomorphisme.
\end{enumerate}
Problème : montrer que la conclusion de la proposition vaut pour $X$.

\medskip\noindent
Remarquons que $W = f^{-1}(U) \subset V$ est un sous-schéma ouvert quasi-compact
étale sur $U$. On peut donc appliquer les
lemmes \ref{spaces-limits-lemma-descend-finite-presentation} et \ref{spaces-limits-lemma-descend-etale}
pour trouver un indice $0 \in I$ et un morphisme étale $W_0 \to U_0$
de présentation finie dont le changement de base à $U$ donne $W$. En posant
$W_i = W_0 \times_{U_0} U_i$, on voit que $W = \lim_{i \geq 0} W_i$. Quitte à
augmenter $0$, on peut supposer que les $W_i$ sont des schémas, voir le
lemme \ref{spaces-limits-lemma-limit-is-scheme}.
En outre, $W_i$ est de type fini sur $\mathbf{Z}$.

\medskip\noindent
Appliquons Limites, lemme \ref{limits-lemma-approximate}, à
$W = \lim_{i \geq 0} W_i$ et à l'inclusion $W \subset V$. Remplaçons $I$
par l'ensemble préordonné filtrant $J$ fourni par ce lemme. Cela permet
d'écrire $V$ comme limite projective filtrante $V = \lim V_i$ de schémas de type fini sur
$\mathbf{Z}$ à morphismes de transition affines, tels que chaque $V_i$ contienne
$W_i$ comme sous-schéma ouvert (de façon compatible avec les morphismes de transition).
Pour chaque $i$, on peut former la somme amalgamée
$$
\xymatrix{
W_i \ar[r] \ar[d]_\Delta & V_i \ar[d] \\
W_i \times_{U_i} W_i \ar[r] & R_i
}
$$
dans la catégorie des schémas. En effet, la flèche verticale de gauche et la flèche horizontale supérieure
sont des immersions ouvertes de schémas. Autrement dit, on peut construire
$R_i$ par recollement de $V_i$ et de $W_i \times_{U_i} W_i$ le long de l'ouvert commun
$W_i$ (voir Schémas, section \ref{schemes-section-glueing-schemes}). Remarquons que
les projections étales $W_i \times_{U_i} W_i \to W_i$ se prolongent
en morphismes étales $s_i, t_i : R_i \to V_i$. Il est clair que le
morphisme $j_i = (t_i, s_i) : R_i \to V_i \times V_i$ est une relation d'équivalence étale
sur $V_i$. Remarquons que $W_i \times_{U_i} W_i$ est
quasi-compact (puisque $U_i$ est quasi-séparé et $W_i$ quasi-compact)
et que $V_i$ est quasi-compact; ainsi $R_i$ est quasi-compact. Pour
$i \geq i'$, le diagramme
\begin{equation}
\label{spaces-limits-equation-cartesian}
\vcenter{
\xymatrix{
R_i \ar[r] \ar[d]_{s_i} & R_{i'} \ar[d]^{s_{i'}} \\
V_i \ar[r] & V_{i'}
}
}
\end{equation}
est cartésien car
$$
(W_{i'} \times_{U_{i'}} W_{i'}) \times_{U_{i'}} U_i =
W_{i'} \times_{U_{i'}} U_i \times_{U_i} U_i \times_{U_{i'}} W_{i'} =
W_i \times_{U_i} W_i.
$$
Considérons l'espace algébrique $X_i = V_i/R_i$ (voir
Espaces, théorème \ref{spaces-theorem-presentation}).
Comme $V_i$ est de type fini sur $\mathbf{Z}$ et que $R_i$ est quasi-compact,
on voit que $X_i$ est quasi-séparé et de type fini sur $\mathbf{Z}$
(voir
Propriétés des espaces, lemme \ref{spaces-properties-lemma-quasi-separated}
et
Morphismes des espaces, lemmes
\ref{spaces-morphisms-lemma-surjection-from-quasi-compact} et
\ref{spaces-morphisms-lemma-finite-type-local}).
Comme la construction de $R_i$ ci-dessus est compatible
avec les morphismes de transition, on obtient des morphismes d'espaces algébriques
$X_i \to X_{i'}$ pour $i \geq i'$. Les diagrammes commutatifs
$$
\xymatrix{
V_i \ar[r] \ar[d] & V_{i'} \ar[d] \\
X_i \ar[r] & X_{i'}
}
$$
sont cartésiens puisque (\ref{spaces-limits-equation-cartesian}) l'est, voir
Groupoïdes, lemme \ref{groupoids-lemma-criterion-fibre-product}.
Comme $V_i \to V_{i'}$ est affine, cela implique que $X_i \to X_{i'}$
est affine, voir
Morphismes des espaces, lemme \ref{spaces-morphisms-lemma-affine-local}.
On peut donc former la limite $X' = \lim X_i$ d'après le
lemme \ref{spaces-limits-lemma-directed-inverse-system-has-limit}.
Nous affirmons que $X \cong X'$, ce qui achève la démonstration de la proposition.

\medskip\noindent
Démonstration de l'affirmation. Posons $R = \lim R_i$.
Par construction, l'espace algébrique $X'$ est
muni d'un morphisme étale surjectif $V \to X'$ tel que
$$
V \times_{X'} V \cong R
$$
(utiliser le lemme \ref{spaces-limits-lemma-directed-inverse-system-has-limit}).
Par construction, $\lim W_i \times_{U_i} W_i = W \times_U W$ et $V = \lim V_i$,
de sorte que $R$ est la réunion de $W \times_U W$ et de $V$, recollés le long de $W$.
La propriété (6) implique que les projections $V \times_X V \to V$ sont des isomorphismes
au-dessus de $f^{-1}(Z) \subset V$. Ainsi, le schéma $V \times_X V$ est la réunion
des ouverts $\Delta_{V/X}(V)$ et $W \times_U W$, dont l'intersection
est $\Delta_{W/X}(W)$. On conclut qu'il existe un unique isomorphisme
$R \cong V \times_X V$ compatible avec les projections vers $V$.
Comme $V \to X$ et $V \to X'$ sont étales et surjectifs, on voit que
$$
X = V/ V \times_X V = V/R = V/V \times_{X'} V = X'
$$
d'après Espaces, lemme \ref{spaces-lemma-space-presentation}, ce qui conclut.
\end{proof}




\section{Applications}
\label{spaces-limits-section-applications}

\noindent
Le lemme suivant se déduit aussi directement de
Espaces décents, lemme
\ref{decent-spaces-lemma-filter-quasi-compact-quasi-separated}
sans passer par l'approximation noethérienne absolue.

\begin{lemma}
\label{spaces-limits-lemma-colimit-finitely-presented}
Soient $S$ un schéma et $X$ un espace algébrique quasi-compact et quasi-séparé
sur $S$. Tout $\mathcal{O}_X$-module quasi-cohérent est une
limite inductive filtrante de $\mathcal{O}_X$-modules de présentation finie.
\end{lemma}

\begin{proof}
On peut considérer $X$ comme un espace algébrique sur $\Spec(\mathbf{Z})$, voir
Espaces, définition \ref{spaces-definition-base-change}, et
Propriétés des espaces, définition \ref{spaces-properties-definition-separated}.
On peut donc appliquer la proposition \ref{spaces-limits-proposition-approximate}
et écrire $X = \lim X_i$, où $X_i$ est de présentation finie sur $\mathbf{Z}$.
Ainsi, $X_i$ est un espace algébrique noethérien, voir
Morphismes des espaces, lemme
\ref{spaces-morphisms-lemma-finite-presentation-noetherian}.
Le morphisme $X \to X_i$ est affine, voir le
lemme \ref{spaces-limits-lemma-directed-inverse-system-has-limit}.
On conclut par
Cohomologie des espaces, lemme
\ref{spaces-cohomology-lemma-direct-colimit-finite-presentation}.
\end{proof}

\noindent
Le reste de cette section consiste en des applications immédiates
du lemme \ref{spaces-limits-lemma-colimit-finitely-presented}.

\begin{lemma}
\label{spaces-limits-lemma-directed-colimit-finite-type}
Soient $S$ un schéma et $X$ un espace algébrique quasi-compact et quasi-séparé
sur $S$.
Soit $\mathcal{F}$ un $\mathcal{O}_X$-module quasi-cohérent.
Alors $\mathcal{F}$ est la limite inductive filtrante de ses sous-modules quasi-cohérents
de type fini.
\end{lemma}

\begin{proof}
Si $\mathcal{G}, \mathcal{H} \subset \mathcal{F}$ sont des
$\mathcal{O}_X$-sous-modules quasi-cohérents de type fini, alors l'image
de $\mathcal{G} \oplus \mathcal{H} \to \mathcal{F}$ est encore un
$\mathcal{O}_X$-sous-module quasi-cohérent de type fini qui les contient
tous deux. On voit ainsi que le système est filtrant.
Pour montrer que $\mathcal{F}$ est la limite inductive de ce système, écrivons
$\mathcal{F} = \colim_i \mathcal{F}_i$ comme
limite inductive filtrante de modules quasi-cohérents de présentation finie, comme dans le
lemme \ref{spaces-limits-lemma-colimit-finitely-presented}.
Alors les images $\mathcal{G}_i = \Im(\mathcal{F}_i \to \mathcal{F})$ sont des
sous-modules quasi-cohérents de type fini de $\mathcal{F}$. Comme
$\mathcal{F}$ est leur limite inductive, le résultat en découle.
\end{proof}

\begin{lemma}
\label{spaces-limits-lemma-finite-directed-colimit-surjective-maps}
Soient $S$ un schéma et $X$ un espace algébrique quasi-compact et quasi-séparé
sur $S$. Soit $\mathcal{F}$ un $\mathcal{O}_X$-module quasi-cohérent
de type fini. On peut alors écrire
$\mathcal{F} = \lim \mathcal{F}_i$, où chaque $\mathcal{F}_i$ est un
$\mathcal{O}_X$-module de présentation finie et tous les morphismes de transition
$\mathcal{F}_i \to \mathcal{F}_{i'}$ sont surjectifs.
\end{lemma}

\begin{proof}
Écrivons $\mathcal{F} = \colim \mathcal{G}_i$ comme limite inductive filtrante de
$\mathcal{O}_X$-modules de présentation finie
(lemme \ref{spaces-limits-lemma-colimit-finitely-presented}).
Nous affirmons que $\mathcal{G}_i \to \mathcal{F}$ est surjectif pour un certain $i$.
En effet, choisissons une surjection étale $U \to X$, où $U$ est un schéma affine.
Choisissons un nombre fini de sections $s_k \in \mathcal{F}(U)$ qui engendrent
$\mathcal{F}|_U$. Comme $U$ est affine, on voit que chaque $s_k$ appartient à l'image
de $\mathcal{G}_i \to \mathcal{F}$ pour $i$ assez grand. Ainsi,
$\mathcal{G}_i \to \mathcal{F}$ est surjectif pour $i$ assez grand.
Fixons un tel $i$ et soit $\mathcal{K} \subset \mathcal{G}_i$ le
noyau de l'application $\mathcal{G}_i \to \mathcal{F}$. Écrivons
$\mathcal{K} = \colim \mathcal{K}_a$
comme limite inductive filtrante de ses sous-modules quasi-cohérents de type fini
(lemme \ref{spaces-limits-lemma-directed-colimit-finite-type}). Alors
$\mathcal{F} = \colim \mathcal{G}_i/\mathcal{K}_a$ répond à la question
posée par le lemme.
\end{proof}

\noindent
Soit $X$ un espace algébrique. Dans le lemme suivant, on utilise la notion
d'une {\it $\mathcal{O}_X$-algèbre quasi-cohérente de présentation finie
$\mathcal{A}$}. Cela signifie que, pour tout schéma affine
$U = \Spec(R)$ étale sur $X$, on a $\mathcal{A}|_U = \widetilde{A}$,
où $A$ est une $R$-algèbre (commutative) de présentation finie
comme $R$-algèbre.

\begin{lemma}
\label{spaces-limits-lemma-algebra-directed-colimit-finite-presentation}
Soient $S$ un schéma et $X$ un espace algébrique quasi-compact et quasi-séparé
sur $S$.
Soit $\mathcal{A}$ une $\mathcal{O}_X$-algèbre quasi-cohérente.
Alors $\mathcal{A}$ est une limite inductive filtrante d'$\mathcal{O}_X$-algèbres
quasi-cohérentes de présentation finie.
\end{lemma}

\begin{proof}
Écrivons d'abord $\mathcal{A} = \colim_i \mathcal{F}_i$ comme limite inductive filtrante
de modules quasi-cohérents de présentation finie, comme dans le
lemme \ref{spaces-limits-lemma-colimit-finitely-presented}.
Pour chaque $i$, soit $\mathcal{B}_i = \text{Sym}(\mathcal{F}_i)$ l'algèbre
symétrique de $\mathcal{F}_i$ sur $\mathcal{O}_X$. Écrivons
$\mathcal{I}_i = \Ker(\mathcal{B}_i \to \mathcal{A})$. Écrivons
$\mathcal{I}_i = \colim_j \mathcal{F}_{i, j}$, où
$\mathcal{F}_{i, j}$ est un sous-module quasi-cohérent de type fini de
$\mathcal{I}_i$, voir le
lemme \ref{spaces-limits-lemma-directed-colimit-finite-type}.
Posons $\mathcal{I}_{i, j} \subset \mathcal{I}_i$
égal à l'idéal de $\mathcal{B}_i$ engendré par $\mathcal{F}_{i, j}$.
Posons $\mathcal{A}_{i, j} = \mathcal{B}_i/\mathcal{I}_{i, j}$.
Alors $\mathcal{A}_{i, j}$ est une $\mathcal{O}_X$-algèbre quasi-cohérente
de présentation finie. Définissons $(i, j) \leq (i', j')$ si
$i \leq i'$ et si l'application $\mathcal{B}_i \to \mathcal{B}_{i'}$
envoie l'idéal $\mathcal{I}_{i, j}$ dans l'idéal $\mathcal{I}_{i', j'}$.
Alors il est clair que $\mathcal{A} = \colim_{i, j} \mathcal{A}_{i, j}$.
\end{proof}

\noindent
Soit $X$ un espace algébrique. Dans le lemme suivant, on utilise la notion
d'une {\it $\mathcal{O}_X$-algèbre quasi-cohérente $\mathcal{A}$
de type fini}. Cela signifie que, pour tout schéma affine
$U = \Spec(R)$ étale sur $X$, on a $\mathcal{A}|_U = \widetilde{A}$,
où $A$ est une $R$-algèbre (commutative) de type fini
comme $R$-algèbre.

\begin{lemma}
\label{spaces-limits-lemma-algebra-directed-colimit-finite-type}
Soient $S$ un schéma et $X$ un espace algébrique quasi-compact et quasi-séparé
sur $S$. Soit $\mathcal{A}$ une $\mathcal{O}_X$-algèbre quasi-cohérente.
Alors $\mathcal{A}$ est la limite inductive filtrante de ses
$\mathcal{O}_X$-sous-algèbres quasi-cohérentes de type fini.
\end{lemma}

\begin{proof}
Omise. Indication : comparer avec la démonstration du
lemme \ref{spaces-limits-lemma-directed-colimit-finite-type}.
\end{proof}

\noindent
Soit $X$ un espace algébrique. Dans le lemme suivant, on utilise la notion
d'une {\it $\mathcal{O}_X$-algèbre quasi-cohérente finie (resp.\ entière)
$\mathcal{A}$}. Cela signifie que, pour tout schéma
affine $U = \Spec(R)$ étale sur $X$, on a
$\mathcal{A}|_U = \widetilde{A}$, où $A$ est une $R$-algèbre (commutative)
finie (resp.\ entière) comme $R$-algèbre.

\begin{lemma}
\label{spaces-limits-lemma-finite-algebra-directed-colimit-finite-finitely-presented}
Soient $S$ un schéma et $X$ un espace algébrique quasi-compact et quasi-séparé
sur $S$. Soit $\mathcal{A}$ une $\mathcal{O}_X$-algèbre quasi-cohérente
finie. Alors $\mathcal{A} = \colim \mathcal{A}_i$
est une limite inductive filtrante d'$\mathcal{O}_X$-algèbres quasi-cohérentes
finies et de présentation finie, à morphismes de transition surjectifs.
\end{lemma}

\begin{proof}
D'après le lemme \ref{spaces-limits-lemma-finite-directed-colimit-surjective-maps},
il existe un $\mathcal{O}_X$-module de présentation finie
$\mathcal{F}$ et une surjection $\mathcal{F} \to \mathcal{A}$.
La structure d'algèbre fournit une surjection
$$
\text{Sym}^*_{\mathcal{O}_X}(\mathcal{F}) \longrightarrow \mathcal{A}
$$
Notons $\mathcal{J}$ le noyau. Écrivons $\mathcal{J} = \colim \mathcal{E}_i$
comme limite inductive filtrante de $\mathcal{O}_X$-sous-modules de type fini
$\mathcal{E}_i$ (lemme \ref{spaces-limits-lemma-directed-colimit-finite-type}). Posons
$$
\mathcal{A}_i = \text{Sym}^*_{\mathcal{O}_X}(\mathcal{F})/(\mathcal{E}_i)
$$
où $(\mathcal{E}_i)$ désigne le faisceau d'idéaux engendré par
l'image de $\mathcal{E}_i \to \text{Sym}^*_{\mathcal{O}_X}(\mathcal{F})$.
Alors chaque $\mathcal{A}_i$ est une $\mathcal{O}_X$-algèbre de présentation finie,
les morphismes de transition sont surjectifs et $\mathcal{A} = \colim \mathcal{A}_i$.
Pour achever la démonstration, il reste
à montrer que $\mathcal{A}_i$ est une $\mathcal{O}_X$-algèbre finie
pour $i$ assez grand. Pour cela, choisissons un morphisme étale surjectif
$U \to X$, où $U$ est un schéma affine. Prenons des générateurs
$f_1, \ldots, f_m \in \Gamma(U, \mathcal{F})$.
Comme $\mathcal{A}(U)$ est une $\mathcal{O}_X(U)$-algèbre finie,
on voit que, pour chaque $j$, il existe un polynôme unitaire
$P_j \in \mathcal{O}(U)[T]$ tel que $P_j(f_j)$ soit nul dans $\mathcal{A}(U)$.
Comme $\mathcal{A} = \colim \mathcal{A}_i$ par construction,
on a $P_j(f_j) = 0$ dans $\mathcal{A}_i(U)$ pour tout $i$ assez grand.
Pour un tel $i$, les algèbres $\mathcal{A}_i$ sont finies.
\end{proof}

\begin{lemma}
\label{spaces-limits-lemma-integral-algebra-directed-colimit-finite}
Soient $S$ un schéma et $X$ un espace algébrique quasi-compact et quasi-séparé
sur $S$. Soit $\mathcal{A}$ une $\mathcal{O}_X$-algèbre quasi-cohérente
entière. Alors
\begin{enumerate}
\item $\mathcal{A}$ est la limite inductive filtrante de ses
$\mathcal{O}_X$-sous-algèbres quasi-cohérentes finies, et
\item $\mathcal{A}$ est une limite inductive filtrante d'$\mathcal{O}_X$-algèbres
finies et de présentation finie.
\end{enumerate}
\end{lemma}

\begin{proof}
D'après le lemme \ref{spaces-limits-lemma-algebra-directed-colimit-finite-type}, on a
$\mathcal{A} = \colim \mathcal{A}_i$, où
$\mathcal{A}_i \subset \mathcal{A}$ parcourt les
$\mathcal{O}_X$-sous-algèbres quasi-cohérentes de type fini.
Toute $\mathcal{O}_X$-sous-algèbre quasi-cohérente de type fini
de $\mathcal{A}$ est finie (utiliser Algèbre, lemme
\ref{algebra-lemma-characterize-finite-in-terms-of-integral},
sur les schémas affines étales sur $X$). Cela prouve (1).

\medskip\noindent
Pour prouver (2), écrivons $\mathcal{A} = \colim \mathcal{F}_i$
comme limite inductive de $\mathcal{O}_X$-modules de présentation finie au moyen du
lemme \ref{spaces-limits-lemma-colimit-finitely-presented}.
Pour chaque $i$, soit $\mathcal{J}_i$ le noyau de l'application
$$
\text{Sym}^*_{\mathcal{O}_X}(\mathcal{F}_i) \longrightarrow \mathcal{A}
$$
Pour $i' \geq i$, il existe une application induite $\mathcal{J}_i \to \mathcal{J}_{i'}$,
et l'on a $\mathcal{A} =
\colim \text{Sym}^*_{\mathcal{O}_X}(\mathcal{F}_i)/\mathcal{J}_i$.
En outre, les $\mathcal{O}_X$-algèbres quasi-cohérentes
$\text{Sym}^*_{\mathcal{O}_X}(\mathcal{F}_i)/\mathcal{J}_i$
sont finies (voir ci-dessus). Écrivons $\mathcal{J}_i = \colim \mathcal{E}_{ik}$
comme limite inductive de $\mathcal{O}_X$-modules de présentation finie.
Étant donnés $i' \geq i$ et $k$, il existe un $k'$ tel qu'on
dispose d'une application $\mathcal{E}_{ik} \to \mathcal{E}_{i'k'}$
rendant
$$
\xymatrix{
\mathcal{J}_i \ar[r] & \mathcal{J}_{i'} \\
\mathcal{E}_{ik} \ar[u] \ar[r] & \mathcal{E}_{i'k'} \ar[u]
}
$$
le diagramme commutatif. Cela résulte de
Cohomologie des espaces, lemme
\ref{spaces-cohomology-lemma-finite-presentation-quasi-compact-colimit}.
Cela induit une application
$$
\mathcal{A}_{ik} =
\text{Sym}^*_{\mathcal{O}_X}(\mathcal{F}_i)/(\mathcal{E}_{ik})
\longrightarrow
\text{Sym}^*_{\mathcal{O}_X}(\mathcal{F}_{i'})/(\mathcal{E}_{i'k'}) =
\mathcal{A}_{i'k'}
$$
où $(\mathcal{E}_{ik})$ désigne l'idéal engendré par $\mathcal{E}_{ik}$.
Les $\mathcal{O}_X$-algèbres quasi-cohérentes $\mathcal{A}_{ki}$
sont de présentation finie et finies pour $k$ assez grand
(voir la démonstration du
lemme \ref{spaces-limits-lemma-finite-algebra-directed-colimit-finite-finitely-presented}).
Enfin, on a
$$
\colim \mathcal{A}_{ik} = \colim \mathcal{A}_i = \mathcal{A}
$$
En effet, la première égalité a été démontrée dans la démonstration du
lemme \ref{spaces-limits-lemma-finite-algebra-directed-colimit-finite-finitely-presented},
et la seconde parce que $\mathcal{A}$ est la limite inductive des
modules $\mathcal{F}_i$.
\end{proof}

\begin{lemma}
\label{spaces-limits-lemma-extend}
Soit $S$ un schéma.
Soit $X$ un espace algébrique quasi-compact et quasi-séparé sur $S$.
Soit $U \subset X$ un sous-espace ouvert quasi-compact.
Soit $\mathcal{F}$ un $\mathcal{O}_X$-module quasi-cohérent.
Soit $\mathcal{G} \subset \mathcal{F}|_U$ un
$\mathcal{O}_U$-sous-module quasi-cohérent de type fini. Alors
il existe un sous-module quasi-cohérent $\mathcal{G}' \subset \mathcal{F}$
de type fini tel que $\mathcal{G}'|_U = \mathcal{G}$.
\end{lemma}

\begin{proof}
Notons $j : U \to X$ le morphisme d'inclusion. Comme $X$ est quasi-séparé
et $U$ quasi-compact, le morphisme $j$ est quasi-compact. Ainsi,
$j_*\mathcal{G} \subset j_*\mathcal{F}|_U$ sont des modules quasi-cohérents
sur $X$ (Morphismes des espaces, lemme
\ref{spaces-morphisms-lemma-pushforward}).
Posons $\mathcal{H} =
\Ker(j_*\mathcal{G} \oplus \mathcal{F} \to j_*\mathcal{F}|_U)$.
Alors $\mathcal{H}|_U = \mathcal{G}$. D'après le
lemme \ref{spaces-limits-lemma-directed-colimit-finite-type},
on peut trouver un sous-module quasi-cohérent de type fini
$\mathcal{H}' \subset \mathcal{H}$ tel que
$\mathcal{H}'|_U = \mathcal{H}|_U = \mathcal{G}$.
Posons $\mathcal{G}' = \Im(\mathcal{H}' \to \mathcal{F})$
pour conclure.
\end{proof}









\section{Approximation relative}
\label{spaces-limits-section-relative-approximation}

\noindent
On examine des variantes de la proposition \ref{spaces-limits-proposition-approximate}
sur une base.

\begin{lemma}
\label{spaces-limits-lemma-approximate-morphism}
Soit $f : X \to Y$ un morphisme d'espaces algébriques quasi-compacts et quasi-séparés
sur $\mathbf{Z}$.
Il existe un ensemble préordonné filtrant $I$ et un système projectif $(f_i : X_i \to Y_i)$
de morphismes d'espaces algébriques indexé par $I$, tel que les morphismes de transition
$X_i \to X_{i'}$ et $Y_i \to Y_{i'}$ soient affines, que $X_i$
et $Y_i$ soient quasi-séparés et de type fini sur $\mathbf{Z}$, et
que $(X \to Y) = \lim (X_i \to Y_i)$.
\end{lemma}

\begin{proof}
Écrivons $X = \lim_{a \in A} X_a$ et $Y = \lim_{b \in B} Y_b$ comme dans la
proposition \ref{spaces-limits-proposition-approximate}, c'est-à-dire avec $X_a$ et $Y_b$
quasi-séparés et de type fini sur $\mathbf{Z}$ et
des morphismes de transition affines.

\medskip\noindent
Fixons $b \in B$.
D'après le lemme \ref{spaces-limits-lemma-better-characterize-relative-limit-preserving}
appliqué à $Y_b$ et à $X = \lim X_a$ sur $\mathbf{Z}$,
il existe un $a \in A$ et un morphisme
$f_{a, b} : X_a \to Y_b$ rendant le diagramme
$$
\xymatrix{
X \ar[d] \ar[r] & Y \ar[d] \\
X_a \ar[r] & Y_b
}
$$
commutatif. Soit $I$ l'ensemble des triplets $(a, b, f_{a, b})$ qu'on
obtient ainsi.

\medskip\noindent
Soient $(a, b, f_{a, b})$ et $(a', b', f_{a', b'})$ dans $I$.
Soit $b'' \leq \min(b, b')$. D'après le
lemme \ref{spaces-limits-lemma-better-characterize-relative-limit-preserving},
il existe encore un $a'' \geq \max(a, a')$ tel que
les composés $X_{a''} \to X_a \to Y_b \to Y_{b''}$ et
$X_{a''} \to X_{a'} \to Y_{b'} \to Y_{b''}$ soient égaux.
Munissons $I$ du préordre
$$
(a, b, f_{a, b}) \geq (a', b', f_{a', b'})
\Leftrightarrow
a \geq a',\ b \geq b',\text{ et }
g_{b, b'} \circ f_{a, b} = f_{a', b'} \circ h_{a, a'}
$$
où $h_{a, a'} : X_a \to X_{a'}$ et $g_{b, b'} : Y_b \to Y_{b'}$
sont les morphismes de transition. Les remarques précédentes montrent que $I$
est filtrant et que les applications $I \to A$, $(a, b, f_{a, b}) \mapsto a$,
et $I \to B$, $(a, b, f_{a, b})$, sont cofinales. Si, pour $i = (a, b, f_{a, b})$,
on pose $X_i = X_a$, $Y_i = Y_b$ et $f_i = f_{a, b}$, on obtient
un système projectif de morphismes indexé par $I$ et l'on a
$$
\lim_{i \in I} X_i = \lim_{a \in A} X_a = X
\quad\text{et}\quad
\lim_{i \in I} S_i = \lim_{b \in B} Y_b = Y
$$
d'après Catégories, lemme \ref{categories-lemma-initial} (rappelons que
les limites sur $I$ sont en réalité les limites sur la catégorie opposée
associée à $I$, de sorte que cofinal devient initial).
Cela achève la démonstration.
\end{proof}

\begin{lemma}
\label{spaces-limits-lemma-relative-approximation}
Soient $S$ un schéma et $f : X \to Y$ un morphisme d'espaces algébriques
sur $S$. Supposons que
\begin{enumerate}
\item $X$ soit quasi-compact et quasi-séparé, et
\item $Y$ soit quasi-séparé.
\end{enumerate}
Alors $X = \lim X_i$ est la limite d'un système projectif filtrant d'espaces algébriques
$X_i$ de présentation finie sur $Y$, à morphismes de transition affines
sur $Y$.
\end{lemma}

\begin{proof}
Comme $|f|(|X|)$ est quasi-compact, on peut remplacer $Y$ par un sous-espace
ouvert quasi-compact dont l'ensemble des points contient $|f|(|X|)$. On peut donc supposer
$Y$ également quasi-compact.
D'après le lemme \ref{spaces-limits-lemma-approximate-morphism}, on peut écrire
$(X \to Y) = \lim (X_i \to Y_i)$ pour un certain système projectif filtrant
de morphismes d'espaces algébriques de type fini sur $\mathbf{Z}$, à
morphismes de transition affines. Comme les limites commutent aux limites
(Catégories, lemme \ref{categories-lemma-colimits-commute}),
on a $X = \lim X_i \times_{Y_i} Y$. Pour $i \geq i'$,
le morphisme de transition
$X_i \times_{Y_i} Y \to X_{i'} \times_{Y_{i'}} Y$ est affine, car c'est
le composé
$$
X_i \times_{Y_i} Y \to X_i \times_{Y_{i'}} Y \to X_{i'} \times_{Y_{i'}} Y
$$
où le premier morphisme est une immersion fermée (d'après
Morphismes des espaces, lemme
\ref{spaces-morphisms-lemma-fibre-product-after-map})
et le second est un changement de base d'un morphisme affine
(Morphismes des espaces, lemme \ref{spaces-morphisms-lemma-base-change-affine}),
et où le composé de morphismes affines est affine
(Morphismes des espaces, lemme \ref{spaces-morphisms-lemma-composition-affine}).
Les morphismes $f_i$ sont de présentation finie
(Morphismes des espaces, lemmes
\ref{spaces-morphisms-lemma-noetherian-finite-type-finite-presentation} et
\ref{spaces-morphisms-lemma-finite-presentation-permanence})
et leurs changements de base $X_i \times_{f_i, Y_i} Y \to Y$
sont donc de présentation finie
(Morphismes des espaces, lemme
\ref{spaces-morphisms-lemma-base-change-finite-presentation}).
\end{proof}






\section{Réalisation d'un morphisme de type fini comme fermé dans un morphisme de présentation finie}
\sectionmark{Type fini comme fermé dans une présentation finie}
\label{spaces-limits-section-finite-type-closed-in-finite-presentation}

\noindent
Cette section est l'analogue de la
section \ref{limits-section-finite-type-closed-in-finite-presentation} de Limites.

\begin{lemma}
\label{spaces-limits-lemma-affine-morphism-is-limit}
Soient $S$ un schéma et $f : X \to Y$ un morphisme affine d'espaces
algébriques sur $S$. Si $Y$ est quasi-compact et
quasi-séparé, alors $X$ est une limite projective filtrante $X = \lim X_i$,
où chaque $X_i$ est affine et de présentation finie sur $Y$.
\end{lemma}

\begin{proof}
Considérons le $\mathcal{O}_Y$-module quasi-cohérent
$\mathcal{A} = f_*\mathcal{O}_X$. D'après le
lemme \ref{spaces-limits-lemma-algebra-directed-colimit-finite-presentation},
on peut écrire $\mathcal{A} = \colim \mathcal{A}_i$ comme une limite
inductive filtrante d'algèbres de présentation finie
sur $\mathcal{O}_Y$, notées $\mathcal{A}_i$.
Posons $X_i = \underline{\Spec}_Y(\mathcal{A}_i)$, voir
Morphismes des espaces, définition
\ref{spaces-morphisms-definition-relative-spec}.
Par construction, $X_i \to Y$ est affine et de présentation finie
et $X = \lim X_i$.
\end{proof}

\begin{lemma}
\label{spaces-limits-lemma-integral-limit-finite-and-finite-presentation}
Soient $S$ un schéma et $f : X \to Y$ un morphisme entier d'espaces
algébriques sur $S$. Supposons $Y$ quasi-compact et quasi-séparé.
Alors $X$ s'écrit comme une limite projective filtrante $X = \lim X_i$,
où les $X_i$ sont finis et de présentation finie sur $Y$.
\end{lemma}

\begin{proof}
Considérons le $\mathcal{O}_Y$-module quasi-cohérent
$\mathcal{A} = f_*\mathcal{O}_X$. D'après le
lemme \ref{spaces-limits-lemma-integral-algebra-directed-colimit-finite},
on peut écrire $\mathcal{A} = \colim \mathcal{A}_i$ comme une limite
inductive filtrante de $\mathcal{O}_Y$-algèbres
$\mathcal{A}_i$ finies et de présentation finie.
Posons $X_i = \underline{\Spec}_Y(\mathcal{A}_i)$, voir
Morphismes des espaces, définition
\ref{spaces-morphisms-definition-relative-spec}.
Par construction, $X_i \to Y$ est fini et de présentation finie et
$X = \lim X_i$.
\end{proof}

\begin{lemma}
\label{spaces-limits-lemma-finite-in-finite-and-finite-presentation}
Soient $S$ un schéma et $f : X \to Y$ un morphisme fini d'espaces
algébriques sur $S$. Supposons $Y$ quasi-compact et quasi-séparé.
Alors $X$ s'écrit comme une limite projective filtrante $X = \lim X_i$,
où les morphismes de transition sont des immersions fermées et les objets
$X_i$ sont finis et de présentation finie sur $Y$.
\end{lemma}

\begin{proof}
Considérons le $\mathcal{O}_Y$-module quasi-cohérent de type fini
$\mathcal{A} = f_*\mathcal{O}_X$. D'après le
lemme \ref{spaces-limits-lemma-finite-algebra-directed-colimit-finite-finitely-presented},
on peut écrire $\mathcal{A} = \colim \mathcal{A}_i$ comme une limite
inductive filtrante de $\mathcal{O}_Y$-algèbres
$\mathcal{A}_i$ finies et de présentation finie, à morphismes de transition surjectifs.
Posons $X_i = \underline{\Spec}_Y(\mathcal{A}_i)$, voir
Morphismes des espaces, définition
\ref{spaces-morphisms-definition-relative-spec}.
Par construction, $X_i \to Y$ est fini et de présentation finie,
les morphismes de transition sont des immersions fermées et $X = \lim X_i$.
\end{proof}

\begin{lemma}
\label{spaces-limits-lemma-closed-is-limit-closed-and-finite-presentation}
\begin{slogan}
Les immersions fermées d'espaces algébriques qcqs peuvent être approchées
par des immersions fermées de présentation finie.
\end{slogan}
Soient $S$ un schéma et $f : X \to Y$ une immersion fermée d'espaces
algébriques sur $S$. Supposons $Y$ quasi-compact et quasi-séparé.
Alors $X$ s'écrit comme une limite projective filtrante $X = \lim X_i$,
où les morphismes de transition sont des immersions fermées et les morphismes
$X_i \to Y$ sont des immersions fermées de présentation finie.
\end{lemma}

\begin{proof}
Soit $\mathcal{I} \subset \mathcal{O}_Y$ le faisceau quasi-cohérent
d'idéaux définissant $X$ comme sous-espace fermé de $Y$. D'après le
lemme \ref{spaces-limits-lemma-directed-colimit-finite-type},
on peut écrire $\mathcal{I} = \colim \mathcal{I}_i$ comme la
limite inductive filtrante de ses sous-modules quasi-cohérents de type fini.
Soit $X_i$ le sous-espace fermé de $X$ défini par $\mathcal{I}_i$.
Alors $X_i \to Y$ est une immersion fermée de présentation finie
et $X = \lim X_i$. Certains détails sont omis.
\end{proof}

\begin{lemma}
\label{spaces-limits-lemma-quasi-affine-closed-in-finite-presentation}
Soient $S$ un schéma et $f : X \to Y$ un morphisme d'espaces algébriques
sur $S$. Supposons que
\begin{enumerate}
\item $f$ soit localement de type fini et quasi-affine, et
\item $Y$ soit quasi-compact et quasi-séparé.
\end{enumerate}
Alors il existe un morphisme de présentation finie
$f' : X' \to Y$ et une immersion fermée $X \to X'$ sur $Y$.
\end{lemma}

\begin{proof}
D'après Morphismes des espaces, lemme
\ref{spaces-morphisms-lemma-characterize-quasi-affine},
il existe une factorisation $X \to Z \to Y$ où
$X \to Z$ est une immersion ouverte quasi-compacte et
$Z \to Y$ est affine. Écrivons $Z = \lim Z_i$, où $Z_i$ est affine et
de présentation finie sur $Y$ (lemme \ref{spaces-limits-lemma-affine-morphism-is-limit}).
Pour un certain $0 \in I$, il existe un ouvert quasi-compact $U_0 \subset Z_0$
tel que $X$ soit isomorphe à l'image réciproque de $U_0$ dans $Z$
(lemme \ref{spaces-limits-lemma-descend-opens}). Soit $U_i$ l'image réciproque de
$U_0$ dans $Z_i$, de sorte que $U = \lim U_i$. D'après le
lemme \ref{spaces-limits-lemma-finite-type-eventually-closed},
le morphisme $X \to U_i$ est une immersion fermée pour un certain $i$ assez grand.
En posant $X' = U_i$, on conclut.
\end{proof}

\begin{lemma}
\label{spaces-limits-lemma-finite-type-closed-in-finite-presentation}
Soient $S$ un schéma et $f : X \to Y$ un morphisme d'espaces algébriques
sur $S$. Supposons :
\begin{enumerate}
\item $f$ soit localement de type fini,
\item $X$ soit quasi-compact et quasi-séparé, et
\item $Y$ soit quasi-compact et quasi-séparé.
\end{enumerate}
Alors il existe un morphisme de présentation finie
$f' : X' \to Y$ et une immersion fermée $X \to X'$ d'espaces
algébriques sur $Y$.
\end{lemma}

\begin{proof}
D'après la proposition \ref{spaces-limits-proposition-approximate},
on peut écrire $X = \lim_i X_i$, où $X_i$ est quasi-séparé et de type fini sur
$\mathbf{Z}$ et les morphismes de transition $f_{ii'} : X_i \to X_{i'}$ sont affines.
Considérons le diagramme commutatif
$$
\xymatrix{
X \ar[r] \ar[rd] & X_{i, Y} \ar[r] \ar[d] & X_i \ar[d] \\
& Y \ar[r] & \Spec(\mathbf{Z})
}
$$
Notons que $X_i$ est de présentation finie sur $\Spec(\mathbf{Z})$, voir
Morphismes des espaces, lemme
\ref{spaces-morphisms-lemma-noetherian-finite-type-finite-presentation}.
Le changement de base $X_{i, Y} \to Y$ est donc de présentation finie d'après
Morphismes des espaces, lemme
\ref{spaces-morphisms-lemma-base-change-finite-presentation}.
Observons que $\lim X_{i, Y} = X \times Y$ et que $X \to X \times Y$ est un
monomorphisme. D'après le lemme \ref{spaces-limits-lemma-finite-type-eventually-closed},
le morphisme $X \to X_{i, Y}$ est un monomorphisme pour $i$ assez grand.
Fixons un tel $i$. Notons que $X \to X_{i, Y}$ est localement de type fini
(Morphismes des espaces, lemme
\ref{spaces-morphisms-lemma-permanence-finite-type})
et est un monomorphisme, donc est séparé et localement quasi-fini
(Morphismes des espaces, lemme
\ref{spaces-morphisms-lemma-monomorphism-loc-finite-type-loc-quasi-finite}).
Le morphisme $X \to X_{i, Y}$ est donc représentable.
Le morphisme $X \to X_{i, Y}$ est alors quasi-affine, car on peut appliquer le
principe énoncé dans Espaces, lemme
\ref{spaces-lemma-representable-transformations-property-implication},
et le résultat pour les morphismes de schémas de Compléments sur les morphismes, lemme
\ref{more-morphisms-lemma-quasi-finite-separated-quasi-affine}.
Ainsi, le lemme \ref{spaces-limits-lemma-quasi-affine-closed-in-finite-presentation}
fournit une factorisation $X \to X' \to X_{i, Y}$,
où $X \to X'$ est une immersion fermée et $X' \to X_{i, Y}$ est de
présentation finie. Enfin, $X' \to Y$ est de présentation finie comme
composé de morphismes de présentation finie
(Morphismes des espaces, lemme
\ref{spaces-morphisms-lemma-composition-finite-presentation}).
\end{proof}

\begin{proposition}
\label{spaces-limits-proposition-separated-closed-in-finite-presentation}
Soient $S$ un schéma et $f : X \to Y$ un morphisme d'espaces algébriques
sur $S$. Supposons que
\begin{enumerate}
\item $f$ soit de type fini et séparé, et
\item $Y$ soit quasi-compact et quasi-séparé.
\end{enumerate}
Alors il existe un morphisme séparé de présentation finie
$f' : X' \to Y$ et une immersion fermée $X \to X'$ sur $Y$.
\end{proposition}

\begin{proof}
D'après le lemme \ref{spaces-limits-lemma-finite-type-closed-in-finite-presentation},
il existe une immersion fermée $X \to Z$, où $Z/Y$ est de
présentation finie. Soit $\mathcal{I} \subset \mathcal{O}_Z$
le faisceau quasi-cohérent d'idéaux définissant $X$ comme
sous-espace fermé de $Y$. D'après le
lemme \ref{spaces-limits-lemma-directed-colimit-finite-type},
on peut écrire $\mathcal{I}$ comme une limite inductive filtrante
$\mathcal{I} = \colim_{a \in A} \mathcal{I}_a$ de ses
faisceaux quasi-cohérents d'idéaux de type fini.
Soit $X_a \subset Z$ le sous-espace fermé défini par $\mathcal{I}_a$.
Ceux-ci forment un système projectif indexé par $A$.
Les morphismes de transition $X_a \to X_{a'}$ sont affines, car
ce sont des immersions fermées. Chaque $X_a$ est quasi-compact et quasi-séparé,
puisqu'il est un sous-espace fermé de $Z$ et que $Z$ est quasi-compact et
quasi-séparé d'après nos hypothèses.
On a $X = \lim_a X_a$, comme il résulte directement de
l'égalité $\mathcal{I} = \colim_{a \in A} \mathcal{I}_a$.
Chacun des morphismes $X_a \to Z$ est de présentation finie, voir
Morphismes, lemme \ref{morphisms-lemma-closed-immersion-finite-presentation}.
Par conséquent, les morphismes $X_a \to Y$ sont de présentation finie.
Il suffit donc de montrer que $X_a \to Y$ est séparé pour un certain
$a \in A$. Cela résulte du lemme \ref{spaces-limits-lemma-eventually-separated},
car on a supposé que $X \to Y$ était séparé.
\end{proof}






\section{Approximation des morphismes propres}
\label{spaces-limits-section-approximate-proper}

\begin{lemma}
\label{spaces-limits-lemma-proper-limit-of-proper-finite-presentation}
Soient $S$ un schéma et $f : X \to Y$ un morphisme propre d'espaces
algébriques sur $S$, avec $Y$ quasi-compact et quasi-séparé. Alors
$X = \lim X_i$ est une limite projective filtrante d'espaces algébriques $X_i$
propres et de présentation finie sur $Y$, dont les morphismes de transition
et les morphismes $X \to X_i$ sont des immersions fermées.
\end{lemma}

\begin{proof}
D'après la proposition \ref{spaces-limits-proposition-separated-closed-in-finite-presentation},
il existe une immersion fermée $X \to X'$, où $X'$ est séparé et de
présentation finie sur $Y$. D'après le
lemme \ref{spaces-limits-lemma-closed-is-limit-closed-and-finite-presentation},
on peut écrire $X = \lim X_i$, où $X_i \to X'$ est une immersion fermée de
présentation finie. Nous affirmons que, pour tout $i$ assez grand,
le morphisme $X_i \to Y$ est propre, ce qui achèvera la démonstration.

\medskip\noindent
Pour le montrer, on peut supposer que $Y$ est un schéma affine, voir
Morphismes des espaces, lemme \ref{spaces-morphisms-lemma-proper-local}.
Ensuite, on applique la version faible du lemme de Chow, voir
Cohomologie des espaces, lemme \ref{spaces-cohomology-lemma-weak-chow},
pour obtenir un diagramme
$$
\xymatrix{
X' \ar[rd] & X'' \ar[d] \ar[l]^\pi \ar[r] & \mathbf{P}^n_Y \ar[dl] \\
& Y &
}
$$
où $X'' \to \mathbf{P}^n_Y$ est une immersion et
$\pi : X'' \to X'$ est propre et surjectif. Notons
$X'_i \subset X''$, resp.\ $\pi^{-1}(X)$, l'image réciproque schématique de
$X_i \subset X'$, resp.\ $X \subset X'$.
Alors $\lim X'_i = \pi^{-1}(X)$. Comme $\pi^{-1}(X) \to Y$ est propre
(Morphismes des espaces, lemme \ref{spaces-morphisms-lemma-composition-proper}),
on voit que $\pi^{-1}(X) \to \mathbf{P}^n_Y$ est une immersion fermée
(Morphismes des espaces, lemmes
\ref{spaces-morphisms-lemma-universally-closed-permanence} et
\ref{spaces-morphisms-lemma-immersion-when-closed}).
Ainsi, pour $i$ assez grand,
$X'_i \to \mathbf{P}^n_Y$ est une immersion fermée d'après le
lemme \ref{spaces-limits-lemma-eventually-closed-immersion}.
Par conséquent, $X'_i$ est propre sur $Y$.
Pour un tel $i$, le morphisme $X_i \to Y$ est propre d'après
Morphismes des espaces, lemme \ref{spaces-morphisms-lemma-image-proper-is-proper}.
\end{proof}

\begin{lemma}
\label{spaces-limits-lemma-proper-limit-of-proper-finite-presentation-noetherian}
Soit $f : X \to Y$ un morphisme propre d'espaces algébriques sur $\mathbf{Z}$,
avec $Y$ quasi-compact et quasi-séparé. Il existe alors un ensemble préordonné
filtrant $I$, un système projectif $(f_i : X_i \to Y_i)$ de morphismes d'espaces
algébriques indexé par $I$, tel que les morphismes de transition $X_i \to X_{i'}$
et $Y_i \to Y_{i'}$ soient affines, que $f_i$ soit propre et de
présentation finie, que $Y_i$ soit de présentation finie sur
$\mathbf{Z}$ et que $(X \to Y) = \lim (X_i \to Y_i)$.
\end{lemma}

\begin{proof}
D'après le lemme \ref{spaces-limits-lemma-proper-limit-of-proper-finite-presentation},
on peut écrire $X = \lim_{k \in K} X_k$, avec $X_k \to Y$ propre et
de présentation finie. Ensuite, par approximation noethérienne absolue
(proposition \ref{spaces-limits-proposition-approximate}), on peut
écrire $Y = \lim_{j \in J} Y_j$, avec $Y_j$ de présentation finie
sur $\mathbf{Z}$.
Pour chaque $k$, il existe un $j$ et un morphisme $X_{k, j} \to Y_j$
de présentation finie tels que $X_k \cong Y \times_{Y_j} X_{k, j}$
comme espaces algébriques sur $Y$, voir le
lemme \ref{spaces-limits-lemma-descend-finite-presentation}.
Quitte à augmenter $j$, on peut supposer $X_{k, j} \to Y_j$ propre, voir le
lemme \ref{spaces-limits-lemma-eventually-proper}. L'ensemble $I$ sera constitué
de ces couples $(k, j)$ et le morphisme correspondant sera $X_{k, j} \to Y_j$.
Pour tout $k' \geq k$, on peut trouver un $j' \geq j$ et un morphisme
$X_{j', k'} \to X_{j, k}$ sur $Y_{j'} \to Y_j$ dont le changement de base à $Y$
donne le morphisme $X_{k'} \to X_k$ (cela résulte encore du
lemme \ref{spaces-limits-lemma-descend-finite-presentation}).
Ces morphismes constituent les morphismes de transition du système. Certains détails
sont omis.
\end{proof}

\noindent
Rappelons le support schématique d'un
module quasi-cohérent de type fini, voir
Morphismes des espaces, définition
\ref{spaces-morphisms-definition-scheme-theoretic-support}.

\begin{lemma}
\label{spaces-limits-lemma-eventually-proper-support}
Hypothèses et notation comme dans la situation \ref{spaces-limits-situation-descent-property}.
Soit $\mathcal{F}_0$ un $\mathcal{O}_{X_0}$-module quasi-cohérent.
Notons $\mathcal{F}$ et $\mathcal{F}_i$ les images réciproques de
$\mathcal{F}_0$ sur $X$ et $X_i$. Supposons que
\begin{enumerate}
\item $f_0$ soit localement de type fini,
\item $\mathcal{F}_0$ soit de type fini,
\item le support schématique de $\mathcal{F}$ soit propre sur $Y$.
\end{enumerate}
Alors le support schématique de $\mathcal{F}_i$ est propre sur $Y_i$
pour un certain $i$.
\end{lemma}

\begin{proof}
On peut remplacer $X_0$ par le support schématique de $\mathcal{F}_0$.
D'après Morphismes des espaces, lemme \ref{spaces-morphisms-lemma-support-finite-type},
cela garantit que $X_i$ est le support de $\mathcal{F}_i$ et que $X$ est le
support de $\mathcal{F}$. Si $Z \subset X$ désigne alors le support
schématique de $\mathcal{F}$, on voit que $Z \to X$ est un homéomorphisme
universel. On en conclut que $X \to Y$ est propre, puisque c'est le cas de
$Z \to Y$ par hypothèse, voir
Morphismes, lemme \ref{morphisms-lemma-image-proper-is-proper}.
D'après le lemme \ref{spaces-limits-lemma-eventually-proper}, le morphisme $X_i \to Y$ est propre
pour un certain $i$. Il s'ensuit que le support schématique $Z_i$ de
$\mathcal{F}_i$ est propre sur $Y$ d'après
Morphismes des espaces, lemmes
\ref{spaces-morphisms-lemma-closed-immersion-proper} et
\ref{spaces-morphisms-lemma-composition-proper}.
\end{proof}








\section{Immersion dans un espace affine}
\label{spaces-limits-section-embedding}

\noindent
Quelques lemmes techniques qui serviront plus loin à démontrer le lemme de Chow.

\begin{lemma}
\label{spaces-limits-lemma-embedding-into-affine-over-ls-qs}
Soient $S$ un schéma et $f : U \to X$ un morphisme d'espaces
algébriques sur $S$. Supposons que $U$ soit un schéma affine, que $f$ soit
localement de type fini, et que $X$ soit quasi-séparé et localement séparé.
Alors il existe une immersion $U \to \mathbf{A}^n_X$ sur $X$.
\end{lemma}

\begin{proof}
Écrivons $U = \Spec(A)$. Écrivons $A = \colim A_i$ comme limite inductive filtrante
de sous-algèbres de type fini sur $\mathbf{Z}$. Pour tout $i$, le morphisme
$U \to U_i = \Spec(A_i)$ induit un morphisme
$$
U \longrightarrow X \times U_i
$$
sur $X$. À la limite, le morphisme $U \to X \times U$ est une immersion
puisque $X$ est localement séparé, voir
Morphismes des espaces, lemme
\ref{spaces-morphisms-lemma-semi-diagonal}.
D'après le lemme \ref{spaces-limits-lemma-finite-type-eventually-closed},
le morphisme $U \to X \times U_i$ est une immersion pour un certain $i$.
Comme $U_i$ est isomorphe à un sous-schéma fermé de
$\mathbf{A}^n_{\mathbf{Z}}$, le lemme en résulte.
\end{proof}

\begin{remark}
\label{spaces-limits-remark-cannot-embed-in-general}
Nous avons vu dans Exemples, section \ref{examples-section-embedding-affines},
que le lemme \ref{spaces-limits-lemma-embedding-into-affine-over-ls-qs}
est faux si l'on omet l'hypothèse que $X$ soit localement séparé.
Cela soulève la question suivante : le
lemme \ref{spaces-limits-lemma-embedding-into-affine-over-ls-qs}
reste-t-il vrai si l'on omet l'hypothèse que $X$ soit quasi-séparé ?
Si vous connaissez la réponse, merci d'écrire à
\href{mailto:stacks.project@gmail.com}{stacks.project@gmail.com}.
\end{remark}

\begin{lemma}
\label{spaces-limits-lemma-embedding-into-affine-over-qs}
Soient $S$ un schéma et $f : Y \to X$ un morphisme d'espaces
algébriques sur $S$. Supposons que $X$ soit noethérien et que $f$ soit de présentation finie.
Alors il existe un ouvert dense $V \subset Y$ et une immersion
$V \to \mathbf{A}^n_X$.
\end{lemma}

\begin{proof}
Les hypothèses impliquent que $Y$ est noethérien
(Morphismes des espaces, lemme
\ref{spaces-morphisms-lemma-finite-presentation-noetherian}).
Alors $Y$ est quasi-séparé et possède donc un sous-schéma ouvert dense
(Propriétés des espaces, proposition
\ref{spaces-properties-proposition-locally-quasi-separated-open-dense-scheme}).
On peut donc supposer que $Y$ est un schéma noethérien.
En retirant les intersections des composantes irréductibles de $Y$
(utiliser Topologie, lemme \ref{topology-lemma-Noetherian} et
Propriétés, lemme \ref{properties-lemma-Noetherian-topology}),
on peut supposer que $Y$ est une réunion disjointe de schémas
noethériens irréductibles. Comme il existe une immersion
$$
\mathbf{A}^n_X \amalg \mathbf{A}^m_X
\longrightarrow
\mathbf{A}^{\max(n, m) + 1}_X
$$
(détails omis), il suffit de démontrer le résultat lorsque
$Y$ est irréductible.

\medskip\noindent
Supposons que $Y$ soit un schéma irréductible. Soit $T \subset |X|$ l'adhérence de
l'image de $f : Y \to X$. Comme $|Y|$ et $|X|$ sont des espaces topologiques sobres
(Propriétés des espaces, lemme
\ref{spaces-properties-lemma-quasi-separated-sober})
$T$ est irréductible et possède un unique point générique $\xi$, lequel est
l'image du point générique $\eta$ de $Y$.
Soit $\mathcal{I} \subset X$ un faisceau quasi-cohérent d'idéaux
définissant sur $T$ la structure d'espace réduit induite
(Propriétés des espaces, définition
\ref{spaces-properties-definition-reduced-induced-space}).
Comme $\mathcal{O}_{Y, \eta}$ est un anneau local artinien, il existe
$n > 0$ tel que $f^{-1}\mathcal{I}^n \mathcal{O}_{Y, \eta} = 0$.
Comme $f^{-1}\mathcal{I}\mathcal{O}_Y$ est un idéal quasi-cohérent de type fini,
on en conclut que $f^{-1}\mathcal{I}^n\mathcal{O}_V = 0$ pour
un ouvert non vide $V \subset Y$. Soit $Z \subset X$ le sous-espace fermé
défini par $\mathcal{I}^n$. Par construction, $V \to Y \to X$ se factorise par
$Z$. Comme $\mathbf{A}^n_Z \to \mathbf{A}^n_X$ est une immersion,
on peut remplacer $X$ par $Z$ et $Y$ par $V$.
Nous sommes ainsi ramenés au cas où $Y$ et $X$ sont irréductibles et où
$Y \to X$ envoie le point générique de $Y$ sur le point générique de $X$.

\medskip\noindent
Supposons que $Y$ et $X$ soient irréductibles, que $Y$ soit un schéma,
et que $Y \to X$ envoie le point générique de
$Y$ sur le point générique de $X$. D'après Propriétés des espaces, proposition
\ref{spaces-properties-proposition-locally-quasi-separated-open-dense-scheme}
$X$ possède un sous-schéma ouvert dense $U \subset X$. Choisissons un ouvert affine
non vide $V \subset Y$ dont l'image dans $X$ est contenue dans $U$. D'après
Morphismes, lemme \ref{morphisms-lemma-quasi-affine-finite-type-over-S},
le morphisme $V \to U$ se factorise en $V \to \mathbf{A}^n_U \to U$. En composant
avec $\mathbf{A}^n_U \to \mathbf{A}^n_X$, on obtient l'immersion voulue.
\end{proof}








\section{sections à support dans une partie fermée}
\label{spaces-limits-section-sections-with-support-in-closed}

\noindent
Cette section est l'analogue de
Propriétés, section \ref{properties-section-sections-with-support-in-closed}.

\begin{lemma}
\label{spaces-limits-lemma-quasi-coherent-finite-type-ideals}
Soit $S$ un schéma.
Soit $X$ un espace algébrique quasi-compact et quasi-séparé.
Soit $U \subset X$ un sous-espace ouvert. Les conditions suivantes sont équivalentes :
\begin{enumerate}
\item $U \to X$ est quasi-compact,
\item $U$ est quasi-compact, et
\item il existe un faisceau quasi-cohérent d'idéaux de type fini
$\mathcal{I} \subset \mathcal{O}_X$ tel que
$|X| \setminus |U| = |V(\mathcal{I})|$.
\end{enumerate}
\end{lemma}

\begin{proof}
Soient $W$ un schéma affine et $\varphi : W \to X$ un morphisme étale
surjectif, voir Propriétés des espaces, lemme
\ref{spaces-properties-lemma-quasi-compact-affine-cover}.
Si (1) est vérifiée, alors $\varphi^{-1}(U) \to W$ est quasi-compact, donc
$\varphi^{-1}(U)$ est quasi-compact, et par conséquent $U$ est quasi-compact
(car $|\varphi^{-1}(U)| \to |U|$ est surjectif). Si (2) est vérifiée, alors
$\varphi^{-1}(U)$ est quasi-compact, car $\varphi$ est quasi-compact
puisque $X$ est quasi-séparé (Morphismes des espaces,
lemme \ref{spaces-morphisms-lemma-quasi-compact-quasi-separated-permanence}).
Ainsi, $\varphi^{-1}(U) \to W$ est un morphisme quasi-compact de schémas d'après
Propriétés, lemme \ref{properties-lemma-quasi-coherent-finite-type-ideals}.
Il en résulte que $U \to X$ est quasi-compact d'après
Morphismes des espaces, lemme \ref{spaces-morphisms-lemma-quasi-compact-local}.
Les conditions (1) et (2) sont donc équivalentes.

\medskip\noindent
Supposons (1) et (2). D'après
Propriétés des espaces, lemme
\ref{spaces-properties-lemma-reduced-closed-subspace},
il existe un unique faisceau quasi-cohérent d'idéaux $\mathcal{J}$ définissant
la structure de sous-espace fermé réduit induite sur $|X| \setminus |U|$.
Remarquons que $\mathcal{J}|_U = \mathcal{O}_U$, qui est un
$\mathcal{O}_U$-module de type fini.
Comme $U$ est quasi-compact, il résulte du
lemme \ref{spaces-limits-lemma-directed-colimit-finite-type}
qu'il existe un sous-faisceau quasi-cohérent
$\mathcal{I} \subset \mathcal{J}$ qui est de type fini
et tel que $\mathcal{I}|_U = \mathcal{J}|_U$.
Alors $|X| \setminus |U| = |V(\mathcal{I})|$, ce qui donne (3). Inversement,
si $\mathcal{I}$ est comme en (3), alors $\varphi^{-1}(U) \subset W$
est un ouvert quasi-compact d'après le lemme pour les schémas
(Propriétés, lemme \ref{properties-lemma-quasi-coherent-finite-type-ideals})
appliqué à $\varphi^{-1}\mathcal{I}$ sur $W$.
La condition (2) est donc vérifiée.
\end{proof}

\begin{lemma}
\label{spaces-limits-lemma-sections-annihilated-by-ideal}
Soient $S$ un schéma et $X$ un espace algébrique sur $S$.
Soit $\mathcal{I} \subset \mathcal{O}_X$ un faisceau quasi-cohérent d'idéaux.
Soit $\mathcal{F}$ un $\mathcal{O}_X$-module quasi-cohérent.
Considérons le faisceau de $\mathcal{O}_X$-modules $\mathcal{F}'$
qui associe à tout objet $U$ de $X_\etale$ le module
$$
\mathcal{F}'(U)
=
\{s \in \mathcal{F}(U) \mid
\mathcal{I}s = 0\}
$$
Supposons que $\mathcal{I}$ soit de type fini. Alors
\begin{enumerate}
\item $\mathcal{F}'$ est un faisceau quasi-cohérent de $\mathcal{O}_X$-modules,
\item pour tout $U$ affine dans $X_\etale$, on a
$\mathcal{F}'(U) = \{s \in \mathcal{F}(U) \mid \mathcal{I}(U)s = 0\}$, et
\item $\mathcal{F}'_x = \{s \in \mathcal{F}_x \mid \mathcal{I}_x s = 0\}$.
\end{enumerate}
\end{lemma}

\begin{proof}
Il est clair que la règle définissant $\mathcal{F}'$ donne un sous-faisceau
de $\mathcal{F}$. On peut donc travailler localement pour la topologie étale sur $X$
afin de vérifier les autres assertions. Le lemme se ramène ainsi au cas des schémas,
qui est donné par
Propriétés, lemme \ref{properties-lemma-sections-annihilated-by-ideal}.
\end{proof}

\begin{definition}
\label{spaces-limits-definition-subsheaf-sections-annihilated-by-ideal}
Soient $S$ un schéma et $X$ un espace algébrique sur $S$.
Soit $\mathcal{I} \subset \mathcal{O}_X$ un faisceau quasi-cohérent
d'idéaux de type fini.
Soit $\mathcal{F}$ un $\mathcal{O}_X$-module quasi-cohérent.
Le sous-faisceau $\mathcal{F}' \subset \mathcal{F}$ défini au
lemme \ref{spaces-limits-lemma-sections-annihilated-by-ideal} ci-dessus est appelé
le {\it sous-faisceau des sections annulées par $\mathcal{I}$}.
\end{definition}

\begin{lemma}
\label{spaces-limits-lemma-push-sections-annihilated-by-ideal}
Soit $S$ un schéma.
Soit $f : X \to Y$ un morphisme quasi-compact et quasi-séparé
d'espaces algébriques sur $S$.
Soit $\mathcal{I} \subset \mathcal{O}_Y$ un faisceau quasi-cohérent
d'idéaux de type fini. Soit $\mathcal{F}$ un
$\mathcal{O}_X$-module quasi-cohérent. Soit $\mathcal{F}' \subset \mathcal{F}$
le sous-faisceau des sections annulées par $f^{-1}\mathcal{I}\mathcal{O}_X$.
Alors $f_*\mathcal{F}' \subset f_*\mathcal{F}$ est le sous-faisceau
des sections annulées par $\mathcal{I}$.
\end{lemma}

\begin{proof}
Démonstration omise. Indication : l'hypothèse que $f$ est quasi-compact et
quasi-séparé implique que $f_*\mathcal{F}$ est quasi-cohérent
(Morphismes des espaces, lemme \ref{spaces-morphisms-lemma-pushforward}),
de sorte que le lemme \ref{spaces-limits-lemma-sections-annihilated-by-ideal} s'applique
à $\mathcal{I}$ et à $f_*\mathcal{F}$.
\end{proof}

\noindent
Passons maintenant au faisceau des sections à support dans une partie fermée.
Là encore, ce n'est pas toujours un faisceau quasi-cohérent ; mais si le complémentaire
de la partie fermée est « rétrocompact » dans l'espace algébrique donné, alors
il l'est.

\begin{lemma}
\label{spaces-limits-lemma-sections-supported-on-closed-subset}
Soient $S$ un schéma et $X$ un espace algébrique sur $S$.
Soient $T \subset |X|$ une partie fermée et $U \subset X$
le sous-espace ouvert tel que $T \amalg |U| = |X|$.
Soit $\mathcal{F}$ un $\mathcal{O}_X$-module quasi-cohérent.
Considérons le faisceau de $\mathcal{O}_X$-modules $\mathcal{F}'$
qui associe à tout objet $\varphi : W \to X$ de
$X_\etale$ le module
$$
\mathcal{F}'(W)
=
\{s \in \mathcal{F}(W) \mid
\text{le support de }s\text{ est contenu dans }|\varphi|^{-1}(T)\}
$$
Si $U \to X$ est quasi-compact, alors
\begin{enumerate}
\item pour $W$ affine, il existe un idéal de type fini
$I \subset \mathcal{O}_X(W)$ tel que $|\varphi|^{-1}(T) = V(I)$,
\item pour $W$ et $I$ comme en (1), on a
$\mathcal{F}'(W) = \{x \in \mathcal{F}(W) \mid
I^nx = 0 \text{ pour un certain } n\}$,
\item $\mathcal{F}'$ est un faisceau quasi-cohérent de $\mathcal{O}_X$-modules.
\end{enumerate}
\end{lemma}

\begin{proof}
Il est clair que la règle définissant $\mathcal{F}'$ donne un sous-faisceau
de $\mathcal{F}$. On peut donc travailler localement pour la topologie étale sur $X$
afin de vérifier les autres assertions. Le lemme se ramène ainsi au cas des schémas,
qui est donné par
Propriétés, lemme \ref{properties-lemma-sections-supported-on-closed-subset}.
\end{proof}

\begin{definition}
\label{spaces-limits-definition-subsheaf-sections-supported-on-closed}
Soient $S$ un schéma et $X$ un espace algébrique sur $S$.
Soit $T \subset |X|$ une partie fermée dont le complémentaire
correspond à un sous-espace ouvert $U \subset X$
dont le morphisme d'inclusion $U \to X$ est quasi-compact.
Soit $\mathcal{F}$ un $\mathcal{O}_X$-module quasi-cohérent.
Le sous-faisceau quasi-cohérent $\mathcal{F}' \subset \mathcal{F}$ défini au
lemme \ref{spaces-limits-lemma-sections-supported-on-closed-subset} ci-dessus est appelé
le {\it sous-faisceau des sections à support dans $T$}.
\end{definition}

\begin{lemma}
\label{spaces-limits-lemma-push-sections-supported-on-closed-subset}
Soit $S$ un schéma.
Soit $f : X \to Y$ un morphisme quasi-compact et quasi-séparé
d'espaces algébriques sur $S$. Soit $T \subset |Y|$ une partie fermée.
Supposons que $|Y| \setminus T$ corresponde à un sous-espace ouvert $V \subset Y$
tel que $V \to Y$ soit quasi-compact. Soit $\mathcal{F}$ un
$\mathcal{O}_X$-module quasi-cohérent. Soit $\mathcal{F}' \subset \mathcal{F}$
le sous-faisceau des sections à support dans $|f|^{-1}T$.
Alors $f_*\mathcal{F}' \subset f_*\mathcal{F}$ est le sous-faisceau
des sections à support dans $T$.
\end{lemma}

\begin{proof}
Démonstration omise. Indications : $|X| \setminus |f|^{-1}T$ est le support du sous-espace ouvert
$U = f^{-1}V \subset X$. Comme $V \to Y$ est quasi-compact, il en est de même de
$U \to X$ (par changement de base). L'hypothèse que $f$ est quasi-compact et
quasi-séparé implique que $f_*\mathcal{F}$ est quasi-cohérent.
Le lemme \ref{spaces-limits-lemma-sections-supported-on-closed-subset}
s'applique donc à $T$ et à $f_*\mathcal{F}$, ainsi qu'à
$|f|^{-1}T$ et à $\mathcal{F}$. L'égalité des modules quasi-cohérents considérés
résulte immédiatement des définitions.
\end{proof}















\section{Caractérisation des espaces affines}
\label{spaces-limits-section-affine}

\noindent
Cette section est l'analogue de Limites, section \ref{limits-section-affine}.

\begin{lemma}
\label{spaces-limits-lemma-affine}
Soient $S$ un schéma et $f : X \to Y$ un morphisme d'espaces algébriques
sur $S$. Supposons que $f$ soit surjectif et fini, et que $X$
soit affine. Alors $Y$ est affine.
\end{lemma}

\begin{proof}
Nous pouvons, et allons, considérer $f : X \to Y$ comme un morphisme d'espaces algébriques sur
$\Spec(\mathbf{Z})$ (voir
Espaces, définition \ref{spaces-definition-base-change}).
Remarquons qu'un morphisme fini est affine et universellement fermé, voir
Morphismes des espaces, lemme
\ref{spaces-morphisms-lemma-integral-universally-closed}.
D'après Morphismes des espaces, lemme
\ref{spaces-morphisms-lemma-image-universally-closed-separated}
on voit que $Y$ est un espace algébrique séparé.
Comme $f$ est surjectif et $X$ quasi-compact, on voit que $Y$ est
quasi-compact.

\medskip\noindent
D'après le lemme \ref{spaces-limits-lemma-finite-in-finite-and-finite-presentation},
on peut écrire $X = \lim X_a$, où chaque $X_a \to Y$ est fini et de
présentation finie. D'après le
lemme \ref{spaces-limits-lemma-limit-is-affine},
on voit que $X_a$ est affine pour $a$ assez grand.
On peut donc supposer, et l'on suppose, que $f : X \to Y$ est fini, surjectif et
de présentation finie.

\medskip\noindent
D'après la proposition \ref{spaces-limits-proposition-approximate}, on peut écrire
$Y = \lim Y_i$ comme limite projective filtrante d'espaces
algébriques de présentation finie sur $\mathbf{Z}$.
D'après le lemme \ref{spaces-limits-lemma-descend-finite-presentation}, on peut
trouver $0 \in I$ et un morphisme $X_0 \to Y_0$ de présentation finie
tels que $X_i = X_0 \times_{Y_0} Y_i$ pour $i \geq 0$
et que $X = \lim_i X_i$. D'après le
lemme \ref{spaces-limits-lemma-descend-finite},
on voit que $X_i \to Y_i$ est fini pour $i$ assez grand.
D'après le lemme \ref{spaces-limits-lemma-descend-surjective},
on voit que $X_i \to Y_i$ est surjectif pour $i$ assez grand.
D'après le lemme \ref{spaces-limits-lemma-limit-is-affine}, on voit que $X_i$ est
affine pour $i$ assez grand. Ainsi, pour $i$ assez grand, on peut appliquer
Cohomologie des espaces, lemme
\ref{spaces-cohomology-lemma-image-affine-finite-morphism-affine-Noetherian}
pour conclure que $Y_i$ est affine. Il s'ensuit que $Y$ est affine, ce qui
achève la démonstration.
\end{proof}

\begin{proposition}
\label{spaces-limits-proposition-affine}
Soient $S$ un schéma et $f : X \to Y$ un morphisme d'espaces algébriques
sur $S$. Supposons que $X$ soit affine et que $f$ soit surjectif et universellement
fermé\footnote{Un morphisme entier est universellement fermé, voir
Morphismes des espaces, lemme
\ref{spaces-morphisms-lemma-integral-universally-closed}.}. Alors $Y$ est affine.
\end{proposition}

\begin{proof}
Nous pouvons, et allons, considérer $f : X \to Y$ comme un morphisme d'espaces algébriques sur
$\Spec(\mathbf{Z})$ (voir
Espaces, définition \ref{spaces-definition-base-change}).
D'après Morphismes des espaces, lemme
\ref{spaces-morphisms-lemma-image-universally-closed-separated}
on voit que $Y$ est un espace algébrique séparé. Puis, d'après
Morphismes des espaces, lemme \ref{spaces-morphisms-lemma-affine-permanence},
on obtient que $f$ est affine. Dès lors, d'après
Morphismes des espaces, lemme
\ref{spaces-morphisms-lemma-integral-universally-closed}
on voit que $f$ est entier.

\medskip\noindent
D'après le paragraphe précédent, on peut supposer que $f : X \to Y$
est surjectif et entier, que $X$ est affine et que $Y$ est séparé.
Comme $f$ est surjectif et $X$ quasi-compact, on en déduit aussi que $Y$ est
quasi-compact.

\medskip\noindent
Considérons le faisceau $\mathcal{A} = f_*\mathcal{O}_X$.
C'est un faisceau quasi-cohérent de $\mathcal{O}_Y$-algèbres, voir
Morphismes des espaces, lemme \ref{spaces-morphisms-lemma-pushforward}.
D'après le lemme \ref{spaces-limits-lemma-colimit-finitely-presented},
on peut écrire $\mathcal{A} = \colim_i \mathcal{F}_i$ comme limite inductive
filtrante de $\mathcal{O}_Y$-modules de type fini. Soit
$\mathcal{A}_i \subset \mathcal{A}$ la sous-$\mathcal{O}_Y$-algèbre
engendrée par $\mathcal{F}_i$. Comme le morphisme d'algèbres
$\mathcal{O}_Y \to \mathcal{A}$ est entier, on voit que chaque $\mathcal{A}_i$
est une $\mathcal{O}_Y$-algèbre quasi-cohérente finie. Ainsi,
$$
X_i = \underline{\Spec}_Y(\mathcal{A}_i) \longrightarrow Y
$$
est un morphisme fini d'espaces algébriques. Ici,
$\underline{\Spec}$ est la construction de Morphismes des espaces, lemme
\ref{spaces-morphisms-lemma-affine-equivalence-algebras}.
Il est clair
que $X = \lim_i X_i$. Par conséquent, d'après le
lemme \ref{spaces-limits-lemma-limit-is-affine},
on voit que, pour $i$ assez grand, le schéma $X_i$ est affine.
De plus, puisque $X \to Y$ se factorise par chaque $X_i$, on voit que
$X_i \to Y$ est surjectif. On conclut donc que $Y$ est affine d'après le
lemme \ref{spaces-limits-lemma-affine}.
\end{proof}

\noindent
Le corollaire suivant du résultat précédent se trouve dans
\cite{CLO}.

\begin{lemma}
\label{spaces-limits-lemma-reduction-scheme}
\begin{reference}
\cite[3.1.12]{CLO}
\end{reference}
Soient $S$ un schéma et $X$ un espace algébrique sur $S$.
Si $X_{red}$ est un schéma, alors $X$ est un schéma.
\end{lemma}

\begin{proof}
Soit $U' \subset X_{red}$ un sous-schéma ouvert affine.
Soit $U \subset X$ le sous-espace ouvert correspondant à l'ouvert
$|U'| \subset |X_{red}| = |X|$. Alors $U' \to U$ est surjectif et
entier. Par conséquent, $U$ est affine d'après la
proposition \ref{spaces-limits-proposition-affine}.
Ainsi, tout point est contenu dans un sous-schéma ouvert de $X$, c'est-à-dire que
$X$ est un schéma.
\end{proof}

\begin{lemma}
\label{spaces-limits-lemma-integral-universally-bijective-scheme}
Soient $S$ un schéma et $f : X \to Y$ un morphisme d'espaces algébriques
sur $S$. Supposons que $f$ soit entier et induise une bijection $|X| \to |Y|$.
Alors $X$ est un schéma si et seulement si $Y$ est un schéma.
\end{lemma}

\begin{proof}
Un morphisme entier est représentable par définition ; si $Y$
est un schéma, il en est donc de même de $X$. Inversement, supposons que $X$ soit un schéma.
Soit $U \subset X$ un ouvert affine. Un morphisme entier est
fermé et $|f|$ est bijectif ; ainsi $|f|(|U|) \subset |Y|$
est ouvert comme complémentaire de $|f|(|X| \setminus |U|)$. Soit
$V \subset Y$ le sous-espace ouvert tel que $|V| = |f|(|U|)$, voir
Propriétés des espaces, lemme \ref{spaces-properties-lemma-open-subspaces}.
Alors $U \to V$ est entier et surjectif ; par conséquent,
$V$ est un schéma affine d'après la proposition \ref{spaces-limits-proposition-affine}.
Cela achève la démonstration.
\end{proof}

\begin{lemma}
\label{spaces-limits-lemma-check-closed-infinitesimally}
Soit $S$ un schéma.
Soient $f : X \to B$ et $B' \to B$ des morphismes d'espaces algébriques sur $S$.
Supposons que
\begin{enumerate}
\item $B' \to B$ soit une immersion fermée,
\item $|B'| \to |B|$ soit bijectif,
\item $X \times_B B' \to B'$ soit une immersion fermée, et
\item $X \to B$ soit de type fini ou $B' \to B$ de présentation finie.
\end{enumerate}
Alors $f : X \to B$ est une immersion fermée.
\end{lemma}

\begin{proof}
Les hypothèses (1) et (2) impliquent que $B_{red} = B'_{red}$.
Posons $X' = X \times_B B'$. Alors $X' \to X$ est une immersion fermée
et $X'_{red} = X_{red}$. Soit $U \to B$ un morphisme étale,
avec $U$ affine. Alors $X' \times_B U \to X \times_B U$ est une
immersion fermée d'espaces algébriques induisant un isomorphisme
sur les espaces réduits sous-jacents. Comme $X' \times_B U$ est un schéma
(puisque $B' \to B$ et $X' \to B'$ sont représentables), il en est de même de
$X \times_B U$ d'après le lemme \ref{spaces-limits-lemma-reduction-scheme}.
Ainsi, $X \to B$ est également représentable. On se ramène donc au
cas des schémas, voir
Morphismes, lemme \ref{morphisms-lemma-check-closed-infinitesimally}.
\end{proof}










\section{Revêtement fini par un schéma}
\label{spaces-limits-section-finite-cover}

\noindent
Comme application des résultats de ce chapitre sur les limites, nous démontrons que, pour
tout espace algébrique quasi-compact et quasi-séparé $X$, il existe un schéma
$Y$ et un morphisme fini surjectif $Y \to X$. Nous utiliserons le résultat déjà
démontré selon lequel on peut trouver un revêtement entier par un schéma,
établi dans
Espaces décents, section \ref{decent-spaces-section-integral-cover}.

\begin{proposition}
\label{spaces-limits-proposition-there-is-a-scheme-finite-over}
Soient $S$ un schéma et $X$ un espace algébrique quasi-compact et quasi-séparé
sur $S$.
\begin{enumerate}
\item Il existe un morphisme fini surjectif $Y \to X$
de présentation finie, où $Y$ est un schéma,
\item étant donné un morphisme étale surjectif $U \to X$, on peut choisir
$Y \to X$ de telle sorte que, pour tout $y \in Y$, il existe un voisinage ouvert
$V \subset Y$ tel que $V \to X$ se factorise par $U$.
\end{enumerate}
\end{proposition}

\begin{proof}
La partie (1) est le cas particulier de (2) où $U = X$.
Soit $Y \to X$ comme dans
Espaces décents, lemme \ref{decent-spaces-lemma-there-is-a-scheme-integral-over}.
Choisissons un recouvrement ouvert affine fini $Y = \bigcup V_j$ tel que $V_j \to X$
se factorise par $U$. On peut écrire $Y = \lim Y_i$, où
$Y_i \to X$ est fini et de présentation finie, voir le
lemme \ref{spaces-limits-lemma-integral-limit-finite-and-finite-presentation}.
Pour $i$ assez grand, l'espace algébrique $Y_i$ est un schéma, voir le
lemme \ref{spaces-limits-lemma-limit-is-scheme}.
Pour $i$ assez grand, on peut trouver des ouverts affines $V_{i, j} \subset Y_i$
dont l'image réciproque dans $Y$ redonne $V_j$, voir le
lemme \ref{spaces-limits-lemma-descend-opens}.
Pour $i$ encore plus grand, les morphismes $V_j \to U$ sur $X$ proviennent
de morphismes $V_{i, j} \to U$ sur $X$, voir la proposition
\ref{spaces-limits-proposition-characterize-locally-finite-presentation}.
Cela achève la démonstration.
\end{proof}

\begin{lemma}
\label{spaces-limits-lemma-integral-limit-finite-and-finite-presentation-refined}
Soient $S$ un schéma et $f : X \to Y$ un morphisme entier d'espaces
algébriques sur $S$. Supposons $Y$ quasi-compact et quasi-séparé.
Soit $V \subset Y$ un sous-espace ouvert quasi-compact tel que
$f^{-1}(V) \to V$ soit fini et de présentation finie.
Alors $X$ peut s'écrire comme limite projective filtrante $X = \lim X_i$,
où les $f_i : X_i \to Y$ sont finis et de présentation finie,
et $f^{-1}(V) \to f_i^{-1}(V)$ est un isomorphisme pour tout $i$.
\end{lemma}

\begin{proof}
Ce lemme est un léger raffinement de la proposition
\ref{spaces-limits-proposition-there-is-a-scheme-finite-over}.
Considérons la $\mathcal{O}_Y$-algèbre quasi-cohérente entière
$\mathcal{A} = f_*\mathcal{O}_X$. Au paragraphe suivant, nous allons
écrire $\mathcal{A} = \colim \mathcal{A}_i$ comme limite inductive
filtrante de $\mathcal{O}_Y$-algèbres finies de présentation finie
$\mathcal{A}_i$ telles que $\mathcal{A}_i|_V = \mathcal{A}|_V$.
Cela fait, posons $X_i = \underline{\Spec}_Y(\mathcal{A}_i)$, voir
Morphismes des espaces, définition
\ref{spaces-morphisms-definition-relative-spec}.
Par construction, $X_i \to Y$ est fini et de présentation finie,
$X = \lim X_i$ et $f_i^{-1}(V) = f^{-1}(V)$.

\medskip\noindent
La démonstration de l'assertion sur les algèbres est analogue à
celle de la partie (2) du
lemme \ref{spaces-limits-lemma-integral-algebra-directed-colimit-finite}.
Écrivons d'abord $\mathcal{A} = \colim \mathcal{F}_i$
comme limite inductive de $\mathcal{O}_Y$-modules de présentation finie, à l'aide du
lemme \ref{spaces-limits-lemma-colimit-finitely-presented}.
Comme $\mathcal{A}|_V$ est un $\mathcal{O}_V$-module de type fini,
on peut supposer, et l'on suppose, que $\mathcal{F}_i|_V \to \mathcal{A}|_V$
est surjectif pour tout $i$.
Pour tout $i$, soit $\mathcal{J}_i$ le noyau du morphisme
$$
\text{Sym}^*_{\mathcal{O}_X}(\mathcal{F}_i) \longrightarrow \mathcal{A}
$$
Pour $i' \geq i$, on obtient un morphisme $\mathcal{J}_i \to \mathcal{J}_{i'}$.
On a $\mathcal{A} =
\colim \text{Sym}^*_{\mathcal{O}_X}(\mathcal{F}_i)/\mathcal{J}_i$.
De plus, les $\mathcal{O}_X$-algèbres quasi-cohérentes
$\text{Sym}^*_{\mathcal{O}_X}(\mathcal{F}_i)/\mathcal{J}_i$
sont finies (comme sous-algèbres quasi-cohérentes de type fini de la
$\mathcal{O}_Y$-algèbre quasi-cohérente entière $\mathcal{A}$ sur $\mathcal{O}_X$).
La restriction de $\text{Sym}^*_{\mathcal{O}_X}(\mathcal{F}_i)/\mathcal{J}_i$
à $V$ est $\mathcal{A}|_V$ par la surjectivité ci-dessus.
Ainsi, $\mathcal{J}_i|_V$ est de type fini comme faisceau d'idéaux de
$\text{Sym}^*_{\mathcal{O}_X}(\mathcal{F}_i)|_V$, du fait que
$\mathcal{A}|_V$ est de présentation finie comme
$\mathcal{O}_Y$-algèbre.
Écrivons $\mathcal{J}_i = \colim \mathcal{E}_{ik}$
comme limite inductive de $\mathcal{O}_X$-modules de présentation finie.
On peut supposer, et l'on suppose, que $\mathcal{E}_{ik}|_V$
engendre $\mathcal{J}_i|_V$ comme faisceau d'idéaux de
$\text{Sym}^*_{\mathcal{O}_X}(\mathcal{F}_i)|_V$, d'après
l'assertion de type fini ci-dessus.
Étant donnés $i' \geq i$ et $k$, il existe $k'$ tel que l'on
dispose d'un morphisme $\mathcal{E}_{ik} \to \mathcal{E}_{i'k'}$
rendant le diagramme
$$
\xymatrix{
\mathcal{J}_i \ar[r] & \mathcal{J}_{i'} \\
\mathcal{E}_{ik} \ar[u] \ar[r] & \mathcal{E}_{i'k'} \ar[u]
}
$$
commutatif. Cela résulte de Cohomologie des espaces, lemme
\ref{spaces-cohomology-lemma-finite-presentation-quasi-compact-colimit}.
On en déduit un morphisme
$$
\mathcal{A}_{ik} =
\text{Sym}^*_{\mathcal{O}_X}(\mathcal{F}_i)/(\mathcal{E}_{ik})
\longrightarrow
\text{Sym}^*_{\mathcal{O}_X}(\mathcal{F}_{i'})/(\mathcal{E}_{i'k'}) =
\mathcal{A}_{i'k'}
$$
où $(\mathcal{E}_{ik})$ désigne l'idéal engendré par $\mathcal{E}_{ik}$.
Les $\mathcal{O}_X$-algèbres quasi-cohérentes $\mathcal{A}_{ki}$
sont de présentation finie et finies pour $k$ assez grand
(voir la démonstration du
lemme \ref{spaces-limits-lemma-finite-algebra-directed-colimit-finite-finitely-presented}).
De plus, on a $\mathcal{A}_{ik}|_V = \mathcal{A}|_V$ par construction.
Enfin, on a
$$
\colim \mathcal{A}_{ik} = \colim \mathcal{A}_i = \mathcal{A}
$$
En effet, la première égalité a été démontrée dans la démonstration du
lemme \ref{spaces-limits-lemma-finite-algebra-directed-colimit-finite-finitely-presented},
et la seconde vaut parce que $\mathcal{A}$ est la limite inductive des
modules $\mathcal{F}_i$.
\end{proof}

\begin{lemma}
\label{spaces-limits-lemma-there-is-a-scheme-finite-over-refined}
Soient $S$ un schéma et $X$ un espace algébrique quasi-compact et quasi-séparé
sur $S$, tel que $|X|$ ait un nombre fini de composantes
irréductibles.
\begin{enumerate}
\item Il existe un morphisme fini surjectif $f : Y \to X$
de présentation finie, où $Y$ est un schéma, tel que $f$
soit fini étale au-dessus d'un ouvert dense quasi-compact $U \subset X$,
\item étant donné un morphisme étale surjectif $V \to X$, on peut choisir
$Y \to X$ de telle sorte que, pour tout $y \in Y$, il existe un voisinage ouvert
$W \subset Y$ tel que $W \to X$ se factorise par $V$.
\end{enumerate}
\end{lemma}

\begin{proof}
La partie (1) est le cas particulier de (2) où $V = X$.

\medskip\noindent
Preuve de (2). Soit $\pi : Y \to X$ comme dans Espaces décents, lemme
\ref{decent-spaces-lemma-there-is-a-scheme-integral-over-refined}
et soit $U \subset X$ un ouvert dense quasi-compact tel que
$\pi^{-1}(U) \to U$ soit fini étale.
Choisissons un recouvrement ouvert affine fini $Y = \bigcup W_j$ tel que
$W_j \to X$ se factorise par $V$. On peut écrire $Y = \lim Y_i$, où
$\pi_i : Y_i \to X$ est fini et de présentation finie et où
$\pi^{-1}(U) \to \pi_i^{-1}(U)$ est un isomorphisme, voir le
lemme \ref{spaces-limits-lemma-integral-limit-finite-and-finite-presentation-refined}.
Pour $i$ assez grand, l'espace algébrique $Y_i$ est un schéma, voir le
lemme \ref{spaces-limits-lemma-limit-is-scheme}.
Pour $i$ assez grand, on peut trouver des ouverts affines $W_{i, j} \subset Y_i$
dont l'image réciproque dans $Y$ redonne $W_j$, voir le
lemme \ref{spaces-limits-lemma-descend-opens}.
Pour $i$ encore plus grand, les morphismes $W_j \to V$ sur $X$ proviennent
de morphismes $W_{i, j} \to U$ sur $X$, voir la proposition
\ref{spaces-limits-proposition-characterize-locally-finite-presentation}.
Cela achève la démonstration.
\end{proof}

\begin{lemma}
\label{spaces-limits-lemma-there-is-a-scheme-finite-over-filtered}
Soient $S$ un schéma et $X$ un espace algébrique quasi-compact et quasi-séparé
sur $S$. Il existe un entier $t \geq 0$ et des
sous-espaces fermés
$$
X \supset Z_0 \supset Z_1 \supset \ldots \supset Z_t = \emptyset
$$
tels que $Z_i \to X$ soit de présentation finie,
que $Z_0 \subset X$ soit un épaississement et que, pour tout $i = 0, \ldots t - 1$,
il existe un schéma $Y_i$ et un morphisme fini surjectif de
présentation finie $Y_i \to Z_i$, qui soit fini étale au-dessus de
$Z_i \setminus Z_{i + 1}$.
\end{lemma}

\begin{proof}
On peut considérer $X$ comme un espace algébrique sur $\Spec(\mathbf{Z})$, voir
Espaces, définition \ref{spaces-definition-base-change}, et
Propriétés des espaces, définition \ref{spaces-properties-definition-separated}.
On peut donc appliquer la proposition \ref{spaces-limits-proposition-approximate}.
Il en résulte que l'on peut trouver un morphisme affine $X \to X_0$,
où $X_0$ est de présentation finie sur $\mathbf{Z}$.
Si l'on peut démontrer le lemme pour $X_0$, on peut alors ramener par changement de base
la stratification et les morphismes à $X$ et obtenir le résultat pour $X$ ;
certains détails sont omis. On se ramène ainsi au cas examiné au
paragraphe suivant.

\medskip\noindent
Supposons $X$ de présentation finie sur $\mathbf{Z}$.
Alors $X$ est noethérien et $|X|$ est un espace topologique
noethérien (ayant un nombre fini de composantes irréductibles), de dimension finie.
On peut donc raisonner par récurrence sur $\dim(|X|)$.
Tout morphisme fini vers $X$ est de présentation finie ;
on peut donc négliger cette condition dans la suite de la démonstration.
D'après le lemme \ref{spaces-limits-lemma-there-is-a-scheme-finite-over-refined},
il existe un morphisme fini surjectif $Y \to X$ qui est
fini étale au-dessus d'un ouvert dense $U \subset X$.
Posons $Z_0 = X$ et soit $Z_1 \subset X$ le sous-espace fermé réduit
tel que $|Z_1| = |X| \setminus |U|$.
Par récurrence, on trouve un entier $t \geq 0$ et une filtration
$$
Z_1 \supset Z_{1, 0} \supset Z_{1, 1} \supset \ldots
\supset Z_{1, t} = \emptyset
$$
par sous-espaces fermés, où $Z_{1, 0} \to Z_1$ est un épaississement
et où il existe des morphismes finis surjectifs $Y_{1, i} \to Z_{1, i}$
qui sont finis étales au-dessus de $Z_{1, i} \setminus Z_{1, i + 1}$.
Comme $Z_1$ est réduit, on a $Z_1 = Z_{1, 0}$.
On peut donc poser $Z_i = Z_{1, i - 1}$ et $Y_i = Y_{1, i - 1}$
pour $i \geq 1$, ce qui démontre le lemme.
\end{proof}











\section{Obtention de schémas}
\label{spaces-limits-section-representable}

\noindent
Voici quelques techniques supplémentaires pour montrer qu'un espace algébrique est un schéma.
La première consiste à montrer qu'il existe un sous-espace fermé minimal
qui n'est pas un schéma.

\begin{lemma}
\label{spaces-limits-lemma-minimal-closed-subspace}
Soient $S$ un schéma et $X$ un espace algébrique quasi-compact et quasi-séparé
sur $S$. Si $X$ n'est pas un schéma, alors il existe
un sous-espace fermé $Z \subset X$ tel que $Z$ ne soit pas un schéma, mais que
tout sous-espace fermé propre $Z' \subset Z$ soit un schéma.
\end{lemma}

\begin{proof}
Nous démontrons cela par le lemme de Zorn. Soit $\mathcal{Z}$ l'ensemble
des sous-espaces fermés $Z$ qui ne sont pas des schémas, ordonné par inclusion.
Par hypothèse, $\mathcal{Z}$ contient $X$ et n'est donc pas vide.
Si les $Z_\alpha$ forment une partie totalement ordonnée de $\mathcal{Z}$, alors
$Z = \bigcap Z_\alpha$ appartient à $\mathcal{Z}$. En effet,
$$
Z = \lim Z_\alpha
$$
et les morphismes de transition sont affines.
On peut donc appliquer le lemme \ref{spaces-limits-lemma-limit-is-scheme} : si $Z$
était un schéma, l'un des $Z_\alpha$ le serait aussi.
(Cela fonctionne même si $Z = \emptyset$, mais remarquons que, d'après le
lemme \ref{spaces-limits-lemma-limit-nonempty}, ce cas ne peut pas se produire.)
Ainsi, $\mathcal{Z}$ possède des éléments minimaux d'après le lemme de Zorn.
\end{proof}

\noindent
Nous pouvons maintenant établir quelques propriétés de ces non-schémas minimaux.

\begin{lemma}
\label{spaces-limits-lemma-minimal-nonscheme}
Soient $S$ un schéma et $X$ un espace algébrique quasi-compact et quasi-séparé
sur $S$. Supposons que tout sous-espace fermé propre
$Z \subset X$ soit un schéma, mais que $X$ ne soit pas un schéma. Alors $X$ est réduit
et irréductible.
\end{lemma}

\begin{proof}
On voit que $X$ est réduit d'après le lemme \ref{spaces-limits-lemma-reduction-scheme}.
Choisissons des parties fermées $T_1 \subset |X|$ et $T_2 \subset |X|$ telles que
$|X| = T_1 \cup T_2$. Si $T_1$ et $T_2$ sont des parties fermées propres,
alors les sous-espaces fermés réduits induits correspondants $Z_1, Z_2 \subset X$
(Propriétés des espaces, définition
\ref{spaces-properties-definition-reduced-induced-space})
sont des schémas, de même que $Z = Z_1 \times_X Z_2 = Z_1 \cap Z_2$, comme sous-schéma
fermé de $Z_1$ ou de $Z_2$. Remarquons que la somme amalgamée
$Z_1 \amalg_Z Z_2$ existe dans la catégorie des schémas, voir
Compléments sur les morphismes, lemme
\ref{more-morphisms-lemma-pushout-along-closed-immersions}.
Une manière de procéder consiste à montrer que $Z_1 \amalg_Z Z_2$ est isomorphe à $X$,
mais nous ne pouvons pas l'utiliser ici, car les sommes amalgamées d'espaces
algébriques ne sont étudiées que plus loin. Nous utiliserons plutôt le
lemme \ref{spaces-limits-lemma-affine} pour trouver un voisinage affine de tout point.
Soit en effet $x \in |X|$. Si $x \not \in Z_1$, alors $x$ possède un voisinage
qui est un schéma, à savoir $X \setminus Z_1$. De même si $x \not \in Z_2$.
Si $x \in Z = Z_1 \cap Z_2$, choisissons un ouvert affine
$U \subset Z_1 \amalg_Z Z_2$ contenant $z$. Alors $U_1 = Z_1 \cap U$
et $U_2 = Z_2 \cap U$ sont des ouverts affines dont les intersections avec
$Z$ coïncident. Comme $|Z_1| = T_1$ et $|Z_2| = T_2$ sont des parties fermées de
$|X|$ dont l'intersection est $|Z|$, on trouve un ouvert $W \subset |X|$
tel que $W \cap T_1 = |U_1|$ et $W \cap T_2 = |U_2|$. Notons encore $W$ le
sous-espace ouvert correspondant de $X$. Alors $x \in |W|$ et le morphisme
$U_1 \amalg U_2 \to W$ est un morphisme fini surjectif dont la source
est un schéma affine. Ainsi, $W$ est un schéma affine d'après le
lemme \ref{spaces-limits-lemma-affine}.
\end{proof}

\noindent
Un point essentiel du lemme suivant est qu'il suffit de vérifier
la condition aux images des points de $X$.

\begin{lemma}
\label{spaces-limits-lemma-enough-local}
Soit $f: X \to S$ un morphisme quasi-compact et quasi-séparé d'un
espace algébrique vers un schéma $S$. Si, pour tout $x \in |X|$ d'image
$s = f(x) \in S$, l'espace algébrique $X \times_S \Spec(\mathcal{O}_{S,s})$
est un schéma, alors $X$ est un schéma.
\end{lemma}

\begin{proof}
Soit $x \in |X|$. Il suffit de trouver un voisinage ouvert $U$ de
$s = f(x)$ tel que $X \times_S U$ soit un schéma.
Comme $X \times_S \Spec(\mathcal{O}_{S, s})$ est un schéma et que
$\mathcal{O}_{S, s} = \colim \mathcal{O}_S(U)$, où la limite inductive porte
sur les voisinages ouverts affines de $s$ dans $S$, on voit que
$$
X \times_S \Spec(\mathcal{O}_{S, s}) = \lim X \times_S U
$$
D'après le lemme \ref{spaces-limits-lemma-limit-is-scheme}, on voit que $X \times_S U$
est un schéma pour un certain $U$.
\end{proof}

\noindent
Au lieu de se restreindre aux anneaux locaux comme dans le lemme \ref{spaces-limits-lemma-enough-local},
on peut se restreindre aux sous-schémas fermés de la base.

\begin{lemma}
\label{spaces-limits-lemma-maximal-ideal}
Soit $\varphi : X \to \Spec(A)$ un morphisme quasi-compact et quasi-séparé
d'un espace algébrique vers un schéma affine.
Si $X$ n'est pas un schéma, alors il existe un idéal $I \subset A$
tel que le changement de base $X_{A/I}$ ne soit pas un schéma, mais que,
pour tout $I \subset I'$ avec $I \not = I'$, le changement de base
$X_{A/I'}$ soit un schéma.
\end{lemma}

\begin{proof}
Nous démontrons cela par le lemme de Zorn. Soit $\mathcal{I}$ l'ensemble
des idéaux $I$ tels que $X_{A/I}$ ne soit pas un schéma. Par
hypothèse, $\mathcal{I}$ contient $(0)$. Si les $I_\alpha$ forment
une chaîne d'idéaux de $\mathcal{I}$, alors
$I = \bigcup I_\alpha$ appartient à $\mathcal{I}$. En effet,
$A/I = \colim A/I_\alpha$, d'où
$$
X_{A/I} = \lim X_{A/I_\alpha}
$$
On peut donc appliquer le lemme \ref{spaces-limits-lemma-limit-is-scheme} : si $X_{A/I}$
était un schéma, l'un des $X_{A/I_\alpha}$ le serait aussi.
Ainsi, $\mathcal{I}$ possède des éléments maximaux d'après le lemme de Zorn.
\end{proof}






\section{Recollement dans les fibres fermées}
\label{spaces-limits-section-change-over-closed-points}

\noindent
En appliquant la théorie précédente au spectre d'un anneau local, on obtient
quelques agréables résultats de recollement pour les espaces algébriques relatifs.
Nous démontrons d'abord un lemme auxiliaire (qui sera considérablement généralisé
dans Amorçage, section \ref{bootstrap-section-applications}).

\begin{lemma}
\label{spaces-limits-lemma-relative-glueing}
Soit $S = U \cup W$ un recouvrement ouvert d'un schéma. Alors le foncteur
$$
FP_S \longrightarrow FP_U \times_{FP_{U \cap W}} FP_W
$$
donné par changement de base est une équivalence, où $FP_T$
est la catégorie des espaces algébriques de présentation finie sur
le schéma $T$.
\end{lemma}

\begin{proof}
Tout d'abord, puisque $S = U \cup W$ est un recouvrement de Zariski, on voit que la
catégorie des faisceaux sur $(\Sch/S)_{fppf}$ est équivalente à la catégorie
des triplets $(\mathcal{F}_U, \mathcal{F}_W, \varphi)$, où
$\mathcal{F}_U$ est un faisceau sur $(\Sch/U)_{fppf}$,
$\mathcal{F}_W$ est un faisceau sur $(\Sch/W)_{fppf}$, et où
$$
\varphi :
\mathcal{F}_U|_{(\Sch/U \cap W)_{fppf}}
\longrightarrow
\mathcal{F}_W|_{(\Sch/U \cap W)_{fppf}}
$$
est un isomorphisme. Voir Sites, lemme \ref{sites-lemma-mapping-property-glue}
(remarquons qu'aucune autre donnée de recollement n'est nécessaire, car
$U \times_S U = U$, $W \times_S W = W$, et que la condition de cocycle
est automatique pour la même raison).
Maintenant, si le faisceau $\mathcal{F}$ sur $(\Sch/S)_{fppf}$
correspond à $(\mathcal{F}_U, \mathcal{F}_W, \varphi)$
par cette équivalence, alors $\mathcal{F}$ est un espace algébrique
si et seulement si $\mathcal{F}_U$ et $\mathcal{F}_W$ sont des espaces algébriques.
Cela résulte immédiatement de
Espaces algébriques, lemme \ref{spaces-lemma-glueing-algebraic-spaces},
car $\mathcal{F}_U \to \mathcal{F}$ et $\mathcal{F}_W \to \mathcal{F}$
sont représentables par des immersions ouvertes et recouvrent $\mathcal{F}$.
Enfin, dans ce cas, l'espace algébrique $\mathcal{F}$ est de
présentation finie sur $S$ si et seulement si $\mathcal{F}_U$ est de présentation finie
sur $U$ et $\mathcal{F}_W$ de présentation finie sur $W$,
d'après Morphismes des espaces, lemmes
\ref{spaces-morphisms-lemma-quasi-compact-local},
\ref{spaces-morphisms-lemma-separated-local}, et
\ref{spaces-morphisms-lemma-finite-presentation-local}.
\end{proof}

\begin{lemma}
\label{spaces-limits-lemma-glueing-near-closed-point}
Soient $S$ un schéma et $s \in S$ un point fermé tel que
$U = S \setminus \{s\} \to S$ soit quasi-compact. En posant
$V = \Spec(\mathcal{O}_{S, s}) \setminus \{s\}$, on a
une équivalence de catégories
$$
FP_S \longrightarrow FP_U \times_{FP_V} FP_{\Spec(\mathcal{O}_{S, s})}
$$
où $FP_T$ est la catégorie des espaces algébriques de présentation finie
sur $T$.
\end{lemma}

\begin{proof}
Soit $W \subset S$ un voisinage ouvert de $s$. Le foncteur
$$
FP_S \to FP_U \times_{FP_{W \setminus \{s\}}} FP_W
$$
est une équivalence de catégories d'après le lemme \ref{spaces-limits-lemma-relative-glueing}.
On a $\mathcal{O}_{S, s} = \colim \mathcal{O}_W(W)$, où
$W$ parcourt les voisinages ouverts affines de $s$.
Ainsi, $\Spec(\mathcal{O}_{S, s}) = \lim W$, où $W$
parcourt les voisinages ouverts affines de $s$.
La catégorie des espaces algébriques de présentation finie
sur $\Spec(\mathcal{O}_{S, s})$ est donc la limite des
catégories d'espaces algébriques de présentation finie sur
$W$, où $W$ parcourt les voisinages ouverts affines
de $s$, voir le
lemme \ref{spaces-limits-lemma-descend-finite-presentation}.
Pour tout ouvert affine $s \in W$, on voit que $U \cap W$
est quasi-compact puisque $U \to S$ est quasi-compact.
Ainsi, $V = \lim W \cap U = \lim W \setminus \{s\}$ est une limite de
schémas quasi-compacts et quasi-séparés (voir
Limites, lemme \ref{limits-lemma-directed-inverse-system-has-limit}).
La catégorie des espaces algébriques de présentation finie
sur $V$ est donc elle aussi la limite des
catégories d'espaces algébriques de présentation finie sur
$W \cap U$, où $W$ parcourt les voisinages ouverts affines
de $s$. Le lemme résulte formellement de la combinaison
de ces résultats.
\end{proof}

\begin{lemma}
\label{spaces-limits-lemma-glueing-near-point}
Soient $S$ un schéma et $U \subset S$ un ouvert rétrocompact.
Soit $s \in S$ un point du complémentaire de $U$. En posant
$V = \Spec(\mathcal{O}_{S, s}) \cap U$, on a
une équivalence de catégories
$$
\colim_{s \in U' \supset U\text{ ouvert}} FP_{U'}
\longrightarrow
FP_U \times_{FP_V} FP_{\Spec(\mathcal{O}_{S, s})}
$$
où $FP_T$ est la catégorie des espaces algébriques de présentation finie
sur $T$.
\end{lemma}

\begin{proof}
Soit $W \subset S$ un voisinage ouvert de $s$. D'après le
lemme \ref{spaces-limits-lemma-relative-glueing}, le foncteur
$$
FP_{U \cup W}
\longrightarrow
FP_U \times_{FP_{U \cap W}} FP_W
$$
est une équivalence de catégories. On a
$\mathcal{O}_{S, s} = \colim \mathcal{O}_W(W)$, où
$W$ parcourt les voisinages ouverts affines de $s$.
Ainsi, $\Spec(\mathcal{O}_{S, s}) = \lim W$, où $W$
parcourt les voisinages ouverts affines de $s$.
La catégorie des espaces algébriques de présentation finie
sur $\Spec(\mathcal{O}_{S, s})$ est donc la limite des
catégories d'espaces algébriques de présentation finie sur
$W$, où $W$ parcourt les voisinages ouverts affines
de $s$, voir le
lemme \ref{spaces-limits-lemma-descend-finite-presentation}.
Pour tout ouvert affine $s \in W$, on voit que $U \cap W$
est quasi-compact puisque $U \to S$ est quasi-compact.
Ainsi, $V = \lim W \cap U$ est une limite de
schémas quasi-compacts et quasi-séparés (voir
Limites, lemme \ref{limits-lemma-directed-inverse-system-has-limit}).
La catégorie des espaces algébriques de présentation finie
sur $V$ est donc elle aussi la limite des
catégories d'espaces algébriques de présentation finie sur
$W \cap U$, où $W$ parcourt les voisinages ouverts affines
de $s$. Le lemme résulte formellement de la combinaison
de ces résultats.
\end{proof}

\begin{lemma}
\label{spaces-limits-lemma-glueing-near-multiple-closed-points}
Soient $S$ un schéma et $s_1, \ldots, s_n \in S$ des
points fermés deux à deux distincts tels que
$U = S \setminus \{s_1, \ldots, s_n\} \to S$ soit quasi-compact. En posant
$S_i = \Spec(\mathcal{O}_{S, s_i})$ et $U_i = S_i \setminus \{s_i\}$,
on a une équivalence de catégories
$$
FP_S \longrightarrow
FP_U \times_{(FP_{U_1} \times \ldots \times FP_{U_n})}
(FP_{S_1} \times \ldots \times FP_{S_n})
$$
où $FP_T$ est la catégorie des espaces algébriques de présentation finie
sur $T$.
\end{lemma}

\begin{proof}
Pour $n = 1$, c'est le lemme \ref{spaces-limits-lemma-glueing-near-closed-point}.
Pour $n > 1$, le lemme se démontre exactement de la même façon ou
peut s'en déduire. Par exemple, supposons que $f_i : X_i \to S_i$
soient des objets de $FP_{S_i}$ et que $f : X \to U$ soit un objet
de $FP_U$, et qu'on se soit donné des isomorphismes $X_i \times_{S_i} U_i = X \times_U U_i$.
D'après le lemme \ref{spaces-limits-lemma-glueing-near-closed-point}, on peut trouver
un morphisme $f' : X' \to U' = S \setminus \{s_1, \ldots, s_{n - 1}\}$
qui est de présentation finie, qui est isomorphe à
$X_i$ au-dessus de $S_i$, qui est isomorphe à $X$ au-dessus de $U$, et
ces isomorphismes sont compatibles avec l'isomorphisme donné
$X_i \times_{S_n} U_n = X \times_U U_n$.
On peut alors appliquer l'hypothèse de récurrence à
$f_i : X_i \to S_i$, $i \leq n - 1$,
$f' : X' \to U'$, et aux
isomorphismes induits $X_i \times_{S_i} U_i = X' \times_{U'} U_i$, $i \leq n - 1$.
Cela démontre la surjectivité essentielle. Nous omettons la preuve de la
pleine fidélité.
\end{proof}







\section{Application aux modifications}
\label{spaces-limits-section-modifications-at-a-point}

\noindent
À l'aide des limites, on peut décrire la catégorie des modifications d'un
espace algébrique décent en un point fermé au moyen de
l'anneau local hensélien.

\begin{lemma}
\label{spaces-limits-lemma-excision-modifications}
Soit $S$ un schéma. Considérons un morphisme séparé et \'etale
$f : V \to W$ d'espaces algébriques sur $S$.
Supposons qu'il existe un
sous-espace fermé $T \subset W$ tel que $f^{-1}T \to T$ soit
un isomorphisme. Alors, en posant $W^0 = W \setminus T$ et
$V^0 = f^{-1}W^0$, le foncteur de changement de base
$$
\left\{
\begin{matrix}
g : X \to W\text{ morphisme d'espaces algébriques} \\
g^{-1}(W^0) \to W^0\text{ est un isomorphisme}
\end{matrix}
\right\}
\longrightarrow
\left\{
\begin{matrix}
h : Y \to V\text{ morphisme d'espaces algébriques} \\
h^{-1}(V^0) \to V^0\text{ est un isomorphisme}
\end{matrix}
\right\}
$$
est une équivalence de catégories.
\end{lemma}

\begin{proof}
Comme $V \to W$ est séparé, on a
$V \times_W V = \Delta(V) \amalg U$ pour un certain sous-espace ouvert et fermé
$U$ de $V \times_W V$. L'hypothèse que $f^{-1}T \to T$ est un
isomorphisme donne $U \times_W T = \emptyset$, c'est-à-dire que les deux
projections $U \to V$ prennent leurs valeurs dans $V^0$.

\medskip\noindent
Étant donné $h : Y \to V$ dans la catégorie de droite, considérons le
foncteur contravariant $X$ sur $(\Sch/S)_{fppf}$ défini par la règle
$$
X(T) = \{(w, y) \mid
w :  T \to W,\ y : T \times_{w, W} V \to Y\text{ morphisme sur }V\}
$$
Notons $g : X \to W$ l'application qui envoie $(w, y) \in X(T)$ sur $w \in W(T)$.
Comme $h^{-1}V^0 \to V^0$ est un isomorphisme, si
$w : T \to W$ prend ses valeurs dans $W^0$, il n'existe qu'un seul choix pour $h$.
Autrement dit, $X \times_{g, W} W^0 = W^0$. D'autre part, considérons
un point à valeurs dans $T$, $(w, y, v)$, de $X \times_{g, W, f} V$.
Alors $w = f \circ v$ et
$$
y : T \times_{f \circ v, W} V \longrightarrow V
$$
est un morphisme sur $V$. Considérons le morphisme
$$
T \times_{f \circ v, W} V \xrightarrow{(v, \text{id}_V)}
V \times_W V = V \amalg U
$$
L'image réciproque de $V$ est $T$, plongé par
$(\text{id}_T, v) : T \to T \times_{f \circ v, W} V$.
Le composé $y' = y \circ (\text{id}_T, v) : T \to Y$
est un morphisme tel que $v = h \circ y'$; il détermine $y$, car la
restriction de $y$ à l'autre composante est uniquement déterminée puisque
$U$ prend ses valeurs dans $V^0$ par la seconde projection. Il s'ensuit que
$X \times_{g, W, f} V \to Y$, $(w, y, v) \mapsto y'$ est un isomorphisme.

\medskip\noindent
Ainsi, il suffit de montrer que $X$ est un espace algébrique.
Puisque $V \to W$ est séparé et \'etale, il est représentable d'après
Morphismes d'espaces, lemme
\ref{spaces-morphisms-lemma-locally-quasi-finite-separated-representable}
(et Morphismes d'espaces, lemme
\ref{spaces-morphisms-lemma-etale-locally-quasi-finite}).
Il est clair que $W^0 \to W$ est représentable et \'etale, puisqu'il s'agit d'une
immersion ouverte. Ainsi
$$
W^0 \amalg Y = X \times_{g, W} W^0 \amalg X \times_{g, W, f} V
= X \times_{g, W} (W^0 \amalg V) \longrightarrow X
$$
est représentable, surjectif et \'etale d'après Espaces, lemmes
\ref{spaces-lemma-base-change-representable-transformations} et
\ref{spaces-lemma-base-change-representable-transformations-property}.
Ainsi, $X$ est un espace
algébrique d'après Espaces, lemme
\ref{spaces-lemma-etale-locally-representable-by-space-gives-space}.
\end{proof}

\begin{lemma}
\label{spaces-limits-lemma-excision-modifications-properties}
Notations et hypothèses comme dans le lemme \ref{spaces-limits-lemma-excision-modifications}.
Soient $g : X \to W$ et $h : Y \to V$ qui se correspondent par l'équivalence.
Alors $g$ est quasi-compact, quasi-séparé, séparé, localement de présentation
finie, de présentation finie, localement de type fini, de type fini,
propre, entier, fini, et ajouter d'autres propriétés ici si et seulement si
$h$ possède la même propriété.
\end{lemma}

\begin{proof}
Si $g$ est quasi-compact, quasi-séparé, séparé, localement de présentation
finie, de présentation finie, localement de type fini, de type fini,
propre, fini, il en est de même de $h$, comme changement de base de $g$, d'après
Morphismes d'espaces, lemmes
\ref{spaces-morphisms-lemma-base-change-quasi-compact},
\ref{spaces-morphisms-lemma-base-change-separated},
\ref{spaces-morphisms-lemma-base-change-finite-presentation},
\ref{spaces-morphisms-lemma-base-change-finite-type},
\ref{spaces-morphisms-lemma-base-change-proper},
\ref{spaces-morphisms-lemma-base-change-integral}.
Réciproquement, soit $P$ une propriété des morphismes d'espaces
algébriques qui est locale pour la topologie \'etale sur la base et qui est vérifiée par
le morphisme identité de tout espace algébrique.
Puisque $\{W^0 \to W, V \to W\}$ est un recouvrement \'etale,
pour prouver que $g$ possède $P$, il suffit de montrer
que $h$ possède $P$. On conclut donc à l'aide de
Morphismes d'espaces, lemmes
\ref{spaces-morphisms-lemma-quasi-compact-local},
\ref{spaces-morphisms-lemma-separated-local},
\ref{spaces-morphisms-lemma-finite-presentation-local},
\ref{spaces-morphisms-lemma-finite-type-local},
\ref{spaces-morphisms-lemma-proper-local},
\ref{spaces-morphisms-lemma-integral-local}.
\end{proof}

\begin{lemma}
\label{spaces-limits-lemma-modifications}
Soit $S$ un schéma. Soit $X$ un espace algébrique décent sur $S$.
Soit $x \in |X|$ un point fermé tel que $U = X \setminus \{x\} \to X$
soit quasi-compact. En posant
$V = \Spec(\mathcal{O}_{X, x}^h) \setminus \{\mathfrak m_x^h\}$,
le foncteur de changement de base
$$
\left\{
\begin{matrix}
f : Y \to X\text{ de présentation finie} \\
f^{-1}(U) \to U\text{ est un isomorphisme}
\end{matrix}
\right\}
\longrightarrow
\left\{
\begin{matrix}
g : Y \to \Spec(\mathcal{O}_{X, x}^h)\text{ de présentation finie} \\
g^{-1}(V) \to V\text{ est un isomorphisme}
\end{matrix}
\right\}
$$
est une équivalence de catégories.
\end{lemma}

\begin{proof}
Soit $a : (W, w) \to (X, x)$ un voisinage \'etale élémentaire de $x$
avec $W$ affine, comme dans
Espaces décents, lemme
\ref{decent-spaces-lemma-decent-space-elementary-etale-neighbourhood}.
Comme $x$ est un point fermé de $X$ et que $w$ est l'unique point de $W$
au-dessus de $x$, on voit que $w$ est un point fermé de $W$. Puisque $a$
est \'etale et identifie les corps résiduels en $x$ et $w$, il
s'ensuit que $a$ induit un isomorphisme $a^{-1}x \to x$ (comme sous-espaces
fermés de $X$ et $W$). On peut donc appliquer
les lemmes \ref{spaces-limits-lemma-excision-modifications} et
\ref{spaces-limits-lemma-excision-modifications-properties}
pour se ramener au cas où $X$ est un schéma affine.

\medskip\noindent
Supposons $X$ affine. Rappelons que $\mathcal{O}_{X, x}^h$
est la limite inductive des $\Gamma(U, \mathcal{O}_U)$ sur les
voisinages \'etales élémentaires affines $(U, u) \to (X, x)$.
Rappelons que la catégorie de ces voisinages est
cofiltrante; voir Espaces décents, lemme
\ref{decent-spaces-lemma-elementary-etale-neighbourhoods} ou
Compléments sur les morphismes, lemme
\ref{more-morphisms-lemma-elementary-etale-neighbourhoods}.
Alors $\Spec(\mathcal{O}_{X, x}^h) = \lim U$ et
$V = \lim U \setminus \{u\}$
(lemme \ref{spaces-limits-lemma-directed-inverse-system-has-limit}),
où les limites sont prises sur la même catégorie. Ainsi, d'après
le lemme \ref{spaces-limits-lemma-descend-finite-presentation},
la catégorie de droite est la limite inductive des catégories
associées aux couples $(U, u)$. D'après le premier
paragraphe, chacune de ces catégories est équivalente à la
catégorie associée au couple $(X, x)$. Cela achève la démonstration.
\end{proof}






\section{Morphismes universellement fermés}
\label{spaces-limits-section-universally-closed}

\noindent
Dans cette section, nous étudions à quelle condition un morphisme quasi-compact
(mais non nécessairement séparé) est universellement fermé. Nous démontrons d'abord un lemme qui
permettra de vérifier cette propriété après un changement de base
localement de présentation finie.

\begin{lemma}
\label{spaces-limits-lemma-separate}
Soit $S$ un schéma. Soient $f : X \to Y$ et $g : Z \to Y$ des
morphismes d'espaces algébriques sur $S$. Soit $z \in |Z|$ et soit
$T \subset |X \times_Y Z|$ une partie fermée
telle que $z \not \in \Im(T \to |Z|)$.
Si $f$ est quasi-compact, il existe
un voisinage \'etale $(V, v) \to (Z, z)$,
un diagramme commutatif
$$
\xymatrix{
V \ar[d] \ar[r]_a & Z' \ar[d]^b \\
Z \ar[r]^g & Y,
}
$$
et une partie fermée $T' \subset |X \times_Y Z'|$ tels que
\begin{enumerate}
\item le morphisme $b : Z' \to Y$ soit localement de présentation finie,
\item en posant $z' = a(v)$, on ait $z' \not \in \Im(T' \to |Z'|)$, et
\item l'image réciproque de $T$ dans $|X \times_Y V|$
ait son image dans $T'$ par $|X \times_Y V| \to |X \times_Y Z'|$.
\end{enumerate}
De plus, on peut supposer que $V$ et $Z'$ sont des schémas affines et, si $Z$
est un schéma, que $V$ est un voisinage ouvert affine de $z$.
\end{lemma}

\begin{proof}
Nous allons déduire ce résultat du résultat analogue pour les morphismes de schémas.
Soit $y \in |Y|$ l'image de $z$. Choisissons d'abord un voisinage \'etale affine
$(U, u) \to (Y, y)$, puis un voisinage \'etale affine
$(V, v) \to (Z, z)$ tel que le morphisme $V \to Y$
se factorise par $U$. On peut alors remplacer
\begin{enumerate}
\item $X \to Y$ par $X \times_Y U \to U$,
\item $Z \to Y$ par $V \to U$,
\item $z$ par $v$, et
\item $T$ par son image réciproque dans
$|(X \times_Y U) \times_U V| = |X \times_Y V|$.
\end{enumerate}
En effet, nous montrerons ci-dessous qu'après avoir remplacé $V$ par un voisinage
ouvert affine de $v$, il existe un morphisme $a : V \to Z'$ vers
un certain $Z' \to U$ de présentation finie et une partie fermée $T'$
de $|(X \times_Y U) \times_U Z'| = |X \times_Y Z'|$ tels que
$T$ ait son image dans $T'$ et que $a(v) \not \in \Im(T' \to |Z'|)$.
On peut donc supposer, et l'on suppose désormais, que $Z$ et $Y$ sont des schémas affines,
étant entendu qu'il faut trouver une solution où $V$
est un voisinage ouvert de $z$.

\medskip\noindent
Puisque $f$ est quasi-compact et $Y$ affine, l'espace algébrique
$X$ est quasi-compact. Choisissons un schéma affine $W$ et un morphisme
\'etale surjectif $W \to X$. Soit $T_W \subset |W \times_Y Z|$
l'image réciproque de $T$. Alors $z$ n'appartient pas à l'image de
$T_W$. D'après le cas des schémas (Limites, lemme \ref{limits-lemma-separate}),
on peut trouver un voisinage ouvert $V \subset Z$ de $z$,
un diagramme commutatif de schémas
$$
\xymatrix{
V \ar[d] \ar[r]_a & Z' \ar[d]^b \\
Z \ar[r]^g & Y,
}
$$
et une partie fermée $T' \subset |W \times_Y Z'|$ tels que
\begin{enumerate}
\item le morphisme $b : Z' \to Y$ soit localement de présentation finie,
\item en posant $z' = a(z)$, on ait $z' \not \in \Im(T' \to Z')$, et
\item $T_1 = T_W \cap |W \times_Y V|$ ait son image dans $T'$ par
$|W \times_Y V| \to |W \times_Y Z'|$.
\end{enumerate}
Le diagramme commutatif
$$
\xymatrix{
W \times_Y Z \ar[d] &
W \times_Y V \ar[l] \ar[rr]_{a_1} \ar[d]_c & &
W \times_Y Z' \ar[d]^q \\
X \times_Y Z &
X \times_Y V \ar[l] \ar[rr]^{a_2} & &
X \times_Y Z'
}
$$
a des carrés cartésiens, et les morphismes verticaux sont surjectifs, \'etales
et, a fortiori, ouverts. En considérant le carré de gauche, on
voit que $T_1 = T_W \cap |W \times_Y V|$ est l'image réciproque de
$T_2 = T \cap |X \times_Y V|$ par $c$. D'après Propriétés des espaces, lemme
\ref{spaces-properties-lemma-points-cartesian}, on a
$a_1(T_1) = q^{-1}(a_2(T_2))$.
D'après Topologie, lemme \ref{topology-lemma-open-morphism-quotient-topology},
on a
$$
q^{-1}\left(\overline{a_2(T_2)}\right) =
\overline{q^{-1}(a_2(T_2))} =
\overline{a_1(T_1)} \subset T'
$$
Comme $q$ est surjectif, l'image de $\overline{a_2(T_2)} \to |Z'|$
ne contient pas $z'$, puisque cela est vrai pour $T'$.
On peut donc prendre le diagramme ci-dessus avec $Z', V, a, b$ et la
partie fermée $\overline{a_2(T_2)} \subset |X \times_Y Z'|$ comme
solution du problème posé par le lemme.
\end{proof}

\begin{lemma}
\label{spaces-limits-lemma-test-universally-closed}
Soit $S$ un schéma.
Soit $f : X \to Y$ un morphisme quasi-compact d'espaces algébriques sur $S$.
Les conditions suivantes sont équivalentes :
\begin{enumerate}
\item $f$ est universellement fermé,
\item pour tout morphisme $Z \to Y$ localement de présentation finie,
l'application $|X \times_Y Z| \to |Z|$ est fermée, et
\item il existe un schéma $V$ et un morphisme \'etale surjectif $V \to Y$
de sorte que l'application $|\mathbf{A}^n \times (X \times_Y V)| \to |\mathbf{A}^n \times V|$
soit fermée pour tout $n \geq 0$.
\end{enumerate}
\end{lemma}

\begin{proof}
Il est clair que (1) implique (2).
Supposons que $|X \times_Y Z| \to |Z|$ ne soit pas fermée pour un certain
morphisme d'espaces algébriques $Z \to Y$ sur $S$. Cela signifie qu'il
existe une partie fermée $T \subset |X \times_Y Z|$
telle que $\Im(T \to |Z|)$ ne soit pas fermée. Choisissons $z \in |Z|$
dans l'adhérence de l'image de $T$, mais non dans cette image.
Appliquons le lemme \ref{spaces-limits-lemma-separate}.
On obtient un voisinage \'etale $(V, v) \to (Z, z)$, un diagramme commutatif
$$
\xymatrix{
V \ar[d] \ar[r]_a & Z' \ar[d]^b \\
Z \ar[r]^g & Y,
}
$$
et une partie fermée $T' \subset |X \times_Y Z'|$ tels que
\begin{enumerate}
\item le morphisme $b : Z' \to Y$ soit localement de présentation finie,
\item en posant $z' = a(v)$, on ait $z' \not \in \Im(T' \to |Z'|)$, et
\item l'image réciproque de $T$ dans $|X \times_Y V|$ ait son image dans $T'$ par
$|X \times_Y V| \to |X \times_Y Z'|$.
\end{enumerate}
Nous affirmons que $z'$ appartient à l'adhérence de $\Im(T' \to |Z'|)$,
ce qui implique que $|X \times_Y Z'| \to |Z'|$ n'est pas fermée.
Cette affirmation montre que (2) implique (1).
Pour la vérifier, considérons
le diagramme commutatif suivant :
$$
\xymatrix{
X \times_Y Z \ar[d] &
X \times_Y V \ar[l] \ar[d] \ar[r] &
X \times_Y Z' \ar[d] \\
Z & V \ar[l] \ar[r]^a & Z'
}
$$
Soit $T_V \subset |X \times_Y V|$ l'image réciproque de $T$.
D'après Propriétés des espaces, lemme \ref{spaces-properties-lemma-points-cartesian},
l'image de $T_V$ dans $|V|$ est l'image réciproque de l'image
de $T$ dans $|Z|$. Puisque $z$ appartient à l'adhérence de l'image de
$T \to |Z|$ et que $|V| \to |Z|$ est ouverte, on voit que $v$ appartient à
l'adhérence de l'image de $T_V \to |V|$. Comme l'image de
$T_V$ dans $|X \times_Y Z'|$ est contenue dans $|T'|$, il s'ensuit
immédiatement que $z' = a(v)$ appartient à l'adhérence de l'image de $T'$.

\medskip\noindent
Il est clair que (1) implique (3). Soit $V \to Y$ comme dans (3).
Si l'on montre que $X \times_Y V \to V$ est universellement fermé,
alors $f$ est universellement fermé d'après
Morphismes d'espaces, lemme
\ref{spaces-morphisms-lemma-universally-closed-local}.
Il suffit donc de montrer que $f : X \to Y$ vérifie (2)
si $f$ est un morphisme quasi-compact d'espaces algébriques,
si $Y$ est un schéma et si $|\mathbf{A}^n \times X| \to |\mathbf{A}^n \times Y|$
est fermée pour tout $n$. Soit $Z \to Y$ localement de présentation finie.
Il faut montrer que l'application $|X \times_Y Z| \to |Z|$ est fermée.
Cette question est locale pour la topologie \'etale sur $Z$; on peut donc supposer $Z$
affine (certains détails sont omis). Puisque $Y$ est un schéma, $Z$ est affine
et $Z \to Y$ est localement de présentation finie, on peut trouver
une immersion $Z \to \mathbf{A}^n \times Y$; voir
Morphismes, lemme \ref{morphisms-lemma-quasi-affine-finite-type-over-S}.
Considérons le diagramme cartésien
$$
\vcenter{
\xymatrix{
X \times_Y Z \ar[d] \ar[r] & \mathbf{A}^n \times X \ar[d] \\
Z \ar[r] & \mathbf{A}^n \times Y
}
}
\quad
\begin{matrix}
\text{qui induit le} \\
\text{carré cartésien}
\end{matrix}
\quad
\vcenter{
\xymatrix{
|X \times_Y Z| \ar[d] \ar[r] & |\mathbf{A}^n \times X| \ar[d] \\
|Z| \ar[r] & |\mathbf{A}^n \times Y|
}
}
$$
d'espaces topologiques, dont les flèches horizontales sont des homéomorphismes
sur des parties localement fermées (Propriétés des espaces, lemme
\ref{spaces-properties-lemma-subspace-induced-topology}).
Ainsi, toute partie fermée $T$
de $|X \times_Y Z|$ est l'image réciproque d'une partie fermée $T'$ de
$|\mathbf{A}^n \times Y|$. Comme l'hypothèse affirme que l'image
de $T'$ dans $|\mathbf{A}^n \times X|$ est fermée, on conclut que
l'image de $T$ dans $|Z|$ est fermée, comme voulu.
\end{proof}

\begin{lemma}
\label{spaces-limits-lemma-limited-base-change}
Soit $S$ un schéma. Soit $f : X \to Y$ un
morphisme d'espaces algébriques sur $S$.
Supposons $f$ séparé et de type fini.
Les conditions suivantes sont équivalentes :
\begin{enumerate}
\item Le morphisme $f$ est propre.
\item Pour tout morphisme $Y \to Z$ localement de présentation finie,
l'application $|X \times_Y Z| \to |Z|$ est fermée, et
\item il existe un schéma $V$ et un morphisme \'etale surjectif $V \to Y$
de sorte que l'application $|\mathbf{A}^n \times (X \times_Y V)| \to |\mathbf{A}^n \times V|$
soit fermée pour tout $n \geq 0$.
\end{enumerate}
\end{lemma}

\begin{proof}
Puisqu'un morphisme propre est exactement un morphisme
séparé, de type fini et universellement fermé, ce
lemme est un cas particulier du lemme \ref{spaces-limits-lemma-test-universally-closed}.
\end{proof}


\section{Critère valuatif noethérien}
\label{spaces-limits-section-Noetherian-valuative-criterion}

\noindent
Nous avons déjà démontré certains résultats dans Cohomologie des espaces, section
\ref{spaces-cohomology-section-Noetherian-valuative-criterion}.
La section correspondante pour les schémas est
Limites, section \ref{limits-section-Noetherian-valuative-criterion}.

\medskip\noindent
Bon nombre des résultats de cette section peuvent (et devraient peut-être)
se démontrer en invoquant le lemme suivant, bien que nous ne l'ayons pas
toujours fait.

\begin{lemma}
\label{spaces-limits-lemma-reach-point-closure-Noetherian}
Soit $S$ un schéma. Soit $f : X \to Y$ un morphisme d'espaces algébriques
sur $S$. Supposons $f$ de type fini et $Y$ localement noethérien.
Soit $y \in |Y|$ un point de l'adhérence de l'image de $|f|$.
Alors il existe un diagramme commutatif
$$
\xymatrix{
\Spec(K) \ar[r] \ar[d] & X \ar[d]^f \\
\Spec(A) \ar[r] & Y
}
$$
où $A$ est un anneau de valuation discrète et $K$ son corps des fractions,
et où le point fermé de $\Spec(A)$ a pour image $y$. En outre, nous pouvons supposer
que le point $x \in |X|$ correspondant à $\Spec(K) \to X$ est un point de
codimension $0$\footnote{Voir la discussion dans
Propriétés des espaces, section \ref{spaces-properties-section-generic-points}.}
et que $K$ est le corps résiduel d'un point
d'un schéma \'etale sur $X$.
\end{lemma}

\begin{proof}
Choisissons un schéma affine $V$, un point $v \in V$ et un morphisme \'etale
$V \to Y$ qui envoie $v$ sur $y$. L'application $|V| \to |Y|$ est ouverte et, d'après
Propriétés des espaces, lemme \ref{spaces-properties-lemma-points-cartesian},
l'image de $|X \times_Y V| \to |V|$ est l'image inverse de
celle de $|f|$. Nous en concluons que le point $v$ appartient à l'adhérence de
l'image de $|X \times_Y V| \to |V|$. Si nous démontrons le lemme pour
$X \times_Y V \to V$ et le point $v$, alors le lemme en résulte pour
$f$ et $y$. Nous nous ramenons ainsi à la situation décrite au
paragraphe suivant.

\medskip\noindent
Supposons donnés $f : X \to Y$ et $y \in |Y|$ comme dans le lemme, où
$Y$ est un schéma affine. Comme $f$ est quasi-compact, nous en déduisons que
$X$ est quasi-compact. Nous pouvons donc choisir un schéma affine $W$ et
un morphisme \'etale surjectif $W \to X$. Alors l'image de
$|f|$ est égale à celle de $W \to Y$. Nous nous ramenons ainsi
au cas des schémas, traité dans
Limites, lemme \ref{limits-lemma-reach-point-closure-Noetherian}.
\end{proof}

\noindent
Énonçons d'abord le résultat concernant
la séparation. Nous utiliserons souvent des diagrammes commutatifs à flèches pleines de
morphismes d'espaces algébriques sur un schéma de base $S$, de la forme suivante :
\begin{equation}
\label{spaces-limits-equation-valuative}
\vcenter{
\xymatrix{
\Spec(K) \ar[r] \ar[d] & X \ar[d] \\
\Spec(A) \ar[r] \ar@{-->}[ru] & Y
}
}
\end{equation}
où $A$ est un anneau de valuation et $K$ son corps des fractions.

\begin{lemma}
\label{spaces-limits-lemma-Noetherian-dvr-valuative-separation}
Soit $S$ un schéma. Soit $f : X \to Y$ un morphisme d'espaces algébriques
sur $S$. Supposons $f$ quasi-séparé et localement de type fini, et
$Y$ localement noethérien. Les conditions suivantes sont équivalentes :
\begin{enumerate}
\item Le morphisme $f$ est séparé.
\item Dans tout diagramme (\ref{spaces-limits-equation-valuative}), il existe au plus
une flèche en tirets.
\item Dans tout diagramme (\ref{spaces-limits-equation-valuative}) où $A$ est un anneau de
valuation discrète, il existe au plus une flèche en tirets.
\item Dans tout diagramme (\ref{spaces-limits-equation-valuative}) où $A$ est un anneau de
valuation discrète et où l'image de $\Spec(K) \to X$ est un point de
codimension $0$ de $X$, il existe au plus une flèche en tirets.
\end{enumerate}
\end{lemma}

\begin{proof}
Nous avons (1) $\Rightarrow$ (2) d'après
Morphismes d'espaces, lemme
\ref{spaces-morphisms-lemma-separated-implies-valuative}.
Les implications (2) $\Rightarrow$ (3) et (3) $\Rightarrow$ (4)
sont immédiates. Il reste à montrer que (4) implique (1).

\medskip\noindent
Supposons (4). Nous devons montrer que la diagonale
$\Delta : X \to X \times_Y X$ est une immersion
fermée. Nous savons déjà que $\Delta$ est représentable, séparé et
localement de type fini, et que c'est un monomorphisme ; voir Morphismes d'espaces, lemme
\ref{spaces-morphisms-lemma-properties-diagonal}.
Choisissons un schéma affine $U$ et un morphisme \'etale
$U \to X \times_Y X$. Posons $V = X \times_{\Delta, X \times_Y X} U$.
Il suffit de montrer que $V \to U$ est une immersion fermée
(Morphismes d'espaces, lemme
\ref{spaces-morphisms-lemma-closed-immersion-local}).
Comme $X \times_Y X$ est localement de type fini sur $Y$, nous voyons que
$U$ est noethérien (utiliser Morphismes d'espaces, lemmes
\ref{spaces-morphisms-lemma-composition-finite-type},
\ref{spaces-morphisms-lemma-base-change-finite-type} et
\ref{spaces-morphisms-lemma-locally-finite-type-locally-noetherian}).
Notons que $V$ est un schéma puisque $\Delta$ est représentable.
De plus, $V$ est quasi-compact parce que $f$ est quasi-séparé.
Ainsi, $V \to U$ est séparé et de type fini. Considérons un diagramme commutatif
$$
\xymatrix{
\Spec(K) \ar[r] \ar[d] & V \ar[d] \\
\Spec(A) \ar[r] \ar@{-->}[ru] & U
}
$$
de morphismes de schémas, où $A$ est un anneau de valuation discrète
de corps des fractions $K$, et où $K$ est le corps résiduel
d'un point générique du schéma noethérien $V$. Comme $V \to X$
est \'etale (en tant que changement de base du morphisme \'etale $U \to X \times_Y X$),
nous voyons que l'image de $\Spec(K) \to V \to X$ est un point de
codimension $0$ ; voir Propriétés des espaces, section
\ref{spaces-properties-section-dimension-local-ring}.
Nous pouvons interpréter le composé $\Spec(A) \to U \to X \times_Y X$
comme une paire de morphismes $a, b : \Spec(A) \to X$ qui coïncident comme morphismes
vers $Y$ et sont égaux après restriction à $\Spec(K)$, leur restriction commune
ayant pour image un point de codimension $0$. Notre hypothèse
(4) garantit donc $a = b$, et nous obtenons la flèche en tirets du diagramme.
D'après Limites, lemme \ref{limits-lemma-Noetherian-dvr-valuative-proper},
nous concluons que $V \to U$ est propre. Autrement dit, $\Delta$ est propre.
Comme $\Delta$ est un monomorphisme, nous en déduisons que $\Delta$ est une
immersion fermée (Morphismes \'etales, lemme
\ref{etale-lemma-characterize-closed-immersion}), comme voulu.
\end{proof}

\begin{lemma}
\label{spaces-limits-lemma-Noetherian-dvr-valuative-proper}
Soit $S$ un schéma. Soit $f : X \to Y$ un morphisme d'espaces algébriques
sur $S$. Supposons $f$ quasi-séparé et de type fini, et $Y$
localement noethérien. Les conditions suivantes sont équivalentes :
\begin{enumerate}
\item $f$ est propre,
\item $f$ satisfait au critère valuatif, voir
Morphismes d'espaces, définition
\ref{spaces-morphisms-definition-valuative-criterion},
\item dans tout diagramme (\ref{spaces-limits-equation-valuative}), il existe exactement
une flèche en tirets,
\item dans tout diagramme (\ref{spaces-limits-equation-valuative}) où $A$ est un anneau de
valuation discrète, il existe exactement une flèche en tirets, et
\item dans tout diagramme (\ref{spaces-limits-equation-valuative}) où $A$ est un anneau de
valuation discrète et où l'image de $\Spec(K) \to X$ est un point de
codimension $0$ de $X$, il existe exactement une flèche en tirets\footnote{Il existe
une formulation plus précise où, pour la partie d'existence, on demande seulement que
la flèche en tirets existe après une extension d'anneaux de valuation discrète.}.
\end{enumerate}
\end{lemma}

\begin{proof}
Nous avons (1) $\Leftrightarrow$ (2) $\Leftrightarrow$ (3) d'après 
Morphismes d'espaces, lemme \ref{spaces-morphisms-lemma-characterize-proper}.
Il est clair que (3) $\Rightarrow$ (4) $\Rightarrow$ (5).
Pour achever la démonstration, montrons que (5) implique (1).

\medskip\noindent
Supposons (5). D'après le lemme \ref{spaces-limits-lemma-Noetherian-dvr-valuative-separation},
nous voyons que $f$ est séparé. Pour achever la démonstration, il suffit de montrer
que $f$ est universellement fermé. Soit $V \to Y$ un morphisme \'etale,
où $V$ est un schéma affine. Il suffit de montrer que le changement de
base $V \times_Y X \to V$ est universellement fermé ; voir
Morphismes d'espaces, lemme
\ref{spaces-morphisms-lemma-universally-closed-local}.
Considérons le diagramme
$$
\xymatrix{
\Spec(K) \ar[r] \ar[d] & V \times_Y X \ar[d] \ar[r] & X \ar[d] \\
\Spec(A) \ar[r] \ar@{-->}[ru] \ar@{..>}[rru] & V \ar[r] & Y
}
$$
ci-dessus, formé de morphismes d'espaces algébriques sur $S$,
qui est commutatif, où $A$ est un anneau de valuation discrète
de corps des fractions $K$ et où $\Spec(K) \to V \times_Y X$
a pour image un point de codimension $0$ de l'espace algébrique
$V \times_Y X$. Puisque $V \times_Y X \to X$ est \'etale, il s'ensuit
que l'image de $\Spec(K) \to X$ est un point de codimension $0$
de $X$. Ainsi, (5) nous donne la plus longue des deux flèches non pleines,
qui est pointillée et s'insère dans le diagramme. Nous obtenons alors bien sûr aussi
la plus courte, qui est en tirets. Il s'ensuit que nos hypothèses valent pour
le morphisme $V \times_Y X \to V$, et nous sommes ramenés au cas traité
au paragraphe suivant.

\medskip\noindent
Supposons que $Y$ soit un schéma affine noethérien. Dans ce cas, $X$ est un
espace algébrique noethérien séparé (nous savons déjà que $f$ est séparé)
de type fini sur $Y$. (En particulier, l'espace algébrique $X$
possède un sous-espace ouvert dense qui est un schéma d'après Propriétés des espaces, proposition
\ref{spaces-properties-proposition-locally-quasi-separated-open-dense-scheme},
bien qu'à strictement parler nous n'en ayons pas besoin.)
Choisissons un schéma quasi-projectif $X'$ sur $Y$ et un morphisme propre surjectif
$X' \to X$ comme dans la forme faible du lemme de Chow
(Cohomologie des espaces, lemme \ref{spaces-cohomology-lemma-weak-chow}).
Nous pouvons remplacer $X'$ par la réunion disjointe des composantes irréductibles
qui dominent une composante irréductible de $X$ ; nous omettons les détails.
En particulier, nous pouvons supposer que les points génériques du schéma $X'$ ont pour image
des points de codimension $0$ de $X$ (dans ce cas, ce sont exactement
les points génériques de $X$). Nous affirmons que $X' \to Y$ est propre.
Cette affirmation implique que $X$ est propre sur $Y$ d'après
Morphismes d'espaces, lemme \ref{spaces-morphisms-lemma-image-proper-is-proper}.
Pour la démontrer, d'après
Limites, lemme \ref{limits-lemma-Noetherian-dvr-valuative-proper},
il suffit de prouver que, dans tout diagramme commutatif à flèches pleines
$$
\xymatrix{
\Spec(K) \ar[r] \ar[d] & X' \ar[r] & X \ar[d] \\
\Spec(A) \ar[rr] \ar@{-->}[ru]^a \ar@{-->}[rru]_b & & Y
}
$$
où $A$ est un anneau de valuation discrète de corps des fractions $K$, et où $K$ est le corps résiduel
d'un point générique de $X'$, nous pouvons trouver la
flèche en tirets $a$ (nous savons déjà qu'elle est unique puisque $X'$ est séparé).
L'hypothèse (5) nous fournit la flèche en tirets $b$.
Alors le morphisme $X' \times_{X, b} \Spec(A) \to \Spec(A)$
est un morphisme propre de schémas et, par le critère valuatif
pour les morphismes de schémas, le relèvement de $b$ fournit le morphisme $a$ recherché.
\end{proof}

\begin{lemma}
\label{spaces-limits-lemma-check-universally-closed-Noetherian}
Soit $S$ un schéma. Soit $f : X \to Y$ un morphisme d'espaces algébriques
sur $S$. Supposons $Y$ localement noethérien et $f$
de type fini. Les conditions suivantes sont équivalentes :
\begin{enumerate}
\item $f$ est universellement fermé,
\item $f$ satisfait à la partie d'existence du critère valuatif,
\item il existe un schéma $V$ et un morphisme \'etale surjectif
$V \to Y$ tels que
$|\mathbf{A}^n \times X \times_Y V| \to |\mathbf{A}^n \times V|$ soit fermée
pour tout $n \geq 0$,
\item dans tout diagramme (\ref{spaces-limits-equation-valuative}) où $A$ est un anneau de
valuation discrète, il existe une extension de corps finie séparable $K'/K$,
un anneau de valuation discrète $A' \subset K'$ qui domine $A$, et
un morphisme $\Spec(A') \to X$ tel que le diagramme suivant soit commutatif :
$$
\xymatrix{
\Spec(K') \ar[r] \ar[d] & \Spec(K) \ar[r] & X \ar[d] \\
\Spec(A') \ar[r] \ar[rru] & \Spec(A) \ar[r] & Y
}
$$
\item dans tout diagramme (\ref{spaces-limits-equation-valuative}) où $A$ est un anneau de
valuation discrète, il existe une extension de corps $K'/K$,
un anneau de valuation $A' \subset K'$ qui domine $A$, et
un morphisme $\Spec(A') \to X$ tel que le diagramme suivant soit commutatif :
$$
\xymatrix{
\Spec(K') \ar[r] \ar[d] & \Spec(K) \ar[r] & X \ar[d] \\
\Spec(A') \ar[r] \ar[rru] & \Spec(A) \ar[r] & Y
}
$$
\end{enumerate}
\end{lemma}

\begin{proof}
Les conditions (1), (2) et (3) sont équivalentes d'après
le lemme \ref{spaces-limits-lemma-test-universally-closed} et
Morphismes d'espaces, lemme
\ref{spaces-morphisms-lemma-quasi-compact-existence-universally-closed}.
Ces conditions équivalentes impliquent (4), car
Morphismes d'espaces, lemme \ref{spaces-morphisms-lemma-finite-separable-enough},
nous dit que nous pouvons toujours choisir l'extension $K'/K$ finie séparable dans
la partie d'existence du critère valuatif, ce qui force automatiquement
$A'$ à être un anneau de valuation discrète par le théorème de Krull-Akizuki
(Algèbre, lemme \ref{algebra-lemma-krull-akizuki}).
L'implication (4) $\Rightarrow$ (5) est immédiate.
Dans le reste de la démonstration, nous montrons que (5) implique (1).

\medskip\noindent
Supposons (5). Choisissons un schéma affine $V$ et un morphisme \'etale $V \to Y$.
Il suffit de montrer que le changement de base de $f$ à $V$ est universellement fermé ;
voir Morphismes d'espaces, lemme
\ref{spaces-morphisms-lemma-universally-closed-local}.
Exactement comme dans la démonstration du lemme \ref{spaces-limits-lemma-Noetherian-dvr-valuative-proper},
nous voyons que l'hypothèse (5) se transmet à ce changement de base ; nous omettons
les détails. Nous sommes ainsi ramenés au cas traité au paragraphe suivant.

\medskip\noindent
Supposons que $Y$ soit un schéma affine noethérien et que (5) soit satisfaite.
Pour prouver que $f$ est universellement fermé, il suffit de montrer que
$|X \times \mathbf{A}^n| \to |Y \times \mathbf{A}^n|$ est fermée
pour tout $n$ (d'après la discussion ci-dessus). Comme l'hypothèse (5)
se transmet au morphisme produit
$X \times \mathbf{A}^n \to Y \times \mathbf{A}^n$ (détails omis),
nous sommes ramenés à prouver que $|X| \to |Y|$ est fermée.

\medskip\noindent
Supposons que $Y$ soit un schéma affine noethérien et que (5) soit satisfaite.
Soit $T \subset |X|$ une partie fermée. Nous devons montrer que
l'image de $T$ dans $|Y|$ est fermée. Nous pouvons remplacer $X$
par la structure réduite de sous-espace fermé induite sur $T$ ; nous
omettons de vérifier que la propriété (5) est préservée par
ce remplacement. Nous sommes ainsi ramenés à prouver que l'image
de $|X| \to |Y|$ est fermée.

\medskip\noindent
Soit $y \in |Y|$ un point de l'adhérence de l'image de
$|X| \to |Y|$. D'après le lemme \ref{spaces-limits-lemma-reach-point-closure-Noetherian},
nous pouvons choisir un diagramme commutatif
$$
\xymatrix{
\Spec(K) \ar[r] \ar[d] & X \ar[d]^f \\
\Spec(A) \ar[r] & Y
}
$$
où $A$ est un anneau de valuation discrète et $K$ son corps des fractions,
et où le point fermé de $\Spec(A)$ a pour image $y$. Il résulte immédiatement
de la propriété (5) que $y$ appartient à l'image de
$|X| \to |Y|$, ce qui achève la démonstration.
\end{proof}


















\section{Critères valuatifs noethériens raffinés}
\label{spaces-limits-section-refined-valuative-criteria}

\noindent
Cette section est l'analogue de
Limites, section \ref{limits-section-refined-valuative-criteria}.
Pour v\'erifier les crit\`eres valuatifs, il n'est en g\'en\'eral pas n\'ecessaire
de consid\'erer tous les diagrammes possibles faisant intervenir des anneaux de valuation.

\begin{lemma}
\label{spaces-limits-lemma-refined-valuative-criterion-proper}
Soit $S$ un sch\'ema. Soient $f : X \to Y$ et $h : U \to X$ des
morphismes d'espaces alg\'ebriques sur $S$. Supposons que $Y$ soit
localement noeth\'erien, que $f$ et $h$ soient de type fini,
que $f$ soit s\'epar\'e et que l'image de $|h| : |U| \to |X|$
soit dense dans $|X|$. Si, pour tout diagramme commutatif de fl\`eches pleines
$$
\xymatrix{
\Spec(K) \ar[r] \ar[d] & U \ar[r]^h & X \ar[d]^f \\
\Spec(A) \ar[rr] \ar@{-->}[rru] & & Y
}
$$
o\`u $A$ est un anneau de valuation discr\`ete de corps des fractions $K$, il
existe une fl\`eche pointill\'ee rendant le diagramme commutatif, alors $f$ est propre.
\end{lemma}

\begin{proof}
Il suffit de montrer que $f$ est universellement ferm\'e.
Soit $V \to Y$ un morphisme \'etale, o\`u $V$ est un sch\'ema affine.
Par Morphismes d'espaces, lemme
\ref{spaces-morphisms-lemma-universally-closed-local}
il suffit de montrer que le changement de base $X \times_Y V \to V$ est
universellement ferm\'e. Par Propri\'et\'es des espaces, lemme
\ref{spaces-properties-lemma-points-cartesian}
l'image $I$ de $|U \times_Y V| \to |X \times_Y V|$
est l'image r\'eciproque de l'image de $|h|$. Puisque
$|X \times_Y V| \to |X|$ est ouverte
(Propri\'et\'es des espaces, lemme \ref{spaces-properties-lemma-etale-open}),
on en d\'eduit que $I$ est dense dans $|X \times_Y V|$.
Les hypoth\`eses du lemme sont donc
satisfaites pour les morphismes $U \times_Y V \to X \times_Y V \to V$.
On peut ainsi supposer que $Y$ est un sch\'ema affine.

\medskip\noindent
Supposons que $Y$ soit un sch\'ema affine. Alors $U$ est quasi-compact. Choisissons
un sch\'ema affine et un morphisme \'etale surjectif $W \to U$.
On remplace alors $U$ par $W$ et l'on suppose que $U$ est affine.
Par la version faible du lemme de Chow
(Cohomologie des espaces, lemme \ref{spaces-cohomology-lemma-weak-chow})
on peut choisir un morphisme propre surjectif $X' \to X$
o\`u $X'$ est un sch\'ema. Alors $U' = X' \times_X U$ est un sch\'ema
et $U' \to X'$ est de type fini. On peut remplacer $X'$ par
l'image sch\'ematique de $h' : U' \to X'$; ainsi $h'(U')$
est dense dans $X'$. Nous affirmons que, pour tout diagramme
$$
\xymatrix{
\Spec(K) \ar[r] \ar[d] & U' \ar[r]^h & X' \ar[d]^{f'} \\
\Spec(A) \ar[rr] \ar@{-->}[rru] & & Y
}
$$
o\`u $A$ est un anneau de valuation discr\`ete de corps des fractions $K$, il
existe une fl\`eche pointill\'ee rendant le diagramme commutatif. En effet, l'hypoth\`ese du
lemme fournit d'abord une fl\`eche $\Spec(A) \to X$, que l'on rel\`eve ensuite
en une fl\`eche $\Spec(A) \to X'$ au moyen du crit\`ere valuatif de propret\'e
(Morphismes d'espaces, lemme \ref{spaces-morphisms-lemma-characterize-proper}).
Le morphisme $X' \to Y$ est s\'epar\'e
comme compos\'e d'un morphisme propre et d'un morphisme s\'epar\'e.
Le cas des sch\'emas montre donc que le morphisme $X' \to Y$ est propre
(Limites, lemme
\ref{limits-lemma-refined-valuative-criterion-proper}).
Par Morphismes d'espaces, lemme
\ref{spaces-morphisms-lemma-image-proper-is-proper}
on en conclut que $X \to Y$ est propre.
\end{proof}

\begin{lemma}
\label{spaces-limits-lemma-refined-valuative-criterion-separated}
Soit $S$ un sch\'ema. Soient $f : X \to Y$ et $h : U \to X$ des
morphismes d'espaces alg\'ebriques sur $S$. Supposons que $Y$ soit
localement noeth\'erien, que $f$ soit localement de type fini et quasi-s\'epar\'e,
que $h$ soit de type fini et que l'image de $|h| : |U| \to |X|$
soit dense dans $|X|$.
Si, pour tout diagramme commutatif de fl\`eches pleines
$$
\xymatrix{
\Spec(K) \ar[r] \ar[d] & U \ar[r]^h & X \ar[d]^f \\
\Spec(A) \ar[rr] \ar@{-->}[rru] & & Y
}
$$
o\`u $A$ est un anneau de valuation discr\`ete de corps des fractions $K$, il
existe au plus une fl\`eche pointill\'ee rendant le diagramme commutatif, alors $f$ est
s\'epar\'e.
\end{lemma}

\begin{proof}
Nous appliquerons le Lemme \ref{spaces-limits-lemma-refined-valuative-criterion-proper}
aux morphismes $U \to X$ et $\Delta : X \to X \times_Y X$.
V\'erifions les conditions. Remarquons que $\Delta$ est quasi-compact puisque
$f$ est quasi-s\'epar\'e. Bien entendu, $\Delta$ est localement de type fini et
s\'epar\'e (comme tout morphisme diagonal).
Enfin, supposons donn\'e un diagramme commutatif de fl\`eches pleines
$$
\xymatrix{
\Spec(K) \ar[r] \ar[d] & U \ar[r]^h & X \ar[d]^\Delta \\
\Spec(A) \ar[rr]^{(a, b)} \ar@{-->}[rru] & & X \times_Y X
}
$$
o\`u $A$ est un anneau de valuation discr\`ete de corps des fractions $K$.
Alors $a$ et $b$ fournissent deux fl\`eches pointill\'ees dans le diagramme du lemme
et doivent \^etre \`egales. On peut donc prendre $a = b$ comme fl\`eche pointill\'ee,
ce qui donne l'existence et ach\`eve la d\'emonstration.
\end{proof}

\begin{lemma}
\label{spaces-limits-lemma-refined-valuative-criterion-universally-closed}
Soit $S$ un sch\'ema.
Soient $f : X \to Y$ et $h : U \to X$ des morphismes d'espaces alg\'ebriques sur $S$.
Supposons que $Y$ soit localement noeth\'erien, que $f$ et $h$ soient de type fini,
que $f$ soit quasi-s\'epar\'e et
que $h(U)$ soit dense dans $X$. Si, pour tout diagramme commutatif de fl\`eches pleines
$$
\xymatrix{
\Spec(K) \ar[r] \ar[d] & U \ar[r]^h & X \ar[d]^f \\
\Spec(A) \ar[rr] \ar@{-->}[rru] & & Y
}
$$
o\`u $A$ est un anneau de valuation discr\`ete de corps des fractions $K$, il
existe une unique fl\`eche pointill\'ee rendant le diagramme commutatif, alors $f$ est propre.
\end{lemma}

\begin{proof}
Combiner les Lemmes \ref{spaces-limits-lemma-refined-valuative-criterion-separated} et
\ref{spaces-limits-lemma-refined-valuative-criterion-proper}.
\end{proof}







\section{Descente des espaces de type fini}
\label{spaces-limits-section-finite-type-quasi-separated}

\noindent
Cette section poursuit le th\`eme de la
section \ref{spaces-limits-section-finite-type-closed-in-finite-presentation}
dans l'esprit des r\'esultats examin\'es dans la
section \ref{spaces-limits-section-descending-relative}.
Elle est aussi l'analogue de
Limites, section \ref{limits-section-finite-type-quasi-separated},
pour les espaces alg\'ebriques.

\begin{situation}
\label{spaces-limits-situation-limit-noetherian}
Soit $S$ un sch\'ema, par exemple $\Spec(\mathbf{Z})$.
Soit $B = \lim_{i \in I} B_i$ la limite
d'un syst\`eme projectif filtrant d'espaces noeth\'eriens sur $S$
\`a morphismes de transition affines
$B_{i'} \to B_i$ pour $i' \geq i$.
\end{situation}

\begin{lemma}
\label{spaces-limits-lemma-good-diagram}
Dans la Situation \ref{spaces-limits-situation-limit-noetherian}.
Soit $X \to B$ un morphisme d'espaces alg\'ebriques
quasi-s\'epar\'e et de type fini.
Alors il existe $i \in I$ et un diagramme
\begin{equation}
\label{spaces-limits-equation-good-diagram}
\vcenter{
\xymatrix{
X \ar[r] \ar[d] & W \ar[d] \\
B \ar[r] & B_i
}
}
\end{equation}
tels que $W \to B_i$ soit de type fini et que
le morphisme induit $X \to B \times_{B_i} W$ soit une immersion
ferm\'ee.
\end{lemma}

\begin{proof}
Par le Lemme \ref{spaces-limits-lemma-finite-type-closed-in-finite-presentation},
il existe une immersion ferm\'ee $X \to X'$
sur $B$, o\`u $X'$ est un espace alg\'ebrique de pr\'esentation finie sur $B$.
Par le Lemme \ref{spaces-limits-lemma-descend-finite-presentation},
il existe $i$ et un morphisme de pr\'esentation finie
$X'_i \to B_i$ dont le changement de base est $X'$. Posons $W = X'_i$.
\end{proof}

\begin{lemma}
\label{spaces-limits-lemma-limit-from-good-diagram}
Dans la Situation \ref{spaces-limits-situation-limit-noetherian}.
Soit $X \to B$ un morphisme d'espaces alg\'ebriques
quasi-s\'epar\'e et de type fini. \`Etant donn\'es $i \in I$ et un diagramme
$$
\vcenter{
\xymatrix{
X \ar[r] \ar[d] & W \ar[d] \\
B \ar[r] & B_i
}
}
$$
comme dans (\ref{spaces-limits-equation-good-diagram}), pour $i' \geq i$, soit
$X_{i'}$ l'image sch\'ematique de $X \to B_{i'} \times_{B_i} W$.
Alors $X = \lim_{i' \geq i} X_{i'}$.
\end{lemma}

\begin{proof}
Puisque $X$ est quasi-compact et quasi-s\'epar\'e, la formation de
l'image sch\'ematique de $X \to B_{i'} \times_{B_i} W$
commute \`a la localisation \'etale
(Morphismes d'espaces, lemme
\ref{spaces-morphisms-lemma-quasi-compact-scheme-theoretic-image}).
On peut donc supposer que $W$ est affine et s'envoie dans un sch\'ema affine
$U_i$ \'etale sur $B_i$. Alors
$$
B_{i'} \times_{B_i} W =
B_{i'} \times_{B_i} U_i \times_{U_i} W =
U_{i'} \times_{U_i} W
$$
o\`u $U_{i'} = B_{i'} \times_{B_i} U_i$ est affine puisque les morphismes de
transition sont affines. Le lemme r\'esulte donc du cas des sch\'emas,
qui est
Limites, lemme \ref{limits-lemma-limit-from-good-diagram}.
\end{proof}

\begin{lemma}
\label{spaces-limits-lemma-morphism-good-diagram}
Dans la Situation \ref{spaces-limits-situation-limit-noetherian}.
Soit $f : X \to Y$ un morphisme d'espaces alg\'ebriques quasi-s\'epar\'es
et de type fini sur $B$. Soient
$$
\vcenter{
\xymatrix{
X \ar[r] \ar[d] & W \ar[d] \\
B \ar[r] & B_{i_1}
}
}
\quad\text{et}\quad
\vcenter{
\xymatrix{
Y \ar[r] \ar[d] & V \ar[d] \\
B \ar[r] & B_{i_2}
}
}
$$
des diagrammes comme dans (\ref{spaces-limits-equation-good-diagram}). Soient
$X = \lim_{i \geq i_1} X_i$ et
$Y = \lim_{i \geq i_2} Y_i$ les descriptions correspondantes
comme limites, donn\'ees par le Lemme \ref{spaces-limits-lemma-limit-from-good-diagram}.
Alors il existe $i_0 \geq \max(i_1, i_2)$ et un morphisme
$$
(f_i)_{i \geq i_0} : (X_i)_{i \geq i_0} \to (Y_i)_{i \geq i_0}
$$
de syst\`emes projectifs sur $(B_i)_{i \geq i_0}$ tel que
$f = \lim_{i \geq i_0} f_i$.
Si $(g_i)_{i \geq i_0} : (X_i)_{i \geq i_0} \to (Y_i)_{i \geq i_0}$
est un second morphisme de syst\`emes projectifs sur $(B_i)_{i \geq i_0}$ tel que
$f = \lim_{i \geq i_0} g_i$,
alors $f_i = g_i$ pour tout $i \gg i_0$.
\end{lemma}

\begin{proof}
Puisque $V \to B_{i_2}$ est de pr\'esentation finie et que
$X = \lim_{i \geq i_1} X_i$, on peut invoquer la Proposition
\ref{spaces-limits-proposition-characterize-locally-finite-presentation}
dans la forme pr\'ecis\'ee par le Lemme \ref{spaces-limits-lemma-better-characterize-relative-limit-preserving}
pour trouver $i_0 \geq \max(i_1, i_2)$ et un morphisme $h : X_{i_0} \to V$
sur $B_{i_2}$ tel que $X \to X_{i_0} \to V$ soit \`egal \`a $X \to Y \to V$.
Pour $i \geq i_0$, on obtient un diagramme commutatif de fl\`eches pleines
$$
\xymatrix{
X \ar[d] \ar[r] &
X_i \ar[r] \ar@{..>}[d] \ar@/_2pc/[dd] |!{[d];[ld]}\hole &
X_{i_0} \ar[d]^h \\
Y \ar[r] \ar[d] & Y_i \ar[r] \ar[d] & V \ar[d] \\
B \ar[r] & B_i \ar[r] & B_{i_0}
}
$$
Puisque l'image de $X \to X_i$ est sch\'ematiquement dense
et que $Y_i$ est l'image sch\'ematique de
$Y \to B_i \times_{B_{i_2}} V$,
le morphisme $X_i \to B_i \times_{B_{i_2}} V$
induit par le diagramme se factorise par $Y_i$
(Morphismes d'espaces, lemme \ref{spaces-morphisms-lemma-factor-factor}).
Ceci prouve l'existence.

\medskip\noindent
Unicit\'e. Soit $E_i \to X_i$ le noyau du couple $f_i$, $g_i$
pour $i \geq i_0$. On a
$E_i = Y_i \times_{\Delta, Y_i \times_{B_i} Y_i, (f_i, g_i)} X_i$.
Ainsi $E_i \to X_i$ est un monomorphisme de pr\'esentation finie, comme changement
de base de la diagonale de $Y_i$ sur $B_i$; voir
Morphismes d'espaces, lemmes \ref{spaces-morphisms-lemma-properties-diagonal} et
\ref{spaces-morphisms-lemma-diagonal-morphism-finite-type}.
Puisque $X_i$ est un sous-espace ferm\'e de $B_i \times_{B_{i_0}} X_{i_0}$
et qu'il en va de m\^eme pour $Y_i$, on a
$$
E_i =
X_i \times_{(B_i \times_{B_{i_0}} X_{i_0})} (B_i \times_{B_{i_0}} E_{i_0}) =
X_i \times_{X_{i_0}} E_{i_0}
$$
De m\^eme, on a $X = X \times_{X_{i_0}} E_{i_0}$. Le Lemme
\ref{spaces-limits-lemma-descend-isomorphism} donne donc $E_i = X_i$ pour $i$ assez grand.
\end{proof}

\begin{remark}
\label{spaces-limits-remark-finite-type-gives-well-defined-system}
Dans la Situation \ref{spaces-limits-situation-limit-noetherian},
les Lemmes \ref{spaces-limits-lemma-good-diagram}, \ref{spaces-limits-lemma-limit-from-good-diagram} et
\ref{spaces-limits-lemma-morphism-good-diagram}
montrent que la cat\'egorie des espaces alg\'ebriques quasi-s\'epar\'es et
de type fini sur $B$ est \`equivalente \`a certains types de
syst\`emes projectifs d'espaces alg\'ebriques sur $(B_i)_{i \in I}$, \`a savoir
ceux obtenus en appliquant le Lemme \ref{spaces-limits-lemma-limit-from-good-diagram}
\`a un diagramme de la forme (\ref{spaces-limits-equation-good-diagram}).
Par exemple, soient $X \to B$ de type fini et quasi-s\'epar\'e,
et deux diagrammes distincts $X \to V_1 \to B_{i_1}$
et $X \to V_2 \to B_{i_2}$ comme dans (\ref{spaces-limits-equation-good-diagram}). En
appliquant le Lemme \ref{spaces-limits-lemma-morphism-good-diagram} \`a $\text{id}_X$
(dans les deux sens),
on voit que les descriptions correspondantes de
$X$ comme limite sont canoniquement isomorphes (quitte \`a restreindre
l'ensemble filtrant $I$). Et ainsi de suite.
\end{remark}

\begin{lemma}
\label{spaces-limits-lemma-morphism-good-diagram-flat}
Notations et hypoth\`eses comme dans le Lemme \ref{spaces-limits-lemma-morphism-good-diagram}.
Si $f$ est plat et de pr\'esentation finie, alors
il existe $i_3 > i_0$ tel que, pour $i \geq i_3$, le morphisme
$f_i$ soit plat, $X_i = Y_i \times_{Y_{i_3}} X_{i_3}$ et
$X = Y \times_{Y_{i_3}} X_{i_3}$.
\end{lemma}

\begin{proof}
Par le Lemme \ref{spaces-limits-lemma-descend-finite-presentation},
on peut choisir $i \geq i_2$ et un morphisme
$U \to Y_i$ de pr\'esentation finie tel que $X = Y \times_{Y_i} U$
(c'est ici que l'on utilise que $f$ est de pr\'esentation finie).
Quitte \`a augmenter $i$, on peut supposer que $U \to Y_i$ est plat; voir le
Lemme \ref{spaces-limits-lemma-descend-flat}.
Comme expliqu\'e dans la Remarque \ref{spaces-limits-remark-finite-type-gives-well-defined-system},
on remplace le diagramme initial ayant servi \`a d\'efinir le syst\`eme
$(X_i)_{i \geq i_1}$ par le syst\`eme correspondant \`a
$X \to U \to B_i$. Ainsi, $X_{i'}$, pour $i' \geq i$, est d\'efini comme
l'image sch\'ematique de $X \to B_{i'} \times_{B_i} U$.

\medskip\noindent
Puisque $U \to Y_i$ est plat (c'est ici que l'on utilise que $f$ est plat),
que $X = Y \times_{Y_i} U$ et
que l'image sch\'ematique de $Y \to Y_i$ est $Y_i$,
on voit que l'image sch\'ematique de $X \to U$ est $U$
(Morphismes d'espaces, lemme
\ref{spaces-morphisms-lemma-flat-base-change-scheme-theoretic-image}).
Remarquons que $Y_{i'} \to B_{i'} \times_{B_i} Y_i$ est une
immersion ferm\'ee pour $i' \geq i$ par construction du syst\`eme des $Y_j$.
Le m\^eme argument que ci-dessus montre alors que l'image sch\'ematique
de $X \to B_{i'} \times_{B_i} U$
est le sous-espace ferm\'e $Y_{i'} \times_{Y_i} U$.
Ainsi $X_{i'} = Y_{i'} \times_{Y_i} U$ pour tout $i' \geq i$,
et le lemme vaut donc avec $i_3 = i$.
\end{proof}

\begin{lemma}
\label{spaces-limits-lemma-morphism-good-diagram-smooth}
Notations et hypoth\`eses comme dans le Lemme \ref{spaces-limits-lemma-morphism-good-diagram}.
Si $f$ est lisse, alors il existe $i_3 > i_0$ tel que, pour
$i \geq i_3$, le morphisme $f_i$ soit lisse.
\end{lemma}

\begin{proof}
Combiner les Lemmes \ref{spaces-limits-lemma-morphism-good-diagram-flat} et
\ref{spaces-limits-lemma-descend-smooth}.
\end{proof}

\begin{lemma}
\label{spaces-limits-lemma-morphism-good-diagram-proper}
Notations et hypoth\`eses comme dans le Lemme \ref{spaces-limits-lemma-morphism-good-diagram}.
Si $f$ est propre, alors il existe $i_3 \geq i_0$ tel que, pour
$i \geq i_3$, le morphisme $f_i$ soit propre.
\end{lemma}

\begin{proof}
Il r\'esulte de la discussion de la 
Remarque \ref{spaces-limits-remark-finite-type-gives-well-defined-system}
que le choix de $i_1$ et de $W$ s'inscrivant dans un diagramme comme
(\ref{spaces-limits-equation-good-diagram}) est sans incidence sur la validit\'e du
lemme. On choisit donc $W$ comme suit.
On choisit d'abord une immersion ferm\'ee $X \to X'$
avec $X' \to Y$ propre et de pr\'esentation finie; voir le
Lemme \ref{spaces-limits-lemma-proper-limit-of-proper-finite-presentation}.
On choisit ensuite $i_3 \geq i_2$ et un morphisme propre $W \to Y_{i_3}$
tel que $X' = Y \times_{Y_{i_3}} W$. Cela est possible puisque
$Y = \lim_{i \geq i_2} Y_i$ et en vertu des
Lemmes \ref{spaces-limits-lemma-relative-approximation} et \ref{spaces-limits-lemma-eventually-proper}.
Avec ce choix de $W$, il r\'esulte imm\'ediatement de la construction que,
pour $i \geq i_3$, l'espace alg\'ebrique $X_i$ est un sous-espace ferm\'e de
$Y_i \times_{Y_{i_3}} W \subset B_i \times_{B_{i_3}} W$
et est donc propre sur $Y_i$.
\end{proof}

\begin{lemma}
\label{spaces-limits-lemma-good-diagram-fibre-product}
Dans la Situation \ref{spaces-limits-situation-limit-noetherian}, supposons donn\'e un
diagramme cart\'esien
$$
\xymatrix{
X^1 \ar[r]_p \ar[d]_q & X^3 \ar[d]^a \\
X^2 \ar[r]^b & X^4
}
$$
d'espaces alg\'ebriques quasi-s\'epar\'es et de type fini sur $B$.
Pour chaque $j = 1, 2, 3, 4$, choisissons $i_j \in I$ et un diagramme
$$
\xymatrix{
X^j \ar[r] \ar[d] & W^j \ar[d] \\
B \ar[r] & B_{i_j}
}
$$
comme dans (\ref{spaces-limits-equation-good-diagram}). Soient
$X^j = \lim_{i \geq i_j} X^j_i$ les descriptions correspondantes comme limites,
comme dans le Lemme \ref{spaces-limits-lemma-morphism-good-diagram}.
Soient $(a_i)_{i \geq i_5}$, $(b_i)_{i \geq i_6}$, $(p_i)_{i \geq i_7}$ et
$(q_i)_{i \geq i_8}$ les morphismes correspondants de syst\`emes projectifs
construits dans le Lemme \ref{spaces-limits-lemma-morphism-good-diagram}. Alors il existe
$i_9 \geq \max(i_5, i_6, i_7, i_8)$ tel que, pour $i \geq i_9$, on ait
$a_i \circ p_i = b_i \circ q_i$ et que
$$
(q_i, p_i) : X^1_i \longrightarrow X^2_i \times_{b_i, X^4_i, a_i} X^3_i
$$
soit une immersion ferm\'ee.
Si $a$ et $b$ sont plats et de pr\'esentation finie, alors il existe
$i_{10} \geq \max(i_5, i_6, i_7, i_8, i_9)$ tel que, pour $i \geq i_{10}$,
le dernier morphisme affich\'e soit un isomorphisme.
\end{lemma}

\begin{proof}
D'apr\`es la discussion de la
Remarque \ref{spaces-limits-remark-finite-type-gives-well-defined-system},
le choix de $W^1$ s'inscrivant dans un diagramme comme
(\ref{spaces-limits-equation-good-diagram}) est sans incidence sur la validit\'e du
lemme. On peut donc choisir $W^1 = W^2 \times_{W^4} W^3$.
Il r\'esulte alors imm\'ediatement de la construction de $X^1_i$ que 
$a_i \circ p_i = b_i \circ q_i$ et que
$$
(q_i, p_i) : X^1_i \longrightarrow X^2_i \times_{b_i, X^4_i, a_i} X^3_i
$$
est une immersion ferm\'ee.

\medskip\noindent
Si $a$ et $b$ sont plats et de pr\'esentation finie, il en va de m\^eme de
$p$ et $q$, changements de base de $a$ et $b$. On peut donc appliquer le
Lemme \ref{spaces-limits-lemma-morphism-good-diagram-flat}
\`a chacun de $a$, $b$, $p$, $q$ et $a \circ p = b \circ q$.
Il en r\'esulte qu'il existe $i_9 \in I$ tel que
$$
(q_i, p_i) : X^1_i \to X^2_i \times_{X^4_i} X^3_i
$$
soit le changement de base de $(q_{i_9}, p_{i_9})$ par le morphisme
$X^4_i \to X^4_{i_9}$ pour tout $i \geq i_9$.
On conclut que $(q_i, p_i)$ est un isomorphisme pour tout $i$ suffisamment
grand par le Lemme \ref{spaces-limits-lemma-descend-isomorphism}.
\end{proof}














 


% OK, start here.
%

\chapter{Diviseurs sur les espaces algébriques}



\phantomsection
\label{spaces-divisors-section-phantom}


\section{Introduction}
\label{spaces-divisors-section-introduction}

\noindent
Dans ce chapitre, nous étudions les diviseurs sur les espaces algébriques et des sujets connexes.
Une référence de base sur les espaces algébriques est \cite{Kn}.









\section{Points associés et points faiblement associés}
\label{spaces-divisors-section-associated}

\noindent
Dans le cas des schémas, nous avons introduit deux notions concurrentes
de points associés : les points associés au sens usuel
(Diviseurs, section \ref{divisors-section-associated})
et les points faiblement associés
(Diviseurs, section \ref{divisors-section-weakly-associated}).
Pour un espace algébrique général, la notion de point associé
est pratiquement inutilisable, et nous ne chercherons même pas à la définir.
Si l'espace algébrique est localement noethérien, nous nous autorisons
à employer l'expression « point associé » à la place de
« point faiblement associé », puisque ces notions coïncident pour les
schémas noethériens (Diviseurs, lemme \ref{divisors-lemma-ass-weakly-ass}).
Avant de donner notre définition, nous avons besoin d'un lemme.

\begin{lemma}
\label{spaces-divisors-lemma-associated}
Soit $S$ un schéma. Soit $X$ un espace algébrique sur $S$.
Soit $\mathcal{F}$ un $\mathcal{O}_X$-module quasi-cohérent.
Soit $x \in |X|$. Les conditions suivantes sont équivalentes :
\begin{enumerate}
\item il existe un morphisme \'etale $f : U \to X$, où $U$ est un schéma,
et un point $u \in U$ d'image $x$ tels que $u$ soit faiblement associé
à $f^*\mathcal{F}$,
\item pour tout morphisme \'etale $f : U \to X$, où $U$ est un schéma,
et tout point $u \in U$ d'image $x$, le point $u$ est faiblement associé
à $f^*\mathcal{F}$,
\item l'idéal maximal de $\mathcal{O}_{X, \overline{x}}$
est un idéal premier faiblement associé au module $\mathcal{F}_{\overline{x}}$.
\end{enumerate}
Si $X$ est localement noethérien, elles sont également équivalentes à :
\begin{enumerate}
\item[(4)] il existe un morphisme \'etale $f : U \to X$, où $U$ est un schéma,
et un point $u \in U$ d'image $x$ tels que $u$ soit associé
à $f^*\mathcal{F}$,
\item[(5)] pour tout morphisme \'etale $f : U \to X$, où $U$ est un schéma,
et tout point $u \in U$ d'image $x$, le point $u$ est associé
à $f^*\mathcal{F}$,
\item[(6)] l'idéal maximal de $\mathcal{O}_{X, \overline{x}}$
est un idéal premier associé au module $\mathcal{F}_{\overline{x}}$.
\end{enumerate}
\end{lemma}

\begin{proof}
Choisissons un schéma $U$, un point $u$ et un morphisme \'etale
$f : U \to X$ qui envoie $u$ sur $x$. Relevons $\overline{x}$ en un point
géométrique de $U$ au-dessus de $u$. Rappelons que
$\mathcal{O}_{X, \overline{x}} = \mathcal{O}_{U, u}^{sh}$
où le hensélisé strict est pris relativement au
relèvement choisi de $\overline{x}$ ; voir
Propriétés des espaces, lemme
\ref{spaces-properties-lemma-describe-etale-local-ring}.
Enfin, nous avons
$$
\mathcal{F}_{\overline{x}} =
(f^*\mathcal{F})_u \otimes_{\mathcal{O}_{U, u}}
\mathcal{O}_{X, \overline{x}} =
(f^*\mathcal{F})_u \otimes_{\mathcal{O}_{U, u}}
\mathcal{O}_{U, u}^{sh}
$$
d'après
Propriétés des espaces, lemme \ref{spaces-properties-lemma-stalk-quasi-coherent}.
L'équivalence de (1), (2) et (3) résulte donc de
Compléments sur la platitude, lemme \ref{flat-lemma-weakly-associated-henselization}.
Si $X$ est localement noethérien,
tout $U$ comme ci-dessus est localement noethérien ;
les conditions (1), resp.\ (2), sont donc équivalentes aux conditions (4), resp.\ (5), d'après
Diviseurs, lemme \ref{divisors-lemma-ass-weakly-ass}.
D'autre part, dans le cas localement noethérien,
l'anneau local $\mathcal{O}_{X, \overline{x}}$ est lui aussi noethérien
(Propriétés des espaces, lemme
\ref{spaces-properties-lemma-Noetherian-local-ring-Noetherian}).
L'équivalence de (3) et (6) résulte donc du même lemme
(ou par Algèbre, lemme \ref{algebra-lemma-ass-weakly-ass}).
\end{proof}

\begin{definition}
\label{spaces-divisors-definition-weakly-associated}
Soit $S$ un schéma. Soit $X$ un espace algébrique sur $S$.
Soit $\mathcal{F}$ un faisceau quasi-cohérent sur $X$.
Soit $x \in |X|$.
\begin{enumerate}
\item On dit que $x$ est {\it faiblement associé} à $\mathcal{F}$
si les conditions équivalentes (1), (2) et (3) du
lemme \ref{spaces-divisors-lemma-associated} sont satisfaites.
\item On note $\text{WeakAss}(\mathcal{F})$ l'ensemble des points
faiblement associés à $\mathcal{F}$.
\item Les {\it points faiblement associés à $X$} sont les points
faiblement associés à $\mathcal{O}_X$.
\end{enumerate}
Si $X$ est localement noethérien, on dira
{\it $x$ est associé à $\mathcal{F}$}
si et seulement si $x$ est faiblement associé à $\mathcal{F}$, et l'on pose
$\text{Ass}(\mathcal{F}) = \text{WeakAss}(\mathcal{F})$.
Enfin (toujours sous l'hypothèse que $X$ est localement noethérien),
on dira que {\it $x$ est un point associé de $X$} si et seulement si
$x$ est un point faiblement associé à $X$.
\end{definition}

\noindent
À ce stade, nous pouvons démontrer les lemmes de rigueur.

\begin{lemma}
\label{spaces-divisors-lemma-weakly-ass-support}
Soit $S$ un schéma. Soit $X$ un espace algébrique sur $S$.
Soit $\mathcal{F}$ un $\mathcal{O}_X$-module quasi-cohérent.
Alors $\text{WeakAss}(\mathcal{F}) \subset \text{Supp}(\mathcal{F})$.
\end{lemma}

\begin{proof}
Cela résulte immédiatement des définitions. Le support d'un faisceau abélien
sur $X$ est défini dans Propriétés des espaces, définition
\ref{spaces-properties-definition-support}.
\end{proof}

\begin{lemma}
\label{spaces-divisors-lemma-ses-weakly-ass}
Soit $S$ un schéma. Soit $X$ un espace algébrique sur $S$.
Soit $0 \to \mathcal{F}_1 \to \mathcal{F}_2 \to \mathcal{F}_3 \to 0$
une suite exacte de faisceaux quasi-cohérents sur $X$.
Alors
$\text{WeakAss}(\mathcal{F}_2) \subset
\text{WeakAss}(\mathcal{F}_1) \cup \text{WeakAss}(\mathcal{F}_3)$
et
$\text{WeakAss}(\mathcal{F}_1) \subset \text{WeakAss}(\mathcal{F}_2)$.
\end{lemma}

\begin{proof}
Pour tout point géométrique $\overline{x} \in X$,
la suite des fibres
$0 \to \mathcal{F}_{1, \overline{x}} \to
\mathcal{F}_{2, \overline{x}} \to
\mathcal{F}_{3, \overline{x}} \to 0$
est une suite exacte de $\mathcal{O}_{X, \overline{x}}$-modules.
Le lemme résulte donc de
Algèbre, lemme \ref{algebra-lemma-weakly-ass}.
\end{proof}

\begin{lemma}
\label{spaces-divisors-lemma-weakly-ass-zero}
Soit $S$ un schéma. Soit $X$ un espace algébrique sur $S$.
Soit $\mathcal{F}$ un $\mathcal{O}_X$-module quasi-cohérent.
Alors
$$
\mathcal{F} = (0) \Leftrightarrow \text{WeakAss}(\mathcal{F}) = \emptyset
$$
\end{lemma}

\begin{proof}
Choisissons un schéma $U$ et un morphisme \'etale surjectif $f : U \to X$.
Alors $\mathcal{F}$ est nul si et seulement si $f^*\mathcal{F}$ est nul.
Le lemme résulte donc de la définition et du lemme dans le
cas des schémas ; voir
Diviseurs, lemme \ref{divisors-lemma-weakly-ass-zero}.
\end{proof}

\begin{lemma}
\label{spaces-divisors-lemma-minimal-support-in-weakly-ass}
Soit $S$ un schéma. Soit $X$ un espace algébrique sur $S$.
Soit $\mathcal{F}$ un $\mathcal{O}_X$-module quasi-cohérent.
Soit $x \in |X|$. Si
\begin{enumerate}
\item $x \in \text{Supp}(\mathcal{F})$
\item $x$ est un point de codimension $0$ de $X$
(Propriétés des espaces, définition
\ref{spaces-properties-definition-dimension-local-ring}).
\end{enumerate}
Alors $x \in \text{WeakAss}(\mathcal{F})$. Si $\mathcal{F}$
est un $\mathcal{O}_X$-module de type fini de support schématique $Z$
(Morphismes d'espaces, définition
\ref{spaces-morphisms-definition-scheme-theoretic-support})
et si $x$ est un point de codimension $0$ de $Z$, alors
$x \in \text{WeakAss}(\mathcal{F})$.
\end{lemma}

\begin{proof}
Puisque $x \in \text{Supp}(\mathcal{F})$, la fibre
$\mathcal{F}_{\overline{x}}$ n'est pas nulle. Ainsi,
$\text{WeakAss}(\mathcal{F}_{\overline{x}})$
est non vide d'après
Algèbre, lemme \ref{algebra-lemma-weakly-ass-zero}.
D'autre part, le spectre de $\mathcal{O}_{X, \overline{x}}$
est réduit à un point. Par définition, $x$ est donc faiblement associé à
$\mathcal{F}$. Pour la dernière assertion, l'homomorphisme
$\mathcal{O}_{X, \overline{x}} \to \mathcal{O}_{Z, \overline{z}}$
est surjectif, le spectre de $\mathcal{O}_{Z, \overline{z}}$
est réduit à un point, et $\mathcal{F}_{\overline{x}}$ est un
module non nul sur $\mathcal{O}_{Z, \overline{z}}$.
\end{proof}

\begin{lemma}
\label{spaces-divisors-lemma-minimal-support-in-weakly-ass-decent}
Soit $S$ un schéma. Soit $X$ un espace algébrique sur $S$.
Soit $\mathcal{F}$ un $\mathcal{O}_X$-module quasi-cohérent.
Soit $x \in |X|$. Si
\begin{enumerate}
\item $X$ est décent (par exemple quasi-séparé ou localement séparé),
\item $x \in \text{Supp}(\mathcal{F})$
\item $x$ n'est pas une spécialisation d'un autre point dans
$\text{Supp}(\mathcal{F})$.
\end{enumerate}
Alors $x \in \text{WeakAss}(\mathcal{F})$.
\end{lemma}

\begin{proof}
(Un espace algébrique quasi-séparé est décent ; voir
Espaces décents, section \ref{decent-spaces-section-reasonable-decent}.
Un espace algébrique localement séparé est décent ; voir
Espaces décents, lemme \ref{decent-spaces-lemma-locally-separated-decent}.)
Choisissons un schéma $U$, un point $u \in U$ et un morphisme \'etale
$f : U \to X$ envoyant $u$ sur $x$. D'après
Espaces décents, lemme
\ref{decent-spaces-lemma-decent-no-specializations-map-to-same-point}
si $u' \leadsto u$ est une spécialisation non triviale, alors
$f(u') \not = x$. Ainsi, $u \in \text{Supp}(f^*\mathcal{F})$
n'est pas une spécialisation d'un autre point de
$\text{Supp}(f^*\mathcal{F})$.
Par conséquent, $u \in \text{WeakAss}(f^*\mathcal{F})$ d'après
Diviseurs, lemme \ref{spaces-divisors-lemma-minimal-support-in-weakly-ass}.
\end{proof}

\begin{lemma}
\label{spaces-divisors-lemma-finite-ass}
Soit $S$ un schéma. Soit $X$ un espace algébrique localement noethérien sur $S$.
Soit $\mathcal{F}$ un $\mathcal{O}_X$-module cohérent.
Alors $\text{Ass}(\mathcal{F}) \cap W$ est fini pour
tout ouvert quasi-compact $W \subset |X|$.
\end{lemma}

\begin{proof}
Choisissons un schéma quasi-compact $U$ et un morphisme \'etale $U \to X$
tel que $W$ soit l'image de $|U| \to |X|$. Alors $U$ est un
schéma noethérien, et l'on conclut en appliquant
Diviseurs, lemme \ref{divisors-lemma-finite-ass}.
\end{proof}

\begin{lemma}
\label{spaces-divisors-lemma-restriction-injective-open-contains-weakly-ass}
Soit $S$ un schéma. Soit $X$ un espace algébrique sur $S$.
Soit $\mathcal{F}$ un $\mathcal{O}_X$-module quasi-cohérent.
Si $U \to X$ est un morphisme \'etale tel que
$\text{WeakAss}(\mathcal{F}) \subset \Im(|U| \to |X|)$, alors
$\Gamma(X, \mathcal{F}) \to \Gamma(U, \mathcal{F})$ est injectif.
\end{lemma}

\begin{proof}
Soit $s \in \Gamma(X, \mathcal{F})$ une section dont la restriction à $U$ est nulle.
Soit $\mathcal{F}' \subset \mathcal{F}$ l'image de l'homomorphisme
$\mathcal{O}_X \to \mathcal{F}$ défini par $s$. Alors $\mathcal{F}'|_U = 0$.
Cela entraîne
$\text{WeakAss}(\mathcal{F}') \cap \Im(|U| \to |X|) = \emptyset$
(par définition des points faiblement associés).
D'autre part,
$\text{WeakAss}(\mathcal{F}') \subset \text{WeakAss}(\mathcal{F})$
par le lemme \ref{spaces-divisors-lemma-ses-weakly-ass}. Nous en déduisons que
$\text{WeakAss}(\mathcal{F}') = \emptyset$.
Ainsi, $\mathcal{F}' = 0$ par le lemme \ref{spaces-divisors-lemma-weakly-ass-zero}.
\end{proof}

\begin{lemma}
\label{spaces-divisors-lemma-weakass-pushforward}
Soit $S$ un schéma. Soit $f : X \to Y$ un morphisme quasi-compact et quasi-séparé
d'espaces algébriques sur $S$. Soit $\mathcal{F}$ un
$\mathcal{O}_X$-module quasi-cohérent. Soit $y \in |Y|$ un point qui n'appartient pas à
l'image de $|f|$. Alors $y$ n'est pas faiblement associé à $f_*\mathcal{F}$.
\end{lemma}

\begin{proof}
D'après Morphismes d'espaces, lemme \ref{spaces-morphisms-lemma-pushforward},
le $\mathcal{O}_Y$-module $f_*\mathcal{F}$ est quasi-cohérent, de sorte que
l'énoncé a un sens.
Choisissons un schéma affine $V$, un point $v \in V$ et un morphisme \'etale
$V \to Y$ envoyant $v$ sur $y$. Nous pouvons remplacer
$f : X \to Y$, $\mathcal{F}$, $y$ par
$X \times_Y V \to V$, $\mathcal{F}|_{X \times_Y V}$, $v$.
Nous pouvons donc supposer que $Y$ est un schéma affine.
Dans ce cas, $X$ est quasi-compact ; nous pouvons donc choisir
un schéma affine $U$ et un morphisme \'etale surjectif $U \to X$.
Notons $g : U \to Y$ le composé.
Alors $f_*\mathcal{F} \subset g_*(\mathcal{F}|_U)$.
Par le lemme \ref{spaces-divisors-lemma-ses-weakly-ass},
on se ramène au cas des schémas, traité dans
Diviseurs, lemme \ref{divisors-lemma-weakass-pushforward}.
\end{proof}

\begin{lemma}
\label{spaces-divisors-lemma-check-injective-on-weakass}
Soit $S$ un schéma. Soit $X$ un espace algébrique sur $S$.
Soit $\varphi : \mathcal{F} \to \mathcal{G}$ un homomorphisme de
$\mathcal{O}_X$-modules quasi-cohérents. Supposons que, pour tout
$x \in |X|$, l'une au moins des conditions suivantes soit satisfaite :
\begin{enumerate}
\item $\mathcal{F}_{\overline{x}} \to \mathcal{G}_{\overline{x}}$
est injectif, ou
\item $x \not \in \text{WeakAss}(\mathcal{F})$.
\end{enumerate}
Alors $\varphi$ est injectif.
\end{lemma}

\begin{proof}
Les hypothèses impliquent que $\text{WeakAss}(\Ker(\varphi)) = \emptyset$
et donc que $\Ker(\varphi) = 0$ par le lemme \ref{spaces-divisors-lemma-weakly-ass-zero}.
\end{proof}

\begin{lemma}
\label{spaces-divisors-lemma-weakass-reduced}
Soit $S$ un schéma. Soit $X$ un espace algébrique réduit sur $S$.
Alors les points faiblement associés à $X$ sont exactement les points de codimension $0$
de $X$.
\end{lemma}

\begin{proof}
En travaillant localement pour la topologie \'etale, cela résulte de
Diviseurs, lemme \ref{divisors-lemma-weakass-reduced}
et
Propriétés des espaces, lemme \ref{spaces-properties-lemma-codimension-0-points}.
\end{proof}










\section{Morphismes et points faiblement associés}
\label{spaces-divisors-section-morphisms-weakly-associated}

\begin{lemma}
\label{spaces-divisors-lemma-weakly-ass-reverse-functorial}
Soit $S$ un schéma.
Soit $f : X \to Y$ un morphisme affine d'espaces algébriques sur $S$.
Soit $\mathcal{F}$ un $\mathcal{O}_X$-module quasi-cohérent.
On a alors
$$
\text{WeakAss}_S(f_*\mathcal{F}) \subset f(\text{WeakAss}_X(\mathcal{F}))
$$
\end{lemma}

\begin{proof}
Choisissons un schéma $V$ et un morphisme \'etale surjectif $V \to Y$.
Posons $U = X \times_Y V$. Alors $U \to V$ est un morphisme affine
de schémas. Par définition des points faiblement associés,
le problème est ramené au morphisme de schémas $U \to V$. Ce cas est
traité dans Diviseurs, lemme \ref{divisors-lemma-weakly-ass-reverse-functorial}.
\end{proof}

\begin{lemma}
\label{spaces-divisors-lemma-ass-functorial-equal}
Soit $S$ un schéma.
Soit $f : X \to Y$ un morphisme affine d'espaces algébriques sur $S$.
Soit $\mathcal{F}$ un $\mathcal{O}_X$-module quasi-cohérent.
Si $X$ est localement noethérien, alors on a
$$
\text{WeakAss}_Y(f_*\mathcal{F}) =
f(\text{WeakAss}_X(\mathcal{F}))
$$
\end{lemma}

\begin{proof}
Choisissons un schéma $V$ et un morphisme \'etale surjectif $V \to Y$.
Posons $U = X \times_Y V$. Alors $U \to V$ est un morphisme affine
et $U$ est localement noethérien.
Par définition des points faiblement associés,
le problème est ramené au morphisme de schémas $U \to V$. Ce cas est
traité dans Diviseurs, lemme \ref{divisors-lemma-ass-functorial-equal}.
\end{proof}

\begin{lemma}
\label{spaces-divisors-lemma-weakly-associated-finite}
Soit $S$ un schéma.
Soit $f : X \to Y$ un morphisme fini d'espaces algébriques sur $S$.
Soit $\mathcal{F}$ un $\mathcal{O}_X$-module quasi-cohérent.
Alors $\text{WeakAss}(f_*\mathcal{F}) = f(\text{WeakAss}(\mathcal{F}))$.
\end{lemma}

\begin{proof}
Choisissons un schéma $V$ et un morphisme \'etale surjectif $V \to Y$.
Posons $U = X \times_Y V$. Alors $U \to V$ est un morphisme fini
de schémas. Par définition des points faiblement associés,
le problème est ramené au morphisme de schémas $U \to V$. Ce cas est
traité dans Diviseurs, lemme \ref{divisors-lemma-weakly-associated-finite}.
\end{proof}

\begin{lemma}
\label{spaces-divisors-lemma-weakly-ass-pullback}
Soit $S$ un schéma. Soit $f : X \to Y$ un morphisme d'espaces algébriques
sur $S$. Soit $\mathcal{G}$ un $\mathcal{O}_Y$-module quasi-cohérent.
Soit $x \in |X|$ et $y = f(x) \in |Y|$. Si
\begin{enumerate}
\item $y \in \text{WeakAss}_S(\mathcal{G})$,
\item $f$ est plat en $x$, et
\item la dimension de l'anneau local de la fibre de $f$ au point $x$
est nulle (Morphismes d'espaces, définition
\ref{spaces-morphisms-definition-dimension-fibre}),
\end{enumerate}
alors $x \in \text{WeakAss}(f^*\mathcal{G})$.
\end{lemma}

\begin{proof}
Choisissons un schéma $V$, un point $v \in V$ et un morphisme \'etale $V \to Y$
envoyant $v$ sur $y$. Choisissons un schéma $U$, un point $u \in U$ et un
morphisme \'etale $U \to V \times_Y X$ envoyant $v$ sur un point situé au-dessus de
$v$ et de $x$. Cela est possible car il existe $t \in |V \times_Y X|$
s'envoyant sur $(v, y)$ d'après Propriétés des espaces, lemme
\ref{spaces-properties-lemma-points-cartesian}.
Par définition, la dimension de $\mathcal{O}_{U_v, u}$ est nulle.
Ainsi, $u$ est un point générique de la fibre $U_v$.
Par définition des points faiblement associés,
le problème est ramené au morphisme de schémas $U \to V$.
Ce cas est traité dans
Diviseurs, lemme \ref{divisors-lemma-weakly-ass-pullback}.
\end{proof}

\begin{lemma}
\label{spaces-divisors-lemma-weakly-ass-change-fields}
Soit $K/k$ une extension de corps. Soit $X$ un espace algébrique sur $k$.
Soit $\mathcal{F}$ un $\mathcal{O}_X$-module quasi-cohérent.
Soit $y \in X_K$ d'image $x \in X$. Si $y$ est un point
faiblement associé à l'image inverse $\mathcal{F}_K$, alors $x$
est un point faiblement associé à $\mathcal{F}$.
\end{lemma}

\begin{proof}
C'est la traduction de
Diviseurs, lemme \ref{divisors-lemma-weakly-ass-change-fields}
dans le langage des espaces algébriques. Nous omettons les détails de cette
traduction.
\end{proof}

\begin{lemma}
\label{spaces-divisors-lemma-finite-flat-weak-assassin-up-down}
Soit $S$ un schéma.
Soit $f : X \to Y$ un morphisme fini et plat d'espaces algébriques.
Soit $\mathcal{G}$ un $\mathcal{O}_Y$-module quasi-cohérent.
Soit $x \in |X|$ un point d'image $y \in |Y|$. Alors
$$
x \in \text{WeakAss}(g^*\mathcal{G})
\Leftrightarrow
y \in \text{WeakAss}(\mathcal{G})
$$
\end{lemma}

\begin{proof}
Cela résulte immédiatement du cas des schémas
(Compléments sur la platitude, lemme \ref{flat-lemma-finite-flat-weak-assassin-up-down})
par localisation \'etale.
\end{proof}

\begin{lemma}
\label{spaces-divisors-lemma-etale-weak-assassin-up-down}
Soit $S$ un schéma.
Soit $f : X \to Y$ un morphisme \'etale d'espaces algébriques.
Soit $\mathcal{G}$ un $\mathcal{O}_Y$-module quasi-cohérent.
Soit $x \in |X|$ un point d'image $y \in |Y|$. Alors
$$
x \in \text{WeakAss}(f^*\mathcal{G})
\Leftrightarrow
y \in \text{WeakAss}(\mathcal{G})
$$
\end{lemma}

\begin{proof}
Cela résulte immédiatement de la définition des points faiblement associés
et, en fait, le lemme correspondant pour les schémas
(Compléments sur la platitude, lemme \ref{flat-lemma-etale-weak-assassin-up-down})
est à la base de notre définition.
\end{proof}







\section{Assassin faible relatif}
\label{spaces-divisors-section-relative-weak-assassin}

\noindent
Nous avons besoin de quelques lemmes pour définir cet objet.

\begin{lemma}
\label{spaces-divisors-lemma-locally-noetherian-fibre}
Soit $S$ un schéma. Soit $f : X \to Y$ un morphisme d'espaces algébriques
sur $S$. Soit $y \in |Y|$. Les conditions suivantes sont équivalentes :
\begin{enumerate}
\item il existe un schéma $V$, un point $v \in V$ et un morphisme \'etale $V \to Y$
envoyant $v$ sur $y$ tels que l'espace algébrique $X_v$ soit localement noethérien,
\item pour tout schéma $V$, tout point $v \in V$ et tout morphisme \'etale $V \to Y$
envoyant $v$ sur $y$, l'espace algébrique $X_v$ est localement noethérien, et
\item il existe un corps $k$ et un morphisme $\Spec(k) \to Y$ représentant
$y$ tels que $X_k$ soit localement noethérien.
\end{enumerate}
S'il existe un corps $k_0$ et un monomorphisme $\Spec(k_0) \to Y$
représentant $y$, ces conditions sont également équivalentes à
\begin{enumerate}
\item[(4)] l'espace algébrique $X_{k_0}$ est localement noethérien.
\end{enumerate}
\end{lemma}

\begin{proof}
Remarquons que $X_v = v \times_Y X = \Spec(\kappa(v)) \times_Y X$.
Les implications (2) $\Rightarrow$ (1) $\Rightarrow$ (3) sont donc claires.
Supposons que $\Spec(k) \to Y$ soit un morphisme depuis le spectre d'un corps
tel que $X_k$ soit localement noethérien. Soit $V \to Y$ un morphisme \'etale
depuis un schéma $V$, et soit $v \in V$ un point envoyé sur $y$.
Alors le schéma $v \times_Y \Spec(k)$ est non vide. Choisissons un
point $w \in v \times_Y \Spec(k)$. Considérons les morphismes
$$
X_v \longleftarrow X_w \longrightarrow X_k
$$
Comme $V \to Y$ est \'etale et que $w$ peut être considéré comme un point de
$V \times_Y \Spec(k)$, l'extension $\kappa(w)/k$
est une extension finie séparable de corps
(Morphismes, lemme \ref{morphisms-lemma-etale-over-field}).
Ainsi, $X_w \to X_k$ est un morphisme fini \'etale obtenu par changement de base de
$w \to \Spec(k)$. Par conséquent, $X_w$ est localement noethérien
(Morphismes des espaces, lemme
\ref{spaces-morphisms-lemma-locally-finite-type-locally-noetherian}).
Le morphisme $X_w \to X_v$ est surjectif, affine et plat,
car il est obtenu par changement de base du morphisme surjectif, affine et plat $w \to v$.
Le fait que $X_w$ soit localement noethérien implique alors que
$X_v$ est localement noethérien. Pour le voir, choisissons un morphisme
\'etale surjectif $U \to X$, puis remarquons que
$U_w \to U_v$ est surjectif, affine et plat. En nous plaçant
localement dans le cas affine sur le schéma $U_v$, nous concluons
que $U_w$ est localement noethérien par
Algèbre, lemme \ref{algebra-lemma-descent-Noetherian}.

\medskip\noindent
Enfin, il suffit de prouver que (3) implique (4)
lorsqu'il existe un monomorphisme $\Spec(k_0) \to Y$ dans la classe de $y$.
Alors $\Spec(k) \to Y$ se factorise en $\Spec(k) \to \Spec(k_0) \to Y$.
L'argument donné ci-dessus montre alors que le fait que $X_k$ soit
localement noethérien implique que $X_{k_0}$ est localement noethérien.
\end{proof}

\begin{definition}
\label{spaces-divisors-definition-locally-Noetherian-fibre}
Soit $S$ un schéma. Soit $f : X \to Y$ un morphisme d'espaces algébriques
sur $S$. Soit $y \in |Y|$. Nous disons que {\it la fibre de $f$ au-dessus de $y$ est
localement noethérienne} si les conditions équivalentes (1), (2) et (3)
du lemme \ref{spaces-divisors-lemma-locally-noetherian-fibre} sont satisfaites.
Nous disons que {\it les fibres de $f$ sont localement noethériennes} si cette propriété
est satisfaite pour tout $y \in |Y|$.
\end{definition}

\noindent
Bien entendu, la manière usuelle de garantir que les fibres sont localement noethériennes consiste
à supposer que le morphisme est localement de type fini.

\begin{lemma}
\label{spaces-divisors-lemma-locally-finite-type-locally-Noetherian-fibres}
Soit $S$ un schéma. Soit $f : X \to Y$ un morphisme d'espaces algébriques
sur $S$. Si $f$ est localement de type fini, alors
les fibres de $f$ sont localement noethériennes.
\end{lemma}

\begin{proof}
Cela résulte de Morphismes d'espaces, lemme
\ref{spaces-morphisms-lemma-locally-finite-type-locally-noetherian}
et du fait que le spectre d'un corps est noethérien.
\end{proof}

\begin{lemma}
\label{spaces-divisors-lemma-relative-assassin}
Soit $S$ un schéma. Soit $f : X \to Y$ un morphisme d'espaces algébriques
sur $S$. Soit $x \in |X|$ et $y = f(x) \in |Y|$.
Soit $\mathcal{F}$ un
$\mathcal{O}_X$-module quasi-cohérent. Considérons les diagrammes commutatifs
$$
\xymatrix{
X \ar[d] & X \times_Y V \ar[d] \ar[l] & X_v \ar[d] \ar[l] \\
Y & V \ar[l] & v \ar[l]
}
\quad
\xymatrix{
X \ar[d] & U \ar[d] \ar[l] & U_v \ar[d] \ar[l] \\
Y & V \ar[l] & v \ar[l]
}
\quad
\xymatrix{
x \ar@{|->}[d] &
x' \ar@{|->}[d] \ar@{|->}[l] &
u \ar@{|->}[ld] \ar@{|->}[l] \\
y &
v \ar@{|->}[l]
}
$$
où $V$ et $U$ sont des schémas, $V \to Y$ et $U \to X \times_Y V$
sont \'etales, et $v \in V$, $x' \in |X_v|$, $u \in U$ sont des points
reliés comme dans le dernier diagramme.
Notons $\mathcal{F}|_{X_v}$ et $\mathcal{F}|_{U_v}$
les images inverses de $\mathcal{F}$.
Les conditions suivantes sont équivalentes :
\begin{enumerate}
\item il existe $V, v, x'$ comme ci-dessus tels que $x'$ soit un point
faiblement associé à $\mathcal{F}|_{X_v}$,
\item pour tous $V \to Y, v, x'$ comme ci-dessus, $x'$ est un point
faiblement associé à $\mathcal{F}|_{X_v}$,
\item il existe $U, V, u, v$ comme ci-dessus tels que $u$ soit un point
faiblement associé à $\mathcal{F}|_{U_v}$,
\item pour tous $U, V, u, v$ comme ci-dessus, $u$ est un point
faiblement associé à $\mathcal{F}|_{U_v}$,
\item il existe un corps $k$, un morphisme $\Spec(k) \to Y$ représentant $y$
et un point $t \in |X_k|$ s'envoyant sur $x$, tels que $t$ soit un point
faiblement associé à $\mathcal{F}|_{X_k}$.
\end{enumerate}
S'il existe un corps $k_0$ et un monomorphisme $\Spec(k_0) \to Y$
représentant $y$, ces conditions sont également équivalentes à
\begin{enumerate}
\item[(6)] $x_0$ est un point faiblement associé à $\mathcal{F}|_{X_{k_0}}$
où $x_0 \in |X_{k_0}|$ est l'unique point s'envoyant sur $x$.
\end{enumerate}
Si la fibre de $f$ au-dessus de $y$ est localement noethérienne, alors, dans les
conditions (1), (2), (3), (4) et (6), nous pouvons remplacer
``faiblement associé'' par ``associé''.
\end{lemma}

\begin{proof}
Remarquons qu'étant donnés $V, v, x'$ comme dans le lemme, nous pouvons trouver
$U \to X \times_Y V$ et $u \in U$ s'envoyant sur $x'$
et qu'alors le morphisme $U_v \to X_v$ est \'etale.
Il est donc clair que (1) et (3) sont équivalentes,
de même que (2) et (4). Chacune de ces conditions implique (5).
Nous allons montrer que (5) implique (2).
Supposons donnés $V, v, x'$, ainsi que $\Spec(k) \to X$ et $t \in |X_k|$,
de sorte que $t$ soit un point faiblement
associé à $\mathcal{F}|_{X_k}$.
Nous pouvons choisir un point $w \in v \times_Y \Spec(k)$.
Nous obtenons alors les morphismes
$$
X_v \longleftarrow X_w \longrightarrow X_k
$$
Comme $V \to Y$ est \'etale et que $w$ peut être considéré comme un point de
$V \times_Y \Spec(k)$, l'extension $\kappa(w)/k$
est une extension finie séparable de corps
(Morphismes, lemme \ref{morphisms-lemma-etale-over-field}).
Ainsi, $X_w \to X_k$ est un morphisme fini \'etale obtenu par changement de base de
$w \to \Spec(k)$. Tout point $x''$ de $X_w$ situé au-dessus de $t$
est donc un point faiblement associé à $\mathcal{F}|_{X_w}$ par le
lemme \ref{spaces-divisors-lemma-etale-weak-assassin-up-down}.
Nous pouvons choisir $x''$ s'envoyant sur $x'$
(Propriétés des espaces, lemme \ref{spaces-properties-lemma-points-cartesian}).
Le lemme \ref{spaces-divisors-lemma-weakly-ass-change-fields}
implique alors que $x'$ est un point faiblement associé
à $\mathcal{F}|_{X_v}$.

\medskip\noindent
Pour achever la preuve, il suffit de montrer que les
conditions équivalentes (1) -- (5) impliquent (6) lorsque l'on se donne
$\Spec(k_0) \to Y$ comme dans (6). Dans ce cas, le morphisme
$\Spec(k) \to Y$ de (5) se factorise de manière unique en $\Spec(k) \to \Spec(k_0) \to Y$.
Alors $x_0$ est l'image de $t$ par le morphisme $X_k \to X_{k_0}$.
Le même lemme que ci-dessus montre donc que (6) est vraie.
\end{proof}

\begin{definition}
\label{spaces-divisors-definition-relative-weak-assassin}
Soit $S$ un schéma. Soit $f : X \to Y$ un morphisme d'espaces algébriques
sur $S$. Soit $\mathcal{F}$ un $\mathcal{O}_X$-module quasi-cohérent.
L'{\it assassin faible relatif de $\mathcal{F}$ dans $X$ sur $Y$}
est l'ensemble $\text{WeakAss}_{X/Y}(\mathcal{F}) \subset |X|$
des $x \in |X|$ pour lesquels les conditions équivalentes du
lemme \ref{spaces-divisors-lemma-relative-assassin} sont satisfaites.
Si les fibres de $f$ sont localement noethériennes
(définition \ref{spaces-divisors-definition-locally-Noetherian-fibre}),
on utilise la notation $\text{Ass}_{X/Y}(\mathcal{F})$.
\end{definition}

\noindent
Avec ces notations, nous pouvons énoncer certains des résultats
déjà démontrés pour les schémas.

\begin{lemma}
\label{spaces-divisors-lemma-bourbaki}
Soit $S$ un schéma.
Soit $f : X \to Y$ un morphisme d'espaces algébriques sur $S$.
Soit $\mathcal{F}$ un $\mathcal{O}_X$-module quasi-cohérent.
Soit $\mathcal{G}$ un $\mathcal{O}_Y$-module quasi-cohérent.
Supposons que
\begin{enumerate}
\item $\mathcal{F}$ est plat sur $Y$,
\item $X$ et $Y$ sont localement noethériens, et
\item les fibres de $f$ sont localement noethériennes.
\end{enumerate}
Alors
$$
\text{Ass}_X(\mathcal{F} \otimes_{\mathcal{O}_X} f^*\mathcal{G}) =
\{x \in \text{Ass}_{X/Y}(\mathcal{F})\text{ tel que }
f(x) \in \text{Ass}_Y(\mathcal{G}) \}
$$
\end{lemma}

\begin{proof}
Par localisation \'etale, l'énoncé résulte immédiatement du résultat
pour les schémas ; voir
Diviseurs, lemme \ref{divisors-lemma-bourbaki}.
Le résultat pour les schémas n'est plus général que parce que
nous n'avons pas défini les points associés pour les
espaces algébriques non noethériens (nous devons donc supposer que $X$
et les fibres de $X \to Y$ sont localement noethériens pour pouvoir
même formuler cet énoncé).
\end{proof}

\begin{lemma}
\label{spaces-divisors-lemma-base-change-relative-assassin}
Soit $S$ un schéma. Soit
$$
\xymatrix{
X' \ar[d]_{f'} \ar[r]_{g'} & X \ar[d]^f \\
Y' \ar[r]^g & Y
}
$$
un diagramme cartésien d'espaces algébriques sur $S$.
Soit $\mathcal{F}$ un $\mathcal{O}_X$-module quasi-cohérent
et posons $\mathcal{F}' = (g')^*\mathcal{F}$.
Si $f$ est localement de type fini, alors
\begin{enumerate}
\item $x' \in \text{Ass}_{X'/Y'}(\mathcal{F}')
\Rightarrow g'(x') \in \text{Ass}_{X/Y}(\mathcal{F})$
\item si $x \in \text{Ass}_{X/Y}(\mathcal{F})$, alors, étant donné
$y' \in |Y'|$ tel que $f(x) = g(y')$, il existe
$x' \in \text{Ass}_{X'/Y'}(\mathcal{F}')$
avec $g'(x') = x$ et $f'(x') = y'$.
\end{enumerate}
\end{lemma}

\begin{proof}
Cela résulte du cas des schémas par localisation \'etale.
Donnons tous les détails. Choisissons un schéma
$V$ et un morphisme \'etale surjectif $V \to Y$.
Choisissons un schéma $U$ et un morphisme
\'etale surjectif $U \to V \times_Y X$. Choisissons un schéma $V'$
et un morphisme \'etale surjectif $V' \to V \times_Y Y'$.
Alors $U' = V' \times_V U$ est un schéma et le morphisme
$U' \to X'$ est surjectif et \'etale.

\medskip\noindent
Preuve de (1). Choisissons $u' \in U'$ s'envoyant sur $x'$.
Notons $v' \in V'$ l'image de $u'$.
La condition $x' \in \text{Ass}_{X'/Y'}(\mathcal{F}')$ équivaut
à $u' \in \text{Ass}(\mathcal{F}|_{U'_{v'}})$
par définition (il est légitime d'écrire $\text{Ass}$ au lieu de $\text{WeakAss}$
car $U'_{v'}$ est localement noethérien).
En appliquant Diviseurs, lemme \ref{divisors-lemma-base-change-relative-assassin},
nous voyons que l'image $u \in U$ de $u'$ appartient à
$\text{Ass}(\mathcal{F}|_{U_v})$ où $v \in V$ est l'image de $u$.
Cela signifie à son tour que $g'(x') \in \text{Ass}_{X/Y}(\mathcal{F})$.

\medskip\noindent
Preuve de (2). Choisissons $u \in U$ s'envoyant sur $x$.
Notons $v \in V$ l'image de $u$.
La condition $x \in \text{Ass}_{X/Y}(\mathcal{F})$ équivaut
à $u \in \text{Ass}(\mathcal{F}|_{U_v})$
par définition. Choisissons un point $v' \in V'$ s'envoyant
sur $y' \in |Y'|$ et sur $v \in V$ (ce qui est possible d'après
Propriétés des espaces, lemme \ref{spaces-properties-lemma-points-cartesian}).
Soit $t \in \Spec(\kappa(v') \otimes_{\kappa(v)} \kappa(u))$
un point générique d'une composante irréductible.
Soit $u' \in U'$ l'image de $t$.
En appliquant Diviseurs, lemme \ref{divisors-lemma-base-change-relative-assassin},
nous voyons que $u' \in \text{Ass}(\mathcal{F}'|_{U'_{v'}})$.
Cela signifie à son tour que $x' \in \text{Ass}_{X'/Y'}(\mathcal{F}')$
où $x' \in |X'|$ est l'image de $u'$.
\end{proof}

\begin{lemma}
\label{spaces-divisors-lemma-base-change-relative-assassin-quasi-finite}
Conservons les notations et les hypothèses du
lemme \ref{spaces-divisors-lemma-base-change-relative-assassin}.
Supposons que $g$ soit localement quasi-fini, ou plus généralement que
pour tout $y' \in |Y'|$ le degré de transcendance de $y'/g(y')$ soit $0$.
Alors $\text{Ass}_{X'/Y'}(\mathcal{F}')$ est l'image réciproque de
$\text{Ass}_{X/Y}(\mathcal{F})$.
\end{lemma}

\begin{proof}
Le degré de transcendance d'un point sur son image est défini dans
Morphismes d'espaces, définition
\ref{spaces-morphisms-definition-dimension-fibre}.
Soit $x' \in |X'|$ d'image $x \in |X|$.
Choisissons un schéma $V$ et un morphisme \'etale surjectif $V \to Y$.
Choisissons un schéma $U$ et un morphisme
\'etale surjectif $U \to V \times_Y X$. Choisissons un schéma $V'$
et un morphisme \'etale surjectif $V' \to V \times_Y Y'$.
Alors $U' = V' \times_V U$ est un schéma et le morphisme
$U' \to X'$ est surjectif et \'etale.
Choisissons $u \in U$ s'envoyant sur $x$.
Notons $v \in V$ l'image de $u$.
La condition $x \in \text{Ass}_{X/Y}(\mathcal{F})$ équivaut
à $u \in \text{Ass}(\mathcal{F}|_{U_v})$
par définition. Choisissons un point $u' \in U'$ s'envoyant
sur $x' \in |X'|$ et sur $u \in U$ (ce qui est possible d'après
Propriétés des espaces, lemme \ref{spaces-properties-lemma-points-cartesian}).
Soit $v' \in V'$ l'image de $u'$.
La condition $x' \in \text{Ass}_{X'/Y'}(\mathcal{F}')$ équivaut
à $u' \in \text{Ass}(\mathcal{F}'|_{U'_{v'}})$
par définition.
Le lemme résulte alors de la discussion de
Diviseurs, remarque \ref{divisors-remark-base-change-relative-assassin},
appliquée à $u' \in \Spec(\kappa(v') \otimes_{\kappa(v)} \kappa(u))$.
\end{proof}

\begin{lemma}
\label{spaces-divisors-lemma-relative-weak-assassin-finite}
Soit $S$ un schéma.
Soit $f : X \to Y$ un morphisme d'espaces algébriques sur $S$.
Soit $i : Z \to X$ un morphisme fini.
Soit $\mathcal{G}$ un $\mathcal{O}_Z$-module quasi-cohérent.
Alors $\text{WeakAss}_{X/Y}(i_*\mathcal{G}) =
i(\text{WeakAss}_{Z/Y}(\mathcal{G}))$.
\end{lemma}

\begin{proof}
Cela résulte du cas des schémas
(Diviseurs, lemme \ref{divisors-lemma-relative-weak-assassin-finite})
par localisation \'etale. Nous omettons les détails.
\end{proof}

\begin{lemma}
\label{spaces-divisors-lemma-relative-assassin-constructible}
Soit $Y$ un schéma. Soit $X$ un espace algébrique de présentation finie
sur $Y$. Soit $\mathcal{F}$ un $\mathcal{O}_X$-module quasi-cohérent
de présentation finie. Soit $U \subset X$ un sous-espace ouvert
tel que $U \to Y$ soit quasi-compact. Alors l'ensemble
$$
E = \{y \in Y \mid \text{Ass}_{X_y}(\mathcal{F}_y) \subset |U_y|\}
$$
est une partie localement constructible de $Y$.
\end{lemma}

\begin{proof}
Remarquons que, puisque $Y$ est un schéma, on peut considérer les fibres
$X_y = \Spec(\kappa(y)) \times_Y X$. (De plus, d'après nos définitions,
l'ensemble $\text{Ass}_{X_y}(\mathcal{F}_y)$ est exactement la fibre de
$\text{Ass}_{X/Y}(\mathcal{F}) \to Y$ au-dessus de $y$, mais nous ne l'utiliserons pas.)
La question est locale sur $Y$ : en effet, il suffit de montrer que
$E$ est constructible lorsque $Y$ est affine.
Dans ce cas, $X$ est quasi-compact. Choisissons un schéma affine $W$
et un morphisme \'etale surjectif $\varphi : W \to X$.
Alors $\text{Ass}_{X_y}(\mathcal{F}_y)$ est l'image de
$\text{Ass}_{W_y}(\varphi^*\mathcal{F}_y)$ pour tous les $y \in Y$.
Par conséquent, le lemme résulte du cas des schémas, appliqué à
l'ouvert $\varphi^{-1}(U) \subset W$ et au morphisme $W \to Y$.
Le cas des schémas est traité dans
Compléments sur les morphismes, lemme
\ref{more-morphisms-lemma-relative-assassin-constructible}.
\end{proof}









\section{Idéaux de Fitting}
\label{spaces-divisors-section-fitting-ideals}

\noindent
Cette section prolonge la discussion de
Diviseurs, section \ref{divisors-section-fitting-ideals}.
Soit $S$ un schéma. Soit $X$ un espace algébrique sur $S$.
Soit $\mathcal{F}$ un $\mathcal{O}_X$-module quasi-cohérent de type fini.
Nous pouvons alors construire les idéaux de Fitting
$$
0 = \text{Fit}_{-1}(\mathcal{F}) \subset \text{Fit}_0(\mathcal{F}) \subset
\text{Fit}_1(\mathcal{F}) \subset \ldots \subset \mathcal{O}_X
$$
comme la suite d'idéaux quasi-cohérents caractérisée par la
propriété suivante : pour tout affine $U = \Spec(A)$ \'etale sur $X$,
si $\mathcal{F}|_U$ correspond au $A$-module $M$, alors
$\text{Fit}_i(\mathcal{F})|_U$
correspond à l'idéal $\text{Fit}_i(M) \subset A$.
Cette définition est indépendante des choix et donne un idéal quasi-cohérent, car
si $A \to B$ est un homomorphisme d'anneaux \'etale, alors le $i$-ième idéal de Fitting
de $M \otimes_A B$ sur $B$ est égal à $\text{Fit}_i(M) B$ d'après
Compléments sur l'algèbre, lemme \ref{more-algebra-lemma-fitting-ideal-basics}, partie (3).
Plus précisément (peut-être), l'existence du faisceau quasi-cohérent d'idéaux
$\text{Fit}_0(\mathcal{O}_X)$ résulte (par exemple) de
la description des faisceaux quasi-cohérents dans
Propriétés des espaces, lemme
\ref{spaces-properties-lemma-characterize-quasi-coherent-small-etale}
et de la propriété d'image inverse donnée dans
Diviseurs, lemme \ref{divisors-lemma-base-change-fitting-ideal}.

\medskip\noindent
L'avantage de cette construction des idéaux de Fitting
est que l'on voit immédiatement que leur formation
commute à la localisation \'etale ; par conséquent, de nombreuses propriétés des
idéaux de Fitting se ramènent immédiatement aux propriétés correspondantes
dans le cas des schémas. Nous utiliserons souvent la
discussion de Propriétés des espaces, section
\ref{spaces-properties-section-properties-modules}
pour passer des propriétés des faisceaux quasi-cohérents
sur les schémas à celles sur les espaces algébriques.

\begin{lemma}
\label{spaces-divisors-lemma-base-change-fitting-ideal}
Soit $S$ un schéma.
Soit $f : X \to Y$ un morphisme d'espaces algébriques sur $S$.
Soit $\mathcal{F}$ un $\mathcal{O}_Y$-module quasi-cohérent de type fini.
Alors
$f^{-1}\text{Fit}_i(\mathcal{F}) \cdot \mathcal{O}_X =
\text{Fit}_i(f^*\mathcal{F})$.
\end{lemma}

\begin{proof}
L'énoncé se ramène à
Diviseurs, lemme \ref{divisors-lemma-base-change-fitting-ideal}
par localisation \'etale.
\end{proof}

\begin{lemma}
\label{spaces-divisors-lemma-fitting-ideal-of-finitely-presented}
Soit $S$ un schéma. Soit $X$ un espace algébrique sur $S$.
Soit $\mathcal{F}$ un $\mathcal{O}_X$-module de présentation finie.
Alors $\text{Fit}_r(\mathcal{F})$ est un idéal quasi-cohérent de type fini.
\end{lemma}

\begin{proof}
L'énoncé se ramène à
Diviseurs, lemme \ref{divisors-lemma-fitting-ideal-of-finitely-presented}
par localisation \'etale.
\end{proof}

\begin{lemma}
\label{spaces-divisors-lemma-on-subscheme-cut-out-by-Fit-0}
Soit $S$ un schéma. Soit $X$ un espace algébrique sur $S$.
Soit $\mathcal{F}$ un $\mathcal{O}_X$-module quasi-cohérent de type fini.
Soit $Z_0 \subset X$ le sous-espace fermé défini par
$\text{Fit}_0(\mathcal{F})$.
Soit $Z \subset X$ le support schématique de $\mathcal{F}$.
Alors
\begin{enumerate}
\item $Z \subset Z_0 \subset X$ comme sous-espaces fermés,
\item $|Z| = |Z_0| = \text{Supp}(\mathcal{F})$ comme parties fermées de $|X|$,
\item il existe un $\mathcal{O}_{Z_0}$-module quasi-cohérent de type fini
$\mathcal{G}_0$ avec
$$
(Z_0 \to X)_*\mathcal{G}_0 = \mathcal{F}.
$$
\end{enumerate}
\end{lemma}

\begin{proof}
Rappelons que la formation de $Z$ commute à la localisation \'etale ; voir
Morphismes d'espaces, définition
\ref{spaces-morphisms-definition-scheme-theoretic-support}
(qui utilise Morphismes d'espaces, lemme
\ref{spaces-morphisms-lemma-scheme-theoretic-support}
pour définir $Z$). Ainsi (1) et (2) découlent du cas des schémas ; voir
Diviseurs, lemme \ref{divisors-lemma-on-subscheme-cut-out-by-Fit-0}.
Pour obtenir $\mathcal{G}_0$ comme dans la partie (3), nous pouvons utiliser
le module $\mathcal{G}$ sur $Z$ fourni par Morphismes d'espaces, lemme
\ref{spaces-morphisms-lemma-scheme-theoretic-support}
et poser $\mathcal{G}_0 = (Z \to Z_0)_*\mathcal{G}$.
\end{proof}

\begin{lemma}
\label{spaces-divisors-lemma-fitting-ideal-generate-locally}
Soit $S$ un schéma. Soit $X$ un espace algébrique sur $S$.
Soit $\mathcal{F}$ un module quasi-cohérent de type fini
sur $\mathcal{O}_X$. Soit $x \in |X|$. Alors $\mathcal{F}$ peut être
engendré par $r$ éléments dans un voisinage \'etale de $x$ si et seulement
si $\text{Fit}_r(\mathcal{F})_{\overline{x}} = \mathcal{O}_{X, \overline{x}}$.
\end{lemma}

\begin{proof}
L'énoncé se ramène à
Diviseurs, lemme \ref{divisors-lemma-fitting-ideal-generate-locally}
par localisation \'etale (ainsi qu'à la description de l'anneau
local dans Propriétés des espaces, section
\ref{spaces-properties-section-stalks-structure-sheaf}
et au fait que le hensélisé strict d'un anneau local
est fidèlement plat sur celui-ci, ce qui montre que l'égalité sur le hensélisé
strict équivaut à l'égalité sur l'anneau local).
\end{proof}

\begin{lemma}
\label{spaces-divisors-lemma-fitting-ideal-finite-locally-free}
Soit $S$ un schéma. Soit $X$ un espace algébrique sur $S$.
Soit $\mathcal{F}$ un module quasi-cohérent de type fini
sur $\mathcal{O}_X$. Soit $r \geq 0$. Les conditions suivantes sont équivalentes :
\begin{enumerate}
\item $\mathcal{F}$ est fini localement libre de rang $r$,
\item $\text{Fit}_{r - 1}(\mathcal{F}) = 0$ et
$\text{Fit}_r(\mathcal{F}) = \mathcal{O}_X$, et
\item $\text{Fit}_k(\mathcal{F}) = 0$ pour $k < r$ et
$\text{Fit}_k(\mathcal{F}) = \mathcal{O}_X$ pour $k \geq r$.
\end{enumerate}
\end{lemma}

\begin{proof}
L'énoncé se ramène à
Diviseurs, lemme \ref{divisors-lemma-fitting-ideal-finite-locally-free}
par localisation \'etale.
\end{proof}

\begin{lemma}
\label{spaces-divisors-lemma-locally-free-rank-r-pullback}
Soit $S$ un schéma. Soit $X$ un espace algébrique sur $S$.
Soit $\mathcal{F}$ un module quasi-cohérent de type fini
sur $\mathcal{O}_X$. Les sous-espaces fermés
$$
X = Z_{-1} \supset Z_0 \supset Z_1 \supset Z_2 \ldots
$$
définis par les idéaux de Fitting de $\mathcal{F}$ ont les propriétés
suivantes :
\begin{enumerate}
\item L'intersection $\bigcap Z_r$ est vide.
\item Le foncteur $(\Sch/X)^{opp} \to \textit{Ens}$ défini par la règle
$$
T \longmapsto
\left\{
\begin{matrix}
\{*\} & \text{si }\mathcal{F}_T\text{ est engendré localement par }
\leq r\text{ sections} \\
\emptyset & \text{sinon}
\end{matrix}
\right.
$$
est représentable par le sous-espace ouvert $X \setminus Z_r$.
\item Le foncteur $F_r : (\Sch/X)^{opp} \to \textit{Ens}$ défini par la règle
$$
T \longmapsto
\left\{
\begin{matrix}
\{*\} & \text{si }\mathcal{F}_T\text{ est localement libre de rang }r\\
\emptyset & \text{sinon}
\end{matrix}
\right.
$$
est représentable par le sous-espace localement fermé $Z_{r - 1} \setminus Z_r$
de $X$.
\end{enumerate}
Si $\mathcal{F}$ est de présentation finie, alors
$Z_r \to X$, $X \setminus Z_r \to X$ et $Z_{r - 1} \setminus Z_r \to X$
sont de présentation finie.
\end{lemma}

\begin{proof}
L'énoncé se ramène à
Diviseurs, lemme \ref{divisors-lemma-locally-free-rank-r-pullback}
par localisation \'etale.
\end{proof}

\begin{lemma}
\label{spaces-divisors-lemma-finite-presentation-module}
Soit $S$ un schéma. Soit $X$ un espace algébrique sur $S$.
Soit $\mathcal{F}$ un $\mathcal{O}_X$-module
de présentation finie. Soit $X = Z_{-1} \subset Z_0 \subset Z_1 \subset \ldots$
comme dans le lemme \ref{spaces-divisors-lemma-locally-free-rank-r-pullback}.
Posons $X_r = Z_{r - 1} \setminus Z_r$.
Alors $X' = \coprod_{r \geq 0} X_r$ représente le foncteur
$$
F_{flat} : \Sch/X \longrightarrow \textit{Ens},\quad\quad
T \longmapsto
\left\{
\begin{matrix}
\{*\} & \text{si }\mathcal{F}_T\text{ est plat sur }T\\
\emptyset & \text{sinon}
\end{matrix}
\right.
$$
De plus, $\mathcal{F}|_{X_r}$ est localement libre de rang $r$ et les
morphismes $X_r \to X$ et $X' \to X$ sont de présentation finie.
\end{lemma}

\begin{proof}
L'énoncé se ramène à
Diviseurs, lemme \ref{divisors-lemma-finite-presentation-module}
par localisation \'etale.
\end{proof}














\section{Diviseurs de Cartier effectifs}
\label{spaces-divisors-section-effective-Cartier-divisors}

\noindent
Pour des raisons de commodité, il semble préférable de définir d'abord la notion de
diviseur de Cartier effectif. Rappelons que dans
Morphismes d'espaces, section \ref{spaces-morphisms-section-closed-immersions},
nous avons étudié la correspondance entre les sous-espaces fermés et les faisceaux
quasi-cohérents d'idéaux. De plus, dans
Propriétés des espaces, section
\ref{spaces-properties-section-properties-modules}, nous avons étudié les propriétés
des modules quasi-cohérents, notamment celle d'être « engendré localement par $1$ élément ».
Ces références montrent que la définition suivante est
compatible avec celle donnée pour les schémas.

\begin{definition}
\label{spaces-divisors-definition-effective-Cartier-divisor}
Soit $S$ un schéma. Soit $X$ un espace algébrique sur $S$.
\begin{enumerate}
\item Un {\it sous-espace fermé localement principal} de $X$ est un sous-espace fermé
dont le faisceau d'idéaux est engendré localement par $1$ élément.
\item Un {\it diviseur de Cartier effectif} sur $X$ est un sous-espace fermé
$D \subset X$ tel que le faisceau d'idéaux $\mathcal{I}_D \subset \mathcal{O}_X$
soit un $\mathcal{O}_X$-module inversible.
\end{enumerate}
\end{definition}

\noindent
Ainsi, un diviseur de Cartier effectif est un sous-espace fermé localement principal,
mais la réciproque n'est pas toujours vraie. Les diviseurs de Cartier effectifs sont des
sous-espaces fermés de codimension pure $1$ au sens le plus fort possible :
ils sont en effet définis localement par un seul élément non diviseur de zéro.
Ils sont en particulier nulle part denses.

\begin{lemma}
\label{spaces-divisors-lemma-characterize-effective-Cartier-divisor}
Soit $S$ un schéma. Soit $X$ un espace algébrique sur $S$.
Soit $D \subset X$ un sous-espace fermé.
Les conditions suivantes sont équivalentes :
\begin{enumerate}
\item Le sous-espace $D$ est un diviseur de Cartier effectif sur $X$.
\item Il existe un schéma $U$ et un morphisme \'etale surjectif $U \to X$ tels que
l'image inverse $D \times_X U$ soit un diviseur de Cartier effectif sur $U$.
\item Pour tout schéma $U$ et tout morphisme \'etale $U \to X$,
l'image inverse $D \times_X U$ est un diviseur de Cartier effectif sur $U$.
\item Pour tout $x \in |D|$, il existe un morphisme \'etale
$(U, u) \to (X, x)$ d'espaces algébriques pointés tel que $U = \Spec(A)$
et $D \times_X U = \Spec(A/(f))$, où $f \in A$ est non diviseur de zéro.
\end{enumerate}
\end{lemma}

\begin{proof}
L'équivalence de (1) -- (3) résulte de la
définition \ref{spaces-divisors-definition-effective-Cartier-divisor}
et des références qui la précèdent.
Supposons (1) et soit $x \in |D|$. Choisissons un schéma $W$ et un
morphisme \'etale surjectif
$W \to X$. Choisissons $w \in D \times_X W$ d'image $x$.
D'après (3), $D \times_X W$ est un diviseur de Cartier
effectif sur $W$. On obtient donc un voisinage \'etale affine $U$
en choisissant un voisinage ouvert affine de $w$ dans $W$, comme dans
Diviseurs, lemme \ref{divisors-lemma-characterize-effective-Cartier-divisor}.

\medskip\noindent
Supposons (4). Alors $\mathcal{I}_D|_U$ est inversible d'après
Diviseurs, lemme \ref{divisors-lemma-characterize-effective-Cartier-divisor}.
Puisque l'on peut trouver un recouvrement \'etale de $X$ par la famille
de tous ces $U$ et de $X \setminus D$, on en déduit que
$\mathcal{I}_D$ est un $\mathcal{O}_X$-module inversible.
\end{proof}

\begin{lemma}
\label{spaces-divisors-lemma-complement-locally-principal-closed-subscheme}
Soit $S$ un schéma. Soit $X$ un espace algébrique sur $S$.
Soit $Z \subset X$ un sous-espace fermé localement
principal. Soit $U = X \setminus Z$. Alors $U \to X$ est un morphisme affine.
\end{lemma}

\begin{proof}
La question est locale sur $X$ pour la topologie \'etale ; voir
Morphismes d'espaces, lemmes \ref{spaces-morphisms-lemma-affine-local}
et le
lemme \ref{spaces-divisors-lemma-characterize-effective-Cartier-divisor}.
L'énoncé découle donc du cas des schémas, traité dans
Diviseurs, lemme
\ref{divisors-lemma-complement-locally-principal-closed-subscheme}.
\end{proof}

\begin{lemma}
\label{spaces-divisors-lemma-complement-effective-Cartier-divisor}
Soit $S$ un schéma. Soit $X$ un espace algébrique sur $S$.
Soit $D \subset X$ un diviseur de Cartier effectif.
Soit $U = X \setminus D$. Alors $U \to X$ est un morphisme affine et $U$
est schématiquement dense dans $X$.
\end{lemma}

\begin{proof}
L'affinité résulte du lemme \ref{spaces-divisors-lemma-complement-locally-principal-closed-subscheme}.
La question de la densité est locale sur $X$ pour la topologie \'etale d'après
Morphismes d'espaces, définition
\ref{spaces-morphisms-definition-scheme-theoretically-dense}.
L'énoncé découle donc du cas des schémas, traité dans
Diviseurs, lemme
\ref{divisors-lemma-complement-effective-Cartier-divisor}.
\end{proof}

\begin{lemma}
\label{spaces-divisors-lemma-effective-Cartier-makes-dimension-drop}
Soit $S$ un schéma. Soit $X$ un espace algébrique sur $S$.
Soit $D \subset X$ un diviseur de Cartier effectif.
Soit $x \in |D|$.
Si $\dim_x(X) < \infty$, alors $\dim_x(D) < \dim_x(X)$.
\end{lemma}

\begin{proof}
La définition d'un diviseur de Cartier effectif et celle de la
dimension d'un espace algébrique en un point
(Propriétés des espaces, définition
\ref{spaces-properties-definition-dimension-at-point})
sont toutes deux locales pour la topologie \'etale. Le lemme découle donc du cas des schémas
traité dans
Diviseurs, lemme \ref{divisors-lemma-effective-Cartier-makes-dimension-drop}.
\end{proof}

\begin{definition}
\label{spaces-divisors-definition-sum-effective-Cartier-divisors}
Soit $S$ un schéma. Soit $X$ un espace algébrique sur $S$.
Étant donnés deux diviseurs de Cartier effectifs
$D_1$, $D_2$ sur $X$, on définit $D = D_1 + D_2$ comme le
sous-espace fermé de $X$ correspondant au faisceau quasi-cohérent
d'idéaux
$\mathcal{I}_{D_1}\mathcal{I}_{D_2} \subset \mathcal{O}_S$.
On appelle ce sous-espace la {\it somme des diviseurs de Cartier effectifs
$D_1$ et $D_2$}.
\end{definition}

\noindent
Il est clair que l'on peut de même définir la somme $\sum n_iD_i$ à partir
d'un nombre fini de diviseurs de Cartier effectifs $D_i$ sur $X$
et d'entiers positifs ou nuls $n_i$.

\begin{lemma}
\label{spaces-divisors-lemma-sum-effective-Cartier-divisors}
La somme de deux diviseurs de Cartier effectifs est un diviseur de
Cartier effectif.
\end{lemma}

\begin{proof}
Démonstration omise. Localement pour la topologie \'etale, l'énoncé se ramène au
fait d'algèbre élémentaire suivant : si $f_1, f_2 \in A$ sont non diviseurs de zéro dans un anneau $A$, alors
$f_1f_2 \in A$ est non diviseur de zéro.
\end{proof}

\begin{lemma}
\label{spaces-divisors-lemma-sum-closed-subschemes-effective-Cartier}
Soit $S$ un schéma. Soit $X$ un espace algébrique sur $S$.
Soient $Z, Y$ deux sous-espaces fermés de $X$
de faisceaux d'idéaux $\mathcal{I}$ et $\mathcal{J}$. Si $\mathcal{I}\mathcal{J}$
définit un diviseur de Cartier effectif $D \subset X$, alors $Z$ et $Y$
sont des diviseurs de Cartier effectifs et $D = Z + Y$.
\end{lemma}

\begin{proof}
D'après le lemme \ref{spaces-divisors-lemma-characterize-effective-Cartier-divisor},
l'énoncé se ramène au cas des schémas, traité dans
Diviseurs, lemme \ref{divisors-lemma-sum-closed-subschemes-effective-Cartier}.
\end{proof}

\noindent
Rappelons que nous avons défini l'image inverse d'un sous-espace fermé
par un morphisme quelconque d'espaces algébriques dans
Morphismes d'espaces, définition
\ref{spaces-morphisms-definition-inverse-image-closed-subspace}.

\begin{lemma}
\label{spaces-divisors-lemma-pullback-locally-principal}
Soit $S$ un schéma.
Soit $f : X' \to X$ un morphisme d'espaces algébriques sur $S$.
Soit $Z \subset X$ un sous-espace fermé localement principal.
Alors l'image inverse $f^{-1}(Z)$ est un sous-espace fermé localement
principal de $X'$.
\end{lemma}

\begin{proof}
Démonstration omise.
\end{proof}

\begin{definition}
\label{spaces-divisors-definition-pullback-effective-Cartier-divisor}
Soit $S$ un schéma.
Soit $f : X' \to X$ un morphisme d'espaces algébriques sur $S$.
Soit $D \subset X$
un diviseur de Cartier effectif. On dit que {\it l'image inverse de
$D$ par $f$ est définie} si le sous-espace fermé $f^{-1}(D) \subset X'$
est un diviseur de Cartier effectif. Dans ce cas, on la note indifféremment
$f^*D$ ou $f^{-1}(D)$ et on l'appelle
{\it l'image inverse du diviseur de Cartier effectif}.
\end{definition}

\noindent
La condition selon laquelle $f^{-1}(D)$ est un diviseur de Cartier effectif
est souvent satisfaite en pratique.

\begin{lemma}
\label{spaces-divisors-lemma-pullback-effective-Cartier-defined}
Soit $S$ un schéma.
Soit $f : X \to Y$ un morphisme d'espaces algébriques sur $S$.
Soit $D \subset Y$ un diviseur de Cartier effectif.
L'image inverse de $D$ par $f$ est définie dans chacun des cas suivants :
\begin{enumerate}
\item $f(x) \not \in |D|$ pour tout point faiblement associé $x$ de $X$,
\item $f$ est plat, et
\item ajouter ici d'autres cas au besoin.
\end{enumerate}
\end{lemma}

\begin{proof}
En travaillant localement pour la topologie \'etale, ce lemme se ramène au cas des schémas ; voir
Diviseurs, lemme \ref{divisors-lemma-pullback-effective-Cartier-defined}.
\end{proof}

\begin{lemma}
\label{spaces-divisors-lemma-pullback-effective-Cartier-divisors-additive}
Soit $S$ un schéma.
Soit $f : X' \to X$ un morphisme d'espaces algébriques sur $S$.
Soient $D_1$, $D_2$ des diviseurs de Cartier effectifs sur $X$.
Si les images inverses de $D_1$ et $D_2$ sont définies, alors celle
de $D = D_1 + D_2$ est définie et
$f^*D = f^*D_1 + f^*D_2$.
\end{lemma}

\begin{proof}
Démonstration omise.
\end{proof}




\section{Diviseurs de Cartier effectifs et faisceaux inversibles}
\label{spaces-divisors-section-effective-Cartier-invertible}

\noindent
Puisqu'un diviseur de Cartier effectif possède un faisceau d'idéaux inversible
(définition \ref{spaces-divisors-definition-effective-Cartier-divisor}), la
définition suivante a un sens.

\begin{definition}
\label{spaces-divisors-definition-invertible-sheaf-effective-Cartier-divisor}
Soit $S$ un schéma. Soit $X$ un espace algébrique sur $S$
et soit $D \subset X$ un diviseur de Cartier effectif dont le faisceau
d'idéaux est $\mathcal{I}_D$.
\begin{enumerate}
\item Le {\it faisceau inversible $\mathcal{O}_X(D)$ associé à $D$}
est défini par
$$
\mathcal{O}_X(D) =
\SheafHom_{\mathcal{O}_X}(\mathcal{I}_D, \mathcal{O}_X) =
\mathcal{I}_D^{\otimes -1}.
$$
\item La section canonique, habituellement notée $1$ ou $1_D$, est la
section globale de $\mathcal{O}_X(D)$ correspondant au
morphisme d'inclusion $\mathcal{I}_D \to \mathcal{O}_X$.
\item On écrit
$\mathcal{O}_X(-D) = \mathcal{O}_X(D)^{\otimes -1} = \mathcal{I}_D$.
\item Étant donné un second diviseur de Cartier effectif $D' \subset X$, on définit
$\mathcal{O}_X(D - D') =
\mathcal{O}_X(D) \otimes_{\mathcal{O}_X} \mathcal{O}_X(-D')$.
\end{enumerate}
\end{definition}

\noindent
Quelques remarques. Nous verrons plus bas que l'application
$D \mapsto \mathcal{O}_X(D)$ transforme l'addition des diviseurs de Cartier
effectifs (définition \ref{spaces-divisors-definition-sum-effective-Cartier-divisors})
en l'addition dans le groupe de Picard de $X$
(lemme \ref{spaces-divisors-lemma-invertible-sheaf-sum-effective-Cartier-divisors}).
Cependant, l'expression $D - D'$ de la définition précédente n'a
aucune signification géométrique. Plus précisément, on peut considérer
l'ensemble des diviseurs de Cartier effectifs sur $X$ comme un monoïde commutatif
$\text{EffCart}(X)$ dont l'élément nul est le diviseur de Cartier effectif vide.
L'application $(D, D') \mapsto \mathcal{O}_X(D - D')$ définit alors
un homomorphisme de groupes
$$
\text{EffCart}(X)^{gp} \longrightarrow \Pic(X)
$$
où le membre de gauche est la complétion en groupe de
$\text{EffCart}(X)$. Autrement dit, lorsqu'on écrit $\mathcal{O}_X(D - D')$,
on peut considérer $D - D'$ comme un élément de $\text{EffCart}(X)^{gp}$.

\begin{lemma}
\label{spaces-divisors-lemma-conormal-effective-Cartier-divisor}
Soit $S$ un schéma. Soit $X$ un espace algébrique sur $S$.
Soit $D \subset X$ un diviseur de Cartier effectif.
Alors le faisceau conormal vérifie $\mathcal{C}_{D/X} = \mathcal{I}_D|D =
\mathcal{O}_X(D)^{\otimes -1}|_D$.
\end{lemma}

\begin{proof}
Démonstration omise.
\end{proof}

\begin{lemma}
\label{spaces-divisors-lemma-invertible-sheaf-sum-effective-Cartier-divisors}
Soit $S$ un schéma. Soit $X$ un espace algébrique sur $S$.
Soient $D_1$, $D_2$ des diviseurs de Cartier effectifs sur $X$.
Soit $D = D_1 + D_2$.
Il existe alors un unique isomorphisme
$$
\mathcal{O}_X(D_1) \otimes_{\mathcal{O}_X} \mathcal{O}_X(D_2)
\longrightarrow
\mathcal{O}_X(D)
$$
qui envoie $1_{D_1} \otimes 1_{D_2}$ sur $1_D$.
\end{lemma}

\begin{proof}
Démonstration omise.
\end{proof}

\begin{definition}
\label{spaces-divisors-definition-regular-section}
Soit $S$ un schéma. Soit $X$ un espace algébrique sur $S$.
Soit $\mathcal{L}$ un faisceau inversible sur $X$.
Une section globale $s \in \Gamma(X, \mathcal{L})$ est appelée
{\it section régulière} si l'application $\mathcal{O}_X \to \mathcal{L}$,
$f \mapsto fs$, est injective.
\end{definition}

\begin{lemma}
\label{spaces-divisors-lemma-regular-section-structure-sheaf}
Soit $S$ un schéma.
Soit $X$ un espace algébrique sur $S$.
Soit $f \in \Gamma(X, \mathcal{O}_X)$.
Les conditions suivantes sont équivalentes :
\begin{enumerate}
\item $f$ est une section régulière, et
\item pour tout $x \in X$, l'image $f \in \mathcal{O}_{X, \overline{x}}$
est non diviseur de zéro.
\item pour tout $U = \Spec(A)$ affine et \'etale sur $X$,
la restriction $f|_U$ est un élément non diviseur de zéro de $A$, et
\item il existe un schéma $U$ et un morphisme \'etale surjectif
$U \to X$ tels que $f|_U$ soit une section régulière de $\mathcal{O}_U$.
\end{enumerate}
\end{lemma}

\begin{proof}
Démonstration omise.
\end{proof}

\noindent
Notons qu'une section globale $s$ d'un $\mathcal{O}_X$-module inversible
$\mathcal{L}$ peut être considérée comme un homomorphisme de $\mathcal{O}_X$-modules
$s : \mathcal{O}_X \to \mathcal{L}$. Son dual est donc un
homomorphisme $s : \mathcal{L}^{\otimes -1} \to \mathcal{O}_X$.
(Voir Modules sur les sites, lemme \ref{sites-modules-lemma-constructions-invertible},
pour le faisceau inversible dual.)

\begin{definition}
\label{spaces-divisors-definition-zero-scheme-s}
Soit $S$ un schéma. Soit $X$ un espace algébrique sur $S$.
Soit $\mathcal{L}$ un faisceau inversible.
Soit $s \in \Gamma(X, \mathcal{L})$.
Le {\it schéma des zéros} de $s$ est le sous-espace fermé $Z(s) \subset X$
défini par le faisceau quasi-cohérent d'idéaux
$\mathcal{I} \subset \mathcal{O}_X$ qui est l'image de
l'homomorphisme $s : \mathcal{L}^{\otimes -1} \to \mathcal{O}_X$.
\end{definition}

\begin{lemma}
\label{spaces-divisors-lemma-zero-scheme}
Soit $S$ un schéma. Soit $X$ un espace algébrique sur $S$.
Soit $\mathcal{L}$ un $\mathcal{O}_X$-module inversible.
Soit $s \in \Gamma(X, \mathcal{L})$.
\begin{enumerate}
\item Considérons les immersions fermées $i : Z \to X$ telles que
$i^*s \in \Gamma(Z, i^*\mathcal{L}))$ soit nulle,
ordonnées par inclusion. Le schéma des zéros $Z(s)$ est l'
élément maximal de cet ensemble ordonné.
\item Pour tout morphisme d'espaces algébriques $f : Y \to X$ sur $S$,
on a $f^*s = 0$ dans $\Gamma(Y, f^*\mathcal{L})$ si et seulement si
$f$ se factorise par $Z(s)$.
\item Le schéma des zéros $Z(s)$ est un sous-espace fermé localement principal de $X$.
\item Le schéma des zéros $Z(s)$ est un diviseur de Cartier effectif sur $X$
si et seulement si $s$ est une section régulière de $\mathcal{L}$.
\end{enumerate}
\end{lemma}

\begin{proof}
Démonstration omise.
\end{proof}

\begin{lemma}
\label{spaces-divisors-lemma-characterize-OD}
Soit $S$ un schéma. Soit $X$ un espace algébrique sur $S$.
\begin{enumerate}
\item Si $D \subset X$ est un diviseur de Cartier effectif, alors
la section canonique $1_D$ de $\mathcal{O}_X(D)$ est régulière.
\item Réciproquement, si $s$ est une section régulière du faisceau
inversible $\mathcal{L}$, alors il existe un unique diviseur de Cartier
effectif $D = Z(s) \subset X$ et un unique isomorphisme
$\mathcal{O}_X(D) \to \mathcal{L}$ qui envoie $1_D$ sur $s$.
\end{enumerate}
Les constructions
$D \mapsto (\mathcal{O}_X(D), 1_D)$ et $(\mathcal{L}, s) \mapsto Z(s)$
définissent des applications réciproques
$$
\left\{
\begin{matrix}
\text{diviseurs de Cartier effectifs sur }X
\end{matrix}
\right\}
\leftrightarrow
\left\{
\begin{matrix}
\text{couples }(\mathcal{L}, s)\text{ formés d'un}\\
\mathcal{O}_X\text{-module inversible et d'une section globale régulière}
\end{matrix}
\right\}
$$
\end{lemma}

\begin{proof}
Démonstration omise.
\end{proof}





\section{Diviseurs de Cartier effectifs sur les espaces noethériens}
\label{spaces-divisors-section-Noetherian-effective-Cartier}

\noindent
Dans le cadre localement noethérien, l'essentiel de l'étude des
diviseurs de Cartier effectifs et des sections régulières se simplifie quelque peu.

\begin{lemma}
\label{spaces-divisors-lemma-effective-Cartier-divisor-Sk}
Soient $S$ un schéma et $X$ un espace algébrique localement noethérien
sur $S$. Soit $D \subset X$ un diviseur de Cartier effectif. Si $X$ vérifie
$(S_k)$, alors $D$ vérifie $(S_{k - 1})$.
\end{lemma}

\begin{proof}
D'après notre définition de la propriété $(S_k)$ pour les espaces algébriques
(Propriétés des espaces, section
\ref{spaces-properties-section-types-properties})
et le
lemme \ref{spaces-divisors-lemma-characterize-effective-Cartier-divisor},
l'énoncé découle du cas des schémas
(Diviseurs, lemme \ref{divisors-lemma-effective-Cartier-divisor-Sk}).
\end{proof}

\begin{lemma}
\label{spaces-divisors-lemma-normal-effective-Cartier-divisor-S1}
Soient $S$ un schéma et $X$ un espace algébrique normal localement noethérien
sur $S$. Soit $D \subset X$ un
diviseur de Cartier effectif. Alors $D$ vérifie $(S_1)$.
\end{lemma}

\begin{proof}
D'après notre définition de la normalité pour les espaces algébriques
(Propriétés des espaces, section
\ref{spaces-properties-section-types-properties})
et le
lemme \ref{spaces-divisors-lemma-characterize-effective-Cartier-divisor},
l'énoncé découle du cas des schémas
(Diviseurs, lemme \ref{divisors-lemma-normal-effective-Cartier-divisor-S1}).
\end{proof}

\noindent
Le lemme suivant permet parfois de construire des diviseurs de
Cartier effectifs.

\begin{lemma}
\label{spaces-divisors-lemma-complement-open-affine-effective-cartier-divisor}
Soit $S$ un schéma. Soit $X$ un espace algébrique régulier, noethérien et séparé
sur $S$. Soit $U \subset X$ un ouvert affine dense. Il existe alors un
diviseur de Cartier effectif $D \subset X$ tel que $U = X \setminus D$.
\end{lemma}

\begin{proof}
Nous affirmons que l'espace algébrique $D$ muni de la structure réduite induite
sur $X \setminus U$ (Propriétés des espaces, définition
\ref{spaces-properties-definition-reduced-induced-space})
est le diviseur de Cartier effectif recherché. La construction
de $D$ commute à la localisation \'etale ; voir la démonstration de
Propriétés des espaces, lemme
\ref{spaces-properties-lemma-reduced-closed-subspace}.
Soit $X' \to X$ un morphisme \'etale surjectif avec $X'$ affine.
Puisque $X$ est séparé, on voit que $U' = X' \times_X U$ est
affine. Puisque l'application $|X'| \to |X|$ est ouverte, on voit que $U'$
est dense dans $X'$. Puisque $D' = X' \times_X D$ est la structure de schéma
réduite induite sur $X' \setminus U'$, on en déduit que
$D'$ est un diviseur de Cartier effectif d'après
Diviseurs, lemme
\ref{divisors-lemma-complement-open-affine-effective-cartier-divisor}
et sa démonstration. C'est ce qu'il fallait démontrer.
\end{proof}

\begin{lemma}
\label{spaces-divisors-lemma-Noetherian-regular-separated-pic-effective-Cartier}
Soit $S$ un schéma. Soit $X$ un espace algébrique régulier, noethérien et séparé
sur $S$. Alors tout $\mathcal{O}_X$-module inversible est isomorphe à
$$
\mathcal{O}_X(D - D') =
\mathcal{O}_X(D) \otimes_{\mathcal{O}_X} \mathcal{O}_X(D')^{\otimes -1}
$$
pour certains diviseurs de Cartier effectifs $D, D'$ sur $X$.
\end{lemma}

\begin{proof}
Soit $\mathcal{L}$ un $\mathcal{O}_X$-module inversible.
Choisissons un ouvert affine dense $U \subset X$ tel que
$\mathcal{L}|_U$ soit trivial. Cela est possible car
$X$ possède un sous-espace ouvert dense qui est un schéma ; voir
Propriétés des espaces, proposition
\ref{spaces-properties-proposition-locally-quasi-separated-open-dense-scheme}.
Notons $s : \mathcal{O}_U \to \mathcal{L}|_U$ la trivialisation.
Le complémentaire de $U$ est un diviseur de Cartier effectif
$D$. Nous affirmons que, pour un certain $n > 0$, l'homomorphisme $s$ se prolonge de manière unique en un homomorphisme
$$
s : \mathcal{O}_X(-nD) \longrightarrow \mathcal{L}
$$
Cette affirmation implique le lemme car elle montre que
$\mathcal{L} \otimes_{\mathcal{O}_X} \mathcal{O}_X(nD)$
possède une section globale régulière, donc est isomorphe à
$\mathcal{O}_X(D')$ pour un certain diviseur de Cartier effectif $D'$
d'après le lemme \ref{spaces-divisors-lemma-characterize-OD}.
Pour démontrer l'affirmation, on peut travailler localement pour la topologie \'etale. On peut donc supposer que
$X$ est un schéma affine noethérien. Puisque
$\mathcal{O}_X(-nD) = \mathcal{I}^n$, où $\mathcal{I} = \mathcal{O}_X(-D)$
est le faisceau d'idéaux de $D$ dans $X$, ce cas résulte de
Cohomologie des schémas, lemme \ref{coherent-lemma-homs-over-open}.
\end{proof}

\noindent
Le lemme suivant relève en réalité d'une autre section.

\begin{lemma}
\label{spaces-divisors-lemma-smooth-over-valuation-ring-effective-Cartier}
Soit $R$ un anneau de valuation de corps des fractions $K$.
Soit $X$ un espace algébrique sur $R$ tel que $X \to \Spec(R)$
soit lisse. Pour tout diviseur de Cartier effectif $D \subset X_K$,
il existe un diviseur de Cartier effectif $D' \subset X$
tel que $D'_K = D$.
\end{lemma}

\begin{proof}
Soit $D' \subset X$ l'image schématique de $D \to X_K \to X$.
Puisque ce morphisme est quasi-compact, la formation de $D'$
commute au changement de base plat ; voir
Morphismes d'espaces, lemme
\ref{spaces-morphisms-lemma-flat-base-change-scheme-theoretic-image}.
On obtient en particulier $D'_K = D$. On peut donc
supposer $X$ affine. Écrivons $X = \Spec(A)$.
Alors $X_K = \Spec(A \otimes_R K)$ et $D$ correspond à
un idéal $I \subset A \otimes_R K$. Il faut montrer que
$J = I \cap A$ définit un diviseur de Cartier effectif dans $X$.
Observons d'abord que $A/J$ est plat sur $R$ (car c'est un
$R$-module sans torsion ; voir Compléments d'algèbre, lemme
\ref{more-algebra-lemma-valuation-ring-torsion-free-flat}),
donc $J$ est de type fini d'après
Compléments d'algèbre, lemme
\ref{more-algebra-lemma-flat-finite-type-valuation-ring-finite-presentation}
et le
lemme d'Algèbre \ref{algebra-lemma-extension}.
Il suffit donc de montrer que $J_\mathfrak q \subset A_\mathfrak q$
est engendré par un seul élément pour tout idéal premier $\mathfrak q \subset A$.
Soit $\mathfrak p = R \cap \mathfrak q$. Alors
$R_\mathfrak p$ est un anneau de valuation
(Algèbre, lemme \ref{algebra-lemma-make-valuation-rings}).
Observons en outre que $A_\mathfrak q/\mathfrak p A_\mathfrak q$
est un anneau régulier d'après Algèbre, lemme
\ref{algebra-lemma-characterize-smooth-over-field}.
On peut donc appliquer Compléments d'algèbre, lemme
\ref{more-algebra-lemma-picard-group-generic-fibre-regular}
pour voir que $I(A_\mathfrak q \otimes_R K)$ est engendré par
un seul élément $f \in A_\mathfrak p \otimes_R K$.
Après avoir chassé les dénominateurs, on peut supposer $f \in A_\mathfrak q$.
Soit $\mathfrak c \subset R_\mathfrak p$ l'idéal de contenu de $f$
(voir Compléments d'algèbre, définition \ref{more-algebra-definition-content-ideal}
et Compléments sur la platitude, lemme
\ref{flat-lemma-flat-finite-type-local-valuation-ring-has-content}).
Puisque $R_\mathfrak p$ est un anneau de valuation et
que $\mathfrak c$ est de type fini
(Compléments d'algèbre, lemme \ref{more-algebra-lemma-content-finitely-generated}),
on a $\mathfrak c = (\pi)$ pour un certain $\pi \in R_\mathfrak p$
(Algèbre, lemme \ref{algebra-lemma-characterize-valuation-ring}).
Après avoir remplacé $f$ par $\pi^{-1}f$, on a $f \in A_\mathfrak q$
et $f \not \in \mathfrak pA_\mathfrak q$.
Affirmation : $I_\mathfrak q = (f)$, ce qui achève la démonstration.
Pour la vérifier, observons que $f \in I_\mathfrak q$.
On a donc une surjection $A_\mathfrak q/(f) \to A_\mathfrak q/I_\mathfrak q$
qui devient un isomorphisme après tensorisation sur $R$ avec $K$.
Il suffit donc de montrer que
$A_\mathfrak q/(f)$ est plat sur $R_\mathfrak p$.
Cela résulte de
Algèbre, lemme \ref{algebra-lemma-grothendieck-general}
et du choix de $f$.
\end{proof}









\section{Diviseurs de Cartier effectifs relatifs}
\label{spaces-divisors-section-effective-Cartier-morphisms}

\noindent
Le lemme suivant montre qu'un diviseur de Cartier effectif qui est
plat sur la base est véritablement une « famille de diviseurs de Cartier effectifs »
sur la base. Par exemple, sa restriction à toute fibre est un diviseur de
Cartier effectif.

\begin{lemma}
\label{spaces-divisors-lemma-relative-Cartier}
Soit $S$ un schéma. Soit $f : X \to Y$ un morphisme
d'espaces algébriques sur $S$. Soit $D \subset X$ un sous-espace fermé.
Supposons que
\begin{enumerate}
\item $D$ soit un diviseur de Cartier effectif, et
\item $D \to Y$ soit un morphisme plat.
\end{enumerate}
Alors, pour tout morphisme de schémas $g : Y' \to Y$, l'image inverse
$(g')^{-1}D$ est un diviseur de Cartier effectif sur $X' = Y' \times_Y X$,
où $g' : X' \to X$ est la projection.
\end{lemma}

\begin{proof}
D'après le lemme \ref{spaces-divisors-lemma-characterize-effective-Cartier-divisor},
la propriété d'être un diviseur de Cartier effectif est locale pour la topologie \'etale.
Le lemme se ramène donc immédiatement au cas des schémas,
qui est Diviseurs, lemme \ref{divisors-lemma-relative-Cartier}.
\end{proof}

\noindent
Ce lemme justifie la définition suivante.

\begin{definition}
\label{spaces-divisors-definition-relative-effective-Cartier-divisor}
Soit $S$ un schéma.
Soit $f : X \to Y$ un morphisme d'espaces algébriques sur $S$.
Un {\it diviseur de Cartier effectif relatif} sur $X/Y$ est un
diviseur de Cartier effectif $D \subset X$ tel que $D \to Y$
soit un morphisme plat d'espaces algébriques.
\end{definition}








\section{Fonctions et sections méromorphes}
\label{spaces-divisors-section-meromorphic-functions}

\noindent
Cette section est l'analogue de
Diviseurs, section \ref{divisors-section-meromorphic-functions}.
Attention : il est encore plus facile de commettre des erreurs sur ce sujet dans le
cas des espaces algébriques que dans celui des schémas !

\medskip\noindent
Soit $S$ un schéma. Soit $X$ un espace algébrique sur $S$.
Pour tout schéma $U$ \'etale sur $X$, nous avons défini l'ensemble
$\mathcal{S}(U) \subset \mathcal{O}_X(U)$
des sections régulières de $\mathcal{O}_X$ sur $U$ ; voir la
définition \ref{spaces-divisors-definition-regular-section}. La restriction
d'une section régulière à $V/U$ \'etale est régulière. Par conséquent,
$\mathcal{S} : U \mapsto \mathcal{S}(U)$ est un sous-faisceau (d'ensembles)
de $\mathcal{O}_X$. On note parfois $\mathcal{S} = \mathcal{S}_X$
pour indiquer la dépendance en $X$.
De plus, $\mathcal{S}(U)$
est une partie multiplicative de l'anneau $\mathcal{O}_X(U)$ pour
tout $U$. On peut donc considérer
le préfaisceau d'anneaux
$$
U \longmapsto \mathcal{S}(U)^{-1} \mathcal{O}_X(U),
$$
sur $X_\etale$ et son faisceau associé ; voir
Modules sur les sites, section \ref{sites-modules-section-localizing-sheaves-rings}.

\begin{definition}
\label{spaces-divisors-definition-sheaf-meromorphic-functions}
Soit $S$ un schéma. Soit $X$ un espace algébrique sur $S$.
Le {\it faisceau des fonctions méromorphes sur $X$} est
le faisceau {\it $\mathcal{K}_X$} sur $X_\etale$ associé au préfaisceau
affiché ci-dessus. Une {\it fonction méromorphe} sur $X$
est une section globale de $\mathcal{K}_X$.
\end{definition}

\noindent
Puisque tout élément de chaque $\mathcal{S}(U)$ est non diviseur de zéro dans
$\mathcal{O}_X(U)$, on voit que l'homomorphisme naturel de faisceaux
d'anneaux $\mathcal{O}_X \to \mathcal{K}_X$ est injectif.
De plus, par compatibilité du passage au faisceau associé et de la prise des fibres,
on voit que
$$
\mathcal{K}_{X, \overline{x}} =
\mathcal{S}_{\overline{x}}^{-1}\mathcal{O}_{X, \overline{x}}
$$
pour tout point géométrique $\overline{x}$ de $X$. L'ensemble
$\mathcal{S}_{\overline{x}}$ est une partie de
l'ensemble des éléments non diviseurs de zéro de $\mathcal{O}_{X, \overline{x}}$, mais
ne lui est pas égal en général.

\begin{lemma}
\label{spaces-divisors-lemma-describe-S}
Soit $S$ un schéma. Soit $X$ un espace algébrique sur $S$.
Pour $U$ affine et \'etale sur $X$, l'ensemble
$\mathcal{S}_X(U)$ est l'ensemble des éléments non diviseurs de zéro de
$\mathcal{O}_X(U)$.
\end{lemma}

\begin{proof}
Cela résulte du lemme \ref{spaces-divisors-lemma-regular-section-structure-sheaf}.
\end{proof}

\noindent
Soit ensuite $\mathcal{F}$ un faisceau de $\mathcal{O}_X$-modules
sur $X_\etale$.
Considérons le préfaisceau $U \mapsto \mathcal{S}(U)^{-1}\mathcal{F}(U)$.
Le faisceau qui lui est associé est
$\mathcal{F} \otimes_{\mathcal{O}_X} \mathcal{K}_X$ ; voir
Modules sur les sites, lemme \ref{sites-modules-lemma-simple-invert-module}.

\begin{definition}
\label{spaces-divisors-definition-meromorphic-section}
Soit $S$ un schéma. Soit $X$ un espace algébrique sur $S$.
Soit $\mathcal{F}$ un faisceau de $\mathcal{O}_X$-modules
sur $X_\etale$.
\begin{enumerate}
\item On note $\mathcal{K}_X(\mathcal{F})$ le faisceau de
$\mathcal{K}_X$-modules associé au préfaisceau
$U \mapsto \mathcal{S}(U)^{-1}\mathcal{F}(U)$. De manière équivalente,
$\mathcal{K}_X(\mathcal{F}) =
\mathcal{F} \otimes_{\mathcal{O}_X} \mathcal{K}_X$ (voir ci-dessus).
\item Une {\it section méromorphe de $\mathcal{F}$}
est une section globale de $\mathcal{K}_X(\mathcal{F})$.
\end{enumerate}
\end{definition}

\noindent
On a en particulier
$$
\mathcal{K}_X(\mathcal{F})_{\overline{x}}
=
\mathcal{F}_{\overline{x}}
\otimes_{\mathcal{O}_{X, \overline{x}}} \mathcal{K}_{X, \overline{x}}
=
\mathcal{S}_{\overline{x}}^{-1}\mathcal{F}_{\overline{x}}
$$
pour tout point géométrique $\overline{x}$ de $X$. Il faut toutefois prendre garde
au fait que $\mathcal{S}_{\overline{x}}$ peut ne pas être l'ensemble des
éléments non diviseurs de zéro de l'anneau local \'etale $\mathcal{O}_{X, \overline{x}}$,
comme nous l'avons signalé plus haut. Les faisceaux de sections méromorphes ne sont pas
des modules quasi-cohérents en général, mais ils ont certaines propriétés communes
avec les modules quasi-cohérents.

\begin{lemma}
\label{spaces-divisors-lemma-meromorphic-quasi-coherent}
Soit $S$ un schéma. Soit $X$ un espace algébrique sur $S$. Supposons que
\begin{enumerate}
\item[(a)] tout point faiblement associé de $X$
soit un point de codimension $0$, et
\item[(b)] $X$ vérifie les conditions équivalentes
du lemme de Morphismes d'espaces
\ref{spaces-morphisms-lemma-prepare-normalization}.
\end{enumerate}
Alors
\begin{enumerate}
\item $\mathcal{K}_X$ est un faisceau quasi-cohérent de $\mathcal{O}_X$-algèbres,
\item pour $U \in X_\etale$ affine, $\mathcal{K}_X(U)$
est l'anneau total des fractions de $\mathcal{O}_X(U)$,
\item pour un point géométrique $\overline{x}$, l'ensemble
$\mathcal{S}_{\overline{x}}$
est l'ensemble des éléments non diviseurs de zéro de $\mathcal{O}_{X, \overline{x}}$, et
\item pour un point géométrique $\overline{x}$, l'anneau
$\mathcal{K}_{X, \overline{x}}$ est l'anneau total des fractions de
$\mathcal{O}_{X, \overline{x}}$.
\end{enumerate}
\end{lemma}

\begin{proof}
Le lemme \ref{spaces-divisors-lemma-regular-section-structure-sheaf}
montre que, pour $U \in X_\etale$ affine,
$\mathcal{S}_X(U) \subset \mathcal{O}_X(U)$
est l'ensemble des éléments non diviseurs de zéro de $\mathcal{O}_X(U)$.
Ainsi, le préfaisceau $\mathcal{S}^{-1}\mathcal{O}_X$ est égal à
$$
U \longmapsto Q(\mathcal{O}_X(U))
$$
sur $X_{affine, \etale}$, avec les notations de
Algèbre, exemple \ref{algebra-example-localize-at-prime}.
Observons que les points de codimension $0$ de $X$ correspondent
aux points génériques de $U$ ; voir
Propriétés des espaces, lemme \ref{spaces-properties-lemma-codimension-0-points}.
Ainsi, si $U = \Spec(A)$, alors $A$ est un anneau ayant un nombre fini d'idéaux premiers minimaux
tel que tout idéal premier faiblement associé de $A$ soit minimal.
Il en est de même de toute extension \'etale de $A$ (car son
spectre est un schéma affine \'etale sur $X$ et peut donc
jouer le rôle de $A$ dans la phrase précédente).
Pour montrer que notre préfaisceau est un faisceau quasi-cohérent,
il suffit de montrer que
$$
Q(A) \otimes_A B \longrightarrow Q(B)
$$
est un isomorphisme lorsque $A \to B$ est un homomorphisme \'etale d'anneaux ; voir
Propriétés des espaces, lemme
\ref{spaces-properties-lemma-characterize-quasi-coherent-small-etale}.
(Pour définir la flèche affichée, observons que, puisque $A \to B$
est plat, il envoie les éléments non diviseurs de zéro sur des éléments non diviseurs de zéro.)
D'après les lemmes d'Algèbre
\ref{algebra-lemma-total-ring-fractions-no-embedded-points} et
\ref{algebra-lemma-weakly-ass-zero-divisors}.
on a
$$
Q(A) = \prod\nolimits_{\mathfrak p \subset A\text{ minimal}} A_\mathfrak p
\quad\text{et}\quad
Q(B) = \prod\nolimits_{\mathfrak q \subset B\text{ minimal}} B_\mathfrak q
$$
Puisque $A \to B$ est \'etale, les idéaux premiers minimaux de $B$ sont exactement les
idéaux premiers de $B$ au-dessus des idéaux premiers minimaux
de $A$ (par exemple d'après Compléments d'algèbre, lemme
\ref{more-algebra-lemma-dimension-etale-extension}).
D'après les lemmes d'Algèbre \ref{algebra-lemma-local-dimension-zero-henselian},
\ref{algebra-lemma-characterize-henselian} (13), et
\ref{algebra-lemma-mop-up}, on voit que $A_\mathfrak p \otimes_A B$
est un produit fini d'anneaux locaux finis \'etales sur $A_\mathfrak p$.
Cela implique clairement que $A_\mathfrak p \otimes_A B =
\prod_{\mathfrak q\text{ au-dessus de }\mathfrak p} B_\mathfrak q$,
comme souhaité.

\medskip\noindent
À ce stade, on sait que (1) et (2) sont satisfaites.
Démonstration de (3). Soit $s \in \mathcal{O}_{X, \overline{x}}$
un élément non diviseur de zéro. On peut trouver un voisinage \'etale
$(U, \overline{u}) \to (X, \overline{x})$
et $f \in \mathcal{O}_X(U)$ d'image $s$.
Soit $u \in U$ le point déterminé par $\overline{u}$.
Puisque $\mathcal{O}_{U, u} \to \mathcal{O}_{X, \overline{x}}$
est fidèlement plat (car c'est un hensélisé strict), on voit que
$f$ s'envoie sur un élément non diviseur de zéro de $\mathcal{O}_{U, u}$.
D'après Diviseurs, lemme \ref{divisors-lemma-meromorphic-weakass-finite},
quitte à rétrécir $U$, on obtient que $f$ est non diviseur de zéro
et donc une section de $\mathcal{S}_X(U)$.
L'assertion (4) résulte de (3) par calcul des fibres.
\end{proof}

\begin{lemma}
\label{spaces-divisors-lemma-meromorphic-quasi-coherent-agree}
Soit $S$ un schéma. Soit $X$ un espace algébrique sur $S$. Supposons que
\begin{enumerate}
\item[(a)] tout point faiblement associé de $X$
soit un point de codimension $0$, et
\item[(b)] $X$ vérifie les conditions équivalentes
du lemme de Morphismes d'espaces
\ref{spaces-morphisms-lemma-prepare-normalization}.
\item[(c)] $X$ soit représentable par un schéma $X_0$
(notation maladroite mais provisoire).
\end{enumerate}
Alors le faisceau des fonctions méromorphes $\mathcal{K}_X$
est le faisceau quasi-cohérent de $\mathcal{O}_X$-algèbres
associé au faisceau quasi-cohérent des fonctions méromorphes
$\mathcal{K}_{X_0}$.
\end{lemma}

\begin{proof}
Pour l'équivalence entre $\QCoh(\mathcal{O}_X)$ et
$\QCoh(\mathcal{O}_{X_0})$, voir
Propriétés des espaces, section \ref{spaces-properties-section-quasi-coherent}.
Le lemme est vrai parce que $\mathcal{K}_X$ et $\mathcal{K}_{X_0}$
sont quasi-cohérents et prennent la même valeur sur les ouverts affines correspondants
de $X$ et $X_0$, d'après le lemme \ref{spaces-divisors-lemma-meromorphic-quasi-coherent} et
Diviseurs, lemme \ref{divisors-lemma-meromorphic-weakass-finite}.
\end{proof}

\begin{definition}
\label{spaces-divisors-definition-pullback-meromorphic-sections}
Soit $S$ un schéma. Soit $f : X \to Y$ un morphisme
d'espaces algébriques sur $S$. On dit que {\it les images inverses des fonctions
méromorphes sont définies pour $f$} si, pour tout diagramme commutatif
$$
\xymatrix{
U \ar[r] \ar[d] & X \ar[d] \\
V \ar[r] & Y
}
$$
où $U \in X_\etale$ et $V \in Y_\etale$, et pour toute
section $s \in \mathcal{S}_Y(V)$, l'image inverse
$f^\sharp(s) \in \mathcal{O}_X(U)$ est un élément
de $\mathcal{S}_X(U)$.
\end{definition}

\noindent
Dans ce cas, il existe un morphisme induit
$f^\sharp : f_{small}^{-1}\mathcal{K}_Y \to \mathcal{K}_X$,
autrement dit, on obtient un diagramme commutatif de morphismes
de topos annelés
$$
\xymatrix{
(\Sh(X_\etale), \mathcal{K}_X) \ar[r] \ar[d]^{f_{small}} &
(\Sh(X_\etale), \mathcal{O}_X) \ar[d]^{f_{small}} \\
(\Sh(Y_\etale), \mathcal{K}_Y) \ar[r] &
(\Sh(Y_\etale), \mathcal{O}_Y)
}
$$
On note parfois $f^*(s) = f^\sharp(s)$ pour une
section $s \in \Gamma(Y, \mathcal{K}_Y)$.

\begin{lemma}
\label{spaces-divisors-lemma-pullback-meromorphic-sections-defined}
Soit $S$ un schéma. Soit $f : X \to Y$ un morphisme d'espaces algébriques
sur $S$. Les images inverses des sections méromorphes sont définies
dans chacun des cas suivants :
\begin{enumerate}
\item les points faiblement associés de $X$ sont envoyés
sur des points de codimension $0$ de $Y$,
\item $f$ est plat,
\item ajouter ici d'autres cas au besoin.
\end{enumerate}
\end{lemma}

\begin{proof}
En travaillant localement pour la topologie \'etale, cela se ramène au cas des schémas ; voir
Diviseurs, lemme \ref{divisors-lemma-pullback-meromorphic-sections-defined}.
Pour effectuer cette réduction, utiliser le lemme \ref{spaces-divisors-lemma-regular-section-structure-sheaf}
(description des sections régulières),
la définition \ref{spaces-divisors-definition-weakly-associated} (définition
des points faiblement associés), et
Propriétés des espaces, lemme \ref{spaces-properties-lemma-codimension-0-points}
(description des points de codimension $0$).
\end{proof}

\begin{lemma}
\label{spaces-divisors-lemma-compute-meromorphic}
Soit $S$ un schéma. Soit $X$ un espace algébrique sur $S$. Supposons que
\begin{enumerate}
\item[(a)] tout point faiblement associé de $X$
soit un point de codimension $0$, et
\item[(b)] $X$ vérifie les conditions équivalentes
du lemme de Morphismes d'espaces
\ref{spaces-morphisms-lemma-prepare-normalization},
\item[(c)] tout point de codimension $0$ de $X$ puisse être représenté
par un monomorphisme $\Spec(k) \to X$.
\end{enumerate}
Soit $X^0 \subset |X|$ l'ensemble des points de codimension $0$ de $X$.
Alors on a
$$
\mathcal{K}_X =
\bigoplus\nolimits_{\eta \in X^0} j_{\eta, *}\mathcal{O}_{X, \eta} =
\prod\nolimits_{\eta \in X^0} j_{\eta, *}\mathcal{O}_{X, \eta}
$$
où $j_\eta : \Spec(\mathcal{O}_{X, \eta}) \to X$ est le morphisme canonique
de Schémas, section \ref{schemes-section-points} ; cela a un sens parce que
$X^0$ est contenu dans le lieu schématique de $X$. De même,
pour tout $\mathcal{O}_X$-module quasi-cohérent $\mathcal{F}$,
on obtient la formule
$$
\mathcal{K}_X(\mathcal{F}) =
\bigoplus\nolimits_{\eta \in X^0} j_{\eta, *}\mathcal{F}_\eta =
\prod\nolimits_{\eta \in X^0} j_{\eta, *}\mathcal{F}_\eta
$$
pour le faisceau des sections méromorphes de $\mathcal{F}$.
Enfin, l'anneau des fonctions rationnelles de $X$ est l'anneau des fonctions
méromorphes sur $X$, autrement dit : $R(X) = \Gamma(X, \mathcal{K}_X)$.
\end{lemma}

\begin{proof}
D'après Espaces décents, lemme \ref{decent-spaces-lemma-get-reasonable} et
section \ref{decent-spaces-section-reasonable-decent},
on voit que $X$ est décent\footnote{Réciproquement, si $X$ est décent,
alors la condition (c) est automatiquement satisfaite.}.
Ainsi, $X^0 \subset |X|$ est l'ensemble des points génériques des composantes
irréductibles
(Espaces décents, lemme \ref{decent-spaces-lemma-decent-generic-points}),
et $X^0$ est localement fini dans $|X|$ d'après (b).
Il en résulte que $X^0$ est contenu dans tout ouvert dense de $|X|$.
En particulier, $X^0$ est contenu dans le lieu schématique
(Espaces décents, théorème \ref{decent-spaces-theorem-decent-open-dense-scheme}).
Ainsi, les anneaux locaux $\mathcal{O}_{X, \eta}$ et les morphismes $j_\eta$
sont définis.

\medskip\noindent
Observons qu'une somme directe localement finie de faisceaux de modules
est égale au produit. Ce fait, joint au fait que $X^0$ est localement
fini dans $|X|$, explique les égalités entre
les sommes directes et les produits dans l'énoncé.
Ensuite, puisque $\mathcal{K}_X(\mathcal{F}) =
\mathcal{F} \otimes_{\mathcal{O}_X} \mathcal{K}_X$,
on voit que la seconde égalité résulte de la première.

\medskip\noindent
Soit
$j : Y = \coprod\nolimits_{\eta \in X^0} \Spec(\mathcal{O}_{X, \eta}) \to X$
le morphisme induit par les morphismes $j_\eta$.
Il faut montrer que $\mathcal{K}_X = j_*\mathcal{O}_Y$.
Observons que $\mathcal{K}_Y = \mathcal{O}_Y$, car $Y$ est une somme
disjointe de spectres d'anneaux locaux de dimension $0$ : dans un anneau local
de dimension zéro, tout élément non diviseur de zéro est inversible.
Ensuite, remarquons que les images inverses des fonctions méromorphes
sont définies pour $j$ d'après le
lemme \ref{spaces-divisors-lemma-pullback-meromorphic-sections-defined}.
Cela donne un morphisme
$$
\mathcal{K}_X \longrightarrow j_*\mathcal{O}_Y.
$$
Soit $U \in X_\etale$ affine. D'après le lemme \ref{spaces-divisors-lemma-meromorphic-quasi-coherent},
le membre de gauche prend pour valeur l'anneau total des fractions de $\mathcal{O}_X(U)$.
D'autre part, le membre de droite est égal au produit des
anneaux locaux de $U$ aux points de codimension $0$, c'est-à-dire aux points génériques
de $U$. Ces deux anneaux sont égaux (comme on l'a déjà vu dans la démonstration
du lemme \ref{spaces-divisors-lemma-meromorphic-quasi-coherent}), d'après les lemmes
d'Algèbre
\ref{algebra-lemma-total-ring-fractions-no-embedded-points} et
\ref{algebra-lemma-weakly-ass-zero-divisors}.
Ainsi, notre morphisme est un isomorphisme.

\medskip\noindent
Enfin, il reste à montrer que $R(X) = \Gamma(X, \mathcal{K}_X)$.
Cela résulte du cas des schémas
(Diviseurs, lemme \ref{divisors-lemma-meromorphic-weakass-finite})
appliqué au lieu schématique $X' \subset X$.
En effet, l'anneau des fonctions rationnelles de $X$ est, par définition,
le même que l'anneau des fonctions rationnelles sur $X'$
car celui-ci est un sous-espace ouvert dense de $X$ (voir ci-dessus).
Bien entendu, $R(X')$ coïncide avec l'anneau des fonctions rationnelles
lorsque $X'$ est considéré comme un schéma.
D'autre part, d'après notre description de $\mathcal{K}_X$ ci-dessus,
et le fait, établi plus haut, que $X^0 \subset |X'|$
est contenu dans tout ouvert dense, on voit que
$\Gamma(X, \mathcal{K}_X) = \Gamma(X', \mathcal{K}_{X'})$.
Enfin, utilisons la compatibilité consignée dans le
lemme \ref{spaces-divisors-lemma-meromorphic-quasi-coherent-agree}.
\end{proof}

\begin{definition}
\label{spaces-divisors-definition-regular-meromorphic-section}
Soit $S$ un schéma.
Soit $X$ un espace algébrique sur $S$.
Soit $\mathcal{L}$ un $\mathcal{O}_X$-module inversible.
Une section méromorphe $s$ de $\mathcal{L}$ est dite {\it régulière}
si l'homomorphisme induit $\mathcal{K}_X \to \mathcal{K}_X(\mathcal{L})$
est injectif.
\end{definition}

\noindent
Précisons quand on peut prendre l'image inverse des sections méromorphes (régulières).

\begin{lemma}
\label{spaces-divisors-lemma-meromorphic-sections-pullback}
Soit $S$ un schéma.
Soit $f : X \to Y$ un morphisme d'espaces algébriques sur $S$.
Supposons que les images inverses des fonctions méromorphes soient définies
pour $f$ (voir la
définition \ref{spaces-divisors-definition-pullback-meromorphic-sections}).
\begin{enumerate}
\item Soit $\mathcal{F}$ un faisceau de $\mathcal{O}_Y$-modules.
Il existe une application canonique d'image inverse
$f^* : \Gamma(Y, \mathcal{K}_Y(\mathcal{F})) \to
\Gamma(X, \mathcal{K}_X(f^*\mathcal{F}))$
pour les sections méromorphes de $\mathcal{F}$.
\item Soit $\mathcal{L}$ un $\mathcal{O}_X$-module inversible.
Une section méromorphe régulière $s$ de $\mathcal{L}$ a pour image inverse
la section méromorphe régulière $f^*s$ de $f^*\mathcal{L}$.
\end{enumerate}
\end{lemma}

\begin{proof}
Omise.
\end{proof}

\begin{lemma}
\label{spaces-divisors-lemma-regular-meromorphic-section-exists}
Soit $S$ un schéma. Soit $X$ un espace algébrique sur $S$
vérifiant les conditions (a), (b) et (c) du
lemme \ref{spaces-divisors-lemma-compute-meromorphic}.
Alors tout $\mathcal{O}_X$-module inversible $\mathcal{L}$ possède
une section méromorphe régulière.
\end{lemma}

\begin{proof}
Avec les notations du lemme \ref{spaces-divisors-lemma-compute-meromorphic},
la fibre $\mathcal{L}_\eta$ de $\mathcal{L}$ est définie pour tout
$\eta \in X^0$ et c'est un module libre de rang $1$ sur $\mathcal{O}_{X, \eta}$.
Choisissons un générateur $s_\eta \in \mathcal{L}_\eta$ pour tout $\eta \in X^0$.
Il résulte immédiatement de la description de
$\mathcal{K}_X$ et $\mathcal{K}_X(\mathcal{L})$
dans le lemme \ref{spaces-divisors-lemma-compute-meromorphic}
que $s = \prod s_\eta$ est une section méromorphe régulière de $\mathcal{L}$.
\end{proof}












\section{Proj relatif}
\label{spaces-divisors-section-relative-proj}

\noindent
Cette section reprend la construction du Proj relatif
dans le cadre des espaces algébriques. Son contenu
correspond à celui de Constructions, section
\ref{constructions-section-relative-proj}
et de Diviseurs, section \ref{divisors-section-relative-proj},
dans le cas des schémas.

\begin{situation}
\label{spaces-divisors-situation-relative-proj}
Ici, $S$ est un schéma, $X$ est un espace algébrique sur $S$ et
$\mathcal{A}$ est une $\mathcal{O}_X$-algèbre graduée quasi-cohérente.
\end{situation}

\noindent
Dans la situation \ref{spaces-divisors-situation-relative-proj}, nous allons définir
un foncteur $F : (\Sch/S)_{fppf}^{opp} \to \textit{Ens}$ qui se révélera
être un espace algébrique. Nous suivrons, mutatis mutandis,
la méthode de
Constructions, section \ref{constructions-section-relative-proj}.
Tout d'abord, pour un schéma $T$ sur $S$, nous appelons
{\it quadruplet au-dessus de $T$} un système
$(d, f : T \to X, \mathcal{L}, \psi)$
\begin{enumerate}
\item $d \geq 1$ est un entier,
\item $f : T \to X$ est un morphisme au-dessus de $S$,
\item $\mathcal{L}$ est un $\mathcal{O}_T$-module inversible, et
\item
$\psi : f^*\mathcal{A}^{(d)} \to \bigoplus_{n \geq 0}\mathcal{L}^{\otimes n}$
est un homomorphisme de $\mathcal{O}_T$-algèbres graduées
tel que $f^*\mathcal{A}_d \to \mathcal{L}$ soit surjectif.
\end{enumerate}
Nous disons que les deux quadruplets $(d, f, \mathcal{L}, \psi)$ et
$(d', f', \mathcal{L}', \psi')$ sont {\it équivalents}\footnote{Cette
définition est motivée par
Constructions, lemme \ref{constructions-lemma-equivalent-relative}.
Son avantage est qu'elle définit manifestement
une relation d'équivalence.}
si et seulement si
$f = f'$ et s'il existe un entier positif $m = ad = a'd'$
ainsi qu'un isomorphisme
$\beta : \mathcal{L}^{\otimes a} \to (\mathcal{L}')^{\otimes a'}$
tel que $\beta \circ \psi|_{f^*\mathcal{A}^{(m)}}$ et
$\psi'|_{f^*\mathcal{A}^{(m)}}$ coïncident
comme homomorphismes d'anneaux gradués
$f^*\mathcal{A}^{(m)} \to \bigoplus_{n \geq 0} (\mathcal{L}')^{\otimes mn}$.
Étant donnés un quadruplet $(d, f, \mathcal{L}, \psi)$
et un morphisme $h : T' \to T$, on dispose du quadruplet image inverse
$(d, f \circ h, h^*\mathcal{L}, h^*\psi)$. L'image inverse respecte
la relation d'équivalence. Enfin, pour un schéma {\it quasi-compact} $T$
sur $S$, on pose
$$
F(T) = \text{l'ensemble des classes d'équivalence de quadruplets au-dessus de }T
$$
et, pour un schéma arbitraire $T$ sur $S$, on pose
$$
F(T)
=
\lim_{V \subset T\text{ ouvert quasi-compact}} F(V).
$$
Autrement dit, un élément $\xi$ de $F(T)$ correspond à un système
compatible de choix d'éléments $\xi_V \in F(V)$, où $V$ parcourt les
ouverts quasi-compacts de $T$. Nous avons ainsi défini notre foncteur
\begin{equation}
\label{spaces-divisors-equation-proj}
F : \Sch^{opp} \longrightarrow \textit{Ens}
\end{equation}
Il existe un morphisme de foncteurs $F \to X$ qui envoie le quadruplet
$(d, f, \mathcal{L}, \psi)$ sur $f$.

\begin{lemma}
\label{spaces-divisors-lemma-relative-proj}
Dans la situation \ref{spaces-divisors-situation-relative-proj}, le foncteur $F$ ci-dessus est un
espace algébrique. Pour tout morphisme $g : Z \to X$, où $Z$ est un schéma,
il existe un isomorphisme canonique
$\underline{\text{Proj}}_Z(g^*\mathcal{A}) = Z \times_X F$
compatible avec tout changement de base ultérieur.
\end{lemma}

\begin{proof}
Il suffit de démontrer la seconde assertion, voir
Espaces, lemme \ref{spaces-lemma-representable-over-space}.
Soit $g : Z \to X$ un morphisme, où $Z$ est un schéma.
Soit $F'$ le foncteur des quadruplets associé
à la $\mathcal{O}_Z$-algèbre graduée quasi-cohérente $g^*\mathcal{A}$.
Il existe alors un isomorphisme canonique $F' = Z \times_X F$ qui envoie
un quadruplet $(d, f : T \to Z, \mathcal{L}, \psi)$ pour $F'$
sur $(d, g \circ f, \mathcal{L}, \psi)$ (détails omis, voir la démonstration de
Constructions, lemme \ref{constructions-lemma-proj-base-change}).
D'après les lemmes
\ref{constructions-lemma-equivalent-relative},
\ref{constructions-lemma-relative-proj} et
\ref{constructions-lemma-glueing-gives-functor-proj} de Constructions, ainsi que la
définition \ref{constructions-definition-relative-proj},
$F'$ est représentable par
$\underline{\text{Proj}}_Z(g^*\mathcal{A})$.
\end{proof}

\noindent
Le lemme ci-dessus montre que la définition suivante a un sens.

\begin{definition}
\label{spaces-divisors-definition-relative-proj}
Soient $S$ un schéma et $X$ un espace algébrique sur $S$.
Soit $\mathcal{A}$ un faisceau quasi-cohérent
d'algèbres graduées sur $\mathcal{O}_X$. Le
{\it spectre homogène relatif de $\mathcal{A}$ sur $X$},
ou le {\it spectre homogène de $\mathcal{A}$ sur $X$}, ou le
{\it Proj relatif de $\mathcal{A}$ sur $X$} est l'espace algébrique
$F$ sur $X$ du lemme \ref{spaces-divisors-lemma-relative-proj}.
On le note $\pi : \underline{\text{Proj}}_X(\mathcal{A}) \to X$.
\end{definition}

\noindent
En particulier, le morphisme structural du Proj relatif est représentable
par construction. On peut aussi décrire le Proj relatif par recollement. Soit
$\varphi : U \to X$ un morphisme \'etale surjectif, où $U$ est un schéma.
Posons $R = U \times_X U$, avec pour morphismes de projection $s, t : R  \to U$.
D'après le lemme \ref{spaces-divisors-lemma-relative-proj}, il existe un isomorphisme canonique
$$
\gamma : 
\underline{\text{Proj}}_U(\varphi^*\mathcal{A})
\longrightarrow
\underline{\text{Proj}}_X(\mathcal{A}) \times_X U
$$
au-dessus de $U$. Soit $\alpha : t^*\varphi^*\mathcal{A} \to s^*\varphi^*\mathcal{A}$
l'isomorphisme canonique fourni par la
proposition de Propriétés des espaces
\ref{spaces-properties-proposition-quasi-coherent}.
Alors le diagramme
$$
\xymatrix{
&
\underline{\text{Proj}}_U(\varphi^*\mathcal{A}) \times_{U, s} R
\ar@{=}[r] &
\underline{\text{Proj}}_R(s^*\varphi^*\mathcal{A})
\ar[dd]_{\text{induit par }\alpha} \\
\underline{\text{Proj}}_X(\mathcal{A}) \times_X R
\ar[ru]_{s^*\gamma} \ar[rd]^{t^*\gamma} \\
&
\underline{\text{Proj}}_U(\varphi^*\mathcal{A}) \times_{U, t} R
\ar@{=}[r] &
\underline{\text{Proj}}_R(t^*\varphi^*\mathcal{A})
}
$$
est commutatif (les signes d'égalité proviennent du
lemme \ref{constructions-lemma-relative-proj-base-change} de Constructions).
Ainsi, si l'on note $\mathcal{A}_U$, $\mathcal{A}_R$
les images inverses de $\mathcal{A}$ sur $U$, $R$, alors
$P = \underline{\text{Proj}}_X(\mathcal{A})$ admet pour revêtement \'etale
le schéma $P_U = \underline{\text{Proj}}_U(\mathcal{A}_U)$, et
$P_U \times_P P_U$ est égal à
$P_R = \underline{\text{Proj}}_R(\mathcal{A}_R)$.
Ces remarques permettent de raisonner de la manière habituelle par
localisation \'etale afin de transporter les résultats sur le Proj relatif du cas
des schémas à celui des espaces algébriques.

\begin{lemma}
\label{spaces-divisors-lemma-twists-of-structure-sheaf}
Dans la situation \ref{spaces-divisors-situation-relative-proj}, le Proj relatif est muni
d'un faisceau quasi-cohérent d'algèbres graduées par $\mathbf{Z}$
$\bigoplus_{n \in \mathbf{Z}}
\mathcal{O}_{\underline{\text{Proj}}_X(\mathcal{A})}(n)$
et d'un homomorphisme canonique d'algèbres graduées
$$
\psi :
\pi^*\mathcal{A}
\longrightarrow
\bigoplus\nolimits_{n \geq 0}
\mathcal{O}_{\underline{\text{Proj}}_X(\mathcal{A})}(n)
$$
dont le changement de base à tout schéma sur $X$ coïncide avec
celui du lemme \ref{constructions-lemma-glue-relative-proj-twists} de Constructions.
\end{lemma}

\begin{proof}
Comme dans la discussion qui suit la définition \ref{spaces-divisors-definition-relative-proj},
choisissons un schéma $U$ et un morphisme \'etale surjectif
$U \to X$, puis posons $R = U \times_X U$, avec projections $s, t : R \to U$,
$\mathcal{A}_U = \mathcal{A}|_U$, $\mathcal{A}_R = \mathcal{A}|_R$,
ainsi que $\pi : P = \underline{\text{Proj}}_X(\mathcal{A}) \to X$,
$\pi_U : P_U = \underline{\text{Proj}}_U(\mathcal{A}_U)$ et
$\pi_R : P_R = \underline{\text{Proj}}_U(\mathcal{A}_R)$.
D'après le
lemme \ref{constructions-lemma-glue-relative-proj-twists} de Constructions,
on dispose d'un faisceau quasi-cohérent d'algèbres graduées par $\mathbf{Z}$
sur $\mathcal{O}_{P_U}$,
$\bigoplus_{n \in \mathbf{Z}} \mathcal{O}_{P_U}(n)$
ainsi que d'un homomorphisme canonique
$\psi_U : \pi_U^*\mathcal{A}_U \to \bigoplus_{n \geq 0} \mathcal{O}_{P_U}(n)$
et de même pour $P_R$. D'après le
lemme \ref{constructions-lemma-relative-proj-base-change} de Constructions,
les images inverses de $\mathcal{O}_{P_U}(n)$ et de $\psi_U$ par l'une ou l'autre projection
$P_R \to P_U$ sont $\mathcal{O}_{P_R}(n)$ et $\psi_R$.
D'après la proposition de Propriétés des espaces
\ref{spaces-properties-proposition-quasi-coherent}
on obtient $\mathcal{O}_{P}(n)$ et $\psi$.
Nous omettons de vérifier la compatibilité avec l'image inverse sur
un schéma arbitraire au-dessus de $X$.
\end{proof}

\noindent
Le Proj relatif étant construit, passons à quelques propriétés
élémentaires.

\begin{lemma}
\label{spaces-divisors-lemma-relative-proj-base-change}
Soit $S$ un schéma. Soit $g : X' \to X$ un morphisme d'espaces algébriques
sur $S$, et soit $\mathcal{A}$ un faisceau quasi-cohérent
d'algèbres graduées sur $\mathcal{O}_X$. Il existe alors un isomorphisme canonique
$$
r :
\underline{\text{Proj}}_{X'}(g^*\mathcal{A})
\longrightarrow
X' \times_X \underline{\text{Proj}}_X(\mathcal{A})
$$
ainsi qu'un isomorphisme correspondant
$$
\theta :
r^*\text{pr}_2^*\left(\bigoplus\nolimits_{d \in \mathbf{Z}}
\mathcal{O}_{\underline{\text{Proj}}_X(\mathcal{A})}(d)\right)
\longrightarrow
\bigoplus\nolimits_{d \in \mathbf{Z}}
\mathcal{O}_{\underline{\text{Proj}}_{X'}(g^*\mathcal{A})}(d)
$$
d'algèbres graduées par $\mathbf{Z}$
sur $\mathcal{O}_{\underline{\text{Proj}}_{X'}(g^*\mathcal{A})}$.
\end{lemma}

\begin{proof}
Soit $F$ le foncteur (\ref{spaces-divisors-equation-proj}) et soit $F'$ le
foncteur correspondant défini à l'aide de $g^*\mathcal{A}$ sur $X'$.
Nous affirmons qu'il existe un isomorphisme canonique $r : F' \to X' \times_X F$
de foncteurs (et, bien sûr, $r$ est l'isomorphisme du lemme).
Il suffit de construire la bijection
$r : F'(T) \to X'(T) \times_{X(T)} F(T)$ pour les schémas quasi-compacts $T$
sur $S$. Tout d'abord, si $\xi = (d', f', \mathcal{L}', \psi')$ est un
quadruplet au-dessus de $T$ pour $F'$, on peut poser
$r(\xi) = (f', (d', g \circ f', \mathcal{L}', \psi'))$. Cette expression a un sens
car $(g \circ f')^*\mathcal{A}^{(d)} = (f')^*(g^*\mathcal{A})^{(d)}$.
L'application inverse envoie la paire $(f', (d, f, \mathcal{L}, \psi))$
sur le quadruplet $(d, f', \mathcal{L}, \psi)$. Nous omettons la démonstration
de la dernière assertion (indication : se ramener au cas des schémas par localisation
\'etale et appliquer le lemme de Constructions
\ref{constructions-lemma-relative-proj-base-change}).
\end{proof}

\begin{lemma}
\label{spaces-divisors-lemma-relative-proj-separated}
Dans la situation \ref{spaces-divisors-situation-relative-proj}, le morphisme
$\pi : \underline{\text{Proj}}_X(\mathcal{A}) \to X$
est séparé.
\end{lemma}

\begin{proof}
D'après le lemme \ref{spaces-morphisms-lemma-separated-local} de Morphismes d'espaces
et la construction du Proj relatif, cela résulte du
cas des schémas, traité dans
Constructions, lemme \ref{constructions-lemma-relative-proj-separated}.
\end{proof}

\begin{lemma}
\label{spaces-divisors-lemma-relative-proj-quasi-compact}
Dans la situation \ref{spaces-divisors-situation-relative-proj}, si l'une des conditions suivantes est satisfaite :
\begin{enumerate}
\item $\mathcal{A}$ est de type fini en tant que faisceau
d'algèbres sur $\mathcal{A}_0$,
\item $\mathcal{A}$ est engendrée par $\mathcal{A}_1$ en tant
qu'algèbre sur $\mathcal{A}_0$, et $\mathcal{A}_1$ est un
$\mathcal{A}_0$-module de type fini,
\item il existe un sous-$\mathcal{A}_0$-module quasi-cohérent de type fini
$\mathcal{F} \subset \mathcal{A}_{+}$ tel que
$\mathcal{A}_{+}/\mathcal{F}\mathcal{A}$ soit un faisceau d'idéaux localement nilpotent
de $\mathcal{A}/\mathcal{F}\mathcal{A}$,
\end{enumerate}
alors $\pi : \underline{\text{Proj}}_X(\mathcal{A}) \to X$ est quasi-compact.
\end{lemma}

\begin{proof}
D'après le lemme \ref{spaces-morphisms-lemma-quasi-compact-local} de Morphismes d'espaces
et la construction du Proj relatif, cela résulte du
cas des schémas, traité dans
Diviseurs, lemme \ref{divisors-lemma-relative-proj-quasi-compact}.
\end{proof}

\begin{lemma}
\label{spaces-divisors-lemma-relative-proj-finite-type}
Dans la situation \ref{spaces-divisors-situation-relative-proj},
si $\mathcal{A}$ est de type fini en tant que faisceau
d'algèbres sur $\mathcal{O}_X$, alors
$\pi : \underline{\text{Proj}}_X(\mathcal{A}) \to X$ est de type fini.
\end{lemma}

\begin{proof}
D'après le lemme \ref{spaces-morphisms-lemma-finite-type-local} de Morphismes d'espaces
et la construction du Proj relatif, cela résulte du
cas des schémas, traité dans
Diviseurs, lemme \ref{divisors-lemma-relative-proj-finite-type}.
\end{proof}

\begin{lemma}
\label{spaces-divisors-lemma-relative-proj-universally-closed}
Dans la situation \ref{spaces-divisors-situation-relative-proj}, si
$\mathcal{O}_X \to \mathcal{A}_0$
est un homomorphisme d'algèbres entier\footnote{Autrement dit, la clôture intégrale
de $\mathcal{O}_X$ dans $\mathcal{A}_0$, voir
Morphismes d'espaces, définition
\ref{spaces-morphisms-definition-integral-closure}, est égale à
$\mathcal{A}_0$.} et si $\mathcal{A}$ est de type fini en tant
qu'algèbre sur $\mathcal{A}_0$, alors
$\pi : \underline{\text{Proj}}_X(\mathcal{A}) \to X$ est universellement fermé.
\end{lemma}

\begin{proof}
D'après le lemme de Morphismes d'espaces
\ref{spaces-morphisms-lemma-universally-closed-local}
et la construction du Proj relatif, cela résulte du
cas des schémas, traité dans
Diviseurs, lemme \ref{divisors-lemma-relative-proj-universally-closed}.
\end{proof}

\begin{lemma}
\label{spaces-divisors-lemma-relative-proj-proper}
Dans la situation \ref{spaces-divisors-situation-relative-proj},
les conditions suivantes sont équivalentes :
\begin{enumerate}
\item $\mathcal{A}_0$ est un $\mathcal{O}_X$-module de type fini
et $\mathcal{A}$ est de type fini en tant qu'algèbre sur $\mathcal{A}_0$,
\item $\mathcal{A}_0$ est un $\mathcal{O}_X$-module de type fini
et $\mathcal{A}$ est de type fini en tant qu'algèbre sur $\mathcal{O}_X$.
\end{enumerate}
Si ces conditions sont satisfaites, alors
$\pi : \underline{\text{Proj}}_X(\mathcal{A}) \to X$
est propre.
\end{lemma}

\begin{proof}
D'après le lemme de Morphismes d'espaces
\ref{spaces-morphisms-lemma-proper-local}
et la construction du Proj relatif, cela résulte du
cas des schémas, traité dans
Diviseurs, lemme \ref{divisors-lemma-relative-proj-universally-closed}.
\end{proof}

\begin{lemma}
\label{spaces-divisors-lemma-relative-proj-generated-in-degree-1}
Soient $S$ un schéma et $X$ un espace algébrique sur $S$.
Soit $\mathcal{A}$ un faisceau quasi-cohérent d'algèbres graduées sur $\mathcal{O}_X$
engendré en tant qu'algèbre sur $\mathcal{A}_0$ par $\mathcal{A}_1$.
Avec $P = \underline{\text{Proj}}_X(\mathcal{A})$, on a
\begin{enumerate}
\item $P$ représente le foncteur $F_1$ qui associe à
$T$ sur $S$ l'ensemble des classes d'isomorphie de
triplets $(f, \mathcal{L}, \psi)$, où $f : T \to X$ est un morphisme
au-dessus de $S$, $\mathcal{L}$ est un $\mathcal{O}_T$-module inversible, et
$\psi : f^*\mathcal{A} \to \bigoplus_{n \geq 0} \mathcal{L}^{\otimes n}$
est un homomorphisme de $\mathcal{O}_T$-algèbres graduées induisant une surjection
$f^*\mathcal{A}_1 \to \mathcal{L}$,
\item l'homomorphisme canonique $\pi^*\mathcal{A}_1 \to \mathcal{O}_P(1)$ est
surjectif, et
\item chaque $\mathcal{O}_P(n)$ est inversible,
et les homomorphismes de multiplication induisent des isomorphismes
$\mathcal{O}_P(n) \otimes_{\mathcal{O}_P} \mathcal{O}_P(m) =
\mathcal{O}_P(n + m)$.
\end{enumerate}
\end{lemma}

\begin{proof}
Omise.
Voir Constructions, lemme \ref{constructions-lemma-apply-relative},
pour le cas des schémas.
\end{proof}







\section{Fonctorialité du Proj relatif}
\label{spaces-divisors-section-functoriality-relative-proj}

\noindent
Cette section est l'analogue de
Constructions, section \ref{constructions-section-functoriality-relative-proj}.

\begin{lemma}
\label{spaces-divisors-lemma-morphism-relative-proj}
Soient $S$ un schéma et $X$ un espace algébrique sur $S$.
Soit $\psi : \mathcal{A} \to \mathcal{B}$ un homomorphisme
d'algèbres graduées quasi-cohérentes sur $\mathcal{O}_X$. Posons
$P = \underline{\text{Proj}}_X(\mathcal{A}) \to X$ et
$Q = \underline{\text{Proj}}_X(\mathcal{B}) \to X$.
Il existe un sous-espace ouvert canonique
$U(\psi) \subset Q$ et un morphisme canonique
d'espaces algébriques
$$
r_\psi :
U(\psi)
\longrightarrow
P
$$
au-dessus de $X$ et un homomorphisme d'algèbres graduées par $\mathbf{Z}$ sur $\mathcal{O}_{U(\psi)}$
$$
\theta = \theta_\psi :
r_\psi^*\left(
\bigoplus\nolimits_{d \in \mathbf{Z}} \mathcal{O}_P(d)
\right)
\longrightarrow
\bigoplus\nolimits_{d \in \mathbf{Z}} \mathcal{O}_{U(\psi)}(d).
$$
Le triplet $(U(\psi), r_\psi, \theta)$ est caractérisé par la propriété
que, pour tout schéma $W$ \'etale sur $X$, le triplet
$$
(U(\psi) \times_X W,\quad
r_\psi|_{U(\psi) \times_X W} : U(\psi) \times_X W \to  P \times_X W,\quad
\theta|_{U(\psi) \times_X W})
$$
est égal au triplet associé à $\psi : \mathcal{A}|_W \to \mathcal{B}|_W$
par Constructions, lemme \ref{constructions-lemma-morphism-relative-proj}.
\end{lemma}

\begin{proof}
Ce lemme résulte de la localisation \'etale et du cas des schémas, voir la
discussion qui suit la
définition \ref{spaces-divisors-definition-relative-proj}. Détails omis.
\end{proof}

\begin{lemma}
\label{spaces-divisors-lemma-morphism-relative-proj-transitive}
Soient $S$ un schéma et $X$ un espace algébrique sur $S$.
Soient $\mathcal{A}$, $\mathcal{B}$ et $\mathcal{C}$ des
algèbres graduées quasi-cohérentes sur $\mathcal{O}_X$.
Posons $P = \underline{\text{Proj}}_X(\mathcal{A})$,
$Q = \underline{\text{Proj}}_X(\mathcal{B})$ et
$R = \underline{\text{Proj}}_X(\mathcal{C})$.
Soient $\varphi : \mathcal{A} \to \mathcal{B}$,
$\psi : \mathcal{B} \to \mathcal{C}$ des homomorphismes d'algèbres graduées sur $\mathcal{O}_X$.
Alors
$$
U(\psi \circ \varphi) = r_\varphi^{-1}(U(\psi))
\quad
\text{et}
\quad
r_{\psi \circ \varphi}
=
r_\varphi \circ r_\psi|_{U(\psi \circ \varphi)}.
$$
De plus,
$$
\theta_\psi \circ r_\psi^*\theta_\varphi
=
\theta_{\psi \circ \varphi}
$$
avec les notations évidentes.
\end{lemma}

\begin{proof}
Omise.
\end{proof}

\begin{lemma}
\label{spaces-divisors-lemma-surjective-graded-rings-map-relative-proj}
Sous les hypothèses et avec les notations du lemme \ref{spaces-divisors-lemma-morphism-relative-proj}
ci-dessus, supposons que $\mathcal{A}_d \to \mathcal{B}_d$ soit surjectif pour
$d \gg 0$. Alors
\begin{enumerate}
\item $U(\psi) = Q$,
\item $r_\psi : Q \to R$ est une immersion fermée, et
\item les homomorphismes $\theta : r_\psi^*\mathcal{O}_P(n) \to \mathcal{O}_Q(n)$
sont surjectifs mais ne sont pas en général des isomorphismes (même si
$\mathcal{A} \to \mathcal{B}$ est surjectif).
\end{enumerate}
\end{lemma}

\begin{proof}
Cela résulte du cas des schémas
(Constructions, lemme
\ref{constructions-lemma-surjective-graded-rings-map-relative-proj})
par localisation \'etale.
\end{proof}

\begin{lemma}
\label{spaces-divisors-lemma-eventual-iso-graded-rings-map-relative-proj}
Sous les hypothèses et avec les notations du lemme \ref{spaces-divisors-lemma-morphism-relative-proj}
ci-dessus, supposons que $\mathcal{A}_d \to \mathcal{B}_d$ soit un isomorphisme pour tout
$d \gg 0$. Alors
\begin{enumerate}
\item $U(\psi) = Q$,
\item $r_\psi : Q \to P$ est un isomorphisme, et
\item les homomorphismes $\theta : r_\psi^*\mathcal{O}_P(n) \to \mathcal{O}_Q(n)$
sont des isomorphismes.
\end{enumerate}
\end{lemma}

\begin{proof}
Cela résulte du cas des schémas
(Constructions, lemme
\ref{constructions-lemma-eventual-iso-graded-rings-map-relative-proj})
par localisation \'etale.
\end{proof}

\begin{lemma}
\label{spaces-divisors-lemma-surjective-generated-degree-1-map-relative-proj}
Sous les hypothèses et avec les notations du lemme \ref{spaces-divisors-lemma-morphism-relative-proj}
ci-dessus, supposons que $\mathcal{A}_d \to \mathcal{B}_d$ soit surjectif pour $d \gg 0$
et que $\mathcal{A}$ soit engendrée par $\mathcal{A}_1$ sur $\mathcal{A}_0$.
Alors
\begin{enumerate}
\item $U(\psi) = Q$,
\item $r_\psi : Q \to P$ est une immersion fermée, et
\item les homomorphismes $\theta : r_\psi^*\mathcal{O}_P(n) \to \mathcal{O}_Q(n)$
sont des isomorphismes.
\end{enumerate}
\end{lemma}

\begin{proof}
Cela résulte du cas des schémas
(Constructions, lemme
\ref{constructions-lemma-surjective-generated-degree-1-map-relative-proj})
par localisation \'etale.
\end{proof}










\section{Faisceaux inversibles et morphismes vers le Proj relatif}
\label{spaces-divisors-section-invertible-relative-proj}

\noindent
Il semble que le lemme suivant puisse être utile quelque part.
La situation est la suivante :
\begin{enumerate}
\item Soient $S$ un schéma et $Y$ un espace algébrique sur $S$.
\item Soit $\mathcal{A}$ une algèbre graduée quasi-cohérente sur $\mathcal{O}_Y$.
\item Notons $\pi : \underline{\text{Proj}}_Y(\mathcal{A}) \to Y$ le Proj
relatif de $\mathcal{A}$ sur $Y$.
\item Soit $f : X \to Y$ un morphisme d'espaces algébriques sur $S$.
\item Soit $\mathcal{L}$ un $\mathcal{O}_X$-module inversible.
\item Soit
$\psi : f^*\mathcal{A} \to \bigoplus_{d \geq 0} \mathcal{L}^{\otimes d}$
un homomorphisme d'algèbres graduées sur $\mathcal{O}_X$.
\end{enumerate}
Ces données étant fixées, soit $U(\psi) \subset X$ le sous-espace ouvert tel que
$$
|U(\psi)| = \bigcup\nolimits_{d \geq 1}
\{\text{lieu où }f^*\mathcal{A}_d \to \mathcal{L}^{\otimes d}
\text{ est surjectif}\}
$$
La formation de $U(\psi) \subset X$ commute au
changement de base par tout morphisme $X' \to X$.

\begin{lemma}
\label{spaces-divisors-lemma-invertible-map-into-relative-proj}
Sous les hypothèses et avec les notations ci-dessus, le morphisme
$\psi$ induit un morphisme canonique d'espaces algébriques sur $Y$
$$
r_{\mathcal{L}, \psi} :
U(\psi) \longrightarrow \underline{\text{Proj}}_Y(\mathcal{A})
$$
ainsi qu'un homomorphisme d'algèbres graduées sur $\mathcal{O}_{U(\psi)}$
$$
\theta :
r_{\mathcal{L}, \psi}^*\left(
\bigoplus\nolimits_{d \geq 0}
\mathcal{O}_{\underline{\text{Proj}}_Y(\mathcal{A})}(d)
\right)
\longrightarrow
\bigoplus\nolimits_{d \geq 0} \mathcal{L}^{\otimes d}|_{U(\psi)}
$$
caractérisé par les propriétés suivantes :
\begin{enumerate}
\item Pour $V \to Y$ \'etale et $d \geq 0$, le diagramme
$$
\xymatrix{
\mathcal{A}_d(V) \ar[d]_{\psi} \ar[r]_{\psi} &
\Gamma(V \times_Y X, \mathcal{L}^{\otimes d}) \ar[d]^{restriction} \\
\Gamma(V \times_Y \underline{\text{Proj}}_Y(\mathcal{A}),
\mathcal{O}_{\underline{\text{Proj}}_Y(\mathcal{A})}(d)) \ar[r]^-\theta &
\Gamma(V \times_Y U(\psi), \mathcal{L}^{\otimes d})
}
$$
est commutatif.
\item Pour tout $d \geq 1$ et tout morphisme $W \to X$, où $W$ est un schéma,
tel que $\psi|_W : f^*\mathcal{A}_d|_W \to \mathcal{L}^{\otimes d}|_W$
soit surjectif, on a (a) $W \to X$ se factorise par $U(\psi)$ et
(b) le composé de $W \to U(\psi)$ avec $r_{\mathcal{L}, \psi}$
coïncide avec le morphisme $W \to \underline{\text{Proj}}_Y(\mathcal{A})$
dont l'existence résulte de la construction de $\underline{\text{Proj}}_Y(\mathcal{A})$,
voir la définition \ref{spaces-divisors-definition-relative-proj}.
\item Considérons un diagramme commutatif
$$
\xymatrix{
X' \ar[r]_{g'} \ar[d]_{f'} & X \ar[d]^f \\
Y' \ar[r]^g & Y
}
$$
où $X'$ et $Y'$ sont des schémas, posons $\mathcal{A}' = g^*\mathcal{A}$
et $\mathcal{L}' = (g')^*\mathcal{L}$ et notons
$\psi' : (f')^*\mathcal{A} \to \bigoplus_{d \geq 0} (\mathcal{L}')^{\otimes d}$
l'image inverse de $\psi$. Notons $U(\psi')$, $r_{\psi', \mathcal{L}'}$
et $\theta'$ le sous-espace ouvert, le morphisme et l'homomorphisme construits
dans Constructions, lemme \ref{spaces-divisors-lemma-invertible-map-into-relative-proj}.
Alors $U(\psi') = (g')^{-1}(U(\psi))$
et $r_{\psi', \mathcal{L}'}$ coïncide avec le changement de base
de $r_{\psi, \mathcal{L}}$ via l'isomorphisme
$\underline{\text{Proj}}_{Y'}(\mathcal{A}') =
Y' \times_Y \underline{\text{Proj}}_Y(\mathcal{A})$
du lemme \ref{spaces-divisors-lemma-relative-proj-base-change}.
De plus, $\theta'$ est l'image inverse de $\theta$.
\end{enumerate}
\end{lemma}

\begin{proof}
Omise. Indications :
Observons d'abord que, pour un schéma quasi-compact
$W$ sur $X$, les conditions suivantes sont équivalentes :
\begin{enumerate}
\item $W \to X$ se factorise par $U(\psi)$, et
\item il existe un $d$ tel que
$\psi|_W : f^*\mathcal{A}_d|_W \to \mathcal{L}^{\otimes d}|_W$
soit surjectif.
\end{enumerate}
On obtient ainsi une description de $U(\psi)$ comme sous-foncteur de $X$
sur notre catégorie de base $(\Sch/S)_{fppf}$. Pour de tels $W$ et $d$,
considérons le quadruplet
$(d, W \to Y, \mathcal{L}|_W, \psi^{(d)}|_W)$.
D'après la définition de $\underline{\text{Proj}}_Y(\mathcal{A})$,
on obtient un morphisme $W \to \underline{\text{Proj}}_Y(\mathcal{A})$.
Notre notion d'équivalence des quadruplets montre que
ce morphisme est indépendant du choix de $d$.
Cela définit manifestement une transformation de foncteurs
$r_{\psi, \mathcal{L}} : U(\psi) \to \underline{\text{Proj}}_Y(\mathcal{A})$,
c'est-à-dire un morphisme d'espaces algébriques.
Par construction, ce morphisme vérifie (2).
Comme le morphisme construit dans
Constructions, lemme \ref{constructions-lemma-invertible-map-into-relative-proj}
vérifie la même propriété, (3) est vraie.

\medskip\noindent
Pour construire $\theta$ et vérifier la compatibilité (1) du
lemme, travaillons localement pour la topologie \'etale sur $Y$ et $X$, en raisonnant comme
dans la discussion qui suit la
définition \ref{spaces-divisors-definition-relative-proj}.
\end{proof}








\section{Faisceaux relativement amples}
\label{spaces-divisors-section-relatively-ample}

\noindent
Cette section est l'analogue de
Morphismes, section \ref{morphisms-section-relatively-ample}
pour les espaces algébriques.
Notre définition d'un faisceau inversible relativement ample est
la suivante.

\begin{definition}
\label{spaces-divisors-definition-relatively-ample}
Soit $S$ un schéma.
Soit $f : X \to Y$ un morphisme d'espaces algébriques sur $S$.
Soit $\mathcal{L}$ un $\mathcal{O}_X$-module inversible.
On dit que $\mathcal{L}$ est {\it relativement ample}, ou {\it ample relativement à $f$},
ou {\it ample sur $X/Y$}, ou {\it $f$-ample}, si $f : X \to Y$
est représentable et si, pour tout morphisme $Z \to Y$
où $Z$ est un schéma, l'image inverse $\mathcal{L}_Z$ de $\mathcal{L}$
sur $X_Z = Z \times_Y X$ est ample sur $X_Z/Z$ au sens de
Morphismes, définition \ref{morphisms-definition-relatively-ample}.
\end{definition}

\noindent
Nous ramènerons presque toujours les questions sur les faisceaux inversibles
relativement amples au cas des schémas. Ainsi, dans cette section, nous n'avons
essentiellement que des vérifications de cohérence.

\begin{lemma}
\label{spaces-divisors-lemma-relatively-ample-sanity-check}
Soit $S$ un schéma.
Soit $f : X \to Y$ un morphisme d'espaces algébriques sur $S$.
Soit $\mathcal{L}$ un $\mathcal{O}_X$-module inversible.
Supposons que $Y$ soit un schéma. Les conditions suivantes sont équivalentes :
\begin{enumerate}
\item $\mathcal{L}$ est ample sur $X/Y$ au sens de la
définition \ref{spaces-divisors-definition-relatively-ample}, et
\item $X$ est un schéma et $\mathcal{L}$ est ample sur $X/Y$
au sens de
Morphismes, définition \ref{morphisms-definition-relatively-ample}.
\end{enumerate}
\end{lemma}

\begin{proof}
Cela résulte des définitions et de
Morphismes, lemme \ref{morphisms-lemma-ample-base-change}
(qui affirme que, pour les schémas, la propriété d'être relativement ample
est préservée par changement de base).
\end{proof}

\begin{lemma}
\label{spaces-divisors-lemma-ample-base-change}
Soit $S$ un schéma.
Soit $f : X \to Y$ un morphisme d'espaces algébriques sur $S$.
Soit $\mathcal{L}$ un $\mathcal{O}_X$-module inversible.
Soit $Y' \to Y$ un morphisme d'espaces algébriques sur $S$.
Soit $f' : X' \to Y'$ le changement de base de $f$ et notons
$\mathcal{L}'$ l'image inverse de $\mathcal{L}$ sur $X'$.
Si $\mathcal{L}$ est $f$-ample, alors $\mathcal{L}'$ est $f'$-ample.
\end{lemma}

\begin{proof}
Cela résulte immédiatement de la définition !
(Indication : transitivité du changement de base.)
\end{proof}

\begin{lemma}
\label{spaces-divisors-lemma-relatively-ample-properties}
Soit $S$ un schéma.
Soit $f : X \to Y$ un morphisme d'espaces algébriques sur $S$.
S'il existe un faisceau inversible $f$-ample, alors
$f$ est représentable, quasi-compact et séparé.
\end{lemma}

\begin{proof}
C'est clair d'après les définitions et
Morphismes, lemme \ref{morphisms-lemma-relatively-ample-separated}.
(En cas de doute, voir le principe exposé dans
Espaces algébriques, lemme
\ref{spaces-lemma-representable-transformations-property-implication}.)
\end{proof}

\begin{lemma}
\label{spaces-divisors-lemma-descend-relatively-ample}
Soit $V \to U$ un morphisme \'etale surjectif de schémas affines.
Soit $X$ un espace algébrique sur $U$.
Soit $\mathcal{L}$ un $\mathcal{O}_X$-module inversible.
Soit $Y = V \times_U X$ et soit $\mathcal{N}$
l'image inverse de $\mathcal{L}$ sur $Y$.
Les conditions suivantes sont équivalentes :
\begin{enumerate}
\item $\mathcal{L}$ est ample sur $X/U$, et
\item $\mathcal{N}$ est ample sur $Y/V$.
\end{enumerate}
\end{lemma}

\begin{proof}
L'implication (1) $\Rightarrow$ (2) résulte du
lemme \ref{spaces-divisors-lemma-ample-base-change}.
Supposons (2). Cela implique que $Y \to V$ est
quasi-compact et séparé (lemme \ref{spaces-divisors-lemma-relatively-ample-properties})
et que $Y$ est un schéma. Il s'ensuit que le morphisme $f : X \to U$ est
quasi-compact et séparé
(Morphismes d'espaces, lemmes
\ref{spaces-morphisms-lemma-quasi-compact-local} et
\ref{spaces-morphisms-lemma-separated-local}).
Posons $\mathcal{A} = \bigoplus_{d \geq 0} f_*\mathcal{L}^{\otimes d}$.
C'est un faisceau quasi-cohérent d'algèbres graduées sur $\mathcal{O}_U$
(Morphismes d'espaces, lemme \ref{spaces-morphisms-lemma-pushforward}).
Par adjonction, nous avons un morphisme
$\psi : f^*\mathcal{A} \to \bigoplus_{d \geq 0} \mathcal{L}^{\otimes d}$.
En appliquant le lemme \ref{spaces-divisors-lemma-invertible-map-into-relative-proj},
nous obtenons un sous-espace ouvert $U(\psi) \subset X$ et un morphisme
$$
r_{\mathcal{L}, \psi} : U(\psi) \to \underline{\text{Proj}}_U(\mathcal{A})
$$
Puisque $h : V \to U$ est \'etale, nous avons
$\mathcal{A}|_V = (Y \to V)_*(\bigoplus_{d \geq 0} \mathcal{N}^{\otimes d})$,
voir Propriétés des espaces, lemme
\ref{spaces-properties-lemma-pushforward-etale-base-change-modules}.
Il s'ensuit que l'image inverse $\psi'$ de $\psi$ sur
$Y$ est le morphisme d'adjonction associé à la situation $(Y \to V, \mathcal{N})$
considérée dans Morphismes, lemme \ref{morphisms-lemma-characterize-relatively-ample}, partie (5).
Comme $\mathcal{N}$ est ample sur $Y/V$, le lemme que nous venons de
citer donne $U(\psi') = Y$ et montre que $r_{\mathcal{N}, \psi'}$
est une immersion ouverte.
Puisque le lemme \ref{spaces-divisors-lemma-invertible-map-into-relative-proj}
nous dit que la formation de $r_{\mathcal{L}, \psi}$
commute au changement de base, nous en concluons que
$U(\psi) = X$ et que nous avons un diagramme commutatif
$$
\xymatrix{
Y \ar[r]_-{r'} \ar[d] &
\underline{\text{Proj}}_V(\mathcal{A}|_V) \ar[d] \ar[r] &
V \ar[d] \\
X \ar[r]^-r &
\underline{\text{Proj}}_U(\mathcal{A}) \ar[r] &
U
}
$$
dont les carrés sont cartésiens. Nous en concluons que $r$ est une
immersion ouverte par
Morphismes d'espaces, lemme \ref{spaces-morphisms-lemma-closed-immersion-local}.
Ainsi, $X$ est un schéma. Nous pouvons alors appliquer
Morphismes, lemme \ref{morphisms-lemma-characterize-relatively-ample}, partie (5),
pour conclure que $\mathcal{L}$ est ample sur $X/U$.
\end{proof}

\begin{lemma}
\label{spaces-divisors-lemma-relatively-ample-local}
Soit $S$ un schéma. Soit $f : X \to Y$ un morphisme d'espaces algébriques
sur $S$. Soit $\mathcal{L}$ un $\mathcal{O}_X$-module inversible.
Les conditions suivantes sont équivalentes :
\begin{enumerate}
\item $\mathcal{L}$ est ample sur $X/Y$,
\item pour tout schéma $Z$ et tout morphisme $Z \to Y$,
l'espace algébrique $X_Z = Z \times_Y X$ est un schéma
et l'image inverse $\mathcal{L}_Z$ est ample sur $X_Z/Z$,
\item pour tout schéma affine $Z$ et tout morphisme $Z \to Y$,
l'espace algébrique $X_Z = Z \times_Y X$ est un schéma
et l'image inverse $\mathcal{L}_Z$ est ample sur $X_Z/Z$,
\item il existe un schéma $V$ et un morphisme \'etale surjectif
$V \to Y$ tels que l'espace algébrique $X_V = V \times_Y X$ soit un schéma
et que l'image inverse $\mathcal{L}_V$ soit ample sur $X_V/V$.
\end{enumerate}
\end{lemma}

\begin{proof}
Les assertions (1) et (2) sont équivalentes par définition.
L'implication (2) $\Rightarrow$ (3) est immédiate.
Si (3) est satisfaite et si $Z \to Y$ est comme dans (2), nous voyons
que la représentabilité de $X_Z \to Z$ se vérifie localement sur les ouverts affines de $Z$.
Ainsi, $X_Z$ est un schéma, par exemple par
Propriétés des espaces, lemme \ref{spaces-properties-lemma-subscheme}.
Il en résulte alors que $\mathcal{L}_Z$ est ample sur $X_Z/Z$, car
cela est vrai localement sur $Z$ et nous pouvons utiliser
Morphismes, lemme \ref{morphisms-lemma-characterize-relatively-ample}.
Ainsi, (1), (2) et (3) sont équivalentes. Ces conditions
impliquent clairement (4).

\medskip\noindent
Supposons (4). Soit $Z \to Y$ un morphisme avec $Z$ affine.
Alors $U = V \times_Y Z \to Z$ est un morphisme \'etale surjectif
tel que l'image inverse de $\mathcal{L}_Z$ par $X_U \to X_Z$
soit relativement ample sur $X_U/U$.
Nous pouvons remplacer $U$ par une somme disjointe finie d'ouverts affines dont les images recouvrent encore la base.
Il s'ensuit que $\mathcal{L}_Z$ est ample sur $X_Z/Z$ par le
lemme \ref{spaces-divisors-lemma-descend-relatively-ample}.
Ainsi (4) $\Rightarrow$ (3), ce qui achève la démonstration.
\end{proof}

\begin{lemma}
\label{spaces-divisors-lemma-quasi-affine-O-ample}
Soit $S$ un schéma. Soit $f : X \to Y$ un morphisme d'espaces algébriques
sur $S$. Alors $f$ est quasi-affine si et seulement si $\mathcal{O}_X$
est ample relativement à $f$.
\end{lemma}

\begin{proof}
Cela résulte du cas des schémas, voir
Morphismes, lemme \ref{morphisms-lemma-quasi-affine-O-ample}.
\end{proof}



\section{Amplitude relative et cohomologie}
\label{spaces-divisors-section-ample-and-proper}

\noindent
Cette section contient quelques résultats apparentés à ceux
de Cohomologie des schémas, sections
\ref{coherent-section-applications-formal-functions} et
\ref{coherent-section-ample-cohomology}.

\medskip\noindent
Le lemme suivant n'est qu'un exemple de ce que l'on peut faire.

\begin{lemma}
\label{spaces-divisors-lemma-vanshing-gives-ample}
Soit $R$ un anneau noethérien. Soit $X$ un espace algébrique sur $R$
tel que le morphisme structural $f : X \to \Spec(R)$ soit propre.
Soit $\mathcal{L}$ un $\mathcal{O}_X$-module inversible.
Les conditions suivantes sont équivalentes :
\begin{enumerate}
\item $\mathcal{L}$ est ample sur $X/R$
(définition \ref{spaces-divisors-definition-relatively-ample}),
\item pour tout $\mathcal{O}_X$-module cohérent $\mathcal{F}$,
il existe un $n_0 \geq 0$ tel que
$H^p(X, \mathcal{F} \otimes \mathcal{L}^{\otimes n}) = 0$
pour tous $n \geq n_0$ et $p > 0$.
\end{enumerate}
\end{lemma}

\begin{proof}
L'implication (1) $\Rightarrow$ (2) résulte de
Cohomologie des schémas, lemme \ref{coherent-lemma-coherent-proper-ample},
car l'hypothèse (1) implique que $X$ est un schéma.
L'implication (2) $\Rightarrow$ (1) est le
lemme de Cohomologie des espaces
\ref{spaces-cohomology-lemma-Noetherian-h1-zero-invertible}.
\end{proof}

\begin{lemma}
\label{spaces-divisors-lemma-ample-on-fibre}
Soit $Y$ un schéma noethérien. Soit $X$ un espace algébrique sur $Y$
tel que le morphisme structural $f : X \to Y$ soit propre.
Soit $\mathcal{L}$ un $\mathcal{O}_X$-module inversible.
Soit $\mathcal{F}$ un $\mathcal{O}_X$-module cohérent.
Soit $y \in Y$ un point tel que $X_y$ soit un schéma et que
$\mathcal{L}_y$ soit ample sur $X_y$.
Il existe alors un $d_0$ tel que, pour tout $d \geq d_0$, on ait
$$
R^pf_*(\mathcal{F} \otimes_{\mathcal{O}_X} \mathcal{L}^{\otimes d})_y = 0
\text{ pour }p > 0
$$
et que l'homomorphisme
$$
f_*(\mathcal{F} \otimes_{\mathcal{O}_X} \mathcal{L}^{\otimes d})_y
\longrightarrow
H^0(X_y, \mathcal{F}_y \otimes_{\mathcal{O}_{X_y}} \mathcal{L}_y^{\otimes d})
$$
soit surjectif.
\end{lemma}

\begin{proof}
Observons que $\mathcal{O}_{Y, y}$ est un anneau local noethérien.
Considérons le morphisme canonique
$c : \Spec(\mathcal{O}_{Y, y}) \to Y$, voir
Schémas, équation (\ref{schemes-equation-canonical-morphism}).
Ce morphisme est plat, puisqu'il identifie les anneaux locaux.
Notons provisoirement $f' : X' \to \Spec(\mathcal{O}_{Y, y})$
le morphisme obtenu de $f$ par ce changement de base. Nous avons
$c^*R^pf_*\mathcal{F} = R^pf'_*\mathcal{F}'$ d'après
Cohomologie des espaces, lemme
\ref{spaces-cohomology-lemma-flat-base-change-cohomology}.
De plus, les fibres $X_y$ et $X'_y$ s'identifient.
Nous pouvons donc supposer que $Y = \Spec(A)$ est le spectre d'un
anneau local noethérien $(A, \mathfrak m, \kappa)$ et que $y \in Y$
correspond à $\mathfrak m$. Dans ce cas,
$R^pf_*(\mathcal{F} \otimes_{\mathcal{O}_X} \mathcal{L}^{\otimes d})_y =
H^p(X, \mathcal{F} \otimes_{\mathcal{O}_X} \mathcal{L}^{\otimes d})$
pour tout $p \geq 0$. Notons $f_y : X_y \to \Spec(\kappa)$ la projection.

\medskip\noindent
Posons $B = \text{Gr}_\mathfrak m(A) =
\bigoplus_{n \geq 0} \mathfrak m^n/\mathfrak m^{n + 1}$.
Considérons le faisceau $\mathcal{B} = f_y^*\widetilde{B}$
d'algèbres graduées quasi-cohérentes sur $\mathcal{O}_{X_y}$.
Nous emploierons les notations de Cohomologie des espaces, section
\ref{spaces-cohomology-section-theorem-formal-functions},
en remplaçant $I$ par $\mathfrak m$.
Puisque $X_y$ est le sous-espace fermé de $X$ défini par
$\mathfrak m\mathcal{O}_X$, nous pouvons considérer
$\mathfrak m^n\mathcal{F}/\mathfrak m^{n + 1}\mathcal{F}$
comme un $\mathcal{O}_{X_y}$-module cohérent, voir
Cohomologie des espaces, lemme \ref{spaces-cohomology-lemma-i-star-equivalence}.
Alors
$\bigoplus_{n \geq 0} \mathfrak m^n\mathcal{F}/\mathfrak m^{n + 1}\mathcal{F}$
est un $\mathcal{B}$-module gradué quasi-cohérent de type fini,
car il est engendré en degré zéro sur $\mathcal{B}$
et sa partie de degré zéro est
$\mathcal{F}_y = \mathcal{F}/\mathfrak m \mathcal{F}$,
qui est un $\mathcal{O}_{X_y}$-module cohérent.
Ainsi, d'après Cohomologie des schémas, lemme
\ref{coherent-lemma-graded-finiteness}, partie (2),
il existe un $d_0$ tel que
$$
H^p(X_y, \mathfrak m^n \mathcal{F}/ \mathfrak m^{n + 1}\mathcal{F}
\otimes_{\mathcal{O}_{X_y}} \mathcal{L}_y^{\otimes d}) = 0
$$
pour tous $p > 0$, $d \geq d_0$ et $n \geq 0$. D'après
Cohomologie des espaces, lemme
\ref{spaces-cohomology-lemma-relative-affine-cohomology},
cela revient à dire que
$
H^p(X, \mathfrak m^n \mathcal{F}/ \mathfrak m^{n + 1}\mathcal{F}
\otimes_{\mathcal{O}_X} \mathcal{L}^{\otimes d}) = 0
$
pour tous $p > 0$, $d \geq d_0$ et $n \geq 0$.

\medskip\noindent
Considérons les suites exactes courtes
$$
0 \to \mathfrak m^n\mathcal{F}/\mathfrak m^{n + 1} \mathcal{F}
\to \mathcal{F}/\mathfrak m^{n + 1} \mathcal{F}
\to \mathcal{F}/\mathfrak m^n \mathcal{F} \to 0
$$
de $\mathcal{O}_X$-modules cohérents. La tensorisation par $\mathcal{L}^{\otimes d}$
est un foncteur exact et nous obtenons des suites exactes courtes
$$
0 \to
\mathfrak m^n\mathcal{F}/\mathfrak m^{n + 1} \mathcal{F}
\otimes_{\mathcal{O}_X} \mathcal{L}^{\otimes d}
\to \mathcal{F}/\mathfrak m^{n + 1} \mathcal{F}
\otimes_{\mathcal{O}_X} \mathcal{L}^{\otimes d}
\to \mathcal{F}/\mathfrak m^n \mathcal{F}
\otimes_{\mathcal{O}_X} \mathcal{L}^{\otimes d} \to 0
$$
En utilisant la suite exacte longue de cohomologie et l'annulation ci-dessus,
nous concluons (par récurrence) que
\begin{enumerate}
\item $H^p(X, \mathcal{F}/\mathfrak m^n \mathcal{F}
\otimes_{\mathcal{O}_X} \mathcal{L}^{\otimes d}) = 0$
pour tous $p > 0$, $d \geq d_0$ et $n \geq 0$, et
\item $H^0(X, \mathcal{F}/\mathfrak m^n \mathcal{F}
\otimes_{\mathcal{O}_X} \mathcal{L}^{\otimes d}) \to
H^0(X_y, \mathcal{F}_y \otimes_{\mathcal{O}_{X_y}} \mathcal{L}_y^{\otimes d})$
est surjectif pour tous $d \geq d_0$ et $n \geq 1$.
\end{enumerate}
D'après le théorème des fonctions formelles
(Cohomologie des espaces, théorème
\ref{spaces-cohomology-theorem-formal-functions}),
nous trouvons que le complété $\mathfrak m$-adique de
$H^p(X, \mathcal{F} \otimes_{\mathcal{O}_X} \mathcal{L}^{\otimes d})$
est nul pour tous $d \geq d_0$ et $p > 0$.
Comme $H^p(X, \mathcal{F} \otimes_{\mathcal{O}_X} \mathcal{L}^{\otimes d})$
est un $A$-module de type fini d'après
Cohomologie des espaces, lemme
\ref{spaces-cohomology-lemma-proper-over-affine-cohomology-finite},
il résulte du lemme de Nakayama (Algèbre, lemme \ref{algebra-lemma-NAK})
que $H^p(X, \mathcal{F} \otimes_{\mathcal{O}_X} \mathcal{L}^{\otimes d})$
est nul pour tous $d \geq d_0$ et $p > 0$.
Pour $p = 0$, nous déduisons du
lemme de Cohomologie des espaces
\ref{spaces-cohomology-lemma-ML-cohomology-powers-ideal}, partie (3),
que $H^0(X, \mathcal{F} \otimes_{\mathcal{O}_X} \mathcal{L}^{\otimes d}) \to
H^0(X_y, \mathcal{F}_y \otimes_{\mathcal{O}_{X_y}} \mathcal{L}_y^{\otimes d})$
est surjectif, ce qui donne la dernière assertion du lemme.
\end{proof}

\begin{lemma}
\label{spaces-divisors-lemma-ample-in-neighbourhood}
(Pour une version plus générale, voir
Descente sur les espaces, lemme \ref{spaces-descent-lemma-ample-in-neighbourhood}).
Soit $Y$ un schéma noethérien. Soit $X$ un espace algébrique sur $Y$
tel que le morphisme structural $f : X \to Y$ soit propre.
Soit $\mathcal{L}$ un $\mathcal{O}_X$-module inversible.
Soit $y \in Y$ un point tel que $X_y$ soit un schéma et que
$\mathcal{L}_y$ soit ample sur $X_y$.
Il existe alors un voisinage ouvert $V \subset Y$
de $y$ tel que $\mathcal{L}|_{f^{-1}(V)}$ soit ample sur $f^{-1}(V)/V$
(au sens de la définition \ref{spaces-divisors-definition-relatively-ample}).
\end{lemma}

\begin{proof}
Choisissons $d_0$ comme dans le lemme \ref{spaces-divisors-lemma-ample-on-fibre} pour
$\mathcal{F} = \mathcal{O}_X$. Choisissons $d \geq d_0$
de sorte qu'il existe un $r \geq 0$ et des sections
$s_{y, 0}, \ldots, s_{y, r} \in H^0(X_y, \mathcal{L}_y^{\otimes d})$
qui définissent une immersion fermée
$$
\varphi_y =
\varphi_{\mathcal{L}_y^{\otimes d}, (s_{y, 0}, \ldots, s_{y, r})} :
X_y \to \mathbf{P}^r_{\kappa(y)}.
$$
C'est possible d'après Morphismes, lemme
\ref{morphisms-lemma-finite-type-over-affine-ample-very-ample},
mais nous utilisons aussi
Morphismes, lemme \ref{morphisms-lemma-image-proper-scheme-closed},
pour voir que $\varphi_y$ est une immersion fermée, et
Constructions, section \ref{constructions-section-projective-space},
pour la description des morphismes vers l'espace projectif
au moyen de faisceaux inversibles et de sections.
Par notre choix de $d_0$, après avoir remplacé $Y$ par un voisinage ouvert
de $y$, nous pouvons choisir
$s_0, \ldots, s_r \in H^0(X, \mathcal{L}^{\otimes d})$
qui s'envoient sur $s_{y, 0}, \ldots, s_{y, r}$.
Soit $X_{s_i} \subset X$ le sous-espace ouvert où $s_i$
est un générateur de $\mathcal{L}^{\otimes d}$. Puisque
les $s_{y, i}$ engendrent $\mathcal{L}_y^{\otimes d}$, nous avons
$|X_y| \subset U = \bigcup |X_{s_i}|$. Comme $X \to Y$ est fermé,
il existe un voisinage ouvert $y \in V \subset Y$
tel que $|f|^{-1}(V) \subset U$. Après avoir remplacé $Y$ par $V$, nous pouvons
supposer que les $s_i$ engendrent $\mathcal{L}^{\otimes d}$. Nous
obtenons ainsi un morphisme
$$
\varphi = \varphi_{\mathcal{L}^{\otimes d}, (s_0, \ldots, s_r)} :
X \longrightarrow \mathbf{P}^r_Y
$$
avec $\mathcal{L}^{\otimes d} \cong \varphi^*\mathcal{O}_{\mathbf{P}^r_Y}(1)$
dont le changement de base en $y$ donne $\varphi_y$ (à strictement parler, il faudrait
rédiger une preuve que la construction des morphismes vers l'espace projectif
donnée dans
Constructions, section \ref{constructions-section-projective-space},
vaut aussi pour décrire les morphismes d'espaces algébriques vers l'espace
projectif ; nous omettons les détails).

\medskip\noindent
Nous terminerons la démonstration par une astuce ; la démonstration « correcte »
consiste à montrer directement que $\varphi$ est une immersion
fermée après changement de base à un voisinage ouvert de $y$.
En effet, d'après
Cohomologie des espaces, lemme
\ref{spaces-cohomology-lemma-proper-finite-fibre-finite-in-neighbourhood},
nous voyons que $\varphi$ est fini au-dessus d'un voisinage ouvert
de la fibre $\mathbf{P}^r_{\kappa(y)}$ de $\mathbf{P}^r_Y \to Y$
au-dessus de $y$. Comme $\mathbf{P}^r_Y \to Y$ est fermé, quitte à
rétrécir $Y$, nous pouvons supposer que $\varphi$ est fini.
En particulier, $X$ est un schéma.
Alors $\mathcal{L}^{\otimes d} \cong \varphi^*\mathcal{O}_{\mathbf{P}^r_Y}(1)$
est ample d'après le résultat très général donné dans
Morphismes, lemme \ref{morphisms-lemma-pullback-ample-tensor-relatively-ample}.
\end{proof}








\section{Sous-espaces fermés du Proj relatif}
\label{spaces-divisors-section-closed-in-relative-proj}

\noindent
Quelques lemmes auxiliaires sur les sous-espaces fermés du Proj relatif.
Cette section est l'analogue de
Diviseurs, section \ref{divisors-section-closed-in-relative-proj}.

\begin{lemma}
\label{spaces-divisors-lemma-closed-subscheme-proj}
Soit $S$ un schéma. Soit $X$ un espace algébrique sur $S$.
Soit $\mathcal{A}$ une $\mathcal{O}_X$-algèbre graduée quasi-cohérente. Soit
$\pi : P = \underline{\text{Proj}}_X(\mathcal{A}) \to X$ le Proj relatif
de $\mathcal{A}$. Soit $i : Z \to P$ un sous-espace fermé. Notons
$\mathcal{I} \subset \mathcal{A}$ le noyau du morphisme canonique
$$
\mathcal{A}
\longrightarrow
\bigoplus\nolimits_{d \geq 0} \pi_*\left((i_*\mathcal{O}_Z)(d)\right)
$$
Si $\pi$ est quasi-compact, il existe un isomorphisme
$Z = \underline{\text{Proj}}_X(\mathcal{A}/\mathcal{I})$.
\end{lemma}

\begin{proof}
Le morphisme $\pi$ est séparé d'après le
lemme \ref{spaces-divisors-lemma-relative-proj-separated}.
Comme $\pi$ est quasi-compact, $\pi_*$ transforme les modules quasi-cohérents
en modules quasi-cohérents, voir
Morphismes d'espaces, lemme \ref{spaces-morphisms-lemma-pushforward}.
Ainsi, $\mathcal{I}$ est un $\mathcal{O}_X$-module quasi-cohérent.
En particulier, $\mathcal{B} = \mathcal{A}/\mathcal{I}$ est une
$\mathcal{O}_X$-algèbre graduée quasi-cohérente. Le morphisme de fonctorialité
$Z' = \underline{\text{Proj}}_X(\mathcal{B}) \to
\underline{\text{Proj}}_X(\mathcal{A})$ est défini partout et est
une immersion fermée, voir le lemme
\ref{spaces-divisors-lemma-surjective-graded-rings-map-relative-proj}.
Il suffit donc de montrer que $Z = Z'$ comme sous-espaces fermés de $P$.

\medskip\noindent
Cela étant, la question est locale pour la topologie \'etale sur la base et l'on
se ramène au cas des schémas
(Diviseurs, lemme \ref{divisors-lemma-closed-subscheme-proj})
par localisation \'etale.
\end{proof}

\noindent
Lorsque le sous-espace fermé est localement défini par un nombre fini
d'équations, nous pouvons le définir par un faisceau d'idéaux de type fini de
$\mathcal{A}$.

\begin{lemma}
\label{spaces-divisors-lemma-closed-subscheme-proj-finite}
Soit $S$ un schéma. Soit $X$ un espace algébrique quasi-compact et quasi-séparé
sur $S$.
Soit $\mathcal{A}$ une $\mathcal{O}_X$-algèbre graduée quasi-cohérente. Soit
$\pi : P = \underline{\text{Proj}}_X(\mathcal{A}) \to X$ le Proj relatif
de $\mathcal{A}$. Soit $i : Z \to P$ un sous-schéma fermé.
Si $\pi$ est quasi-compact et si $i$ est de présentation finie, alors il existe
un $d > 0$ et un sous-$\mathcal{O}_X$-module quasi-cohérent de type fini
$\mathcal{F} \subset \mathcal{A}_d$ tel que
$Z = \underline{\text{Proj}}_X(\mathcal{A}/\mathcal{F}\mathcal{A})$.
\end{lemma}

\begin{proof}
Le lecteur peut reprendre les arguments employés dans le cas des schémas. Nous
allons toutefois montrer, par une astuce, que le lemme résulte de ce cas.
Soit $\mathcal{I} \subset \mathcal{A}$ l'idéal gradué quasi-cohérent
qui définit $Z$ dans le lemme \ref{spaces-divisors-lemma-closed-subscheme-proj}.
Choisissons un schéma affine $U$ et un morphisme \'etale surjectif
$U \to X$, voir Propriétés des espaces, lemme
\ref{spaces-properties-lemma-quasi-compact-affine-cover}.
D'après le cas des schémas
(Diviseurs, lemme \ref{divisors-lemma-closed-subscheme-proj-finite}),
il existe un $d > 0$ et un
sous-$\mathcal{O}_U$-module quasi-cohérent de type fini
$\mathcal{F}' \subset \mathcal{I}_d|_U \subset \mathcal{A}_d|_U$,
tel que $Z \times_X U$ soit égal à
$\underline{\text{Proj}}_U(\mathcal{A}|_U/\mathcal{F}'\mathcal{A}|_U)$.
D'après Limites des espaces, lemme
\ref{spaces-limits-lemma-directed-colimit-finite-type},
nous pouvons trouver un sous-module quasi-cohérent de type fini
$\mathcal{F} \subset \mathcal{I}_d$ tel que
$\mathcal{F}' \subset \mathcal{F}|_U$. Posons
$Z' = \underline{\text{Proj}}_X(\mathcal{A}/\mathcal{F}\mathcal{A})$.
Alors $Z' \to P$ est une immersion fermée
(lemme \ref{spaces-divisors-lemma-surjective-generated-degree-1-map-relative-proj})
et $Z \subset Z'$ puisque $\mathcal{F}\mathcal{A} \subset \mathcal{I}$.
D'autre part, $Z' \times_X U \subset Z \times_X U$ d'après notre
choix de $\mathcal{F}$. Ainsi, $Z = Z'$, comme voulu.
\end{proof}

\begin{lemma}
\label{spaces-divisors-lemma-closed-subscheme-proj-finite-type}
Soit $S$ un schéma. Soit $X$ un espace algébrique quasi-compact et quasi-séparé
sur $S$.
Soit $\mathcal{A}$ une $\mathcal{O}_X$-algèbre graduée quasi-cohérente.
Soit $\pi : P = \underline{\text{Proj}}_X(\mathcal{A}) \to X$ le Proj relatif
de $\mathcal{A}$. Soit $i : Z \to X$ un sous-espace fermé.
Soit $U \subset X$ un ouvert. Supposons que
\begin{enumerate}
\item $\pi$ soit quasi-compact,
\item $i$ soit de présentation finie,
\item $|U| \cap |\pi|(|i|(|Z|)) = \emptyset$,
\item $U$ soit quasi-compact,
\item $\mathcal{A}_n$ soit un $\mathcal{O}_X$-module de type fini pour tout $n$.
\end{enumerate}
Alors il existe un $d > 0$ et un
sous-$\mathcal{O}_X$-module quasi-cohérent de type fini $\mathcal{F} \subset \mathcal{A}_d$ tel que (a)
$Z = \underline{\text{Proj}}_X(\mathcal{A}/\mathcal{F}\mathcal{A})$
et (b) le support de $\mathcal{A}_d/\mathcal{F}$ soit disjoint de $U$.
\end{lemma}

\begin{proof}
Nous employons la même astuce que dans la démonstration du
lemme \ref{spaces-divisors-lemma-closed-subscheme-proj-finite}
pour nous ramener au cas des schémas.
Soit $\mathcal{I} \subset \mathcal{A}$ l'idéal gradué quasi-cohérent
qui définit $Z$ dans le lemme \ref{spaces-divisors-lemma-closed-subscheme-proj}.
Choisissons un schéma affine $W$ et un morphisme \'etale surjectif
$W \to X$, voir Propriétés des espaces, lemme
\ref{spaces-properties-lemma-quasi-compact-affine-cover}.
D'après le cas des schémas
(Diviseurs, lemme \ref{divisors-lemma-closed-subscheme-proj-finite-type}),
il existe un $d > 0$ et un
sous-$\mathcal{O}_W$-module quasi-cohérent de type fini
$\mathcal{F}' \subset \mathcal{I}_d|_W \subset \mathcal{A}_d|_W$,
tel que (a) $Z \times_X W$ soit égal à
$\underline{\text{Proj}}_W(\mathcal{A}|_W/\mathcal{F}'\mathcal{A}|_W)$
et que (b) le support de $\mathcal{A}_d|_W/\mathcal{F}'$ soit disjoint de
$U \times_X W$. D'après Limites des espaces, lemme
\ref{spaces-limits-lemma-directed-colimit-finite-type},
nous pouvons trouver un sous-module quasi-cohérent de type fini
$\mathcal{F} \subset \mathcal{I}_d$ tel que
$\mathcal{F}' \subset \mathcal{F}|_W$. Posons
$Z' = \underline{\text{Proj}}_X(\mathcal{A}/\mathcal{F}\mathcal{A})$.
Alors $Z' \to P$ est une immersion fermée
(lemme \ref{spaces-divisors-lemma-surjective-generated-degree-1-map-relative-proj})
et $Z \subset Z'$ puisque $\mathcal{F}\mathcal{A} \subset \mathcal{I}$.
D'autre part, $Z' \times_X W \subset Z \times_X W$ d'après notre
choix de $\mathcal{F}$. Ainsi, $Z = Z'$.
Enfin, le support de $\mathcal{A}_d/\mathcal{F}$ est contenu dans
$X \setminus U$, car $\mathcal{A}_d|_W/\mathcal{F}|_W$ est un quotient
de $\mathcal{A}_d|_W/\mathcal{F}'$, dont le support est contenu dans
$W \setminus U \times_X W$. Le lemme en résulte.
\end{proof}

\begin{lemma}
\label{spaces-divisors-lemma-conormal-sheaf-section-projective-bundle}
Soit $S$ un schéma et soit $X$ un espace algébrique sur $S$.
Soit $\mathcal{E}$ un $\mathcal{O}_X$-module quasi-cohérent.
Il existe une bijection
$$
\left\{
\begin{matrix}
\text{sections }\sigma\text{ du } \\
\text{morphisme } \mathbf{P}(\mathcal{E}) \to X
\end{matrix}
\right\}
\leftrightarrow
\left\{
\begin{matrix}
\text{surjections }\mathcal{E} \to \mathcal{L}\text{ où} \\
\mathcal{L}\text{ est un }\mathcal{O}_X\text{-module inversible}
\end{matrix}
\right\}
$$
Dans ce cas, $\sigma$ est une immersion fermée et il existe un
isomorphisme canonique
$$
\Ker(\mathcal{E} \to \mathcal{L})
\otimes_{\mathcal{O}_X} \mathcal{L}^{\otimes -1}
\longrightarrow
\mathcal{C}_{\sigma(X)/\mathbf{P}(\mathcal{E})}
$$
La bijection et l'isomorphisme sont tous deux compatibles avec le changement de base.
\end{lemma}

\begin{proof}
Puisque les constructions sont compatibles avec le changement de base, il suffit de
vérifier l'assertion localement pour la topologie \'etale sur $X$. Nous pouvons donc supposer que $X$ est
un schéma, et le résultat est alors
Diviseurs, lemme \ref{divisors-lemma-conormal-sheaf-section-projective-bundle}.
\end{proof}






\section{Éclatements}
\label{spaces-divisors-section-blowing-up}

\noindent
L'éclatement est un outil important en géométrie algébrique.

\begin{definition}
\label{spaces-divisors-definition-blow-up}
Soit $S$ un schéma. Soit $X$ un espace algébrique sur $S$.
Soit $\mathcal{I} \subset \mathcal{O}_X$ un faisceau quasi-cohérent
d'idéaux, et soit $Z \subset X$ le sous-espace fermé correspondant
à $\mathcal{I}$
(Morphismes d'espaces, lemme
\ref{spaces-morphisms-lemma-closed-immersion-ideals}).
L'{\it éclatement de $X$ le long de $Z$}, ou l'
{\it éclatement de $X$ suivant le faisceau d'idéaux $\mathcal{I}$} est
le morphisme
$$
b :
\underline{\text{Proj}}_X
\left(\bigoplus\nolimits_{n \geq 0} \mathcal{I}^n\right)
\longrightarrow
X
$$
Le {\it diviseur exceptionnel} de l'éclatement est l'image inverse
$b^{-1}(Z)$. Le sous-espace $Z$ est parfois appelé le {\it centre} de l'éclatement.
\end{definition}

\noindent
Nous verrons plus loin que le diviseur exceptionnel est un diviseur de Cartier
effectif. De plus, l'éclatement se caractérise comme le « plus petit »
espace algébrique au-dessus de $X$ tel que l'image inverse de $Z$ soit un
diviseur de Cartier effectif.

\medskip\noindent
Si $b : X' \to X$ est l'éclatement de $X$ le long de $Z$, on note souvent
$\mathcal{O}_{X'}(n)$ les faisceaux tordus du faisceau structural. Notons que ceux-ci
sont des $\mathcal{O}_{X'}$-modules inversibles et que
$\mathcal{O}_{X'}(n) = \mathcal{O}_{X'}(1)^{\otimes n}$,
car $X'$ est le Proj relatif d'une algèbre graduée quasi-cohérente
sur $\mathcal{O}_X$ qui est engendrée en degré $1$, voir le
lemme \ref{spaces-divisors-lemma-relative-proj-generated-in-degree-1}.

\begin{lemma}
\label{spaces-divisors-lemma-blowing-up-affine}
Soit $S$ un schéma. Soit $X$ un espace algébrique sur $S$.
Soit $\mathcal{I} \subset \mathcal{O}_X$ un
faisceau quasi-cohérent d'idéaux. Soit $U = \Spec(A)$ un schéma affine
\'etale sur $X$, et soit $I \subset A$ l'idéal correspondant à
$\mathcal{I}|_U$. Si $X' \to X$ est l'éclatement de $X$ suivant $\mathcal{I}$,
il existe alors un isomorphisme canonique
$$
U \times_X X' = \text{Proj}(\bigoplus\nolimits_{d \geq 0} I^d)
$$
de schémas au-dessus de $U$, où le membre de droite est
le spectre homogène de l'algèbre de Rees de $I$ dans $A$.
De plus, $U \times_X X'$ possède un recouvrement ouvert affine par les
spectres des algèbres d'éclatement affines $A[\frac{I}{a}]$.
\end{lemma}

\begin{proof}
Notons que la restriction $\mathcal{I}|_U$ est égale à l'image inverse
de $\mathcal{I}$ par le morphisme $U \to X$, voir
Propriétés des espaces, section \ref{spaces-properties-section-modules}.
Le lemme résulte donc de la combinaison du lemme \ref{spaces-divisors-lemma-relative-proj} et de
Diviseurs, lemme \ref{divisors-lemma-blowing-up-affine}.
\end{proof}

\begin{lemma}
\label{spaces-divisors-lemma-flat-base-change-blowing-up}
Soit $S$ un schéma.
Soit $X_1 \to X_2$ un morphisme plat d'espaces algébriques sur $S$.
Soit $Z_2 \subset X_2$ un sous-espace fermé.
Soit $Z_1$ l'image inverse de $Z_2$ dans $X_1$.
Soit $X'_i$ l'éclatement de centre $Z_i$ dans $X_i$. Il existe alors un diagramme
cartésien
$$
\xymatrix{
X_1' \ar[r] \ar[d] & X_2' \ar[d] \\
X_1 \ar[r] & X_2
}
$$
d'espaces algébriques sur $S$.
\end{lemma}

\begin{proof}
Soit $\mathcal{I}_2$ le faisceau d'idéaux de $Z_2$ dans $X_2$.
Notons $g : X_1 \to X_2$ le morphisme donné. Le faisceau d'idéaux
$\mathcal{I}_1$ de $Z_1$ est alors l'image de
$g^*\mathcal{I}_2 \to \mathcal{O}_{X_1}$
(voir Morphismes d'espaces, définition
\ref{spaces-morphisms-definition-inverse-image-closed-subspace}
et la discussion qui suit cette définition).
D'après le lemme \ref{spaces-divisors-lemma-relative-proj-base-change},
$X_1 \times_{X_2} X_2'$ est le Proj relatif de
$\bigoplus_{n \geq 0} g^*\mathcal{I}_2^n$. Comme $g$ est plat, le morphisme
$g^*\mathcal{I}_2^n \to \mathcal{O}_{X_1}$ est injectif et son image est
$\mathcal{I}_1^n$. Ainsi, $X_1 \times_{X_2} X_2' = X_1'$.
\end{proof}

\begin{lemma}
\label{spaces-divisors-lemma-blowing-up-gives-effective-Cartier-divisor}
Soit $S$ un schéma. Soit $X$ un espace algébrique sur $S$.
Soit $Z \subset X$ un sous-espace fermé.
L'éclatement $b : X' \to X$ de centre $Z$ dans $X$
possède les propriétés suivantes :
\begin{enumerate}
\item $b|_{b^{-1}(X \setminus Z)} : b^{-1}(X \setminus Z) \to X \setminus Z$
est un isomorphisme,
\item le diviseur exceptionnel $E = b^{-1}(Z)$ est un diviseur de Cartier effectif
sur $X'$,
\item il existe un isomorphisme canonique
$\mathcal{O}_{X'}(-1) = \mathcal{O}_{X'}(E)$
\end{enumerate}
\end{lemma}

\begin{proof}
Soit $U$ un schéma et soit $U \to X$ un morphisme \'etale surjectif.
Comme la formation de l'éclatement commute au changement de base plat
(lemme \ref{spaces-divisors-lemma-flat-base-change-blowing-up}),
nous pouvons démontrer chacune de ces assertions après changement de base à $U$.
Cela nous ramène au cas des schémas.
Dans ce cas, le résultat est
Diviseurs, lemme
\ref{divisors-lemma-blowing-up-gives-effective-Cartier-divisor}.
\end{proof}

\begin{lemma}[Propriété universelle de l'éclatement]
\label{spaces-divisors-lemma-universal-property-blowing-up}
\begin{slogan}
Éclater un fermé pour en faire un diviseur de Cartier.
\end{slogan}
Soit $S$ un schéma.
Soit $X$ un espace algébrique sur $S$.
Soit $Z \subset X$ un sous-espace fermé.
Soit $\mathcal{C}$ la sous-catégorie pleine de $(\textit{Espaces}/X)$ composée
des $Y \to X$ tels que l'image inverse de $Z$ soit un
diviseur de Cartier effectif sur $Y$. Alors l'éclatement $b : X' \to X$ de centre $Z$ dans $X$
est un objet final de $\mathcal{C}$.
\end{lemma}

\begin{proof}
Le morphisme $b : X' \to X$ est un objet de $\mathcal{C}$ d'après le
lemme \ref{spaces-divisors-lemma-blowing-up-gives-effective-Cartier-divisor}.
Soit $f : Y \to X$ un objet de $\mathcal{C}$. Nous devons montrer qu'il existe
un unique morphisme $Y \to X'$ au-dessus de $X$. Posons $D = f^{-1}(Z)$.
Soit $\mathcal{I} \subset \mathcal{O}_X$ le faisceau d'idéaux de $Z$
et soit $\mathcal{I}_D$ le faisceau d'idéaux de $D$. Alors
$f^*\mathcal{I} \to \mathcal{I}_D$ est une surjection
vers un $\mathcal{O}_Y$-module inversible. Elle se prolonge en un morphisme
$\psi : \bigoplus f^*\mathcal{I}^d \to \bigoplus \mathcal{I}_D^d$
d'algèbres graduées sur $\mathcal{O}_Y$. (Nous observons que
$\mathcal{I}_D^d = \mathcal{I}_D^{\otimes d}$ puisque $D$ est un
diviseur de Cartier effectif.) D'après le
lemme \ref{spaces-divisors-lemma-relative-proj-generated-in-degree-1},
le triplet $(f : Y \to X, \mathcal{I}_D, \psi)$ définit un
morphisme $Y \to X'$ au-dessus de $X$. La restriction
$$
Y \setminus D \longrightarrow X' \setminus b^{-1}(Z) = X \setminus Z
$$
est unique. L'ouvert $Y \setminus D$ est schématiquement dense dans $Y$
d'après le lemme \ref{spaces-divisors-lemma-complement-effective-Cartier-divisor}. 
Ainsi, le morphisme $Y \to X'$ est unique d'après
Morphismes d'espaces, lemme \ref{spaces-morphisms-lemma-equality-of-morphisms}
(de plus, $b$ est séparé d'après le lemme
\ref{spaces-divisors-lemma-relative-proj-separated}).
\end{proof}

\begin{lemma}
\label{spaces-divisors-lemma-blow-up-effective-Cartier-divisor}
Soit $S$ un schéma. Soit $X$ un espace algébrique sur $S$.
Soit $Z \subset X$ un diviseur de Cartier effectif.
L'éclatement de $X$ le long de $Z$ est le morphisme identité de $X$.
\end{lemma}

\begin{proof}
Cela résulte immédiatement de la propriété universelle des éclatements
(lemme \ref{spaces-divisors-lemma-universal-property-blowing-up}).
\end{proof}

\begin{lemma}
\label{spaces-divisors-lemma-blow-up-reduced-space}
Soit $S$ un schéma. Soit $X$ un espace algébrique sur $S$.
Soit $\mathcal{I} \subset \mathcal{O}_X$ un
faisceau quasi-cohérent d'idéaux. Si $X$ est réduit, alors
l'éclatement $X'$ de $X$ suivant $\mathcal{I}$ est réduit.
\end{lemma}

\begin{proof}
Soit $U$ un schéma et soit $U \to X$ un morphisme \'etale surjectif.
Comme la formation de l'éclatement commute au changement de base plat
(lemme \ref{spaces-divisors-lemma-flat-base-change-blowing-up}),
nous pouvons démontrer l'assertion après changement de base à $U$.
Cela nous ramène au cas des schémas.
Dans ce cas, le résultat est
Diviseurs, lemme \ref{divisors-lemma-blow-up-reduced-scheme}.
\end{proof}

\begin{lemma}
\label{spaces-divisors-lemma-blowup-finite-nr-irreducibles}
Soit $S$ un schéma. Soit $X$ un espace algébrique sur $S$. Soit
$b : X' \to X$ l'éclatement de $X$ le long d'un sous-espace fermé. Si
$X$ vérifie les conditions équivalentes de
Morphismes d'espaces, lemme \ref{spaces-morphisms-lemma-prepare-normalization},
alors il en est de même de $X'$.
\end{lemma}

\begin{proof}
Cela résulte immédiatement du lemme cité dans l'énoncé,
de la description locale pour la topologie \'etale des éclatements donnée au
lemme \ref{spaces-divisors-lemma-blowing-up-affine}, et de
Diviseurs, lemme \ref{divisors-lemma-blow-up-and-irreducible-components}.
\end{proof}

\begin{lemma}
\label{spaces-divisors-lemma-blow-up-pullback-effective-Cartier}
Soit $S$ un schéma. Soit $X$ un espace algébrique sur $S$.
Soit $b : X' \to X$ un éclatement de $X$ le long d'un sous-espace fermé.
Pour tout diviseur de Cartier effectif $D$ sur $X$, l'image inverse
$b^{-1}D$ est définie (voir la définition
\ref{spaces-divisors-definition-pullback-effective-Cartier-divisor}).
\end{lemma}

\begin{proof}
D'après les lemmes \ref{spaces-divisors-lemma-blowing-up-affine} et
\ref{spaces-divisors-lemma-characterize-effective-Cartier-divisor},
cela se ramène au fait algébrique suivant :
soient $A$ un anneau, $I \subset A$ un idéal, $a \in I$, et $x \in A$
un non-diviseur de zéro. L'image de $x$ dans $A[\frac{I}{a}]$ est alors un
non-diviseur de zéro. En effet, supposons que $x (y/a^n) = 0$ dans $A[\frac{I}{a}]$.
Alors $a^mxy = 0$ dans $A$ pour un certain $m$. Donc $a^my = 0$ puisque $x$ est un
non-diviseur de zéro. Ainsi $y/a^n$ est nul dans $A[\frac{I}{a}]$, comme voulu.
\end{proof}

\begin{lemma}
\label{spaces-divisors-lemma-blowing-up-two-ideals}
Soit $S$ un schéma. Soit $X$ un espace algébrique sur $S$.
Soient $\mathcal{I} \subset \mathcal{O}_X$ et $\mathcal{J}$ des
faisceaux quasi-cohérents d'idéaux. Soit $b : X' \to X$ l'éclatement
de $X$ suivant $\mathcal{I}$. Soit $b' : X'' \to X'$ l'éclatement de
$X'$ suivant $b^{-1}\mathcal{J} \mathcal{O}_{X'}$. Alors $X'' \to X$
est canoniquement isomorphe à l'éclatement de $X$ suivant $\mathcal{I}\mathcal{J}$.
\end{lemma}

\begin{proof}
Soit $E \subset X'$ le diviseur exceptionnel de $b$, qui est un diviseur de
Cartier effectif d'après le
lemme \ref{spaces-divisors-lemma-blowing-up-gives-effective-Cartier-divisor}.
Alors $(b')^{-1}E$ est un diviseur de Cartier effectif sur $X''$ d'après le
lemme \ref{spaces-divisors-lemma-blow-up-pullback-effective-Cartier}.
Soit $E' \subset X''$ le diviseur exceptionnel de $b'$ (lui aussi un diviseur de
Cartier effectif). Considérons le diviseur de Cartier effectif
$E'' = E' + (b')^{-1}E$. Par construction, l'idéal de $E''$ est
$(b \circ b')^{-1}\mathcal{I} (b \circ b')^{-1}\mathcal{J} \mathcal{O}_{X''}$.
D'après le lemme \ref{spaces-divisors-lemma-universal-property-blowing-up},
il existe donc un morphisme canonique de $X''$ vers l'éclatement $c : Y \to X$
de $X$ le long de $\mathcal{I}\mathcal{J}$. Inversement, puisque $\mathcal{I}\mathcal{J}$
a pour image inverse un idéal inversible, on voit que
$c^{-1}\mathcal{I}\mathcal{O}_Y$ définit
un diviseur de Cartier effectif, voir le
lemme \ref{spaces-divisors-lemma-sum-closed-subschemes-effective-Cartier}.
Il existe donc un morphisme $c' : Y \to X'$ au-dessus de $X$ par le
lemme \ref{spaces-divisors-lemma-universal-property-blowing-up}.
Alors $(c')^{-1}b^{-1}\mathcal{J}\mathcal{O}_Y = c^{-1}\mathcal{J}\mathcal{O}_Y$,
qui définit lui aussi un diviseur de Cartier effectif. Il existe donc un morphisme
$c'' : Y \to X''$ au-dessus de $X'$. Nous omettons de vérifier que ce
morphisme est l'inverse du morphisme $X'' \to Y$ construit précédemment.
\end{proof}

\begin{lemma}
\label{spaces-divisors-lemma-blowing-up-projective}
Soit $S$ un schéma. Soit $X$ un espace algébrique sur $S$.
Soit $\mathcal{I} \subset \mathcal{O}_X$ un faisceau quasi-cohérent
d'idéaux. Soit $b : X' \to X$ l'éclatement de $X$
suivant le faisceau d'idéaux $\mathcal{I}$. Si $\mathcal{I}$ est de type fini, alors
$b : X' \to X$ est un morphisme propre.
\end{lemma}

\begin{proof}
Soit $U$ un schéma et soit $U \to X$ un morphisme \'etale surjectif.
Comme la formation de l'éclatement commute au changement de base plat
(lemme \ref{spaces-divisors-lemma-flat-base-change-blowing-up}),
nous pouvons démontrer l'assertion après changement de base à $U$
(voir Morphismes d'espaces, lemme
\ref{spaces-morphisms-lemma-proper-local}).
Cela nous ramène au cas des schémas.
Dans ce cas, le morphisme $b$ est projectif d'après
Diviseurs, lemme \ref{divisors-lemma-blowing-up-projective},
donc propre d'après
Morphismes, lemme \ref{morphisms-lemma-locally-projective-proper}.
\end{proof}

\begin{lemma}
\label{spaces-divisors-lemma-composition-finite-type-blowups}
Soit $S$ un schéma et soit $X$ un espace algébrique sur $S$.
Supposons que $X$ soit quasi-compact et quasi-séparé.
Soit $Z \subset X$ un sous-espace fermé de présentation finie.
Soit $b : X' \to X$ l'éclatement de centre $Z$.
Soit $Z' \subset X'$ un sous-espace fermé de présentation finie.
Soit $X'' \to X'$ l'éclatement de centre $Z'$.
Il existe un sous-espace fermé $Y \subset X$ de présentation finie
tel que
\begin{enumerate}
\item $|Y| = |Z| \cup |b|(|Z'|)$, et
\item la composée $X'' \to X$ soit isomorphe à l'éclatement
de $X$ le long de $Y$.
\end{enumerate}
\end{lemma}

\begin{proof}
La condition que $Z \to X$ soit de présentation finie signifie que
$Z$ est défini par un faisceau d'idéaux quasi-cohérent de type fini
$\mathcal{I} \subset \mathcal{O}_X$, voir
Morphismes d'espaces, lemme
\ref{spaces-morphisms-lemma-closed-immersion-finite-presentation}.
Écrivons $\mathcal{A} = \bigoplus_{n \geq 0} \mathcal{I}^n$, de sorte que
$X' = \underline{\text{Proj}}(\mathcal{A})$.
Notons que $X \setminus Z$ est un sous-espace ouvert quasi-compact de $X$ d'après
Limites des espaces, lemme
\ref{spaces-limits-lemma-quasi-coherent-finite-type-ideals}.
Puisque $b^{-1}(X \setminus Z) \to X \setminus Z$ est un isomorphisme
(lemme \ref{spaces-divisors-lemma-blowing-up-gives-effective-Cartier-divisor}), le même
résultat montre que
$b^{-1}(X \setminus Z) \setminus Z'$ est un sous-espace ouvert quasi-compact de $X'$.
Ainsi, $U = X \setminus (Z \cup b(Z'))$ est un sous-espace ouvert quasi-compact de $X$.
D'après le lemme \ref{spaces-divisors-lemma-closed-subscheme-proj-finite-type},
il existe un $d > 0$ et un
sous-$\mathcal{O}_X$-module de type fini $\mathcal{F} \subset \mathcal{I}^d$ tel
que $Z' = \underline{\text{Proj}}(\mathcal{A}/\mathcal{F}\mathcal{A})$
et que le support de $\mathcal{I}^d/\mathcal{F}$ soit contenu
dans $X \setminus U$.

\medskip\noindent
Comme $\mathcal{F} \subset \mathcal{I}^d$ est un sous-$\mathcal{O}_X$-module,
nous pouvons considérer $\mathcal{F} \subset \mathcal{I}^d \subset \mathcal{O}_X$
comme un faisceau d'idéaux quasi-cohérent de type fini sur $X$. Notons-le
$\mathcal{J} \subset \mathcal{O}_X$ pour éviter toute confusion. Puisque
$\mathcal{I}^d / \mathcal{J}$ et $\mathcal{O}/\mathcal{I}^d$ ont
leur support dans $|X| \setminus |U|$, on voit que $|V(\mathcal{J})|$ est contenu
dans $|X| \setminus |U|$. Réciproquement, comme $\mathcal{J} \subset \mathcal{I}^d$,
on voit que $|Z| \subset |V(\mathcal{J})|$. Sur
$X \setminus Z \cong X' \setminus b^{-1}(Z)$, le faisceau d'idéaux
$\mathcal{J}$ définit $Z'$ (voir la formule affichée ci-dessous). Ainsi,
$|V(\mathcal{J})|$ est égal à $|Z| \cup |b|(|Z'|)$. Par suite,
$|V(\mathcal{I}\mathcal{J})| = |Z| \cup |b|(|Z'|)$. De plus,
$\mathcal{I}\mathcal{J}$ est un idéal de type fini, comme produit de deux tels idéaux.
Nous affirmons que $X'' \to X$ est isomorphe à l'éclatement de $X$ suivant
$\mathcal{I}\mathcal{J}$, ce qui achève la démonstration du lemme en posant
$Y = V(\mathcal{I}\mathcal{J})$.

\medskip\noindent
Rappelons d'abord que l'éclatement de $X$ suivant $\mathcal{I}\mathcal{J}$
est le même que l'éclatement de $X'$ suivant $b^{-1}\mathcal{J} \mathcal{O}_{X'}$,
voir le lemme \ref{spaces-divisors-lemma-blowing-up-two-ideals}.
Il suffit donc de montrer que l'éclatement de $X'$ suivant
$b^{-1}\mathcal{J} \mathcal{O}_{X'}$ coïncide avec l'éclatement de $X'$
le long de $Z'$. Nous allons montrer que
$$
b^{-1}\mathcal{J} \mathcal{O}_{X'} = \mathcal{I}_E^d \mathcal{I}_{Z'}
$$
comme faisceaux d'idéaux sur $X''$. Cela donnera le résultat, puisque
$\mathcal{I}_E^d$ définit le diviseur de Cartier effectif $dE$
et que nous pouvons utiliser les lemmes \ref{spaces-divisors-lemma-blow-up-effective-Cartier-divisor} et
\ref{spaces-divisors-lemma-blowing-up-two-ideals}.

\medskip\noindent
Pour établir l'égalité affichée des idéaux, nous pouvons travailler localement.
Avec les notations $A$, $I$, $a \in I$ du lemme \ref{spaces-divisors-lemma-blowing-up-affine},
on voit que $\mathcal{F}$ correspond à un sous-$R$-module $M \subset I^d$
qui s'envoie isomorphiquement sur un idéal $J \subset R$. La condition
$Z' = \underline{\text{Proj}}(\mathcal{A}/\mathcal{F}\mathcal{A})$
signifie que $Z' \cap \Spec(A[\frac{I}{a}])$ est défini par l'idéal
engendré par les éléments $m/a^d$, où $m \in M$. Supposons que $m \in M$
corresponde à la fonction $f \in J$. Alors, dans l'algèbre d'éclatement affine
$A' = A[\frac{I}{a}]$, on a $f = (a^dm)/a^d = a^d (m/a^d)$.
L'égalité en résulte.
\end{proof}









\section{Transformée stricte}
\label{spaces-divisors-section-strict-transform}

\noindent
Cette section est l'analogue de
Diviseurs, section \ref{divisors-section-strict-transform}.
Soient $S$ un schéma, $B$ un espace algébrique sur $S$, et
$Z \subset B$ un sous-espace fermé.
Soit $b : B' \to B$ l'éclatement de $B$ le long de $Z$ et désignons par $E \subset B'$
le diviseur exceptionnel $E = b^{-1}Z$. Dans la suite, nous considérerons
souvent un espace algébrique $X$ sur $B$ et formerons le diagramme cartésien
$$
\xymatrix{
\text{pr}_{B'}^{-1}E \ar[r] \ar[d] &
X \times_B B' \ar[r]_-{\text{pr}_X} \ar[d]_{\text{pr}_{B'}} &
X \ar[d]^f \\
E \ar[r] & B' \ar[r] & B
}
$$
Puisque $E$ est un diviseur de Cartier effectif
(lemme \ref{spaces-divisors-lemma-blowing-up-gives-effective-Cartier-divisor}),
on voit que $\text{pr}_{B'}^{-1}E \subset X \times_B B'$
est localement principal
(lemme \ref{spaces-divisors-lemma-pullback-locally-principal}).
Ainsi, le morphisme d'inclusion du complémentaire de
$\text{pr}_{B'}^{-1}E$ dans $X \times_B B'$
est affine et, en particulier, quasi-compact
(lemme \ref{spaces-divisors-lemma-complement-locally-principal-closed-subscheme}).
Par conséquent, pour un $\mathcal{O}_{X \times_B B'}$-module quasi-cohérent
$\mathcal{G}$, le sous-faisceau des sections à support dans $|\text{pr}_{B'}^{-1}E|$
est un sous-module quasi-cohérent ; voir
Limites d'espaces, définition
\ref{spaces-limits-definition-subsheaf-sections-supported-on-closed}.
Si $\mathcal{G}$ est un faisceau quasi-cohérent d'algèbres, par exemple
$\mathcal{G} = \mathcal{O}_{X \times_B B'}$, ce sous-faisceau est alors un idéal
de $\mathcal{G}$.

\begin{definition}
\label{spaces-divisors-definition-strict-transform}
Reprenons $Z \subset B$ et $f : X \to B$ comme ci-dessus.
\begin{enumerate}
\item Étant donné un $\mathcal{O}_X$-module quasi-cohérent $\mathcal{F}$,
la {\it transformée stricte} de $\mathcal{F}$ relativement à l'éclatement
de $B$ le long de $Z$ est le quotient $\mathcal{F}'$ de $\text{pr}_X^*\mathcal{F}$
par le sous-module des sections à support dans $|\text{pr}_{B'}^{-1}E|$.
\item La {\it transformée stricte} de $X$ est le sous-espace fermé
$X' \subset X \times_B B'$ défini par l'idéal quasi-cohérent des
sections de $\mathcal{O}_{X \times_B B'}$ à support dans
$|\text{pr}_{B'}^{-1}E|$.
\end{enumerate}
\end{definition}

\noindent
Notons que la transformée stricte par un éclatement dépend du
sous-espace fermé qui sert de centre à cet éclatement
(et non du seul morphisme $B' \to B$).

\begin{lemma}[Localisation étale et transformée stricte]
\label{spaces-divisors-lemma-strict-transform-local}
Dans la situation de la définition \ref{spaces-divisors-definition-strict-transform}.
Soit
$$
\xymatrix{
U \ar[r] \ar[d] & X \ar[d] \\
V \ar[r] & B
}
$$
un diagramme commutatif de morphismes, où $U$ et $V$ sont des schémas et où
les flèches horizontales sont étales. Soit $V' \to V$ l'éclatement de $V$
le long de $Z \times_B V$. Alors
\begin{enumerate}
\item $V' = V \times_B B'$ et les morphismes
$V' \to B'$ et $U \times_V V' \to X \times_B B'$ sont étales,
\item la transformée stricte $U'$ de $U$ relativement à $V' \to V$
est égale à $X' \times_X U$, où $X'$ est la transformée stricte de $X$
relativement à $B' \to B$, et
\item pour un $\mathcal{O}_X$-module quasi-cohérent $\mathcal{F}$, la
restriction de la transformée stricte $\mathcal{F}'$ à
$U \times_V V'$ est la transformée stricte de $\mathcal{F}|_U$ relativement
à $V' \to V$.
\end{enumerate}
\end{lemma}

\begin{proof}
L'assertion (1) résulte de la commutation de l'éclatement au changement de base
plat (lemme \ref{spaces-divisors-lemma-flat-base-change-blowing-up}), de la platitude des
morphismes étales et de leur stabilité par changement de base. L'assertion (3)
résulte alors de ce que la formation du faisceau des sections à support dans
un fermé commute à l'image inverse
par les morphismes étales ; voir Limites d'espaces, lemme
\ref{spaces-limits-lemma-sections-supported-on-closed-subset}.
L'assertion (2) résulte de (3), appliquée à $\mathcal{F} = \mathcal{O}_X$.
\end{proof}

\begin{lemma}
\label{spaces-divisors-lemma-strict-transform}
Dans la situation de la définition \ref{spaces-divisors-definition-strict-transform}.
\begin{enumerate}
\item La transformée stricte $X'$ de $X$ est l'éclatement de $X$ le long du
sous-espace fermé $f^{-1}Z$ de $X$.
\item Pour un $\mathcal{O}_X$-module quasi-cohérent $\mathcal{F}$, la
transformée stricte $\mathcal{F}'$ est canoniquement isomorphe à
l'image directe par $X' \to X \times_B B'$ de la transformée stricte de
$\mathcal{F}$ relativement à l'éclatement $X' \to X$.
\end{enumerate}
\end{lemma}

\begin{proof}
Soit $X'' \to X$ l'éclatement de $X$ le long de $f^{-1}Z$. Par la propriété
universelle de l'éclatement (lemme \ref{spaces-divisors-lemma-universal-property-blowing-up}),
il existe un diagramme commutatif
$$
\xymatrix{
X'' \ar[r] \ar[d] & X \ar[d] \\
B' \ar[r] & B
}
$$
d'où un morphisme $i : X'' \to X \times_B B'$. La première assertion
du lemme affirme que $i$ est une immersion fermée d'image $X'$.
La seconde affirme que $\mathcal{F}' = i_*\mathcal{F}''$,
où $\mathcal{F}''$ est la transformée stricte de $\mathcal{F}$
relativement à l'éclatement $X'' \to X$. Ces assertions se vérifient
localement pour la topologie étale sur $X$ ; on se ramène donc au cas des schémas
(Diviseurs, lemme \ref{divisors-lemma-strict-transform}).
Certains détails sont omis.
\end{proof}

\begin{lemma}
\label{spaces-divisors-lemma-strict-transform-flat}
Dans la situation de la définition \ref{spaces-divisors-definition-strict-transform}.
\begin{enumerate}
\item Si $X$ est plat sur $B$ en tout point situé au-dessus de $Z$, alors
la transformée stricte de $X$ coïncide avec le changement de base $X \times_B B'$.
\item Soit $\mathcal{F}$ un $\mathcal{O}_X$-module quasi-cohérent.
Si $\mathcal{F}$ est plat sur $B$ en tout point situé au-dessus de $Z$, alors
la transformée stricte $\mathcal{F}'$ de $\mathcal{F}$ est égale à
l'image inverse $\text{pr}_X^*\mathcal{F}$.
\end{enumerate}
\end{lemma}

\begin{proof}
Omis. Indication : cela résulte du cas des schémas
(Diviseurs, lemme \ref{divisors-lemma-strict-transform-flat})
par localisation étale
(lemme \ref{spaces-divisors-lemma-strict-transform-local}).
\end{proof}

\begin{lemma}
\label{spaces-divisors-lemma-strict-transform-affine}
Soit $S$ un schéma. Soit $B$ un espace algébrique sur $S$.
Soit $Z \subset B$ un sous-espace fermé.
Soit $b : B' \to B$ l'éclatement de centre $Z$ dans $B$. Soit
$g : X \to Y$ un morphisme affine d'espaces algébriques sur $B$.
Soit $\mathcal{F}$ un faisceau quasi-cohérent sur $X$.
Soit $g' : X \times_B B' \to Y \times_B B'$ le changement de base
de $g$. Soit $\mathcal{F}'$ la transformée stricte de $\mathcal{F}$
relativement à $b$. Alors $g'_*\mathcal{F}'$ est la transformée stricte
de $g_*\mathcal{F}$.
\end{lemma}

\begin{proof}
Omis. Indication : cela résulte du cas des schémas
(Diviseurs, lemme \ref{divisors-lemma-strict-transform-affine})
par localisation étale (lemme \ref{spaces-divisors-lemma-strict-transform-local}).
\end{proof}

\begin{lemma}
\label{spaces-divisors-lemma-strict-transform-different-centers}
Soit $S$ un schéma. Soit $B$ un espace algébrique sur $S$.
Soit $Z \subset B$ un sous-espace fermé.
Soit $D \subset B$ un diviseur de Cartier effectif.
Soit $Z' \subset B$ le sous-espace fermé défini par le produit
des faisceaux d'idéaux de $Z$ et de $D$.
Soit $B' \to B$ l'éclatement de $B$ le long de $Z$.
\begin{enumerate}
\item L'éclatement de $B$ le long de $Z'$ est isomorphe au morphisme $B' \to B$.
\item Soient $f : X \to B$ un morphisme d'espaces algébriques et $\mathcal{F}$
un $\mathcal{O}_X$-module quasi-cohérent. Si le sous-faisceau de $\mathcal{F}$ des
sections à support dans $|f^{-1}D|$ est nul, alors la
transformée stricte de $\mathcal{F}$ relativement à l'éclatement
le long de $Z$ coïncide avec la transformée stricte de $\mathcal{F}$ relativement
à l'éclatement de $B$ le long de $Z'$.
\end{enumerate}
\end{lemma}

\begin{proof}
Omis. Indication : cela résulte du cas des schémas
(Diviseurs, lemme \ref{divisors-lemma-strict-transform-different-centers})
par localisation étale (lemme \ref{spaces-divisors-lemma-strict-transform-local}).
\end{proof}

\begin{lemma}
\label{spaces-divisors-lemma-strict-transform-composition-blowups}
Soit $S$ un schéma. Soit $B$ un espace algébrique sur $S$.
Soit $Z \subset B$ un sous-espace fermé.
Soit $b : B' \to B$ l'éclatement de centre $Z$.
Soit $Z' \subset B'$ un sous-espace fermé.
Soit $B'' \to B'$ l'éclatement de centre $Z'$.
Soit $Y \subset B$ un sous-schéma fermé tel que
$|Y| = |Z| \cup |b|(|Z'|)$ et que la composée $B'' \to B$
soit isomorphe à l'éclatement de $B$ le long de $Y$.
Dans cette situation, pour tout schéma $X$ sur $B$ et tout
$\mathcal{F} \in \QCoh(\mathcal{O}_X)$, on a
\begin{enumerate}
\item la transformée stricte de $\mathcal{F}$ relativement à l'éclatement
de $B$ le long de $Y$ est égale à la transformée stricte, relativement à
l'éclatement $B'' \to B'$ le long de $Z'$, de la transformée stricte de $\mathcal{F}$
relativement à l'éclatement $B' \to B$ de $B$ le long de $Z$, et
\item la transformée stricte de $X$ relativement à l'éclatement
de $B$ le long de $Y$ est égale à la transformée stricte, relativement à
l'éclatement $B'' \to B'$ le long de $Z'$, de la transformée stricte de $X$
relativement à l'éclatement $B' \to B$ de $B$ le long de $Z$.
\end{enumerate}
\end{lemma}

\begin{proof}
Omis. Indication : cela résulte du cas des schémas
(Diviseurs, lemme \ref{divisors-lemma-strict-transform-composition-blowups})
par localisation étale (lemme \ref{spaces-divisors-lemma-strict-transform-local}).
\end{proof}

\begin{lemma}
\label{spaces-divisors-lemma-strict-transform-universally-injective}
Dans la situation de la définition \ref{spaces-divisors-definition-strict-transform}.
Supposons que
$$
0 \to \mathcal{F}_1 \to \mathcal{F}_2 \to \mathcal{F}_3 \to 0
$$
soit une suite exacte de faisceaux quasi-cohérents sur $X$ qui reste
exacte après tout changement de base $T \to B$. Alors les transformées strictes
$\mathcal{F}_i'$ relativement à tout éclatement $B' \to B$
forment elles aussi la suite exacte
$0 \to \mathcal{F}'_1 \to \mathcal{F}'_2 \to \mathcal{F}'_3 \to 0$.
\end{lemma}

\begin{proof}
Omis. Indication : cela résulte du cas des schémas
(Diviseurs, lemme \ref{divisors-lemma-strict-transform-universally-injective})
par localisation étale (lemme \ref{spaces-divisors-lemma-strict-transform-local}).
\end{proof}

\begin{lemma}
\label{spaces-divisors-lemma-strict-transform-blowup-fitting-ideal}
Soit $S$ un schéma. Soit $B$ un espace algébrique sur $S$.
Soit $\mathcal{F}$ un $\mathcal{O}_B$-module quasi-cohérent de type fini.
Soit $Z_k \subset S$ le sous-schéma fermé défini par
$\text{Fit}_k(\mathcal{F})$ ; voir la section \ref{spaces-divisors-section-fitting-ideals}.
Soit $B' \to B$ l'éclatement de $B$ le long de $Z_k$ et soit
$\mathcal{F}'$ la transformée stricte de $\mathcal{F}$.
Alors $\mathcal{F}'$ peut être engendré localement par $\leq k$
sections.
\end{lemma}

\begin{proof}
Omis. Cela résulte du cas des schémas
(Diviseurs, lemme \ref{divisors-lemma-strict-transform-blowup-fitting-ideal})
par localisation étale (lemme \ref{spaces-divisors-lemma-strict-transform-local}).
\end{proof}

\begin{lemma}
\label{spaces-divisors-lemma-strict-transform-blowup-fitting-ideal-locally-free}
Soit $S$ un schéma. Soit $B$ un espace algébrique sur $S$.
Soit $\mathcal{F}$ un $\mathcal{O}_B$-module quasi-cohérent de type fini.
Soit $Z_k \subset S$ le sous-schéma fermé défini par
$\text{Fit}_k(\mathcal{F})$ ; voir la section \ref{spaces-divisors-section-fitting-ideals}.
Supposons que $\mathcal{F}$ soit localement libre de rang $k$ sur $B \setminus Z_k$.
Soit $B' \to B$ l'éclatement de $B$ le long de $Z_k$ et soit
$\mathcal{F}'$ la transformée stricte de $\mathcal{F}$.
Alors $\mathcal{F}'$ est localement libre de rang $k$.
\end{lemma}

\begin{proof}
Omis. Cela résulte du cas des schémas
(Diviseurs, lemme
\ref{divisors-lemma-strict-transform-blowup-fitting-ideal-locally-free})
par localisation étale (lemme \ref{spaces-divisors-lemma-strict-transform-local}).
\end{proof}












\section{Éclatements admissibles}
\label{spaces-divisors-section-admissible-blowups}

\noindent
Afin de mieux maîtriser les éclatements considérés, nous introduisons la
terminologie standard suivante.

\begin{definition}
\label{spaces-divisors-definition-admissible-blowup}
Soit $S$ un schéma. Soit $X$ un espace algébrique sur $S$.
Soit $U \subset X$ un sous-espace ouvert. Un morphisme
$X' \to X$ est appelé {\it éclatement $U$-admissible} s'il existe une
immersion fermée $Z \to X$ de présentation finie, avec $Z$ disjoint de
$U$, telle que $X'$ soit isomorphe à l'éclatement de $X$ le long de $Z$.
\end{definition}

\noindent
Rappelons que $Z \to X$ est de présentation finie si et seulement si le
faisceau d'idéaux $\mathcal{I}_Z \subset \mathcal{O}_X$ est de type fini ; voir
Morphismes d'espaces, lemme
\ref{spaces-morphisms-lemma-closed-immersion-finite-presentation}.
En particulier, un éclatement $U$-admissible est un morphisme propre ; voir le
lemme \ref{spaces-divisors-lemma-blowing-up-projective}.
Notons que plusieurs centres peuvent donner lieu au même morphisme.
La condition exige donc seulement l'existence d'un centre disjoint de
$U$ qui produise $X'$.
Enfin, comme le morphisme $b : X' \to X$ est un isomorphisme au-dessus de $U$ (voir le
lemme \ref{spaces-divisors-lemma-blowing-up-gives-effective-Cartier-divisor}), nous identifierons souvent,
par abus de notation, $U$ à un sous-espace ouvert de $X'$.

\begin{lemma}
\label{spaces-divisors-lemma-composition-admissible-blowups}
Soit $S$ un schéma.
Soit $X$ un espace algébrique quasi-compact et quasi-séparé sur $S$.
Soit $U \subset X$ un sous-espace ouvert quasi-compact.
Soit $b : X' \to X$ un éclatement $U$-admissible.
Soit $X'' \to X'$ un éclatement $U$-admissible.
Alors la composée $X'' \to X$ est un éclatement $U$-admissible.
\end{lemma}

\begin{proof}
Cela résulte immédiatement du
lemme plus précis \ref{spaces-divisors-lemma-composition-finite-type-blowups}.
\end{proof}

\begin{lemma}
\label{spaces-divisors-lemma-extend-admissible-blowups}
Soit $S$ un schéma.
Soit $X$ un espace algébrique quasi-compact et quasi-séparé.
Soient $U, V \subset X$ des sous-espaces ouverts quasi-compacts.
Soit $b : V' \to V$ un éclatement $U \cap V$-admissible.
Alors il existe un éclatement $U$-admissible $X' \to X$
dont la restriction à $V$ est $V'$.
\end{lemma}

\begin{proof}
Soit $\mathcal{I} \subset \mathcal{O}_V$ le faisceau quasi-cohérent d'idéaux
de type fini tel que $V(\mathcal{I})$ soit
disjoint de $U \cap V$ et que $V'$ soit isomorphe à
l'éclatement de $V$ le long de $\mathcal{I}$. Soit
$\mathcal{I}' \subset \mathcal{O}_{U \cup V}$ le faisceau quasi-cohérent
d'idéaux dont la restriction à $U$ est $\mathcal{O}_U$ et
dont la restriction à $V$ est $\mathcal{I}$.
Par Limites d'espaces, lemme \ref{spaces-limits-lemma-extend},
il existe un faisceau quasi-cohérent d'idéaux de type fini
$\mathcal{J} \subset \mathcal{O}_X$ dont la restriction à $U \cup V$ est
$\mathcal{I}'$. Le lemme en résulte.
\end{proof}

\begin{lemma}
\label{spaces-divisors-lemma-dominate-admissible-blowups}
Soit $S$ un schéma.
Soit $X$ un espace algébrique quasi-compact et quasi-séparé sur $S$.
Soit $U \subset X$ un sous-espace ouvert quasi-compact.
Soient $b_i : X_i \to X$, $i = 1, \ldots, n$, des éclatements $U$-admissibles.
Il existe un éclatement $U$-admissible $b : X' \to X$ tel que
(a) $b$ se factorise sous la forme $X' \to X_i \to X$ pour $i = 1, \ldots, n$, et
(b) chacun des morphismes $X' \to X_i$ soit un éclatement $U$-admissible.
\end{lemma}

\begin{proof}
Soit $\mathcal{I}_i \subset \mathcal{O}_X$ le faisceau quasi-cohérent d'idéaux
de type fini tel que $V(\mathcal{I}_i)$ soit
disjoint de $U$ et que $X_i$ soit isomorphe à
l'éclatement de $X$ le long de $\mathcal{I}_i$. Posons
$\mathcal{I} = \mathcal{I}_1 \cdot \ldots \cdot \mathcal{I}_n$
et soit $X'$ l'éclatement de $X$ le long de $\mathcal{I}$. Alors
$X' \to X$ se factorise par $b_i$ d'après le lemme \ref{spaces-divisors-lemma-blowing-up-two-ideals}.
\end{proof}

\begin{lemma}
\label{spaces-divisors-lemma-separate-disjoint-opens-by-blowing-up}
Soit $S$ un schéma.
Soit $X$ un espace algébrique quasi-compact et quasi-séparé sur $S$.
Soient $U, V$ des sous-espaces ouverts quasi-compacts disjoints de $X$.
Il existe alors un éclatement $U \cup V$-admissible $b : X' \to X$
tel que $X'$ soit une réunion disjointe de sous-espaces ouverts
$X' = X'_1 \amalg X'_2$, avec $b^{-1}(U) \subset X'_1$ et
$b^{-1}(V) \subset X'_2$.
\end{lemma}

\begin{proof}
Choisissons des faisceaux quasi-cohérents d'idéaux de type fini $\mathcal{I}$
et, respectivement, $\mathcal{J}$ tels que $X \setminus U = V(\mathcal{I})$
et, respectivement, $X \setminus V = V(\mathcal{J})$ ; voir
Limites d'espaces, lemme
\ref{spaces-limits-lemma-quasi-coherent-finite-type-ideals}.
Alors $|V(\mathcal{I}\mathcal{J})| = |X|$. Par suite,
$\mathcal{I}\mathcal{J}$ est un faisceau d'idéaux localement nilpotent.
Comme $\mathcal{I}$ et $\mathcal{J}$ sont de type fini et $X$
est quasi-compact, il existe un $n > 0$ tel que
$\mathcal{I}^n \mathcal{J}^n = 0$. On peut donc remplacer, et l'on remplace, $\mathcal{I}$
par $\mathcal{I}^n$ et $\mathcal{J}$ par $\mathcal{J}^n$. On a alors
$\mathcal{I} \mathcal{J} = 0$. Soit $b : X' \to X$ l'éclatement
le long de $\mathcal{I} + \mathcal{J}$. Cet éclatement est $U \cup V$-admissible
car $|V(\mathcal{I} + \mathcal{J})| = |X| \setminus |U| \cup |V|$.
Nous allons montrer que $X'$ est une réunion disjointe de sous-espaces ouverts
$X' = X'_1 \amalg X'_2$, comme dans l'énoncé du lemme.

\medskip\noindent
Comme $|V(\mathcal{I} + \mathcal{J})|$ est le complémentaire de
$|U \cup V|$, on en déduit que $V \cup U$ est schématiquement
dense dans $X'$ ; voir les
lemmes \ref{spaces-divisors-lemma-blowing-up-gives-effective-Cartier-divisor} et
\ref{spaces-divisors-lemma-complement-effective-Cartier-divisor}.
Ainsi, si une telle décomposition $X' = X'_1 \amalg X'_2$
en sous-espaces à la fois ouverts et fermés existe, $X'_1$ est l'adhérence
schématique de $U$ dans $X'$, et, de même, $X'_2$ est
l'adhérence schématique de $V$ dans $X'$. Comme $U \to X'$
et $V \to X'$ sont quasi-compacts, la formation des adhérences
schématiques commute à la localisation étale (Morphismes d'espaces,
lemme \ref{spaces-morphisms-lemma-quasi-compact-scheme-theoretic-image}).
Pour vérifier l'existence de $X'_1$ et $X'_2$, on peut donc travailler localement
pour la topologie étale sur $X$. On se ramène ainsi au cas des schémas, traité
dans la démonstration de Diviseurs, lemme
\ref{divisors-lemma-separate-disjoint-opens-by-blowing-up}.
\end{proof}
























% OK, start here.
%

\chapter{Espaces algébriques sur les corps}



\phantomsection
\label{spaces-over-fields-section-phantom}


\section{Introduction}
\label{spaces-over-fields-section-introduction}

\noindent
Ce chapitre est l'analogue du chapitre sur les variétés dans le cadre
des espaces algébriques. Une référence pour les espaces algébriques est
\cite{Kn}.


\section{Conventions}
\label{spaces-over-fields-section-conventions}

\noindent
Nous faisons l'hypothèse permanente que tous les schémas appartiennent à
un gros site fppf $\Sch_{fppf}$. De plus, tout anneau $A$ considéré
est tel que $\Spec(A)$ est (isomorphe à) un
objet de ce grand site.

\medskip\noindent
Soit $S$ un schéma et soit $X$ un espace algébrique sur $S$.
Dans ce chapitre et le suivant, nous écrirons $X \times_S X$
pour le produit de $X$ par lui-même (dans la catégorie des espaces algébriques
sur $S$), au lieu de $X \times X$.



\section{Morphismes génériquement finis}
\label{spaces-over-fields-section-generically-finite}

\noindent
Cette section prolonge la discussion de
Espaces décents, section \ref{decent-spaces-section-generically-finite},
ainsi que son analogue, pour les morphismes d'espaces algébriques, dans
Variétés, section \ref{varieties-section-generically-finite}.

\begin{lemma}
\label{spaces-over-fields-lemma-quasi-finite-in-codim-1}
Soit $S$ un schéma. Soit $f : X \to Y$ un morphisme d'espaces algébriques
sur $S$. Supposons que $f$ soit localement de type fini et que $Y$ soit
localement noethérien. Soit $y \in |Y|$ un point de codimension $\leq 1$ de $Y$.
Soit $X^0 \subset |X|$ l'ensemble des points de codimension $0$ de $X$.
Supposons de plus que l'une des conditions suivantes soit satisfaite
\begin{enumerate}
\item pour tout $x \in X^0$, le degré de transcendance de $x/f(x)$ est $0$,
\item pour tout $x \in X^0$ tel que $f(x) \leadsto y$,
le degré de transcendance de $x/f(x)$ est $0$,
\item $f$ est quasi-fini en tout $x \in X^0$,
\item $f$ est quasi-fini en tout point d'une partie dense de $|X|$,
\item en ajouter d'autres ici.
\end{enumerate}
Alors $f$ est quasi-fini en tout point de $X$ situé au-dessus de $y$.
\end{lemma}

\begin{proof}
Nous voulons ramener la preuve au cas des schémas. Pour ce faire, nous
choisissons un diagramme commutatif
$$
\xymatrix{
U \ar[r] \ar[d]_g & X \ar[d]^f \\
V \ar[r] & Y
}
$$
où $U$ et $V$ sont des schémas et où les flèches horizontales sont \'etales
et surjectives. Choisissons $v \in V$ d'image $y$. Observons que
$V$ est localement noethérien et que $\dim(\mathcal{O}_{V, v}) \leq 1$
(voir Propriétés des espaces, Définitions
\ref{spaces-properties-definition-dimension-local-ring} et
Remarque \ref{spaces-properties-remark-list-properties-local-etale-topology}).
La fibre $U_v$ de $U \to V$ au-dessus de $v$ s'envoie surjectivement sur
$f^{-1}(\{y\}) \subset |X|$. L'image réciproque de $X^0$ dans $U$
est exactement l'ensemble des
points génériques des composantes irréductibles de $U$ (Propriétés des espaces, Lemme
\ref{spaces-properties-lemma-codimension-0-points}).
Si $\eta \in U$ est un tel point, d'image $x \in X^0$, alors
le degré de transcendance de $x / f(x)$ est le degré de transcendance de
$\kappa(\eta)$ sur $\kappa(g(\eta))$
(Morphismes d'espaces, Définition
\ref{spaces-morphisms-definition-dimension-fibre}).
Observons que $U \to V$ est quasi-fini en $u \in U$ si et seulement si
$f$ est quasi-fini en l'image de $u$ dans $X$.

\medskip\noindent
Cas (1). Ici le cas (1) de
Variétés, Lemme \ref{varieties-lemma-quasi-finite-in-codim-1} s'applique
et nous concluons que $U \to V$ est quasi-fini en tout point de $U_v$.
Donc $f$ est quasi-fini en tout point situé au-dessus de $y$.

\medskip\noindent
Cas (2). Soit $u \in U$ un point générique d'une composante irréductible
dont l'image dans $V$ se spécialise en $v$. Alors l'image $x \in X^0$ de
$u$ vérifie $f(x) \leadsto y$. Par conséquent, nous voyons que
le cas (2) de
Variétés, Lemme \ref{varieties-lemma-quasi-finite-in-codim-1} s'applique
et nous concluons comme précédemment.

\medskip\noindent
Le cas (3) découle du cas (3) de
Variétés, Lemme \ref{varieties-lemma-quasi-finite-in-codim-1}.

\medskip\noindent
Dans le cas (4), puisque l'application $|U| \to |X|$ est ouverte, nous voyons que
l'ensemble des points où $U \to V$ est quasi-fini est également dense.
Ainsi le cas (4) de
Variétés, Lemme \ref{varieties-lemma-quasi-finite-in-codim-1} s'applique.
\end{proof}

\begin{lemma}
\label{spaces-over-fields-lemma-finite-in-codim-1}
Soit $S$ un schéma. Soit $f : X \to Y$ un morphisme d'espaces algébriques
sur $S$. Supposons que $f$ soit propre et que $Y$ soit localement
noethérien. Soit $y \in Y$ un point de codimension $\leq 1$ de $Y$.
Soit $X^0 \subset |X|$ l'ensemble des points de codimension $0$ de $X$.
Supposons de plus que l'une des
conditions suivantes soit satisfaite
\begin{enumerate}
\item pour tout $x \in X^0$, le degré de transcendance de $x/f(x)$ est $0$,
\item pour tout $x \in X^0$ tel que $f(x) \leadsto y$, le degré de transcendance
de $x/f(x)$ est $0$,
\item $f$ est quasi-fini en tout $x \in X^0$,
\item $f$ est quasi-fini en tout point d'une partie dense de $|X|$,
\item en ajouter d'autres ici.
\end{enumerate}
Alors il existe un sous-espace ouvert $Y' \subset Y$ contenant $y$ tel que
$Y' \times_Y X \to Y'$ soit fini.
\end{lemma}

\begin{proof}
D'après le lemme \ref{spaces-over-fields-lemma-quasi-finite-in-codim-1}, le morphisme $f$ est
quasi-fini en tout point situé au-dessus de $y$. Soit $\overline{y} : \Spec(k) \to Y$
un point géométrique situé au-dessus de $y$. Alors $|X_{\overline{y}}|$ est un espace
discret (Espaces décents, Lemme
\ref{decent-spaces-lemma-conditions-on-fibre-and-qf}).
L'espace $X_{\overline{y}}$ étant quasi-compact puisque $f$ est propre, nous concluons
que $|X_{\overline{y}}|$ est fini.
Nous pouvons donc appliquer Cohomologie des espaces, Lemme
\ref{spaces-cohomology-lemma-proper-finite-fibre-finite-in-neighbourhood}
pour conclure.
\end{proof}

\begin{lemma}
\label{spaces-over-fields-lemma-modification-normal-iso-over-codimension-1}
Soit $S$ un schéma. Soit $X$ un espace algébrique noethérien sur $S$.
Soit $f : Y \to X$ un morphisme propre birationnel d'espaces algébriques
tel que $Y$ soit réduit.
Soit $U \subset X$ l'ouvert maximal sur lequel $f$ est un isomorphisme.
Alors $U$ contient
\begin{enumerate}
\item tout point de codimension $0$ de $X$,
\item tout $x \in |X|$ de codimension $1$ de $X$ tel que l'anneau local
de $X$ en $x$ soit normal (Propriétés des espaces, Remarque
\ref{spaces-properties-remark-list-properties-local-ring-local-etale-topology}),
et
\item tout $x \in |X|$ tel que la fibre de $|Y| \to |X|$ au-dessus de $x$ soit
finie et tel que l'anneau local de $X$ en $x$ soit normal.
\end{enumerate}
\end{lemma}

\begin{proof}
La partie (1) résulte du lemme
\ref{decent-spaces-lemma-birational-isomorphism-over-dense-open} du chapitre Espaces décents
(et du fait que les espaces algébriques noethériens $X$ et $Y$
sont quasi-séparés et donc décents).
La partie (2) résulte de la partie (3) et du lemme \ref{spaces-over-fields-lemma-finite-in-codim-1}
(et du fait que les morphismes finis ont des fibres finies).
Soit $x \in |X|$ satisfaisant (3). D'après
Cohomologie des espaces, Lemme
\ref{spaces-cohomology-lemma-proper-finite-fibre-finite-in-neighbourhood}
(qui s'applique d'après le lemme
\ref{decent-spaces-lemma-conditions-on-fibre-and-qf} du chapitre Espaces décents),
nous pouvons supposer que $f$ est fini. Choisissons un schéma affine $X'$ et
un morphisme \'etale $X' \to X$ et un point $x' \in X$ s'envoyant sur $x$.
Il suffit de montrer qu'il existe un voisinage ouvert $U'$ de $x' \in X'$
tel que $Y \times_X X' \to X'$ soit un isomorphisme au-dessus de $U'$
(en effet, $U$ contient alors l'image de $U'$ dans $X$, voir Espaces, Lemme
\ref{spaces-lemma-descent-representable-transformations-property}).
Alors $Y \times_X X' \to X$ est un morphisme fini et birationnel
(Espaces décents, Lemme \ref{decent-spaces-lemma-birational-etale-localization}).
Comme un morphisme fini est affine, nous nous ramenons
au cas d'un morphisme fini birationnel de schémas affines noethériens
$Y \to X$ et d'un point $x \in X$ tel que $\mathcal{O}_{X, x}$ soit un
anneau intègre normal. Ceci est traité dans Variétés, Lemme
\ref{varieties-lemma-modification-normal-iso-over-codimension-1}.
\end{proof}






\section{Espaces algébriques intègres}
\label{spaces-over-fields-section-integral-spaces}

\noindent
Nous n'avons pas encore défini la notion d'espace algébrique intègre. La
difficulté vient de ce que la propriété d'être intègre n'est pas locale pour la topologie \'etale des schémas.
Nous pourrions prendre comme définition le fait que $X$ soit réduit et que $|X|$ soit irréductible,
caractérisation donnée dans Propriétés, Lemme \ref{properties-lemma-characterize-integral}.
Dans ce cas, l'espace algébrique
décrit dans Espaces, Exemple \ref{spaces-example-infinite-product}
serait intègre, ce qui ne paraît pas souhaitable.
Pour éviter ce type de pathologie, nous supposerons de plus que $X$ est
un espace algébrique décent, bien qu'une hypothèse plus faible puisse peut-être suffire.

\begin{definition}
\label{spaces-over-fields-definition-integral-algebraic-space}
Soit $S$ un schéma. On dit qu'un espace algébrique $X$ sur $S$ est
{\it intègre} s'il est réduit et décent, et si $|X|$ est irréductible.
\end{definition}

\noindent
Dans ce cas, l'espace topologique irréductible $|X|$ est sobre
(Espaces décents, Proposition \ref{decent-spaces-proposition-reasonable-sober}).
Il admet donc un unique point générique $x$. En fait, le lemme
\ref{decent-spaces-lemma-finitely-many-irreducible-components} du chapitre
Espaces décents caractérise les espaces algébriques décents
ayant un nombre fini de
composantes irréductibles. Il en résulte qu'un
espace algébrique $X$ est intègre s'il est
réduit, s'il admet un sous-schéma ouvert dense et irréductible $X'$,
de point générique $x'$, et si le morphisme $x' \to X$ est quasi-compact.

\begin{lemma}
\label{spaces-over-fields-lemma-integral-algebraic-space-rational-functions}
Soit $S$ un schéma. Soit $X$ un espace algébrique intègre sur $S$.
Soit $\eta \in |X|$ le point générique de $X$.
On a des identifications canoniques
$$
R(X) = \mathcal{O}_{X, \eta}^h = \kappa(\eta)
$$
où $R(X)$ désigne l'anneau des fonctions rationnelles défini dans
Morphismes d'espaces, Définition
\ref{spaces-morphisms-definition-ring-of-rational-functions},
$\kappa(\eta)$ est le corps résiduel défini dans
Espaces décents, Définition \ref{decent-spaces-definition-residue-field},
et $\mathcal{O}_{X, \eta}^h$ est l'anneau local hensélien défini dans
Espaces décents, Définition
\ref{decent-spaces-definition-elemenary-etale-neighbourhood}.
En particulier, ces anneaux sont des corps.
\end{lemma}

\begin{proof}
Comme $X$ est un schéma dans un voisinage ouvert de $\eta$ (voir ci-dessus),
l'assertion découle immédiatement du résultat correspondant pour
les schémas ; voir Morphismes, Lemme
\ref{morphisms-lemma-integral-scheme-rational-functions}.
Nous utilisons aussi le fait que l'hensélisation d'un corps est ce corps lui-même,
ainsi que la compatibilité de nos définitions de ces objets
pour les espaces algébriques avec celles qui valent pour les schémas.
Nous omettons les détails.
\end{proof}

\noindent
Cela conduit à la définition suivante.

\begin{definition}
\label{spaces-over-fields-definition-function-field}
Soit $S$ un schéma. Soit $X$ un espace algébrique intègre sur $S$.
Le {\it corps des fonctions}, ou le {\it corps des fonctions rationnelles},
de $X$ est le corps $R(X)$ du
lemme \ref{spaces-over-fields-lemma-integral-algebraic-space-rational-functions}.
\end{definition}

\noindent
Nous noterons parfois ce corps $k(X)$ au lieu de $R(X)$.

\begin{lemma}
\label{spaces-over-fields-lemma-integral-sections}
Soit $S$ un schéma. Soit $X$ un espace algébrique intègre sur $S$.
Alors $\Gamma(X, \mathcal{O}_X)$ est un anneau intègre.
\end{lemma}

\begin{proof}
Posons $R = \Gamma(X, \mathcal{O}_X)$. Si $f, g \in R$ sont non nuls et
$fg = 0$, alors $X = V(f) \cup V(g)$, où $V(f)$ désigne le sous-espace fermé
de $X$ défini par $f$. Comme $X$ est irréductible, on a soit
$V(f) = X$, soit $V(g) = X$. Il s'ensuit que soit $f = 0$, soit $g = 0$, d'après
Propriétés des espaces, Lemme \ref{spaces-properties-lemma-reduced-space}.
\end{proof}

\noindent
Voici un lemme sur les espaces algébriques intègres et normaux.

\begin{lemma}
\label{spaces-over-fields-lemma-normal-integral-cover-by-affines}
Soit $S$ un schéma. Soit $X$ un espace algébrique intègre et normal sur $S$.
Pour tout $x \in |X|$, il existe un schéma affine intègre et normal $U$
et un morphisme \'etale $U \to X$ dont l'image contient $x$.
\end{lemma}

\begin{proof}
Choisissons un schéma affine $U$ et un morphisme \'etale $U \to X$ dont
l'image contient $x$. Soient $u_i$, $i \in I$, les points génériques des
composantes irréductibles de $U$. Alors chaque $u_i$ a pour image le point générique de $X$
(Espaces décents, Lemme \ref{decent-spaces-lemma-decent-generic-points}). D'après
notre définition d'un espace décent
(Espaces décents, Définition \ref{decent-spaces-definition-very-reasonable}),
on voit que $I$ est fini. Ainsi, $U = \Spec(A)$, où $A$ est un anneau normal
ayant un nombre fini d'idéaux premiers minimaux.
Par conséquent, $A = \prod_{i \in I} A_i$ est un produit d'anneaux intègres normaux d'après
Algèbre, Lemme \ref{algebra-lemma-characterize-reduced-ring-normal}.
Alors $U = \coprod U_i$ avec $U_i = \Spec(A_i)$ et $x$ appartient à l'image de
$U_i \to X$ pour un certain $i$. Cela démontre le lemme.
\end{proof}

\begin{lemma}
\label{spaces-over-fields-lemma-normal-integral-sections}
Soit $S$ un schéma. Soit $X$ un espace algébrique intègre et normal sur $S$.
Alors $\Gamma(X, \mathcal{O}_X)$ est un anneau intègre normal.
\end{lemma}

\begin{proof}
Posons $R = \Gamma(X, \mathcal{O}_X)$. Alors $R$ est un anneau intègre d'après
Lemme \ref{spaces-over-fields-lemma-integral-sections}.
Soit $f = a/b$ un élément du corps des fractions de $R$
qui est entier sur $R$.
Pour tout $U \to X$ \'etale, où $U$ est un schéma, il existe au plus un
$f_U \in \Gamma(U, \mathcal{O}_U)$ tel que $b|_U f_U = a|_U$.
En effet, $U$ est réduit et les points génériques de $U$ s'envoient sur
le point générique de $X$, ce qui implique que $b|_U$ est un
non diviseur de zéro.
Pour tout $x \in |X|$, choisissons $U \to X$ comme dans
Lemme \ref{spaces-over-fields-lemma-normal-integral-cover-by-affines}.
Il existe alors un unique $f_U \in \Gamma(U, \mathcal{O}_U)$
tel que $b|_U f_U = a|_U$, puisque
$\Gamma(U, \mathcal{O}_U)$ est un anneau intègre normal d'après
Propriétés, Lemme \ref{properties-lemma-normal-integral-sections}.
Par l'unicité ci-dessus, ces $f_U$
se recollent en une section globale $f$ du faisceau
structural, c'est-à-dire en un élément de $R$.
\end{proof}

\begin{lemma}
\label{spaces-over-fields-lemma-decent-irreducible-closed}
Soit $S$ un schéma. Soit $X$ un espace algébrique décent sur $S$.
On a des bijections canoniques entre les ensembles suivants :
\begin{enumerate}
\item l'ensemble des points de $X$, c'est-à-dire $|X|$,
\item l'ensemble des parties fermées irréductibles de $|X|$,
\item l'ensemble des sous-espaces fermés intègres de $X$.
\end{enumerate}
La bijection de (1) sur (2) envoie $x$ sur $\overline{\{x\}}$.
La bijection de (3) sur (2) envoie $Z$ sur $|Z|$.
\end{lemma}

\begin{proof}
L'application ainsi définie est une bijection entre (1) et (2), puisque $|X|$ est
sobre d'après
Espaces décents, Proposition \ref{decent-spaces-proposition-reasonable-sober}.
Soit $T \subset |X|$ une partie fermée irréductible. Il existe un
unique sous-espace fermé réduit $Z \subset X$ tel que
$|Z| = T$ : $Z$ est le sous-espace porté par $T$ et muni de la structure réduite
induite ; voir Propriétés des espaces, Définition
\ref{spaces-properties-definition-reduced-induced-space}.
C'est un espace algébrique intègre car il est décent,
réduit et irréductible.
\end{proof}






\section{Morphismes entre espaces algébriques intègres}
\label{spaces-over-fields-section-morphisms-between-integral-spaces}

\noindent
Le lemme suivant caractérise les morphismes dominants de degré fini
entre espaces algébriques intègres.

\begin{lemma}
\label{spaces-over-fields-lemma-finite-degree}
Soit $S$ un schéma. Soient $X$, $Y$ des espaces algébriques intègres sur $S$.
Soient $x \in |X|$ et $y \in |Y|$ les points génériques. Soit $f : X \to Y$
un morphisme localement de type fini. Supposons que $f$ soit dominant
(Morphismes d'espaces, Définition \ref{spaces-morphisms-definition-dominant}).
Les conditions suivantes sont équivalentes :
\begin{enumerate}
\item le degré de transcendance de $x/y$ est $0$,
\item l'extension $\kappa(x)/\kappa(y)$ (voir la démonstration) est finie,
\item il existe des ouverts affines non vides $U \subset X$ et $V \subset Y$
tels que $f(U) \subset V$ et que $f|_U : U \to V$ soit fini,
\item $f$ est quasi-fini en $x$, et
\item $x$ est le seul point de $|X|$ ayant pour image $y$.
\end{enumerate}
Si $f$ est séparé ou si $f$ est quasi-compact, alors ces conditions sont
également équivalentes à
\begin{enumerate}
\item[(6)] il existe un ouvert affine non vide $V \subset Y$ tel
que $f^{-1}(V) \to V$ soit fini.
\end{enumerate}
\end{lemma}

\begin{proof}
Un argument de topologie élémentaire montre que $f(x) = y$, puisque $f$ est dominant.
Soit $Y' \subset Y$ le lieu schématique de $Y$ et soit
$X' \subset f^{-1}(Y')$ le lieu schématique de $f^{-1}(Y')$.
D'après la discussion ci-dessus, en utilisant
Espaces décents, Proposition \ref{decent-spaces-proposition-reasonable-sober} et
Théorème \ref{decent-spaces-theorem-decent-open-dense-scheme},
on voit que $x \in |X'|$ et $y \in |Y'|$.
Alors $f|_{X'} : X' \to Y'$ est un morphisme de schémas intègres
qui est localement de type fini. On en déduit que les conditions (1), (2) et (3)
sont équivalentes d'après Morphismes, Lemme \ref{morphisms-lemma-finite-degree}.

\medskip\noindent
La condition (4) implique la condition (1) d'après
Morphismes d'espaces, Lemme \ref{spaces-morphisms-lemma-compare-tr-deg}
appliqué à $X \to Y \to Y$.
Réciproquement, la condition (3) implique la condition (4), car tout morphisme
fini est quasi-fini et $x \in U$, puisque $x$
est le point générique. Ainsi (1) -- (4) sont équivalentes.

\medskip\noindent
Supposons satisfaites les conditions équivalentes (1) -- (4). Supposons que
$x' \mapsto y$. Alors $x \leadsto x'$ est une spécialisation dans la
fibre de $|X| \to |Y|$ au-dessus de $y$. Si $x' \not = x$, alors $f$ n'est pas
quasi-fini en $x$ d'après Espaces décents, Lemme
\ref{decent-spaces-lemma-conditions-on-point-in-fibre-and-qf}.
Donc $x = x'$ et (5) est vérifiée. Réciproquement, si (5) est vérifiée, alors
(5) est vérifiée pour le morphisme de schémas $X' \to Y'$ (voir ci-dessus)
et nous pouvons utiliser
Morphismes, Lemme \ref{morphisms-lemma-finite-degree}
pour voir que (1) est vérifiée.

\medskip\noindent
Remarquons que (6) implique les conditions équivalentes (1) -- (5)
sans hypothèse supplémentaire sur $f$. Pour achever la démonstration,
il faut montrer que les conditions équivalentes (1) -- (5) impliquent (6).
Cela résulte du lemme
\ref{decent-spaces-lemma-finite-over-dense-open} du chapitre Espaces décents.
\end{proof}

\begin{definition}
\label{spaces-over-fields-definition-degree}
Soit $S$ un schéma.
Soient $X$ et $Y$ des espaces algébriques intègres sur $S$.
Soit $f : X \to Y$ un morphisme localement de type fini et dominant.
Supposons que l'une des conditions équivalentes (1) -- (5) du
lemme \ref{spaces-over-fields-lemma-finite-degree} soit satisfaite. Soient $x \in |X|$ et $y \in |Y|$
les points génériques. Alors l'entier positif
$$
\deg(X/Y) = [\kappa(x) : \kappa(y)]
$$
est appelé le {\it degré de $X$ sur $Y$}.
\end{definition}

\begin{lemma}
\label{spaces-over-fields-lemma-degree-composition}
Soit $S$ un schéma.
Soient $X$, $Y$, $Z$ des espaces algébriques intègres sur $S$.
Soient $f : X \to Y$ et $g : Y \to Z$ des morphismes dominants et localement
de type fini. Supposons que l'une des conditions équivalentes
(1) -- (5) du lemme \ref{spaces-over-fields-lemma-finite-degree} soit satisfaite pour $f$ et $g$. Alors
$$
\deg(X/Z) = \deg(X/Y) \deg(Y/Z).
$$
\end{lemma}

\begin{proof}
Cela découle de la multiplicativité des degrés dans les tours
d'extensions finies de corps ; voir
Corps, Lemme \ref{fields-lemma-multiplicativity-degrees}.
\end{proof}










\section{Diviseurs de Weil}
\label{spaces-over-fields-section-Weil-divisors}

\noindent
Cette section est l'analogue de la section
\ref{divisors-section-Weil-divisors} de Diviseurs.

\medskip\noindent
Nous allons introduire les diviseurs de Weil et l'équivalence rationnelle des
diviseurs de Weil pour les espaces algébriques intègres localement noethériens.
Comme nous ne supposons pas que nos espaces algébriques sont quasi-compacts, nous devons
être un peu prudents lors de la définition des diviseurs de Weil. Nous devons permettre
des sommes infinies de diviseurs premiers, car une fonction rationnelle peut avoir,
par exemple, une infinité de pôles. Dans le cas quasi-compact, nos diviseurs
de Weil sont des sommes finies comme d'habitude. Voici un lemme fondamental que nous
utiliserons souvent pour montrer que des familles de sous-espaces fermés sont localement finies.

\begin{lemma}
\label{spaces-over-fields-lemma-components-locally-finite}
Soit $S$ un schéma et soit $X$ un espace algébrique localement noethérien
sur $S$. Si $T \subset |X|$ est une partie fermée,
alors la famille des composantes irréductibles de $T$ est localement finie.
\end{lemma}

\begin{proof}
L'espace topologique $|X|$ est localement noethérien
(Propriétés des espaces, lemme \ref{spaces-properties-lemma-Noetherian-topology}).
Un espace topologique noethérien possède un nombre fini de
composantes irréductibles et tout sous-espace d'un espace noethérien est noethérien
(Topologie, lemme \ref{topology-lemma-Noetherian}).
Le lemme découle donc de la définition d'une famille localement finie
(Topologie, définition \ref{topology-definition-locally-finite}).
\end{proof}

\noindent
Soit $S$ un schéma. Soit $X$ un espace algébrique décent sur $S$.
Soit $Z$ un sous-espace fermé intègre de $X$ et soit
$\xi \in |Z|$ son point générique. Alors la codimension de
$|Z|$ dans $|X|$ est égale à la dimension de l'anneau local
de $X$ en $\xi$, d'après
le lemme d'Espaces décents \ref{decent-spaces-lemma-codimension-local-ring}.
Rappelons que nous indiquons également cela en disant que
{\it $\xi$ est un point de codimension $1$ sur $X$}, voir la
définition de Propriétés des espaces
\ref{spaces-properties-definition-dimension-local-ring}.

\begin{definition}
\label{spaces-over-fields-definition-Weil-divisor}
Soit $S$ un schéma.
Soit $X$ un espace algébrique intègre localement noethérien sur $S$.
\begin{enumerate}
\item Un {\it diviseur premier} est un sous-espace fermé intègre $Z \subset X$
de codimension $1$, c'est-à-dire que le point générique de $|Z|$ est un point
de codimension $1$ sur $X$.
\item Un {\it diviseur de Weil} est une somme formelle $D = \sum n_Z Z$ où
la somme porte sur les diviseurs premiers de $X$ et la famille
$\{|Z| : n_Z \not = 0\}$ est localement finie dans $|X|$
(Topologie, définition \ref{topology-definition-locally-finite}).
\end{enumerate}
Le groupe de tous les diviseurs de Weil sur $X$ est noté $\text{Div}(X)$.
\end{definition}

\noindent
Notre prochaine tâche consiste à définir le diviseur de Weil associé à une fonction
rationnelle. Pour cela, nous devons définir l'ordre d'annulation d'une
fonction rationnelle sur un espace algébrique intègre localement noethérien $X$
le long d'un diviseur premier $Z$. Soit $\xi \in |Z|$ le point générique.
Nous rencontrons ici le problème que l'anneau local $\mathcal{O}_{X, \xi}$
n'existe pas et que l'anneau local hensélien $\mathcal{O}_{X, \xi}^h$
peut ne pas être un anneau intègre, voir l'exemple \ref{spaces-over-fields-example-not-a-domain}.
Pour contourner cela, nous utilisons le lemme suivant.

\begin{lemma}
\label{spaces-over-fields-lemma-order-vanishing}
Soit $S$ un schéma. Soit $X$ un espace algébrique intègre localement noethérien
sur $S$. Soit $Z \subset X$ un diviseur premier et soit $\xi \in |Z|$
le point générique. Alors l'anneau local hensélien $\mathcal{O}_{X, \xi}^h$
est un anneau local noethérien réduit de dimension $1$ et
il existe une application canonique injective
$$
R(X) \longrightarrow Q(\mathcal{O}_{X, \xi}^h)
$$
du corps des fonctions $R(X)$ de $X$ dans l'anneau total des fractions.
\end{lemma}

\begin{proof}
Nous utiliserons les résultats d'Espaces décents, section
\ref{decent-spaces-section-residue-fields-henselian-local-rings}.
Soit $(U, u) \to (X, \xi)$ un voisinage \'etale élémentaire.
Observons que $U$ est localement noethérien et réduit.
Ainsi, $\mathcal{O}_{U, u}$ est un anneau de dimension $1$
(par notre définition des diviseurs premiers),
réduit et noethérien.
Après avoir remplacé $U$ par un voisinage ouvert affine de $u$
nous pouvons supposer que $U$ est noethérien et affine.
Après avoir remplacé $U$ par un ouvert plus petit, nous pouvons supposer que toute
composante irréductible de $U$ passe par $u$.
Comme $U \to X$ est ouvert et que $X$ est irréductible, $U \to X$ est dominant.
Nous obtenons donc une application d'anneaux $R(X) \to R(U)$ en composant des applications rationnelles, voir
Morphismes d'espaces, section \ref{spaces-morphisms-section-rational-maps}.
Comme $R(X)$ est un corps, cette application est injective.
Par notre choix de $U$, nous voyons que $R(U)$ est l'anneau total des fractions
$Q(\mathcal{O}_{U, u})$, voir
Morphismes, lemme \ref{morphisms-lemma-integral-scheme-rational-functions}
et
Algèbre, lemme \ref{algebra-lemma-total-ring-fractions-no-embedded-points}.

\medskip\noindent
À ce stade, nous avons démontré toutes les assertions du lemme avec
$\mathcal{O}_{U, u}$ à la place de $\mathcal{O}_{X, \xi}^h$.
Cependant, $\mathcal{O}_{X, \xi}^h$ est l'hensélisation de
$\mathcal{O}_{U, u}$. Ainsi, $\mathcal{O}_{X, \xi}^h$ est un anneau
réduit noethérien de dimension $1$, voir
Compléments d'algèbre, lemmes
\ref{more-algebra-lemma-henselization-reduced},
\ref{more-algebra-lemma-henselization-dimension}, et
\ref{more-algebra-lemma-henselization-noetherian}.
Comme $\mathcal{O}_{U, u} \to \mathcal{O}_{X, \xi}^h$ est
fidèlement plat d'après Compléments d'algèbre, lemme
\ref{more-algebra-lemma-dumb-properties-henselization}
elle envoie les non diviseurs de zéro sur des non diviseurs de zéro.
Nous obtenons donc une application canonique
$Q(\mathcal{O}_{U, u}) \to Q(\mathcal{O}_{X, \xi}^h)$
et nous obtenons notre application.
Nous omettons la vérification que l'application est indépendante du
choix de $(U, u) \to (X, x)$ ; une approche légèrement meilleure
consisterait d'abord à observer que
$\colim Q(\mathcal{O}_{U, u}) = Q(\mathcal{O}_{X, \xi}^h)$.
\end{proof}

\begin{definition}
\label{spaces-over-fields-definition-order-vanishing}
Soit $S$ un schéma. Soit $X$ un espace algébrique intègre localement noethérien
sur $S$. Soit $f \in R(X)^*$. Pour tout diviseur premier
$Z \subset X$, nous définissons l'{\it ordre d'annulation de $f$ le long de $Z$}
comme l'entier
$$
\text{ord}_Z(f) =
\text{longueur}_{\mathcal{O}_{X, \xi}^h}
(\mathcal{O}_{X, \xi}^h/a \mathcal{O}_{X, \xi}^h) -
\text{longueur}_{\mathcal{O}_{X, \xi}^h}
(\mathcal{O}_{X, \xi}^h/b \mathcal{O}_{X, \xi}^h)
$$
où $a, b \in \mathcal{O}_{X, \xi}^h$ sont des non diviseurs de zéro
tels que l'image de $f$ dans $Q(\mathcal{O}_{X, \xi}^h)$
(lemme \ref{spaces-over-fields-lemma-order-vanishing}) soit égale à $a/b$.
Ceci est bien défini d'après
Algèbre, lemme \ref{algebra-lemma-ord-additive}.
\end{definition}

\noindent
Si $\mathcal{O}_{X, \xi}^h$ est un anneau intègre, nous obtenons
$$
\text{ord}_Z(f) = \text{ord}_{\mathcal{O}_{X, \xi}^h}(f)
$$
où le membre de droite est la notion de
la définition d'Algèbre \ref{algebra-definition-ord}.
Notons que pour $f, g \in R(X)^*$ nous avons
$$
\text{ord}_Z(fg) = \text{ord}_Z(f) + \text{ord}_Z(g).
$$
Bien sûr, il peut arriver que $\text{ord}_Z(f) < 0$.
Dans ce cas, nous disons que $f$ possède un {\it pôle} le long de $Z$
et que $-\text{ord}_Z(f) > 0$ est l'
{\it ordre du pôle de $f$ le long de $Z$}. Il est important de remarquer que
la condition $\text{ord}_Z(f) \geq 0$ n'est {\bf pas} équivalente
à la condition $f \in \mathcal{O}_{X, \xi}^h$, sauf si l'anneau local
$\mathcal{O}_{X, \xi}$ est un anneau de valuation discrète.

\begin{lemma}
\label{spaces-over-fields-lemma-order-vanishing-agrees}
Soit $S$ un schéma. Soit $X$ un espace algébrique intègre localement noethérien
sur $S$. Soit $f \in R(X)^*$. Si le diviseur premier
$Z \subset X$ rencontre le lieu schématique de $X$, alors l'ordre
d'annulation $\text{ord}_Z(f)$ de la définition \ref{spaces-over-fields-definition-order-vanishing}
coïncide avec l'ordre d'annulation de
Diviseurs, définition \ref{divisors-definition-order-vanishing}.
\end{lemma}

\begin{proof}
Après avoir rétréci $X$, nous pouvons supposer que $X$ est un schéma intègre noethérien.
Si $\xi \in Z$ désigne le point générique, alors
$\mathcal{O}_{X, \xi}^h$ est l'hensélisation de $\mathcal{O}_{X, \xi}$
(Espaces décents, lemme \ref{decent-spaces-lemma-describe-henselian-local-ring}).
Pour démontrer le lemme, il suffit et il est nécessaire de montrer que
$$
\text{longueur}_{\mathcal{O}_{X, \xi}}
(\mathcal{O}_{X, \xi}/a \mathcal{O}_{X, \xi}) =
\text{longueur}_{\mathcal{O}_{X, \xi}^h}
(\mathcal{O}_{X, \xi}^h/a \mathcal{O}_{X, \xi}^h)
$$
Ceci découle immédiatement de
Algèbre, lemme \ref{algebra-lemma-pullback-module}
(et du fait que $\mathcal{O}_{X, \xi} \to \mathcal{O}_{X, \xi}^h$
est un homomorphisme local plat d'anneaux locaux noethériens).
\end{proof}

\begin{lemma}
\label{spaces-over-fields-lemma-divisor-locally-finite}
Soit $S$ un schéma. Soit $X$ un espace algébrique intègre localement noethérien
sur $S$. Soit $f \in R(X)^*$. Alors les familles
$$
\{Z \subset X \mid Z\text{ est un diviseur premier de point générique }\xi
\text{ et }f\text{ n'appartient pas à }\mathcal{O}_{X, \xi}\}
$$
et
$$
\{Z \subset X \mid Z \text{ est un diviseur premier et }\text{ord}_Z(f) \not = 0\}
$$
sont localement finies dans $X$.
\end{lemma}

\begin{proof}
Il existe un sous-espace ouvert non vide $U \subset X$
tel que $f$ corresponde à une section de $\Gamma(U, \mathcal{O}_X^*)$.
Par conséquent, les diviseurs premiers qui peuvent apparaître dans les ensembles du lemme
correspondent tous à des composantes irréductibles de $|X| \setminus |U|$.
Le lemme \ref{spaces-over-fields-lemma-components-locally-finite} donne donc le résultat voulu.
\end{proof}

\noindent
Ce lemme nous permet de donner la définition suivante.

\begin{definition}
\label{spaces-over-fields-definition-principal-divisor}
Soit $S$ un schéma.
Soit $X$ un espace algébrique intègre localement noethérien sur $S$.
Soit $f \in R(X)^*$.
Le {\it diviseur de Weil principal associé à $f$} est le diviseur de Weil
$$
\text{div}(f) = \text{div}_X(f) = \sum \text{ord}_Z(f) [Z]
$$
où la somme porte sur les diviseurs premiers et où $\text{ord}_Z(f)$ est comme dans
la définition \ref{spaces-over-fields-definition-order-vanishing}. Ceci a un sens
d'après le lemme \ref{spaces-over-fields-lemma-divisor-locally-finite}.
\end{definition}

\begin{lemma}
\label{spaces-over-fields-lemma-div-additive}
Soit $S$ un schéma.
Soit $X$ un espace algébrique intègre localement noethérien sur $S$.
Soient $f, g \in R(X)^*$. Alors
$$
\text{div}_X(fg) = \text{div}_X(f) + \text{div}_X(g)
$$
en tant que diviseurs de Weil sur $X$.
\end{lemma}

\begin{proof}
Ceci est clair par additivité des fonctions $\text{ord}$.
\end{proof}

\noindent
Nous voyons d'après le lemme ci-dessus que la famille des diviseurs de Weil principaux
forme un sous-groupe du groupe de tous les diviseurs de Weil. Cela conduit à la
définition suivante.

\begin{definition}
\label{spaces-over-fields-definition-class-group}
Soit $S$ un schéma.
Soit $X$ un espace algébrique intègre localement noethérien sur $S$. Le
{\it groupe des classes de diviseurs de Weil} de $X$ est le quotient du
groupe des diviseurs de Weil par le sous-groupe des diviseurs de Weil principaux.
Notation : $\text{Cl}(X)$.
\end{definition}

\noindent
Par construction, nous obtenons un complexe exact
\begin{equation}
\label{spaces-over-fields-equation-Weil-divisor-class}
R(X)^* \xrightarrow{\text{div}} \text{Div}(X) \to \text{Cl}(X) \to 0
\end{equation}
que nous pouvons considérer comme une présentation de $\text{Cl}(X)$. Notre prochaine tâche
est de relier le groupe des classes de diviseurs de Weil au groupe de Picard.

\begin{example}
\label{spaces-over-fields-example-Z-mod-2}
Cet exemple prolonge celui de Morphismes d'espaces,
exemple \ref{spaces-morphisms-example-universal-homeomorphism}.
Considérons l'espace algébrique
$X = \mathbf{A}^1_k/\{t \sim -t \mid t \not = 0\}$.
C'est un espace algébrique lisse sur le corps $k$.
Il existe un homéomorphisme universel
$$
X \longrightarrow \mathbf{A}^1_k = \Spec(k[t])
$$
qui est un isomorphisme au-dessus de $\mathbf{A}^1_k \setminus \{0\}$.
Nous en concluons que $X$ est noethérien et intègre.
Comme $\dim(X) = 1$, les diviseurs premiers de $X$ sont
les points fermés de $X$.
Considérons l'unique point fermé $x \in |X|$ situé au-dessus de $0 \in \mathbf{A}^1_k$.
Comme $X \setminus \{x\}$ s'envoie isomorphiquement sur $\mathbf{A}^1 \setminus \{0\}$,
nous voyons que les classes dans $\text{Cl}(X)$ des points fermés distincts
de $x$ sont nulles. Cependant, le diviseur de $t$ sur $X$ est $2[x]$.
Nous concluons que $\text{Cl}(X) = \mathbf{Z}/2\mathbf{Z}$.
\end{example}

\begin{example}
\label{spaces-over-fields-example-not-a-domain}
Soit $k$ un corps. Soit
$$
U = \Spec(k[x, y]/(xy))
$$
qui est l'union des axes de coordonnées dans $\mathbf{A}^2_k$.
Notons $\Delta : U \to U \times_k U$ la diagonale et
$\Delta' : U \to U \times_k U$ l'application $u \mapsto (u, \sigma(u))$
où $\sigma : U \to U$, $(x, y) \mapsto (y, x)$
est l'automorphisme qui échange les axes de coordonnées. Posons
$$
R = \Delta(U) \amalg \Delta'(U \setminus \{0_U\})
$$
où $0_U \in U$ est l'origine. Il est facile de voir que
$R$ est une relation d'équivalence \'etale sur $U$. Le quotient $X = U/R$
est un espace algébrique. Le morphisme $U \to \mathbf{A}^1_k$,
$(x, y) \mapsto x + y$ est invariant par $R$ et définit donc
un morphisme
$$
X \longrightarrow \mathbf{A}^1_k
$$
Ce morphisme est un homéomorphisme universel et un isomorphisme au-dessus de
$\mathbf{A}^1_k \setminus \{0\}$. Il s'ensuit que $X$ est intègre
et noethérien. Exactement comme dans l'exemple \ref{spaces-over-fields-example-Z-mod-2}
le lecteur montre que $\text{Cl}(X) = \mathbf{Z}/2\mathbf{Z}$
avec pour générateur la classe correspondant à l'unique point fermé $x \in |X|$
situé au-dessus de $0 \in \mathbf{A}^1_k$. Cependant, dans ce cas, l'anneau local hensélien
de $X$ en $x$ n'est pas un anneau intègre, car c'est l'hensélisation
de $\mathcal{O}_{U, 0_U}$.
\end{example}














\section{Classe de diviseur de Weil associée à un module inversible}
\label{spaces-over-fields-section-c1}

\noindent
Dans cette section, nous suivons exactement la même progression que dans la
section \ref{spaces-over-fields-section-Weil-divisors} pour définir une application canonique
$\Pic(X) \to \text{Cl}(X)$
sur un espace algébrique intègre localement noethérien.

\medskip\noindent
Soit $S$ un schéma. Soit $X$ un espace algébrique intègre localement noethérien
sur $S$. Soit $\mathcal{L}$ un $\mathcal{O}_X$-module inversible.
D'après Diviseurs sur les espaces, lemme
\ref{spaces-divisors-lemma-regular-meromorphic-section-exists}
il existe une section méromorphe régulière
$s \in \Gamma(X, \mathcal{K}_X(\mathcal{L}))$.
En fait, d'après Diviseurs sur les espaces, lemme
\ref{spaces-divisors-lemma-compute-meromorphic}
c'est la même chose qu'un élément non nul de
$\mathcal{L}_\eta$, où $\eta \in |X|$ est le point générique.
Le même lemme nous dit que si $\mathcal{L} = \mathcal{O}_X$,
alors $s$ est la même chose qu'une fonction rationnelle non nulle sur $X$
(ainsi ce que nous ferons ci-dessous coïncide avec la construction de la
section \ref{spaces-over-fields-section-Weil-divisors}).

\medskip\noindent
Soit $Z \subset X$ un diviseur premier et soit $\xi \in |Z|$ le point générique.
Nous allons définir l'ordre d'annulation de $s$ le long de $Z$.
Considérons le morphisme canonique
$$
c_\xi : \Spec(\mathcal{O}_{X, \xi}^h) \longrightarrow X
$$
dont la source est le spectre de l'anneau local hensélien de $X$ en $\xi$
(Espaces décents, définition \ref{decent-spaces-definition-henselian-local-ring}).
Le tiré en arrière $\mathcal{L}_\xi = c_\xi^*\mathcal{L}$
est un module inversible, donc trivial ; choisissons un générateur $s_\xi$ de
$\mathcal{L}_\xi$. Comme $c_\xi$ est plat, les tirés en arrière des fonctions méromorphes
et des sections (régulières) sont définis pour $c_\xi$, voir
Diviseurs sur les espaces, définition
\ref{spaces-divisors-definition-pullback-meromorphic-sections} et
les lemmes \ref{spaces-divisors-lemma-pullback-meromorphic-sections-defined} et
\ref{spaces-divisors-lemma-meromorphic-sections-pullback}.
Nous obtenons ainsi
$$
c_\xi^*(s) = f s_\xi
$$
pour un élément non diviseur de zéro $f \in Q(\mathcal{O}_{X, \xi}^h)$.
Nous utilisons ici le lemme de Diviseurs \ref{divisors-lemma-locally-Noetherian-K}
pour identifier l'espace des sections méromorphes de
$\mathcal{L}_\xi \cong \mathcal{O}_{\Spec(\mathcal{O}_{X, \xi}^h)}$
en termes de l'anneau total des fractions de $\mathcal{O}_{X, \xi}^h$.
Convenons de noter cet élément
$$
s/s_\xi = f \in  Q(\mathcal{O}_{X, \xi}^h)
$$
Observons que $f = s/s_\xi$
est remplacé par $uf$ où $u \in \mathcal{O}_{X, \xi}^h$
est une unité lorsque nous changeons le choix de $s_\xi$.

\begin{definition}
\label{spaces-over-fields-definition-order-vanishing-meromorphic}
Soit $S$ un schéma. Soit $X$ un espace algébrique intègre localement noethérien
sur $S$. Soit $\mathcal{L}$ un
$\mathcal{O}_X$-module inversible.
Soit $s \in \Gamma(X, \mathcal{K}_X(\mathcal{L}))$
une section méromorphe régulière de $\mathcal{L}$.
Pour tout diviseur premier $Z \subset X$ dont le point générique est $\xi \in |Z|$
nous définissons l'
{\it ordre d'annulation de $s$ le long de $Z$}
comme l'entier
$$
\text{ord}_{Z, \mathcal{L}}(s) =
\text{longueur}_{\mathcal{O}_{X, \xi}^h}
(\mathcal{O}_{X, \xi}^h/a \mathcal{O}_{X, \xi}^h) -
\text{longueur}_{\mathcal{O}_{X, \xi}^h}
(\mathcal{O}_{X, \xi}^h/b \mathcal{O}_{X, \xi}^h)
$$
où $a, b \in \mathcal{O}_{X, \xi}^h$ sont des non diviseurs de zéro
tels que l'élément $s/s_\xi$ de $Q(\mathcal{O}_{X, \xi}^h)$
construit ci-dessus soit égal à $a/b$.
Ceci est bien défini d'après ce qui précède et
le lemme d'Algèbre \ref{algebra-lemma-ord-additive}.
\end{definition}

\noindent
Comme expliqué ci-dessus, une section méromorphe régulière $s$ de $\mathcal{O}_X$
peut s'écrire $s = f \cdot 1$ où $f$
est une fonction rationnelle non nulle sur $X$ et nous avons
$\text{ord}_Z(f) = \text{ord}_{Z, \mathcal{O}_X}(s)$.
Comme dans le cas des diviseurs principaux, nous avons le lemme suivant.

\begin{lemma}
\label{spaces-over-fields-lemma-divisor-meromorphic-locally-finite}
Soit $S$ un schéma. Soit $X$ un espace algébrique intègre localement noethérien
sur $S$. Soit $\mathcal{L}$ un $\mathcal{O}_X$-module inversible.
Soit $s \in \mathcal{K}_X(\mathcal{L})$ une
section méromorphe régulière (c'est-à-dire non nulle) de $\mathcal{L}$. Alors les ensembles
$$
\{Z \subset X \mid Z \text{ est un diviseur premier de point générique }\xi
\text{ et }s\text{ n'appartient pas à }\mathcal{L}_\xi\}
$$
et
$$
\{Z \subset X \mid Z \text{ est un diviseur premier et }
\text{ord}_{Z, \mathcal{L}}(s) \not = 0\}
$$
sont localement finis dans $X$.
\end{lemma}

\begin{proof}
Il existe un sous-espace ouvert non vide $U \subset X$ tel que $s$
corresponde à une section de $\Gamma(U, \mathcal{L})$ qui engendre
$\mathcal{L}$ sur $U$. Par conséquent, les diviseurs premiers qui peuvent apparaître
dans les ensembles du lemme correspondent tous à des composantes irréductibles de
$|X| \setminus |U|$. Le lemme \ref{spaces-over-fields-lemma-components-locally-finite}
donne le résultat voulu.
\end{proof}

\begin{lemma}
\label{spaces-over-fields-lemma-divisor-meromorphic-well-defined}
Soit $S$ un schéma. Soit $X$ un espace algébrique intègre localement noethérien
sur $S$. Soit $\mathcal{L}$ un $\mathcal{O}_X$-module inversible.
Soient $s, s' \in \mathcal{K}_X(\mathcal{L})$ des sections méromorphes
non nulles de $\mathcal{L}$. Alors $f = s/s'$
est un élément de $R(X)^*$ et nous avons
$$
\sum \text{ord}_{Z, \mathcal{L}}(s)[Z]
=
\sum \text{ord}_{Z, \mathcal{L}}(s')[Z]
+
\text{div}(f)
$$
en tant que diviseurs de Weil.
\end{lemma}

\begin{proof}
C'est clair d'après les définitions.
Notons que le lemme \ref{spaces-over-fields-lemma-divisor-meromorphic-locally-finite}
garantit que les sommes sont bien des diviseurs de Weil.
\end{proof}

\begin{definition}
\label{spaces-over-fields-definition-divisor-invertible-sheaf}
Soit $S$ un schéma. Soit $X$ un espace algébrique intègre localement noethérien
sur $S$. Soit $\mathcal{L}$ un $\mathcal{O}_X$-module inversible.
\begin{enumerate}
\item Pour toute section méromorphe non nulle $s$ de $\mathcal{L}$
nous définissons le {\it diviseur de Weil associé à $s$} comme
$$
\text{div}_\mathcal{L}(s) =
\sum \text{ord}_{Z, \mathcal{L}}(s) [Z] \in \text{Div}(X)
$$
où la somme porte sur les diviseurs premiers. Ceci est bien défini d'après
le lemme \ref{spaces-over-fields-lemma-divisor-meromorphic-locally-finite}.
\item Nous définissons la {\it classe de diviseur de Weil associée à $\mathcal{L}$}
comme l'image de $\text{div}_\mathcal{L}(s)$ dans $\text{Cl}(X)$
où $s$ est une quelconque section méromorphe non nulle de $\mathcal{L}$ sur $X$.
Ceci est bien défini d'après
le lemme \ref{spaces-over-fields-lemma-divisor-meromorphic-well-defined}.
\end{enumerate}
\end{definition}

\noindent
Comme prévu, cette construction est additive par rapport au module inversible.

\begin{lemma}
\label{spaces-over-fields-lemma-c1-additive}
Soit $S$ un schéma. Soit $X$ un espace algébrique intègre localement noethérien
sur $S$. Soient $\mathcal{L}$ et $\mathcal{N}$ deux $\mathcal{O}_X$-modules inversibles.
Soit $s$, resp.\ $t$, une section méromorphe non nulle de $\mathcal{L}$,
resp.\ de $\mathcal{N}$. Alors $st$ est une section méromorphe non nulle de
$\mathcal{L} \otimes_{\mathcal{O}_X} \mathcal{N}$ et
$$
\text{div}_{\mathcal{L} \otimes \mathcal{N}}(st)
=
\text{div}_\mathcal{L}(s) + \text{div}_\mathcal{N}(t)
$$
dans $\text{Div}(X)$. En particulier, la classe de diviseur de Weil de
$\mathcal{L} \otimes_{\mathcal{O}_X} \mathcal{N}$ est la somme
des classes de diviseurs de Weil de $\mathcal{L}$ et de $\mathcal{N}$.
\end{lemma}

\begin{proof}
Soit $s$, resp.\ $t$, une section méromorphe non nulle de
$\mathcal{L}$, resp.\ de $\mathcal{N}$. Alors $st$ est une section
méromorphe non nulle de $\mathcal{L} \otimes \mathcal{N}$.
Soit $Z \subset X$ un diviseur premier. Soit $\xi \in |Z|$ son point
générique. Choisissons des générateurs $s_\xi \in \mathcal{L}_\xi$ et
$t_\xi \in \mathcal{N}_\xi$, avec les notations introduites plus haut
dans cette section. Alors $s_\xi \otimes t_\xi$ est un générateur
de $(\mathcal{L} \otimes \mathcal{N})_\xi$.
Ainsi, $st/(s_\xi t_\xi) = (s/s_\xi)(t/t_\xi)$ dans
$Q(\mathcal{O}_{X, \xi}^h)$. En appliquant l'additivité du
lemme d'Algèbre \ref{algebra-lemma-ord-additive},
nous concluons que
$$
\text{div}_{\mathcal{L} \otimes \mathcal{N}, Z}(st)
=
\text{div}_{\mathcal{L}, Z}(s) + \text{div}_{\mathcal{N}, Z}(t)
$$
Nous omettons certains détails.
\end{proof}

\noindent
Soit $S$ un schéma. Soit $X$ un espace algébrique intègre localement noethérien
sur $S$. Les constructions et les lemmes ci-dessus donnent un homomorphisme
de groupes abéliens
\begin{equation}
\label{spaces-over-fields-equation-c1}
\Pic(X) \longrightarrow \text{Cl}(X)
\end{equation}
qui associe à tout module inversible sa classe de diviseur de Weil.

\begin{lemma}
\label{spaces-over-fields-lemma-normal-c1-injective}
Soit $S$ un schéma. Soit $X$ un espace algébrique intègre localement noethérien
sur $S$. Si $X$ est normal, alors l'application (\ref{spaces-over-fields-equation-c1})
$\Pic(X) \to \text{Cl}(X)$ est injective.
\end{lemma}

\begin{proof}
Soit $\mathcal{L}$ un $\mathcal{O}_X$-module inversible dont la
classe de diviseur de Weil associée est triviale. Soit $s$ une section
méromorphe régulière de $\mathcal{L}$. L'hypothèse signifie que
$\text{div}_\mathcal{L}(s) = \text{div}(f)$ pour un certain
$f \in R(X)^*$. Alors $t = f^{-1}s$ est une section
méromorphe régulière de $\mathcal{L}$ telle que
$\text{div}_\mathcal{L}(t) = 0$, voir le
lemme \ref{spaces-over-fields-lemma-divisor-meromorphic-well-defined}.
Nous affirmons que $t$ définit une trivialisation de $\mathcal{L}$.
Cette assertion achève la démonstration du lemme.
La démonstration qui suit est un peu maladroite, car nous ne
disposons pas encore d'une théorie très développée ; nous suggérons
au lecteur de la passer.

\medskip\noindent
Nous pouvons vérifier cette assertion localement pour la topologie \'etale. Soit
$U \in X_\etale$ affine tel que $\mathcal{L}|_U$ soit trivial. Soit
$s_U \in \Gamma(U, \mathcal{L}|_U)$ une trivialisation. D'après le
lemme de Propriétés \ref{properties-lemma-normal-locally-finite-nr-irreducibles},
nous pouvons aussi supposer $U$ intègre. Écrivons $U = \Spec(A)$, où $A$ est
un anneau intègre normal noethérien de corps des fractions $K$.
Nous pouvons écrire $t|_U = f s_U$ pour un élément $f$ de $K$, voir par exemple le
lemme de Diviseurs sur les espaces \ref{spaces-divisors-lemma-meromorphic-quasi-coherent}.
Soit $\mathfrak p \subset A$ un idéal premier de hauteur un
correspondant à un point de codimension $1$, $u \in U$, dont l'image
est un point de codimension $1$, $\xi \in |X|$. Choisissons une trivialisation
$s_\xi$ de $\mathcal{L}_\xi$ comme au début de cette section.
Choisissons un point géométrique $\overline{u}$ de $U$ situé au-dessus de $u$.
Alors
$$
(\mathcal{O}_{X, \xi}^h)^{sh} =
\mathcal{O}_{X, \overline{u}} =
\mathcal{O}_{U, u}^{sh} = (A_\mathfrak p)^{sh}
$$
voir les lemmes d'Espaces décents
\ref{decent-spaces-lemma-henselian-local-ring-strict}
et le lemme de
Propriétés des espaces
\ref{spaces-properties-lemma-describe-etale-local-ring}.
La normalité de $X$ montre que tous ces anneaux sont des
anneaux de valuation discrète. Les trivialisations $s_U$ et $s_\xi$
diffèrent par une unité comme sections du tiré en arrière de $\mathcal{L}$ sur
$\Spec(\mathcal{O}_{X, \overline{u}})$.
Écrivons $t = f_\xi s_\xi$ avec $f_\xi \in Q(\mathcal{O}_{X, \xi}^h)$.
Nous en concluons que $f_\xi$ et $f$ diffèrent par une unité
dans $Q(\mathcal{O}_{X, \overline{u}})$.
Si $Z \subset X$ désigne le diviseur premier correspondant à $\xi$
(lemme \ref{spaces-over-fields-lemma-decent-irreducible-closed}), alors
$0 = \text{ord}_{Z, \mathcal{L}}(t) =
\text{ord}_{\mathcal{O}_{X, \xi}^h}(f_\xi)$
et, comme $\mathcal{O}_{X, \xi}^h$ est un anneau de valuation discrète,
nous voyons que $f_\xi$ est une unité.
Ainsi, $f$ est une unité dans $\mathcal{O}_{X, \overline{u}}$
et, en particulier, $f \in A_\mathfrak p$.
Cela implique $f \in A$ d'après le lemme d'Algèbre
\ref{algebra-lemma-normal-domain-intersection-localizations-height-1}.
Nous concluons que $t \in \Gamma(X, \mathcal{L})$.
En répétant le raisonnement avec $t^{-1}$ considéré comme section méromorphe
de $\mathcal{L}^{\otimes -1}$, on achève la démonstration.
\end{proof}

















\section{Modifications et altérations}
\label{spaces-over-fields-section-modifications-alterations}

\noindent
Notre notion d'espace algébrique intègre permet de définir une modification
comme suit.

\begin{definition}
\label{spaces-over-fields-definition-modification}
Soit $S$ un schéma. Soit $X$ un espace algébrique intègre sur $S$. Une
{\it modification de $X$} est un morphisme propre birationnel
$f : X' \to X$ d'espaces algébriques sur $S$, où $X'$ est intègre.
\end{definition}

\noindent
Pour les morphismes birationnels d'espaces algébriques, voir la
définition d'Espaces décents \ref{decent-spaces-definition-birational}.

\begin{lemma}
\label{spaces-over-fields-lemma-modification-iso-over-open}
Soit $f : X' \to X$ une modification comme dans la
définition \ref{spaces-over-fields-definition-modification}.
Il existe un ouvert non vide $U \subset X$ tel que $f^{-1}(U) \to U$
soit un isomorphisme.
\end{lemma}

\begin{proof}
D'après le
lemme \ref{spaces-over-fields-lemma-finite-degree}, il existe un ouvert non vide $U \subset X$ tel
que $f^{-1}(U) \to U$ soit fini. Par platitude générique
(Morphismes d'espaces, proposition
\ref{spaces-morphisms-proposition-generic-flatness-reduced}),
nous pouvons supposer $f^{-1}(U) \to U$ plat et de présentation finie.
Ainsi, $f^{-1}(U) \to U$ est fini localement libre
(Morphismes d'espaces, lemme \ref{spaces-morphisms-lemma-finite-flat}).
Comme $f$ est birationnel, le degré de $X'$ sur $X$ est $1$.
Par conséquent, $f^{-1}(U) \to U$ est fini localement libre de degré $1$,
autrement dit, c'est un isomorphisme.
\end{proof}

\begin{definition}
\label{spaces-over-fields-definition-alteration}
Soit $S$ un schéma. Soit $X$ un espace algébrique intègre sur $S$.
Une {\it altération de $X$} est un morphisme propre dominant $f : Y \to X$
d'espaces algébriques sur $S$, avec $Y$ intègre, tel que $f^{-1}(U) \to U$
soit fini pour un certain ouvert non vide $U \subset X$.
\end{definition}

\noindent
Si $f : Y \to X$ est un morphisme dominant et propre entre espaces
algébriques intègres, alors c'est une altération dès que l'extension induite
des corps résiduels aux points génériques est finie. Voici l'énoncé
précis.

\begin{lemma}
\label{spaces-over-fields-lemma-alteration-generically-finite}
Soit $S$ un schéma. Soit $f : X \to Y$ un morphisme propre dominant
d'espaces algébriques intègres sur $S$. Alors $f$ est une altération
si et seulement si l'une des conditions équivalentes (1) -- (6) du
lemme \ref{spaces-over-fields-lemma-finite-degree} est satisfaite.
\end{lemma}

\begin{proof}
C'est une conséquence immédiate du lemme cité dans l'énoncé.
\end{proof}

\begin{lemma}
\label{spaces-over-fields-lemma-alteration-contained-in}
Soit $S$ un schéma. Soit $f : X \to Y$ un morphisme propre surjectif
d'espaces algébriques sur $S$. Supposons $Y$ intègre. Alors
il existe un sous-espace fermé intègre $X' \subset X$ tel que
$f' = f|_{X'} : X' \to Y$ soit une altération.
\end{lemma}

\begin{proof}
Soit $V \subset Y$ un ouvert affine non vide
(Espaces décents, théorème \ref{decent-spaces-theorem-decent-open-dense-scheme}).
Soit $\eta \in V$ le point générique. Alors
$X_\eta$ est un espace algébrique propre non vide sur $\eta$.
Choisissons un point fermé $x \in |X_\eta|$
(il en existe, car $|X_\eta|$ est un espace topologique quasi-compact
et sobre ; voir Espaces décents, proposition
\ref{decent-spaces-proposition-reasonable-sober},
et Topologie, lemme \ref{topology-lemma-quasi-compact-closed-point}).
Soit $X'$ le sous-espace fermé, muni de la structure réduite induite, porté par
$\overline{\{x\}} \subset |X|$ (Propriétés des espaces, définition
\ref{spaces-properties-definition-reduced-induced-space}).
Alors $f' : X' \to Y$ est surjectif, puisque son image contient $\eta$.
De plus, $f'$ est propre comme composé d'une immersion fermée
et d'un morphisme propre. Enfin, la fibre $X'_\eta$ ne possède qu'un
seul point ; pour le voir, appliquons le
lemme d'Espaces décents \ref{decent-spaces-lemma-topology-fibre}
à $X \to Y$ et à $X' \to Y$, au point $\eta$.
Comme $Y$ est décent et que $X' \to Y$ est séparé, $X'$ est décent
(Espaces décents, lemmes
\ref{decent-spaces-lemma-properties-trivial-implications} et
\ref{decent-spaces-lemma-property-over-property}).
Ainsi, $f'$ est une altération d'après le
lemme \ref{spaces-over-fields-lemma-alteration-generically-finite}.
\end{proof}








\section{Lieu schématique}
\label{spaces-over-fields-section-schematic}

\noindent
Nous avons déjà démontré plusieurs résultats sur le lieu schématique
d'un espace algébrique. En voici une liste de références :
\begin{enumerate}
\item Propriétés des espaces, sections
\ref{spaces-properties-section-schematic} et
\ref{spaces-properties-section-getting-a-scheme},
\item Espaces décents, section \ref{decent-spaces-section-schematic},
\item Propriétés des espaces, lemme
\ref{spaces-properties-lemma-point-like-spaces}
$\Leftarrow$
Espaces décents, lemme \ref{decent-spaces-lemma-decent-point-like-spaces}
$\Leftarrow$
Espaces décents, lemme \ref{decent-spaces-lemma-when-field},
\item Limites d'espaces, section \ref{spaces-limits-section-affine}, et
\item Limites d'espaces, section \ref{spaces-limits-section-representable}.
\end{enumerate}
Dans certains cas, certains types de morphismes d'espaces algébriques
sont automatiquement représentables, par exemple les morphismes
séparés et localement quasi-finis (Morphismes d'espaces, lemme
\ref{spaces-morphisms-lemma-locally-quasi-finite-separated-representable})
et les monomorphismes plats (Compléments sur les morphismes d'espaces, lemme
\ref{spaces-more-morphisms-lemma-flat-case}).
Dans la section \ref{spaces-over-fields-section-schematic-and-field-extension},
nous étudierons le comportement du lieu
schématique par extension du corps de base.

\begin{lemma}
\label{spaces-over-fields-lemma-locally-finite-type-dim-zero}
Soit $S$ un schéma. Soit $X$ un espace algébrique sur $S$.
Supposons que $X$ satisfasse au moins l'une des conditions suivantes :
\begin{enumerate}
\item $X$ est quasi-séparé et $\dim(X) = 0$,
\item $X$ est localement de type fini sur un corps $k$ et $\dim(X) = 0$,
\item $X$ est noethérien et $\dim(X) = 0$, ou
\item en ajouter d'autres ici.
\end{enumerate}
Alors $X$ est un schéma séparé et tout ouvert quasi-compact de $X$ est affine.
\end{lemma}

\begin{proof}
Si nous démontrons que tout ouvert quasi-compact de $X$ est affine, alors
$X$ est un schéma séparé. Nous pouvons donc supposer $X$ quasi-compact
et cherchons à montrer que $X$ est affine.
Les cas (2) et (3) découlent immédiatement du cas (1), mais nous donnerons des
démonstrations séparées de (2) et (3), qui utilisent beaucoup moins de théorie.

\medskip\noindent
Démonstration de (3). Soit $U$ un schéma affine et soit $U \to X$ un
morphisme \'etale. Posons $R = U \times_X U$. Les deux morphismes de
projection $s, t : R \to U$ sont des morphismes étales de schémas. D'après la
définition de Propriétés des espaces \ref{spaces-properties-definition-dimension},
nous avons $\dim(U) = 0$ et $\dim(R) = 0$.
Comme $R$ est un schéma localement noethérien de dimension $0$,
$R$ est une réunion disjointe de spectres
d'anneaux locaux artiniens
(Propriétés, lemme \ref{properties-lemma-locally-Noetherian-dimension-0}).
Puisque nous avons supposé $X$ noethérien (donc quasi-séparé), nous
concluons que $R$ est quasi-compact. Ainsi, $R$ est un schéma affine
(utiliser Schémas, lemme \ref{schemes-lemma-disjoint-union-affines}).
Les morphismes étales $s, t : R \to U$ induisent des extensions finies
des corps résiduels. Le lemme d'Algèbre
\ref{algebra-lemma-essentially-of-finite-type-into-artinian-local}
montre donc que $s$ et $t$ sont finis
(nous omettons un détail mineur).
La proposition de Groupoïdes
\ref{groupoids-proposition-finite-flat-equivalence}
montre alors que $X = U/R$ est un schéma affine.

\medskip\noindent
Démonstration de (2), presque identique à celle de (3).
Soit $U$ un schéma affine et soit $U \to X$ un morphisme étale surjectif.
Posons $R = U \times_X U$. Les deux morphismes de projection
$s, t : R \to U$ sont des morphismes étales de schémas. D'après la
définition de Propriétés des espaces \ref{spaces-properties-definition-dimension},
nous avons $\dim(U) = 0$ et, de même, $\dim(R) = 0$.
D'autre part, le morphisme $U \to \Spec(k)$ est localement de type
fini comme composé du morphisme étale $U \to X$ et de
$X \to \Spec(k)$ ; voir
Morphismes d'espaces,
les lemmes \ref{spaces-morphisms-lemma-composition-finite-type} et
\ref{spaces-morphisms-lemma-etale-locally-finite-type}.
De même, $R \to \Spec(k)$ est localement de type fini.
Le lemme de Variétés
\ref{varieties-lemma-algebraic-scheme-dim-0}
montre donc que $U$ et $R$ sont des réunions disjointes de spectres de
$k$-algèbres locales artiniennes finies sur $k$. Il en va donc
de même pour $U \times_{\Spec(k)} U$. Comme
$$
R = U \times_X U \longrightarrow U \times_{\Spec(k)} U
$$
est un monomorphisme, $R$ est une réunion finie (!) de spectres de
$k$-algèbres finies. Il s'ensuit que $R$ est affine ; voir
Schémas, lemme \ref{schemes-lemma-disjoint-union-affines}.
En appliquant de nouveau le
lemme de Variétés \ref{varieties-lemma-algebraic-scheme-dim-0},
nous voyons que $R$ est fini sur $k$. Ainsi, $s, t$
sont finis ; voir
Morphismes, lemme \ref{morphisms-lemma-finite-permanence}.
La proposition de Groupoïdes
\ref{groupoids-proposition-finite-flat-equivalence}
montre alors que $X = U/R$ est un schéma affine.

\medskip\noindent
Démonstration cohomologique de (1). D'après le lemme de Cohomologie des espaces
\ref{spaces-cohomology-lemma-vanishing-above-dimension},
les groupes de cohomologie supérieure s'annulent pour tout
faisceau quasi-cohérent $\mathcal{F}$ sur $X$. Par conséquent, $X$
est affine (donc, en particulier, un schéma) d'après la
proposition de Cohomologie des espaces
\ref{spaces-cohomology-proposition-vanishing-affine}.

\medskip\noindent
Démonstration géométrique de (1). Choisissons une stratification
$$
\emptyset = U_{n + 1} \subset
U_n \subset U_{n - 1} \subset \ldots \subset U_1 = X
$$
et des morphismes étales $f_p : V_p \to U_p$ comme dans le
lemme d'Espaces décents
\ref{decent-spaces-lemma-filter-quasi-compact-quasi-separated}
(nous utiliserons ci-dessous toutes leurs propriétés).
Alors $\dim(V_p) = 0$ par notre définition de la dimension des espaces
algébriques. Le lemme de Propriétés \ref{properties-lemma-dimension-zero}
s'applique donc à chaque $V_p$. Puis $f_p^{-1}(U_{p + 1}) \subset V_p$
est un ouvert quasi-compact, donc à la fois affine et fermé.
Il s'ensuit que $|T_p| \subset |U_p|$ (voir la référence citée)
est à la fois ouvert et fermé. Ainsi, $X$ est une réunion disjointe
de sous-espaces ouverts et fermés dont les structures réduites sont des schémas.
Il s'ensuit que $X$ est un schéma
(Limites d'espaces, lemme \ref{spaces-limits-lemma-reduction-scheme}).
La démonstration s'achève alors par le cas des schémas
déjà cité ci-dessus.
\end{proof}

\noindent
Le lemme suivant affirme qu'un espace algébrique quasi-séparé
est un schéma hors d'une partie de codimension au moins $1$.

\begin{lemma}
\label{spaces-over-fields-lemma-generic-point-in-schematic-locus}
Soit $S$ un schéma. Soit $X$ un espace algébrique quasi-séparé sur $S$.
Soit $x \in |X|$. Les conditions suivantes sont équivalentes :
\begin{enumerate}
\item $x$ est un point de codimension $0$ sur $X$ ;
\item l'anneau local de $X$ en $x$ est de dimension $0$ ;
\item $x$ est un point générique d'une composante irréductible de $|X|$.
\end{enumerate}
Si ces conditions sont satisfaites, il existe un sous-espace ouvert de $X$
qui contient $x$ et qui est un schéma.
\end{lemma}

\begin{proof}
L'équivalence de (1), (2) et (3) découle du
lemme d'Espaces décents \ref{decent-spaces-lemma-decent-generic-points}
et du fait qu'un espace algébrique quasi-séparé est décent
(Espaces décents, section \ref{decent-spaces-section-reasonable-decent}).
Nous donnerons toutefois, dans le paragraphe suivant, une démonstration plus
élémentaire de cette équivalence.

\medskip\noindent
Notons que (1) et (2) sont équivalents par définition
(Propriétés des espaces, définition
\ref{spaces-properties-definition-dimension-local-ring}).
Pour démontrer l'équivalence de (1) et (3), nous pouvons supposer $X$ quasi-compact.
Choisissons
$$
\emptyset = U_{n + 1} \subset
U_n \subset U_{n - 1} \subset \ldots \subset U_1 = X
$$
et $f_i : V_i \to U_i$ comme dans le lemme d'Espaces décents
\ref{decent-spaces-lemma-filter-quasi-compact-quasi-separated}.
Supposons $x \in U_i$, $x \not \in U_{i + 1}$. Alors $x = f_i(y)$ pour
un unique $y \in V_i$. Si (1) est satisfaite, $y$ est un point générique
d'une composante irréductible de $V_i$ (Propriétés des espaces, lemme
\ref{spaces-properties-lemma-codimension-0-points}).
Comme $f_i^{-1}(U_{i + 1})$ est un ouvert quasi-compact de $V_i$
qui ne contient pas $y$, il existe un voisinage ouvert $W \subset V_i$
de $y$ disjoint de $f_i^{-1}(V_i)$
(voir
Propriétés, lemme \ref{properties-lemma-generic-point-in-constructible},
ou, plus simplement, Algèbre, lemme
\ref{algebra-lemma-standard-open-containing-maximal-point}).
Alors $f_i|_W : W \to X$ est un isomorphisme sur son image, et donc
$x = f_i(y)$ est un point générique de $|X|$. Réciproquement, supposons (3).
Alors $f_i$ envoie $\overline{\{y\}}$ sur la composante irréductible
$\overline{\{x\}}$ de $|U_i|$. Comme $|f_i|$ est bijectif au-dessus de
$\overline{\{x\}}$, il s'ensuit que $\overline{\{y\}}$
est une composante irréductible de $U_i$. Ainsi, $x$ est un point de
codimension $0$.

\medskip\noindent
La dernière assertion du lemme est la
proposition de Propriétés des espaces
\ref{spaces-properties-proposition-locally-quasi-separated-open-dense-scheme}.
\end{proof}

\noindent
Le lemme suivant affirme qu'un espace algébrique séparé localement noethérien
est un schéma en codimension $1$, c'est-à-dire hors de la codimension $2$.

\begin{lemma}
\label{spaces-over-fields-lemma-codim-1-point-in-schematic-locus}
\begin{slogan}
Les espaces algébriques séparés sont des schémas en codimension 1.
\end{slogan}
Soit $S$ un schéma. Soit $X$ un espace algébrique sur $S$.
Soit $x \in |X|$. Si $X$ est séparé et localement noethérien, et si
la dimension de l'anneau local de $X$ en $x$ est $\leq 1$
(Propriétés des espaces, définition
\ref{spaces-properties-definition-dimension-local-ring}),
alors il existe un sous-espace ouvert de $X$ contenant $x$ qui est un schéma.
\end{lemma}

\begin{proof}
(Voir la remarque ci-dessous pour une autre approche qui évite les résultats sur
les groupoïdes finis.) Nous pouvons remplacer $X$ par un voisinage quasi-compact
de $x$ ; nous pouvons donc supposer $X$ quasi-compact, séparé et noethérien.
Il existe un schéma $U$ et un morphisme fini surjectif $U \to X$ ;
voir Limites d'espaces, proposition
\ref{spaces-limits-proposition-there-is-a-scheme-finite-over}.
Soit $R = U \times_X U$. Alors $j : R \to U \times_S U$ définit une relation
d'équivalence et nous obtenons un schéma en groupoïdes $(U, R, s, t, c)$ sur $S$,
avec $s, t$ finis et $U$ noethérien et séparé.
Soit $\{u_1, \ldots, u_n\} \subset U$ l'ensemble des points ayant pour image $x$.
Alors $\dim(\mathcal{O}_{U, u_i}) \leq 1$ d'après le
lemme d'Espaces décents
\ref{decent-spaces-lemma-dimension-local-ring-quasi-finite}.

\medskip\noindent
D'après le lemme de Compléments sur les groupoïdes
\ref{more-groupoids-lemma-find-affine-codimension-1},
il existe un ouvert affine $R$-invariant $W \subset U$ qui contient
l'orbite $\{u_1, \ldots, u_n\}$. Comme $U \to X$ est fini surjectif,
l'application continue $|U| \to |X|$ est fermée et surjective, donc
submersive d'après le lemme de Topologie
\ref{topology-lemma-closed-morphism-quotient-topology}.
Ainsi, $f(W)$ est ouvert et il existe un sous-espace ouvert $X' \subset X$
tel que $f : W \to X'$ soit un morphisme fini surjectif.
Alors $X'$ est un schéma affine d'après le
lemme de Cohomologie des espaces
\ref{spaces-cohomology-lemma-image-affine-finite-morphism-affine-Noetherian},
ce qui achève la démonstration.
\end{proof}

\begin{remark}
\label{spaces-over-fields-remark-alternate-proof-scheme-codim-1}
Voici une esquisse de démonstration du
lemme \ref{spaces-over-fields-lemma-codim-1-point-in-schematic-locus}
qui n'utilise pas le
lemme de Compléments sur les groupoïdes
\ref{more-groupoids-lemma-find-affine-codimension-1}.

\medskip\noindent
Étape 1. Nous pouvons supposer que $X$ est un espace algébrique réduit,
noethérien et séparé (par exemple d'après Cohomologie des espaces, lemme
\ref{spaces-cohomology-lemma-image-affine-finite-morphism-affine-Noetherian},
ou d'après
Limites d'espaces, lemme \ref{spaces-limits-lemma-reduction-scheme}),
et nous pouvons choisir un morphisme fini surjectif
$Y \to X$, où $Y$ est un schéma noethérien (d'après
Limites d'espaces, proposition
\ref{spaces-limits-proposition-there-is-a-scheme-finite-over}).

\medskip\noindent
Étape 2. Après avoir remplacé $X$ par un voisinage ouvert de $x$, il
existe un morphisme fini birationnel $X' \to X$ et un sous-schéma fermé
$Y' \subset X' \times_X Y$ tels que $Y' \to X'$ soit surjectif,
fini et localement libre. En effet, puisque $X$ est réduit, il existe un
sous-espace ouvert dense $U \subset X$ au-dessus duquel $Y$ est plat (Morphismes d'espaces,
proposition \ref{spaces-morphisms-proposition-generic-flatness-reduced}).
Nous pouvons alors choisir un éclatement $U$-admissible $b : \tilde X \to X$ tel
que la transformée stricte $\tilde Y$ de $Y$ soit plate sur $\tilde X$ ; voir
Compléments sur les morphismes d'espaces, lemme
\ref{spaces-more-morphisms-lemma-flat-after-blowing-up}.
(On peut aussi utiliser les schémas de Hilbert si l'on veut éviter
le résultat sur les éclatements.)
Soit alors $X' \subset \tilde X$ l'adhérence
schématique de $b^{-1}(U)$ et $Y' = X' \times_{\tilde X} \tilde Y$.
Comme $x$ est un point de codimension $1$, $X' \to X$ est fini au-dessus d'un
voisinage de $x$ (lemme \ref{spaces-over-fields-lemma-finite-in-codim-1}).

\medskip\noindent
Étape 3. Après avoir restreint $X$ à un voisinage plus petit de $x$, nous
obtenons que $X'$ est un schéma. En effet, $Y'$ est un schéma, $Y' \to X'$
est fini localement libre, et tout ensemble fini de points de codimension $1$
de $Y'$ est contenu dans un ouvert affine. Utiliser la
proposition de Propriétés des espaces
\ref{spaces-properties-proposition-finite-flat-equivalence-global}
et la
proposition de Variétés
\ref{varieties-proposition-finite-set-of-points-of-codim-1-in-affine}.

\medskip\noindent
Étape 4. Il existe un ouvert affine $W' \subset X'$ qui contient tous les points
situés au-dessus de $x$ et qui est l'image réciproque d'un sous-espace ouvert de $X$.
Pour le démontrer, soit $Z \subset X$ l'adhérence de l'ensemble des points
où $X' \to X$ n'est pas un isomorphisme. Nous pouvons supposer $x \in Z$, sinon
nous avons déjà terminé. Alors $x$ est un point générique d'une composante
irréductible de $Z$ et, après restriction de $X$, nous pouvons supposer $Z$ affine
(lemme \ref{spaces-over-fields-lemma-generic-point-in-schematic-locus}).
L'image réciproque $Z' \subset X'$ est alors elle aussi un schéma affine.
Soient $x_1, \ldots, x_n \in Z'$ les points ayant pour image $x$.
Nous pouvons trouver un ouvert affine $W'$ de $X'$ dont l'intersection avec
$Z'$ est l'image réciproque d'un ouvert principal de $Z$ contenant $x$.
En effet, choisissons d'abord un ouvert affine $W' \subset X'$ contenant
$x_1, \ldots, x_n$ grâce à la proposition de Variétés
\ref{varieties-proposition-finite-set-of-points-of-codim-1-in-affine}.
Choisissons ensuite un ouvert principal $D(f) \subset Z$ contenant $x$
dont l'image réciproque $D(f|_{Z'})$ est contenue dans $W' \cap Z'$.
Choisissons enfin $f' \in \Gamma(W', \mathcal{O}_{W'})$ dont la restriction
est $f|_{Z'}$, puis remplaçons $W'$ par $D(f') \subset W'$.
Comme $X' \to X$ est un isomorphisme hors de $Z' \to Z$, le choix
de $W'$ garantit que l'image $W \subset X$ de $W'$ est ouverte
et qu'elle a pour image réciproque $W'$ dans $X'$.

\medskip\noindent
Étape 5. Alors $W' \to W$ est un morphisme fini surjectif et $W$ est un schéma d'après le
lemme de Cohomologie des espaces
\ref{spaces-cohomology-lemma-image-affine-finite-morphism-affine-Noetherian},
ce qui achève la démonstration.
\end{remark}





\section{Lieu schématique et extension de corps}
\label{spaces-over-fields-section-schematic-and-field-extension}

\noindent
Il peut arriver qu'un espace algébrique non représentable sur un corps $k$
devienne représentable (c'est-à-dire un schéma) après changement de base à une extension
de $k$. Voir Espaces, exemple \ref{spaces-example-non-representable-descent}.
Dans cette section, nous étudions cette question.

\begin{lemma}
\label{spaces-over-fields-lemma-scheme-after-purely-inseparable-base-change}
Soit $k$ un corps. Soit $X$ un espace algébrique sur $k$.
S'il existe une extension radicielle de corps $k'/k$
telle que $X_{k'}$ soit un schéma, alors $X$ est un schéma.
\end{lemma}

\begin{proof}
Le morphisme $X_{k'} \to X$ est entier, surjectif et
universellement injectif. Le lemme découle donc du
lemme de Limites d'espaces
\ref{spaces-limits-lemma-integral-universally-bijective-scheme}.
\end{proof}

\begin{lemma}
\label{spaces-over-fields-lemma-when-scheme-after-base-change}
Soit $k$ un corps muni d'une clôture algébrique $\overline{k}$.
Soit $X$ un espace algébrique quasi-séparé sur $k$.
\begin{enumerate}
\item S'il existe une extension de corps $K/k$ telle que
$X_K$ soit un schéma, alors $X_{\overline{k}}$ est un schéma.
\item Si $X$ est quasi-compact et s'il existe une extension de corps
$K/k$ telle que $X_K$ soit un schéma, alors $X_{k'}$
est un schéma pour une certaine extension finie séparable $k'$ de $k$.
\end{enumerate}
\end{lemma}

\begin{proof}
Puisque tout espace algébrique est la réunion de ses sous-espaces ouverts
quasi-compacts, la première partie du lemme découle de
la seconde (nous omettons certains détails). Supposons donc $X$ quasi-compact
et donnée une extension $K/k$ pour laquelle $X_K$ est représentable.
Écrivons $K = \bigcup A$ comme la limite inductive des sous-$k$-algèbres de type fini
$A$. D'après le lemme de Limites d'espaces \ref{spaces-limits-lemma-limit-is-scheme},
$X_A$ est un schéma pour une certaine $A$. Choisissons un idéal maximal
$\mathfrak m \subset A$. D'après le théorème des zéros de Hilbert
(Algèbre, théorème \ref{algebra-theorem-nullstellensatz}),
le corps résiduel $k' = A/\mathfrak m$ est une extension finie de $k$.
Ainsi, $X_{k'}$ est un schéma. Si $k' \supset k$ n'est pas
séparable, soit $k'/k''/k$ la sous-extension
fournie par le lemme de Corps \ref{fields-lemma-separable-first}.
Comme $k'/k''$ est radicielle, le
lemme \ref{spaces-over-fields-lemma-scheme-after-purely-inseparable-base-change}
montre que l'espace algébrique $X_{k''}$ est un schéma. Puisque $k''|k$ est séparable,
la démonstration est achevée.
\end{proof}

\begin{lemma}
\label{spaces-over-fields-lemma-base-change-by-Galois}
Soit $k'/k$ une extension galoisienne finie de groupe de Galois $G$.
Soit $X$ un espace algébrique sur $k$. Alors $G$ agit librement sur
l'espace algébrique $X_{k'}$ et $X = X_{k'}/G$ au sens du
lemme de Propriétés des espaces \ref{spaces-properties-lemma-quotient}.
\end{lemma}

\begin{proof}
Démonstration omise. Indications : montrer d'abord que $\Spec(k) = \Spec(k')/G$.
Utiliser ensuite la compatibilité des quotients avec le changement de base.
\end{proof}

\begin{lemma}
\label{spaces-over-fields-lemma-when-quotient-scheme-at-point}
Soit $S$ un schéma. Soit $X$ un espace algébrique sur $S$ et
soit $G$ un groupe fini agissant librement sur $X$. Posons $Y = X/G$ comme
dans le lemme de Propriétés des espaces \ref{spaces-properties-lemma-quotient}.
Pour $y \in |Y|$, les conditions suivantes sont équivalentes :
\begin{enumerate}
\item $y$ appartient au lieu schématique de $Y$ ;
\item il existe un ouvert affine $U \subset X$
qui contient l'image réciproque de $y$.
\end{enumerate}
\end{lemma}

\begin{proof}
Il découle de la construction de $Y = X/G$ dans le
lemme de Propriétés des espaces \ref{spaces-properties-lemma-quotient}
que le morphisme $X \to Y$ est surjectif et \'etale.
Nous avons bien sûr $X \times_Y X = X \times G$ ; le morphisme
$X \to Y$ est donc même fini étale. Il est aussi surjectif.
Le lemme découle ainsi du
lemme d'Espaces décents \ref{decent-spaces-lemma-when-quotient-scheme-at-point}.
\end{proof}

\begin{lemma}
\label{spaces-over-fields-lemma-scheme-after-purely-transcendental-base-change}
Soit $k$ un corps. Soit $X$ un espace algébrique
quasi-séparé sur $k$. S'il existe une extension transcendante pure
$K/k$ telle que $X_K$ soit un schéma, alors
$X$ est un schéma.
\end{lemma}

\begin{proof}
Puisque tout espace algébrique est la réunion de ses sous-espaces ouverts
quasi-compacts, nous pouvons supposer $X$ quasi-compact (certains détails sont omis).
Rappelons (Corps, définition \ref{fields-definition-transcendence})
que l'hypothèse sur l'extension $K/k$ signifie que
$K$ est le corps des fractions d'un anneau de polynômes (en un nombre éventuellement infini
de variables) sur $k$. Ainsi, $K = \bigcup A$ est la réunion de sous-algèbres
dont chacune est une localisation d'une algèbre polynomiale en un nombre fini de variables sur $k$.
D'après le lemme de Limites d'espaces \ref{spaces-limits-lemma-limit-is-scheme},
$X_A$ est un schéma pour une certaine $A$. Écrivons
$$
A = k[x_1, \ldots, x_n][1/f]
$$
pour un certain élément non nul $f \in k[x_1, \ldots, x_n]$.

\medskip\noindent
Si $k$ est infini, nous pouvons achever la démonstration comme suit : choisissons
$a_1, \ldots, a_n \in k$ tels que $f(a_1, \ldots, a_n) \not = 0$.
Alors $(a_1, \ldots, a_n)$ définit un homomorphisme de $k$-algèbres $A \to k$
qui envoie $x_i$ sur $a_i$ et $1/f$ sur $1/f(a_1, \ldots, a_n)$.
Ainsi, le changement de base $X_A \times_{\Spec(A)} \Spec(k) \cong X$ est un
schéma, comme voulu.

\medskip\noindent
Dans ce paragraphe, achevons la démonstration lorsque $k$ est fini. Dans ce
cas, écrivons $X = \lim X_i$, où les $X_i$ sont de présentation finie sur $k$
et les morphismes de transition sont affines
(Limites d'espaces, lemme \ref{spaces-limits-lemma-relative-approximation}).
En utilisant le lemme de Limites d'espaces \ref{spaces-limits-lemma-limit-is-scheme},
nous voyons que $X_{i, A}$ est un schéma pour un certain $i$. Nous pouvons donc supposer
$X \to \Spec(k)$ de présentation finie. Soit $x \in |X|$ un point fermé.
Nous pouvons représenter $x$ par une immersion fermée
$\Spec(\kappa) \to X$
(Espaces décents, lemme \ref{decent-spaces-lemma-decent-space-closed-point}).
Alors $\Spec(\kappa) \to \Spec(k)$ est de type fini, donc $\kappa$
est une extension finie de $k$ (d'après le théorème des zéros de Hilbert ; voir
Algèbre, théorème \ref{algebra-theorem-nullstellensatz} ;
nous omettons certains détails). Posons $[\kappa : k] = d$. Choisissons un entier
$n \gg 0$ premier à $d$ et soit $k'/k$ l'extension
de degré $n$. Alors $k'/k$ est galoisienne, avec $G = \text{Aut}(k'/k)$
cyclique d'ordre $n$. Si $n$ est assez grand, il existe un homomorphisme de $k$-algèbres
$A \to k'$ pour la même raison que ci-dessus.
Alors $X_{k'}$ est un schéma et $X = X_{k'}/G$
(lemme \ref{spaces-over-fields-lemma-base-change-by-Galois}).
D'autre part, comme $n$ et $d$ sont premiers entre eux, nous avons
$$
\Spec(\kappa) \times_{X} X_{k'} =
\Spec(\kappa) \times_{\Spec(k)} \Spec(k') =
\Spec(\kappa \otimes_k k')
$$
est le spectre d'un corps. Autrement dit, la fibre de $X_{k'} \to X$
au-dessus de $x$ est réduite à un seul point. Le
lemme \ref{spaces-over-fields-lemma-when-quotient-scheme-at-point}
montre donc que $x$ appartient au lieu schématique de $X$, comme voulu.
\end{proof}

\begin{remark}
\label{spaces-over-fields-remark-when-does-the-argument-work}
Soit $k$ un corps fini. Soit $K/k$ une extension de corps géométriquement
irréductible. Alors $K$ est la limite inductive de $k$-algèbres de type fini
géométriquement irréductibles $A$. Pour une telle $A$, les estimations
de Lang et Weil \cite{LW} montrent que, pour $n \gg 0$, il existe
un homomorphisme de $k$-algèbres $A \to k'$ avec $k'/k$ de degré $n$.
L'analyse de l'argument donné dans la démonstration du
lemme \ref{spaces-over-fields-lemma-scheme-after-purely-transcendental-base-change}
montre que, si $X$ est un espace algébrique quasi-séparé sur $k$
et si $X_K$ est un schéma, alors $X$ est un schéma. Si nous avons un jour besoin de ce
résultat, nous le formulerons précisément et le démontrerons ici.
\end{remark}

\begin{lemma}
\label{spaces-over-fields-lemma-scheme-over-algebraic-closure-enough-affines}
Soit $k$ un corps muni d'une clôture algébrique $\overline{k}$. Soit $X$
un espace algébrique sur $k$ tel que
\begin{enumerate}
\item $X$ est décent et localement de type fini sur $k$ ;
\item $X_{\overline{k}}$ est un schéma ;
\item tout ensemble fini de points $\overline{k}$-rationnels de $X_{\overline{k}}$
est contenu dans un ouvert affine.
\end{enumerate}
Alors $X$ est un schéma.
\end{lemma}

\begin{proof}
Si $K/k$ est une extension, le changement de base $X_K$ est
décent (Espaces décents, lemme
\ref{decent-spaces-lemma-representable-named-properties})
et localement de type fini
sur $K$ (Morphismes d'espaces, lemme
\ref{spaces-morphisms-lemma-base-change-finite-type}).
D'après le lemme \ref{spaces-over-fields-lemma-scheme-after-purely-inseparable-base-change},
il suffit de démontrer que $X$ devient un schéma après changement de base à
la perfection de $k$ ; nous pouvons donc supposer que $k$ est un corps parfait
(cette étape n'est pas strictement nécessaire, mais facilite
la compréhension des autres arguments).
En recouvrant $X$ par des ouverts quasi-compacts, il suffit de démontrer
le lemme lorsque $X$ est quasi-compact (nous omettons un détail mineur).
Dans ce cas, $|X|$ est un espace topologique sobre
(Espaces décents, proposition
\ref{decent-spaces-proposition-reasonable-sober}).
Il suffit donc de montrer que tout point fermé de $|X|$
appartient au lieu schématique de $X$
(utiliser Propriétés des espaces, lemme \ref{spaces-properties-lemma-subscheme}, et
Topologie, lemme \ref{topology-lemma-quasi-compact-closed-point}).

\medskip\noindent
Soit $x \in |X|$ un point fermé. D'après le lemme d'Espaces décents
\ref{decent-spaces-lemma-decent-space-closed-point},
il existe une immersion fermée $\Spec(l) \to X$ représentant $x$.
Alors $\Spec(l) \to \Spec(k)$ est de type fini (Morphismes d'espaces,
lemme \ref{spaces-morphisms-lemma-composition-finite-type}) et nous
concluons que $l$ est une extension finie de $k$
d'après le théorème des zéros de Hilbert (Algèbre, théorème
\ref{algebra-theorem-nullstellensatz}). Elle est séparable, car
$k$ est parfait. Ainsi, le schéma
$$
\Spec(l) \times_X X_{\overline{k}} =
\Spec(l) \times_{\Spec(k)} \Spec(\overline{k}) =
\Spec(l \otimes_k \overline{k})
$$
est la réunion disjointe d'un nombre fini de points $\overline{k}$-rationnels.
D'après l'hypothèse (3), il existe un ouvert affine $W \subset X_{\overline{k}}$
qui contient ces points.

\medskip\noindent
D'après le lemme \ref{spaces-over-fields-lemma-when-scheme-after-base-change}, $X_{k'}$
est un schéma pour une certaine extension finie $k'/k$. Après avoir agrandi
$k'$, nous pouvons supposer qu'il existe un ouvert affine $U' \subset X_{k'}$
dont le changement de base à $\overline{k}$ redonne $W$
(utiliser que $X_{\overline{k}}$ est la limite des schémas $X_{k''}$
pour les extensions finies $k' \subset k'' \subset \overline{k}$, ainsi que les
lemmes de Limites \ref{limits-lemma-descend-opens} et
\ref{limits-lemma-limit-affine}). Nous pouvons supposer
que $k'/k$ est une extension galoisienne (prendre la clôture normale ;
Corps, lemme \ref{fields-lemma-normal-closure}, et utiliser
que $k$ est parfait). Posons $G = \text{Gal}(k'/k)$.
Par construction, le sous-schéma fermé stable par $G$
$\Spec(l) \times_X X_{k'}$ est contenu dans $U'$.
Ainsi, $x$ appartient au lieu schématique d'après les
lemmes \ref{spaces-over-fields-lemma-base-change-by-Galois} et
\ref{spaces-over-fields-lemma-when-quotient-scheme-at-point}.
\end{proof}

\noindent
Les deux lemmes suivants devraient figurer ailleurs.
Comparer le lemme suivant au
lemme d'Espaces décents \ref{decent-spaces-lemma-conditions-on-space-over-field}.

\begin{lemma}
\label{spaces-over-fields-lemma-locally-quasi-finite-over-field}
Soit $k$ un corps. Soit $X$ un espace algébrique sur $k$.
Les conditions suivantes sont équivalentes :
\begin{enumerate}
\item $X$ est localement quasi-fini sur $k$ ;
\item $X$ est localement de type fini sur $k$ et de dimension $0$ ;
\item $X$ est un schéma localement quasi-fini sur $k$ ;
\item $X$ est un schéma localement de type fini sur $k$ et de
dimension $0$ ;
\item $X$ est une réunion disjointe de spectres de $k$-algèbres locales artiniennes
$A$ finies sur $k$ telles que $\dim_k(A) < \infty$.
\end{enumerate}
\end{lemma}

\begin{proof}
Comme nous sommes sur un corps, la dimension relative de $X/k$ est égale à
la dimension de $X$. Le
lemme de Morphismes d'espaces
\ref{spaces-morphisms-lemma-locally-quasi-finite-rel-dimension-0}
montre donc que (1) et (2) sont équivalents. Il découle alors du
lemme \ref{spaces-over-fields-lemma-locally-finite-type-dim-zero}
(et d'implications immédiates) que (1) -- (4) sont équivalents.
Enfin, le
lemme de Variétés \ref{varieties-lemma-algebraic-scheme-dim-0}
montre que (1) -- (4) sont équivalents à (5).
\end{proof}

\begin{lemma}
\label{spaces-over-fields-lemma-mono-towards-locally-quasi-finite-over-field}
Soit $k$ un corps. Soit $f : X \to Y$ un monomorphisme d'espaces algébriques
sur $k$. Si $Y$ est localement quasi-fini sur $k$, alors $X$ l'est aussi.
\end{lemma}

\begin{proof}
Supposons $Y$ localement quasi-fini sur $k$. D'après le
lemme \ref{spaces-over-fields-lemma-locally-quasi-finite-over-field},
nous avons $Y = \coprod \Spec(A_i)$, où chaque $A_i$ est un
anneau local artinien fini sur $k$. D'après le
lemme d'Espaces décents
\ref{decent-spaces-lemma-monomorphism-toward-disjoint-union-dim-0-rings},
$X$ est un schéma. Considérons $X_i = f^{-1}(\Spec(A_i))$.
Alors $X_i$ possède soit un point, soit aucun. Si $X_i$ n'a aucun point, il
n'y a rien à démontrer. Si $X_i$ possède un point, alors
$X_i = \Spec(B_i)$, où $B_i$ est un anneau local de dimension zéro,
et $A_i \to B_i$ est un épimorphisme d'anneaux. En particulier,
$A_i/\mathfrak m_{A_i} = B_i/\mathfrak m_{A_i}B_i$, et
$A_i \to B_i$ est surjectif d'après le lemme de Nakayama,
Algèbre, lemme \ref{algebra-lemma-NAK}
(car $\mathfrak m_{A_i}$ est un idéal nilpotent !).
Ainsi, $B_i$ est une $k$-algèbre locale finie, et le
lemme \ref{spaces-over-fields-lemma-locally-quasi-finite-over-field}
montre que $X \to \Spec(k)$ est localement quasi-fini.
\end{proof}











\section{Espaces algébriques géométriquement réduits}
\label{spaces-over-fields-section-geometrically-reduced}

\noindent
Si $X$ est un espace algébrique réduit sur un corps, il peut arriver que $X$
cesse d'être réduit après extension du corps de base. Cela ne se produit pas
pour les espaces algébriques géométriquement réduits.

\begin{definition}
\label{spaces-over-fields-definition-geometrically-reduced}
Soit $k$ un corps.
Soit $X$ un espace algébrique sur $k$.
\begin{enumerate}
\item Soit $x \in |X|$ un point. Nous disons que $X$ est
{\it géométriquement réduit en $x$} si $\mathcal{O}_{X, \overline{x}}$
est géométriquement réduit sur $k$.
\item Nous disons que $X$ est {\it géométriquement réduit} sur $k$
si $X$ est géométriquement réduit en tout point de $X$.
\end{enumerate}
\end{definition}

\noindent
Observons que, si $X$ est géométriquement réduit en $x$, alors l'anneau local
de $X$ en $x$ est réduit (Propriétés des espaces, lemme
\ref{spaces-properties-lemma-reduced-local-ring}).
De même, si $X$ est géométriquement réduit sur $k$, alors $X$ est réduit
(d'après Propriétés des espaces, lemme \ref{spaces-properties-lemma-reduced-space}).
Le lemme suivant implique notamment que cette définition est compatible
avec la propriété correspondante pour les schémas sur un corps.

\begin{lemma}
\label{spaces-over-fields-lemma-geometrically-reduced-at-point}
Soit $k$ un corps. Soit $X$ un espace algébrique sur $k$.
Soit $x \in |X|$. Les conditions suivantes sont équivalentes :
\begin{enumerate}
\item $X$ est géométriquement réduit en $x$ ;
\item il existe un voisinage \'etale $(U, u) \to (X, x)$
tel que $U$ soit un schéma et que $U$ soit géométriquement réduit en $u$ ;
\item pour tout voisinage \'etale $(U, u) \to (X, x)$
tel que $U$ soit un schéma, $U$ est géométriquement réduit en $u$.
\end{enumerate}
\end{lemma}

\begin{proof}
Rappelons que l'anneau local $\mathcal{O}_{X, \overline{x}}$
est l'hensélisé strict de $\mathcal{O}_{U, u}$ ; voir le
lemme de Propriétés des espaces
\ref{spaces-properties-lemma-describe-etale-local-ring}.
D'après le lemme de Variétés \ref{varieties-lemma-geometrically-reduced-at-point},
$U$ est géométriquement réduit en $u$ si et seulement
si $\mathcal{O}_{U, u}$ est géométriquement réduit sur $k$.
Il nous reste donc à montrer ceci : si $A$ est une $k$-algèbre locale,
$A$ est géométriquement réduite sur $k$ si et seulement si
$A^{sh}$ est géométriquement réduite sur $k$.
Vérifions-le à partir de la définition des algèbres
géométriquement réduites (Algèbre, définition \ref{algebra-definition-geometrically-reduced}).
Soit $K/k$ une extension de corps.
Comme $A \to A^{sh}$ est fidèlement plat
(Compléments d'algèbre, lemme \ref{more-algebra-lemma-dumb-properties-henselization}),
le morphisme
$A \otimes_k K \to A^{sh} \otimes_k K$ est fidèlement plat
(Algèbre, lemme \ref{algebra-lemma-flat-base-change}).
Par conséquent, si $A^{sh} \otimes_k K$ est réduit, alors
$A \otimes_k K$ l'est aussi, d'après le lemme d'Algèbre \ref{algebra-lemma-descent-reduced}.
Réciproquement, rappelons que $A^{sh}$ est une limite inductive
d'$A$-algèbres \'etales ; voir le
lemme d'Algèbre \ref{algebra-lemma-strict-henselization}.
Ainsi, $A^{sh} \otimes_k K$ est une limite inductive
filtrante d'algèbres \'etales sur $A \otimes_k K$.
Nous concluons par le lemme d'Algèbre \ref{algebra-lemma-reduced-goes-up}.
\end{proof}

\begin{lemma}
\label{spaces-over-fields-lemma-geometrically-reduced}
Soit $k$ un corps. Soit $X$ un espace algébrique sur $k$.
Les conditions suivantes sont équivalentes :
\begin{enumerate}
\item $X$ est géométriquement réduit ;
\item il existe un morphisme \'etale surjectif $U \to X$, où $U$
est un schéma, et $U$ est géométriquement réduit ;
\item pour tout morphisme \'etale $U \to X$
où $U$ est un schéma, $U$ est géométriquement réduit.
\end{enumerate}
\end{lemma}

\begin{proof}
Cela résulte immédiatement des définitions et du
lemme \ref{spaces-over-fields-lemma-geometrically-reduced-at-point}.
\end{proof}

\noindent
Cette notion ne présente pas d'intérêt en caractéristique nulle.

\begin{lemma}
\label{spaces-over-fields-lemma-perfect-reduced}
Soit $X$ un espace algébrique sur un corps parfait $k$ (par exemple,
$k$ est de caractéristique nulle).
\begin{enumerate}
\item Pour $x \in |X|$, si $\mathcal{O}_{X, \overline{x}}$ est
réduit, alors $X$ est géométriquement réduit en $x$.
\item Si $X$ est réduit, alors $X$ est géométriquement réduit sur $k$.
\end{enumerate}
\end{lemma}

\begin{proof}
La première assertion résulte du
lemme d'Algèbre \ref{algebra-lemma-separable-extension-preserves-reducedness}
et de la définition d'un corps parfait
(Algèbre, définition \ref{algebra-definition-perfect}).
La seconde assertion résulte de la première.
\end{proof}

\begin{lemma}
\label{spaces-over-fields-lemma-geometrically-reduced-positive-characteristic}
Soit $k$ un corps de caractéristique $p > 0$. Soit $X$ un espace algébrique
sur $k$. Les conditions suivantes sont équivalentes :
\begin{enumerate}
\item $X$ est géométriquement réduit sur $k$ ;
\item $X_{k'}$ est réduit pour toute extension de corps $k'/k$ ;
\item $X_{k'}$ est réduit pour toute
extension radicielle finie de corps $k'/k$ ;
\item $X_{k^{1/p}}$ est réduit ;
\item $X_{k^{perf}}$ est réduit ;
\item $X_{\bar k}$ est réduit.
\end{enumerate}
\end{lemma}

\begin{proof}
Choisissons un morphisme \'etale surjectif $U \to X$, où $U$ est un schéma.
Le lemme \ref{spaces-over-fields-lemma-geometrically-reduced} ramène l'énoncé au
résultat pour $U$ sur $k$.
Voir le lemme de Variétés \ref{varieties-lemma-geometrically-reduced}.
\end{proof}

\begin{lemma}
\label{spaces-over-fields-lemma-geometrically-reduced-upstairs}
Soit $k$ un corps. Soit $X$ un espace algébrique sur $k$.
Soit $k'/k$ une extension de corps. Soit $x \in |X|$ un point et soit
$x' \in |X_{k'}|$ un point au-dessus de $x$.
Les conditions suivantes sont équivalentes :
\begin{enumerate}
\item $X$ est géométriquement réduit en $x$ ;
\item $X_{k'}$ est géométriquement réduit en $x'$.
\end{enumerate}
En particulier, $X$ est géométriquement réduit sur $k$ si et seulement si
$X_{k'}$ est géométriquement réduit sur $k'$.
\end{lemma}

\begin{proof}
Choisissons un morphisme \'etale $U \to X$, où $U$ est un schéma,
et un point $u \in U$ dont l'image est $x \in |X|$.
D'après le lemme de Propriétés des espaces \ref{spaces-properties-lemma-points-cartesian},
nous pouvons choisir un point $u' \in U_{k'} = U \times_X X_{k'}$ dont les images
sont $u$ et $x'$. D'après le lemme \ref{spaces-over-fields-lemma-geometrically-reduced-at-point},
l'énoncé se ramène au lemme pour $U, u, u'$, à savoir au
lemme de Variétés \ref{varieties-lemma-geometrically-reduced-upstairs}.
\end{proof}

\begin{lemma}
\label{spaces-over-fields-lemma-geometrically-reduced-etale-local}
Soit $k$ un corps. Soit $f : X \to Y$ un morphisme
d'espaces algébriques sur $k$. Soit $x \in |X|$ un point
dont l'image est $y \in |Y|$.
\begin{enumerate}
\item si $f$ est \'etale en $x$, alors
$X$ est géométriquement réduit en $x$ $\Leftrightarrow$
$Y$ est géométriquement réduit en $y$ ;
\item si $f$ est \'etale surjectif, alors
$X$ est géométriquement réduit $\Leftrightarrow$
$Y$ est géométriquement réduit.
\end{enumerate}
\end{lemma}

\begin{proof}
L'assertion (1) est claire, car
$\mathcal{O}_{X, \overline{x}} = \mathcal{O}_{Y, \overline{y}}$
lorsque $f$ est \'etale en $x$.
L'assertion (2) découle immédiatement de l'assertion (1).
\end{proof}






\section{Espaces algébriques géométriquement connexes}
\label{spaces-over-fields-section-geometrically-connected}

\noindent
Si $X$ est un espace algébrique connexe sur un corps, il peut arriver que
$X$ devienne non connexe après extension du corps de base. Cela ne se
produit pas pour les espaces algébriques géométriquement connexes.

\begin{definition}
\label{spaces-over-fields-definition-geometrically-connected}
Soit $X$ un espace algébrique sur le corps $k$. Nous disons que $X$ est
{\it géométriquement connexe} sur $k$ si le changement de base $X_{k'}$
est connexe pour toute extension de corps $k'$ de $k$.
\end{definition}

\noindent
Par convention, un espace topologique connexe est non vide ; à plus forte raison,
les espaces algébriques géométriquement connexes sont non vides.

\begin{lemma}
\label{spaces-over-fields-lemma-geometrically-connected-check-after-extension}
Soit $X$ un espace algébrique sur le corps $k$.
Soit $k'/k$ une extension de corps.
Alors $X$ est géométriquement connexe sur $k$ si et seulement si
$X_{k'}$ est géométriquement connexe sur $k'$.
\end{lemma}

\begin{proof}
Si $X$ est géométriquement connexe sur $k$, il est clair que
$X_{k'}$ est géométriquement connexe sur $k'$. Réciproquement,
pour toute extension de corps $k''/k$, il existe une extension de corps commune
$k'''/k$ de $k'/k$ et $k''/k$ ; voir le
lemme de Corps \ref{fields-lemma-common-extension-field}. Comme le
morphisme $X_{k'''} \to X_{k''}$ est surjectif (c'est un changement de base
d'un morphisme surjectif entre spectres de corps), la
connexité de $X_{k'''}$ entraîne celle de $X_{k''}$.
Ainsi, si $X_{k'}$ est géométriquement connexe sur $k'$, alors
$X$ est géométriquement connexe sur $k$.
\end{proof}

\begin{lemma}
\label{spaces-over-fields-lemma-bijection-connected-components}
Soit $k$ un corps. Soient $X$, $Y$ des espaces algébriques sur $k$.
Supposons $X$ géométriquement connexe sur $k$.
Alors le morphisme de projection
$$
p : X \times_k Y \longrightarrow Y
$$
induit une bijection entre les composantes connexes.
\end{lemma}

\begin{proof}
Soit $y \in |Y|$ représenté par un morphisme $\Spec(K) \to Y$,
où $K$ est un corps. La fibre de $|X \times_k Y| \to |Y|$ au-dessus de $y$
est l'image de $|X_K| \to |X \times_k Y|$ d'après le
lemme de Propriétés des espaces \ref{spaces-properties-lemma-points-cartesian}.
Ces fibres sont donc connexes puisque $X$ est, par hypothèse,
géométriquement connexe. D'après le
lemme de Morphismes d'espaces
\ref{spaces-morphisms-lemma-space-over-field-universally-open}
l'application $|p|$ est ouverte.
Nous pouvons donc appliquer le lemme de Topologie
\ref{topology-lemma-connected-fibres-connected-components}
pour conclure.
\end{proof}

\begin{lemma}
\label{spaces-over-fields-lemma-separably-closed-field-connected-components}
Soit $k'/k$ une extension de corps. Soit $X$ un espace algébrique
sur $k$. Supposons $k$ séparablement clos. Alors le morphisme
$X_{k'} \to X$ induit une bijection entre les composantes connexes. En particulier,
$X$ est géométriquement connexe sur $k$ si et seulement si $X$ est connexe.
\end{lemma}

\begin{proof}
Comme $k$ est séparablement clos, nous voyons que
$k'$ est géométriquement connexe sur $k$ ; voir le
lemme d'Algèbre
\ref{algebra-lemma-separably-closed-connected-implies-geometric}.
Par conséquent, $Z = \Spec(k')$ est géométriquement connexe sur $k$ d'après le
lemme de Variétés \ref{varieties-lemma-affine-geometrically-connected}.
Comme $X_{k'} = Z \times_k X$, le résultat est un cas particulier du
lemme \ref{spaces-over-fields-lemma-bijection-connected-components}.
\end{proof}

\begin{lemma}
\label{spaces-over-fields-lemma-characterize-geometrically-connected}
Soit $k$ un corps. Soit $X$ un espace algébrique sur $k$.
Soit $\overline{k}$ une clôture séparable de $k$.
Alors $X$ est géométriquement connexe si et seulement si le changement de base
$X_{\overline{k}}$ est connexe.
\end{lemma}

\begin{proof}
Supposons $X_{\overline{k}}$ connexe. Soit $k'/k$ une extension de
corps. Il existe une extension de corps $\overline{k}'/\overline{k}$
dans laquelle $k'$ se plonge dans $\overline{k}'$ par-dessus $k$.
D'après le lemme \ref{spaces-over-fields-lemma-separably-closed-field-connected-components},
$X_{\overline{k}'}$ est connexe.
Comme $X_{\overline{k}'} \to X_{k'}$ est surjectif, nous concluons
que $X_{k'}$ est connexe, comme voulu.
\end{proof}

\noindent
Soit $k$ un corps. Soit $\overline{k}/k$ une extension galoisienne
(éventuellement infinie). Par exemple, $\overline{k}$ pourrait être la
clôture séparable de $k$.
À tout $\sigma \in \text{Gal}(\overline{k}/k)$ correspond un
automorphisme
$
\Spec(\sigma) :
\Spec(\overline{k})
\longrightarrow
\Spec(\overline{k})
$.
Remarquons que
$\Spec(\sigma) \circ \Spec(\tau) = \Spec(\tau \circ \sigma)$.
Nous obtenons donc une action
$$
\text{Gal}(\overline{k}/k)^{opp} \times \Spec(\overline{k})
\longrightarrow
\Spec(\overline{k})
$$
du groupe opposé sur le schéma $\Spec(\overline{k})$.
Soit $X$ un espace algébrique sur $k$. Comme
$X_{\overline{k}} =
\Spec(\overline{k}) \times_{\Spec(k)} X$
par définition, l'action ci-dessus induit une action canonique
\begin{equation}
\label{spaces-over-fields-equation-galois-action-base-change-kbar}
\text{Gal}(\overline{k}/k)^{opp} \times X_{\overline{k}}
\longrightarrow
X_{\overline{k}}.
\end{equation}

\begin{lemma}
\label{spaces-over-fields-lemma-Galois-action-quasi-compact-open}
Soit $k$ un corps. Soit $X$ un espace algébrique sur $k$.
Soit $\overline{k}$ une extension galoisienne (éventuellement infinie) de $k$.
Soit $V \subset X_{\overline{k}}$ un ouvert quasi-compact.
Alors :
\begin{enumerate}
\item il existe une sous-extension finie $\overline{k}/k'/k$
et un ouvert quasi-compact $V' \subset X_{k'}$ tels que
$V = (V')_{\overline{k}}$,
\item il existe un sous-groupe ouvert $H \subset \text{Gal}(\overline{k}/k)$
tel que $\sigma(V) = V$ pour tout $\sigma \in H$.
\end{enumerate}
\end{lemma}

\begin{proof}
Choisissons un schéma $U$ et un morphisme \'etale surjectif $U \to X$.
Choisissons un ouvert quasi-compact $W \subset U_{\overline{k}}$ dont
l'image dans $X_{\overline{k}}$ est $V$. Cela est possible, car
$|U_{\overline{k}}| \to |X_{\overline{k}}|$ est continue et
$|U_{\overline{k}}|$ possède une base d'ouverts quasi-compacts. Nous pouvons appliquer le
lemme de Variétés
\ref{varieties-lemma-Galois-action-quasi-compact-open}
à $W \subset U_{\overline{k}}$, ce qui donne le lemme.
\end{proof}

\begin{lemma}
\label{spaces-over-fields-lemma-closed-fixed-by-Galois}
Soit $k$ un corps. Soit $\overline{k}/k$ une extension galoisienne
(éventuellement infinie). Soit $X$ un espace algébrique sur $k$. Soit
$\overline{T} \subset |X_{\overline{k}}|$ vérifiant les propriétés suivantes :
\begin{enumerate}
\item $\overline{T}$ est une partie fermée de $|X_{\overline{k}}|$ ;
\item pour tout $\sigma \in \text{Gal}(\overline{k}/k)$,
nous avons $\sigma(\overline{T}) = \overline{T}$.
\end{enumerate}
Alors il existe une partie fermée $T \subset |X|$ dont l'image réciproque
dans $|X_{k'}|$ est $\overline{T}$.
\end{lemma}

\begin{proof}
Soit $T \subset |X|$ l'image de $\overline{T}$.
Comme $|X_{\overline{k}}| \to |X|$ est surjectif, l'assertion signifie
que $T$ est fermé et que son image réciproque est $\overline{T}$.
Choisissons un schéma $U$ et un morphisme \'etale surjectif $U \to X$.
Dans le cas des schémas
(voir le lemme de Variétés \ref{varieties-lemma-closed-fixed-by-Galois}),
il existe une partie fermée $T' \subset |U|$ dont l'image réciproque
dans $|U_{\overline{k}}|$ est l'image réciproque de $\overline{T}$.
Comme $|U_{\overline{k}}| \to |X_{\overline{k}}|$ est surjectif,
$T'$ est l'image réciproque de $T$ par $|U| \to |X|$.
Par notre construction de la topologie sur $|X|$, cela signifie que $T$ est
fermé. De la même manière, on voit que $\overline{T}$ est l'image
réciproque de $T$.
\end{proof}

\begin{lemma}
\label{spaces-over-fields-lemma-characterize-geometrically-disconnected}
Soit $k$ un corps. Soit $X$ un espace algébrique sur $k$.
Les conditions suivantes sont équivalentes :
\begin{enumerate}
\item $X$ est géométriquement connexe ;
\item pour toute extension de corps finie séparable $k'/k$,
l'espace algébrique $X_{k'}$ est connexe.
\end{enumerate}
\end{lemma}

\begin{proof}
Cette démonstration est identique à celle du
lemme de Variétés \ref{varieties-lemma-characterize-geometrically-disconnected},
si ce n'est que
nous remplaçons le
lemme de Variétés \ref{varieties-lemma-characterize-geometrically-connected}
par le lemme \ref{spaces-over-fields-lemma-characterize-geometrically-connected},
nous remplaçons le
lemme de Variétés \ref{varieties-lemma-Galois-action-quasi-compact-open}
par le lemme \ref{spaces-over-fields-lemma-Galois-action-quasi-compact-open}, et
nous remplaçons le
lemme de Variétés \ref{varieties-lemma-closed-fixed-by-Galois}
par le lemme \ref{spaces-over-fields-lemma-closed-fixed-by-Galois}.
Nous invitons le lecteur à lire cette démonstration plutôt que celle-ci.

\medskip\noindent
Il découle immédiatement de la définition que (1) implique (2).
Supposons que $X$ ne soit pas géométriquement connexe.
Soit $k \subset \overline{k}$ une clôture séparable
de $k$. D'après le
lemme \ref{spaces-over-fields-lemma-characterize-geometrically-connected},
$X_{\overline{k}}$ est non connexe.
Écrivons $X_{\overline{k}} = \overline{U} \amalg \overline{V}$,
où $\overline{U}$ et $\overline{V}$ sont des sous-espaces algébriques ouverts, fermés et non vides
de $X_{\overline{k}}$.

\medskip\noindent
Soit $W \subset X$ un sous-espace ouvert quasi-compact quelconque.
Alors $W_{\overline{k}} \cap \overline{U}$ et
$W_{\overline{k}} \cap \overline{V}$ sont des sous-espaces ouverts et fermés de
$W_{\overline{k}}$. En particulier, $W_{\overline{k}} \cap \overline{U}$ et
$W_{\overline{k}} \cap \overline{V}$ sont quasi-compacts et, d'après le
lemme \ref{spaces-over-fields-lemma-Galois-action-quasi-compact-open},
$W_{\overline{k}} \cap \overline{U}$ et
$W_{\overline{k}} \cap \overline{V}$
sont tous deux définis sur une sous-extension finie et stables par un
sous-groupe ouvert de $\text{Gal}(\overline{k}/k)$.
Nous l'utiliserons désormais sans autre mention.

\medskip\noindent
Choisissons un sous-espace ouvert quasi-compact $W_0 \subset X$ tel que
$W_{0, \overline{k}} \cap \overline{U}$ et
$W_{0, \overline{k}} \cap \overline{V}$ soient tous deux non vides.
Choisissons une sous-extension finie $\overline{k}/k'/k$
et une décomposition $W_{0, k'} = U_0' \amalg V_0'$ en parties ouvertes et fermées
telle que
$W_{0, \overline{k}} \cap \overline{U} = (U'_0)_{\overline{k}}$ et
$W_{0, \overline{k}} \cap \overline{V} = (V'_0)_{\overline{k}}$.
Posons $H = \text{Gal}(\overline{k}/k') \subset \text{Gal}(\overline{k}/k)$.
En particulier,
$\sigma(W_{0, \overline{k}} \cap \overline{U}) =
W_{0, \overline{k}} \cap \overline{U}$, et de même pour
$\overline{V}$.

\medskip\noindent
$W_0$, $k'$ étant choisis comme ci-dessus, pour tout sous-espace ouvert quasi-compact
$W \subset X$, posons
$$
U_W =
\bigcap\nolimits_{\sigma \in H} \sigma(W_{\overline{k}} \cap \overline{U}),
\quad
V_W =
\bigcup\nolimits_{\sigma \in H} \sigma(W_{\overline{k}} \cap \overline{V}).
$$
Comme $W_{\overline{k}} \cap \overline{U}$ et
$W_{\overline{k}} \cap \overline{V}$ sont fixés par un sous-groupe ouvert de
$\text{Gal}(\overline{k}/k)$, la réunion et l'intersection
ci-dessus sont finies. Ainsi, $U_W$ et $V_W$ sont des sous-espaces ouverts et fermés.
De plus, par construction, $W_{\bar k} = U_W \amalg V_W$.

\medskip\noindent
Affirmons que, si $W \subset W' \subset X$ sont des sous-espaces ouverts
quasi-compacts, alors $W_{\overline{k}} \cap U_{W'} = U_W$ et
$W_{\overline{k}} \cap V_{W'} = V_W$. Vérification omise.
Il s'ensuit qu'en posant $U = \bigcup_{W \subset X} U_W$
et $V = \bigcup_{W \subset X} V_W$, nous obtenons
$X_{\overline{k}} = U \amalg V$ comme réunion disjointe de parties ouvertes
et fermées.
Il est clair que $V$ est non vide, puisqu'il est construit en prenant des
réunions (localement). D'autre part, $U$ est non vide puisqu'il contient
$W_0 \cap \overline{U}$ par construction. Enfin, $U, V \subset X_{\bar k}$
sont fermés et stables par $H$ par construction. Donc, d'après le
lemme \ref{spaces-over-fields-lemma-closed-fixed-by-Galois},
nous avons $U = (U')_{\bar k}$ et $V = (V')_{\bar k}$ pour certaines
parties fermées $U', V' \subset X_{k'}$. Manifestement, $X_{k'} = U' \amalg V'$,
et $X_{k'}$ est non connexe, comme voulu.
\end{proof}






\section{Espaces algébriques géométriquement irréductibles}
\label{spaces-over-fields-section-geometrically-irreducible}

\noindent
L'exemple d'Espaces \ref{spaces-example-infinite-product} montre qu'il
vaut mieux ne pas considérer les espaces algébriques irréductibles en toute
généralité\footnote{Bien entendu, si nous disons « l'espace algébrique $X$
est irréductible », nous voulons probablement dire « l'espace topologique $|X|$
est irréductible ».}. Pour les espaces algébriques décents (par exemple quasi-séparés),
ce type de désastre ne se produit pas. Nous posons donc la
définition suivante seulement sous l'hypothèse que notre espace algébrique est décent.

\begin{definition}
\label{spaces-over-fields-definition-geometrically-irreducible}
Soit $k$ un corps.
Soit $X$ un espace algébrique décent sur $k$.
Nous disons que $X$ est {\it géométriquement irréductible} si l'espace
topologique $|X_{k'}|$ est
irréductible\footnote{Un espace irréductible est non vide.}
pour toute extension de corps $k'$ de $k$.
\end{definition}

\noindent
Observons que $X_{k'}$ est un espace algébrique décent
(Espaces décents, lemme
\ref{decent-spaces-lemma-representable-named-properties}).
Par conséquent, l'espace topologique $|X_{k'}|$ est sobre,
d'après la proposition d'Espaces décents \ref{decent-spaces-proposition-reasonable-sober}.





\section{Espaces algébriques géométriquement intègres}
\label{spaces-over-fields-section-geometrically-integral}

\noindent
Rappelons que les espaces algébriques intègres sont décents par définition ; voir
la section \ref{spaces-over-fields-section-integral-spaces}.

\begin{definition}
\label{spaces-over-fields-definition-geometrically-integral}
Soit $X$ un espace algébrique sur le corps $k$. Nous disons que $X$ est
{\it géométriquement intègre} sur $k$ si l'espace algébrique
$X_{k'}$ est intègre (définition \ref{spaces-over-fields-definition-integral-algebraic-space})
pour toute extension de corps $k'$ de $k$.
\end{definition}

\noindent
En particulier, $X$ est un espace algébrique décent.
On peut relier cette propriété au fait d'être géométriquement réduit et
géométriquement irréductible comme suit.

\begin{lemma}
\label{spaces-over-fields-lemma-geometrically-integral}
Soit $k$ un corps. Soit $X$ un espace algébrique décent sur $k$.
Alors $X$ est géométriquement intègre sur $k$ si et seulement si
$X$ est à la fois géométriquement réduit et géométriquement irréductible
sur $k$.
\end{lemma}

\begin{proof}
C'est une conséquence immédiate des définitions, car notre
notion d'intégrité (sous l'hypothèse de décence) équivaut à
être réduit et irréductible.
\end{proof}

\begin{lemma}
\label{spaces-over-fields-lemma-proper-geometrically-reduced-global-sections}
Soit $k$ un corps. Soit $X$ un espace algébrique propre sur $k$.
\begin{enumerate}
\item $A = H^0(X, \mathcal{O}_X)$ est une $k$-algèbre de dimension finie ;
\item $A = \prod_{i = 1, \ldots, n} A_i$ est un produit de
$k$-algèbres locales artiniennes, avec un facteur par composante connexe de $|X|$ ;
\item si $X$ est réduit, alors $A = \prod_{i = 1, \ldots, n} k_i$
est un produit de corps, chacun étant une extension finie de $k$ ;
\item si $X$ est géométriquement réduit, alors $k_i$ est fini séparable
sur $k$ ;
\item si $X$ est géométriquement connexe, alors $A$ est géométriquement
irréductible sur $k$ ;
\item si $X$ est géométriquement irréductible, alors $A$ est géométriquement
irréductible sur $k$ ;
\item si $X$ est géométriquement réduit et connexe, alors $A = k$ ;
\item si $X$ est géométriquement intègre, alors $A = k$.
\end{enumerate}
\end{lemma}

\begin{proof}
D'après le lemme de Cohomologie des espaces
\ref{spaces-cohomology-lemma-proper-over-affine-cohomology-finite}
$A = H^0(X, \mathcal{O}_X)$ est une
$k$-algèbre de dimension finie. Cela démontre (1).

\medskip\noindent
Alors $A$ est un produit d'anneaux locaux d'après le
lemme d'Algèbre \ref{algebra-lemma-finite-dimensional-algebra} et la
proposition d'Algèbre \ref{algebra-proposition-dimension-zero-ring}.
Si $X = Y \amalg Z$, avec $Y$ et $Z$ des sous-espaces ouverts de $X$, nous obtenons
un idempotent $e \in A$ en prenant la section de $\mathcal{O}_X$
qui vaut $1$ sur $Y$ et $0$ sur $Z$. Réciproquement, si $e \in A$
est un idempotent, nous obtenons une décomposition correspondante de $|X|$.
Enfin, comme $|X|$ est un espace topologique noethérien
(d'après le lemme de Morphismes d'espaces
\ref{spaces-morphisms-lemma-finite-presentation-noetherian} et
le lemme de Propriétés des espaces
\ref{spaces-properties-lemma-Noetherian-topology})
ses composantes connexes sont ouvertes. Les composantes connexes
de $|X|$ correspondent donc $1$ pour $1$ aux idempotents primitifs de $A$.
Cela démontre (2).

\medskip\noindent
Si $X$ est réduit, alors $A$ est réduit
(Propriétés des espaces, lemme \ref{spaces-properties-lemma-reduced-space}).
Les anneaux locaux $A_i = k_i$ sont donc réduits et sont par conséquent des corps
(par exemple d'après le lemme d'Algèbre \ref{algebra-lemma-minimal-prime-reduced-ring}).
Cela démontre (3).

\medskip\noindent
Si $X$ est géométriquement réduit, la même propriété vaut pour
$A \otimes_k \overline{k} =
H^0(X_{\overline{k}}, \mathcal{O}_{X_{\overline{k}}})$
(voir le lemme de Cohomologie des espaces
\ref{spaces-cohomology-lemma-flat-base-change-cohomology} pour l'égalité).
Il s'ensuit que $k_i \otimes_k \overline{k}$ est un produit
de corps et que $k_i/k$ est donc séparable, par exemple d'après les
lemmes d'Algèbre
\ref{algebra-lemma-characterize-separable-field-extensions} et
\ref{algebra-lemma-geometrically-reduced-finite-purely-inseparable-extension}.
Cela démontre (4).

\medskip\noindent
Si $X$ est géométriquement connexe, alors $A \otimes_k \overline{k} =
H^0(X_{\overline{k}}, \mathcal{O}_{X_{\overline{k}}})$
est un anneau local de dimension zéro d'après (2) ; son
spectre n'a donc qu'un point et il est en particulier irréductible.
Ainsi, $A$ est géométriquement irréductible. Cela démontre (5).
Bien entendu, (5) entraîne (6).

\medskip\noindent
Si $X$ est géométriquement réduit et connexe, alors
$A = k_1$ est un corps et l'extension $k_1/k$ est finie séparable et
géométriquement irréductible. Mais $k_1 \otimes_k \overline{k}$
est alors un produit de $[k_1 : k]$ copies de $\overline{k}$, d'où
$k_1 = k$. Cela démontre (7). Bien entendu, (7) entraîne (8).
\end{proof}

\begin{lemma}
\label{spaces-over-fields-lemma-characterize-trivial-pic-integral}
Soit $k$ un corps. Soit $X$ un espace algébrique propre intègre sur $k$.
Soit $\mathcal{L}$ un $\mathcal{O}_X$-module inversible.
Si $H^0(X, \mathcal{L})$ et $H^0(X, \mathcal{L}^{\otimes - 1})$
sont tous deux non nuls, alors $\mathcal{L} \cong \mathcal{O}_X$.
\end{lemma}

\begin{proof}
Soient $s \in H^0(X, \mathcal{L})$ et $t \in H^0(X, \mathcal{L}^{\otimes - 1})$
des sections non nulles. Soit $x \in |X|$ un point appartenant au support de $s$.
Choisissons un voisinage affine \'etale $(U, u) \to (X, x)$ tel que
$\mathcal{L}|_U \cong \mathcal{O}_U$. Alors $s|_U$ correspond à une fonction régulière non nulle
sur un schéma réduit (car $X$ est réduit), à savoir $U$, et elle ne
s'annule donc pas en un point générique d'une composante irréductible de $U$. D'après le
lemme d'Espaces décents \ref{decent-spaces-lemma-decent-generic-points},
le point générique $\eta$ de $|X|$ appartient au support de $s$.
Il en va de même pour $t$. Alors, bien sûr, $st$ est non nul, car
l'anneau local de $X$ en $\eta$ est un corps (d'après le lemme précité,
l'anneau local est de dimension zéro ; comme $X$ est réduit, l'anneau local est
réduit ; appliquer le lemme d'Algèbre \ref{algebra-lemma-minimal-prime-reduced-ring}).
Or, nous avons vu que $K = H^0(X, \mathcal{O}_X)$
est un corps au lemme \ref{spaces-over-fields-lemma-proper-geometrically-reduced-global-sections}.
Ainsi, $st$ est partout non nul, et
$s : \mathcal{O}_X \to \mathcal{L}$ est un isomorphisme.
\end{proof}






\section{Dimension}
\label{spaces-over-fields-section-dimension}

\noindent
Dans cette section, nous poursuivons l'étude de la dimension.
Voici une liste des résultats antérieurs :
\begin{enumerate}
\item la dimension est définie dans la
section de Propriétés des espaces
\ref{spaces-properties-section-dimension},
\item la dimension de l'anneau local est définie dans la
section de Propriétés des espaces
\ref{spaces-properties-section-dimension-local-ring},
\item quelques résultats figurent dans les lemmes de Propriétés des espaces
\ref{spaces-properties-lemma-dimension-local-ring} et
\ref{spaces-properties-lemma-dimension-decent-invariant-under-etale},
\item la dimension relative est définie dans la
section de Morphismes d'espaces \ref{spaces-morphisms-section-relative-dimension},
\item des résultats sur la dimension des fibres figurent dans la
section de Morphismes d'espaces \ref{spaces-morphisms-section-dimension-fibres},
\item une forme faible de la formule de dimension figure dans la
section de Morphismes d'espaces \ref{spaces-morphisms-section-dimension-formula},
\item un résultat sur la lissité et la dimension figure dans le lemme de Morphismes d'espaces
\ref{spaces-morphisms-lemma-smoothness-dimension-spaces},
\item pour les espaces décents, la dimension est $\dim(|X|)$, d'après le
lemme d'Espaces décents \ref{decent-spaces-lemma-dimension-decent-space},
\item les morphismes quasi-finis et la dimension sont étudiés dans les
lemmes d'Espaces décents
\ref{decent-spaces-lemma-dimension-local-ring-quasi-finite} et
\ref{decent-spaces-lemma-dimension-quasi-finite}.
\end{enumerate}
Dans la section de Compléments sur les morphismes d'espaces
\ref{spaces-more-morphisms-section-dimension}
nous étudierons les sauts de
dimension dans les fibres d'un morphisme de type fini.

\begin{lemma}
\label{spaces-over-fields-lemma-integral-dimension}
Soit $S$ un schéma. Soit $f : X \to Y$ un morphisme entier
d'espaces algébriques. Alors $\dim(X) \leq \dim(Y)$.
Si $f$ est surjectif, alors $\dim(X) = \dim(Y)$.
\end{lemma}

\begin{proof}
Choisissons $V \to Y$ \'etale surjectif, avec $V$ un schéma.
Alors $U = X \times_Y V$ est un schéma et $U \to V$ est entier
(et surjectif si $f$ est surjectif).
D'après le lemme de Propriétés des espaces
\ref{spaces-properties-lemma-dimension-decent-invariant-under-etale}
nous avons $\dim(X) = \dim(U)$ et $\dim(Y) = \dim(V)$.
Le résultat découle donc du cas des schémas,
qui est le lemme de Morphismes \ref{morphisms-lemma-integral-dimension}.
\end{proof}

\begin{lemma}
\label{spaces-over-fields-lemma-alteration-dimension}
Soit $S$ un schéma. Soit $f : X \to Y$ un morphisme d'espaces algébriques
sur $S$. Supposons que :
\begin{enumerate}
\item $Y$ est localement noethérien ;
\item $X$ et $Y$ sont des espaces algébriques intègres ;
\item $f$ est dominant ;
\item $f$ est localement de type fini.
\end{enumerate}
Si $x \in |X|$ et $y \in |Y|$ sont les points génériques, alors
$$
\dim(X) \leq \dim(Y) + \text{degré de transcendance de }x/y.
$$
Si $f$ est propre, alors il y a égalité.
\end{lemma}

\begin{proof}
Rappelons que $|X|$ et $|Y|$ sont des espaces topologiques irréductibles et sobres ; voir la
discussion qui suit la définition \ref{spaces-over-fields-definition-integral-algebraic-space}.
Ainsi, la dominance de $f$ signifie que $|f|$ envoie $x$ sur $y$.
De plus, $x \in |X|$ est l'unique point auquel
l'anneau local de $X$ est de dimension $0$ ; voir le
lemme d'Espaces décents \ref{decent-spaces-lemma-decent-generic-points}.
D'après le lemme de Morphismes d'espaces
\ref{spaces-morphisms-lemma-dimension-formula-general}
la dimension de l'anneau local de $X$ en
tout point $x' \in |X|$ est au plus égale à celle de l'anneau local
de $Y$ en $y' = f(x')$, augmentée du degré de transcendance de $x/y$.
Comme la dimension de $X$, resp.\ celle de $Y$, est le
supremum des dimensions des anneaux locaux en $x'$, resp.\ en $y'$
(Propriétés des espaces, lemme \ref{spaces-properties-lemma-dimension}),
nous obtenons l'inégalité.

\medskip\noindent
Supposons $f$ propre.
Soit $V \subset Y$ un sous-espace ouvert quasi-compact non vide.
Si nous pouvons démontrer l'égalité pour le morphisme $f^{-1}(V) \to V$,
nous l'obtenons pour $X \to Y$. Nous pouvons donc supposer que
$X$ et $Y$ sont quasi-compacts.
Observons que $X$ est quasi-séparé, car c'est
un espace algébrique décent localement noethérien ; voir le
lemme d'Espaces décents
\ref{decent-spaces-lemma-locally-Noetherian-decent-quasi-separated}.
Nous pouvons donc choisir $Y' \to Y$ fini surjectif, où $Y'$
est un schéma ; voir la proposition de Limites d'espaces
\ref{spaces-limits-proposition-there-is-a-scheme-finite-over}.
Après avoir remplacé $Y'$ par un sous-schéma fermé convenable, nous
pouvons supposer $Y'$ intègre ; voir, par exemple, le lemme plus général
\ref{spaces-over-fields-lemma-alteration-contained-in}.
D'après le même lemme, nous pouvons choisir un sous-espace fermé
$X' \subset X \times_Y Y'$ tel que $X'$ soit intègre
et que $X' \to X$ soit fini surjectif.
Alors $X'$ est aussi localement noethérien
(Morphismes d'espaces, lemme
\ref{spaces-morphisms-lemma-locally-finite-type-locally-noetherian})
et nous pouvons utiliser la proposition de Limites d'espaces
\ref{spaces-limits-proposition-there-is-a-scheme-finite-over}
une nouvelle fois pour choisir un morphisme fini surjectif $X'' \to X'$
avec $X''$ un schéma. Comme précédemment, nous pouvons supposer $X''$
intègre. On obtient le diagramme
$$
\xymatrix{
X'' \ar[d] \ar[r] & X \ar[d]^f \\
Y' \ar[r] & Y
}
$$
D'après le lemme \ref{spaces-over-fields-lemma-integral-dimension}, nous avons $\dim(X'') = \dim(X)$
et $\dim(Y') = \dim(Y)$. Comme $X$ et $Y$ possèdent des voisinages ouverts
de $x$, resp.\ de $y$, qui sont des schémas, nous voyons aussitôt que les points génériques
$x'' \in X''$, resp.\ $y' \in Y'$, sont les uniques points envoyés
sur $x$, resp.\ sur $y$, et que les extensions de corps résiduels
$\kappa(x'')/\kappa(x)$ et $\kappa(y')/\kappa(y)$ sont finies.
Il s'ensuit que le degré de transcendance de $x''/y'$ est le
même que celui de $x/y$. L'égalité découle donc
du cas des schémas, qui est le
lemme de Morphismes \ref{morphisms-lemma-alteration-dimension}.
\end{proof}





\section{Espaces lisses sur les corps}
\label{spaces-over-fields-section-smooth}

\noindent
Cette section est l'analogue de la
section de Variétés \ref{varieties-section-smooth}.

\begin{lemma}
\label{spaces-over-fields-lemma-smooth-regular}
Soit $k$ un corps.
Soit $X$ un espace algébrique lisse sur $k$.
Alors $X$ est un espace algébrique régulier.
\end{lemma}

\begin{proof}
Choisissons un schéma $U$ et un morphisme \'etale surjectif $U \to X$.
Le morphisme $U \to \Spec(k)$ est lisse, car il est composé
d'un morphisme \'etale (donc lisse) et d'un morphisme lisse (voir les
lemmes de Morphismes d'espaces \ref{spaces-morphisms-lemma-etale-smooth}
et \ref{spaces-morphisms-lemma-composition-smooth}).
Par conséquent, $U$ est régulier d'après le
lemme de Variétés \ref{varieties-lemma-smooth-regular}.
D'après la
définition de Propriétés des espaces
\ref{spaces-properties-definition-type-property}
cela signifie que $X$ est régulier.
\end{proof}

\begin{lemma}
\label{spaces-over-fields-lemma-smooth-separable-closed-points-dense}
Soit $k$ un corps. Soit $X$ un espace algébrique lisse sur $\Spec(k)$.
L'ensemble des $x \in |X|$ qui sont images de morphismes $\Spec(k') \to X$
où $k' \supset k$ est fini séparable est dense dans $|X|$.
\end{lemma}

\begin{proof}
Choisissons un schéma $U$ et un morphisme \'etale surjectif $U \to X$.
Le morphisme $U \to \Spec(k)$ est lisse, car il est composé
d'un morphisme \'etale (donc lisse) et d'un morphisme lisse (voir les
lemmes de Morphismes d'espaces \ref{spaces-morphisms-lemma-etale-smooth}
et \ref{spaces-morphisms-lemma-composition-smooth}).
Nous pouvons donc appliquer le lemme de Variétés
\ref{varieties-lemma-smooth-separable-closed-points-dense}, selon lequel
les points fermés de $U$ dont les corps résiduels sont finis séparables sur
$k$ sont denses. Le lemme en résulte par notre définition de la
topologie sur $|X|$.
\end{proof}







\section{Caractéristiques d'Euler-Poincaré}
\label{spaces-over-fields-section-euler}

\noindent
Dans cette section, nous démontrons quelques propriétés élémentaires des caractéristiques
d'Euler-Poincaré des faisceaux cohérents sur les espaces algébriques propres sur des corps.

\begin{definition}
\label{spaces-over-fields-definition-euler-characteristic}
Soit $k$ un corps. Soit $X$ un espace algébrique propre sur $k$. Soit $\mathcal{F}$
un $\mathcal{O}_X$-module cohérent. Dans cette situation, la
{\it caractéristique d'Euler-Poincaré de $\mathcal{F}$} est l'entier
$$
\chi(X, \mathcal{F}) = \sum\nolimits_i (-1)^i \dim_k H^i(X, \mathcal{F}).
$$
La justification de cette formule est donnée ci-dessous.
\end{definition}

\noindent
Dans la situation de la définition, les espaces vectoriels
$H^i(X, \mathcal{F})$ non nuls sont en nombre fini (lemme de Cohomologie des espaces
\ref{spaces-cohomology-lemma-vanishing-quasi-separated})
et ils sont tous de dimension finie
(lemme de Cohomologie des espaces
\ref{spaces-cohomology-lemma-proper-over-affine-cohomology-finite}). Ainsi,
$\chi(X, \mathcal{F}) \in \mathbf{Z}$ est bien défini. Remarquons que
cette définition dépend du corps $k$ et non pas seulement de la paire
$(X, \mathcal{F})$.

\begin{lemma}
\label{spaces-over-fields-lemma-euler-characteristic-additive}
Soit $k$ un corps. Soit $X$ un espace algébrique propre sur $k$.
Soit $0 \to \mathcal{F}_1 \to \mathcal{F}_2 \to \mathcal{F}_3 \to 0$
une suite exacte courte de modules cohérents sur $X$. Alors
$$
\chi(X, \mathcal{F}_2) = \chi(X, \mathcal{F}_1) + \chi(X, \mathcal{F}_3)
$$
\end{lemma}

\begin{proof}
Considérons la suite exacte longue de cohomologie
$$
0 \to H^0(X, \mathcal{F}_1) \to H^0(X, \mathcal{F}_2) \to
H^0(X, \mathcal{F}_3) \to H^1(X, \mathcal{F}_1) \to \ldots
$$
associée à la suite exacte courte du lemme. Le théorème du rang
en algèbre linéaire montre que
$$
0 = \dim H^0(X, \mathcal{F}_1) - \dim H^0(X, \mathcal{F}_2)
+ \dim H^0(X, \mathcal{F}_3) - \dim H^1(X, \mathcal{F}_1) + \ldots
$$
Cela implique immédiatement le lemme.
\end{proof}

\begin{lemma}
\label{spaces-over-fields-lemma-euler-characteristic-morphism}
Soit $k$ un corps. Soit $f : Y \to X$ un morphisme d'
espaces algébriques propres sur $k$. Soit $\mathcal{G}$ un
$\mathcal{O}_Y$-module cohérent. Alors
$$
\chi(Y, \mathcal{G}) = \sum (-1)^i \chi(X, R^if_*\mathcal{G})
$$
\end{lemma}

\begin{proof}
La formule a un sens : les faisceaux $R^if_*\mathcal{G}$ sont cohérents
et ceux qui sont non nuls sont en nombre fini ; voir les
lemmes de Cohomologie des espaces
\ref{spaces-cohomology-lemma-proper-pushforward-coherent} et
\ref{spaces-cohomology-lemma-vanishing-higher-direct-images}.
D'après le lemme de Cohomologie sur les sites \ref{sites-cohomology-lemma-Leray},
il existe une suite spectrale telle que
$$
E_2^{p, q} = H^p(X, R^qf_*\mathcal{G})
$$
qui aboutit à $H^{p + q}(Y, \mathcal{G})$. La finitude de la cohomologie
sur $X$ montre que les $E_2^{p, q}$ non nuls sont en nombre fini
et que chaque $E_2^{p, q}$ est un espace vectoriel de dimension finie. Il s'ensuit
qu'il en va de même pour $E_r^{p, q}$ lorsque $r \geq 2$ et que
$$
\sum (-1)^{p + q} \dim_k E_r^{p, q}
$$
est indépendante de $r$. Comme, pour $r$ assez grand, nous avons
$E_r^{p, q} = E_\infty^{p, q}$ et que la convergence signifie qu'il
existe sur $H^n(Y, \mathcal{G})$ une filtration dont les quotients gradués sont
les $E_\infty^{p, q}$ avec $p + 1 = n$ (c'est le sens de la convergence
de la suite spectrale), nous concluons.
\end{proof}








\section{Nombres d'intersection}
\label{spaces-over-fields-section-num}

\noindent
Dans cette section, nous étudions la caractéristique d'Euler-Poincaré des
faisceaux cohérents sur les espaces algébriques propres afin d'obtenir des
nombres d'intersection pour les modules inversibles. Notre outil principal sera le
lemme suivant.

\begin{lemma}
\label{spaces-over-fields-lemma-numerical-polynomial-from-euler}
Soit $k$ un corps. Soit $X$ un espace algébrique propre sur $k$.
Soit $\mathcal{F}$ un $\mathcal{O}_X$-module cohérent. Soient
$\mathcal{L}_1, \ldots, \mathcal{L}_r$ des $\mathcal{O}_X$-modules inversibles.
L'application
$$
(n_1, \ldots, n_r) \longmapsto
\chi(X, \mathcal{F} \otimes
\mathcal{L}_1^{\otimes n_1} \otimes \ldots \otimes
\mathcal{L}_r^{\otimes n_r})
$$
est un polynôme numérique en $n_1, \ldots, n_r$ de degré total au
plus égal à la dimension du support schématique de $\mathcal{F}$.
\end{lemma}

\begin{proof}
Soit $Z \subset X$ le support schématique de $\mathcal{F}$.
Alors $\mathcal{F} = i_*\mathcal{G}$ pour un
$\mathcal{O}_Z$-module cohérent $\mathcal{G}$
(lemme de Cohomologie des espaces
\ref{spaces-cohomology-lemma-coherent-support-closed})
et nous avons
$$
\chi(X, \mathcal{F} \otimes
\mathcal{L}_1^{\otimes n_1} \otimes \ldots \otimes
\mathcal{L}_r^{\otimes n_r}) =
\chi(Z, \mathcal{G} \otimes
i^*\mathcal{L}_1^{\otimes n_1} \otimes \ldots \otimes
i^*\mathcal{L}_r^{\otimes n_r})
$$
d'après la formule de projection
(lemme de Cohomologie sur les sites \ref{sites-cohomology-lemma-projection-formula})
et le lemme de Cohomologie des espaces
\ref{spaces-cohomology-lemma-relative-affine-cohomology}.
Comme $|Z| = \text{Supp}(\mathcal{F})$, il suffit
de montrer que
$$
P_\mathcal{F}(n_1, \ldots, n_r) :
(n_1, \ldots, n_r)
\longmapsto
\chi(X, \mathcal{F} \otimes
\mathcal{L}_1^{\otimes n_1} \otimes \ldots \otimes
\mathcal{L}_r^{\otimes n_r})
$$
est un polynôme numérique en $n_1, \ldots, n_r$ de degré total au
plus égal à $\dim(X)$. Disons que la propriété $\mathcal{P}$ est vérifiée par le
$\mathcal{O}_X$-module cohérent $\mathcal{F}$ si l'assertion ci-dessus est vraie.

\medskip\noindent
Nous démontrerons cet énoncé par dévissage ; plus précisément, nous
vérifierons que les conditions (1), (2) et (3) du
lemme de Cohomologie des espaces
\ref{spaces-cohomology-lemma-property-higher-rank-cohomological-variant}
sont satisfaites.

\medskip\noindent
Vérification de la condition (1). Soit
$$
0 \to \mathcal{F}_1 \to \mathcal{F}_2 \to \mathcal{F}_3 \to 0
$$
une suite exacte courte de faisceaux cohérents sur $X$.
D'après le lemme \ref{spaces-over-fields-lemma-euler-characteristic-additive}, nous avons
$$
P_{\mathcal{F}_2}(n_1, \ldots, n_r) =
P_{\mathcal{F}_1}(n_1, \ldots, n_r) +
P_{\mathcal{F}_3}(n_1, \ldots, n_r)
$$
Il est alors clair que, si deux des trois faisceaux $\mathcal{F}_i$
vérifient la propriété $\mathcal{P}$, le troisième la vérifie aussi.

\medskip\noindent
La condition (2) résulte de l'égalité
$P_{\mathcal{F}^{\oplus m}}(n_1, \ldots, n_r) =
mP_\mathcal{F}(n_1, \ldots, n_r)$.

\medskip\noindent
Démonstration de (3). Soit $i : Z \to X$ un sous-espace fermé réduit tel que
$|Z|$ soit irréductible. Nous devons trouver un module cohérent $\mathcal{G}$
sur $X$, de support $Z$, tel que $\mathcal{P}$ soit vérifiée pour $\mathcal{G}$.
Nous donnerons deux constructions : l'une au moyen du lemme de Chow, l'autre
au moyen d'un revêtement fini par un schéma.

\medskip\noindent
Existence de $\mathcal{G}$ au moyen d'un revêtement fini par un schéma.
Choisissons $\pi : Z' \to Z$ fini surjectif, où $Z'$ est un schéma ; voir la
proposition de Limites d'espaces
\ref{spaces-limits-proposition-there-is-a-scheme-finite-over}.
Posons $\mathcal{G} = i_*\pi_*\mathcal{O}_{Z'} = (i \circ \pi)_*\mathcal{O}_{Z'}$.
Remarquons que $Z'$ est propre sur $k$ et que le support de $\mathcal{G}$ est $Y$
(détails omis). Nous avons
$$
R(\pi \circ i)_*(\mathcal{O}_{Z'}) = \mathcal{G}
\quad\text{et}\quad
R(\pi \circ i)_*(\pi^*i^*(\mathcal{L}_1^{\otimes n_1} \otimes \ldots \otimes
\mathcal{L}_r^{\otimes n_r})
) = \mathcal{G} \otimes \mathcal{L}_1^{\otimes n_1} \otimes \ldots \otimes
\mathcal{L}_r^{\otimes n_r}
$$
La première égalité vaut parce que $i \circ \pi$ est affine
(lemme de Cohomologie des espaces
\ref{spaces-cohomology-lemma-affine-vanishing-higher-direct-images})
et la seconde résulte de la première et de la formule de projection
(lemme de Cohomologie sur les sites \ref{sites-cohomology-lemma-projection-formula}).
En utilisant Leray
(lemme de Cohomologie sur les sites \ref{sites-cohomology-lemma-apply-Leray}),
nous obtenons
$$
P_\mathcal{G}(n_1, \ldots, n_r) =
\chi(Z', \pi^*i^*(\mathcal{L}_1^{\otimes n_1} \otimes \ldots \otimes
\mathcal{L}_r^{\otimes n_r}))
$$
D'après le cas des schémas
(lemme de Variétés \ref{varieties-lemma-numerical-polynomial-from-euler}),
il s'agit d'un polynôme numérique en
$n_1, \ldots, n_r$ de degré au plus égal à $\dim(Z')$.
Nous concluons puisque $\dim(Z') \leq \dim(Z) \leq \dim(X)$.
La première inégalité résulte du
lemme d'Espaces décents \ref{decent-spaces-lemma-dimension-quasi-finite}.

\medskip\noindent
Existence de $\mathcal{G}$ au moyen du lemme de Chow. Nous appliquons le
lemme de Cohomologie des espaces \ref{spaces-cohomology-lemma-weak-chow}
au morphisme $Z \to \Spec(k)$. Nous obtenons ainsi un
morphisme propre surjectif $f : Y \to Z$ sur $\Spec(k)$,
où $Y$ est un sous-schéma fermé de $\mathbf{P}^m_k$ pour un certain $m$.
Après remplacement de $Y$ par un sous-schéma fermé, nous pouvons supposer que $Y$
est intègre et que $f : Y \to Z$ est une altération ; voir le
lemme \ref{spaces-over-fields-lemma-alteration-contained-in}.
Notons $\mathcal{O}_Y(n)$ le tiré en arrière de $\mathcal{O}_{\mathbf{P}^m_k}(n)$.
Choisissons $n > 0$ tel que $R^pf_*\mathcal{O}_Y(n) = 0$
pour $p > 0$ ; voir le lemme de Cohomologie des espaces
\ref{spaces-cohomology-lemma-kill-by-twisting}.
Nous affirmons que $\mathcal{G} = i_*f_*\mathcal{O}_Y(n)$ vérifie $\mathcal{P}$.
En effet, le cas des schémas
(lemme de Variétés \ref{varieties-lemma-numerical-polynomial-from-euler})
montre que
$$
(n_1, \ldots, n_r)
\longmapsto
\chi(Y, \mathcal{O}_Y(n) \otimes
f^*i^*(\mathcal{L}_1^{\otimes n_1} \otimes \ldots \otimes
\mathcal{L}_r^{\otimes n_r}))
$$
est un polynôme numérique en $n_1, \ldots, n_r$ de degré total au
plus égal à $\dim(Y)$. D'autre part, d'après la formule de projection
(lemme de Cohomologie sur les sites \ref{sites-cohomology-lemma-projection-formula}),
\begin{align*}
i_*Rf_*\left(
\mathcal{O}_Y(n) \otimes
f^*i^*(\mathcal{L}_1^{\otimes n_1} \otimes \ldots \otimes
\mathcal{L}_r^{\otimes n_r})\right)
& =
i_*Rf_*\mathcal{O}_Y(n) \otimes
\mathcal{L}_1^{\otimes n_1} \otimes \ldots \otimes
\mathcal{L}_r^{\otimes n_r} \\
& =
\mathcal{G} \otimes \mathcal{L}_1^{\otimes n_1} \otimes \ldots \otimes
\mathcal{L}_r^{\otimes n_r}
\end{align*}
la dernière égalité résultant de notre choix de $n$. D'après
Leray (lemme de Cohomologie sur les sites \ref{sites-cohomology-lemma-apply-Leray}),
nous obtenons
$$
\chi(Y, \mathcal{O}_Y(n) \otimes
f^*i^*(\mathcal{L}_1^{\otimes n_1} \otimes \ldots \otimes
\mathcal{L}_r^{\otimes n_r})) =
P_\mathcal{G}(n_1, \ldots, n_r)
$$
et nous concluons puisque $\dim(Y) \leq \dim(Z) \leq \dim(X)$.
La première inégalité résulte du
lemme de Morphismes d'espaces
\ref{spaces-morphisms-lemma-alteration-dimension-general}
et du fait que $Y \to Z$ est une altération (de sorte que
l'extension induite des corps résiduels aux points génériques est finie).
\end{proof}

\noindent
Le lemme suivant montre, en gros, que le coefficient dominant ne dépend que
de la longueur du module cohérent aux points génériques de son
support.

\begin{lemma}
\label{spaces-over-fields-lemma-numerical-polynomial-leading-term}
Soit $k$ un corps. Soit $X$ un espace algébrique propre sur $k$. Soit
$\mathcal{F}$ un $\mathcal{O}_X$-module cohérent. Soient
$\mathcal{L}_1, \ldots, \mathcal{L}_r$ des $\mathcal{O}_X$-modules inversibles.
Posons $d = \dim(\text{Supp}(\mathcal{F}))$.
Soient $Z_i \subset X$ les composantes irréductibles
de $\text{Supp}(\mathcal{F})$ de dimension $d$. Soit $\overline{x}_i$
un point générique géométrique de $Z_i$ et posons
$m_i = \text{longueur}_{\mathcal{O}_{X, \overline{x}_i}}
(\mathcal{F}_{\overline{x}_i})$.
Alors
$$
\chi(X, \mathcal{F} \otimes \mathcal{L}_1^{\otimes n_1} \otimes \ldots \otimes
\mathcal{L}_r^{\otimes n_r}) -
\sum\nolimits_i
m_i\ \chi(Z_i, \mathcal{L}_1^{\otimes n_1} \otimes \ldots \otimes
\mathcal{L}_r^{\otimes n_r}|_{Z_i})
$$
est un polynôme numérique en $n_1, \ldots, n_r$ de degré total $< d$.
\end{lemma}

\begin{proof}
Démontrons d'abord un énoncé légèrement plus faible. Posons
$\dim(X) = N$ et soient $X_i \subset X$ les composantes
irréductibles de dimension $N$. Soit $\overline{x}_i$ un point
générique géométrique de $X_i$. L'anneau local \'etale
$\mathcal{O}_{X, \overline{x}_i}$ est noethérien de
dimension $0$ ; ainsi, pour tout $\mathcal{O}_X$-module cohérent
$\mathcal{F}$, la longueur
$$
m_i(\mathcal{F}) = \text{longueur}_{\mathcal{O}_{X, \overline{x}_i}}
(\mathcal{F}_{\overline{x}_i})
$$
est un entier $\geq 0$. Nous affirmons que
$$
E(\mathcal{F}) =
\chi(X, \mathcal{F} \otimes \mathcal{L}_1^{\otimes n_1} \otimes \ldots \otimes
\mathcal{L}_r^{\otimes n_r}) -
\sum\nolimits_i
m_i(\mathcal{F})\ \chi(Z_i, \mathcal{L}_1^{\otimes n_1} \otimes \ldots \otimes
\mathcal{L}_r^{\otimes n_r}|_{Z_i})
$$
est un polynôme numérique en $n_1, \ldots, n_r$ de degré total $< N$.
Nous le démontrerons au moyen du lemme de Cohomologie des espaces
\ref{spaces-cohomology-lemma-property-higher-rank-cohomological-variant}.
Pour toute suite exacte courte $0 \to \mathcal{F}' \to \mathcal{F} \to
\mathcal{F}'' \to 0$, nous avons
$E(\mathcal{F}) = E(\mathcal{F}') + E(\mathcal{F}'')$.
Cela résulte de l'additivité des caractéristiques d'Euler-Poincaré
(lemme \ref{spaces-over-fields-lemma-euler-characteristic-additive})
et de l'additivité des longueurs
(lemme d'Algèbre \ref{algebra-lemma-length-additive}).
Cela implique immédiatement les propriétés (1) et (2) du lemme de Cohomologie des espaces
\ref{spaces-cohomology-lemma-property-higher-rank-cohomological-variant}.
Enfin, la propriété (3) est vérifiée avec $\mathcal{G} = \mathcal{O}_Z$
pour tout sous-espace fermé réduit et irréductible $Z \subset X$.
En effet, si $Z = Z_{i_0}$ pour un certain $i_0$, alors
$m_i(\mathcal{G}) = \delta_{i_0i}$ et nous concluons que $E(\mathcal{G}) = 0$.
Si $Z \not = Z_i$ pour tout $i$, alors $m_i(\mathcal{G}) = 0$ pour tout $i$,
$\dim(Z) < N$ et le résultat découle du
lemme \ref{spaces-over-fields-lemma-numerical-polynomial-from-euler}.

\medskip\noindent
Démonstration de l'énoncé du lemme.
Soit $Z \subset X$ le support schématique de $\mathcal{F}$.
Alors $\mathcal{F} = i_*\mathcal{G}$ pour un
$\mathcal{O}_Z$-module cohérent $\mathcal{G}$
(lemme de Cohomologie des espaces
\ref{spaces-cohomology-lemma-coherent-support-closed})
et nous avons
$$
\chi(X, \mathcal{F} \otimes
\mathcal{L}_1^{\otimes n_1} \otimes \ldots \otimes
\mathcal{L}_r^{\otimes n_r}) =
\chi(Z, \mathcal{G} \otimes
i^*\mathcal{L}_1^{\otimes n_1} \otimes \ldots \otimes
i^*\mathcal{L}_r^{\otimes n_r})
$$
d'après la formule de projection
(lemme de Cohomologie sur les sites \ref{sites-cohomology-lemma-projection-formula})
et le lemme de Cohomologie des espaces
\ref{spaces-cohomology-lemma-relative-affine-cohomology}.
Comme $|Z| = \text{Supp}(\mathcal{F})$, nous avons $Z_i \subset Z$
pour tout $i$, et ce sont les composantes irréductibles
de $Z$ de dimension $d$. Nous considérons désormais $\overline{x}_i$
comme un point géométrique de $Z$. L'application
$i^\sharp : \mathcal{O}_X \to i_*\mathcal{O}_Z$
détermine une surjection
$$
\mathcal{O}_{X, \overline{x}_i} \to \mathcal{O}_{Z, \overline{x}_i}
$$
Cette application induit un isomorphisme de modules
$\mathcal{G}_{\overline{x}_i} = \mathcal{F}_{\overline{x}_i}$
puisque $\mathcal{F} = i_*\mathcal{G}$. Il en résulte que
$$
m_i = \text{longueur}_{\mathcal{O}_{X, \overline{x}_i}}
(\mathcal{F}_{\overline{x}_i}) =
\text{longueur}_{\mathcal{O}_{Z, \overline{x}_i}}
(\mathcal{G}_{\overline{x}_i})
$$
Nous voyons donc que l'expression du lemme est égale à
$$
\chi(Z, \mathcal{G} \otimes \mathcal{L}_1^{\otimes n_1} \otimes \ldots \otimes
\mathcal{L}_r^{\otimes n_r}) -
\sum\nolimits_i
m_i\ \chi(Z_i, \mathcal{L}_1^{\otimes n_1} \otimes \ldots \otimes
\mathcal{L}_r^{\otimes n_r}|_{Z_i})
$$
et le résultat découle de la discussion du premier paragraphe
(appliquée à $Z$ plutôt qu'à $X$).
\end{proof}

\begin{definition}
\label{spaces-over-fields-definition-intersection-number}
Soit $k$ un corps. Soit $X$ un espace algébrique propre sur $k$. Soit
$i : Z \to X$ un sous-espace fermé de dimension $d$. Soient
$\mathcal{L}_1, \ldots, \mathcal{L}_d$ des
$\mathcal{O}_X$-modules inversibles. Nous définissons le {\it nombre d'intersection}
$(\mathcal{L}_1 \cdots \mathcal{L}_d \cdot Z)$
comme le coefficient de $n_1 \ldots n_d$ dans le polynôme numérique
$$
\chi(X, i_*\mathcal{O}_Z \otimes \mathcal{L}_1^{\otimes n_1} \otimes
\ldots \otimes \mathcal{L}_d^{\otimes n_d}) =
\chi(Z, \mathcal{L}_1^{\otimes n_1} \otimes
\ldots \otimes \mathcal{L}_d^{\otimes n_d}|_Z)
$$
Dans le cas particulier
où $\mathcal{L}_1 = \ldots = \mathcal{L}_d = \mathcal{L}$,
nous écrivons $(\mathcal{L}^d \cdot Z)$.
\end{definition}

\noindent
L'égalité affichée dans la définition résulte de
la formule de projection
(Cohomologie, section \ref{cohomology-section-projection-formula}) et du
lemme de Cohomologie des schémas
\ref{coherent-lemma-relative-affine-cohomology}.
Démontrons quelques lemmes sur ces nombres d'intersection.

\begin{lemma}
\label{spaces-over-fields-lemma-intersection-number-integer}
Dans la situation de la définition \ref{spaces-over-fields-definition-intersection-number},
le nombre d'intersection
$(\mathcal{L}_1 \cdots \mathcal{L}_d \cdot Z)$
est un entier.
\end{lemma}

\begin{proof}
Tout polynôme numérique de degré $e$ en $n_1, \ldots, n_d$ s'écrit
de manière unique comme combinaison $\mathbf{Z}$-linéaire des fonctions
${n_1 \choose k_1}{n_2 \choose k_2} \ldots {n_d \choose k_d}$ avec
$k_1 + \ldots + k_d \leq e$. Appliquons ceci avec $e = d$.
La vérification est laissée en exercice.
\end{proof}

\begin{lemma}
\label{spaces-over-fields-lemma-intersection-number-additive}
Dans la situation de la définition \ref{spaces-over-fields-definition-intersection-number},
le nombre d'intersection
$(\mathcal{L}_1 \cdots \mathcal{L}_d \cdot Z)$
est additif : si $\mathcal{L}_i = \mathcal{L}_i' \otimes \mathcal{L}_i''$,
alors nous avons
$$
(\mathcal{L}_1 \cdots \mathcal{L}_i \cdots \mathcal{L}_d \cdot Z) =
(\mathcal{L}_1 \cdots \mathcal{L}_i' \cdots \mathcal{L}_d \cdot Z) +
(\mathcal{L}_1 \cdots \mathcal{L}_i'' \cdots \mathcal{L}_d \cdot Z)
$$
\end{lemma}

\begin{proof}
Cela résulte du lemme \ref{spaces-over-fields-lemma-numerical-polynomial-from-euler}, car
la fonction
$$
(n_1, \ldots, n_{i - 1}, n_i', n_i'', n_{i + 1}, \ldots, n_d)
\mapsto
\chi(Z, \mathcal{L}_1^{\otimes n_1} \otimes
\ldots \otimes (\mathcal{L}_i')^{\otimes n_i'} \otimes
(\mathcal{L}_i'')^{\otimes n_i''} \otimes \ldots \otimes
\mathcal{L}_d^{\otimes n_d}|_Z)
$$
est un polynôme numérique de degré total au plus égal à $d$ en $d + 1$ variables.
\end{proof}

\begin{lemma}
\label{spaces-over-fields-lemma-intersection-number-in-terms-of-components}
Dans la situation de la définition \ref{spaces-over-fields-definition-intersection-number},
soient $Z_i \subset Z$ les composantes irréductibles de dimension $d$. Soit
$m_i = \text{longueur}_{\mathcal{O}_{X, \overline{x}_i}}
(\mathcal{O}_{Z, \overline{x}_i})$
où $\overline{x}_i$ est un point générique géométrique de $Z_i$. Alors
$$
(\mathcal{L}_1 \cdots \mathcal{L}_d \cdot Z) =
\sum m_i(\mathcal{L}_1 \cdots \mathcal{L}_d \cdot Z_i)
$$
\end{lemma}

\begin{proof}
Cela résulte immédiatement du lemme \ref{spaces-over-fields-lemma-numerical-polynomial-leading-term}
et des définitions.
\end{proof}

\begin{lemma}
\label{spaces-over-fields-lemma-intersection-number-and-pullback}
Soit $k$ un corps. Soit $f : Y \to X$ un morphisme d'
espaces algébriques propres sur $k$.
Soit $Z \subset Y$ un sous-espace fermé intègre de dimension $d$ et soient
$\mathcal{L}_1, \ldots, \mathcal{L}_d$ des $\mathcal{O}_X$-modules inversibles.
Alors
$$
(f^*\mathcal{L}_1 \cdots f^*\mathcal{L}_d \cdot Z) =
\deg(f|_Z : Z \to f(Z)) (\mathcal{L}_1 \cdots \mathcal{L}_d \cdot f(Z))
$$
où $\deg(Z \to f(Z))$ est défini comme dans la définition \ref{spaces-over-fields-definition-degree},
ou vaut $0$ si $\dim(f(Z)) < d$.
\end{lemma}

\begin{proof}
Dans l'énoncé, $f(Z) \subset X$ est l'image schématique de $f$
et c'est aussi l'espace algébrique muni de la structure réduite induite sur la
partie fermée $f(|Z|) \subset X$ ; voir le lemme de Morphismes d'espaces
\ref{spaces-morphisms-lemma-scheme-theoretic-image-reduced}.
Alors $Z$ et $f(Z)$ sont des espaces algébriques réduits, propres (donc décents)
sur $k$, et par suite intègres
(définition \ref{spaces-over-fields-definition-integral-algebraic-space}).
Le membre de gauche se calcule comme le coefficient de $n_1 \ldots n_d$
dans la fonction
$$
\chi(Y, \mathcal{O}_Z \otimes f^*\mathcal{L}_1^{\otimes n_1} \otimes
\ldots \otimes f^*\mathcal{L}_d^{\otimes n_d}) =
\sum (-1)^i
\chi(X, R^if_*\mathcal{O}_Z \otimes
\mathcal{L}_1^{\otimes n_1} \otimes \ldots \otimes
\mathcal{L}_d^{\otimes n_d})
$$
L'égalité résulte du lemme \ref{spaces-over-fields-lemma-euler-characteristic-morphism}
et de la formule de projection
(lemme de Cohomologie \ref{cohomology-lemma-projection-formula}).
Si $f(Z)$ est de dimension $< d$, alors le membre de droite
est un polynôme de degré total $< d$ d'après le
lemme \ref{spaces-over-fields-lemma-numerical-polynomial-from-euler},
et le résultat est établi. Supposons $\dim(f(Z)) = d$. Alors
la théorie de la dimension (lemme \ref{spaces-over-fields-lemma-alteration-dimension}) montre
que les conditions équivalentes (1) -- (5) du
lemme \ref{spaces-over-fields-lemma-finite-degree} sont satisfaites. Ainsi,
$\deg(Z \to f(Z))$ est bien défini.
D'après le lemme \ref{spaces-over-fields-lemma-finite-degree}, déjà utilisé,
$f : Z \to f(Z)$ est fini au-dessus d'un ouvert non vide
$V$ de $f(Z)$ ; après avoir éventuellement rétréci $V$, nous pouvons supposer que
$V$ est un schéma. Soit $\xi \in V$ le point générique.
Ainsi, $\deg(f : Z \to f(Z))$ est la longueur du germe de
$f_*\mathcal{O}_Z$ en $\xi$ sur $\mathcal{O}_{X, \xi}$,
et le germe de $R^if_*\mathcal{O}_X$ en $\xi$ est nul pour $i > 0$
(par exemple d'après le lemme de Cohomologie des espaces
\ref{spaces-cohomology-lemma-finite-higher-direct-image-zero}).
Ainsi, les termes $\chi(X, R^if_*\mathcal{O}_Z \otimes
\mathcal{L}_1^{\otimes n_1} \otimes \ldots \otimes
\mathcal{L}_d^{\otimes n_d})$ pour $i > 0$
ont un degré total $< d$ et
$$
\chi(X, f_*\mathcal{O}_Z \otimes
\mathcal{L}_1^{\otimes n_1} \otimes \ldots \otimes
\mathcal{L}_d^{\otimes n_d})
=
\deg(f : Z \to f(Z)) \chi(f(Z),
\mathcal{L}_1^{\otimes n_1} \otimes \ldots \otimes
\mathcal{L}_d^{\otimes n_d}|_{f(Z)})
$$
à un polynôme de degré total $< d$ près, d'après le
lemme \ref{spaces-over-fields-lemma-numerical-polynomial-leading-term}.
Le résultat voulu en découle.
\end{proof}

\begin{lemma}
\label{spaces-over-fields-lemma-numerical-intersection-effective-Cartier-divisor}
Soit $k$ un corps. Soit $X$ un espace algébrique propre sur $k$.
Soit $Z \subset X$ un sous-espace fermé de dimension $d$.
Soient $\mathcal{L}_1, \ldots, \mathcal{L}_d$
des $\mathcal{O}_X$-modules inversibles. Supposons qu'il existe un
diviseur de Cartier effectif $D \subset Z$ tel que
$\mathcal{L}_1|_Z \cong \mathcal{O}_Z(D)$. Alors
$$
(\mathcal{L}_1 \cdots \mathcal{L}_d \cdot Z) =
(\mathcal{L}_2 \cdots \mathcal{L}_d \cdot D)
$$
\end{lemma}

\begin{proof}
Nous pouvons remplacer $X$ par $Z$ et $\mathcal{L}_i$ par $\mathcal{L}_i|_Z$.
Ainsi, nous pouvons supposer $X = Z$ et $\mathcal{L}_1 = \mathcal{O}_X(D)$.
Alors $\mathcal{L}_1^{-1}$ est le faisceau d'idéaux de $D$ et nous pouvons
considérer la suite exacte courte
$$
0 \to \mathcal{L}_1^{\otimes -1} \to \mathcal{O}_X \to \mathcal{O}_D \to 0
$$
Posons
$P(n_1, \ldots, n_d) =
\chi(X, \mathcal{L}_1^{\otimes n_1} \otimes \ldots \otimes
\mathcal{L}_d^{\otimes n_d})$
et
$Q(n_1, \ldots, n_d) =
\chi(D, \mathcal{L}_1^{\otimes n_1} \otimes \ldots \otimes
\mathcal{L}_d^{\otimes n_d}|_D)$.
L'additivité
(lemme \ref{spaces-over-fields-lemma-euler-characteristic-additive})
donne
$$
P(n_1, \ldots, n_d) - P(n_1 - 1, n_2, \ldots, n_d) =
Q(n_1, \ldots, n_d)
$$
Comme le degré total de $P$ est au plus égal à $d$, nous voyons que
le coefficient de $n_1 \ldots n_d$ dans $P$ est égal au coefficient
de $n_2 \ldots n_d$ dans $Q$.
\end{proof}












% OK, start here.
%

\chapter{Topologies sur les espaces algébriques}


\phantomsection
\label{spaces-topologies-section-phantom}





\section{Introduction}
\label{spaces-topologies-section-introduction}

\noindent
Dans ce chapitre, nous introduisons quelques topologies sur la
catégorie des espaces algébriques. Comparer avec le contenu de \cite{SGA1},
\cite{Ner}, \cite{LM-B} et \cite{Kn}.
Avant cela, signalons que
de nombreux choix de sites (au sens de
Sites, définition \ref{sites-definition-site}) donnent
la même notion de faisceau sur la catégorie sous-jacente. Ainsi,
nos choix peuvent différer légèrement de ceux des références,
mais conduisent finalement aux mêmes groupes de cohomologie, etc.




\section{Le procédé général}
\label{spaces-topologies-section-procedure}

\noindent
Dans cette section, nous expliquons un procédé général pour construire les
sites avec lesquels nous travaillerons. Cette discussion n'aura guère de
sens si le lecteur n'a pas lu
Topologies, section \ref{topologies-section-procedure}.

\medskip\noindent
Soit $S$ un schéma de base.
Prenons une catégorie $\Sch_\alpha$ construite comme dans
Ensembles, lemme \ref{sets-lemma-construct-category} en partant de
$S$ et de tout ensemble de schémas sur $S$ que l'on souhaite inclure.
Choisissons un ensemble quelconque de
recouvrements $\text{Cov}_{fppf}$ sur $\Sch_\alpha$, comme dans
Ensembles, lemme \ref{sets-lemma-coverings-site},
en partant de la catégorie $\Sch_\alpha$ et de la classe des recouvrements fppf.
Notons $\Sch_{fppf}$ le grand site fppf ainsi
obtenu, et notons $(\Sch/S)_{fppf}$ le
grand site fppf de $S$. (Tout ceci est exactement prescrit dans Topologies,
section \ref{topologies-section-fppf}).

\medskip\noindent
Pour les choix ci-dessus, la catégorie des espaces algébriques sur $S$
possède un ensemble de classes d'isomorphisme. Pour le voir, on peut utiliser le
fait que tout espace algébrique sur $S$ est de la forme $U/R$ pour
une certaine relation d'équivalence \'etale $j : R \to U \times_S U$ avec
$U, R \in \Ob((\Sch/S)_{fppf})$, voir
Espaces, lemme \ref{spaces-lemma-space-presentation}.
On peut donc trouver une sous-catégorie pleine $\textit{Espaces}/S$ de la catégorie des
espaces algébriques sur $S$ qui possède un ensemble d'objets
tel que tout espace algébrique soit isomorphe à un objet de
$\textit{Espaces}/S$. Nous fixons une telle catégorie.

\medskip\noindent
Dans les sections suivantes, pour une topologie $\tau$ donnée, le grand site
$(\textit{Espaces}/S)_\tau$ (resp.\ le grand site $(\textit{Espaces}/X)_\tau$
d'un espace algébrique $X$ sur $S$)
a pour catégorie sous-jacente la catégorie $\textit{Espaces}/S$
(resp.\ la sous-catégorie $\textit{Espaces}/X$ de $\textit{Espaces}/S$, voir
Catégories, exemple \ref{categories-example-category-over-X}).
Le procédé pour en faire un site consiste, comme d'habitude, à définir une
classe de recouvrements $\tau$ et à utiliser
Ensembles, lemme \ref{sets-lemma-coverings-site},
pour choisir un ensemble suffisamment grand de recouvrements qui définit la topologie.

\medskip\noindent
Signalons que le {\it petit site \'etale $X_\etale$
d'un espace algébrique $X$} a déjà été défini dans
Propriétés des espaces, définition
\ref{spaces-properties-definition-etale-site}.
Ses objets sont les schémas \'etales sur $X$, qui sont nombreux
par définition d'un espace algébrique. Cependant,
un site plus naturel du point de vue de ce chapitre (comparer
Topologies, définition \ref{topologies-definition-big-small-etale})
est le site $X_{spaces, \etale}$ de
Propriétés des espaces, définition
\ref{spaces-properties-definition-spaces-etale-site}.
Ces deux sites définissent le même topos ; voir
Propriétés des espaces, lemme \ref{spaces-properties-lemma-compare-etale-sites}.
Nous ne les redéfinirons pas dans ce chapitre ; nous les
utiliserons simplement.






\section{Topologie de Zariski}
\label{spaces-topologies-section-zariski}

\noindent
Dans
Espaces, section \ref{spaces-section-Zariski}
nous avons introduit la notion de recouvrement de Zariski d'un espace algébrique par
des sous-espaces ouverts. Voici la notion correspondante, où les sous-espaces ouverts
sont remplacés par des immersions ouvertes.

\begin{definition}
\label{spaces-topologies-definition-zariski-covering}
Soit $S$ un schéma et soit $X$ un espace algébrique sur $S$.
Un {\it recouvrement de Zariski de $X$} est une famille de morphismes
$\{f_i : X_i \to X\}_{i \in I}$ d'espaces algébriques sur $S$
telle que chaque $f_i$ soit une immersion ouverte
et telle que
$$
|X| = \bigcup\nolimits_{i \in I} |f_i|(|X_i|),
$$
c'est-à-dire que cette famille est surjective.
\end{definition}

\noindent
Bien que les recouvrements de Zariski soient parfois utiles, la topologie correspondante
sur la catégorie des espaces algébriques est beaucoup trop grossière et peu
utile. Elle définit néanmoins un site.

\begin{lemma}
\label{spaces-topologies-lemma-zariski}
Soit $S$ un schéma.
Soit $X$ un espace algébrique sur $S$.
\begin{enumerate}
\item Si $X' \to X$ est un isomorphisme, alors $\{X' \to X\}$
est un recouvrement de Zariski de $X$.
\item Si $\{X_i \to X\}_{i\in I}$ est un recouvrement de Zariski et si, pour chaque
$i$, nous avons un recouvrement de Zariski $\{X_{ij} \to X_i\}_{j\in J_i}$, alors
$\{X_{ij} \to X\}_{i \in I, j\in J_i}$ est un recouvrement de Zariski.
\item Si $\{X_i \to X\}_{i\in I}$ est un recouvrement de Zariski
et si $X' \to X$ est un morphisme d'espaces algébriques, alors
$\{X' \times_X X_i \to X'\}_{i\in I}$ est un recouvrement de Zariski.
\end{enumerate}
\end{lemma}

\begin{proof}
Démonstration omise.
\end{proof}










\section{Topologie \'etale}
\label{spaces-topologies-section-etale}

\noindent
Dans cette section, nous étudions la notion de recouvrement \'etale des
espaces algébriques et définissons le grand site \'etale d'un
espace algébrique. Comparer avec
Topologies, section \ref{topologies-section-etale}.

\begin{definition}
\label{spaces-topologies-definition-etale-covering}
Soit $S$ un schéma et soit $X$ un espace algébrique sur $S$.
Un {\it recouvrement \'etale de $X$} est une famille de morphismes
$\{f_i : X_i \to X\}_{i \in I}$ d'espaces algébriques sur $S$
telle que chaque $f_i$ soit \'etale
et telle que
$$
|X| = \bigcup\nolimits_{i \in I} |f_i|(|X_i|),
$$
c'est-à-dire que la famille des morphismes est surjective.
\end{definition}

\noindent
C'est exactement la même notion que celle de
Topologies, définition \ref{topologies-definition-etale-covering}.
En particulier, si $X$ et tous les $X_i$ sont des schémas, on retrouve la
notion usuelle de recouvrement \'etale de schémas.

\begin{lemma}
\label{spaces-topologies-lemma-zariski-etale}
Tout recouvrement de Zariski est un recouvrement \'etale.
\end{lemma}

\begin{proof}
Cela découle immédiatement des définitions et du fait qu'une immersion ouverte
est un morphisme \'etale (cela résulte de
Morphismes, lemme \ref{morphisms-lemma-open-immersion-etale}, via
Espaces, lemme
\ref{spaces-lemma-representable-transformations-property-implication}
puisque les immersions sont représentables).
\end{proof}

\begin{lemma}
\label{spaces-topologies-lemma-etale}
Soit $S$ un schéma.
Soit $X$ un espace algébrique sur $S$.
\begin{enumerate}
\item Si $X' \to X$ est un isomorphisme, alors $\{X' \to X\}$
est un recouvrement \'etale de $X$.
\item Si $\{X_i \to X\}_{i\in I}$ est un recouvrement \'etale et si, pour chaque
$i$, nous avons un recouvrement \'etale $\{X_{ij} \to X_i\}_{j\in J_i}$, alors
$\{X_{ij} \to X\}_{i \in I, j\in J_i}$ est un recouvrement \'etale.
\item Si $\{X_i \to X\}_{i\in I}$ est un recouvrement \'etale
et si $X' \to X$ est un morphisme d'espaces algébriques, alors
$\{X' \times_X X_i \to X'\}_{i\in I}$ est un recouvrement \'etale.
\end{enumerate}
\end{lemma}

\begin{proof}
Démonstration omise.
\end{proof}

\noindent
Le lemme suivant montre que les sites
$(\textit{Espaces}/X)_\etale$ et $(\textit{Espaces}/X)_{smooth}$
ont la même catégorie de faisceaux.

\begin{lemma}
\label{spaces-topologies-lemma-etale-dominates-smooth}
Soit $S$ un schéma et soit $X$ un espace algébrique sur $S$.
Soit $\{X_i \to X\}_{i \in I}$ un recouvrement lisse de $X$.
Alors il existe un recouvrement \'etale $\{U_j \to X\}_{j \in J}$
de $X$ qui raffine $\{X_i \to X\}_{i \in I}$.
\end{lemma}

\begin{proof}
Choisissons d'abord un schéma $U$ et un morphisme \'etale surjectif $U \to X$.
Pour chaque $i$, choisissons un schéma $W_i$ et un morphisme \'etale surjectif
$W_i \to X_i$. Alors $\{W_i \to X\}_{i \in I}$ est un recouvrement lisse
qui raffine $\{X_i \to X\}_{i \in I}$. Ainsi,
$\{W_i \times_X U \to U\}_{i \in I}$ est un recouvrement lisse de schémas.
Par le lemme \ref{more-morphisms-lemma-etale-dominates-smooth} de Compléments sur les morphismes,
nous pouvons choisir un recouvrement \'etale $\{U_j \to U\}$ qui raffine
$\{W_i \times_X U \to U\}$. Alors $\{U_j \to X\}_{j \in J}$
est un recouvrement \'etale qui raffine $\{X_i \to X\}_{i \in I}$.
\end{proof}

\begin{definition}
\label{spaces-topologies-definition-big-etale-site}
Soit $S$ un schéma. On appelle grand site \'etale
{\it $(\textit{Espaces}/S)_\etale$} tout site construit comme suit :
\begin{enumerate}
\item Choisissons un grand site \'etale $(\Sch/S)_\etale$ comme dans
Topologies, section \ref{topologies-section-etale}.
\item Prenons pour catégorie sous-jacente la catégorie $\textit{Espaces}/S$
des espaces algébriques sur $S$ (voir la discussion de la
section \ref{spaces-topologies-section-procedure} expliquant qu'il s'agit d'un ensemble).
\item Choisissons un ensemble quelconque de recouvrements comme dans
Ensembles, lemme \ref{sets-lemma-coverings-site}, en partant de la
catégorie $\textit{Espaces}/S$ et de la classe des recouvrements \'etales
de la définition \ref{spaces-topologies-definition-etale-covering}.
\end{enumerate}
\end{definition}

\noindent
Après cette définition, on peut localiser pour obtenir le site \'etale
d'un espace algébrique.

\begin{definition}
\label{spaces-topologies-definition-big-small-etale}
Soit $S$ un schéma. Soit $(\textit{Espaces}/S)_\etale$ comme dans
la définition \ref{spaces-topologies-definition-big-etale-site}.
Soit $X$ un espace algébrique sur $S$, c'est-à-dire un objet de
$(\textit{Espaces}/S)_\etale$. Alors le grand site \'etale
{\it $(\textit{Espaces}/X)_\etale$} de $X$
est la localisation du site $(\textit{Espaces}/S)_\etale$
en $X$, introduite dans Sites, section \ref{sites-section-localize}.
\end{definition}

\noindent
Rappelons que, pour un espace algébrique $X$ sur $S$ comme dans
la définition, nous avons déjà défini les petits sites \'etales
$X_{spaces, \etale}$ et $X_\etale$, voir
Propriétés des espaces, section \ref{spaces-properties-section-etale-site}.
Nous identifierons tacitement les topos correspondants au moyen du
foncteur d'inclusion $X_\etale \subset X_{spaces, \etale}$
(Propriétés des espaces, lemme \ref{spaces-properties-lemma-compare-etale-sites})
et l'appellerons le petit topos \'etale de $X$.
Nous établissons ensuite des relations entre les topos
associés à ces sites.

\begin{lemma}
\label{spaces-topologies-lemma-put-in-T-etale}
Soit $S$ un schéma. Soit $f : Y \to X$ un morphisme dans
$(\textit{Espaces}/S)_\etale$. Le foncteur d'inclusion
$Y_{spaces, \etale} \to (\textit{Espaces}/X)_\etale$
est cocontinu et induit un morphisme de topos
$$
i_f :
\Sh(Y_\etale)
\longrightarrow
\Sh((\textit{Espaces}/X)_\etale)
$$
Pour un faisceau $\mathcal{G}$ sur $(\textit{Espaces}/X)_\etale$
nous avons la formule $(i_f^{-1}\mathcal{G})(U/Y) = \mathcal{G}(U/X)$.
Le foncteur $i_f^{-1}$ possède également un adjoint à gauche $i_{f, !}$ qui commute
aux produits fibrés et aux égalisateurs.
\end{lemma}

\begin{proof}
Notons $u : Y_{spaces, \etale} \to (\textit{Espaces}/X)_\etale$ le foncteur.
Autrement dit, étant donné un morphisme \'etale $j : U \to Y$ correspondant
à un objet de $Y_{spaces, \etale}$, posons
$u(U \to Y) = (f \circ j : U \to X)$.
La catégorie $Y_{spaces, \etale}$
possède des produits fibrés et des égalisateurs, et $u$ commute à ces constructions.
Il est immédiat que $u$ est cocontinu.
Le foncteur $u$ est également continu, car $u$ transforme les recouvrements en recouvrements et
commute aux produits fibrés. Le lemme découle donc de
Sites, lemmes \ref{sites-lemma-when-shriek}
et \ref{sites-lemma-preserve-equalizers}.
\end{proof}

\begin{lemma}
\label{spaces-topologies-lemma-at-the-bottom-etale}
Soit $S$ un schéma. Soit $X$ un objet de $(\textit{Espaces}/S)_\etale$.
Le foncteur d'inclusion $X_{spaces, \etale} \to (\textit{Espaces}/X)_\etale$
satisfait aux hypothèses de Sites, lemme \ref{sites-lemma-bigger-site},
et induit donc un morphisme de sites
$$
\pi_X :
(\textit{Espaces}/X)_\etale
\longrightarrow
X_{spaces, \etale}
$$
et un morphisme de topos
$$
i_X :
\Sh(X_\etale)
\longrightarrow
\Sh((\textit{Espaces}/X)_\etale)
$$
tel que $\pi_X \circ i_X = \text{id}$. De plus, $i_X = i_{\text{id}_X}$
avec $i_{\text{id}_X}$ comme dans le lemme \ref{spaces-topologies-lemma-put-in-T-etale}.
En particulier, le foncteur $i_X^{-1} = \pi_{X, *}$ est décrit par la règle
$i_X^{-1}(\mathcal{G})(U/X) = \mathcal{G}(U/X)$.
\end{lemma}

\begin{proof}
Dans ce cas, le foncteur
$u : X_{spaces, \etale} \to (\textit{Espaces}/X)_\etale$,
en plus des propriétés établies dans la démonstration du
lemme \ref{spaces-topologies-lemma-put-in-T-etale} ci-dessus, est également pleinement fidèle
et transforme l'objet final en objet final.
Le lemme découle de Sites, lemme \ref{sites-lemma-bigger-site}.
\end{proof}

\begin{definition}
\label{spaces-topologies-definition-restriction-small-etale}
Dans la situation du lemme \ref{spaces-topologies-lemma-at-the-bottom-etale},
le foncteur $i_X^{-1} = \pi_{X, *}$ est souvent
appelé la {\it restriction au petit site \'etale}, et, pour un faisceau
$\mathcal{F}$ sur le grand site \'etale, nous notons souvent
$\mathcal{F}|_{X_\etale}$ cette restriction.
\end{definition}

\noindent
Avec cette notation, nous avons, pour un faisceau $\mathcal{F}$ sur le
grand site et un faisceau $\mathcal{G}$ sur le petit site,
\begin{align*}
\Mor_{\Sh(X_\etale)}(
\mathcal{F}|_{X_\etale},
\mathcal{G})
& =
\Mor_{\Sh((\textit{Espaces}/X)_\etale)}(
\mathcal{F},
i_{X, *}\mathcal{G}) \\
\Mor_{\Sh(X_\etale)}(
\mathcal{G},
\mathcal{F}|_{X_\etale})
& =
\Mor_{\Sh((\textit{Espaces}/X)_\etale)}(
\pi_X^{-1}\mathcal{G},
\mathcal{F})
\end{align*}
De plus, $(i_{X, *}\mathcal{G})|_{X_\etale} = \mathcal{G}$
et $(\pi_X^{-1}\mathcal{G})|_{X_\etale} = \mathcal{G}$.

\begin{lemma}
\label{spaces-topologies-lemma-morphism-big-etale}
Soit $S$ un schéma. Soit $f : Y \to X$ un morphisme dans
$(\textit{Espaces}/S)_\etale$. Le foncteur
$$
u :
(\textit{Espaces}/Y)_\etale
\longrightarrow
(\textit{Espaces}/X)_\etale,
\quad
V/Y \longmapsto V/X
$$
est cocontinu et admet un adjoint à droite continu
$$
v :
(\textit{Espaces}/X)_\etale
\longrightarrow
(\textit{Espaces}/Y)_\etale,
\quad
(U \to X) \longmapsto (U \times_X Y \to Y).
$$
Ils induisent le même morphisme de topos
$$
f_{big} :
\Sh((\textit{Espaces}/Y)_\etale)
\longrightarrow
\Sh((\textit{Espaces}/X)_\etale)
$$
Nous avons $f_{big}^{-1}(\mathcal{G})(U/Y) = \mathcal{G}(U/X)$.
Nous avons $f_{big, *}(\mathcal{F})(U/X) = \mathcal{F}(U \times_X Y/Y)$.
De plus, $f_{big}^{-1}$ possède un adjoint à gauche $f_{big!}$ qui commute aux
produits fibrés et aux égalisateurs.
\end{lemma}

\begin{proof}
Le foncteur $u$ est cocontinu, continu et commute aux produits fibrés
et aux égalisateurs (détails omis ; comparer avec la démonstration du
lemme \ref{spaces-topologies-lemma-put-in-T-etale}).
Par conséquent,
les lemmes de Sites \ref{sites-lemma-when-shriek} et
\ref{sites-lemma-preserve-equalizers}
s'appliquent et nous obtenons la formule
pour $f_{big}^{-1}$ ainsi que l'existence de $f_{big!}$. De plus,
le foncteur $v$ est adjoint à droite du foncteur précédent car, pour $U/Y$ et $V/X$,
nous avons $\Mor_X(u(U), V) = \Mor_Y(U, V \times_X Y)$
comme souhaité. Nous pouvons donc appliquer
les lemmes de Sites \ref{sites-lemma-have-functor-other-way} et
\ref{sites-lemma-have-functor-other-way-morphism} pour obtenir la
formule de $f_{big, *}$.
\end{proof}

\begin{lemma}
\label{spaces-topologies-lemma-morphism-big-small-etale}
Soit $S$ un schéma. Soit $f : Y \to X$ un morphisme dans
$(\textit{Espaces}/S)_\etale$.
\begin{enumerate}
\item Nous avons $i_f = f_{big} \circ i_Y$, avec $i_f$ comme dans
le lemme \ref{spaces-topologies-lemma-put-in-T-etale} et $i_Y$ comme dans
le lemme \ref{spaces-topologies-lemma-at-the-bottom-etale}.
\item Le foncteur $X_{spaces, \etale} \to Y_{spaces, \etale}$,
$(U \to X) \mapsto (U \times_X Y \to Y)$ est continu et induit
un morphisme de sites
$$
f_{spaces, \etale} : Y_{spaces, \etale} \longrightarrow X_{spaces, \etale}
$$
Le morphisme correspondant de petits topos \'etales est noté
$$
f_{small} : \Sh(Y_\etale) \to \Sh(X_\etale)
$$
Nous avons $f_{small, *}(\mathcal{F})(U/X) = \mathcal{F}(U \times_X Y/Y)$.
\item Nous avons un diagramme commutatif de morphismes de sites
$$
\xymatrix{
Y_{spaces, \etale} \ar[d]_{f_{spaces, \etale}} &
(\textit{Espaces}/Y)_\etale \ar[d]^{f_{big}} \ar[l]^-{\pi_Y}\\
X_{spaces, \etale} &
(\textit{Espaces}/X)_\etale \ar[l]_-{\pi_X}
}
$$
de sorte que $f_{small} \circ \pi_Y = \pi_X \circ f_{big}$ comme morphismes de topos.
\item Nous avons $f_{small} = \pi_X \circ f_{big} \circ i_Y = \pi_X \circ i_f$.
\end{enumerate}
\end{lemma}

\begin{proof}
L'égalité $i_f = f_{big} \circ i_Y$ découle de l'égalité
$i_f^{-1} = i_Y^{-1} \circ f_{big}^{-1}$, qui est claire d'après
les descriptions de ces foncteurs ci-dessus.
Nous obtenons ainsi (1).

\medskip\noindent
Le foncteur $u : X_{spaces, \etale} \to Y_{spaces, \etale}$,
$u(U \to X) = (U \times_X Y \to Y)$ induit, comme on l'a montré,
un morphisme de sites et le morphisme correspondant de petits
topos \'etales ; voir
Propriétés des espaces, lemme
\ref{spaces-properties-lemma-functoriality-etale-site}. La description
de l'image directe est claire.

\medskip\noindent
La partie (3) découle du fait que $\pi_X$ et $\pi_Y$ sont donnés par les
foncteurs d'inclusion, tandis que $f_{spaces, \etale}$ et $f_{big}$ sont donnés par les
foncteurs de changement de base $U \mapsto U \times_X Y$.

\medskip\noindent
L'énoncé (4) découle de (3) par précomposition avec $i_Y$.
\end{proof}

\noindent
Dans la situation du lemme, en utilisant la terminologie de
la définition \ref{spaces-topologies-definition-restriction-small-etale},
nous avons, pour $\mathcal{F}$ un faisceau sur le grand site \'etale de $Y$,
$$
(f_{big, *}\mathcal{F})|_{X_\etale} =
f_{small, *}(\mathcal{F}|_{Y_\etale}),
$$
Cette égalité découle de la commutativité du diagramme de
sites du lemme, puisque la restriction au petit site \'etale de
$Y$, resp.\ $X$, est donnée par $\pi_{Y, *}$, resp.\ $\pi_{X, *}$.
Une formule analogue faisant intervenir les images inverses et les restrictions est fausse.

\begin{lemma}
\label{spaces-topologies-lemma-composition-etale}
Soit $S$ un schéma. Étant donnés des morphismes $f : X \to Y$, $g : Y \to Z$
dans $(\textit{Espaces}/S)_\etale$, nous avons
$g_{big} \circ f_{big} = (g \circ f)_{big}$ et
$g_{small} \circ f_{small} = (g \circ f)_{small}$.
\end{lemma}

\begin{proof}
Cela découle de la description explicite des foncteurs image directe
et image inverse sur les grands sites donnée par le
lemme \ref{spaces-topologies-lemma-morphism-big-etale}. Pour les foncteurs
sur les petits sites, cela découle de la description des
foncteurs image directe dans le lemme \ref{spaces-topologies-lemma-morphism-big-small-etale}.
\end{proof}

\begin{lemma}
\label{spaces-topologies-lemma-morphism-big-small-cartesian-diagram-etale}
Soit $S$ un schéma. Considérons un diagramme cartésien
$$
\xymatrix{
Y' \ar[r]_{g'} \ar[d]_{f'} & Y \ar[d]^f \\
X' \ar[r]^g & X
}
$$
dans $(\textit{Espaces}/S)_\etale$.
$i_g^{-1} \circ f_{big, *} = f'_{small, *} \circ (i_{g'})^{-1}$
et $g_{big}^{-1} \circ f_{big, *} = f'_{big, *} \circ (g'_{big})^{-1}$.
\end{lemma}

\begin{proof}
Puisque le diagramme est cartésien, nous avons, pour $U'/X'$,
$U' \times_{X'} Y' = U' \times_X Y$. Ainsi, les deux foncteurs
$i_g^{-1} \circ f_{big, *}$ et $f'_{small, *} \circ (i_{g'})^{-1}$
associent à tout faisceau $\mathcal{F}$ sur $(\textit{Espaces}/Y)_\etale$ le faisceau
$U' \mapsto \mathcal{F}(U' \times_{X'} Y')$ sur $X'_\etale$
(voir les lemmes \ref{spaces-topologies-lemma-put-in-T-etale} et \ref{spaces-topologies-lemma-morphism-big-etale}).
La seconde égalité se démontre de la même manière ou peut être
déduite du résultat très général de
Sites, lemme \ref{sites-lemma-localize-morphism}.
\end{proof}

\begin{remark}
\label{spaces-topologies-remark-change-topologies-ringed}
Les sites $(\textit{Espaces}/X)_\etale$ et $X_{spaces, \etale}$
sont munis de faisceaux structuraux. Pour le petit site \'etale, nous
l'avons vu dans Propriétés des espaces, section
\ref{spaces-properties-section-structure-sheaf}.
Le faisceau structural $\mathcal{O}$ sur le grand site \'etale
$(\textit{Espaces}/X)_\etale$ est défini en associant à un objet
$U$ les sections globales du faisceau structural de $U$.
Cela a un sens puisque $U$ est lui-même un espace algébrique
et possède donc un faisceau structural. Comme $\mathcal{O}_U$
est un faisceau sur le site \'etale de $U$, le préfaisceau $\mathcal{O}$
ainsi défini satisfait à la condition de faisceau pour les recouvrements de $U$, c'est-à-dire que
$\mathcal{O}$ est un faisceau.
Nous pouvons munir les morphismes $i_f$, $\pi_X$, $i_X$, $f_{small}$ et
$f_{big}$ définis ci-dessus de structures de morphismes de sites annelés, respectivement de topos annelés.
Examinons-les successivement.
\begin{enumerate}
\item Dans le lemme \ref{spaces-topologies-lemma-put-in-T-etale}, notons $\mathcal{O}$
le faisceau structural sur $(\textit{Espaces}/X)_\etale$.
Nous avons $(i_f^{-1}\mathcal{O})(U/Y) = \mathcal{O}_U(U) = \mathcal{O}_Y(U)$
par construction.
On obtient ainsi un isomorphisme $i_f^\sharp : i_f^{-1}\mathcal{O} \to \mathcal{O}_Y$.
\item Dans le lemme \ref{spaces-topologies-lemma-at-the-bottom-etale}, il a été noté
que $i_X$ est un cas particulier de $i_f$ avec $f = \text{id}_X$ ;
nous sommes donc ramenés au cas (1).
\item Dans le lemme \ref{spaces-topologies-lemma-at-the-bottom-etale}, le morphisme
$\pi_X$ vérifie $(\pi_{X, *}\mathcal{O})(U) = \mathcal{O}(U) =
\mathcal{O}_X(U)$. Cette égalité permet de définir
$\pi_X^\sharp : \mathcal{O}_X \to \pi_{X, *}\mathcal{O}$.
\item Dans le lemme \ref{spaces-topologies-lemma-morphism-big-small-etale},
l'extension de $f_{small}$ en un morphisme de topos annelés
a été discutée dans Propriétés des espaces, lemme
\ref{spaces-properties-lemma-morphism-ringed-topoi}.
\item Dans le lemme \ref{spaces-topologies-lemma-morphism-big-small-etale},
le foncteur $f_{big}^{-1}$ est simplement la restriction
via le foncteur d'inclusion
$(\textit{Espaces}/Y)_\etale \to (\textit{Espaces}/X)_\etale$.
Soit $\mathcal{O}_1$ le faisceau structural sur $(\textit{Espaces}/X)_\etale$
et soit $\mathcal{O}_2$ le faisceau structural sur $(\textit{Espaces}/Y)_\etale$.
Nous obtenons un isomorphisme canonique
$f_{big}^\sharp : f_{big}^{-1}\mathcal{O}_1 \to \mathcal{O}_2$.
\end{enumerate}
De plus, ces définitions sont également compatibles avec les compositions.
Nous omettons l'énoncé et la démonstration détaillés.
\end{remark}









\section{Topologie lisse}
\label{spaces-topologies-section-smooth}

\noindent
Dans cette section, nous étudions la notion de recouvrement lisse
d'espaces algébriques et définissons le grand site lisse d'un
espace algébrique. Comparer avec
Topologies, section \ref{topologies-section-smooth}.

\begin{definition}
\label{spaces-topologies-definition-smooth-covering}
Soit $S$ un schéma et soit $X$ un espace algébrique sur $S$.
Un {\it recouvrement lisse de $X$} est une famille de morphismes
$\{f_i : X_i \to X\}_{i \in I}$ d'espaces algébriques sur $S$
telle que chaque $f_i$ soit lisse
et telle que
$$
|X| = \bigcup\nolimits_{i \in I} |f_i|(|X_i|),
$$
c'est-à-dire que cette famille est surjective.
\end{definition}

\noindent
C'est exactement la même notion que celle de
Topologies, définition \ref{topologies-definition-smooth-covering}.
En particulier, si $X$ et tous les $X_i$ sont des schémas, on retrouve la
notion usuelle de recouvrement lisse de schémas.

\begin{lemma}
\label{spaces-topologies-lemma-zariski-etale-smooth}
Tout recouvrement \'etale est un recouvrement lisse, et a fortiori,
tout recouvrement de Zariski est un recouvrement lisse.
\end{lemma}

\begin{proof}
Cela découle immédiatement des définitions, du fait qu'un morphisme \'etale
est lisse (Morphismes des espaces, lemme
\ref{spaces-morphisms-lemma-etale-smooth}) et du lemme
\ref{spaces-topologies-lemma-zariski-etale}.
\end{proof}

\begin{lemma}
\label{spaces-topologies-lemma-smooth}
Soit $S$ un schéma.
Soit $X$ un espace algébrique sur $S$.
\begin{enumerate}
\item Si $X' \to X$ est un isomorphisme, alors $\{X' \to X\}$
est un recouvrement lisse de $X$.
\item Si $\{X_i \to X\}_{i\in I}$ est un recouvrement lisse et si, pour chaque
$i$, nous avons un recouvrement lisse $\{X_{ij} \to X_i\}_{j\in J_i}$, alors
$\{X_{ij} \to X\}_{i \in I, j\in J_i}$ est un recouvrement lisse.
\item Si $\{X_i \to X\}_{i\in I}$ est un recouvrement lisse
et si $X' \to X$ est un morphisme d'espaces algébriques, alors
$\{X' \times_X X_i \to X'\}_{i\in I}$ est un recouvrement lisse.
\end{enumerate}
\end{lemma}

\begin{proof}
Démonstration omise.
\end{proof}

\noindent
À suivre...









\section{Topologie syntomique}
\label{spaces-topologies-section-syntomic}

\noindent
Dans cette section, nous étudions la notion de recouvrement syntomique
d'espaces algébriques et définissons le grand site syntomique d'un
espace algébrique. Comparer avec
Topologies, section \ref{topologies-section-syntomic}.

\begin{definition}
\label{spaces-topologies-definition-syntomic-covering}
Soit $S$ un schéma et soit $X$ un espace algébrique sur $S$.
Un {\it recouvrement syntomique de $X$} est une famille de morphismes
$\{f_i : X_i \to X\}_{i \in I}$ d'espaces algébriques sur $S$
telle que chaque $f_i$ soit syntomique
et telle que
$$
|X| = \bigcup\nolimits_{i \in I} |f_i|(|X_i|),
$$
c'est-à-dire que cette famille est surjective.
\end{definition}

\noindent
C'est exactement la même notion que celle de
Topologies, définition \ref{topologies-definition-syntomic-covering}.
En particulier, si $X$ et tous les $X_i$ sont des schémas, on retrouve la
notion usuelle de recouvrement syntomique de schémas.

\begin{lemma}
\label{spaces-topologies-lemma-zariski-etale-smooth-syntomic}
Tout recouvrement lisse est un recouvrement syntomique, et a fortiori,
tout recouvrement \'etale ou de Zariski est un recouvrement syntomique.
\end{lemma}

\begin{proof}
Cela découle immédiatement des définitions, du fait qu'un morphisme lisse
est syntomique (Morphismes des espaces, lemme
\ref{spaces-morphisms-lemma-smooth-syntomic}) et du lemme
\ref{spaces-topologies-lemma-zariski-etale-smooth}.
\end{proof}

\begin{lemma}
\label{spaces-topologies-lemma-syntomic}
Soit $S$ un schéma.
Soit $X$ un espace algébrique sur $S$.
\begin{enumerate}
\item Si $X' \to X$ est un isomorphisme, alors $\{X' \to X\}$
est un recouvrement syntomique de $X$.
\item Si $\{X_i \to X\}_{i\in I}$ est un recouvrement syntomique et si, pour chaque
$i$, nous avons un recouvrement syntomique $\{X_{ij} \to X_i\}_{j\in J_i}$, alors
$\{X_{ij} \to X\}_{i \in I, j\in J_i}$ est un recouvrement syntomique.
\item Si $\{X_i \to X\}_{i\in I}$ est un recouvrement syntomique
et si $X' \to X$ est un morphisme d'espaces algébriques, alors
$\{X' \times_X X_i \to X'\}_{i\in I}$ est un recouvrement syntomique.
\end{enumerate}
\end{lemma}

\begin{proof}
Démonstration omise.
\end{proof}

\noindent
À suivre...










\section{Topologie fppf}
\label{spaces-topologies-section-fppf}

\noindent
Dans cette section, nous étudions la notion de recouvrement fppf d'espaces algébriques,
et définissons le grand site fppf d'un espace algébrique. Comparer avec
Topologies, section \ref{topologies-section-fppf}.

\begin{definition}
\label{spaces-topologies-definition-fppf-covering}
Soit $S$ un schéma et soit $X$ un espace algébrique sur $S$.
Un {\it recouvrement fppf de $X$} est une famille de morphismes
$\{f_i : X_i \to X\}_{i \in I}$ d'espaces algébriques sur $S$
telle que chaque $f_i$ soit plat et localement de présentation finie
et telle que
$$
|X| = \bigcup\nolimits_{i \in I} |f_i|(|X_i|),
$$
c'est-à-dire que cette famille est surjective.
\end{definition}

\noindent
C'est exactement la même notion que celle de
Topologies, définition \ref{topologies-definition-fppf-covering}.
En particulier, si $X$ et tous les $X_i$ sont des schémas, on retrouve la notion usuelle
de recouvrement fppf de schémas.

\begin{lemma}
\label{spaces-topologies-lemma-zariski-etale-smooth-syntomic-fppf}
Tout recouvrement syntomique est un recouvrement fppf, et a fortiori,
tout recouvrement lisse, \'etale ou de Zariski est un recouvrement fppf.
\end{lemma}

\begin{proof}
Cela découle des définitions et du fait qu'un morphisme syntomique
est plat et localement de présentation finie
(Morphismes des espaces, lemmes
\ref{spaces-morphisms-lemma-syntomic-locally-finite-presentation} et
\ref{spaces-morphisms-lemma-syntomic-flat}) et
du lemme \ref{spaces-topologies-lemma-zariski-etale-smooth-syntomic}.
\end{proof}

\begin{lemma}
\label{spaces-topologies-lemma-fppf}
Soit $S$ un schéma.
Soit $X$ un espace algébrique sur $S$.
\begin{enumerate}
\item Si $X' \to X$ est un isomorphisme, alors $\{X' \to X\}$
est un recouvrement fppf de $X$.
\item Si $\{X_i \to X\}_{i\in I}$ est un recouvrement fppf et si, pour chaque
$i$, nous avons un recouvrement fppf $\{X_{ij} \to X_i\}_{j\in J_i}$, alors
$\{X_{ij} \to X\}_{i \in I, j\in J_i}$ est un recouvrement fppf.
\item Si $\{X_i \to X\}_{i\in I}$ est un recouvrement fppf
et si $X' \to X$ est un morphisme d'espaces algébriques, alors
$\{X' \times_X X_i \to X'\}_{i\in I}$ est un recouvrement fppf.
\end{enumerate}
\end{lemma}

\begin{proof}
Démonstration omise.
\end{proof}

\begin{lemma}
\label{spaces-topologies-lemma-refine-fppf-schemes}
Soit $S$ un schéma et soit $X$ un espace algébrique sur $S$.
Supposons que $\mathcal{U} = \{f_i : X_i \to X\}_{i \in I}$ soit un
recouvrement fppf de $X$. Alors il existe un raffinement
$\mathcal{V} = \{g_i : T_i \to X\}$ de $\mathcal{U}$ qui est un
recouvrement fppf tel que chaque $T_i$ soit un schéma.
\end{lemma}

\begin{proof}
Démonstration omise. Indication : pour chaque $i$, choisir un schéma $T_i$ et un morphisme \'etale surjectif
$T_i \to X_i$. Vérifier ensuite que $\{T_i \to X\}$ est un recouvrement fppf.
\end{proof}

\begin{lemma}
\label{spaces-topologies-lemma-fppf-covering-surjective}
Soit $S$ un schéma.
Soit $\{f_i : X_i \to X\}_{i \in I}$ un recouvrement fppf d'espaces algébriques
sur $S$. Alors le morphisme de faisceaux
$$
\coprod X_i \longrightarrow X
$$
est surjectif.
\end{lemma}

\begin{proof}
Cela découle du lemme
\ref{spaces-lemma-surjective-flat-locally-finite-presentation} du chapitre Espaces.
Voir également la remarque
\ref{spaces-remark-warning} du chapitre Espaces
si la signification de ce lemme vous semble obscure.
\end{proof}

\begin{definition}
\label{spaces-topologies-definition-big-fppf-site}
Soit $S$ un schéma. On appelle grand site fppf
{\it $(\textit{Espaces}/S)_{fppf}$} tout site construit comme suit :
\begin{enumerate}
\item Choisissons un grand site fppf $(\Sch/S)_{fppf}$ comme dans
Topologies, section \ref{topologies-section-fppf}.
\item Prenons pour catégorie sous-jacente la catégorie $\textit{Espaces}/S$
des espaces algébriques sur $S$ (voir la discussion de la
section \ref{spaces-topologies-section-procedure} expliquant qu'il s'agit d'un ensemble).
\item Choisissons un ensemble quelconque de recouvrements comme dans
Ensembles, lemme \ref{sets-lemma-coverings-site}, en partant de la
catégorie $\textit{Espaces}/S$ et de la classe des recouvrements fppf
de la définition \ref{spaces-topologies-definition-fppf-covering}.
\end{enumerate}
\end{definition}

\noindent
Après cette définition, on peut localiser pour obtenir le site fppf
d'un espace algébrique.

\begin{definition}
\label{spaces-topologies-definition-big-small-fppf}
Soit $S$ un schéma. Soit $(\textit{Espaces}/S)_{fppf}$ comme dans
la définition \ref{spaces-topologies-definition-big-fppf-site}.
Soit $X$ un espace algébrique sur $S$, c'est-à-dire un objet de
$(\textit{Espaces}/S)_{fppf}$. Alors le grand site fppf
{\it $(\textit{Espaces}/X)_{fppf}$} de $X$
est la localisation du site $(\textit{Espaces}/S)_{fppf}$
en $X$, introduite dans Sites, section \ref{sites-section-localize}.
\end{definition}

\noindent
Nous établissons ensuite quelques relations entre les topos
associés à ces sites.

\begin{lemma}
\label{spaces-topologies-lemma-morphism-big-fppf}
Soit $S$ un schéma.
Soit $f : Y \to X$ un morphisme d'espaces algébriques sur $S$.
Le foncteur
$$
u : (\textit{Espaces}/Y)_{fppf} \longrightarrow (\textit{Espaces}/X)_{fppf},
\quad
V/Y \longmapsto V/X
$$
est cocontinu et admet un adjoint à droite continu
$$
v : (\textit{Espaces}/X)_{fppf} \longrightarrow (\textit{Espaces}/Y)_{fppf},
\quad
(U \to X) \longmapsto (U \times_X Y \to Y).
$$
Ils induisent le même morphisme de topos
$$
f_{big} :
\Sh((\textit{Espaces}/Y)_{fppf})
\longrightarrow
\Sh((\textit{Espaces}/X)_{fppf})
$$
Nous avons $f_{big}^{-1}(\mathcal{G})(U/Y) = \mathcal{G}(U/X)$.
Nous avons $f_{big, *}(\mathcal{F})(U/X) = \mathcal{F}(U \times_X Y/Y)$.
De plus, $f_{big}^{-1}$ possède un adjoint à gauche $f_{big!}$ qui commute aux
produits fibrés et aux égalisateurs.
\end{lemma}

\begin{proof}
Le foncteur $u$ est cocontinu, continu et commute aux produits fibrés
et aux égalisateurs. Par conséquent,
les lemmes de Sites \ref{sites-lemma-when-shriek} et
\ref{sites-lemma-preserve-equalizers}
s'appliquent et nous obtenons la formule
pour $f_{big}^{-1}$ ainsi que l'existence de $f_{big!}$. De plus,
le foncteur $v$ est adjoint à droite du foncteur précédent car, pour $U/Y$ et $V/X$,
nous avons $\Mor_X(u(U), V) = \Mor_Y(U, V \times_X Y)$
comme souhaité. Nous pouvons donc appliquer
les lemmes de Sites \ref{sites-lemma-have-functor-other-way} et
\ref{sites-lemma-have-functor-other-way-morphism} pour obtenir la
formule de $f_{big, *}$.
\end{proof}

\begin{lemma}
\label{spaces-topologies-lemma-composition-fppf}
Soit $S$ un schéma. Étant donnés des morphismes $f : X \to Y$, $g : Y \to Z$
d'espaces algébriques sur $S$, nous avons
$g_{big} \circ f_{big} = (g \circ f)_{big}$.
\end{lemma}

\begin{proof}
Cela découle de la description explicite des foncteurs image directe
et image inverse sur les grands sites donnée par le
lemme \ref{spaces-topologies-lemma-morphism-big-fppf}.
\end{proof}








\section{La topologie ph}
\label{spaces-topologies-section-ph}

\noindent
Dans cette section, nous définissons la topologie ph. Il s'agit de la topologie
engendrée par les recouvrements \'etales et les morphismes propres surjectifs ; voir
le lemme \ref{spaces-topologies-lemma-characterize-sheaf}.

\begin{definition}
\label{spaces-topologies-definition-ph-covering}
Soit $S$ un schéma et soit $X$ un espace algébrique sur $S$.
Un {\it recouvrement ph de $X$} est une famille
de morphismes $\{f_i : X_i \to X\}_{i \in I}$ d'espaces algébriques sur $S$
telle que chaque $f_i$ soit localement de type fini et telle que, pour tout
$h : U \to X$ avec $U$ affine, il existe un recouvrement ph standard
$\{g_j : U_j \to U\}_{j = 1, \ldots, m}$ qui raffine la famille
$\{X_i \times_X U \to U\}_{i \in I}$.
\end{definition}

\noindent
Autrement dit, il existe des indices $i_1, \ldots, i_m \in I$ et des
morphismes $h_j : U_j \to X_{i_j}$ tels que
$f_{i_j} \circ h_j = h \circ g_j$. Notons que, si $X$ et tous les $X_i$ sont
représentables, il s'agit de la même notion qu'un recouvrement ph de schémas défini dans
Topologies, définition \ref{topologies-definition-ph-covering}.

\begin{lemma}
\label{spaces-topologies-lemma-zariski-etale-smooth-syntomic-fppf-ph}
Tout recouvrement fppf est un recouvrement ph, et a fortiori,
tout recouvrement syntomique, lisse, \'etale ou de Zariski est un recouvrement ph.
\end{lemma}

\begin{proof}
Montrons qu'un recouvrement fppf est un recouvrement ph ; le
reste découle alors du lemme \ref{spaces-topologies-lemma-zariski-etale-smooth-syntomic-fppf}.
Soit $\{X_i \to X\}_{i \in I}$ un recouvrement fppf d'espaces algébriques
sur un schéma de base $S$. Soit $U$ un schéma affine et soit
$U \to X$ un morphisme. Nous pouvons raffiner le recouvrement fppf
$\{X_i \times_X U \to U\}_{i \in I}$ par un recouvrement fppf
$\{T_i \to U\}_{i \in I}$ où chaque $T_i$ est un schéma
(lemme \ref{spaces-topologies-lemma-refine-fppf-schemes}).
Nous pouvons alors trouver un recouvrement ph standard $\{U_j \to U\}_{j = 1, \ldots, m}$
raffinant $\{T_i \to U\}_{i \in I}$, par le
lemme \ref{more-morphisms-lemma-fppf-ph} de Compléments sur les morphismes
(et par la définition des recouvrements ph pour les schémas).
Ainsi, $\{X_i \to X\}_{i \in I}$ est un recouvrement ph par définition.
\end{proof}

\begin{lemma}
\label{spaces-topologies-lemma-surjective-proper-ph}
Soit $S$ un schéma. Soit $f : Y \to X$ un morphisme propre surjectif
d'espaces algébriques sur $S$. Alors $\{Y \to X\}$ est un recouvrement ph.
\end{lemma}

\begin{proof}
Soit $U \to X$ un morphisme, où $U$ est affine.
Par le lemme de Chow (sous la forme faible donnée dans
Cohomologie des espaces, lemme \ref{spaces-cohomology-lemma-weak-chow}),
il existe un morphisme propre surjectif de schémas
$V \to U$ qui se factorise par $Y \times_X U \to U$.
En choisissant un recouvrement fini de $V$ par des ouverts affines, nous obtenons un
recouvrement ph standard de $U$ qui raffine $\{Y \times_X U \to U\}$,
comme souhaité.
\end{proof}

\begin{lemma}
\label{spaces-topologies-lemma-ph}
Soit $S$ un schéma. Soit $X$ un espace algébrique sur $S$.
\begin{enumerate}
\item Si $X' \to X$ est un isomorphisme, alors $\{X' \to X\}$
est un recouvrement ph de $X$.
\item Si $\{X_i \to X\}_{i\in I}$ est un recouvrement ph et si, pour tout
$i$, on a un recouvrement ph $\{X_{ij} \to X_i\}_{j\in J_i}$, alors
$\{X_{ij} \to X\}_{i \in I, j\in J_i}$ est un recouvrement ph.
\item Si $\{X_i \to X\}_{i\in I}$ est un recouvrement ph
et si $X' \to X$ est un morphisme d'espaces algébriques, alors
$\{X' \times_X X_i \to X'\}_{i\in I}$ est un recouvrement ph.
\end{enumerate}
\end{lemma}

\begin{proof}
La partie (1) est claire. Considérons $g : X' \to X$ et
un recouvrement ph $\{X_i \to X\}_{i\in I}$ comme dans (3). Par le
lemme \ref{spaces-morphisms-lemma-base-change-finite-type} de Morphismes des espaces,
les morphismes $X' \times_X X_i \to X'$ sont localement de type fini.
Si $h' : Z \to X'$ est un morphisme d'un schéma affine
vers $X'$, posons $h = g \circ h' : Z \to X$. L'hypothèse
sur $\{X_i \to X\}_{i\in I}$ signifie qu'il existe un recouvrement ph standard
$\{Z_j \to Z\}_{j = 1, \ldots, n}$ et des morphismes $Z_j \to X_{i(j)}$ au-dessus de
$h$, pour certains $i(j) \in I$. Par la propriété universelle du produit fibré,
nous obtenons aussi des morphismes $Z_j \to X' \times_X X_{i(j)}$ au-dessus de $h'$.
Ainsi, $\{X' \times_X X_i \to X'\}_{i\in I}$ est un recouvrement ph.
Ceci démontre (3).

\medskip\noindent
Soient $\{X_i \to X\}_{i\in I}$ et $\{X_{ij} \to X_i\}_{j\in J_i}$ comme
dans (2). Soit $h : Z \to X$ un morphisme d'un schéma affine vers $X$.
Par hypothèse, il existe un recouvrement ph standard
$\{Z_j \to Z\}_{j = 1, \ldots, n}$ et des morphismes $h_j : Z_j \to X_{i(j)}$
au-dessus de $h$, pour certains indices $i(j) \in I$. Par hypothèse, il existe des
recouvrements ph standard
$\{Z_{j, l} \to Z_j\}_{l = 1, \ldots, n(j)}$
et des morphismes $Z_{j, l} \to X_{i(j)j(l)}$ au-dessus de
$h_j$, pour certains indices $j(l) \in J_{i(j)}$. Par le
lemme \ref{topologies-lemma-refine-by-standard-ph} de Topologies,
la famille $\{Z_{j, l} \to Z\}$ peut être raffinée par un recouvrement ph standard.
Nous en concluons que $\{X_{ij} \to X\}_{i \in I, j\in J_i}$
est un recouvrement ph.
\end{proof}

\begin{definition}
\label{spaces-topologies-definition-big-ph-site}
Soit $S$ un schéma. On appelle grand site ph {\it $(\textit{Espaces}/S)_{ph}$}
tout site construit comme suit :
\begin{enumerate}
\item Choisissons un grand site ph $(\Sch/S)_{ph}$ comme dans
Topologies, section \ref{topologies-section-ph}.
\item Prenons pour catégorie sous-jacente la catégorie $\textit{Espaces}/S$
des espaces algébriques sur $S$ (voir la discussion de la
section \ref{spaces-topologies-section-procedure}, qui explique pourquoi il s'agit d'un ensemble).
\item Choisissons un ensemble quelconque de recouvrements comme dans le
lemme \ref{sets-lemma-coverings-site} d'Ensembles, en partant de la
catégorie $\textit{Espaces}/S$ et de la classe des recouvrements ph
de la définition \ref{spaces-topologies-definition-ph-covering}.
\end{enumerate}
\end{definition}

\noindent
Après cette définition, on peut localiser pour obtenir le site ph
d'un espace algébrique.

\begin{definition}
\label{spaces-topologies-definition-big-small-ph}
Soit $S$ un schéma. Soit $(\textit{Espaces}/S)_{ph}$ comme dans la
définition \ref{spaces-topologies-definition-big-ph-site}.
Soit $X$ un espace algébrique sur $S$, c'est-à-dire un objet de
$(\textit{Espaces}/S)_{ph}$. Alors le grand site ph
{\it $(\textit{Espaces}/X)_{ph}$} de $X$
est la localisation du site $(\textit{Espaces}/S)_{ph}$
en $X$, introduite dans Sites, section \ref{sites-section-localize}.
\end{definition}

\noindent
Voici la caractérisation annoncée des faisceaux ph.

\begin{lemma}
\label{spaces-topologies-lemma-characterize-sheaf}
Soit $S$ un schéma. Soit $X$ un espace algébrique sur $S$.
Soit $\mathcal{F}$ un préfaisceau sur $(\textit{Espaces}/X)_{ph}$.
Alors $\mathcal{F}$ est un faisceau si et seulement si
\begin{enumerate}
\item $\mathcal{F}$ satisfait à la condition de faisceau pour les recouvrements \'etales, et
\item si $f : V \to U$ est un morphisme propre surjectif de
$(\textit{Espaces}/X)_{ph}$, alors
$\mathcal{F}(U)$ s'envoie bijectivement sur l'égalisateur
des deux applications $\mathcal{F}(V) \to \mathcal{F}(V \times_U V)$.
\end{enumerate}
\end{lemma}

\begin{proof}
Montrons que, si (1) et (2) sont satisfaites, alors $\mathcal{F}$ est un faisceau.
Soit $\{T_i \to T\}$ un recouvrement ph, c'est-à-dire un recouvrement dans
$(\textit{Espaces}/X)_{ph}$.
Vérifions la condition de faisceau pour ce recouvrement.
Soient $s_i \in \mathcal{F}(T_i)$ des sections dont les restrictions sont égales
sur $T_i \times_T T_{i'}$. Montrons qu'il existe une
unique section $s \in \mathcal{F}$ dont la restriction est $s_i$ sur $T_i$.
Soit $\{U_j \to T\}$ un recouvrement \'etale avec les $U_j$ affines.
D'après la propriété (1), il suffit de produire des sections $s_j \in \mathcal{F}(U_j)$
qui coïncident sur $U_j \cap U_{j'}$ afin de construire $s$.
Considérons les recouvrements ph $\{T_i \times_T U_j \to U_j\}$.
Alors les $s_{ji} = s_i|_{T_i \times_T U_j}$ sont des sections qui coïncident
sur $(T_i \times_T U_j) \times_{U_j} (T_{i'} \times_T U_j)$.
Choisissons un morphisme propre surjectif $V_j \to U_j$ et un recouvrement
fini par des ouverts affines, $V_j = \bigcup V_{jk}$, tels que le recouvrement ph standard
$\{V_{jk} \to U_j\}$ raffine $\{T_i \times_T U_j \to U_j\}$.
Si $s_{jk} \in \mathcal{F}(V_{jk})$
désigne l'image inverse de $s_{ji}$ sur $V_{jk}$ par les
morphismes implicites, alors les $s_{jk}$ se recollent en une section
$s'_j \in \mathcal{F}(V_j)$. En utilisant encore une fois leur coïncidence sur les intersections,
nous voyons que $s'_j$ appartient à l'égalisateur des deux
applications $\mathcal{F}(V_j) \to \mathcal{F}(V_j \times_{U_j} V_j)$.
Par (2), la section $s'_j$ provient donc d'une unique section
$s_j \in \mathcal{F}(U_j)$. Nous omettons de vérifier que ces
sections $s_j$ possèdent toutes les propriétés voulues.
\end{proof}

\noindent
Nous établissons ensuite quelques relations entre les topos
associés à ces sites.

\begin{lemma}
\label{spaces-topologies-lemma-morphism-big-ph}
Soit $S$ un schéma.
Soit $f : Y \to X$ un morphisme d'espaces algébriques sur $S$.
Le foncteur
$$
u : (\textit{Espaces}/Y)_{ph} \longrightarrow (\textit{Espaces}/X)_{ph},
\quad
V/Y \longmapsto V/X
$$
est cocontinu et possède un adjoint à droite continu
$$
v : (\textit{Espaces}/X)_{ph} \longrightarrow (\textit{Espaces}/Y)_{ph},
\quad
(U \to X) \longmapsto (U \times_X Y \to Y).
$$
Ils induisent le même morphisme de topos
$$
f_{big} :
\Sh((\textit{Espaces}/Y)_{ph})
\longrightarrow
\Sh((\textit{Espaces}/X)_{ph})
$$
Nous avons $f_{big}^{-1}(\mathcal{G})(U/Y) = \mathcal{G}(U/X)$.
Nous avons $f_{big, *}(\mathcal{F})(U/X) = \mathcal{F}(U \times_X Y/Y)$.
De plus, $f_{big}^{-1}$ possède un adjoint à gauche $f_{big!}$ qui commute aux
produits fibrés et aux égalisateurs.
\end{lemma}

\begin{proof}
Le foncteur $u$ est cocontinu, continu et commute aux produits fibrés
et aux égalisateurs. Par conséquent,
les lemmes de Sites \ref{sites-lemma-when-shriek} et
\ref{sites-lemma-preserve-equalizers}
s'appliquent et nous en déduisons la formule
pour $f_{big}^{-1}$ ainsi que l'existence de $f_{big!}$. De plus,
le foncteur $v$ est un adjoint à droite car, pour $U/Y$ et $V/X$,
nous avons $\Mor_X(u(U), V) = \Mor_Y(U, V \times_X Y)$,
comme souhaité. Nous pouvons donc appliquer
les lemmes de Sites \ref{sites-lemma-have-functor-other-way} et
\ref{sites-lemma-have-functor-other-way-morphism} pour obtenir la
formule de $f_{big, *}$.
\end{proof}

\begin{lemma}
\label{spaces-topologies-lemma-composition-ph}
Soit $S$ un schéma. Étant donnés des morphismes $f : X \to Y$, $g : Y \to Z$
d'espaces algébriques sur $S$, nous avons
$g_{big} \circ f_{big} = (g \circ f)_{big}$.
\end{lemma}

\begin{proof}
Cela découle de la description simple des foncteurs image directe
et image inverse sur les grands sites donnée par le
lemme \ref{spaces-topologies-lemma-morphism-big-ph}.
\end{proof}

\begin{lemma}
\label{spaces-topologies-lemma-cech-enough}
Soit $S$ un schéma. Soit $X$ un espace algébrique sur $S$.
Soit $P$ une propriété des objets de $(\textit{Espaces}/X)_{fppf}$
telle que, pour tout recouvrement $\{U_i \to U\}$ de
$(\textit{Espaces}/X)_{fppf}$, l'implication suivante soit vérifiée :
$$
P(U_{i_0} \times_U \ldots \times_U U_{i_p})
\text{ pour tous }
p \geq 0,\ i_0, \ldots, i_p \in I
\Rightarrow P(U)
$$
Si $P(U)$ est vraie pour tout $U$ affine, plat et localement de présentation finie sur $X$,
alors $P(X)$.
\end{lemma}

\begin{proof}
Soit $U$ un espace algébrique séparé, plat et localement de présentation finie sur $X$.
Nous pouvons choisir un recouvrement \'etale $\{U_i \to U\}_{i \in I}$ avec $U_i$
affine. Comme $U$ est séparé, nous en concluons que
$U_{i_0} \times_U \ldots \times_U U_{i_p}$ est toujours affine.
Ainsi, $P(U_{i_0} \times_U \ldots \times_U U_{i_p})$ est toujours vraie.
Par conséquent, $P(U)$ est vraie. Choisissons un schéma $U$ qui est une réunion disjointe de
schémas affines et un morphisme \'etale surjectif $U \to X$.
Alors $U \times_X \ldots \times_X U$ (avec $p + 1$ facteurs)
est un espace algébrique séparé
et \'etale sur $X$. D'après ce qui précède, $P(U \times_X \ldots \times_X U)$ est donc vraie.
Nous en concluons que $P(X)$ est vraie.
\end{proof}












\section{Topologie fpqc}
\label{spaces-topologies-section-fpqc}

\noindent
Nous étudions brièvement la notion de recouvrement fpqc d'espaces algébriques.
Comparer avec
Topologies, section \ref{topologies-section-fpqc}.
Nous montrerons dans
Descente sur les espaces,
proposition \ref{spaces-descent-proposition-fpqc-descent-quasi-coherent},
que les faisceaux quasi-cohérents satisfont à la descente relativement à de tels recouvrements.

\begin{definition}
\label{spaces-topologies-definition-fpqc-covering}
Soit $S$ un schéma et soit $X$ un espace algébrique sur $S$.
Un {\it recouvrement fpqc de $X$} est une famille de morphismes
$\{f_i : X_i \to X\}_{i \in I}$ d'espaces algébriques
telle que chaque $f_i$ soit plat et que, pour tout schéma affine
$Z$ et tout morphisme $h : Z \to X$, il existe un recouvrement fpqc standard
$\{g_j : Z_j \to Z\}_{j = 1, \ldots, m}$ qui raffine la famille
$\{X_i \times_X Z \to Z\}_{i \in I}$.
\end{definition}

\noindent
Autrement dit, il existe des indices $i_1, \ldots, i_m \in I$ et des
morphismes $h_j : Z_j \to X_{i_j}$ tels que
$f_{i_j} \circ h_j = h \circ g_j$. Remarquons que, si $X$ et tous les $X_i$ sont
représentables, cela revient à un recouvrement fpqc de schémas d'après le
lemme \ref{topologies-lemma-fpqc-covering-affines-mapping-in} de Topologies.

\begin{lemma}
\label{spaces-topologies-lemma-zariski-etale-smooth-syntomic-fppf-fpqc}
Tout recouvrement fppf est un recouvrement fpqc et a fortiori,
tout recouvrement syntomique, lisse, \'etale ou de Zariski est un recouvrement fpqc.
\end{lemma}

\begin{proof}
Montrons qu'un recouvrement fppf est un recouvrement fpqc ; le
reste découle alors du
lemme \ref{spaces-topologies-lemma-zariski-etale-smooth-syntomic-fppf}.
Soit $\{f_i : U_i \to U\}_{i \in I}$ un
recouvrement fppf d'espaces algébriques sur $S$.
Par définition, les $f_i$ sont plats, ce qui vérifie la première
condition de la définition \ref{spaces-topologies-definition-fpqc-covering}. Pour vérifier la
seconde, soit $V \to U$ un morphisme avec $V$ affine.
Nous pouvons choisir un recouvrement \'etale $\{V_{ij} \to V \times_U U_i\}$
avec les $V_{ij}$ affines. Alors les composés
$f_{ij} : V_{ij} \to V \times_U U_i \to V$ sont plats et localement de
présentation finie, chacun étant composé de morphismes possédant ces propriétés
(Morphismes des espaces, lemmes
\ref{spaces-morphisms-lemma-composition-finite-presentation},
\ref{spaces-morphisms-lemma-composition-flat},
\ref{spaces-morphisms-lemma-etale-flat}, et
\ref{spaces-morphisms-lemma-etale-locally-finite-presentation}).
Ces morphismes sont donc ouverts
(Morphismes des espaces, lemme \ref{spaces-morphisms-lemma-fppf-open}),
et nous voyons que
$|V| = \bigcup_{i \in I} \bigcup_{j \in J_i} f_{ij}(|V_{ij}|)$
est un recouvrement ouvert de $|V|$.
Comme $|V|$ est quasi-compact, ce recouvrement possède un
raffinement fini.
Disons que $V_{i_1j_1}, \ldots, V_{i_Nj_N}$ conviennent.
Alors $\{V_{i_kj_k} \to V\}_{k = 1, \ldots, N}$ est
un recouvrement fpqc standard de $V$ qui raffine la famille
$\{U_i \times_U V \to V\}$.
Ceci achève la démonstration.
\end{proof}

\begin{lemma}
\label{spaces-topologies-lemma-fpqc}
Soit $S$ un schéma.
Soit $X$ un espace algébrique sur $S$.
\begin{enumerate}
\item Si $X' \to X$ est un isomorphisme, alors $\{X' \to X\}$
est un recouvrement fpqc de $X$.
\item Si $\{X_i \to X\}_{i\in I}$ est un recouvrement fpqc et si, pour tout
$i$, on a un recouvrement fpqc $\{X_{ij} \to X_i\}_{j\in J_i}$, alors
$\{X_{ij} \to X\}_{i \in I, j\in J_i}$ est un recouvrement fpqc.
\item Si $\{X_i \to X\}_{i\in I}$ est un recouvrement fpqc
et si $X' \to X$ est un morphisme d'espaces algébriques, alors
$\{X' \times_X X_i \to X'\}_{i\in I}$ est un recouvrement fpqc.
\end{enumerate}
\end{lemma}

\begin{proof}
La partie (1) est claire. Considérons $g : X' \to X$ et
un recouvrement fpqc $\{X_i \to X\}_{i\in I}$ comme dans (3). Par le
lemme \ref{spaces-morphisms-lemma-base-change-flat} de Morphismes des espaces,
les morphismes $X' \times_X X_i \to X'$
sont plats. Si $h' : Z \to X'$ est un morphisme d'un schéma affine
vers $X'$, posons $h = g \circ h' : Z \to X$. L'hypothèse
sur $\{X_i \to X\}_{i\in I}$ signifie qu'il existe un recouvrement fpqc standard
$\{Z_j \to Z\}_{j = 1, \ldots, n}$ et des morphismes $Z_j \to X_{i(j)}$ au-dessus de
$h$, pour certains $i(j) \in I$. Par la propriété universelle du produit fibré,
nous obtenons aussi des morphismes $Z_j \to X' \times_X X_{i(j)}$ au-dessus de $h'$.
Ainsi, $\{X' \times_X X_i \to X'\}_{i\in I}$ est un recouvrement fpqc.
Ceci démontre (3).

\medskip\noindent
Soient $\{X_i \to X\}_{i\in I}$ et $\{X_{ij} \to X_i\}_{j\in J_i}$ comme
dans (2). Soit $h : Z \to X$ un morphisme d'un schéma affine vers $X$.
Par hypothèse, il existe un recouvrement fpqc standard
$\{Z_j \to Z\}_{j = 1, \ldots, n}$ et des morphismes $h_j : Z_j \to X_{i(j)}$
au-dessus de $h$, pour certains indices $i(j) \in I$. Par hypothèse, il existe des
recouvrements fpqc standard
$\{Z_{j, l} \to Z_j\}_{l = 1, \ldots, n(j)}$
et des morphismes $Z_{j, l} \to X_{i(j)j(l)}$ au-dessus de
$h_j$, pour certains indices $j(l) \in J_{i(j)}$. Par le
lemme \ref{topologies-lemma-fpqc-affine-axioms} de Topologies,
la famille $\{Z_{j, l} \to Z\}$ est un recouvrement fpqc standard.
Nous en concluons que $\{X_{ij} \to X\}_{i \in I, j\in J_i}$
est un recouvrement fpqc.
\end{proof}

\begin{lemma}
\label{spaces-topologies-lemma-recognize-fpqc-covering}
Soit $S$ un schéma et soit $X$ un espace algébrique sur $S$.
Supposons que $\{f_i : X_i \to X\}_{i \in I}$ soit une famille de morphismes
d'espaces algébriques de but $X$. Soit $U \to X$ un morphisme
\'etale surjectif d'un schéma vers $X$. Alors
$\{f_i : X_i \to X\}_{i \in I}$ est un recouvrement fpqc de $X$ si et seulement
si $\{U \times_X X_i \to U\}_{i \in I}$ est un recouvrement fpqc de $U$.
\end{lemma}

\begin{proof}
Si $\{X_i \to X\}_{i \in I}$ est un recouvrement fpqc, il en va de même de
$\{U \times_X X_i \to U\}_{i \in I}$ par le lemme \ref{spaces-topologies-lemma-fpqc}.
Supposons que $\{U \times_X X_i \to U\}_{i \in I}$ soit un recouvrement fpqc.
Soit $h : Z \to X$ un morphisme d'un schéma affine vers $X$.
Alors $U \times_X Z \to Z$ est un morphisme \'etale surjectif
de schémas, donc en particulier ouvert. Nous pouvons donc trouver un nombre fini d'ouverts affines
$W_1, \ldots, W_t$ de $U \times_X Z$ dont les images recouvrent $Z$.
Pour chaque $j$, nous pouvons appliquer la condition selon laquelle
$\{U \times_X X_i \to U\}_{i \in I}$ est un recouvrement fpqc
au morphisme $W_j \to U$, et obtenir un recouvrement fpqc standard
$\{W_{jl} \to W_j\}$ qui raffine $\{W_j \times_X X_i \to W_j\}_{i \in I}$.
Ainsi, $\{W_{jl} \to Z\}$ est un recouvrement fpqc standard de $Z$
(voir le
lemme \ref{topologies-lemma-fpqc-affine-axioms} de Topologies)
qui raffine $\{Z \times_X X_i \to Z\}$, ce qui achève la démonstration.
\end{proof}

\begin{lemma}
\label{spaces-topologies-lemma-refine-fpqc-schemes}
Soit $S$ un schéma et soit $X$ un espace algébrique sur $S$.
Supposons que $\mathcal{U} = \{f_i : X_i \to X\}_{i \in I}$ soit un
recouvrement fpqc de $X$. Il existe alors un raffinement
$\mathcal{V} = \{g_i : T_i \to X\}$ de $\mathcal{U}$ qui est un
recouvrement fpqc tel que chaque $T_i$ soit un schéma.
\end{lemma}

\begin{proof}
Démonstration omise. Indication : pour chaque $i$, choisir un schéma $T_i$ et un morphisme \'etale
surjectif $T_i \to X_i$. Vérifier ensuite que $\{T_i \to X\}$ est un recouvrement fpqc.
\end{proof}

\noindent
À suivre...










% OK, start here.
%

\chapter{Descente et espaces algébriques}


\phantomsection
\label{spaces-descent-section-phantom}


\section{Introduction}
\label{spaces-descent-section-introduction}

\noindent
Dans le chapitre consacré aux topologies sur les espaces algébriques (voir
Topologies sur les espaces, section \ref{spaces-topologies-section-introduction})
nous avons introduit les recouvrements \'etales, fppf, lisses, syntomiques et fpqc des
espaces algébriques.
Dans ce chapitre, nous étudions les structures sur les espaces algébriques
dont la descente est possible relativement à de tels recouvrements.
Voir par exemple \cite{Gr-I}, \cite{Gr-II}, \cite{Gr-III},
\cite{Gr-IV}, \cite{Gr-V} et \cite{Gr-VI}.



\section{Conventions}
\label{spaces-descent-section-conventions}

\noindent
L'hypothèse permanente est que tous les schémas appartiennent à
un grand site fppf $\Sch_{fppf}$. De plus, tous les anneaux $A$ considérés
sont tels que $\Spec(A)$ est isomorphe à un
objet de ce grand site.

\medskip\noindent
Soit $S$ un schéma et soit $X$ un espace algébrique sur $S$.
Dans ce chapitre et dans le suivant, nous écrirons $X \times_S X$
pour le produit de $X$ avec lui-même (dans la catégorie des espaces algébriques
sur $S$), au lieu de $X \times X$.






\section{Données de descente pour les faisceaux quasi-cohérents}
\label{spaces-descent-section-equivalence}

\noindent
Cette section est l'analogue de
Descente, section \ref{descent-section-equivalence}
pour les espaces algébriques.
Il est naturel de lire d'abord cette dernière.

\begin{definition}
\label{spaces-descent-definition-descent-datum-quasi-coherent}
Soit $S$ un schéma. Soit $\{f_i : X_i \to X\}_{i \in I}$ une famille
de morphismes d'espaces algébriques sur $S$ à cible fixée $X$.
\begin{enumerate}
\item Une {\it donnée de descente $(\mathcal{F}_i, \varphi_{ij})$
pour les faisceaux quasi-cohérents} relativement à la famille donnée
consiste en la donnée d'un faisceau quasi-cohérent $\mathcal{F}_i$ sur $X_i$ pour
chaque $i \in I$, et d'un isomorphisme de
$\mathcal{O}_{X_i \times_X X_j}$-modules quasi-cohérents
$\varphi_{ij} : \text{pr}_0^*\mathcal{F}_i \to \text{pr}_1^*\mathcal{F}_j$
pour chaque paire $(i, j) \in I^2$
tel que, pour tout triplet d'indices $(i, j, k) \in I^3$, le
diagramme
$$
\xymatrix{
\text{pr}_0^*\mathcal{F}_i \ar[rd]_{\text{pr}_{01}^*\varphi_{ij}}
\ar[rr]_{\text{pr}_{02}^*\varphi_{ik}} & &
\text{pr}_2^*\mathcal{F}_k \\
& \text{pr}_1^*\mathcal{F}_j \ar[ru]_{\text{pr}_{12}^*\varphi_{jk}} &
}
$$
de $\mathcal{O}_{X_i \times_X X_j \times_X X_k}$-modules
est commutatif. C'est la {\it condition de cocycle}.
\item Un {\it morphisme $\psi : (\mathcal{F}_i, \varphi_{ij}) \to
(\mathcal{F}'_i, \varphi'_{ij})$ de données de descente} est donné
par une famille $\psi = (\psi_i)_{i\in I}$ de morphismes de
$\mathcal{O}_{X_i}$-modules $\psi_i : \mathcal{F}_i \to \mathcal{F}'_i$
telle que tous les diagrammes
$$
\xymatrix{
\text{pr}_0^*\mathcal{F}_i \ar[r]_{\varphi_{ij}} \ar[d]_{\text{pr}_0^*\psi_i}
& \text{pr}_1^*\mathcal{F}_j \ar[d]^{\text{pr}_1^*\psi_j} \\
\text{pr}_0^*\mathcal{F}'_i \ar[r]^{\varphi'_{ij}} &
\text{pr}_1^*\mathcal{F}'_j \\
}
$$
soient commutatifs.
\end{enumerate}
\end{definition}

\begin{lemma}
\label{spaces-descent-lemma-map-families}
Soit $S$ un schéma.
Soient $\mathcal{U} = \{U_i \to U\}_{i \in I}$ et
$\mathcal{V} = \{V_j \to V\}_{j \in J}$
des familles de morphismes d'espaces algébriques sur $S$ à cibles fixées.
Soit $(g, \alpha : I \to J, (g_i)) : \mathcal{U} \to \mathcal{V}$
un morphisme de familles de morphismes à cible fixée, voir
Sites, définition \ref{sites-definition-morphism-coverings}.
Soit $(\mathcal{F}_j, \varphi_{jj'})$ une donnée de descente
pour les faisceaux quasi-cohérents relativement à la
famille $\{V_j \to V\}_{j \in J}$. Alors
\begin{enumerate}
\item Le système
$$
\left(g_i^*\mathcal{F}_{\alpha(i)},
(g_i \times g_{i'})^*\varphi_{\alpha(i)\alpha(i')}\right)
$$
est une donnée de descente relativement à la famille $\{U_i \to U\}_{i \in I}$.
\item Cette construction est fonctorielle en la donnée de descente
$(\mathcal{F}_j, \varphi_{jj'})$.
\item Étant donné un second morphisme
$(g', \alpha' : I \to J, (g'_i))$
de familles de morphismes à cible fixée
tel que $g = g'$, il existe un isomorphisme fonctoriel de données de descente
$$
(g_i^*\mathcal{F}_{\alpha(i)},
(g_i \times g_{i'})^*\varphi_{\alpha(i)\alpha(i')})
\cong
((g'_i)^*\mathcal{F}_{\alpha'(i)},
(g'_i \times g'_{i'})^*\varphi_{\alpha'(i)\alpha'(i')}).
$$
\end{enumerate}
\end{lemma}

\begin{proof}
Démonstration omise. Indication : les morphismes
$g_i^*\mathcal{F}_{\alpha(i)} \to (g'_i)^*\mathcal{F}_{\alpha'(i)}$
qui donnent l'isomorphisme de données de descente en (3)
sont les images réciproques des morphismes $\varphi_{\alpha(i)\alpha'(i)}$ par les
morphismes $(g_i, g'_i) : U_i \to V_{\alpha(i)} \times_V V_{\alpha'(i)}$.
\end{proof}

\noindent
Soit $g : U \to V$ un morphisme d'espaces algébriques.
Le lemme ci-dessus montre qu'il existe un foncteur image réciproque bien défini
entre les catégories de données de descente relatives à des familles de
morphismes ayant respectivement $V$ et $U$ pour buts, pourvu qu'il existe entre ces
familles un morphisme qui ``vit au-dessus de $g$''.

\begin{definition}
\label{spaces-descent-definition-descent-datum-effective-quasi-coherent}
Soit $S$ un schéma.
Soit $\{U_i \to U\}_{i \in I}$ une famille de morphismes d'espaces algébriques
sur $S$ à cible fixée.
\begin{enumerate}
\item Soit $\mathcal{F}$ un $\mathcal{O}_U$-module quasi-cohérent.
L'unique donnée de descente portée par $\mathcal{F}$ relativement au recouvrement
$\{U \to U\}$ est appelée la {\it donnée de descente triviale}.
\item L'image réciproque de la donnée de descente triviale relativement à
la famille $\{U_i \to U\}$ est appelée la {\it donnée de descente canonique}.
Notation : $(\mathcal{F}|_{U_i}, can)$.
\item Une donnée de descente $(\mathcal{F}_i, \varphi_{ij})$
pour les faisceaux quasi-cohérents relativement à la famille donnée
est dite {\it effective} s'il existe un faisceau quasi-cohérent
$\mathcal{F}$ sur $U$ tel que $(\mathcal{F}_i, \varphi_{ij})$
soit isomorphe à $(\mathcal{F}|_{U_i}, can)$.
\end{enumerate}
\end{definition}

\begin{lemma}
\label{spaces-descent-lemma-zariski-descent-effective}
Soit $S$ un schéma. Soit $U$ un espace algébrique sur $S$.
Soit $\{U_i \to U\}$ un recouvrement de Zariski de $U$, voir
Topologies sur les espaces,
définition \ref{spaces-topologies-definition-zariski-covering}.
Toute donnée de descente pour les faisceaux quasi-cohérents
relativement à la famille $\mathcal{U} = \{U_i \to U\}$ est
effective. De plus, le foncteur de la catégorie des
$\mathcal{O}_U$-modules quasi-cohérents vers la catégorie
des données de descente relativement à $\{U_i \to U\}$ est pleinement fidèle.
\end{lemma}

\begin{proof}
Démonstration omise.
\end{proof}












\section{Descente fpqc des faisceaux quasi-cohérents}
\label{spaces-descent-section-fpqc-descent-quasi-coherent}

\noindent
La principale application de la descente plate des modules est
l'énoncé de descente correspondant pour les faisceaux quasi-cohérents
relativement aux recouvrements fpqc.

\begin{proposition}
\label{spaces-descent-proposition-fpqc-descent-quasi-coherent}
Soit $S$ un schéma.
Soit $\{X_i \to X\}$ un recouvrement fpqc d'espaces algébriques sur $S$, voir
Topologies sur les espaces,
définition \ref{spaces-topologies-definition-fpqc-covering}.
Toute donnée de descente pour les faisceaux quasi-cohérents
relativement à $\{X_i \to X\}$ est effective.
De plus, le foncteur de la catégorie des
$\mathcal{O}_X$-modules quasi-cohérents vers la catégorie
des données de descente relativement à $\{X_i \to X\}$ est pleinement fidèle.
\end{proposition}

\begin{proof}
C'est, à peu de chose près, une conséquence formelle du
résultat correspondant pour les schémas, voir
Descente, proposition \ref{descent-proposition-fpqc-descent-quasi-coherent}.
Voici une stratégie de démonstration :
\begin{enumerate}
\item Le fait que $\{X_i \to X\}$ soit un raffinement du recouvrement trivial
$\{X \to X\}$ donne, via
le lemme \ref{spaces-descent-lemma-map-families},
un foncteur $\QCoh(\mathcal{O}_X) \to DD(\{X_i \to X\})$ de la
catégorie des $\mathcal{O}_X$-modules quasi-cohérents vers la catégorie des
données de descente pour la famille donnée.
\item Afin de démontrer la proposition, nous allons construire un
foncteur quasi-inverse
$back : DD(\{X_i \to X\}) \to \QCoh(\mathcal{O}_X)$.
\item En appliquant à nouveau
le lemme \ref{spaces-descent-lemma-map-families}
nous voyons qu'il existe un foncteur
$DD(\{X_i \to X\}) \to DD(\{T_j \to X\})$
si $\{T_j \to X\}$ est un raffinement de la famille donnée.
Ainsi, pour construire le foncteur $back$, nous pouvons supposer que
chaque $X_i$ est un schéma, voir
Topologies sur les espaces,
lemme \ref{spaces-topologies-lemma-refine-fpqc-schemes}.
Cela nous ramène au cas où tous les $X_i$ sont des schémas.
\item Un faisceau quasi-cohérent sur $X$ est, par définition, un
$\mathcal{O}_X$-module quasi-cohérent sur $X_\etale$. Pour tout
$U \in \Ob(X_\etale)$, on obtient alors un recouvrement fppf
$\{U_i \times_X X_i \to U\}$ par des schémas et un morphisme
$g : \{U_i \times_X X_i \to U\} \to \{X_i \to X\}$ de recouvrements
au-dessus de $U \to X$. Étant donnée une donnée de descente
$\xi = (\mathcal{F}_i, \varphi_{ij})$, on obtient un
$\mathcal{O}_U$-module quasi-cohérent $\mathcal{F}_{\xi, U}$ correspondant
à l'image réciproque $g^*\xi$ fournie par le
lemme \ref{spaces-descent-lemma-map-families}
sur le recouvrement de $U$, puis on applique l'effectivité de la descente
pour les recouvrements fppf de schémas, voir Descente, proposition \ref{descent-proposition-fpqc-descent-quasi-coherent}.
\item Vérifier que $\xi \mapsto \mathcal{F}_{\xi, U}$ est fonctorielle en $\xi$.
Vérification omise.
\item Vérifier que $\xi \mapsto \mathcal{F}_{\xi, U}$ est compatible
avec les morphismes $U \to U'$ du site $X_\etale$, de sorte que
le système de faisceaux $\mathcal{F}_{\xi, U}$ corresponde à un faisceau quasi-cohérent
$\mathcal{F}_\xi$ sur $X_\etale$, voir
Propriétés des espaces,
lemme \ref{spaces-properties-lemma-characterize-quasi-coherent-small-etale}.
Détails omis.
\item Vérifier que $back : \xi \mapsto \mathcal{F}_\xi$ est quasi-inverse
du foncteur construit en (1). Vérification omise.
\end{enumerate}
La démonstration est ainsi achevée.
\end{proof}









\section{Modules quasi-cohérents et schémas affines}
\label{spaces-descent-section-alternative-quasi-coherent}

\noindent
Soit $S$ un schéma. Soit $X$ un espace algébrique sur $S$.
Rappelons que $X_{affine, \etale}$ est la sous-catégorie pleine de $X_\etale$
dont les objets sont les schémas affines; on la munit d'une structure de site en prenant
pour recouvrements les recouvrements \'etales standards. Voir Propriétés des espaces, définition
\ref{spaces-properties-definition-affine-etale-site}.
D'après Propriétés des espaces, lemme \ref{spaces-properties-lemma-alternative},
on a une équivalence de topos
$g : \Sh(X_{affine, \etale}) \to \Sh(X_\etale)$
dont le foncteur image réciproque est donné par la restriction.
Rappelons que $\mathcal{O}_X$ désigne le faisceau structural sur
$X_\etale$. Nous obtenons alors une équivalence
\begin{equation}
\label{spaces-descent-equation-alternative-small-ringed}
(\Sh(X_{affine, \etale}), \mathcal{O}_X|_{X_{affine, \etale}})
\longrightarrow
(\Sh(X_\etale), \mathcal{O}_X)
\end{equation}
de topos annelés. Nous écrirons souvent $\mathcal{O}_X$
au lieu de $\mathcal{O}_X|_{X_{affine, \etale}}$.
Cela permet également de comparer les modules quasi-cohérents.

\begin{lemma}
\label{spaces-descent-lemma-quasi-coherent-alternative-small}
Soit $S$ un schéma. Soit $X$ un espace algébrique sur $S$.
Soit $\mathcal{F}$ un préfaisceau de $\mathcal{O}_X$-modules
sur $X_{affine, \etale}$. Les conditions suivantes sont équivalentes :
\begin{enumerate}
\item pour tout morphisme $U \to U'$ de $X_{affine, \etale}$ la flèche
$\mathcal{F}(U') \otimes_{\mathcal{O}_X(U')} \mathcal{O}_X(U)
\to \mathcal{F}(U)$ est un isomorphisme,
\item $\mathcal{F}$ est un module quasi-cohérent sur le site annelé
$(X_{affine, \etale}, \mathcal{O}_X)$ au sens de
Modules sur les sites, définition \ref{sites-modules-definition-site-local},
\item $\mathcal{F}$ correspond à un module quasi-cohérent sur $X$
via l'équivalence (\ref{spaces-descent-equation-alternative-small-ringed}),
\end{enumerate}
\end{lemma}

\begin{proof}
Supposons que (1) soit satisfaite. Pour montrer que $\mathcal{F}$ est un faisceau, soit
$\mathcal{U} = \{U_i \to U\}_{i = 1, \ldots, n}$ un recouvrement
de $X_{affine, \etale}$. La condition de faisceau pour $\mathcal{F}$
et $\mathcal{U}$, par notre hypothèse sur $\mathcal{F}$,
se réduit à montrer que
$$
0 \to \mathcal{F}(U) \to
\prod \mathcal{F}(U) \otimes_{\mathcal{O}_X(U)} \mathcal{O}_X(U_i) \to
\prod \mathcal{F}(U) \otimes_{\mathcal{O}_X(U)} \mathcal{O}_X(U_i \times_U U_j)
$$
Cette suite est exacte, car $\mathcal{O}_X(U) \to \prod \mathcal{O}_X(U_i)$
est fidèlement plat
(d'après Descente, lemme \ref{descent-lemma-standard-covering-Cech}, et parce que
les recouvrements de $X_{affine, \etale}$ sont les recouvrements \'etales
standards); on peut alors appliquer Descente, lemme \ref{descent-lemma-ff-exact}.
Ensuite, nous montrons que $\mathcal{F}$ est quasi-cohérent sur $X_{affine, \etale}$.
En effet, pour $U$ dans $X_{affine, \etale}$, posons $R = \mathcal{O}_X(U)$
et choisissons une présentation
$$
\bigoplus\nolimits_{k \in K} R
\longrightarrow
\bigoplus\nolimits_{l \in L} R
\longrightarrow
\mathcal{F}(U)
\longrightarrow 0
$$
par des $R$-modules libres. D'après la propriété (1) et l'exactitude à droite du produit tensoriel,
nous voyons que, pour tout morphisme $U' \to U$ dans $X_{affine, \etale}$,
on obtient une présentation
$$
\bigoplus\nolimits_{k \in K} \mathcal{O}_X(U')
\longrightarrow
\bigoplus\nolimits_{l \in L} \mathcal{O}_X(U')
\longrightarrow
\mathcal{F}(U')
\longrightarrow 0
$$
Autrement dit, la restriction de $\mathcal{F}$
à la catégorie localisée $X_{affine, etale}/U$ possède une présentation
$$
\bigoplus\nolimits_{k \in K} \mathcal{O}_X|_{X_{affine, \etale}/U}
\longrightarrow
\bigoplus\nolimits_{l \in L} \mathcal{O}_X|_{X_{affine, \etale}/U}
\longrightarrow
\mathcal{F}|_{X_{affine, \etale}/U}
\longrightarrow 0
$$
ce qui est la présentation requise pour montrer que $\mathcal{F}$ est quasi-cohérent.
Toutes nos excuses pour cette horrible notation; cela achève la démonstration
que (1) implique (2).

\medskip\noindent
Puisque la notion de module quasi-cohérent est intrinsèque
(Modules sur les sites, lemme \ref{sites-modules-lemma-special-locally-free})
nous voyons que l'équivalence (\ref{spaces-descent-equation-alternative-small-ringed})
induit une équivalence entre les catégories de modules quasi-cohérents.
Ainsi, (2) et (3) sont équivalentes.

\medskip\noindent
Supposons (3) et démontrons (1). Soit
$\mathcal{G}$ un module quasi-cohérent sur $X$ correspondant à $\mathcal{F}$.
Soit $h : U \to U' \to X$ un morphisme de $X_{affine, \etale}$.
Notons $f : U \to X$ et $f' : U' \to X$ les morphismes structuraux,
de sorte que $f = f' \circ h$. Nous avons
$\mathcal{F}(U') = \Gamma(U', (f')^*\mathcal{G})$ et
$\mathcal{F}(U) = \Gamma(U, f^*\mathcal{G}) = \Gamma(U, h^*(f')^*\mathcal{G})$.
Ainsi, (1) résulte de Schémas, lemme \ref{schemes-lemma-widetilde-pullback}.
\end{proof}
% preserved blank authority line 391
% preserved blank authority line 392
% preserved blank authority line 393
% preserved blank authority line 394
% preserved blank authority line 395
\section{Descente des propriétés de finitude des modules}
\label{spaces-descent-section-descent-finiteness}

\noindent
Cette section est, pour les espaces algébriques, l'analogue de
Descente, section \ref{descent-section-descent-finiteness}.
Nous voulons montrer que certaines conditions de finitude portant sur un module
quasi-cohérent peuvent se vérifier sur les membres d'un
recouvrement. Nous utiliserons à plusieurs reprises le schéma de démonstration suivant.
Supposons que $X$ soit un espace algébrique et que $\{X_i \to X\}$
soit un recouvrement fppf (resp.\ fpqc). Soit $U \to X$ un morphisme \'etale
surjectif tel que $U$ soit un schéma. Il existe alors un
recouvrement fppf (resp.\ fpqc) $\{Y_j \to X\}$ tel que
\begin{enumerate}
\item $\{Y_j \to X\}$ est un raffinement de $\{X_i \to X\}$,
\item chaque $Y_j$ est un schéma, et
\item chaque morphisme $Y_j \to X$ se factorise par $U$, et
\item $\{Y_j \to U\}$ est un recouvrement fppf (resp.\ fpqc) de $U$.
\end{enumerate}
En effet, raffinons d'abord $\{X_i \to X\}$ par un recouvrement fppf (resp.\ fpqc)
tel que chaque $X_i$ soit un schéma, voir
Topologies sur les espaces, lemme \ref{spaces-topologies-lemma-refine-fppf-schemes},
resp.\ le lemme \ref{spaces-topologies-lemma-refine-fpqc-schemes}.
Posons alors $Y_i = U \times_X X_i$. Un module quasi-cohérent
sur $\mathcal{O}_X$, noté $\mathcal{F}$, est de type fini, de
présentation finie, etc., si et seulement si le module quasi-cohérent
sur $\mathcal{O}_U$, noté $\mathcal{F}|_U$, est de type fini, de
présentation finie, etc. L'existence du
raffinement $\{Y_j \to X\}$ permet donc de ramener la démonstration des
lemmes suivants au cas des schémas. Nous l'indiquerons en disant
que « {\it le résultat se déduit du cas des schémas par
localisation \'etale} ».

\begin{lemma}
\label{spaces-descent-lemma-finite-type-descends}
Soit $X$ un espace algébrique sur un schéma $S$.
Soit $\mathcal{F}$ un $\mathcal{O}_X$-module quasi-cohérent.
Soit $\{f_i : X_i \to X\}_{i \in I}$ un recouvrement fpqc tel que
chaque $f_i^*\mathcal{F}$ soit un $\mathcal{O}_{X_i}$-module de type fini.
Alors $\mathcal{F}$ est un $\mathcal{O}_X$-module de type fini.
\end{lemma}

\begin{proof}
Ce résultat se déduit du cas des schémas, voir
Descente, lemme \ref{descent-lemma-finite-type-descends},
par localisation \'etale.
\end{proof}

\begin{lemma}
\label{spaces-descent-lemma-finite-presentation-descends}
Soit $X$ un espace algébrique sur un schéma $S$.
Soit $\mathcal{F}$ un $\mathcal{O}_X$-module quasi-cohérent.
Soit $\{f_i : X_i \to X\}_{i \in I}$ un recouvrement fpqc tel que
chaque $f_i^*\mathcal{F}$ soit un $\mathcal{O}_{X_i}$-module de présentation
finie. Alors $\mathcal{F}$ est un $\mathcal{O}_X$-module
de présentation finie.
\end{lemma}

\begin{proof}
Ce résultat se déduit du cas des schémas, voir
Descente, lemme \ref{descent-lemma-finite-presentation-descends},
par localisation \'etale.
\end{proof}

\begin{lemma}
\label{spaces-descent-lemma-flat-descends}
Soit $X$ un espace algébrique sur un schéma $S$.
Soit $\mathcal{F}$ un $\mathcal{O}_X$-module quasi-cohérent.
Soit $\{f_i : X_i \to X\}_{i \in I}$ un recouvrement fpqc tel que
chaque $f_i^*\mathcal{F}$ soit un $\mathcal{O}_{X_i}$-module plat.
Alors $\mathcal{F}$ est un $\mathcal{O}_X$-module plat.
\end{lemma}

\begin{proof}
Ce résultat se déduit du cas des schémas, voir
Descente, lemme \ref{descent-lemma-flat-descends},
par localisation \'etale.
\end{proof}

\begin{lemma}
\label{spaces-descent-lemma-finite-locally-free-descends}
Soit $X$ un espace algébrique sur un schéma $S$.
Soit $\mathcal{F}$ un $\mathcal{O}_X$-module quasi-cohérent.
Soit $\{f_i : X_i \to X\}_{i \in I}$ un recouvrement fpqc tel que
chaque $f_i^*\mathcal{F}$ soit un $\mathcal{O}_{X_i}$-module localement libre de type fini.
Alors $\mathcal{F}$ est un $\mathcal{O}_X$-module localement libre de type fini.
\end{lemma}

\begin{proof}
Ce résultat se déduit du cas des schémas, voir
Descente, lemme \ref{descent-lemma-finite-locally-free-descends},
par localisation \'etale.
\end{proof}

\noindent
La définition d'un faisceau quasi-cohérent localement projectif se trouve dans
Propriétés des espaces, section
\ref{spaces-properties-section-locally-projective}.
Il y est également démontré que cette notion est stable par image réciproque.

\begin{lemma}
\label{spaces-descent-lemma-locally-projective-descends}
Soit $X$ un espace algébrique sur un schéma $S$.
Soit $\mathcal{F}$ un $\mathcal{O}_X$-module quasi-cohérent.
Soit $\{f_i : X_i \to X\}_{i \in I}$ un recouvrement fpqc tel que
chaque $f_i^*\mathcal{F}$ soit un $\mathcal{O}_{X_i}$-module localement projectif.
Alors $\mathcal{F}$ est un $\mathcal{O}_X$-module localement projectif.
\end{lemma}

\begin{proof}
Ce résultat se déduit du cas des schémas, voir
Descente, lemme \ref{descent-lemma-locally-projective-descends},
par localisation \'etale.
\end{proof}

\noindent
Ajoutons ici deux résultats liés aux précédents, mais
de nature légèrement différente.

\begin{lemma}
\label{spaces-descent-lemma-finite-over-finite-module}
Soit $S$ un schéma.
Soit $f : X \to Y$ un morphisme d'espaces algébriques sur $S$.
Soit $\mathcal{F}$ un $\mathcal{O}_X$-module quasi-cohérent.
Supposons que $f$ soit fini.
Alors $\mathcal{F}$ est un $\mathcal{O}_X$-module de type fini
si et seulement si $f_*\mathcal{F}$ est un $\mathcal{O}_Y$-module de type
fini.
\end{lemma}

\begin{proof}
Puisque $f$ est fini, il est représentable. Choisissons un schéma $V$ et un
morphisme \'etale surjectif $V \to Y$. Alors $U = V \times_Y X$ est un schéma muni d'un
morphisme \'etale surjectif vers $X$ et d'un morphisme fini
$\psi : U \to V$ (le changement de base de $f$). Puisque
$\psi_*(\mathcal{F}|_U) = f_*\mathcal{F}|_V$,
l'assertion se déduit immédiatement de la version pour les schémas,
c'est-à-dire du
lemme \ref{descent-lemma-finite-over-finite-module} de Descente.
\end{proof}

\begin{lemma}
\label{spaces-descent-lemma-finite-finitely-presented-module}
Soit $S$ un schéma.
Soit $f : X \to Y$ un morphisme d'espaces algébriques sur $S$.
Soit $\mathcal{F}$ un $\mathcal{O}_X$-module quasi-cohérent.
Supposons que $f$ soit fini et de présentation finie.
Alors $\mathcal{F}$ est un $\mathcal{O}_X$-module de présentation finie
si et seulement si $f_*\mathcal{F}$ est un $\mathcal{O}_Y$-module de présentation
finie.
\end{lemma}

\begin{proof}
Puisque $f$ est fini, il est représentable. Choisissons un schéma $V$ et un
morphisme \'etale surjectif $V \to Y$. Alors $U = V \times_Y X$ est un schéma muni d'un
morphisme \'etale surjectif vers $X$ et d'un morphisme fini
$\psi : U \to V$ (le changement de base de $f$). Puisque
$\psi_*(\mathcal{F}|_U) = f_*\mathcal{F}|_V$,
l'assertion se déduit immédiatement de la version pour les schémas,
c'est-à-dire du
lemme \ref{descent-lemma-finite-finitely-presented-module} de Descente.
\end{proof}






\section{Recouvrements fpqc}
\label{spaces-descent-section-fpqc}

\noindent
Cette section est l'analogue de la
section \ref{descent-section-fpqc-universal-effective-epimorphisms} de Descente.
Nous ne savons pas actuellement si tous les résultats valables pour les
recouvrements fpqc de schémas le sont aussi pour les espaces algébriques.

\begin{lemma}
\label{spaces-descent-lemma-open-fpqc-covering}
Soit $S$ un schéma.
Soit $\{f_i : T_i \to T\}_{i \in I}$ un recouvrement fpqc
d'espaces algébriques sur $S$.
Pour chaque $i$, soit $W_i \subset T_i$ un sous-espace ouvert,
et supposons que, pour tous $i, j \in I$,
$\text{pr}_0^{-1}(W_i) = \text{pr}_1^{-1}(W_j)$ en tant que sous-espaces ouverts
de $T_i \times_T T_j$. Alors il existe un unique sous-espace ouvert
$W \subset T$ tel que $W_i = f_i^{-1}(W)$ pour chaque $i$.
\end{lemma}

\begin{proof}
D'après le
lemme \ref{spaces-topologies-lemma-refine-fpqc-schemes} de Topologies sur les espaces,
nous pouvons supposer que chaque $T_i$ est un schéma.
Choisissons un schéma $U$ et un morphisme \'etale surjectif $U \to T$.
Alors $\{T_i \times_T U \to U\}$ est un recouvrement fpqc de $U$
et $T_i \times_T U$ est un schéma pour chaque $i$. La famille
des ouverts $W_i \times_T U$ provient donc d'un unique
sous-schéma ouvert $W' \subset U$, d'après le
lemme \ref{descent-lemma-open-fpqc-covering} de Descente.
Comme $U \to X$ est ouvert, nous pouvons définir $W \subset X$ comme l'ouvert de Zariski
image de $W'$; voir
Propriétés des espaces, section \ref{spaces-properties-section-points}.
Nous omettons de vérifier que cette construction convient, c'est-à-dire que
$W_i$ est l'image réciproque de $W$ pour chaque $i$.
\end{proof}

\begin{lemma}
\label{spaces-descent-lemma-fpqc-universal-effective-epimorphisms}
Soit $S$ un schéma. Soit $\{T_i \to T\}$ un recouvrement fpqc d'espaces algébriques
sur $S$; voir Topologies sur les espaces, définition
\ref{spaces-topologies-definition-fpqc-covering}.
Alors, pour tout espace algébrique $B$ sur $S$, la suite
$$
\xymatrix{
\Mor_S(T, B) \ar[r] &
\prod\nolimits_i \Mor_S(T_i, B) \ar@<1ex>[r] \ar@<-1ex>[r] &
\prod\nolimits_{i, j} \Mor_S(T_i \times_T T_j, B)
}
$$
est un diagramme égalisateur.
Autrement dit, tout foncteur représentable sur la catégorie des
espaces algébriques sur $S$ vérifie la condition de faisceau pour les
recouvrements fpqc.
\end{lemma}

\begin{proof}
Nous savons que l'assertion est vraie si $\{T_i \to T\}$ est un recouvrement fpqc de
schémas, voir Propriétés des espaces, proposition
\ref{spaces-properties-proposition-sheaf-fpqc}.
C'est le point essentiel; nous invitons le lecteur à omettre la suite
de cette démonstration, qui est formelle. Choisissons un schéma $U$ et un
morphisme \'etale surjectif
$U \to T$. Choisissons un schéma $U_i$ et un morphisme $U_i \to T_i \times_T U$
\'etale surjectif. Alors $\{U_i \to U\}$ est un
recouvrement fpqc. Cela découle de
Topologies sur les espaces, lemmes \ref{spaces-topologies-lemma-fpqc} et
\ref{spaces-topologies-lemma-recognize-fpqc-covering}.
Ce qui précède donne le résultat pour $\{U_i \to U\}$.

\medskip\noindent
Concrètement, supposons que les $b_i : T_i \to B$
forment une famille de morphismes telle que
$b_i \circ \text{pr}_0 = b_j \circ \text{pr}_1$ comme morphismes
$T_i \times_T T_j \to B$. Définissons $a_i : U_i \to B$ comme le
composé obtenu en faisant suivre $U_i \to T_i$ de $b_i$. D'après ce qui précède,
il existe un unique morphisme $a : U \to B$ tel que
$a_i$ soit le composé de $a$ avec $U_i \to U$.
L'unicité assure que $a \circ \text{pr}_0 = a \circ \text{pr}_1$
comme morphismes $U \times_T U \to B$. Puisque $T = U/(U \times_T U)$
comme faisceau, $a$ provient d'un unique morphisme $b : T \to B$.
Une chasse au diagramme montre que $b$ est le morphisme recherché.
\end{proof}










\section{Descente des propriétés de finitude et de lissité des morphismes}
\sectionmark{Descente : finitude et lissité}
\label{spaces-descent-section-descent-finiteness-morphisms}

\noindent
Un lemme du type suivant est parfois utile.

\begin{lemma}
\label{spaces-descent-lemma-curiosity}
Soit $S$ un schéma. Soient $X \to Y \to Z$ des morphismes d'espaces algébriques.
Soit $P$ l'une des propriétés suivantes pour un morphisme d'espaces algébriques
sur $S$ :
être plat, être localement de type fini ou être localement de présentation finie.
Supposons que $X \to Z$ vérifie $P$ et que
$X \to Y$ soit une surjection de faisceaux sur $(\Sch/S)_{fppf}$.
Alors $Y \to Z$ vérifie $P$.
\end{lemma}

\begin{proof}
Choisissons un schéma $W$ et un morphisme \'etale surjectif $W \to Z$.
Choisissons un schéma $V$ et un morphisme \'etale surjectif $V \to W \times_Z Y$.
Choisissons un schéma $U$ et un morphisme \'etale surjectif $U \to V \times_Y X$.
Par hypothèse, il existe un recouvrement fppf $\{V_i \to V\}$ et des
relèvements $V_i \to X$ du morphisme $V_i \to Y$. Comme $U \to X$ est surjectif
\'etale, nous voyons que, sur les membres du recouvrement fppf
$\{V_i \times_X U \to V\}$, il existe des relèvements dans $U$. Par conséquent, $U \to V$ induit
une surjection de faisceaux sur $(\Sch/S)_{fppf}$.
D'après la définition de la propriété $P$ pour un
morphisme d'espaces algébriques (voir
Morphismes d'espaces,
définition \ref{spaces-morphisms-definition-flat},
définition \ref{spaces-morphisms-definition-locally-finite-type}, et
définition \ref{spaces-morphisms-definition-locally-finite-presentation})
nous savons que $U \to W$ vérifie $P$, et il reste à montrer que $V \to W$ vérifie $P$.
La question est ainsi ramenée au cas des morphismes de schémas
traité dans le
lemme \ref{descent-lemma-curiosity} de Descente.
\end{proof}

\noindent
Un cas plus usuel du lemme précédent est le suivant.
(La version « plate » découle de
Morphismes d'espaces, lemme \ref{spaces-morphisms-lemma-flat-permanence}.)

\begin{lemma}
\label{spaces-descent-lemma-flat-finitely-presented-permanence}
Soit $S$ un schéma. Considérons le diagramme suivant
$$
\xymatrix{
X \ar[rr]_f \ar[rd]_p & &
Y \ar[dl]^q \\
& B
}
$$
de morphismes d'espaces algébriques sur $S$, supposé commutatif.
Supposons que $f$ soit surjectif, plat et localement de présentation finie
et que $p$ soit localement de présentation finie (resp.\ localement
de type fini). Alors $q$ est localement de présentation finie
(resp.\ localement de type fini).
\end{lemma}

\begin{proof}
Puisque $\{X \to Y\}$ est un recouvrement fppf, il induit une surjection de
faisceaux fppf (Topologies sur les espaces, lemme
\ref{spaces-topologies-lemma-fppf-covering-surjective}); le
lemme est donc un cas particulier du lemme \ref{spaces-descent-lemma-curiosity}.
On peut aussi le déduire plus directement de
l'analogue pour les schémas. La question est en effet locale pour la topologie \'etale
sur $B$ et sur $Y$ (Morphismes d'espaces, lemmes
\ref{spaces-morphisms-lemma-finite-type-local} et
\ref{spaces-morphisms-lemma-finite-presentation-local}).
Nous pouvons donc supposer que $B$ et $Y$ sont des
schémas affines. Comme l'application $|X| \to |Y|$ est ouverte
(Morphismes d'espaces, lemme \ref{spaces-morphisms-lemma-fppf-open}),
nous pouvons choisir un
schéma affine $U$ et un morphisme \'etale $U \to X$ tel que le
morphisme composé $U \to Y$ soit surjectif. Dans ce cas, le résultat
se déduit de Descente, lemme
\ref{descent-lemma-flat-finitely-presented-permanence}.
\end{proof}

\begin{lemma}
\label{spaces-descent-lemma-syntomic-smooth-etale-permanence}
Soit $S$ un schéma. Considérons le diagramme suivant
$$
\xymatrix{
X \ar[rr]_f \ar[rd]_p & &
Y \ar[dl]^q \\
& B
}
$$
de morphismes d'espaces algébriques sur $S$, supposé commutatif.
Supposons que
\begin{enumerate}
\item $f$ soit surjectif et syntomique (resp.\ lisse, resp.\ \'etale),
\item $p$ soit syntomique (resp.\ lisse, resp.\ \'etale).
\end{enumerate}
Alors $q$ est syntomique (resp.\ lisse, resp.\ \'etale).
\end{lemma}

\begin{proof}
Le résultat se déduit de l'analogue pour les schémas.
La question est en effet locale pour la topologie \'etale sur $B$ et sur $Y$
(Morphismes d'espaces, lemmes
\ref{spaces-morphisms-lemma-syntomic-local},
\ref{spaces-morphisms-lemma-smooth-local}, et
\ref{spaces-morphisms-lemma-etale-local}).
Nous pouvons donc supposer que $B$ et $Y$ sont des
schémas affines. Comme l'application $|X| \to |Y|$ est ouverte
(Morphismes d'espaces, lemme \ref{spaces-morphisms-lemma-fppf-open}),
nous pouvons choisir un
schéma affine $U$ et un morphisme \'etale $U \to X$ tel que le
morphisme composé $U \to Y$ soit surjectif. Dans ce cas, le résultat
se déduit de Descente, lemme
\ref{descent-lemma-syntomic-smooth-etale-permanence}.
\end{proof}

\noindent
En fait, ce résultat peut être renforcé comme suit.

\begin{lemma}
\label{spaces-descent-lemma-smooth-permanence}
Soit $S$ un schéma. Considérons le diagramme suivant
$$
\xymatrix{
X \ar[rr]_f \ar[rd]_p & &
Y \ar[dl]^q \\
& B
}
$$
de morphismes d'espaces algébriques sur $S$, supposé commutatif. Supposons que
\begin{enumerate}
\item $f$ soit surjectif, plat et localement de présentation finie,
\item $p$ soit lisse (resp.\ \'etale).
\end{enumerate}
Alors $q$ est lisse (resp.\ \'etale).
\end{lemma}

\begin{proof}
Le résultat se déduit de l'analogue pour les schémas.
La question est en effet locale pour la topologie \'etale sur $B$ et sur $Y$
(Morphismes d'espaces, lemmes
\ref{spaces-morphisms-lemma-smooth-local} et
\ref{spaces-morphisms-lemma-etale-local}).
Nous pouvons donc supposer que $B$ et $Y$ sont des
schémas affines. Comme l'application $|X| \to |Y|$ est ouverte
(Morphismes d'espaces, lemme \ref{spaces-morphisms-lemma-fppf-open}),
nous pouvons choisir un
schéma affine $U$ et un morphisme \'etale $U \to X$ tel que le
morphisme composé $U \to Y$ soit surjectif. Dans ce cas, le résultat
se déduit de Descente, lemme
\ref{descent-lemma-smooth-permanence}.
\end{proof}

\begin{lemma}
\label{spaces-descent-lemma-syntomic-permanence}
Soit $S$ un schéma. Considérons le diagramme suivant
$$
\xymatrix{
X \ar[rr]_f \ar[rd]_p & &
Y \ar[dl]^q \\
& B
}
$$
de morphismes d'espaces algébriques sur $S$, supposé commutatif. Supposons que
\begin{enumerate}
\item $f$ soit surjectif, plat et localement de présentation finie,
\item $p$ soit syntomique.
\end{enumerate}
Alors $q$ et $f$ sont tous deux syntomiques.
\end{lemma}

\begin{proof}
Le résultat se déduit de l'analogue pour les schémas.
La question est en effet locale pour la topologie \'etale sur $B$ et sur $Y$
(Morphismes d'espaces, lemme
\ref{spaces-morphisms-lemma-syntomic-local}).
Nous pouvons donc supposer que $B$ et $Y$ sont des
schémas affines. Comme l'application $|X| \to |Y|$ est ouverte
(Morphismes d'espaces, lemme \ref{spaces-morphisms-lemma-fppf-open}),
nous pouvons choisir un
schéma affine $U$ et un morphisme \'etale $U \to X$ tel que le
morphisme composé $U \to Y$ soit surjectif. Dans ce cas, le résultat
se déduit de Descente, lemme
\ref{descent-lemma-syntomic-permanence}.
\end{proof}
% preserved blank authority line 843
% preserved blank authority line 844
% preserved blank authority line 845
% preserved blank authority line 846
% preserved blank authority line 847
% preserved blank authority line 848




\section{Descente de propriétés des espaces}
\label{spaces-descent-section-descending-properties-spaces}

\noindent
Dans cette section, nous rassemblons quelques résultats du type suivant.

\begin{lemma}
\label{spaces-descent-lemma-descend-unibranch}
Soit $S$ un schéma.
Soit $f : X \to Y$ un morphisme d'espaces algébriques sur $S$.
Soit $x \in |X|$.
Si $f$ est plat en $x$ et si $X$ est géométriquement unibranche en $x$, alors $Y$ est
géométriquement unibranche en $f(x)$.
\end{lemma}

\begin{proof}
Considérons l'homomorphisme d'anneaux strictement locaux
$\mathcal{O}_{Y, f(\overline{x})} \to \mathcal{O}_{X, \overline{x}}$.
Par
Morphismes d'espaces, lemme
\ref{spaces-morphisms-lemma-flat-at-point-etale-local-rings}
celui-ci est plat. Ainsi, si $\mathcal{O}_{X, \overline{x}}$ possède un unique idéal premier minimal,
il en va de même de $\mathcal{O}_{Y, f(\overline{x})}$ (par le théorème de descente des idéaux premiers, voir
Algèbre, lemme \ref{algebra-lemma-flat-going-down}).
\end{proof}

\begin{lemma}
\label{spaces-descent-lemma-descend-reduced}
\begin{slogan}
Un morphisme plat et surjectif d'espaces algébriques dont la source est réduite
a un but réduit.
\end{slogan}
Soit $S$ un schéma.
Soit $f : X \to Y$ un morphisme d'espaces algébriques sur $S$.
Si $f$ est plat et surjectif et si $X$ est réduit, alors $Y$ est réduit.
\end{lemma}

\begin{proof}
Choisissons un schéma $V$ et un morphisme \'etale surjectif $V \to Y$.
Choisissons un schéma $U$ et un morphisme \'etale surjectif
$U \to X \times_Y V$. Comme $f$ est surjectif et plat, le morphisme de
schémas $U \to V$ est surjectif et plat. Nous nous ramenons ainsi
au cas des schémas (le caractère réduit de $X$ et de $Y$ étant défini
en termes du caractère réduit de $U$ et de $V$, voir
Propriétés des espaces,
section \ref{spaces-properties-section-types-properties}).
Le cas des schémas est traité dans
Descente, lemme \ref{descent-lemma-descend-reduced}.
\end{proof}

\begin{lemma}
\label{spaces-descent-lemma-descend-locally-Noetherian}
Soit $f : X \to Y$ un morphisme d'espaces algébriques.
Si $f$ est localement de présentation finie, plat et surjectif et si
$X$ est localement noethérien, alors $Y$ est localement noethérien.
\end{lemma}

\begin{proof}
Choisissons un schéma $V$ et un morphisme \'etale surjectif $V \to Y$.
Choisissons un schéma $U$ et un morphisme \'etale surjectif
$U \to X \times_Y V$. Comme $f$ est surjectif, plat et localement de
présentation finie, le morphisme de schémas $U \to V$ est surjectif, plat et
localement de présentation finie. Nous nous ramenons ainsi
au cas des schémas (la propriété, pour $X$ et $Y$, d'être localement noethériens
étant définie par la même propriété pour $U$ et $V$, voir
Propriétés des espaces,
section \ref{spaces-properties-section-types-properties}).
Dans le cas des schémas, le résultat découle de
Descente, lemme \ref{descent-lemma-Noetherian-local-fppf}.
\end{proof}

\begin{lemma}
\label{spaces-descent-lemma-descend-regular}
Soit $f : X \to Y$ un morphisme d'espaces algébriques.
Si $f$ est localement de présentation finie, plat et surjectif et si
$X$ est régulier, alors $Y$ est régulier.
\end{lemma}

\begin{proof}
Par
le lemme \ref{spaces-descent-lemma-descend-locally-Noetherian}
nous savons que $Y$ est localement noethérien.
Choisissons un schéma $V$ et un morphisme \'etale surjectif $V \to Y$.
Il suffit de montrer que tous les anneaux locaux de $V$ sont réguliers ;
voir
Propriétés, lemme \ref{properties-lemma-characterize-regular}.
Choisissons un schéma $U$ et un morphisme \'etale surjectif
$U \to X \times_Y V$. Comme $f$ est surjectif et plat, le morphisme de schémas
$U \to V$ est surjectif et plat. Par hypothèse, $U$ est un schéma régulier
et, en particulier, tous ses anneaux locaux sont réguliers (d'après le lemme ci-dessus).
Le lemme découle donc de
Algèbre, lemme \ref{algebra-lemma-flat-under-regular}.
\end{proof}

\begin{lemma}
\label{spaces-descent-lemma-reduced-local-smooth}
Soit $f : X \to Y$ un morphisme lisse d'espaces algébriques.
Si $Y$ est réduit, alors $X$ est réduit. Si $f$ est surjectif
et si $X$ est réduit, alors $Y$ est réduit.
\end{lemma}

\begin{proof}
Choisissons un diagramme commutatif
$$
\xymatrix{
U \ar[d] \ar[r] & V \ar[d] \\
X \ar[r] & Y
}
$$
où $U$ et $V$ sont des schémas, les flèches verticales sont surjectives et
\'etales, et $U \to X \times_Y V$ est surjectif et \'etale. Observons que $X$ est
un espace algébrique réduit si et seulement si $U$ est un schéma réduit
par notre définition des espaces algébriques réduits dans
Propriétés des espaces, section \ref{spaces-properties-section-types-properties}.
De même pour $Y$ et $V$. 
Le morphisme $U \to V$ est un morphisme lisse de schémas, voir
Morphismes d'espaces, lemme \ref{spaces-morphisms-lemma-smooth-local}.
Puisque le caractère réduit est local pour la topologie lisse dans la catégorie des
schémas (Descente, lemme \ref{descent-lemma-reduced-local-smooth}),
il s'ensuit que $U$ est réduit si $V$ est réduit.
D'autre part, si $X \to Y$ est surjectif, alors $U \to V$ est
surjectif et, dans ce cas, si $U$ est réduit, alors $V$ est réduit.
\end{proof}







\section{Descente de propriétés des morphismes}
\label{spaces-descent-section-descending-properties-morphisms}

\noindent
Dans cette section, nous définissons ce que signifie, pour une propriété des morphismes
d'espaces algébriques, être locale sur le but pour une topologie. Comparer avec
Descente, section \ref{descent-section-descending-properties-morphisms}.

\begin{definition}
\label{spaces-descent-definition-property-morphisms-local}
Soit $S$ un schéma.
Soit $\mathcal{P}$ une propriété des morphismes d'espaces algébriques sur $S$.
Soit $\tau \in \{fpqc, fppf, syntomic, smooth, \etale\}$.
On dit que $\mathcal{P}$ est {\it $\tau$-locale sur la base}, ou
{\it $\tau$-locale sur le but}, ou
{\it locale sur la base pour la topologie $\tau$} si, pour tout
$\tau$-recouvrement $\{Y_i \to Y\}_{i \in I}$ d'espaces algébriques
et tout morphisme d'espaces algébriques $f : X \to Y$, on
a
$$
f \text{ vérifie }\mathcal{P}
\Leftrightarrow
\text{chaque }Y_i \times_Y X \to Y_i\text{ vérifie }\mathcal{P}.
$$
\end{definition}

\noindent
Bien entendu, puisque les isomorphismes sont toujours des recouvrements,
on voit (ou l'on exige) que la propriété $\mathcal{P}$ est vérifiée par $X \to Y$
si et seulement si toute flèche $X' \to Y'$ isomorphe à $X \to Y$ la vérifie.
Si une propriété est $\tau$-locale sur le but, elle est préservée
par changement de base le long de morphismes figurant dans des $\tau$-recouvrements. Voici
un énoncé formel.

\begin{lemma}
\label{spaces-descent-lemma-pullback-property-local-target}
Soit $S$ un schéma.
Soit $\tau \in \{fpqc, fppf, syntomic, smooth, \etale\}$.
Soit $\mathcal{P}$ une propriété des morphismes d'espaces algébriques sur $S$
qui est $\tau$-locale sur le but. Soit $f : X \to Y$ un morphisme vérifiant la propriété
$\mathcal{P}$. Pour tout morphisme $Y' \to Y$ qui est
plat, resp.\ plat et localement de présentation finie, resp.\ syntomique,
resp.\ \'etale, le changement de base
$f' : Y' \times_Y X \to Y'$ de $f$ vérifie la propriété $\mathcal{P}$.
\end{lemma}

\begin{proof}
Cela résulte de ce que l'on peut compléter $Y' \to Y$ en une famille de
morphismes qui constitue un $\tau$-recouvrement.
\end{proof}

\noindent
Une conséquence simple et souvent utilisée de ce qui précède est la suivante. Si
$f : X \to Y$ vérifie la propriété $\mathcal{P}$, laquelle est $\tau$-locale
sur le but, et si $f(X) \subset V$
pour un certain sous-espace ouvert $V \subset Y$, alors le morphisme induit
$X \to V$ vérifie encore $\mathcal{P}$. Démonstration : le changement de base de
$f$ par $V \to Y$ donne $X \to V$.

\begin{lemma}
\label{spaces-descent-lemma-largest-open-of-the-base}
Soit $S$ un schéma.
Soit $\tau \in \{fppf, syntomic, smooth, \etale\}$.
Soit $\mathcal{P}$ une propriété des morphismes d'espaces algébriques sur $S$
qui est $\tau$-locale sur le but. Pour tout morphisme d'espaces algébriques
$f : X \to Y$ sur $S$, il existe un plus grand sous-espace ouvert
$W(f) \subset Y$ tel que la restriction $X_{W(f)} \to W(f)$ vérifie
$\mathcal{P}$. En outre,
\begin{enumerate}
\item si $g : Y' \to Y$ est un morphisme d'espaces algébriques qui est
plat et localement de présentation finie, syntomique, lisse ou \'etale,
et si le changement de base $f' : X_{Y'} \to Y'$ vérifie $\mathcal{P}$, alors
$g$ se factorise par $W(f)$,
\item si $g : Y' \to Y$ est plat et localement de présentation finie,
syntomique, lisse ou \'etale, alors $W(f') = g^{-1}(W(f))$, et
\item si $\{g_i : Y_i \to Y\}$ est un $\tau$-recouvrement, alors
$g_i^{-1}(W(f)) = W(f_i)$, où $f_i$ est le changement de base de $f$
par $Y_i \to Y$.
\end{enumerate}
\end{lemma}

\begin{proof}
Considérons la réunion $W_{set} \subset |Y|$ des images
$g(|Y'|) \subset |Y|$ des morphismes $g : Y' \to Y$ satisfaisant aux conditions suivantes :
\begin{enumerate}
\item $g$ est plat et localement de présentation finie, syntomique,
lisse ou \'etale, et
\item le changement de base $Y' \times_{g, Y} X \to Y'$ vérifie la propriété
$\mathcal{P}$.
\end{enumerate}
Comme un tel morphisme $g$ est ouvert (voir
Morphismes d'espaces, lemme \ref{spaces-morphisms-lemma-fppf-open}),
on voit que $W_{set}$ est une partie ouverte de $|Y|$. Notons $W \subset Y$
le sous-espace ouvert dont l'ensemble de points est $W_{set}$, voir
Propriétés des espaces, lemme \ref{spaces-properties-lemma-open-subspaces}.
Comme $\mathcal{P}$ est locale pour la topologie $\tau$, la restriction
$X_W \to W$ vérifie la propriété $\mathcal{P}$, car on dispose d'un recouvrement
$\{Y' \to W\}$ de $W$ tel que les changements de base vérifient $\mathcal{P}$.
Cela établit l'existence et montre que $W(f)$ satisfait à (1).
Pour établir (2), remarquons que $W(f') \supset g^{-1}(W(f))$, car
$\mathcal{P}$ est stable par changement de base le long de morphismes plats et localement de
présentation finie, syntomiques, lisses ou \'etales ; voir
le lemme \ref{spaces-descent-lemma-pullback-property-local-target}.
Réciproquement, si $Y'' \subset Y'$ est un ouvert tel que
$X_{Y''} \to Y''$ vérifie la propriété $\mathcal{P}$, alors $Y'' \to Y$ se factorise
par $W$ en vertu de la construction, c'est-à-dire $Y'' \subset g^{-1}(W(f))$. Cela
prouve (2). L'assertion (3) découle de (2), car chaque morphisme
$Y_i \to Y$ est plat et localement de présentation finie, syntomique,
lisse ou \'etale d'après notre définition d'un $\tau$-recouvrement.
\end{proof}

\begin{lemma}
\label{spaces-descent-lemma-descending-properties-morphisms}
Soit $S$ un schéma. Soit $\mathcal{P}$ une propriété des morphismes
d'espaces algébriques sur $S$. Supposons que
\begin{enumerate}
\item si les morphismes $X_i \to Y_i$, $i = 1, 2$, vérifient la propriété $\mathcal{P}$, alors
$X_1 \amalg X_2 \to Y_1 \amalg Y_2$ la vérifie aussi,
\item un morphisme d'espaces algébriques $f : X \to Y$ vérifie la propriété
$\mathcal{P}$ si et seulement si, pour tout schéma affine $Z$ et tout
morphisme $Z \to Y$, le changement de base $Z \times_Y X \to Z$ de $f$
vérifie la propriété $\mathcal{P}$, et
\item pour tout morphisme plat et surjectif de schémas affines
$Z' \to Z$ sur $S$ et tout morphisme $f : X \to Z$ d'un espace algébrique
vers $Z$, on a
$$
f' : Z' \times_Z X \to Z'\text{ vérifie }\mathcal{P}
\Rightarrow
f\text{ vérifie }\mathcal{P}.
$$
\end{enumerate}
Alors $\mathcal{P}$ est locale sur la base pour la topologie fpqc.
\end{lemma}

\begin{proof}
Si $\mathcal{P}$ satisfait à la condition (2), elle est automatiquement
stable par tout changement de base. L'implication directe de la
définition \ref{spaces-descent-definition-property-morphisms-local} est donc acquise.

\medskip\noindent
Soit $\{Y_i \to Y\}_{i \in I}$ un recouvrement fpqc d'espaces algébriques sur $S$.
Soit $f : X \to Y$ un morphisme d'espaces algébriques sur $S$.
Supposons que chaque changement de base $f_i : Y_i \times_Y X \to Y_i$ vérifie la propriété
$\mathcal{P}$. Nous voulons montrer que $f$ vérifie $\mathcal{P}$.
Soit $Z$ un schéma affine et soit $Z \to Y$ un morphisme.
D'après (2), il suffit de montrer que le morphisme d'espaces algébriques
$Z \times_Y X \to Z$ vérifie $\mathcal{P}$.
Puisque $\{Y_i \to Y\}_{i \in I}$ est un recouvrement fpqc, il existe
un recouvrement fpqc standard $\{Z_j \to Z\}_{j = 1, \ldots , n}$
et des morphismes $Z_j \to Y_{i_j}$ au-dessus de $Y$, pour des indices convenables $i_j \in I$.
Puisque $f_{i_j}$ vérifie $\mathcal{P}$, on voit que
$$
Z_j \times_Y X
=
Z_j \times_{Y_{i_j}} (Y_{i_j} \times_Y X)
\longrightarrow
Z_j
$$
vérifie $\mathcal{P}$, car c'est un changement de base de $f_{i_j}$ (voir la première remarque de la
démonstration). Posons $Z' = \coprod_{j = 1, \ldots, n} Z_j$ ; alors $Z' \to Z$ est
un morphisme plat et surjectif de schémas affines sur $S$. D'après (1),
on conclut que $Z' \times_Y X \to Z'$ vérifie la propriété $\mathcal{P}$.
Comme il s'agit du changement de base du morphisme $Z \times_Y X \to Z$
par le morphisme $Z' \to Z$, on conclut que
$Z \times_Y X \to Z$ vérifie la propriété $\mathcal{P}$, comme souhaité.
\end{proof}
	
	
\section{Descente de propriétés des morphismes pour la topologie fpqc}
\sectionmark{Descente de propriétés pour la topologie fpqc}
\label{spaces-descent-section-descending-properties-morphisms-fpqc}


\noindent
Dans cette section, nous montrons qu'un grand nombre de propriétés
de morphismes d'espaces algébriques sont locales sur la base
pour la topologie fpqc. Comparer avec
Descente, section \ref{descent-section-descending-properties-morphisms-fpqc}
pour le cas des morphismes de schémas.

\begin{lemma}
\label{spaces-descent-lemma-descending-property-quasi-compact}
Soit $S$ un schéma.
La propriété $\mathcal{P}(f) =$``$f$ est quasi-compact''
est locale sur la base pour la topologie fpqc dans la catégorie des espaces algébriques sur $S$.
\end{lemma}

\begin{proof}
Nous appliquerons le
lemme \ref{spaces-descent-lemma-descending-properties-morphisms}.
Ses hypothèses (1) et (2) résultent de
Morphismes des espaces,
Lemme \ref{spaces-morphisms-lemma-quasi-compact-local}.
Soit $Z' \to Z$ un morphisme plat et surjectif de schémas affines sur $S$.
Soit $f : X \to Z$ un morphisme d'espaces algébriques, et supposons
que le changement de base $f' : Z' \times_Z X \to Z'$ soit quasi-compact. Il faut
montrer que $f$ est quasi-compact. Il résulte de nouveau de
Morphismes des espaces,
Lemme \ref{spaces-morphisms-lemma-quasi-compact-local},
qu'il suffit de montrer que, pour tout schéma affine $Y$ et tout
morphisme $Y \to Z$, le produit fibré $Y \times_Z X$ est quasi-compact.
Voici le diagramme correspondant :
\begin{equation}
\label{spaces-descent-equation-cube}
\vcenter{
\xymatrix{
Y \times_Z Z' \times_Z X \ar[dd] \ar[rr] \ar[rd] & &
Z' \times_Z X \ar'[d][dd]^{f'} \ar[rd] \\
& Y \times_Z X \ar[dd] \ar[rr] & & X \ar[dd]^f \\
Y \times_Z Z' \ar'[r][rr] \ar[rd] & & Z' \ar[rd] \\
& Y \ar[rr] & & Z
}
}
\end{equation}
Tous les carrés sont cartésiens, et le carré inférieur est formé
de schémas affines. L'hypothèse selon laquelle $f'$ est quasi-compact, jointe au
fait que $Y \times_Z Z'$ est affine, implique que
$Y \times_Z Z' \times_Z X$ est quasi-compact. Comme
$$
Y \times_Z Z' \times_Z X \longrightarrow Y \times_Z X
$$
est surjectif, puisqu'il s'obtient par changement de base à partir de $Z' \to Z$,
on en conclut que $Y \times_Z X$ est quasi-compact ; voir
Morphismes des espaces,
Lemme \ref{spaces-morphisms-lemma-surjection-from-quasi-compact}.
Cela achève la démonstration.
\end{proof}

\begin{lemma}
\label{spaces-descent-lemma-descending-property-quasi-separated}
Soit $S$ un schéma.
La propriété $\mathcal{P}(f) =$``$f$ est quasi-séparé''
est locale sur la base pour la topologie fpqc dans la catégorie des espaces algébriques sur $S$.
\end{lemma}

\begin{proof}
Tout changement de base d'un morphisme quasi-séparé est quasi-séparé ; voir
Morphismes des espaces,
Lemme \ref{spaces-morphisms-lemma-base-change-separated}.
On obtient donc le sens direct de la
définition \ref{spaces-descent-definition-property-morphisms-local}.

\medskip\noindent
Soit $\{Y_i \to Y\}_{i \in I}$ un recouvrement fpqc d'espaces algébriques sur $S$.
Soit $f : X \to Y$ un morphisme d'espaces algébriques sur $S$.
Supposons que chaque changement de base $X_i := Y_i \times_Y X \to Y_i$ soit quasi-séparé.
Cela signifie que chacun des morphismes
$$
\Delta_i :
X_i
\longrightarrow
X_i \times_{Y_i} X_i = Y_i \times_Y (X \times_Y X)
$$
est quasi-compact. Le changement de base d'un recouvrement fpqc est un recouvrement fpqc ; voir
Topologies sur les espaces, Lemme \ref{spaces-topologies-lemma-fpqc} ;
ainsi, $\{Y_i \times_Y (X \times_Y X) \to X \times_Y X\}$
est un recouvrement fpqc d'espaces algébriques. De plus, chaque
$\Delta_i$ est le changement de base du morphisme
$\Delta : X \to X \times_Y X$. Il résulte donc du
lemme \ref{spaces-descent-lemma-descending-property-quasi-compact}
que $\Delta$ est quasi-compact, c'est-à-dire que $f$ est quasi-séparé.
\end{proof}

\begin{lemma}
\label{spaces-descent-lemma-descending-property-universally-closed}
Soit $S$ un schéma.
La propriété $\mathcal{P}(f) =$``$f$ est universellement fermé''
est locale sur la base pour la topologie fpqc dans la catégorie des espaces algébriques sur $S$.
\end{lemma}

\begin{proof}
Nous appliquerons le
lemme \ref{spaces-descent-lemma-descending-properties-morphisms}.
Ses hypothèses (1) et (2) résultent de
Morphismes des espaces,
Lemme \ref{spaces-morphisms-lemma-universally-closed-local}.
Soit $Z' \to Z$ un morphisme plat et surjectif de schémas affines sur $S$.
Soit $f : X \to Z$ un morphisme d'espaces algébriques, et supposons
que le changement de base $f' : Z' \times_Z X \to Z'$ soit universellement fermé.
Il faut montrer que $f$ est universellement fermé. Il résulte de nouveau de
Morphismes des espaces,
Lemme \ref{spaces-morphisms-lemma-universally-closed-local},
qu'il suffit de montrer que, pour tout schéma affine $Y$ et tout
morphisme $Y \to Z$, l'application $|Y \times_Z X| \to |Y|$ est fermée.
Considérons le cube (\ref{spaces-descent-equation-cube}).
L'hypothèse selon laquelle $f'$ est universellement fermé implique que
$|Y \times_Z Z' \times_Z X| \to |Y \times_Z Z'|$ est fermée.
Comme $Y \times_Z Z' \to Y$ est quasi-compact, surjectif et plat,
puisqu'il s'obtient par changement de base à partir de $Z' \to Z$,
l'application $|Y \times_Z Z'| \to |Y|$ est submersive ; voir
Morphismes, Lemme \ref{morphisms-lemma-fpqc-quotient-topology}.
De plus, l'application
$$
|Y \times_Z Z' \times_Z X|
\longrightarrow
|Y \times_Z Z'| \times_{|Y|} |Y \times_Z X|
$$
est surjective ; voir
Propriétés des espaces, Lemme \ref{spaces-properties-lemma-points-cartesian}.
Il résulte de la topologie élémentaire que $|Y \times_Z X| \to |Y|$ est fermée.
\end{proof}

\begin{lemma}
\label{spaces-descent-lemma-descending-property-universally-open}
Soit $S$ un schéma.
La propriété $\mathcal{P}(f) =$``$f$ est universellement ouvert''
est locale sur la base pour la topologie fpqc dans la catégorie des espaces algébriques sur $S$.
\end{lemma}

\begin{proof}
La démonstration est identique à celle du
lemme \ref{spaces-descent-lemma-descending-property-universally-closed}.
\end{proof}

\begin{lemma}
\label{spaces-descent-lemma-descending-property-universally-submersive}
La propriété $\mathcal{P}(f) =$``$f$ est universellement submersif''
est locale sur la base pour la topologie fpqc.
\end{lemma}

\begin{proof}
La démonstration est identique à celle du
lemme \ref{spaces-descent-lemma-descending-property-universally-closed}.
\end{proof}

\begin{lemma}
\label{spaces-descent-lemma-descending-property-surjective}
La propriété $\mathcal{P}(f) =$``$f$ est surjectif''
est locale sur la base pour la topologie fpqc.
\end{lemma}

\begin{proof}
Démonstration omise. (Indication : utiliser
Propriétés des espaces, Lemme \ref{spaces-properties-lemma-points-cartesian}.)
\end{proof}

\begin{lemma}
\label{spaces-descent-lemma-descending-property-universally-injective}
La propriété $\mathcal{P}(f) =$``$f$ est universellement injectif''
est locale sur la base pour la topologie fpqc.
\end{lemma}

\begin{proof}
Nous appliquerons le
lemme \ref{spaces-descent-lemma-descending-properties-morphisms}.
Ses hypothèses (1) et (2) résultent de
Morphismes des espaces,
Lemme \ref{spaces-morphisms-lemma-universally-closed-local}.
Soit $Z' \to Z$ un morphisme plat et surjectif de schémas affines
sur $S$, et soit $f : X \to Z$ un morphisme d'un espace algébrique vers $Z$.
Supposons que le changement de base $f' : X' \to Z'$ soit universellement injectif.
Soit $K$ un corps, et soient $a, b : \Spec(K) \to X$
deux morphismes tels que $f \circ a = f \circ b$.
Comme $Z' \to Z$ est surjectif, il existe une extension de corps
$K'/K$ et un morphisme
$\Spec(K') \to Z'$
de sorte que le diagramme suivant, où les flèches pleines sont données, soit commutatif :
$$
\xymatrix{
\Spec(K') \ar[rrd] \ar@{-->}[rd]_{a', b'} \ar[dd] \\
 &
X' \ar[r] \ar[d] &
Z' \ar[d] \\
\Spec(K) \ar[r]^{a, b} &
X \ar[r] &
Z
}
$$
Comme le carré est cartésien, on obtient les deux flèches pointillées $a'$, $b'$ qui rendent le
diagramme commutatif. Puisque $X' \to Z'$ est universellement injectif, on a $a' = b'$.
Cela entraîne $a = b$, car $\{\Spec(K') \to \Spec(K)\}$
est un recouvrement fpqc ; voir
Propriétés des espaces, Proposition
\ref{spaces-properties-proposition-sheaf-fpqc}.
Ainsi, $f$ est universellement injectif, ce qu'il fallait démontrer.
\end{proof}

\begin{lemma}
\label{spaces-descent-lemma-descending-property-universal-homeomorphism}
La propriété $\mathcal{P}(f) =$``$f$ est un homéomorphisme universel''
est locale sur la base pour la topologie fpqc.
\end{lemma}

\begin{proof}
On peut le démontrer exactement de la même manière que dans le
lemme \ref{spaces-descent-lemma-descending-property-universally-closed}.
On peut aussi utiliser le fait qu'une application d'espaces topologiques est un
homéomorphisme si et seulement si elle est injective, surjective et ouverte.
Ainsi, un homéomorphisme universel équivaut à un morphisme
surjectif, universellement injectif et universellement ouvert.
Voir Morphismes des espaces, Lemme
\ref{spaces-morphisms-lemma-base-change-surjective} et
Morphismes des espaces, Définitions
\ref{spaces-morphisms-definition-universally-injective},
\ref{spaces-morphisms-definition-surjective},
\ref{spaces-morphisms-definition-open},
\ref{spaces-morphisms-definition-universal-homeomorphism}.
Le lemme résulte donc des
lemmes \ref{spaces-descent-lemma-descending-property-surjective},
\ref{spaces-descent-lemma-descending-property-universally-injective} et
\ref{spaces-descent-lemma-descending-property-universally-open}.
\end{proof}

\begin{lemma}
\label{spaces-descent-lemma-descending-property-locally-finite-type}
La propriété $\mathcal{P}(f) =$``$f$ est localement de type fini''
est locale sur la base pour la topologie fpqc.
\end{lemma}

\begin{proof}
Nous appliquerons le
lemme \ref{spaces-descent-lemma-descending-properties-morphisms}.
Ses hypothèses (1) et (2) résultent de
Morphismes des espaces,
Lemme \ref{spaces-morphisms-lemma-finite-type-local}.
Soit $Z' \to Z$ un morphisme plat et surjectif de schémas affines sur $S$.
Soit $f : X \to Z$ un morphisme d'espaces algébriques, et supposons
que le changement de base $f' : Z' \times_Z X \to Z'$ soit localement de type fini.
Il faut montrer que $f$ est localement de type fini. Soit $U$ un schéma
et soit $U \to X$ un morphisme étale surjectif. D'après
Morphismes des espaces,
Lemme \ref{spaces-morphisms-lemma-finite-type-local},
il suffit encore de montrer que $U \to Z$ est localement de type fini.
Puisque $f'$ est localement de type fini et que $Z' \times_Z U$ est un
schéma étale sur $Z' \times_Z X$, on déduit encore du même lemme que
$Z' \times_Z U \to Z'$ est localement de type fini.
Comme $\{Z' \to Z\}$ est un recouvrement fpqc, on en conclut que
$U \to Z$ est localement de type fini d'après
Descente, Lemme \ref{descent-lemma-descending-property-locally-finite-type},
ce qu'il fallait démontrer.
\end{proof}

\begin{lemma}
\label{spaces-descent-lemma-descending-property-locally-finite-presentation}
La propriété $\mathcal{P}(f) =$``$f$ est localement de présentation finie''
est locale sur la base pour la topologie fpqc.
\end{lemma}

\begin{proof}
Nous appliquerons le
lemme \ref{spaces-descent-lemma-descending-properties-morphisms}.
Ses hypothèses (1) et (2) résultent de
Morphismes des espaces,
Lemme \ref{spaces-morphisms-lemma-finite-presentation-local}.
Soit $Z' \to Z$ un morphisme plat et surjectif de schémas affines sur $S$.
Soit $f : X \to Z$ un morphisme d'espaces algébriques, et supposons
que le changement de base $f' : Z' \times_Z X \to Z'$ soit localement de
présentation finie.
Il faut montrer que $f$ est localement de présentation finie. Soit $U$ un schéma
et soit $U \to X$ un morphisme étale surjectif. D'après
Morphismes des espaces,
Lemme \ref{spaces-morphisms-lemma-finite-presentation-local},
il suffit encore de montrer que $U \to Z$ est localement de présentation finie.
Puisque $f'$ est localement de présentation finie et que $Z' \times_Z U$ est un
schéma étale sur $Z' \times_Z X$, on déduit encore du même lemme que
$Z' \times_Z U \to Z'$ est localement de présentation finie.
Comme $\{Z' \to Z\}$ est un recouvrement fpqc, on en conclut que
$U \to Z$ est localement de présentation finie d'après
Descente,
Lemme \ref{descent-lemma-descending-property-locally-finite-presentation},
ce qu'il fallait démontrer.
\end{proof}

\begin{lemma}
\label{spaces-descent-lemma-descending-property-finite-type}
La propriété $\mathcal{P}(f) =$``$f$ est de type fini''
est locale sur la base pour la topologie fpqc.
\end{lemma}

\begin{proof}
Il suffit de combiner les lemmes \ref{spaces-descent-lemma-descending-property-quasi-compact}
et \ref{spaces-descent-lemma-descending-property-locally-finite-type}.
\end{proof}

\begin{lemma}
\label{spaces-descent-lemma-descending-property-finite-presentation}
La propriété $\mathcal{P}(f) =$``$f$ est de présentation finie''
est locale sur la base pour la topologie fpqc.
\end{lemma}

\begin{proof}
Il suffit de combiner les lemmes \ref{spaces-descent-lemma-descending-property-quasi-compact},
\ref{spaces-descent-lemma-descending-property-quasi-separated} et
\ref{spaces-descent-lemma-descending-property-locally-finite-presentation}.
\end{proof}

\begin{lemma}
\label{spaces-descent-lemma-descending-property-flat}
La propriété $\mathcal{P}(f) =$``$f$ est plat''
est locale sur la base pour la topologie fpqc.
\end{lemma}

\begin{proof}
Nous appliquerons le
lemme \ref{spaces-descent-lemma-descending-properties-morphisms}.
Ses hypothèses (1) et (2) résultent de
Morphismes des espaces,
Lemme \ref{spaces-morphisms-lemma-flat-local}.
Soit $Z' \to Z$ un morphisme plat et surjectif de schémas affines sur $S$.
Soit $f : X \to Z$ un morphisme d'espaces algébriques, et supposons
que le changement de base $f' : Z' \times_Z X \to Z'$ soit plat.
Il faut montrer que $f$ est plat. Soit $U$ un schéma
et soit $U \to X$ un morphisme étale surjectif. D'après
Morphismes des espaces,
Lemme \ref{spaces-morphisms-lemma-flat-local},
il suffit encore de montrer que $U \to Z$ est plat.
Puisque $f'$ est plat et que $Z' \times_Z U$ est un
schéma étale sur $Z' \times_Z X$, on déduit encore du même lemme que
$Z' \times_Z U \to Z'$ est plat.
Comme $\{Z' \to Z\}$ est un recouvrement fpqc, on en conclut que
$U \to Z$ est plat d'après
Descente, Lemme \ref{descent-lemma-descending-property-flat},
ce qu'il fallait démontrer.
\end{proof}

\begin{lemma}
\label{spaces-descent-lemma-descending-property-open-immersion}
La propriété $\mathcal{P}(f) =$``$f$ est une immersion ouverte''
est locale sur la base pour la topologie fpqc.
\end{lemma}

\begin{proof}
Nous appliquerons le
lemme \ref{spaces-descent-lemma-descending-properties-morphisms}.
Ses hypothèses (1) et (2) résultent de
Morphismes des espaces,
Lemme \ref{spaces-morphisms-lemma-closed-immersion-local}.
Considérons un diagramme cartésien
$$
\xymatrix{
X' \ar[r] \ar[d] & X \ar[d] \\
Z' \ar[r] & Z
}
$$
d'espaces algébriques sur $S$
où $Z' \to Z$ est un morphisme plat et surjectif de schémas affines
et $X' \to Z'$ une immersion ouverte. Il faut montrer que $X \to Z$
est une immersion ouverte. L'ouvert $|X'| \subset |Z'|$ correspond à un
sous-schéma ouvert $U' \subset Z'$ (isomorphe à $X'$)
tel que $\text{pr}_0^{-1}(U') = \text{pr}_1^{-1}(U')$
comme sous-schémas ouverts de $Z' \times_Z Z'$. Il existe donc un
sous-schéma ouvert $U \subset Z$ tel que $X' = (Z' \to Z)^{-1}(U)$ ; voir
Descente, Lemme \ref{descent-lemma-open-fpqc-covering}.
D'après Propriétés des espaces,
Proposition \ref{spaces-properties-proposition-sheaf-fpqc},
on voit que $X$ vérifie la condition de faisceau pour la topologie fpqc.
Nous avons maintenant le recouvrement fpqc $\mathcal{U} = \{U' \to U\}$
et l'élément $U' \to X' \to X \in \check{H}^0(\mathcal{U}, X)$.
La condition de faisceau fournit un morphisme $U \to X$ tel que
$$
\xymatrix{
U' \ar[r] \ar[d]^{\cong} \ar@/_3ex/[dd] & U \ar[d] \ar@/^3ex/[dd] \\
X' \ar[r] \ar[d] & X \ar[d] \\
Z' \ar[r] & Z
}
$$
soit commutatif. D'autre part, pour tout schéma $T$ sur $S$
et tout point à valeurs dans $T$, c'est-à-dire tout morphisme $T \to X$, le composé $T \to X \to Z$ est un
morphisme tel que $Z' \times_Z T \to Z'$ se factorise par $U'$. Cela
signifie que $T \to Z$ se factorise par $U$. Autrement dit, le morphisme
de faisceaux $U \to X$ est bijectif, ce qui achève la démonstration.
\end{proof}

\begin{lemma}
\label{spaces-descent-lemma-descending-property-isomorphism}
La propriété $\mathcal{P}(f) =$``$f$ est un isomorphisme''
est locale sur la base pour la topologie fpqc.
\end{lemma}

\begin{proof}
Il suffit de combiner les lemmes \ref{spaces-descent-lemma-descending-property-surjective}
et \ref{spaces-descent-lemma-descending-property-open-immersion}.
\end{proof}

\begin{lemma}
\label{spaces-descent-lemma-descending-property-affine}
La propriété $\mathcal{P}(f) =$``$f$ est affine''
est locale sur la base pour la topologie fpqc.
\end{lemma}

\begin{proof}
Nous appliquerons le
lemme \ref{spaces-descent-lemma-descending-properties-morphisms}.
Ses hypothèses (1) et (2) résultent de
Morphismes des espaces,
Lemme \ref{spaces-morphisms-lemma-affine-local}.
Soit $Z' \to Z$ un morphisme plat et surjectif de schémas affines sur $S$.
Soit $f : X \to Z$ un morphisme d'espaces algébriques, et supposons
que le changement de base $f' : Z' \times_Z X \to Z'$ soit affine.
Soit $X'$ un schéma représentant $Z' \times_Z X$.
On obtient un isomorphisme canonique
$$
\varphi : X' \times_Z Z' \longrightarrow Z' \times_Z X'
$$
car les deux schémas représentent l'espace algébrique $Z' \times_Z Z' \times_Z X$.
C'est une donnée de descente pour $X'/Z'/Z$ ; voir
Descente, Définition \ref{descent-definition-descent-datum}
(vérification omise ; comparer avec
Descente, Lemme \ref{descent-lemma-descent-data-sheaves}).
Puisque $X' \to Z'$ est affine, cette donnée de descente est effective ; voir
Descente, Lemme \ref{descent-lemma-affine}.
Il existe donc un schéma $Y \to Z$ sur $Z$, ainsi qu'un
isomorphisme $\psi : Z' \times_Z Y \to X'$ compatible avec les données de descente.
Bien entendu, $Y \to Z$ est affine (par construction ou d'après
Descente, Lemme \ref{descent-lemma-descending-property-affine}).
Remarquons que $\mathcal{Y} = \{Z' \times_Z Y \to Y\}$ est un
recouvrement fpqc. En interprétant $\psi$ comme un élément de
$X(Z' \times_Z Y)$, on voit que $\psi \in \check{H}^0(\mathcal{Y}, X)$.
La condition de faisceau pour $X$ relativement à ce recouvrement (voir
Propriétés des espaces, Proposition
\ref{spaces-properties-proposition-sheaf-fpqc})
fournit un morphisme $Y \to X$.
Par construction, le changement de base de ce morphisme par $Z'$ est un isomorphisme ; le morphisme lui-même est donc
un isomorphisme d'après le
lemme \ref{spaces-descent-lemma-descending-property-isomorphism}.
Cela prouve que $X$ est représentable par un schéma affine, ce qui achève la démonstration.
\end{proof}

\begin{lemma}
\label{spaces-descent-lemma-descending-property-closed-immersion}
La propriété $\mathcal{P}(f) =$``$f$ est une immersion fermée''
est locale sur la base pour la topologie fpqc.
\end{lemma}

\begin{proof}
Nous appliquerons le
lemme \ref{spaces-descent-lemma-descending-properties-morphisms}.
Ses hypothèses (1) et (2) résultent de
Morphismes des espaces,
Lemme \ref{spaces-morphisms-lemma-closed-immersion-local}.
Considérons un diagramme cartésien
$$
\xymatrix{
X' \ar[r] \ar[d] & X \ar[d] \\
Z' \ar[r] & Z
}
$$
d'espaces algébriques sur $S$
où $Z' \to Z$ est un morphisme plat et surjectif de schémas affines
et $X' \to Z'$ une immersion fermée. Il faut montrer que $X \to Z$
est une immersion fermée. Le morphisme $X' \to Z'$ est affine. D'après le
lemme \ref{spaces-descent-lemma-descending-property-affine},
on voit que $X$ est un schéma et que $X \to Z$ est affine.
Il résulte de
Descente, Lemme \ref{descent-lemma-descending-property-closed-immersion},
que $X \to Z$ est une immersion fermée, ce qu'il fallait démontrer.
\end{proof}

\begin{lemma}
\label{spaces-descent-lemma-descending-property-separated}
La propriété $\mathcal{P}(f) =$``$f$ est séparé''
est locale sur la base pour la topologie fpqc.
\end{lemma}

\begin{proof}
Tout changement de base d'un morphisme séparé est séparé ; voir
Morphismes des espaces,
Lemme \ref{spaces-morphisms-lemma-base-change-separated}.
Cela donne donc l'implication directe de la
définition \ref{spaces-descent-definition-property-morphisms-local}.

\medskip\noindent
Soit $\{Y_i \to Y\}_{i \in I}$ un recouvrement fpqc d'espaces algébriques sur $S$.
Soit $f : X \to Y$ un morphisme d'espaces algébriques sur $S$.
Supposons que chaque changement de base $X_i := Y_i \times_Y X \to Y_i$ soit séparé.
Cela signifie que chacun des morphismes
$$
\Delta_i :
X_i
\longrightarrow
X_i \times_{Y_i} X_i = Y_i \times_Y (X \times_Y X)
$$
est une immersion fermée. Le changement de base d'un recouvrement fpqc est un
recouvrement fpqc ; voir
Topologies sur les espaces, Lemme \ref{spaces-topologies-lemma-fpqc} ;
ainsi, $\{Y_i \times_Y (X \times_Y X) \to X \times_Y X\}$
est un recouvrement fpqc d'espaces algébriques. De plus, chaque
$\Delta_i$ est le changement de base du morphisme
$\Delta : X \to X \times_Y X$. Il résulte donc du
lemme \ref{spaces-descent-lemma-descending-property-closed-immersion}
que $\Delta$ est une immersion fermée, c'est-à-dire que $f$ est séparé.
\end{proof}

\begin{lemma}
\label{spaces-descent-lemma-descending-property-proper}
La propriété $\mathcal{P}(f) =$``$f$ est propre''
est locale sur la base pour la topologie fpqc.
\end{lemma}

\begin{proof}
Le lemme résulte de la combinaison des
lemmes \ref{spaces-descent-lemma-descending-property-universally-closed},
\ref{spaces-descent-lemma-descending-property-separated}
et \ref{spaces-descent-lemma-descending-property-finite-type}.
\end{proof}

\begin{lemma}
\label{spaces-descent-lemma-descending-property-quasi-affine}
La propriété $\mathcal{P}(f) =$``$f$ est quasi-affine''
est locale sur la base pour la topologie fpqc.
\end{lemma}

\begin{proof}
Nous appliquerons le
lemme \ref{spaces-descent-lemma-descending-properties-morphisms}.
Ses hypothèses (1) et (2) résultent de
Morphismes des espaces,
Lemme \ref{spaces-morphisms-lemma-quasi-affine-local}.
Soit $Z' \to Z$ un morphisme plat et surjectif de schémas affines sur $S$.
Soit $f : X \to Z$ un morphisme d'espaces algébriques, et supposons
que le changement de base $f' : Z' \times_Z X \to Z'$ soit quasi-affine.
Soit $X'$ un schéma représentant $Z' \times_Z X$.
On obtient un isomorphisme canonique
$$
\varphi : X' \times_Z Z' \longrightarrow Z' \times_Z X'
$$
car les deux schémas représentent l'espace algébrique $Z' \times_Z Z' \times_Z X$.
C'est une donnée de descente pour $X'/Z'/Z$ ; voir
Descente, Définition \ref{descent-definition-descent-datum}
(vérification omise ; comparer avec
Descente, Lemme \ref{descent-lemma-descent-data-sheaves}).
Puisque $X' \to Z'$ est quasi-affine, cette donnée de descente est effective ; voir
Descente, Lemme \ref{descent-lemma-quasi-affine}.
Il existe donc un schéma $Y \to Z$ sur $Z$, ainsi qu'un
isomorphisme $\psi : Z' \times_Z Y \to X'$ compatible avec les données de descente.
Bien entendu, $Y \to Z$ est quasi-affine (par construction ou d'après
Descente, Lemme \ref{descent-lemma-descending-property-quasi-affine}).
Remarquons que $\mathcal{Y} = \{Z' \times_Z Y \to Y\}$ est un
recouvrement fpqc. En interprétant $\psi$ comme un élément de
$X(Z' \times_Z Y)$, on voit que $\psi \in \check{H}^0(\mathcal{Y}, X)$.
La condition de faisceau pour $X$ (voir
Propriétés des espaces, Proposition
\ref{spaces-properties-proposition-sheaf-fpqc})
fournit un morphisme $Y \to X$.
Par construction, le changement de base de ce morphisme par $Z'$ est un isomorphisme ; il est donc
un isomorphisme d'après le
lemme \ref{spaces-descent-lemma-descending-property-isomorphism}.
Cela prouve que $X$ est représentable par un schéma quasi-affine, ce qui achève la démonstration.
\end{proof}

\begin{lemma}
\label{spaces-descent-lemma-descending-property-quasi-compact-immersion}
La propriété $\mathcal{P}(f) =$``$f$ est une immersion quasi-compacte''
est locale sur la base pour la topologie fpqc.
\end{lemma}

\begin{proof}
Nous appliquerons le
lemme \ref{spaces-descent-lemma-descending-properties-morphisms}.
Ses hypothèses (1) et (2) résultent de
Morphismes des espaces,
Lemmes \ref{spaces-morphisms-lemma-closed-immersion-local} et
\ref{spaces-morphisms-lemma-quasi-compact-local}.
Considérons un diagramme cartésien
$$
\xymatrix{
X' \ar[r] \ar[d] & X \ar[d] \\
Z' \ar[r] & Z
}
$$
d'espaces algébriques sur $S$
où $Z' \to Z$ est un morphisme plat et surjectif de schémas affines
et $X' \to Z'$ une immersion quasi-compacte. Il faut montrer que $X \to Z$
est une immersion quasi-compacte. Le morphisme $X' \to Z'$ est quasi-affine. D'après le
lemme \ref{spaces-descent-lemma-descending-property-quasi-affine},
on voit que $X$ est un schéma et que $X \to Z$ est quasi-affine.
Il résulte de
Descente, Lemme \ref{descent-lemma-descending-property-quasi-compact-immersion},
que $X \to Z$ est une immersion quasi-compacte, ce qu'il fallait démontrer.
\end{proof}

\begin{lemma}
\label{spaces-descent-lemma-descending-property-integral}
La propriété $\mathcal{P}(f) =$``$f$ est entier''
est locale sur la base pour la topologie fpqc.
\end{lemma}

\begin{proof}
Un morphisme est entier si et seulement s'il est affine et
universellement fermé. Voir
Morphismes des espaces,
Lemme \ref{spaces-morphisms-lemma-integral-universally-closed}.
Le lemme résulte donc de la combinaison des
lemmes \ref{spaces-descent-lemma-descending-property-universally-closed}
et \ref{spaces-descent-lemma-descending-property-affine}.
\end{proof}

\begin{lemma}
\label{spaces-descent-lemma-descending-property-finite}
La propriété $\mathcal{P}(f) =$``$f$ est fini''
est locale sur la base pour la topologie fpqc.
\end{lemma}

\begin{proof}
Un morphisme est fini si et seulement s'il est entier
et localement de type fini. Voir
Morphismes des espaces, Lemme \ref{spaces-morphisms-lemma-finite-integral}.
Le lemme résulte donc de la combinaison des
lemmes \ref{spaces-descent-lemma-descending-property-locally-finite-type}
et \ref{spaces-descent-lemma-descending-property-integral}.
\end{proof}

\begin{lemma}
\label{spaces-descent-lemma-descending-property-quasi-finite}
Les propriétés
$\mathcal{P}(f) =$``$f$ est localement quasi-fini''
et
$\mathcal{P}(f) =$``$f$ est quasi-fini''
sont locales sur la base pour la topologie fpqc.
\end{lemma}

\begin{proof}
Nous avons déjà vu que la propriété ``quasi-compact'' est locale sur la base pour la topologie fpqc ; voir
Lemme \ref{spaces-descent-lemma-descending-property-quasi-compact}. Il suffit donc
de démontrer le lemme pour ``localement quasi-fini''. Nous appliquerons le
lemme \ref{spaces-descent-lemma-descending-properties-morphisms}.
Ses hypothèses (1) et (2) résultent de
Morphismes des espaces,
Lemme \ref{spaces-morphisms-lemma-quasi-finite-local}.
Soit $Z' \to Z$ un morphisme plat et surjectif de schémas affines sur $S$.
Soit $f : X \to Z$ un morphisme d'espaces algébriques, et supposons
que le changement de base $f' : Z' \times_Z X \to Z'$ soit localement quasi-fini.
Il faut montrer que $f$ est localement quasi-fini. Soit $U$ un schéma
et soit $U \to X$ un morphisme étale surjectif. D'après
Morphismes des espaces,
Lemme \ref{spaces-morphisms-lemma-quasi-finite-local},
il suffit encore de montrer que $U \to Z$ est localement quasi-fini.
Puisque $f'$ est localement quasi-fini et que $Z' \times_Z U$ est un
schéma étale sur $Z' \times_Z X$, on déduit encore du même lemme que
$Z' \times_Z U \to Z'$ est localement quasi-fini.
Comme $\{Z' \to Z\}$ est un recouvrement fpqc, on en conclut que
$U \to Z$ est localement quasi-fini d'après
Descente, Lemme \ref{descent-lemma-descending-property-quasi-finite},
ce qu'il fallait démontrer.
\end{proof}

\begin{lemma}
\label{spaces-descent-lemma-descending-property-syntomic}
La propriété $\mathcal{P}(f) =$``$f$ est syntomique''
est locale sur la base pour la topologie fpqc.
\end{lemma}

\begin{proof}
Nous appliquerons le
lemme \ref{spaces-descent-lemma-descending-properties-morphisms}.
Ses hypothèses (1) et (2) résultent de
Morphismes des espaces,
Lemme \ref{spaces-morphisms-lemma-syntomic-local}.
Soit $Z' \to Z$ un morphisme plat et surjectif de schémas affines sur $S$.
Soit $f : X \to Z$ un morphisme d'espaces algébriques, et supposons
que le changement de base $f' : Z' \times_Z X \to Z'$ soit syntomique.
Il faut montrer que $f$ est syntomique. Soit $U$ un schéma
et soit $U \to X$ un morphisme étale surjectif. D'après
Morphismes des espaces,
Lemme \ref{spaces-morphisms-lemma-syntomic-local},
il suffit encore de montrer que $U \to Z$ est syntomique.
Puisque $f'$ est syntomique et que $Z' \times_Z U$ est un
schéma étale sur $Z' \times_Z X$, on déduit encore du même lemme que
$Z' \times_Z U \to Z'$ est syntomique.
Comme $\{Z' \to Z\}$ est un recouvrement fpqc, on en conclut que
$U \to Z$ est syntomique d'après
Descente, Lemme \ref{descent-lemma-descending-property-syntomic},
ce qu'il fallait démontrer.
\end{proof}

\begin{lemma}
\label{spaces-descent-lemma-descending-property-smooth}
La propriété $\mathcal{P}(f) =$``$f$ est lisse''
est locale sur la base pour la topologie fpqc.
\end{lemma}

\begin{proof}
Nous appliquerons le
lemme \ref{spaces-descent-lemma-descending-properties-morphisms}.
Ses hypothèses (1) et (2) résultent de
Morphismes des espaces,
Lemme \ref{spaces-morphisms-lemma-smooth-local}.
Soit $Z' \to Z$ un morphisme plat et surjectif de schémas affines sur $S$.
Soit $f : X \to Z$ un morphisme d'espaces algébriques, et supposons
que le changement de base $f' : Z' \times_Z X \to Z'$ soit lisse.
Il faut montrer que $f$ est lisse. Soit $U$ un schéma
et soit $U \to X$ un morphisme étale surjectif. D'après
Morphismes des espaces,
Lemme \ref{spaces-morphisms-lemma-smooth-local},
il suffit encore de montrer que $U \to Z$ est lisse.
Puisque $f'$ est lisse et que $Z' \times_Z U$ est un
schéma étale sur $Z' \times_Z X$, on déduit encore du même lemme que
$Z' \times_Z U \to Z'$ est lisse.
Comme $\{Z' \to Z\}$ est un recouvrement fpqc, on en conclut que
$U \to Z$ est lisse d'après
Descente, Lemme \ref{descent-lemma-descending-property-smooth},
ce qu'il fallait démontrer.
\end{proof}

\begin{lemma}
\label{spaces-descent-lemma-descending-property-unramified}
La propriété $\mathcal{P}(f) =$``$f$ est non ramifié''
est locale sur la base pour la topologie fpqc.
\end{lemma}

\begin{proof}
Nous appliquerons le
lemme \ref{spaces-descent-lemma-descending-properties-morphisms}.
Ses hypothèses (1) et (2) résultent de
Morphismes des espaces,
Lemme \ref{spaces-morphisms-lemma-unramified-local}.
Soit $Z' \to Z$ un morphisme plat et surjectif de schémas affines sur $S$.
Soit $f : X \to Z$ un morphisme d'espaces algébriques, et supposons
que le changement de base $f' : Z' \times_Z X \to Z'$ soit non ramifié.
Il faut montrer que $f$ est non ramifié. Soit $U$ un schéma
et soit $U \to X$ un morphisme étale surjectif. D'après
Morphismes des espaces,
Lemme \ref{spaces-morphisms-lemma-unramified-local},
il suffit encore de montrer que $U \to Z$ est non ramifié.
Puisque $f'$ est non ramifié et que $Z' \times_Z U$ est un
schéma étale sur $Z' \times_Z X$, on déduit encore du même lemme que
$Z' \times_Z U \to Z'$ est non ramifié.
Comme $\{Z' \to Z\}$ est un recouvrement fpqc, on en conclut que
$U \to Z$ est non ramifié d'après
Descente, Lemme \ref{descent-lemma-descending-property-unramified},
ce qu'il fallait démontrer.
\end{proof}

\begin{lemma}
\label{spaces-descent-lemma-descending-property-etale}
La propriété $\mathcal{P}(f) =$``$f$ est étale''
est locale sur la base pour la topologie fpqc.
\end{lemma}

\begin{proof}
Nous appliquerons le
lemme \ref{spaces-descent-lemma-descending-properties-morphisms}.
Ses hypothèses (1) et (2) résultent de
Morphismes des espaces,
Lemme \ref{spaces-morphisms-lemma-etale-local}.
Soit $Z' \to Z$ un morphisme plat et surjectif de schémas affines sur $S$.
Soit $f : X \to Z$ un morphisme d'espaces algébriques, et supposons
que le changement de base $f' : Z' \times_Z X \to Z'$ soit étale.
Il faut montrer que $f$ est étale. Soit $U$ un schéma
et soit $U \to X$ un morphisme étale surjectif. D'après
Morphismes des espaces,
Lemme \ref{spaces-morphisms-lemma-etale-local},
il suffit encore de montrer que $U \to Z$ est étale.
Puisque $f'$ est étale et que $Z' \times_Z U$ est un
schéma étale sur $Z' \times_Z X$, on déduit encore du même lemme que
$Z' \times_Z U \to Z'$ est étale.
Comme $\{Z' \to Z\}$ est un recouvrement fpqc, on en conclut que
$U \to Z$ est étale d'après
Descente, Lemme \ref{descent-lemma-descending-property-etale},
ce qu'il fallait démontrer.
\end{proof}

\begin{lemma}
\label{spaces-descent-lemma-descending-property-finite-locally-free}
La propriété $\mathcal{P}(f) =$``$f$ est fini localement libre''
est locale sur la base pour la topologie fpqc.
\end{lemma}

\begin{proof}
Un morphisme est fini localement libre si et seulement s'il est
fini, plat et localement de présentation finie
(Morphismes des espaces, Lemme \ref{spaces-morphisms-lemma-finite-flat}).
Le résultat découle donc des lemmes
\ref{spaces-descent-lemma-descending-property-finite},
\ref{spaces-descent-lemma-descending-property-flat} et
\ref{spaces-descent-lemma-descending-property-locally-finite-presentation}.
\end{proof}

\begin{lemma}
\label{spaces-descent-lemma-descending-property-monomorphism}
La propriété $\mathcal{P}(f) =$``$f$ est un monomorphisme''
est locale sur la base pour la topologie fpqc.
\end{lemma}

\begin{proof}
Soit $f : X \to Y$ un morphisme d'espaces algébriques.
Soit $\{Y_i \to Y\}$ un recouvrement fpqc, et supposons
que chacun des changements de base $f_i : X_i \to Y_i$ de $f$ soit
un monomorphisme. Il faut montrer que $f$ est un monomorphisme.

\medskip\noindent
Première démonstration. Le morphisme $f$ est un monomorphisme si et seulement si
$\Delta : X \to X \times_Y X$ est un isomorphisme. En appliquant ce critère aux
$f_i$, on voit que chacun des morphismes
$$
\Delta_i :
X_i
\longrightarrow
X_i \times_{Y_i} X_i = Y_i \times_Y (X \times_Y X)
$$
est un isomorphisme. Le changement de base d'un recouvrement fpqc est un recouvrement fpqc ; voir
Topologies sur les espaces, Lemme \ref{spaces-topologies-lemma-fpqc} ;
ainsi, $\{Y_i \times_Y (X \times_Y X) \to X \times_Y X\}$
est un recouvrement fpqc d'espaces algébriques. De plus, chaque
$\Delta_i$ est le changement de base du morphisme
$\Delta : X \to X \times_Y X$. Il résulte donc du
lemme \ref{spaces-descent-lemma-descending-property-isomorphism}
que $\Delta$ est un isomorphisme, c'est-à-dire que $f$ est un monomorphisme.

\medskip\noindent
Seconde démonstration.
Soit $V$ un schéma et soit $V \to Y$ un morphisme étale surjectif.
Si l'on peut montrer que $V \times_Y X \to V$ est un monomorphisme, alors
$X \to Y$ est un monomorphisme. En effet, dans tout
diagramme cartésien de faisceaux
$$
\vcenter{
\xymatrix{
\mathcal{F} \ar[r]_a \ar[d]_b & \mathcal{G} \ar[d]^c \\
\mathcal{H} \ar[r]^d & \mathcal{I}
}
}
\quad
\quad
\mathcal{F} = \mathcal{H} \times_\mathcal{I} \mathcal{G}
$$
si $c$ est un morphisme surjectif de faisceaux et si $a$ est injectif, alors
$d$ est lui aussi injectif. Cela ramène le problème au cas où $Y$ est
un schéma. En outre, dans ce cas, on peut supposer que les espaces algébriques
$Y_i$ sont des schémas, puisqu'on peut toujours raffiner le recouvrement afin de se
ramener à cette situation ; voir
Topologies sur les espaces, Lemme \ref{spaces-topologies-lemma-refine-fpqc-schemes}.

\medskip\noindent
Supposons que $\{Y_i \to Y\}$ soit un recouvrement fpqc de schémas.
Soient $a, b : T \to X$ deux morphismes
tels que $f \circ a = f \circ b$. Il faut montrer que $a = b$.
Puisque $f_i$ est un monomorphisme, on a $a_i = b_i$, où
$a_i, b_i : Y_i \times_Y T \to X_i$ sont
les changements de base. En particulier, les composés
$Y_i \times_Y T \to T \to X$ sont égaux.
Comme $\{Y_i \times_Y T \to T\}$ est un recouvrement fpqc,
on déduit que $a = b$ d'après Propriétés des espaces, Proposition
\ref{spaces-properties-proposition-sheaf-fpqc}.
\end{proof}



% preserved blank authority line 2020
\section{Descente des propriétés des morphismes pour la topologie fppf}
\label{spaces-descent-section-descending-properties-morphisms-fppf}

\noindent
Nous établissons dans cette section certaines propriétés des morphismes d'espaces
algébriques dont nous n'avons pas (encore) pu montrer qu'elles sont locales sur
la base pour la topologie fpqc, mais qui sont locales sur la base
pour la topologie fppf.

\begin{lemma}
\label{spaces-descent-lemma-descending-fppf-property-immersion}
La propriété $\mathcal{P}(f) =$``$f$ est une immersion''
est locale sur la base pour la topologie fppf.
\end{lemma}

\begin{proof}
Soit $f : X \to Y$ un morphisme d'espaces algébriques.
Soit $\{Y_i \to Y\}_{i \in I}$ un recouvrement fppf de $Y$.
Soit $f_i : X_i \to Y_i$ le changement de base de $f$.

\medskip\noindent
Si $f$ est une immersion, alors chaque $f_i$ est une immersion d'après
Espaces, lemme \ref{spaces-lemma-base-change-immersions}.
Cela démontre l'implication directe de la
définition \ref{spaces-descent-definition-property-morphisms-local}.

\medskip\noindent
Réciproquement, supposons que chaque $f_i$ soit une immersion. D'après
Morphismes d'espaces,
lemme \ref{spaces-morphisms-lemma-immersions-monomorphisms},
chaque $f_i$ est alors séparé. D'après
Morphismes d'espaces,
lemme \ref{spaces-morphisms-lemma-immersion-quasi-finite},
chaque $f_i$ est alors localement quasi-fini.
Les lemmes \ref{spaces-descent-lemma-descending-property-separated}
et \ref{spaces-descent-lemma-descending-property-quasi-finite} montrent donc que
$f$ est localement quasi-fini et séparé.
D'après
Morphismes d'espaces, lemme
\ref{spaces-morphisms-lemma-locally-quasi-finite-separated-representable},
il s'ensuit que $f$ est représentable !

\medskip\noindent
D'après
Morphismes d'espaces, lemme \ref{spaces-morphisms-lemma-closed-immersion-local},
il suffit de montrer que, pour tout schéma $Z$ et tout morphisme $Z \to Y$,
le changement de base $Z \times_Y X \to Z$ est une immersion. D'après
Topologies sur les espaces, lemme \ref{spaces-topologies-lemma-refine-fppf-schemes},
on peut trouver un recouvrement fppf $\{Z_i \to Z\}$ par des schémas qui raffine
le changement de base du recouvrement $\{Y_i \to Y\}$ à $Z$.
Par conséquent, $Z \times_Y X \to Z$ (qui est un morphisme de schémas
d'après le résultat du paragraphe précédent) devient une immersion
après changement de base par chacun des morphismes d'un recouvrement fppf de $Z$ formé de schémas.
Ainsi, $Z \times_Y X \to Z$ est une immersion d'après le résultat pour les schémas ; voir
Descente, lemme \ref{descent-lemma-descending-fppf-property-immersion}.
\end{proof}

\begin{lemma}
\label{spaces-descent-lemma-descending-fppf-property-locally-separated}
La propriété $\mathcal{P}(f) =$``$f$ est localement séparé''
est locale sur la base pour la topologie fppf.
\end{lemma}

\begin{proof}
Tout changement de base d'un morphisme localement séparé est localement séparé ; voir
Morphismes d'espaces,
lemme \ref{spaces-morphisms-lemma-base-change-separated}.
Cela établit l'implication directe de la
définition \ref{spaces-descent-definition-property-morphisms-local}.

\medskip\noindent
Soit $\{Y_i \to Y\}_{i \in I}$ un recouvrement fppf d'espaces algébriques sur $S$.
Soit $f : X \to Y$ un morphisme d'espaces algébriques sur $S$.
Supposons que chaque changement de base $X_i := Y_i \times_Y X \to Y_i$ soit localement séparé.
Cela signifie que chacun des morphismes
$$
\Delta_i :
X_i
\longrightarrow
X_i \times_{Y_i} X_i = Y_i \times_Y (X \times_Y X)
$$
est une immersion. Tout changement de base d'un recouvrement fppf est un
recouvrement fppf ; voir
Topologies sur les espaces, lemme \ref{spaces-topologies-lemma-fppf}.
Par conséquent, $\{Y_i \times_Y (X \times_Y X) \to X \times_Y X\}$
est un recouvrement fppf d'espaces algébriques. De plus, chaque
$\Delta_i$ est le changement de base du morphisme
$\Delta : X \to X \times_Y X$. Il résulte donc du
lemme \ref{spaces-descent-lemma-descending-fppf-property-immersion}
que $\Delta$ est une immersion, c'est-à-dire que $f$ est localement séparé.
\end{proof}





\section{Application de la descente des propriétés des morphismes}
\label{spaces-descent-section-application-descending-properties-morphisms}

\noindent
Cette section est l'analogue de Descente, section
\ref{descent-section-application-descending-properties-morphisms}.

\begin{lemma}
\label{spaces-descent-lemma-descending-property-ample}
Soit $S$ un schéma.
Soit $f : X \to Y$ un morphisme d'espaces algébriques sur $S$.
Soit $\mathcal{L}$ un $\mathcal{O}_X$-module inversible.
Soit $\{g_i : Y_i \to Y\}_{i \in I}$ un recouvrement fpqc.
Soit $f_i : X_i \to Y_i$ le changement de base de $f$, et soit $\mathcal{L}_i$
l'image inverse de $\mathcal{L}$ sur $X_i$.
Les conditions suivantes sont équivalentes :
\begin{enumerate}
\item $\mathcal{L}$ est ample sur $X/Y$ ; et
\item $\mathcal{L}_i$ est ample sur $X_i/Y_i$
pour tout $i \in I$.
\end{enumerate}
\end{lemma}

\begin{proof}
L'implication (1) $\Rightarrow$ (2) résulte de
Diviseurs sur les espaces, lemme \ref{spaces-divisors-lemma-ample-base-change}.
Supposons (2). Pour vérifier que $\mathcal{L}$ est ample sur $X/Y$, on peut
travailler localement pour la topologie \'etale sur $Y$ ; voir
Diviseurs sur les espaces, lemme \ref{spaces-divisors-lemma-relatively-ample-local}.
On peut donc supposer que $Y$ est un schéma, puis que
chaque $Y_i$ est également un schéma ; voir
Topologies sur les espaces, lemme \ref{spaces-topologies-lemma-refine-fpqc-schemes}.
Autrement dit, on peut supposer que
$\{Y_i \to Y\}$ est un recouvrement fpqc de schémas.

\medskip\noindent
D'après Diviseurs sur les espaces, lemme
\ref{spaces-divisors-lemma-relatively-ample-properties},
$X_i \to Y_i$ est représentable (c'est-à-dire que $X_i$ est un schéma),
quasi-compact et séparé. Les lemmes
\ref{spaces-descent-lemma-descending-property-quasi-compact} et
\ref{spaces-descent-lemma-descending-property-separated} montrent donc que $f$ est quasi-compact et séparé.
Cela signifie que
$\mathcal{A} = \bigoplus_{d \geq 0} f_*\mathcal{L}^{\otimes d}$
est une algèbre graduée quasi-cohérente sur $\mathcal{O}_Y$
(Morphismes d'espaces, lemme \ref{spaces-morphisms-lemma-pushforward}).
De plus, la formation de $\mathcal{A}$ commute aux
changements de base plats d'après
Cohomologie des espaces, lemme
\ref{spaces-cohomology-lemma-flat-base-change-cohomology}.
En particulier, si l'on pose
$\mathcal{A}_i = \bigoplus_{d \geq 0} f_{i, *}\mathcal{L}_i^{\otimes d}$,
alors $\mathcal{A}_i = g_i^*\mathcal{A}$.
Il s'ensuit que les morphismes naturels
$\psi_d : f^*\mathcal{A}_d \to \mathcal{L}^{\otimes d}$
de $\mathcal{O}_X$-modules
deviennent, après image inverse, les morphismes naturels
$\psi_{i, d} : f_i^*(\mathcal{A}_i)_d \to \mathcal{L}_i^{\otimes d}$
de $\mathcal{O}_{X_i}$-modules. Puisque $\mathcal{L}_i$ est ample sur $X_i/Y_i$,
pour tout point $x_i \in X_i$, il existe un $d \geq 1$
tel que $f_i^*(\mathcal{A}_i)_d \to \mathcal{L}_i^{\otimes d}$
soit surjectif sur la fibre en $x_i$. Cela résulte soit directement
de la définition d'un module relativement ample, soit de
Morphismes, lemme \ref{morphisms-lemma-characterize-relatively-ample}.
Si $x \in |X|$, on peut choisir un $i$ et un $x_i \in X_i$
dont l'image est $x$. Comme
$\mathcal{O}_{X, \overline{x}} \to \mathcal{O}_{X_i, \overline{x}_i}$
est plat, donc fidèlement plat, on conclut que, pour tout $x \in |X|$,
il existe un $d \geq 1$ tel que
$f^*\mathcal{A}_d \to \mathcal{L}^{\otimes d}$
soit surjectif sur la fibre en $x$.
Il s'ensuit que l'ouvert $U(\psi) \subset X$ du
lemme de Diviseurs sur les espaces
\ref{spaces-divisors-lemma-invertible-map-into-relative-proj}
correspondant au morphisme
$\psi : f^*\mathcal{A} \to \bigoplus_{d \geq 0} \mathcal{L}^{\otimes d}$
d'algèbres graduées sur $\mathcal{O}_X$
est égal à $X$. Considérons le morphisme correspondant
$$
r_{\mathcal{L}, \psi} : X \longrightarrow \underline{\text{Proj}}_Y(\mathcal{A})
$$
Il ressort clairement de ce qui précède que le changement de base de
$r_{\mathcal{L}, \psi}$ à $Y_i$ est le morphisme
$r_{\mathcal{L}_i, \psi_i}$, qui est une immersion ouverte d'après
Morphismes, lemme \ref{morphisms-lemma-characterize-relatively-ample}.
Ainsi, $r_{\mathcal{L}, \psi}$ est une immersion ouverte
d'après le lemme \ref{spaces-descent-lemma-descending-property-open-immersion}.
Par conséquent, $X$ est un schéma,
et l'on conclut que $\mathcal{L}$ est ample sur $X/Y$ d'après
Morphismes, lemme \ref{morphisms-lemma-characterize-relatively-ample}.
\end{proof}

\begin{lemma}
\label{spaces-descent-lemma-ample-in-neighbourhood}
Soit $S$ un schéma.
Soit $f : X \to Y$ un morphisme propre d'espaces algébriques sur $S$.
Soit $\mathcal{L}$ un $\mathcal{O}_X$-module inversible.
Il existe un sous-espace ouvert $V \subset Y$ caractérisé par
la propriété suivante :
Un morphisme $Y' \to Y$ d'espaces algébriques se factorise
par $V$ si et seulement si l'image inverse $\mathcal{L}'$
de $\mathcal{L}$ sur $X' = Y' \times_Y X$ est ample sur $X'/Y'$
(au sens de Diviseurs sur les espaces, définition
\ref{spaces-divisors-definition-relatively-ample}).
\end{lemma}

\begin{proof}
Supposons le lemme établi lorsque $Y$ est un schéma.
Soit $U$ un schéma et soit $U \to Y$ un morphisme \'etale surjectif.
Soit $R = U \times_Y U$, avec les projections $t, s : R \to U$.
Notons $X_U = U \times_Y X$ et désignons par $\mathcal{L}_U$ l'image inverse correspondante.
On obtient alors un sous-schéma ouvert $V' \subset U$ comme dans le lemme
pour $(X_U \to U, \mathcal{L}_U)$. La caractérisation fonctorielle donne
$s^{-1}(V') = t^{-1}(V')$.
Il existe donc un sous-espace ouvert $V \subset Y$ tel que
$V'$ soit l'image inverse de $V$ dans $U$.
En particulier, $V' \to V$ est \'etale et surjectif, et l'on
conclut que $\mathcal{L}_V$ est ample sur $X_V/V$
(Diviseurs sur les espaces, lemme \ref{spaces-divisors-lemma-relatively-ample-local}).
Maintenant, si $Y' \to Y$ est un morphisme tel que
$\mathcal{L}'$ soit ample sur $X'/Y'$, alors
$U \times_Y Y' \to Y'$ doit se factoriser par $V'$,
et l'on conclut que $Y' \to Y$ se factorise par $V$.
Ainsi, $V \subset Y$ possède la propriété énoncée dans le lemme.
On est ainsi ramené au cas traité au
paragraphe suivant.

\medskip\noindent
Supposons que $Y$ soit un schéma. La question étant locale sur $Y$,
on peut supposer que $Y$ est un schéma affine. Montrons
l'assertion suivante :
\begin{enumerate}
\item[(A)] Si $\Spec(k) \to Y$ est un morphisme tel que
$\mathcal{L}_k$ soit ample sur $X_k/k$, alors il existe un
voisinage ouvert $V \subset Y$ de l'image de $\Spec(k) \to Y$
tel que $\mathcal{L}_V$ soit ample sur $X_V/V$.
\end{enumerate}
Il est clair que (A) entraîne le lemme.

\medskip\noindent
Soient $X \to Y$, $\mathcal{L}$ et $\Spec(k) \to Y$ comme dans (A).
D'après le lemme \ref{spaces-descent-lemma-descending-property-ample},
on peut supposer que $k = \kappa(y)$ est le corps résiduel d'un point $y$
de $Y$.

\medskip\noindent
Comme $Y$ est affine, on peut trouver un ensemble préordonné filtrant $I$ et un
système projectif de morphismes $X_i \to Y_i$ d'espaces algébriques,
où $Y_i$ est de présentation finie sur $\mathbf{Z}$, où les
morphismes de transition $X_i \to X_{i'}$ et $Y_i \to Y_{i'}$ sont affines,
où $X_i \to Y_i$ est propre et de présentation finie, et tel que
$X \to Y = \lim (X_i \to Y_i)$. Voir Limites d'espaces, lemme
\ref{spaces-limits-lemma-proper-limit-of-proper-finite-presentation-noetherian}.
Quitte à remplacer $I$ par une partie cofinale, on peut supposer que $Y_i$ est un schéma (affine) pour tout $i$ ;
voir Limites d'espaces, lemme \ref{spaces-limits-lemma-limit-is-affine}.
Quitte encore à remplacer $I$ par une partie cofinale, on peut supposer disposer d'un système compatible de
$\mathcal{O}_{X_i}$-modules inversibles $\mathcal{L}_i$
dont l'image inverse est $\mathcal{L}$ ; voir
Limites d'espaces, lemme \ref{spaces-limits-lemma-descend-invertible-modules}.
Soit $y_i \in Y_i$ l'image de $y$.
Alors $\kappa(y) = \colim \kappa(y_i)$.
Par conséquent, $X_y = \lim X_{i, y_i}$ et, quitte à remplacer $I$ par une partie cofinale,
on peut supposer que $X_{i, y_i}$ est un schéma pour tout $i$ ; voir
Limites d'espaces, lemme \ref{spaces-limits-lemma-limit-is-scheme}.
Il existe donc un $i$ tel que $\mathcal{L}_{i, y_i}$
soit ample sur $X_{i, y_i}$, d'après
Limites, lemme \ref{limits-lemma-limit-ample}.
D'après Diviseurs sur les espaces, lemme \ref{spaces-divisors-lemma-ample-in-neighbourhood},
on trouve un voisinage ouvert
$V_i \subset Y_i$ de $y_i$ tel que
la restriction de $\mathcal{L}_i$ à $f_i^{-1}(V_i)$
soit ample relativement à $V_i$.
On achève la démonstration en prenant pour $V \subset Y$ l'image inverse de
$V_i$ (indications : utiliser
Morphismes, lemme \ref{morphisms-lemma-ample-base-change},
le fait que $X \to Y \times_{Y_i} X_i$ est affine,
et le fait que l'image inverse d'un
faisceau inversible ample par un morphisme affine est ample, d'après
Morphismes, lemme \ref{morphisms-lemma-pullback-ample-tensor-relatively-ample}).
\end{proof}













\section{Propriétés des morphismes locales sur la source}
\label{spaces-descent-section-properties-morphisms-local-source}

\noindent
Dans cette section, nous définissons ce que signifie, pour une propriété des morphismes
d'espaces algébriques, être locale sur la source. Comparer avec
Descente, section \ref{descent-section-properties-morphisms-local-source}.

\begin{definition}
\label{spaces-descent-definition-property-morphisms-local-source}
Soit $S$ un schéma.
Soit $\mathcal{P}$ une propriété des morphismes d'espaces algébriques sur $S$.
Soit $\tau \in \{fpqc, \linebreak[0] fppf, \linebreak[0] syntomic, \linebreak[0]
smooth, \linebreak[0] \etale\}$. Nous disons que $\mathcal{P}$ est
{\it $\tau$-locale sur la source}, ou
{\it locale sur la source pour la topologie $\tau$}, si, pour
tout morphisme $f : X \to Y$ d'espaces algébriques sur $S$ et tout
$\tau$-recouvrement $\{X_i \to X\}_{i \in I}$ d'espaces algébriques, on a
$$
f \text{ vérifie }\mathcal{P}
\Leftrightarrow
\text{chaque }X_i \to Y\text{ vérifie }\mathcal{P}.
$$
\end{definition}

\noindent
Précisons que, puisque les isomorphismes sont toujours couvrants,
on voit (ou l'on exige) que la propriété $\mathcal{P}$ est vérifiée par $X \to Y$
si et seulement si elle est vérifiée par toute flèche $X' \to Y'$ isomorphe à $X \to Y$.
Si une propriété est $\tau$-locale sur la source, elle reste vérifiée après
précomposition par tout morphisme figurant dans un $\tau$-recouvrement. Voici
un énoncé formel.

\begin{lemma}
\label{spaces-descent-lemma-precompose-property-local-source}
Soit $S$ un schéma.
Soit $\tau \in \{fpqc, \linebreak[0] fppf, \linebreak[0] syntomic, \linebreak[0]
smooth, \linebreak[0] \etale\}$.
Soit $\mathcal{P}$ une propriété des morphismes d'espaces algébriques sur $S$
qui est $\tau$-locale sur la source. Supposons que $f : X \to Y$ vérifie la propriété
$\mathcal{P}$. Pour tout morphisme $a : X' \to X$ qui est respectivement
plat, plat et localement de présentation finie, syntomique,
lisse ou \'etale, la composée $f \circ a : X' \to Y$ vérifie
la propriété $\mathcal{P}$.
\end{lemma}

\begin{proof}
Cela résulte du fait que l'on peut compléter $X' \to X$ en une famille de
morphismes formant un recouvrement pour la topologie $\tau$.
\end{proof}

\begin{lemma}
\label{spaces-descent-lemma-transfer-from-schemes}
Soit $S$ un schéma.
Soit $\tau \in \{fpqc, \linebreak[0] fppf, \linebreak[0] syntomic, \linebreak[0]
smooth, \linebreak[0] \etale\}$.
Supposons que $\mathcal{P}$ soit une propriété des morphismes de schémas sur $S$
qui soit locale sur la source et le but pour la topologie \'etale. Notons $\mathcal{P}_{spaces}$
la propriété correspondante des morphismes d'espaces algébriques sur $S$ ; voir
Morphismes d'espaces, définition \ref{spaces-morphisms-definition-P}.
Si $\mathcal{P}$ est locale sur la source pour la topologie $\tau$, alors
$\mathcal{P}_{spaces}$ est locale sur la source pour la topologie $\tau$.
\end{lemma}

\begin{proof}
Soit $f : X \to Y$ un morphisme d'espaces algébriques sur $S$.
Soit $\{X_i \to X\}_{i \in I}$ un recouvrement d'espaces algébriques pour la topologie $\tau$.
Choisissons un schéma $V$ et un morphisme \'etale surjectif $V \to Y$.
Choisissons un schéma $U$ et un morphisme \'etale surjectif $U \to X \times_Y V$.
Pour chaque $i$, choisissons un schéma $U_i$ et un morphisme \'etale surjectif
$U_i \to X_i \times_X U$.

\medskip\noindent
Remarquons que $\{X_i \times_X U \to U\}_{i \in I}$ est un recouvrement pour la topologie $\tau$.
Remarquons que chacune des familles $\{U_i \to X_i \times_X U\}$ est un recouvrement \'etale,
donc un recouvrement pour la topologie $\tau$. Ainsi, $\{U_i \to U\}_{i \in I}$ est un
recouvrement pour la topologie $\tau$ dans la catégorie des espaces algébriques sur $S$. Mais puisque $U$ et chaque $U_i$
sont des schémas, $\{U_i \to U\}_{i \in I}$ est un
recouvrement pour la topologie $\tau$ dans la catégorie des schémas sur $S$.

\medskip\noindent
On a maintenant
\begin{align*}
f \text{ vérifie }\mathcal{P}_{spaces}
& \Leftrightarrow
U \to V \text{ vérifie }\mathcal{P} \\
& \Leftrightarrow
\text{chaque }U_i \to V \text{ vérifie }\mathcal{P} \\
& \Leftrightarrow
\text{chaque }X_i \to Y\text{ vérifie }\mathcal{P}_{spaces}.
\end{align*}
La première et la dernière équivalence résultent de la définition de
$\mathcal{P}_{spaces}$, tandis que l'équivalence intermédiaire résulte de l'hypothèse
selon laquelle $\mathcal{P}$ est locale sur la source pour la topologie $\tau$.
\end{proof}
% preserved blank authority line 2405
% preserved blank authority line 2406
% preserved blank authority line 2407
% preserved blank authority line 2408
% preserved blank authority line 2409
% preserved blank authority line 2410
\section{Propriétés des morphismes locales sur la source pour la topologie fpqc}
\label{spaces-descent-section-fpqc-local-source}

\noindent
Voici quelques propriétés des morphismes qui sont locales sur la source pour la topologie fpqc.

\begin{lemma}
\label{spaces-descent-lemma-flat-fpqc-local-source}
La propriété $\mathcal{P}(f)=$``$f$ est plat'' est locale sur la source pour la topologie fpqc.
\end{lemma}

\begin{proof}
L'assertion résulte du
lemme \ref{spaces-descent-lemma-transfer-from-schemes}
en utilisant
Morphismes d'espaces, définition \ref{spaces-morphisms-definition-flat}
et
Descente, lemme \ref{descent-lemma-flat-fpqc-local-source}.
\end{proof}










\section{Propriétés des morphismes locales sur la source pour la topologie fppf}
\label{spaces-descent-section-fppf-local-source}

\noindent
Voici quelques propriétés des morphismes qui sont locales sur la source pour la topologie fppf.

\begin{lemma}
\label{spaces-descent-lemma-locally-finite-presentation-fppf-local-source}
La propriété $\mathcal{P}(f)=$``$f$ est localement de présentation finie''
est locale sur la source pour la topologie fppf.
\end{lemma}

\begin{proof}
L'assertion résulte du
lemme \ref{spaces-descent-lemma-transfer-from-schemes}
en utilisant
Morphismes d'espaces,
définition \ref{spaces-morphisms-definition-locally-finite-presentation}
et
Descente,
lemme \ref{descent-lemma-locally-finite-presentation-fppf-local-source}.
\end{proof}

\begin{lemma}
\label{spaces-descent-lemma-locally-finite-type-fppf-local-source}
La propriété $\mathcal{P}(f)=$``$f$ est localement de type fini''
est locale sur la source pour la topologie fppf.
\end{lemma}

\begin{proof}
L'assertion résulte du
lemme \ref{spaces-descent-lemma-transfer-from-schemes}
en utilisant
Morphismes d'espaces,
définition \ref{spaces-morphisms-definition-locally-finite-type}
et
Descente, lemme \ref{descent-lemma-locally-finite-type-fppf-local-source}.
\end{proof}

\begin{lemma}
\label{spaces-descent-lemma-open-fppf-local-source}
La propriété $\mathcal{P}(f)=$``$f$ est ouvert''
est locale sur la source pour la topologie fppf.
\end{lemma}

\begin{proof}
L'assertion résulte du
lemme \ref{spaces-descent-lemma-transfer-from-schemes}
en utilisant
Morphismes d'espaces, définition \ref{spaces-morphisms-definition-open}
et
Descente, lemme \ref{descent-lemma-open-fppf-local-source}.
\end{proof}

\begin{lemma}
\label{spaces-descent-lemma-universally-open-fppf-local-source}
La propriété $\mathcal{P}(f)=$``$f$ est universellement ouvert''
est locale sur la source pour la topologie fppf.
\end{lemma}

\begin{proof}
L'assertion résulte du
lemme \ref{spaces-descent-lemma-transfer-from-schemes}
en utilisant
Morphismes d'espaces, définition \ref{spaces-morphisms-definition-open}
et
Descente, lemme \ref{descent-lemma-universally-open-fppf-local-source}.
\end{proof}



\section{Propriétés des morphismes locales sur la source pour la topologie syntomique}
\label{spaces-descent-section-syntomic-local-source}

\noindent
Voici quelques propriétés des morphismes qui sont locales sur la source pour la topologie syntomique.

\begin{lemma}
\label{spaces-descent-lemma-syntomic-syntomic-local-source}
La propriété $\mathcal{P}(f)=$``$f$ est syntomique''
est locale sur la source pour la topologie syntomique.
\end{lemma}

\begin{proof}
L'assertion résulte du
lemme \ref{spaces-descent-lemma-transfer-from-schemes}
en utilisant
Morphismes d'espaces, définition \ref{spaces-morphisms-definition-syntomic}
et
Descente, lemme \ref{descent-lemma-syntomic-syntomic-local-source}.
\end{proof}




\section{Propriétés des morphismes locales sur la source pour la topologie lisse}
\label{spaces-descent-section-smooth-local-source}

\noindent
Voici quelques propriétés des morphismes qui sont locales sur la source pour la topologie lisse.

\begin{lemma}
\label{spaces-descent-lemma-smooth-smooth-local-source}
La propriété $\mathcal{P}(f)=$``$f$ est lisse''
est locale sur la source pour la topologie lisse.
\end{lemma}

\begin{proof}
L'assertion résulte du
lemme \ref{spaces-descent-lemma-transfer-from-schemes}
en utilisant
Morphismes d'espaces, définition \ref{spaces-morphisms-definition-smooth}
et
Descente, lemme \ref{descent-lemma-smooth-smooth-local-source}.
\end{proof}



\section{Propriétés des morphismes locales sur la source pour la topologie \'etale}
\label{spaces-descent-section-etale-local-source}

\noindent
Voici quelques propriétés des morphismes qui sont locales sur la source pour la topologie \'etale.

\begin{lemma}
\label{spaces-descent-lemma-etale-etale-local-source}
La propriété $\mathcal{P}(f)=$``$f$ est \'etale''
est locale sur la source pour la topologie \'etale.
\end{lemma}

\begin{proof}
L'assertion résulte du
lemme \ref{spaces-descent-lemma-transfer-from-schemes}
en utilisant
Morphismes d'espaces,
définition \ref{spaces-morphisms-definition-etale}
et
Descente, lemme \ref{descent-lemma-etale-etale-local-source}.
\end{proof}

\begin{lemma}
\label{spaces-descent-lemma-locally-quasi-finite-etale-local-source}
La propriété $\mathcal{P}(f)=$``$f$ est localement quasi-fini''
est locale sur la source pour la topologie \'etale.
\end{lemma}

\begin{proof}
L'assertion résulte du
lemme \ref{spaces-descent-lemma-transfer-from-schemes}
en utilisant
Morphismes d'espaces,
définition \ref{spaces-morphisms-definition-locally-quasi-finite}
et
Descente, lemme \ref{descent-lemma-locally-quasi-finite-etale-local-source}.
\end{proof}

\begin{lemma}
\label{spaces-descent-lemma-unramified-etale-local-source}
La propriété $\mathcal{P}(f)=$``$f$ est non ramifié''
est locale sur la source pour la topologie \'etale.
\end{lemma}

\begin{proof}
L'assertion résulte du
lemme \ref{spaces-descent-lemma-transfer-from-schemes}
en utilisant
Morphismes d'espaces, définition \ref{spaces-morphisms-definition-unramified}
et
Descente, lemme \ref{descent-lemma-unramified-etale-local-source}.
\end{proof}


% preserved blank authority line 2612
% preserved blank authority line 2613
\section{Propriétés des morphismes locales pour la topologie lisse sur la source et le but}
\sectionmark{Propriétés locales pour la topologie lisse}
\label{spaces-descent-section-properties-local-source-target}

\noindent
Soit $\mathcal{P}$ une propriété des morphismes d'espaces algébriques. La phrase
« $\mathcal{P}$ est locale pour la topologie lisse sur la
source et le but » possède un sens intuitif. Toutefois, cette notion ne revient
pas à demander que $\mathcal{P}$ soit à la fois locale pour la topologie lisse
sur la source et locale pour la topologie lisse sur le but.
Nous avons étudié très en détail un phénomène analogue (pour la topologie \'etale
et la catégorie des schémas) dans
Descente, section \ref{descent-section-properties-etale-local-source-target}
(pour un aperçu rapide, voir
Descente, remarque \ref{descent-remark-compare-definitions}).
Il existe toutefois une différence importante entre le cas de la topologie lisse
et celui de la topologie \'etale. Pour la saisir, nous invitons le lecteur
à comparer
Descente, lemme \ref{descent-lemma-local-source-target-implies}
et
le lemme \ref{spaces-descent-lemma-local-source-target-implies},
ainsi qu'à comparer
Descente, lemme \ref{descent-lemma-local-source-target-characterize}
et
le lemme \ref{spaces-descent-lemma-local-source-target-characterize}.
En effet, dans le cadre étale, le choix du « recouvrement » étale du
but est indifférent, tandis qu'il ne l'est pas dans le cadre lisse.

\begin{definition}
\label{spaces-descent-definition-local-source-target}
Soit $S$ un schéma.
Soit $\mathcal{P}$ une propriété des morphismes d'espaces algébriques sur $S$.
Nous disons que $\mathcal{P}$ est {\it locale pour la topologie lisse sur la source et le but} si
\begin{enumerate}
\item (stabilité par précomposition par des morphismes lisses)
si $f : X \to Y$ est lisse et si $g : Y \to Z$ vérifie $\mathcal{P}$,
alors $g \circ f$ vérifie $\mathcal{P}$,
\item (stabilité par changement de base lisse)
si $f : X \to Y$ vérifie $\mathcal{P}$ et si $Y' \to Y$ est lisse, alors
le changement de base $f' : Y' \times_Y X \to Y'$ vérifie $\mathcal{P}$, et
\item (caractère local) pour un morphisme $f : X \to Y$, les conditions suivantes sont
équivalentes :
\begin{enumerate}
\item $f$ vérifie $\mathcal{P}$,
\item pour tout $x \in |X|$, il existe un diagramme commutatif
$$
\xymatrix{
U \ar[d]_a \ar[r]_h & V \ar[d]^b \\
X \ar[r]^f & Y
}
$$
dont les flèches verticales sont lisses et pour lequel il existe $u \in |U|$ tel que $a(u) = x$ et que
$h$ vérifie $\mathcal{P}$.
\end{enumerate}
\end{enumerate}
\end{definition}

\noindent
Ce qui précède nous sert de définition. Dans les lemmes ci-dessous, nous montrerons que
cela équivaut à ce que $\mathcal{P}$ soit locale pour la topologie lisse sur le but,
locale pour la topologie lisse sur la source et stable par postcomposition par des morphismes lisses.

\begin{lemma}
\label{spaces-descent-lemma-local-source-target-implies}
Soit $S$ un schéma.
Soit $\mathcal{P}$ une propriété des morphismes d'espaces algébriques sur $S$
qui est locale pour la topologie lisse sur la source et le but. Alors :
\begin{enumerate}
\item $\mathcal{P}$ est locale pour la topologie lisse sur la source,
\item $\mathcal{P}$ est locale pour la topologie lisse sur le but,
\item $\mathcal{P}$ est stable par postcomposition par des morphismes lisses :
si $f : X \to Y$ vérifie $\mathcal{P}$ et si $g : Y \to Z$ est lisse, alors
$g \circ f$ vérifie $\mathcal{P}$.
\end{enumerate}
\end{lemma}

\begin{proof}
Écrivons tous les détails.

\medskip\noindent
Preuve de (1). Soit $f : X \to Y$ un morphisme d'espaces algébriques sur $S$.
Soit $\{X_i \to X\}_{i \in I}$ un recouvrement lisse de $X$. Si chaque composé
$h_i : X_i \to Y$ vérifie $\mathcal{P}$, alors, pour tout $|x| \in X$, on peut trouver
un indice $i \in I$ et un point $x_i \in |X_i|$ qui s'envoie sur $x$. Alors
$(X_i, x_i) \to (X, x)$ est un morphisme lisse de couples,
$\text{id}_Y : Y \to Y$ est un morphisme lisse et $h_i$ est comme dans la partie (3) de la
définition \ref{spaces-descent-definition-local-source-target}.
On voit donc que $f$ vérifie $\mathcal{P}$.
Réciproquement, si $f$ vérifie $\mathcal{P}$, alors chaque $X_i \to Y$ vérifie
$\mathcal{P}$ d'après la
partie (1) de la définition \ref{spaces-descent-definition-local-source-target}.

\medskip\noindent
Preuve de (2). Soit $f : X \to Y$ un morphisme d'espaces algébriques sur $S$.
Soit $\{Y_i \to Y\}_{i \in I}$ un recouvrement lisse de $Y$.
Notons $X_i = Y_i \times_Y X$ et $h_i : X_i \to Y_i$ le changement de base
de $f$. Si chaque $h_i : X_i \to Y_i$ vérifie $\mathcal{P}$, alors, pour tout
$x \in |X|$, choisissons un indice $i \in I$ et un point $x_i \in |X_i|$ qui s'envoie sur $x$.
Alors $(X_i, x_i) \to (X, x)$ est un morphisme lisse de couples, $Y_i \to Y$ est
lisse et $h_i$ est comme dans la partie (3) de la
définition \ref{spaces-descent-definition-local-source-target}.
On voit donc que $f$ vérifie $\mathcal{P}$.
Réciproquement, si $f$ vérifie $\mathcal{P}$, alors chaque $X_i \to Y_i$ vérifie
$\mathcal{P}$ d'après la
partie (2) de la définition \ref{spaces-descent-definition-local-source-target}.

\medskip\noindent
Preuve de (3). Supposons que $f : X \to Y$ vérifie $\mathcal{P}$ et que $g : Y \to Z$ soit
lisse. Pour tout $x \in |X|$, considérons $(X, x) \to (X, x)$ comme un
morphisme lisse de couples ; $Y \to Z$ est un morphisme lisse, et $h = f$ est comme
dans la partie (3) de la
définition \ref{spaces-descent-definition-local-source-target}.
On voit donc que $g \circ f$ vérifie $\mathcal{P}$.
\end{proof}

\noindent
Le lemme suivant est l'analogue de
Morphismes, lemme \ref{morphisms-lemma-locally-P-characterize}.

\begin{lemma}
\label{spaces-descent-lemma-local-source-target-characterize}
Soit $S$ un schéma.
Soit $\mathcal{P}$ une propriété des morphismes d'espaces algébriques sur $S$
qui est locale pour la topologie lisse sur la source et le but. Soit $f : X \to Y$ un morphisme
d'espaces algébriques sur $S$. Les conditions suivantes sont équivalentes :
\begin{enumerate}
\item[(a)] $f$ vérifie la propriété $\mathcal{P}$,
\item[(b)] pour tout $x \in |X|$, il existe un morphisme lisse de couples
$a : (U, u) \to (X, x)$, un morphisme lisse $b : V \to Y$ et
un morphisme $h : U \to V$ tel que $f \circ a = b \circ h$ et que
$h$ vérifie $\mathcal{P}$,
\item[(c)] pour un certain diagramme commutatif
$$
\xymatrix{
U \ar[d]_a \ar[r]_h & V \ar[d]^b \\
X \ar[r]^f & Y
}
$$
où $a$ et $b$ sont lisses et $a$ est surjectif, le morphisme $h$ vérifie $\mathcal{P}$,
\item[(d)] pour tout diagramme commutatif
$$
\xymatrix{
U \ar[d]_a \ar[r]_h & V \ar[d]^b \\
X \ar[r]^f & Y
}
$$
où $b$ est lisse et $U \to X \times_Y V$ est lisse,
le morphisme $h$ vérifie $\mathcal{P}$,
\item[(e)] il existe un recouvrement lisse $\{Y_i \to Y\}_{i \in I}$ tel
que chaque changement de base $Y_i \times_Y X \to Y_i$ vérifie $\mathcal{P}$,
\item[(f)] il existe un recouvrement lisse $\{X_i \to X\}_{i \in I}$ tel
que chaque composé $X_i \to Y$ vérifie $\mathcal{P}$,
\item[(g)] il existe un recouvrement lisse $\{Y_i \to Y\}_{i \in I}$ et,
pour chaque $i \in I$, un recouvrement lisse
$\{X_{ij} \to Y_i \times_Y X\}_{j \in J_i}$ tel que chaque morphisme
$X_{ij} \to Y_i$ vérifie $\mathcal{P}$.
\end{enumerate}
\end{lemma}

\begin{proof}
L'équivalence de (a) et (b) fait partie de la
définition \ref{spaces-descent-definition-local-source-target}.
L'équivalence de (a) et (e) est donnée par la
partie (2) du lemme \ref{spaces-descent-lemma-local-source-target-implies}.
L'équivalence de (a) et (f) est donnée par la
partie (1) du lemme \ref{spaces-descent-lemma-local-source-target-implies}.
Puisque (a) est maintenant équivalente à (e) et à (f), il s'ensuit que
(a) est équivalente à (g).

\medskip\noindent
Il est clair que (c) implique (b). Si (b) est vérifiée, alors, pour tout
$x \in |X|$, nous pouvons choisir un morphisme lisse de couples
$a_x : (U_x, u_x) \to (X, x)$, un morphisme lisse $b_x : V_x \to Y$ et
un morphisme $h_x : U_x \to V_x$ tel que $f \circ a_x = b_x \circ h_x$ et que
$h_x$ vérifie $\mathcal{P}$. Alors $h = \coprod h_x : \coprod U_x \to \coprod V_x$
forme, avec $a = \coprod a_x$ et $b = \coprod b_x$, un diagramme comme en (c).
(Remarquons que $h$ vérifie la propriété $\mathcal{P}$, puisque $\{V_x \to \coprod V_x\}$
est un recouvrement lisse et que $\mathcal{P}$ est locale pour la topologie lisse sur le but.)
Ainsi, (b) et (c) sont équivalentes.

\medskip\noindent
Nous savons maintenant que (a), (b), (c), (e), (f) et (g) sont équivalentes.
Supposons (a) vérifiée. Soient $U, V, a, b, h$ comme en (d). Alors
$X \times_Y V \to V$ vérifie $\mathcal{P}$ puisque $\mathcal{P}$ est stable par
changement de base lisse ; il s'ensuit que $U \to V$ vérifie $\mathcal{P}$ puisque $\mathcal{P}$
est stable par précomposition par des morphismes lisses. Réciproquement, si (d)
est vérifiée, prendre $U = X$ et $V = Y$ montre que $f$ vérifie $\mathcal{P}$.
\end{proof}

\begin{lemma}
\label{spaces-descent-lemma-smooth-local-source-target}
Soit $S$ un schéma.
Soit $\mathcal{P}$ une propriété des morphismes d'espaces algébriques sur $S$.
Supposons que
\begin{enumerate}
\item $\mathcal{P}$ soit locale pour la topologie lisse sur la source,
\item $\mathcal{P}$ soit locale pour la topologie lisse sur le but, et
\item $\mathcal{P}$ soit stable par postcomposition par des morphismes lisses :
si $f : X \to Y$ vérifie $\mathcal{P}$ et si $Y \to Z$ est un morphisme lisse,
alors $X \to Z$ vérifie $\mathcal{P}$.
\end{enumerate}
Alors $\mathcal{P}$ est locale pour la topologie lisse sur la source et le but.
\end{lemma}

\begin{proof}
Soit $\mathcal{P}$ une propriété des morphismes d'espaces algébriques qui
satisfait aux conditions (1), (2) et (3) du lemme. D'après le
lemme \ref{spaces-descent-lemma-precompose-property-local-source},
on voit que $\mathcal{P}$ est stable par précomposition par des
morphismes lisses. D'après le
lemme \ref{spaces-descent-lemma-pullback-property-local-target},
on voit que $\mathcal{P}$ est stable par changement de base lisse.
Il suffit donc de prouver que la partie (3) de la
définition \ref{spaces-descent-definition-local-source-target}
est vérifiée.

\medskip\noindent
Plus précisément, supposons que $f : X \to Y$ soit un morphisme
d'espaces algébriques sur $S$ qui vérifie la
partie (3)(b) de la définition \ref{spaces-descent-definition-local-source-target}.
Autrement dit, pour tout $x \in X$, il existe un morphisme lisse
$a_x : U_x \to X$, un point $u_x \in |U_x|$ qui s'envoie sur $x$,
un morphisme lisse $b_x : V_x \to Y$ et un morphisme $h_x : U_x \to V_x$
tels que $f \circ a_x = b_x \circ h_x$ et que $h_x$ vérifie $\mathcal{P}$.
Pour achever la preuve du lemme, il suffit de montrer que $f$ vérifie $\mathcal{P}$.
Posons $U = \coprod U_x$, $a = \coprod a_x$, $V = \coprod V_x$,
$b = \coprod b_x$ et $h = \coprod h_x$. Nous obtenons un
diagramme commutatif
$$
\xymatrix{
U \ar[d]_a \ar[r]_h & V \ar[d]^b \\
X \ar[r]^f & Y
}
$$
où $a$ et $b$ sont lisses et $a$ est surjectif. Remarquons que $h$ vérifie $\mathcal{P}$
puisque chaque $h_x$ la vérifie et que $\mathcal{P}$ est locale pour la topologie lisse sur le but.
Comme $a$ est surjectif et que $\mathcal{P}$ est locale pour la topologie lisse sur la source,
il suffit de prouver que $b \circ h$ vérifie $\mathcal{P}$.
Cela résulte de l'hypothèse selon laquelle $\mathcal{P}$ est stable par
postcomposition par un morphisme lisse et du fait que $b$ est lisse.
\end{proof}

\begin{remark}
\label{spaces-descent-remark-list-local-source-target}
À l'aide du
lemme \ref{spaces-descent-lemma-smooth-local-source-target}
et des résultats établis dans les sections précédentes de ce chapitre, on dresse aisément
une liste de types de morphismes dont la propriété est locale pour la topologie lisse sur la
source et le but. Dans chaque cas, nous indiquons le lemme qui implique que
la propriété est locale pour la topologie lisse sur la source et celui qui implique qu'elle
est locale pour la topologie lisse sur le but. Dans chaque cas, la troisième hypothèse
du
lemme \ref{spaces-descent-lemma-smooth-local-source-target}
se vérifie immédiatement ; nous l'omettons. Voici la liste :
\begin{enumerate}
\item plat, voir les
lemmes \ref{spaces-descent-lemma-flat-fpqc-local-source} et
\ref{spaces-descent-lemma-descending-property-flat},
\item localement de présentation finie, voir les
lemmes \ref{spaces-descent-lemma-locally-finite-presentation-fppf-local-source} et
\ref{spaces-descent-lemma-descending-property-locally-finite-presentation},
\item localement de type fini, voir les
lemmes \ref{spaces-descent-lemma-locally-finite-type-fppf-local-source} et
\ref{spaces-descent-lemma-descending-property-locally-finite-type},
\item universellement ouvert, voir les
lemmes \ref{spaces-descent-lemma-universally-open-fppf-local-source} et
\ref{spaces-descent-lemma-descending-property-universally-open},
\item syntomique, voir les
lemmes \ref{spaces-descent-lemma-syntomic-syntomic-local-source} et
\ref{spaces-descent-lemma-descending-property-syntomic},
\item lisse, voir les
lemmes \ref{spaces-descent-lemma-smooth-smooth-local-source} et
\ref{spaces-descent-lemma-descending-property-smooth},
\item ajouter ici d'autres cas au besoin.
\end{enumerate}
\end{remark}






% CH074 slice boundary 2895
\section{Propriétés des morphismes locales pour la topologie \'etale sur la source et la topologie lisse sur le but}
\sectionmark{Topologie \'etale sur la source, lisse sur le but}
\label{spaces-descent-section-properties-etale-smooth-local-source-target}

\noindent
Cette section est l'analogue de la
section \ref{spaces-descent-section-properties-local-source-target}
pour les propriétés des morphismes qui sont locales pour la topologie \'etale
sur la source et pour la topologie lisse sur le but.
Nous donnons à cette propriété un nom ridiculement long
afin de ne pas avoir à l'employer trop souvent.

\begin{definition}
\label{spaces-descent-definition-etale-smooth-local-source-target}
Soit $S$ un schéma.
Soit $\mathcal{P}$ une propriété des morphismes d'espaces algébriques sur $S$.
Nous disons que $\mathcal{P}$ est {\it locale pour la topologie \'etale sur la source et pour la topologie lisse sur le but} si
\begin{enumerate}
\item (stabilité par précomposition par des morphismes \'etales)
si $f : X \to Y$ est \'etale et si $g : Y \to Z$ vérifie $\mathcal{P}$,
alors $g \circ f$ vérifie $\mathcal{P}$,
\item (stabilité par changement de base lisse)
si $f : X \to Y$ vérifie $\mathcal{P}$ et si $Y' \to Y$ est lisse, alors
le changement de base $f' : Y' \times_Y X \to Y'$ vérifie $\mathcal{P}$, et
\item (caractère local) pour un morphisme $f : X \to Y$, les conditions suivantes sont
équivalentes :
\begin{enumerate}
\item $f$ vérifie $\mathcal{P}$,
\item pour tout $x \in |X|$, il existe un diagramme commutatif
$$
\xymatrix{
U \ar[d]_a \ar[r]_h & V \ar[d]^b \\
X \ar[r]^f & Y
}
$$
où $b$ est lisse, où $U \to X \times_Y V$ est \'etale
et où il existe $u \in |U|$ tel que $a(u) = x$ et que
$h$ vérifie $\mathcal{P}$.
\end{enumerate}
\end{enumerate}
\end{definition}

\noindent
Ce qui précède nous sert de définition. Dans les lemmes ci-dessous, nous montrerons que
cela équivaut à ce que $\mathcal{P}$ soit locale pour la topologie \'etale sur le but,
locale pour la topologie lisse sur la source et stable par postcomposition par des
morphismes \'etales.

\begin{lemma}
\label{spaces-descent-lemma-etale-smooth-local-source-target-implies}
Soit $S$ un schéma.
Soit $\mathcal{P}$ une propriété des morphismes d'espaces algébriques sur $S$
qui est locale pour la topologie \'etale sur la source et pour la topologie lisse sur le but. Alors :
\begin{enumerate}
\item $\mathcal{P}$ est locale pour la topologie \'etale sur la source,
\item $\mathcal{P}$ est locale pour la topologie lisse sur le but,
\item $\mathcal{P}$ est stable par postcomposition par des morphismes \'etales :
si $f : X \to Y$ vérifie $\mathcal{P}$ et si $g : Y \to Z$ est \'etale, alors
$g \circ f$ vérifie $\mathcal{P}$, et
\item $\mathcal{P}$ possède la propriété de permanence suivante : si $f : X \to Y$ et
$g : Y \to Z$ est \'etale et si $g \circ f$ vérifie $\mathcal{P}$, alors
$f$ vérifie $\mathcal{P}$.
\end{enumerate}
\end{lemma}

\begin{proof}
Écrivons tous les détails.

\medskip\noindent
Preuve de (1). Soit $f : X \to Y$ un morphisme d'espaces algébriques sur $S$.
Soit $\{X_i \to X\}_{i \in I}$ un recouvrement \'etale de $X$. Si chaque composé
$h_i : X_i \to Y$ vérifie $\mathcal{P}$, alors, pour tout $|x| \in X$, on peut trouver
un indice $i \in I$ et un point $x_i \in |X_i|$ qui s'envoie sur $x$. Alors
$(X_i, x_i) \to (X, x)$ est un morphisme \'etale de couples,
$\text{id}_Y : Y \to Y$ est un morphisme lisse et $h_i$ est comme dans la partie (3) de la
définition \ref{spaces-descent-definition-etale-smooth-local-source-target}.
On voit donc que $f$ vérifie $\mathcal{P}$.
Réciproquement, si $f$ vérifie $\mathcal{P}$, alors chaque $X_i \to Y$ vérifie
$\mathcal{P}$ d'après la
partie (1) de la définition \ref{spaces-descent-definition-etale-smooth-local-source-target}.

\medskip\noindent
Preuve de (2). Soit $f : X \to Y$ un morphisme d'espaces algébriques sur $S$.
Soit $\{Y_i \to Y\}_{i \in I}$ un recouvrement lisse de $Y$.
Notons $X_i = Y_i \times_Y X$ et $h_i : X_i \to Y_i$ le changement de base
de $f$. Si chaque $h_i : X_i \to Y_i$ vérifie $\mathcal{P}$, alors, pour tout
$x \in |X|$, choisissons un indice $i \in I$ et un point $x_i \in |X_i|$ qui s'envoie sur $x$.
Alors $X_i \to X \times_Y Y_i$ est un morphisme \'etale
(puisque c'est un isomorphisme), $Y_i \to Y$ est
lisse et $h_i$ est comme dans la partie (3) de la
définition \ref{spaces-descent-definition-local-source-target}.
On voit donc que $f$ vérifie $\mathcal{P}$.
Réciproquement, si $f$ vérifie $\mathcal{P}$, alors chaque $X_i \to Y_i$ vérifie
$\mathcal{P}$ d'après la
partie (2) de la définition \ref{spaces-descent-definition-local-source-target}.

\medskip\noindent
Preuve de (3). Supposons que $f : X \to Y$ vérifie $\mathcal{P}$ et que $g : Y \to Z$ soit
\'etale. Le morphisme $X \to Y \times_Z X$ est \'etale en tant que morphisme
entre des espaces algébriques \'etales sur $X$ (
Propriétés des espaces, lemme \ref{spaces-properties-lemma-etale-permanence}).
De plus, $Y \to Z$ est \'etale, donc lisse.
Ainsi, le diagramme
$$
\xymatrix{
X \ar[d] \ar[r]_f & Y \ar[d] \\
X \ar[r]^{g \circ f} & Z
}
$$
convient, pour tout $x \in |X|$, dans la partie (3) de la
définition \ref{spaces-descent-definition-local-source-target},
et nous concluons que $g \circ f$ vérifie $\mathcal{P}$.

\medskip\noindent
Preuve de (4). Soient $f : X \to Y$ un morphisme et $g : Y \to Z$ \'etale
tels que $g \circ f$ vérifie $\mathcal{P}$. Alors, d'après la
partie (2) de la définition \ref{spaces-descent-definition-etale-smooth-local-source-target},
on voit que $\text{pr}_Y : Y \times_Z X \to Y$ vérifie $\mathcal{P}$. Mais
le morphisme $(f, 1) : X \to Y \times_Z X$ est \'etale, car c'est une section de la
projection \'etale $\text{pr}_X : Y \times_Z X \to X$ ; voir
Morphismes d'espaces, lemme \ref{spaces-morphisms-lemma-etale-permanence}.
Ainsi, $f = \text{pr}_Y \circ (f, 1)$ vérifie $\mathcal{P}$ d'après la
partie (1) de la définition \ref{spaces-descent-definition-etale-smooth-local-source-target}.
\end{proof}

\begin{lemma}
\label{spaces-descent-lemma-etale-smooth-local-source-target-characterize}
Soit $S$ un schéma.
Soit $\mathcal{P}$ une propriété des morphismes d'espaces algébriques sur $S$
qui est locale pour la topologie étale sur la source et pour la topologie lisse sur le but.
Soit $f : X \to Y$ un morphisme
d'espaces algébriques sur $S$. Les conditions suivantes sont équivalentes :
\begin{enumerate}
\item[(a)] $f$ vérifie la propriété $\mathcal{P}$,
\item[(b)] pour tout $x \in |X|$, il existe un morphisme lisse $b : V \to Y$,
un morphisme \'etale $a : U \to V \times_Y X$ et un point $u \in |U|$
qui s'envoie sur $x$, de telle sorte que $U \to V$ vérifie $\mathcal{P}$,
\item[(c)] pour un certain diagramme commutatif
$$
\xymatrix{
U \ar[d]_a \ar[r]_h & V \ar[d]^b \\
X \ar[r]^f & Y
}
$$
où $b$ est lisse, $U \to V \times_Y X$ est \'etale et $a$ est surjectif,
le morphisme $h$ vérifie $\mathcal{P}$,
\item[(d)] pour tout diagramme commutatif
$$
\xymatrix{
U \ar[d]_a \ar[r]_h & V \ar[d]^b \\
X \ar[r]^f & Y
}
$$
où $b$ est lisse et $U \to X \times_Y V$ est \'etale, le morphisme $h$
vérifie $\mathcal{P}$,
\item[(e)] il existe un recouvrement lisse $\{Y_i \to Y\}_{i \in I}$ tel
que chaque changement de base $Y_i \times_Y X \to Y_i$ vérifie $\mathcal{P}$,
\item[(f)] il existe un recouvrement \'etale $\{X_i \to X\}_{i \in I}$ tel
que chaque composé $X_i \to Y$ vérifie $\mathcal{P}$,
\item[(g)] il existe un recouvrement lisse $\{Y_i \to Y\}_{i \in I}$ et,
pour chaque $i \in I$, un recouvrement \'etale
$\{X_{ij} \to Y_i \times_Y X\}_{j \in J_i}$ tel que chaque morphisme
$X_{ij} \to Y_i$ vérifie $\mathcal{P}$.
\end{enumerate}
\end{lemma}

\begin{proof}
L'équivalence de (a) et (b) fait partie de la
définition \ref{spaces-descent-definition-etale-smooth-local-source-target}.
L'équivalence de (a) et (e) est donnée par la
partie (2) du lemme \ref{spaces-descent-lemma-etale-smooth-local-source-target-implies}.
L'équivalence de (a) et (f) est donnée par la
partie (1) du lemme \ref{spaces-descent-lemma-etale-smooth-local-source-target-implies}.
Puisque (a) est maintenant équivalente à (e) et à (f), il s'ensuit que
(a) est équivalente à (g).

\medskip\noindent
Il est clair que (c) implique (b). Si (b) est vérifiée, alors, pour tout
$x \in |X|$, nous pouvons choisir un morphisme lisse
$b_x : V_x \to Y$, un morphisme \'etale $U_x \to V_x \times_Y X$,
et un point $u_x \in |U_x|$ qui s'envoie sur $x$,
de telle sorte que $U_x \to V_x$ vérifie $\mathcal{P}$.
Alors $h = \coprod h_x : \coprod U_x \to \coprod V_x$
forme, avec $a = \coprod a_x$ et $b = \coprod b_x$, un diagramme comme en (c).
(Remarquons que $h$ vérifie la propriété $\mathcal{P}$, puisque $\{V_x \to \coprod V_x\}$
est un recouvrement lisse et que $\mathcal{P}$ est locale pour la topologie lisse sur le but.)
Ainsi, (b) et (c) sont équivalentes.

\medskip\noindent
Nous savons maintenant que (a), (b), (c), (e), (f) et (g) sont équivalentes.
Supposons (a) vérifiée. Soient $U, V, a, b, h$ comme en (d). Alors
$X \times_Y V \to V$ vérifie $\mathcal{P}$ puisque $\mathcal{P}$ est stable par
changement de base lisse ; il s'ensuit que $U \to V$ vérifie $\mathcal{P}$ puisque $\mathcal{P}$
est stable par précomposition par des morphismes \'etales. Réciproquement, si (d)
est vérifiée, prendre $U = X$ et $V = Y$ montre que $f$ vérifie $\mathcal{P}$.
\end{proof}

\begin{lemma}
\label{spaces-descent-lemma-etale-smooth-local-source-target}
Soit $S$ un schéma.
Soit $\mathcal{P}$ une propriété des morphismes d'espaces algébriques sur $S$.
Supposons que
\begin{enumerate}
\item $\mathcal{P}$ soit locale pour la topologie \'etale sur la source,
\item $\mathcal{P}$ soit locale pour la topologie lisse sur le but, et
\item $\mathcal{P}$ soit stable par postcomposition par des immersions ouvertes :
si $f : X \to Y$ vérifie $\mathcal{P}$ et si $Y \subset Z$ est une immersion ouverte,
alors $X \to Z$ vérifie $\mathcal{P}$.
\end{enumerate}
Alors $\mathcal{P}$ est locale pour la topologie \'etale sur la source et pour la topologie lisse sur le but.
\end{lemma}

\begin{proof}
Soit $\mathcal{P}$ une propriété des morphismes d'espaces algébriques qui
satisfait aux conditions (1), (2) et (3) du lemme. D'après le
lemme \ref{spaces-descent-lemma-precompose-property-local-source},
on voit que $\mathcal{P}$ est stable par précomposition par des
morphismes \'etales. D'après le
lemme \ref{spaces-descent-lemma-pullback-property-local-target},
on voit que $\mathcal{P}$ est stable par changement de base lisse.
Il suffit donc de prouver que la partie (3) de la
définition \ref{spaces-descent-definition-local-source-target}
est vérifiée.

\medskip\noindent
Plus précisément, supposons que $f : X \to Y$ soit un morphisme
d'espaces algébriques sur $S$ qui vérifie la
partie (3)(b) de la définition \ref{spaces-descent-definition-local-source-target}.
Autrement dit, pour tout $x \in X$, il existe
un morphisme lisse $b_x : V_x \to Y$,
un morphisme \'etale $U_x \to V_x \times_Y X$ et
un point $u_x \in |U_x|$ qui s'envoie sur $x$
de telle sorte que $h_x : U_x \to V_x$ vérifie $\mathcal{P}$.
Pour achever la preuve du lemme, il suffit de montrer que $f$ vérifie $\mathcal{P}$.

\medskip\noindent
Soit $a_x : U_x \to X$ le composé $U_x \to V_x \times_Y X \to X$.
Posons $U = \coprod U_x$, $a = \coprod a_x$, $V = \coprod V_x$,
$b = \coprod b_x$ et $h = \coprod h_x$. Nous obtenons un
diagramme commutatif
$$
\xymatrix{
U \ar[d]_a \ar[r]_h & V \ar[d]^b \\
X \ar[r]^f & Y
}
$$
où $b$ est lisse, $U \to V \times_Y X$ est \'etale et $a$ est surjectif.
Remarquons que $h$ vérifie $\mathcal{P}$ puisque chaque $h_x$ la vérifie et que $\mathcal{P}$
est locale pour la topologie lisse sur le but. Au paragraphe suivant, nous montrons
que l'on peut supposer que $U, V, X, Y$ sont des schémas ; nous invitons le lecteur
à passer ce paragraphe.

\medskip\noindent
Soient $X, Y, U, V, a, b, f, h$ comme au paragraphe précédent. Il faut
montrer que $f$ vérifie $\mathcal{P}$. Soit $X' \to X$ un morphisme \'etale surjectif
avec $X_i$ un schéma. Posons $U' = X' \times_X U$. Alors
$U' \to X'$ est surjectif et $U' \to X' \times_Y V$ est \'etale.
Comme $\mathcal{P}$ est locale pour la topologie \'etale sur la source, on voit que
$U' \to V$ vérifie $\mathcal{P}$ et qu'il suffit de montrer que
$X' \to Y$ vérifie $\mathcal{P}$. Autrement dit, on peut supposer
que $X$ est un schéma. Choisissons ensuite un morphisme \'etale surjectif
$Y' \to Y$, où $Y'$ est un schéma. Posons $V' = V \times_Y Y'$,
$X' = X \times_Y Y'$ et $U' = U \times_Y Y'$. Alors
$U' \to X'$ est surjectif et $U' \to X' \times_{Y'} V'$ est \'etale.
Comme $\mathcal{P}$ est locale pour la topologie lisse sur le but, on voit que $U' \to V'$ vérifie
$\mathcal{P}$ et qu'il suffit de prouver que $X' \to Y'$ vérifie $\mathcal{P}$.
Ainsi, on peut supposer que $X$ et $Y$ sont tous deux des schémas.
Choisissons un morphisme \'etale surjectif $V' \to V$
avec $V'$ un schéma. Posons $U' = U \times_V V'$.
Alors $U' \to X$ est surjectif et $U' \to X \times_Y V'$ est \'etale.
Comme $\mathcal{P}$ est locale pour la topologie lisse sur la source, on voit que
$U' \to V'$ vérifie $\mathcal{P}$. Nous pouvons donc remplacer $U, V$ par
$U', V'$ et supposer que $X, Y, V$ sont des schémas.
Enfin, remplaçons $U$ par un schéma muni d'un morphisme \'etale surjectif vers $U$ ;
on voit alors que l'on peut supposer que $U, V, X, Y$ sont tous des schémas.

\medskip\noindent
Si $U, V, X, Y$ sont des schémas, alors $f$ vérifie $\mathcal{P}$
d'après Descente, lemme \ref{descent-lemma-etale-tau-local-source-target}.
\end{proof}

\begin{remark}
\label{spaces-descent-remark-list-etale-smooth-local-source-target}
À l'aide du lemme \ref{spaces-descent-lemma-etale-smooth-local-source-target}
et des résultats établis dans les sections précédentes de ce chapitre, on dresse aisément
une liste de types de morphismes dont la propriété est locale pour la topologie lisse sur la
source et le but. Dans chaque cas, nous indiquons le lemme qui implique que
la propriété est locale pour la topologie étale sur la source et celui qui implique qu'elle
est locale pour la topologie lisse sur le but. Dans chaque cas, la troisième hypothèse
du lemme \ref{spaces-descent-lemma-etale-smooth-local-source-target}
se vérifie immédiatement ; nous l'omettons. Voici la liste :
\begin{enumerate}
\item \'etale, voir les
lemmes \ref{spaces-descent-lemma-etale-etale-local-source} et
\ref{spaces-descent-lemma-descending-property-etale},
\item localement quasi-fini, voir les
lemmes \ref{spaces-descent-lemma-locally-quasi-finite-etale-local-source} et
\ref{spaces-descent-lemma-descending-property-quasi-finite},
\item non ramifié, voir les
lemmes \ref{spaces-descent-lemma-unramified-etale-local-source} et
\ref{spaces-descent-lemma-descending-property-unramified}, et
\item ajouter ici d'autres cas au besoin.
\end{enumerate}
Bien entendu, toute propriété énumérée dans la
remarque \ref{spaces-descent-remark-list-local-source-target}
est a fortiori un exemple que l'on pourrait faire figurer ici.
\end{remark}






% CH074 slice boundary 3208
\section{Données de descente pour des espaces au-dessus d'espaces}
\label{spaces-descent-section-descent-datum}

\noindent
Cette section est l'analogue de Descente, section
\ref{descent-section-descent-datum}, dans le cas des espaces algébriques.
La plupart des arguments de cette section sont formels et ne reposent
que sur la définition d'une donnée de descente.

\begin{definition}
\label{spaces-descent-definition-descent-datum}
Soit $S$ un schéma. Soit $f : Y \to X$ un morphisme d'espaces algébriques
au-dessus de $S$.
\begin{enumerate}
\item Soit $V \to Y$ un morphisme d'espaces algébriques.
Une {\it donnée de descente pour $V/Y/X$} est un isomorphisme
$\varphi : V \times_X Y \to Y \times_X V$ d'espaces algébriques au-dessus de
$Y \times_X Y$ satisfaisant à la {\it condition de cocycle} suivante : le diagramme
$$
\xymatrix{
V \times_X Y \times_X Y \ar[rd]^{\varphi_{01}} \ar[rr]_{\varphi_{02}} &
&
Y \times_X Y \times_X V\\
&
Y \times_X V \times_X Y \ar[ru]^{\varphi_{12}}
}
$$
est commutatif (avec les notations évidentes).
\item On dit également que le couple $(V/Y, \varphi)$ est
une {\it donnée de descente relativement à $Y \to X$}.
\item Un {\it morphisme $f : (V/Y, \varphi) \to (V'/Y, \varphi')$ de données
de descente relativement à $Y \to X$} est un morphisme
$f : V \to V'$ d'espaces algébriques au-dessus de $Y$ tel que
le diagramme
$$
\xymatrix{
V \times_X Y \ar[r]_{\varphi} \ar[d]_{f \times \text{id}_Y} &
Y \times_X V \ar[d]^{\text{id}_Y \times f} \\
V' \times_X Y \ar[r]^{\varphi'} & Y \times_X V'
}
$$
soit commutatif.
\end{enumerate}
\end{definition}

\begin{remark}
\label{spaces-descent-remark-easier}
Soit $S$ un schéma.
Soit $Y \to X$ un morphisme d'espaces algébriques au-dessus de $S$.
Soit $(V/Y, \varphi)$ une donnée de descente relativement à $Y \to X$.
On peut considérer l'isomorphisme $\varphi$ comme un isomorphisme
$$
(Y \times_X Y) \times_{\text{pr}_0, Y} V
\longrightarrow
(Y \times_X Y) \times_{\text{pr}_1, Y} V
$$
d'espaces algébriques au-dessus de $Y \times_X Y$. Ainsi, de manière informelle, on peut
considérer $\varphi$ comme une application
$\varphi : \text{pr}_0^*V \to \text{pr}_1^*V$\footnote{Malheureusement,
nous avons choisi ici le « mauvais » sens pour notre flèche. Dans les
définitions \ref{spaces-descent-definition-descent-datum} et
\ref{spaces-descent-definition-descent-datum-for-family-of-morphisms}
il faudrait prendre le sens opposé à celui adopté dans la
définition \ref{spaces-descent-definition-descent-datum-quasi-coherent},
selon le principe général que les « fonctions » et les « espaces » sont duaux.}.
La condition de cocycle exprime alors que
$\text{pr}_{02}^*\varphi =
\text{pr}_{12}^*\varphi \circ \text{pr}_{01}^*\varphi$.
Elle est ainsi très semblable à celle d'une donnée de descente sur des
faisceaux quasi-cohérents.
\end{remark}

\noindent
Voici la définition lorsqu'on dispose d'une famille de morphismes
à cible fixée.

\begin{definition}
\label{spaces-descent-definition-descent-datum-for-family-of-morphisms}
Soit $S$ un schéma.
Soit $\{X_i \to X\}_{i \in I}$ une famille de morphismes
 d'espaces algébriques au-dessus de $S$, à cible fixée $X$.
\begin{enumerate}
\item Une {\it donnée de descente $(V_i, \varphi_{ij})$ relativement à la
famille $\{X_i \to X\}$} est constituée d'un espace algébrique $V_i$ au-dessus de $X_i$
pour chaque $i \in I$, et d'un isomorphisme
$\varphi_{ij} : V_i \times_X X_j \to X_i \times_X V_j$
d'espaces algébriques au-dessus de $X_i \times_X X_j$, pour chaque couple $(i, j) \in I^2$,
tel que, pour tout triplet d'indices $(i, j, k) \in I^3$,
le diagramme
$$
\xymatrix{
V_i \times_X X_j \times_X X_k
\ar[rd]^{\text{pr}_{01}^*\varphi_{ij}}
\ar[rr]_{\text{pr}_{02}^*\varphi_{ik}} &
&
X_i \times_X X_j \times_X V_k\\
&
X_i \times_X V_j \times_X X_k
\ar[ru]^{\text{pr}_{12}^*\varphi_{jk}}
}
$$
d'espaces algébriques au-dessus de $X_i \times_X X_j \times_X X_k$ soit commutatif
(avec les notations évidentes).
\item Un {\it morphisme
$\psi : (V_i, \varphi_{ij}) \to (V'_i, \varphi'_{ij})$
de données de descente} est constitué d'une famille $\psi = (\psi_i)_{i \in I}$
de morphismes $\psi_i : V_i \to V'_i$ d'espaces algébriques au-dessus de $X_i$
telle que tous les diagrammes
$$
\xymatrix{
V_i \times_X X_j \ar[r]_{\varphi_{ij}} \ar[d]_{\psi_i \times \text{id}} &
X_i \times_X V_j \ar[d]^{\text{id} \times \psi_j} \\
V'_i \times_X X_j \ar[r]^{\varphi'_{ij}} & X_i \times_X V'_j
}
$$
soient commutatifs.
\end{enumerate}
\end{definition}

\begin{remark}
\label{spaces-descent-remark-easier-family}
Soit $S$ un schéma.
Soit $\{X_i \to X\}_{i \in I}$ une famille de morphismes
 d'espaces algébriques au-dessus de $S$, à cible fixée $X$.
Soit $(V_i, \varphi_{ij})$ une donnée de descente relativement à
$\{X_i \to X\}$. On peut considérer les isomorphismes $\varphi_{ij}$
comme des isomorphismes
$$
(X_i \times_X X_j) \times_{\text{pr}_0, X_i} V_i
\longrightarrow
(X_i \times_X X_j) \times_{\text{pr}_1, X_j} V_j
$$
d'espaces algébriques au-dessus de $X_i \times_X X_j$. Ainsi, de manière informelle, on peut
considérer $\varphi_{ij}$ comme un isomorphisme
$\text{pr}_0^*V_i \to \text{pr}_1^*V_j$ au-dessus de $X_i \times_X X_j$.
La condition de cocycle exprime alors que
$\text{pr}_{02}^*\varphi_{ik} =
\text{pr}_{12}^*\varphi_{jk} \circ \text{pr}_{01}^*\varphi_{ij}$.
Elle est ainsi très semblable à celle d'une donnée de descente sur des
faisceaux quasi-cohérents.
\end{remark}

\noindent
La raison pour laquelle nous travaillerons habituellement avec une famille réduite
à un seul morphisme est le lemme suivant.

\begin{lemma}
\label{spaces-descent-lemma-family-is-one}
Soit $S$ un schéma.
Soit $\{X_i \to X\}_{i \in I}$ une famille de morphismes
 d'espaces algébriques au-dessus de $S$, à cible fixée $X$.
Posons $Y = \coprod_{i \in I} X_i$.
Il existe une équivalence canonique de catégories
$$
\begin{matrix}
\text{catégorie des données de descente } \\
\text{relativement à la famille } \{X_i \to X\}_{i \in I}
\end{matrix}
\longrightarrow
\begin{matrix}
\text{ catégorie des données de descente} \\
\text{ relativement à } Y/X
\end{matrix}
$$
qui envoie $(V_i, \varphi_{ij})$ sur $(V, \varphi)$, où
$V = \coprod_{i\in I} V_i$ et $\varphi = \coprod \varphi_{ij}$.
\end{lemma}

\begin{proof}
Remarquons que $Y \times_X Y = \coprod_{ij} X_i \times_X X_j$
et de même pour les produits fibrés itérés.
Se donner un morphisme $V \to Y$ revient exactement à
se donner une famille $V_i \to X_i$. De même, se donner une donnée de descente
$\varphi$ revient exactement à se donner une famille $\varphi_{ij}$.
\end{proof}

\begin{lemma}
\label{spaces-descent-lemma-pullback}
Changement de base des données de descente. Soit $S$ un schéma.
\begin{enumerate}
\item Soit
$$
\xymatrix{
Y' \ar[r]_f \ar[d]_{a'} & Y \ar[d]^a \\
X' \ar[r]^h & X
}
$$
soit un diagramme commutatif d'espaces algébriques au-dessus de $S$.
La construction
$$
(V \to Y, \varphi) \longmapsto f^*(V \to Y, \varphi) = (V' \to Y', \varphi')
$$
où $V' = Y' \times_Y V$ et où
$\varphi'$ est définie comme la composée
$$
\xymatrix{
V' \times_{X'} Y' \ar@{=}[r] &
(Y' \times_Y V) \times_{X'} Y' \ar@{=}[r] &
(Y' \times_{X'} Y') \times_{Y \times_X Y} (V \times_X Y)
\ar[d]^{\text{id} \times \varphi} \\
Y' \times_{X'} V' \ar@{=}[r] &
Y' \times_{X'} (Y' \times_Y V) &
(Y' \times_X Y') \times_{Y \times_X Y} (Y \times_X V) \ar@{=}[l]
}
$$
définit un foncteur de la catégorie des données de descente
relativement à $Y \to X$ vers la catégorie des données de descente
relativement à $Y' \to X'$.
\item Étant donnés deux morphismes $f_i : Y' \to Y$, $i = 0, 1$, rendant le
diagramme commutatif, les foncteurs $f_0^*$ et $f_1^*$ sont
canoniquement isomorphes.
\end{enumerate}
\end{lemma}

\begin{proof}
Nous omettons la démonstration de (1), mais remarquons que le morphisme
$\varphi'$ est le morphisme $(f \times f)^*\varphi$ dans les
notations introduites dans la remarque \ref{spaces-descent-remark-easier}.
Pour (2), indiquons quel morphisme
$f_0^*V \to f_1^*V$ fournit l'isomorphisme fonctoriel. En effet,
puisque $f_0$ et $f_1$ s'insèrent tous deux dans le diagramme commutatif,
il existe un unique morphisme $r : Y' \to Y \times_X Y$
tel que $f_i = \text{pr}_i \circ r$. Prenons alors
\begin{eqnarray*}
f_0^*V & = &
Y' \times_{f_0, Y} V \\
& = &
Y' \times_{\text{pr}_0 \circ r, Y} V \\
& = &
Y' \times_{r, Y \times_X Y} (Y \times_X Y) \times_{\text{pr}_0, Y} V \\
& \xrightarrow{\varphi} &
Y' \times_{r, Y \times_X Y} (Y \times_X Y) \times_{\text{pr}_1, Y} V \\
& = &
Y' \times_{\text{pr}_1 \circ r, Y} V \\
& = &
Y' \times_{f_1, Y} V \\
& = & f_1^*V
\end{eqnarray*}
Nous omettons la vérification.
\end{proof}

\begin{definition}
\label{spaces-descent-definition-pullback-functor}
Avec $S, X, X', Y, Y', f, a, a', h$ comme dans le lemme \ref{spaces-descent-lemma-pullback},
le foncteur
$$
(V, \varphi) \longmapsto f^*(V, \varphi)
$$
construit dans ce lemme est appelé le {\it foncteur de changement de base} des données de descente.
\end{definition}

\begin{lemma}
\label{spaces-descent-lemma-pullback-family}
Soit $S$ un schéma. Soient $\mathcal{U}' = \{X'_i \to X'\}_{i \in I'}$ et
$\mathcal{U} = \{X_j \to X\}_{i \in I}$ des familles de morphismes à
cible fixée. Soient $\alpha : I' \to I$, $g : X' \to X$ et les morphismes
$g_i : X'_i \to X_{\alpha(i)}$ les composantes d'un morphisme de familles
de morphismes à cible fixée, voir
Sites, définition \ref{sites-definition-morphism-coverings}.
\begin{enumerate}
\item Soit $(V_i, \varphi_{ij})$ une donnée de descente relativement à la
famille $\mathcal{U}$. Le système
$$
\left(
g_i^*V_{\alpha(i)}, (g_i \times g_j)^*\varphi_{\alpha(i) \alpha(j)}
\right)
$$
(avec les notations de la remarque \ref{spaces-descent-remark-easier-family})
est une donnée de descente relativement à $\mathcal{U}'$.
\item Cette construction définit un foncteur de la catégorie des
données de descente relativement à $\mathcal{U}$ vers la catégorie des
données de descente relativement à $\mathcal{U}'$.
\item Soient aussi $\beta : I' \to I$, $h : X' \to X$ et les morphismes
$h'_i : X'_i \to X_{\beta(i)}$ les composantes d'un second morphisme de familles
de morphismes à cible fixée. Si $g = h$, alors les deux foncteurs obtenus
entre les catégories de données de descente sont canoniquement isomorphes.
\item Via le lemme \ref{spaces-descent-lemma-family-is-one}, ces foncteurs coïncident
avec les foncteurs de changement de base construits dans le lemme \ref{spaces-descent-lemma-pullback}.
\end{enumerate}
\end{lemma}

\begin{proof}
L'assertion résulte du lemme \ref{spaces-descent-lemma-pullback} au moyen de la
correspondance du lemme \ref{spaces-descent-lemma-family-is-one}.
\end{proof}

\begin{definition}
\label{spaces-descent-definition-pullback-functor-family}
Avec $\mathcal{U}' = \{X'_i \to X'\}_{i \in I'}$,
$\mathcal{U} = \{X_i \to X\}_{i \in I}$, $\alpha : I' \to I$,
$g : X' \to X$ et $g_i : X'_i \to X_{\alpha(i)}$ comme dans le
lemme \ref{spaces-descent-lemma-pullback-family}, le foncteur
$$
(V_i, \varphi_{ij}) \longmapsto
(g_i^*V_{\alpha(i)}, (g_i \times g_j)^*\varphi_{\alpha(i) \alpha(j)})
$$
construit dans ce lemme
est appelé le {\it foncteur de changement de base} des données de descente.
\end{definition}

\noindent
Si $\mathcal{U}$ et $\mathcal{U}'$ ont la même cible $X$,
et si $\mathcal{U}'$ raffine $\mathcal{U}$ (voir
Sites, définition \ref{sites-definition-morphism-coverings}),
mais qu'aucun couple explicite $(\alpha, g_i)$ n'est donné, nous pouvons néanmoins
parler du foncteur de changement de base, puisque le
lemme \ref{spaces-descent-lemma-pullback-family} montre que le choix du couple n'importe pas
(à isomorphisme canonique près).

\begin{definition}
\label{spaces-descent-definition-effective}
Soit $S$ un schéma. Soit $f : Y \to X$ un morphisme d'espaces algébriques au-dessus de
$S$.
\begin{enumerate}
\item À tout espace algébrique $U$ au-dessus de $X$ est associée la
{\it donnée de descente triviale} de $U$ relativement à $\text{id} : X \to X$, à savoir
le morphisme identité de $U$.
\item Le lemme \ref{spaces-descent-lemma-pullback} fournit alors une
{\it donnée de descente canonique} sur $Y \times_X U$
relativement à $Y \to X$, obtenue par changement de base de la donnée de descente
triviale par $f$. Cette donnée de descente est souvent
notée $(Y \times_X U, can)$.
\item Une donnée de descente $(V, \varphi)$ relativement à $Y/X$
est dite {\it effective} si $(V, \varphi)$
est isomorphe à la donnée de descente canonique
$(Y \times_X U, can)$ pour un espace algébrique $U$ au-dessus de $X$.
\end{enumerate}
\end{definition}

\noindent
Ainsi, l'effectivité signifie qu'il existe un espace algébrique $U$
au-dessus de $X$ et un isomorphisme $\psi : V \to Y \times_X U$
au-dessus de $Y$ tels que $\varphi$ soit égale à la composée
$$
V \times_X Y \xrightarrow{\psi \times \text{id}_Y}
Y \times_X U \times_S Y =
Y \times_X Y \times_X U
\xrightarrow{\text{id}_Y \times \psi^{-1}}
Y \times_X V
$$
Il existe ici une légère difficulté : cette définition
(dans son esprit) entre en conflit avec celle donnée dans
Descente, définition \ref{descent-definition-effective},
lorsque $Y$ et $X$ sont des schémas. Cependant, le contexte permettra toujours
de savoir laquelle des deux est visée.

\begin{definition}
\label{spaces-descent-definition-effective-family}
Soit $S$ un schéma.
Soit $\{X_i \to X\}$ une famille de morphismes d'espaces algébriques au-dessus de $S$,
à cible fixée $X$.
\begin{enumerate}
\item Étant donné un espace algébrique $U$ au-dessus de $X$,
il existe une {\it donnée de descente canonique} sur la famille des
espaces algébriques $X_i \times_X U$, obtenue par changement de base de la donnée de descente
triviale de $U$ relativement à $\{\text{id} : S \to S\}$.
Nous notons cette donnée de descente $(X_i \times_X U, can)$.
\item Une donnée de descente $(V_i, \varphi_{ij})$
relativement à $\{X_i \to S\}$ est dite {\it effective}
s'il existe un espace algébrique $U$ au-dessus de $X$ tel que
$(V_i, \varphi_{ij})$ soit isomorphe à $(X_i \times_X U, can)$.
\end{enumerate}
\end{definition}






\section{Données de descente en termes de faisceaux}
\label{spaces-descent-section-descent-data-sheaves}

\noindent
Cette section est l'analogue de Descente, section
\ref{descent-section-descent-data-sheaves}.
Elle diffère légèrement, puisque les espaces algébriques sont déjà des faisceaux.

\begin{lemma}
\label{spaces-descent-lemma-descent-data-sheaves}
Soit $S$ un schéma. Soit $\{X_i \to X\}_{i \in I}$ un recouvrement fppf
d'espaces algébriques au-dessus de $S$ (Topologies sur les espaces,
définition \ref{spaces-topologies-definition-fppf-covering}).
Il existe une équivalence de catégories
$$
\left\{
\begin{matrix}
\text{données de descente }(V_i, \varphi_{ij})\\
\text{relativement à }\{X_i \to X\}
\end{matrix}
\right\}
\leftrightarrow
\left\{
\begin{matrix}
\text{faisceaux }F\text{ sur }(\Sch/S)_{fppf}\text{ munis}\\
\text{d'un morphisme }F \to X\text{ tel que chaque}\\
X_i \times_X F\text{ soit un espace algébrique}
\end{matrix}
\right\}.
$$
De plus,
\begin{enumerate}
\item l'espace algébrique $X_i \times_X F$ du membre de droite
correspond à $V_i$ du membre de gauche, et
\item le faisceau $F$ est un espace algébrique\footnote{Nous verrons
plus loin que c'est toujours le cas si $I$ n'est pas trop grand, voir
Amorçage, lemme \ref{bootstrap-lemma-descend-algebraic-space}.}
si et seulement si la
donnée de descente $(X_i, \varphi_{ij})$ correspondante est effective.
\end{enumerate}
\end{lemma}

\begin{proof}
Construisons le foncteur de droite à gauche.
Soit $F \to X$ un morphisme de faisceaux sur $(\Sch/S)_{fppf}$ tel que chaque
$V_i = X_i \times_X F$ soit un espace algébrique. Nous disposons de la
projection $V_i \to X_i$.
Alors $V_i \times_X X_j$ et $X_i \times_X V_j$
représentent tous deux le faisceau $X_i \times_X F \times_X X_j$ ;
nous obtenons donc un isomorphisme
$$
\varphi_{ii'} : V_i \times_X X_j \to X_i \times_X V_j
$$
Il est immédiat que les applications $\varphi_{ij}$
sont des morphismes au-dessus de $X_i \times_X X_j$ et satisfont à la
condition de cocycle. Le foncteur de droite à gauche est donné
par cette construction $F \mapsto (V_i, \varphi_{ij})$.

\medskip\noindent
Construisons maintenant un foncteur de gauche à droite.
Les isomorphismes $\varphi_{ij}$ fournissent des isomorphismes
$$
\varphi_{ij} : V_i \times_X X_j \longrightarrow X_i \times_X V_j
$$
au-dessus de $X_i \times X_j$. Posons $F$ égal au coégalisateur du
diagramme suivant :
$$
\xymatrix{
\coprod_{i, i'} V_i \times_X X_j
\ar@<1ex>[rr]^-{\text{pr}_0}
\ar@<-1ex>[rr]_-{\text{pr}_1 \circ \varphi_{ij}}
& &
\coprod_i V_i \ar[r]
&
F
}
$$
La condition de cocycle assure que $F$ est muni d'un morphisme
$F \to X$ et que $X_i \times_X F$ est isomorphe à $V_i$.
Le foncteur de gauche à droite est donné
par cette construction $(V_i, \varphi_{ij}) \mapsto F$.

\medskip\noindent
Nous omettons de vérifier que ces constructions
sont des foncteurs quasi-inverses l'un de l'autre. Les derniers énoncés
(1) et (2) résultent des constructions.
\end{proof}









% OK, start here.
%

\chapter{Catégories dérivées des espaces}



\phantomsection
\label{spaces-perfect-section-phantom}


\section{Introduction}
\label{spaces-perfect-section-introduction}

\noindent
Dans ce chapitre, nous étudions les catégories dérivées de modules sur les espaces algébriques.
Il ne semble pas exister de bonnes références introductives consacrées à ce sujet;
la littérature le traite en renvoyant à des articles portant sur les catégories dérivées
de modules sur les champs algébriques; voir par exemple
\cite{olsson_sheaves}.



\section{Conventions}
\label{spaces-perfect-section-conventions}

\noindent
Si $\mathcal{A}$ est une catégorie abélienne et $M$ un objet
de $\mathcal{A}$, nous notons également $M$ l'objet de $K(\mathcal{A})$
et/ou de $D(\mathcal{A})$ correspondant au complexe qui a
$M$ en degré $0$ et qui est nul dans tous les autres degrés.

\medskip\noindent
Si $A$ est un anneau, $K(A)$ désigne la catégorie homotopique
des complexes de $A$-modules et $D(A)$ la catégorie dérivée associée.
De même, si $(X, \mathcal{O}_X)$ est un espace annelé, le symbole
$K(\mathcal{O}_X)$ désigne la catégorie homotopique des complexes de
$\mathcal{O}_X$-modules et $D(\mathcal{O}_X)$ la catégorie dérivée
associée.





\section{Généralités}
\label{spaces-perfect-section-generalities}

\noindent
Dans cette section, nous rassemblons quelques résultats généraux sur la cohomologie des
complexes non bornés de modules sur les espaces algébriques.

\begin{lemma}
\label{spaces-perfect-lemma-restrict-direct-image-open}
Soit $S$ un schéma. Soit $f : X \to Y$ un morphisme d'espaces algébriques
sur $S$. Étant donné un morphisme \'etale $V \to Y$, posons $U = V \times_Y X$
et notons $g : U \to V$ le morphisme de projection. Alors
$(Rf_*E)|_V = Rg_*(E|_U)$ pour $E$ dans $D(\mathcal{O}_X)$.
\end{lemma}

\begin{proof}
Représentons $E$ par un complexe K-injectif $\mathcal{I}^\bullet$ de
$\mathcal{O}_X$-modules. Alors $Rf_*(E) = f_*\mathcal{I}^\bullet$
et $Rg_*(E|_U) = g_*(\mathcal{I}^\bullet|_U)$ d'après
Cohomologie sur les sites, lemme
\ref{sites-cohomology-lemma-restrict-K-injective-to-open}.
Le résultat découle donc de
Propriétés des espaces,
lemme \ref{spaces-properties-lemma-pushforward-etale-base-change-modules}.
\end{proof}

\begin{definition}
\label{spaces-perfect-definition-supported-on}
Soit $S$ un schéma. Soit $X$ un espace algébrique sur $S$.
Soit $E$ un objet de $D(\mathcal{O}_X)$.
Soit $T \subset |X|$ un sous-ensemble fermé.
Nous disons que $E$ est {\it supporté sur $T$} si les
faisceaux de cohomologie $H^i(E)$ sont supportés sur $T$.
\end{definition}








\section{Catégorie dérivée des modules quasi-cohérents sur le petit site \'etale}
\label{spaces-perfect-section-derived-quasi-coherent-etale}

\noindent
Soit $X$ un schéma. Dans cette section, nous montrons que
$D_\QCoh(\mathcal{O}_X)$
peut être définie au moyen du petit site \'etale $X_\etale$ de $X$.
Notons $\mathcal{O}_\etale$ le faisceau structural sur
$X_\etale$. 
Considérons le morphisme de sites annelés
\begin{equation}
\label{spaces-perfect-equation-epsilon}
\epsilon :
(X_\etale, \mathcal{O}_\etale)
\longrightarrow
(X_{Zar}, \mathcal{O}_X).
\end{equation}
noté $\text{id}_{small, \etale, Zar}$ dans
Descente, lemme \ref{descent-lemma-compare-sites}.

\begin{lemma}
\label{spaces-perfect-lemma-epsilon-flat}
Le morphisme $\epsilon$ de (\ref{spaces-perfect-equation-epsilon})
est un morphisme plat de sites annelés. En particulier, le foncteur
$\epsilon^* : \textit{Mod}(\mathcal{O}_X) \to
\textit{Mod}(\mathcal{O}_\etale)$ est exact.
De plus, si $\epsilon^*\mathcal{F} = 0$, alors $\mathcal{F} = 0$.
\end{lemma}

\begin{proof}
La platitude du morphisme $\epsilon$ résulte de
Descente, lemme \ref{descent-lemma-compare-etale-zariski-flat}.
Voici une autre démonstration. Il faut montrer que
$\mathcal{O}_\etale$ est un $\epsilon^{-1}\mathcal{O}_X$-module plat.
Pour cela, il suffit de vérifier que
$\mathcal{O}_{X, x} \to \mathcal{O}_{\etale, \overline{x}}$
est plat pour tout point géométrique $\overline{x}$ de $X$, voir
Modules sur les sites, lemme
\ref{sites-modules-lemma-check-flat-stalks},
Sites, lemme
\ref{sites-lemma-point-morphism-sites},
et
Cohomologie étale, remarques
\ref{etale-cohomology-remarks-enough-points}.
D'après Cohomologie étale, lemme
\ref{etale-cohomology-lemma-describe-etale-local-ring}
nous voyons que $\mathcal{O}_{\etale, \overline{x}}$ est la
hensélisation stricte de $\mathcal{O}_{X, x}$. Ainsi
$\mathcal{O}_{X, x} \to \mathcal{O}_{\etale, \overline{x}}$
est fidèlement plat d'après Compléments d'algèbre,
lemme \ref{more-algebra-lemma-dumb-properties-henselization}.

\medskip\noindent
L'exactitude de $\epsilon^*$ découle de la platitude
de $\epsilon$ par
Modules sur les sites, lemme \ref{sites-modules-lemma-flat-pullback-exact}.

\medskip\noindent
Soit $\mathcal{F}$ un $\mathcal{O}_X$-module. Si
$\epsilon^*\mathcal{F} = 0$, alors, avec les notations ci-dessus,
$$
0 = \epsilon^*\mathcal{F}_{\overline{x}} =
\mathcal{F}_x \otimes_{\mathcal{O}_{X, x}} \mathcal{O}_{\etale, \overline{x}}
$$
(Modules sur les sites, lemme \ref{sites-modules-lemma-pullback-stalk})
pour tout point géométrique $\overline{x}$. Par fidélité de la platitude de
$\mathcal{O}_{X, x} \to \mathcal{O}_{\etale, \overline{x}}$
nous concluons que $\mathcal{F}_x = 0$ pour tout $x \in X$.
\end{proof}

\noindent
Soit $X$ un schéma. Avec les notations de (\ref{spaces-perfect-equation-epsilon}),
rappelons que $\epsilon^* : \QCoh(\mathcal{O}_X)
\to \QCoh(\mathcal{O}_\etale)$
est une équivalence d'après
Descente, proposition \ref{descent-proposition-equivalence-quasi-coherent} et
remarque \ref{descent-remark-change-topologies-ringed-sites}.
De plus, $\QCoh(\mathcal{O}_\etale)$ est une
sous-catégorie de Serre de
$\textit{Mod}(\mathcal{O}_\etale)$ d'après
Descente, lemme \ref{descent-lemma-equivalence-quasi-coherent-limits}.
Nous pouvons donc définir $D_\QCoh(\mathcal{O}_\etale)$ comme la sous-catégorie triangulée
de $D(\mathcal{O}_\etale)$ dont les objets sont les
complexes dont les faisceaux de cohomologie sont quasi-cohérents, voir
Catégories dérivées, section \ref{derived-section-triangulated-sub}.
Le foncteur $\epsilon^*$ est exact (lemme \ref{spaces-perfect-lemma-epsilon-flat});
il induit donc
$\epsilon^* :  D(\mathcal{O}_X) \to D(\mathcal{O}_\etale)$
et, comme les images inverses de modules quasi-cohérents sont quasi-cohérentes,
également $\epsilon^* : D_\QCoh(\mathcal{O}_X) \to
D_\QCoh(\mathcal{O}_\etale)$.

\begin{lemma}
\label{spaces-perfect-lemma-derived-quasi-coherent-small-etale-site}
Soit $X$ un schéma. Le foncteur
$\epsilon^* : D_\QCoh(\mathcal{O}_X) \to
D_\QCoh(\mathcal{O}_\etale)$
défini ci-dessus est une équivalence.
\end{lemma}

\begin{proof}
Nous démontrerons ceci en montrant que le foncteur
$R\epsilon_* : D(\mathcal{O}_\etale) \to D(\mathcal{O}_X)$
induit un quasi-inverse. Nous utiliserons librement le fait que $\epsilon_*$
est donné par la restriction à $X_{Zar} \subset X_\etale$ et la description de
$\epsilon^* = \text{id}_{small, \etale, Zar}^*$
dans Descente, lemme \ref{descent-lemma-compare-sites}.

\medskip\noindent
Pour un $\mathcal{O}_X$-module quasi-cohérent $\mathcal{F}$, le morphisme d'adjonction
$\mathcal{F} \to \epsilon_*\epsilon^*\mathcal{F}$ est un isomorphisme, car
$\mathcal{F}^a$
(Descente, définition \ref{descent-definition-structure-sheaf})
est un faisceau, comme démontré dans
Descente, lemme \ref{descent-lemma-sheaf-condition-holds}.
Réciproquement, tout $\mathcal{O}_\etale$-module quasi-cohérent
$\mathcal{H}$ est de la forme $\epsilon^*\mathcal{F}$ pour un certain
$\mathcal{O}_X$-module quasi-cohérent $\mathcal{F}$, voir
Descente, proposition \ref{descent-proposition-equivalence-quasi-coherent}.
Alors $\mathcal{F} = \epsilon_*\mathcal{H}$ d'après ce qui précède et
nous concluons que le morphisme d'adjonction
$\epsilon^*\epsilon_*\mathcal{H} \to \mathcal{H}$ est un isomorphisme pour tout
$\mathcal{O}_\etale$-module quasi-cohérent $\mathcal{H}$.

\medskip\noindent
Soit $E$ un objet de $D_\QCoh(\mathcal{O}_\etale)$
et notons $\mathcal{H}^q = H^q(E)$ son faisceau de cohomologie d'indice $q$.
Soit $\mathcal{B}$ l'ensemble des objets affines de $X_\etale$.
Alors $H^p(U, \mathcal{H}^q) = 0$ pour tout $p > 0$, tout $q \in \mathbf{Z}$,
et tout $U \in \mathcal{B}$, voir
Descente, proposition \ref{descent-proposition-same-cohomology-quasi-coherent}
et
Cohomologie des schémas, lemme
\ref{coherent-lemma-quasi-coherent-affine-cohomology-zero}.
D'après Cohomologie sur les sites, lemme
\ref{sites-cohomology-lemma-cohomology-over-U-trivial}
cela signifie que
$$
H^q(U, E) = H^0(U, \mathcal{H}^q)
$$
pour tout $U \in \mathcal{B}$. En particulier, ceci vaut
pour les ouverts affines $U \subset X$. Il s'ensuit que la cohomologie d'indice $q$ de
$R\epsilon_*E$ au-dessus de $U$ est la valeur du faisceau $\epsilon_*\mathcal{H}^q$
sur $U$. En faisceautisant, nous obtenons
$$
H^q(R\epsilon_*E) = \epsilon_*\mathcal{H}^q
$$
ce qui montre en particulier que $R\epsilon_*$ induit un foncteur
$D_\QCoh(\mathcal{O}_\etale) \to D_\QCoh(\mathcal{O}_X)$.
Comme $\epsilon^*$ est exact, nous obtenons alors
$H^q(\epsilon^*R\epsilon_*E) = \epsilon^*\epsilon_*\mathcal{H}^q =
\mathcal{H}^q$ (par ce qui précède). Ainsi, le morphisme d'adjonction
$\epsilon^*R\epsilon_*E \to E$ est un isomorphisme.

\medskip\noindent
Réciproquement, pour $F \in D_\QCoh(\mathcal{O}_X)$, le
morphisme d'adjonction $F \to R\epsilon_*\epsilon^*F$
est un isomorphisme pour la même raison, à savoir parce que
les faisceaux de cohomologie de $R\epsilon_*\epsilon^*F$
sont isomorphes à
$\epsilon_*H^m(\epsilon^*F) = \epsilon_*\epsilon^*H^m(F) = H^m(F)$.
\end{proof}










\section{Catégorie dérivée des modules quasi-cohérents}
\label{spaces-perfect-section-derived-quasi-coherent}

\noindent
Soit $S$ un schéma. Le lemme
\ref{spaces-perfect-lemma-derived-quasi-coherent-small-etale-site}
montre que la catégorie $D_\QCoh(\mathcal{O}_S)$ peut être définie
au moyen de complexes de $\mathcal{O}_S$-modules sur le schéma $S$
ou de complexes de $\mathcal{O}$-modules sur le petit site \'etale
de $S$. La définition suivante est donc compatible avec celle
du cas des schémas.

\begin{definition}
\label{spaces-perfect-definition-derived-quasi-coherent}
Soit $S$ un schéma. Soit $X$ un espace algébrique sur $S$.
La {\it catégorie dérivée des $\mathcal{O}_X$-modules dont les
faisceaux de cohomologie sont quasi-cohérents} est notée
$D_\QCoh(\mathcal{O}_X)$.
\end{definition}

\noindent
Cette définition a un sens d'après
Propriétés des espaces, lemme
\ref{spaces-properties-lemma-properties-quasi-coherent}
et
Catégories dérivées, section \ref{derived-section-triangulated-sub}.
Nous obtenons ainsi un foncteur canonique
\begin{equation}
\label{spaces-perfect-equation-compare}
D(\QCoh(\mathcal{O}_X))
\longrightarrow
D_\QCoh(\mathcal{O}_X)
\end{equation}
voir Catégories dérivées, équation (\ref{derived-equation-compare}).

\medskip\noindent
Observons qu'un morphisme plat $f : Y \to X$ d'espaces algébriques
induit un foncteur exact
$f^* : \textit{Mod}(\mathcal{O}_X) \to \textit{Mod}(\mathcal{O}_Y)$,
voir
Morphismes des espaces, lemme \ref{spaces-morphisms-lemma-flat-morphism-sites}
et
Modules sur les sites, lemme \ref{sites-modules-lemma-flat-pullback-exact}.
En particulier, $Lf^* : D(\mathcal{O}_X) \to D(\mathcal{O}_Y)$
se calcule sur tout complexe représentant
(Catégories dérivées, lemme \ref{derived-lemma-right-derived-exact-functor}).
Nous écrirons $Lf^* = f^*$ lorsque $f$ est plat; dans ce cas,
$H^i(f^*E) = f^*H^i(E)$ pour $E$ dans $D(\mathcal{O}_X)$.
Nous utiliserons souvent ceci lorsque $f$ est \'etale. Dans ce cas, le
foncteur image inverse est simplement la restriction à $Y_\etale$;
voir Propriétés des espaces, équation
(\ref{spaces-properties-equation-restrict-modules}).

\begin{lemma}
\label{spaces-perfect-lemma-check-quasi-coherence-on-covering}
Soit $S$ un schéma. Soit $X$ un espace algébrique sur $S$.
Soit $E$ un objet de $D(\mathcal{O}_X)$. Les conditions suivantes sont équivalentes:
\begin{enumerate}
\item $E$ est en $D_\QCoh(\mathcal{O}_X)$,
\item pour tout morphisme \'etale $\varphi : U \to X$ où $U$ est un
schéma affine, $\varphi^*E$ est un objet de
$D_\QCoh(\mathcal{O}_U)$,
\item pour tout morphisme \'etale $\varphi : U \to X$ où $U$ est un schéma,
$\varphi^*E$ est un objet de
$D_\QCoh(\mathcal{O}_U)$,
\item il existe un morphisme \'etale surjectif $\varphi : U \to X$
où $U$ est un schéma tel que $\varphi^*E$ soit un objet de
$D_\QCoh(\mathcal{O}_U)$, et
\item il existe un morphisme \'etale surjectif d'espaces algébriques
$f : Y \to X$ tel que $Lf^*E$ soit un objet de
$D_\QCoh(\mathcal{O}_Y)$.
\end{enumerate}
\end{lemma}

\begin{proof}
Ceci découle immédiatement de la discussion qui précède le lemme et de
Propriétés des espaces, lemme
\ref{spaces-properties-lemma-characterize-quasi-coherent}.
\end{proof}

\begin{lemma}
\label{spaces-perfect-lemma-quasi-coherence-direct-sums}
Soit $S$ un schéma. Soit $X$ un espace algébrique sur $S$.
Alors $D_\QCoh(\mathcal{O}_X)$ admet des sommes directes.
\end{lemma}

\begin{proof}
D'après Injectifs, lemme \ref{injectives-lemma-derived-products},
la catégorie dérivée $D(\mathcal{O}_X)$ a des sommes directes et
celles-ci se calculent en prenant les sommes directes terme à terme de représentants quelconques.
Ainsi, il est clair que le faisceau de cohomologie d'une somme directe est le
somme directe des faisceaux de cohomologie, car prendre des sommes directes est
un foncteur exact (dans toute catégorie abélienne de Grothendieck).
Le lemme découle alors du fait que la somme directe de faisceaux quasi-cohérents est quasi-cohérente, voir
Propriétés des espaces, lemme
\ref{spaces-properties-lemma-properties-quasi-coherent}.
\end{proof}

\noindent
Nous aurons besoin de quelques informations sur les limites dérivées.
Nous avertissons le lecteur que, dans le lemme ci-dessous, la limite dérivée n'est en général pas
un objet de $D_\QCoh$.

\begin{lemma}
\label{spaces-perfect-lemma-Rlim-quasi-coherent}
Soit $S$ un schéma. Soit $X$ un espace algébrique sur $S$.
Soit $(K_n)$ un système inverse d'objets de
$D_\QCoh(\mathcal{O}_X)$, de limite dérivée
$K = R\lim K_n$ dans $D(\mathcal{O}_X)$. Supposons que $H^q(K_{n + 1}) \to H^q(K_n)$
soit surjectif pour tout $q \in \mathbf{Z}$ et tout $n \geq 1$.
Alors
\begin{enumerate}
\item $H^q(K) = \lim H^q(K_n)$,
\item $R\lim H^q(K_n) = \lim H^q(K_n)$, et
\item pour tout ouvert affine $U \subset X$ nous avons
$H^p(U, \lim H^q(K_n)) = 0$ pour $p > 0$.
\end{enumerate}
\end{lemma}

\begin{proof}
Soit $\mathcal{B} \subset \Ob(X_\etale)$ l'ensemble des objets affines.
Comme $H^q(K_n)$ est quasi-cohérent, on a $H^p(U, H^q(K_n)) = 0$
pour $U \in \mathcal{B}$, d'après la discussion dans
Cohomologie des espaces, section
\ref{spaces-cohomology-section-higher-direct-image}
et
Cohomologie des schémas, lemme
\ref{coherent-lemma-quasi-coherent-affine-cohomology-zero}.
De plus, les applications $H^0(U, H^q(K_{n + 1})) \to H^0(U, H^q(K_n))$
sont surjectives pour $U \in \mathcal{B}$, par un raisonnement analogue.
La partie (1) découle de Cohomologie sur les sites, lemme
\ref{sites-cohomology-lemma-derived-limit-suitable-system}
dont nous venons de vérifier les conditions.
Les parties (2) et (3) découlent de
Cohomologie sur les sites, lemme
\ref{sites-cohomology-lemma-inverse-limit-is-derived-limit}.
\end{proof}

\begin{lemma}
\label{spaces-perfect-lemma-quasi-coherence-pullback}
Soit $S$ un schéma.
Soit $f : Y \to X$ un morphisme d'espaces algébriques sur $S$.
Le foncteur $Lf^*$ envoie $D_\QCoh(\mathcal{O}_X)$
dans $D_\QCoh(\mathcal{O}_Y)$.
\end{lemma}

\begin{proof}
Choisissons un diagramme
$$
\xymatrix{
U \ar[d]_a \ar[r]_h & V \ar[d]^b \\
X \ar[r]^f & Y
}
$$
où $U$ et $V$ sont des schémas, les flèches verticales sont \'etale, et
$a$ est surjectif. Comme $a^* \circ Lf^* = Lh^* \circ b^*$, le résultat
découle du
lemme \ref{spaces-perfect-lemma-check-quasi-coherence-on-covering}
et du cas des schémas, qui est
Catégories dérivées des schémas, lemme
\ref{perfect-lemma-quasi-coherence-pullback}.
\end{proof}

\begin{lemma}
\label{spaces-perfect-lemma-quasi-coherence-tensor-product}
Soit $S$ un schéma. Soit $X$ un espace algébrique sur $S$.
Pour $K, L$ objets de $D_\QCoh(\mathcal{O}_X)$,
le produit tensoriel dérivé $K \otimes^\mathbf{L} L$ appartient à
$D_\QCoh(\mathcal{O}_X)$.
\end{lemma}

\begin{proof}
Soit $\varphi : U \to X$ un morphisme \'etale surjectif provenant d'un schéma $U$.
Comme
$\varphi^*(K \otimes_{\mathcal{O}_X}^\mathbf{L} L) =
\varphi^*K \otimes_{\mathcal{O}_U}^\mathbf{L} \varphi^*L$
nous voyons par le
lemme \ref{spaces-perfect-lemma-check-quasi-coherence-on-covering}
que ceci découle du cas des schémas, qui est
Catégories dérivées des schémas, lemme
\ref{perfect-lemma-quasi-coherence-tensor-product}.
\end{proof}

\noindent
Le lemme suivant nous aidera à « calculer » un foncteur dérivé à droite
sur un objet de $D_\QCoh(\mathcal{O}_X)$.

\begin{lemma}
\label{spaces-perfect-lemma-nice-K-injective}
Soit $S$ un schéma. Soit $X$ un espace algébrique sur $S$. Soit $E$ un
objet de $D_\QCoh(\mathcal{O}_X)$. Alors le morphisme
$E \to R\lim \tau_{\geq -n}E$ de
Catégories dérivées, remarque
\ref{derived-remark-map-into-derived-limit-truncations}
est un isomorphisme\footnote{En particulier,
$E$ admet un représentant K-injectif, voir
Catégories dérivées, lemme \ref{derived-lemma-difficulty-K-injectives}.}.
\end{lemma}

\begin{proof}
Notons $\mathcal{H}^i = H^i(E)$ le faisceau de cohomologie d'indice $i$ de $E$.
Soit $\mathcal{B}$ l'ensemble des objets affines de $X_\etale$.
Alors $H^p(U, \mathcal{H}^i) = 0$ pour tout $p > 0$, tout $i \in \mathbf{Z}$,
et tout $U \in \mathcal{B}$, puisque $U$ est un schéma affine.
Voir la discussion dans
Cohomologie des espaces, section
\ref{spaces-cohomology-section-higher-direct-image}
et
Cohomologie des schémas, lemme
\ref{coherent-lemma-quasi-coherent-affine-cohomology-zero}.
Ainsi, le lemme découle de
Cohomologie sur les sites, lemme \ref{sites-cohomology-lemma-is-limit-dimension}
avec $d = 0$.
\end{proof}

\begin{lemma}
\label{spaces-perfect-lemma-application-nice-K-injective}
Soit $S$ un schéma. Soit $X$ un espace algébrique sur $S$.
Soit $F : \textit{Mod}(\mathcal{O}_X) \to \textit{Ab}$
un foncteur et $N \geq 0$ un entier. Supposons que
\begin{enumerate}
\item $F$ est exact à gauche,
\item $F$ commute aux produits directs dénombrables,
\item $R^pF(\mathcal{F}) = 0$ pour tout $p \geq N$ et tout $\mathcal{F}$
quasi-cohérent.
\end{enumerate}
Alors, pour $E \in D_\QCoh(\mathcal{O}_X)$,
\begin{enumerate}
\item $H^i(RF(\tau_{\leq a}E) \to H^i(RF(E))$ est un isomorphisme
pour $i \leq a$,
\item $H^i(RF(E)) \to H^i(RF(\tau_{\geq b - N + 1}E))$ est un isomorphisme
pour $i \geq b$,
\item si $H^i(E) = 0$ pour $i \not \in [a, b]$ pour certains
$-\infty \leq a \leq b \leq \infty$, alors $H^i(RF(E)) = 0$
pour $i \not \in [a, b + N - 1]$.
\end{enumerate}
\end{lemma}

\begin{proof}
L'énoncé (1) est
Catégories dérivées, lemme \ref{derived-lemma-negative-vanishing}.

\medskip\noindent
Démonstration de (2). Écrivons $E_n = \tau_{\geq -n}E$. On a
$E = R\lim E_n$, voir le lemme \ref{spaces-perfect-lemma-nice-K-injective}. Ainsi,
$RF(E) = R\lim RF(E_n)$ dans $D(\textit{Ab})$ d'après Injectifs, lemme
\ref{injectives-lemma-RF-commutes-with-Rlim}. Ainsi, pour tout $i \in \mathbf{Z}$,
on a une suite exacte courte
$$
0 \to R^1\lim H^{i - 1}(RF(E_n)) \to H^i(RF(E)) \to \lim H^i(RF(E_n)) \to 0
$$
voir Compléments d'algèbre, remarque
\ref{more-algebra-remark-compare-derived-limit}.
Pour démontrer (2), nous montrerons que le terme de gauche est nul
et que le terme de droite est égal à $H^i(RF(E_{-b + N - 1})$
pour tout $b$ tel que $i \geq b$.

\medskip\noindent
Pour tout $n$, on a un triangle distingué
$$
H^{-n}(E)[n] \to E_n \to E_{n - 1} \to H^{-n}(E)[n + 1]
$$
(Catégories dérivées, remarque
\ref{derived-remark-truncation-distinguished-triangle})
dans $D(\mathcal{O}_X)$. Comme $H^{-n}(E)$ est quasi-cohérent, on a
$$
H^i(RF(H^{-n}(E)[n])) = R^{i + n}F(H^{-n}(E)) = 0
$$
pour $i + n \geq N$, et
$$
H^i(RF(H^{-n}(E)[n + 1])) = R^{i + n + 1}F(H^{-n}(E)) = 0
$$
pour $i + n + 1 \geq N$. Nous en déduisons que
$$
H^i(RF(E_n)) \to H^i(RF(E_{n - 1}))
$$
est un isomorphisme pour $n \geq N - i$. Ainsi, les systèmes $H^i(RF(E_n))$
satisfont tous la condition ML et le terme $R^1\lim$ de notre suite exacte courte
est nul (voir la discussion dans
Compléments d'algèbre, section \ref{more-algebra-section-Rlim}).
De plus, le système $H^i(RF(E_n))$ est constant à partir
de $n = N - i - 1$, comme souhaité.

\medskip\noindent
Démonstration de (3). Sous l'hypothèse sur $E$, on a
$\tau_{\leq a - 1}E = 0$ et l'on obtient l'annulation
de $H^i(RF(E))$ pour $i \leq a - 1$ par (1).
De même, $\tau_{\geq b + 1}E = 0$ et l'on obtient donc
l'annulation de $H^i(RF(E))$ pour $i \geq b + n$ d'après
la partie (2).
\end{proof}









\section{Image directe totale}
\label{spaces-perfect-section-total-direct-image}

\noindent
Le lemme suivant est l'analogue du
lemme de Cohomologie des espaces
\ref{spaces-cohomology-lemma-vanishing-higher-direct-images}.

\begin{lemma}
\label{spaces-perfect-lemma-quasi-coherence-direct-image}
Soit $S$ un schéma. Soit $f : X \to Y$ un morphisme quasi-séparé et quasi-compact
d'espaces algébriques sur $S$.
\begin{enumerate}
\item Le foncteur $Rf_*$ envoie $D_\QCoh(\mathcal{O}_X)$
dans $D_\QCoh(\mathcal{O}_Y)$.
\item Si $Y$ est quasi-compact, il existe un entier $N = N(X, Y, f)$
tel que, pour un objet $E$ de $D_\QCoh(\mathcal{O}_X)$
dont $H^m(E) = 0$ pour $m > 0$, on ait
$H^m(Rf_*E) = 0$ pour $m \geq N$.
\item En fait, si $Y$ est quasi-compact, on peut trouver $N = N(X, Y, f)$
tel que, pour tout morphisme d'espaces algébriques $Y' \to Y$,
la même conclusion vaille pour le foncteur $R(f')_*$,
où $f' : X' \to Y'$ est le changement de base de $f$.
\end{enumerate}
\end{lemma}

\begin{proof}
Soit $E$ un objet de $D_\QCoh(\mathcal{O}_X)$.
Pour démontrer (1), il faut montrer que $Rf_*E$ a des
faisceaux de cohomologie quasi-cohérents. Cette question est locale sur $Y$; on peut donc
supposer $Y$ quasi-compact. Choisissons $N = N(X, Y, f)$ comme dans
Cohomologie des espaces, lemme
\ref{spaces-cohomology-lemma-vanishing-higher-direct-images}.
Ainsi $R^pf_*\mathcal{F} = 0$ pour tout module quasi-cohérent $\mathcal{O}_X$-
module $\mathcal{F}$ et tout $p \geq N$. De plus, $R^pf_*\mathcal{F}$
est quasi-cohérent pour tout $p$, d'après
Cohomologie des espaces, lemme \ref{spaces-cohomology-lemma-higher-direct-image}.
Ces énoncés restent vrais après changement de base.

\medskip\noindent
Commençons par supposer que $E$ est borné inférieurement. Nous montrerons que (1), (2) et (3) sont valables
pour un tel $E$ avec notre choix de $N$. Dans ce cas, on peut par exemple utiliser la
suite spectrale
$$
R^pf_*H^q(E) \Rightarrow R^{p + q}f_*E
$$
(Catégories dérivées, lemme \ref{derived-lemma-two-ss-complex-functor}),
la quasi-cohérence de $R^pf_*H^q(E)$ et l'annulation de $R^pf_*H^q(E)$
pour $p \geq N$ pour voir que (1), (2) et (3) sont valables dans ce cas.

\medskip\noindent
Prouvons maintenant (2) et (3). Supposons que $H^m(E) = 0$ pour $m > 0$.
Soit $V$ un objet affine de $Y_\etale$.
On a $H^p(V \times_Y X, \mathcal{F}) = 0$ pour $p \geq N$, voir
Cohomologie des espaces, lemme
\ref{spaces-cohomology-lemma-quasi-coherence-higher-direct-images-application}.
Nous pouvons donc appliquer le lemme \ref{spaces-perfect-lemma-application-nice-K-injective}
au foncteur $\Gamma(V \times_Y X, -)$ pour obtenir
$$
R\Gamma(V, Rf_*E) = R\Gamma(V \times_Y X, E)
$$
dont la cohomologie est nulle en degrés $\geq N$. Comme ceci vaut pour
tout $V$ affine de $Y_\etale$, on conclut que $H^m(Rf_*E) = 0$
pour $m \geq N$.

\medskip\noindent
Passons à la démonstration de (1) dans le cas général. Rappelons qu'il existe un
triangle distingué
$$
\tau_{\leq -n - 1}E \to E \to \tau_{\geq -n}E \to
(\tau_{\leq -n - 1}E)[1]
$$
dans $D(\mathcal{O}_X)$, voir Catégories dérivées, remarque
\ref{derived-remark-truncation-distinguished-triangle}.
D'après (2), $Rf_*\tau_{\leq -n - 1}E$
a des faisceaux de cohomologie nuls en degrés $\geq -n + N$.
Ainsi, pour un entier $q$ donné, $R^qf_*E$ est égal
à $R^qf_*\tau_{\geq -n}E$ pour un certain $n$, et le résultat
ci-dessus s'applique.
\end{proof}

\begin{lemma}
\label{spaces-perfect-lemma-quasi-coherence-pushforward-direct-sums}
Soit $S$ un schéma. Soit $f : X \to Y$ un morphisme quasi-séparé et
quasi-compact d'espaces algébriques sur $S$. Alors
$Rf_* : D_\QCoh(\mathcal{O}_X) \to D_\QCoh(\mathcal{O}_Y)$
commute aux sommes directes.
\end{lemma}

\begin{proof}
Soit $E_i$ une famille d'objets de $D_\QCoh(\mathcal{O}_X)$
et posons $E = \bigoplus E_i$. Nous voulons montrer que le morphisme
$$
\bigoplus Rf_*E_i \longrightarrow Rf_*E
$$
est un isomorphisme. Nous montrerons qu'il induit un isomorphisme sur
les faisceaux de cohomologie en degré $0$, ce qui démontrera le lemme.
Choisissons un entier $N$ comme dans le lemme \ref{spaces-perfect-lemma-quasi-coherence-direct-image}.
Alors $R^0f_*E = R^0f_*\tau_{\geq -N}E$ et
$R^0f_*E_i = R^0f_*\tau_{\geq -N}E_i$ d'après le lemme cité. Observons que
$\tau_{\geq -N}E = \bigoplus \tau_{\geq -N}E_i$.
Nous pouvons donc supposer que tous les $E_i$ ont des faisceaux de cohomologie nuls
en degrés $< -N$. Nous utilisons ensuite les suites spectrales
$$
R^pf_*H^q(E) \Rightarrow R^{p + q}f_*E
\quad\text{et}\quad
R^pf_*H^q(E_i) \Rightarrow R^{p + q}f_*E_i
$$
(Catégories dérivées, lemme \ref{derived-lemma-two-ss-complex-functor})
pour nous ramener au cas d'une somme directe de faisceaux quasi-cohérents.
Ce cas est traité par
Cohomologie des espaces, lemme \ref{spaces-cohomology-lemma-colimit-cohomology}.
\end{proof}

\begin{remark}
\label{spaces-perfect-remark-match-total-direct-images}
Soit $S$ un schéma. Soit $f : X \to Y$ un morphisme d'espaces algébriques
représentables $X$ et $Y$ sur $S$. Soit $f_0 : X_0 \to Y_0$ un
morphisme de schémas représentant $f$ (notation maladroite mais temporaire).
Alors le diagramme
$$
\xymatrix{
D_\QCoh(\mathcal{O}_{X_0})
\ar@{=}[rrrrrr]_{\text{lemme
\ref{spaces-perfect-lemma-derived-quasi-coherent-small-etale-site}}}
& & & & & &
D_\QCoh(\mathcal{O}_X) \\
D_\QCoh(\mathcal{O}_{Y_0})
\ar[u]^{Lf^*_0}
\ar@{=}[rrrrrr]^{\text{lemme
\ref{spaces-perfect-lemma-derived-quasi-coherent-small-etale-site}}}
& & & & & &
D_\QCoh(\mathcal{O}_Y) \ar[u]_{Lf^*}
}
$$
(lemme \ref{spaces-perfect-lemma-quasi-coherence-pullback} et
Catégories dérivées des schémas, lemme
\ref{perfect-lemma-quasi-coherence-pullback})
est commutatif. En effet, les
équivalences
$D_\QCoh(\mathcal{O}_{X_0}) \to D_\QCoh(\mathcal{O}_X)$
et
$D_\QCoh(\mathcal{O}_{Y_0}) \to D_\QCoh(\mathcal{O}_Y)$
du lemme \ref{spaces-perfect-lemma-derived-quasi-coherent-small-etale-site}
proviennent de l'image inverse par les morphismes (plats) de sites annelés
$\epsilon : X_\etale \to X_{0, Zar}$ et
$\epsilon : Y_\etale \to Y_{0, Zar}$
et le diagramme de sites annelés
$$
\xymatrix{
X_{0, Zar} \ar[d]_{f_0} & X_\etale \ar[l]^\epsilon \ar[d]^f \\
Y_{0, Zar} & Y_\etale \ar[l]_\epsilon
}
$$
est commutatif (détails omis). Si $f$ est quasi-compact et
quasi-séparé, ou de manière équivalente si $f_0$ est quasi-compact et
quasi-séparé, alors on affirme que
$$
\xymatrix{
D_\QCoh(\mathcal{O}_{X_0})
\ar[d]_{Rf_{0, *}} \ar@{=}[rrrrrr]_{\text{lemme
\ref{spaces-perfect-lemma-derived-quasi-coherent-small-etale-site}}}
& & & & & &
D_\QCoh(\mathcal{O}_X) \ar[d]^{Rf_*} \\
D_\QCoh(\mathcal{O}_{Y_0})
\ar@{=}[rrrrrr]^{\text{lemme
\ref{spaces-perfect-lemma-derived-quasi-coherent-small-etale-site}}}
& & & & & &
D_\QCoh(\mathcal{O}_Y)
}
$$
(lemme \ref{spaces-perfect-lemma-quasi-coherence-direct-image} et
Catégories dérivées des schémas, lemme
\ref{perfect-lemma-quasi-coherence-direct-image})
est également commutatif. Cela découle encore du
diagramme de sites affiché ci-dessus, car la démonstration du lemme
\ref{spaces-perfect-lemma-derived-quasi-coherent-small-etale-site}
montre que le foncteur $R\epsilon_*$ fournit les équivalences
$D_\QCoh(\mathcal{O}_X) \to D_\QCoh(\mathcal{O}_{X_0})$
et
$D_\QCoh(\mathcal{O}_Y) \to D_\QCoh(\mathcal{O}_{Y_0})$.
\end{remark}

\begin{lemma}
\label{spaces-perfect-lemma-affine-morphism}
Soit $S$ un schéma. Soit $f : X \to Y$ un morphisme affine d'espaces algébriques
sur $S$. Alors
$Rf_* : D_\QCoh(\mathcal{O}_X) \to D_\QCoh(\mathcal{O}_Y)$
réfléchit les isomorphismes.
\end{lemma}

\begin{proof}
L'énoncé signifie qu'un morphisme $\alpha : E \to F$ de
$D_\QCoh(\mathcal{O}_X)$ est un isomorphisme si
$Rf_*\alpha$ est un isomorphisme. On peut le vérifier sur les faisceaux de cohomologie.
En particulier, la question est locale pour la topologie \'etale sur $Y$. On peut donc supposer
$Y$ et donc $X$ affines. Dans ce cas, le problème se ramène au
cas des schémas
(Catégories dérivées des schémas, lemme \ref{perfect-lemma-affine-morphism})
au moyen du lemme \ref{spaces-perfect-lemma-derived-quasi-coherent-small-etale-site} et de la
Remarque \ref{spaces-perfect-remark-match-total-direct-images}.
\end{proof}

\begin{lemma}
\label{spaces-perfect-lemma-affine-morphism-pull-push}
Soit $S$ un schéma. Soit $f : X \to Y$ un morphisme affine d'espaces algébriques
sur $S$. Pour $E$ dans $D_\QCoh(\mathcal{O}_Y)$, on a
$Rf_* Lf^* E = E \otimes^\mathbf{L}_{\mathcal{O}_Y} f_*\mathcal{O}_X$.
\end{lemma}

\begin{proof}
Comme $f$ est affine, le morphisme $f_*\mathcal{O}_X \to Rf_*\mathcal{O}_X$
est un isomorphisme (Cohomologie des espaces, lemme
\ref{spaces-cohomology-lemma-affine-vanishing-higher-direct-images}).
Il existe un morphisme canonique
$E \otimes^\mathbf{L} f_*\mathcal{O}_X =
E \otimes^\mathbf{L} Rf_*\mathcal{O}_X \to Rf_* Lf^* E$
adjoint au morphisme
$$
Lf^*(E \otimes^\mathbf{L} Rf_*\mathcal{O}_X) =
Lf^*E \otimes^\mathbf{L} Lf^*Rf_*\mathcal{O}_X \longrightarrow
Lf^* E \otimes^\mathbf{L} \mathcal{O}_X = Lf^* E
$$
provenant de $1 : Lf^*E \to Lf^*E$ et du morphisme canonique
$Lf^*Rf_*\mathcal{O}_X \to \mathcal{O}_X$. Pour vérifier que le morphisme ainsi construit
est un isomorphisme, on peut travailler localement sur $Y$. On peut donc supposer
$Y$ et donc $X$ affines. Dans ce cas, le problème se ramène au
cas des schémas
(Catégories dérivées des schémas, lemme
\ref{perfect-lemma-affine-morphism-pull-push})
au moyen du lemme \ref{spaces-perfect-lemma-derived-quasi-coherent-small-etale-site} et de la
Remarque \ref{spaces-perfect-remark-match-total-direct-images}.
\end{proof}








\section{Être propre au-dessus d'une base}
\label{spaces-perfect-section-proper-over-base}

\noindent
Cette section est l'analogue de la section de Cohomologie des schémas
\ref{coherent-section-proper-over-base}.
Comme toujours pour les résultats concernant la topologie des ensembles de points,
il faut être prudent en transposant le matériel aux espaces algébriques.

\begin{lemma}
\label{spaces-perfect-lemma-closed-proper-over-base}
Soit $S$ un schéma. Soit $f : X \to Y$ un morphisme d'espaces algébriques
sur $S$, localement de type fini. Soit $T \subset |X|$ un sous-ensemble fermé.
Les conditions suivantes sont équivalentes
\begin{enumerate}
\item le morphisme $Z \to Y$ est propre si $Z$ est la structure d'espace algébrique réduite
induite sur $T$
(Propriétés des espaces, définition
\ref{spaces-properties-definition-reduced-induced-space}),
\item pour un sous-espace fermé $Z \subset X$ tel que $|Z| = T$,
le morphisme $Z \to Y$ est propre, et
\item pour tout sous-espace fermé $Z \subset X$ tel que $|Z| = T$, le morphisme
$Z \to Y$ est propre.
\end{enumerate}
\end{lemma}

\begin{proof}
Les implications (3) $\Rightarrow$ (1) et (1) $\Rightarrow$ (2)
sont immédiates. Il suffit donc de démontrer que (2) implique (3).
Nous invitons le lecteur à trouver sa propre démonstration de ce fait.
Soient $Z'$ et $Z''$ des sous-espaces fermés tels que $T = |Z'| = |Z''|$
et tels que $Z' \to Y$ soit un morphisme propre d'espaces algébriques.
Nous devons montrer que $Z'' \to Y$ est également propre.
Soit $Z''' = Z' \cup Z''$ l'union schématique, voir
Morphismes des espaces, définition
\ref{spaces-morphisms-definition-scheme-theoretic-intersection-union}.
Alors $Z'''$ est un autre sous-espace fermé tel que $|Z'''| = T$.
Cela découle par exemple de la description des unions schématiques dans
Morphismes des espaces, lemme \ref{spaces-morphisms-lemma-scheme-theoretic-union}.
Comme $Z'' \to Z'''$ est une immersion fermée, il suffit de démontrer
que $Z''' \to Y$ est propre (voir
Morphismes des espaces, lemmes
\ref{spaces-morphisms-lemma-closed-immersion-proper} et
\ref{spaces-morphisms-lemma-composition-proper}).
Le morphisme $Z' \to Z'''$ est une immersion fermée bijective
et donc, en particulier, surjectif et universellement fermé.
Le fait que $Z' \to Y$ soit séparé implique alors que
$Z''' \to Y$ est séparé, voir
Morphismes des espaces, lemme
\ref{spaces-morphisms-lemma-image-universally-closed-separated}.
De plus, $Z''' \to Y$ est localement de type fini
puisque $X \to Y$ est localement de type fini
(Morphismes des espaces, lemmes
\ref{spaces-morphisms-lemma-immersion-locally-finite-type} et
\ref{spaces-morphisms-lemma-composition-finite-type}).
Comme $Z' \to Y$ est quasi-compact et $Z' \to Z'''$ est un
homéomorphisme universel, $Z''' \to Y$ est quasi-compact.
Enfin, puisque $Z' \to Y$ est universellement fermé,
il en va de même pour $Z''' \to Y$ d'après
Morphismes des espaces, lemme \ref{spaces-morphisms-lemma-image-proper-is-proper}.
Ceci achève la démonstration.
\end{proof}

\begin{definition}
\label{spaces-perfect-definition-proper-over-base}
Soit $S$ un schéma.
Soit $f : X \to Y$ un morphisme d'espaces algébriques sur $S$
localement de type fini.
Soit $T \subset |X|$ un sous-ensemble fermé.
Nous disons que {\it $T$ est propre au-dessus de $Y$}
si les conditions équivalentes du lemme \ref{spaces-perfect-lemma-closed-proper-over-base}
sont satisfaites.
\end{definition}

\noindent
Le lemme utilisé dans la définition ci-dessus est faux si le morphisme
$f : X \to Y$ n'est pas localement de type fini. Nous invitons donc
le lecteur à ne pas utiliser cette terminologie si $f$ n'est pas de
type fini localement.

\begin{lemma}
\label{spaces-perfect-lemma-closed-closed-proper-over-base}
Soit $S$ un schéma.
Soit $f : X \to Y$ un morphisme d'espaces algébriques sur $S$
localement de type fini.
Soit $T' \subset T \subset |X|$ des sous-ensembles fermés.
Si $T$ est propre au-dessus de $Y$, il en va de même pour $T'$.
\end{lemma}

\begin{proof}
Omis.
\end{proof}

\begin{lemma}
\label{spaces-perfect-lemma-base-change-closed-proper-over-base}
Soit $S$ un schéma.
Considérons un diagramme cartésien d'espaces algébriques sur $S$
$$
\xymatrix{
X' \ar[d]_{f'} \ar[r]_{g'} & X \ar[d]^f \\
Y' \ar[r]^g & Y
}
$$
tel que $f$ soit localement de type fini.
Si $T$ est un sous-ensemble fermé de $|X|$ propre au-dessus de $Y$, alors
$|g'|^{-1}(T)$ est un sous-ensemble fermé de $|X'|$ propre au-dessus de $Y'$.
\end{lemma}

\begin{proof}
Observons que l'énoncé a un sens puisque $f'$ est localement de
type fini, d'après Morphismes des espaces, lemme
\ref{spaces-morphisms-lemma-base-change-finite-type}.
Soit $Z \subset X$ la structure de sous-espace fermé réduite induite sur $T$.
Notons $Z' = (g')^{-1}(Z)$ l'image inverse schématique.
Alors $Z' = X' \times_X Z = (Y' \times_Y X) \times_X Z = Y' \times_Y Z$
est propre au-dessus de $Y'$ comme changement de base de $Z$ au-dessus de $Y$
(Morphismes des espaces, lemme \ref{spaces-morphisms-lemma-base-change-proper}).
D'autre part, $T' = |Z'|$. Le lemme en découle.
\end{proof}

\begin{lemma}
\label{spaces-perfect-lemma-functoriality-closed-proper-over-base}
Soit $S$ un schéma. Soit $B$ un espace algébrique sur $S$.
Soit $f : X \to Y$ un morphisme d'espaces algébriques qui
sont localement de type fini sur $B$.
\begin{enumerate}
\item Si $Y$ est séparé sur $B$ et si $T \subset |X|$ est un sous-ensemble fermé
propre au-dessus de $B$, alors $|f|(T)$ est un sous-ensemble fermé de $|Y|$ propre au-dessus de $B$.
\item Si $f$ est universellement fermé et si $T \subset |X|$ est un
sous-ensemble fermé propre au-dessus de $B$, alors $|f|(T)$ est un sous-ensemble fermé
de $Y$ propre au-dessus de $B$.
\item Si $f$ est propre et si $T \subset |Y|$ est un sous-ensemble fermé
propre au-dessus de $B$, alors $|f|^{-1}(T)$ est un sous-ensemble fermé de $|X|$
propre au-dessus de $B$.
\end{enumerate}
\end{lemma}

\begin{proof}
Démonstration de (1). Supposons que $Y$ soit séparé sur $B$ et que $T \subset |X|$
soit un sous-ensemble fermé propre au-dessus de $B$. Soit $Z$ la structure réduite induite
de sous-espace fermé sur $T$ et appliquons
Morphismes des espaces, lemme
\ref{spaces-morphisms-lemma-scheme-theoretic-image-is-proper}
à $Z \to Y$ au-dessus de $B$ pour conclure.

\medskip\noindent
Démonstration de (2). Supposons que $f$ soit universellement fermé et que $T \subset |X|$ soit un
sous-ensemble fermé propre au-dessus de $B$. Soit $Z$ la structure réduite induite
de sous-espace fermé sur $T$ et soit $Z'$ la structure réduite
de sous-espace fermé induite sur $|f|(T)$. Nous obtenons un morphisme
induit $Z \to Z'$.
Notons $Z'' = f^{-1}(Z')$ l'image inverse schématique.
Alors $Z'' \to Z'$ est universellement fermé comme changement de base de $f$
(Morphismes des espaces, lemme \ref{spaces-morphisms-lemma-base-change-proper}).
Par conséquent, $Z \to Z'$ est universellement fermé comme composée de
l'immersion fermée $Z \to Z''$ et de $Z'' \to Z'$
(Morphismes des espaces, lemmes
\ref{spaces-morphisms-lemma-closed-immersion-proper} et
\ref{spaces-morphisms-lemma-composition-proper}).
Nous concluons que $Z' \to B$ est séparé d'après
Morphismes des espaces, lemme
\ref{spaces-morphisms-lemma-image-universally-closed-separated}.
Comme $Z \to B$ est quasi-compact et $Z \to Z'$ surjectif,
on voit que $Z' \to B$ est quasi-compact.
Comme $Z' \to B$ est la composée de $Z' \to Y$ et de $Y \to B$,
on voit que $Z' \to B$ est localement de type fini
(Morphismes des espaces, lemmes
\ref{spaces-morphisms-lemma-immersion-locally-finite-type} et
\ref{spaces-morphisms-lemma-composition-finite-type}).
Enfin, puisque $Z \to B$ est universellement fermé,
il en va de même pour $Z' \to B$ d'après
Morphismes des espaces, lemme
\ref{spaces-morphisms-lemma-image-proper-is-proper}.
Ceci achève la démonstration.

\medskip\noindent
Démonstration de (3). Supposons $f$ propre et $T \subset |Y|$ un sous-ensemble fermé
propre au-dessus de $B$. Soit $Z$ la structure de sous-espace fermé réduite induite
sur $T$. Notons $Z' = f^{-1}(Z)$ l'image inverse schématique.
Alors $Z' \to Z$ est propre comme changement de base de $f$
(Morphismes des espaces, lemme \ref{spaces-morphisms-lemma-base-change-proper}).
Par conséquent, $Z' \to B$ est propre comme composée de $Z' \to Z$
et de $Z \to B$
(Morphismes des espaces, lemme \ref{spaces-morphisms-lemma-composition-proper}).
Ceci achève la démonstration.
\end{proof}

\begin{lemma}
\label{spaces-perfect-lemma-union-closed-proper-over-base}
Soit $S$ un schéma.
Soit $f : X \to Y$ un morphisme d'espaces algébriques sur $S$
localement de type fini.
Soient $T_i \subset |X|$, $i = 1, \ldots, n$, des sous-ensembles fermés.
Si les $T_i$, $i = 1, \ldots, n$, sont propres au-dessus de $Y$, il en va de même
pour $T_1 \cup \ldots \cup T_n$.
\end{lemma}

\begin{proof}
Soit $Z_i$ la structure de sous-schéma fermé réduite induite sur $T_i$.
Le morphisme
$$
Z_1 \amalg \ldots \amalg Z_n \longrightarrow X
$$
est fini d'après Morphismes des espaces, lemmes
\ref{spaces-morphisms-lemma-closed-immersion-finite} et
\ref{spaces-morphisms-lemma-finite-union-finite}.
Comme les morphismes finis sont universellement fermés
(Morphismes des espaces, lemme \ref{spaces-morphisms-lemma-finite-proper})
et que $Z_1 \amalg \ldots \amalg Z_n$ est propre au-dessus de $S$,
nous concluons par la partie (2) du lemme
\ref{spaces-perfect-lemma-functoriality-closed-proper-over-base}
que l'image $Z_1 \cup \ldots \cup Z_n$ est propre au-dessus de $S$.
\end{proof}

\noindent
Soit $S$ un schéma.
Soit $f : X \to Y$ un morphisme d'espaces algébriques au-dessus de $S$
localement de type fini. Soit $\mathcal{F}$ un
$\mathcal{O}_X$-module de type fini et quasi-cohérent. Alors le support
$\text{Supp}(\mathcal{F})$ de $\mathcal{F}$ est un sous-ensemble fermé de $|X|$, voir
Morphismes des espaces, lemme \ref{spaces-morphisms-lemma-support-finite-type}.
Il est donc sensé de dire que
« le support de $\mathcal{F}$ est propre au-dessus de $Y$ ».

\begin{lemma}
\label{spaces-perfect-lemma-module-support-proper-over-base}
Soit $S$ un schéma. Soit $f : X \to Y$ un morphisme d'espaces
algébriques au-dessus de $S$, localement de type fini.
Soit $\mathcal{F}$ un $\mathcal{O}_X$-module de type fini et
quasi-cohérent. Les conditions suivantes sont équivalentes :
\begin{enumerate}
\item le support de $\mathcal{F}$ est propre au-dessus de $Y$ ;
\item le support schématique de $\mathcal{F}$
(Morphismes des espaces, définition
\ref{spaces-morphisms-definition-scheme-theoretic-support})
est propre au-dessus de $Y$ ;
\item il existe un sous-espace fermé $Z \subset X$ et un
$\mathcal{O}_Z$-module $\mathcal{G}$ de type fini et quasi-cohérent
tel que (a) $Z \to Y$ soit propre et (b)
$(Z \to X)_*\mathcal{G} = \mathcal{F}$.
\end{enumerate}
\end{lemma}

\begin{proof}
Le support $\text{Supp}(\mathcal{F})$ de $\mathcal{F}$ est un sous-ensemble fermé
de $|X|$, voir Morphismes des espaces, lemme
\ref{spaces-morphisms-lemma-support-finite-type}.
Nous pouvons donc appliquer la définition \ref{spaces-perfect-definition-proper-over-base}.
Comme le support schématique de $\mathcal{F}$ est un sous-espace fermé
dont le sous-ensemble fermé sous-jacent est $\text{Supp}(\mathcal{F})$,
les conditions (1) et (2) sont équivalentes d'après la
définition \ref{spaces-perfect-definition-proper-over-base}.
Il est clair que (2) implique (3).
Réciproquement, si (3) est vraie, alors
$\text{Supp}(\mathcal{F}) \subset |Z|$
et donc $\text{Supp}(\mathcal{F})$
est propre au-dessus de $Y$, par exemple d'après le
lemme \ref{spaces-perfect-lemma-closed-closed-proper-over-base}.
\end{proof}

\begin{lemma}
\label{spaces-perfect-lemma-base-change-module-support-proper-over-base}
Soit $S$ un schéma.
Considérons un diagramme cartésien d'espaces algébriques au-dessus de $S$
$$
\xymatrix{
X' \ar[d]_{f'} \ar[r]_{g'} & X \ar[d]^f \\
Y' \ar[r]^g & Y
}
$$
tel que $f$ soit localement de type fini. Soit $\mathcal{F}$ un
$\mathcal{O}_X$-module de type fini et quasi-cohérent.
Si le support de $\mathcal{F}$ est propre au-dessus de $Y$, alors
le support de $(g')^*\mathcal{F}$ est propre au-dessus de $Y'$.
\end{lemma}

\begin{proof}
Observons que l'énoncé a un sens puisque
$(g')^*\mathcal{F}$ est de type fini d'après
Modules sur les sites, lemme \ref{sites-modules-lemma-local-pullback}.
On a $\text{Supp}((g')^*\mathcal{F}) = |g'|^{-1}(\text{Supp}(\mathcal{F}))$
d'après Morphismes des espaces, lemme \ref{spaces-morphisms-lemma-support-finite-type}.
Le lemme découle donc du
lemme \ref{spaces-perfect-lemma-base-change-closed-proper-over-base}.
\end{proof}

\begin{lemma}
\label{spaces-perfect-lemma-cat-module-support-proper-over-base}
Soit $S$ un schéma.
Soit $f : X \to Y$ un morphisme d'espaces algébriques au-dessus de $S$
localement de type fini. Soient $\mathcal{F}$ et $\mathcal{G}$
des $\mathcal{O}_X$-modules de type fini et quasi-cohérents.
\begin{enumerate}
\item Si les supports de $\mathcal{F}$ et $\mathcal{G}$
sont propres au-dessus de $Y$, il en va de même
pour $\mathcal{F} \oplus \mathcal{G}$, toute extension
de $\mathcal{G}$ par $\mathcal{F}$, $\Im(u)$ et $\Coker(u)$
pour tout morphisme de $\mathcal{O}_X$-modules $u : \mathcal{F} \to \mathcal{G}$,
ainsi que pour tout quotient quasi-cohérent de $\mathcal{F}$ ou de $\mathcal{G}$.
\item Si $Y$ est localement noethérien, la catégorie des
$\mathcal{O}_X$-modules cohérents dont le support est propre au-dessus de
$Y$ est une sous-catégorie de Serre (Homologie, définition
\ref{homology-definition-serre-subcategory})
de la catégorie abélienne des $\mathcal{O}_X$-modules cohérents.
\end{enumerate}
\end{lemma}

\begin{proof}
Preuve de (1). Soient $T$ et $T'$ les supports de $\mathcal{F}$
et de $\mathcal{G}$. Tous les faisceaux mentionnés en (1)
ont alors un support contenu dans $T \cup T'$. L'assertion découle donc
des lemmes \ref{spaces-perfect-lemma-closed-closed-proper-over-base} et
\ref{spaces-perfect-lemma-union-closed-proper-over-base}
pourvu que nous vérifiions que ces faisceaux sont de type fini
et quasi-cohérents.
Pour la quasi-cohérence, nous renvoyons le lecteur à
Propriétés des espaces, section \ref{spaces-properties-section-quasi-coherent}.
Pour la propriété « de type fini », nous renvoyons le lecteur à
Propriétés des espaces, section
\ref{spaces-properties-section-properties-modules}.

\medskip\noindent
Preuve de (2). La preuve est la même que celle de (1). Notons que
les assertions ont un sens puisque $X$ est localement noethérien d'après
Morphismes des espaces, lemme
\ref{spaces-morphisms-lemma-locally-finite-type-locally-noetherian}
et d'après la description de la catégorie des modules cohérents dans
Cohomologie des espaces, section \ref{spaces-cohomology-section-coherent}.
\end{proof}

\begin{lemma}
\label{spaces-perfect-lemma-support-proper-over-base-pushforward}
Soit $S$ un schéma. Soit $f : X \to Y$ un morphisme d'espaces algébriques
au-dessus de $S$. Supposons $f$ localement de type fini et $Y$ localement noethérien.
Soit $\mathcal{F}$ un $\mathcal{O}_X$-module cohérent dont le support
est propre au-dessus de $Y$. Alors $R^pf_*\mathcal{F}$ est un
$\mathcal{O}_Y$-module cohérent pour tout $p \geq 0$.
\end{lemma}

\begin{proof}
D'après le lemme \ref{spaces-perfect-lemma-module-support-proper-over-base},
il existe une immersion fermée $i : Z \to X$ telle que
$g = f \circ i : Z \to Y$ soit propre et
$\mathcal{F} = i_*\mathcal{G}$ pour un certain module cohérent $\mathcal{G}$
sur $Z$. On voit que $R^pg_*\mathcal{G}$
est cohérent sur $S$ d'après Cohomologie des espaces, lemme
\ref{spaces-cohomology-lemma-proper-pushforward-coherent}.
D'autre part, $R^qi_*\mathcal{G} = 0$ pour $q > 0$
(Cohomologie des espaces, lemme
\ref{spaces-cohomology-lemma-finite-pushforward-coherent}).
Par Cohomologie sur les sites, lemme \ref{sites-cohomology-lemma-relative-Leray},
nous obtenons $R^pf_*\mathcal{F} = R^pg_*\mathcal{G}$, d'où le lemme.
\end{proof}






\section{Catégorie dérivée des modules cohérents}
\label{spaces-perfect-section-derived-coherent}

\noindent
Soit $S$ un schéma. Soit $X$ un espace algébrique localement noethérien au-dessus de $S$.
Dans ce cas, la catégorie
$\textit{Coh}(\mathcal{O}_X) \subset \textit{Mod}(\mathcal{O}_X)$
des $\mathcal{O}_X$-modules cohérents est une sous-catégorie de Serre faible, voir
Homologie, section \ref{homology-section-serre-subcategories}
et
Cohomologie des espaces, lemme
\ref{spaces-cohomology-lemma-coherent-abelian-Noetherian}.
Notons
$$
D_{\textit{Coh}}(\mathcal{O}_X) \subset D(\mathcal{O}_X)
$$
la sous-catégorie des complexes dont les faisceaux de cohomologie sont cohérents, voir
Catégories dérivées, section \ref{derived-section-triangulated-sub}.
Nous obtenons ainsi un foncteur canonique
\begin{equation}
\label{spaces-perfect-equation-compare-coherent}
D(\textit{Coh}(\mathcal{O}_X))
\longrightarrow
D_{\textit{Coh}}(\mathcal{O}_X)
\end{equation}
voir Catégories dérivées, équation (\ref{derived-equation-compare}).

\begin{lemma}
\label{spaces-perfect-lemma-direct-image-coherent}
Soit $S$ un schéma. Soit $f : X \to Y$ un morphisme d'espaces algébriques
au-dessus de $S$. Supposons $f$ localement de type fini et $Y$ noethérien.
Soit $E$ un objet de $D^b_{\textit{Coh}}(\mathcal{O}_X)$ tel que le
support de $H^i(E)$ soit propre au-dessus de $Y$ pour tout $i$.
Alors $Rf_*E$ est un objet de $D^b_{\textit{Coh}}(\mathcal{O}_Y)$.
\end{lemma}

\begin{proof}
Considérons la suite spectrale
$$
R^pf_*H^q(E) \Rightarrow R^{p + q}f_*E
$$
voir Catégories dérivées, lemme \ref{derived-lemma-two-ss-complex-functor}.
D'après l'hypothèse et le lemme \ref{spaces-perfect-lemma-support-proper-over-base-pushforward},
les faisceaux $R^pf_*H^q(E)$ sont cohérents. Donc
$R^{p + q}f_*E$ est cohérent, c'est-à-dire $Rf_*E \in D_{\textit{Coh}}(\mathcal{O}_Y)$.
La bornitude inférieure est triviale. La bornitude supérieure
découle de
Cohomologie des espaces, lemme
\ref{spaces-cohomology-lemma-vanishing-higher-direct-images}
ou du
lemme \ref{spaces-perfect-lemma-quasi-coherence-direct-image}.
\end{proof}

\begin{lemma}
\label{spaces-perfect-lemma-direct-image-coherent-bdd-below}
Soit $S$ un schéma. Soit $f : X \to Y$ un morphisme d'espaces algébriques
au-dessus de $S$. Supposons $f$ localement de type fini et $Y$ noethérien.
Soit $E$ un objet de $D^+_{\textit{Coh}}(\mathcal{O}_X)$ tel que le support de $H^i(E)$
soit propre au-dessus de $Y$ pour tout $i$.
Alors $Rf_*E$ est un objet de $D^+_{\textit{Coh}}(\mathcal{O}_Y)$.
\end{lemma}

\begin{proof}
La preuve est la même que celle du
lemme \ref{spaces-perfect-lemma-direct-image-coherent}.
On peut aussi le déduire du
lemme \ref{spaces-perfect-lemma-direct-image-coherent}
en considérant l'action du foncteur exact $Rf_*$ sur les
triangles distingués
$\tau_{\leq a}E \to E \to \tau_{\geq a + 1}E \to \tau_{\leq a}E[1]$.
\end{proof}

\begin{lemma}
\label{spaces-perfect-lemma-coherent-internal-hom}
Soit $S$ un schéma. Soit $X$ un espace algébrique localement noethérien au-dessus de $S$.
Si $L$ appartient à $D^+_{\textit{Coh}}(\mathcal{O}_X)$
et $K$ à $D^-_{\textit{Coh}}(\mathcal{O}_X)$, alors
$R\SheafHom(K, L)$ appartient à $D^+_{\textit{Coh}}(\mathcal{O}_X)$.
\end{lemma}

\begin{proof}
Nous pouvons vérifier qu'un objet de $D(\mathcal{O}_X)$ appartient à
$D_{\textit{Coh}}(\mathcal{O}_X)$ localement pour la topologie étale sur $X$, voir
Cohomologie des espaces, lemme \ref{spaces-cohomology-lemma-coherent-Noetherian}.
On se ramène donc au cas des schémas, voir
Catégories dérivées des schémas, lemme \ref{perfect-lemma-coherent-internal-hom}.
\end{proof}

\begin{lemma}
\label{spaces-perfect-lemma-ext-finite}
Soit $A$ un anneau noethérien. Soit $X$ un espace algébrique propre au-dessus de $A$.
Pour $L$ dans $D^+_{\textit{Coh}}(\mathcal{O}_X)$ et $K$ dans
$D^-_{\textit{Coh}}(\mathcal{O}_X)$, les $A$-modules
$\Ext_{\mathcal{O}_X}^n(K, L)$ sont de type fini.
\end{lemma}

\begin{proof}
Rappelons que
$$
\Ext_{\mathcal{O}_X}^n(K, L) =
H^n(X, R\SheafHom_{\mathcal{O}_X}(K, L)) =
H^n(\Spec(A), Rf_*R\SheafHom_{\mathcal{O}_X}(K, L))
$$
voir Cohomologie sur les sites, lemme \ref{sites-cohomology-lemma-section-RHom-over-U}
et Cohomologie sur les sites, section \ref{sites-cohomology-section-leray}.
Le résultat découle donc des lemmes \ref{spaces-perfect-lemma-coherent-internal-hom} et
\ref{spaces-perfect-lemma-direct-image-coherent-bdd-below}.
\end{proof}




\section{Principe de récurrence}
\label{spaces-perfect-section-induction}

\noindent
Dans cette section, nous discutons d'un principe de récurrence pour les espaces algébriques,
analogue à celui qui est établi dans
Cohomologie des schémas, lemme \ref{coherent-lemma-induction-principle}
pour les schémas. Pour le formuler, nous introduisons la notion de
{\it carré distingué élémentaire} ; cette terminologie est empruntée
à \cite{MV}.
Le principe sous cette forme est implicite dans l'article \cite{GruRay}
de Raynaud et Gruson.
Un principe apparenté pour les champs algébriques est donné dans
\cite[Théorème D]{rydh_etale_devissage} par David Rydh.

\begin{definition}
\label{spaces-perfect-definition-elementary-distinguished-square}
Soit $S$ un schéma. Un diagramme commutatif
$$
\xymatrix{
U \times_W V \ar[r] \ar[d] & V \ar[d]^f \\
U \ar[r]^j & W
}
$$
d'espaces algébriques au-dessus de $S$ est appelé un {\it carré distingué élémentaire}
si
\begin{enumerate}
\item $U$ est un sous-espace ouvert de $W$ et $j$ est le morphisme d'inclusion ;
\item $f$ est étale ;
\item en posant $T = W \setminus U$ (muni de la structure réduite induite),
le morphisme $f^{-1}(T) \to T$ est un isomorphisme.
\end{enumerate}
Nous l'indiquerons en disant : « Soit $(U \subset W, f : V \to W)$
un carré distingué élémentaire. »
\end{definition}

\noindent
Notons que si $(U \subset W, f : V \to W)$ est un carré distingué
élémentaire, alors $W = U \cup f(V)$. Ainsi, $\{U \to W, V \to W\}$ est
un recouvrement étale de $W$. Ces recouvrements étales possèdent
de bonnes propriétés et sont, en un certain sens,
« suffisamment nombreux ».

\begin{lemma}
\label{spaces-perfect-lemma-make-more-elementary-distinguished-squares}
Soit $S$ un schéma. Soit $(U \subset W, f : V \to W)$ un carré
distingué élémentaire d'espaces algébriques au-dessus de $S$.
\begin{enumerate}
\item Si $V' \subset V$ et si
$U \subset U' \subset W$ sont des sous-espaces ouverts et $W' = U' \cup f(V')$,
alors $(U' \subset W', f|_{V'} : V' \to W')$ est un carré distingué
élémentaire.
\item Si $p : W' \to W$ est un morphisme d'espaces algébriques, alors
$(p^{-1}(U) \subset W', V \times_W W' \to W')$ est un carré distingué
élémentaire.
\item Si $S' \to S$ est un morphisme de schémas, alors
$(S' \times_S U \subset S' \times_S W, S' \times_S V \to S' \times_S W)$
est un carré distingué élémentaire.
\end{enumerate}
\end{lemma}

\begin{proof}
Omis.
\end{proof}

\begin{lemma}
\label{spaces-perfect-lemma-induction-principle}
Soit $S$ un schéma. Soit $X$ un espace algébrique quasi-compact et quasi-séparé
au-dessus de $S$. Soit $P$ une propriété des objets quasi-compacts
et quasi-séparés de $X_{spaces, \etale}$. Supposons que
\begin{enumerate}
\item tout objet affine de $X_{spaces, \etale}$ v\'erifie $P$ ;
\item pour tout carré distingué élémentaire $(U \subset W, f : V \to W)$
tel que
\begin{enumerate}
\item $W$ soit un objet quasi-compact et quasi-séparé de
$X_{spaces, \etale}$ ;
\item $U$ soit quasi-compact ;
\item $V$ soit affine ;
\item $U$, $V$ et $U \times_W V$ v\'erifient $P$,
\end{enumerate}
alors $W$ v\'erifie $P$.
\end{enumerate}
Alors tout objet quasi-compact et quasi-séparé de
$X_{spaces, \etale}$ v\'erifie $P$, et en particulier $X$.
\end{lemma}

\begin{proof}
Nous affirmons d'abord que tout objet représentable, quasi-compact et
quasi-séparé de $X_{spaces, \etale}$ v\'erifie $P$.
Supposons en effet que $U \to X$ soit étale et que $U$ soit un schéma
quasi-compact et quasi-séparé. Par l'hypothèse (1), tout ouvert affine de
$U$ v\'erifie $P$. De plus, si $W, V \subset U$ sont des ouverts
quasi-compacts, avec $V$ affine, et si $W$, $V$ et $W \cap V$ v\'erifient
$P$, alors $W \cup V$ v\'erifie $P$ d'après (2)
(car le couple $(W \subset W \cup V, V \to W \cup V)$ est un carré
distingué élémentaire). Le principe de récurrence pour les schémas montre donc
que $U$ v\'erifie $P$; voir
Cohomologie des schémas, lemme \ref{coherent-lemma-induction-principle}.

\medskip\noindent
Pour achever la preuve, il suffit de montrer que $X$ v\'erifie $P$
(car nous pouvons remplacer $X$ par n'importe quel objet quasi-compact et quasi-séparé
de $X_{spaces, \etale}$ auquel nous voulons appliquer le résultat).
Nous utiliserons la filtration
$$
\emptyset = U_{n + 1} \subset
U_n \subset U_{n - 1} \subset \ldots \subset U_1 = X
$$
et les morphismes $f_p : V_p \to U_p$ du
Espaces décents, lemme
\ref{decent-spaces-lemma-filter-quasi-compact-quasi-separated}.
Nous allons montrer par récurrence descendante sur $p$ que $U_p$ v\'erifie
$P$. Notons que $U_{n + 1}$ v\'erifie $P$ d'après (1), puisque l'espace
algébrique vide est affine. Supposons que $U_{p + 1}$ v\'erifie $P$.
Le couple $(U_{p + 1} \subset U_p, f_p : V_p \to U_p)$ est un carré distingué
élémentaire, mais (2) ne s'applique pas nécessairement, car $V_p$ peut ne pas être affine.
Toutefois, puisque $V_p$ est un schéma quasi-compact, nous pouvons choisir un recouvrement
fini par des ouverts affines $V_p = V_{p, 1} \cup \ldots \cup V_{p, m}$.
Posons $W_{p, 0} = U_{p + 1}$ et
$$
W_{p, i} = U_{p + 1} \cup f_p(V_{p, 1} \cup \ldots \cup V_{p, i})
$$
pour $i = 1, \ldots, m$. Ce sont des sous-espaces ouverts quasi-compacts de $X$.
Nous avons alors
$$
U_{p + 1} = W_{p, 0} \subset
W_{p, 1} \subset \ldots \subset
W_{p, m} = U_p
$$
et les couples
$$
(W_{p, 0} \subset W_{p, 1}, f_p|_{V_{p, 1}}),
(W_{p, 1} \subset W_{p, 2}, f_p|_{V_{p, 2}}),\ldots,
(W_{p, m - 1} \subset W_{p, m}, f_p|_{V_{p, m}})
$$
sont des carrés distingués élémentaires d'après le
lemme \ref{spaces-perfect-lemma-make-more-elementary-distinguished-squares}.
Notons que chaque $V_{p, i}$ v\'erifie $P$ (c'est un schéma affine), de même
que $W_{p, i} \times_{W_{p, i + 1}} V_{p, i + 1}$, car celui-ci est un ouvert
quasi-compact de $V_{p, i + 1}$ et vérifie donc $P$ d'après le premier
paragraphe de cette preuve. Ainsi, (2) s'applique à chacun de ces carrés et,
par récurrence, on conclut que
$W_{p, 1}, \ldots, W_{p, m} = U_p$ vérifient $P$.
\end{proof}

\begin{lemma}
\label{spaces-perfect-lemma-induction-principle-separated}
Soit $S$ un schéma. Soit $X$ un espace algébrique quasi-compact et quasi-séparé
au-dessus de $S$. Soit
$\mathcal{B} \subset \Ob(X_{spaces, \etale})$.
Soit $P$ une propriété des éléments de $\mathcal{B}$.
Supposons que
\begin{enumerate}
\item tout $W \in \mathcal{B}$ soit quasi-compact et quasi-séparé ;
\item si $W \in \mathcal{B}$ et si $U \subset W$ est un ouvert quasi-compact, alors
$U \in \mathcal{B}$ ;
\item si $V \in \Ob(X_{spaces, \etale})$ est affine, alors
(a) $V \in \mathcal{B}$ et (b) $V$ v\'erifie $P$ ;
\item pour tout carré distingué élémentaire $(U \subset W, f : V \to W)$
tel que
\begin{enumerate}
\item $W \in \mathcal{B}$ ;
\item $U$ soit quasi-compact ;
\item $V$ soit affine ;
\item $U$, $V$ et $U \times_W V$ v\'erifient $P$,
\end{enumerate}
alors $W$ v\'erifie $P$.
\end{enumerate}
Alors tout $W \in \mathcal{B}$ v\'erifie $P$.
\end{lemma}

\begin{proof}
La preuve est exactement la même que celle du
lemme \ref{spaces-perfect-lemma-induction-principle}.
(Notons que (4)(d) a un sens puisque $U \times_W V$ est un ouvert
quasi-compact de $V$, donc un élément de $\mathcal{B}$ d'après les conditions
(2) et (3).)
\end{proof}

\begin{remark}
\label{spaces-perfect-remark-how-to}
Comment choisir la collection $\mathcal{B}$ dans le
lemme \ref{spaces-perfect-lemma-induction-principle-separated} ?
Voici quelques exemples :
\begin{enumerate}
\item Si $X$ est quasi-compact et séparé, nous pouvons choisir pour
$\mathcal{B}$ l'ensemble des objets quasi-compacts et séparés
de $X_{spaces, \etale}$. Alors $X \in \mathcal{B}$ et $\mathcal{B}$
satisfait (1), (2) et (3)(a). Avec ce choix de $\mathcal{B}$, le
lemme \ref{spaces-perfect-lemma-induction-principle-separated} redonne le
lemme \ref{spaces-perfect-lemma-induction-principle}.
\item Si $X$ est quasi-compact et à diagonale affine
sur $\mathbf{Z}$ (comme dans Propriétés des espaces, définition
\ref{spaces-properties-definition-separated}), nous pouvons choisir pour
$\mathcal{B}$ l'ensemble des objets de $X_{spaces, \etale}$ qui sont
quasi-compacts et dont la diagonale sur $\mathbf{Z}$ est affine. Là encore,
$X \in \mathcal{B}$ et $\mathcal{B}$ satisfait (1), (2) et (3)(a).
\item Si $X$ est quasi-compact et quasi-séparé, le plus petit sous-ensemble
$\mathcal{B}$ qui contient $X$ et satisfait (1), (2) et (3)(a) est donné par
la règle suivante : $W \in \mathcal{B}$ si et seulement si $W$ est soit un
sous-espace ouvert quasi-compact de $X$, soit un ouvert quasi-compact d'un
objet affine de $X_{spaces, \etale}$.
\end{enumerate}
\end{remark}

\noindent
Voici une variante dans laquelle on propage la validit\'e d'une propri\'et\'e
d'un ouvert \`a des ouverts plus grands.

\begin{lemma}
\label{spaces-perfect-lemma-induction-principle-enlarge}
Soit $S$ un sch\'ema. Soit $X$ un espace alg\'ebrique quasi-compact et
quasi-s\'epar\'e sur $S$. Soit $W \subset X$ un sous-espace ouvert
quasi-compact. Soit $P$ une propri\'et\'e des sous-espaces ouverts
quasi-compacts de $X$. Supposons que
\begin{enumerate}
\item $W$ v\'erifie $P$, et
\item pour tout carr\'e distingu\'e \`el\'ementaire
$(W_1 \subset W_2, f : V \to W_2)$ satisfaisant aux conditions suivantes :
\begin{enumerate}
\item $W_1$ et $W_2$ soient des sous-espaces ouverts quasi-compacts de $X$,
\item $W \subset W_1$,
\item $V$ soit affine, et
\item $W_1$ v\'erifie $P$,
\end{enumerate}
alors $W_2$ v\'erifie $P$.
\end{enumerate}
Alors $X$ v\'erifie $P$.
\end{lemma}

\begin{proof}
On peut d\'eduire ce r\'esultat du
lemme \ref{spaces-perfect-lemma-induction-principle-separated}, mais nous allons plut\^ot
donner un argument direct en reprenant explicitement la preuve du
lemme \ref{spaces-perfect-lemma-induction-principle}. Nous utiliserons la filtration
$$
\emptyset = U_{n + 1} \subset
U_n \subset U_{n - 1} \subset \ldots \subset U_1 = X
$$
et les morphismes $f_p : V_p \to U_p$ du lemme
\ref{decent-spaces-lemma-filter-quasi-compact-quasi-separated} de
Espaces décents. Nous allons montrer, par r\'ecurrence descendante sur $p$,
que $W_p = W \cup U_p$ v\'erifie $P$. Cela ach\`evera la preuve puisque
$W_1 = X$. Notons que $W_{n + 1} = W \cup U_{n + 1} = W$ v\'erifie $P$
d'apr\`es (1). Supposons que $W_{p + 1}$ v\'erifie $P$. Le sous-espace
$W_p \setminus W_{p + 1}$, muni de la structure r\'eduite induite, est un
sous-espace ferm\'e de $U_p \setminus U_{p + 1}$.
Comme $(U_{p + 1} \subset U_p, f_p : V_p \to U_p)$ est un carr\'e distingu\'e
\'el\'ementaire, il en est de m\^eme de
$(W_{p + 1} \subset W_p, f_p : V_p \to W_p)$.
Il se peut toutefois que (2) ne s'applique pas, car $V_p$ n'est pas
n\'ecessairement affine. Puisque $V_p$ est un sch\'ema quasi-compact, nous
pouvons choisir un recouvrement ouvert affine fini
$V_p = V_{p, 1} \cup \ldots \cup V_{p, m}$. Posons
$W_{p, 0} = W_{p + 1}$ et
$$
W_{p, i} = W_{p + 1} \cup f_p(V_{p, 1} \cup \ldots \cup V_{p, i})
$$
pour $i = 1, \ldots, m$. Ce sont des sous-espaces ouverts quasi-compacts
de $X$ contenant $W$. On a alors
$$
W_{p + 1} = W_{p, 0} \subset
W_{p, 1} \subset \ldots \subset
W_{p, m} = W_p
$$
et les couples
$$
(W_{p, 0} \subset W_{p, 1}, f_p|_{V_{p, 1}}),
(W_{p, 1} \subset W_{p, 2}, f_p|_{V_{p, 2}}),\ldots,
(W_{p, m - 1} \subset W_{p, m}, f_p|_{V_{p, m}})
$$
sont des carr\'es distingu\'es \`el\'ementaires d'apr\`es le
lemme \ref{spaces-perfect-lemma-make-more-elementary-distinguished-squares}.
La condition (2) s'applique donc \`a chacun d'eux et, par r\'ecurrence sur
$i$, on conclut que $W_{p, 1}, \ldots, W_{p, m} = W_p$ v\'erifient $P$.
\end{proof}



\section{Mayer--Vietoris}
\label{spaces-perfect-section-mayer-vietoris}

\noindent
Dans cette section, nous montrons qu'un carr\'e distingu\'e \`el\'ementaire
donne naissance \`a diverses suites de Mayer--Vietoris.

\medskip\noindent
Soit $S$ un sch\'ema et soit $U \to X$ un morphisme \'etale d'espaces
alg\'ebriques sur $S$. Dans la section
\ref{spaces-properties-section-localize} de Propri\'et\'es des espaces,
on a montr\'e que
$U_{spaces, \etale} = X_{spaces, \etale}/U$
de fa\c{c}on compatible avec les faisceaux structuraux. Dans cette situation,
on consid\`ere donc souvent le morphisme $j_U : U \to X$ comme un morphisme
de localisation (voir Modules sur les sites, d\'efinition
\ref{sites-modules-definition-localize-ringed-site}).
En particulier, on consid\`ere l'image inverse $j_U^*$ comme la restriction
\`a $U$ et on la note souvent ${}|_U$; cette convention est compatible avec
l'\'equation
(\ref{spaces-properties-equation-restrict-modules})
de Propri\'et\'es des espaces. On voit notamment que
\begin{equation}
\label{spaces-perfect-equation-stalk-restriction}
(\mathcal{F}|_U)_{\overline{u}} = \mathcal{F}_{\overline{x}}
\end{equation}
si $\overline{u}$ est un point g\'eom\'etrique de $U$ et $\overline{x}$ son
image dans $X$. En outre, le foncteur de restriction admet un adjoint \`a
gauche exact $j_{U!}$; voir Modules sur les sites, lemmes
\ref{sites-modules-lemma-extension-by-zero} et
\ref{sites-modules-lemma-extension-by-zero-exact}.
Enfin, rappelons que, si $\mathcal{G}$ est un $\mathcal{O}_X$-module, alors
\begin{equation}
\label{spaces-perfect-equation-stalk-j-shriek}
(j_{U!}\mathcal{G})_{\overline{x}} =
\bigoplus\nolimits_{\overline{u}} \mathcal{G}_{\overline{u}}
\end{equation}
pour tout point g\'eom\'etrique $\overline{x} : \Spec(k) \to X$, o\`u la
somme directe porte sur les morphismes $\overline{u} : \Spec(k) \to U$
tels que $j_U \circ \overline{u} = \overline{x}$; voir Modules sur les sites,
lemme \ref{sites-modules-lemma-stalk-j-shriek}, et Propri\'et\'es des espaces,
lemme
\ref{spaces-properties-lemma-points-small-etale-site}.

\begin{lemma}
\label{spaces-perfect-lemma-exact-sequence-lower-shriek}
Soit $S$ un sch\'ema et soit $(U \subset X, V \to X)$ un carr\'e distingu\'e
\'el\'ementaire d'espaces alg\'ebriques sur $S$.
\begin{enumerate}
\item Pour tout faisceau de $\mathcal{O}_X$-modules $\mathcal{F}$, on a une
suite exacte courte
$$
0 \to j_{U \times_X V!}\mathcal{F}|_{U \times_X V} \to
j_{U!}\mathcal{F}|_U \oplus j_{V!}\mathcal{F}|_V \to \mathcal{F} \to 0
$$
\item Pour tout objet $E$ de $D(\mathcal{O}_X)$, on a un triangle distingu\'e
$$
j_{U \times_X V!}E|_{U \times_X V} \to
j_{U!}E|_U \oplus j_{V!}E|_V \to E \to 
j_{U \times_X V!}E|_{U \times_X V}[1]
$$
dans $D(\mathcal{O}_X)$.
\end{enumerate}
\end{lemma}

\begin{proof}
Pour montrer que la suite de (1) est exacte, il suffit de le v\'erifier sur
les fibres aux points g\'eom\'etriques, d'apr\`es Propri\'et\'es des espaces,
th\'eor\`eme
\ref{spaces-properties-theorem-exactness-stalks}.
Soit $\overline{x}$ un point g\'eom\'etrique de $X$. En prenant les fibres en
$\overline{x}$ et en utilisant les \`equations
(\ref{spaces-perfect-equation-stalk-restriction}) et (\ref{spaces-perfect-equation-stalk-j-shriek}), on
obtient la suite
$$
0 \to
\bigoplus\nolimits_{(\overline{u}, \overline{v})} \mathcal{F}_{\overline{x}}
\to
\bigoplus\nolimits_{\overline{u}} \mathcal{F}_{\overline{x}}
\oplus
\bigoplus\nolimits_{\overline{v}} \mathcal{F}_{\overline{x}}
\to
\mathcal{F}_{\overline{x}} \to 0
$$
Cette suite est exacte car, pour tout $\overline{x}$, ou bien il existe
exactement un $\overline{u}$ qui s'envoie sur $\overline{x}$, ou bien il
n'en existe aucun et il existe exactement un $\overline{v}$ qui s'envoie
sur $\overline{x}$.

\medskip\noindent
Preuve de (2). Nous avons vu dans Cohomologie sur les sites, section
\ref{sites-cohomology-section-properties-K-injective}, que les foncteurs de
restriction et de prolongement par z\'ero entre cat\'egories d\'eriv\'ees se
calculent en appliquant simplement le foncteur terme \`a terme \`a un
complexe quelconque. Soit $\mathcal{E}^\bullet$ un complexe de
$\mathcal{O}_X$-modules repr\'esentant $E$. Le triangle distingu\'e du lemme
est celui associ\'e, d'apr\`es Cat\'egories d\'eriv\'ees, section
\ref{derived-section-canonical-delta-functor}, et plus particuli\`erement le
lemme \ref{derived-lemma-derived-canonical-delta-functor}, \`a la suite de
complexes de $\mathcal{O}_X$-modules
$$
0 \to j_{U \times_X V!}\mathcal{E}^\bullet|_{U \times_X V} \to
j_{U!}\mathcal{E}^\bullet|_U \oplus j_{V!}\mathcal{E}^\bullet|_V
\to \mathcal{E}^\bullet \to 0
$$
qui est une suite exacte courte d'apr\`es (1).
\end{proof}

\begin{lemma}
\label{spaces-perfect-lemma-exact-sequence-j-star}
Soit $S$ un sch\'ema et soit $(U \subset X, V \to X)$ un carr\'e distingu\'e
\'el\'ementaire d'espaces alg\'ebriques sur $S$.
\begin{enumerate}
\item Pour tout faisceau de $\mathcal{O}_X$-modules $\mathcal{F}$, on a une
suite exacte courte
$$
0 \to \mathcal{F} \to
j_{U, *}\mathcal{F}|_U \oplus j_{V, *}\mathcal{F}|_V \to
j_{U \times_X V, *}\mathcal{F}|_{U \times_X V} \to 0
$$
\item Pour tout objet $E$ de $D(\mathcal{O}_X)$, on a un triangle distingu\'e
$$
E \to 
Rj_{U, *}E|_U \oplus Rj_{V, *}E|_V \to
Rj_{U \times_X V, *}E|_{U \times_X V} \to
E[1]
$$
dans $D(\mathcal{O}_X)$.
\end{enumerate}
\end{lemma}

\begin{proof}
Soit $W$ un objet de $X_\etale$. Nous affirmons que la suite
$$
0 \to
\mathcal{F}(W) \to
\mathcal{F}(W \times_X U) \oplus \mathcal{F}(W \times_X V) \to
\mathcal{F}(W \times_X U \times_X V)
$$
est exacte et que, localement sur $W$, tout \`el\'ement du dernier groupe
se rel\`eve dans celui du milieu. D'apr\`es le
lemme \ref{spaces-perfect-lemma-make-more-elementary-distinguished-squares}, les donn\'ees
$(W \times_X U \subset W, V \times_X W \to W)$ forment un carr\'e distingu\'e
\'el\'ementaire. Nous pouvons donc supposer que $W = X$, et il suffit
d'\'etablir le m\^eme fait pour
$$
0 \to
\mathcal{F}(X) \to
\mathcal{F}(U) \oplus \mathcal{F}(V) \to
\mathcal{F}(U \times_X V)
$$
Nous avons vu que
$$
0 \to j_{U \times_X V!}\mathcal{O}_{U \times_X V}
\to j_{U!}\mathcal{O}_U \oplus
j_{V!}\mathcal{O}_V \to
\mathcal{O}_X \to 0
$$
est une suite exacte de $\mathcal{O}_X$-modules, d'apr\`es le
lemme \ref{spaces-perfect-lemma-exact-sequence-lower-shriek}. L'application du foncteur
contravariant $\Hom_{\mathcal{O}_X}(- , \mathcal{F})$ donne alors la suite
ci-dessus. Cela signifie aussi que l'obstruction au rel\`evement de
$s \in \mathcal{F}(U \times_X V)$ en un \`el\'ement de
$\mathcal{F}(U) \oplus \mathcal{F}(V)$ appartient \`a
$\Ext^1_{\mathcal{O}_X}(\mathcal{O}_X, \mathcal{F}) =
H^1(X, \mathcal{F})$. Par localit\'e de la cohomologie (Cohomologie sur les
sites, lemme
\ref{sites-cohomology-lemma-kill-cohomology-class-on-covering})
cette obstruction s'annule localement pour la topologie \'etale sur $X$, ce
qui ach\`eve la preuve de (1).

\medskip\noindent
Preuve de (2). Choisissons un complexe K-injectif
$\mathcal{I}^\bullet$ repr\'esentant $E$ dont les termes $\mathcal{I}^n$
sont des objets injectifs de $\textit{Mod}(\mathcal{O}_X)$; voir Injectifs,
th\'eor\`eme
\ref{injectives-theorem-K-injective-embedding-grothendieck}.
Alors $\mathcal{I}^\bullet|U$ est un complexe K-injectif (Cohomologie sur
les sites, lemme
\ref{sites-cohomology-lemma-restrict-K-injective-to-open}).
Ainsi $Rj_{U, *}E|_U$ est repr\'esent\'e par
$j_{U, *}\mathcal{I}^\bullet|_U$, et de m\^eme pour $V$ et $U \times_X V$.
Le triangle distingu\'e du lemme est donc celui associ\'e, d'apr\`es
Cat\'egories d\'eriv\'ees, section
\ref{derived-section-canonical-delta-functor}, et plus particuli\`erement le
lemme \ref{derived-lemma-derived-canonical-delta-functor}, \`a la suite exacte
courte de complexes
$$
0 \to
\mathcal{I}^\bullet \to
j_{U, *}\mathcal{I}^\bullet|_U \oplus j_{V, *}\mathcal{I}^\bullet|_V \to
j_{U \times_X V, *}\mathcal{I}^\bullet|_{U \times_X V} \to
0.
$$
Cette suite est exacte d'apr\`es (1).
\end{proof}

\begin{lemma}
\label{spaces-perfect-lemma-unbounded-relative-mayer-vietoris}
Soit $S$ un sch\'ema. Soit $f : X \to Y$ un morphisme d'espaces alg\'ebriques
sur $S$, et soit $(U \subset X, V \to X)$ un carr\'e distingu\'e
\'el\'ementaire. Notons $a = f|_U : U \to Y$, $b = f|_V : V \to Y$ et
$c = f|_{U \times_X V} : U \times_X V \to Y$ les restrictions.
Pour tout objet $E$ de $D(\mathcal{O}_X)$, il existe un triangle distingu\'e
$$
Rf_*E \to
Ra_*(E|_U) \oplus Rb_*(E|_V) \to
Rc_*(E|_{U \times_X V}) \to
Rf_*E[1]
$$
dans $D(\mathcal{O}_Y)$. Ce triangle est fonctoriel en $E$.
\end{lemma}

\begin{proof}
Choisissons un complexe K-injectif $\mathcal{I}^\bullet$ repr\'esentant $E$.
Nous pouvons supposer que $\mathcal{I}^n$ est un objet injectif de
$\textit{Mod}(\mathcal{O}_X)$ pour tout $n$; voir Injectifs, th\'eor\`eme
\ref{injectives-theorem-K-injective-embedding-grothendieck}.
Alors $Rf_*E$ est calcul\'e par $f_*\mathcal{I}^\bullet$. Il en est de m\^eme
sur $U$, $V$ et $U \times_X V$, d'apr\`es Cohomologie sur les sites, lemme
\ref{sites-cohomology-lemma-restrict-K-injective-to-open}. Le triangle
distingu\'e du lemme est donc celui associ\'e, d'apr\`es Cat\'egories
d\'eriv\'ees, section \ref{derived-section-canonical-delta-functor}, et plus
particuli\`erement le lemme
\ref{derived-lemma-derived-canonical-delta-functor}, \`a la suite exacte
courte de complexes
$$
0 \to
f_*\mathcal{I}^\bullet \to
a_*\mathcal{I}^\bullet|_U \oplus b_*\mathcal{I}^\bullet|_V \to
c_*\mathcal{I}^\bullet|_{U \times_X V} \to
0.
$$
Pour voir qu'il s'agit d'une suite exacte courte de complexes, raisonnons
comme suit. Prenons un objet injectif $\mathcal{I}$ de
$\textit{Mod}(\mathcal{O}_X)$. Appliquons $f_*$ \`a la suite exacte courte
$$
0 \to \mathcal{I} \to
j_{U, *}\mathcal{I}|_U \oplus j_{V, *}\mathcal{I}|_V \to
j_{U \times_X V, *}\mathcal{I}|_{U \times_X V} \to 0
$$
du lemme \ref{spaces-perfect-lemma-exact-sequence-j-star}. Comme
$R^1f_*\mathcal{I} = 0$, on obtient la suite exacte courte
$$
0 \to f_*\mathcal{I} \to
f_*j_{U, *}\mathcal{I}|_U \oplus f_*j_{V, *}\mathcal{I}|_V \to
f_*j_{U \times_X V, *}\mathcal{I}|_{U \times_X V} \to 0
$$
Pour conclure, il suffit d'observer que $a_* = f_*j_{U, *}$, et de m\^eme
pour $b_*$ et $c_*$.
\end{proof}

\begin{lemma}
\label{spaces-perfect-lemma-mayer-vietoris-hom}
Soit $S$ un sch\'ema et soit $(U \subset X, V \to X)$ un carr\'e distingu\'e
\'el\'ementaire d'espaces alg\'ebriques sur $S$. Pour tous objets $E$ et $F$
de $D(\mathcal{O}_X)$, on a une suite de Mayer--Vietoris
$$
\xymatrix{
& \ldots \ar[r] &
\Ext^{-1}(E_{U \times_X V}, F_{U \times_X V}) \ar[lld] \\
\Hom(E, F) \ar[r] &
\Hom(E_U, F_U) \oplus
\Hom(E_V, F_V) \ar[r] &
\Hom(E_{U \times_X V}, F_{U \times_X V})
}
$$
o\`u les indices d\'esignent les restrictions aux ouverts correspondants et
les $\Hom$ sont pris dans les cat\'egories d\'eriv\'ees correspondantes.
\end{lemma}

\begin{proof}
Le triangle distingu\'e du lemme
\ref{spaces-perfect-lemma-exact-sequence-lower-shriek} donne une suite exacte longue de
$\Hom$ (Cat\'egories d\'eriv\'ees, lemme
\ref{derived-lemma-representable-homological}). On utilise ensuite l'\'egalit\'e
$\Hom(j_{U!}E|_U, F) = \Hom(E|_U, F|_U)$, qui r\'esulte de Cohomologie sur les
sites, lemme
\ref{sites-cohomology-lemma-adjoint-lower-shriek-restrict}.
\end{proof}

\begin{lemma}
\label{spaces-perfect-lemma-unbounded-mayer-vietoris}
Soit $S$ un sch\'ema et soit $(U \subset X, V \to X)$ un carr\'e distingu\'e
\'el\'ementaire d'espaces alg\'ebriques sur $S$. Pour tout objet $E$ de
$D(\mathcal{O}_X)$, on a un triangle distingu\'e
$$
R\Gamma(X, E) \to R\Gamma(U, E) \oplus R\Gamma(V, E) \to
R\Gamma(U \times_X V, E) \to R\Gamma(X, E)[1]
$$
et, en particulier, une suite exacte longue de cohomologie
$$
\ldots \to
H^n(X, E) \to
H^n(U, E) \oplus H^n(V, E) \to
H^n(U \times_X V, E) \to
H^{n + 1}(X, E) \to \ldots
$$
La construction du triangle distingu\'e et de la suite exacte longue est
fonctorielle en $E$.
\end{lemma}

\begin{proof}
Choisissons un complexe K-injectif $\mathcal{I}^\bullet$ repr\'esentant $E$
dont les termes $\mathcal{I}^n$ sont des objets injectifs de
$\textit{Mod}(\mathcal{O}_X)$; voir Injectifs, th\'eor\`eme
\ref{injectives-theorem-K-injective-embedding-grothendieck}.
Dans la preuve du lemme \ref{spaces-perfect-lemma-exact-sequence-j-star}, nous avons obtenu
une suite exacte courte de complexes
$$
0 \to \mathcal{I}^\bullet \to
j_{U, *}\mathcal{I}^\bullet|_U \oplus j_{V, *}\mathcal{I}^\bullet|_V \to
j_{U \times_X V, *}\mathcal{I}^\bullet|_{U \times_X V} \to 0
$$
Comme $H^1(X, \mathcal{I}^n) = 0$, le passage aux sections globales donne
une suite exacte de complexes
$$
0 \to \Gamma(X, \mathcal{I}^\bullet) \to
\Gamma(U, \mathcal{I}^\bullet) \oplus
\Gamma(V, \mathcal{I}^\bullet) \to
\Gamma(U \times_X V, \mathcal{I}^\bullet) \to 0
$$
Puisque ces complexes repr\'esentent $R\Gamma(X, E)$, $R\Gamma(U, E)$,
$R\Gamma(V, E)$ et $R\Gamma(U \times_X V, E)$, on obtient un triangle
distingu\'e d'apr\`es Cat\'egories d\'eriv\'ees, section
\ref{derived-section-canonical-delta-functor}, et plus particuli\`erement le
lemme \ref{derived-lemma-derived-canonical-delta-functor}.
\end{proof}

\begin{lemma}
\label{spaces-perfect-lemma-restrict-lower-shriek}
Soit $S$ un sch\'ema. Soit $j : U \to X$ un morphisme \'etale d'espaces
alg\'ebriques sur $S$. Pour tout morphisme \'etale $V \to X$, posons
$W = V \times_X U$ et notons $j_W : W \to V$ le morphisme de projection.
Alors $(j_!E)|_V = j_{W!}(E|_W)$ pour tout $E$ de $D(\mathcal{O}_U)$.
\end{lemma}

\begin{proof}
Cela r\'esulte de l'\'egalit\'e
$(j_!\mathcal{F})|_V = j_{W!}(\mathcal{F}|_W)$ pour tout
$\mathcal{O}_X$-module $\mathcal{F}$, laquelle d\'ecoule imm\'ediatement de la
construction des foncteurs $j_!$ et $j_{W!}$; voir Modules sur les sites,
lemme \ref{sites-modules-lemma-extension-by-zero}.
\end{proof}

\begin{lemma}
\label{spaces-perfect-lemma-pushforward-with-support-in-open}
Soit $S$ un sch\'ema et soit $(U \subset X, j : V \to X)$ un carr\'e
distingu\'e \`el\'ementaire d'espaces alg\'ebriques sur $S$. Posons
$T = |X| \setminus |U|$.
\begin{enumerate}
\item Si $E$ est un objet de $D(\mathcal{O}_X)$ \`a support dans $T$, alors
(a) $E \to Rj_*(E|_V)$ et (b) $j_!(E|_V) \to E$ sont des isomorphismes.
\item Si $F$ est un objet de $D(\mathcal{O}_V)$ \`a support dans $j^{-1}T$,
alors (a) $F \to (j_!F)|_V$, (b) $(Rj_*F)|_V \to F$ et (c)
$j_!F \to Rj_*F$ sont des isomorphismes.
\end{enumerate}
\end{lemma}

\begin{proof}
Soit $E$ un objet de $D(\mathcal{O}_X)$ dont les faisceaux de cohomologie
sont \`a support dans $T$. Alors $E|_U = 0$ et
$E|_{U \times_X V} = 0$, puisque $T$ ne rencontre pas $U$ et que $j^{-1}T$
ne rencontre pas $U \times_X V$. L'assertion (1)(a) r\'esulte donc du
lemme \ref{spaces-perfect-lemma-exact-sequence-j-star}. De la m\^eme mani\`ere, (1)(b)
r\'esulte du lemme \ref{spaces-perfect-lemma-exact-sequence-lower-shriek}.

\medskip\noindent
Soit $F$ un objet de $D(\mathcal{O}_V)$ dont les faisceaux de cohomologie
sont \`a support dans $j^{-1}T$. D'apr\`es le
lemme \ref{spaces-perfect-lemma-restrict-direct-image-open}, on a
$(Rj_*F)|_U = Rj_{W, *}(F|_W) = 0$, puisque $F|_W = 0$ par hypoth\`ese.
De m\^eme, $(j_!F)|_U = j_{W!}(F|_W) = 0$ d'apr\`es le
lemme \ref{spaces-perfect-lemma-restrict-lower-shriek}. Ainsi $j_!F$ et $Rj_*F$ sont \`a
support dans $T$, tandis que $(j_!F)|_V$ et $(Rj_*F)|_V$ sont \`a support
dans $j^{-1}(T)$. Pour v\'erifier que les morphismes (2)(a), (b) et (c) sont
des isomorphismes dans la cat\'egorie d\'eriv\'ee, il suffit de v\'erifier
qu'ils induisent des isomorphismes sur les fibres des faisceaux de
cohomologie aux points g\'eom\'etriques de $T$ et de $j^{-1}(T)$, d'apr\`es
Propri\'et\'es des espaces, th\'eor\`eme
\ref{spaces-properties-theorem-exactness-stalks}.
Nous pouvons le faire apr\`es avoir remplac\'e $X$ par $V$, $U$ par
$U \times_X V$, $V$ par $V \times_X V$ et $F$ par $F|_{V \times_X V}$
(restriction suivant la premi\`ere projection); voir les lemmes
\ref{spaces-perfect-lemma-restrict-direct-image-open},
\ref{spaces-perfect-lemma-restrict-lower-shriek} et
\ref{spaces-perfect-lemma-make-more-elementary-distinguished-squares}.
Comme $V \times_X V \to V$ admet une section, cela ram\`ene (2) au cas o\`u
$j : V \to X$ admet une section.

\medskip\noindent
Supposons que $j$ admette une section $\sigma : X \to V$. Posons
$V' = \sigma(X)$; c'est un sous-espace ouvert de $V$. Posons aussi
$U' = j^{-1}(U)$; c'est un autre sous-espace ouvert de $V$. Alors
$(U' \subset V, V' \to V)$ est un carr\'e distingu\'e \`el\'ementaire.
Comme $F$ est \`a support dans $j^{-1}(T)$, on a $F|_{U'} = 0$ et
$F|_{V' \cap U'} = 0$. Notons $j' : V' \to V$ l'immersion ouverte et
$j_{V'} : V' \to X$ le compos\'e $V' \to V \to X$, qui est l'inverse de
$\sigma$. Posons $F' = \sigma^*F$. Les triangles distingu\'es des lemmes
\ref{spaces-perfect-lemma-exact-sequence-lower-shriek} et
\ref{spaces-perfect-lemma-exact-sequence-j-star} montrent que
$F = j'_!(F|_{V'})$ et $F = Rj'_*(F|_{V'})$. Il s'ensuit que
$j_!F = j_!j'_!(F|_{V'}) = j_{V'!}F = F'$, car
$j_{V'} : V' \to X$ est un isomorphisme inverse de $\sigma$. De m\^eme,
$Rj_*F = Rj_*Rj'_*F = Rj_{V', *}F = F'$. Cela prouve (2)(c). Pour prouver
(2)(a) et (2)(b), il suffit de montrer que $F = F'|_V$. C'est clair, car
$F$ et $F'|_V$ ont tous deux pour restriction z\'ero sur $U'$ et sur
$U' \cap V'$, et le m\^eme objet sur $V'$.
\end{proof}

\noindent
On peut recoller des complexes !

\begin{lemma}
\label{spaces-perfect-lemma-glue}
Soit $S$ un sch\'ema et soit $(U \subset X, V \to X)$ un carr\'e distingu\'e
\'el\'ementaire d'espaces alg\'ebriques sur $S$. Supposons donn\'es
\begin{enumerate}
\item un objet $A$ de $D(\mathcal{O}_U)$,
\item un objet $B$ de $D(\mathcal{O}_V)$, et
\item un isomorphisme $c : A|_{U \times_X V} \to B|_{U \times_X V}$.
\end{enumerate}
Alors il existe un objet $F$ de $D(\mathcal{O}_X)$ et des isomorphismes
$f : F|_U \to A$ et $g : F|_V \to B$ tels que
$c = g|_{U \times_X V} \circ f^{-1}|_{U \times_X V}$. En outre, soient
\begin{enumerate}
\item un objet $E$ de $D(\mathcal{O}_X)$,
\item un morphisme $a : A \to E|_U$ dans $D(\mathcal{O}_U)$,
\item un morphisme $b : B \to E|_V$ dans $D(\mathcal{O}_V)$,
\end{enumerate}
tels que
$$
a|_{U \times_X V}  = b|_{U \times_X V} \circ c.
$$
Il existe alors un morphisme $F \to E$ dans $D(\mathcal{O}_X)$ dont la
restriction \`a $U$ est $a \circ f$ et la restriction \`a $V$ est
$b \circ g$.
\end{lemma}

\begin{proof}
Notons $j_U$, $j_V$ et $j_{U \times_X V}$ les morphismes correspondants
vers $X$. Choisissons un triangle distingu\'e
$$
F \to Rj_{U, *}A \oplus Rj_{V, *}B \to
Rj_{U \times_X V, *}(B|_{U \times_X V}) \to F[1]
$$
Ici, le morphisme
$Rj_{V, *}B \to Rj_{U \times_X V, *}(B|_{U \times_X V})$ est le morphisme
naturel. Le morphisme
$Rj_{U, *}A \to Rj_{U \times_X V, *}(B|_{U \times_X V})$
est le compos\'e de
$Rj_{U, *}A \to Rj_{U \times_X V, *}(A|_{U \times_X V})$
avec $Rj_{U \times_X V, *}c$. En restreignant \`a $U$, on obtient
$$
F|_U \to A \oplus (Rj_{V, *}B)|_U \to
(Rj_{U \times_X V, *}(B|_{U \times_X V}))|_U \to F|_U[1]
$$
Notons $j : U \times_X V \to U$. La compatibilit\'e entre restriction et
image directe totale (lemme \ref{spaces-perfect-lemma-restrict-direct-image-open}) montre
que $(Rj_{V, *}B)|_U$ et
$(Rj_{U \times_X V, *}(B|_{U \times_X V}))|_U$
sont tous deux canoniquement isomorphes \`a
$Rj_*(B|_{U \times_X V})$. Le second morphisme de la derni\`ere formule
affich\'ee admet donc une section, d'o\`u l'on conclut que
$F|_U \to A$ est un isomorphisme.

\medskip\noindent
Pour voir que $F|_V \to B$ est un isomorphisme, nous emploierons l'astuce
suivante. Choisissons un triangle distingu\'e
$$
F|_V \to B \to B' \to F[1]|_V
$$
dans $D(\mathcal{O}_V)$. Puisque $F|_U \to A$ est un isomorphisme et que
l'on dispose de l'isomorphisme
$c : A|_{U \times_X V} \to B|_{U \times_X V}$, la restriction de
$F|_V \to B$ \`a $U \times_X V$ est un isomorphisme. Ainsi $B'$ est \`a
support dans $j_V^{-1}(T)$, o\`u $T = |X| \setminus |U|$. D'autre part, on a
un morphisme de triangles distingu\'es
$$
\xymatrix{
F \ar[r] \ar[d] &
Rj_{U, *}F|_U \oplus Rj_{V, *}F|_V \ar[r] \ar[d] &
Rj_{U \times_X V, *}F|_{U \times_X V} \ar[r] \ar[d] &
F[1] \ar[d] \\
F \ar[r] &
Rj_{U, *}A \oplus Rj_{V, *}B \ar[r] &
Rj_{U \times_X V, *}(B|_{U \times_X V}) \ar[r] &
F[1]
}
$$
Tous les morphismes verticaux de ce diagramme, sauf peut-\^etre
$Rj_{V, *}F|_V \to Rj_{V, *}B$, sont des isomorphismes; ce dernier en est
donc aussi un (Cat\'egories d\'eriv\'ees, lemme
\ref{derived-lemma-third-isomorphism-triangle}). Il en r\'esulte que
$Rj_{V, *}B' = 0$, puis que $B' = 0$ d'apr\`es le
lemme \ref{spaces-perfect-lemma-pushforward-with-support-in-open}.

\medskip\noindent
L'existence du morphisme $F \to E$ r\'esulte de la suite de Mayer--Vietoris
pour $\Hom$; voir le lemme \ref{spaces-perfect-lemma-mayer-vietoris-hom}.
\end{proof}







\section{Le coh\'erateur}
\label{spaces-perfect-section-coherator}

\noindent
Soit $S$ un sch\'ema et soit $X$ un espace alg\'ebrique sur $S$.
Le {\it coh\'erateur} est un foncteur
$$
Q_X :
\textit{Mod}(\mathcal{O}_X)
\longrightarrow
\QCoh(\mathcal{O}_X)
$$
adjoint \`a droite au foncteur d'inclusion
$\QCoh(\mathcal{O}_X) \to \textit{Mod}(\mathcal{O}_X)$.
Il existe pour tout espace alg\'ebrique $X$ et, de plus, le morphisme
d'adjonction $Q_X(\mathcal{F}) \to \mathcal{F}$ est un isomorphisme pour tout
module quasi-coh\'erent $\mathcal{F}$; voir Propri\'et\'es des espaces,
proposition
\ref{spaces-properties-proposition-coherator}.
Comme $Q_X$ est exact \`a gauche (en tant qu'adjoint \`a droite), on peut
consid\'erer son extension d\'eriv\'ee \`a droite
$$
RQ_X :
D(\mathcal{O}_X)
\longrightarrow
D(\QCoh(\mathcal{O}_X)).
$$
Puisque $Q_X$ est adjoint \`a droite au foncteur d'inclusion
$\QCoh(\mathcal{O}_X) \to \textit{Mod}(\mathcal{O}_X)$, le foncteur $RQ_X$
est adjoint \`a droite au foncteur canonique
$D(\QCoh(\mathcal{O}_X)) \to D(\mathcal{O}_X)$, d'apr\`es Cat\'egories
d\'eriv\'ees, lemme \ref{derived-lemma-derived-adjoint-functors}.

\medskip\noindent
Dans cette section, nous \'etudierons le foncteur $RQ_X$. Dans la
section \ref{spaces-perfect-section-better-coherator}, nous \'etudierons l'adjoint \`a droite,
\'etroitement apparent\'e, du foncteur d'inclusion
$D_\QCoh(\mathcal{O}_X) \to D(\mathcal{O}_X)$, lorsqu'il existe.

\begin{lemma}
\label{spaces-perfect-lemma-affine-pushforward}
Soit $S$ un sch\'ema et soit $f : X \to Y$ un morphisme affine d'espaces
alg\'ebriques sur $S$. Alors $f_*$ d\'efinit un foncteur d\'eriv\'e
$f_* : D(\QCoh(\mathcal{O}_X)) \to D(\QCoh(\mathcal{O}_Y))$ tel que le
diagramme
$$
\xymatrix{
D(\QCoh(\mathcal{O}_X)) \ar[d]_{f_*} \ar[r] &
D_\QCoh(\mathcal{O}_X) \ar[d]^{Rf_*} \\
D(\QCoh(\mathcal{O}_Y)) \ar[r] &
D_\QCoh(\mathcal{O}_Y)
}
$$
est commutatif.
\end{lemma}

\begin{proof}
Le foncteur
$f_* : \QCoh(\mathcal{O}_X) \to \QCoh(\mathcal{O}_Y)$
est exact; voir Cohomologie des espaces, lemme
\ref{spaces-cohomology-lemma-affine-vanishing-higher-direct-images}.
Ainsi $f_*$ d\'efinit un foncteur d\'eriv\'e
$f_* : D(\QCoh(\mathcal{O}_X)) \to D(\QCoh(\mathcal{O}_Y))$ en appliquant
simplement $f_*$ \`a un complexe repr\'esentant l'objet consid\'er\'e; voir
Cat\'egories d\'eriv\'ees, lemme
\ref{derived-lemma-right-derived-exact-functor}. Pour tout complexe
$\mathcal{F}^\bullet$ de $\mathcal{O}_X$-modules, on dispose d'un morphisme
canonique
$f_*\mathcal{F}^\bullet \to Rf_*\mathcal{F}^\bullet$.
Pour achever la preuve, montrons que c'est un quasi-isomorphisme lorsque
chaque terme $\mathcal{F}^n$ du complexe $\mathcal{F}^\bullet$ est
quasi-coh\'erent. L'assertion est locale sur $Y$ pour la topologie \'etale;
nous pouvons donc supposer $Y$ affine. Tout morphisme affine \'etant
repr\'esentable, la compatibilit\'e de la
remarque \ref{spaces-perfect-remark-match-total-direct-images} nous ram\`ene au cas des
sch\'emas. Ce cas est trait\'e dans Cat\'egories d\'eriv\'ees des sch\'emas,
lemme \ref{perfect-lemma-affine-pushforward}.
\end{proof}

\begin{lemma}
\label{spaces-perfect-lemma-flat-pushforward-coherator}
Soit $S$ un sch\'ema et soit $f : X \to Y$ un morphisme d'espaces alg\'ebriques
sur $S$. Supposons $f$ quasi-compact, quasi-s\'epar\'e et plat. Si l'on note
$$
\Phi : D(\QCoh(\mathcal{O}_X)) \to D(\QCoh(\mathcal{O}_Y))
$$
le foncteur d\'eriv\'e droit de
$f_* : \QCoh(\mathcal{O}_X) \to \QCoh(\mathcal{O}_Y)$
alors $RQ_Y \circ Rf_* = \Phi \circ RQ_X$.
\end{lemma}

\begin{proof}
Nous allons le prouver en montrant que $RQ_Y \circ Rf_*$ et
$\Phi \circ RQ_X$ sont adjoints \`a droite au m\^eme foncteur
$D(\QCoh(\mathcal{O}_Y)) \to D(\mathcal{O}_X)$.

\medskip\noindent
Comme $f$ est quasi-compact et quasi-s\'epar\'e, $f_*$ pr\'eserve la
quasi-coh\'erence; voir Morphismes d'espaces, lemme
\ref{spaces-morphisms-lemma-pushforward}. Rappelons que
$\QCoh(\mathcal{O}_X)$ est une cat\'egorie ab\'elienne de Grothendieck
(Propri\'et\'es des espaces, proposition
\ref{spaces-properties-proposition-coherator}).
Ainsi tout $K$ de $D(\QCoh(\mathcal{O}_X))$ peut \^etre repr\'esent\'e par un
complexe K-injectif $\mathcal{I}^\bullet$ de $\QCoh(\mathcal{O}_X)$; voir
Injectifs, th\'eor\`eme
\ref{injectives-theorem-K-injective-embedding-grothendieck}.
On peut alors poser $\Phi(K) = f_*\mathcal{I}^\bullet$.

\medskip\noindent
Comme $f$ est plat, le foncteur $f^*$ est exact. Il d\'efinit donc les
foncteurs
$f^* : D(\mathcal{O}_Y) \to D(\mathcal{O}_X)$ et
$f^* : D(\QCoh(\mathcal{O}_Y)) \to D(\QCoh(\mathcal{O}_X))$.
Le foncteur $f^* = Lf^* : D(\mathcal{O}_Y) \to D(\mathcal{O}_X)$ est adjoint
\`a gauche \`a
$Rf_* : D(\mathcal{O}_X) \to D(\mathcal{O}_Y)$,
voir Cohomologie sur les sites, lemme
\ref{sites-cohomology-lemma-adjoint}. De m\^eme, le foncteur
$f^* : D(\QCoh(\mathcal{O}_Y)) \to D(\QCoh(\mathcal{O}_X))$
est adjoint \`a gauche \`a
$\Phi : D(\QCoh(\mathcal{O}_X)) \to D(\QCoh(\mathcal{O}_Y))$, d'apr\`es
Cat\'egories d\'eriv\'ees, lemme
\ref{derived-lemma-derived-adjoint-functors}.

\medskip\noindent
Soient $A$ un objet de $D(\QCoh(\mathcal{O}_Y))$ et $E$ un objet de
$D(\mathcal{O}_X)$. Alors
\begin{align*}
\Hom_{D(\QCoh(\mathcal{O}_Y))}(A, RQ_Y(Rf_*E))
& =
\Hom_{D(\mathcal{O}_Y)}(A, Rf_*E) \\
& =
\Hom_{D(\mathcal{O}_X)}(f^*A, E) \\
& =
\Hom_{D(\QCoh(\mathcal{O}_X))}(f^*A, RQ_X(E)) \\
& =
\Hom_{D(\QCoh(\mathcal{O}_Y))}(A, \Phi(RQ_X(E)))
\end{align*}
Cela donne le r\'esultat voulu.
\end{proof}

\begin{lemma}
\label{spaces-perfect-lemma-affine-coherator}
Soit $S$ un sch\'ema et soit $X$ un espace alg\'ebrique affine sur $S$.
Posons $A = \Gamma(X, \mathcal{O}_X)$. Alors
\begin{enumerate}
\item $Q_X : \textit{Mod}(\mathcal{O}_X) \to \QCoh(\mathcal{O}_X)$
est le foncteur qui associe \`a $\mathcal{F}$ le
$\mathcal{O}_X$-module quasi-coh\'erent associ\'e au $A$-module
$\Gamma(X, \mathcal{F})$,
\item $RQ_X : D(\mathcal{O}_X) \to D(\QCoh(\mathcal{O}_X))$
est le foncteur qui associe \`a $E$ le complexe de
$\mathcal{O}_X$-modules quasi-coh\'erents associ\'e \`a l'objet
$R\Gamma(X, E)$ de $D(A)$,
\item la restriction de $RQ_X$ \`a $D_\QCoh(\mathcal{O}_X)$ d\'efinit un
quasi-inverse de (\ref{spaces-perfect-equation-compare}).
\end{enumerate}
\end{lemma}

\begin{proof}
Soit $X_0 = \Spec(A)$ le sch\'ema affine qui repr\'esente $X$. Rappelons qu'il
existe un morphisme de sites annel\'es
$\epsilon : X_\etale \to X_{0, Zar}$
qui induit des \'equivalences
$$
\xymatrix{
\QCoh(\mathcal{O}_X) \ar@<1ex>[r]^{{\epsilon_*}} &
\QCoh(\mathcal{O}_{X_0}) \ar@<1ex>[l]^{{\epsilon^*}}
}
$$
voir le lemme
\ref{spaces-perfect-lemma-derived-quasi-coherent-small-etale-site}.
Par unicit\'e des foncteurs adjoints, on a donc
$Q_X = \epsilon^* \circ Q_{X_0} \circ \epsilon_*$. L'assertion (1) r\'esulte
alors de la description de $Q_{X_0}$ dans Cat\'egories d\'eriv\'ees des
sch\'emas, lemme \ref{perfect-lemma-affine-coherator}, et de l'\'egalit\'e
$\Gamma(X_0, \epsilon_*\mathcal{F}) = \Gamma(X, \mathcal{F})$.
L'assertion (2) r\'esulte de (1) et de l'exactitude du foncteur des
$A$-modules vers les $\mathcal{O}_X$-modules quasi-coh\'erents. La troisi\`eme
assertion d\'ecoule alors du r\'esultat pour les sch\'emas (Cat\'egories
d\'eriv\'ees des sch\'emas, lemme \ref{perfect-lemma-affine-coherator}) et du
lemme
\ref{spaces-perfect-lemma-derived-quasi-coherent-small-etale-site}.
\end{proof}

\noindent
Nous \'etablissons maintenant un crit\`ere assurant que le foncteur
$D(\QCoh(\mathcal{O}_X)) \to D_\QCoh(\mathcal{O}_X)$
est une \'equivalence.

\begin{lemma}
\label{spaces-perfect-lemma-argument-proves}
Soit $S$ un sch\'ema et soit $X$ un espace alg\'ebrique quasi-compact et
quasi-s\'epar\'e sur $S$. Supposons que, pour tout morphisme \'etale
$j : V \to W$, o\`u $W \subset X$ est ouvert quasi-compact et $V$ affine,
le foncteur d\'eriv\'e droit
$$
\Phi : D(\QCoh(\mathcal{O}_V)) \to D(\QCoh(\mathcal{O}_W))
$$
du foncteur exact \`a gauche
$j_* : \QCoh(\mathcal{O}_V) \to \QCoh(\mathcal{O}_W)$
s'inscrive dans le diagramme commutatif
$$
\xymatrix{
D(\QCoh(\mathcal{O}_V)) \ar[d]_\Phi \ar[r]_{i_V} &
D_\QCoh(\mathcal{O}_V) \ar[d]^{Rj_*} \\
D(\QCoh(\mathcal{O}_W)) \ar[r]^{i_W} &
D_\QCoh(\mathcal{O}_W)
}
$$
Alors le foncteur (\ref{spaces-perfect-equation-compare})
$$
D(\QCoh(\mathcal{O}_X))
\longrightarrow
D_\QCoh(\mathcal{O}_X)
$$
est une \'equivalence, dont $RQ_X$ est un quasi-inverse.
\end{lemma}

\begin{proof}
Nous utilisons d'abord le principe de r\'ecurrence pour montrer que $i_X$ est
pleinement fid\`ele. Plus pr\'ecis\'ement, nous appliquerons le
lemme \ref{spaces-perfect-lemma-induction-principle-enlarge}. Soit
$(U \subset W, V \to W)$ un carr\'e distingu\'e \`el\'ementaire, avec $V$
affine et $U,W$ ouverts quasi-compacts de $X$. Supposons $i_U$ pleinement
fid\`ele. Il faut montrer que $i_W$ est pleinement fid\`ele. Nous pouvons
remplacer $X$ par $W$, c'est-\`a-dire supposer que $W = X$; cette op\'eration
ne sert qu'\`a simplifier les notations et pr\'eserve la condition de l'\'enonc\'e.

\medskip\noindent
Soient $A$ et $B$ des objets de $D(\QCoh(\mathcal{O}_X))$. Nous voulons
montrer que
$$
\Hom_{D(\QCoh(\mathcal{O}_X))}(A, B)
\longrightarrow
\Hom_{D(\mathcal{O}_X)}(i_X(A), i_X(B))
$$
est bijectif. Posons $T = |X| \setminus |U|$.

\medskip\noindent
Supposons d'abord que $i_X(B)$ soit \`a support dans $T$. Dans ce cas, le
morphisme
$$
i_X(B) \to Rj_{V, *}(i_X(B)|_V) = Rj_{V, *}(i_V(B|_V))
$$
est un quasi-isomorphisme (lemme
\ref{spaces-perfect-lemma-pushforward-with-support-in-open}). Par hypoth\`ese, on a dans
$D(\mathcal{O}_X)$ un isomorphisme
$i_X(\Phi(B|_V)) \to Rj_{V, *}(i_V(B|_V))$. En outre, $\Phi$ et ${-}|_V$
sont des foncteurs adjoints entre les cat\'egories d\'eriv\'ees de modules
quasi-coh\'erents (Cat\'egories d\'eriv\'ees, lemme
\ref{derived-lemma-derived-adjoint-functors}). Le morphisme d'adjonction
$B \to \Phi(B|_V)$ devient un isomorphisme apr\`es application de $i_X$;
c'est donc un isomorphisme dans $D(\QCoh(\mathcal{O}_X))$. Par cons\'equent,
\begin{align*}
\Mor_{D(\QCoh(\mathcal{O}_X))}(A, B)
& =
\Mor_{D(\QCoh(\mathcal{O}_X))}(A, \Phi(B|_V)) \\
& =
\Mor_{D(\QCoh(\mathcal{O}_V))}(A|_V, B|_V) \\
& =
\Mor_{D(\mathcal{O}_V)}(i_V(A|_V), i_V(B|_V)) \\
& =
\Mor_{D(\mathcal{O}_X)}(i_X(A), Rj_{V, *}(i_V(B|_V))) \\
& =
\Mor_{D(\mathcal{O}_X)}(i_X(A), i_X(B))
\end{align*}
comme voulu. Nous avons utilis\'e ici le fait que $i_V$ est pleinement fid\`ele
(lemme \ref{spaces-perfect-lemma-affine-coherator}).

\medskip\noindent
Dans le cas g\'en\'eral, choisissons un complexe quelconque
$\mathcal{B}^\bullet$ de $\mathcal{O}_X$-modules quasi-coh\'erents qui
repr\'esente $B$, puis un quasi-isomorphisme quelconque
$s : \mathcal{B}^\bullet|_U \to \mathcal{C}^\bullet$ de complexes de modules
quasi-coh\'erents sur $U$. Comme $j_U : U \to X$ est quasi-compact et
quasi-s\'epar\'e, le foncteur $j_{U, *}$ envoie les modules quasi-coh\'erents
sur des modules quasi-coh\'erents (Morphismes d'espaces, lemme
\ref{spaces-morphisms-lemma-pushforward}). On dispose donc d'un morphisme
canonique
$\mathcal{B}^\bullet \to j_{U, *}(\mathcal{B}^\bullet|_U) \to
j_{U, *}\mathcal{C}^\bullet$
de complexes de modules quasi-coh\'erents sur $X$. Posons
$B'' = j_{U, *}\mathcal{C}^\bullet$ dans $D(\QCoh(\mathcal{O}_X))$ et
choisissons un triangle distingu\'e
$$
B \to B'' \to B' \to B[1]
$$
dans $D(\QCoh(\mathcal{O}_X))$. Comme le premier morphisme du triangle se
restreint en un isomorphisme sur $U$, l'objet $B'$ est \`a support dans $T$.
Dans le diagramme
$$
\xymatrix{
\Hom_{D(\QCoh(\mathcal{O}_X))}(A, B'[-1]) \ar[r] \ar[d] &
\Hom_{D(\mathcal{O}_X)}(i_X(A), i_X(B')[-1]) \ar[d] \\
\Hom_{D(\QCoh(\mathcal{O}_X))}(A, B) \ar[r] \ar[d] &
\Hom_{D(\mathcal{O}_X)}(i_X(A), i_X(B)) \ar[d] \\
\Hom_{D(\QCoh(\mathcal{O}_X))}(A, B'') \ar[r] \ar[d] &
\Hom_{D(\mathcal{O}_X)}(i_X(A), i_X(B'')) \ar[d] \\
\Hom_{D(\QCoh(\mathcal{O}_X))}(A, B') \ar[r] &
\Hom_{D(\mathcal{O}_X)}(i_X(A), i_X(B'))
}
$$
les colonnes sont exactes et les fl\`eches horizontales sup\'erieure et
inf\'erieure sont bijectives. Choisissons enfin un complexe
$\mathcal{A}^\bullet$ de modules quasi-coh\'erents repr\'esentant $A$.

\medskip\noindent
Soit $\alpha : i_X(A) \to i_X(B)$ un morphisme dans
$D(\mathcal{O}_X)$. Puisque $i_U$ est pleinement fid\`ele, la restriction
$\alpha|_U$ provient d'un morphisme de $D(\QCoh(\mathcal{O}_U))$. On peut
donc choisir $s : \mathcal{B}^\bullet|_U \to \mathcal{C}^\bullet$ comme
ci-dessus de telle sorte que $\alpha|_U$ soit repr\'esent\'e par un v\'eritable
morphisme de complexes
$\mathcal{A}^\bullet|_U \to \mathcal{C}^\bullet$. Celui-ci correspond \`a un
morphisme de complexes
$\mathcal{A}^\bullet \to j_{U, *}\mathcal{C}^\bullet$. Autrement dit, l'image
de $\alpha$ dans $\Hom_{D(\mathcal{O}_X)}(i_X(A), i_X(B''))$ provient d'un
\'el\'ement de $\Hom_{D(\QCoh(\mathcal{O}_X))}(A, B'')$. Une chasse au
diagramme montre alors que $\alpha$ provient d'un morphisme $A \to B$ dans
$D(\QCoh(\mathcal{O}_X))$.

Supposons enfin que $a : A \to B$ soit un morphisme de
$D(\QCoh(\mathcal{O}_X))$ qui devienne nul dans $D(\mathcal{O}_X)$. En
choisissant convenablement $\mathcal{B}^\bullet$, nous pouvons supposer que
$a$ est repr\'esent\'e par un morphisme de complexes
$a^\bullet : \mathcal{A}^\bullet \to \mathcal{B}^\bullet$. Comme $i_U$ est
pleinement fid\`ele, la restriction $a^\bullet|_U$ est nulle dans
$D(\QCoh(\mathcal{O}_U))$. On peut donc choisir $s$ de sorte que
$s \circ a^\bullet|_U : \mathcal{A}^\bullet|_U \to \mathcal{C}^\bullet$
soit homotope \`a z\'ero. En appliquant le foncteur $j_{U, *}$, on en d\'eduit
que $\mathcal{A}^\bullet \to j_{U, *}\mathcal{C}^\bullet$ est homotope \`a
z\'ero. Ainsi $a$ s'envoie sur z\'ero dans
$\Hom_{D(\QCoh(\mathcal{O}_X))}(A, B'')$.
Nous pouvons donc supposer que $a$ est l'image d'un \`el\'ement
$b \in \Hom_{D(\QCoh(\mathcal{O}_X))}(A, B'[-1])$. L'image de $b$ dans
$\Hom_{D(\mathcal{O}_X)}(i_X(A), i_X(B')[-1])$ provient d'un
$\gamma \in \Hom_{D(\mathcal{O}_X)}(A, B''[-1])$, puisque $a$ s'envoie sur
z\'ero dans le groupe de droite. Les fl\`eches horizontales \'etant
surjectives, comme nous l'avons vu plus haut, $\gamma$ provient d'un
$c$ de
$\Hom_{D(\QCoh(\mathcal{O}_X))}(A, B''[-1])$
et il en r\'esulte que $a = 0$, comme voulu.

\medskip\noindent
Nous savons maintenant que $i_X$ est pleinement fid\`ele pour l'espace $X$
initial. Comme $RQ_X$ est son adjoint \`a droite, on a
$RQ_X \circ i_X = \text{id}$ (Cat\'egories, lemme
\ref{categories-lemma-adjoint-fully-faithful}).
Pour achever la preuve, montrons que, pour tout
$E$ de $D_\QCoh(\mathcal{O}_X)$, le morphisme
$i_X(RQ_X(E)) \to E$ est un isomorphisme. Choisissons un triangle distingu\'e
$$
i_X(RQ_X(E)) \to E \to E' \to i_X(RQ_X(E))[1]
$$
dans $D_\QCoh(\mathcal{O}_X)$. Un argument formel fond\'e sur ce qui pr\'ec\`ede
montre que $i_X(RQ_X(E')) = 0$. Il suffit donc de prouver que, pour
$E \in D_\QCoh(\mathcal{O}_X)$, la condition $i_X(RQ_X(E)) = 0$ entra\^ine
$E = 0$. Consid\'erons un morphisme \'etale $j : V \to X$, avec $V$ affine.
D'apr\`es les lemmes \ref{spaces-perfect-lemma-affine-coherator} et
\ref{spaces-perfect-lemma-flat-pushforward-coherator}, ainsi que notre hypoth\`ese, on a
$$
Rj_*(E|_V) = Rj_*(i_V(RQ_V(E|_V))) = i_X(\Phi(RQ_V(E|_V))) =
i_X(RQ_X(Rj_*(E|_V)))
$$
Choisissons un triangle distingu\'e
$$
E \to Rj_*(E|_V) \to E' \to E[1]
$$
En appliquant $RQ_X$, on obtient un triangle distingu\'e
$$
0 \to RQ_X(Rj_*(E|_V)) \to RQ_X(E') \to 0[1]
$$
c'est-\`a-dire que le morphisme du milieu est un isomorphisme. Avec la suite
d'\'egalit\'es ci-dessus, on voit que notre premier triangle distingu\'e devient
un triangle distingu\'e
$$
E \to i_X(RQ_X(E')) \to E' \to E[1]
$$
o\`u le morphisme du milieu est le morphisme d'adjonction. Or le compos\'e
$E \to E'$ est nul; par adjonction, $E \to i_X(RQ_X(E'))$ est donc nul.
Ce morphisme \'etant isomorphe au morphisme
$E \to Rj_*(E|_V)$ adjoint \`a $\text{id} : E|_V \to E|_V$, on en d\'eduit
que $E|_V$ est nul. Comme cela vaut pour tout $V$ affine et \'etale sur $X$,
on conclut que $E$ est nul, comme voulu.
\end{proof}

\begin{proposition}
\label{spaces-perfect-proposition-quasi-compact-affine-diagonal}
Soit $S$ un sch\'ema et soit $X$ un espace alg\'ebrique quasi-compact sur $S$
dont la diagonale sur $\mathbf{Z}$ est affine (au sens de Propri\'et\'es des
espaces, d\'efinition \ref{spaces-properties-definition-separated}). Alors le
foncteur
(\ref{spaces-perfect-equation-compare})
$$
D(\QCoh(\mathcal{O}_X))
\longrightarrow
D_\QCoh(\mathcal{O}_X)
$$
est une \'equivalence, dont $RQ_X$ est un quasi-inverse.
\end{proposition}

\begin{proof}
Soit $V \to W$ un morphisme \'etale, avec $V$ affine et $W$ un sous-espace
ouvert quasi-compact de $X$. Le morphisme $V \to W$ est affine, car la
diagonale de $W$ sur $\mathbf{Z}$ est affine et $V$ est affine (Morphismes
d'espaces, lemme \ref{spaces-morphisms-lemma-affine-permanence}). Le
lemme \ref{spaces-perfect-lemma-affine-pushforward} assure alors que l'hypoth\`ese du
lemme \ref{spaces-perfect-lemma-argument-proves} est satisfaite, ce qui permet de conclure.
\end{proof}

\begin{lemma}
\label{spaces-perfect-lemma-direct-image-coherator}
Soit $S$ un sch\'ema et soit $f : X \to Y$ un morphisme d'espaces alg\'ebriques
sur $S$. Supposons $X$ et $Y$ quasi-compacts et \`a diagonale affine sur
$\mathbf{Z}$ (au sens de Propri\'et\'es des espaces, d\'efinition
\ref{spaces-properties-definition-separated}). Si l'on note
$$
\Phi : D(\QCoh(\mathcal{O}_X)) \to D(\QCoh(\mathcal{O}_Y))
$$
le foncteur d\'eriv\'e droit de
$f_* : \QCoh(\mathcal{O}_X) \to \QCoh(\mathcal{O}_Y)$
alors le diagramme
$$
\xymatrix{
D(\QCoh(\mathcal{O}_X)) \ar[d]_\Phi \ar[r] &
D_\QCoh(\mathcal{O}_X) \ar[d]^{Rf_*} \\
D(\QCoh(\mathcal{O}_Y)) \ar[r] &
D_\QCoh(\mathcal{O}_Y)
}
$$
est commutatif.
\end{lemma}

\begin{proof}
Les fl\`eches horizontales du diagramme sont des \'equivalences de cat\'egories
d'apr\`es la proposition
\ref{spaces-perfect-proposition-quasi-compact-affine-diagonal}. Nous pouvons donc identifier
ces cat\'egories, et faire de m\^eme pour tout autre espace alg\'ebrique
quasi-compact \`a diagonale affine. L'\'enonc\'e du lemme affirme alors que le
morphisme canonique $\Phi(K) \to Rf_*(K)$ est un isomorphisme pour tout
$K$ de $D(\QCoh(\mathcal{O}_X))$. Remarquons que, si
$K_1 \to K_2 \to K_3 \to K_1[1]$ est un triangle distingu\'e dans
$D(\QCoh(\mathcal{O}_X))$ et que l'assertion est vraie pour deux des trois
objets, elle est vraie pour le troisi\`eme.

\medskip\noindent
Soit $\mathcal{B} \subset \Ob(X_{spaces, \etale})$ l'ensemble des objets
quasi-compacts \`a diagonale affine. Pour $U \in \mathcal{B}$ et tout
morphisme $g : U \to Z$, o\`u $Z$ est un espace alg\'ebrique quasi-compact
sur $S$ \`a diagonale affine, notons
$$
\Phi_g : D(\QCoh(\mathcal{O}_U)) \to D(\QCoh(\mathcal{O}_Z))
$$
le foncteur d\'eriv\'e droit de $g_*$. D\'efinissons $P(U)$ comme l'assertion :
``pour tout $K$ de $D(\QCoh(\mathcal{O}_U))$ et tout $g : U \to Z$ comme
ci-dessus, le morphisme $\Phi_g(K) \to Rg_*K$ est un isomorphisme''.
D'apr\`es la remarque \ref{spaces-perfect-remark-how-to}, les conditions (1), (2) et (3)(a)
du lemme \ref{spaces-perfect-lemma-induction-principle-separated} sont satisfaites; il reste
\`a prouver (3)(b) et (4).

\medskip\noindent
V\'erifions la condition (3)(b). Soit $U$ un sch\'ema affine \'etale sur $X$
et soit $g : U \to Z$ comme ci-dessus. Puisque la diagonale de $Z$ est
affine, le morphisme $g : U \to Z$ est affine (Morphismes d'espaces, lemme
\ref{spaces-morphisms-lemma-affine-permanence}).
Le lemme \ref{spaces-perfect-lemma-affine-pushforward} montre donc que $U$ v\'erifie $P$.

\medskip\noindent
V\'erifions la condition (4). Soit $(U \subset W, V \to W)$ un carr\'e
distingu\'e \`el\'ementaire de $X_{spaces, \etale}$, avec
$U,W,V \in \mathcal{B}$ et $V$ affine. Supposons que $U$, $V$ et
$U \times_W V$ v\'erifient $P$. Il faut montrer que $W$ v\'erifie $P$.
Soit $g : W \to Z$ un morphisme vers un espace alg\'ebrique quasi-compact
\`a diagonale affine, et soit $K$ un objet de
$D(\QCoh(\mathcal{O}_W))$. Consid\'erons le triangle distingu\'e
$$
K \to Rj_{U, *}K|_U \oplus Rj_{V, *}K|_V \to
Rj_{U \times_W V, *}K|_{U \times_W V} \to K[1]
$$
dans $D(\mathcal{O}_W)$. Par la propri\'et\'e des deux sur trois rappel\'ee
ci-dessus, il suffit de montrer que
$\Phi_g(Rj_{U, *}K|_U) \to Rg_*(Rj_{U, *}K|_U)$ est un isomorphisme, et de
m\^eme pour $V$ et $U \times_W V$. C'est l'objet du paragraphe suivant.

\medskip\noindent
Soit $j : U \to W$ un morphisme de $X_{spaces, \etale}$, avec
$U,W \in \mathcal{B}$ et $U$ v\'erifiant $P$. Soit $g : W \to Z$ un
morphisme vers un espace alg\'ebrique quasi-compact \`a diagonale affine.
Pour achever la preuve, il faut montrer que
$\Phi_g(Rj_*K) \to Rg_*(Rj_*K)$
est un isomorphisme pour tout $K$ de $D(\QCoh(\mathcal{O}_U))$.
Soit $\mathcal{I}^\bullet$ un complexe K-injectif de
$\QCoh(\mathcal{O}_U)$ repr\'esentant $K$. En appliquant $P(U)$ \`a $j$, on
voit que $j_*\mathcal{I}^\bullet$ repr\'esente $Rj_*K$. Comme
$j_* : \QCoh(\mathcal{O}_U) \to \QCoh(\mathcal{O}_W)$ admet l'adjoint \`a
gauche exact $j^* : \QCoh(\mathcal{O}_W) \to \QCoh(\mathcal{O}_U)$, le
complexe $j_*\mathcal{I}^\bullet$ est K-injectif dans
$\QCoh(\mathcal{O}_W)$; voir Cat\'egories d\'eriv\'ees, lemme
\ref{derived-lemma-adjoint-preserve-K-injectives}. Ainsi
$\Phi_g(Rj_*K)$ est repr\'esent\'e par
$g_*j_*\mathcal{I}^\bullet = (g \circ j)_*\mathcal{I}^\bullet$.
En appliquant $P(U)$ \`a $g \circ j$, on voit que ce complexe repr\'esente
$R_{g \circ j, *}(K) = Rg_*(Rj_*K)$, ce qui ach\`eve la preuve.
\end{proof}









\section{Le coh\'erateur pour les espaces noeth\'eriens}
\label{spaces-perfect-section-coherator-Noetherian}

\noindent
Il nous faut encore quelques r\'esultats sur les modules injectifs pour traiter
le cas d'un espace alg\'ebrique noeth\'erien.

\begin{lemma}
\label{spaces-perfect-lemma-affine-injective-colimit-direct-sum-pushforwards-artin}
Soit $S$ un sch\'ema affine noeth\'erien. Tout objet injectif de
$\QCoh(\mathcal{O}_S)$ est une limite inductive filtrante
$\colim_i \mathcal{F}_i$ de faisceaux quasi-coh\'erents de la forme
$$
\mathcal{F}_i = (Z_i \to S)_*\mathcal{G}_i
$$
o\`u $Z_i$ est le spectre d'un anneau artinien et $\mathcal{G}_i$ un module
coh\'erent sur $Z_i$.
\end{lemma}

\begin{proof}
Posons $S = \Spec(A)$. Soit $\mathcal{J}$ un objet injectif de
$\QCoh(\mathcal{O}_S)$. Comme $\QCoh(\mathcal{O}_S)$ est
\'equivalente \`a la cat\'egorie des $A$-modules, on a
$\mathcal{J}=\widetilde{J}$ pour un $A$-module injectif $J$.
D'apr\`es Complexes dualisants, proposition
\ref{dualizing-proposition-structure-injectives-noetherian}
on peut \'ecrire $J = \bigoplus E_\alpha$, avec $E_\alpha$ ind\'ecomposable,
donc isomorphe \`a l'enveloppe injective du corps r\'esiduel en un point.
Ainsi (puisqu'une somme disjointe finie de sch\'emas artiniens est
artinienne), on peut supposer que $J$ est l'enveloppe injective de
$\kappa(\mathfrak p)$ pour un id\'eal premier $\mathfrak p$ de $A$.
Alors $J = \bigcup J[\mathfrak p^n]$, o\`u $J[\mathfrak p^n]$ est
l'enveloppe injective de $\kappa(\mathfrak p)$ sur
$A_\mathfrak p/\mathfrak p^nA_\mathfrak p$; voir
Complexes dualisants, lemme \ref{dualizing-lemma-union-artinian}.
Ainsi $\widetilde{J}$ est la limite inductive des faisceaux
$(Z_n \to S)_*\mathcal{G}_n$, o\`u
$Z_n = \Spec(A_\mathfrak p/\mathfrak p^nA_\mathfrak p)$ et
$\mathcal{G}_n$ est le faisceau coh\'erent associ\'e au
$A_\mathfrak p/\mathfrak p^nA_\mathfrak p$-module fini
$J[\mathfrak p^n]$. La finitude r\'esulte de
Complexes dualisants, lemme \ref{dualizing-lemma-finite}.
\end{proof}

\begin{lemma}
\label{spaces-perfect-lemma-injective-colimit-direct-sum-pushforwards-artin}
Soit $S$ un sch\'ema affine et soit $X$ un espace alg\'ebrique noeth\'erien
sur $S$. Tout objet injectif de $\QCoh(\mathcal{O}_X)$ est facteur direct
d'une limite inductive filtrante $\colim_i \mathcal{F}_i$ de faisceaux
quasi-coh\'erents de la forme
$$
\mathcal{F}_i = (Z_i \to X)_*\mathcal{G}_i
$$
o\`u $Z_i$ est le spectre d'un anneau artinien et $\mathcal{G}_i$ un module
coh\'erent sur $Z_i$.
\end{lemma}

\begin{proof}
Choisissons un sch\'ema affine $U$ et un morphisme \'etale surjectif
$j : U \to X$ (Propri\'et\'es des espaces, lemme
\ref{spaces-properties-lemma-quasi-compact-affine-cover}).
Alors $U$ est un sch\'ema affine noeth\'erien. Choisissons un objet injectif
$\mathcal{J}'$ de $\QCoh(\mathcal{O}_U)$ et une injection
$\mathcal{J}|_U \to \mathcal{J}'$. Le morphisme
$$
\mathcal{J} \to j_*\mathcal{J}'
$$
est alors injectif dans $\QCoh(\mathcal{O}_X)$; il identifie donc
$\mathcal{J}$ \`a un facteur direct de $j_*\mathcal{J}'$.
Le r\'esultat d\'ecoule donc du r\'esultat correspondant pour
$\mathcal{J}'$, d\'emontr\'e dans le lemme
\ref{spaces-perfect-lemma-affine-injective-colimit-direct-sum-pushforwards-artin}.
\end{proof}

\begin{lemma}
\label{spaces-perfect-lemma-flat-pullback-injective-quasi-coherent}
Soit $S$ un sch\'ema. Soit $f : X \to Y$ un morphisme plat,
quasi-compact et quasi-s\'epar\'e d'espaces alg\'ebriques sur $S$. Si
$\mathcal{J}$ est un objet injectif de $\QCoh(\mathcal{O}_X)$,
alors $f_*\mathcal{J}$ est un objet injectif de
$\QCoh(\mathcal{O}_Y)$.
\end{lemma}

\begin{proof}
Puisque $f$ est quasi-compact et quasi-s\'epar\'e, le foncteur
$f_*$ transforme les faisceaux quasi-coh\'erents en faisceaux
quasi-coh\'erents (Morphismes d'espaces, lemme
\ref{spaces-morphisms-lemma-pushforward}). Le foncteur $f^*$ est un adjoint
\`a gauche de $f_*$ qui transforme les injections en injections.
Le r\'esultat d\'ecoule donc d'Homologie, lemme
\ref{homology-lemma-adjoint-preserve-injectives}.
\end{proof}

\begin{lemma}
\label{spaces-perfect-lemma-injective-pushforward}
Soit $S$ un sch\'ema et soit $X$ un espace alg\'ebrique noeth\'erien sur $S$.
Si $\mathcal{J}$ est un objet injectif de $\QCoh(\mathcal{O}_X)$, alors
\begin{enumerate}
\item $H^p(U, \mathcal{J}|_U) = 0$ pour $p > 0$ et pour tout espace
alg\'ebrique $U$ quasi-compact et quasi-s\'epar\'e, \'etale sur $X$;
\item pour tout morphisme $f : X \to Y$ d'espaces alg\'ebriques sur $S$,
avec $Y$ quasi-s\'epar\'e, on a $R^pf_*\mathcal{J} = 0$ pour $p > 0$.
\end{enumerate}
\end{lemma}

\begin{proof}
D\'emonstration de (1). \'Ecrivons $\mathcal{J}$ comme facteur direct de
$\colim \mathcal{F}_i$, avec
$\mathcal{F}_i = (Z_i \to X)_*\mathcal{G}_i$, comme dans le lemme
\ref{spaces-perfect-lemma-injective-colimit-direct-sum-pushforwards-artin}.
Il suffit manifestement de d\'emontrer l'annulation pour
$\colim \mathcal{F}_i$. Comme l'image inverse commute aux limites inductives
et que $U$ est quasi-compact et quasi-s\'epar\'e, il suffit de d\'emontrer
$H^p(U, \mathcal{F}_i|_U) = 0$ pour $p > 0$; voir Cohomologie des espaces,
lemme \ref{spaces-cohomology-lemma-colimits}. Observons que
$Z_i \to X$ est un morphisme affine; voir Morphismes d'espaces, lemme
\ref{spaces-morphisms-lemma-Artinian-affine}. Par suite,
$$
\mathcal{F}_i|_U = (Z_i \times_X U \to U)_*\mathcal{G}'_i =
R(Z_i \times_X U \to U)_*\mathcal{G}'_i
$$
o\`u $\mathcal{G}'_i$ est l'image inverse de $\mathcal{G}_i$ sur
$Z_i \times_X U$; voir Cohomologie des espaces, lemme
\ref{spaces-cohomology-lemma-affine-base-change}.
Comme $Z_i \times_X U$ est affine, on en conclut que $\mathcal{G}'_i$
n'a pas de cohomologie sup\'erieure sur $Z_i \times_X U$.
La suite spectrale de Leray donne le m\^eme r\'esultat pour
$\mathcal{F}_i|_U$ (Cohomologie sur les sites, lemme
\ref{sites-cohomology-lemma-apply-Leray}).

\medskip\noindent
D\'emonstration de (2). Soit $f : X \to Y$ un morphisme d'espaces
alg\'ebriques sur $S$, et soit $V \to Y$ un morphisme \'etale avec $V$
affine. Alors $V \times_Y X \to X$ est un morphisme \'etale et
$V \times_Y X$ est un espace alg\'ebrique quasi-compact et quasi-s\'epar\'e,
\'etale sur $X$ (d\'etails omis). Par (1),
$H^p(V \times_Y X, \mathcal{J})$ est donc nul. Comme
$R^pf_*\mathcal{J}$ est le faisceau associ\'e au pr\'efaisceau
$V \mapsto H^p(V \times_Y X, \mathcal{J})$, la d\'emonstration est achev\'ee.
\end{proof}

\begin{lemma}
\label{spaces-perfect-lemma-Noetherian-pushforward}
Soit $S$ un sch\'ema et soit $f : X \to Y$ un morphisme d'espaces
alg\'ebriques noeth\'eriens sur $S$. Le foncteur $f_*$ sur les faisceaux
quasi-coh\'erents admet alors une extension d\'eriv\'ee \`a droite
$\Phi : D(\QCoh(\mathcal{O}_X)) \to D(\QCoh(\mathcal{O}_Y))$
telle que le diagramme
$$
\xymatrix{
D(\QCoh(\mathcal{O}_X)) \ar[d]_{\Phi} \ar[r] &
D_\QCoh(\mathcal{O}_X) \ar[d]^{Rf_*} \\
D(\QCoh(\mathcal{O}_Y)) \ar[r] &
D_\QCoh(\mathcal{O}_Y)
}
$$
soit commutatif.
\end{lemma}

\begin{proof}
Comme $X$ et $Y$ sont noeth\'eriens, le morphisme est quasi-compact et
quasi-s\'epar\'e (voir Morphismes d'espaces, lemme
\ref{spaces-morphisms-lemma-quasi-compact-quasi-separated-permanence}).
Ainsi $f_*$ pr\'eserve la quasi-coh\'erence; voir Morphismes d'espaces,
lemme \ref{spaces-morphisms-lemma-pushforward}.
Soit maintenant $K$ un objet de $D(\QCoh(\mathcal{O}_X))$.
Comme $\QCoh(\mathcal{O}_X)$ est une cat\'egorie ab\'elienne de Grothendieck
(Propri\'et\'es des espaces, proposition
\ref{spaces-properties-proposition-coherator}), on peut repr\'esenter $K$
par un complexe K-injectif $\mathcal{I}^\bullet$ tel que chaque
$\mathcal{I}^n$ soit un objet injectif de $\QCoh(\mathcal{O}_X)$; voir
Injectifs, th\'eor\`eme
\ref{injectives-theorem-K-injective-embedding-grothendieck}.
On d\'efinit alors le foncteur $\Phi$ en posant
$$
\Phi(K) = f_*\mathcal{I}^\bullet
$$
o\`u le membre de droite est consid\'er\'e comme un objet de
$D(\QCoh(\mathcal{O}_Y))$. Pour achever la preuve du lemme, il suffit de
montrer que le morphisme canonique
$$
f_*\mathcal{I}^\bullet \longrightarrow Rf_*\mathcal{I}^\bullet
$$
est un isomorphisme dans $D(\mathcal{O}_Y)$. Pour cela, il suffit de montrer
qu'il induit un isomorphisme sur les faisceaux de cohomologie. Fixons
$m \in \mathbf{Z}$ et soit $N = N(X, Y, f)$ comme dans le lemme
\ref{spaces-perfect-lemma-quasi-coherence-direct-image}. Consid\'erons la suite exacte courte
$$
0 \to \sigma_{\geq m - N - 1}\mathcal{I}^\bullet \to
\mathcal{I}^\bullet \to \sigma_{\leq m - N - 2}\mathcal{I}^\bullet \to 0
$$
de complexes de faisceaux quasi-coh\'erents sur $X$. D'apr\`es le lemme
\ref{spaces-perfect-lemma-quasi-coherence-direct-image}, les faisceaux de cohomologie de
$Rf_*\sigma_{\leq m - N - 2}\mathcal{I}^\bullet$
sont nuls en degr\'es $\geq m - 1$. On en d\'eduit que
$R^mf_*\mathcal{I}^\bullet$ est isomorphe \`a
$R^mf_*\sigma_{\geq m - N - 1}\mathcal{I}^\bullet$.
Autrement dit, on peut supposer que $\mathcal{I}^\bullet$ est un complexe
born\'e inf\'erieurement d'objets injectifs de $\QCoh(\mathcal{O}_X)$.
Ce cas r\'esulte du lemme d'acyclicit\'e de Leray (Cat\'egories d\'eriv\'ees,
lemme \ref{derived-lemma-leray-acyclicity}), l'annulation requise \'etant
fournie par le lemme \ref{spaces-perfect-lemma-injective-pushforward}.
\end{proof}

\begin{proposition}
\label{spaces-perfect-proposition-Noetherian}
Soit $S$ un sch\'ema et soit $X$ un espace alg\'ebrique noeth\'erien sur $S$.
Alors le foncteur (\ref{spaces-perfect-equation-compare})
$$
D(\QCoh(\mathcal{O}_X))
\longrightarrow
D_\QCoh(\mathcal{O}_X)
$$
est une \'equivalence dont un quasi-inverse est donn\'e par $RQ_X$.
\end{proposition}

\begin{proof}
Cela r\'esulte imm\'ediatement des lemmes
\ref{spaces-perfect-lemma-Noetherian-pushforward} et \ref{spaces-perfect-lemma-argument-proves}.
\end{proof}













\section{Complexes pseudo-coh\'erents et parfaits}
\label{spaces-perfect-section-spell-out}

\noindent
Dans cette section, nous \'etudions les notions g\'en\'erales d\'efinies dans
Cohomologie sur les sites, sections
\ref{sites-cohomology-section-strictly-perfect},
\ref{sites-cohomology-section-pseudo-coherent},
\ref{sites-cohomology-section-tor}, et
\ref{sites-cohomology-section-perfect}
pour le site \'etale d'un espace alg\'ebrique. Nous les comparons en particulier
aux notions correspondantes pour les sch\'emas.

\medskip\noindent
Comparons tout d'abord la notion de complexe pseudo-coh\'erent sur un sch\'ema
et sur le petit site \'etale qui lui est associ\'e.

\begin{lemma}
\label{spaces-perfect-lemma-descend-finite-type}
Soit $X$ un sch\'ema et soit $\mathcal{F}$ un $\mathcal{O}_X$-module.
Les conditions suivantes sont \'equivalentes :
\begin{enumerate}
\item $\mathcal{F}$ est de type fini comme $\mathcal{O}_X$-module;
\item $\epsilon^*\mathcal{F}$ est de type fini comme
$\mathcal{O}_\etale$-module sur le petit site \'etale de $X$.
\end{enumerate}
Ici $\epsilon$ est comme dans (\ref{spaces-perfect-equation-epsilon}).
\end{lemma}

\begin{proof}
L'implication (1) $\Rightarrow$ (2) est un fait g\'en\'eral; voir Modules sur
les sites, lemme \ref{sites-modules-lemma-local-pullback}.
Supposons (2). Il existe par hypoth\`ese une famille couvrante
$\{f_i : X_i \to X\}$ pour la topologie \'etale telle que
$\epsilon^*\mathcal{F}|_{(X_i)_\etale}$ soit engendr\'e par un nombre fini de
sections. Soit $x \in X$. Nous allons montrer que $\mathcal{F}$ est engendr\'e
par un nombre fini de sections dans un voisinage de $x$.
Choisissons $i$ tel que $x$ appartienne \`a l'image de $X_i \to X$ et posons
$X' = X_i$. Soient
$s_1, \ldots, s_n \in
\Gamma(X', \epsilon^*\mathcal{F}|_{X'_\etale})$
des sections g\'en\'eratrices. Comme
$\epsilon^*\mathcal{F} =
\epsilon^{-1}\mathcal{F} \otimes_{\epsilon^{-1}\mathcal{O}_X}
\mathcal{O}_\etale$
on peut trouver un morphisme \'etale $X'' \to X'$ tel que $x$ appartienne
\`a l'image de $X'' \to X$ et que
$s_i|_{X''} = \sum s_{ij} \otimes a_{ij}$, pour des sections
$s_{ij} \in \epsilon^{-1}\mathcal{F}(X'')$ et
$a_{ij} \in \mathcal{O}_\etale(X'')$. Notons $U \subset X$ l'image de
$X'' \to X$. C'est un sous-sch\'ema ouvert, car $f'' : X'' \to X$ est \'etale
(Morphismes, lemme \ref{morphisms-lemma-etale-open}). Quitte \`a r\'etr\'ecir
encore $X''$, on peut supposer que les $s_{ij}$ proviennent d'\'el\'ements
$t_{ij} \in \mathcal{F}(U)$, d'apr\`es la construction du foncteur image
inverse $\epsilon^{-1}$. Nous affirmons que les $t_{ij}$ engendrent
$\mathcal{F}|_U$, ce qui ach\`eve la preuve du lemme. En effet, l'image inverse
sur $X''$ du morphisme correspondant
$\mathcal{O}_U^{\oplus N} \to \mathcal{F}|_U$ est surjective. Comme
$f'' : X'' \to U$ est un morphisme plat surjectif de sch\'emas, on en d\'eduit
que $\mathcal{O}_U^{\oplus N} \to \mathcal{F}|_U$ est surjectif en passant
aux fibres et en utilisant le fait que
$\mathcal{O}_{U, f''(z)} \to \mathcal{O}_{X'', z}$
est fid\`element plat pour tout $z \in X''$.
\end{proof}

\noindent
Dans la situation ci-dessus, le morphisme de sites $\epsilon$ est plat et
d\'efinit donc une image inverse sur les complexes de modules.

\begin{lemma}
\label{spaces-perfect-lemma-descend-pseudo-coherent}
Soit $X$ un sch\'ema et soit $E$ un objet de $D(\mathcal{O}_X)$.
Les conditions suivantes sont \'equivalentes :
\begin{enumerate}
\item $E$ est $m$-pseudo-coh\'erent;
\item $\epsilon^*E$ est $m$-pseudo-coh\'erent sur le petit site \'etale de $X$.
\end{enumerate}
Ici $\epsilon$ est comme dans (\ref{spaces-perfect-equation-epsilon}).
\end{lemma}

\begin{proof}
L'implication (1) $\Rightarrow$ (2) est un fait g\'en\'eral; voir Cohomologie
sur les sites, lemme
\ref{sites-cohomology-lemma-pseudo-coherent-pullback}.
Supposons $\epsilon^*E$ $m$-pseudo-coh\'erent. Nous utiliserons sans plus le
mentionner le fait que $\epsilon^*$ est un foncteur exact et que, par suite,
$$
\epsilon^*H^i(E) = H^i(\epsilon^*E).
$$
Pour montrer que $E$ est $m$-pseudo-coh\'erent, on peut travailler localement
sur $X$ et donc supposer $X$ quasi-compact (par exemple affine).
Comme $X$ est quasi-compact, toute famille couvrante
$\{U_i \to X\}$ pour la topologie \'etale admet un raffinement fini. Ainsi
$\epsilon^*E$ est un objet de
$D^{-}(\mathcal{O}_\etale)$; voir les commentaires qui suivent Cohomologie
sur les sites, d\'efinition
\ref{sites-cohomology-definition-pseudo-coherent}.
D'apr\`es le lemme \ref{spaces-perfect-lemma-epsilon-flat}, il s'ensuit que $E$ est un objet
de $D^-(\mathcal{O}_X)$.

\medskip\noindent
Soit $n \in \mathbf{Z}$ le plus grand entier tel que $H^n(E)$ soit non nul;
c'est aussi le plus grand entier tel que $H^n(\epsilon^*E)$ soit non nul.
Nous d\'emontrerons le lemme par r\'ecurrence sur $n-m$.
Si $n < m$, le lemme est clair. Si $n \geq m$, alors
$H^n(\epsilon^*E)$ est un $\mathcal{O}_\etale$-module de type fini; voir
Cohomologie sur les sites, lemme
\ref{sites-cohomology-lemma-finite-cohomology}.
Le $\mathcal{O}_X$-module $H^n(E)$ est donc de type fini d'apr\`es le lemme
\ref{spaces-perfect-lemma-descend-finite-type}. En rempla\c{c}ant $X$ par les membres d'un
recouvrement ouvert, on peut supposer qu'il existe une surjection
$\mathcal{O}_X^{\oplus t} \to H^n(E)$. Localement sur $X$, on peut la relever
en un morphisme de complexes
$\alpha : \mathcal{O}_X^{\oplus t}[-n] \to E$ (d\'etails omis).
Choisissons un triangle distingu\'e
$$
\mathcal{O}_X^{\oplus t}[-n] \to E \to C \to \mathcal{O}_X^{\oplus t}[-n + 1]
$$
La cohomologie de $C$ est alors nulle en degr\'es $\geq n$. D'autre part, le
complexe $\epsilon^*C$ est $m$-pseudo-coh\'erent; voir Cohomologie sur les
sites, lemme \ref{sites-cohomology-lemma-cone-pseudo-coherent}.
Par l'hypoth\`ese de r\'ecurrence, $C$ est donc $m$-pseudo-coh\'erent. Une
nouvelle application de Cohomologie sur les sites, lemme
\ref{sites-cohomology-lemma-cone-pseudo-coherent}, permet de conclure.
\end{proof}

\begin{lemma}
\label{spaces-perfect-lemma-descend-tor-amplitude}
Soit $X$ un sch\'ema et soit $E$ un objet de $D(\mathcal{O}_X)$. Alors :
\begin{enumerate}
\item $E$ est de Tor-amplitude contenue dans $[a,b]$ si et seulement si
$\epsilon^*E$ est de Tor-amplitude contenue dans $[a,b]$;
\item $E$ est de Tor-dimension finie si et seulement si $\epsilon^*E$ est de
Tor-dimension finie.
\end{enumerate}
Ici $\epsilon$ est comme dans (\ref{spaces-perfect-equation-epsilon}).
\end{lemma}

\begin{proof}
L'implication directe r\'esulte de Cohomologie sur les sites, lemme
\ref{sites-cohomology-lemma-tor-amplitude-pullback}.
R\'eciproquement, supposons $\epsilon^*E$ de Tor-amplitude contenue dans
$[a,b]$.
Soit $\mathcal{F}$ un $\mathcal{O}_X$-module. Comme $\epsilon$ est un
morphisme plat de sites annel\'es (lemme \ref{spaces-perfect-lemma-epsilon-flat}), on a
$$
\epsilon^*(E \otimes^\mathbf{L}_{\mathcal{O}_X} \mathcal{F})
=
\epsilon^*E
\otimes^\mathbf{L}_{\mathcal{O}_\etale}
\epsilon^*\mathcal{F}
$$
L'annulation suppos\'ee des faisceaux de cohomologie du membre de droite
entra\^ine donc, par le lemme \ref{spaces-perfect-lemma-epsilon-flat}, l'annulation voulue
des faisceaux de cohomologie de
$E \otimes^\mathbf{L}_{\mathcal{O}_X} \mathcal{F}$.
\end{proof}

\begin{lemma}
\label{spaces-perfect-lemma-tor-dimension-rel}
Soit $f : X \to Y$ un morphisme de sch\'emas et soit $E$ un objet de
$D(\mathcal{O}_X)$. Alors :
\begin{enumerate}
\item $E$, consid\'er\'e comme objet de $D(f^{-1}\mathcal{O}_Y)$, est
de Tor-amplitude contenue dans $[a,b]$ si et seulement si $\epsilon^*E$,
consid\'er\'e
comme objet de $D(f_{small}^{-1}\mathcal{O}_{Y_\etale})$, est de Tor-amplitude
contenue dans $[a,b]$;
\item $E$, consid\'er\'e comme objet de $D(f^{-1}\mathcal{O}_Y)$, est
localement de Tor-dimension finie si et seulement si $\epsilon^*E$,
consid\'er\'e comme objet de
$D(f_{small}^{-1}\mathcal{O}_{Y_\etale})$, est localement de Tor-dimension
finie.
\end{enumerate}
Ici $\epsilon$ est comme dans (\ref{spaces-perfect-equation-epsilon}).
\end{lemma}

\begin{proof}
L'implication directe de (1) r\'esulte de Cohomologie sur les sites, lemme
\ref{sites-cohomology-lemma-tor-amplitude-pullback}.
Soit $x \in X$ un point et soit $\overline{x}$ un point g\'eom\'etrique
au-dessus de $x$. Posons $y=f(x)$ et notons $\overline{y}$ le point
g\'eom\'etrique de $Y$ provenant de $\overline{x}$. Alors
$(f^{-1}\mathcal{O}_Y)_x = \mathcal{O}_{Y,y}$ (Faisceaux, lemme
\ref{sheaves-lemma-stalk-pullback}) et
$$
(f_{small}^{-1}\mathcal{O}_{Y_\etale})_{\overline{x}} =
\mathcal{O}_{Y_\etale, \overline{y}} =
\mathcal{O}_{Y, y}^{sh}
$$
est le hens\'elis\'e strict (d'apr\`es Cohomologie \'etale, lemmes
\ref{etale-cohomology-lemma-stalk-pullback} et
\ref{etale-cohomology-lemma-describe-etale-local-ring}).
Comme la fibre de $\mathcal{O}_{X_\etale}$ en $\overline{x}$ est
$\mathcal{O}_{X,x}^{sh}$, on obtient
$$
(\epsilon^*E)_{\overline{x}} =
E_x \otimes_{\mathcal{O}_{X, x}}^\mathbf{L} \mathcal{O}_{X, x}^{sh}
$$
par transitivit\'e des images inverses. Si $\epsilon^*E$ est de Tor-amplitude
contenue dans $[a,b]$ comme complexe de
$f_{small}^{-1}\mathcal{O}_{Y_\etale}$-modules, alors
$(\epsilon^*E)_{\overline{x}}$ est de Tor-amplitude contenue dans $[a,b]$
comme
complexe de $\mathcal{O}_{Y,y}^{sh}$-modules (car le passage aux fibres est
une image inverse, d'apr\`es le lemme cit\'e ci-dessus). Par Autour de la
platitude, lemme \ref{flat-lemma-tor-amplitude-up-down-henselization},
la Tor-amplitude de
$(\epsilon^*E)_{\overline{x}}$
comme complexe de $\mathcal{O}_{Y,y}$-modules est contenue dans $[a,b]$.
Comme $\mathcal{O}_{X,x} \to \mathcal{O}_{X,x}^{sh}$ est fid\`element plat
(Autour de l'alg\`ebre, lemme
\ref{more-algebra-lemma-dumb-properties-henselization}) et que
$(\epsilon^*E)_{\overline{x}} =
E_x \otimes_{\mathcal{O}_{X, x}}^\mathbf{L} \mathcal{O}_{X, x}^{sh}$
on peut appliquer Autour de l'alg\`ebre, lemme
\ref{more-algebra-lemma-no-change-tor-amplitude}, et conclure que la
Tor-amplitude de $E_x$, comme complexe de
$\mathcal{O}_{Y,y}$-modules, est contenue
dans $[a,b]$. D'apr\`es Cohomologie, lemme
\ref{cohomology-lemma-tor-amplitude-stalk}, $E$, consid\'er\'e comme objet de
$D(f^{-1}\mathcal{O}_Y)$, est donc de Tor-amplitude contenue dans $[a,b]$.
Ceci d\'emontre l'implication r\'eciproque de (1). L'assertion (2) d\'ecoule
formellement de (1).
\end{proof}

\begin{lemma}
\label{spaces-perfect-lemma-descend-perfect}
Soit $X$ un sch\'ema et soit $E$ un objet de $D(\mathcal{O}_X)$. Alors $E$ est
un objet parfait de $D(\mathcal{O}_X)$ si et seulement si $\epsilon^*E$ est
un objet parfait de $D(\mathcal{O}_\etale)$. Ici $\epsilon$ est comme dans
(\ref{spaces-perfect-equation-epsilon}).
\end{lemma}

\begin{proof}
L'implication directe r\'esulte du r\'esultat g\'en\'eral de Cohomologie sur les
sites, lemme \ref{sites-cohomology-lemma-perfect-pullback}.
Pour la r\'eciproque, on utilise l'\'equivalence de Cohomologie sur les sites,
lemme \ref{sites-cohomology-lemma-perfect}, et les r\'esultats correspondants
pour les complexes pseudo-coh\'erents et les complexes de Tor-dimension finie,
\`a savoir les lemmes \ref{spaces-perfect-lemma-descend-pseudo-coherent} et
\ref{spaces-perfect-lemma-descend-tor-amplitude}. Quelques d\'etails sont omis.
\end{proof}

\begin{lemma}
\label{spaces-perfect-lemma-pseudo-coherent}
Soit $S$ un sch\'ema et soit $X$ un espace alg\'ebrique sur $S$.
Si $E$ est un objet $m$-pseudo-coh\'erent de $D(\mathcal{O}_X)$, alors
$H^i(E)$ est un $\mathcal{O}_X$-module quasi-coh\'erent pour $i>m$.
Si $E$ est pseudo-coh\'erent, alors $E$ est un objet de
$D_\QCoh(\mathcal{O}_X)$.
\end{lemma}

\begin{proof}
Localement, $H^i(E)$ est isomorphe \`a $H^i(\mathcal{E}^\bullet)$, o\`u
$\mathcal{E}^\bullet$ est strictement parfait. Les faisceaux
$\mathcal{E}^i$ sont facteurs directs de modules libres de type fini, donc
quasi-coh\'erents. Le lemme en r\'esulte.
\end{proof}

\begin{lemma}
\label{spaces-perfect-lemma-identify-pseudo-coherent-noetherian}
Soit $S$ un sch\'ema, soit $X$ un espace alg\'ebrique noeth\'erien sur $S$ et
soit $E$ un objet de $D_\QCoh(\mathcal{O}_X)$. Pour
$m \in \mathbf{Z}$, les conditions suivantes sont \'equivalentes :
\begin{enumerate}
\item $H^i(E)$ est coh\'erent pour $i \geq m$ et nul pour $i \gg 0$;
\item $E$ est $m$-pseudo-coh\'erent.
\end{enumerate}
En particulier, $E$ est pseudo-coh\'erent si et seulement si $E$ est un objet
de $D^-_{\textit{Coh}}(\mathcal{O}_X)$.
\end{lemma}

\begin{proof}
Comme $X$ est quasi-compact, on peut trouver un sch\'ema affine $U$ et un
morphisme \'etale surjectif $U \to X$ (Propri\'et\'es des espaces, lemme
\ref{spaces-properties-lemma-quasi-compact-affine-cover}).
Observons que $U$ est noeth\'erien. L'objet $E$ est $m$-pseudo-coh\'erent si et
seulement si $E|_U$ l'est (cela r\'esulte de la d\'efinition ou de Cohomologie
sur les sites, lemme
\ref{sites-cohomology-lemma-pseudo-coherent-independent-representative}).
De m\^eme, $H^i(E)$ est coh\'erent si et seulement si
$H^i(E)|_U=H^i(E|_U)$ est coh\'erent (voir Cohomologie des espaces, lemme
\ref{spaces-cohomology-lemma-coherent-Noetherian}).
On peut donc supposer $X$ repr\'esentable.

\medskip\noindent
Si $X$ est repr\'esent\'e par un sch\'ema $X_0$, on peut (lemme
\ref{spaces-perfect-lemma-derived-quasi-coherent-small-etale-site}) \'ecrire
$E=\epsilon^*E_0$, o\`u $E_0$ est un objet de
$D_\QCoh(\mathcal{O}_{X_0})$ et o\`u
$\epsilon : X_\etale \to (X_0)_{Zar}$ est comme dans
(\ref{spaces-perfect-equation-epsilon}). Dans ce cas, le lemme
\ref{spaces-perfect-lemma-descend-pseudo-coherent} montre que $E$ est
$m$-pseudo-coh\'erent si et seulement si $E_0$ l'est. De m\^eme, d'apr\`es le
lemme \ref{spaces-perfect-lemma-descend-finite-type}, $H^i(E_0)$ est de type fini (c'est-\`a-dire
coh\'erent) si et seulement si $H^i(E)$ l'est. Enfin, le lemme
\ref{spaces-perfect-lemma-epsilon-flat} montre que $H^i(E_0)=0$ si et seulement si
$H^i(E)=0$. On est donc ramen\'e au cas des sch\'emas, qui est le lemme de
Cat\'egories d\'eriv\'ees des sch\'emas
\ref{perfect-lemma-identify-pseudo-coherent-noetherian}.
\end{proof}

\begin{lemma}
\label{spaces-perfect-lemma-tor-qc-qs}
Soit $S$ un sch\'ema, soit $X$ un espace alg\'ebrique quasi-s\'epar\'e sur $S$,
soit $E$ un objet de $D_\QCoh(\mathcal{O}_X)$ et soient $a \leq b$.
Les conditions suivantes sont \'equivalentes :
\begin{enumerate}
\item $E$ est de Tor-amplitude contenue dans $[a,b]$;
\item pour tout $\mathcal{F}$ de $\QCoh(\mathcal{O}_X)$, on a
$H^i(E \otimes_{\mathcal{O}_X}^\mathbf{L} \mathcal{F})=0$ pour
$i \notin [a,b]$.
\end{enumerate}
\end{lemma}

\begin{proof}
Il est clair que (1) implique (2). Supposons (2). Soit $j : U \to X$ un
morphisme \'etale avec $U$ affine. Comme $X$ est quasi-s\'epar\'e, $j$ est \`a
la fois quasi-compact et s\'epar\'e. Par cons\'equent, $j_*$ transforme les
modules quasi-coh\'erents en modules quasi-coh\'erents (Morphismes d'espaces,
lemme
\ref{spaces-morphisms-lemma-pushforward}).
Soit $\mathcal{G}$ un module quasi-coh\'erent quelconque sur $U$.
Par hypoth\`ese,
$H^i(E \otimes_{\mathcal{O}_X}^\mathbf{L} j_*\mathcal{G})$ est nul pour
$i \notin [a,b]$. En prenant l'image inverse par le morphisme plat $j$, on
voit que
$H^i(E|_U \otimes_{\mathcal{O}_U}^\mathbf{L} j^*j_*\mathcal{G})$ est nul
pour $i \notin [a,b]$. D'apr\`es Cohomologie des espaces, lemme
\ref{spaces-cohomology-lemma-etale-pull-push-split}, $\mathcal{G}$ est un
facteur direct de $j^*j_*\mathcal{G}$. La condition (2) entra\^ine donc
l'annulation de
$H^i(E|_U \otimes_{\mathcal{O}_U}^\mathbf{L} \mathcal{G})$
pour $i \notin [a,b]$ et pour tout $\mathcal{O}_U$-module quasi-coh\'erent
$\mathcal{G}$. Puisqu'il suffit de d\'emontrer que $E|_U$ est de Tor-amplitude
contenue dans $[a,b]$, on est ramen\'e au cas o\`u $X$ est repr\'esentable.

\medskip\noindent
Si $X$ est repr\'esent\'e par un sch\'ema $X_0$, on peut (lemme
\ref{spaces-perfect-lemma-derived-quasi-coherent-small-etale-site}) \'ecrire
$E=\epsilon^*E_0$, o\`u $E_0$ est un objet de
$D_\QCoh(\mathcal{O}_{X_0})$ et o\`u
$\epsilon : X_\etale \to (X_0)_{Zar}$ est comme dans
(\ref{spaces-perfect-equation-epsilon}). Pour tout module quasi-coh\'erent
$\mathcal{F}_0$ sur $X_0$, le module $\epsilon^*\mathcal{F}_0$ est
quasi-coh\'erent sur $X$ et
$$
H^i(E \otimes_{\mathcal{O}_X}^\mathbf{L} \epsilon^*\mathcal{F}_0)
=
\epsilon^*H^i(E_0 \otimes_{\mathcal{O}_{X_0}}^\mathbf{L} \mathcal{F}_0)
$$
car $\epsilon$ est plat (lemme \ref{spaces-perfect-lemma-epsilon-flat}). En outre,
l'annulation de ces faisceaux pour $i \notin [a,b]$ entra\^ine celle de
$H^i(E_0 \otimes_{\mathcal{O}_{X_0}}^\mathbf{L} \mathcal{F}_0)$
par le m\^eme lemme. Le probl\`eme est donc ramen\'e au cas des sch\'emas,
trait\'e dans Cat\'egories d\'eriv\'ees des sch\'emas, lemme
\ref{perfect-lemma-tor-qc-qs}.
\end{proof}

\begin{lemma}
\label{spaces-perfect-lemma-descend-RHom}
Soit $X$ un sch\'ema et soient $E,F$ des objets de $D(\mathcal{O}_X)$.
Supposons que l'une des conditions suivantes soit satisfaite :
\begin{enumerate}
\item $E$ est pseudo-coh\'erent et $F$ appartient \`a
$D^+(\mathcal{O}_X)$;
\item $E$ est parfait et $F$ est quelconque.
\end{enumerate}
Alors il existe un isomorphisme canonique
$$
\epsilon^*R\SheafHom(E, F) \longrightarrow R\SheafHom(\epsilon^*E, \epsilon^*F)
$$
Ici $\epsilon$ est comme dans (\ref{spaces-perfect-equation-epsilon}).
\end{lemma}

\begin{proof}
Rappelons que $\epsilon$ est plat (lemme \ref{spaces-perfect-lemma-epsilon-flat}), donc que
$\epsilon^*=L\epsilon^*$. Il existe un morphisme canonique du membre de gauche
vers celui de droite d'apr\`es Cohomologie sur les sites, remarque
\ref{sites-cohomology-remark-prepare-fancy-base-change}.
Pour voir que c'est un isomorphisme, on peut travailler localement, c'est-\`a-dire
supposer que $X$ est un sch\'ema affine.

\medskip\noindent
Dans le cas (1), on peut repr\'esenter $E$ par un complexe born\'e
sup\'erieurement $\mathcal{E}^\bullet$ de $\mathcal{O}_X$-modules libres de
type fini; voir Cat\'egories d\'eriv\'ees des sch\'emas, lemme
\ref{perfect-lemma-lift-pseudo-coherent}. On peut aussi repr\'esenter $F$ par
un complexe born\'e inf\'erieurement $\mathcal{F}^\bullet$ de
$\mathcal{O}_X$-modules. En appliquant Cohomologie, lemme
\ref{cohomology-lemma-Rhom-complex-of-direct-summands-finite-free}
on voit que $R\SheafHom(E,F)$ est repr\'esent\'e par le complexe de termes
$$
\bigoplus\nolimits_{n = - p + q}
\SheafHom_{\mathcal{O}_X}(\mathcal{E}^p, \mathcal{F}^q)
$$
En appliquant Cohomologie sur les sites, lemme
\ref{sites-cohomology-lemma-Rhom-complex-of-direct-summands-finite-free}
on voit que $R\SheafHom(\epsilon^*E,\epsilon^*F)$ est repr\'esent\'e par le
complexe de termes
$$
\bigoplus\nolimits_{n = - p + q}
\SheafHom_{\mathcal{O}_\etale}
(\epsilon^*\mathcal{E}^p, \epsilon^*\mathcal{F}^q)
$$
L'assertion du lemme se ram\`ene donc au fait, vrai, que le morphisme canonique
$$
\epsilon^*\SheafHom_{\mathcal{O}_X}(\mathcal{E}, \mathcal{F})
\longrightarrow
\SheafHom_{\mathcal{O}_\etale}
(\epsilon^*\mathcal{E}, \epsilon^*\mathcal{F})
$$
est un isomorphisme pour tout $\mathcal{O}_X$-module $\mathcal{F}$ et tout
$\mathcal{O}_X$-module libre de type fini $\mathcal{E}$.

\medskip\noindent
Dans le cas (2), on peut repr\'esenter $E$ par un complexe strictement parfait
$\mathcal{E}^\bullet$ de $\mathcal{O}_X$-modules; on utilise Cat\'egories
d\'eriv\'ees des sch\'emas, lemmes
\ref{perfect-lemma-affine-compare-bounded} et
\ref{perfect-lemma-perfect-affine}, ainsi que le fait qu'un complexe parfait
de modules est repr\'esent\'e par un complexe fini de modules projectifs de
type fini. On peut donc reprendre exactement la d\'emonstration ci-dessus en
rempla\c{c}ant la r\'ef\'erence \`a Cohomologie, lemme
\ref{cohomology-lemma-Rhom-complex-of-direct-summands-finite-free}
par une r\'ef\'erence \`a Cohomologie, lemme
\ref{cohomology-lemma-Rhom-strictly-perfect}.
\end{proof}

\begin{lemma}
\label{spaces-perfect-lemma-quasi-coherence-internal-hom}
Soit $S$ un sch\'ema, soit $X$ un espace alg\'ebrique sur $S$ et soient $L,K$
des objets de $D(\mathcal{O}_X)$. Si l'une des conditions suivantes est
satisfaite :
\begin{enumerate}
\item $L$ appartient \`a $D^+_\QCoh(\mathcal{O}_X)$ et $K$ est
pseudo-coh\'erent;
\item $L$ appartient \`a $D_\QCoh(\mathcal{O}_X)$ et $K$ est parfait;
\end{enumerate}
alors $R\SheafHom(K,L)$ appartient \`a $D_\QCoh(\mathcal{O}_X)$.
\end{lemma}

\begin{proof}
Cela r\'esulte de l'analogue pour les sch\'emas (Cat\'egories d\'eriv\'ees des
sch\'emas, lemme \ref{perfect-lemma-quasi-coherence-internal-hom}), du crit\`ere
du lemme \ref{spaces-perfect-lemma-check-quasi-coherence-on-covering}, des crit\`eres des
lemmes \ref{spaces-perfect-lemma-descend-pseudo-coherent} et \ref{spaces-perfect-lemma-descend-perfect},
et du r\'esultat du lemme \ref{spaces-perfect-lemma-descend-RHom}.
\end{proof}

\begin{lemma}
\label{spaces-perfect-lemma-internal-hom-evaluate-tensor-isomorphism}
Soit $S$ un sch\'ema, soit $X$ un espace alg\'ebrique sur $S$ et soient
$K,L,M$ des objets de $D_\QCoh(\mathcal{O}_X)$. Le morphisme
$$
K \otimes_{\mathcal{O}_X}^\mathbf{L} R\SheafHom(M, L)
\longrightarrow
R\SheafHom(M, K \otimes_{\mathcal{O}_X}^\mathbf{L} L)
$$
de Cohomologie sur les sites, lemme
\ref{sites-cohomology-lemma-internal-hom-diagonal-better}
est un isomorphisme dans les cas suivants :
\begin{enumerate}
\item $M$ est parfait;
\item $K$ est parfait;
\item $M$ est pseudo-coh\'erent, $L \in D^+(\mathcal{O}_X)$ et $K$ est de
Tor-dimension finie.
\end{enumerate}
\end{lemma}

\begin{proof}
La propri\'et\'e pour ce morphisme d'\^etre un isomorphisme se v\'erifie
localement pour la topologie \'etale; on peut donc supposer que $X$ est un
sch\'ema affine. D'apr\`es le lemme
\ref{spaces-perfect-lemma-derived-quasi-coherent-small-etale-site}, on peut alors \'ecrire
$K,L,M$ sous la forme $\epsilon^*K_0,\epsilon^*L_0,\epsilon^*M_0$, pour des
$K_0,L_0,M_0$ de $D_\QCoh(\mathcal{O}_X)$. Le lemme
\ref{spaces-perfect-lemma-descend-RHom} ram\`ene les cas (1) et (3) au cas des sch\'emas,
qui est Cat\'egories d\'eriv\'ees des sch\'emas, lemme
\ref{perfect-lemma-internal-hom-evaluate-tensor-isomorphism}.
Si $K$ est parfait mais qu'aucune autre hypoth\`ese n'est faite, on ne sait pas
si l'un ou l'autre membre de la fl\`eche appartient \`a
$D_\QCoh(\mathcal{O}_X)$; le r\'esultat reste toutefois vrai, car, apr\`es une
localisation suppl\'ementaire, $K$ sera repr\'esent\'e par un complexe fini de
modules libres de type fini, auquel cas l'assertion est claire.
\end{proof}









\section{Approximation par des complexes parfaits}
\label{spaces-perfect-section-approximation}

\noindent
Dans cette section, nous poursuivons l'\'etude commenc\'ee dans Cat\'egories
d\'eriv\'ees des sch\'emas, section \ref{perfect-section-approximation}.

\begin{definition}
\label{spaces-perfect-definition-approximation-holds}
Soit $S$ un sch\'ema et soit $X$ un espace alg\'ebrique sur $S$.
Consid\'erons les triplets $(T,E,m)$ tels que
\begin{enumerate}
\item $T \subset |X|$ est une partie ferm\'ee;
\item $E$ est un objet de $D_\QCoh(\mathcal{O}_X)$;
\item $m \in \mathbf{Z}$.
\end{enumerate}
On dit que {\it la propri\'et\'e d'approximation est satisfaite pour le triplet}
$(T,E,m)$ s'il existe un objet parfait $P$ de $D(\mathcal{O}_X)$ \`a support
dans $T$ et un morphisme $\alpha : P \to E$ qui induit des isomorphismes
$H^i(P) \to H^i(E)$ pour $i>m$ et une surjection
$H^m(P) \to H^m(E)$.
\end{definition}

\noindent
La propri\'et\'e d'approximation ne peut pas \^etre satisfaite pour tout
triplet. Les remarques qui suivent Cat\'egories d\'eriv\'ees des sch\'emas,
d\'efinition \ref{perfect-definition-approximation-holds}, en expliquent la
raison.

\begin{definition}
\label{spaces-perfect-definition-approximation}
Soit $S$ un sch\'ema et soit $X$ un espace alg\'ebrique sur $S$.
On dit que {\it la propri\'et\'e d'approximation par des complexes parfaits est
satisfaite} sur $X$ si, pour toute partie ferm\'ee $T \subset |X|$ telle que le
morphisme $X \setminus T \to X$ soit quasi-compact, il existe un entier $r$
tel que, pour tout triplet $(T,E,m)$ comme dans la d\'efinition
\ref{spaces-perfect-definition-approximation-holds} et v\'erifiant
\begin{enumerate}
\item $E$ est $(m-r)$-pseudo-coh\'erent;
\item $H^i(E)$ est \`a support dans $T$ pour $i \geq m-r$,
\end{enumerate}
la propri\'et\'e d'approximation soit satisfaite.
\end{definition}

\begin{lemma}
\label{spaces-perfect-lemma-pushforward-perfect}
Soit $S$ un sch\'ema et soit $(U \subset X, j : V \to X)$ un carr\'e
distingu\'e \'el\'ementaire d'espaces alg\'ebriques sur $S$.
Soit $E$ un objet parfait de $D(\mathcal{O}_V)$ \`a support dans
$j^{-1}(T)$, o\`u $T=|X|\setminus|U|$. Alors $Rj_*E$ est un objet parfait de
$D(\mathcal{O}_X)$.
\end{lemma}

\begin{proof}
La propri\'et\'e d'\^etre parfait est locale sur $X_\etale$. Il suffit donc de
v\'erifier que les restrictions de $Rj_*E$ \`a $U$ et \`a $V$ sont parfaites.
On a $Rj_*E|_V=E$ d'apr\`es le lemme
\ref{spaces-perfect-lemma-pushforward-with-support-in-open}, et cet objet est parfait.
On a $Rj_*E|_U=0$, car $E|_{V\setminus j^{-1}(T)}=0$ (utiliser le lemme
\ref{spaces-perfect-lemma-restrict-direct-image-open}).
\end{proof}

\begin{lemma}
\label{spaces-perfect-lemma-open}
Soit $S$ un sch\'ema et soit $(U \subset X,j : V \to X)$ un carr\'e distingu\'e
\'el\'ementaire d'espaces alg\'ebriques sur $S$. Soit $T$ une partie ferm\'ee de
$|X|\setminus|U|$ et soit $(T,E,m)$ un triplet comme dans la d\'efinition
\ref{spaces-perfect-definition-approximation-holds}. Si
\begin{enumerate}
\item la propri\'et\'e d'approximation est satisfaite pour
$(j^{-1}T,E|_V,m)$;
\item les faisceaux $H^i(E)$ sont \`a support dans $T$ pour $i\geq m$,
\end{enumerate}
alors la propri\'et\'e d'approximation est satisfaite pour $(T,E,m)$.
\end{lemma}

\begin{proof}
Soit $P \to E|_V$ une approximation du triplet $(j^{-1}T,E|_V,m)$ sur $V$.
D'apr\`es le lemme \ref{spaces-perfect-lemma-pushforward-perfect}, $Rj_*P$ est un objet
parfait de $D(\mathcal{O}_X)$. D'autre part, le lemme
\ref{spaces-perfect-lemma-pushforward-with-support-in-open} donne $Rj_*P=j_!P$.
L'objet $j_!P$ est \`a support dans $T$, comme le montre par exemple
(\ref{spaces-perfect-equation-stalk-j-shriek}). On obtient donc une approximation
$Rj_*P=j_!P \to j_!(E|_V) \to E$.
\end{proof}

\begin{lemma}
\label{spaces-perfect-lemma-approximation-affine}
Soit $S$ un sch\'ema et soit $X$ un espace alg\'ebrique sur $S$ repr\'esentable
par un sch\'ema affine. Alors la propri\'et\'e d'approximation est satisfaite
pour tout triplet $(T,E,m)$ comme dans la d\'efinition
\ref{spaces-perfect-definition-approximation-holds} pour lequel il existe un entier
$r\geq0$ tel que
\begin{enumerate}
\item $E$ est $m$-pseudo-coh\'erent;
\item $H^i(E)$ est \`a support dans $T$ pour $i\geq m-r+1$;
\item $X\setminus T$ est la r\'eunion de $r$ ouverts affines.
\end{enumerate}
En particulier, la propri\'et\'e d'approximation par des complexes parfaits est
satisfaite pour les sch\'emas affines.
\end{lemma}

\begin{proof}
Soit $X_0$ un sch\'ema affine repr\'esentant $X$ et soit $T_0\subset X_0$ la
partie ferm\'ee correspondant \`a $T$. Soit
$\epsilon : X_\etale \to X_{0,Zar}$ le morphisme
(\ref{spaces-perfect-equation-epsilon}). On peut \'ecrire $E=\epsilon^*E_0$ pour un objet
$E_0$ de $D_\QCoh(\mathcal{O}_{X_0})$; voir le lemme
\ref{spaces-perfect-lemma-derived-quasi-coherent-small-etale-site}.
Le lemme \ref{spaces-perfect-lemma-descend-pseudo-coherent} montre que $E_0$ est
$m$-pseudo-coh\'erent. En comparant les fibres des faisceaux de cohomologie
(voir la preuve du lemme \ref{spaces-perfect-lemma-epsilon-flat}), on voit que $H^i(E_0)$ est
\`a support dans $T_0$ pour $i\geq m-r+1$. D'apr\`es Cat\'egories d\'eriv\'ees
des sch\'emas, lemme \ref{perfect-lemma-approximation-affine}, il existe une
approximation $P_0\to E_0$ de $(T_0,E_0,m)$. Le lemme
\ref{spaces-perfect-lemma-descend-perfect} montre que $P=\epsilon^*P_0$ est un objet parfait
de $D(\mathcal{O}_X)$. En prenant l'image inverse, on obtient l'approximation
voulue $P=\epsilon^*P_0\to\epsilon^*E_0=E$.
\end{proof}

\begin{lemma}
\label{spaces-perfect-lemma-induction-step}
Soit $S$ un sch\'ema et soit $(U \subset X,j : V \to X)$ un carr\'e distingu\'e
\'el\'ementaire d'espaces alg\'ebriques sur $S$. Supposons $U$ quasi-compact,
$V$ affine et $U\times_XV$ quasi-compact. Si la propri\'et\'e d'approximation
par des complexes parfaits est satisfaite sur $U$, elle l'est sur $X$.
\end{lemma}

\begin{proof}
Soit $T\subset|X|$ une partie ferm\'ee telle que
$X\setminus T\to X$ soit quasi-compact. Soit $r_U$ l'entier de la d\'efinition
\ref{spaces-perfect-definition-approximation} associ\'e au couple $(U,T\cap|U|)$.
Posons $T'=T\setminus|U|$ et munissons $T'$ de la structure r\'eduite induite
de sous-espace. Comme $|T'|$ est contenu dans $|X|\setminus|U|$, le morphisme
$j^{-1}(T')\to T'$ est un isomorphisme. De plus,
$V\setminus j^{-1}(T')$ est quasi-compact : c'est le produit fibr\'e de
$U\times_XV$ et de $X\setminus T$ au-dessus de $X$, et nous avons suppos\'e
$U\times_XV$ quasi-compact ainsi que $X\setminus T\to X$ quasi-compact.
Soit $r'$ le nombre d'ouverts affines n\'ecessaires pour recouvrir
$V\setminus j^{-1}(T')$. Nous affirmons que $r=\max(r_U,r')$ convient au
couple $(X,T)$.

\medskip\noindent
Pour le voir, choisissons un triplet $(T,E,m)$ tel que $E$ soit
$(m-r)$-pseudo-coh\'erent et que $H^i(E)$ soit \`a support dans $T$ pour
$i\geq m-r$. Soit $t$ le plus grand entier tel que $H^t(E)|_U$ soit non nul.
(Un tel entier existe, car $U$ est quasi-compact et $E|_U$ est
$(m-r)$-pseudo-coh\'erent.) Nous montrerons par r\'ecurrence sur $t$ que $E$
admet une approximation.

\medskip\noindent
Initialisation : $t\leq m-r'$. Cela signifie que $H^i(E)$ est \`a support
dans $T'$ pour $i\geq m-r'$. Le lemme
\ref{spaces-perfect-lemma-approximation-affine} assure donc l'existence sur $V$ d'une
approximation $P\to E|_V$ de $(T',E|_V,m)$. En appliquant le lemme
\ref{spaces-perfect-lemma-open}, on voit que $(T',E,m)$ admet une approximation. Celle-ci est
aussi une approximation de $(T,E,m)$.

\medskip\noindent
H\'er\'edit\'e. Choisissons une approximation $P\to E|_U$ de
$(T\cap|U|,E|_U,m)$. Elle fournit en particulier une surjection
$H^t(P)\to H^t(E|_U)$. Dans la suite de la preuve, nous utiliserons
l'\'equivalence du lemme
\ref{spaces-perfect-lemma-derived-quasi-coherent-small-etale-site} (et les compatibilit\'es de
la remarque \ref{spaces-perfect-remark-match-total-direct-images}) pour les espaces
alg\'ebriques repr\'esentables $V$ et $U\times_XV$. Nous utiliserons aussi le
fait que les notions de $(m-r)$-pseudo-coh\'erence, respectivement de
perfection, co\"incident sur les sites de Zariski et \'etale; voir les lemmes
\ref{spaces-perfect-lemma-descend-pseudo-coherent} et \ref{spaces-perfect-lemma-descend-perfect}.
On peut donc appliquer les r\'esultats de Cat\'egories d\'eriv\'ees des sch\'emas,
section \ref{perfect-section-lift}, \`a l'immersion ouverte
$U\times_XV\subset V$. Ainsi, Cat\'egories d\'eriv\'ees des sch\'emas, lemme
\ref{perfect-lemma-lift-perfect-complex-plus-shift-support}, montre qu'il
existe un objet parfait $Q$ de $D(\mathcal{O}_V)$, \`a support dans
$j^{-1}(T)$, et un isomorphisme
$Q|_{U\times_XV}\to(P\oplus P[1])|_{U\times_XV}$. D'apr\`es Cat\'egories
d\'eriv\'ees des sch\'emas, lemme \ref{perfect-lemma-lift-map}, on peut
remplacer $Q$ par $Q\otimes^\mathbf{L}I$ et supposer que le morphisme
$$
Q|_{U \times_X V} \longrightarrow
(P \oplus P[1])|_{U \times_X V} \longrightarrow
P|_{U \times_X V} \longrightarrow E|_{U \times_X V}
$$
se rel\`eve en $Q\to E|_V$. Le lemme \ref{spaces-perfect-lemma-glue} fournit alors un
morphisme $a:R\to E$ de $D(\mathcal{O}_X)$ tel que $a|_U$ soit isomorphe \`a
$P\oplus P[1]\to E|_U$ et que $a|_V$ soit isomorphe \`a $Q\to E|_V$.
Ainsi $R$ est parfait et \`a support dans $T$, et le morphisme
$H^t(R)\to H^t(E)$ est surjectif apr\`es restriction \`a $U$.
Choisissons un triangle distingu\'e
$$
R \to E \to E' \to R[1]
$$
Alors $E'$ est $(m-r)$-pseudo-coh\'erent (Cohomologie sur les sites, lemme
\ref{sites-cohomology-lemma-cone-pseudo-coherent}), $H^i(E')|_U=0$ pour
$i\geq t$, et $H^i(E')$ est \`a support dans $T$ pour $i\geq m-r$.
Par l'hypoth\`ese de r\'ecurrence, on trouve une approximation $R'\to E'$ de
$(T,E',m)$. Ins\'erons le compos\'e $R'\to E'\to R[1]$ dans un triangle
distingu\'e $R\to R''\to R'\to R[1]$, puis prolongeons les morphismes
$R'\to E'$ et $R[1]\to R[1]$ en un morphisme de triangles distingu\'es
$$
\xymatrix{
R \ar[r] \ar[d] & R'' \ar[d] \ar[r] & R' \ar[d] \ar[r] & R[1] \ar[d] \\
R \ar[r] &  E \ar[r] & E' \ar[r] & R[1]
}
$$
en utilisant TR3. Alors $R''$ est un complexe parfait (Cohomologie sur les
sites, lemme \ref{sites-cohomology-lemma-two-out-of-three-perfect}) \`a support
dans $T$. Une chasse au diagramme imm\'ediate montre que $R''\to E$ est
l'approximation cherch\'ee.
\end{proof}

\begin{theorem}
\label{spaces-perfect-theorem-approximation}
Soit $S$ un sch\'ema et soit $X$ un espace alg\'ebrique quasi-compact et
quasi-s\'epar\'e sur $S$. Alors la propri\'et\'e d'approximation par des
complexes parfaits est satisfaite sur $X$.
\end{theorem}

\begin{proof}
Cela r\'esulte du principe de r\'ecurrence du lemme
\ref{spaces-perfect-lemma-induction-principle} et des lemmes \ref{spaces-perfect-lemma-induction-step} et
\ref{spaces-perfect-lemma-approximation-affine}.
\end{proof}






\section{Engendrement des cat\'egories d\'eriv\'ees}
\label{spaces-perfect-section-generating}

\noindent
Cette section est l'analogue de Cat\'egories d\'eriv\'ees des sch\'emas,
section \ref{perfect-section-generating}. Nous commen\c{c}ons toutefois par
d\'emontrer le lemme suivant, analogue de Cat\'egories d\'eriv\'ees des
sch\'emas, lemme
\ref{perfect-lemma-lift-map-from-perfect-complex-with-support}.

\begin{lemma}
\label{spaces-perfect-lemma-lift-map-from-perfect-complex-with-support}
Soit $S$ un sch\'ema. Soit $X$ un espace alg\'ebrique quasi-compact et
quasi-s\'epar\'e sur $S$. Soit $W\subset X$ un ouvert quasi-compact. Soit
$T\subset|X|$ une partie ferm\'ee telle que le morphisme
$X\setminus T\to X$ soit quasi-compact. Soit $E$ un objet de
$D_\QCoh(\mathcal{O}_X)$. Soit $\alpha:P\to E|_W$ un morphisme, o\`u $P$ est
un objet parfait de $D(\mathcal{O}_W)$ \`a support dans $T\cap W$. Il existe
alors un morphisme $\beta:R\to E$, o\`u $R$ est un objet parfait de
$D(\mathcal{O}_X)$ \`a support dans $T$, tel que $P$ soit un facteur direct
de $R|_W$ dans $D(\mathcal{O}_W)$, de fa\c{c}on compatible avec $\alpha$ et
$\beta|_W$.
\end{lemma}

\begin{proof}
Nous allons utiliser le principe de r\'ecurrence du lemme
\ref{spaces-perfect-lemma-induction-principle-enlarge}. On se ram\`ene donc imm\'ediatement
au cas d'un carr\'e distingu\'e \'el\'ementaire
$(W\subset X,f:V\to X)$, avec $V$ affine et $P\to E|_W$ comme dans l'\'enonc\'e
du lemme. Dans la suite de la d\'emonstration, nous utiliserons le lemme
\ref{spaces-perfect-lemma-derived-quasi-coherent-small-etale-site} (ainsi que les
compatibilit\'es de la remarque \ref{spaces-perfect-remark-match-total-direct-images}) pour
les espaces alg\'ebriques repr\'esentables $V$ et $W\times_XV$. Nous
utiliserons aussi le fait que les notions d'objet parfait sur les sites de
Zariski et \'etale co\"incident; voir le lemme \ref{spaces-perfect-lemma-descend-perfect}.

\medskip\noindent
D'apr\`es Cat\'egories d\'eriv\'ees des sch\'emas, lemme
\ref{perfect-lemma-lift-perfect-complex-plus-shift-support}, on peut choisir un
objet parfait $Q$ de $D(\mathcal{O}_V)$, \`a support dans $f^{-1}T$, et un
isomorphisme
$Q|_{W\times_XV}\to(P\oplus P[1])|_{W\times_XV}$. D'apr\`es Cat\'egories
d\'eriv\'ees des sch\'emas, lemme \ref{perfect-lemma-lift-map}, on peut
remplacer $Q$ par $Q\otimes^\mathbf{L}I$ (qui est encore \`a support dans
$f^{-1}T$) et supposer que le morphisme
$$
Q|_{W \times_X V} \to (P \oplus P[1])|_{W \times_X V}
\longrightarrow P|_{W \times_X V}
\longrightarrow
E|_{W \times_X V}
$$
se rel\`eve en $Q\to E|_V$. Le lemme \ref{spaces-perfect-lemma-glue} fournit un morphisme
$a:R\to E$ de $D(\mathcal{O}_X)$ tel que $a|_W$ soit isomorphe \`a
$P\oplus P[1]\to E|_W$ et que $a|_V$ soit isomorphe \`a $Q\to E|_V$.
Ainsi, $R$ est parfait et \`a support dans $T$, comme souhait\'e.
\end{proof}

\begin{remark}
\label{spaces-perfect-remark-addendum}
La d\'emonstration du lemme
\ref{spaces-perfect-lemma-lift-map-from-perfect-complex-with-support} montre que
$$
R|_W = P \oplus P^{\oplus n_1}[1] \oplus \ldots \oplus P^{\oplus n_m}[m]
$$
pour certains $m\geq0$ et $n_j\geq0$. Le faisceau de cohomologie de plus haut
degr\'e de $R|_W$ est donc \`egal \`a celui de $P$. En r\'ep\'etant la
construction pour le morphisme
$P^{\oplus n_1}[1]\oplus\ldots\oplus P^{\oplus n_m}[m]\to R|_W$, en prenant
des c\^ones et en raisonnant par r\'ecurrence, on peut obtenir l'\'egalit\'e des
faisceaux de cohomologie de $R|_W$ et de $P$ au-dessus de tout degr\'e fix\'e.
\end{remark}

\begin{lemma}
\label{spaces-perfect-lemma-direct-summand-of-a-restriction}
Soit $S$ un sch\'ema. Soit $X$ un espace alg\'ebrique quasi-compact et
quasi-s\'epar\'e sur $S$. Soit $W$ un sous-espace ouvert quasi-compact de $X$.
Soit $P$ un objet parfait de $D(\mathcal{O}_W)$. Alors $P$ est un facteur
direct de la restriction d'un objet parfait de $D(\mathcal{O}_X)$.
\end{lemma}

\begin{proof}
C'est un cas particulier du lemme
\ref{spaces-perfect-lemma-lift-map-from-perfect-complex-with-support}.
\end{proof}

\begin{theorem}
\label{spaces-perfect-theorem-bondal-van-den-Bergh}
Soit $S$ un sch\'ema. Soit $X$ un espace alg\'ebrique quasi-compact et
quasi-s\'epar\'e sur $S$. La cat\'egorie $D_\QCoh(\mathcal{O}_X)$ peut \^etre
engendr\'ee par un seul objet parfait. Plus pr\'ecis\'ement, il existe un objet
parfait $P$ de $D(\mathcal{O}_X)$ tel que, pour
$E\in D_\QCoh(\mathcal{O}_X)$, les conditions suivantes soient \'equivalentes :
\begin{enumerate}
\item $E = 0$;
\item $\Hom_{D(\mathcal{O}_X)}(P[n], E) = 0$ pour tout $n \in \mathbf{Z}$.
\end{enumerate}
\end{theorem}

\begin{proof}
Nous allons le d\'emontrer au moyen du principe de r\'ecurrence du lemme
\ref{spaces-perfect-lemma-induction-principle}.

\medskip\noindent
Si $X$ est affine, alors $\mathcal{O}_X$ est un g\'en\'erateur parfait. Cela
r\'esulte du lemme \ref{spaces-perfect-lemma-derived-quasi-coherent-small-etale-site} et de
Cat\'egories d\'eriv\'ees des sch\'emas, lemme
\ref{perfect-lemma-affine-compare-bounded}.

\medskip\noindent
Supposons que $(U\subset X,f:V\to X)$ soit un carr\'e distingu\'e
\'el\'ementaire, avec $U$ quasi-compact, tel que le th\'eor\`eme soit vrai pour
$U$ et que $V$ soit un sch\'ema affine. Soit $P$ un objet parfait de
$D(\mathcal{O}_U)$ qui engendre $D_\QCoh(\mathcal{O}_U)$. En utilisant le
lemme \ref{spaces-perfect-lemma-direct-summand-of-a-restriction}, on peut choisir un objet
parfait $Q$ de $D(\mathcal{O}_X)$ dont la restriction \`a $U$ est une somme
directe ayant $P$ pour facteur direct. Posons $V=\Spec(A)$. Soit $Z\subset V$
le sous-sch\'ema ferm\'e r\'eduit, image inverse de $X\setminus U$, qui s'envoie
isomorphiquement sur celui-ci (voir la d\'efinition
\ref{spaces-perfect-definition-elementary-distinguished-square}). C'est une partie ferm\'ee
r\'etrocompacte de $V$. Choisissons $f_1,\ldots,f_r\in A$ tels que
$Z=V(f_1,\ldots,f_r)$. Soit $K\in D(\mathcal{O}_V)$ l'objet parfait
correspondant au complexe de Koszul associ\'e \`a $f_1,\ldots,f_r$ sur $A$.
Comme $K$ est \`a support dans $Z$, son image directe $K'=Rf_*K$ est un objet
parfait de $D(\mathcal{O}_X)$ dont la restriction \`a $V$ est $K$ (voir les
lemmes \ref{spaces-perfect-lemma-pushforward-perfect} et
\ref{spaces-perfect-lemma-pushforward-with-support-in-open}). Nous affirmons que
$Q\oplus K'$ engendre $D_\QCoh(\mathcal{O}_X)$.

\medskip\noindent
Soit $E$ un objet de $D_\QCoh(\mathcal{O}_X)$ tel qu'il n'existe aucun
morphisme non nul d'un translat\'e de $Q\oplus K'$ vers $E$. D'apr\`es le lemme
\ref{spaces-perfect-lemma-pushforward-with-support-in-open}, on a $K'=f_!K$, donc
$$
\Hom_{D(\mathcal{O}_X)}(K'[n], E) = \Hom_{D(\mathcal{O}_V)}(K[n], E|_V)
$$
Par suite, Cat\'egories d\'eriv\'ees des sch\'emas, lemme
\ref{perfect-lemma-orthogonal-koszul-complex} (en utilisant aussi le lemme
\ref{spaces-perfect-lemma-derived-quasi-coherent-small-etale-site}), montre que l'annulation
de ces groupes entra\^ine que $E|_V$ est isomorphe \`a
$R(U\times_XV\to V)_*E|_{U\times_XV}$. Il en r\'esulte que
$E=R(U\to X)_*E|_U$ (nous omettons un d\'etail mineur). Dans ce cas,
$$
\Hom_{D(\mathcal{O}_X)}(Q[n], E) = \Hom_{D(\mathcal{O}_U)}(Q|_U[n], E|_U)
$$
contient $\Hom_{D(\mathcal{O}_U)}(P[n],E|_U)$ comme facteur direct. Le choix
de $P$ montre donc que l'annulation de ces groupes entra\^ine $E|_U=0$, d'o\`u
$E=0$.
\end{proof}

\begin{remark}
\label{spaces-perfect-remark-pullback-generator}
Soit $S$ un sch\'ema. Soit $f:X\to Y$ un morphisme d'espaces alg\'ebriques
quasi-compacts et quasi-s\'epar\'es sur $S$. Soit
$E\in D_\QCoh(\mathcal{O}_Y)$ un g\'en\'erateur (voir le th\'eor\`eme
\ref{spaces-perfect-theorem-bondal-van-den-Bergh}). Les conditions suivantes sont alors
\'equivalentes :
\begin{enumerate}
\item pour $K\in D_\QCoh(\mathcal{O}_X)$, on a $Rf_*K=0$ si et seulement si
$K=0$;
\item $Rf_* : D_\QCoh(\mathcal{O}_X) \to D_\QCoh(\mathcal{O}_Y)$ est
conservatif;
\item $Lf^*E$ engendre $D_\QCoh(\mathcal{O}_X)$.
\end{enumerate}
L'\'equivalence de (1) et (2) est une cons\'equence formelle du fait que
$Rf_*:D_\QCoh(\mathcal{O}_X)\to D_\QCoh(\mathcal{O}_Y)$ est un foncteur exact
de cat\'egories triangul\'ees. De m\^eme, l'\'equivalence de (1) et (3) r\'esulte
formellement du fait que $Lf^*$ est adjoint \`a gauche de $Rf_*$. Ces
conditions sont satisfaites si $f$ est affine (lemme
\ref{spaces-perfect-lemma-affine-morphism}), si $f$ est une immersion ouverte, ou si $f$ est
une compos\'ee de morphismes de ces types.
\end{remark}

\noindent
Le r\'esultat suivant renforce le th\'eor\`eme
\ref{spaces-perfect-theorem-bondal-van-den-Bergh}; il se d\'emontre par les m\^emes m\'ethodes.
Soit $T\subset|X|$ une partie ferm\'ee, o\`u $X$ est un espace alg\'ebrique.
Notons $D_T(\mathcal{O}_X)$ la sous-cat\'egorie triangul\'ee strictement pleine
et satur\'ee form\'ee des complexes dont les faisceaux de cohomologie sont \`a
support dans $T$.

\begin{lemma}
\label{spaces-perfect-lemma-generator-with-support}
Soit $S$ un sch\'ema. Soit $X$ un espace alg\'ebrique quasi-compact et
quasi-s\'epar\'e sur $S$. Soit $T\subset|X|$ une partie ferm\'ee telle que
$|X|\setminus T$ soit quasi-compact. Avec les notations ci-dessus, la
cat\'egorie $D_{\QCoh,T}(\mathcal{O}_X)$ est engendr\'ee par un seul objet
parfait.
\end{lemma}

\begin{proof}
Nous allons le d\'emontrer au moyen du principe de r\'ecurrence du lemme
\ref{spaces-perfect-lemma-induction-principle}. D'apr\`es Cat\'egories d\'eriv\'ees des
sch\'emas, lemme \ref{perfect-lemma-generator-with-support}, la propri\'et\'e est
vraie pour les objets repr\'esentables, quasi-compacts et quasi-s\'epar\'es du
site $X_{spaces,\etale}$.

\medskip\noindent
Supposons que $(U\subset X,f:V\to X)$ soit un carr\'e distingu\'e
\'el\'ementaire tel que le lemme soit vrai pour $U$ et que $V$ soit affine.
Pour achever la d\'emonstration, il faut \'etablir le r\'esultat pour $X$. Soit
$P$ un objet parfait de $D(\mathcal{O}_U)$, \`a support dans $T\cap U$, qui
engendre $D_{\QCoh,T\cap U}(\mathcal{O}_U)$. D'apr\`es le lemme
\ref{spaces-perfect-lemma-lift-map-from-perfect-complex-with-support}, on peut choisir un
objet parfait $Q$ de $D(\mathcal{O}_X)$, \`a support dans $T$, dont la
restriction \`a $U$ est une somme directe ayant $P$ pour facteur direct.
\'Ecrivons $V=\Spec(B)$ et posons $Z=X\setminus U$. Alors $f^{-1}Z$ est une
partie ferm\'ee de $V$ telle que $V\setminus f^{-1}Z$ soit quasi-compact.
Comme $X$ est quasi-s\'epar\'e, il s'ensuit que
$f^{-1}Z\cap f^{-1}T=f^{-1}(Z\cap T)$ est une partie ferm\'ee de $V$ telle que
$W=V\setminus f^{-1}(Z\cap T)$ soit quasi-compact. On peut donc choisir
$g_1,\ldots,g_s\in B$ tels que
$f^{-1}(Z\cap T)=V(g_1,\ldots,g_s)$. Soit $K\in D(\mathcal{O}_V)$ l'objet
parfait correspondant au complexe de Koszul associ\'e \`a
$g_1,\ldots,g_s$ sur $B$. Comme $K$ est \`a support dans la partie ferm\'ee
$f^{-1}(Z\cap T)\subset V$, son image directe $K'=R(V\to X)_*K$ est un objet
parfait de $D(\mathcal{O}_X)$ dont la restriction \`a $V$ est $K$ (voir les
lemmes \ref{spaces-perfect-lemma-pushforward-perfect} et
\ref{spaces-perfect-lemma-pushforward-with-support-in-open}). Nous affirmons que
$Q\oplus K'$ engendre $D_{\QCoh,T}(\mathcal{O}_X)$.

\medskip\noindent
Soit $E$ un objet de $D_{\QCoh,T}(\mathcal{O}_X)$ tel qu'il n'existe aucun
morphisme non nul d'un translat\'e de $Q\oplus K'$ vers $E$. D'apr\`es le lemme
\ref{spaces-perfect-lemma-pushforward-with-support-in-open}, on a
$K'=R(V\to X)_!K$, donc
$$
\Hom_{D(\mathcal{O}_X)}(K'[n], E) = \Hom_{D(\mathcal{O}_V)}(K[n], E|_V)
$$
Cat\'egories d\'eriv\'ees des sch\'emas, lemme
\ref{perfect-lemma-orthogonal-koszul-complex}, donne donc
$E|_V=Rj_*E|_W$, o\`u $j:W\to V$ est l'inclusion. On a le diagramme
$$
\xymatrix{
W \ar[r]_j & V & f^{-1}(Z \cap T) \ar[l] \ar[d] \\
V \setminus f^{-1}Z \ar[u]^{j'} \ar[ru]_{j''} & & f^{-1}Z \ar[lu]
}
$$
Comme $E$ est \`a support dans $T$, on voit que $E|_W$ est \`a support dans
$f^{-1}T\cap W=f^{-1}T\cap(V\setminus f^{-1}Z)$, qui est ferm\'e dans $W$.
On en conclut que
$$
E|_V = Rj_*(E|_W) = Rj_*(Rj'_*(E|_{U \times_X V})) =
Rj''_*(E|_{U \times_X V})
$$
La seconde \'egalit\'e est l'assertion (1) de Cohomologie, lemme
\ref{cohomology-lemma-pushforward-restriction}. Celui-ci s'applique parce que
$V$ est un sch\'ema et que $E$ a des faisceaux de cohomologie quasi-coh\'erents;
l'image directe par l'immersion ouverte quasi-compacte $j'$ co\"incide donc
avec l'image directe sur les sch\'emas sous-jacents; voir la remarque
\ref{spaces-perfect-remark-match-total-direct-images}. Il en r\'esulte que
$E=R(U\to X)_*E|_U$ (nous omettons un d\'etail mineur). Dans ce cas,
$$
\Hom_{D(\mathcal{O}_X)}(Q[n], E) = \Hom_{D(\mathcal{O}_U)}(Q|_U[n], E|_U)
$$
contient $\Hom_{D(\mathcal{O}_U)}(P[n],E|_U)$ comme facteur direct. Le choix
de $P$ montre donc que l'annulation de ces groupes entra\^ine $E|_U=0$, d'o\`u
$E=0$.
\end{proof}










\section{Objets compacts et parfaits}
\label{spaces-perfect-section-compact}

\noindent
Cette section est l'analogue de
Cat\'egories d\'eriv\'ees des sch\'emas, section \ref{perfect-section-compact}.

\begin{proposition}
\label{spaces-perfect-proposition-compact-is-perfect}
Soit $S$ un sch\'ema.
Soit $X$ un espace alg\'ebrique quasi-compact et quasi-s\'epar\'e sur $S$.
Un objet de $D_\QCoh(\mathcal{O}_X)$ est compact
si et seulement s'il est parfait.
\end{proposition}

\begin{proof}
Si $K$ est un objet parfait de $D(\mathcal{O}_X)$, de dual
$K^\vee$ (Cohomologie sur les sites, lemme
\ref{sites-cohomology-lemma-dual-perfect-complex}),
on a
$$
\Hom_{D(\mathcal{O}_X)}(K, M) =
H^0(X, K^\vee \otimes_{\mathcal{O}_X}^\mathbf{L} M)
$$
fonctoriellement en $M$. Comme $K^\vee \otimes_{\mathcal{O}_X}^\mathbf{L} -$
commute aux sommes directes et que $H^0(X, -)$ commute aux sommes
directes dans $D_\QCoh(\mathcal{O}_X)$ d'apr\`es le
lemme \ref{spaces-perfect-lemma-quasi-coherence-pushforward-direct-sums},
on en conclut que $K$ est compact dans $D_\QCoh(\mathcal{O}_X)$.

\medskip\noindent
R\'eciproquement, soit $K$ un objet compact de $D_\QCoh(\mathcal{O}_X)$.
Pour montrer que $K$ est parfait, il suffit de montrer que
$K|_U$ est parfait pour tout sch\'ema affine $U$ \'etale sur $X$; voir
Cohomologie sur les sites, lemme
\ref{sites-cohomology-lemma-perfect-independent-representative}.
Observons que $j : U \to X$ est un morphisme quasi-compact et s\'epar\'e.
Par cons\'equent,
$Rj_* : D_\QCoh(\mathcal{O}_U) \to D_\QCoh(\mathcal{O}_X)$
commute aux sommes directes; voir le
lemme \ref{spaces-perfect-lemma-quasi-coherence-pushforward-direct-sums}.
L'adjonction entre la restriction \`a $U$ et $Rj_*$ montre donc que
$K|_U$ est un objet compact de $D_\QCoh(\mathcal{O}_U)$.
On se ram\`ene ainsi au cas o\`u $X$ est affine, et en particulier \`a celui
d'un sch\'ema quasi-compact et quasi-s\'epar\'e. Au moyen des
lemmes \ref{spaces-perfect-lemma-derived-quasi-coherent-small-etale-site} et
\ref{spaces-perfect-lemma-descend-perfect},
on se ram\`ene au cas des sch\'emas, c'est-\`a-dire \`a
Cat\'egories d\'eriv\'ees des sch\'emas, proposition
\ref{perfect-proposition-compact-is-perfect}.
\end{proof}

\begin{remark}
\label{spaces-perfect-remark-classical-generator}
Soit $S$ un sch\'ema.
Soit $X$ un espace alg\'ebrique quasi-compact et quasi-s\'epar\'e sur $S$.
Soit $G$ un
objet parfait de $D(\mathcal{O}_X)$ qui engendre
$D_\QCoh(\mathcal{O}_X)$. D'apr\`es le th\'eor\`eme
\ref{spaces-perfect-theorem-bondal-van-den-Bergh}, il en existe au moins un. En combinant
le lemme \ref{spaces-perfect-lemma-quasi-coherence-direct-sums}, la
proposition \ref{spaces-perfect-proposition-compact-is-perfect} et
Cat\'egories d\'eriv\'ees, proposition
\ref{derived-proposition-generator-versus-classical-generator},
on voit que $G$ est un g\'en\'erateur classique de
$D_{perf}(\mathcal{O}_X)$.
\end{remark}

\noindent
Le r\'esultat suivant renforce la
proposition \ref{spaces-perfect-proposition-compact-is-perfect}.
Soit $T \subset |X|$ une partie ferm\'ee, o\`u $X$ est un espace alg\'ebrique.
Comme pr\'ec\'edemment, $D_T(\mathcal{O}_X)$ d\'esigne la sous-cat\'egorie
triangul\'ee strictement pleine et satur\'ee form\'ee des complexes dont les
faisceaux de cohomologie sont \`a support dans $T$. Comme la formation des
sommes directes commute \`a celle des faisceaux de cohomologie, il s'ensuit
que $D_T(\mathcal{O}_X)$ admet des sommes directes et que celles-ci co\"incident
avec les sommes directes dans $D(\mathcal{O}_X)$.

\begin{lemma}
\label{spaces-perfect-lemma-compact-is-perfect-with-support}
Soit $S$ un sch\'ema.
Soit $X$ un espace alg\'ebrique quasi-compact et quasi-s\'epar\'e sur $S$.
Soit $T \subset |X|$ une partie ferm\'ee telle que $|X| \setminus T$
soit quasi-compact. Un objet de $D_{\QCoh, T}(\mathcal{O}_X)$ est compact
si et seulement s'il est parfait en tant qu'objet de $D(\mathcal{O}_X)$.
\end{lemma}

\begin{proof}
La remarque qui pr\'ec\`ede le lemme montre que
$D_{\QCoh, T}(\mathcal{O}_X)$ est une cat\'egorie triangul\'ee
admettant des sommes directes.
D'apr\`es la proposition \ref{spaces-perfect-proposition-compact-is-perfect},
les objets parfaits donnent des objets compacts de
$D_\QCoh(\mathcal{O}_X)$, et donc a fortiori de toute sous-cat\'egorie
stable par sommes directes. Pour la r\'eciproque, nous utiliserons le fait
qu'il existe un g\'en\'erateur
$E \in D_{\QCoh, T}(\mathcal{O}_X)$ qui est un complexe parfait
de $\mathcal{O}_X$-modules; voir le
lemme \ref{spaces-perfect-lemma-generator-with-support}.
D'apr\`es ce qui pr\'ec\`ede, $E$ est donc compact. Il r\'esulte alors de
Cat\'egories d\'eriv\'ees, proposition
\ref{derived-proposition-generator-versus-classical-generator},
que $E$ est un g\'en\'erateur classique de la sous-cat\'egorie pleine
des objets compacts de $D_{\QCoh, T}(\mathcal{O}_X)$.
Ainsi, tout objet compact se construit \`a partir de $E$ par
une suite finie d'op\'erations consistant \`a
(a) prendre des translat\'es, (b) prendre des sommes directes finies,
(c) prendre des c\^ones et
(d) prendre des facteurs directs. Chacune de ces op\'erations pr\'eserve
la propri\'et\'e d'\^etre parfait, ce qui prouve le r\'esultat.
\end{proof}

\begin{remark}
\label{spaces-perfect-remark-classical-generator-with-support}
Soit $S$ un sch\'ema.
Soit $X$ un espace alg\'ebrique quasi-compact et quasi-s\'epar\'e sur $S$.
Soit $T \subset |X|$ une partie ferm\'ee telle que $|X| \setminus T$
soit quasi-compact. Soit $G$ un
objet parfait de $D_{\QCoh, T}(\mathcal{O}_X)$ qui engendre
$D_{\QCoh, T}(\mathcal{O}_X)$. D'apr\`es le lemme
\ref{spaces-perfect-lemma-generator-with-support}, il en existe au moins un. En combinant
le fait que $D_{\QCoh, T}(\mathcal{O}_X)$ admet des sommes directes, le
lemme \ref{spaces-perfect-lemma-compact-is-perfect-with-support} et
Cat\'egories d\'eriv\'ees, proposition
\ref{derived-proposition-generator-versus-classical-generator},
on voit que $G$ est un g\'en\'erateur classique de
$D_{perf, T}(\mathcal{O}_X)$.
\end{remark}

\noindent
Le lemme suivant applique les id\'ees qui interviennent dans
la d\'emonstration du lemme pr\'ec\'edent.

\begin{lemma}
\label{spaces-perfect-lemma-map-from-pseudo-coherent-to-complex-with-support}
Soit $S$ un sch\'ema. Soit $X$ un espace alg\'ebrique quasi-compact et
quasi-s\'epar\'e sur $S$. Soit $T \subset |X|$
une partie ferm\'ee dont le compl\'ementaire $U \subset X$ est quasi-compact.
Soit $\alpha : P \to E$ un morphisme de $D_\QCoh(\mathcal{O}_X)$ se trouvant
dans l'un des deux cas suivants :
\begin{enumerate}
\item $P$ est parfait et $E$ est \`a support dans $T$;
\item $P$ est pseudo-coh\'erent, $E$ est \`a support dans $T$ et $E$ est
born\'e inf\'erieurement.
\end{enumerate}
Alors il existe un complexe parfait de $\mathcal{O}_X$-modules $I$
et un morphisme $I \to \mathcal{O}_X[0]$ tels que
$I \otimes^\mathbf{L} P \to E$ soit nul et que
$I|_U \to \mathcal{O}_U[0]$ soit un
isomorphisme.
\end{lemma}

\begin{proof}
Posons $\mathcal{D} = D_{\QCoh, T}(\mathcal{O}_X)$. Dans les deux cas, le complexe
$K = R\SheafHom(P, E)$ est un objet de $\mathcal{D}$. Voir le
lemme \ref{spaces-perfect-lemma-quasi-coherence-internal-hom} pour la quasi-coh\'erence.
Il est clair que $K$ est \`a support dans $T$, puisque la formation de
$R\SheafHom$ commute aux restrictions aux ouverts.
Le morphisme $\alpha$ d\'efinit un \'el\'ement de
$H^0(K) = \Hom_{D(\mathcal{O}_X)}(\mathcal{O}_X[0], K)$.
Il suffit donc de d\'emontrer le r\'esultat pour le morphisme
$\alpha : \mathcal{O}_X[0] \to K$.

\medskip\noindent
Soit $E \in \mathcal{D}$ un g\'en\'erateur parfait; voir le
lemme \ref{spaces-perfect-lemma-generator-with-support}. \'Ecrivons
$$
K = \text{hocolim} K_n
$$
comme dans Cat\'egories d\'eriv\'ees, lemme
\ref{derived-lemma-write-as-colimit}, en utilisant le g\'en\'erateur $E$.
Comme le foncteur $\mathcal{D} \to D(\mathcal{O}_X)$ commute aux sommes
directes, l'\'egalit\'e $K = \text{hocolim} K_n$ vaut aussi dans
$D(\mathcal{O}_X)$. Puisque $\mathcal{O}_X$ est un objet compact de
$D(\mathcal{O}_X)$, il existe un entier $n$ et un morphisme
$\alpha_n : \mathcal{O}_X \to K_n$ qui induit $\alpha$; voir
Cat\'egories d\'eriv\'ees, lemme
\ref{derived-lemma-commutes-with-countable-sums}.
En appliquant Cat\'egories d\'eriv\'ees, lemme
\ref{derived-lemma-factor-through}, au morphisme
$\mathcal{O}_X[0] \to K_n$ de la cat\'egorie ambiante
$D(\mathcal{O}_X)$, on voit que $\alpha_n$ se factorise sous la forme
$\mathcal{O}_X[0] \to Q \to K_n$, o\`u $Q$ est un objet
de $\langle E \rangle$. On en conclut que $Q$ est un complexe parfait
\`a support dans $T$.

\medskip\noindent
Choisissons un triangle distingu\'e
$$
I \to \mathcal{O}_X[0] \to Q \to I[1]
$$
Par construction, $I$ est parfait, le morphisme $I \to \mathcal{O}_X[0]$
se restreint \`a un isomorphisme sur $U$, et le compos\'e
$I \to K$ est nul puisque $\alpha$ se factorise par $Q$.
Cela prouve le lemme.
\end{proof}

\section{Les cat\'egories d\'eriv\'ees comme cat\'egories de modules}
\label{spaces-perfect-section-derived-is-dga}

\noindent
Cette section est l'analogue de Cat\'egories d\'eriv\'ees des sch\'emas,
section \ref{perfect-section-derived-is-dga}.

\begin{lemma}
\label{spaces-perfect-lemma-tensor-with-QCoh-complex}
Soit $S$ un sch\'ema et soit $X$ un espace alg\'ebrique sur $S$. Soit
$K^\bullet$ un complexe de $\mathcal{O}_X$-modules dont les faisceaux de
cohomologie sont quasi-coh\'erents. Soit
$(E,d)=\Hom_{\text{Comp}^{dg}(\mathcal{O}_X)}(K^\bullet,K^\bullet)$ l'alg\`ebre
diff\'erentielle gradu\'ee des endomorphismes. Alors le foncteur
$$
- \otimes_E^\mathbf{L} K^\bullet :
D(E, \text{d}) \longrightarrow D(\mathcal{O}_X)
$$
du chapitre Alg\`ebres diff\'erentielles gradu\'ees, lemme
\ref{dga-lemma-tensor-with-complex-derived}, prend ses valeurs dans
$D_\QCoh(\mathcal{O}_X)$.
\end{lemma}

\begin{proof}
Soit $P$ un $E$-module diff\'erentiel gradu\'e poss\'edant la propri\'et\'e P.
Soit $F_\bullet$ une filtration de $P$ comme dans Alg\`ebres diff\'erentielles
gradu\'ees, section \ref{dga-section-P-resolutions}. On a alors
$$
P \otimes_E K^\bullet = \text{hocolim}\ F_iP \otimes_E K^\bullet.
$$
Chacun des $F_iP$ poss\`ede une filtration finie dont les gradu\'es associ\'es
sont des sommes directes de $E[k]$. Le r\'esultat s'en d\'eduit ais\'ement.
\end{proof}

\noindent
On peut renforcer le lemme suivant : l'annulation est uniforme lorsque $L$
parcourt les complexes dont les faisceaux de cohomologie non nuls sont
concentr\'es dans un intervalle fix\'e.

\begin{lemma}
\label{spaces-perfect-lemma-ext-from-perfect-into-bounded-QCoh}
Soit $S$ un sch\'ema. Soit $X$ un espace alg\'ebrique quasi-compact et
quasi-s\'epar\'e sur $S$. Soit $K$ un objet parfait de
$D(\mathcal{O}_X)$. Alors
\begin{enumerate}
\item il existe des entiers $a\leq b$ tels que
$\Hom_{D(\mathcal{O}_X)}(K,L)=0$ pour tout
$L\in D_\QCoh(\mathcal{O}_X)$ tel que $H^i(L)=0$ pour $i\in[a,b]$;
\item si $L$ est born\'e, alors
$\Ext^n_{D(\mathcal{O}_X)}(K,L)$ est nul sauf pour un nombre fini de valeurs de
$n$.
\end{enumerate}
\end{lemma}

\begin{proof}
L'assertion (2) r\'esulte de (1), puisque
$\Ext^n_{D(\mathcal{O}_X)}(K,L)=
\Hom_{D(\mathcal{O}_X)}(K,L[n])$. D\'emontrons (1). Comme $K$ est parfait, on a
$$
\Ext^i_{D(\mathcal{O}_X)}(K, L) =
H^i(X, K^\vee \otimes_{\mathcal{O}_X}^\mathbf{L} L),
$$
o\`u $K^\vee$ est le complexe parfait dual de $K$; voir Cohomologie sur les
sites, lemme \ref{sites-cohomology-lemma-dual-perfect-complex}. Remarquons que
$P=K^\vee\otimes_{\mathcal{O}_X}^\mathbf{L}L$ appartient \`a $D_\QCoh(X)$
d'apr\`es les lemmes \ref{spaces-perfect-lemma-quasi-coherence-tensor-product} et
\ref{spaces-perfect-lemma-pseudo-coherent} (ce dernier montre qu'un complexe parfait a des
faisceaux de cohomologie quasi-coh\'erents). Supposons que $K^\vee$ soit de
Tor-amplitude contenue dans $[a,b]$. La suite spectrale
$$
E_1^{p, q} = H^p(K^\vee \otimes_{\mathcal{O}_X}^\mathbf{L} H^q(L))
\Rightarrow
H^{p + q}(K^\vee \otimes_{\mathcal{O}_X}^\mathbf{L} L)
$$
montre que $H^j(K^\vee\otimes_{\mathcal{O}_X}^\mathbf{L}L)$ est nul si
$H^q(L)=0$ pour $q\in[j-b,j-a]$. Soit $N$ l'entier
$\max(d_p+p)$ du lemme de Cohomologie des espaces
\ref{spaces-cohomology-lemma-vanishing-quasi-separated}. Alors
$H^0(X,K^\vee\otimes_{\mathcal{O}_X}^\mathbf{L}L)$ s'annule si les faisceaux de
cohomologie
$$
H^{-N}(K^\vee \otimes_{\mathcal{O}_X}^\mathbf{L} L),
\ H^{-N + 1}(K^\vee \otimes_{\mathcal{O}_X}^\mathbf{L} L),
\ \ldots,
\ H^0(K^\vee \otimes_{\mathcal{O}_X}^\mathbf{L} L)
$$
sont nuls. En effet, le lemme cit\'e et le lemme
\ref{spaces-perfect-lemma-application-nice-K-injective} donnent
$$
H^0(X, K^\vee \otimes_{\mathcal{O}_X}^\mathbf{L} L) =
H^0(X, \tau_{\geq -N}(K^\vee \otimes_{\mathcal{O}_X}^\mathbf{L} L));
$$
l'annulation des faisceaux de cohomologie montre que ce dernier groupe est
\'egal \`a
$H^0(X,\tau_{\geq1}(K^\vee\otimes_{\mathcal{O}_X}^\mathbf{L}L))$, lequel est
nul d'apr\`es Cat\'egories d\'eriv\'ees, lemme
\ref{derived-lemma-negative-vanishing}. Par cons\'equent,
$\Hom_{D(\mathcal{O}_X)}(K,L)$ est nul si $H^i(L)=0$ pour
$i\in[-b-N,-a]$.
\end{proof}

\noindent
Le th\'eor\`eme suivant est l'analogue de Cat\'egories d\'eriv\'ees des sch\'emas,
th\'eor\`eme \ref{perfect-theorem-DQCoh-is-Ddga}.

\begin{theorem}
\label{spaces-perfect-theorem-DQCoh-is-Ddga}
Soit $S$ un sch\'ema. Soit $X$ un espace alg\'ebrique quasi-compact et
quasi-s\'epar\'e sur $S$. Il existe alors une alg\`ebre diff\'erentielle gradu\'ee
$(E,\text{d})$ dont seuls un nombre fini de groupes de cohomologie $H^i(E)$
sont non nuls, et telle que $D_\QCoh(\mathcal{O}_X)$ soit \'equivalente \`a
$D(E,\text{d})$.
\end{theorem}

\begin{proof}
Soit $K^\bullet$ un complexe K-injectif de $\mathcal{O}$-modules qui soit
parfait et qui engendre $D_\QCoh(\mathcal{O}_X)$. Un tel complexe existe par le
th\'eor\`eme \ref{spaces-perfect-theorem-bondal-van-den-Bergh} et par l'existence des
r\'esolutions K-injectives. Nous allons montrer que le th\'eor\`eme est vrai pour
$$
(E, \text{d}) = \Hom_{\text{Comp}^{dg}(\mathcal{O}_X)}(K^\bullet, K^\bullet),
$$
o\`u $\text{Comp}^{dg}(\mathcal{O}_X)$ est la cat\'egorie diff\'erentielle
gradu\'ee des complexes de $\mathcal{O}$-modules. Voir Alg\`ebres
diff\'erentielles gradu\'ees, section
\ref{dga-section-variant-base-change}. Comme $K^\bullet$ est K-injectif, on a
\begin{equation}
\label{spaces-perfect-equation-E-is-OK}
H^n(E) = \Ext^n_{D(\mathcal{O}_X)}(K^\bullet, K^\bullet)
\end{equation}
pour tout $n\in\mathbf{Z}$. Ces groupes Ext ne sont non nuls
que pour un nombre fini d'indices, d'apr\`es le lemme \ref{spaces-perfect-lemma-ext-from-perfect-into-bounded-QCoh}.
Consid\'erons le foncteur
$$
- \otimes_E^\mathbf{L} K^\bullet :
D(E, \text{d}) \longrightarrow D(\mathcal{O}_X)
$$
d'Alg\`ebres diff\'erentielles gradu\'ees, lemme
\ref{dga-lemma-tensor-with-complex-derived}. Comme
$K^\bullet$ est parfait, il d\'efinit un objet compact de
$D_\QCoh(\mathcal{O}_X)$; voir la proposition \ref{spaces-perfect-proposition-compact-is-perfect}.
Avec (\ref{spaces-perfect-equation-E-is-OK}), l'argument du lemme d'Alg\`ebres diff\'erentielles
gradu\'ees \ref{dga-lemma-fully-faithful-in-compact-case} montre que le foncteur
ci-dessus \`a valeurs dans $D_\QCoh(\mathcal{O}_X)$ est pleinement fid\`ele; son adjoint \`a droite est
$$
R\Hom(K^\bullet, - ) : D_\QCoh(\mathcal{O}_X) \longrightarrow D(E, \text{d})
$$
d'apr\`es Alg\`ebres diff\'erentielles gradu\'ees, lemme
\ref{dga-lemma-tensor-with-complex-hom-adjoint}; par les propri\'et\'es g\'en\'erales
des foncteurs adjoints, cet adjoint est un quasi-inverse \`a gauche. Le lemme
\ref{spaces-perfect-lemma-tensor-with-QCoh-complex} donne pr\'ecis\'ement ce foncteur
$$
- \otimes_E^\mathbf{L} K^\bullet :
D(E, \text{d}) \longrightarrow D_\QCoh(\mathcal{O}_X),
$$
et, puisque $K^\bullet$ engendre $D_\QCoh(\mathcal{O}_X)$, le noyau de
l'adjoint restreint \`a $D_\QCoh(\mathcal{O}_X)$ est nul. Un argument formel
fournit alors l'\'equivalence cherch\'ee; voir Cat\'egories d\'eriv\'ees, lemme
\ref{derived-lemma-fully-faithful-adjoint-kernel-zero}.
\end{proof}

\begin{remark}[Variante avec support]
\label{spaces-perfect-remark-DQCoh-is-Ddga-with-support}
Soit $S$ un sch\'ema. Soit $X$ un espace alg\'ebrique quasi-compact et
quasi-s\'epar\'e sur $S$. Soit $T\subset|X|$ une partie ferm\'ee telle que
$|X|\setminus T$ soit quasi-compact. L'analogue du th\'eor\`eme
\ref{spaces-perfect-theorem-DQCoh-is-Ddga} est vrai pour
$D_{\QCoh,T}(\mathcal{O}_X)$. Cela r\'esulte exactement du m\^eme argument que
dans la d\'emonstration du th\'eor\`eme, en utilisant les lemmes
\ref{spaces-perfect-lemma-generator-with-support} et
\ref{spaces-perfect-lemma-compact-is-perfect-with-support}, ainsi qu'une variante avec
supports du lemme \ref{spaces-perfect-lemma-tensor-with-QCoh-complex}. Si ce r\'esultat devient
n\'ecessaire, nous en donnerons ici un \'enonc\'e pr\'ecis et une d\'emonstration
d\'etaill\'ee.
\end{remark}

\begin{remark}[Unicit\'e de l'alg\`ebre diff\'erentielle gradu\'ee]
\label{spaces-perfect-remark-independence-choice}
Soit $X$ un espace alg\'ebrique quasi-compact et quasi-s\'epar\'e sur un anneau
$R$. La construction de la d\'emonstration du th\'eor\`eme
\ref{spaces-perfect-theorem-DQCoh-is-Ddga} fournit une alg\`ebre diff\'erentielle gradu\'ee
$(A,\text{d})$ sur $R$ telle que $D_\QCoh(X)$ soit $R$-lin\'eairement
\'equivalente \`a $D(A,\text{d})$ comme cat\'egorie triangul\'ee. On peut se
demander dans quelle mesure $(A,\text{d})$ est unique. La r\'eponse est
seulement un peu plus pr\'ecise que l'affirmation selon laquelle
$(A,\text{d})$ est d\'efinie \`a \'equivalence d\'eriv\'ee pr\`es. Supposons en
effet que $(B,\text{d})$ soit une seconde paire de ce type. On a alors
$$
(A, \text{d}) = \Hom_{\text{Comp}^{dg}(\mathcal{O}_X)}(K^\bullet, K^\bullet)
$$
et
$$
(B, \text{d}) = \Hom_{\text{Comp}^{dg}(\mathcal{O}_X)}(L^\bullet, L^\bullet)
$$
pour des complexes K-injectifs $K^\bullet$ et $L^\bullet$ de
$\mathcal{O}_X$-modules correspondant \`a des g\'en\'erateurs parfaits de
$D_\QCoh(\mathcal{O}_X)$. Posons
$$
\Omega = \Hom_{\text{Comp}^{dg}(\mathcal{O}_X)}(K^\bullet, L^\bullet)
\quad
\Omega' = \Hom_{\text{Comp}^{dg}(\mathcal{O}_X)}(L^\bullet, K^\bullet).
$$
Alors $\Omega$ est un $B^{opp}\otimes_RA$-module diff\'erentiel gradu\'e, et
$\Omega'$ un $A^{opp}\otimes_RB$-module diff\'erentiel gradu\'e. De plus,
l'\'equivalence
$$
D(A, \text{d}) \to D_\QCoh(\mathcal{O}_X) \to
D(B, \text{d})
$$
est donn\'ee par le foncteur $-\otimes_A^\mathbf{L}\Omega'$, et de m\^eme pour
le quasi-inverse. Nous sommes donc dans la situation d'Alg\`ebres
diff\'erentielles gradu\'ees, remarque
\ref{dga-remark-hochschild-cohomology}. Si cette remarque devient n\'ecessaire,
nous en donnerons ici un \'enonc\'e pr\'ecis et une d\'emonstration d\'etaill\'ee.
\end{remark}

\section{Caract\'erisation des complexes pseudo-coh\'erents, I}
\label{spaces-perfect-section-pseudo-coherent-hocolim}

\noindent
Cette \'etude se poursuivra dans Compl\'ements sur les morphismes d'espaces,
section
\ref{spaces-more-morphisms-section-characterize-pseudo-coherent}.
On peut caract\'eriser les objets pseudo-coh\'erents
comme des colimites homotopiques d'approximations par des objets parfaits.

\begin{lemma}
\label{spaces-perfect-lemma-pseudo-coherent-hocolim}
Soit $S$ un sch\'ema.
Soit $X$ un espace alg\'ebrique quasi-compact et quasi-s\'epar\'e sur $S$.
Soit $K \in D(\mathcal{O}_X)$. Les conditions suivantes sont \'equivalentes :
\begin{enumerate}
\item $K$ est pseudo-coh\'erent;
\item $K = \text{hocolim} K_n$, o\`u
$K_n$ est parfait et $\tau_{\geq -n}K_n \to \tau_{\geq -n}K$
est un isomorphisme pour tout $n$.
\end{enumerate}
\end{lemma}

\begin{proof}
L'implication (2) $\Rightarrow$ (1) est valable sur tout site annel\'e.
Supposons en effet que (2) soit satisfaite. Rappelons qu'un objet parfait de la
cat\'egorie d\'eriv\'ee est pseudo-coh\'erent; voir
Cohomologie sur les sites, lemme \ref{sites-cohomology-lemma-perfect}.
Il r\'esulte alors des d\'efinitions que
$\tau_{\geq -n}K_n$ est $(-n + 1)$-pseudo-coh\'erent,
donc que $\tau_{\geq -n}K$ est $(-n + 1)$-pseudo-coh\'erent,
et par suite que $K$ est $(-n + 1)$-pseudo-coh\'erent. Comme ceci vaut pour
tout $n$, le complexe $K$ est pseudo-coh\'erent; voir
Cohomologie sur les sites, d\'efinition
\ref{sites-cohomology-definition-pseudo-coherent}.

\medskip\noindent
Supposons (1). Commen\c{c}ons par choisir une approximation
$K_1 \to K$ de $(X, K, -2)$ par un complexe parfait $K_1$; voir les
d\'efinitions \ref{spaces-perfect-definition-approximation-holds} et
\ref{spaces-perfect-definition-approximation}, ainsi que le
th\'eor\`eme \ref{spaces-perfect-theorem-approximation}.
Supposons par r\'ecurrence que l'on ait
$$
K_1 \to K_2 \to \ldots \to K_n \to K
$$
avec les $K_i$ parfaits et tels que
$\tau_{\geq -i}K_i \to \tau_{\geq -i}K$ soit un isomorphisme
pour tout $1 \leq i \leq n$. Choisissons alors $a \leq b$ comme dans le
lemme \ref{spaces-perfect-lemma-ext-from-perfect-into-bounded-QCoh}
pour l'objet parfait $K_n$. Choisissons une approximation
$K_{n + 1} \to K$ de $(X, K, \min(a - 1, -n - 1))$.
Choisissons un triangle distingu\'e
$$
K_{n + 1} \to K \to C \to K_{n + 1}[1]
$$
On voit alors que $C \in D_\QCoh(\mathcal{O}_X)$ v\'erifie
$H^i(C) = 0$ pour $i \geq a$. Le choix de $a, b$ donne donc
$\Hom_{D(\mathcal{O}_X)}(K_n, C) = 0$.
Le compos\'e $K_n \to K \to C$ est par cons\'equent nul. D'apr\`es
Cat\'egories d\'eriv\'ees, lemme
\ref{derived-lemma-representable-homological},
le morphisme $K_n \to K$ se factorise par $K_{n + 1}$,
ce qui \'etablit l'h\'er\'edit\'e.

\medskip\noindent
Il reste \`a d\'emontrer que $K = \text{hocolim} K_n$.
Cela r\'esulte de l'application de
Cat\'egories d\'eriv\'ees, lemme \ref{derived-lemma-cohomology-of-hocolim},
aux foncteurs
$H^i( - ) : D(\mathcal{O}_X) \to \textit{Mod}(\mathcal{O}_X)$
et du choix des $K_n$.
\end{proof}

\begin{lemma}
\label{spaces-perfect-lemma-pseudo-coherent-hocolim-with-support}
Soit $X$ un sch\'ema quasi-compact et quasi-s\'epar\'e.
Soit $T \subset X$ une partie ferm\'ee telle que $X \setminus T$
soit quasi-compact. Soit $K \in D(\mathcal{O}_X)$ \`a support dans $T$.
Les conditions suivantes sont \'equivalentes :
\begin{enumerate}
\item $K$ est pseudo-coh\'erent;
\item $K = \text{hocolim} K_n$, o\`u
$K_n$ est parfait, \`a support dans $T$, et
$\tau_{\geq -n}K_n \to \tau_{\geq -n}K$ est un isomorphisme pour tout $n$.
\end{enumerate}
\end{lemma}

\begin{proof}
La d\'emonstration est exactement la m\^eme que celle du
lemme \ref{spaces-perfect-lemma-pseudo-coherent-hocolim},
si ce n'est que, pour choisir les approximations, on utilise
les triplets $(T, K, m)$.
\end{proof}

\section{Retour sur le cohérateur}
\label{spaces-perfect-section-better-coherator}

\noindent
Dans la section \ref{spaces-perfect-section-coherator}, nous avons construit et étudié
l'adjoint à droite $RQ_X$ du foncteur canonique
$D(\QCoh(\mathcal{O}_X)) \to D(\mathcal{O}_X)$.
On l'a construit comme le foncteur dérivé droit du cohérateur
$Q_X : \textit{Mod}(\mathcal{O}_X) \to \QCoh(\mathcal{O}_X)$.
Dans la présente section, nous étudions les conditions sous lesquelles le foncteur d'inclusion
$$
D_\QCoh(\mathcal{O}_X) \longrightarrow D(\mathcal{O}_X)
$$
admet un adjoint à droite. S'il existe, nous le
noterons\footnote{Cette notation n'est probablement pas standard. Nous avons toutefois déjà
utilisé $Q_X$ pour le cohérateur et $RQ_X$ pour son foncteur dérivé.}
$$
DQ_X :
D(\mathcal{O}_X) \longrightarrow D_\QCoh(\mathcal{O}_X)
$$
Il s'avère que les espaces algébriques quasi-compacts et quasi-séparés
admettent un tel adjoint à droite.

\begin{lemma}
\label{spaces-perfect-lemma-better-coherator}
Soit $S$ un schéma.
Soit $X$ un espace algébrique quasi-compact et quasi-séparé sur $S$.
Le foncteur d'inclusion $D_\QCoh(\mathcal{O}_X) \to D(\mathcal{O}_X)$
admet un adjoint à droite.
\end{lemma}

\begin{proof}[Première démonstration]
Nous allons le démontrer à l'aide du principe de récurrence du
lemme \ref{spaces-perfect-lemma-induction-principle}.
Si $D(\QCoh(\mathcal{O}_X)) \to D_\QCoh(\mathcal{O}_X)$
est une équivalence, le lemme est vrai, car le foncteur
$RQ_X$ de la section \ref{spaces-perfect-section-coherator} est adjoint à droite au foncteur
$D(\QCoh(\mathcal{O}_X)) \to D(\mathcal{O}_X)$.
En particulier, notre lemme est vrai pour les espaces algébriques affines ; voir le
lemme \ref{spaces-perfect-lemma-affine-coherator}.
Il suffit donc de montrer ceci : si $(U \subset X, f : V \to X)$
est un carré distingué élémentaire, avec $U$ quasi-compact
et $V$ affine, et si le lemme vaut pour $U$, $V$ et $U \times_X V$,
alors il vaut pour $X$.

\medskip\noindent
L'adjoint existe si et seulement si, pour tout objet $E$ de
$D(\mathcal{O}_X)$, on peut trouver un triangle distingué
$$
E' \to E \to K \to E'[1]
$$
dans $D(\mathcal{O}_X)$
tel que $E'$ appartienne à $D_\QCoh(\mathcal{O}_X)$ et que
$\Hom(M, K) = 0$ pour tout $M$ de $D_\QCoh(\mathcal{O}_X)$. Voir
Catégories dérivées, lemme \ref{derived-lemma-right-adjoint}.
Considérons le triangle distingué
$$
E \to Rj_{U, *}E|_U \oplus Rj_{V, *}E|_V \to
Rj_{U \times_X V, *}E|_{U \times_X V} \to E[1]
$$
dans $D(\mathcal{O}_X)$ du lemme \ref{spaces-perfect-lemma-exact-sequence-j-star}.
D'après Catégories dérivées, lemme \ref{derived-lemma-prepare-adjoint},
il suffit de construire les triangles distingués voulus
pour $Rj_{U, *}E|_U$, $Rj_{V, *}E|_V$ et
$Rj_{U \times_X V, *}E|_{U \times_X V}$. Nous sommes ainsi ramenés à l'assertion
étudiée au paragraphe suivant.

\medskip\noindent
Soit $j : U \to X$ un morphisme étale, où
$U$ est quasi-compact et quasi-séparé et où le lemme est vrai pour $U$.
Soit $L$ un objet de $D(\mathcal{O}_U)$.
Il existe alors un triangle distingué
$$
E' \to Rj_*L \to K \to E'[1]
$$
dans $D(\mathcal{O}_X)$
tel que $E'$ appartienne à $D_\QCoh(\mathcal{O}_X)$ et que
$\Hom(M, K) = 0$ pour tout $M$ de $D_\QCoh(\mathcal{O}_X)$.
Pour le voir, choisissons un triangle distingué
$$
L' \to L \to Q \to L'[1]
$$
dans $D(\mathcal{O}_U)$ tel que $L'$ appartienne à $D_\QCoh(\mathcal{O}_U)$
et que $\Hom(N, Q) = 0$ pour tout $N$ de $D_\QCoh(\mathcal{O}_U)$.
Cela est possible, car les conditions du lemme
\ref{derived-lemma-right-adjoint} de Catégories dérivées
sont équivalentes.
Nous obtenons un triangle distingué
$$
Rj_*L' \to Rj_*L \to Rj_*Q \to Rj_*L'[1]
$$
dans $D(\mathcal{O}_X)$. Notons que $Rj_*L'$ appartient à $D_\QCoh(\mathcal{O}_X)$
d'après le lemme \ref{spaces-perfect-lemma-quasi-coherence-direct-image}.
D'autre part, si $M$ appartient à $D_\QCoh(\mathcal{O}_X)$, alors
$$
\Hom(M, Rj_*Q) = \Hom(Lj^*M, Q) = 0
$$
car $Lj^*M$ appartient à $D_\QCoh(\mathcal{O}_U)$ d'après le
lemme \ref{spaces-perfect-lemma-quasi-coherence-pullback}.
Cela achève la démonstration.
\end{proof}

\begin{proof}[Seconde démonstration]
L'adjoint existe d'après Catégories dérivées, proposition
\ref{derived-proposition-brown}. Les hypothèses sont satisfaites :
d'abord, $D_\QCoh(\mathcal{O}_X)$ admet les sommes directes,
et le foncteur d'inclusion commute aux sommes directes
(lemme \ref{spaces-perfect-lemma-quasi-coherence-direct-sums}).
D'autre part, $D_\QCoh(\mathcal{O}_X)$
est compactement engendrée, puisqu'elle possède un
générateur parfait d'après le théorème \ref{spaces-perfect-theorem-bondal-van-den-Bergh},
et puisque les objets parfaits sont compacts d'après la
proposition \ref{spaces-perfect-proposition-compact-is-perfect}.
\end{proof}

\begin{lemma}
\label{spaces-perfect-lemma-pushforward-better-coherator}
Soit $S$ un schéma.
Soit $f : X \to Y$ un morphisme quasi-compact et quasi-séparé
d'espaces algébriques sur $S$.
Si les adjoints à droite $DQ_X$ et $DQ_Y$
des foncteurs d'inclusion $D_\QCoh \to D$ existent pour $X$ et $Y$, alors
$$
Rf_* \circ DQ_X = DQ_Y \circ Rf_*
$$
\end{lemma}

\begin{proof}
L'énoncé a un sens parce que $Rf_*$ envoie
$D_\QCoh(\mathcal{O}_X)$ dans $D_\QCoh(\mathcal{O}_Y)$ d'après le
lemme \ref{spaces-perfect-lemma-quasi-coherence-direct-image}.
Il est vrai parce que $Lf^*$ envoie de même
$D_\QCoh(\mathcal{O}_Y)$ dans $D_\QCoh(\mathcal{O}_X)$
(lemme \ref{spaces-perfect-lemma-quasi-coherence-pullback}) ;
ainsi, $Rf_* \circ DQ_X$ et $DQ_Y \circ Rf_*$ sont tous deux
adjoints à droite de $Lf^* : D_\QCoh(\mathcal{O}_Y) \to D(\mathcal{O}_X)$.
\end{proof}

\begin{remark}
\label{spaces-perfect-remark-explain-consequence}
Soit $S$ un schéma. Soit $(U \subset X, f : V \to X)$ un
carré distingué élémentaire d'espaces algébriques sur $S$.
Supposons $X$, $U$ et $V$ quasi-compacts et quasi-séparés.
D'après le lemme \ref{spaces-perfect-lemma-better-coherator}, les foncteurs
$DQ_X$, $DQ_U$, $DQ_V$, $DQ_{U \times_X V}$ existent. De plus, on a un
triangle distingué canonique
$$
DQ_X(K) \to Rj_{U, *}DQ_U(K|_U) \oplus Rj_{V, *}DQ_V(K|_V)
\to Rj_{U \times_X V, *}DQ_{U \times_X V}(K|_{U \times_X V})
\to DQ_X(K)[1]
$$
pour tout $K \in D(\mathcal{O}_X)$. Cela résulte de l'application du
foncteur exact $DQ_X$ au triangle distingué du
lemme \ref{spaces-perfect-lemma-exact-sequence-j-star}
et de trois applications du lemme \ref{spaces-perfect-lemma-pushforward-better-coherator}.
\end{remark}

\begin{lemma}
\label{spaces-perfect-lemma-boundedness-better-coherator}
Soit $S$ un schéma.
Soit $X$ un espace algébrique quasi-compact et quasi-séparé sur $S$.
Le foncteur $DQ_X$ du lemme \ref{spaces-perfect-lemma-better-coherator}
possède la propriété de bornitude suivante :
il existe un entier $N = N(X)$ tel que, si
$K$ appartient à $D(\mathcal{O}_X)$ et si
$H^i(U, K) = 0$ pour tout $U$ affine étale sur $X$ et tout $i \not \in [a, b]$, alors
les faisceaux de cohomologie $H^i(DQ_X(K))$ sont nuls pour
$i \not \in [a, b + N]$.
\end{lemma}

\begin{proof}
Nous allons le démontrer à l'aide du principe de récurrence du
lemme \ref{spaces-perfect-lemma-induction-principle}.

\medskip\noindent
Si $X$ est affine, le lemme est vrai avec $N = 0$, car
$RQ_X = DQ_X$ s'obtient alors en prenant le complexe de
faisceaux quasi-cohérents associé à $R\Gamma(X, K)$.
Voir le lemme \ref{spaces-perfect-lemma-affine-coherator}.

\medskip\noindent
Soit $(U \subset W, f : V \to W)$ un carré distingué élémentaire,
où $W$ est quasi-compact et quasi-séparé, $U \subset W$
est un ouvert quasi-compact et $V$ est affine, tel que
le lemme vaille pour $U$, $V$ et $U \times_W V$.
Notons $N(U)$, $N(V)$ et $N(U \times_W V)$ les entiers correspondants.
Supposons maintenant que $K$ appartienne à $D(\mathcal{O}_W)$ et que
$H^i(T, K) = 0$ pour tout $T$ affine étale sur $W$ et tout $i \not \in [a, b]$.
Alors $K|_U$, $K|_V$ et $K|_{U \times_W V}$ ont la même propriété.
Il en résulte que $DQ_U(K|_U)$, $DQ_V(K|_V)$ et
$DQ_{U \times_W V}(K|_{U \times_W V})$ ont leurs faisceaux de cohomologie nuls
hors de l'intervalle $[a, b + \max(N(U), N(V), N(U \times_W V))]$.
Puisque les foncteurs $Rj_{U, *}$, $Rj_{V, *}$ et $Rj_{U \times_W V, *}$
sont de dimension cohomologique finie sur $D_\QCoh$ d'après le
lemme \ref{spaces-perfect-lemma-quasi-coherence-direct-image},
il existe un entier $N$ tel que
$Rj_{U, *}DQ_U(K|_U)$, $Rj_{V, *}DQ_V(K|_V)$ et
$Rj_{U \times_W V, *}DQ_{U \times_W V}(K|_{U \times_W V})$ aient leurs
faisceaux de cohomologie nuls hors de l'intervalle $[a, b + N]$.
On conclut enfin au moyen du triangle distingué de la
remarque \ref{spaces-perfect-remark-explain-consequence}.
\end{proof}

\begin{example}
\label{spaces-perfect-example-inverse-limit-quasi-coherent}
Soit $S$ un schéma.
Soit $X$ un espace algébrique quasi-compact et quasi-séparé sur $S$.
Soit $(\mathcal{F}_n)$
un système projectif de faisceaux quasi-cohérents sur $X$.
Comme $DQ_X$ est adjoint à droite, il commute aux
produits, donc aux limites projectives dérivées.
Il en résulte que
$$
DQ_X(R\lim \mathcal{F}_n) =
(R\lim\text{ dans }D_\QCoh(\mathcal{O}_X))(\mathcal{F}_n)
$$
où la première $R\lim$ est calculée dans $D(\mathcal{O}_X)$.
En fait, posons $K = R\lim \mathcal{F}_n$.
Pour tout $U$ affine étale sur $X$, on a
$$
H^i(U, K) =
H^i(R\Gamma(U, R\lim \mathcal{F}_n)) =
H^i(R\lim R\Gamma(U, \mathcal{F}_n)) =
H^i(R\lim \Gamma(U, \mathcal{F}_n))
$$
En effet, la cohomologie commute aux limites projectives dérivées et les
faisceaux quasi-cohérents
$\mathcal{F}_n$ n'ont pas de cohomologie supérieure sur les espaces affines.
Le calcul de $R\lim$ dans la catégorie des
groupes abéliens montre que $H^i(U, K) = 0$
sauf si $i \in [0, 1]$. On en conclut que
le $R\lim$ dans $D_\QCoh(\mathcal{O}_X)$, qui est
$DQ_X(K)$ d'après ce qui précède, appartient à $D^b_\QCoh(\mathcal{O}_X)$
et que ses faisceaux de cohomologie sont nuls en degrés négatifs
d'après le lemme \ref{spaces-perfect-lemma-boundedness-better-coherator}.
\end{example}

\section{Cohomologie et changement de base, IV}
\label{spaces-perfect-section-cohomology-and-base-change-perfect}

\noindent
Cette section est l'analogue de Cat\'egories d\'eriv\'ees des sch\'emas,
section \ref{perfect-section-cohomology-and-base-change-perfect}.

\begin{lemma}
\label{spaces-perfect-lemma-cohomology-base-change}
Soit $S$ un sch\'ema. Soit $f:X\to Y$ un morphisme quasi-compact et
quasi-s\'epar\'e d'espaces alg\'ebriques sur $S$. Pour
$E\in D_\QCoh(\mathcal{O}_X)$ et $K\in D_\QCoh(\mathcal{O}_Y)$, on a
$$
Rf_*(E) \otimes_{\mathcal{O}_Y}^\mathbf{L} K =
Rf_*(E \otimes_{\mathcal{O}_X}^\mathbf{L} Lf^*K).
$$
\end{lemma}

\begin{proof}
Sans aucune hypoth\`ese, il existe un morphisme
$Rf_*(E)\otimes_{\mathcal{O}_Y}^\mathbf{L}K\to
Rf_*(E\otimes_{\mathcal{O}_X}^\mathbf{L}Lf^*K)$. C'est en effet l'adjoint du
morphisme canonique
$$
Lf^*(Rf_*(E) \otimes_{\mathcal{O}_Y}^\mathbf{L} K) =
Lf^*(Rf_*(E)) \otimes_{\mathcal{O}_X}^\mathbf{L} Lf^*K
\longrightarrow
E \otimes_{\mathcal{O}_X}^\mathbf{L} Lf^*K
$$
provenant du morphisme $Lf^*Rf_*E\to E$; voir Cohomologie sur les sites,
lemmes \ref{sites-cohomology-lemma-pullback-tensor-product} et
\ref{sites-cohomology-lemma-adjoint}. Pour v\'erifier qu'il s'agit d'un
isomorphisme, on peut travailler localement pour la topologie \'etale sur $Y$.
On se ram\`ene donc au cas o\`u $Y$ est un sch\'ema affine.

\medskip\noindent
Supposons que $K=\bigoplus K_i$ soit une somme directe de complexes
$K_i\in D_\QCoh(\mathcal{O}_Y)$. Si l'assertion est vraie pour chacun des
$K_i$, elle l'est pour $K$. En effet, les foncteurs $Lf^*$ et
$\otimes^\mathbf{L}$ commutent aux sommes directes par construction, et $Rf_*$
commute aux sommes directes pour les complexes \`a faisceaux de cohomologie
quasi-coh\'erents, d'apr\`es le lemme
\ref{spaces-perfect-lemma-quasi-coherence-pushforward-direct-sums}. De plus, si
$K\to L\to M\to K[1]$ est un triangle distingu\'e de $D_\QCoh(Y)$ et si
l'assertion du lemme est vraie pour deux des trois objets $K,L,M$, elle l'est
pour le troisi\`eme, car les foncteurs en jeu sont des foncteurs exacts de
cat\'egories triangul\'ees.

\medskip\noindent
Supposons $Y$ affine, disons $Y=\Spec(A)$. Le foncteur
$\widetilde{\ }:D(A)\to D_\QCoh(\mathcal{O}_Y)$ est une \'equivalence d'apr\`es
le lemme \ref{spaces-perfect-lemma-derived-quasi-coherent-small-etale-site} et Cat\'egories
d\'eriv\'ees des sch\'emas, lemme
\ref{perfect-lemma-affine-compare-bounded}. Notons $T$ la propri\'et\'e, pour
$K\in D(A)$, selon laquelle l'assertion du lemme est vraie pour
$\widetilde K$. La discussion pr\'ec\'edente et Compl\'ements d'alg\`ebre,
remarque \ref{more-algebra-remark-P-resolution}, montrent qu'il suffit de
prouver $T$ pour $A[k]$. Cela ach\`eve la d\'emonstration, puisque l'assertion
du lemme est claire pour les translat\'es du faisceau structural.
\end{proof}

\begin{definition}
\label{spaces-perfect-definition-tor-independent}
Soit $S$ un sch\'ema. Soit $B$ un espace alg\'ebrique sur $S$. Soient $X$ et
$Y$ des espaces alg\'ebriques sur $B$. On dit que $X$ et $Y$ sont
{\it Tor-ind\'ependants sur $B$} si, pour tout diagramme commutatif
$$
\xymatrix{
\Spec(k) \ar[d]_{\overline{y}} \ar[dr]_{\overline{b}} \ar[r]_-{\overline{x}} &
X \ar[d] \\
Y \ar[r] & B
}
$$
de points g\'eom\'etriques, les anneaux
$\mathcal{O}_{X,\overline{x}}$ et $\mathcal{O}_{Y,\overline{y}}$ sont
Tor-ind\'ependants sur $\mathcal{O}_{B,\overline{b}}$ (voir Compl\'ements
d'alg\`ebre, d\'efinition \ref{more-algebra-definition-tor-independent}).
\end{definition}

\noindent
Le lemme suivant montre en particulier que cette d\'efinition co\"incide avec
la d\'efinition donn\'ee pour les espaces alg\'ebriques repr\'esentables.

\begin{lemma}
\label{spaces-perfect-lemma-tor-independent}
Soit $S$ un sch\'ema. Soit $B$ un espace alg\'ebrique sur $S$. Soient $X$ et
$Y$ des espaces alg\'ebriques sur $B$. Les conditions suivantes sont
\'equivalentes :
\begin{enumerate}
\item $X$ et $Y$ sont Tor-ind\'ependants sur $B$;
\item pour tout diagramme commutatif
$$
\xymatrix{
U \ar[d] \ar[r] & W \ar[d] & V \ar[d] \ar[l] \\
X \ar[r] & B & Y \ar[l]
}
$$
dont les fl\`eches verticales sont \'etales, $U$ et $V$ sont
Tor-ind\'ependants sur $W$;
\item il existe un diagramme commutatif comme en (2) tel que (a) $W\to B$
soit \'etale surjectif, (b) $U\to X\times_BW$ soit \'etale surjectif,
(c) $V\to Y\times_BW$ soit \'etale surjectif, et tel que $U$ et $V$ soient
Tor-ind\'ependants sur $W$;
\item il existe un diagramme commutatif comme en (3), avec $U,V,W$ des
sch\'emas, tel que $U$ et $V$ soient Tor-ind\'ependants sur $W$ au sens de
Cat\'egories d\'eriv\'ees des sch\'emas, d\'efinition
\ref{perfect-definition-tor-independent}.
\end{enumerate}
\end{lemma}

\begin{proof}
Pour un morphisme \'etale $\varphi:U\to X$ d'espaces alg\'ebriques et un point
g\'eom\'etrique $\overline u$, le morphisme d'anneaux locaux
$\mathcal{O}_{X,\varphi(\overline u)}\to\mathcal{O}_{U,\overline u}$ est un
isomorphisme. On en d\'eduit l'\'equivalence de (1) et (2), ainsi que
l'implication (1) $\Rightarrow$ (3). Supposons (3) et choisissons un diagramme
de points g\'eom\'etriques comme dans la d\'efinition
\ref{spaces-perfect-definition-tor-independent}. Les hypoth\`eses permettent d'abord de
relever $\overline b$ en un point g\'eom\'etrique $\overline w$ de $W$, puis de
relever le point g\'eom\'etrique $(\overline x,\overline b)$ en un point
g\'eom\'etrique $\overline u$ de $U$, et enfin de relever
$(\overline y,\overline b)$ en un point g\'eom\'etrique $\overline v$ de $V$.
Pour trouver ces rel\`evements, on utilise Propri\'et\'es des espaces, lemme
\ref{spaces-properties-lemma-geometric-lift-to-usual}. La remarque pr\'ec\'edente
sur les anneaux locaux montre alors que la condition de la d\'efinition est
satisfaite pour le diagramme donn\'e.

\medskip\noindent
Ces points pr\'eliminaires \'etant fix\'es, l'assertion (4) revient 
manifestement \`a dire que la d\'efinition \ref{spaces-perfect-definition-tor-independent}
co\"incide avec Cat\'egories d\'eriv\'ees des sch\'emas, d\'efinition
\ref{perfect-definition-tor-independent}, lorsque $X,Y,B$ sont des sch\'emas.

\medskip\noindent
Soient $\overline x,\overline b,\overline y$ comme dans la d\'efinition
\ref{spaces-perfect-definition-tor-independent}, au-dessus des points $x,y,b$. Rappelons que
$\mathcal{O}_{X,\overline{x}}=\mathcal{O}_{X,x}^{sh}$ (Propri\'et\'es des
espaces, lemme \ref{spaces-properties-lemma-describe-etale-local-ring}), et de
m\^eme pour les deux autres. D'apr\`es Alg\`ebre, lemme
\ref{algebra-lemma-strictly-henselian-functorial-improve},
$\mathcal{O}_{X,\overline x}$ est un hens\'elis\'e strict de
$\mathcal{O}_{X,x}\otimes_{\mathcal{O}_{B,b}}\mathcal{O}_{B,\overline b}$.
En particulier, le morphisme d'anneaux
$$
\mathcal{O}_{X, x} \otimes_{\mathcal{O}_{B, b}} \mathcal{O}_{B, \overline{b}}
\longrightarrow
\mathcal{O}_{X, \overline{x}}
$$
est plat (Compl\'ements d'alg\`ebre, lemme
\ref{more-algebra-lemma-dumb-properties-henselization}). D'apr\`es Compl\'ements
d'alg\`ebre, lemme \ref{more-algebra-lemma-tor-independent-flat}, on a
$$
\text{Tor}_i^{\mathcal{O}_{B, b}}(\mathcal{O}_{X, x}, \mathcal{O}_{Y, y})
\otimes_{\mathcal{O}_{X, x} \otimes_{\mathcal{O}_{B, b}} \mathcal{O}_{Y, y}}
(\mathcal{O}_{X, \overline{x}} \otimes_{\mathcal{O}_{B, \overline{b}}}
\mathcal{O}_{Y, \overline y})
=
\text{Tor}_i^{\mathcal{O}_{B, \overline{b}}}(
\mathcal{O}_{X, \overline{x}}, \mathcal{O}_{Y, \overline{y}}).
$$
Par cons\'equent, si $X$ et $Y$ sont Tor-ind\'ependants sur $B$ comme
sch\'emas, ils le sont aussi comme espaces alg\'ebriques sur $B$.

\medskip\noindent
R\'eciproquement, on peut supposer $X,Y,B$ affines. Le morphisme d'anneaux
$$
\mathcal{O}_{X, x} \otimes_{\mathcal{O}_{B, b}} \mathcal{O}_{Y, y}
\longrightarrow
\mathcal{O}_{X, \overline{x}} \otimes_{\mathcal{O}_{B, \overline{b}}}
\mathcal{O}_{Y, \overline y}
$$
est plat d'apr\`es les observations pr\'ec\'edentes. De plus, l'image du
morphisme sur les spectres contient tous les id\'eaux premiers
$\mathfrak s\subset
\mathcal{O}_{X,x}\otimes_{\mathcal{O}_{B,b}}\mathcal{O}_{Y,y}$ situ\'es
au-dessus de $\mathfrak m_x$ et $\mathfrak m_y$. Cette remarque et la formule
pr\'ec\'edente pour les Tor montrent que, si $X$ et $Y$ sont
Tor-ind\'ependants sur $B$ comme espaces alg\'ebriques, alors
$$
\text{Tor}_i^{\mathcal{O}_{B, b}}
(\mathcal{O}_{X, x}, \mathcal{O}_{Y, y})_\mathfrak s = 0
$$
pour tout $i>0$ et tout $\mathfrak s$ comme ci-dessus. Le lemme de Compl\'ements
d'alg\`ebre \ref{more-algebra-lemma-tor-independent}, appliqu\'e aux morphismes
d'anneaux
$\Gamma(B,\mathcal{O}_B)\to\Gamma(X,\mathcal{O}_X)$ et
$\Gamma(B,\mathcal{O}_B)\to\Gamma(Y,\mathcal{O}_Y)$, montre alors que $X$ et
$Y$ sont Tor-ind\'ependants sur $B$.
\end{proof}

\begin{lemma}
\label{spaces-perfect-lemma-compare-base-change}
Soit $S$ un sch\'ema. Soit $g:Y'\to Y$ un morphisme d'espaces alg\'ebriques
sur $S$. Soit $f:X\to Y$ un morphisme quasi-compact et quasi-s\'epar\'e
d'espaces alg\'ebriques sur $S$. Consid\'erons le diagramme de changement de
base
$$
\xymatrix{
X' \ar[r]_{g'} \ar[d]_{f'} &
X \ar[d]^f \\
Y' \ar[r]^g &
Y
}
$$
Si $X$ et $Y'$ sont Tor-ind\'ependants sur $Y$, alors, pour tout
$E\in D_\QCoh(\mathcal{O}_X)$, on a
$Rf'_*L(g')^*E=Lg^*Rf_*E$.
\end{lemma}

\begin{proof}
Pour tout objet $E$ de $D(\mathcal{O}_X)$, la remarque de Cohomologie sur les
sites \ref{sites-cohomology-remark-base-change} fournit un morphisme canonique
de changement de base $Lg^*Rf_*E\to Rf'_*L(g')^*E$. Pour v\'erifier que c'est
un isomorphisme, on peut travailler localement pour la topologie \'etale sur
$Y'$. On peut donc supposer que $g:Y'\to Y$ est un morphisme de sch\'emas
affines. En particulier, $g$ est affine et il suffit de montrer que
$$
Rg_*Lg^*Rf_*E \to Rg_*Rf'_*L(g')^*E = Rf_*(Rg'_* L(g')^* E)
$$
est un isomorphisme; voir le lemme \ref{spaces-perfect-lemma-affine-morphism} (et utiliser les
lemmes \ref{spaces-perfect-lemma-quasi-coherence-pullback},
\ref{spaces-perfect-lemma-quasi-coherence-tensor-product} et
\ref{spaces-perfect-lemma-quasi-coherence-direct-image} pour voir que les objets
$Rf'_*L(g')^*E$ et $Lg^*Rf_*E$ ont des faisceaux de cohomologie
quasi-coh\'erents). Le morphisme $g'$ est lui aussi affine (Morphismes
d'espaces, lemme \ref{spaces-morphisms-lemma-base-change-affine}). D'apr\`es le
lemme \ref{spaces-perfect-lemma-affine-morphism-pull-push}, le morphisme devient
$$
Rf_*E \otimes_{\mathcal{O}_Y}^\mathbf{L} g_*\mathcal{O}_{Y'}
\longrightarrow
Rf_*(E \otimes_{\mathcal{O}_X}^\mathbf{L} g'_*\mathcal{O}_{X'}).
$$
On a $g'_*\mathcal{O}_{X'}=f^*g_*\mathcal{O}_{Y'}$. Le lemme
\ref{spaces-perfect-lemma-cohomology-base-change} ram\`ene donc la d\'emonstration \`a
$Lf^*g_*\mathcal{O}_{Y'}=f^*g_*\mathcal{O}_{Y'}$. Cette \'egalit\'e r\'esulte de
l'hypoth\`ese selon laquelle $X$ et $Y'$ sont Tor-ind\'ependants sur $Y$. Pour
la v\'erifier, on peut travailler localement pour la topologie \'etale sur $X$
et supposer aussi $X$ affine. \'Ecrivons $X=\Spec(A)$, $Y=\Spec(R)$ et
$Y'=\Spec(R')$. L'hypoth\`ese entra\^ine que $A$ et $R'$ sont
Tor-ind\'ependants sur $R$ (voir le lemme \ref{spaces-perfect-lemma-tor-independent} et
Compl\'ements d'alg\`ebre, lemme \ref{more-algebra-lemma-tor-independent}),
c'est-\`a-dire que $\text{Tor}_i^R(A,R')=0$ pour $i>0$. Autrement dit,
$A\otimes_R^\mathbf{L}R'=A\otimes_RR'$, ce qui signifie exactement que
$Lf^*g_*\mathcal{O}_{Y'}=f^*g_*\mathcal{O}_{Y'}$.
\end{proof}

\noindent
Le lemme suivant sera utilis\'e dans le chapitre consacr\'e aux complexes
dualisants.

\begin{lemma}
\label{spaces-perfect-lemma-affine-morphism-and-hom-out-of-perfect}
Soit $g:S'\to S$ un morphisme de sch\'emas affines. Consid\'erons un carr\'e
cart\'esien
$$
\xymatrix{
X' \ar[r]_{g'} \ar[d]_{f'} & X \ar[d]^f \\
S' \ar[r]^g & S
}
$$
d'espaces alg\'ebriques quasi-compacts et quasi-s\'epar\'es. Supposons $g$ et
$f$ Tor-ind\'ependants. \'Ecrivons $S=\Spec(R)$ et $S'=\Spec(R')$. Pour
$M,K\in D(\mathcal{O}_X)$, le morphisme canonique
$$
R\Hom_X(M, K) \otimes^\mathbf{L}_R R'
\longrightarrow
R\Hom_{X'}(L(g')^*M, L(g')^*K)
$$
de $D(R')$ est un isomorphisme dans les deux cas suivants :
\begin{enumerate}
\item $M\in D(\mathcal{O}_X)$ est parfait et $K\in D_\QCoh(X)$;
\item $M\in D(\mathcal{O}_X)$ est pseudo-coh\'erent,
$K\in D_\QCoh^+(X)$ et $R'$ est de Tor-dimension finie sur $R$.
\end{enumerate}
\end{lemma}

\begin{proof}
Il existe dans $D(\Gamma(X,\mathcal{O}_X))$ un morphisme canonique
$R\Hom_X(M,K)\to R\Hom_{X'}(L(g')^*M,L(g')^*K)$ entre complexes
d'homomorphismes globaux; voir Cohomologie sur les sites, section
\ref{sites-cohomology-section-global-RHom}. Par restriction des scalaires, on
peut le regarder comme un morphisme de $D(R)$. L'adjonction entre le foncteur de
restriction des scalaires et $-\otimes_R^\mathbf{L}R'$ donne alors le morphisme
affich\'e dans l'\'enonc\'e. Une fois ce morphisme d\'efini, il suffit de montrer qu'il est un
isomorphisme dans la cat\'egorie d\'eriv\'ee des groupes ab\'eliens.

\medskip\noindent
Le membre de droite est \'egal \`a
$$
R\Hom_X(M, R(g')_*L(g')^*K) =
R\Hom_X(M, K \otimes_{\mathcal{O}_X}^\mathbf{L} g'_*\mathcal{O}_{X'})
$$
d'apr\`es le lemme \ref{spaces-perfect-lemma-affine-morphism-pull-push}. Dans les deux cas,
le complexe $R\SheafHom(M,K)$ est un objet de $D_\QCoh(\mathcal{O}_X)$
d'apr\`es le lemme \ref{spaces-perfect-lemma-quasi-coherence-internal-hom}. Il existe un
morphisme naturel
$$
R\SheafHom(M, K) \otimes_{\mathcal{O}_X}^\mathbf{L} g'_*\mathcal{O}_{X'}
\longrightarrow
R\SheafHom(M, K \otimes_{\mathcal{O}_X}^\mathbf{L} g'_*\mathcal{O}_{X'})
$$
qui est un isomorphisme dans les deux cas, d'apr\`es le lemme
\ref{spaces-perfect-lemma-internal-hom-evaluate-tensor-isomorphism}. Pour voir que ce lemme
s'applique dans le cas (2), remarquons que
$g'_*\mathcal{O}_{X'}=Rg'_*\mathcal{O}_{X'}=
Lf^*g_*\mathcal{O}_{S'}$, la seconde \'egalit\'e r\'esultant du lemme
\ref{spaces-perfect-lemma-compare-base-change}. Les lemmes Cat\'egories d\'eriv\'ees des
sch\'emas \ref{perfect-lemma-tor-dimension-affine},
\ref{spaces-perfect-lemma-descend-tor-amplitude} et Cohomologie sur les sites
\ref{sites-cohomology-lemma-tor-amplitude-pullback} montrent alors que
$g'_*\mathcal{O}_{X'}$ est de Tor-dimension finie. Dans les deux cas, en
rempla\c{c}ant $K$ par $R\SheafHom(M,K)$, on se ram\`ene donc \`a montrer que
$$
R\Gamma(X, K) \otimes^\mathbf{L}_R R' \longrightarrow
R\Gamma(X, K \otimes^\mathbf{L}_{\mathcal{O}_X} g'_*\mathcal{O}_{X'})
$$
est un isomorphisme. Le membre de gauche est \'egal \`a
$R\Gamma(X',L(g')^*K)$ d'apr\`es le lemme
\ref{spaces-perfect-lemma-affine-morphism-pull-push}. Le r\'esultat d\'ecoule donc du lemme
\ref{spaces-perfect-lemma-compare-base-change}.
\end{proof}

\begin{remark}
\label{spaces-perfect-remark-multiplication-map}
Avec les notations du lemme
\ref{spaces-perfect-lemma-affine-morphism-and-hom-out-of-perfect}, le diagramme
$$
\xymatrix{
R\Hom_X(M, Rg'_*L) \otimes_R^\mathbf{L} R' \ar[r] \ar[d]_\mu &
R\Hom_{X'}(L(g')^*M, L(g')^*Rg'_*L) \ar[d]^a \\
R\Hom_X(M, R(g')_*L) \ar@{=}[r] &
R\Hom_{X'}(L(g')^*M, L)
}
$$
est commutatif; la fl\`eche horizontale sup\'erieure est le morphisme du lemme,
$\mu$ est le morphisme de multiplication, et $a$ provient du morphisme
d'adjonction $L(g')^*Rg'_*L\to L$. Pour tout $K'\in D(R')$, le morphisme de
multiplication est le morphisme d'adjonction
$K'\otimes_R^\mathbf{L}R'\to K'$.
\end{remark}

\begin{lemma}
\label{spaces-perfect-lemma-tor-independence-and-tor-amplitude}
Soit $S$ un sch\'ema. Consid\'erons un carr\'e cart\'esien d'espaces alg\'ebriques
$$
\xymatrix{
X' \ar[r]_{g'} \ar[d]_{f'} & X \ar[d]^f \\
Y' \ar[r]^g & Y
}
$$
sur $S$. Supposons $g$ et $f$ Tor-ind\'ependants.
\begin{enumerate}
\item Si $E\in D(\mathcal{O}_X)$ est de Tor-amplitude contenue dans $[a,b]$
comme complexe de $f^{-1}\mathcal{O}_Y$-modules, alors $L(g')^*E$ est de
Tor-amplitude contenue dans $[a,b]$ comme complexe de
$(f')^{-1}\mathcal{O}_{Y'}$-modules.
\item Si $\mathcal{G}$ est un $\mathcal{O}_X$-module plat sur $Y$, alors
$L(g')^*\mathcal{G}=(g')^*\mathcal{G}$.
\end{enumerate}
\end{lemma}

\begin{proof}
On peut calculer la Tor-dimension sur les fibres; voir Cohomologie sur les
sites, lemme \ref{sites-cohomology-lemma-tor-amplitude-stalk}, et Propri\'et\'es
des espaces, th\'eor\`eme \ref{spaces-properties-theorem-exactness-stalks}.
Si $\overline{x}'$ est un point g\'eom\'etrique de $X'$ d'image $\overline x$
dans $X$, alors
$$
(L(g')^*E)_{\overline{x}'} =
E_{\overline{x}}
\otimes_{\mathcal{O}_{X, \overline{x}}}^\mathbf{L}
\mathcal{O}_{X', \overline{x}'}.
$$
Soient $\overline y'$ dans $Y'$ et $\overline y$ dans $Y$ les images de
$\overline x'$ et $\overline x$. Comme $X$ et $Y'$ sont Tor-ind\'ependants sur
$Y$, Compl\'ements d'alg\`ebre, lemme
\ref{more-algebra-lemma-base-change-comparison}, montre que le membre de droite
de la formule pr\'ec\'edente est \'egal \`a
$E_{\overline{x}}
\otimes_{\mathcal{O}_{Y, \overline{y}}}^\mathbf{L}
\mathcal{O}_{Y', \overline{y}'}$
dans $D(\mathcal{O}_{Y',\overline y'})$. L'assertion (1) r\'esulte donc de
Compl\'ements d'alg\`ebre, lemme
\ref{more-algebra-lemma-pull-tor-amplitude}. Pour (2), il suffit d'observer que
la platitude de $\mathcal G$ \'equivaut \`a ce que $\mathcal G[0]$ soit de
Tor-amplitude contenue dans $[0,0]$, puis d'appliquer (1).
\end{proof}

\section{Cohomologie et changement de base, V}
\label{spaces-perfect-section-cohomology-base-change}

\noindent
Cette section est l'analogue de Catégories dérivées des schémas, section
\ref{perfect-section-cohomology-base-change}. Dans la
section \ref{spaces-perfect-section-cohomology-and-base-change-perfect},
nous avons établi un théorème de changement de base sous l'hypothèse que les
morphismes sont Tor-indépendants. Même dans le cas affine, aucun théorème de
changement de base n'est possible sans une telle condition ; voir
Compléments d'algèbre, section \ref{more-algebra-section-tor-independence}.
Nous étudions ici les conditions sous lesquelles on peut obtenir un résultat
de changement de base « complexe par complexe ».

\medskip\noindent
Pour mettre cela en œuvre, soit $S$ un schéma de base et
supposons donné un diagramme commutatif
$$
\xymatrix{
X' \ar[r]_{g'} \ar[d]_{f'} &
X \ar[d]^f \\
Y' \ar[r]^g &
Y
}
$$
d'espaces algébriques sur $S$ (nous supposerons généralement qu'il est cartésien).
Soit $K \in D_\QCoh(\mathcal{O}_X)$
et soit $L(g')^*K \to K'$ un morphisme dans $D_\QCoh(\mathcal{O}_{X'})$.
Pour un point géométrique $\overline{x}'$ de $X'$, considérons les points
géométriques
$\overline{x} = g'(\overline{x}')$,
$\overline{y}' = f'(\overline{x}')$,
$\overline{y} = f(\overline{x}) = g(\overline{y}')$
de $X$, $Y'$ et $Y$.
Nous pouvons alors considérer les morphismes
$$
K_{\overline{x}}
\otimes_{\mathcal{O}_{Y, \overline{y}}}^\mathbf{L}
\mathcal{O}_{Y', \overline{y}'}
\to
K_{\overline{x}}
\otimes_{\mathcal{O}_{X, \overline{x}}}^\mathbf{L}
\mathcal{O}_{X', \overline{x}'} \to
K'_{\overline{x}'}
$$
où le premier morphisme est celui de Compléments d'algèbre,
équation (\ref{more-algebra-equation-comparison-map}),
et le second provient de
$(L(g')^*K)_{\overline{x}'} =
K_{\overline{x}} \otimes_{\mathcal{O}_{X, \overline{x}}}^\mathbf{L}
\mathcal{O}_{X', \overline{x}'}$
et du morphisme donné $L(g')^*K \to K'$.
Pour tout $i \in \mathbf{Z}$, on obtient une structure de
$\mathcal{O}_{X, \overline{x}} \otimes_{\mathcal{O}_{Y, \overline{y}}}
\mathcal{O}_{Y', \overline{y}'}$-module sur
$H^i(K_{\overline{x}}
\otimes_{\mathcal{O}_{Y, \overline{y}}}^\mathbf{L}
\mathcal{O}_{Y', \overline{y}'})$.
En regroupant ces données, on obtient des morphismes canoniques
\begin{equation}
\label{spaces-perfect-equation-bc}
H^i(K_{\overline{x}}
\otimes_{\mathcal{O}_{Y, \overline{y}}}^\mathbf{L}
\mathcal{O}_{Y', \overline{y}'})
\otimes_{(\mathcal{O}_{X, \overline{x}}
\otimes_{\mathcal{O}_{Y, \overline{y}}}
\mathcal{O}_{Y', \overline{y}'})}
\mathcal{O}_{X', \overline{x}'}
\longrightarrow
H^i(K'_{\overline{x}'})
\end{equation}
de $\mathcal{O}_{X', \overline{x}'}$-modules.

\begin{lemma}
\label{spaces-perfect-lemma-single-complex-base-change-condition}
Soit $S$ un schéma. Soit
$$
\xymatrix{
X' \ar[r]_{g'} \ar[d]_{f'} &
X \ar[d]^f \\
Y' \ar[r]^g &
Y
}
$$
un diagramme cartésien d'espaces algébriques sur $S$.
Soit $K \in D_\QCoh(\mathcal{O}_X)$ et soit $L(g')^*K \to K'$
un morphisme dans $D_\QCoh(\mathcal{O}_{X'})$. Les conditions suivantes sont équivalentes :
\begin{enumerate}
\item pour tout $x' \in X'$ et tout $i \in \mathbf{Z}$, le morphisme (\ref{spaces-perfect-equation-bc})
est un isomorphisme ;
\item pour tout diagramme commutatif
$$
\xymatrix{
& U \ar[d] \ar[rd]^a \\
V' \ar[r] \ar[rd]^c & V \ar[rd]^b & X \ar[d]^f \\
& Y' \ar[r]^g & Y
}
$$
où $a, b, c$ sont étales, où $U, V, V'$ sont des schémas et où $U' = V' \times_V U$,
les conditions équivalentes de
Catégories dérivées des schémas, lemme
\ref{perfect-lemma-single-complex-base-change-condition},
sont satisfaites pour $(U \to X)^*K$ et $(U' \to X')^*K'$ ;
\item il existe un diagramme comme en (2) tel que $U' \to X'$ soit surjectif.
\end{enumerate}
\end{lemma}

\begin{proof}
Notons que la condition (1) est locale pour la topologie étale sur $X'$.
En explicitant les implications formelles des résultats connus, on voit qu'il suffit de prouver
que la condition (1) du présent lemme équivaut à la condition
(1) de Catégories dérivées des schémas, lemme
\ref{perfect-lemma-single-complex-base-change-condition},
lorsque $X, Y, Y', X'$ sont représentables par des schémas
$X_0, Y_0, Y'_0, X'_0$. Notons $f_0, g_0, g'_0, f'_0$ les
morphismes entre ces schémas qui correspondent à $f, g, g', f'$.
On peut supposer que $K = \epsilon^*K_0$ et $K' = \epsilon^*K'_0$
pour certains objets $K_0 \in D_\QCoh(\mathcal{O}_{X_0})$
et $K'_0 \in D_\QCoh(\mathcal{O}_{X'_0})$ ; voir le
lemme \ref{spaces-perfect-lemma-derived-quasi-coherent-small-etale-site}.
De plus, le morphisme $L(g')^*K \to K'$ est l'image inverse
d'un morphisme $L(g_0)^*K_0 \to K'_0$, avec les notations de la
remarque \ref{spaces-perfect-remark-match-total-direct-images}.
Rappelons que $\mathcal{O}_{X, \overline{x}}$ est l'hensélisé
strict de $\mathcal{O}_{X, x}$
(Propriétés des espaces, lemme
\ref{spaces-properties-lemma-describe-etale-local-ring})
et que l'on a
$$
K_{\overline{x}} =
K_{0, x} \otimes_{\mathcal{O}_{X, x}}^\mathbf{L} \mathcal{O}_{X, \overline{x}}
\quad\text{et}\quad
K'_{\overline{x}'} =
K'_{0, x'} \otimes_{\mathcal{O}_{X', x'}}^\mathbf{L}
\mathcal{O}_{X', \overline{x}'}
$$
(de façon analogue à Propriétés des espaces, lemme
\ref{spaces-properties-lemma-stalk-quasi-coherent}).
Considérons le diagramme commutatif
$$
\xymatrix{
H^i(K_{\overline{x}}
\otimes_{\mathcal{O}_{Y, \overline{y}}}^\mathbf{L}
\mathcal{O}_{Y', \overline{y}'})
\otimes_{(\mathcal{O}_{X, \overline{x}}
\otimes_{\mathcal{O}_{Y, \overline{y}}}
\mathcal{O}_{Y', \overline{y}'})}
\mathcal{O}_{X', \overline{x}'}
\ar[r] &
H^i(K'_{\overline{x}'}) \\
H^i(K_{0, x} \otimes_{\mathcal{O}_{Y, y}}^\mathbf{L} \mathcal{O}_{Y', y'})
\otimes_{(\mathcal{O}_{X, x} \otimes_{\mathcal{O}_{Y, y}} \mathcal{O}_{Y', y'})}
\mathcal{O}_{X', x'}
\ar[u] \ar[r] &
H^i(K'_{0, x'}) \ar[u]
}
$$
Il faut montrer que le morphisme horizontal inférieur est un isomorphisme si et seulement
si le morphisme horizontal supérieur en est un. Comme
$\mathcal{O}_{X', x'} \to \mathcal{O}_{X', \overline{x}'}$
est fidèlement plat (Compléments d'algèbre, lemme
\ref{more-algebra-lemma-dumb-properties-henselization}),
il suffit de montrer que le morphisme supérieur s'obtient du morphisme
inférieur par le changement de base donné par ce morphisme.
Cela résulte immédiatement des relations entre les fibres
données plus haut pour les objets de droite.
Pour les objets de gauche, il suffit de montrer que
\begin{align*}
&
H^i\left(
(K_{0, x}
\otimes_{\mathcal{O}_{X, x}}^\mathbf{L} \mathcal{O}_{X, \overline{x}})
\otimes_{\mathcal{O}_{Y, \overline{y}}}^\mathbf{L}
\mathcal{O}_{Y', \overline{y}'}\right) \\
& =
H^i(K_{0, x} \otimes_{\mathcal{O}_{Y, y}}^\mathbf{L} \mathcal{O}_{Y', y'})
\otimes_{(\mathcal{O}_{X, x} \otimes_{\mathcal{O}_{Y, y}} \mathcal{O}_{Y', y'})}
(\mathcal{O}_{X, \overline{x}}
\otimes_{\mathcal{O}_{Y, \overline{y}}}
\mathcal{O}_{Y', \overline{y}'})
\end{align*}
Cela résulte de Compléments d'algèbre, lemme
\ref{more-algebra-lemma-lemma-tor-independent-flat-compare}.
Les hypothèses de platitude de ce lemme découlent de ce qui précède
ainsi que d'Algèbre, lemme
\ref{algebra-lemma-strictly-henselian-functorial-improve},
qui entraîne que
$\mathcal{O}_{X, \overline{x}}$ est l'hensélisé strict de
$\mathcal{O}_{X, x} \otimes_{\mathcal{O}_{Y, y}}
\mathcal{O}_{Y, \overline{y}}$
et que
$\mathcal{O}_{Y', \overline{y}'}$ est l'hensélisé strict de
$\mathcal{O}_{Y', y'} \otimes_{\mathcal{O}_{Y, y}}
\mathcal{O}_{Y, \overline{y}}$.
\end{proof}

\begin{lemma}
\label{spaces-perfect-lemma-single-complex-base-change}
Soit $S$ un schéma. Soit
$$
\xymatrix{
X' \ar[r]_{g'} \ar[d]_{f'} &
X \ar[d]^f \\
Y' \ar[r]^g &
Y
}
$$
un diagramme cartésien d'espaces algébriques sur $S$.
Soit $K \in D_\QCoh(\mathcal{O}_X)$ et soit $L(g')^*K \to K'$
un morphisme dans $D_\QCoh(\mathcal{O}_{X'})$. Si
\begin{enumerate}
\item les conditions équivalentes du
lemme \ref{spaces-perfect-lemma-single-complex-base-change-condition} sont satisfaites ; et
\item $f$ est quasi-compact et quasi-séparé,
\end{enumerate}
alors le composé $Lg^*Rf_*K \to Rf'_*L(g')^*K \to Rf'_*K'$
est un isomorphisme.
\end{lemma}

\begin{proof}
Pour vérifier que le morphisme est un isomorphisme, on peut raisonner localement pour la topologie étale sur $Y'$.
On peut donc supposer que $g : Y' \to Y$ est un morphisme de schémas affines.
Dans ce cas, nous utiliserons le principe de récurrence du
lemme \ref{spaces-perfect-lemma-induction-principle}
pour montrer que, pour tout espace algébrique quasi-compact et quasi-séparé
$U$ étale sur $X$,
le morphisme construit de même
$Lg^*R(U \to Y)_*K|_U \to R(U' \to Y')_*K'|_{U'}$
est un isomorphisme. Ici, $U' = X' \times_{g', X} U = Y' \times_{g, Y} U$.

\medskip\noindent
Si $U$ est un schéma (par exemple s'il est affine), le résultat est vrai.
En effet, $Y, Y', U, U'$ sont alors des schémas, et $K$ et $K'$ proviennent
d'objets des catégories dérivées des schémas sous-jacents d'après le
lemme \ref{spaces-perfect-lemma-derived-quasi-coherent-small-etale-site},
et la condition de
Catégories dérivées des schémas,
lemme \ref{perfect-lemma-single-complex-base-change-condition},
est satisfaite pour ces complexes d'après le
lemme \ref{spaces-perfect-lemma-single-complex-base-change-condition}.
Ainsi, en vertu des compatibilités expliquées dans la
remarque \ref{spaces-perfect-remark-match-total-direct-images},
on peut appliquer le résultat dans le cas des schémas, à savoir
Catégories dérivées des schémas, lemme
\ref{perfect-lemma-single-complex-base-change}.

\medskip\noindent
Passons à l'étape de récurrence. Soit $(U \subset W, V \to W)$ un carré
distingué élémentaire, où $W$ est un espace algébrique quasi-compact
et quasi-séparé, étale sur $X$, où $U$ est quasi-compact, où $V$ est affine,
et pour lequel le résultat vaut pour $U$, $V$ et $U \times_W V$.
Pour simplifier les notations, remplaçons $W$ par $X$ (ce qui est permis à ce stade).
Notons $a : U \to Y$, $b : V \to Y$ et $c : U \times_X V \to Y$
les morphismes évidents. Soient $a' : U' \to Y'$, $b' : V' \to Y'$
et $c' : U' \times_{X'} V' \to Y'$ les changements de base de $a$, $b$ et $c$.
À l'aide des triangles distingués de Mayer--Vietoris relatif
(lemme \ref{spaces-perfect-lemma-unbounded-relative-mayer-vietoris}),
on obtient un diagramme commutatif
$$
\xymatrix{
Lg^*Rf_*K \ar[r] \ar[d] & Rf'_*K' \ar[d] \\
Lg^*Ra_*K|_U \oplus Lg^*Rb_*K|_V \ar[r] \ar[d] &
Ra'_* K'|_{U'} \oplus Rb'_* K'|_{V'} \ar[d] \\
Lg^*Rc_*K|_{U \times_X V} \ar[r] \ar[d] &
Rc'_*K'|_{U' \times_{X'} V'} \ar[d] \\
Lg^*Rf_* K[1] \ar[r] &
Rf'_* K'[1]
}
$$
Comme les deuxième et troisième morphismes horizontaux sont des isomorphismes, le premier
l'est aussi (Catégories dérivées, lemme \ref{derived-lemma-third-isomorphism-triangle}),
ce qui achève la démonstration du lemme.
\end{proof}

\begin{lemma}
\label{spaces-perfect-lemma-single-complex-base-change-condition-inherited}
Soit $S$ un schéma. Soit
$$
\xymatrix{
X' \ar[r]_{g'} \ar[d]_{f'} &
X \ar[d]^f \\
S' \ar[r]^g &
S
}
$$
un diagramme cartésien d'espaces algébriques sur $S$.
Soit $K \in D_\QCoh(\mathcal{O}_X)$ et soit $L(g')^*K \to K'$
un morphisme dans $D_\QCoh(\mathcal{O}_{X'})$. Si les conditions équivalentes du
lemme \ref{spaces-perfect-lemma-single-complex-base-change-condition} sont satisfaites, alors :
\begin{enumerate}
\item pour $E \in D_\QCoh(\mathcal{O}_X)$, les conditions
équivalentes du lemme \ref{spaces-perfect-lemma-single-complex-base-change-condition} sont satisfaites
pour $L(g')^*(E \otimes^\mathbf{L} K) \to L(g')^*E \otimes^\mathbf{L} K'$ ;
\item si $E$ est un objet parfait de $D(\mathcal{O}_X)$, les conditions équivalentes du
lemme \ref{spaces-perfect-lemma-single-complex-base-change-condition} sont satisfaites pour
$L(g')^*R\SheafHom(E, K) \to R\SheafHom(L(g')^*E, K')$ ;
\item si $K$ est borné inférieurement et si $E$ est un objet pseudo-cohérent de
$D(\mathcal{O}_X)$, les conditions équivalentes du
lemme \ref{spaces-perfect-lemma-single-complex-base-change-condition} sont satisfaites pour
$L(g')^*R\SheafHom(E, K) \to R\SheafHom(L(g')^*E, K')$.
\end{enumerate}
\end{lemma}

\begin{proof}
L'énoncé a un sens, car les complexes en jeu ont des faisceaux de cohomologie
quasi-cohérents d'après les lemmes
\ref{spaces-perfect-lemma-quasi-coherence-pullback},
\ref{spaces-perfect-lemma-quasi-coherence-tensor-product} et
\ref{spaces-perfect-lemma-quasi-coherence-internal-hom}, ainsi que
Cohomologie sur les sites, lemmes
\ref{sites-cohomology-lemma-pseudo-coherent-pullback} et
\ref{sites-cohomology-lemma-perfect-pullback}.
Cela étant dit, dans le cas (1), on peut vérifier que les morphismes (\ref{spaces-perfect-equation-bc})
sont des isomorphismes en calculant la source et le but
de (\ref{spaces-perfect-equation-bc}) à l'aide de l'associativité du produit tensoriel ; voir
Compléments d'algèbre, lemme \ref{more-algebra-lemma-triple-tensor-product}.
Le morphisme des cas (2) et (3) est le composé
$$
L(g')^*R\SheafHom(E, K) \to R\SheafHom(L(g')^*E, L(g')^*K)
\to R\SheafHom(L(g')^*E, K')
$$
où le premier morphisme est celui de
Cohomologie sur les sites, remarque
\ref{sites-cohomology-remark-prepare-fancy-base-change},
et où le second provient du morphisme donné $L(g')^*K \to K'$.
Pour prouver que les morphismes (\ref{spaces-perfect-equation-bc}) sont des isomorphismes, on représente
$E_x$ par un complexe borné de $\mathcal{O}_{X, x}$-modules projectifs de type fini
dans le cas (2), ou par un complexe borné supérieurement de modules libres de type fini dans le cas (3),
puis on calcule la source et le but du morphisme.
Nous omettons quelques détails.
\end{proof}

\begin{lemma}
\label{spaces-perfect-lemma-base-change-tensor}
Soit $S$ un schéma. Soit $f : X \to Y$ un morphisme quasi-compact et
quasi-séparé d'espaces algébriques sur $S$.
Soit $E \in D_\QCoh(\mathcal{O}_X)$.
Soit $\mathcal{G}^\bullet$ un complexe borné supérieurement de
$\mathcal{O}_X$-modules quasi-cohérents plats sur $Y$.
Alors la formation de
$$
Rf_*(E \otimes^\mathbf{L}_{\mathcal{O}_X} \mathcal{G}^\bullet)
$$
commute à tout changement de base (voir la démonstration pour l'énoncé précis).
\end{lemma}

\begin{proof}
L'énoncé signifie ce qui suit. Soit $g : Y' \to Y$ un morphisme d'espaces
algébriques et considérons le diagramme de changement de base
$$
\xymatrix{
X' \ar[r]_{g'} \ar[d]_{f'} &
X \ar[d]^f \\
Y' \ar[r]^g &
Y
}
$$
autrement dit, $X' = Y' \times_Y X$. Le lemme affirme que
$$
Lg^*Rf_*(E \otimes^\mathbf{L}_{\mathcal{O}_X} \mathcal{G}^\bullet)
\longrightarrow
Rf'_*(L(g')^*E \otimes^\mathbf{L}_{\mathcal{O}_{X'}} (g')^*\mathcal{G}^\bullet)
$$
est un isomorphisme. Notons qu'au second membre nous n'utilisons {\bf pas}
l'image inverse dérivée de $\mathcal{G}^\bullet$.
Pour le démontrer, appliquons les lemmes \ref{spaces-perfect-lemma-single-complex-base-change} et
\ref{spaces-perfect-lemma-single-complex-base-change-condition-inherited} : il
suffit de prouver que le morphisme canonique
$$
L(g')^*\mathcal{G}^\bullet \to (g')^*\mathcal{G}^\bullet
$$
satisfait aux conditions équivalentes du
lemme \ref{spaces-perfect-lemma-single-complex-base-change-condition}.
Cela résulte de la vérification de la condition sur les fibres, où le résultat
découle immédiatement du fait que
$\mathcal{G}^\bullet_{\overline{x}}
\otimes_{\mathcal{O}_{Y, \overline{y}}}
\mathcal{O}_{Y', \overline{y}'}$
calcule le produit tensoriel dérivé en vertu de nos hypothèses sur le complexe
$\mathcal{G}^\bullet$.
\end{proof}

\begin{lemma}
\label{spaces-perfect-lemma-base-change-RHom}
Soit $S$ un schéma. Soit $f : X \to Y$ un morphisme quasi-compact et
quasi-séparé d'espaces algébriques sur $S$. Soit $E$
un objet de $D(\mathcal{O}_X)$. Soit $\mathcal{G}^\bullet$
un complexe de $\mathcal{O}_X$-modules quasi-cohérents. Si
\begin{enumerate}
\item $E$ est parfait, $\mathcal{G}^\bullet$ est borné supérieurement
et $\mathcal{G}^n$ est plat sur $Y$ ; ou si
\item $E$ est pseudo-cohérent, $\mathcal{G}^\bullet$ est borné
et $\mathcal{G}^n$ est plat sur $Y$,
\end{enumerate}
alors la formation de
$$
Rf_*R\SheafHom(E, \mathcal{G}^\bullet)
$$
commute à tout changement de base (voir la démonstration pour l'énoncé précis).
\end{lemma}

\begin{proof}
L'énoncé signifie ce qui suit. Soit $g : Y' \to Y$ un morphisme d'espaces
algébriques et considérons le diagramme de changement de base
$$
\xymatrix{
X' \ar[r]_{g'} \ar[d]_{f'} &
X \ar[d]^f \\
Y' \ar[r]^g &
Y
}
$$
autrement dit, $X' = Y' \times_Y X$. Le lemme affirme que
$$
Lg^*Rf_*R\SheafHom(E, \mathcal{G}^\bullet)
\longrightarrow
R(f')_*R\SheafHom(L(g')^*E, (g')^*\mathcal{G}^\bullet)
$$
est un isomorphisme. Notons qu'au second membre nous n'utilisons {\bf pas}
l'image inverse dérivée de $\mathcal{G}^\bullet$. Pour le démontrer, appliquons les
lemmes \ref{spaces-perfect-lemma-single-complex-base-change} et
\ref{spaces-perfect-lemma-single-complex-base-change-condition-inherited} : il
suffit de prouver que le morphisme canonique
$$
L(g')^*\mathcal{G}^\bullet \to (g')^*\mathcal{G}^\bullet
$$
satisfait aux conditions équivalentes du
lemme \ref{spaces-perfect-lemma-single-complex-base-change-condition}.
Cela a été montré dans la démonstration du lemme \ref{spaces-perfect-lemma-base-change-tensor}.
\end{proof}

\section{Construction de complexes parfaits}
\label{spaces-perfect-section-producing-perfect}

\noindent
Le lemme suivant est notre principal outil technique pour construire des
complexes parfaits. Les variantes ult\'erieures de ce r\'esultat se ram\`eneront
\`a ce lemme par approximation noeth\'erienne.

\begin{lemma}
\label{spaces-perfect-lemma-perfect-direct-image}
Soit $S$ un sch\'ema. Soit $Y$ un espace alg\'ebrique noeth\'erien sur $S$.
Soit $f:X\to Y$ un morphisme d'espaces alg\'ebriques localement de type fini
et quasi-s\'epar\'e. Soit $E\in D(\mathcal{O}_X)$ tel que
\begin{enumerate}
\item $E\in D^b_{\textit{Coh}}(\mathcal{O}_X)$ ;
\item le support de $H^i(E)$ est propre sur $Y$ pour tout $i$ ;
\item $E$ est de Tor-dimension finie comme objet de
$D(f^{-1}\mathcal{O}_Y)$.
\end{enumerate}
Alors $Rf_*E$ est un objet parfait de $D(\mathcal{O}_Y)$.
\end{lemma}

\begin{proof}
D'apr\`es le lemme \ref{spaces-perfect-lemma-direct-image-coherent}, $Rf_*E$ est un objet de
$D^b_{\textit{Coh}}(\mathcal{O}_Y)$. Il est donc pseudo-coh\'erent
(lemme \ref{spaces-perfect-lemma-identify-pseudo-coherent-noetherian}). Il suffit par
cons\'equent de montrer que $Rf_*E$ est de Tor-dimension finie ; voir
Cohomologie sur les sites, lemme \ref{sites-cohomology-lemma-perfect}.
D'apr\`es le lemme \ref{spaces-perfect-lemma-tor-qc-qs}, il suffit de v\'erifier que
$Rf_*(E)\otimes_{\mathcal{O}_Y}^\mathbf{L}\mathcal{F}$ a une cohomologie
uniform\'ement born\'ee pour tout faisceau quasi-coh\'erent $\mathcal{F}$ sur
$Y$. Le fait que ce complexe soit born\'e sup\'erieurement est imm\'ediat, puisque
$Rf_*(E)$ l'est. Soit $T\subset |X|$ la r\'eunion des supports de $H^i(E)$ pour
tout $i$. Les hypoth\`eses (1) et (2), ainsi que le lemme
\ref{spaces-perfect-lemma-union-closed-proper-over-base}, montrent que $T$ est propre
sur $Y$. Il existe donc un sous-espace ouvert quasi-compact
$X'\subset X$ contenant $T$. En posant $f'=f|_{X'}$, on a
$Rf_*(E)=Rf'_*(E|_{X'})$, car la restriction de $E$ \`a $X\setminus T$ est
nulle. On peut ainsi remplacer $X$ par $X'$ et supposer $f$ quasi-compact.
Nous avons suppos\'e $f$ quasi-s\'epar\'e. Par cons\'equent,
$$
Rf_*(E) \otimes_{\mathcal{O}_Y}^\mathbf{L} \mathcal{F} =
Rf_*\left(E \otimes_{\mathcal{O}_X}^\mathbf{L} Lf^*\mathcal{F}\right) =
Rf_*\left(E \otimes_{f^{-1}\mathcal{O}_Y}^\mathbf{L} f^{-1}\mathcal{F}\right)
$$
d'apr\`es le lemme \ref{spaces-perfect-lemma-cohomology-base-change} et Cohomologie sur les
sites, lemme \ref{sites-cohomology-lemma-variant-derived-pullback}.
Par l'hypoth\`ese (3), le complexe
$E\otimes_{f^{-1}\mathcal{O}_Y}^\mathbf{L}f^{-1}\mathcal{F}$ a ses faisceaux
de cohomologie dans un intervalle fini fix\'e, disons $[a,b]$. Son image par
$Rf_*$ a donc sa cohomologie dans $[a,\infty)$, ce qui ach\`eve la d\'emonstration.
\end{proof}

\begin{lemma}
\label{spaces-perfect-lemma-tensor-perfect}
Soit $S$ un sch\'ema. Soit $B$ un espace alg\'ebrique noeth\'erien sur $S$.
Soit $f:X\to B$ un morphisme d'espaces alg\'ebriques localement de type fini
et quasi-s\'epar\'e. Soit $E\in D(\mathcal{O}_X)$ parfait. Soit
$\mathcal{G}^\bullet$ un complexe born\'e de $\mathcal{O}_X$-modules coh\'erents,
plats sur $B$ et dont le support est propre sur $B$. Alors
$K=Rf_*(E\otimes^\mathbf{L}_{\mathcal{O}_X}\mathcal{G}^\bullet)$ est un objet
parfait de $D(\mathcal{O}_B)$.
\end{lemma}

\begin{proof}
L'objet $K$ est parfait d'apr\`es le lemme \ref{spaces-perfect-lemma-perfect-direct-image}.
V\'erifions que les hypoth\`eses de ce lemme sont satisfaites. Localement, $E$
est isomorphe \`a un complexe fini de $\mathcal{O}_X$-modules libres de type
fini. Ainsi, localement,
$E\otimes^\mathbf{L}_{\mathcal{O}_X}\mathcal{G}^\bullet$ est isomorphe \`a un
complexe fini dont les termes sont de la forme
$$
\bigoplus\nolimits_{i = a, \ldots, b} (\mathcal{G}^i)^{\oplus r_i}
$$
pour certains entiers $a,b,r_a,\ldots,r_b$. Il en r\'esulte imm\'ediatement que
les faisceaux de cohomologie
$H^i(E\otimes^\mathbf{L}_{\mathcal{O}_X}\mathcal{G}^\bullet)$ sont coh\'erents.
L'hypoth\`ese sur la Tor-dimension en d\'ecoule elle aussi, puisque
$\mathcal{G}^i$ est plat sur $f^{-1}\mathcal{O}_B$.
\end{proof}

\begin{lemma}
\label{spaces-perfect-lemma-ext-perfect}
Soit $S$ un sch\'ema. Soit $B$ un espace alg\'ebrique noeth\'erien sur $S$.
Soit $f:X\to B$ un morphisme d'espaces alg\'ebriques localement de type fini
et quasi-s\'epar\'e. Soit $E\in D(\mathcal{O}_X)$ parfait. Soit
$\mathcal{G}^\bullet$ un complexe born\'e de $\mathcal{O}_X$-modules coh\'erents,
plats sur $B$ et dont le support est propre sur $B$. Alors
$K=Rf_*R\SheafHom(E,\mathcal{G}^\bullet)$ est un objet parfait de
$D(\mathcal{O}_B)$.
\end{lemma}

\begin{proof}
Comme $E$ est un complexe parfait, il existe un complexe parfait dual
$E^\vee$ ; voir Cohomologie sur les sites, lemme
\ref{sites-cohomology-lemma-dual-perfect-complex}. On a
$R\SheafHom(E,\mathcal{G}^\bullet)=
E^\vee\otimes^\mathbf{L}_{\mathcal{O}_X}\mathcal{G}^\bullet$.
La perfection de $K$ r\'esulte donc du lemme \ref{spaces-perfect-lemma-tensor-perfect}.
\end{proof}

\section{Une formule de projection pour Ext}
\label{spaces-perfect-section-ext}

\noindent
Le lemme \ref{spaces-perfect-lemma-compute-ext} (ou des r\'esultats analogues de la
litt\'erature) est parfois utile pour v\'erifier les propri\'et\'es d'une th\'eorie
des obstructions n\'ecessaires \`a la v\'erification de l'un des crit\`eres d'Artin
pour les foncteurs Quot, les sch\'emas de Hilbert et d'autres probl\`emes de
modules. Supposons que $f:X\to Y$ soit un morphisme propre, plat et de
pr\'esentation finie d'espaces alg\'ebriques, et que
$E\in D(\mathcal{O}_X)$ soit parfait. Le lemme affirme alors que
$$
\Ext^i_X(E, f^*\mathcal{F}) =
\Ext^i_Y((Rf_*E^\vee)^\vee, \mathcal{F})
$$
pour tout faisceau quasi-coh\'erent $\mathcal{F}$ sur $Y$. Sous cette forme,
l'identit\'e prend l'allure d'une formule de projection pour Ext ; de fait, le
r\'esultat d\'ecoule assez facilement du lemme \ref{spaces-perfect-lemma-cohomology-base-change}.

\begin{lemma}
\label{spaces-perfect-lemma-compute-tensor-perfect}
Conservons les hypoth\`eses et les notations du lemme
\ref{spaces-perfect-lemma-tensor-perfect}. Il existe alors des isomorphismes fonctoriels
$$
H^i(B, K \otimes^\mathbf{L}_{\mathcal{O}_B} \mathcal{F})
\longrightarrow
H^i(X, E \otimes^\mathbf{L}_{\mathcal{O}_X}
(\mathcal{G}^\bullet \otimes_{\mathcal{O}_X} f^*\mathcal{F}))
$$
pour $\mathcal{F}$ quasi-coh\'erent sur $B$, compatibles avec les homomorphismes
cobord (voir la d\'emonstration).
\end{lemma}

\begin{proof}
On a
$$
\mathcal{G}^\bullet \otimes_{\mathcal{O}_X}^\mathbf{L} Lf^*\mathcal{F} =
\mathcal{G}^\bullet \otimes_{f^{-1}\mathcal{O}_B}^\mathbf{L} f^{-1}\mathcal{F} =
\mathcal{G}^\bullet \otimes_{f^{-1}\mathcal{O}_B} f^{-1}\mathcal{F} =
\mathcal{G}^\bullet \otimes_{\mathcal{O}_X} f^*\mathcal{F}.
$$
La premi\`ere \`egalit\'e r\'esulte de Cohomologie sur les sites, lemme
\ref{sites-cohomology-lemma-variant-derived-pullback}, la deuxi\`eme du fait que
$\mathcal{G}^n$ est un $f^{-1}\mathcal{O}_B$-module plat, et la troisi\`eme de
la d\'efinition des images inverses. On obtient donc
\begin{align*}
H^i(X, E \otimes^\mathbf{L}_{\mathcal{O}_X}
(\mathcal{G}^\bullet \otimes_{\mathcal{O}_X} f^*\mathcal{F}))
& =
H^i(X, E \otimes^\mathbf{L}_{\mathcal{O}_X} \mathcal{G}^\bullet
\otimes_{\mathcal{O}_X}^\mathbf{L} Lf^*\mathcal{F}) \\
& =
H^i(B,
Rf_*(E \otimes^\mathbf{L}_{\mathcal{O}_X} \mathcal{G}^\bullet
\otimes^\mathbf{L}_{\mathcal{O}_X} Lf^*\mathcal{F})) \\
& =
H^i(B,
Rf_*(E \otimes^\mathbf{L}_{\mathcal{O}_X} \mathcal{G}^\bullet)
\otimes^\mathbf{L}_{\mathcal{O}_B} \mathcal{F}) \\
& =
H^i(B, K \otimes^\mathbf{L}_{\mathcal{O}_B} \mathcal{F}).
\end{align*}
La premi\`ere \`egalit\'e vient de ce qui pr\'ec\`ede, la deuxi\`eme de la suite
spectrale de Leray (Cohomologie sur les sites, remarque
\ref{sites-cohomology-remark-before-Leray}), et la troisi\`eme du lemme
\ref{spaces-perfect-lemma-cohomology-base-change}. La compatibilit\'e avec les homomorphismes
cobord signifie ceci : pour toute suite exacte courte
$0\to\mathcal{F}_1\to\mathcal{F}_2\to\mathcal{F}_3\to0$, les isomorphismes
s'ins\`erent dans des diagrammes commutatifs
$$
\xymatrix{
H^i(B, K \otimes^\mathbf{L}_{\mathcal{O}_B} \mathcal{F}_3)
\ar[r] \ar[d]_\delta &
H^i(X, E \otimes^\mathbf{L}_{\mathcal{O}_X}
(\mathcal{G}^\bullet \otimes_{\mathcal{O}_X} f^*\mathcal{F}_3)) \ar[d]^\delta \\
H^{i + 1}(B, K \otimes^\mathbf{L}_{\mathcal{O}_B} \mathcal{F}_1)
\ar[r] &
H^{i + 1}(X, E \otimes^\mathbf{L}_{\mathcal{O}_X}
(\mathcal{G}^\bullet \otimes_{\mathcal{O}_X} f^*\mathcal{F}_1))
}
$$
o\`u les homomorphismes cobord proviennent du triangle distingu\'e
$$
K \otimes^\mathbf{L}_{\mathcal{O}_B} \mathcal{F}_1 \to
K \otimes^\mathbf{L}_{\mathcal{O}_B} \mathcal{F}_2 \to
K \otimes^\mathbf{L}_{\mathcal{O}_B} \mathcal{F}_3 \to
K \otimes^\mathbf{L}_{\mathcal{O}_B} \mathcal{F}_1[1]
$$
et du triangle distingu\'e de $D(\mathcal{O}_X)$ associ\'e \`a la suite exacte
courte de complexes
$$
0 \to
\mathcal{G}^\bullet \otimes_{\mathcal{O}_X} f^*\mathcal{F}_1 \to
\mathcal{G}^\bullet \otimes_{\mathcal{O}_X} f^*\mathcal{F}_2 \to
\mathcal{G}^\bullet \otimes_{\mathcal{O}_X} f^*\mathcal{F}_3 \to 0.
$$
Cette suite est exacte parce que $\mathcal{G}^n$ est plat sur $B$.
Nous omettons la v\'erification de la commutativit\'e du diagramme affich\'e.
\end{proof}

\begin{lemma}
\label{spaces-perfect-lemma-compute-ext-perfect}
Conservons les hypoth\`eses et les notations du lemme
\ref{spaces-perfect-lemma-ext-perfect}. Il existe alors des isomorphismes fonctoriels
$$
H^i(B, K \otimes^\mathbf{L}_{\mathcal{O}_B} \mathcal{F})
\longrightarrow
\Ext^i_{\mathcal{O}_X}(E,
\mathcal{G}^\bullet \otimes_{\mathcal{O}_X} f^*\mathcal{F})
$$
pour $\mathcal{F}$ quasi-coh\'erent sur $B$, compatibles avec les homomorphismes
cobord (voir la d\'emonstration).
\end{lemma}

\begin{proof}
Comme dans la d\'emonstration du lemme \ref{spaces-perfect-lemma-ext-perfect}, soit $E^\vee$
le complexe parfait dual, et rappelons que
$K=Rf_*(E^\vee\otimes_{\mathcal{O}_X}^\mathbf{L}\mathcal{G}^\bullet)$.
Comme on a aussi
$$
\Ext^i_{\mathcal{O}_X}(E,
\mathcal{G}^\bullet \otimes_{\mathcal{O}_X} f^*\mathcal{F})
=
H^i(X, E^\vee \otimes^\mathbf{L}_{\mathcal{O}_X}
(\mathcal{G}^\bullet \otimes_{\mathcal{O}_X} f^*\mathcal{F})),
$$
par construction de $E^\vee$, l'existence des isomorphismes r\'esulte du
lemme \ref{spaces-perfect-lemma-compute-tensor-perfect} appliqu\'e \`a $E^\vee$ et
$\mathcal{G}^\bullet$. La compatibilit\'e avec les homomorphismes cobord signifie
ceci : pour toute suite exacte courte
$0\to\mathcal{F}_1\to\mathcal{F}_2\to\mathcal{F}_3\to0$, les isomorphismes
s'ins\`erent dans des diagrammes commutatifs
$$
\xymatrix{
H^i(B, K \otimes^\mathbf{L}_{\mathcal{O}_B} \mathcal{F}_3)
\ar[r] \ar[d]_\delta &
\Ext^i_{\mathcal{O}_X}(E,
\mathcal{G}^\bullet \otimes_{\mathcal{O}_X} f^*\mathcal{F}_3) \ar[d]^\delta \\
H^{i + 1}(B, K \otimes^\mathbf{L}_{\mathcal{O}_B} \mathcal{F}_1)
\ar[r] &
\Ext^{i + 1}_{\mathcal{O}_X}(E,
\mathcal{G}^\bullet \otimes_{\mathcal{O}_X} f^*\mathcal{F}_1)
}
$$
o\`u les homomorphismes cobord proviennent du triangle distingu\'e
$$
K \otimes^\mathbf{L}_{\mathcal{O}_B} \mathcal{F}_1 \to
K \otimes^\mathbf{L}_{\mathcal{O}_B} \mathcal{F}_2 \to
K \otimes^\mathbf{L}_{\mathcal{O}_B} \mathcal{F}_3 \to
K \otimes^\mathbf{L}_{\mathcal{O}_B} \mathcal{F}_1[1]
$$
et du triangle distingu\'e de $D(\mathcal{O}_X)$ associ\'e \`a la suite exacte
courte de complexes
$$
0 \to
\mathcal{G}^\bullet \otimes_{\mathcal{O}_X} f^*\mathcal{F}_1 \to
\mathcal{G}^\bullet \otimes_{\mathcal{O}_X} f^*\mathcal{F}_2 \to
\mathcal{G}^\bullet \otimes_{\mathcal{O}_X} f^*\mathcal{F}_3 \to 0.
$$
Cette suite est exacte parce que $\mathcal{G}^n$ est plat sur $B$.
Nous omettons la v\'erification de la commutativit\'e du diagramme affich\'e.
\end{proof}

\begin{lemma}
\label{spaces-perfect-lemma-compute-ext}
Soit $S$ un sch\'ema. Soit $f:X\to B$ un morphisme d'espaces alg\'ebriques sur
$S$, soit $E\in D(\mathcal{O}_X)$, et soit $\mathcal{G}^\bullet$ un complexe de
$\mathcal{O}_X$-modules. Supposons que
\begin{enumerate}
\item $B$ soit noeth\'erien ;
\item $f$ soit localement de type fini et quasi-s\'epar\'e ;
\item $E\in D^-_{\textit{Coh}}(\mathcal{O}_X)$ ;
\item $\mathcal{G}^\bullet$ soit un complexe born\'e de
$\mathcal{O}_X$-modules coh\'erents, plats sur $B$ et dont le support est propre
sur $B$.
\end{enumerate}
Alors les deux assertions suivantes sont vraies :
\begin{enumerate}
\item[(A)] pour tout $m\in\mathbf{Z}$, il existe un objet parfait $K$ de
$D(\mathcal{O}_B)$ et des homomorphismes fonctoriels
$$
\alpha^i_\mathcal{F} :
\Ext^i_{\mathcal{O}_X}(E,
\mathcal{G}^\bullet \otimes_{\mathcal{O}_X} f^*\mathcal{F})
\longrightarrow
H^i(B, K \otimes^\mathbf{L}_{\mathcal{O}_B} \mathcal{F})
$$
pour $\mathcal{F}$ quasi-coh\'erent sur $B$, compatibles avec les homomorphismes
cobord (voir la d\'emonstration), tels que $\alpha^i_\mathcal{F}$ soit un
isomorphisme pour $i\leq m$ ;
\item[(B)] il existe $L\in D(\mathcal{O}_B)$ pseudo-coh\'erent et des
isomorphismes fonctoriels
$$
\Ext^i_{\mathcal{O}_B}(L, \mathcal{F}) \longrightarrow
\Ext^i_{\mathcal{O}_X}(E,
\mathcal{G}^\bullet \otimes_{\mathcal{O}_X} f^*\mathcal{F})
$$
pour $\mathcal{F}$ quasi-coh\'erent sur $B$, compatibles avec les homomorphismes
cobord.
\end{enumerate}
\end{lemma}

\begin{proof}
D\'emonstration de (A). Supposons $\mathcal{G}^i$ non nul seulement pour
$i\in[a,b]$. On peut remplacer $X$ par un voisinage ouvert quasi-compact de la
r\'eunion des supports des $\mathcal{G}^i$. On peut donc supposer $X$
noeth\'erien. Dans ce cas, $X$ et $f$ sont quasi-compacts et quasi-s\'epar\'es.
Choisissons une approximation $P\to E$ par un complexe parfait $P$ de
$(X,E,-m-1+a)$ (ce qui est possible par le th\'eor\`eme
\ref{spaces-perfect-theorem-approximation}). Le morphisme induit
$$
\Ext^i_{\mathcal{O}_X}(E,
\mathcal{G}^\bullet \otimes_{\mathcal{O}_X} f^*\mathcal{F})
\longrightarrow
\Ext^i_{\mathcal{O}_X}(P,
\mathcal{G}^\bullet \otimes_{\mathcal{O}_X} f^*\mathcal{F})
$$
est alors un isomorphisme pour $i\leq m$. En effet, le noyau, resp. le
conoyau, de ce morphisme est un quotient, resp. un sous-module, de
$$
\Ext^i_{\mathcal{O}_X}(C,
\mathcal{G}^\bullet \otimes_{\mathcal{O}_X} f^*\mathcal{F})
\quad\text{resp.}\quad
\Ext^{i + 1}_{\mathcal{O}_X}(C,
\mathcal{G}^\bullet \otimes_{\mathcal{O}_X} f^*\mathcal{F}),
$$
o\`u $C$ est le c\^one de $P\to E$. Comme les faisceaux de cohomologie de $C$
sont nuls en degr\'es $\geq-m-1+a$, ces groupes d'Ext sont nuls pour
$i\leq m+1$ d'apr\`es Cat\'egories d\'eriv\'ees, lemme
\ref{derived-lemma-negative-exts}. On est ainsi ramen\'e au cas o\`u $E$ est un
complexe parfait, qui est le lemme \ref{spaces-perfect-lemma-compute-ext-perfect}.
La compatibilit\'e avec les homomorphismes cobord est expliqu\'ee dans la
d\'emonstration du lemme \ref{spaces-perfect-lemma-compute-ext-perfect}.

\medskip\noindent
D\'emonstration de (B). Comme dans la d\'emonstration de (A), on peut supposer
$X$ noeth\'erien. Le lemme \ref{spaces-perfect-lemma-identify-pseudo-coherent-noetherian}
montre que $E$ est pseudo-coh\'erent. D'apr\`es le lemme
\ref{spaces-perfect-lemma-pseudo-coherent-hocolim}, on peut \`ecrire
$E=\text{hocolim}\,E_n$, o\`u les $E_n$ sont parfaits et o\`u $E_n\to E$ induit
un isomorphisme sur les troncations $\tau_{\geq-n}$. Soit $E_n^\vee$ le
complexe parfait dual (Cohomologie sur les sites, lemme
\ref{sites-cohomology-lemma-dual-perfect-complex}). On obtient un syst\`eme
projectif $\ldots\to E_3^\vee\to E_2^\vee\to E_1^\vee$ d'objets parfaits.
Celui-ci donne \`a son tour un syst\`eme projectif
$$
\ldots \to K_3 \to K_2 \to K_1\quad\text{avec}\quad
K_n = Rf_*(E_n^\vee \otimes_{\mathcal{O}_X}^\mathbf{L} \mathcal{G}^\bullet)
$$
d'objets parfaits sur $B$ ; voir le lemme \ref{spaces-perfect-lemma-tensor-perfect}.
D'apr\`es le lemme \ref{spaces-perfect-lemma-compute-ext-perfect} et sa d\'emonstration, ainsi
que les arguments du paragraphe pr\'ec\'edent (avec $P=E_n$), pour tout
$\mathcal{F}$ quasi-coh\'erent sur $B$, on dispose d'homomorphismes canoniques
fonctoriels
$$
\xymatrix{
& \Ext^i_{\mathcal{O}_X}(E,
\mathcal{G}^\bullet \otimes_{\mathcal{O}_X} f^*\mathcal{F})
\ar[ld] \ar[rd] \\
H^i(B, K_{n + 1} \otimes_{\mathcal{O}_B}^\mathbf{L} \mathcal{F})
\ar[rr] & &
H^i(B, K_n \otimes_{\mathcal{O}_B}^\mathbf{L} \mathcal{F})
}
$$
qui sont des isomorphismes pour $i\leq n+a$. Soit $L_n=K_n^\vee$ le complexe
parfait dual. On voit alors que
$L_1\to L_2\to L_3\to\ldots$ est un syst\`eme inductif d'objets parfaits de
$D(\mathcal{O}_B)$ tel que, pour tout $\mathcal{F}$ quasi-coh\'erent sur $B$,
les morphismes
$$
\Ext^i_{\mathcal{O}_B}(L_{n + 1}, \mathcal{F})
\longrightarrow
\Ext^i_{\mathcal{O}_B}(L_n, \mathcal{F})
$$
sont des isomorphismes pour $i\leq n+a-1$. Il en r\'esulte que
$L_n\to L_{n+1}$ induit un isomorphisme sur les troncations
$\tau_{\geq-n-a+2}$ (indication : prendre le c\^one de $L_n\to L_{n+1}$ et
consid\'erer son dernier faisceau de cohomologie non nul). Ainsi,
$L=\text{hocolim}\,L_n$ est pseudo-coh\'erent ; voir le lemme
\ref{spaces-perfect-lemma-pseudo-coherent-hocolim}. La propri\'et\'e universelle des colimites
homotopiques donne
$\Ext^i_{\mathcal{O}_B}(L,\mathcal{F})=
\Ext^i_{\mathcal{O}_B}(L_n,\mathcal{F})$ pour $i\leq n+a-3$, ce qui ach\`eve
la d\'emonstration.
\end{proof}

\begin{remark}
\label{spaces-perfect-remark-base-change-of-L}
Le complexe pseudo-coh\'erent $L$ de la partie (B) du lemme
\ref{spaces-perfect-lemma-compute-ext} est canoniquement associ\'e \`a la situation. Par
exemple, la construction de $L$ dans (B) est compatible avec le changement de
base. Autrement dit, pour tout diagramme cart\'esien de sch\'emas
$$
\xymatrix{
X' \ar[r]_{g'} \ar[d]_{f'} &
X \ar[d]^f \\
Y' \ar[r]^g &
Y
}
$$
on a des isomorphismes canoniques fonctoriels
$$
\Ext^i_{\mathcal{O}_{Y'}}(Lg^*L, \mathcal{F}') \longrightarrow
\Ext^i_{\mathcal{O}_{X'}}(L(g')^*E,
(g')^*\mathcal{G}^\bullet \otimes_{\mathcal{O}_{X'}} (f')^*\mathcal{F}')
$$
pour $\mathcal{F}'$ quasi-coh\'erent sur $Y'$. Notons que nous n'utilisons
{\bf pas} l'image inverse d\'eriv\'ee de $\mathcal{G}^\bullet$ au membre de
droite. Si ce r\'esultat nous est un jour n\'ecessaire, nous en formulerons ici
un \'enonc\'e pr\'ecis et en donnerons une d\'emonstration d\'etaill\'ee.
\end{remark}

\section{Limites et cat\'egories d\'eriv\'ees}
\label{spaces-perfect-section-limits}

\noindent
Nous rassemblons dans cette section quelques r\'esultats sur la cat\'egorie
d\'eriv\'ee d'un espace alg\'ebrique qui est la limite projective d'un syst\`eme
projectif d'espaces alg\'ebriques. Plus pr\'ecis\'ement, nous nous placerons dans
la situation suivante.

\begin{situation}
\label{spaces-perfect-situation-descent}
Soit $S$ un sch\'ema. Soit $X=\lim_{i\in I}X_i$ la limite projective d'un syst\`eme
projectif d'espaces alg\'ebriques sur $S$, dont les morphismes de transition
affines sont $f_{i'i}:X_{i'}\to X_i$. Notons $f_i:X\to X_i$ la projection.
Nous supposons $X_i$ quasi-compact et quasi-s\'epar\'e pour tout $i\in I$.
Nous choisissons en outre un \`el\'ement $0\in I$.
\end{situation}

\begin{lemma}
\label{spaces-perfect-lemma-descend-homomorphisms}
Pla\c{c}ons-nous dans la situation \ref{spaces-perfect-situation-descent}. Soient $E_0$ et
$K_0$ des objets de $D(\mathcal{O}_{X_0})$. Posons
$E_i=Lf_{i0}^*E_0$ et $K_i=Lf_{i0}^*K_0$ pour $i\geq0$, ainsi que
$E=Lf_0^*E_0$ et $K=Lf_0^*K_0$. Alors le morphisme
$$
\colim_{i \geq 0} \Hom_{D(\mathcal{O}_{X_i})}(E_i, K_i)
\longrightarrow
\Hom_{D(\mathcal{O}_X)}(E, K)
$$
est un isomorphisme dans chacun des cas suivants :
\begin{enumerate}
\item $E_0$ est parfait et $K_0\in D_\QCoh(\mathcal{O}_{X_0})$ ;
\item $E_0$ est pseudo-coh\'erent et
$K_0\in D_\QCoh(\mathcal{O}_{X_0})$ est de Tor-dimension finie.
\end{enumerate}
\end{lemma}

\begin{proof}
Pour tout objet quasi-compact et quasi-s\'epar\'e $U_0$ de
$(X_0)_{spaces,\etale}$, consid\'erons la propri\'et\'e $P$ selon laquelle le
morphisme canonique
$$
\colim_{i \geq 0} \Hom_{D(\mathcal{O}_{U_i})}(E_i|_{U_i}, K_i|_{U_i})
\longrightarrow
\Hom_{D(\mathcal{O}_U)}(E|_U, K|_U)
$$
est un isomorphisme, o\`u $U=X\times_{X_0}U_0$ et
$U_i=X_i\times_{X_0}U_0$. Nous allons montrer, par le principe de r\'ecurrence
du lemme \ref{spaces-perfect-lemma-induction-principle}, que chaque $U_0$ v\'erifie $P$.
La condition (2) de ce lemme r\'esulte imm\'ediatement de Mayer--Vietoris pour
les Hom de la cat\'egorie d\'eriv\'ee ; voir le lemme
\ref{spaces-perfect-lemma-mayer-vietoris-hom}. Il suffit donc de d\'emontrer le lemme lorsque
$X_0$ est affine.

\medskip\noindent
Si $X_0$ est affine, le r\'esultat d\'ecoule du cas des sch\'emas ; voir
Cat\'egories d\'eriv\'ees des sch\'emas, lemme
\ref{perfect-lemma-descend-homomorphisms}. Pour le constater, utilisons
l'\'equivalence du lemme \ref{spaces-perfect-lemma-derived-quasi-coherent-small-etale-site}
et la correspondance des propri\'et\'es expliqu\'ee dans les lemmes
\ref{spaces-perfect-lemma-descend-pseudo-coherent},
\ref{spaces-perfect-lemma-descend-tor-amplitude} et
\ref{spaces-perfect-lemma-descend-perfect}.
\end{proof}

\begin{lemma}
\label{spaces-perfect-lemma-perfect-on-limit}
Dans la situation \ref{spaces-perfect-situation-descent}, la cat\'egorie des objets parfaits
de $D(\mathcal{O}_X)$ est la limite inductive des cat\'egories d'objets
parfaits de $D(\mathcal{O}_{X_i})$.
\end{lemma}

\begin{proof}
Pour tout objet quasi-compact et quasi-s\'epar\'e $U_0$ de
$(X_0)_{spaces,\etale}$, consid\'erons la propri\'et\'e $P$ selon laquelle le
foncteur
$$
\colim_{i \geq 0} D_{perf}(\mathcal{O}_{U_i})
\longrightarrow
D_{perf}(\mathcal{O}_U)
$$
est une \`equivalence, o\`u l'indice ${}_{perf}$ d\'esigne la sous-cat\'egorie
pleine des objets parfaits, et o\`u $U=X\times_{X_0}U_0$ et
$U_i=X_i\times_{X_0}U_0$. Nous allons montrer, par le principe de r\'ecurrence
du lemme \ref{spaces-perfect-lemma-induction-principle}, que chaque $U_0$ v\'erifie $P$.
D'apr\`es le lemme \ref{spaces-perfect-lemma-descend-homomorphisms}, nous savons d\'ej\`a que
le foncteur est pleinement fid\`ele. Il suffit donc d'\'etablir sa surjectivit\'e
essentielle.

\medskip\noindent
V\'erifions d'abord la condition (2) du principe de r\'ecurrence. Supposons donc
donn\'e un carr\'e distingu\'e \`el\'ementaire
$(U_0\subset X_0,V_0\to X_0)$, et supposons que $U_0$, $V_0$ et
$U_0\times_{X_0}V_0$ v\'erifient $P$. Soit $E$ un objet parfait de
$D(\mathcal{O}_X)$. On peut trouver un indice $i\geq0$ et des objets parfaits
$E_{U,i}$ sur $U_i$ et $E_{V,i}$ sur $V_i$, dont les images inverses sur $U$
et $V$ sont respectivement isomorphes \`a $E|_U$ et $E|_V$. Notons
$$
a : E_{U, i} \to (R(X \to X_i)_*E)|_{U_i}
\quad\text{et}\quad
b : E_{V, i} \to (R(X \to X_i)_*E)|_{V_i}
$$
les morphismes adjoints aux isomorphismes
$L(U\to U_i)^*E_{U,i}\to E|_U$ et
$L(V\to V_i)^*E_{V,i}\to E|_V$. Par pleine fid\'elit\'e, quitte \`a augmenter
$i$, on peut trouver un isomorphisme
$c:E_{U,i}|_{U_i\times_{X_i}V_i}\to
E_{V,i}|_{U_i\times_{X_i}V_i}$ dont l'image inverse donne les identifications
$$
L(U \to U_i)^*E_{U, i}|_{U \times_X V} \to E|_{U \times_X V} \to
L(V \to V_i)^*E_{V, i}|_{U \times_X V}.
$$
Appliquons le lemme \ref{spaces-perfect-lemma-glue}. On obtient un objet $E_i$ sur $X_i$ et
un morphisme $d:E_i\to R(X\to X_i)_*E$ dont les restrictions sur $U_i$ et
$V_i$ sont les morphismes $a$ et $b$. Il est alors clair que $E_i$ est parfait
et que $d$ est adjoint \`a un isomorphisme $L(X\to X_i)^*E_i\to E$.

\medskip\noindent
V\'erifions enfin la condition (1) du principe de r\'ecurrence, autrement dit
le cas o\`u $X_0$ est affine. Celui-ci r\'esulte du cas des sch\'emas ; voir
Cat\'egories d\'eriv\'ees des sch\'emas, lemme
\ref{perfect-lemma-descend-perfect}. Pour le constater, utilisons
l'\'equivalence du lemme \ref{spaces-perfect-lemma-derived-quasi-coherent-small-etale-site}
et la correspondance du lemme \ref{spaces-perfect-lemma-descend-perfect}.
\end{proof}

\section{Cohomologie et changement de base, VI}
\label{spaces-perfect-section-cohomology-and-base-change-final}

\noindent
Cette dernière section sur la cohomologie et le changement de base poursuit
la discussion des sections
\ref{spaces-perfect-section-cohomology-and-base-change-perfect},
\ref{spaces-perfect-section-cohomology-base-change} et
\ref{spaces-perfect-section-producing-perfect}.
Un cas particulier facile à saisir est donné dans
la remarque \ref{spaces-perfect-remark-explain-perfect-direct-image}.

\begin{lemma}
\label{spaces-perfect-lemma-base-change-tensor-perfect}
Soit $S$ un schéma. Soit $f : X \to Y$ un morphisme de présentation finie
entre des espaces algébriques sur $S$. Soit $E \in D(\mathcal{O}_X)$ un objet
parfait. Soit $\mathcal{G}^\bullet$ un complexe borné de
$\mathcal{O}_X$-modules de présentation finie dont les termes sont plats sur
$Y$ et dont le support est propre sur $Y$. Alors
$$
K = Rf_*(E \otimes_{\mathcal{O}_X}^\mathbf{L} \mathcal{G}^\bullet)
$$
est un objet parfait de $D(\mathcal{O}_Y)$ et sa formation
commute à tout changement de base.
\end{lemma}

\begin{proof}
L'assertion concernant le changement de base est le
lemme \ref{spaces-perfect-lemma-base-change-tensor}.
Il suffit donc de montrer que $K$ est un objet parfait. Si $Y$ est
noethérien, cela résulte du
lemme \ref{spaces-perfect-lemma-tensor-perfect}.
Nous nous ramènerons à ce cas par approximation noethérienne.
Nous conseillons au lecteur de ne pas lire la suite de cette démonstration.

\medskip\noindent
La question étant locale sur $Y$, nous pouvons supposer $Y$ affine.
Écrivons $Y = \Spec(R)$. Écrivons $R = \colim R_i$ comme limite inductive
filtrante d'anneaux noethériens $R_i$. D'après Limites d'espaces, lemme
\ref{spaces-limits-lemma-descend-finite-presentation},
il existe un $i$ et un espace algébrique $X_i$ de présentation finie
sur $R_i$ dont le changement de base à $R$ est $X$. D'après
Limites d'espaces, lemme
\ref{spaces-limits-lemma-descend-modules-finite-presentation},
quitte à remplacer $i$ par un indice plus grand, nous pouvons supposer
qu'il existe un complexe borné de
$\mathcal{O}_{X_i}$-modules de présentation finie $\mathcal{G}_i^\bullet$ dont
l'image inverse sur $X$ est $\mathcal{G}^\bullet$. Quitte à remplacer $i$
par un indice plus grand,
nous pouvons supposer que $\mathcal{G}_i^n$ est plat sur $R_i$ ; voir
Limites d'espaces, lemme
\ref{spaces-limits-lemma-descend-flat}.
Quitte à remplacer $i$ par un indice plus grand, nous pouvons supposer que
le support de $\mathcal{G}_i^n$ est propre sur $R_i$ ; voir
Limites d'espaces, lemme \ref{spaces-limits-lemma-eventually-proper-support}.
Enfin, d'après le lemme \ref{spaces-perfect-lemma-perfect-on-limit},
quitte à remplacer $i$ par un indice plus grand, nous pouvons supposer
qu'il existe un objet parfait $E_i$ de $D(\mathcal{O}_{X_i})$ dont l'image
inverse sur $X$ est $E$. D'après le lemme \ref{spaces-perfect-lemma-tensor-perfect},
$K_i =
Rf_{i, *}(E_i \otimes_{\mathcal{O}_{X_i}}^\mathbf{L} \mathcal{G}_i^\bullet)$
est parfait sur $\Spec(R_i)$, où $f_i : X_i \to \Spec(R_i)$ est
le morphisme structural. D'après le résultat de changement de base
(lemme \ref{spaces-perfect-lemma-base-change-tensor}), l'image inverse de $K_i$
sur $Y = \Spec(R)$ est $K$, ce qui permet de conclure.
\end{proof}

\begin{remark}
\label{spaces-perfect-remark-explain-perfect-direct-image}
Soit $R$ un anneau. Soit $X$ un espace algébrique de présentation finie sur
$R$. Soit $\mathcal{G}$ un $\mathcal{O}_X$-module de présentation finie,
plat sur $R$ et dont le support est propre sur $R$. D'après le
lemme \ref{spaces-perfect-lemma-base-change-tensor-perfect},
il existe un complexe fini de $R$-modules projectifs de type fini
$M^\bullet$ tel que l'on ait
$$
R\Gamma(X_{R'}, \mathcal{G}_{R'}) = M^\bullet \otimes_R R'
$$
fonctoriellement en la $R$-algèbre $R'$.
\end{remark}

\begin{lemma}
\label{spaces-perfect-lemma-base-change-tensor-pseudo-coherent}
Soit $S$ un schéma.
Soit $f : X \to Y$ un morphisme de présentation finie
entre des espaces algébriques sur $S$.
Soit $E \in D(\mathcal{O}_X)$ un objet pseudo-cohérent.
Soit $\mathcal{G}^\bullet$ un complexe borné supérieurement de
$\mathcal{O}_X$-modules de présentation finie dont les termes sont
plats sur $Y$ et dont le support est propre sur $Y$. Alors
$$
K = Rf_*(E \otimes_{\mathcal{O}_X}^\mathbf{L} \mathcal{G}^\bullet)
$$
est un objet pseudo-cohérent de $D(\mathcal{O}_Y)$ et sa formation
commute à tout changement de base.
\end{lemma}

\begin{proof}
L'assertion concernant le changement de base est le
lemme \ref{spaces-perfect-lemma-base-change-tensor}.
Il suffit donc de montrer que $K$ est un objet pseudo-cohérent.
Cela résultera du lemme \ref{spaces-perfect-lemma-base-change-tensor-perfect}
grâce à une approximation par des complexes parfaits. Nous conseillons au
lecteur de ne pas lire la suite de cette démonstration.

\medskip\noindent
La question est locale pour la topologie étale sur $Y$ ; nous pouvons donc
supposer $Y$ affine. Alors $X$ est quasi-compact et quasi-séparé. De plus, il
existe un entier $N$ tel que l'image directe totale
$Rf_* : D_\QCoh(\mathcal{O}_X) \to D_\QCoh(\mathcal{O}_Y)$
soit de dimension cohomologique $N$, comme l'explique le
lemme \ref{spaces-perfect-lemma-quasi-coherence-direct-image}.
Choisissons un entier $b$ tel que $\mathcal{G}^i = 0$ pour $i > b$.
Il suffit de montrer que $K$ est $m$-pseudo-cohérent pour
tout $m$. Choisissons une approximation $P \to E$ par un complexe parfait $P$
de $(X, E, m - N - 1 - b)$. Cela est possible d'après le
théorème \ref{spaces-perfect-theorem-approximation}.
Choisissons un triangle distingué
$$
P \to E \to C \to P[1]
$$
dans $D_\QCoh(\mathcal{O}_X)$. Les faisceaux de cohomologie de $C$ sont nuls
en degrés $\geq m - N - 1 - b$. Par conséquent, les faisceaux de cohomologie
de $C \otimes^\mathbf{L} \mathcal{G}^\bullet$ sont nuls en degrés
$\geq m - N - 1$.
Ainsi, les faisceaux de cohomologie de
$Rf_*(C \otimes^\mathbf{L} \mathcal{G}^\bullet)$ sont nuls en degrés $\geq m - 1$.
Par conséquent,
$$
Rf_*(P \otimes^\mathbf{L} \mathcal{G}^\bullet) \to
Rf_*(E \otimes^\mathbf{L} \mathcal{G}^\bullet)
$$
est un isomorphisme sur les faisceaux de cohomologie en degrés $\geq m$.
Supposons ensuite que $H^i(P) = 0$ pour $i > a$. Alors
$
P \otimes^\mathbf{L} \sigma_{\geq m - N - 1 - a}\mathcal{G}^\bullet
\longrightarrow
P \otimes^\mathbf{L} \mathcal{G}^\bullet
$
est un isomorphisme sur les faisceaux de cohomologie en degrés
$\geq m - N - 1$.
Nous trouvons donc de nouveau que
$$
Rf_*(P \otimes^\mathbf{L} \sigma_{\geq m - N - 1 - a}\mathcal{G}^\bullet) \to
Rf_*(P \otimes^\mathbf{L} \mathcal{G}^\bullet)
$$
est un isomorphisme sur les faisceaux de cohomologie en degrés $\geq m$.
D'après le lemme \ref{spaces-perfect-lemma-base-change-tensor-perfect}, la source
est un complexe parfait.
Nous en concluons que $K$ est $m$-pseudo-cohérent, comme voulu.
\end{proof}

\begin{lemma}
\label{spaces-perfect-lemma-flat-proper-perfect-direct-image-general}
Soit $S$ un schéma. Soit $f : X \to Y$ un morphisme propre
et de présentation finie entre des espaces algébriques sur $S$.
\begin{enumerate}
\item Soit $E \in D(\mathcal{O}_X)$ un objet parfait et supposons $f$ plat.
Alors
$Rf_*E$ est un objet parfait de $D(\mathcal{O}_Y)$ et sa formation
commute à tout changement de base.
\item Soit $\mathcal{G}$ un $\mathcal{O}_X$-module de présentation finie,
plat sur $Y$. Alors $Rf_*\mathcal{G}$ est un objet parfait de
$D(\mathcal{O}_Y)$ et sa formation commute à tout changement de base.
\end{enumerate}
\end{lemma}

\begin{proof}
Ce sont des cas particuliers du
lemme \ref{spaces-perfect-lemma-base-change-tensor-perfect}, obtenus en prenant
(1) $\mathcal{G}^\bullet$ égal au complexe $\mathcal{O}_X$ placé en degré $0$
et (2) $E = \mathcal{O}_X$ et $\mathcal{G}^\bullet$ égal au complexe constitué
de $\mathcal{G}$ placé en degré $0$.
\end{proof}

\begin{lemma}
\label{spaces-perfect-lemma-flat-proper-pseudo-coherent-direct-image-general}
Soit $S$ un sch\'ema. Soit $f : X \to Y$ un morphisme propre, plat et
de pr\'esentation finie entre des espaces alg\'ebriques sur $S$.
Soit $E \in D(\mathcal{O}_X)$ un objet
pseudo-coh\'erent. Alors $Rf_*E$ est un objet pseudo-coh\'erent de
$D(\mathcal{O}_Y)$ et sa formation commute \`a tout changement de base.
\end{lemma}

\noindent
Plus g\'en\'eralement, si $f : X \to Y$ est propre et si $E$ sur $X$ est
pseudo-coh\'erent relativement \`a $Y$ (Compl\'ements sur les morphismes d'espaces, d\'efinition
\ref{spaces-more-morphisms-definition-relative-pseudo-coherence}),
alors $Rf_*E$ est pseudo-coh\'erent
(mais sa formation ne commute pas au changement de base \`a ce degr\'e de g\'en\'eralit\'e).
Pour les sch\'emas, cet \`enonc\'e est d\'emontr\'e dans \cite{Kiehl}.

\begin{proof}
C'est le cas particulier du
lemme \ref{spaces-perfect-lemma-base-change-tensor-pseudo-coherent} obtenu en prenant
$\mathcal{G} = \mathcal{O}_X$.
\end{proof}

\begin{lemma}
\label{spaces-perfect-lemma-pullback-and-limits}
Soit $R$ un anneau. Soit $X$ un espace alg\'ebrique et soit
$f : X \to \Spec(R)$ un morphisme propre, plat et
de pr\'esentation finie. Soit $(M_n)$ un syst\`eme projectif
de $R$-modules dont les morphismes de transition sont surjectifs.
Alors le morphisme canonique
$$
\mathcal{O}_X \otimes_R (\lim M_n)
\longrightarrow
\lim \mathcal{O}_X \otimes_R M_n
$$
induit un isomorphisme de la source sur l'objet obtenu en appliquant $DQ_X$
\`a la cible.
\end{lemma}

\begin{proof}
L'\'enonc\'e signifie que, pour tout objet $E$ de
$D_\QCoh(\mathcal{O}_X)$, le morphisme induit
$$
\Hom(E, \mathcal{O}_X \otimes_R (\lim M_n))
\longrightarrow
\Hom(E, \lim \mathcal{O}_X \otimes_R M_n)
$$
est un isomorphisme. Comme $D_\QCoh(\mathcal{O}_X)$ admet
un g\'en\'erateur parfait (th\'eor\`eme \ref{spaces-perfect-theorem-bondal-van-den-Bergh}),
il suffit de le v\'erifier pour $E$ parfait.
D'apr\`es le lemme \ref{spaces-perfect-lemma-Rlim-quasi-coherent}, on a
$\lim \mathcal{O}_X \otimes_R M_n = R\lim \mathcal{O}_X \otimes_R M_n$.
Le foncteur exact
$R\Hom_X(E, -) : D_\QCoh(\mathcal{O}_X) \to D(R)$,
consid\'er\'e dans Cohomologie sur les sites, section
\ref{sites-cohomology-section-global-RHom}, commute aux produits, donc aux
limites projectives d\'eriv\'ees ; par cons\'equent,
$$
R\Hom_X(E, \lim \mathcal{O}_X \otimes_R M_n) =
R\lim R\Hom_X(E, \mathcal{O}_X \otimes_R M_n)
$$
Soit $E^\vee$ le complexe parfait dual ; voir
Cohomologie sur les sites, lemme \ref{sites-cohomology-lemma-dual-perfect-complex}.
On a
$$
R\Hom_X(E, \mathcal{O}_X \otimes_R M_n) =
R\Gamma(X, E^\vee \otimes_{\mathcal{O}_X}^\mathbf{L} Lf^*M_n) =
R\Gamma(X, E^\vee) \otimes_R^\mathbf{L} M_n
$$
d'apr\`es le lemme \ref{spaces-perfect-lemma-cohomology-base-change}.
Le lemme \ref{spaces-perfect-lemma-flat-proper-perfect-direct-image-general} montre que
$R\Gamma(X, E^\vee)$ est un complexe parfait de $R$-modules.
En particulier, c'est un complexe pseudo-coh\'erent ; le
lemme \ref{more-algebra-lemma-pseudo-coherent-tensor-limit} de Compl\'ements
d'alg\`ebre donne donc
$$
R\lim R\Gamma(X, E^\vee) \otimes_R^\mathbf{L} M_n =
R\Gamma(X, E^\vee) \otimes_R^\mathbf{L} \lim M_n
$$
ce qu'il fallait d\'emontrer.
\end{proof}

\begin{lemma}
\label{spaces-perfect-lemma-perfect-enough}
Soit $A$ un anneau. Soit $X$ un espace alg\'ebrique sur $A$, quasi-compact
et quasi-s\'epar\'e. Soit $K \in D^-_\QCoh(\mathcal{O}_X)$.
Si $R\Gamma(X, E \otimes^\mathbf{L} K)$ est pseudo-coh\'erent
dans $D(A)$ pour tout $E$ parfait de $D(\mathcal{O}_X)$,
alors $R\Gamma(X, E \otimes^\mathbf{L} K)$ est pseudo-coh\'erent
dans $D(A)$ pour tout $E$ pseudo-coh\'erent de $D(\mathcal{O}_X)$.
\end{lemma}

\begin{proof}
Il existe un entier $N$ tel que
$R\Gamma(X, -) : D_\QCoh(\mathcal{O}_X) \to D(A)$
soit de dimension cohomologique $N$, comme expliqu\'e dans le
lemme \ref{spaces-perfect-lemma-quasi-coherence-direct-image}.
Soit $b \in \mathbf{Z}$ tel que $H^i(K) = 0$ pour $i > b$.
Soit $E$ pseudo-coh\'erent sur $X$.
Il suffit de montrer que $R\Gamma(X, E \otimes^\mathbf{L} K)$
est $m$-pseudo-coh\'erent pour tout $m$.
Choisissons une approximation $P \to E$ par un complexe parfait $P$
de $(X, E, m - N - 1 - b)$. C'est possible d'apr\`es le
th\'eor\`eme \ref{spaces-perfect-theorem-approximation}.
Choisissons un triangle distingu\'e
$$
P \to E \to C \to P[1]
$$
dans $D_\QCoh(\mathcal{O}_X)$. Les faisceaux de cohomologie de $C$ sont nuls
en degr\'es $\geq m - N - 1 - b$. Par suite, les faisceaux de cohomologie
de $C \otimes^\mathbf{L} K$ sont nuls en degr\'es $\geq m - N - 1$.
Les groupes de cohomologie de $R\Gamma(X, C \otimes^\mathbf{L} K)$
sont donc nuls en degr\'es $\geq m - 1$. Ainsi,
$$
R\Gamma(X, P \otimes^\mathbf{L} K) \to R\Gamma(X, E \otimes^\mathbf{L} K)
$$
induit un isomorphisme sur la cohomologie en degr\'es $\geq m$.
Par hypoth\`ese, la source est pseudo-coh\'erente.
On en d\'eduit que $R\Gamma(X, E \otimes^\mathbf{L} K)$
est $m$-pseudo-coh\'erent, comme voulu.
\end{proof}

\begin{lemma}
\label{spaces-perfect-lemma-base-change-RHom-perfect}
Soit $S$ un sch\'ema.
Soit $f : X \to Y$ un morphisme de pr\'esentation finie
entre des espaces alg\'ebriques sur $S$.
Soit $E \in D(\mathcal{O}_X)$ un objet parfait. Soit $\mathcal{G}^\bullet$
un complexe born\'e de $\mathcal{O}_X$-modules de pr\'esentation finie,
dont les termes sont plats sur $Y$ et dont le support est propre sur $Y$. Alors
$$
K = Rf_*R\SheafHom(E, \mathcal{G}^\bullet)
$$
est un objet parfait de $D(\mathcal{O}_Y)$ et sa formation
commute \`a tout changement de base.
\end{lemma}

\begin{proof}
L'assertion sur le changement de base est le lemme \ref{spaces-perfect-lemma-base-change-RHom}.
Il suffit donc de montrer que $K$ est un objet parfait. Si $Y$ est
noeth\'erien, cela r\'esulte du lemme \ref{spaces-perfect-lemma-ext-perfect}.
Nous allons nous ramener \`a ce cas par approximation noeth\'erienne.
Nous conseillons au lecteur de ne pas lire la suite de cette d\'emonstration.

\medskip\noindent
La question \'etant locale sur $Y$, on peut donc supposer $Y$ affine.
Écrivons $Y = \Spec(R)$. Écrivons $R = \colim R_i$ comme limite inductive
filtrante d'anneaux noeth\'eriens $R_i$. D'apr\`es Limites d'espaces, lemme
\ref{spaces-limits-lemma-descend-finite-presentation},
il existe un $i$ et un espace alg\'ebrique $X_i$ de pr\'esentation finie
sur $R_i$ dont le changement de base \`a $R$ est $X$. D'apr\`es
Limites d'espaces, lemme
\ref{spaces-limits-lemma-descend-modules-finite-presentation},
on peut supposer, quitte \`a remplacer $i$ par un indice plus grand,
qu'il existe un complexe born\'e de $\mathcal{O}_{X_i}$-modules de pr\'esentation
finie $\mathcal{G}_i^\bullet$ dont
l'image inverse sur $X$ est $\mathcal{G}^\bullet$. Quitte \`a remplacer $i$
par un indice plus grand,
on peut supposer que $\mathcal{G}_i^n$ est plat sur $R_i$ ; voir
Limites d'espaces, lemme
\ref{spaces-limits-lemma-descend-flat}.
Quitte \`a remplacer $i$ par un indice plus grand, on peut supposer que
le support de $\mathcal{G}_i^n$ est propre sur $R_i$ ; voir
Limites d'espaces, lemme \ref{spaces-limits-lemma-eventually-proper-support}.
Enfin, d'apr\`es le lemme \ref{spaces-perfect-lemma-descend-perfect},
on peut, quitte \`a remplacer $i$ par un indice plus grand, supposer qu'il
existe un objet parfait $E_i$ de $D(\mathcal{O}_{X_i})$ dont l'image inverse
sur $X$ est $E$. En appliquant le lemme \ref{spaces-perfect-lemma-compute-ext-perfect}
\`a $X_i \to \Spec(R_i)$, $E_i$, $\mathcal{G}_i^\bullet$ et en utilisant la
propri\'et\'e de changement de base d\'ej\`a d\'emontr\'ee, on obtient le r\'esultat.
\end{proof}







\section{Complexes parfaits}
\label{spaces-perfect-section-perfect-complexes}

\noindent
Nous commençons par étudier les lieux de saut des nombres de Betti des complexes
parfaits. Définissons d'abord les nombres de Betti.

\medskip\noindent
Soit $S$ un sch\'ema. Soit $X$ un espace alg\'ebrique sur $S$.
Soit $E$ un objet de $D(\mathcal{O}_X)$.
Soit $x\in|X|$. Nous voulons d\'efinir
$\beta_i(x)\in\{0,1,2,\ldots\}\cup\{\infty\}$.
Pour cela, choisissons un morphisme $f:\Spec(k)\to X$ représentant $x$.
Alors $Lf^*E$ est un objet de
$D(\Spec(k)_\etale,\mathcal{O})$.
D'apr\`es Cohomologie \'etale, lemme
\ref{etale-cohomology-lemma-all-modules-quasi-coherent}, et le th\'eor\`eme
\ref{etale-cohomology-theorem-quasi-coherent}, on a
$D(\Spec(k)_\etale,\mathcal{O})=D(k)$, la cat\'egorie d\'eriv\'ee des
$k$-espaces vectoriels. Ainsi, $Lf^*E$ est un complexe de $k$-espaces
vectoriels et l'on peut poser
$\beta_i(x)=\dim_k H^i(Lf^*E)$. On voit facilement que cette valeur ne d\'epend
pas du choix du repr\'esentant de $x$. De plus, si $X$ est un sch\'ema, cette
notion co\"incide avec celle employ\'ee dans Cat\'egories d\'eriv\'ees des
sch\'emas, section \ref{perfect-section-perfect-complexes}.

\begin{lemma}
\label{spaces-perfect-lemma-jump-loci}
Soit $S$ un sch\'ema. Soit $X$ un espace alg\'ebrique sur $S$.
Soit $E\in D(\mathcal{O}_X)$ pseudo-coh\'erent (par exemple parfait).
Pour tout $i\in\mathbf{Z}$, consid\'erons la fonction
$$
\beta_i : |X| \longrightarrow \{0, 1, 2, \ldots\}
$$
d\'efinie ci-dessus. Alors
\begin{enumerate}
\item la formation de $\beta_i$ commute \`a tout changement de base ;
\item les fonctions $\beta_i$ sont semi-continues sup\'erieurement ;
\item les ensembles de niveau de $\beta_i$ sont localement constructibles
pour la topologie \'etale.
\end{enumerate}
\end{lemma}

\begin{proof}
Choisissons un sch\'ema $U$ et un morphisme \'etale surjectif
$\varphi:U\to X$. Alors $L\varphi^*E$ est un complexe pseudo-coh\'erent sur le
sch\'ema $U$ (utiliser le lemme \ref{spaces-perfect-lemma-descend-pseudo-coherent}), et l'on
peut appliquer le r\'esultat pour les sch\'emas ; voir Cat\'egories d\'eriv\'ees
des sch\'emas, lemme \ref{perfect-lemma-jump-loci}. L'assertion (3) signifie
que les images inverses des ensembles de niveau sur $U$ sont localement
constructibles ; voir Propri\'et\'es des espaces, d\'efinition
\ref{spaces-properties-definition-locally-constructible}.
\end{proof}

\begin{lemma}
\label{spaces-perfect-lemma-jump-loci-geometric}
Soit $Y$ un sch\'ema et soit $X$ un espace alg\'ebrique sur $Y$ tel que le
morphisme structural $f:X\to Y$ soit plat, propre et de pr\'esentation finie.
Soit $\mathcal{F}$ un $\mathcal{O}_X$-module de pr\'esentation finie, plat sur
$Y$. Pour $i\in\mathbf{Z}$ fix\'e, consid\'erons la fonction
$$
\beta_i : |Y| \to \{0, 1, 2, \ldots\},\quad
y \longmapsto \dim_{\kappa(y)} H^i(X_y, \mathcal{F}_y).
$$
Alors
\begin{enumerate}
\item la formation de $\beta_i$ commute \`a tout changement de base ;
\item les fonctions $\beta_i$ sont semi-continues sup\'erieurement ;
\item les ensembles de niveau de $\beta_i$ sont localement constructibles
dans $Y$.
\end{enumerate}
\end{lemma}

\begin{proof}
D'apr\`es le th\'eor\`eme de cohomologie et changement de base (plus
pr\'ecis\'ement, le lemme
\ref{spaces-perfect-lemma-flat-proper-perfect-direct-image-general}), l'objet
$K=Rf_*\mathcal{F}$ est un objet parfait de la cat\'egorie d\'eriv\'ee de $Y$
dont la formation commute \`a tout changement de base. En particulier,
$$
H^i(X_y, \mathcal{F}_y) = H^i(K \otimes_{\mathcal{O}_Y}^\mathbf{L} \kappa(y)).
$$
Le lemme r\'esulte donc du lemme \ref{spaces-perfect-lemma-jump-loci}.
\end{proof}

\begin{lemma}
\label{spaces-perfect-lemma-chi-locally-constant}
Soit $S$ un sch\'ema. Soit $X$ un espace alg\'ebrique sur $S$.
Soit $E\in D(\mathcal{O}_X)$ parfait. La fonction
$$
\chi_E : |X| \longrightarrow \mathbf{Z},\quad
x \longmapsto \sum (-1)^i \beta_i(x)
$$
est localement constante sur $X$.
\end{lemma}

\begin{proof}
D\'emonstration omise. Indications :
on se ram\`ene au cas des sch\'emas par localisation \'etale. Voir Cat\'egories
d\'eriv\'ees des sch\'emas, lemme
\ref{perfect-lemma-chi-locally-constant}.
\end{proof}

\begin{lemma}
\label{spaces-perfect-lemma-open-where-cohomology-in-degree-i-rank-r}
Soit $S$ un sch\'ema. Soit $X$ un espace alg\'ebrique sur $S$.
Soit $E\in D(\mathcal{O}_X)$ parfait.
\'Etant donn\'es $i,r\in\mathbf{Z}$, il existe un sous-espace ouvert
$U\subset X$ caract\'eris\'e par les propri\'et\'es suivantes :
\begin{enumerate}
\item $E|_U\cong H^i(E|_U)[-i]$ et $H^i(E|_U)$ est un
$\mathcal{O}_U$-module localement libre de rang $r$ ;
\item un morphisme $f:Y\to X$ se factorise par $U$ si et seulement si
$Lf^*E$ est isomorphe \`a un module localement libre de rang $r$ plac\'e en
degr\'e $i$.
\end{enumerate}
\end{lemma}

\begin{proof}
D\'emonstration omise. Indications :
on se ram\`ene au cas des sch\'emas par localisation \'etale. Voir Cat\'egories
d\'eriv\'ees des sch\'emas, lemme
\ref{perfect-lemma-open-where-cohomology-in-degree-i-rank-r}.
\end{proof}

\begin{lemma}
\label{spaces-perfect-lemma-open-where-cohomology-in-degree-i-rank-r-geometric}
Soit $S$ un sch\'ema.
Soit $f:X\to Y$ un morphisme d'espaces alg\'ebriques sur $S$, propre, plat et
de pr\'esentation finie.
Soit $\mathcal{F}$ un $\mathcal{O}_X$-module de pr\'esentation finie, plat sur
$Y$. Fixons $i,r\in\mathbf{Z}$.
Il existe alors un sous-espace ouvert $V\subset Y$ poss\'edant la propri\'et\'e
suivante : un morphisme $T\to Y$ se factorise par $V$ si et seulement si
$Rf_{T,*}\mathcal{F}_T$ est isomorphe \`a un module localement libre de type
fini et de rang $r$, plac\'e en degr\'e $i$.
\end{lemma}

\begin{proof}
D'apr\`es le th\'eor\`eme de cohomologie et changement de base
(lemme \ref{spaces-perfect-lemma-flat-proper-perfect-direct-image-general}), l'objet
$K=Rf_*\mathcal{F}$ est un objet parfait de la cat\'egorie d\'eriv\'ee de $Y$
dont la formation commute \`a tout changement de base. Ce lemme r\'esulte donc
imm\'ediatement du lemme
\ref{spaces-perfect-lemma-open-where-cohomology-in-degree-i-rank-r}.
\end{proof}

\begin{lemma}
\label{spaces-perfect-lemma-locally-closed-where-H0-locally-free}
Soit $S$ un sch\'ema. Soit $X$ un espace alg\'ebrique sur $S$.
Soit $E\in D(\mathcal{O}_X)$ parfait, de Tor-amplitude contenue dans
$[a,b]$ pour certains $a,b\in\mathbf{Z}$.
Soit $r\geq0$.
Il existe alors un sous-espace localement ferm\'e $j:Z\to X$ caract\'eris\'e
par les propri\'et\'es suivantes :
\begin{enumerate}
\item $H^a(Lj^*E)$ est un $\mathcal{O}_Z$-module localement libre de rang $r$ ;
\item un morphisme $f:Y\to X$ se factorise par $Z$ si et seulement si, pour
tout morphisme $g:Y'\to Y$, le $\mathcal{O}_{Y'}$-module
$H^a(L(f\circ g)^*E)$ est localement libre de rang $r$.
\end{enumerate}
De plus, $j:Z\to X$ est de pr\'esentation finie et
\begin{enumerate}
\item[(3)] si $f:Y\to X$ se factorise sous la forme
$Y\xrightarrow{g}Z\to X$, alors
$H^a(Lf^*E)=g^*H^a(Lj^*E)$ ;
\item[(4)] si $\beta_a(x)\leq r$ pour tout $x\in|X|$, alors $j$ est une
immersion ferm\'ee et, pour tout $f:Y\to X$, les conditions suivantes sont
\'equivalentes :
\begin{enumerate}
\item $f:Y\to X$ se factorise par $Z$ ;
\item $H^a(Lf^*E)$ est un $\mathcal{O}_Y$-module localement libre de rang $r$ ;
\end{enumerate}
et, si $r=1$, elles sont encore \'equivalentes \`a
\begin{enumerate}
\item[(c)] $\mathcal{O}_Y\to
\SheafHom_{\mathcal{O}_Y}(H^a(Lf^*E),H^a(Lf^*E))$ est injectif.
\end{enumerate}
\end{enumerate}
\end{lemma}

\begin{proof}
D\'emonstration omise. Indications :
on se ram\`ene au cas des sch\'emas par localisation \'etale. Voir Cat\'egories
d\'eriv\'ees des sch\'emas, lemme
\ref{perfect-lemma-locally-closed-where-H0-locally-free}.
\end{proof}

\begin{lemma}
\label{spaces-perfect-lemma-proper-flat-h0}
Soit $S$ un sch\'ema.
Soit $f:X\to Y$ un morphisme d'espaces alg\'ebriques sur $S$. Supposons que
\begin{enumerate}
\item $f$ soit propre, plat et de pr\'esentation finie ;
\item pour tout morphisme $\Spec(k)\to Y$, o\`u $k$ est un corps, on ait
$k=H^0(X_k,\mathcal{O}_{X_k})$.
\end{enumerate}
Alors
\begin{enumerate}
\item[(a)] $f_*\mathcal{O}_X=\mathcal{O}_Y$, et cette \'egalit\'e subsiste apr\`es
tout changement de base ;
\item[(b)] localement pour la topologie \'etale sur $Y$, on a
$$
Rf_*\mathcal{O}_X = \mathcal{O}_Y \oplus P
$$
dans $D(\mathcal{O}_Y)$, o\`u $P$ est parfait et de Tor-amplitude dans
$[1,\infty)$.
\end{enumerate}
\end{lemma}

\begin{proof}
Il suffit de d\'emontrer (a) et (b) localement pour la topologie \'etale sur
$Y$ ; nous pouvons donc supposer, ce que nous ferons, que $Y$ est un sch\'ema
affine. D'apr\`es le th\'eor\`eme de cohomologie et changement de base
(lemme \ref{spaces-perfect-lemma-flat-proper-perfect-direct-image-general}), le complexe
$E=Rf_*\mathcal{O}_X$ est parfait et sa formation commute \`a tout changement
de base. En particulier, pour $y\in Y$, on a
$H^0(E\otimes^\mathbf{L}\kappa(y))=
H^0(X_y,\mathcal{O}_{X_y})=\kappa(y)$.
Ainsi, $\beta_0(y)\leq1$ pour tout $y\in Y$, avec les notations du lemme
\ref{spaces-perfect-lemma-jump-loci}. Appliquons le lemme
\ref{spaces-perfect-lemma-locally-closed-where-H0-locally-free} avec $a=0$ et $r=1$.
On obtient un sous-sch\'ema ferm\'e universel $j:Z\to Y$ tel que
$H^0(Lj^*E)$ soit inversible, caract\'eris\'e par l'\'equivalence de (4)(a),
(b) et (c) du lemme. Puisque la formation de $E$ commute au changement de
base, on a
$$
Lf^*E = R\text{pr}_{1, *}\mathcal{O}_{X \times_Y X}.
$$
Le morphisme $\text{pr}_1:X\times_YX\to X$ poss\`ede une section, \`a savoir le
morphisme diagonal $\Delta$ de $X$ relativement à $Y$. On obtient donc des morphismes
$$
\mathcal{O}_X \longrightarrow R\text{pr}_{1, *}\mathcal{O}_{X \times_Y X}
\longrightarrow \mathcal{O}_X
$$
dans $D(\mathcal{O}_X)$, dont le compos\'e est l'identit\'e. Par cons\'equent,
$R\text{pr}_{1,*}\mathcal{O}_{X\times_YX}=\mathcal{O}_X\oplus E'$ dans
$D(\mathcal{O}_X)$. Ainsi, $\mathcal{O}_X$ est un facteur direct de
$H^0(Lf^*E)$, et l'\'equivalence de (4)(c) et (4)(a) dans le lemme pr\'ecit\'e
montre que $X\to Y$ se factorise par $Z$. Puisque $\{X\to Y\}$ est un
recouvrement fppf, on a $Z=Y$. Ainsi, $f_*\mathcal{O}_X$ est un
$\mathcal{O}_Y$-module inversible. On conclut que
$\mathcal{O}_Y\to f_*\mathcal{O}_X$ est un isomorphisme, car un morphisme
d'anneaux $A\to B$ tel que $B$ soit inversible comme $A$-module est un
isomorphisme. Les hypoth\`eses \'etant stables par changement de base, (a) est
d\'emontr\'e.

\medskip\noindent
D\'emonstration de (b). Nous avons vu ci-dessus que, pour tout $y\in Y$, le
morphisme $\mathcal{O}_Y\to H^0(E\otimes^\mathbf{L}\kappa(y))$ est surjectif.
Le lemme \ref{more-algebra-lemma-better-cut-complex-in-two} de Compl\'ements
d'alg\`ebre montre donc que, dans un voisinage ouvert de $y$, on a une
d\'ecomposition $Rf_*\mathcal{O}_X=\mathcal{O}_Y\oplus P$.
\end{proof}

\begin{lemma}
\label{spaces-perfect-lemma-proper-flat-geom-red-connected}
Soit $S$ un sch\'ema.
Soit $f:X\to Y$ un morphisme d'espaces alg\'ebriques sur $S$. Supposons que
\begin{enumerate}
\item $f$ soit propre, plat et de pr\'esentation finie ;
\item les fibres g\'eom\'etriques de $f$ soient r\'eduites et connexes.
\end{enumerate}
Alors $f_*\mathcal{O}_X=\mathcal{O}_Y$, et cette \'egalit\'e subsiste apr\`es
tout changement de base.
\end{lemma}

\begin{proof}
D'apr\`es le lemme \ref{spaces-perfect-lemma-proper-flat-h0}, il suffit de montrer que
$k=H^0(X_k,\mathcal{O}_{X_k})$ pour tout morphisme $\Spec(k)\to Y$, o\`u $k$
est un corps. Cela r\'esulte du lemme
\ref{spaces-over-fields-lemma-proper-geometrically-reduced-global-sections}
du chapitre Espaces sur un corps, et du fait que $X_k$ est g\'eom\'etriquement connexe
et g\'eom\'etriquement r\'eduit.
\end{proof}






\section{Autres applications}
\label{spaces-perfect-section-other-applications}

\noindent
Nous \'enon\c{c}ons et d\'emontrons dans cette section quelques r\'esultats qui
se d\'eduisent de la th\'eorie d\'evelopp\'ee ci-dessus.

\begin{lemma}
\label{spaces-perfect-lemma-countable-cohomology}
Soit $S$ un sch\'ema. Soit $X$ un espace alg\'ebrique quasi-compact et
quasi-s\'epar\'e sur $S$. Soit $K$ un objet de
$D_\QCoh(\mathcal{O}_X)$ tel que, pour tout sch\'ema affine \'etale sur $X$,
les ensembles de sections des faisceaux de cohomologie $H^i(K)$ soient
d\'enombrables. Alors, pour tout morphisme \'etale quasi-compact et quasi-s\'epar\'e
$U\to X$ et tout objet parfait $E$ de $D(\mathcal{O}_X)$, les ensembles
$$
H^i(U, K \otimes^\mathbf{L} E),\quad \Ext^i(E|_U, K|_U)
$$
sont d\'enombrables.
\end{lemma}

\begin{proof}
D'apr\`es Cohomologie sur les sites, lemme
\ref{sites-cohomology-lemma-dual-perfect-complex}, il suffit de d\'emontrer le
r\'esultat pour les groupes $H^i(U,K\otimes^\mathbf{L}E)$.
Nous allons employer le principe de r\'ecurrence du lemme
\ref{spaces-perfect-lemma-induction-principle}.

\medskip\noindent
Lorsque $U=\Spec(A)$ est affine, le r\'esultat d\'ecoule du cas des sch\'emas ;
voir Cat\'egories d\'eriv\'ees des sch\'emas, lemme
\ref{perfect-lemma-countable-cohomology}.

\medskip\noindent
Pour achever la d\'emonstration, il suffit d'\'etablir ceci : si
$(U\subset W,V\to W)$ est un carr\'e distingu\'e \'el\'ementaire et si le
r\'esultat vaut pour $U$, $V$ et $U\times_WV$, alors il vaut pour $W$.
C'est une cons\'equence imm\'ediate de la suite de Mayer--Vietoris ; voir le
lemme \ref{spaces-perfect-lemma-unbounded-mayer-vietoris}.
\end{proof}

\begin{lemma}
\label{spaces-perfect-lemma-countable}
Soit $S$ un sch\'ema.
Soit $X$ un espace alg\'ebrique quasi-compact et quasi-s\'epar\'e sur $S$.
Supposons d\'enombrables les ensembles de sections de $\mathcal{O}_X$ sur les
sch\'emas affines \'etales sur $X$. Soit $K$ un objet de
$D_\QCoh(\mathcal{O}_X)$. Les conditions suivantes sont \'equivalentes :
\begin{enumerate}
\item $K=\text{hocolim}\ E_n$, o\`u les $E_n$ sont des objets parfaits de
$D(\mathcal{O}_X)$ ;
\item les faisceaux de cohomologie $H^i(K)$ ont des ensembles d\'enombrables
de sections sur les sch\'emas affines \'etales sur $X$.
\end{enumerate}
\end{lemma}

\begin{proof}
Si (1) est vraie, alors (2) l'est aussi, car la prise des faisceaux de
cohomologie commute aux colimites homotopiques (d'apr\`es Cat\'egories
d\'eriv\'ees, lemme \ref{derived-lemma-cohomology-of-hocolim}) et parce qu'un
complexe parfait est localement isomorphe \`a un complexe fini de
$\mathcal{O}_X$-modules libres de type fini ; il v\'erifie donc (2) par
l'hypoth\`ese faite sur $X$.

\medskip\noindent
Supposons (2).
Choisissons un complexe K-injectif $\mathcal{K}^\bullet$ repr\'esentant $K$.
Choisissons un g\'en\'erateur parfait $E$ de
$D_\QCoh(\mathcal{O}_X)$ et repr\'esentons-le par un complexe K-injectif
$\mathcal{I}^\bullet$. D'apr\`es le th\'eor\`eme
\ref{spaces-perfect-theorem-DQCoh-is-Ddga} et sa d\'emonstration, il existe une \'equivalence
de cat\'egories triangul\'ees
$F:D_\QCoh(\mathcal{O}_X)\to D(A,\text{d})$, o\`u $(A,\text{d})$ est
l'alg\`ebre diff\'erentielle gradu\'ee
$$
(A, \text{d}) =
\Hom_{\text{Comp}^{dg}(\mathcal{O}_X)}
(\mathcal{I}^\bullet, \mathcal{I}^\bullet),
$$
qui envoie $K$ sur le module diff\'erentiel gradu\'e
$$
M = \Hom_{\text{Comp}^{dg}(\mathcal{O}_X)}
(\mathcal{I}^\bullet, \mathcal{K}^\bullet).
$$
Notons que $H^i(A)=\Ext^i(E,E)$ et que
$H^i(M)=\Ext^i(E,K)$. De plus, puisque $F$ est une \'equivalence, $F$ et son
quasi-inverse commutent aux colimites homotopiques. Il suffit donc d'\'ecrire
$M$ comme colimite homotopique d'objets compacts de $D(A,\text{d})$.
D'apr\`es Alg\`ebre diff\'erentielle gradu\'ee, lemme
\ref{dga-lemma-countable}, il suffit que $\Ext^i(E,E)$ et $\Ext^i(E,K)$ soient
d\'enombrables pour tout $i$. Cela r\'esulte du lemme
\ref{spaces-perfect-lemma-countable-cohomology}.
\end{proof}

\begin{lemma}
\label{spaces-perfect-lemma-computing-sections-as-colim}
Soit $A$ un anneau. Soit $f:U\to X$ un morphisme plat d'espaces alg\'ebriques
de pr\'esentation finie sur $A$. Alors
\begin{enumerate}
\item il existe un syst\`eme projectif d'objets parfaits $L_n$ de
$D(\mathcal{O}_X)$ tel que
$$
R\Gamma(U, Lf^*K) = \text{hocolim}\ R\Hom_X(L_n, K)
$$
dans $D(A)$, fonctoriellement en $K\in D_\QCoh(\mathcal{O}_X)$ ;
\item il existe un syst\`eme inductif d'objets parfaits $E_n$ de
$D(\mathcal{O}_X)$ tel que
$$
R\Gamma(U, Lf^*K) = \text{hocolim}\ R\Gamma(X, E_n \otimes^\mathbf{L} K)
$$
dans $D(A)$, fonctoriellement en $K\in D_\QCoh(\mathcal{O}_X)$.
\end{enumerate}
\end{lemma}

\begin{proof}
D'apr\`es le lemme \ref{spaces-perfect-lemma-cohomology-base-change}, on a, fonctoriellement
en $K$,
$$
R\Gamma(U, Lf^*K) = R\Gamma(X, Rf_*\mathcal{O}_U \otimes^\mathbf{L} K).
$$
Notons que $R\Gamma(X,-)$ commute aux colimites homotopiques, puisqu'il commute
aux sommes directes d'apr\`es le lemme
\ref{spaces-perfect-lemma-quasi-coherence-pushforward-direct-sums}. De m\^eme,
$-\otimes^\mathbf{L}K$ commute aux colimites d\'eriv\'ees, puisqu'il commute aux
sommes directes (les sommes directes dans $D(\mathcal{O}_X)$ \'etant donn\'ees
par les sommes directes des complexes qui les repr\'esentent). Pour d\'emontrer
(2), il suffit donc d'\'ecrire
$Rf_*\mathcal{O}_U=\text{hocolim}\ E_n$ pour un syst\`eme inductif d'objets
parfaits $E_n$ de $D(\mathcal{O}_X)$. Cela fait, on obtient (1) en posant
$L_n=E_n^\vee$ ; voir Cohomologie sur les sites, lemme
\ref{sites-cohomology-lemma-dual-perfect-complex}.

\medskip\noindent
\'Ecrivons $A=\colim A_i$, avec les $A_i$ de type fini sur $\mathbf{Z}$.
D'apr\`es Limites d'espaces, lemme
\ref{spaces-limits-lemma-descend-finite-presentation}, il existe un $i$ et des
morphismes $U_i\to X_i\to\Spec(A_i)$ de pr\'esentation finie dont le changement
de base \`a $\Spec(A)$ redonne $U\to X\to\Spec(A)$.
Quitte \`a remplacer $i$ par un indice plus grand, on peut supposer $f_i:U_i\to X_i$ plat ; voir
Limites d'espaces, lemme \ref{spaces-limits-lemma-descend-flat}.
D'apr\`es le lemme \ref{spaces-perfect-lemma-compare-base-change}, l'image inverse d\'eriv\'ee
de $Rf_{i,*}\mathcal{O}_{U_i}$ par $g:X\to X_i$ est \'egale \`a
$Rf_*\mathcal{O}_U$. Puisque $Lg^*$ commute aux colimites d\'eriv\'ees, il
suffit de d\'emontrer le r\'esultat pour $f_i$. Nous pouvons donc supposer $U$
et $X$ de type fini sur $\mathbf{Z}$.

\medskip\noindent
Supposons que $f:U\to X$ soit un morphisme d'espaces alg\'ebriques de type fini
sur $\mathbf{Z}$. Pour achever la d\'emonstration, nous allons montrer que
$Rf_*\mathcal{O}_U$ est une colimite homotopique de complexes parfaits.
Appliquons pour cela le lemme \ref{spaces-perfect-lemma-countable}. Il suffit donc de montrer
que les faisceaux $R^if_*\mathcal{O}_U$ ont des ensembles d\'enombrables de
sections sur les sch\'emas affines \'etales sur $X$. Cela r\'esulte du lemme
\ref{spaces-perfect-lemma-countable-cohomology} appliqu\'e au faisceau structural.
\end{proof}









\section{La propri\'et\'e de r\'esolution}
\label{spaces-perfect-section-resolution-property}

\noindent
Cette section est, pour les espaces alg\'ebriques, l'analogue de Cat\'egories
d\'eriv\'ees des sch\'emas, section
\ref{perfect-section-resolution-property} ; nous invitons \`a lire d'abord
cette derni\`ere. On ne sait pas actuellement si tout espace alg\'ebrique lisse
et propre sur un corps poss\`ede la propri\'et\'e de r\'esolution, ou si cette
assertion est fausse. Si vous connaissez la r\'eponse, veuillez \'ecrire \`a
\href{mailto:stacks.project@gmail.com}{stacks.project@gmail.com}.

\medskip\noindent
Nous pouvons donner la d\'efinition suivante, bien que sa consid\'eration pour
des espaces alg\'ebriques g\'en\'eraux n'ait gu\`ere de sens.

\begin{definition}
\label{spaces-perfect-definition-resolution-property}
Soit $S$ un sch\'ema. Soit $X$ un espace alg\'ebrique sur $S$.
On dit que $X$ poss\`ede la {\it propri\'et\'e de r\'esolution} si tout
$\mathcal{O}_X$-module quasi-coh\'erent de type fini est quotient d'un
$\mathcal{O}_X$-module localement libre de type fini.
\end{definition}

\noindent
Si $X$ est un espace alg\'ebrique quasi-compact et quasi-s\'epar\'e, il suffit de
v\'erifier que tout $\mathcal{O}_X$-module de pr\'esentation finie
(automatiquement quasi-coh\'erent) est quotient d'un $\mathcal{O}_X$-module
localement libre de type fini ; voir Limites d'espaces, lemme
\ref{spaces-limits-lemma-finite-directed-colimit-surjective-maps}.
Si $X$ est un espace alg\'ebrique noeth\'erien, les modules quasi-coh\'erents de
type fini sont exactement les $\mathcal{O}_X$-modules coh\'erents ; voir
Cohomologie des espaces, lemme
\ref{spaces-cohomology-lemma-coherent-Noetherian}.

\begin{lemma}
\label{spaces-perfect-lemma-resolution-property-ample-relative}
Soit $S$ un sch\'ema.
Soit $f:X\to Y$ un morphisme d'espaces alg\'ebriques sur $S$. Supposons que
\begin{enumerate}
\item $Y$ soit quasi-compact et quasi-s\'epar\'e et poss\`ede la propri\'et\'e de
r\'esolution ;
\item il existe sur $X$ un module inversible $f$-ample (Diviseurs sur les
espaces, d\'efinition
\ref{spaces-divisors-definition-relatively-ample}).
\end{enumerate}
Alors $X$ poss\`ede la propri\'et\'e de r\'esolution.
\end{lemma}

\begin{proof}
Soit $\mathcal{F}$ un $\mathcal{O}_X$-module quasi-coh\'erent de type fini.
Soit $\mathcal{L}$ un module inversible $f$-ample.
Choisissons un sch\'ema affine $V$ et un morphisme \'etale surjectif
$V\to Y$. Posons $U=V\times_YX$. Alors $\mathcal{L}|_U$ est ample sur $U$.
D'apr\`es Propri\'et\'es, proposition
\ref{properties-proposition-characterize-ample}, il existe un nombre fini de
morphismes
$s_i:\mathcal{L}^{\otimes n_i}|_U\to\mathcal{F}|_U$ conjointement surjectifs.
Consid\'erons les $\mathcal{O}_Y$-modules quasi-coh\'erents
$$
\mathcal{H}_n =
f_*(\mathcal{F} \otimes_{\mathcal{O}_X} \mathcal{L}^{\otimes n}).
$$
On peut regarder $s_i$ comme une section sur $V$ du faisceau
$\mathcal{H}_{-n_i}$. Supposons que l'on puisse trouver des
$\mathcal{O}_Y$-modules localement libres de type fini $\mathcal{E}_i$ et des
morphismes $\mathcal{E}_i\to\mathcal{H}_{-n_i}$ dont l'image contienne $s_i$.
Alors les morphismes correspondants
$$
f^*\mathcal{E}_i
\otimes_{\mathcal{O}_X}
\mathcal{L}^{\otimes n_i} \longrightarrow \mathcal{F}
$$
sont conjointement surjectifs, ce qui d\'emontre le lemme. D'apr\`es Limites
d'espaces, lemme \ref{spaces-limits-lemma-directed-colimit-finite-type}, il
existe pour chaque $i$ un sous-module quasi-coh\'erent de type fini
$\mathcal{H}'_i\subset\mathcal{H}_{-n_i}$ qui contient la section $s_i$ sur
$V$. La propri\'et\'e de r\'esolution de $Y$ fournit alors des surjections
$\mathcal{E}_i\to\mathcal{H}'_i$, d'o\`u la conclusion.
\end{proof}

\begin{lemma}
\label{spaces-perfect-lemma-resolution-property-goes-up-affine}
Soit $S$ un sch\'ema.
Soit $f:X\to Y$ un morphisme affine ou quasi-affine d'espaces alg\'ebriques
sur $S$, avec $Y$ quasi-compact et quasi-s\'epar\'e.
Si $Y$ poss\`ede la propri\'et\'e de r\'esolution, alors $X$ la poss\`ede aussi.
\end{lemma}

\begin{proof}
D'apr\`es Diviseurs sur les espaces, lemme
\ref{spaces-divisors-lemma-quasi-affine-O-ample}, c'est un cas particulier du
lemme \ref{spaces-perfect-lemma-resolution-property-ample-relative}.
\end{proof}

\noindent
Voici un cas o\`u l'on peut faire descendre la propri\'et\'e de r\'esolution.

\begin{lemma}
\label{spaces-perfect-lemma-resolution-property-goes-down-finite-flat}
Soit $S$ un sch\'ema.
Soit $f:X\to Y$ un morphisme surjectif, fini et localement libre d'espaces
alg\'ebriques sur $S$.
Si $X$ poss\`ede la propri\'et\'e de r\'esolution, alors $Y$ la poss\`ede aussi.
\end{lemma}

\begin{proof}
La condition signifie que $f$ est affine et que $f_*\mathcal{O}_X$ est un
$\mathcal{O}_Y$-module localement libre de type fini et de rang positif.
Soit $\mathcal{G}$ un $\mathcal{O}_Y$-module quasi-coh\'erent de type fini.
Par hypoth\`ese, il existe une surjection
$\mathcal{E}\to f^*\mathcal{G}$, o\`u $\mathcal{E}$ est un
$\mathcal{O}_X$-module localement libre de type fini. Puisque $f_*$ est exact
(Cohomologie des espaces, section
\ref{spaces-cohomology-section-finite-morphisms}), on obtient une surjection
$$
f_*\mathcal{E}
\longrightarrow
f_*f^*\mathcal{G} = \mathcal{G} \otimes_{\mathcal{O}_Y} f_*\mathcal{O}_X.
$$
En tensorisant cette surjection par le dual de $f_*\mathcal{O}_X$, puis en composant avec l'évaluation, on obtient une surjection
$$
f_*\mathcal{E}
\otimes_{\mathcal{O}_Y}
\SheafHom_{\mathcal{O}_Y}(f_*\mathcal{O}_X, \mathcal{O}_Y)
\longrightarrow
\mathcal{G}.
$$
Comme $f_*\mathcal{E}$ est localement libre de type fini, on conclut.
\end{proof}

\noindent
Pour d'autres r\'esultats sur la propri\'et\'e de r\'esolution des espaces
alg\'ebriques, voir Compl\'ements sur les morphismes d'espaces, section
\ref{spaces-more-morphisms-section-resolution-property}.












\section{D\'etection de la bornitude}
\label{spaces-perfect-section-detecting-boundedness}

\noindent
Nous montrons dans cette section que les g\'en\'erateurs compacts de $D_\QCoh$
d'un espace alg\'ebrique quasi-compact et quasi-s\'epar\'e, construits dans la
section \ref{spaces-perfect-section-generating}, poss\`edent une propri\'et\'e particuli\`ere.
Nous conseillons de lire d'abord cette section, tr\`es proche de celle-ci.

\begin{lemma}
\label{spaces-perfect-lemma-bounded-truncation}
Soit $S$ un sch\'ema.
Soit $X$ un espace alg\'ebrique quasi-compact et quasi-s\'epar\'e sur $S$.
Soient $P\in D_{perf}(\mathcal{O}_X)$ et
$E\in D_{\QCoh}(\mathcal{O}_X)$. Soit $a\in\mathbf{Z}$.
Les conditions suivantes sont \'equivalentes :
\begin{enumerate}
\item $\Hom_{D(\mathcal{O}_X)}(P[-i],E)=0$ pour $i\gg0$ ;
\item $\Hom_{D(\mathcal{O}_X)}(P[-i],\tau_{\geq a}E)=0$ pour $i\gg0$.
\end{enumerate}
\end{lemma}

\begin{proof}
En utilisant le triangle $\tau_{<a}E\to E\to\tau_{\geq a}E\to$, on voit
qu'il suffit, pour obtenir l'\'equivalence, de montrer que
$$
\Hom_{D(\mathcal{O}_X)}(P[-i], \tau_{< a} E) =
\Hom_{D(\mathcal{O}_X)}(P, (\tau_{< a} E)[i]) = 0
$$
pour $i\gg0$. Comme $P$ est parfait, cela r\'esulte du lemme
\ref{spaces-perfect-lemma-ext-from-perfect-into-bounded-QCoh}.
\end{proof}

\begin{lemma}
\label{spaces-perfect-lemma-bounded-below-truncation}
Soit $S$ un sch\'ema.
Soit $X$ un espace alg\'ebrique quasi-compact et quasi-s\'epar\'e sur $S$.
Soient $P\in D_{perf}(\mathcal{O}_X)$ et
$E\in D_{\QCoh}(\mathcal{O}_X)$. Soit $a\in\mathbf{Z}$.
Les conditions suivantes sont \'equivalentes :
\begin{enumerate}
\item $\Hom_{D(\mathcal{O}_X)}(P[-i],E)=0$ pour $i\ll0$ ;
\item $\Hom_{D(\mathcal{O}_X)}(P[-i],\tau_{\leq a}E)=0$ pour $i\ll0$.
\end{enumerate}
\end{lemma}

\begin{proof}
En utilisant le triangle $\tau_{\leq a}E\to E\to\tau_{>a}E\to$, on voit
qu'il suffit, pour obtenir l'\'equivalence, de montrer que
$$
\Hom_{D(\mathcal{O}_X)}(P[-i], \tau_{> a} E) =
\Hom_{D(\mathcal{O}_X)}(P, (\tau_{> a} E)[i]) = 0
$$
pour $i\ll0$. Comme $P$ est parfait, cela r\'esulte du lemme
\ref{spaces-perfect-lemma-ext-from-perfect-into-bounded-QCoh}.
\end{proof}

\begin{proposition}
\label{spaces-perfect-proposition-detecting-bounded-above}
Soit $S$ un sch\'ema.
Soit $X$ un espace alg\'ebrique quasi-compact et quasi-s\'epar\'e sur $S$.
Soit $G\in D_{perf}(\mathcal{O}_X)$ un complexe parfait qui engendre
$D_\QCoh(\mathcal{O}_X)$. Soit $E\in D_\QCoh(\mathcal{O}_X)$.
Les conditions suivantes sont \'equivalentes :
\begin{enumerate}
\item $E\in D^-_\QCoh(\mathcal{O}_X)$ ;
\item $\Hom_{D(\mathcal{O}_X)}(G[-i],E)=0$ pour $i\gg0$ ;
\item $\Ext^i_X(G,E)=0$ pour $i\gg0$ ;
\item $R\Hom_X(G,E)$ appartient \`a $D^-(\mathbf{Z})$ ;
\item $H^i(X,G^\vee\otimes_{\mathcal{O}_X}^\mathbf{L}E)=0$ pour $i\gg0$ ;
\item $R\Gamma(X,G^\vee\otimes_{\mathcal{O}_X}^\mathbf{L}E)$ appartient \`a
$D^-(\mathbf{Z})$ ;
\item pour tout objet parfait $P$ de $D(\mathcal{O}_X)$,
\begin{enumerate}
\item les assertions (2), (3) et (4) valent apr\`es remplacement de $G$ par
$P$ ;
\item $H^i(X,P\otimes_{\mathcal{O}_X}^\mathbf{L}E)=0$ pour $i\gg0$ ;
\item $R\Gamma(X,P\otimes_{\mathcal{O}_X}^\mathbf{L}E)$ appartient \`a
$D^-(\mathbf{Z})$.
\end{enumerate}
\end{enumerate}
\end{proposition}

\begin{proof}
Supposons (1). Comme
$\Hom_{D(\mathcal{O}_X)}(G[-i],E)=\Hom_{D(\mathcal{O}_X)}(G,E[i])$, ce groupe
est nul pour $i\gg0$ d'apr\`es le lemme
\ref{spaces-perfect-lemma-ext-from-perfect-into-bounded-QCoh}. Ainsi, (1) implique (2).

\medskip\noindent
Les assertions (2), (3) et (4) sont \'equivalentes d'apr\`es Cohomologie sur
les sites, section \ref{sites-cohomology-section-global-RHom}.
Les assertions (5) et (6) sont \'equivalentes puisque, par d\'efinition,
$H^i(X,-)=H^i(R\Gamma(X,-))$. D'apr\`es le lemme
\ref{sites-cohomology-lemma-dual-perfect-complex} de Cohomologie sur les sites
et l'\'egalit\'e $(G[-i])^\vee=G^\vee[i]$, elles sont \'equivalentes aux
assertions (2), (3) et (4).

\medskip\noindent
Il est clair que (7) implique (2). R\'eciproquement, montrons que les conditions
\'equivalentes (2)--(6) impliquent (7). Rappelons que $G$ est un g\'en\'erateur
classique de $D_{perf}(\mathcal{O}_X)$ d'apr\`es la remarque
\ref{spaces-perfect-remark-classical-generator}. Pour $P\in D_{perf}(\mathcal{O}_X)$, notons
$T(P)$ l'assertion selon laquelle $R\Hom_X(P,E)$ appartient \`a
$D^-(\mathbf{Z})$. La classe des $P$ tels que $T(P)$ soit vraie est stable par
sommes directes, facteurs directs et d\'ecalages, et v\'erifie la propri\'et\'e des
deux sur trois pour les triangles distingu\'es. La remarque
\ref{derived-remark-check-on-generator} de Cat\'egories d\'eriv\'ees montre donc
que (4) implique $T(P)$ pour tout $P\in D_{perf}(\mathcal{O}_X)$. Le m\^eme
raisonnement s'applique \`a toutes les autres propri\'et\'es ; pour (7)(b) et
(7)(c), on utilise aussi que $P\mapsto P^\vee$ est une auto-\'equivalence de
$D_{perf}(\mathcal{O}_X)$. Nous omettons ce petit d\'etail.

\medskip\noindent
Nous allons montrer que les conditions \'equivalentes (2)--(7) impliquent (1)
en employant le principe de r\'ecurrence du lemme
\ref{spaces-perfect-lemma-induction-principle}.

\medskip\noindent
Supposons d'abord $X$ affine. L'implication (2)--(7) $\Rightarrow$ (1) d\'ecoule
du cas des sch\'emas ; voir Cat\'egories d\'eriv\'ees des sch\'emas, proposition
\ref{perfect-proposition-detecting-bounded-above}.

\medskip\noindent
Supposons maintenant que $(U\subset X,j:V\to X)$ soit un carr\'e distingu\'e
\'el\'ementaire d'espaces alg\'ebriques quasi-compacts et quasi-s\'epar\'es sur
$S$, et que l'implication (2)--(7) $\Rightarrow$ (1) soit connue pour $U$, $V$
et $U\times_XV$. Il reste \`a la d\'emontrer pour $X$. Supposons donc que
$E\in D_\QCoh(\mathcal{O}_X)$ v\'erifie (2)--(7). D'apr\`es le lemme
\ref{spaces-perfect-lemma-direct-summand-of-a-restriction} et le th\'eor\`eme
\ref{spaces-perfect-theorem-bondal-van-den-Bergh}, il existe sur $X$ un complexe parfait $Q$
tel que $Q|_U$ engendre $D_\QCoh(\mathcal{O}_U)$.

\medskip\noindent
\'Ecrivons $V=\Spec(A)$. Soit $Z\subset V$ le sous-sch\'ema ferm\'e r\'eduit,
image inverse de $X\setminus U$, qui s'envoie isomorphiquement sur celui-ci
(voir la d\'efinition \ref{spaces-perfect-definition-elementary-distinguished-square}).
C'est une partie ferm\'ee r\'etrocompacte de $V$. Choisissons
$f_1,\ldots,f_r\in A$ tels que $Z=V(f_1,\ldots,f_r)$. Soit
$K\in D(\mathcal{O}_V)$ l'objet parfait correspondant au complexe de Koszul
associ\'e \`a $f_1,\ldots,f_r$ sur $A$. Puisque $K$ est \`a support dans $Z$,
son image directe $K'=Rj_*K$ est un objet parfait de $D(\mathcal{O}_X)$ dont
la restriction \`a $V$ est $K$ (voir les lemmes
\ref{spaces-perfect-lemma-pushforward-perfect} et
\ref{spaces-perfect-lemma-pushforward-with-support-in-open}). Par hypoth\`ese,
$R\Hom_{\mathcal{O}_X}(Q,E)$ et $R\Hom_{\mathcal{O}_X}(K',E)$ sont born\'es
sup\'erieurement.

\medskip\noindent
D'apr\`es le lemme \ref{spaces-perfect-lemma-pushforward-with-support-in-open}, on a
$K'=j_!K$, et donc
$$
\Hom_{D(\mathcal{O}_X)}(K'[-i], E) = \Hom_{D(\mathcal{O}_V)}(K[-i], E|_V) = 0
$$
pour $i\gg0$. On peut donc appliquer \`a $E|_V$ le lemme
\ref{perfect-lemma-orthogonal-koszul-first-variant} de Cat\'egories d\'eriv\'ees
des sch\'emas. Il existe un entier $a$ tel que
$\tau_{\geq a}(E|_V)=
\tau_{\geq a}R(U\times_XV\to V)_*(E|_{U\times_XV})$.
Alors $\tau_{\geq a}E=\tau_{\geq a}R(U\to X)_*(E|_U)$ ; on le v\'erifie en
restreignant le morphisme canonique \`a $U$ et \`a $V$. En appliquant deux fois
le lemme \ref{spaces-perfect-lemma-bounded-truncation}, on obtient
$$
\Hom_{D(\mathcal{O}_U)}(Q|_U [-i], E|_U) =
\Hom_{D(\mathcal{O}_X)}(Q[-i], R (U \to X)_* (E|_U)) = 0
$$
pour $i\gg0$. Puisque la proposition vaut pour $U$ et le g\'en\'erateur
$Q|_U$, on a $E|_U\in D^-_\QCoh(\mathcal{O}_U)$. Le foncteur
$R(U\to X)_*$ pr\'eserve $D^-_\QCoh$ d'apr\`es le lemme
\ref{spaces-perfect-lemma-quasi-coherence-direct-image} ; ainsi,
$\tau_{\geq a}E\in D^-_\QCoh(\mathcal{O}_X)$, puis
$E\in D^-_\QCoh(\mathcal{O}_X)$.
\end{proof}

\begin{proposition}
\label{spaces-perfect-proposition-detecting-bounded-below}
Soit $S$ un sch\'ema.
Soit $X$ un espace alg\'ebrique quasi-compact et quasi-s\'epar\'e sur $S$.
Soit $G\in D_{perf}(\mathcal{O}_X)$ un complexe parfait qui engendre
$D_\QCoh(\mathcal{O}_X)$. Soit $E\in D_\QCoh(\mathcal{O}_X)$.
Les conditions suivantes sont \'equivalentes :
\begin{enumerate}
\item $E\in D^+_\QCoh(\mathcal{O}_X)$ ;
\item $\Hom_{D(\mathcal{O}_X)}(G[-i],E)=0$ pour $i\ll0$ ;
\item $\Ext^i_X(G,E)=0$ pour $i\ll0$ ;
\item $R\Hom_X(G,E)$ appartient \`a $D^+(\mathbf{Z})$ ;
\item $H^i(X,G^\vee\otimes_{\mathcal{O}_X}^\mathbf{L}E)=0$ pour $i\ll0$ ;
\item $R\Gamma(X,G^\vee\otimes_{\mathcal{O}_X}^\mathbf{L}E)$ appartient \`a
$D^+(\mathbf{Z})$ ;
\item pour tout objet parfait $P$ de $D(\mathcal{O}_X)$,
\begin{enumerate}
\item les assertions (2), (3) et (4) valent apr\`es remplacement de $G$ par
$P$ ;
\item $H^i(X,P\otimes_{\mathcal{O}_X}^\mathbf{L}E)=0$ pour $i\ll0$ ;
\item $R\Gamma(X,P\otimes_{\mathcal{O}_X}^\mathbf{L}E)$ appartient \`a
$D^+(\mathbf{Z})$.
\end{enumerate}
\end{enumerate}
\end{proposition}

\begin{proof}
Supposons (1). Comme
$\Hom_{D(\mathcal{O}_X)}(G[-i],E)=\Hom_{D(\mathcal{O}_X)}(G,E[i])$, ce groupe
est nul pour $i\ll0$ d'apr\`es le lemme
\ref{spaces-perfect-lemma-ext-from-perfect-into-bounded-QCoh}. Ainsi, (1) implique (2).

\medskip\noindent
Les assertions (2), (3) et (4) sont \'equivalentes d'apr\`es Cohomologie sur
les sites, section \ref{sites-cohomology-section-global-RHom}.
Les assertions (5) et (6) sont \'equivalentes puisque, par d\'efinition,
$H^i(X,-)=H^i(R\Gamma(X,-))$. D'apr\`es le lemme
\ref{sites-cohomology-lemma-dual-perfect-complex} de Cohomologie sur les sites
et l'\'egalit\'e $(G[-i])^\vee=G^\vee[i]$, elles sont \'equivalentes aux
assertions (2), (3) et (4).

\medskip\noindent
Il est clair que (7) implique (2). R\'eciproquement, montrons que les conditions
\'equivalentes (2)--(6) impliquent (7). Rappelons que $G$ est un g\'en\'erateur
classique de $D_{perf}(\mathcal{O}_X)$ d'apr\`es la remarque
\ref{spaces-perfect-remark-classical-generator}. Pour $P\in D_{perf}(\mathcal{O}_X)$, notons
$T(P)$ l'assertion selon laquelle $R\Hom_X(P,E)$ appartient \`a
$D^+(\mathbf{Z})$. La classe des $P$ tels que $T(P)$ soit vraie est stable par
sommes directes, facteurs directs et d\'ecalages, et v\'erifie la propri\'et\'e des
deux sur trois pour les triangles distingu\'es. La remarque
\ref{derived-remark-check-on-generator} de Cat\'egories d\'eriv\'ees montre donc
que (4) implique $T(P)$ pour tout $P\in D_{perf}(\mathcal{O}_X)$. Le m\^eme
raisonnement s'applique \`a toutes les autres propri\'et\'es ; pour (7)(b) et
(7)(c), on utilise aussi que $P\mapsto P^\vee$ est une auto-\'equivalence de
$D_{perf}(\mathcal{O}_X)$. Nous omettons ce petit d\'etail.

\medskip\noindent
Nous allons montrer que les conditions \'equivalentes (2)--(7) impliquent (1)
en employant le principe de r\'ecurrence du lemme
\ref{spaces-perfect-lemma-induction-principle}.

\medskip\noindent
Supposons d'abord $X$ affine. L'implication (2)--(7) $\Rightarrow$ (1) d\'ecoule
du cas des sch\'emas ; voir Cat\'egories d\'eriv\'ees des sch\'emas, proposition
\ref{perfect-proposition-detecting-bounded-below}.

\medskip\noindent
Supposons maintenant que $(U\subset X,j:V\to X)$ soit un carr\'e distingu\'e
\'el\'ementaire d'espaces alg\'ebriques quasi-compacts et quasi-s\'epar\'es sur
$S$, et que l'implication (2)--(7) $\Rightarrow$ (1) soit connue pour $U$, $V$
et $U\times_XV$. Il reste \`a la d\'emontrer pour $X$. Supposons donc que
$E\in D_\QCoh(\mathcal{O}_X)$ v\'erifie (2)--(7). D'apr\`es le lemme
\ref{spaces-perfect-lemma-direct-summand-of-a-restriction} et le th\'eor\`eme
\ref{spaces-perfect-theorem-bondal-van-den-Bergh}, il existe sur $X$ un complexe parfait $Q$
tel que $Q|_U$ engendre $D_\QCoh(\mathcal{O}_U)$.

\medskip\noindent
\'Ecrivons $V=\Spec(A)$. Soit $Z\subset V$ le sous-sch\'ema ferm\'e r\'eduit,
image inverse de $X\setminus U$, qui s'envoie isomorphiquement sur celui-ci
(voir la d\'efinition \ref{spaces-perfect-definition-elementary-distinguished-square}).
C'est une partie ferm\'ee r\'etrocompacte de $V$. Choisissons
$f_1,\ldots,f_r\in A$ tels que $Z=V(f_1,\ldots,f_r)$. Soit
$K\in D(\mathcal{O}_V)$ l'objet parfait correspondant au complexe de Koszul
associ\'e \`a $f_1,\ldots,f_r$ sur $A$. Puisque $K$ est \`a support dans $Z$,
son image directe $K'=Rj_*K$ est un objet parfait de $D(\mathcal{O}_X)$ dont
la restriction \`a $V$ est $K$ (voir les lemmes
\ref{spaces-perfect-lemma-pushforward-perfect} et
\ref{spaces-perfect-lemma-pushforward-with-support-in-open}). Par hypoth\`ese,
$R\Hom_{\mathcal{O}_X}(Q,E)$ et $R\Hom_{\mathcal{O}_X}(K',E)$ sont born\'es
inf\'erieurement.

\medskip\noindent
D'apr\`es le lemme \ref{spaces-perfect-lemma-pushforward-with-support-in-open}, on a
$K'=j_!K$, et donc
$$
\Hom_{D(\mathcal{O}_X)}(K'[-i], E) = \Hom_{D(\mathcal{O}_V)}(K[-i], E|_V) = 0
$$
pour $i\ll0$. On peut donc appliquer \`a $E|_V$ le lemme
\ref{perfect-lemma-orthogonal-koszul-second-variant} de Cat\'egories d\'eriv\'ees
des sch\'emas. Il existe un entier $a$ tel que
$\tau_{\leq a}(E|_V)=
\tau_{\leq a}R(U\times_XV\to V)_*(E|_{U\times_XV})$.
Alors $\tau_{\leq a}E=\tau_{\leq a}R(U\to X)_*(E|_U)$ ; on le v\'erifie en
restreignant le morphisme canonique \`a $U$ et \`a $V$. En appliquant deux fois
le lemme \ref{spaces-perfect-lemma-bounded-below-truncation}, on obtient
$$
\Hom_{D(\mathcal{O}_U)}(Q|_U [-i], E|_U) =
\Hom_{D(\mathcal{O}_X)}(Q[-i], R (U \to X)_* (E|_U)) = 0
$$
pour $i\ll0$. Puisque la proposition vaut pour $U$ et le g\'en\'erateur
$Q|_U$, on a $E|_U\in D^+_\QCoh(\mathcal{O}_U)$. Le foncteur
$R(U\to X)_*$ pr\'eserve $D^+_\QCoh$ d'apr\`es le lemme
\ref{spaces-perfect-lemma-quasi-coherence-direct-image} ; ainsi,
$\tau_{\leq a}E\in D^+_\QCoh(\mathcal{O}_X)$, puis
$E\in D^+_\QCoh(\mathcal{O}_X)$.
\end{proof}










\section{Objets quasi-coh\'erents dans la cat\'egorie d\'eriv\'ee}
\label{spaces-perfect-section-QC}

\noindent
Soit $S$ un sch\'ema. Soit $X$ un espace alg\'ebrique sur $S$.
Rappelons que $X_{affine,\etale}$ d\'esigne la cat\'egorie des objets affines de
$X_\etale$, munie de la topologie d\'efinie par les recouvrements \'etales
standard ; voir Propri\'et\'es des espaces, d\'efinition
\ref{spaces-properties-definition-affine-etale-site}.
Rappelons aussi que le topos de $X_{affine,\etale}$ est le petit topos \'etale
de $X$ ; voir Propri\'et\'es des espaces, lemme
\ref{spaces-properties-lemma-alternative}. Le site $X_\etale$ est muni d'un
faisceau structural $\mathcal{O}_X$, dont nous notons encore
$\mathcal{O}_X$ la restriction \`a $X_{affine,\etale}$. Il existe alors une
\'equivalence de topos annel\'es
$$
(\Sh(X_{affine, \etale}), \mathcal{O}_X)
\longrightarrow
(\Sh(X_\etale), \mathcal{O}_X).
$$
Voir Descente sur les espaces, \'equation
(\ref{spaces-descent-equation-alternative-small-ringed}), ainsi que la discussion de
Descente sur les espaces, section
\ref{spaces-descent-section-alternative-quasi-coherent}.

\medskip\noindent
Dans cette section, nous notons $X_{affine}$ la cat\'egorie sous-jacente \`a
$X_{affine,\etale}$, munie de la topologie chaotique, c'est-\`a-dire telle que
les faisceaux s'identifient aux pr\'efaisceaux. En particulier, le faisceau
structural $\mathcal{O}_X$ devient aussi un faisceau sur $X_{affine}$.
On obtient un morphisme de sites annel\'es
$$
\epsilon :
(X_{affine, \etale}, \mathcal{O}_X)
\longrightarrow
(X_{affine}, \mathcal{O}_X)
$$
comme dans Cohomologie sur les sites, section
\ref{sites-cohomology-section-compare}. Nous identifierons dans cette section
$D_\QCoh(\mathcal{O}_X)$ \`a la cat\'egorie
$\mathit{QC}(X_{affine},\mathcal{O}_X)$ introduite dans Cohomologie sur les
sites, section \ref{sites-cohomology-section-modules-cohomology}.

\begin{lemma}
\label{spaces-perfect-lemma-DQCoh-alternative-small}
Dans la situation ci-dessus, il existe des \'equivalences exactes canoniques
entre les cat\'egories triangul\'ees suivantes :
\begin{enumerate}
\item $D_\QCoh(\mathcal{O}_X)$ ;
\item $D_\QCoh(X_{affine, \etale}, \mathcal{O}_X)$ ;
\item $D_\QCoh(X_{affine}, \mathcal{O}_X)$ ;
\item $\mathit{QC}(X_{affine}, \mathcal{O}_X)$.
\end{enumerate}
\end{lemma}

\begin{proof}
Si $U\to V\to X$ sont des morphismes \'etales et si $U$ et $V$ sont affines,
alors le morphisme d'anneaux
$\mathcal{O}_X(V)\to\mathcal{O}_X(U)$ est plat. L'\'equivalence de (3) et (4)
est donc un cas particulier de Cohomologie sur les sites, lemme
\ref{sites-cohomology-lemma-cartesian-flat-quasi-coherent}, dont la
d\'emonstration \'eclaire aussi l'\'enonc\'e.

\medskip\noindent
La discussion qui pr\'ec\`ede le lemme montre qu'il existe une \'equivalence de
topos annel\'es
$(\Sh(X_{affine,\etale}),\mathcal{O}_X)\to
(\Sh(X_\etale),\mathcal{O}_X)$, donc une \'equivalence entre les cat\'egories
ab\'eliennes de modules. Puisque la notion de module quasi-coh\'erent est
intrins\`eque (Modules sur les sites, lemme
\ref{sites-modules-lemma-special-locally-free}), cette \'equivalence pr\'eserve
les sous-cat\'egories de modules quasi-coh\'erents. On obtient ainsi une
\'equivalence exacte canonique entre les cat\'egories triangul\'ees de (1) et
(2).

\medskip\noindent
Pour obtenir une \'equivalence exacte entre les cat\'egories triangul\'ees de
(2) et (3), appliquons Cohomologie sur les sites, lemme
\ref{sites-cohomology-lemma-compare-topologies-derived-adequate-modules}, au
morphisme
$\epsilon:(X_{affine,\etale},\mathcal{O}_X)\to
(X_{affine},\mathcal{O}_X)$ ci-dessus. Prenons
$\mathcal{B}=\Ob(X_{affine})$ et prenons pour
$\mathcal{A}\subset\textit{PMod}(X_{affine},\mathcal{O})$ la sous-cat\'egorie
pleine des pr\'efaisceaux $\mathcal{F}$ tels que
$\mathcal{F}(V)\otimes_{\mathcal{O}_X(V)}\mathcal{O}_X(U)\to\mathcal{F}(U)$
soit un isomorphisme. D'apr\`es Descente sur les espaces, lemme
\ref{spaces-descent-lemma-quasi-coherent-alternative-small}, les objets de
$\mathcal{A}$ sont exactement les faisceaux pour la topologie \'etale qui sont
des modules quasi-coh\'erents sur
$(X_{affine,\etale},\mathcal{O}_X)$. D'autre part, d'apr\`es Modules sur les
sites, lemme \ref{sites-modules-lemma-chaotic-quasi-coherent}, les objets de
$\mathcal{A}$ sont exactement les modules quasi-coh\'erents sur
$(X_{affine},\mathcal{O}_X)$, c'est-\`a-dire pour la topologie chaotique. Si les
hypoth\`eses du lemme
\ref{sites-cohomology-lemma-compare-topologies-derived-adequate-modules} de
Cohomologie sur les sites sont satisfaites, $\epsilon^*$ et $R\epsilon_*$
fournissent donc bien l'\'equivalence canonique entre les cat\'egories de (2) et
(3).

\medskip\noindent
Il reste \`a v\'erifier quatre conditions :
\begin{enumerate}
\item Tout objet de $\mathcal{A}$ est un faisceau pour la topologie \'etale.
Nous l'avons vu ci-dessus.
\item $\mathcal{A}$ est une sous-cat\'egorie de Serre faible de
$\textit{Mod}(X_{affine,\etale},\mathcal{O}_X)$. Nous avons vu ci-dessus que
$\mathcal{A}=\QCoh(X_{affine,\etale},\mathcal{O}_X)$ et que, par
l'\'equivalence
$\textit{Mod}(X_{affine,\etale},\mathcal{O})=
\textit{Mod}(X_\etale,\mathcal{O}_X)$, ses objets correspondent aux modules
quasi-coh\'erents sur $X$. L'assertion r\'esulte donc du lemme
\ref{spaces-properties-lemma-properties-quasi-coherent} de Propri\'et\'es des
espaces et du lemme \ref{homology-lemma-characterize-weak-serre-subcategory}
d'Homologie.
\item Tout objet de $X_{affine}$ poss\`ede, pour la topologie chaotique, un
recouvrement dont les membres appartiennent \`a $\mathcal{B}$. Cela tient au
fait que $\mathcal{B}$ contient tous les objets.
\item Pour tout objet $U$ de $X_{affine}$ et tout $\mathcal{F}$ de
$\mathcal{A}$, on a $H^p_{Zar}(U,\mathcal{F})=0$ pour $p>0$. Cela r\'esulte de
l'annulation de la cohomologie des modules quasi-coh\'erents sur les sch\'emas
affines ; voir la discussion de Cohomologie des espaces, section
\ref{spaces-cohomology-section-higher-direct-image}, et Cohomologie des
sch\'emas, lemme
\ref{coherent-lemma-quasi-coherent-affine-cohomology-zero}.
\end{enumerate}
La d\'emonstration est achev\'ee.
\end{proof}

\begin{remark}
\label{spaces-perfect-remark-QC-compare}
Soit $S$ un sch\'ema. Soit $X$ un espace alg\'ebrique sur $S$.
Nous montrerons plus loin que
$\mathit{QC}((\textit{Aff}/X),\mathcal{O})$ est aussi canoniquement
\'equivalente \`a $D_\QCoh(\mathcal{O}_X)$. Voir Faisceaux sur les champs,
proposition \ref{stacks-sheaves-proposition-QC-compare}.
\end{remark}










% OK, start here.
%

\chapter{Compléments sur les morphismes d'espaces algébriques}



\phantomsection
\label{spaces-more-morphisms-section-phantom}


\section{Introduction}
\label{spaces-more-morphisms-section-introduction}

\noindent
Dans ce chapitre, nous poursuivons notre étude des propriétés des morphismes d'espaces algébriques.
Une référence fondamentale est \cite{Kn}.





\section{Conventions}
\label{spaces-more-morphisms-section-conventions}

\noindent
Nous supposons une fois pour toutes que tous les schémas appartiennent à
un grand site fppf $\Sch_{fppf}$. De plus, tous les anneaux $A$ considérés
ont la propriété que $\Spec(A)$ est (isomorphe) à un
objet de ce grand site.

\medskip\noindent
Soit $S$ un schéma et soit $X$ un espace algébrique sur $S$.
Dans ce chapitre et les suivants, nous noterons $X \times_S X$
le produit de $X$ par lui-même (dans la catégorie des espaces algébriques
sur $S$), au lieu de $X \times X$.








\section{Morphismes radiciels}
\label{spaces-more-morphisms-section-radicial}

\noindent
Il s'avère qu'un morphisme radiciel n'est pas la même chose qu'un morphisme
universellement injectif, contrairement à ce qui se passe pour les
morphismes de schémas. En fait, il est un peu plus fort.

\begin{definition}
\label{spaces-more-morphisms-definition-radicial}
Soit $S$ un schéma. Soit $f : X \to Y$ un morphisme d'espaces algébriques
au-dessus de $S$. On dit que $f$ est {\it radiciel} si, pour tout morphisme
$\Spec(K) \to Y$ où $K$ est un corps, la réduction
$(\Spec(K) \times_Y X)_{red}$ est soit vide, soit
est représentable par le spectre d'une extension de corps purement inséparable de $K$.
\end{definition}

\begin{lemma}
\label{spaces-more-morphisms-lemma-radicial-implies-universally-injective}
Un morphisme radiciel d'espaces algébriques est universellement injectif.
\end{lemma}

\begin{proof}
Soit $S$ un schéma. Soit $f : X \to Y$ un morphisme radiciel
d'espaces algébriques au-dessus de $S$.
Il est clair d'après la définition que, pour un morphisme
$\Spec(K) \to Y$, il existe au plus un relèvement de ce morphisme
en un morphisme vers $X$. Nous concluons donc que $f$ est universellement
injectif par
Morphismes d'espaces,
lemme \ref{spaces-morphisms-lemma-universally-injective}.
\end{proof}

\begin{example}
\label{spaces-more-morphisms-example-universally-injective-not-radicial}
Il n'est plus vrai que « universellement injectif » soit équivalent à « radiciel ».
Par exemple, le morphisme
$$
X = [\Spec(\overline{\mathbf{Q}})/
\text{Gal}(\overline{\mathbf{Q}}/\mathbf{Q})]
\longrightarrow
S = \Spec(\mathbf{Q})
$$
d'Espaces,
exemple \ref{spaces-example-Qbar}
est universellement injectif, mais n'est pas radiciel au sens ci-dessus.
\end{example}

\noindent
Néanmoins, il arrive souvent que l'implication réciproque soit vraie.

\begin{lemma}
\label{spaces-more-morphisms-lemma-when-universally-injective-radicial}
Soit $S$ un schéma. Soit $f : X \to Y$ un morphisme universellement injectif
d'espaces algébriques au-dessus de $S$.
\begin{enumerate}
\item Si $f$ est décent, alors $f$ est radiciel.
\item Si $f$ est quasi-séparé, alors $f$ est radiciel.
\item Si $f$ est localement séparé, alors $f$ est radiciel.
\end{enumerate}
\end{lemma}

\begin{proof}
Soit $\mathcal{P}$ une propriété des morphismes d'espaces algébriques
qui est stable par changement de base et composition et qui est satisfaite pour
les immersions fermées. Supposons que $f : X \to Y$ possède $\mathcal{P}$ et
soit universellement injectif. Alors, dans la situation de la
définition \ref{spaces-more-morphisms-definition-radicial},
le morphisme $(\Spec(K) \times_Y X)_{red} \to \Spec(K)$
est universellement injectif et possède $\mathcal{P}$. Cela ramène le
problème de démontrer
$$
\mathcal{P} + \text{universellement injectif}
\Rightarrow
\text{radicial}
$$
au problème de démontrer que tout espace algébrique réduit non vide $X$
sur un corps, dont le morphisme structural $X \to \Spec(K)$ est universellement
injectif et possède $\mathcal{P}$, est représentable par le spectre d'un corps.
En effet, alors $X \to \Spec(K)$ sera un morphisme de schémas et
nous concluons par l'équivalence entre radiciel et universellement injectif pour
les morphismes de schémas, voir
Morphismes, lemme \ref{morphisms-lemma-universally-injective}.

\medskip\noindent
Démontrons (1). Supposons que $f$ soit décent et universellement injectif. D'après
Espaces décents,
les lemmes \ref{decent-spaces-lemma-base-change-relative-conditions},
\ref{decent-spaces-lemma-composition-relative-conditions} et
\ref{decent-spaces-lemma-properties-trivial-implications}
(pour voir qu'une immersion est décente), nous voyons que la discussion du
premier paragraphe s'applique.
Soit $X$ un espace algébrique réduit décent non vide
universellement injectif sur un corps $K$. En particulier, nous voyons que $|X|$
est un singleton. D'après
Espaces décents, lemme \ref{decent-spaces-lemma-when-field},
nous concluons que $X \cong \Spec(L)$ pour une extension
$K \subset L$, comme souhaité.

\medskip\noindent
Un morphisme quasi-séparé est décent, voir
Espaces décents,
lemme \ref{decent-spaces-lemma-properties-trivial-implications}.
Ainsi, (1) implique (2).

\medskip\noindent
Démontrons (3).
Rappelons que les axiomes de séparation sont stables par changement de base
et composition et que les immersions fermées sont séparées, voir
Morphismes d'espaces,
les lemmes \ref{spaces-morphisms-lemma-base-change-separated},
\ref{spaces-morphisms-lemma-composition-separated} et
\ref{spaces-morphisms-lemma-immersions-monomorphisms}.
La discussion du premier paragraphe de la démonstration s'applique donc.
Soit $X$ un espace algébrique réduit, universellement injectif et
localement séparé sur un corps $K$.
En particulier, $|X|$ est un singleton, donc $X$ est quasi-compact, voir
Propriétés des espaces, lemme \ref{spaces-properties-lemma-quasi-compact-space}.
Nous pouvons trouver un morphisme \'etale surjectif $U \to X$ avec $U$ affine, voir
Propriétés des espaces,
lemme \ref{spaces-properties-lemma-quasi-compact-affine-cover}.
Considérons le morphisme de schémas
$$
j :
U \times_X U
\longrightarrow
U \times_{\Spec(K)} U
$$
Comme $X \to \Spec(K)$ est universellement injectif, $j$ est surjectif,
et comme $X \to \Spec(K)$ est localement séparé, $j$ est une immersion.
Une immersion surjective est une immersion fermée, voir
Schémas, lemme \ref{schemes-lemma-immersion-when-closed}.
Ainsi $R = U \times_X U$ est affine comme sous-schéma fermé d'un schéma affine.
En particulier, $R$ est quasi-compact.
Il s'ensuit que $X = U/R$ est quasi-séparé, et le résultat découle de (2).
\end{proof}

\begin{remark}
\label{spaces-more-morphisms-remark-weakly-radicial}
Soit $X \to Y$ un morphisme d'espaces algébriques.
Pour certaines applications (des morphismes radiciels),
il suffit d'exiger que, pour tout
$\Spec(K) \to Y$ où $K$ est un corps,
\begin{enumerate}
\item l'espace $|\Spec(K) \times_Y X|$ soit un singleton,
\item il existe un monomorphisme
$\Spec(L) \to \Spec(K) \times_Y X$, et
\item $K \subset L$ soit purement inséparable.
\end{enumerate}
Si nécessaire, nous appellerons ultérieurement un tel morphisme {\it faiblement radiciel}.
Par exemple, si $X \to Y$ est un morphisme faiblement radiciel surjectif,
alors $X(k) \to Y(k)$ est surjectif pour tout corps algébriquement clos $k$.
Remarquons que le changement de base
$X_{\overline{\mathbf{Q}}} \to \Spec(\overline{\mathbf{Q}})$
du morphisme de
l'exemple \ref{spaces-more-morphisms-example-universally-injective-not-radicial}
est faiblement radiciel, mais pas radiciel. L'analogue du
lemme \ref{spaces-more-morphisms-lemma-when-universally-injective-radicial}
est que, si $X \to Y$ possède la propriété ($\beta$) et est universellement
injectif, alors il est faiblement radiciel (démonstration omise).
\end{remark}

\begin{lemma}
\label{spaces-more-morphisms-lemma-check-universally-injective}
Soit $S$ un schéma. Soit $f : X \to Y$ un morphisme d'espaces algébriques
au-dessus de $S$. Supposons que
\begin{enumerate}
\item $f$ soit localement de type fini,
\item pour tout morphisme \'etale $V \to Y$, l'application $|X \times_Y V| \to |V|$
soit injective.
\end{enumerate}
Alors $f$ est universellement injectif.
\end{lemma}

\begin{proof}
La question est locale pour la topologie \'etale sur $Y$, d'après
Morphismes d'espaces, lemme
\ref{spaces-morphisms-lemma-universally-injective-local}.
Nous pouvons donc supposer que $Y$ est un schéma.
Alors $Y$ est en particulier décent et, d'après Espaces décents, lemme
\ref{decent-spaces-lemma-conditions-on-point-in-fibre-and-qf}
nous voyons que $f$ est localement quasi-fini.
Soit $y \in Y$ un point et soit $X_y$ la fibre schématique.
Supposons que $X_y$ ne soit pas vide. D'après Espaces sur les corps, lemme
\ref{spaces-over-fields-lemma-locally-quasi-finite-over-field}
nous voyons que $X_y$ est un schéma localement quasi-fini sur
$\kappa(y)$. Comme $|X_y| \subset |X|$ est la fibre de $|X| \to |Y|$
au-dessus de $y$, nous voyons que $X_y$ possède un unique point $x$. Il en est de même
pour $X_y \times_{\Spec(\kappa(y))} \Spec(k)$ pour toute
extension finie séparable $k/\kappa(y)$,
car nous pouvons réaliser $k$ comme corps résiduel en un point
au-dessus de $y$ dans un schéma \'etale au-dessus de $Y$,
voir Compléments sur les morphismes, lemme
\ref{more-morphisms-lemma-realize-prescribed-residue-field-extension-etale}.
Ainsi $X_y$ est géométriquement connexe, voir
Variétés, lemme \ref{varieties-lemma-characterize-geometrically-disconnected}.
Cela implique que l'extension finie $\kappa(x)/\kappa(y)$
est purement inséparable.

\medskip\noindent
Nous concluons (dans le cas où $Y$ est un schéma)
que, pour tout $y \in Y$, la fibre $X_y$ est soit vide,
soit $(X_y)_{red} = \Spec(\kappa(x))$ avec
$\kappa(y) \subset \kappa(x)$ purement inséparable.
Ainsi $f$ est radiciel (certains détails sont omis), donc universellement injectif par
le lemme \ref{spaces-more-morphisms-lemma-radicial-implies-universally-injective}.
\end{proof}




\section{Monomorphismes}
\label{spaces-more-morphisms-section-monomorphisms}

\noindent
Cette section fait suite à
Morphismes d'espaces, section \ref{spaces-morphisms-section-monomorphisms}.
Nous voudrions savoir si tout monomorphisme d'espaces algébriques est représentable.
Si vous pouvez démontrer que c'est vrai ou fournir un
contre-exemple, veuillez envoyer un courriel à
\href{mailto:stacks.project@gmail.com}{stacks.project@gmail.com}.
Pour le moment, on connaît les cas suivants
\begin{enumerate}
\item pour les monomorphismes qui sont localement de type fini
(plus généralement, tout morphisme séparé et localement quasi-fini
est représentable par Morphismes d'espaces, lemme
\ref{spaces-morphisms-lemma-locally-quasi-finite-separated-representable}
et un monomorphisme localement de type fini est
localement quasi-fini par Morphismes d'espaces, lemme
\ref{spaces-morphisms-lemma-monomorphism-loc-finite-type-loc-quasi-finite}),
\item si le but est une union disjointe de spectres d'anneaux locaux de dimension zéro
(Espaces décents, lemme
\ref{decent-spaces-lemma-monomorphism-toward-disjoint-union-dim-0-rings}), et
\item pour les monomorphismes plats (voir ci-dessous).
\end{enumerate}

\begin{lemma}[David Rydh]
\label{spaces-more-morphisms-lemma-flat-case}
Un monomorphisme plat d'espaces algébriques est représentable par des schémas.
\end{lemma}

\begin{proof}
Soit $f : X \to Y$ un monomorphisme plat d'espaces algébriques.
Pour démontrer que $f$ est représentable, il faut montrer que
$X \times_Y V$ est un schéma pour tout schéma $V$ muni d'un morphisme vers $Y$.
Comme la propriété d'être un schéma est locale (Propriétés des espaces,
lemme \ref{spaces-properties-lemma-subscheme}), nous pouvons
supposer que $V$ est affine. Nous pouvons donc supposer que $Y = \Spec(B)$
est un schéma affine. Ensuite, nous pouvons supposer que $X$ est quasi-compact
en remplaçant $X$ par un ouvert quasi-compact. L'espace $X$ est
séparé puisque $X \to X \times_{\Spec(B)} X$ est un isomorphisme.
En appliquant Limites d'espaces, lemme \ref{spaces-limits-lemma-enough-local},
nous nous ramenons au cas où $B$ est local, où $X \to \Spec(B)$ est un
monomorphisme plat et où
il existe un point $x \in X$ s'envoyant sur le point fermé de $\Spec(B)$.
Alors $X \to \Spec(B)$ est surjectif puisque les générisations
se relèvent le long des morphismes plats d'espaces algébriques séparés, voir
Espaces décents, lemme \ref{decent-spaces-lemma-generalizations-lift-flat}.
Ainsi, $\{X \to \Spec(B)\}$ est un recouvrement fpqc.
Alors $X \to \Spec(B)$ est un morphisme qui devient un isomorphisme
après changement de base par $X \to \Spec(B)$. Il est donc un isomorphisme par
descente fpqc, voir Descente sur les espaces, lemme
\ref{spaces-descent-lemma-descending-property-isomorphism}.
\end{proof}

\noindent
Le résultat suivant est (en un certain sens) une variante du lemme ci-dessus.

\begin{lemma}
\label{spaces-more-morphisms-lemma-ui-case}
Soit $S$ un schéma. Soit $f : X \to Y$ un monomorphisme quasi-compact
d'espaces algébriques tel que, pour tout $T \to Y$, l'application
$$
\mathcal{O}_T \to f_{T,*}\mathcal{O}_{X \times_Y T}
$$
soit injective. Alors $f$ est un isomorphisme (et donc représentable par des schémas).
\end{lemma}

\begin{proof}
La question est locale pour la topologie \'etale sur $Y$; nous pouvons donc supposer que $Y = \Spec(A)$
est affine. Alors $X$ est quasi-compact et nous pouvons choisir un schéma affine
$U = \Spec(B)$ et un morphisme \'etale surjectif $U \to X$
(Propriétés des espaces, lemme
\ref{spaces-properties-lemma-quasi-compact-affine-cover}).
Remarquons que $U \times_X U = \Spec(B \otimes_A B)$. Ainsi, la catégorie des
faisceaux quasi-cohérents de $\mathcal{O}_X$-modules est équivalente à la
catégorie $DD_{B/A}$ des données de descente sur les modules pour $A \to B$.
Voir Propriétés des espaces, proposition
\ref{spaces-properties-proposition-quasi-coherent},
Descente, définition \ref{descent-definition-descent-datum-modules}, et
Descente, sous-section \ref{descent-subsection-descent-modules-morphisms}.
D'autre part,
$$
A \to B
$$
est une application d'anneaux universellement injective. En effet, étant donné un
$A$-module $M$, nous voyons que $A \oplus M \to B \otimes_A (A \oplus M)$
est injective par l'hypothèse du lemme. Ainsi,
$DD_{B/A}$ est équivalente à la catégorie des $A$-modules par
Descente, théorème \ref{descent-theorem-descent}. Ainsi, le tiré en arrière le long de
$f : X \to \Spec(A)$ détermine une équivalence de catégories des
modules quasi-cohérents. En particulier, $f^*$ est exact sur les
modules quasi-cohérents et nous voyons que $f$ est plat
(un petit détail est omis). De plus, il est clair que $f$ est surjectif
(par exemple parce que $\Spec(B) \to \Spec(A)$ est surjectif).
Ainsi, $\{X \to \Spec(A)\}$ est un recouvrement fpqc.
Alors $X \to \Spec(A)$ est un morphisme qui devient un isomorphisme
après changement de base par $X \to \Spec(A)$. Il est donc un isomorphisme par
descente fpqc, voir Descente sur les espaces, lemme
\ref{spaces-descent-lemma-descending-property-isomorphism}.
\end{proof}

\begin{lemma}
\label{spaces-more-morphisms-lemma-flat-surjective-monomorphism}
Un monomorphisme quasi-compact, plat et surjectif d'espaces algébriques
est un isomorphisme.
\end{lemma}

\begin{proof}
Un tel morphisme satisfait les hypothèses du lemme \ref{spaces-more-morphisms-lemma-ui-case}.
\end{proof}







\section{Faisceau conormal d'une immersion}
\label{spaces-more-morphisms-section-conormal-sheaf}

\noindent
Soit $S$ un schéma. Soit $i : Z \to X$ une immersion fermée d'espaces algébriques
au-dessus de $S$. Soit $\mathcal{I} \subset \mathcal{O}_X$ le
faisceau quasi-cohérent d'idéaux correspondant, voir
Morphismes d'espaces,
lemme \ref{spaces-morphisms-lemma-closed-immersion-ideals}.
Considérons la suite exacte courte
$$
0 \to \mathcal{I}^2 \to \mathcal{I} \to \mathcal{I}/\mathcal{I}^2 \to 0
$$
de faisceaux quasi-cohérents sur $X$. Comme le faisceau $\mathcal{I}/\mathcal{I}^2$
est annulé par $\mathcal{I}$, il correspond à un faisceau sur $Z$ par
Morphismes d'espaces, lemme \ref{spaces-morphisms-lemma-i-star-equivalence}.
Ce $\mathcal{O}_Z$-module quasi-cohérent est le
{\it faisceau conormal de $Z$ dans $X$} et est souvent noté
$\mathcal{I}/\mathcal{I}^2$, par l'abus de notation mentionné dans
Morphismes d'espaces,
section \ref{spaces-morphisms-section-closed-immersions-quasi-coherent}.

\medskip\noindent
Dans le cas où $i : Z \to X$ est une immersion (localement fermée), nous définissons le
faisceau conormal de $i$ comme le faisceau conormal de l'immersion fermée
$i : Z \to X \setminus \partial Z$, voir
Morphismes d'espaces, remarque \ref{spaces-morphisms-remark-immersion}.
Il est souvent noté
$\mathcal{I}/\mathcal{I}^2$ où $\mathcal{I}$ est le faisceau d'idéaux
de l'immersion fermée $i : Z \to X \setminus \partial Z$.

\begin{definition}
\label{spaces-more-morphisms-definition-conormal-sheaf}
Soit $i : Z \to X$ une immersion. Le {\it faisceau conormal
$\mathcal{C}_{Z/X}$ de $Z$ dans $X$} ou le {\it faisceau conormal de $i$}
est le $\mathcal{O}_Z$-module quasi-cohérent $\mathcal{I}/\mathcal{I}^2$
décrit ci-dessus.
\end{definition}

\noindent
Dans \cite[IV, définition 16.1.2]{EGA}, ce faisceau est noté
$\mathcal{N}_{Z/X}$. Nous ne suivrons pas cette convention, car nous
souhaitons réserver la notation $\mathcal{N}_{Z/X}$
au {\it faisceau normal de l'immersion}. Il est défini par
$$
\mathcal{N}_{Z/X} =
\SheafHom_{\mathcal{O}_Z}(\mathcal{C}_{Z/X}, \mathcal{O}_Z) =
\SheafHom_{\mathcal{O}_Z}(\mathcal{I}/\mathcal{I}^2, \mathcal{O}_Z)
$$
pourvu que le faisceau conormal soit de présentation finie (sinon le
faisceau normal peut même ne pas être quasi-cohérent). Nous reviendrons sur le
faisceau normal ultérieurement (insérer ici une référence future).

\begin{lemma}
\label{spaces-more-morphisms-lemma-etale-conormal}
Soit $S$ un schéma. Soit $i : Z \to X$ une immersion.
Soit $\varphi : U \to X$ un morphisme \'etale où $U$ est un schéma.
Posons $Z_U = U \times_X Z$, qui est un sous-schéma localement fermé de $U$.
Alors
$$
\mathcal{C}_{Z/X}|_{Z_U} = \mathcal{C}_{Z_U/U}
$$
canoniquement et fonctoriellement en $U$.
\end{lemma}

\begin{proof}
Soit $T \subset X$ un sous-espace fermé tel que $i$ définisse une
immersion fermée dans $X \setminus T$.
Soit $\mathcal{I}$ le faisceau quasi-cohérent d'idéaux sur
$X \setminus T$ définissant $Z$. Le lemme affirme simplement que
$\mathcal{I}|_{U \setminus \varphi^{-1}(T)}$ est le faisceau d'idéaux de
l'immersion $Z_U \to U \setminus \varphi^{-1}(T)$.
C'est clair d'après la construction de $\mathcal{I}$ dans
Morphismes d'espaces, lemme \ref{spaces-morphisms-lemma-closed-immersion-ideals}.
\end{proof}

\begin{lemma}
\label{spaces-more-morphisms-lemma-conormal-functorial}
Soit $S$ un schéma. Soit
$$
\xymatrix{
Z \ar[r]_i \ar[d]_f & X \ar[d]^g \\
Z' \ar[r]^{i'} & X'
}
$$
un diagramme commutatif d'espaces algébriques au-dessus de $S$.
Supposons que $i$ et $i'$ soient des immersions. Il existe une application canonique
de $\mathcal{O}_Z$-modules
$$
f^*\mathcal{C}_{Z'/X'}
\longrightarrow
\mathcal{C}_{Z/X}
$$
\end{lemma}

\begin{proof}
D'abord, trouvons des sous-espaces ouverts $U' \subset X'$ et $U \subset X$ tels que
$g(U) \subset U'$ et tels que $i(Z) \subset U$ et $i(Z') \subset U'$
soient fermés (l'existence est omise). En remplaçant $X$ par $U$ et $X'$ par
$U'$, nous pouvons supposer que $i$ et $i'$ sont des immersions fermées.
Soient $\mathcal{I}' \subset \mathcal{O}_{X'}$ et
$\mathcal{I} \subset \mathcal{O}_X$ les faisceaux quasi-cohérents d'idéaux
associés à $i'$ et $i$, voir
Morphismes d'espaces, lemme \ref{spaces-morphisms-lemma-closed-immersion-ideals}.
Considérons la composée
$$
g^{-1}\mathcal{I}' \to g^{-1}\mathcal{O}_{X'}
\xrightarrow{g^\sharp} \mathcal{O}_X \to
\mathcal{O}_X/\mathcal{I} = i_*\mathcal{O}_Z
$$
Comme $g(i(Z)) \subset Z'$, nous concluons que cette composée est nulle (voir
l'énoncé sur les factorisations dans
Morphismes d'espaces,
lemme \ref{spaces-morphisms-lemma-closed-immersion-ideals}).
Nous obtenons ainsi un diagramme commutatif
$$
\xymatrix{
0 \ar[r] &
\mathcal{I} \ar[r] &
\mathcal{O}_X \ar[r] &
i_*\mathcal{O}_Z \ar[r] &
0 \\
0 \ar[r] &
g^{-1}\mathcal{I}' \ar[r] \ar[u] &
g^{-1}\mathcal{O}_{X'} \ar[r] \ar[u] &
g^{-1}i'_*\mathcal{O}_{Z'} \ar[r] \ar[u] &
0
}
$$
La ligne inférieure est exacte puisque $g^{-1}$ est un foncteur exact.
Par exactitude, nous voyons aussi que
$(g^{-1}\mathcal{I}')^2 = g^{-1}((\mathcal{I}')^2)$.
Le diagramme induit donc une application
$g^{-1}(\mathcal{I}'/(\mathcal{I}')^2) \to \mathcal{I}/\mathcal{I}^2$.
En tirant en arrière (par exemple à l'aide de $i^{-1}$) vers $Z$, nous obtenons
$i^{-1}g^{-1}(\mathcal{I}'/(\mathcal{I}')^2) \to \mathcal{C}_{Z/X}$.
Comme $i^{-1}g^{-1} = f^{-1}(i')^{-1}$, cela donne une application
$f^{-1}\mathcal{C}_{Z'/X'} \to \mathcal{C}_{Z/X}$, qui induit
l'application cherchée.
\end{proof}

\begin{lemma}
\label{spaces-more-morphisms-lemma-conormal-functorial-more}
Soit $S$ un schéma. Le faisceau conormal de la
définition \ref{spaces-more-morphisms-definition-conormal-sheaf} et sa fonctorialité du
lemme \ref{spaces-more-morphisms-lemma-conormal-functorial}
satisfont les propriétés suivantes :
\begin{enumerate}
\item Si $Z \to X$ est une immersion de schémas au-dessus de $S$, alors le faisceau conormal
coïncide avec celui de
Morphismes, définition \ref{morphisms-definition-conormal-sheaf}.
\item Si, dans le
lemme \ref{spaces-more-morphisms-lemma-conormal-functorial},
tous les espaces sont des schémas, alors l'application
$f^*\mathcal{C}_{Z'/X'} \to \mathcal{C}_{Z/X}$ est la même
que celle construite dans
Morphismes, lemme \ref{morphisms-lemma-conormal-functorial}.
\item Étant donné un diagramme commutatif
$$
\xymatrix{
Z \ar[r]_i \ar[d]_f & X \ar[d]^g \\
Z' \ar[r]^{i'} \ar[d]_{f'} & X' \ar[d]^{g'} \\
Z'' \ar[r]^{i''} & X''
}
$$
alors l'application $(f' \circ f)^*\mathcal{C}_{Z''/X''} \to \mathcal{C}_{Z/X}$
est la même que la composée de
$f^*\mathcal{C}_{Z'/X'} \to \mathcal{C}_{Z/X}$
avec le tiré en arrière par $f$ de
$(f')^*\mathcal{C}_{Z''/X''} \to \mathcal{C}_{Z'/X'}$.
\end{enumerate}
\end{lemma}

\begin{proof}
Démonstration omise. Remarquons que la partie (1) est un cas particulier du
lemme \ref{spaces-more-morphisms-lemma-etale-conormal}.
\end{proof}

\begin{lemma}
\label{spaces-more-morphisms-lemma-conormal-functorial-flat}
Soit $S$ un schéma. Soit
$$
\xymatrix{
Z \ar[r]_i \ar[d]_f & X \ar[d]^g \\
Z' \ar[r]^{i'} & X'
}
$$
un diagramme de produit fibré d'espaces algébriques au-dessus de $S$. Supposons que
$i$ et $i'$ soient des immersions. Alors l'application canonique
$f^*\mathcal{C}_{Z'/X'} \to \mathcal{C}_{Z/X}$ du
lemme \ref{spaces-more-morphisms-lemma-conormal-functorial}
est surjective. Si $g$ est plat, c'est un isomorphisme.
\end{lemma}

\begin{proof}
Choisissons un diagramme commutatif
$$
\xymatrix{
U \ar[r] \ar[d] & X \ar[d] \\
U' \ar[r] & X'
}
$$
où $U$ et $U'$ sont des schémas et où les flèches horizontales sont surjectives
et \'etales, voir
Espaces, lemme \ref{spaces-lemma-lift-morphism-presentations}.
En utilisant alors
les lemmes \ref{spaces-more-morphisms-lemma-etale-conormal} et \ref{spaces-more-morphisms-lemma-conormal-functorial-more},
nous voyons que la question se ramène au cas d'un morphisme de schémas.
Dans le cas des schémas, il s'agit du
lemme \ref{morphisms-lemma-conormal-functorial-flat} de Morphismes.
\end{proof}

\begin{lemma}
\label{spaces-more-morphisms-lemma-transitivity-conormal}
Soit $S$ un schéma.
Soient $Z \to Y \to X$ des immersions d'espaces algébriques.
Alors il existe une suite exacte canonique
$$
i^*\mathcal{C}_{Y/X} \to
\mathcal{C}_{Z/X} \to
\mathcal{C}_{Z/Y} \to 0
$$
où les applications proviennent du
lemme \ref{spaces-more-morphisms-lemma-conormal-functorial}
et où $i : Z \to Y$ est le premier morphisme.
\end{lemma}

\begin{proof}
Soit $U$ un schéma et soit $U \to X$ un morphisme \'etale surjectif. Par
les lemmes \ref{spaces-more-morphisms-lemma-etale-conormal} et \ref{spaces-more-morphisms-lemma-conormal-functorial-more},
l'exactitude de la suite se traduit immédiatement par
l'exactitude de la suite correspondante pour les immersions de schémas
$Z \times_X U \to Y \times_X U \to U$. Le lemme découle donc du
lemme \ref{morphisms-lemma-transitivity-conormal} de Morphismes.
\end{proof}








\section{Cône normal d'une immersion}
\label{spaces-more-morphisms-section-normal-cone}

\noindent
Soit $S$ un schéma. Soit $i : Z \to X$ une immersion fermée d'espaces algébriques
au-dessus de $S$. Soit $\mathcal{I} \subset \mathcal{O}_X$ le
faisceau quasi-cohérent d'idéaux correspondant, voir
Morphismes d'espaces, lemme
\ref{spaces-morphisms-lemma-closed-immersion-ideals}.
Considérons le faisceau quasi-cohérent d'algèbres graduées de $\mathcal{O}_X$
$\bigoplus_{n \geq 0} \mathcal{I}^n/\mathcal{I}^{n + 1}$.
Comme les faisceaux $\mathcal{I}^n/\mathcal{I}^{n + 1}$
sont tous annulés par $\mathcal{I}$, cette algèbre graduée
correspond à un faisceau quasi-cohérent d'algèbres graduées de $\mathcal{O}_Z$ par
Morphismes d'espaces, lemme \ref{spaces-morphisms-lemma-i-star-equivalence}.
Cette algèbre graduée quasi-cohérente de $\mathcal{O}_Z$ est appelée
{\it algèbre conormale de $Z$ dans $X$} et est souvent simplement notée
$\bigoplus_{n \geq 0} \mathcal{I}^n/\mathcal{I}^{n + 1}$
par l'abus de notation mentionné dans
Morphismes d'espaces, section
\ref{spaces-morphisms-section-closed-immersions-quasi-coherent}.

\medskip\noindent
Dans le cas où $i : Z \to X$ est une immersion (localement fermée), nous définissons l'algèbre conormale
de $i$ comme l'algèbre conormale de l'immersion fermée
$i : Z \to X \setminus \partial Z$, voir Morphismes d'espaces, remarque
\ref{spaces-morphisms-remark-immersion}.
Elle est souvent notée
$\bigoplus_{n \geq 0} \mathcal{I}^n/\mathcal{I}^{n + 1}$
où $\mathcal{I}$ est le faisceau d'idéaux
de l'immersion fermée $i : Z \to X \setminus \partial Z$.

\begin{definition}
\label{spaces-more-morphisms-definition-conormal-algebra}
Soit $i : Z \to X$ une immersion. L'{\it algèbre conormale
$\mathcal{C}_{Z/X, *}$ de $Z$ dans $X$} ou l'{\it algèbre conormale de $i$}
est le faisceau quasi-cohérent d'algèbres graduées de $\mathcal{O}_Z$
$\bigoplus_{n \geq 0} \mathcal{I}^n/\mathcal{I}^{n + 1}$ décrit ci-dessus.
\end{definition}

\noindent
Ainsi $\mathcal{C}_{Z/X, 1} = \mathcal{C}_{Z/X}$ est le faisceau conormal
de l'immersion. De plus, $\mathcal{C}_{Z/X, 0} = \mathcal{O}_Z$ et
$\mathcal{C}_{Z/X, n}$ est un $\mathcal{O}_Z$-module quasi-cohérent
caractérisé par la propriété
\begin{equation}
\label{spaces-more-morphisms-equation-conormal-in-degree-n}
i_*\mathcal{C}_{Z/X, n} = \mathcal{I}^n/\mathcal{I}^{n + 1}
\end{equation}
où $i : Z \to X \setminus \partial Z$ et $\mathcal{I}$ est le faisceau
d'idéaux de $i$ comme ci-dessus. Enfin, remarquons qu'il existe une application canonique surjective
\begin{equation}
\label{spaces-more-morphisms-equation-conormal-algebra-quotient}
\text{Sym}^*(\mathcal{C}_{Z/X}) \longrightarrow \mathcal{C}_{Z/X, *}
\end{equation}
d'algèbres graduées quasi-cohérentes de $\mathcal{O}_Z$, qui est un isomorphisme
en degrés $0$ et $1$.

\begin{lemma}
\label{spaces-more-morphisms-lemma-etale-conormal-algebra}
Soit $S$ un schéma. Soit $i : Z \to X$ une immersion d'espaces algébriques
au-dessus de $S$. Soit $\varphi : U \to X$ un morphisme \'etale où $U$ est un
schéma. Posons $Z_U = U \times_X Z$, qui est un sous-schéma localement fermé de $U$.
Alors
$$
\mathcal{C}_{Z/X, *}|_{Z_U} = \mathcal{C}_{Z_U/U, *}
$$
canoniquement et fonctoriellement en $U$.
\end{lemma}

\begin{proof}
Soit $T \subset X$ un sous-espace fermé tel que $i$ définisse une immersion
fermée dans $X \setminus T$. Soit $\mathcal{I}$ le faisceau quasi-cohérent
d'idéaux sur $X \setminus T$ définissant $Z$. Le lemme découle
du fait que
$\mathcal{I}|_{U \setminus \varphi^{-1}(T)}$ est le faisceau d'idéaux de
l'immersion $Z_U \to U \setminus \varphi^{-1}(T)$.
C'est clair d'après la construction de $\mathcal{I}$ dans
Morphismes d'espaces, lemme
\ref{spaces-morphisms-lemma-closed-immersion-ideals}.
\end{proof}

\begin{lemma}
\label{spaces-more-morphisms-lemma-conormal-algebra-functorial}
Soit $S$ un schéma. Soit
$$
\xymatrix{
Z \ar[r]_i \ar[d]_f & X \ar[d]^g \\
Z' \ar[r]^{i'} & X'
}
$$
un diagramme commutatif d'espaces algébriques au-dessus de $S$.
Supposons que $i$ et $i'$ soient des immersions. Il existe une application canonique
d'algèbres graduées de $\mathcal{O}_Z$
$$
f^*\mathcal{C}_{Z'/X', *}
\longrightarrow
\mathcal{C}_{Z/X, *}
$$
\end{lemma}

\begin{proof}
Commençons par trouver des sous-espaces ouverts $U' \subset X'$ et $U \subset X$ tels que
$g(U) \subset U'$ et tels que $i(Z) \subset U$ et $i(Z') \subset U'$
soient fermés (l'existence est omise). En remplaçant $X$ par $U$ et $X'$ par
$U'$, nous pouvons supposer que $i$ et $i'$ sont des immersions fermées.
Soient $\mathcal{I}' \subset \mathcal{O}_{X'}$ et
$\mathcal{I} \subset \mathcal{O}_X$ les faisceaux quasi-cohérents d'idéaux
associés à $i'$ et $i$, voir
Morphismes d'espaces, lemme \ref{spaces-morphisms-lemma-closed-immersion-ideals}.
Considérons la composée
$$
g^{-1}\mathcal{I}' \to g^{-1}\mathcal{O}_{X'}
\xrightarrow{g^\sharp} \mathcal{O}_X \to
\mathcal{O}_X/\mathcal{I} = i_*\mathcal{O}_Z
$$
Comme $g(i(Z)) \subset Z'$, nous concluons que cette composée est nulle (voir
l'énoncé sur les factorisations dans
Morphismes d'espaces,
lemme \ref{spaces-morphisms-lemma-closed-immersion-ideals}).
Nous obtenons ainsi un diagramme commutatif
$$
\xymatrix{
0 \ar[r] &
\mathcal{I} \ar[r] &
\mathcal{O}_X \ar[r] &
i_*\mathcal{O}_Z \ar[r] &
0 \\
0 \ar[r] &
g^{-1}\mathcal{I}' \ar[r] \ar[u] &
g^{-1}\mathcal{O}_{X'} \ar[r] \ar[u] &
g^{-1}i'_*\mathcal{O}_{Z'} \ar[r] \ar[u] &
0
}
$$
La ligne inférieure est exacte puisque $g^{-1}$ est un foncteur exact.
Par exactitude, nous voyons aussi que
$(g^{-1}\mathcal{I}')^n = g^{-1}((\mathcal{I}')^n)$ pour tout $n \geq 1$.
Le diagramme induit donc une application
$g^{-1}((\mathcal{I}')^n/(\mathcal{I}')^{n + 1}) \to
\mathcal{I}^n/\mathcal{I}^{n + 1}$.
En tirant en arrière (par exemple à l'aide de $i^{-1}$) vers $Z$, nous obtenons
$i^{-1}g^{-1}((\mathcal{I}')^n/(\mathcal{I}')^{n + 1}) \to
\mathcal{C}_{Z/X, n}$.
Comme $i^{-1}g^{-1} = f^{-1}(i')^{-1}$, cela donne des applications
$f^{-1}\mathcal{C}_{Z'/X', n} \to \mathcal{C}_{Z/X, n}$, qui induisent
l'application cherchée.
\end{proof}

\begin{lemma}
\label{spaces-more-morphisms-lemma-conormal-algebra-functorial-flat}
Soit $S$ un schéma. Soit
$$
\xymatrix{
Z \ar[r]_i \ar[d]_f & X \ar[d]^g \\
Z' \ar[r]^{i'} & X'
}
$$
un carré cartésien d'espaces algébriques au-dessus de $S$ avec
$i$ et $i'$ immersions. Alors l'application canonique
$f^*\mathcal{C}_{Z'/X', *} \to \mathcal{C}_{Z/X, *}$ du
lemme \ref{spaces-more-morphisms-lemma-conormal-algebra-functorial}
est surjective. Si $g$ est plat, c'est un isomorphisme.
\end{lemma}

\begin{proof}
Nous pouvons vérifier l'assertion après localisation \'etale de $X'$.
Dans ce cas, nous pouvons supposer que $X' \to X$ est un morphisme de schémas,
d'où $Z$ et $Z'$ sont des schémas, et le résultat découle
du cas des schémas, voir
Diviseurs, lemme \ref{divisors-lemma-conormal-algebra-functorial-flat}.
\end{proof}

\noindent
Nous utilisons les mêmes conventions pour les cônes et les fibrés vectoriels sur
les espaces algébriques que pour les schémas (où nous utilisons
les conventions de l'EGA), voir
Constructions, sections \ref{constructions-section-cone} et
\ref{constructions-section-vector-bundle}.
En particulier, un fibré vectoriel est un objet très général
(et n'est pas localement isomorphe à un fibré en espaces affines).

\begin{definition}
\label{spaces-more-morphisms-definition-normal-cone}
Soit $S$ un schéma. Soit $i : Z \to X$ une immersion d'espaces algébriques
au-dessus de $S$. Le {\it cône normal $C_ZX$} de $Z$ dans $X$ est
$$
C_ZX = \underline{\Spec}_Z(\mathcal{C}_{Z/X, *})
$$
voir Morphismes d'espaces,
définition \ref{spaces-morphisms-definition-relative-spec}. Le
{\it fibré normal} de $Z$ dans $X$ est le fibré vectoriel
$$
N_ZX = \underline{\Spec}_Z(\text{Sym}(\mathcal{C}_{Z/X}))
$$
\end{definition}

\noindent
Ainsi $C_ZX \to Z$ est un cône au-dessus de $Z$ et $N_ZX \to Z$ est un fibré vectoriel
au-dessus de $Z$. De plus, la surjection canonique
(\ref{spaces-more-morphisms-equation-conormal-algebra-quotient}) d'algèbres graduées
définit une immersion fermée canonique
\begin{equation}
\label{spaces-more-morphisms-equation-normal-cone-in-normal-bundle}
C_ZX \longrightarrow N_ZX
\end{equation}
de cônes au-dessus de $Z$.










\section{Faisceau des différentielles d'un morphisme}
\label{spaces-more-morphisms-section-sheaf-differentials}

\noindent
Nous invitons le lecteur à consulter la section correspondante
dans le chapitre d'algèbre commutative
(Algèbre, section \ref{algebra-section-differentials}),
la section correspondante dans le chapitre sur les morphismes de schémas
(Morphismes, section \ref{morphisms-section-sheaf-differentials})
ainsi que
Modules sur les sites, section \ref{sites-modules-section-differentials}.
Nous montrons d'abord que la notion de faisceau des différentielles pour un
morphisme de schémas coïncide avec la notion correspondante pour les
petits sites \'etale (annelés).

\medskip\noindent
Pour énoncer clairement le lemme suivant, nous revenons temporairement à la notation
où $\mathcal{F}^a$ désigne le faisceau de $\mathcal{O}_{X_\etale}$-modules
associé à un $\mathcal{O}_X$-module quasi-cohérent $\mathcal{F}$
sur le schéma $X$, voir
Descente, définition \ref{descent-definition-structure-sheaf}.

\begin{lemma}
\label{spaces-more-morphisms-lemma-match-modules-differentials}
Soit $f : X \to Y$ un morphisme de schémas. Soit
$f_{small} : X_\etale \to Y_\etale$ le morphisme associé
des petits sites \'etale, voir
Descente, remarque \ref{descent-remark-change-topologies-ringed}.
Alors il existe un isomorphisme canonique
$$
(\Omega_{X/Y})^a = \Omega_{X_\etale/Y_\etale}
$$
compatible avec les dérivations universelles. Ici, le premier module
est le faisceau sur $X_\etale$ associé
au $\mathcal{O}_X$-module quasi-cohérent $\Omega_{X/Y}$, voir
Morphismes, définition \ref{morphisms-definition-sheaf-differentials},
et le second module est celui de
Modules sur les sites,
définition \ref{sites-modules-definition-module-differentials}.
\end{lemma}

\begin{proof}
Soit $h : U \to X$ un morphisme \'etale. Dans ce cas, l'application naturelle
$h^*\Omega_{X/Y} \to \Omega_{U/Y}$ est un isomorphisme, voir
Compléments sur les morphismes,
lemme \ref{more-morphisms-lemma-sheaf-differentials-etale-localization}.
Cela signifie qu'il existe une $\mathcal{O}_{Y_\etale}$-dérivation naturelle
$$
\text{d}^a : \mathcal{O}_{X_\etale} \longrightarrow (\Omega_{X/Y})^a
$$
puisque nous venons de voir que la valeur de $(\Omega_{X/Y})^a$ sur tout objet
$U$ de $X_\etale$ s'identifie canoniquement à
$\Gamma(U, \Omega_{U/Y})$. Par la propriété universelle de
$\text{d}_{X/Y} :
\mathcal{O}_{X_\etale}
\to
\Omega_{X_\etale/Y_\etale}$
il existe une unique application $\mathcal{O}_{X_\etale}$-linéaire
$c : \Omega_{X_\etale/Y_\etale} \to (\Omega_{X/Y})^a$
telle que
$\text{d}^a = c \circ \text{d}_{X/Y}$.

\medskip\noindent
Réciproquement, supposons que $\mathcal{F}$ soit un
$\mathcal{O}_{X_\etale}$-module
et que $D : \mathcal{O}_{X_\etale} \to \mathcal{F}$ soit une
$\mathcal{O}_{Y_\etale}$-dérivation. Nous pouvons simplement restreindre
$D$ au petit site de Zariski $X_{Zar}$ de $X$. Comme les faisceaux sur $X_{Zar}$
coïncident avec les faisceaux sur $X$, voir
Descente, remarque \ref{descent-remark-Zariski-site-space},
nous voyons que $D|_{X_{Zar}} : \mathcal{O}_X \to \mathcal{F}|_{X_{Zar}}$
est simplement une dérivation $Y$ « usuelle ». Nous obtenons donc une application
$\psi : \Omega_{X/Y} \longrightarrow \mathcal{F}|_{X_{Zar}}$
telle que $D|_{X_{Zar}} = \psi \circ \text{d}$. En particulier, si nous
appliquons ceci avec $\mathcal{F} = \Omega_{X_\etale/Y_\etale}$
nous obtenons une application
$$
c' :
\Omega_{X/Y}
\longrightarrow
\Omega_{X_\etale/Y_\etale}|_{X_{Zar}}
$$
Considérons le morphisme de sites annelés
$\text{id}_{small, \etale, Zar} : X_\etale \to X_{Zar}$
décrit dans
Descente, remarque \ref{descent-remark-change-topologies-ringed} et
lemme \ref{descent-lemma-compare-sites}.
Comme le foncteur de restriction $\mathcal{F} \mapsto \mathcal{F}|_{X_{Zar}}$
est égal à $\text{id}_{small, \etale, Zar, *}$, comme
$\text{id}_{small, \etale, Zar}^*$ est adjoint à gauche de
$\text{id}_{small, \etale, Zar, *}$ et comme
$(\Omega_{X/Y})^a = \text{id}_{small, \etale, Zar}^*\Omega_{X/Y}$,
nous voyons que $c'$ est adjoint à une application
$$
c'' :
(\Omega_{X/Y})^a
\longrightarrow
\Omega_{X_\etale/Y_\etale}.
$$
Nous affirmons que $c''$ et $c'$ sont inverses l'un de l'autre.
Cette assertion achève la démonstration du lemme.
Pour le voir, il suffit de montrer que $c''(\text{d}(f)) = \text{d}_{X/Y}(f)$
et $c(\text{d}_{X/Y}(f)) = \text{d}(f)$ lorsque $f$ est une section locale de
$\mathcal{O}_X$ sur un ouvert de $X$. Nous omettons la vérification.
\end{proof}

\noindent
Cette observation ouvre la voie à la définition suivante. Pour une autre possibilité, voir
la remarque \ref{spaces-more-morphisms-remark-alternative}.

\begin{definition}
\label{spaces-more-morphisms-definition-sheaf-differentials}
Soit $S$ un schéma. Soit $f : X \to Y$ un morphisme d'espaces algébriques
au-dessus de $S$. Le {\it faisceau des différentielles $\Omega_{X/Y}$ de $X$ sur $Y$}
est le faisceau des différentielles
(Modules sur les sites,
définition \ref{sites-modules-definition-sheaf-differentials})
pour le morphisme de topoi annelés
$$
(f_{small}, f^\sharp) :
(X_\etale, \mathcal{O}_X)
\to
(Y_\etale, \mathcal{O}_Y)
$$
de
Propriétés des espaces,
lemme \ref{spaces-properties-lemma-morphism-ringed-topoi}.
La {\it dérivation universelle $Y$} sera notée
$\text{d}_{X/Y} : \mathcal{O}_X \to \Omega_{X/Y}$.
\end{definition}

\noindent
D'après
le lemme \ref{spaces-more-morphisms-lemma-match-modules-differentials},
cela n'entre pas en conflit avec la notion déjà existante
dans le cas où $X$ et $Y$ sont représentables. Désormais, si $X$ et $Y$
sont représentables, nous ne faisons plus de distinction entre le faisceau des différentielles
défini ci-dessus et celui défini dans
Morphismes, définition \ref{morphisms-definition-sheaf-differentials}.
Nous voulons relier ceci aux modules usuels des différentielles pour les
morphismes de schémas. Voici le lemme clé.

\begin{lemma}
\label{spaces-more-morphisms-lemma-localize-differentials}
Soit $S$ un schéma. Soit $f : X \to Y$ un morphisme d'espaces algébriques
au-dessus de $S$. Considérons un diagramme commutatif quelconque
$$
\xymatrix{
U \ar[d]_a \ar[r]_\psi & V \ar[d]^b \\
X \ar[r]^f & Y
}
$$
où les flèches verticales sont des morphismes \'etales d'espaces algébriques. Alors
$$
\Omega_{X/Y}|_{U_\etale} = \Omega_{U/V}
$$
En particulier, si $U$ et $V$ sont des schémas, ceci est égal au
faisceau usuel des différentielles du morphisme de schémas $U \to V$.
\end{lemma}

\begin{proof}
D'après
Propriétés des espaces, lemme \ref{spaces-properties-lemma-etale-morphism-topoi}
et l'Équation (\ref{spaces-properties-equation-restrict}),
on peut considérer la restriction d'un faisceau sur $X_\etale$ à
$U_\etale$ comme l'image inverse par $a_{small}$. De même pour $b$. Par
Modules sur les sites, Lemme \ref{sites-modules-lemma-localize-differentials},
on a
$$
\Omega_{X/Y}|_{U_\etale} =
\Omega_{\mathcal{O}_{U_\etale}/
a_{small}^{-1}f_{small}^{-1}\mathcal{O}_{Y_\etale}}
$$
Comme $a_{small}^{-1}f_{small}^{-1}\mathcal{O}_{Y_\etale}
= \psi_{small}^{-1}b_{small}^{-1}\mathcal{O}_{Y_\etale}
= \psi_{small}^{-1}\mathcal{O}_{V_\etale}$, le lemme en résulte.
\end{proof}

\begin{lemma}
\label{spaces-more-morphisms-lemma-module-differentials-quasi-coherent}
Soit $S$ un schéma. Soit $f : X \to Y$ un morphisme d'espaces algébriques
sur $S$. Alors $\Omega_{X/Y}$ est un $\mathcal{O}_X$-module quasi-cohérent.
\end{lemma}

\begin{proof}
Choisissons un diagramme comme dans le
Lemme \ref{spaces-more-morphisms-lemma-localize-differentials},
avec $a$ et $b$ surjectifs et $U$ et $V$ des schémas.
On voit alors que $\Omega_{X/Y}|_U = \Omega_{U/V}$, qui est
quasi-cohérent (par exemple, d'après
Morphismes, Lemme \ref{morphisms-lemma-differentials-diagonal}).
Nous concluons donc que $\Omega_{X/Y}$ est quasi-cohérent par
Propriétés des espaces,
Lemme \ref{spaces-properties-lemma-characterize-quasi-coherent}.
\end{proof}

\begin{remark}
\label{spaces-more-morphisms-remark-alternative}
Maintenant que nous savons que $\Omega_{X/Y}$ est quasi-cohérent, nous pouvons tenter
de le construire d'une autre manière. On peut par exemple utiliser le résultat de
Propriétés des espaces,
section \ref{spaces-properties-section-quasi-coherent-presentation},
pour construire le faisceau des différentielles par recollement.
Par exemple, si $Y$ est un schéma et si $U \to X$ est un morphisme étale surjectif
d'un schéma vers $X$, alors $\Omega_{U/Y}$ est
un $\mathcal{O}_U$-module quasi-cohérent et, comme $s, t : R \to U$
sont étales, on obtient un isomorphisme
$$
\alpha : s^*\Omega_{U/Y} \to \Omega_{R/Y} \to t^*\Omega_{U/Y}
$$
en utilisant
Morphismes, Lemme \ref{morphisms-lemma-triangle-differentials-smooth}.
On vérifie qu'il satisfait la condition de cocycle, et c'est terminé.
Si $Y$ n'est pas un schéma, on définit $\Omega_{U/Y}$ comme le conoyau
du morphisme $(U \to Y)^*\Omega_{Y/S} \to \Omega_{U/S}$, puis on procède comme
précédemment. Cette procédure en deux étapes est quelque peu disgracieuse. Une autre possibilité
consiste à recoller les faisceaux $\Omega_{U/V}$ pour tout diagramme comme dans le
Lemme \ref{spaces-more-morphisms-lemma-localize-differentials},
mais cela n'est pas non plus très élégant. Les deux méthodes fonctionnent néanmoins et
donnent une construction légèrement plus élémentaire du faisceau des
différentielles.
\end{remark}

\begin{lemma}
\label{spaces-more-morphisms-lemma-functoriality-differentials}
Soit $S$ un schéma. Soit
$$
\xymatrix{
X' \ar[d] \ar[r]_f & X \ar[d] \\
Y' \ar[r] & Y
}
$$
un diagramme commutatif d'espaces algébriques. Le morphisme
$f^\sharp : \mathcal{O}_X \to f_*\mathcal{O}_{X'}$, composé avec le morphisme
$f_*\text{d}_{X'/Y'} : f_*\mathcal{O}_{X'} \to f_*\Omega_{X'/Y'}$, est une
$Y$-dérivation. On obtient donc un morphisme canonique de $\mathcal{O}_X$-modules
$\Omega_{X/Y} \to f_*\Omega_{X'/Y'}$ et, par
l'adjonction entre $f_*$ et $f^*$, un
homomorphisme canonique de $\mathcal{O}_{X'}$-modules
$$
c_f : f^*\Omega_{X/Y} \longrightarrow \Omega_{X'/Y'}.
$$
Il est uniquement caractérisé par la propriété suivante :
$f^*\text{d}_{X/Y}(t)$ est envoyé sur $\text{d}_{X'/Y'}(f^* t)$
pour toute section locale $t$ de $\mathcal{O}_X$.
\end{lemma}

\begin{proof}
Il s'agit d'un cas particulier de
Modules sur les sites, Lemme
\ref{sites-modules-lemma-functoriality-differentials-ringed-topoi}.
\end{proof}

\begin{lemma}
\label{spaces-more-morphisms-lemma-check-functoriality-differentials}
Soit $S$ un schéma. Soit
$$
\xymatrix{
X'' \ar[d] \ar[r]_g & X' \ar[d] \ar[r]_f & X \ar[d] \\
Y'' \ar[r] & Y' \ar[r] & Y
}
$$
un diagramme commutatif d'espaces algébriques sur $S$. Alors on a
$$
c_{f \circ g} = c_g \circ g^* c_f
$$
comme morphismes $(f \circ g)^*\Omega_{X/Y} \to \Omega_{X''/Y''}$.
\end{lemma}

\begin{proof}
Démonstration omise. Indication : utiliser la caractérisation de $c_f, c_g, c_{f \circ g}$
par l'effet de ces morphismes sur les sections locales.
\end{proof}

\begin{lemma}
\label{spaces-more-morphisms-lemma-triangle-differentials}
Soit $S$ un schéma.
Soient $f : X \to Y$ et $g : Y \to B$ des morphismes d'espaces algébriques sur $S$.
Alors il existe une suite exacte canonique
$$
f^*\Omega_{Y/B} \to \Omega_{X/B} \to \Omega_{X/Y} \to 0
$$
où les morphismes proviennent d'applications du
Lemme \ref{spaces-more-morphisms-lemma-functoriality-differentials}.
\end{lemma}

\begin{proof}
Cela découle de la version de ce résultat pour les schémas, voir
Morphismes, Lemme \ref{morphisms-lemma-triangle-differentials},
par localisation étale, voir le
Lemme \ref{spaces-more-morphisms-lemma-localize-differentials}.
\end{proof}

\begin{lemma}
\label{spaces-more-morphisms-lemma-immersion-differentials}
Soit $S$ un schéma. Si $X \to Y$ est une immersion
d'espaces algébriques sur $S$, alors $\Omega_{X/S}$ est nul.
\end{lemma}

\begin{proof}
Cela découle de la version de ce résultat pour les schémas, voir
Morphismes, Lemme \ref{morphisms-lemma-immersion-differentials},
par localisation étale, voir le
Lemme \ref{spaces-more-morphisms-lemma-localize-differentials}.
\end{proof}

\begin{lemma}
\label{spaces-more-morphisms-lemma-differentials-relative-immersion}
Soit $S$ un schéma. Soit $B$ un espace algébrique sur $S$.
Soit $i : Z \to X$ une immersion d'espaces algébriques sur $B$.
Il existe une suite exacte canonique
$$
\mathcal{C}_{Z/X} \to i^*\Omega_{X/B} \to \Omega_{Z/B} \to 0
$$
où la première flèche est induite par $\text{d}_{X/B}$
et où la seconde flèche provient du
Lemme \ref{spaces-more-morphisms-lemma-functoriality-differentials}.
\end{lemma}

\begin{proof}
Il s'agit de la version pour les espaces algébriques de
Morphismes, Lemme \ref{morphisms-lemma-differentials-relative-immersion},
et elle découlera de ce lemme par
localisation étale, voir les
Lemmes \ref{spaces-more-morphisms-lemma-localize-differentials} et
\ref{spaces-more-morphisms-lemma-etale-conormal}.
Il faut toutefois s'assurer que l'on peut définir globalement la première flèche.
Expliquons donc ici le sens de « induite par $\text{d}_{X/B}$ ».
On peut en effet supposer que $i$ est une immersion fermée après avoir remplacé $X$
par un sous-espace ouvert. Soit $\mathcal{I} \subset \mathcal{O}_X$
le faisceau quasi-cohérent d'idéaux correspondant à $Z \subset X$.
Alors $\text{d}_{X/S} : \mathcal{I} \to \Omega_{X/S}$
envoie le sous-faisceau $\mathcal{I}^2 \subset \mathcal{I}$ dans
$\mathcal{I}\Omega_{X/S}$. Il induit donc un morphisme
$\mathcal{I}/\mathcal{I}^2 \to \Omega_{X/S}/\mathcal{I}\Omega_{X/S}$
qui est $\mathcal{O}_X/\mathcal{I}$-linéaire.
D'après
Morphismes d'espaces, Lemme \ref{spaces-morphisms-lemma-i-star-equivalence},
cela correspond au morphisme $\mathcal{C}_{Z/X} \to i^*\Omega_{X/S}$ voulu.
\end{proof}

\begin{lemma}
\label{spaces-more-morphisms-lemma-differentials-relative-immersion-section}
Soit $S$ un schéma. Soit $B$ un espace algébrique sur $S$.
Soit $i : Z \to X$ une immersion d'espaces algébriques sur $B$, et
supposons que $i$ admette (localement pour la topologie étale) un inverse à gauche. Alors la
suite canonique
$$
0 \to \mathcal{C}_{Z/X} \to i^*\Omega_{X/B} \to \Omega_{Z/B} \to 0
$$
du
Lemme \ref{spaces-more-morphisms-lemma-differentials-relative-immersion}
est (localement pour la topologie étale) exacte scindée.
\end{lemma}

\begin{proof}
Précision : nous affirmons que, si $g : X \to Z$ est un inverse à gauche de $i$
sur $B$, alors $i^*c_g$ est un inverse à droite du morphisme
$i^*\Omega_{X/B} \to \Omega_{Z/B}$.
Ceci étant dit, le résultat découle du résultat correspondant pour les
morphismes de schémas par localisation étale, voir les
Lemmes \ref{spaces-more-morphisms-lemma-localize-differentials} et
\ref{spaces-more-morphisms-lemma-etale-conormal}.
\end{proof}

\begin{lemma}
\label{spaces-more-morphisms-lemma-base-change-differentials}
Soit $S$ un schéma.
Soit $X \to Y$ un morphisme d'espaces algébriques sur $S$.
Soit $g : Y' \to Y$ un morphisme d'espaces algébriques sur $S$.
Soit $X' = X_{Y'}$ le changement de base de $X$.
Notons $g' : X' \to X$ la projection.
Alors le morphisme
$$
(g')^*\Omega_{X/Y} \to \Omega_{X'/Y'}
$$
du
Lemme \ref{spaces-more-morphisms-lemma-functoriality-differentials}
est un isomorphisme.
\end{lemma}

\begin{proof}
Cela découle de la version pour les schémas, voir
Morphismes, Lemme \ref{morphisms-lemma-base-change-differentials},
et de la localisation étale, voir le
Lemme \ref{spaces-more-morphisms-lemma-localize-differentials}.
\end{proof}

\begin{lemma}
\label{spaces-more-morphisms-lemma-differential-product}
Soit $S$ un schéma.
Soient $f : X \to B$ et $g : Y \to B$ des morphismes d'espaces algébriques
sur $S$ ayant le même but.
Soient $p : X \times_B Y \to X$ et $q : X \times_B Y \to Y$ les
morphismes de projection. Les morphismes du
Lemme \ref{spaces-more-morphisms-lemma-functoriality-differentials}
$$
p^*\Omega_{X/B} \oplus q^*\Omega_{Y/B}
\longrightarrow
\Omega_{X \times_B Y/B}
$$
donnent un isomorphisme.
\end{lemma}

\begin{proof}
Cela découle de la version pour les schémas, voir
Morphismes, Lemme \ref{morphisms-lemma-differential-product},
et de la localisation étale, voir le
Lemme \ref{spaces-more-morphisms-lemma-localize-differentials}.
\end{proof}

\begin{lemma}
\label{spaces-more-morphisms-lemma-finite-type-differentials}
Soit $S$ un schéma.
Soit $f : X \to Y$ un morphisme d'espaces algébriques sur $S$.
Si $f$ est localement de type fini, alors $\Omega_{X/Y}$ est
un $\mathcal{O}_X$-module de type fini.
\end{lemma}

\begin{proof}
Cela découle de la version pour les schémas, voir
Morphismes, Lemme \ref{morphisms-lemma-finite-type-differentials},
et de la localisation étale, voir le
Lemme \ref{spaces-more-morphisms-lemma-localize-differentials}.
\end{proof}

\begin{lemma}
\label{spaces-more-morphisms-lemma-finite-presentation-differentials}
Soit $S$ un schéma.
Soit $f : X \to Y$ un morphisme d'espaces algébriques sur $S$.
Si $f$ est localement de présentation finie, alors $\Omega_{X/Y}$ est
un $\mathcal{O}_X$-module de présentation finie.
\end{lemma}

\begin{proof}
Cela découle de la version pour les schémas, voir
Morphismes, Lemme \ref{morphisms-lemma-finite-presentation-differentials},
et de la localisation étale, voir le
Lemme \ref{spaces-more-morphisms-lemma-localize-differentials}.
\end{proof}

\begin{lemma}
\label{spaces-more-morphisms-lemma-smooth-omega-finite-locally-free}
Soit $S$ un schéma.
Soit $f : X \to Y$ un morphisme lisse d'espaces algébriques sur $S$.
Alors le module des différentielles $\Omega_{X/Y}$
est localement libre de type fini.
\end{lemma}

\begin{proof}
L'énoncé est local pour la topologie étale sur $X$ et $Y$ par le
Lemme \ref{spaces-more-morphisms-lemma-localize-differentials}.
Il découle donc du cas des schémas, voir
Morphismes, Lemme \ref{morphisms-lemma-smooth-omega-finite-locally-free}.
\end{proof}













\section{Invariance topologique du site étale}
\label{spaces-more-morphisms-section-topological-invariance}

\noindent
Nous montrons que le site $X_{spaces, \etale}$ est un « invariant
topologique ». Il s'ensuit que $X_\etale$, qui
est formé des objets représentables de $X_{spaces, \etale}$,
est lui aussi un invariant topologique, voir le Lemme \ref{spaces-more-morphisms-lemma-topological-invariance}.

\begin{theorem}
\label{spaces-more-morphisms-theorem-topological-invariance}
Soit $S$ un schéma.
Soit $f : X \to Y$ un morphisme d'espaces algébriques sur $S$.
Supposons que $f$ soit entier, universellement injectif et surjectif.
Le foncteur
$$
V \longmapsto V_X = X \times_Y V
$$
définit une équivalence de catégories
$Y_{spaces, \etale} \to X_{spaces, \etale}$.
\end{theorem}

\begin{proof}
Le morphisme $f$ est représentable et constitue un homéomorphisme universel, voir
Morphismes d'espaces,
section \ref{spaces-morphisms-section-universal-homeomorphisms}.

\medskip\noindent
Nous montrons d'abord que le foncteur est fidèle.
Supposons que $V', V$ soient des objets de $Y_{spaces, \etale}$ et
que $a, b : V' \to V$ soient des morphismes distincts sur $Y$.
Comme $V', V$ sont étales sur $Y$, l'égalisateur
$$
E =  V' \times_{(a, b), V \times_Y V, \Delta_{V/Y}} V
$$
de $a, b$ est lui aussi étale sur $Y$. Ainsi, $E \to V'$ est un monomorphisme étale
(c'est-à-dire une immersion ouverte) qui est un isomorphisme si et seulement s'il est
surjectif. Comme $X \to Y$ est un homéomorphisme universel, nous voyons que c'est
le cas si et seulement si $E_X = V'_X$, c'est-à-dire si et seulement si $a_X = b_X$.

\medskip\noindent
Nous montrons ensuite que le foncteur est pleinement fidèle.
Supposons que $V', V$ soient des objets de $Y_{spaces, \etale}$ et
que $c : V'_X \to V_X$ soit un morphisme sur $X$. Nous voulons construire
un morphisme $a : V' \to V$ sur $Y$ tel que $a_X = c$.
Soit $a' : V'' \to V'$ un morphisme étale surjectif tel que $V''$ soit
un espace algébrique séparé. Si l'on peut construire un morphisme
$a'' : V'' \to V$ tel que $a''_X = c \circ a'_X$, alors les deux composées
$$
V'' \times_{V'} V'' \xrightarrow{\text{pr}_i} V'' \xrightarrow{a''} V
$$
sont égales par la fidélité du foncteur établie au premier
paragraphe. Ainsi, $a''$ se factorise par un unique morphisme
$a : V' \to V$, puisque $V'$ est (en tant que faisceau) le quotient de $V''$ par
la relation d'équivalence $V'' \times_{V'} V''$. Nous pouvons donc supposer que
$V'$ est séparé. Dans ce cas, le graphe
$$
\Gamma_c \subset (V' \times_Y V)_X
$$
est ouvert et fermé (détails omis). Comme $X \to Y$ est un homéomorphisme
universel, il existe un sous-espace ouvert et fermé
$\Gamma \subset V' \times_Y V$ tel que $\Gamma_X = \Gamma_c$.
La projection $\Gamma \to V'$ est un morphisme étale dont le changement
de base à $X$ est un isomorphisme. Ainsi, $\Gamma \to V'$ est étale,
universellement injectif et surjectif, donc est un isomorphisme d'après
Morphismes d'espaces,
Lemme \ref{spaces-morphisms-lemma-etale-universally-injective-open}.
Ainsi, $\Gamma$ est le graphe d'un morphisme $a : V' \to V$, comme souhaité.

\medskip\noindent
Enfin, nous montrons que le foncteur est essentiellement surjectif.
Supposons que $U$ soit un objet de $X_{spaces, \etale}$.
Il faut trouver un objet $V$ de $Y_{spaces, \etale}$
tel que $V_X \cong U$. Soit $U' \to U$ un morphisme étale surjectif
tel que $U' \cong V'_X$ et $U' \times_U U' \cong V''_X$
pour certains objets $V'', V'$ de $Y_{spaces, \etale}$.
Alors, par la pleine fidélité du foncteur, on obtient des morphismes
$s, t : V'' \to V'$ tels que $t_X = \text{pr}_0$ et $s_X = \text{pr}_1$
comme morphismes $U' \times_U U' \to U'$. Puisque
$(\text{pr}_0, \text{pr}_1) : U' \times_U U' \to U' \times_S U'$
est une relation d'équivalence étale et que $U' \to V'$ et
$U' \times_U U' \to V''$ sont universellement injectifs et surjectifs,
on en déduit que
$(t, s) : V'' \to V' \times_S V'$ est une relation d'équivalence étale.
Alors le quotient $V = V'/V''$ (voir
Espaces, Théorème \ref{spaces-theorem-presentation})
est l'espace algébrique $V$ sur $Y$. Il existe un morphisme
$V' \to V$ tel que $V'' = V' \times_V V'$. On obtient ainsi un morphisme
$V \to Y$ (voir
Descente sur les espaces, Lemme
\ref{spaces-descent-lemma-fpqc-universal-effective-epimorphisms}).
Après changement de base à $X$, on voit que l'on a un morphisme $U' \to V_X$
et un isomorphisme compatible $U' \times_{V_X} U' = U' \times_U U'$, ce qui
implique que $V_X \cong U$ (en appliquant encore une fois le lemme qui vient d'être cité).

\medskip\noindent
Choisissons un schéma $W$ et un morphisme étale surjectif $W \to Y$.
Choisissons un schéma $U'$ et un morphisme étale surjectif $U' \to U \times_X W_X$.
Notons que $U'$ et $U' \times_U U'$ sont des schémas étales sur $X$ dont le
morphisme structural vers $X$ se factorise par le schéma $W_X$.
Ainsi, d'après
Cohomologie étale,
Théorème \ref{etale-cohomology-theorem-topological-invariance},
il existe des schémas $V', V''$ étales sur $W$ dont les changements de base à
$W_X$ sont respectivement isomorphes à $U'$ et $U' \times_U U'$.
Cela achève la démonstration.
\end{proof}

\begin{lemma}
\label{spaces-more-morphisms-lemma-topological-invariance}
Sous les hypothèses et avec les notations du
Théorème \ref{spaces-more-morphisms-theorem-topological-invariance},
l'équivalence de catégories
$Y_{spaces, \etale} \to X_{spaces, \etale}$
induit par restriction des équivalences de catégories
$Y_\etale \to X_\etale$ et $Y_{affine, \etale} \to X_{affine, \etale}$.
\end{lemma}

\begin{proof}
Il s'agit simplement de dire que, pour tout objet
$V \in Y_{spaces, \etale}$, $V$ est un schéma (respectivement un schéma affine) si et
seulement si $V \times_Y X$ est un schéma (respectivement un schéma affine). Comme $V \times_Y X \to V$
est entier, universellement injectif et surjectif (en tant que changement
de base de $X \to Y$), cela
découle du Lemme
\ref{spaces-limits-lemma-integral-universally-bijective-scheme} de Limites d'espaces et de la
Proposition \ref{spaces-limits-proposition-affine}.
\end{proof}

\begin{remark}
\label{spaces-more-morphisms-remark-topological-invariance-etale-site}
\begin{reference}
Courriel de Lenny Taelman daté du 1er mai 2016.
\end{reference}
Un homéomorphisme universel d'espaces algébriques n'est pas nécessairement représentable, voir
Morphismes d'espaces,
Exemple \ref{spaces-morphisms-example-universal-homeomorphism}.
En fait, le Théorème \ref{spaces-more-morphisms-theorem-topological-invariance} n'est
pas valable pour tous les homéomorphismes universels. Pour le voir, soit $k$ un
corps algébriquement clos de caractéristique $0$ et soit
$$
\mathbf{A}^1 \to X \to \mathbf{A}^1
$$
comme dans Morphismes d'espaces,
Exemple \ref{spaces-morphisms-example-universal-homeomorphism}.
Rappelons que le premier morphisme est étale et identifie
$t$ à $-t$ pour $t \in \mathbf{A}^1_k \setminus \{0\}$
et que le second morphisme est notre homéomorphisme universel.
Comme $\mathbf{A}^1_k$ n'admet aucun
revêtement fini étale connexe non trivial
(car $k$ est algébriquement clos de caractéristique zéro ; détails omis),
il suffit de construire un revêtement fini étale connexe non trivial
$Y \to X$. Pour cela, soit $Y$ la droite affine
à origine dédoublée
(Schémas, Exemple \ref{schemes-example-affine-space-zero-doubled}).
Alors $Y = Y_1 \cup Y_2$, avec $Y_i = \mathbf{A}^1_k$, est obtenu par recollement
le long de $\mathbf{A}^1_k \setminus \{0\}$.
Pour définir le morphisme $Y \to X$, nous utilisons les morphismes
$$
Y_1 \xrightarrow{1} \mathbf{A}^1_k \to X
\quad\text{et}\quad
Y_2 \xrightarrow{-1} \mathbf{A}^1_k \to X.
$$
Ces morphismes se recollent sur $Y_1 \cap Y_2$ par construction de $X$ et
définissent donc un morphisme $Y \to X$. En fait, nous affirmons que
$$
\xymatrix{
Y \ar[d] & Y_1 \amalg Y_2 \ar[l] \ar[d] \\
X & \mathbf{A}^1_k \ar[l]
}
$$
est un carré cartésien. Nous omettons les détails ; on peut utiliser par exemple
Groupoïdes, Lemme \ref{groupoids-lemma-criterion-fibre-product}.
Comme $\mathbf{A}^1_k \to X$ est étale et
surjectif, cela montre que $Y \to X$
est fini étale de degré $2$, ce qui fournit l'exemple voulu.

\medskip\noindent
Plus simplement, on peut raisonner comme suit. Le schéma $Y$ est muni de l'action du groupe
$G = \{+1, -1\}$, qui opère librement, l'élément $-1$ agissant en échangeant
$Y_1$ et $Y_2$ et en changeant le signe de la coordonnée. Alors
$X = Y/G$ (voir Espaces, Définition \ref{spaces-definition-quotient})
et donc $Y \to X$ est fini étale. On peut aussi montrer directement
qu'il existe un homéomorphisme universel $X \to \mathbf{A}^1_k$
en utilisant $t \mapsto t^2$ sur les espaces affines. En fait, ce $X$ est
le même que le $X$ ci-dessus.
\end{remark}








\section{Épaississements}
\label{spaces-more-morphisms-section-thickenings}

\noindent
La terminologie suivante n'est peut-être pas tout à fait standard, mais elle est commode.

\begin{definition}
\label{spaces-more-morphisms-definition-thickening}
Épaississements. Soit $S$ un schéma.
\begin{enumerate}
\item Nous disons qu'un espace algébrique $X'$ est un {\it épaississement} d'un espace
algébrique $X$ si $X$ est un sous-espace fermé de $X'$ et si les espaces topologiques
associés sont égaux.
\item Nous disons que $X'$ est un {\it épaississement du premier ordre} de $X$ si
$X$ est un sous-espace fermé de $X'$ et si le faisceau quasi-cohérent d'idéaux
$\mathcal{I} \subset \mathcal{O}_{X'}$ qui définit $X$ est de carré nul.
\item Nous disons que $X'$ est un {\it épaississement d'ordre fini} de $X$
si $X$ est un sous-espace fermé de $X$ et si le faisceau quasi-cohérent d'idéaux
$\mathcal{I} \subset \mathcal{O}_{X'}$ qui définit $X$ est nilpotent, c'est-à-dire s'il
existe un entier $n \geq 0$ tel que $\mathcal{I}^{n + 1} = 0$.
\item Nous disons que $X'$ est un {\it épaississement d'ordre $n$} de $X$
si $X$ est un sous-espace fermé de $X$ et si $\mathcal{I}^{n + 1} = 0$,
où $\mathcal{I} \subset \mathcal{O}_{X'}$ est le
faisceau quasi-cohérent d'idéaux qui définit $X$.
\item Étant donnés deux épaississements $X \subset X'$ et $Y \subset Y'$, un
{\it morphisme d'épaississements} est un morphisme $f' : X' \to Y'$ tel que
$f(X) \subset Y$, c'est-à-dire tel que $f'|_X$ se factorise par le sous-espace
fermé $Y$. Dans cette situation, nous posons $f = f'|_X : X \to Y$ et disons
que $(f, f') : (X \subset X') \to (Y \subset Y')$ est un morphisme
d'épaississements.
\item Soit $B$ un espace algébrique. Nous définissons de même les
{\it épaississements sur $B$} et les
{\it morphismes d'épaississements sur $B$}. Cela signifie que les espaces
$X, X', Y, Y'$ ci-dessus sont des espaces algébriques munis d'un morphisme
structural vers $B$ et que les morphismes
$X \to X'$, $Y \to Y'$ et $f' : X' \to Y'$ sont des morphismes sur $B$.
\end{enumerate}
\end{definition}

\noindent
L'équivalence fondamentale.
Notons que si $X \subset X'$ est un épaississement, alors $X \to X'$
est entier et universellement bijectif. Il en résulte que
\begin{equation}
\label{spaces-more-morphisms-equation-equivalence-etale-spaces}
X_{spaces, \etale} = X'_{spaces, \etale}
\end{equation}
par le foncteur d'image inverse, voir le
Théorème \ref{spaces-more-morphisms-theorem-topological-invariance}.
Nous pouvons donc considérer $\mathcal{O}_{X'}$ comme un faisceau sur
$X_{spaces, \etale}$. Nous obtenons ainsi une équivalence canonique
de topos localement annelés
\begin{equation}
\label{spaces-more-morphisms-equation-fundamental-equivalence}
(\Sh(X'_{spaces, \etale}), \mathcal{O}_{X'})
\cong
(\Sh(X_{spaces, \etale}), \mathcal{O}_{X'})
\end{equation}
Dans la suite, nous la combinerons fréquemment avec le résultat de pleine fidélité de
Propriétés des espaces, Théorème \ref{spaces-properties-theorem-fully-faithful}.
Par exemple, à l'immersion fermée $i_X : X \to X'$ correspond
le morphisme surjectif $i_X^\sharp : \mathcal{O}_{X'} \to \mathcal{O}_X$.

\medskip\noindent
Soit $S$ un schéma et soit $B$ un espace algébrique sur $S$.
Soit $(f, f') : (X \subset X') \to (Y \subset Y')$ un morphisme
d'épaississements sur $B$. Notons que le diagramme de foncteurs continus
$$
\xymatrix{
X_{spaces, \etale} &
Y_{spaces, \etale} \ar[l] \\
X'_{spaces, \etale} \ar[u] &
Y'_{spaces, \etale} \ar[u] \ar[l]
}
$$
est commutatif et que les flèches verticales sont des équivalences. Ainsi,
$f_{spaces, \etale}$, $f_{small}$,
$f'_{spaces, \etale}$, et $f'_{small}$
définissent tous le même morphisme de topos. Nous pouvons donc considérer
$$
(f')^\sharp :
f_{spaces, \etale}^{-1}\mathcal{O}_{Y'}
\longrightarrow
\mathcal{O}_{X'}
$$
comme un morphisme de faisceaux de $\mathcal{O}_B$-algèbres qui s'insère dans le
diagramme commutatif
$$
\xymatrix{
f_{spaces, \etale}^{-1}\mathcal{O}_Y \ar[r]_-{f^\sharp} \ar[r] &
\mathcal{O}_X \\
f_{spaces, \etale}^{-1}\mathcal{O}_{Y'} \ar[r]^-{(f')^\sharp}
\ar[u]^{i_Y^\sharp} &
\mathcal{O}_{X'} \ar[u]_{i_X^\sharp}
}
$$
Ici, $i_X : X \to X'$ et $i_Y : Y \to Y'$ désignent les
immersions fermées données.

\begin{lemma}
\label{spaces-more-morphisms-lemma-first-order-thickening-maps}
Soit $S$ un schéma. Soit $B$ un espace algébrique sur $S$.
Soient $X \subset X'$ et $Y \subset Y'$ des épaississements
d'espaces algébriques sur $B$. Soit $f : X \to Y$ un morphisme d'espaces
algébriques sur $B$. Étant donné un morphisme de $\mathcal{O}_B$-algèbres
$$
\alpha : f_{spaces, \etale}^{-1}\mathcal{O}_{Y'} \to \mathcal{O}_{X'}
$$
tel que le diagramme
$$
\xymatrix{
f_{spaces, \etale}^{-1}\mathcal{O}_Y \ar[r]_-{f^\sharp} \ar[r] &
\mathcal{O}_X \\
f_{spaces, \etale}^{-1}\mathcal{O}_{Y'} \ar[r]^-\alpha
\ar[u]^{i_Y^\sharp} &
\mathcal{O}_{X'} \ar[u]_{i_X^\sharp}
}
$$
soit commutatif, il existe un unique morphisme $(f, f')$
d'épaississements sur $B$ tel que $\alpha = (f')^\sharp$.
\end{lemma}

\begin{proof}
Pour trouver $f'$, d'après
Propriétés des espaces, Théorème \ref{spaces-properties-theorem-fully-faithful},
il suffit de montrer que le morphisme de topos annelés
$$
(f_{spaces, \etale}, \alpha) :
(\Sh(X_{spaces, \etale}), \mathcal{O}_{X'})
\longrightarrow
(\Sh(Y_{spaces, \etale}), \mathcal{O}_{Y'})
$$
est un morphisme de topos localement annelés. Cela découle directement
de la définition des morphismes de topos localement annelés
(Modules sur les sites,
Définition \ref{sites-modules-definition-morphism-locally-ringed-topoi}),
du fait que $(f, f^\sharp)$ est un morphisme de topos localement annelés
(Propriétés des espaces,
Lemme \ref{spaces-properties-lemma-morphism-locally-ringed}),
du fait que $\alpha$ s'insère dans le diagramme commutatif donné et
du fait que les noyaux de $i_X^\sharp$ et de $i_Y^\sharp$ sont
localement nilpotents. Enfin, l'égalité $f' \circ i_X = i_Y \circ f$
découle de la commutativité du diagramme et d'une nouvelle application de
Propriétés des espaces, Théorème \ref{spaces-properties-theorem-fully-faithful}.
Nous omettons la vérification que $f'$ est un morphisme sur $B$.
\end{proof}

\begin{lemma}
\label{spaces-more-morphisms-lemma-open-subspace-thickening}
Soit $S$ un schéma. Soit $X \subset X'$ un épaississement
d'espaces algébriques sur $S$. Pour tout sous-espace ouvert $U \subset X$, il
existe un unique sous-espace ouvert $U' \subset X'$ tel que
$U = X \times_{X'} U'$.
\end{lemma}

\begin{proof}
Soit $U' \to X'$ l'objet de $X'_{spaces, \etale}$
qui correspond à l'objet $U \to X$ de $X_{spaces, \etale}$
par (\ref{spaces-more-morphisms-equation-equivalence-etale-spaces}). Le morphisme
$U' \to X'$ est étale et universellement injectif, donc une immersion ouverte, voir
Morphismes d'espaces,
Lemme \ref{spaces-morphisms-lemma-etale-universally-injective-open}.
\end{proof}

\noindent
Soit $i_X : X \to X'$ un épaississement d'espaces algébriques. Toute section locale
du noyau $\mathcal{I} = \Ker(i_X^\sharp) \subset \mathcal{O}_{X'}$ est
localement nilpotente. Cela donne un certain contrôle sur le faisceau structural
$\mathcal{O}_{X'}$, mais, techniquement, la classe des épaississements d'ordre fini
$X \subset X'$ est beaucoup plus facile à manier.
En effet, si $X'$ est un épaississement d'ordre $n$, nous avons une filtration
$$
0 \subset \mathcal{I}^n \subset \mathcal{I}^{n - 1} \subset
\ldots \subset \mathcal{I} \subset \mathcal{O}_{X'}
$$
et nous voyons que $X'$ est muni d'une filtration par des sous-espaces fermés
$$
X = X_0 \subset X_1 \subset \ldots \subset X_{n - 1} \subset X_{n + 1} = X'
$$
telle que chaque paire $X_i \subset X_{i + 1}$ est un épaississement du premier ordre
sur $B$. De simples arguments de récurrence permettent de reformuler nombre de résultats démontrés pour les
épaississements du premier ordre en des résultats sur les épaississements d'ordre fini.

\begin{lemma}
\label{spaces-more-morphisms-lemma-first-order-thickening-surjective}
Soit $S$ un schéma. Soit $X \subset X'$ un épaississement
d'espaces algébriques sur $S$. Soit $U$ un objet affine de
$X_{spaces, \etale}$. Alors
$$
\Gamma(U, \mathcal{O}_{X'}) \to \Gamma(U, \mathcal{O}_X)
$$
est surjectif, où l'on considère $\mathcal{O}_{X'}$ comme un faisceau sur
$X_{spaces, \etale}$ au moyen de (\ref{spaces-more-morphisms-equation-fundamental-equivalence}).
\end{lemma}

\begin{proof}
Soit $U' \to X'$ le morphisme étale d'espaces algébriques tel que
$U = X \times_{X'} U'$, voir le Théorème \ref{spaces-more-morphisms-theorem-topological-invariance}.
D'après Limites d'espaces, Lemme \ref{spaces-limits-lemma-affine}, nous voyons
que $U'$ est un schéma affine. Par conséquent,
$\Gamma(U, \mathcal{O}_{X'}) = \Gamma(U', \mathcal{O}_{U'}) \to
\Gamma(U, \mathcal{O}_U)$
est surjectif, puisque $U \to U'$ est une immersion fermée de schémas affines.
Nous donnons ci-dessous une démonstration directe pour les épaississements d'ordre fini,
cas le plus utilisé en pratique.
\end{proof}

\begin{proof}[Démonstration pour les épaississements d'ordre fini]
Nous pouvons supposer que $X \subset X'$ est un épaississement du premier ordre, par le
principe expliqué ci-dessus. Notons $\mathcal{I}$ le noyau de la surjection
$\mathcal{O}_{X'} \to \mathcal{O}_X$. Comme $\mathcal{I}$ est un
$\mathcal{O}_{X'}$-module quasi-cohérent et que $\mathcal{I}^2 = 0$ par définition
d'un épaississement du premier ordre, nous pouvons appliquer
Morphismes d'espaces, Lemme
\ref{spaces-morphisms-lemma-i-star-equivalence}
pour voir que $\mathcal{I}$ est un $\mathcal{O}_X$-module quasi-cohérent.
Le lemme découle donc de la suite exacte longue de cohomologie
associée à la suite exacte courte
$$
0 \to \mathcal{I} \to \mathcal{O}_{X'} \to \mathcal{O}_X \to 0
$$
et du fait que $H^1_\etale(U, \mathcal{I}) = 0$, puisque
$\mathcal{I}$ est quasi-cohérent, voir
Descente, Proposition \ref{descent-proposition-same-cohomology-quasi-coherent},
et Cohomologie des schémas, Lemme
\ref{coherent-lemma-quasi-coherent-affine-cohomology-zero}.
\end{proof}

\begin{lemma}
\label{spaces-more-morphisms-lemma-thickening-scheme}
Soit $S$ un schéma. Soit $X \subset X'$ un épaississement d'espaces algébriques
sur $S$. Si $X$ est (représentable par) un schéma, alors $X'$ l'est aussi.
\end{lemma}

\begin{proof}
Notons que $X'_{red} = X_{red}$. Par conséquent, si $X$ est un schéma, alors
$X'_{red}$ est un schéma. Le résultat découle donc de
Limites d'espaces, Lemme
\ref{spaces-limits-lemma-reduction-scheme}.
Nous donnons ci-dessous une démonstration directe pour les épaississements d'ordre fini,
cas le plus souvent utilisé en pratique.
\end{proof}

\begin{proof}[Démonstration pour les épaississements d'ordre fini]
Il suffit de le démontrer lorsque $X'$ est un épaississement du premier ordre de $X$. D'après
Propriétés des espaces, Lemme \ref{spaces-properties-lemma-subscheme},
il existe un plus grand sous-espace ouvert de $X'$ qui soit un schéma. Nous devons donc
montrer que tout point $x$ de $|X'| = |X|$ est contenu dans un sous-espace ouvert de
$X'$ qui soit un schéma. En utilisant le
Lemme \ref{spaces-more-morphisms-lemma-open-subspace-thickening},
nous pouvons remplacer $X \subset X'$ par $U \subset U'$, avec $x \in U$ et $U$
un schéma affine. Nous pouvons donc supposer que $X$ est affine.
Nous sommes ainsi ramenés au cas traité dans le paragraphe suivant.

\medskip\noindent
Supposons que $X \subset X'$ soit un épaississement du premier ordre, où $X$ est un schéma
affine. Posons $A = \Gamma(X, \mathcal{O}_X)$ et
$A' = \Gamma(X', \mathcal{O}_{X'})$. D'après le
Lemme \ref{spaces-more-morphisms-lemma-first-order-thickening-surjective},
le morphisme $A' \to A$ est surjectif. Son noyau $I$ est un idéal de carré nul. D'après
Propriétés des espaces,
Lemme \ref{spaces-properties-lemma-morphism-to-affine-scheme},
nous obtenons un morphisme canonique $f : X' \to \Spec(A')$ qui s'insère
dans le diagramme commutatif suivant
$$
\xymatrix{
X \ar@{=}[d] \ar[r] &  X' \ar[d]^f \\
\Spec(A) \ar[r] & \Spec(A')
}
$$
Comme les flèches horizontales sont des épaississements, il est clair que $f$ est
universellement injectif et surjectif. Il suffit donc de montrer que
$f$ est étale, car alors
Morphismes d'espaces,
Lemme \ref{spaces-morphisms-lemma-etale-universally-injective-open},
entraînera que $f$ est un isomorphisme.

\medskip\noindent
Pour démontrer que $f$ est étale, choisissons un schéma affine $U'$ et un
morphisme étale $U' \to X'$. Il suffit de montrer que
$U' \to X' \to \Spec(A')$ est étale, voir
Propriétés des espaces, Définition \ref{spaces-properties-definition-etale}.
Écrivons $U' = \Spec(B')$. Posons $U = X \times_{X'} U'$. Comme $U$
est un sous-espace fermé de $U'$, c'est un sous-schéma fermé, et donc
$U = \Spec(B)$ avec $B' \to B$ surjectif. Notons
$J = \Ker(B' \to B)$ et remarquons que $J = \Gamma(U, \mathcal{I})$,
où $\mathcal{I} = \Ker(\mathcal{O}_{X'} \to \mathcal{O}_X)$
sur $X_{spaces, \etale}$, comme dans la démonstration du
Lemme \ref{spaces-more-morphisms-lemma-first-order-thickening-surjective}.
Le morphisme $U' \to X' \to \Spec(A')$ induit un diagramme
commutatif
$$
\xymatrix{
0 \ar[r] &
J \ar[r] &
B' \ar[r] &
B \ar[r] & 0 \\
0 \ar[r] &
I \ar[r] \ar[u] &
A' \ar[r] \ar[u] &
A \ar[r] \ar[u] & 0
}
$$
Maintenant, comme $\mathcal{I}$ est un $\mathcal{O}_X$-module quasi-cohérent,
nous avons $\mathcal{I} = (\widetilde I)^a$, voir
Descente, Définition \ref{descent-definition-structure-sheaf},
pour les notations, et
Descente, Proposition \ref{descent-proposition-equivalence-quasi-coherent},
pour la justification. Nous voyons donc que $J = I \otimes_A B$.
Enfin, notons que $A \to B$ est étale, car $U \to X$ est étale en tant que
changement de base du morphisme étale $U' \to X'$.
Nous concluons que $A' \to B'$ est étale par
Algèbre, Lemme \ref{algebra-lemma-lift-etale-infinitesimal}.
\end{proof}

\begin{lemma}
\label{spaces-more-morphisms-lemma-thickening-equivalence}
Soit $S$ un schéma. Soit $X \subset X'$ un épaississement d'espaces algébriques
sur $S$. Le foncteur
$$
V' \longmapsto V = X \times_{X'} V'
$$
définit une équivalence de catégories
$X'_\etale \to X_\etale$.
\end{lemma}

\begin{proof}
Le foncteur $V' \mapsto V$ définit une équivalence de catégories
$X'_{spaces, \etale} \to X_{spaces, \etale}$, voir le
Théorème \ref{spaces-more-morphisms-theorem-topological-invariance}.
Il suffit donc de montrer que $V$ est un schéma si et seulement si $V'$ est
un schéma. C'est le contenu du
Lemme \ref{spaces-more-morphisms-lemma-thickening-scheme}.
\end{proof}

\noindent
Les épaississements du premier ordre se décrivent comme suit.

\begin{lemma}
\label{spaces-more-morphisms-lemma-first-order-thickening}
Soit $S$ un schéma.
Soit $f : X \to B$ un morphisme d'espaces algébriques sur $S$.
Considérons une suite exacte courte
$$
0 \to \mathcal{I} \to \mathcal{A} \to \mathcal{O}_X \to 0
$$
de faisceaux sur $X_\etale$, où $\mathcal{A}$ est un faisceau de
$f^{-1}\mathcal{O}_B$-algèbres, $\mathcal{A} \to \mathcal{O}_X$ est une surjection
de faisceaux de $f^{-1}\mathcal{O}_B$-algèbres, et $\mathcal{I}$ son noyau.
Si
\begin{enumerate}
\item $\mathcal{I}$ est un idéal de carré nul dans $\mathcal{A}$, et
\item $\mathcal{I}$ est quasi-cohérent en tant que $\mathcal{O}_X$-module,
\end{enumerate}
alors il existe un épaississement du premier ordre
$X \subset X'$ sur $B$ et un isomorphisme
$\mathcal{O}_{X'} \to \mathcal{A}$ de $f^{-1}\mathcal{O}_B$-algèbres
compatible aux surjections vers $\mathcal{O}_X$.
\end{lemma}

\begin{proof}
Dans cette démonstration, nous reprenons certains des arguments employés dans les
démonstrations des
Lemmes \ref{spaces-more-morphisms-lemma-first-order-thickening-surjective} et
\ref{spaces-more-morphisms-lemma-thickening-scheme}.
Nous traitons d'abord le cas $B = S = \Spec(\mathbf{Z})$.
Soit $U$ un schéma affine et soit $U \to X$ étale.
Alors
$$
0 \to \mathcal{I}(U) \to \mathcal{A}(U) \to \mathcal{O}_X(U) \to 0
$$
est exacte, car $H^1(U_\etale, \mathcal{I}) = 0$ puisque
$\mathcal{I}$ est quasi-cohérent, voir
Descente, Proposition \ref{descent-proposition-same-cohomology-quasi-coherent},
et Cohomologie des schémas, Lemme
\ref{coherent-lemma-quasi-coherent-affine-cohomology-zero}.
Si $V \to U$ est un morphisme d'objets affines de $X_{spaces, \etale}$,
alors
$$
\mathcal{I}(V) = \mathcal{I}(U) \otimes_{\mathcal{O}_X(U)} \mathcal{O}_X(V)
$$
puisque $\mathcal{I}$ est un $\mathcal{O}_X$-module quasi-cohérent, voir
Descente, Proposition \ref{descent-proposition-equivalence-quasi-coherent}.
Ainsi $\mathcal{A}(U) \to \mathcal{A}(V)$ est un homomorphisme d'anneaux
étale, voir
Algèbre, Lemme \ref{algebra-lemma-lift-etale-infinitesimal}.
Nous voyons donc que
$$
U \longmapsto U' = \Spec(\mathcal{A}(U))
$$
est un foncteur de $X_{affine, \etale}$ vers la catégorie des schémas
affines et des morphismes étales. En fait, nous affirmons que ce foncteur peut
s'étendre en un foncteur $U \mapsto U'$ sur tout $X_\etale$.
Pour le voir, si $U$ est un objet de $X_\etale$, notons que
$$
0 \to \mathcal{I}|_{U_{Zar}} \to \mathcal{A}|_{U_{Zar}} \to
\mathcal{O}_X|_{U_{Zar}} \to 0
$$
et que $\mathcal{I}|_{U_{Zar}}$ est un faisceau quasi-cohérent sur $U$, voir
Descente,
Proposition \ref{descent-proposition-equivalence-quasi-coherent-functorial}.
Par conséquent, d'après
Compléments sur les morphismes, Lemme \ref{more-morphisms-lemma-first-order-thickening},
nous obtenons un épaississement du premier ordre de schémas $U \subset U'$ tel que
$\mathcal{O}_{U'}$ soit isomorphe à $\mathcal{A}|_{U_{Zar}}$. Il est clair que
cette construction est compatible avec la construction affine ci-dessus.

\medskip\noindent
Choisissons une présentation $X = U/R$, voir
Espaces, Définition \ref{spaces-definition-presentation},
de sorte que $s, t : R \to U$ définissent une relation d'équivalence étale.
En appliquant le foncteur ci-dessus, nous obtenons une relation d'équivalence
étale $s', t' : R' \to U'$ dans les schémas. Considérons l'espace algébrique
$X' = U'/R'$ (voir
Espaces, Théorème \ref{spaces-theorem-presentation}).
Le morphisme $X = U/R \to U'/R' = X'$ est un épaississement du premier ordre.
Considérons $\mathcal{O}_{X'}$ comme un faisceau sur $X_\etale$.
Par construction, nous avons un isomorphisme
$$
\gamma :
\mathcal{O}_{X'}|_{U_\etale}
\longrightarrow
\mathcal{A}|_{U_\etale}
$$
tel que $s^{-1}\gamma$ coïncide avec $t^{-1}\gamma$ sur $R_\etale$.
Par conséquent, d'après
Propriétés des espaces, Lemme \ref{spaces-properties-lemma-descent-sheaf},
cela implique que $\gamma$ provient d'un unique isomorphisme
$\mathcal{O}_{X'} \to \mathcal{A}$, comme voulu.

\medskip\noindent
Pour traiter le cas d'un espace algébrique de base quelconque $B$, nous
construisons d'abord $X'$ comme espace algébrique sur $\mathbf{Z}$, ainsi que ci-dessus.
Nous utilisons ensuite l'isomorphisme $\mathcal{O}_{X'} \to \mathcal{A}$ pour
définir $f^{-1}\mathcal{O}_B \to \mathcal{O}_{X'}$. D'après le
Lemme \ref{spaces-more-morphisms-lemma-first-order-thickening-maps},
cela définit un morphisme $X' \to B$ compatible avec le morphisme donné
$X \to B$, ce qui achève la démonstration.
\end{proof}

\begin{lemma}
\label{spaces-more-morphisms-lemma-base-change-thickening}
Soit $S$ un schéma. Soit $Y \subset Y'$ un épaississement d'espaces algébriques
sur $S$. Soit $X' \to Y'$ un morphisme et posons $X = Y \times_{Y'} X'$.
Alors $(X \subset X') \to (Y \subset Y')$
est un morphisme d'épaississements. Si $Y \subset Y'$ est un épaississement du premier
ordre (respectivement d'ordre fini), alors $X \subset X'$ est un épaississement du premier
ordre (respectivement d'ordre fini).
\end{lemma}

\begin{proof}
Omis.
\end{proof}

\begin{lemma}
\label{spaces-more-morphisms-lemma-composition-thickening}
Soit $S$ un schéma. Si $X \subset X'$ et $X' \subset X''$ sont des
épaississements d'espaces algébriques sur $S$, alors $X \subset X''$ en est un également.
\end{lemma}

\begin{proof}
Omis.
\end{proof}

\begin{lemma}
\label{spaces-more-morphisms-lemma-descending-property-thickening}
La propriété d'être un épaississement est locale pour la topologie fpqc.
Il en va de même pour les épaississements du premier ordre.
\end{lemma}

\begin{proof}
L'énoncé signifie ceci : soit $S$ un schéma et soit
$X \to X'$ un morphisme d'espaces algébriques sur $S$.
Soit $\{g_i : X'_i \to X'\}$ un recouvrement fpqc d'espaces algébriques
tel que le changement de base $X_i \to X'_i$ soit un épaississement pour tout $i$.
Alors $X \to X'$ est un épaississement. Comme les morphismes $g_i$ sont conjointement
surjectifs, nous concluons que $X \to X'$ est surjectif. D'après
Descente sur les espaces, Lemme
\ref{spaces-descent-lemma-descending-property-closed-immersion},
nous concluons que $X \to X'$ est une immersion fermée.
Ainsi $X \to X'$ est un épaississement. Nous omettons la démonstration dans le
cas des épaississements du premier ordre.
\end{proof}








\section{Morphismes d'épaississements}
\label{spaces-more-morphisms-section-morphisms-thickenings}

\noindent
Si $(f, f') : (X \subset X') \to (Y \subset Y')$ est un morphisme
d'épaississements d'espaces algébriques, les propriétés de $f$
se transmettent souvent à $f'$. Il existe plusieurs variantes.

\begin{lemma}
\label{spaces-more-morphisms-lemma-thicken-property-morphisms}
Soit $S$ un schéma. Soit $(f, f') : (X \subset X') \to (Y \subset Y')$
un morphisme d'épaississements d'espaces algébriques sur $S$. Alors
\begin{enumerate}
\item $f$ est un morphisme affine si et seulement si $f'$ est un morphisme affine,
\item $f$ est un morphisme surjectif si et seulement si $f'$ est un morphisme surjectif,
\item $f$ est quasi-compact si et seulement si $f'$ est quasi-compact,
\item $f$ est universellement fermé si et seulement si $f'$ est universellement fermé,
\item $f$ est entier si et seulement si $f'$ est entier,
\item $f$ est (quasi-)séparé si et seulement si $f'$ est (quasi-)séparé,
\item $f$ est universellement injectif si et seulement si $f'$ est universellement injectif,
\item $f$ est universellement ouvert si et seulement si $f'$ est universellement ouvert,
\item $f$ est représentable si et seulement si $f'$ est représentable, et
\item ajouter d'autres propriétés ici.
\end{enumerate}
\end{lemma}

\begin{proof}
Observons que $Y \to Y'$ et $X \to X'$ sont des morphismes entiers et des
homéomorphismes universels. Cela entraîne immédiatement les assertions
(2), (3), (4), (7) et (8).
L'assertion (1) découle de
Limites d'espaces, Proposition \ref{spaces-limits-proposition-affine},
qui nous dit qu'il existe une correspondance bijective entre les
schémas affines étales sur $X$ et sur $X'$, et entre les schémas affines
étales sur $Y$ et sur $Y'$.
L'assertion (5) découle de (1) et (4) par
Morphismes d'espaces, Lemme
\ref{spaces-morphisms-lemma-integral-universally-closed}.
Enfin, notons que
$$
X \times_Y X = X \times_{Y'} X \to X \times_{Y'} X' \to X' \times_{Y'} X'
$$
est un épaississement (les deux flèches sont des épaississements par le
Lemme \ref{spaces-more-morphisms-lemma-base-change-thickening}).
Ainsi, en appliquant (3) et (4) au morphisme
$(X \subset X') \to (X \times_Y X \to X' \times_{Y'} X')$
nous obtenons (6). Enfin, l'assertion (9) découle du fait qu'un espace
algébrique qui est un épaississement d'un schéma est encore un schéma, voir le
Lemme \ref{spaces-more-morphisms-lemma-thickening-scheme}.
\end{proof}

\begin{lemma}
\label{spaces-more-morphisms-lemma-thicken-property-morphisms-cartesian}
Soit $S$ un schéma. Soit $(f, f') : (X \subset X') \to (Y \subset Y')$ un
morphisme d'épaississements d'espaces algébriques sur $S$ tel que
$X = Y \times_{Y'} X'$. Si $X \subset X'$ est un épaississement d'ordre fini, alors
\begin{enumerate}
\item $f$ est une immersion fermée si et seulement si $f'$ est une immersion fermée,
\item $f$ est localement de type fini si et seulement si $f'$ est
localement de type fini,
\item $f$ est localement quasi-fini si et seulement si $f'$ est localement
quasi-fini,
\item $f$ est localement de type fini de dimension relative $d$ si et
seulement si $f'$ est localement de type fini de dimension relative $d$,
\item $\Omega_{X/Y} = 0$ si et seulement si $\Omega_{X'/Y'} = 0$,
\item $f$ est non ramifié si et seulement si $f'$ est non ramifié,
\item $f$ est propre si et seulement si $f'$ est propre,
\item $f$ est un morphisme fini si et seulement si $f'$ est un morphisme fini,
\item $f$ est un monomorphisme si et seulement si $f'$ est un monomorphisme,
\item $f$ est une immersion si et seulement si $f'$ est une immersion, et
\item ajouter d'autres propriétés ici.
\end{enumerate}
\end{lemma}

\begin{proof}
Choisissons un schéma $V'$ et un morphisme étale surjectif $V' \to Y'$.
Choisissons un schéma $U'$ et un morphisme étale surjectif
$U' \to X' \times_{Y'} V'$. Posons $V = Y \times_{Y'} V'$ et
$U = X \times_{X'} U'$. Pour les propriétés des morphismes qui sont locales
pour la topologie étale, nous pouvons alors nous ramener au morphisme d'épaississements de schémas
$(U \subset U') \to (V \subset V')$ et appliquer Compléments sur les morphismes, Lemme
\ref{more-morphisms-lemma-thicken-property-morphisms-cartesian}.
Cela démontre (2), (3), (4), (5) et (6).

\medskip\noindent
Les propriétés des morphismes figurant en (1), (7), (8), (9) et (10) sont stables
par changement de base ; par conséquent, si $f'$ possède la propriété $\mathcal{P}$, alors
$f$ la possède également. Voir
Espaces, Lemme \ref{spaces-lemma-base-change-immersions},
et
Morphismes d'espaces, Lemmes
\ref{spaces-morphisms-lemma-base-change-proper},
\ref{spaces-morphisms-lemma-base-change-integral}, et
\ref{spaces-morphisms-lemma-base-change-monomorphism}.

\medskip\noindent
Le sens intéressant de (1), (7), (8), (9) et (10) consiste à supposer
que $f$ possède la propriété et à en déduire que $f'$ la possède aussi.
Par récurrence sur l'ordre de l'épaississement, nous pouvons
supposer que $Y \subset Y'$ est un épaississement du premier ordre, voir la
discussion ci-dessus sur les épaississements d'ordre fini.

\medskip\noindent
Démonstration de (1). Choisissons un schéma $V'$ et un morphisme étale surjectif
$V' \to Y'$. Posons $V = Y \times_{Y'} V'$, $U' = X' \times_{Y'} V'$ et
$U = X \times_Y V$. Alors $U \to V$ est une immersion fermée, ce qui
implique que $U$ est un schéma, et entraîne à son tour que $U'$ est
un schéma (Lemme \ref{spaces-more-morphisms-lemma-thickening-scheme}). Nous pouvons donc appliquer
le lemme dans le cas des schémas
(Compléments sur les morphismes, Lemme
\ref{more-morphisms-lemma-thicken-property-morphisms-cartesian})
à $(U \subset U') \to (V \subset V')$ pour conclure.

\medskip\noindent
Démonstration de (7). Elle résulte de la combinaison de (2) avec les
résultats du Lemme \ref{spaces-more-morphisms-lemma-thicken-property-morphisms}
et du fait qu'être propre équivaut à être quasi-compact $+$
séparé $+$ localement de type fini $+$ universellement fermé.

\medskip\noindent
Démonstration de (8). Elle résulte de la combinaison de (2) avec les
résultats du Lemme \ref{spaces-more-morphisms-lemma-thicken-property-morphisms}
et du fait qu'être fini équivaut à être entier $+$ localement
de type fini (Morphismes, Lemme \ref{morphisms-lemma-finite-integral}).

\medskip\noindent
Démonstration de (9). Comme $f$ est un monomorphisme, nous avons $X = X \times_Y X$.
Nous pouvons appliquer les résultats déjà démontrés au morphisme
d'épaississements $(X \subset X') \to (X \times_Y X \subset X' \times_{Y'} X')$.
Nous concluons par (1) que $X' \to X' \times_{Y'} X'$ est une immersion fermée.
En fait, c'est un épaississement du premier ordre, car l'idéal qui définit
l'immersion fermée $X' \to X' \times_{Y'} X'$ est contenu dans l'image inverse
de l'idéal $\mathcal{I} \subset \mathcal{O}_{Y'}$ qui définit $Y$ dans $Y'$.
En effet, $X = X \times_Y X = (X' \times_{Y'} X') \times_{Y'} Y$ est contenu
dans $X'$. Le faisceau conormal de l'immersion fermée
$\Delta : X' \to X' \times_{Y'} X'$ est égal à $\Omega_{X'/Y'}$
(c'est l'analogue de
Morphismes, Lemme \ref{morphisms-lemma-differentials-diagonal},
pour les espaces algébriques, et cela découle soit d'une localisation étale,
soit de la combinaison des
Lemmes \ref{spaces-more-morphisms-lemma-differentials-relative-immersion-section} et
\ref{spaces-more-morphisms-lemma-differential-product} ; certains détails sont omis).
Il suffit donc de montrer que $\Omega_{X'/Y'} = 0$, ce qui découle de (5)
et de l'énoncé correspondant pour $X/Y$.

\medskip\noindent
Démonstration de (10). Si $f : X \to Y$ est une immersion, elle se factorise sous la forme
$X \to V \to Y$, où $V \to Y$ est un sous-espace ouvert et $X \to V$ une
immersion fermée, voir
Morphismes d'espaces, Remarque \ref{spaces-morphisms-remark-immersion}.
Soit $V' \subset Y'$ le sous-espace ouvert dont
l'espace topologique sous-jacent $|V'|$ est égal à $|V| \subset |Y| = |Y'|$.
Alors $X' \to Y'$ se factorise par $V'$, et nous concluons par (1) que $X' \to V'$
est une immersion fermée. Cela achève la démonstration.
\end{proof}

\noindent
Le lemme suivant est une variante du précédent. Au lieu de supposer
que les épaississements en jeu sont d'ordre fini (ce qui permet de transférer
de $f$ à $f'$ la propriété d'être localement de type fini),
nous supposons plutôt que $f$ et $f'$ sont tous deux localement
de type fini.

\begin{lemma}
\label{spaces-more-morphisms-lemma-properties-that-extend-over-thickenings}
Soit $S$ un schéma. Soit $(f, f') : (X \subset X') \to (Y \to Y')$ un
morphisme d'épaississements d'espaces algébriques sur $S$. Supposons que $f$ et $f'$
soient localement de type fini et que $X = Y \times_{Y'} X'$. Alors
\begin{enumerate}
\item $f$ est localement quasi-fini si et seulement si $f'$ est localement quasi-fini,
\item $f$ est fini si et seulement si $f'$ est fini,
\item $f$ est une immersion fermée si et seulement si $f'$ est une immersion fermée,
\item $\Omega_{X/Y} = 0$ si et seulement si $\Omega_{X'/Y'} = 0$,
\item $f$ est non ramifié si et seulement si $f'$ est non ramifié,
\item $f$ est un monomorphisme si et seulement si $f'$ est un monomorphisme,
\item $f$ est une immersion si et seulement si $f'$ est une immersion,
\item $f$ est propre si et seulement si $f'$ est propre, et
\item ajouter d'autres propriétés ici.
\end{enumerate}
\end{lemma}

\begin{proof}
Choisissons un schéma $V'$ et un morphisme étale surjectif $V' \to Y'$.
Choisissons un schéma $U'$ et un morphisme étale surjectif
$U' \to X' \times_{Y'} V'$. Posons $V = Y \times_{Y'} V'$ et
$U = X \times_{X'} U'$. Pour les propriétés des morphismes qui sont locales
pour la topologie étale, nous pouvons alors nous ramener au morphisme d'épaississements de schémas
$(U \subset U') \to (V \subset V')$ et appliquer
Compléments sur les morphismes, Lemme
\ref{more-morphisms-lemma-properties-that-extend-over-thickenings}.
Cela démontre (1), (4) et (5).

\medskip\noindent
Les propriétés figurant en (2), (3), (6), (7) et (8) sont stables
par changement de base ; par conséquent, si $f'$ possède la propriété $\mathcal{P}$, alors
$f$ la possède également. Voir Espaces, Lemme \ref{spaces-lemma-base-change-immersions},
et
Morphismes d'espaces, Lemmes
\ref{spaces-morphisms-lemma-base-change-proper},
\ref{spaces-morphisms-lemma-base-change-integral}, et
\ref{spaces-morphisms-lemma-base-change-monomorphism}.
Ainsi, dans chaque cas, il suffit de démontrer que, si $f$ possède
la propriété voulue, alors $f'$ la possède aussi.

\medskip\noindent
Le cas (2) découle du cas (5) du Lemme \ref{spaces-more-morphisms-lemma-thicken-property-morphisms}
et du fait que les morphismes finis sont exactement
les morphismes entiers qui sont localement de type fini
(Morphismes d'espaces, Lemme \ref{spaces-morphisms-lemma-finite-integral}).

\medskip\noindent
Cas (3). Cela découle immédiatement de
Limites d'espaces, Lemme
\ref{spaces-limits-lemma-check-closed-infinitesimally}.

\medskip\noindent
Démonstration de (6). Comme $f$ est un monomorphisme, nous avons $X = X \times_Y X$.
Nous pouvons appliquer les résultats déjà démontrés au morphisme d'épaississements
$(X \subset X') \to (X \times_Y X \subset X' \times_{Y'} X')$.
Nous concluons par (3) que $\Delta_{X'/Y'} : X' \to X' \times_{Y'} X'$
est une immersion fermée. En fait, $\Delta_{X'/Y'}$ induit une bijection
$|X'| \to |X' \times_{Y'} X'|$, donc $\Delta_{X'/Y'}$ est un épaississement.
D'autre part, $\Delta_{X'/Y'}$ est localement de présentation finie d'après
Morphismes d'espaces, Lemme
\ref{spaces-morphisms-lemma-diagonal-morphism-finite-type}.
Autrement dit, $\Delta_{X'/Y'}(X')$ est défini par
un faisceau quasi-cohérent d'idéaux
$\mathcal{J} \subset \mathcal{O}_{X' \times_{Y'} X'}$ de type fini.
Comme $\Omega_{X'/Y'} = 0$ par (5), nous voyons que
le faisceau conormal de $X' \to X' \times_{Y'} X'$ est nul.
(Le faisceau conormal de l'immersion fermée $\Delta_{X'/Y'}$ est égal à
$\Omega_{X'/Y'}$ ; c'est l'analogue de
Morphismes, Lemme \ref{morphisms-lemma-differentials-diagonal},
pour les espaces algébriques, et cela découle soit d'une localisation étale,
soit de la combinaison des
Lemmes \ref{spaces-more-morphisms-lemma-differentials-relative-immersion-section} et
\ref{spaces-more-morphisms-lemma-differential-product} ; certains détails sont omis.)
Autrement dit, $\mathcal{J}/\mathcal{J}^2 = 0$.
Cela implique que $\Delta_{X'/Y'}$ est un isomorphisme, par exemple
par Algèbre, Lemme \ref{algebra-lemma-ideal-is-squared-union-connected}.

\medskip\noindent
Démonstration de (7). Si $f : X \to Y$ est une immersion, elle se factorise sous la forme
$X \to V \to Y$, où $V \to Y$ est un sous-espace ouvert et $X \to V$ une
immersion fermée, voir
Morphismes d'espaces, Remarque \ref{spaces-morphisms-remark-immersion}.
Soit $V' \subset Y'$ le sous-espace ouvert dont
l'espace topologique sous-jacent $|V'|$ est égal à $|V| \subset |Y| = |Y'|$.
Alors $X' \to Y'$ se factorise par $V'$, et nous concluons par (3) que $X' \to V'$
est une immersion fermée.

\medskip\noindent
Le cas (8) découle du Lemme \ref{spaces-more-morphisms-lemma-thicken-property-morphisms}
et de la définition des morphismes propres comme morphismes quasi-compacts,
universellement fermés et séparés qui sont localement de type fini.
\end{proof}










\section{Groupes de Picard des épaississements}
\label{spaces-more-morphisms-section-picard-group-thickening}

\noindent
Quelques éléments sur les groupes de Picard des épaississements.

\begin{lemma}
\label{spaces-more-morphisms-lemma-picard-group-first-order-thickening}
Soit $S$ un schéma. Soit $X \subset X'$ un épaississement du premier ordre
d'espaces algébriques sur $S$, dont le faisceau d'idéaux est $\mathcal{I}$.
Il existe alors une suite exacte canonique
$$
\xymatrix{
0 \ar[r] &
H^0(X, \mathcal{I}) \ar[r] &
H^0(X', \mathcal{O}_{X'}^*) \ar[r] &
H^0(X, \mathcal{O}^*_X) \ar `r[d] `d[l] `l[llld] `d[dll] [dll] \\
& H^1(X, \mathcal{I}) \ar[r] &
\Pic(X') \ar[r] &
\Pic(X) \ar `r[d] `d[l] `l[llld] `d[dll] [dll] \\
& H^2(X, \mathcal{I}) \ar[r] & \ldots \ar[r] & \ldots
}
$$
de groupes abéliens.
\end{lemma}

\begin{proof}
Rappelons que $X_\etale = X'_\etale$, voir le
Lemme \ref{spaces-more-morphisms-lemma-thickening-equivalence} et, plus généralement, la
discussion de la section \ref{spaces-more-morphisms-section-thickenings}.
La suite du lemme est la suite exacte longue de cohomologie
associée à la suite exacte courte de faisceaux de groupes abéliens
$$
0 \to \mathcal{I} \to \mathcal{O}_{X'}^* \to \mathcal{O}_X^* \to 0
$$
sur $X_\etale$, où le premier morphisme envoie une section locale $f$ de $\mathcal{I}$
sur la section inversible $1 + f$ de $\mathcal{O}_{X'}$.
Nous utilisons aussi l'identification du groupe de Picard d'un
site annelé au premier groupe de cohomologie du faisceau
des éléments inversibles, voir
Cohomologie sur les sites, Lemme \ref{sites-cohomology-lemma-h1-invertible}.
\end{proof}








\section{Voisinages infinitésimaux}
\label{spaces-more-morphisms-section-first-order-infinitesimal-neighbourhood}

\noindent
Une construction naturelle d'épaississements d'ordre fini est la suivante.
Supposons que $i : Z \to X$ soit une immersion d'espaces algébriques. Choisissons un
sous-espace ouvert $U \subset X$ tel que $i$ identifie $Z$ au sous-espace
fermé $Z \subset U$ (voir
Morphismes d'espaces, Remarque \ref{spaces-morphisms-remark-immersion}).
Soit $\mathcal{I} \subset \mathcal{O}_U$ le
faisceau quasi-cohérent d'idéaux qui définit $Z$ dans $U$, voir
Morphismes d'espaces,
Lemme \ref{spaces-morphisms-lemma-closed-immersion-ideals}.
Pour $n \geq 1$, nous pouvons considérer le sous-espace fermé $Z_n \subset U$
défini par le faisceau quasi-cohérent d'idéaux $\mathcal{I}^{n + 1}$.

\begin{definition}
\label{spaces-more-morphisms-definition-first-order-infinitesimal-neighbourhood}
Soit $i : Z \to X$ une immersion d'espaces algébriques.
\begin{enumerate}
\item Le {\it voisinage infinitésimal du premier ordre} de $Z$ dans $X$ est
l'épaississement du premier ordre $Z \subset Z_1$ sur $X$ décrit ci-dessus.
\item Le {\it voisinage infinitésimal d'ordre $n$} de $Z$ dans $X$ est
l'épaississement d'ordre $n$ $Z \subset Z_n$ sur $X$ décrit ci-dessus.
\end{enumerate}
\end{definition}

\noindent
Cet épaississement possède la propriété universelle suivante (ce qui dissipera
toute crainte que la construction ci-dessus ne dépende du choix de l'ouvert
$U$).

\begin{lemma}
\label{spaces-more-morphisms-lemma-first-order-infinitesimal-neighbourhood}
Soit $i : Z \to X$ une immersion d'espaces algébriques.
\begin{enumerate}
\item Le voisinage infinitésimal du premier ordre $Z'$ de $Z$ dans $X$
possède la propriété universelle suivante :
pour tout diagramme commutatif
$$
\xymatrix{
Z \ar[d]_i & T \ar[l]^a \ar[d] \\
X & T' \ar[l]_b
}
$$
où $T \subset T'$ est un épaississement du premier ordre sur $X$, il existe
un unique morphisme $(a', a) : (T \subset T') \to (Z \subset Z')$
d'épaississements sur $X$.
\item Pour $n \geq 1$, le voisinage infinitésimal d'ordre $n$
$Z_n$ de $Z$ dans $X$ possède la propriété universelle suivante :
pour tout diagramme commutatif
$$
\xymatrix{
Z \ar[d]_i & T \ar[l]^a \ar[d] \\
X & T' \ar[l]_b
}
$$
où $T \subset T'$ est un épaississement d'ordre $n$ sur $X$, il existe
un unique morphisme $(a', a) : (T \subset T') \to (Z \subset Z_n)$
d'épaississements sur $X$.
\end{enumerate}
\end{lemma}

\begin{proof}
Nous ne démontrerons que (1).
Soit $U \subset X$ le sous-espace ouvert utilisé dans la construction de $Z'$,
c'est-à-dire un ouvert tel que $Z$ soit identifié à un sous-espace fermé de $U$
défini par le faisceau quasi-cohérent d'idéaux $\mathcal{I}$.
Comme $|T| = |T'|$, nous voyons que $|b|(|T'|) \subset |U|$. Nous pouvons donc
considérer $b$ comme un morphisme à valeurs dans $U$, voir
Propriétés des espaces,
Lemme \ref{spaces-properties-lemma-factor-through-open-subspace}.
Soit $\mathcal{J} \subset \mathcal{O}_{T'}$
le faisceau quasi-cohérent d'idéaux de carré nul qui définit $T$.
La commutativité du diagramme donne $b|_T = i \circ a$, où
$i : Z \to U$ est l'immersion fermée. Nous en déduisons que
$b^\sharp(b^{-1}\mathcal{I}) \subset \mathcal{J}$ d'après
Morphismes d'espaces,
Lemme \ref{spaces-morphisms-lemma-closed-immersion-ideals}.
Comme $T'$ est un épaississement du premier ordre de $T$, nous avons $\mathcal{J}^2 = 0$,
donc $b^\sharp(b^{-1}(\mathcal{I}^2)) = 0$. D'après
Morphismes d'espaces, Lemme \ref{spaces-morphisms-lemma-closed-immersion-ideals},
il s'ensuit que $b$ se factorise par $Z'$. En notant $a' : T' \to Z'$
cette factorisation, nous concluons.
\end{proof}

\begin{lemma}
\label{spaces-more-morphisms-lemma-infinitesimal-neighbourhood-conormal}
Soit $i : Z \to X$ une immersion d'espaces algébriques.
Soit $Z \subset Z'$ le voisinage infinitésimal du premier ordre
de $Z$ dans $X$. Alors le diagramme
$$
\xymatrix{
Z \ar[r] \ar[d] & Z' \ar[d] \\
Z \ar[r] & X
}
$$
induit un morphisme de faisceaux conormaux
$\mathcal{C}_{Z/X} \to \mathcal{C}_{Z/Z'}$ d'après le
Lemme \ref{spaces-more-morphisms-lemma-conormal-functorial}.
Ce morphisme est un isomorphisme.
\end{lemma}

\begin{proof}
Cela résulte immédiatement de la construction de $Z'$ ci-dessus.
\end{proof}





\section{Transformations formellement lisses, étales et non ramifiées}
\sectionmark{Transformations lisses, étales et non ramifiées}
\label{spaces-more-morphisms-section-formally-smooth-etale-unramified}

\noindent
Rappelons qu'un homomorphisme d'anneaux $R \to A$ est dit
{\it formellement lisse}, resp.\ {\it formellement étale},
resp.\ {\it formellement non ramifié}
(voir Algèbre, Définition \ref{algebra-definition-formally-smooth},
resp.\ Définition \ref{algebra-definition-formally-etale},
resp.\ Définition \ref{algebra-definition-formally-unramified})
si, pour tout diagramme commutatif de flèches pleines
$$
\xymatrix{
A \ar[r] \ar@{-->}[rd] & B/I \\
R \ar[r] \ar[u] & B \ar[u]
}
$$
où $I \subset B$ est un idéal de carré nul, il
existe une flèche pointillée, resp.\ il en existe une unique, resp.\ il en existe
au plus une, qui rend le diagramme commutatif. Cela motive
l'analogue suivant pour les morphismes d'espaces algébriques et, plus
généralement, pour les foncteurs.

\begin{definition}
\label{spaces-more-morphisms-definition-formally-smooth-etale-unramified}
Soit $S$ un schéma.
Soit $a : F \to G$ une transformation de foncteurs
$F, G : (\Sch/S)_{fppf}^{opp} \to \textit{Ens}$.
Considérons les diagrammes commutatifs de flèches pleines de la forme
$$
\xymatrix{
F \ar[d]_a & T \ar[d]^i \ar[l] \\
G & T' \ar[l] \ar@{-->}[lu]
}
$$
où $T$ et $T'$ sont des schémas affines et $i$ une immersion fermée
définie par un idéal de carré nul.
\begin{enumerate}
\item On dit que $a$ est {\it formellement lisse} si, pour tout
diagramme de flèches pleines comme ci-dessus, il existe une flèche pointillée qui rend le diagramme
commutatif\footnote{Ce n'est là qu'une des définitions possibles.
Une condition légèrement plus faible consisterait à exiger que
la flèche pointillée existe localement pour la topologie fppf sur $T'$. Cette notion plus faible
possède, en un certain sens, de meilleures propriétés formelles.}.
\item On dit que $a$ est {\it formellement étale} si, pour tout
diagramme de flèches pleines comme ci-dessus, il existe exactement une flèche pointillée qui rend le diagramme
commutatif.
\item On dit que $a$ est {\it formellement non ramifiée} si, pour tout
diagramme de flèches pleines comme ci-dessus, il existe au plus une flèche pointillée qui rend le diagramme
commutatif.
\end{enumerate}
\end{definition}

\begin{lemma}
\label{spaces-more-morphisms-lemma-formally-etale-is-combination}
Soit $S$ un schéma.
Soit $a : F \to G$ une transformation de foncteurs
$F, G : (\Sch/S)_{fppf}^{opp} \to \textit{Ens}$.
Alors $a$ est formellement étale si et seulement si $a$ est à la fois formellement
lisse et formellement non ramifiée.
\end{lemma}

\begin{proof}
Cela résulte formellement de la définition.
\end{proof}

\begin{lemma}
\label{spaces-more-morphisms-lemma-composition-formally-smooth-etale-unramified}
Composition.
\begin{enumerate}
\item Un composé de transformations de foncteurs formellement lisses est formellement
lisse.
\item Un composé de transformations de foncteurs formellement étales est formellement
étale.
\item Un composé de transformations de foncteurs formellement non ramifiées est
formellement non ramifié.
\end{enumerate}
\end{lemma}

\begin{proof}
C'est formel.
\end{proof}

\begin{lemma}
\label{spaces-more-morphisms-lemma-base-change-formally-smooth-etale-unramified}
Soit $S$ un schéma appartenant à $\Sch_{fppf}$.
Soient $F, G, H : (\Sch/S)_{fppf}^{opp} \to \textit{Ens}$.
Soient $a : F \to G$, $b : H \to G$ des transformations de foncteurs.
Considérons le diagramme de produit fibré
$$
\xymatrix{
H \times_{b, G, a} F \ar[r]_-{b'} \ar[d]_{a'} & F \ar[d]^a \\
H \ar[r]^b & G
}
$$
\begin{enumerate}
\item Si $a$ est formellement lisse, alors le changement de base $a'$ est
formellement lisse.
\item Si $a$ est formellement étale, alors le changement de base $a'$ est
formellement étale.
\item Si $a$ est formellement non ramifiée, alors le changement de base $a'$ est
formellement non ramifiée.
\end{enumerate}
\end{lemma}

\begin{proof}
C'est formel.
\end{proof}

\begin{lemma}
\label{spaces-more-morphisms-lemma-representable-property-formally-property}
Soit $S$ un schéma.
Soient $F, G : (\Sch/S)_{fppf}^{opp} \to \textit{Ens}$.
Soit $a : F \to G$ une transformation de foncteurs représentable.
\begin{enumerate}
\item Si $a$ est lisse, alors $a$ est formellement lisse.
\item Si $a$ est étale, alors $a$ est formellement étale.
\item Si $a$ est non ramifiée, alors $a$ est formellement non ramifiée.
\end{enumerate}
\end{lemma}

\begin{proof}
Considérons un diagramme commutatif de flèches pleines
$$
\xymatrix{
F \ar[d]_a & T \ar[d]^i \ar[l] \\
G & T' \ar[l] \ar@{-->}[lu]
}
$$
comme dans la
Définition \ref{spaces-more-morphisms-definition-formally-smooth-etale-unramified}.
Alors $F \times_G T'$ est un schéma lisse (resp.\ étale, resp.\ non ramifié)
sur $T'$. Par conséquent, d'après
Compléments sur les morphismes, Lemme \ref{more-morphisms-lemma-smooth-formally-smooth}
(resp.\ Lemme \ref{more-morphisms-lemma-etale-formally-etale},
resp.\ Lemme \ref{more-morphisms-lemma-unramified-formally-unramified}),
nous pouvons tracer (resp.\ tracer de manière unique, resp.\ tracer d'au plus
une manière) la flèche pointillée du diagramme
$$
\xymatrix{
F \times_G T' \ar[d] & T \ar[d]^i \ar[l] \\
T' & T' \ar[l] \ar@{-->}[lu]
}
$$
et nous obtenons donc aussi l'assertion correspondante dans le premier diagramme.
\end{proof}

\begin{lemma}
\label{spaces-more-morphisms-lemma-etale-on-top}
Soit $S$ un schéma appartenant à $\Sch_{fppf}$.
Soient $F, G, H : (\Sch/S)_{fppf}^{opp} \to \textit{Ens}$.
Soient $a : F \to G$, $b : G \to H$ des transformations de foncteurs.
Supposons que $a$ soit représentable, surjective et étale.
\begin{enumerate}
\item Si $b$ est formellement lisse, alors $b \circ a$ est formellement lisse.
\item Si $b$ est formellement étale, alors $b \circ a$ est formellement étale.
\item Si $b$ est formellement non ramifiée, alors $b \circ a$ est formellement non ramifiée.
\end{enumerate}
Réciproquement, considérons un diagramme commutatif de flèches pleines
$$
\xymatrix{
G \ar[d]_b & T \ar[d]^i \ar[l] \\
H & T' \ar[l] \ar@{-->}[lu]
}
$$
où $T'$ est un schéma affine sur $S$
et $i : T \to T'$ une immersion fermée définie par un idéal de carré nul.
\begin{enumerate}
\item[(4)] Si $b \circ a$ est formellement lisse, alors, pour tout $t \in T$,
il existe un morphisme étale de schémas affines $U' \to T'$ et un morphisme
$U' \to G$ tels que
$$
\xymatrix{
G \ar[d]_b & T \ar[l] & T \times_{T'} U' \ar[d] \ar[l]\\
H & T' \ar[l] & U' \ar[llu] \ar[l]
}
$$
le diagramme soit commutatif et $t$ appartienne à l'image de $U' \to T'$.
\item[(5)] Si $b \circ a$ est formellement non ramifiée, alors il existe au plus
une flèche pointillée dans le diagramme ci-dessus, c'est-à-dire que $b$ est formellement non ramifiée.
\item[(6)] Si $b \circ a$ est formellement étale, alors il existe exactement une
flèche pointillée dans le diagramme ci-dessus, c'est-à-dire que $b$ est formellement étale.
\end{enumerate}
\end{lemma}

\begin{proof}
Supposons que $b$ soit formellement lisse (resp.\ formellement étale,
resp.\ formellement non ramifiée). Puisqu'un morphisme étale est à la fois
lisse et non ramifié, nous voyons que $a$ est représentable et lisse
(resp.\ étale, resp. non ramifiée). Les assertions (1), (2) et (3)
résultent donc de la combinaison du
Lemme \ref{spaces-more-morphisms-lemma-representable-property-formally-property}
et du
Lemme \ref{spaces-more-morphisms-lemma-composition-formally-smooth-etale-unramified}.

\medskip\noindent
Supposons que $b \circ a$ soit formellement lisse. Considérons un diagramme
comme dans l'énoncé du lemme. Posons $W = F \times_G T$.
Par hypothèse, $W$ est un schéma étale surjectif sur $T$. D'après
Morphismes étales, Théorème \ref{etale-theorem-remarkable-equivalence},
il existe un schéma $W'$ étale sur $T'$ tel que $W = T \times_{T'} W'$.
Choisissons un sous-schéma ouvert affine $U' \subset W'$ tel que $t$ appartienne à
l'image de $U' \to T'$. Puisque $b \circ a$ est formellement
lisse, il existe un morphisme $U' \to F$ tel que
$$
\xymatrix{
F \ar[d]_{b \circ a} & W \ar[l] & T \times_{T'} U' \ar[d] \ar[l]\\
H & T' \ar[l] & U' \ar[llu] \ar[l]
}
$$
le diagramme soit commutatif. Le composé $U' \to F \to G$ fournit un
morphisme comme dans l'assertion (5) du lemme.

\medskip\noindent
Supposons que $f, g : T' \to G$ soient deux flèches pointillées qui complètent le
diagramme du lemme. Posons $W = F \times_G T$.
Par hypothèse, $W$ est un schéma étale surjectif sur $T$. D'après
Morphismes étales, Théorème \ref{etale-theorem-remarkable-equivalence},
il existe un schéma $W'$ étale sur $T'$ tel que $W = T \times_{T'} W'$.
Comme $a$ est formellement étale, les composés
$$
W' \to T' \xrightarrow{f} G
\quad\text{et}\quad
W' \to T' \xrightarrow{g} G
$$
se relèvent en des morphismes $f', g' : W' \to F$ (on relève sur les ouverts affines et l'on recolle par
unicité). Si maintenant $b \circ a : F \to H$ est formellement non ramifiée, alors
$f' = g'$ et donc $f = g$, puisque $W' \to T'$ est un recouvrement étale. Cela démontre
l'assertion (6) du lemme.

\medskip\noindent
Supposons que $b \circ a$ soit formellement étale. D'après l'assertion (4), nous
pouvons trouver, localement pour la topologie étale sur $T'$, une flèche pointillée qui complète le diagramme,
et cette flèche est unique d'après l'assertion (5). Nous pouvons donc recoller les
solutions locales pour obtenir l'assertion (6). Certains détails sont omis.
\end{proof}

\begin{remark}
\label{spaces-more-morphisms-remark-tempting}
On pourrait être tenté de penser que, dans la situation du
Lemme \ref{spaces-more-morphisms-lemma-etale-on-top},
nous avons
« $b$ formellement lisse » $\Leftrightarrow$ « $b \circ a$ formellement lisse ».
Cependant, cela est vraisemblablement faux en général.
\end{remark}

\begin{lemma}
\label{spaces-more-morphisms-lemma-formally-permanence}
Soit $S$ un schéma.
Soient $F, G, H : (\Sch/S)_{fppf}^{opp} \to \textit{Ens}$.
Soient $a : F \to G$, $b : G \to H$ des transformations de foncteurs.
Supposons que $b$ soit formellement non ramifiée.
\begin{enumerate}
\item Si $b \circ a$ est formellement non ramifiée, alors $a$ est formellement non ramifiée.
\item Si $b \circ a$ est formellement étale, alors $a$ est formellement étale.
\item Si $b \circ a$ est formellement lisse, alors $a$ est formellement lisse.
\end{enumerate}
\end{lemma}

\begin{proof}
Soit $T \subset T'$ une immersion fermée de schémas affines définie par un idéal
de carré nul. Soient $g' : T' \to G$ et $f : T \to F$ tels que
$g'|_T = a \circ f$. Comme $b$ est formellement non ramifiée, il existe une correspondance
bijective entre
$$
\{f' : T' \to F \mid f = f'|_T\text{ et }a \circ f' = g'\}
$$
et
$$
\{f' : T' \to F \mid f = f'|_T\text{ et }b \circ a \circ f' = b \circ g'\}.
$$
Le lemme en résulte formellement.
\end{proof}








\section{Morphismes formellement non ramifiés}
\label{spaces-more-morphisms-section-formally-unramified}

\noindent
Dans cette section, nous précisons ce que signifie, pour un morphisme d'espaces algébriques,
d'être formellement non ramifié.

\begin{definition}
\label{spaces-more-morphisms-definition-formally-unramified}
Soit $S$ un schéma. Un morphisme $f : X \to Y$ d'espaces algébriques sur $S$
est dit {\it formellement non ramifié} s'il est formellement non ramifié en tant que
transformation de foncteurs au sens de la
Définition \ref{spaces-more-morphisms-definition-formally-smooth-etale-unramified}.
\end{definition}

\noindent
Nous ne répéterons pas les résultats démontrés dans le cadre plus général des
transformations de foncteurs formellement non ramifiées à la
section \ref{spaces-more-morphisms-section-formally-smooth-etale-unramified}.
On peut caractériser cette propriété par l'annulation du
module des différentielles relatives, voir le
Lemme \ref{spaces-more-morphisms-lemma-characterize-formally-unramified}.

\begin{lemma}
\label{spaces-more-morphisms-lemma-formally-unramified}
Soit $S$ un schéma. Soit $f : X \to Y$ un morphisme d'espaces algébriques sur
$S$. Les assertions suivantes sont équivalentes :
\begin{enumerate}
\item $f$ est formellement non ramifié,
\item pour tout diagramme
$$
\xymatrix{
U \ar[d] \ar[r]_\psi & V \ar[d] \\
X \ar[r]^f & Y
}
$$
où $U$ et $V$ sont des schémas et les flèches verticales sont étales,
le morphisme de schémas $\psi$ est formellement non ramifié (au sens de
Compléments sur les morphismes,
Définition \ref{more-morphisms-definition-formally-unramified}), et
\item pour un tel diagramme dont les flèches verticales sont surjectives, le morphisme
$\psi$ est formellement non ramifié.
\end{enumerate}
\end{lemma}

\begin{proof}
Supposons $f$ formellement non ramifié. D'après le
Lemme \ref{spaces-more-morphisms-lemma-representable-property-formally-property},
les morphismes $U \to X$ et $V \to Y$ sont formellement non ramifiés. D'après le
Lemme \ref{spaces-more-morphisms-lemma-composition-formally-smooth-etale-unramified},
le composé $U \to Y$ est donc formellement non ramifié. Il résulte alors du
Lemme \ref{spaces-more-morphisms-lemma-formally-permanence}
que $U \to V$ est formellement non ramifié. Ainsi, (1) implique (2), et (2)
implique trivialement (3).

\medskip\noindent
Supposons donné un diagramme comme en (3). D'après le
Lemme \ref{spaces-more-morphisms-lemma-representable-property-formally-property},
le morphisme $V \to Y$ est formellement non ramifié. D'après le
Lemme \ref{spaces-more-morphisms-lemma-composition-formally-smooth-etale-unramified},
le composé $U \to Y$ est donc formellement non ramifié. Il résulte alors du
Lemme \ref{spaces-more-morphisms-lemma-etale-on-top}
que $X \to Y$ est formellement non ramifié, c'est-à-dire que (1) est vérifiée.
\end{proof}

\begin{lemma}
\label{spaces-more-morphisms-lemma-formally-unramified-not-affine}
Soit $S$ un schéma.
Si $f : X \to Y$ est un morphisme formellement non ramifié d'espaces algébriques
sur $S$, alors, pour tout diagramme commutatif de flèches pleines
$$
\xymatrix{
X \ar[d]_f & T \ar[d]^i \ar[l] \\
Y & T' \ar[l] \ar@{-->}[lu]
}
$$
où $T \subset T'$ est un épaississement du premier ordre d'espaces algébriques
sur $S$, il existe au plus une flèche pointillée qui rend le diagramme commutatif.
Autrement dit, dans la
Définition \ref{spaces-more-morphisms-definition-formally-unramified},
on peut supprimer la condition que $T$ soit un schéma affine.
\end{lemma}

\begin{proof}
Cela tient au fait qu'il existe un morphisme étale surjectif
$U' \to T'$, où $U'$ est une réunion disjointe de schémas affines (voir
Propriétés des espaces, Lemme
\ref{spaces-properties-lemma-cover-by-union-affines}),
et qu'un morphisme $T' \to X$ est déterminé par sa restriction à $U'$.
\end{proof}

\begin{lemma}
\label{spaces-more-morphisms-lemma-composition-formally-unramified}
Un composé de morphismes formellement non ramifiés est formellement non ramifié.
\end{lemma}

\begin{proof}
C'est formel.
\end{proof}

\begin{lemma}
\label{spaces-more-morphisms-lemma-base-change-formally-unramified}
Un changement de base d'un morphisme formellement non ramifié est formellement non ramifié.
\end{lemma}

\begin{proof}
C'est formel.
\end{proof}

\begin{lemma}
\label{spaces-more-morphisms-lemma-characterize-formally-unramified}
Soit $S$ un schéma. Soit $f : X \to Y$ un morphisme d'espaces algébriques sur
$S$. Les assertions suivantes sont équivalentes :
\begin{enumerate}
\item $f$ est formellement non ramifié, et
\item $\Omega_{X/Y} = 0$.
\end{enumerate}
\end{lemma}

\begin{proof}
Cela résulte de la combinaison du
Lemme \ref{spaces-more-morphisms-lemma-formally-unramified},
de Compléments sur les morphismes,
Lemme \ref{more-morphisms-lemma-formally-unramified-differentials},
et du
Lemme \ref{spaces-more-morphisms-lemma-localize-differentials}.
\end{proof}

\begin{lemma}
\label{spaces-more-morphisms-lemma-unramified-formally-unramified}
Soit $S$ un schéma.
Soit $f : X \to Y$ un morphisme d'espaces algébriques sur $S$.
Les assertions suivantes sont équivalentes :
\begin{enumerate}
\item le morphisme $f$ est non ramifié,
\item le morphisme $f$ est localement de type fini et $\Omega_{X/Y} = 0$, et
\item le morphisme $f$ est localement de type fini et formellement non ramifié.
\end{enumerate}
\end{lemma}

\begin{proof}
Choisissons un diagramme
$$
\xymatrix{
U \ar[d] \ar[r]_\psi & V \ar[d] \\
X \ar[r]^f & Y
}
$$
où $U$ et $V$ sont des schémas et les flèches verticales sont étales et
surjectives. Nous avons alors
\begin{align*}
f\text{ non ramifié}
& \Leftrightarrow
\psi\text{ non ramifié} \\
& \Leftrightarrow
\psi\text{ localement de type fini et }\Omega_{U/V} = 0 \\
& \Leftrightarrow
f\text{ localement de type fini et }\Omega_{X/Y} = 0 \\
& \Leftrightarrow
f\text{ localement de type fini et formellement non ramifié}
\end{align*}
Nous avons utilisé ici
Morphismes, Lemme \ref{morphisms-lemma-unramified-omega-zero}, et le
Lemme \ref{spaces-more-morphisms-lemma-characterize-formally-unramified}.
\end{proof}

\begin{lemma}
\label{spaces-more-morphisms-lemma-universally-injective-unramified}
Soit $S$ un schéma.
Soit $f : X \to Y$ un morphisme d'espaces algébriques sur $S$.
Les assertions suivantes sont équivalentes :
\begin{enumerate}
\item $f$ est non ramifié et est un monomorphisme,
\item $f$ est non ramifié et universellement injectif,
\item $f$ est localement de type fini et est un monomorphisme,
\item $f$ est universellement injectif, localement de type fini et
formellement non ramifié.
\end{enumerate}
De plus, dans ce cas, $f$ est aussi représentable, séparé et
localement quasi-fini.
\end{lemma}

\begin{proof}
Nous avons vu au
Lemme \ref{spaces-more-morphisms-lemma-unramified-formally-unramified}
qu'être formellement non ramifié et localement de type fini équivaut
à être non ramifié.
Ainsi, (4) est équivalente à (2).
Un monomorphisme est certainement formellement non ramifié ; ainsi, (3) implique (4).
Il est clair que (1) implique (3). Enfin, si (2) est vérifiée, alors
$\Delta : X \to X \times_Y X$ est à la fois une immersion ouverte
(Morphismes d'espaces, Lemme
\ref{spaces-morphisms-lemma-diagonal-unramified-morphism})
et surjective
(Morphismes d'espaces, Lemme
\ref{spaces-morphisms-lemma-universally-injective}) ;
c'est donc un isomorphisme, c'est-à-dire que $f$ est un monomorphisme. Nous voyons ainsi que
(2) implique (1).
Enfin, nous voyons que $f$ est représentable, séparé et localement
quasi-fini d'après
Morphismes d'espaces, Lemmes
\ref{spaces-morphisms-lemma-monomorphism-loc-finite-type-loc-quasi-finite} et
\ref{spaces-morphisms-lemma-locally-quasi-finite-separated-representable}.
\end{proof}

\begin{lemma}
\label{spaces-more-morphisms-lemma-characterize-closed-immersion}
Soit $S$ un schéma.
Soit $f : X \to Y$ un morphisme d'espaces algébriques sur $S$.
Les assertions suivantes sont équivalentes :
\begin{enumerate}
\item $f$ est une immersion fermée,
\item $f$ est universellement fermé, non ramifié et est un monomorphisme,
\item $f$ est universellement fermé, non ramifié et universellement injectif,
\item $f$ est universellement fermé, localement de type fini et est un monomorphisme,
\item $f$ est universellement fermé, universellement injectif, localement de
type fini et formellement non ramifié.
\end{enumerate}
\end{lemma}

\begin{proof}
L'équivalence de (2) -- (5) résulte immédiatement du
Lemme \ref{spaces-more-morphisms-lemma-universally-injective-unramified}.
De plus, si (2) -- (5) sont vérifiées, alors $f$ est représentable.
De même, si (1) est vérifiée, alors $f$ est représentable.
Le résultat découle donc du cas des schémas, voir
Morphismes étales, Lemme \ref{etale-lemma-characterize-closed-immersion}.
\end{proof}





\section{Épaississements universels du premier ordre}
\label{spaces-more-morphisms-section-universal-thickening}

\noindent
Soit $S$ un schéma.
Soit $h : Z \to X$ un morphisme d'espaces algébriques sur $S$.
Un {\it épaississement universel du premier ordre} de $Z$ sur $X$ est un
épaississement du premier ordre $Z \subset Z'$ sur $X$ tel que, pour
tout épaississement du premier ordre $T \subset T'$
sur $X$ et tout diagramme commutatif de flèches pleines
\begin{equation}
\label{spaces-more-morphisms-equation-universal-first-order-thickening}
\vcenter{
\xymatrix{
& Z \ar[ld] & & T \ar[rd] \ar[ll]^a \\
Z' \ar[rrd] & & & & T' \ar@{..>}[llll]_{a'} \ar[lld]^b \\
 & & X
}
}
\end{equation}
il existe une unique flèche pointillée qui rend le diagramme commutatif.
Notons que, dans cette situation, $(a, a') : (T \subset T') \to (Z \subset Z')$
est un morphisme d'épaississements sur $X$. Ainsi, s'il existe un épaississement universel
du premier ordre, celui-ci est unique à isomorphisme unique près.
En général, un épaississement universel du premier ordre
n'existe pas, mais il existe lorsque $h$ est formellement non ramifié.
Avant de le démontrer, montrons qu'un épaississement universel du premier ordre
dans la catégorie des schémas l'est également dans la
catégorie des espaces algébriques.

\begin{lemma}
\label{spaces-more-morphisms-lemma-check-universal-first-order-thickening}
Soit $S$ un schéma.
Soit $h : Z \to X$ un morphisme d'espaces algébriques sur $S$.
Soit $Z \subset Z'$ un épaississement du premier ordre sur $X$.
Les assertions suivantes sont équivalentes :
\begin{enumerate}
\item $Z \subset Z'$ est un épaississement universel du premier ordre,
\item pour tout diagramme (\ref{spaces-more-morphisms-equation-universal-first-order-thickening})
où $T'$ est un schéma, il existe une unique flèche pointillée qui rend le diagramme commutatif, et
\item pour tout diagramme (\ref{spaces-more-morphisms-equation-universal-first-order-thickening})
où $T'$ est un schéma affine, il existe une unique flèche pointillée qui rend le
diagramme commutatif.
\end{enumerate}
\end{lemma}

\begin{proof}
Les implications (1) $\Rightarrow$ (2) $\Rightarrow$ (3) sont formelles.
Supposons (3) et considérons un diagramme arbitraire
(\ref{spaces-more-morphisms-equation-universal-first-order-thickening}).
Choisissons une présentation $T' = U'/R'$, voir
Espaces, Définition \ref{spaces-definition-presentation}.
Nous pouvons supposer que $U' = \coprod U'_i$ est une réunion disjointe
de schémas affines, de sorte que $R' = U' \times_{T'} U' = \coprod_{i, j} U'_i \times_T' U'_j$.
Pour chaque paire $(i, j)$, choisissons un recouvrement ouvert affine
$U'_i \times_T' U'_j = \bigcup_k R'_{ijk}$. Notons $U_i, R_{ijk}$
les produits fibrés avec $T$ au-dessus de $T'$. Alors chacun des couples
$U_i \subset U'_i$ et $R_{ijk} \subset R'_{ijk}$
est un épaississement du premier ordre de schémas affines.
Notons $a_i : U_i \to Z$, resp.\ $a_{ijk} : R_{ijk} \to Z$,
le composé de $a : T \to Z$ avec le morphisme
$U_i \to T$, resp.\ $R_{ijk} \to T$.
En appliquant (3) à $a_i : U_i \to Z$,
nous obtenons des morphismes uniques $a'_i : U'_i \to Z'$.
En appliquant (3) à $a_{ijk}$, nous voyons que les deux composés
$R'_{ijk} \to R'_i \to Z'$ et $R'_{ijk} \to R'_j \to Z'$
sont égaux. Ainsi, $a' = \coprod a'_i : U' = \coprod U'_i \to Z'$
descend au faisceau quotient $T' = U'/R'$, ce qui conclut.
\end{proof}

\begin{lemma}
\label{spaces-more-morphisms-lemma-universal-thickening-over-formally-etale}
Soit $S$ un schéma.
Soient $Z \to Y \to X$ des morphismes d'espaces algébriques sur $S$.
Si $Z \subset Z'$ est un épaississement universel du premier ordre de
$Z$ sur $Y$ et si $Y \to X$ est formellement étale, alors $Z \subset Z'$ est
un épaississement universel du premier ordre de $Z$ sur $X$.
\end{lemma}

\begin{proof}
Cela est formel. En effet, d'après le
Lemme \ref{spaces-more-morphisms-lemma-check-universal-first-order-thickening},
il suffit de considérer les diagrammes commutatifs de flèches pleines
(\ref{spaces-more-morphisms-equation-universal-first-order-thickening})
où $T'$ est un schéma affine. Le composé
$T \to Z \to Y$ se relève de manière unique en $T' \to Y$, puisque $Y \to X$ est
supposé formellement étale. Par conséquent, le fait que
$Z \subset Z'$ soit un épaississement universel du premier ordre sur $Y$
fournit le morphisme voulu $a' : T' \to Z'$.
\end{proof}

\begin{lemma}
\label{spaces-more-morphisms-lemma-etale-morphism-of-universal-thickenings}
Soit $S$ un schéma.
Soient $Z \to Y \to X$ des morphismes d'espaces algébriques sur $S$.
Supposons que $Z \to Y$ soit étale.
\begin{enumerate}
\item Si $Y \subset Y'$ est un épaississement universel du premier ordre de
$Y$ sur $X$, alors l'unique morphisme étale $Z' \to Y'$ tel
que $Z = Y \times_{Y'} Z'$ (voir le
Théorème \ref{spaces-more-morphisms-theorem-topological-invariance})
est un épaississement universel du premier ordre de $Z$ sur $X$.
\item Si $Z \to Y$ est surjectif et si
$(Z \subset Z') \to (Y \subset Y')$ est un morphisme étale
d'épaississements du premier ordre sur $X$, tandis que $Z'$ est un épaississement universel du
premier ordre de $Z$ sur $X$, alors $Y'$ est un épaississement universel du
premier ordre de $Y$ sur $X$.
\end{enumerate}
\end{lemma}

\begin{proof}
Démonstration de (1). D'après le
Lemme \ref{spaces-more-morphisms-lemma-check-universal-first-order-thickening},
il suffit de considérer les diagrammes commutatifs de flèches pleines
(\ref{spaces-more-morphisms-equation-universal-first-order-thickening})
où $T'$ est un schéma affine. Le composé
$T \to Z \to Y$ se relève de manière unique en $T' \to Y'$, puisque $Y'$ est
l'épaississement universel du premier ordre. Le fait que
$Z' \to Y'$ soit étale implique alors (voir le
Lemme \ref{spaces-more-morphisms-lemma-representable-property-formally-property})
que $T' \to Y'$ se relève en le
morphisme voulu $a' : T' \to Z'$.

\medskip\noindent
Démonstration de (2). Soit $T \subset T'$ un épaississement du premier ordre sur
$X$ et soit $a : T \to Y$ un morphisme. Posons $W = T \times_Y Z$
et notons $c : W \to Z$ la projection.
Soit $W' \to T'$ l'unique morphisme étale tel que
$W = T \times_{T'} W'$, voir le
Théorème \ref{spaces-more-morphisms-theorem-topological-invariance}.
Notons que $W' \to T'$ est surjectif puisque $Z \to Y$ est surjectif.
Par hypothèse, nous obtenons un unique morphisme $c' : W' \to Z'$
sur $X$ et dont $c$ est la restriction à $W$. Par unicité, les deux restrictions
de $c'$ à $W' \times_{T'} W'$ sont égales (puisque les deux restrictions de
$c$ à $W \times_T W$ sont égales). Ainsi, $c'$ descend en un unique
morphisme $a' : T' \to Y'$, ce qui conclut.
\end{proof}

\begin{lemma}
\label{spaces-more-morphisms-lemma-universal-thickening}
Soit $S$ un schéma.
Soit $h : Z \to X$ un morphisme formellement non ramifié d'espaces
algébriques sur $S$.
Il existe un épaississement universel du premier ordre $Z \subset Z'$ de
$Z$ sur $X$.
\end{lemma}

\begin{proof}
Choisissons un diagramme commutatif quelconque
$$
\xymatrix{
V \ar[d] \ar[r] & U \ar[d] \\
Z \ar[r] & X
}
$$
où $V$ et $U$ sont des schémas et les flèches verticales sont étales.
Notons que $V \to U$ est un morphisme formellement non ramifié de schémas, voir le
Lemme \ref{spaces-more-morphisms-lemma-formally-unramified}.
En combinant le
Lemme \ref{spaces-more-morphisms-lemma-check-universal-first-order-thickening}
et
Compléments sur les morphismes, Lemme \ref{more-morphisms-lemma-universal-thickening},
nous voyons qu'il existe un épaississement universel du premier ordre $V \subset V'$
de $V$ sur $U$. D'après le
Lemme \ref{spaces-more-morphisms-lemma-universal-thickening-over-formally-etale}, partie (1),
$V'$ est un épaississement universel du premier ordre de $V$ sur $X$.

\medskip\noindent
Fixons un schéma $U$ et un morphisme étale surjectif $U \to X$.
L'argument ci-dessus montre que, pour tout $V \to Z$ étale, où $V$ est
un schéma tel que $V \to Z \to X$ se factorise par $U$, il existe un
épaississement universel du premier ordre $V \subset V'$ de $V$ sur $X$
(qui ne dépend toutefois pas de la factorisation choisie de $V \to X$
par $U$). Nous pouvons maintenant choisir $V$ tel que $V \to Z$ soit étale
surjectif (voir
Espaces, Lemme \ref{spaces-lemma-lift-morphism-presentations}).
Alors $R = V \times_Z V$ est un schéma étale sur $Z$ tel que
$R \to X$ se factorise lui aussi par $U$.
Nous obtenons donc des épaississements universels du premier ordre
$V \subset V'$ et $R \subset R'$ sur $X$.
Puisque $V \subset V'$ est un épaississement universel du premier ordre,
les deux projections $s, t : R \to V$ se relèvent en des morphismes
$s', t': R' \to V'$. D'après le
Lemme \ref{spaces-more-morphisms-lemma-etale-morphism-of-universal-thickenings},
puisque $R'$ est l'épaississement universel du premier ordre de $R$ sur $X$,
ces morphismes sont étales.
Alors $(t', s') : R' \to V'$ est une relation d'équivalence étale,
et nous pouvons poser $Z' = V'/R'$. Puisque $V' \to Z'$ est étale surjectif
et que $v'$ est l'épaississement universel du premier ordre de $V$ sur $X$,
nous déduisons du
Lemme \ref{spaces-more-morphisms-lemma-universal-thickening-over-formally-etale}, partie (2),
que $Z'$ est un épaississement universel du premier ordre de $Z$ sur $X$.
\end{proof}

\begin{definition}
\label{spaces-more-morphisms-definition-universal-thickening}
Soit $S$ un schéma.
Soit $h : Z \to X$ un morphisme formellement non ramifié
d'espaces algébriques sur $S$.
\begin{enumerate}
\item L'{\it épaississement universel du premier ordre} de $Z$ sur $X$
est l'épaississement $Z \subset Z'$ construit au
Lemme \ref{spaces-more-morphisms-lemma-universal-thickening}.
\item Le {\it faisceau conormal de $Z$ sur $X$} est le faisceau conormal
de $Z$ dans son épaississement universel du premier ordre $Z'$ sur $X$.
\end{enumerate}
Dans cette situation, nous notons souvent le faisceau conormal $\mathcal{C}_{Z/X}$.
\end{definition}

\noindent
Nous voyons ainsi qu'il existe une suite exacte courte de faisceaux
$$
0 \to \mathcal{C}_{Z/X} \to \mathcal{O}_{Z'} \to \mathcal{O}_Z \to 0
$$
sur $Z_\etale$, et $\mathcal{C}_{Z/X}$ est un $\mathcal{O}_Z$-module
quasi-cohérent. Le lemme suivant montre qu'il n'y a pas de
conflit entre cette définition et celle du cas où $Z \to X$
est une immersion.

\begin{lemma}
\label{spaces-more-morphisms-lemma-immersion-universal-thickening}
Soit $S$ un schéma.
Soit $i : Z \to X$ une immersion d'espaces algébriques sur $S$. Alors
\begin{enumerate}
\item $i$ est formellement non ramifié,
\item l'épaississement universel du premier ordre de $Z$ sur $X$ est le voisinage
infinitésimal du premier ordre de $Z$ dans $X$ de la
Définition \ref{spaces-more-morphisms-definition-first-order-infinitesimal-neighbourhood},
\item le faisceau conormal de $i$ au sens de la
Définition \ref{spaces-more-morphisms-definition-conormal-sheaf}
coïncide avec le faisceau conormal de $i$ au sens de la
Définition \ref{spaces-more-morphisms-definition-universal-thickening}.
\end{enumerate}
\end{lemma}

\begin{proof}
Une immersion d'espaces algébriques est, par définition, un morphisme représentable.
Par conséquent, d'après
Morphismes, Lemmes \ref{morphisms-lemma-open-immersion-unramified} et
\ref{morphisms-lemma-closed-immersion-unramified},
une immersion est non ramifiée (par le principe abstrait du chapitre
Espaces, Lemme
\ref{spaces-lemma-representable-transformations-property-implication}).
Elle est donc formellement non ramifiée d'après le
Lemme \ref{spaces-more-morphisms-lemma-unramified-formally-unramified}.
Les autres assertions résultent de la combinaison des
Lemmes \ref{spaces-more-morphisms-lemma-first-order-infinitesimal-neighbourhood} et
\ref{spaces-more-morphisms-lemma-infinitesimal-neighbourhood-conormal}
avec les définitions.
\end{proof}

\begin{lemma}
\label{spaces-more-morphisms-lemma-universal-thickening-unramified}
Soit $S$ un schéma.
Soit $Z \to X$ un morphisme formellement non ramifié d'espaces algébriques sur $S$.
Alors l'épaississement universel du premier ordre $Z'$ est formellement
non ramifié sur $X$.
\end{lemma}

\begin{proof}
Soit $T \subset T'$ un épaississement du premier ordre de schémas affines sur $X$.
Soit
$$
\xymatrix{
Z' \ar[d] & T \ar[l]^c \ar[d] \\
X & T' \ar[l] \ar[lu]^{a, b}
}
$$
un diagramme commutatif. Posons $T_0 = c^{-1}(Z) \subset T$ et
$T'_a = a^{-1}(Z)$ (au sens schématique).
Comme $Z'$ est un épaississement du premier ordre de $Z$, nous voyons que $T'$
est un épaississement du premier ordre de $T'_a$. De plus, puisque $c = a|_T$, nous avons
$T_0 = T \cap T'_a$ (au sens schématique). Comme $T'$ est un épaississement du premier ordre
de $T$, il s'ensuit que $T'_a$
est un épaississement du premier ordre de $T_0$. Or $a|_{T'_a}$ et $b|_{T'_a}$
sont des morphismes de $T'_a$ vers $Z'$ sur $X$ qui coïncident sur $T_0$ comme
morphismes vers $Z$. La propriété universelle de $Z'$ donne donc
$a|_{T'_a} = b|_{T'_a}$. Ainsi, $a$ et $b$ sont des morphismes depuis
l'épaississement du premier ordre $T'$ de $T'_a$ dont les restrictions à
$T'_a$ coïncident comme morphismes vers $Z$. En utilisant une nouvelle fois la propriété universelle de
$Z'$, nous obtenons $a = b$. Autrement dit, la propriété qui définit
un morphisme formellement non ramifié est vérifiée par $Z' \to X$, comme souhaité.
\end{proof}

\begin{lemma}
\label{spaces-more-morphisms-lemma-universal-thickening-functorial}
Soit $S$ un schéma.
Considérons un diagramme commutatif d'espaces algébriques sur $S$
$$
\xymatrix{
Z \ar[r]_h \ar[d]_f & X \ar[d]^g \\
W \ar[r]^{h'} & Y
}
$$
où $h$ et $h'$ sont formellement non ramifiés. Soit $Z \subset Z'$ l'épaississement universel
du premier ordre de $Z$ sur $X$. Soit $W \subset W'$ l'épaississement universel
du premier ordre de $W$ sur $Y$. Il existe un morphisme canonique
$(f, f') : (Z, Z') \to (W, W')$ d'épaississements sur $Y$ qui s'insère dans
le diagramme commutatif suivant
$$
\xymatrix{
& & & Z' \ar[ld] \ar[d]^{f'} \\
Z \ar[rr] \ar[d]_f \ar[rrru] & & X \ar[d] & W' \ar[ld] \\
W \ar[rrru]|!{[rr];[rruu]}\hole \ar[rr] & & Y
}
$$
En particulier, le morphisme $(f, f')$ d'épaississements induit un morphisme
de faisceaux conormaux $f^*\mathcal{C}_{W/Y} \to \mathcal{C}_{Z/X}$.
\end{lemma}

\begin{proof}
La première assertion résulte immédiatement de la propriété universelle de $W'$.
Le morphisme induit sur les faisceaux conormaux est celui du
Lemme \ref{spaces-more-morphisms-lemma-conormal-functorial}
appliqué à $(Z \subset Z') \to (W \subset W')$.
\end{proof}

\begin{lemma}
\label{spaces-more-morphisms-lemma-universal-thickening-fibre-product}
Soit $S$ un schéma. Soit
$$
\xymatrix{
Z \ar[r]_h \ar[d]_f & X \ar[d]^g \\
W \ar[r]^{h'} & Y
}
$$
un diagramme de produit fibré d'espaces algébriques sur $S$, où
$h'$ est formellement non ramifié. Alors $h$ est formellement non ramifié et, si
$W \subset W'$ est l'épaississement universel du premier ordre de $W$ sur $Y$,
alors $Z = X \times_Y W \subset X \times_Y W'$ est l'épaississement universel
du premier ordre de $Z$ sur $X$. En particulier, le morphisme canonique
$f^*\mathcal{C}_{W/Y} \to \mathcal{C}_{Z/X}$ du
Lemme \ref{spaces-more-morphisms-lemma-universal-thickening-functorial}
est surjectif.
\end{lemma}

\begin{proof}
Le morphisme $h$ est formellement non ramifié d'après le
Lemme \ref{spaces-more-morphisms-lemma-base-change-formally-unramified}.
Il est clair que $X \times_Y W'$ est un épaississement du premier ordre.
On vérifie immédiatement qu'il possède la propriété universelle,
car $W'$ la possède (d'après les propriétés universelles des
produits fibrés). Voir le
Lemme \ref{spaces-more-morphisms-lemma-conormal-functorial-flat}
pour comprendre pourquoi cela implique la surjectivité du morphisme de faisceaux conormaux.
\end{proof}

\begin{lemma}
\label{spaces-more-morphisms-lemma-universal-thickening-fibre-product-flat}
Soit $S$ un schéma. Soit
$$
\xymatrix{
Z \ar[r]_h \ar[d]_f & X \ar[d]^g \\
W \ar[r]^{h'} & Y
}
$$
un diagramme de produit fibré d'espaces algébriques sur $S$, où
$h'$ est formellement non ramifié et $g$ plat. Dans ce cas, le morphisme correspondant
$Z' \to W'$ entre les épaississements universels du premier ordre est plat, et
$f^*\mathcal{C}_{W/Y} \to \mathcal{C}_{Z/X}$ est un isomorphisme.
\end{lemma}

\begin{proof}
La platitude est préservée par changement de base, voir
Morphismes d'espaces, Lemme \ref{spaces-morphisms-lemma-base-change-flat}.
La première assertion découle donc de la description de $W'$ au
Lemme \ref{spaces-more-morphisms-lemma-universal-thickening-fibre-product}.
Il est clair que $X \times_Y W'$ est un épaississement du premier ordre.
On vérifie immédiatement qu'il possède la propriété universelle,
car $W'$ la possède (d'après les propriétés universelles des
produits fibrés). Voir le
Lemme \ref{spaces-more-morphisms-lemma-conormal-functorial-flat}
pour comprendre pourquoi cela implique que le morphisme de faisceaux conormaux est un isomorphisme.
\end{proof}

\begin{lemma}
\label{spaces-more-morphisms-lemma-universal-thickening-localize}
La formation des épaississements universels du premier ordre commute à la
localisation étale. Plus précisément, soit $h : Z \to X$ un morphisme formellement non ramifié
d'espaces algébriques sur un schéma de base $S$.
Considérons
$$
\xymatrix{
V \ar[d] \ar[r] & U \ar[d] \\
Z \ar[r] & X
}
$$
un diagramme commutatif dont les flèches verticales sont étales.
Soit $Z'$ l'épaississement universel du premier ordre de $Z$ sur $X$.
Alors $V \to U$ est formellement non ramifié, et l'épaississement universel du premier
ordre $V'$ de $V$ sur $U$ est étale sur $Z'$.
En particulier, $\mathcal{C}_{Z/X}|_V = \mathcal{C}_{V/U}$.
\end{lemma}

\begin{proof}
La première assertion résulte du
Lemme \ref{spaces-more-morphisms-lemma-formally-unramified}.
La compatibilité des épaississements universels du premier ordre
résulte des
Lemmes \ref{spaces-more-morphisms-lemma-universal-thickening-over-formally-etale} et
\ref{spaces-more-morphisms-lemma-etale-morphism-of-universal-thickenings}.
\end{proof}

\begin{lemma}
\label{spaces-more-morphisms-lemma-differentials-universally-unramified}
Soit $S$ un schéma. Soit $B$ un espace algébrique sur $S$.
Soit $h : Z \to X$ un morphisme formellement non ramifié d'espaces algébriques
sur $B$. Soit $Z \subset Z'$ l'épaississement universel du premier ordre de $Z$
sur $X$, de morphisme structural $h' : Z' \to X$. Le morphisme canonique
$$
\text{d}h' : (h')^*\Omega_{X/B} \to \Omega_{Z'/B}
$$
induit un isomorphisme
$h^*\Omega_{X/B} \to \Omega_{Z'/B} \otimes \mathcal{O}_Z$.
\end{lemma}

\begin{proof}
Le morphisme $c_{h'}$ est celui qui est défini au
Lemme \ref{spaces-more-morphisms-lemma-functoriality-differentials}.
Si $i : Z \to Z'$ est l'immersion fermée donnée, alors
$i^*c_{h'}$ est un morphisme
$h^*\Omega_{X/S} \to \Omega_{Z'/S} \otimes \mathcal{O}_Z$.
Pour vérifier qu'il s'agit d'un isomorphisme, on se ramène au cas des schémas
par localisation étale, voir le
Lemme \ref{spaces-more-morphisms-lemma-universal-thickening-localize}
et le
Lemme \ref{spaces-more-morphisms-lemma-localize-differentials}.
Dans ce cas, le résultat est
Compléments sur les morphismes,
Lemme \ref{more-morphisms-lemma-differentials-universally-unramified}.
\end{proof}

\begin{lemma}
\label{spaces-more-morphisms-lemma-universally-unramified-differentials-sequence}
Soit $S$ un schéma. Soit $B$ un espace algébrique sur $S$.
Soit $h : Z \to X$ un morphisme formellement non ramifié d'espaces
algébriques sur $B$.
Il existe une suite exacte canonique
$$
\mathcal{C}_{Z/X} \to h^*\Omega_{X/B} \to \Omega_{Z/B} \to 0.
$$
La première flèche est induite par $\text{d}_{Z'/B}$, où
$Z'$ est le voisinage universel du premier ordre de $Z$ sur $X$.
\end{lemma}

\begin{proof}
Nous savons qu'il existe une suite exacte canonique
$$
\mathcal{C}_{Z/Z'} \to
\Omega_{Z'/S} \otimes \mathcal{O}_Z \to
\Omega_{Z/S} \to 0.
$$
Cela résulte du
Lemme \ref{spaces-more-morphisms-lemma-differentials-relative-immersion}.
Le résultat s'obtient donc en appliquant le
Lemme \ref{spaces-more-morphisms-lemma-differentials-universally-unramified}.
\end{proof}

\begin{lemma}
\label{spaces-more-morphisms-lemma-two-unramified-morphisms}
Soit $S$ un schéma. Soit
$$
\xymatrix{
Z \ar[r]_i \ar[rd]_j & X \ar[d] \\
& Y
}
$$
un diagramme commutatif d'espaces algébriques sur $S$,
où $i$ et $j$ sont formellement non ramifiés. Il existe alors une
suite exacte canonique
$$
\mathcal{C}_{Z/Y} \to
\mathcal{C}_{Z/X} \to
i^*\Omega_{X/Y} \to 0
$$
où la première flèche provient du
Lemme \ref{spaces-more-morphisms-lemma-universal-thickening-functorial}
et la seconde du
Lemme \ref{spaces-more-morphisms-lemma-universally-unramified-differentials-sequence}.
\end{lemma}

\begin{proof}
Les morphismes étant définis, la vérification de l'exactitude de la suite
se ramène au cas des schémas par localisation étale, voir le
Lemme \ref{spaces-more-morphisms-lemma-universal-thickening-localize}
et le
Lemme \ref{spaces-more-morphisms-lemma-localize-differentials}.
Dans ce cas, le résultat est
Compléments sur les morphismes,
Lemme \ref{more-morphisms-lemma-two-unramified-morphisms}.
\end{proof}

\begin{lemma}
\label{spaces-more-morphisms-lemma-transitivity-conormal-unramified}
Soit $S$ un schéma.
Soient $Z \to Y \to X$ des morphismes formellement non ramifiés
d'espaces algébriques sur $S$.
\begin{enumerate}
\item Si $Z \subset Z'$ est l'épaississement universel du premier ordre de $Z$
sur $X$ et si $Y \subset Y'$ est l'épaississement universel du premier ordre de $Y$
sur $X$, alors il existe un morphisme $Z' \to Y'$ et $Y \times_{Y'} Z'$ est
l'épaississement universel du premier ordre de $Z$ sur $Y$.
\item Il existe une suite exacte canonique
$$
i^*\mathcal{C}_{Y/X} \to
\mathcal{C}_{Z/X} \to
\mathcal{C}_{Z/Y} \to 0
$$
où les morphismes proviennent du
Lemme \ref{spaces-more-morphisms-lemma-universal-thickening-functorial}
et $i : Z \to Y$ est le premier morphisme.
\end{enumerate}
\end{lemma}

\begin{proof}
Le morphisme $h : Z' \to Y'$ de (1) provient du
Lemme \ref{spaces-more-morphisms-lemma-universal-thickening-functorial}.
L'assertion selon laquelle $Y \times_{Y'} Z'$ est l'épaississement universel du premier ordre
de $Z$ sur $Y$ résulte immédiatement des propriétés universelles
de $Z'$ et de $Y'$. D'après le
Lemme \ref{spaces-more-morphisms-lemma-transitivity-conormal},
nous avons une suite exacte
$$
(i')^*\mathcal{C}_{Y \times_{Y'} Z'/Z'} \to
\mathcal{C}_{Z/Z'} \to
\mathcal{C}_{Z/Y \times_{Y'} Z'} \to 0
$$
où $i' : Z \to Y \times_{Y'} Z'$ est le morphisme donné. D'après le
Lemme \ref{spaces-more-morphisms-lemma-conormal-functorial-flat},
il existe une surjection
$h^*\mathcal{C}_{Y/Y'} \to \mathcal{C}_{Y \times_{Y'} Z'/Z'}$.
Combinée aux égalités
$\mathcal{C}_{Y/Y'} = \mathcal{C}_{Y/X}$,
$\mathcal{C}_{Z/Z'} = \mathcal{C}_{Z/X}$, et
$\mathcal{C}_{Z/Y \times_{Y'} Z'} = \mathcal{C}_{Z/Y}$,
cette surjection démontre le lemme.
\end{proof}













\section{Morphismes formellement étales}
\label{spaces-more-morphisms-section-formally-etale}

\noindent
Dans cette section, nous explicitons ce que signifie pour un morphisme d'espaces algébriques
d'être formellement étale.

\begin{definition}
\label{spaces-more-morphisms-definition-formally-etale}
Soit $S$ un schéma. Un morphisme $f : X \to Y$ d'espaces algébriques sur $S$
est dit {\it formellement étale} s'il est formellement étale en tant que
transformation de foncteurs au sens de la
Définition \ref{spaces-more-morphisms-definition-formally-smooth-etale-unramified}.
\end{definition}

\noindent
Nous ne redonnerons pas les résultats démontrés dans le cadre plus général des
transformations de foncteurs formellement étales dans la
section \ref{spaces-more-morphisms-section-formally-smooth-etale-unramified}.

\begin{lemma}
\label{spaces-more-morphisms-lemma-formally-etale}
Soit $S$ un schéma. Soit $f : X \to Y$ un morphisme d'espaces algébriques sur
$S$. Les assertions suivantes sont équivalentes :
\begin{enumerate}
\item $f$ est formellement étale,
\item pour tout diagramme
$$
\xymatrix{
U \ar[d] \ar[r]_\psi & V \ar[d] \\
X \ar[r]^f & Y
}
$$
où $U$ et $V$ sont des schémas et où les flèches verticales sont étales,
le morphisme de schémas $\psi$ est formellement étale (au sens de
Compléments sur les morphismes,
Définition \ref{more-morphisms-definition-formally-etale}), et
\item il existe un tel diagramme dont les flèches verticales sont surjectives et dans lequel
$\psi$ est formellement étale.
\end{enumerate}
\end{lemma}

\begin{proof}
Supposons $f$ formellement étale. D'après le
Lemme \ref{spaces-more-morphisms-lemma-representable-property-formally-property},
les morphismes $U \to X$ et $V \to Y$ sont formellement étales. Ainsi, d'après le
Lemme \ref{spaces-more-morphisms-lemma-composition-formally-smooth-etale-unramified},
la composée $U \to Y$ est formellement étale. Il résulte alors du
Lemme \ref{spaces-more-morphisms-lemma-formally-permanence}
que $U \to V$ est formellement étale. Ainsi, (1) implique (2), et (2)
implique trivialement (3).

\medskip\noindent
Supposons donné un diagramme comme en (3). D'après le
Lemme \ref{spaces-more-morphisms-lemma-representable-property-formally-property},
le morphisme $V \to Y$ est formellement étale. Ainsi, d'après le
Lemme \ref{spaces-more-morphisms-lemma-composition-formally-smooth-etale-unramified},
la composée $U \to Y$ est formellement étale. Il résulte alors du
Lemme \ref{spaces-more-morphisms-lemma-etale-on-top}
que $X \to Y$ est formellement étale, c'est-à-dire que (1) est vérifiée.
\end{proof}

\begin{lemma}
\label{spaces-more-morphisms-lemma-formally-etale-not-affine}
Soit $S$ un schéma.
Soit $f : X \to Y$ un morphisme formellement étale d'espaces algébriques sur $S$.
Alors, pour tout diagramme commutatif de flèches pleines
$$
\xymatrix{
X \ar[d]_f & T \ar[d]^i \ar[l]_a \\
Y & T' \ar[l] \ar@{-->}[lu]
}
$$
où $T \subset T'$ est un épaississement du premier ordre d'espaces algébriques
sur $Y$, il existe une unique flèche pointillée qui rend le diagramme commutatif.
Autrement dit, dans la
Définition \ref{spaces-more-morphisms-definition-formally-etale},
on peut supprimer la condition que $T$ soit affine.
\end{lemma}

\begin{proof}
Soit $U' \to T'$ un morphisme étale surjectif, où $U' = \coprod U'_i$
est une somme disjointe de schémas affines. Posons
$U_i = T \times_{T'} U'_i$. On obtient alors des morphismes
$a'_i : U'_i \to X$ tels que $a'_i|_{U_i}$ soit égal à la composée
$U_i \to T \to X$. Par unicité (voir le
Lemme \ref{spaces-more-morphisms-lemma-formally-unramified-not-affine}),
on voit que $a'_i$ et $a'_j$ coïncident sur le produit fibré
$U'_i \times_{T'} U'_j$. Par suite, $\coprod a'_i : U' \to X$
descend en un unique morphisme $a' : T' \to X$.
\end{proof}

\begin{lemma}
\label{spaces-more-morphisms-lemma-composition-formally-etale}
La composée de morphismes formellement étales est formellement étale.
\end{lemma}

\begin{proof}
C'est formel.
\end{proof}

\begin{lemma}
\label{spaces-more-morphisms-lemma-base-change-formally-etale}
Un changement de base d'un morphisme formellement étale est formellement étale.
\end{lemma}

\begin{proof}
C'est formel.
\end{proof}

\begin{lemma}
\label{spaces-more-morphisms-lemma-characterize-formally-etale}
Soit $S$ un schéma.
Soit $f : X \to Y$ un morphisme d'espaces algébriques sur $S$.
Les assertions suivantes sont équivalentes :
\begin{enumerate}
\item $f$ est formellement étale,
\item $f$ est formellement non ramifié et l'épaississement universel du premier ordre
de $X$ sur $Y$ est égal à $X$,
\item $f$ est formellement non ramifié et $\mathcal{C}_{X/Y} = 0$, et
\item $\Omega_{X/Y} = 0$ et $\mathcal{C}_{X/Y} = 0$.
\end{enumerate}
\end{lemma}

\begin{proof}
En réalité, la dernière assertion n'a de sens que parce que $\Omega_{X/Y} = 0$
implique que $\mathcal{C}_{X/Y}$ est défini par le
Lemme \ref{spaces-more-morphisms-lemma-characterize-formally-unramified}
et la
Définition \ref{spaces-more-morphisms-definition-universal-thickening}.
Cela montre aussi que (3) et (4) sont équivalentes.

\medskip\noindent
Chacune des hypothèses (1), (2) et (3) implique que $f$ est formellement
non ramifié. On peut donc supposer $f$ formellement non ramifié. L'équivalence
de (1), (2) et (3) résulte de la propriété universelle de l'épaississement
universel du premier ordre $X'$ de $X$ sur $S$ et du fait que
$X = X' \Leftrightarrow \mathcal{C}_{X/Y} = 0$, puisque,
par définition, $\mathcal{C}_{X/Y} = \mathcal{C}_{X/X'}$
est le faisceau d'idéaux de $X$ dans $X'$.
\end{proof}

\begin{lemma}
\label{spaces-more-morphisms-lemma-unramified-flat-formally-etale}
Un morphisme plat et non ramifié est formellement étale.
\end{lemma}

\begin{proof}
Cela résulte du cas des schémas, voir
Compléments sur les morphismes,
Lemme \ref{more-morphisms-lemma-unramified-flat-formally-etale},
et de la localisation étale, voir les
Lemmes \ref{spaces-more-morphisms-lemma-formally-unramified} et \ref{spaces-more-morphisms-lemma-formally-etale},
ainsi que
Morphismes d'espaces, Lemme \ref{spaces-morphisms-lemma-flat-local}.
\end{proof}

\begin{lemma}
\label{spaces-more-morphisms-lemma-etale-formally-etale}
Soit $S$ un schéma.
Soit $f : X \to Y$ un morphisme d'espaces algébriques sur $S$.
Les assertions suivantes sont équivalentes :
\begin{enumerate}
\item le morphisme $f$ est étale, et
\item le morphisme $f$ est localement de présentation finie et
formellement étale.
\end{enumerate}
\end{lemma}

\begin{proof}
Cela résulte du cas des schémas, voir
Compléments sur les morphismes,
Lemme \ref{more-morphisms-lemma-etale-formally-etale},
et de la localisation étale, voir le
Lemme \ref{spaces-more-morphisms-lemma-formally-etale},
ainsi que
Morphismes d'espaces,
Lemmes \ref{spaces-morphisms-lemma-finite-presentation-local} et
\ref{spaces-morphisms-lemma-etale-local}.
\end{proof}















\section{Déformations infinitésimales de morphismes}
\label{spaces-more-morphisms-section-action-by-derivations}

\noindent
Dans cette section, nous expliquons comment une dérivation peut servir à
déformer infinitésimalement un morphisme. Dans toute la section, nous utilisons le fait qu'un
faisceau sur un épaississement $X'$ de $X$ peut être considéré comme un faisceau sur $X$, voir les
Équations (\ref{spaces-more-morphisms-equation-equivalence-etale-spaces}) et
(\ref{spaces-more-morphisms-equation-fundamental-equivalence}).

\begin{lemma}
\label{spaces-more-morphisms-lemma-difference-derivation}
Soit $S$ un schéma. Soit $B$ un espace algébrique sur $S$.
Soient $X \subset X'$ et $Y \subset Y'$ deux épaississements du premier ordre
d'espaces algébriques sur $B$.
Soient $(a, a'), (b, b') : (X \subset X') \to (Y \subset Y')$
deux morphismes d'épaississements sur $B$. Supposons que
\begin{enumerate}
\item $a = b$, et
\item les deux morphismes $a^*\mathcal{C}_{Y/Y'} \to \mathcal{C}_{X/X'}$
(Lemme \ref{spaces-more-morphisms-lemma-conormal-functorial})
soient égaux.
\end{enumerate}
Alors le morphisme $(a')^\sharp - (b')^\sharp$ se factorise sous la forme
$$
\mathcal{O}_{Y'} \to \mathcal{O}_Y \xrightarrow{D}
a_*\mathcal{C}_{X/X'} \to a_*\mathcal{O}_{X'}
$$
où $D$ est une $\mathcal{O}_B$-dérivation.
\end{lemma}

\begin{proof}
Au lieu de travailler sur $Y$, travaillons sur $X$. L'avantage est que le foncteur d'image inverse
$a^{-1}$ est exact. En utilisant (1) et (2), on obtient un diagramme commutatif
à lignes exactes
$$
\xymatrix{
0 \ar[r] &
\mathcal{C}_{X/X'} \ar[r] &
\mathcal{O}_{X'} \ar[r] &
\mathcal{O}_X \ar[r] & 0 \\
0 \ar[r] &
a^{-1}\mathcal{C}_{Y/Y'} \ar[r] \ar[u] &
a^{-1}\mathcal{O}_{Y'}
\ar[r] \ar@<1ex>[u]^{(a')^\sharp} \ar@<-1ex>[u]_{(b')^\sharp} &
a^{-1}\mathcal{O}_Y \ar[r] \ar[u] & 0
}
$$
Or il est général que, dans une telle situation, la différence des
morphismes de $\mathcal{O}_B$-algèbres $(a')^\sharp$ et $(b')^\sharp$ est une
$\mathcal{O}_B$-dérivation de $a^{-1}\mathcal{O}_Y$ dans $\mathcal{C}_{X/X'}$.
Par adjonction entre les foncteurs $a^{-1}$ et $a_*$, cela revient
à une $\mathcal{O}_B$-dérivation de
$\mathcal{O}_Y$ dans $a_*\mathcal{C}_{X/X'}$. Nous omettons certains détails.
\end{proof}

\noindent
Notons que, dans la situation du lemme ci-dessus, on peut écrire
$D$ sous la forme
\begin{equation}
\label{spaces-more-morphisms-equation-D}
D = \text{d}_{Y/B} \circ \theta
\end{equation}
où $\theta$ est une application $\mathcal{O}_Y$-linéaire
$\theta : \Omega_{Y/B} \to a_*\mathcal{C}_{X/X'}$.
Bien entendu, par adjonction encore, on peut alors considérer $\theta$ comme une
application $\mathcal{O}_X$-linéaire
$\theta : a^*\Omega_{Y/B} \to \mathcal{C}_{X/X'}$.

\begin{lemma}
\label{spaces-more-morphisms-lemma-action-by-derivations}
Soit $S$ un schéma. Soit $B$ un espace algébrique sur $S$.
Soit $(a, a') : (X \subset X') \to (Y \subset Y')$
un morphisme d'épaississements du premier ordre sur $B$.
Soit
$$
\theta : a^*\Omega_{Y/B} \to \mathcal{C}_{X/X'}
$$
une application $\mathcal{O}_X$-linéaire. Il existe alors un unique morphisme d'épaississements
$(b, b') : (X \subset X') \to (Y \subset Y')$ tel que
(1) et (2) du
Lemme \ref{spaces-more-morphisms-lemma-difference-derivation}
soient vérifiées et que la dérivation $D$ et $\theta$ soient reliées par
l'Équation (\ref{spaces-more-morphisms-equation-D}).
\end{lemma}

\begin{proof}
Considérons le morphisme
$$
\alpha = (a')^\sharp + D : a^{-1}\mathcal{O}_{Y'} \to \mathcal{O}_{X'}
$$
où $D$ est comme dans l'Équation (\ref{spaces-more-morphisms-equation-D}). Puisque $D$ est une
$\mathcal{O}_B$-dérivation, il en résulte que $\alpha$ est un morphisme de
faisceaux de $\mathcal{O}_B$-algèbres. Par construction, on a
$i_X^\sharp \circ \alpha = a^\sharp \circ i_Y^\sharp$, où
$i_X : X \to X'$ et $i_Y : Y \to Y'$ sont les immersions fermées données. D'après le
Lemme \ref{spaces-more-morphisms-lemma-first-order-thickening-maps},
on obtient un unique morphisme
$(a, b') : (X  \subset X') \to (Y \subset Y')$ d'épaississements
sur $B$ tel que $\alpha = (b')^\sharp$. En posant $b = a$,
le résultat est démontré.
\end{proof}

\begin{remark}
\label{spaces-more-morphisms-remark-action-by-derivations}
Mêmes hypothèses et notations qu'au Lemme \ref{spaces-more-morphisms-lemma-action-by-derivations}.
L'action d'une section locale $\theta$ sur $a'$ est parfois notée
$\theta \cdot a'$. Notons que cela signifie seulement que
$(a')^\sharp$ et $(\theta \cdot a')^\sharp$ diffèrent d'une dérivation
$D$ reliée à $\theta$ par l'Équation (\ref{spaces-more-morphisms-equation-D}).
\end{remark}

\begin{lemma}
\label{spaces-more-morphisms-lemma-sheaf}
Soit $S$ un schéma. Soit $B$ un espace algébrique sur $S$.
Soient $X \subset X'$ et $Y \subset Y'$ des épaississements du premier ordre
sur $B$. Supposons donnés un morphisme $a : X \to Y$ et un morphisme
$A : a^*\mathcal{C}_{Y/Y'} \to \mathcal{C}_{X/X'}$ de
$\mathcal{O}_X$-modules. Pour un objet $U'$ de
$(X')_{spaces, \etale}$ avec $U = X \times_{X'} U'$
considérons les morphismes $a' : U' \to Y'$ tels que
\begin{enumerate}
\item $a'$ soit un morphisme sur $B$,
\item $a'|_U = a|_U$, et
\item le morphisme induit
$a^*\mathcal{C}_{Y/Y'}|_U \to \mathcal{C}_{X/X'}|_U$
soit la restriction de $A$ à $U$.
\end{enumerate}
Alors la règle
\begin{equation}
\label{spaces-more-morphisms-equation-sheaf}
U' \mapsto
\{a' : U' \to Y'\text{ tels que (1), (2), (3) soient satisfaites.}\}
\end{equation}
définit un faisceau d'ensembles sur $(X')_{spaces, \etale}$.
\end{lemma}

\begin{proof}
Notons $\mathcal{F}$ la règle du lemme.
Le morphisme de restriction $\mathcal{F}(U') \to \mathcal{F}(V')$ associé à
$V' \subset U' \subset X'$
pour la règle $\mathcal{F}$ est bien le morphisme de restriction $a' \mapsto a'|_{V'}$.
Avec cette définition, il est clair que $\mathcal{F}$ est un
faisceau, car les morphismes d'espaces algébriques vérifient la descente étale, voir
Descente sur les espaces,
Lemme \ref{spaces-descent-lemma-fpqc-universal-effective-epimorphisms}.
\end{proof}

\begin{lemma}
\label{spaces-more-morphisms-lemma-action-sheaf}
Mêmes notations et hypothèses qu'au Lemme \ref{spaces-more-morphisms-lemma-sheaf}.
Nous identifions les faisceaux sur $X$ et $X'$ au moyen de
(\ref{spaces-more-morphisms-equation-equivalence-etale-spaces}).
Le faisceau
$$
\SheafHom_{\mathcal{O}_X}(a^*\Omega_{Y/B}, \mathcal{C}_{X/X'})
$$
opère sur le faisceau (\ref{spaces-more-morphisms-equation-sheaf}). En outre, cette action
est simplement transitive pour tout objet $U'$ de $(X')_{spaces, \etale}$
au-dessus duquel le faisceau (\ref{spaces-more-morphisms-equation-sheaf}) possède une section.
\end{lemma}

\begin{proof}
C'est la combinaison des
Lemmes \ref{spaces-more-morphisms-lemma-difference-derivation},
\ref{spaces-more-morphisms-lemma-action-by-derivations},
et \ref{spaces-more-morphisms-lemma-sheaf}.
\end{proof}

\begin{remark}
\label{spaces-more-morphisms-remark-special-case}
Un cas particulier des
Lemmes \ref{spaces-more-morphisms-lemma-difference-derivation},
\ref{spaces-more-morphisms-lemma-action-by-derivations},
\ref{spaces-more-morphisms-lemma-sheaf}, et
\ref{spaces-more-morphisms-lemma-action-sheaf}
est celui où $Y = Y'$. Dans ce cas, le morphisme $A$ est toujours nul.
Le faisceau du
Lemme \ref{spaces-more-morphisms-lemma-sheaf}
est simplement donné par la règle
$$
U' \mapsto
\{a' : U' \to Y\text{ au-dessus de }B\text{ tel que } a'|_U = a|_U\}
$$
sur lequel opère le faisceau
$\SheafHom_{\mathcal{O}_X}(a^*\Omega_{Y/B}, \mathcal{C}_{X/X'})$.
\end{remark}

\begin{remark}
\label{spaces-more-morphisms-remark-another-special-case}
Un autre cas particulier des
Lemmes \ref{spaces-more-morphisms-lemma-difference-derivation},
\ref{spaces-more-morphisms-lemma-action-by-derivations},
\ref{spaces-more-morphisms-lemma-sheaf}, et
\ref{spaces-more-morphisms-lemma-action-sheaf}
est celui où $B$ est lui-même l'espace ambiant d'un épaississement $Z \subset Z' = B$
et où $Y = Z \times_{Z'} Y'$. Représentons la situation par
$$
\xymatrix{
(X \subset X') \ar@{..>}[rr]_{(a, ?)} \ar[rd]_{(g, g')} & &
(Y \subset Y') \ar[ld]^{(h, h')} \\
& (Z \subset Z')
}
$$
Dans ce cas, le morphisme $A : a^*\mathcal{C}_{Y/Y'} \to \mathcal{C}_{X/X'}$
est déterminé par $a$ : le morphisme
$h^*\mathcal{C}_{Z/Z'} \to \mathcal{C}_{Y/Y'}$ est surjectif (car nous avons
supposé $Y = Z \times_{Z'} Y'$) ; par conséquent, son image inverse
$g^*\mathcal{C}_{Z/Z'} = a^*h^*\mathcal{C}_{Z/Z'} \to
a^*\mathcal{C}_{Y/Y'}$ est surjective, et la composée
$g^*\mathcal{C}_{Z/Z'} \to a^*\mathcal{C}_{Y/Y'} \to \mathcal{C}_{X/X'}$
doit être le morphisme canonique induit par $g'$. Ainsi, le faisceau du
Lemme \ref{spaces-more-morphisms-lemma-sheaf}
est simplement donné par la règle
$$
U' \mapsto
\{a' : U' \to Y'\text{ au-dessus de }Z'\text{ tel que } a'|_U = a|_U\}
$$
sur lequel opère le faisceau
$\SheafHom_{\mathcal{O}_X}(a^*\Omega_{Y/Z}, \mathcal{C}_{X/X'})$.
\end{remark}

\begin{lemma}
\label{spaces-more-morphisms-lemma-action-by-derivations-etale-localization}
Soit $S$ un schéma. Considérons un diagramme commutatif d'épaississements
du premier ordre
$$
\vcenter{
\xymatrix{
(T_2 \subset T_2') \ar[d]_{(h, h')} \ar[rr]_{(a_2, a_2')} & &
(X_2 \subset X_2') \ar[d]^{(f, f')} \\
(T_1 \subset T_1') \ar[rr]^{(a_1, a_1')} & &
(X_1 \subset X_1')
}
}
\quad
\begin{matrix}
\text{et un} \\
\text{diagramme commutatif}
\end{matrix}
\quad
\vcenter{
\xymatrix{
X_2' \ar[r] \ar[d] & B_2 \ar[d] \\
X_1' \ar[r] & B_1
}
}
$$
d'espaces algébriques sur $S$,
où $X_2 \to X_1$ et $B_2 \to B_1$ sont étales.
Pour toute application $\mathcal{O}_{T_1}$-linéaire
$\theta_1 : a_1^*\Omega_{X_1/B_1} \to \mathcal{C}_{T_1/T'_1}$, soit
$\theta_2$ la composée
$$
\xymatrix{
a_2^*\Omega_{X_2/B_2} \ar@{=}[r] &
h^*a_1^*\Omega_{X_1/B_1} \ar[r]^-{h^*\theta_1} &
h^*\mathcal{C}_{T_1/T'_1} \ar[r] &
\mathcal{C}_{T_2/T'_2}
}
$$
(le signe d'égalité est expliqué dans la démonstration). Alors le diagramme
$$
\xymatrix{
T_2' \ar[rr]_{\theta_2 \cdot a_2'} \ar[d] & & X'_2 \ar[d] \\
T_1' \ar[rr]^{\theta_1 \cdot a_1'} & & X'_1
}
$$
est commutatif, où les actions $\theta_2 \cdot a_2'$ et $\theta_1 \cdot a_1'$
sont celles de la Remarque \ref{spaces-more-morphisms-remark-action-by-derivations}.
\end{lemma}

\begin{proof}
Le signe d'égalité provient de l'identification
$f^*\Omega_{X_1/S_1} = \Omega_{X_2/S_2}$ que l'on obtient
parce que la construction du faisceau des différentielles est
compatible avec la localisation étale (tant à la source qu'au but), voir le
Lemme \ref{spaces-more-morphisms-lemma-localize-differentials}.
En effet, on a ainsi
$a_2^*\Omega_{X_2/S_2} = a_2^*f^*\Omega_{X_1/S_1} =
h^*a_1^*\Omega_{X_1/S_1}$ puisque $f \circ a_2 = a_1 \circ h$.
Cela étant, la commutativité du diagramme se vérifie
localement pour la topologie étale. On peut donc supposer que $T'_i$, $X'_i$,
$B_2$ et $B_1$ sont des schémas ; dans ce cas, le lemme
résulte de
Compléments sur les morphismes, Lemme
\ref{more-morphisms-lemma-action-by-derivations-etale-localization}.
Autre démonstration : d'après le Lemme \ref{spaces-more-morphisms-lemma-first-order-thickening-maps},
il suffit de montrer qu'un certain diagramme de faisceaux
d'anneaux sur $X_1'$ est commutatif ; on raisonne alors exactement
comme dans la démonstration du lemme précité de
Compléments sur les morphismes, Lemme
\ref{more-morphisms-lemma-action-by-derivations-etale-localization}
pour voir que c'est bien le cas.
\end{proof}










\section{Déformations infinitésimales d'espaces algébriques}
\label{spaces-more-morphisms-section-deform}

\noindent
Le lemme élémentaire suivant est souvent un outil commode pour vérifier si
une déformation infinitésimale d'un morphisme est plate.

\begin{lemma}
\label{spaces-more-morphisms-lemma-deform}
Soit $S$ un schéma. Soit $(f, f') : (X \subset X') \to (Y \subset Y')$ un
morphisme d'épaississements du premier ordre d'espaces algébriques sur $S$. Supposons que
$f$ soit plat. Alors les assertions suivantes sont équivalentes :
\begin{enumerate}
\item $f'$ est plat et $X = Y \times_{Y'} X'$, et
\item le morphisme canonique $f^*\mathcal{C}_{Y/Y'} \to \mathcal{C}_{X/X'}$
est un isomorphisme.
\end{enumerate}
\end{lemma}

\begin{proof}
Choisissons un schéma $V'$ et un morphisme étale surjectif $V' \to Y'$.
Choisissons un schéma $U'$ et un morphisme étale surjectif
$U' \to X' \times_{Y'} V'$. Posons $U = X \times_{X'} U'$ et
$V = Y \times_{Y'} V'$. D'après notre définition d'un morphisme plat
d'espaces algébriques, le morphisme induit $g : U \to V$ est un morphisme plat
de schémas, et $f'$ est plat si et seulement si le morphisme correspondant
$g' : U' \to V'$ est plat. De plus, $X = Y \times_{Y'} X'$ si et seulement
si $U = V \times_{V'} V'$. Enfin, le morphisme
$f^*\mathcal{C}_{Y/Y'} \to \mathcal{C}_{X/X'}$
est un isomorphisme si et seulement si
$g^*\mathcal{C}_{V/V'} \to \mathcal{C}_{U/U'}$ est un isomorphisme.
Le lemme résulte donc de son analogue pour les morphismes de schémas, voir
Compléments sur les morphismes, Lemme \ref{more-morphisms-lemma-deform}.
\end{proof}

\noindent
Le lemme suivant est la version « nilpotente » du
« critère de platitude par fibres », voir la
section \ref{spaces-more-morphisms-section-criterion-flat-fibres}.

\begin{lemma}
\label{spaces-more-morphisms-lemma-flatness-morphism-thickenings}
Soit $S$ un schéma. Considérons un diagramme commutatif
$$
\xymatrix{
(X \subset X') \ar[rr]_{(f, f')} \ar[rd] & & (Y \subset Y') \ar[ld] \\
& (B \subset B')
}
$$
d'épaississements d'espaces algébriques sur $S$. Supposons que
\begin{enumerate}
\item $X'$ soit plat sur $B'$,
\item $f$ soit plat,
\item $B \subset B'$ soit un épaississement d'ordre fini, et
\item $X = B \times_{B'} X'$ et $Y = B \times_{B'} Y'$.
\end{enumerate}
Alors $f'$ est plat et $Y'$ est plat sur $B'$ en tout point de
l'image de $f'$.
\end{lemma}

\begin{proof}
Choisissons un schéma $U'$ et un morphisme étale surjectif $U' \to B'$.
Choisissons un schéma $V'$ et un morphisme étale surjectif
$V' \to U' \times_{B'} Y'$.
Choisissons un schéma $W'$ et un morphisme étale surjectif
$W' \to V' \times_{Y'} X'$. Soient $U, V, W$ les changements de base
de $U', V', W'$ par $B \to B'$. Alors la platitude de $f'$ est
équivalente à celle de $W' \to V'$, et nous savons
que $W \to V$ est plat. Nous pouvons donc appliquer le lemme
dans le cas des schémas au diagramme
$$
\xymatrix{
(W \subset W') \ar[rr] \ar[rd] & & (V \subset V') \ar[ld] \\
& (U \subset U')
}
$$
d'épaississements de schémas. Voir
Compléments sur les morphismes, Lemme
\ref{more-morphisms-lemma-flatness-morphism-thickenings}.
L'assertion relative à la platitude de $Y'/B'$ aux points de
l'image de $f'$ se démontre de la même manière.
\end{proof}

\noindent
De nombreuses propriétés des morphismes de schémas sont préservées par les
déformations plates.

\begin{lemma}
\label{spaces-more-morphisms-lemma-deform-property}
Soit $S$ un schéma. Considérons un diagramme commutatif
$$
\xymatrix{
(X \subset X') \ar[rr]_{(f, f')} \ar[rd] & & (Y \subset Y') \ar[ld] \\
& (B \subset B')
}
$$
d'épaississements d'espaces algébriques sur $S$. Supposons que $B \subset B'$
soit un épaississement d'ordre fini, que $X'$ soit plat sur $B'$, que $X = B \times_{B'} X'$,
et que $Y = B \times_{B'} Y'$. Alors
\begin{enumerate}
\item $f$ est représentable si et seulement si $f'$ est représentable,
\label{spaces-more-morphisms-item-representable}
\item $f$ est plat si et seulement si $f'$ est plat,
\label{spaces-more-morphisms-item-flat}
\item $f$ est un isomorphisme si et seulement si $f'$ est un isomorphisme,
\label{spaces-more-morphisms-item-isomorphism}
\item $f$ est une immersion ouverte si et seulement si $f'$ est une immersion ouverte,
\label{spaces-more-morphisms-item-open-immersion}
\item $f$ est quasi-compact si et seulement si $f'$ est quasi-compact,
\label{spaces-more-morphisms-item-quasi-compact}
\item $f$ est universellement fermé si et seulement si $f'$ est universellement fermé,
\label{spaces-more-morphisms-item-universally-closed}
\item $f$ est (quasi-)séparé si et seulement si $f'$ est (quasi-)séparé,
\label{spaces-more-morphisms-item-separated}
\item $f$ est un monomorphisme si et seulement si $f'$ est un monomorphisme,
\label{spaces-more-morphisms-item-monomorphism}
\item $f$ est surjectif si et seulement si $f'$ est surjectif,
\label{spaces-more-morphisms-item-surjective}
\item $f$ est universellement injectif si et seulement si $f'$ est universellement injectif,
\label{spaces-more-morphisms-item-universally-injective}
\item $f$ est affine si et seulement si $f'$ est affine,
\label{spaces-more-morphisms-item-affine}
\item
\label{spaces-more-morphisms-item-finite-type}
$f$ est localement de type fini si et seulement si $f'$ est localement de type fini,
\item $f$ est localement quasi-fini si et seulement si $f'$ est localement quasi-fini,
\label{spaces-more-morphisms-item-quasi-finite}
\item
\label{spaces-more-morphisms-item-finite-presentation}
$f$ est localement de présentation finie si et seulement si $f'$ est localement de
présentation finie,
\item
\label{spaces-more-morphisms-item-relative-dimension-d}
$f$ est localement de type fini de dimension relative $d$ si et seulement si
$f'$ est localement de type fini de dimension relative $d$,
\item $f$ est universellement ouvert si et seulement si $f'$ est universellement ouvert,
\label{spaces-more-morphisms-item-universally-open}
\item $f$ est syntomique si et seulement si $f'$ est syntomique,
\label{spaces-more-morphisms-item-syntomic}
\item $f$ est lisse si et seulement si $f'$ est lisse,
\label{spaces-more-morphisms-item-smooth}
\item $f$ est non ramifié si et seulement si $f'$ est non ramifié,
\label{spaces-more-morphisms-item-unramified}
\item $f$ est étale si et seulement si $f'$ est étale,
\label{spaces-more-morphisms-item-etale}
\item $f$ est propre si et seulement si $f'$ est propre,
\label{spaces-more-morphisms-item-proper}
\item $f$ est entier si et seulement si $f'$ est entier,
\label{spaces-more-morphisms-item-integral}
\item $f$ est fini si et seulement si $f'$ est fini,
\label{spaces-more-morphisms-item-finite}
\item
\label{spaces-more-morphisms-item-finite-locally-free}
$f$ est fini localement libre (de rang $d$) si et seulement si $f'$
est fini localement libre (de rang $d$), et
\item ajouter d'autres cas ici.
\end{enumerate}
\end{lemma}

\begin{proof}
Le cas (\ref{spaces-more-morphisms-item-representable}) résulte du
Lemme \ref{spaces-more-morphisms-lemma-thicken-property-morphisms}.

\medskip\noindent
Choisissons un schéma $U'$ et un morphisme étale surjectif $U' \to B'$.
Choisissons un schéma $V'$ et un morphisme étale surjectif
$V' \to U' \times_{B'} Y'$.
Choisissons un schéma $W'$ et un morphisme étale surjectif
$W' \to V' \times_{Y'} X'$. Soient $U, V, W$ les changements de base
de $U', V', W'$ par $B \to B'$. Considérons le diagramme
$$
\xymatrix{
(W \subset W') \ar[rr] \ar[rd] & & (V \subset V') \ar[ld] \\
& (U \subset U')
}
$$
d'épaississements de schémas. Pour chacune des propriétés qui sont
locales pour la topologie étale à la source et au but, le résultat découle immédiatement
du résultat correspondant pour les morphismes d'épaississements de schémas,
appliqué au diagramme ci-dessus. Ainsi, les cas
(\ref{spaces-more-morphisms-item-flat}), (\ref{spaces-more-morphisms-item-finite-type}),
(\ref{spaces-more-morphisms-item-quasi-finite}), (\ref{spaces-more-morphisms-item-finite-presentation}),
(\ref{spaces-more-morphisms-item-relative-dimension-d}), (\ref{spaces-more-morphisms-item-syntomic}),
(\ref{spaces-more-morphisms-item-smooth}), (\ref{spaces-more-morphisms-item-unramified}), (\ref{spaces-more-morphisms-item-etale})
résultent des cas correspondants de
Compléments sur les morphismes, Lemme \ref{more-morphisms-lemma-deform-property}.

\medskip\noindent
Puisque $X \to X'$ et $Y \to Y'$ sont des homéomorphismes universels,
toute question portant sur la topologie des morphismes
$X \to Y$ et $X' \to Y'$ reçoit la même réponse. Ainsi,
les cas (\ref{spaces-more-morphisms-item-quasi-compact}), (\ref{spaces-more-morphisms-item-universally-closed}),
(\ref{spaces-more-morphisms-item-surjective}), (\ref{spaces-more-morphisms-item-universally-injective}), et
(\ref{spaces-more-morphisms-item-universally-open}) sont vérifiés.

\medskip\noindent
Dans chacun des cas restants, nous démontrons seulement l'implication
$f\text{ possède }P \Rightarrow f'\text{ possède }P$, car l'autre
implication résulte du fait que $P$ est stable par
changement de base, voir
Espaces, Lemme \ref{spaces-lemma-base-change-immersions} et
Morphismes d'espaces, Lemmes
\ref{spaces-morphisms-lemma-base-change-separated},
\ref{spaces-morphisms-lemma-base-change-monomorphism},
\ref{spaces-morphisms-lemma-base-change-affine},
\ref{spaces-morphisms-lemma-base-change-proper},
\ref{spaces-morphisms-lemma-base-change-integral}, et
\ref{spaces-morphisms-lemma-base-change-finite-locally-free}.

\medskip\noindent
Le cas (\ref{spaces-more-morphisms-item-open-immersion}). Supposons que $f$ soit une immersion ouverte.
Alors $f'$ est étale d'après (\ref{spaces-more-morphisms-item-etale}) et universellement injectif
d'après (\ref{spaces-more-morphisms-item-universally-injective}) ;
ainsi, $f'$ est une immersion ouverte, voir
Morphismes d'espaces, Lemme
\ref{spaces-morphisms-lemma-etale-universally-injective-open}.
On peut éviter d'utiliser ce lemme en
utilisant d'abord (\ref{spaces-more-morphisms-item-representable}) pour se ramener au cas des schémas.

\medskip\noindent
Le cas (\ref{spaces-more-morphisms-item-isomorphism}) résulte des cas
(\ref{spaces-more-morphisms-item-open-immersion}) et (\ref{spaces-more-morphisms-item-surjective}).

\medskip\noindent
Le cas (\ref{spaces-more-morphisms-item-separated}). Voir le
Lemme \ref{spaces-more-morphisms-lemma-thicken-property-morphisms}.

\medskip\noindent
Le cas (\ref{spaces-more-morphisms-item-monomorphism}). Supposons que $f$ soit un monomorphisme.
Considérons le morphisme diagonal $\Delta_{X'/Y'} : X' \to X' \times_{Y'} X'$.
Le changement de base de $\Delta_{X'/Y'}$ par $B \to B'$ est $\Delta_{X/Y}$,
qui est un isomorphisme par hypothèse. D'après (\ref{spaces-more-morphisms-item-isomorphism}),
nous concluons que $\Delta_{X'/Y'}$ est un isomorphisme et donc que
$f'$ est un monomorphisme.

\medskip\noindent
Le cas (\ref{spaces-more-morphisms-item-affine}). Voir le Lemme \ref{spaces-more-morphisms-lemma-thicken-property-morphisms}.

\medskip\noindent
Le cas (\ref{spaces-more-morphisms-item-proper}). Voir le
Lemme \ref{spaces-more-morphisms-lemma-thicken-property-morphisms-cartesian}.

\medskip\noindent
Le cas (\ref{spaces-more-morphisms-item-integral}). Voir le
Lemme \ref{spaces-more-morphisms-lemma-thicken-property-morphisms}.

\medskip\noindent
Le cas (\ref{spaces-more-morphisms-item-finite}). Voir le
Lemme \ref{spaces-more-morphisms-lemma-thicken-property-morphisms-cartesian}.

\medskip\noindent
Le cas (\ref{spaces-more-morphisms-item-finite-locally-free}). Supposons $f$ fini localement libre.
D'après (\ref{spaces-more-morphisms-item-finite}), $f'$ est fini.
D'après (\ref{spaces-more-morphisms-item-flat}), $f'$ est plat.
D'après (\ref{spaces-more-morphisms-item-finite-presentation}), $f'$ est localement de
présentation finie. Ainsi, $f'$ est fini localement libre d'après
Morphismes d'espaces, Lemme \ref{spaces-morphisms-lemma-finite-flat}.
\end{proof}

\noindent
Le lemme suivant est la version « localement nilpotente » du
« critère de platitude par fibres », voir la
section \ref{spaces-more-morphisms-section-criterion-flat-fibres}.

\begin{lemma}
\label{spaces-more-morphisms-lemma-flatness-morphism-thickenings-fp-over-ft}
Soit $S$ un schéma. Considérons un diagramme commutatif
$$
\xymatrix{
(X \subset X') \ar[rr]_{(f, f')} \ar[rd] & & (Y \subset Y') \ar[ld] \\
& (B \subset B')
}
$$
d'épaississements d'espaces algébriques sur $S$. Supposons que
\begin{enumerate}
\item $Y' \to B'$ soit localement de type fini,
\item $X' \to B'$ soit plat et localement de présentation finie,
\item $f$ soit plat, et
\item $X = B \times_{B'} X'$ et $Y = B \times_{B'} Y'$.
\end{enumerate}
Alors $f'$ est plat et, pour tout $y' \in |Y'|$ dans l'image de $|f'|$,
le morphisme $Y' \to B'$ est plat en $y'$.
\end{lemma}

\begin{proof}
Choisissons un schéma $U'$ et un morphisme étale surjectif $U' \to B'$.
Choisissons un schéma $V'$ et un morphisme étale surjectif
$V' \to U' \times_{B'} Y'$.
Choisissons un schéma $W'$ et un morphisme étale surjectif
$W' \to V' \times_{Y'} X'$. Soient $U, V, W$ les changements de base
de $U', V', W'$ par $B \to B'$. Alors la platitude de $f'$ est
équivalente à celle de $W' \to V'$, et nous savons
que $W \to V$ est plat. Nous pouvons donc appliquer le lemme
dans le cas des schémas au diagramme
$$
\xymatrix{
(W \subset W') \ar[rr] \ar[rd] & & (V \subset V') \ar[ld] \\
& (U \subset U')
}
$$
d'épaississements de schémas. Voir
Compléments sur les morphismes, Lemme
\ref{more-morphisms-lemma-flatness-morphism-thickenings-fp-over-ft}.
L'assertion relative à la platitude de $Y'/B'$ aux points de
l'image de $f'$ se démontre de la même manière.
\end{proof}

\noindent
De nombreuses propriétés des morphismes de schémas sont préservées par les
déformations plates comme dans le lemme ci-dessus.

\begin{lemma}
\label{spaces-more-morphisms-lemma-deform-property-fp-over-ft}
Soit $S$ un schéma. Considérons un diagramme commutatif
$$
\xymatrix{
(X \subset X') \ar[rr]_{(f, f')} \ar[rd] & & (Y \subset Y') \ar[ld] \\
& (B \subset B')
}
$$
d'épaississements d'espaces algébriques sur $S$.
Supposons que $Y' \to B'$ soit localement de type fini,
que $X' \to B'$ soit plat et localement de présentation finie,
que $X = B \times_{B'} X'$ et que $Y = B \times_{B'} Y'$. Alors
\begin{enumerate}
\item $f$ est représentable si et seulement si $f'$ est représentable,
\label{spaces-more-morphisms-item-representable-fp-over-ft}
\item $f$ est plat si et seulement si $f'$ est plat,
\label{spaces-more-morphisms-item-flat-fp-over-ft}
\item $f$ est un isomorphisme si et seulement si $f'$ est un isomorphisme,
\label{spaces-more-morphisms-item-isomorphism-fp-over-ft}
\item $f$ est une immersion ouverte si et seulement si $f'$ est une immersion ouverte,
\label{spaces-more-morphisms-item-open-immersion-fp-over-ft}
\item $f$ est quasi-compact si et seulement si $f'$ est quasi-compact,
\label{spaces-more-morphisms-item-quasi-compact-fp-over-ft}
\item $f$ est universellement fermé si et seulement si $f'$ est universellement fermé,
\label{spaces-more-morphisms-item-universally-closed-fp-over-ft}
\item $f$ est (quasi-)séparé si et seulement si $f'$ est (quasi-)séparé,
\label{spaces-more-morphisms-item-separated-fp-over-ft}
\item $f$ est un monomorphisme si et seulement si $f'$ est un monomorphisme,
\label{spaces-more-morphisms-item-monomorphism-fp-over-ft}
\item $f$ est surjectif si et seulement si $f'$ est surjectif,
\label{spaces-more-morphisms-item-surjective-fp-over-ft}
\item $f$ est universellement injectif si et seulement si $f'$ est universellement injectif,
\label{spaces-more-morphisms-item-universally-injective-fp-over-ft}
\item $f$ est affine si et seulement si $f'$ est affine,
\label{spaces-more-morphisms-item-affine-fp-over-ft}
\item $f$ est localement quasi-fini si et seulement si $f'$ est localement quasi-fini,
\label{spaces-more-morphisms-item-quasi-finite-fp-over-ft}
\item
\label{spaces-more-morphisms-item-relative-dimension-d-fp-over-ft}
$f$ est localement de type fini de dimension relative $d$ si et seulement si
$f'$ est localement de type fini de dimension relative $d$,
\item $f$ est universellement ouvert si et seulement si $f'$ est universellement ouvert,
\label{spaces-more-morphisms-item-universally-open-fp-over-ft}
\item $f$ est syntomique si et seulement si $f'$ est syntomique,
\label{spaces-more-morphisms-item-syntomic-fp-over-ft}
\item $f$ est lisse si et seulement si $f'$ est lisse,
\label{spaces-more-morphisms-item-smooth-fp-over-ft}
\item $f$ est non ramifié si et seulement si $f'$ est non ramifié,
\label{spaces-more-morphisms-item-unramified-fp-over-ft}
\item $f$ est étale si et seulement si $f'$ est étale,
\label{spaces-more-morphisms-item-etale-fp-over-ft}
\item $f$ est propre si et seulement si $f'$ est propre,
\label{spaces-more-morphisms-item-proper-fp-over-ft}
\item $f$ est fini si et seulement si $f'$ est fini,
\label{spaces-more-morphisms-item-finite-fp-over-ft}
\item
\label{spaces-more-morphisms-item-finite-locally-free-fp-over-ft}
$f$ est fini localement libre (de rang $d$) si et seulement si $f'$
est fini localement libre (de rang $d$), et
\item ajouter d'autres cas ici.
\end{enumerate}
\end{lemma}

\begin{proof}
Le cas (\ref{spaces-more-morphisms-item-representable-fp-over-ft}) résulte du
Lemme \ref{spaces-more-morphisms-lemma-thicken-property-morphisms}.

\medskip\noindent
Choisissons un schéma $U'$ et un morphisme étale surjectif $U' \to B'$.
Choisissons un schéma $V'$ et un morphisme étale surjectif
$V' \to U' \times_{B'} Y'$.
Choisissons un schéma $W'$ et un morphisme étale surjectif
$W' \to V' \times_{Y'} X'$. Soient $U, V, W$ les changements de base
de $U', V', W'$ par $B \to B'$. Considérons le diagramme
$$
\xymatrix{
(W \subset W') \ar[rr] \ar[rd] & & (V \subset V') \ar[ld] \\
& (U \subset U')
}
$$
d'épaississements de schémas. Pour chacune des propriétés qui sont
locales pour la topologie étale à la source et au but, le résultat découle immédiatement
du résultat correspondant pour les morphismes d'épaississements de schémas,
appliqué au diagramme ci-dessus. Ainsi, les cas
(\ref{spaces-more-morphisms-item-flat-fp-over-ft}),
(\ref{spaces-more-morphisms-item-quasi-finite-fp-over-ft}),
(\ref{spaces-more-morphisms-item-relative-dimension-d-fp-over-ft}),
(\ref{spaces-more-morphisms-item-syntomic-fp-over-ft}),
(\ref{spaces-more-morphisms-item-smooth-fp-over-ft}),
(\ref{spaces-more-morphisms-item-unramified-fp-over-ft}),
(\ref{spaces-more-morphisms-item-etale-fp-over-ft})
résultent des cas correspondants de
Compléments sur les morphismes, Lemme \ref{more-morphisms-lemma-deform-property-fp-over-ft}.

\medskip\noindent
Puisque $X \to X'$ et $Y \to Y'$ sont des homéomorphismes universels,
toute question portant sur la topologie des morphismes
$X \to Y$ et $X' \to Y'$ reçoit la même réponse. Ainsi,
les cas (\ref{spaces-more-morphisms-item-quasi-compact-fp-over-ft}),
(\ref{spaces-more-morphisms-item-universally-closed-fp-over-ft}),
(\ref{spaces-more-morphisms-item-surjective-fp-over-ft}),
(\ref{spaces-more-morphisms-item-universally-injective-fp-over-ft}), et
(\ref{spaces-more-morphisms-item-universally-open-fp-over-ft}) sont vérifiés.

\medskip\noindent
Dans chacun des cas restants, nous démontrons seulement l'implication
$f\text{ possède }P \Rightarrow f'\text{ possède }P$, car l'autre
implication résulte du fait que $P$ est stable par
changement de base, voir
Espaces, Lemme \ref{spaces-lemma-base-change-immersions} et
Morphismes d'espaces, Lemmes
\ref{spaces-morphisms-lemma-base-change-separated},
\ref{spaces-morphisms-lemma-base-change-monomorphism},
\ref{spaces-morphisms-lemma-base-change-affine},
\ref{spaces-morphisms-lemma-base-change-proper},
\ref{spaces-morphisms-lemma-base-change-integral}, et
\ref{spaces-morphisms-lemma-base-change-finite-locally-free}.

\medskip\noindent
Le cas (\ref{spaces-more-morphisms-item-open-immersion-fp-over-ft}).
Supposons que $f$ soit une immersion ouverte.
Alors $f'$ est étale d'après (\ref{spaces-more-morphisms-item-etale-fp-over-ft}) et universellement injectif
d'après (\ref{spaces-more-morphisms-item-universally-injective-fp-over-ft}) ;
ainsi, $f'$ est une immersion ouverte, voir
Morphismes d'espaces, Lemme
\ref{spaces-morphisms-lemma-etale-universally-injective-open}.
On peut éviter d'utiliser ce lemme en
utilisant d'abord (\ref{spaces-more-morphisms-item-representable-fp-over-ft}) pour se ramener au cas des schémas.

\medskip\noindent
Le cas (\ref{spaces-more-morphisms-item-isomorphism-fp-over-ft}) résulte des cas
(\ref{spaces-more-morphisms-item-open-immersion-fp-over-ft}) et (\ref{spaces-more-morphisms-item-surjective-fp-over-ft}).

\medskip\noindent
Le cas (\ref{spaces-more-morphisms-item-separated-fp-over-ft}). Voir le
Lemme \ref{spaces-more-morphisms-lemma-thicken-property-morphisms}.

\medskip\noindent
Le cas (\ref{spaces-more-morphisms-item-monomorphism-fp-over-ft}). Supposons que $f$ soit un monomorphisme.
Considérons le morphisme diagonal $\Delta_{X'/Y'} : X' \to X' \times_{Y'} X'$.
Le changement de base de $\Delta_{X'/Y'}$ par $B \to B'$ est $\Delta_{X/Y}$,
qui est un isomorphisme par hypothèse. D'après (\ref{spaces-more-morphisms-item-isomorphism-fp-over-ft}),
nous concluons que $\Delta_{X'/Y'}$ est un isomorphisme et donc que
$f'$ est un monomorphisme.

\medskip\noindent
Le cas (\ref{spaces-more-morphisms-item-affine-fp-over-ft}).
Voir le Lemme \ref{spaces-more-morphisms-lemma-thicken-property-morphisms}.

\medskip\noindent
Le cas (\ref{spaces-more-morphisms-item-proper-fp-over-ft}). Voir le
Lemme \ref{spaces-more-morphisms-lemma-properties-that-extend-over-thickenings}.

\medskip\noindent
Le cas (\ref{spaces-more-morphisms-item-finite-fp-over-ft}). Voir le
Lemme \ref{spaces-more-morphisms-lemma-properties-that-extend-over-thickenings}.

\medskip\noindent
Le cas (\ref{spaces-more-morphisms-item-finite-locally-free-fp-over-ft}).
Supposons $f$ fini localement libre.
D'après (\ref{spaces-more-morphisms-item-finite-fp-over-ft}), $f'$ est fini.
D'après (\ref{spaces-more-morphisms-item-flat-fp-over-ft}), $f'$ est plat.
De plus, $f'$ est localement de présentation finie d'après
Morphismes d'espaces, Lemme
\ref{spaces-morphisms-lemma-finite-presentation-permanence}.
Ainsi, $f'$ est fini localement libre d'après
Morphismes d'espaces, Lemme \ref{spaces-morphisms-lemma-finite-flat}.
\end{proof}













\section{Morphismes formellement lisses}
\label{spaces-more-morphisms-section-formally-smooth}

\noindent
Dans cette section, nous introduisons la notion de morphisme formellement lisse
$X \to Y$ d'espaces algébriques. Un tel morphisme est
caractérisé par la propriété que, pour $T$, les points de $X$ à valeurs dans ce schéma se relèvent
aux épaississements infinitésimaux de $T$, pourvu que $T$ soit affine.
Le résultat principal affirme qu'un morphisme formellement lisse et
localement de présentation finie est lisse, voir le
Lemme \ref{spaces-more-morphisms-lemma-smooth-formally-smooth}.
Il se trouve que ce critère est souvent plus facile à utiliser que le
critère jacobien.

\begin{definition}
\label{spaces-more-morphisms-definition-formally-smooth}
Soit $S$ un schéma. Un morphisme $f : X \to Y$ d'espaces algébriques sur $S$
est dit {\it formellement lisse} s'il est formellement lisse comme
transformation de foncteurs au sens de la
Définition \ref{spaces-more-morphisms-definition-formally-smooth-etale-unramified}.
\end{definition}

\noindent
Dans le cas des morphismes formellement non ramifiés et formellement étales,
on pourrait omettre l'hypothèse que $T'$ soit affine, voir les
Lemmes \ref{spaces-more-morphisms-lemma-formally-unramified-not-affine} et
\ref{spaces-more-morphisms-lemma-formally-etale-not-affine}.
Ce n'est plus vrai dans le cas des morphismes formellement lisses.
En fait, une condition un peu plus naturelle consisterait à pouvoir
compléter la flèche pointillée localement pour la topologie étale sur $T'$.
En effet, l'analyse de la démonstration du
Lemme \ref{spaces-more-morphisms-lemma-smooth-formally-smooth}
montre que cela serait équivalent à la définition sous sa forme
actuelle. Il est aussi vrai qu'exiger l'existence de la flèche
pointillée localement pour la topologie fppf sur $T'$ suffirait, mais cela est un peu
plus difficile à démontrer.

\medskip\noindent
Nous ne reformulerons pas les résultats démontrés dans le cadre plus général des
transformations de foncteurs formellement lisses à la
section \ref{spaces-more-morphisms-section-formally-smooth-etale-unramified}.

\begin{lemma}
\label{spaces-more-morphisms-lemma-composition-formally-smooth}
Une composée de morphismes formellement lisses est formellement lisse.
\end{lemma}

\begin{proof}
Démonstration omise.
\end{proof}

\begin{lemma}
\label{spaces-more-morphisms-lemma-base-change-formally-smooth}
Un changement de base d'un morphisme formellement lisse est formellement lisse.
\end{lemma}

\begin{proof}
Démonstration omise, mais voir
Algèbre, Lemme \ref{algebra-lemma-base-change-fs}
pour la version algébrique.
\end{proof}

\begin{lemma}
\label{spaces-more-morphisms-lemma-formally-etale-unramified-smooth}
Soit $f : X \to S$ un morphisme de schémas.
Alors $f$ est formellement étale si et seulement si
$f$ est formellement lisse et formellement non ramifié.
\end{lemma}

\begin{proof}
Démonstration omise.
\end{proof}

\noindent
Voici un lemme auxiliaire qui sera englobé par le
Lemme \ref{spaces-more-morphisms-lemma-formally-smooth}.

\begin{lemma}
\label{spaces-more-morphisms-lemma-helper-formally-smooth}
Soit $S$ un schéma. Soit
$$
\xymatrix{
U \ar[d] \ar[r]_\psi & V \ar[d] \\
X \ar[r]^f & Y
}
$$
un diagramme commutatif de morphismes d'espaces algébriques sur $S$.
Si les flèches verticales sont étales et si $f$ est formellement lisse, alors
$\psi$ est formellement lisse.
\end{lemma}

\begin{proof}
D'après le
Lemme \ref{spaces-more-morphisms-lemma-representable-property-formally-property},
les morphismes $U \to X$ et $V \to Y$ sont formellement étales. D'après le
Lemme \ref{spaces-more-morphisms-lemma-composition-formally-smooth-etale-unramified},
la composée $U \to Y$ est formellement lisse. D'après le
Lemme \ref{spaces-more-morphisms-lemma-formally-permanence},
nous voyons que $\psi : U \to V$ est formellement lisse.
\end{proof}

\noindent
Le lemme suivant est le résultat principal de cette section.
Combiné avec la
Proposition de Limites d'espaces
\ref{spaces-limits-proposition-characterize-locally-finite-presentation},
il implique que l'on peut reconnaître si un morphisme d'espaces algébriques
$f : X \to Y$ est lisse au moyen de propriétés « simples » de la
transformation de foncteurs $X \to Y$.

\begin{lemma}[Critère infinitésimal de relèvement]
\label{spaces-more-morphisms-lemma-smooth-formally-smooth}
Soit $S$ un schéma.
Soit $f : X \to Y$ un morphisme d'espaces algébriques sur $S$.
Les assertions suivantes sont équivalentes :
\begin{enumerate}
\item Le morphisme $f$ est lisse.
\item Le morphisme $f$ est localement de présentation finie et
formellement lisse.
\end{enumerate}
\end{lemma}

\begin{proof}
Supposons que $f : X \to S$ soit localement de présentation finie et formellement lisse.
Considérons un diagramme commutatif
$$
\xymatrix{
U \ar[d] \ar[r]_\psi & V \ar[d] \\
X \ar[r]^f & Y
}
$$
où $U$ et $V$ sont des schémas et où les flèches verticales sont étales et
surjectives. D'après le
Lemme \ref{spaces-more-morphisms-lemma-helper-formally-smooth},
nous voyons que $\psi : U \to V$ est formellement lisse. D'après
Morphismes d'espaces, Lemme
\ref{spaces-morphisms-lemma-finite-presentation-local}
le morphisme $\psi$ est localement de présentation finie.
Par conséquent, dans le cas des schémas, le morphisme
$\psi$ est lisse, voir
Compléments sur les morphismes, Lemme \ref{more-morphisms-lemma-smooth-formally-smooth}.
Ainsi, $f$ est lisse, voir
Morphismes d'espaces, Lemme
\ref{spaces-morphisms-lemma-smooth-local}.

\medskip\noindent
Réciproquement, supposons que $f : X \to Y$ soit lisse.
Considérons un diagramme commutatif de flèches pleines
$$
\xymatrix{
X \ar[d]_f & T \ar[d]^i \ar[l]^a \\
Y & T' \ar[l] \ar@{-->}[lu]
}
$$
comme dans la Définition \ref{spaces-more-morphisms-definition-formally-smooth}. Nous allons montrer que la
flèche pointillée existe, ce qui démontrera que $f$ est formellement lisse.
Soit $\mathcal{F}$ le faisceau d'ensembles sur $(T')_{spaces, \etale}$ du
Lemme \ref{spaces-more-morphisms-lemma-sheaf}, dans le cas particulier examiné à la
Remarque \ref{spaces-more-morphisms-remark-special-case}. Soit
$$
\mathcal{H} =
\SheafHom_{\mathcal{O}_T}(a^*\Omega_{X/Y}, \mathcal{C}_{T/T'})
$$
le faisceau de $\mathcal{O}_T$-modules sur $T_{spaces, \etale}$
muni de l'action $\mathcal{H} \times \mathcal{F} \to \mathcal{F}$
du Lemme \ref{spaces-more-morphisms-lemma-action-sheaf}.
L'action $\mathcal{H} \times \mathcal{F} \to \mathcal{F}$
fait de $\mathcal{F}$ un pseudo-torseur sous $\mathcal{H}$, voir
Cohomologie sur les sites, Définition \ref{sites-cohomology-definition-torsor}.
Notre but est de montrer que $\mathcal{F}$ est un torseur trivial sous $\mathcal{H}$.
Il y a deux étapes : (I) Pour montrer que $\mathcal{F}$ est un torseur,
nous devons montrer que $\mathcal{F}$ possède localement pour la topologie étale
une section. (II) Pour montrer que $\mathcal{F}$ est le torseur trivial,
il suffit de montrer que $H^1(T_\etale, \mathcal{H}) = 0$, voir
Cohomologie sur les sites, Lemme \ref{sites-cohomology-lemma-torsors-h1}.

\medskip\noindent
Démontrons d'abord (I). Pour cela, choisissons un diagramme commutatif
$$
\xymatrix{
U \ar[d] \ar[r]_\psi & V \ar[d] \\
X \ar[r]^f & Y
}
$$
où $U$ et $V$ sont des schémas et où les flèches verticales sont étales et
surjectives. Comme $f$ est supposé lisse, nous voyons que $\psi$ est lisse et
donc formellement lisse d'après le
Lemme \ref{spaces-more-morphisms-lemma-representable-property-formally-property}.
D'après le même lemme, le morphisme $V \to Y$ est formellement étale. Ainsi, d'après le
Lemme \ref{spaces-more-morphisms-lemma-composition-formally-smooth-etale-unramified},
la composée $U \to Y$ est formellement lisse. L'étape (I) résulte alors du
Lemme \ref{spaces-more-morphisms-lemma-etale-on-top}, partie (4).

\medskip\noindent
Démontrons enfin (II). D'après le
Lemme \ref{spaces-more-morphisms-lemma-finite-presentation-differentials},
nous voyons que $\Omega_{X/S}$ est de présentation finie.
Ainsi, $a^*\Omega_{X/S}$ est de présentation finie (voir
Propriétés d'espaces,
section \ref{spaces-properties-section-properties-modules}).
Par conséquent, le faisceau
$\mathcal{H} =
\SheafHom_{\mathcal{O}_T}(a^*\Omega_{X/Y}, \mathcal{C}_{T/T'})$
est quasi-cohérent d'après
Propriétés d'espaces,
Lemme \ref{spaces-properties-lemma-properties-quasi-coherent}.
Ainsi, d'après
Descente, Proposition \ref{descent-proposition-same-cohomology-quasi-coherent}
et Cohomologie des schémas, Lemme
\ref{coherent-lemma-quasi-coherent-affine-cohomology-zero}
nous avons
$$
H^1(T_{spaces, \etale}, \mathcal{H}) =
H^1(T_\etale, \mathcal{H}) =
H^1(T, \mathcal{H}) = 0
$$
comme voulu.
\end{proof}

\noindent
Les morphismes lisses satisfont une forte propriété locale de relèvement, voir le
Lemme \ref{spaces-more-morphisms-lemma-smooth-strong-lift}. Si, dans le lemme, nous
supposons que $T'$ est affine, alors
nous ignorons s'il est nécessaire de prendre un recouvrement étale.
Plus précisément, si nous avons un diagramme commutatif
$$
\xymatrix{
X \ar[d] & T \ar[l] \ar[d] \\
Y & T' \ar[l] \ar@{..>}[lu]
}
$$
d'espaces algébriques où $X \to Y$ est lisse
et où $T \to T'$ est un épaississement de schémas affines,
existe-t-il toujours une flèche pointillée rendant le diagramme
commutatif ? Si vous connaissez la réponse ou une référence, veuillez écrire à
\href{mailto:stacks.project@gmail.com}{stacks.project@gmail.com}.

\begin{lemma}
\label{spaces-more-morphisms-lemma-smooth-strong-lift}
Soit $S$ un schéma. Considérons un diagramme commutatif
$$
\xymatrix{
X \ar[d] & T \ar[l] \ar[d] \\
Y & T' \ar[l]
}
$$
d'espaces algébriques sur $S$, où $X \to Y$ est lisse
et où $T \to T'$ est un épaississement. Alors il existe un
recouvrement étale $\{T'_i \to T'\}$ tel que nous
puissions trouver la flèche pointillée dans
$$
\xymatrix{
X \ar[d] & T \ar[l] \ar[d] & T \times_{T'} T'_i \ar[l] \ar[d] \\
Y & T' \ar[l] & T'_i \ar[l] \ar@{..>}[llu]
}
$$
qui rend le diagramme commutatif (pour tout $i$).
\end{lemma}

\begin{proof}
Choisissons un recouvrement étale $\{Y_i \to Y\}$ où chaque $Y_i$ est affine.
Après avoir remplacé $T'$ par le recouvrement étale induit, nous pouvons supposer
que $Y$ est affine.

\medskip\noindent
Supposons que $Y$ soit affine. Choisissons un recouvrement étale $\{X_i \to X\}$.
Il donne naissance à un recouvrement étale de $T$. Ce recouvrement étale de $T$
provient d'un recouvrement étale de $T'$
(d'après le Théorème \ref{spaces-more-morphisms-theorem-topological-invariance}, voir la
discussion de la section \ref{spaces-more-morphisms-section-thickenings}).
Nous pouvons donc supposer que $X$ est affine.

\medskip\noindent
Supposons que $X$ et $Y$ soient affines. Nous pouvons encore remplacer
$T'$ par un recouvrement étale et supposer que $T'$ est affine. Dans ce cas, le lemme résulte de
Algèbre, Lemme \ref{algebra-lemma-smooth-strong-lift}.
\end{proof}







\noindent
Nous poursuivons un peu pour montrer que la propriété d'être formellement lisse est locale
pour la topologie étale à la source. Pour commencer, nous montrons qu'un morphisme formellement lisse possède un
faisceau de différentielles bien comporté. La notion de module quasi-cohérent
localement projectif est définie dans
Propriétés d'espaces,
section \ref{spaces-properties-section-locally-projective}.

\begin{lemma}
\label{spaces-more-morphisms-lemma-formally-smooth-sheaf-differentials}
Soit $S$ un schéma.
Soit $f : X \to Y$ un morphisme formellement lisse d'espaces algébriques sur $S$.
Alors $\Omega_{X/Y}$ est localement projectif sur $X$.
\end{lemma}

\begin{proof}
Choisissons un diagramme
$$
\xymatrix{
U \ar[d] \ar[r]_\psi & V \ar[d] \\
X \ar[r]^f & Y
}
$$
où $U$ et $V$ sont des schémas affines (!) et où les flèches verticales sont étales.
D'après le
Lemme \ref{spaces-more-morphisms-lemma-helper-formally-smooth},
nous voyons que $\psi : U \to V$ est formellement lisse. Ainsi,
$\Gamma(V, \mathcal{O}_V) \to \Gamma(U, \mathcal{O}_U)$ est
un homomorphisme d'anneaux formellement lisse, voir
Compléments sur les morphismes, Lemme \ref{more-morphisms-lemma-affine-formally-smooth}.
Par conséquent, d'après
Algèbre, Lemme \ref{algebra-lemma-characterize-formally-smooth-again},
le $\Gamma(U, \mathcal{O}_U)$-module
$\Omega_{\Gamma(U, \mathcal{O}_U)/\Gamma(V, \mathcal{O}_V)}$
est projectif. Ainsi, $\Omega_{U/V}$ est localement projectif, voir
Propriétés, section \ref{properties-section-locally-projective}.
Puisque $\Omega_{X/Y}|_U = \Omega_{U/V}$, nous voyons que $\Omega_{X/Y}$ est
lui aussi localement projectif. (En effet, nous pouvons trouver un recouvrement étale de
$X$ par des $U$ affines s'insérant dans des diagrammes comme ci-dessus --- les détails sont
omis.)
\end{proof}

\begin{lemma}
\label{spaces-more-morphisms-lemma-h1-is-zero}
Soit $T$ un schéma affine.
Soient $\mathcal{F}$ et $\mathcal{G}$ des
$\mathcal{O}_T$-modules quasi-cohérents sur $T_\etale$.
Considérons le faisceau Hom interne
$\mathcal{H} = \SheafHom_{\mathcal{O}_T}(\mathcal{F}, \mathcal{G})$
sur $T_\etale$.
Si $\mathcal{F}$ est localement projectif, alors
$H^1(T_\etale, \mathcal{H}) = 0$.
\end{lemma}

\begin{proof}
D'après la définition d'un faisceau localement projectif sur un espace algébrique (voir
Propriétés d'espaces,
Définition \ref{spaces-properties-definition-locally-projective}),
nous voyons que $\mathcal{F}_{Zar} = \mathcal{F}|_{T_{Zar}}$ est un faisceau
localement projectif sur le schéma $T$. Ainsi, $\mathcal{F}_{Zar}$ est un
facteur direct d'un $\mathcal{O}_{T_{Zar}}$-module libre. Nous
concluons alors (puisque $\mathcal{F} = (\mathcal{F}_{Zar})^a$, voir
Descente, Proposition \ref{descent-proposition-equivalence-quasi-coherent})
que $\mathcal{F}$ est un facteur direct d'un $\mathcal{O}_T$-module libre
sur $T_\etale$. Nous pouvons donc supposer que
$\mathcal{F} = \bigoplus_{i \in I} \mathcal{O}_T$ est un module libre.
Dans ce cas, $\mathcal{H} = \prod_{i \in I} \mathcal{G}$ est
un produit de modules quasi-cohérents. D'après
Cohomologie sur les sites, Lemme \ref{sites-cohomology-lemma-cohomology-products},
nous concluons que $H^1 = 0$, car la cohomologie d'un faisceau quasi-cohérent
sur un schéma affine est nulle, voir
Descente, Proposition \ref{descent-proposition-same-cohomology-quasi-coherent}
et Cohomologie des schémas, Lemme
\ref{coherent-lemma-quasi-coherent-affine-cohomology-zero}.
\end{proof}

\begin{lemma}
\label{spaces-more-morphisms-lemma-formally-smooth}
Soit $S$ un schéma.
Soit $f : X \to Y$ un morphisme d'espaces algébriques sur
$S$. Les assertions suivantes sont équivalentes :
\begin{enumerate}
\item $f$ est formellement lisse,
\item pour tout diagramme
$$
\xymatrix{
U \ar[d] \ar[r]_\psi & V \ar[d] \\
X \ar[r]^f & Y
}
$$
où $U$ et $V$ sont des schémas et où les flèches verticales sont étales,
le morphisme de schémas $\psi$ est formellement lisse (au sens de
Compléments sur les morphismes,
Définition \ref{more-morphisms-definition-formally-unramified}), et
\item pour un tel diagramme à flèches verticales surjectives, le morphisme
$\psi$ est formellement lisse.
\end{enumerate}
\end{lemma}

\begin{proof}
Nous avons vu que (1) implique (2) et (3) dans le
Lemme \ref{spaces-more-morphisms-lemma-helper-formally-smooth}.
Supposons (3).
La démonstration du fait que $f$ est formellement lisse est entièrement semblable à
la démonstration de (1) $\Rightarrow$ (2) du
Lemme \ref{spaces-more-morphisms-lemma-smooth-formally-smooth}.

\medskip\noindent
Considérons un diagramme commutatif de flèches pleines
$$
\xymatrix{
X \ar[d]_f & T \ar[d]^i \ar[l]^a \\
Y & T' \ar[l] \ar@{-->}[lu]
}
$$
comme dans la Définition \ref{spaces-more-morphisms-definition-formally-smooth}.
Nous allons montrer que la flèche pointillée existe, ce qui
démontrera que $f$ est formellement lisse.
Soit $\mathcal{F}$ le faisceau d'ensembles sur $(T')_{spaces, \etale}$ du
Lemme \ref{spaces-more-morphisms-lemma-sheaf}, dans le cas particulier examiné à la
Remarque \ref{spaces-more-morphisms-remark-special-case}.
Soit
$$
\mathcal{H} =
\SheafHom_{\mathcal{O}_T}(a^*\Omega_{X/Y}, \mathcal{C}_{T/T'})
$$
le faisceau de $\mathcal{O}_T$-modules sur $T_{spaces, \etale}$
muni de l'action $\mathcal{H} \times \mathcal{F} \to \mathcal{F}$ comme dans le
Lemme \ref{spaces-more-morphisms-lemma-action-sheaf}.
L'action $\mathcal{H} \times \mathcal{F} \to \mathcal{F}$
fait de $\mathcal{F}$ un pseudo-torseur sous $\mathcal{H}$, voir
Cohomologie sur les sites, Définition \ref{sites-cohomology-definition-torsor}.
Notre but est de montrer que $\mathcal{F}$ est un torseur trivial sous $\mathcal{H}$.
Il y a deux étapes : (I) Pour montrer que $\mathcal{F}$ est un torseur,
nous devons montrer que $\mathcal{F}$ possède localement pour la topologie étale
une section. (II) Pour montrer que $\mathcal{F}$ est le torseur trivial,
il suffit de montrer que $H^1(T_\etale, \mathcal{H}) = 0$, voir
Cohomologie sur les sites, Lemme \ref{sites-cohomology-lemma-torsors-h1}.

\medskip\noindent
Démontrons d'abord (I). Pour cela, considérons un diagramme
(qui existe puisque nous supposons (3))
$$
\xymatrix{
U \ar[d] \ar[r]_\psi & V \ar[d] \\
X \ar[r]^f & Y
}
$$
où $U$ et $V$ sont des schémas, où les flèches verticales sont étales et
surjectives, et où $\psi$ est formellement lisse. D'après le
Lemme \ref{spaces-more-morphisms-lemma-representable-property-formally-property},
le morphisme $V \to Y$ est formellement étale. Ainsi, d'après le
Lemme \ref{spaces-more-morphisms-lemma-composition-formally-smooth-etale-unramified},
la composée $U \to Y$ est formellement lisse. L'étape (I) résulte alors du
Lemme \ref{spaces-more-morphisms-lemma-etale-on-top}, partie (4).

\medskip\noindent
Démontrons enfin (II). D'après le
Lemme \ref{spaces-more-morphisms-lemma-formally-smooth-sheaf-differentials},
nous voyons que $\Omega_{U/V}$ est localement projectif.
Ainsi, $\Omega_{X/Y}$ est localement projectif, voir
Descente sur les espaces,
Lemme \ref{spaces-descent-lemma-locally-projective-descends}.
Par conséquent, $a^*\Omega_{X/Y}$ est localement projectif, voir
Propriétés d'espaces, Lemme
\ref{spaces-properties-lemma-locally-projective-pullback}.
Ainsi,
$$
H^1(T_\etale, \mathcal{H}) =
H^1(T_\etale,
\SheafHom_{\mathcal{O}_T}(a^*\Omega_{X/Y}, \mathcal{C}_{T/T'}) = 0
$$
d'après le
Lemme \ref{spaces-more-morphisms-lemma-h1-is-zero},
comme voulu.
\end{proof}

\begin{lemma}
\label{spaces-more-morphisms-lemma-descending-property-formally-smooth}
La propriété $\mathcal{P}(f) =$ « $f$ est formellement lisse »
est locale pour la topologie fpqc sur la base.
\end{lemma}

\begin{proof}
Soit $f : X \to Y$ un morphisme d'espaces algébriques sur un schéma $S$.
Choisissons un ensemble d'indices $I$ et des diagrammes
$$
\xymatrix{
U_i \ar[d] \ar[r]_{\psi_i} & V_i \ar[d] \\
X \ar[r]^f & Y
}
$$
à flèches verticales étales, où $U_i$ et $V_i$ sont des schémas affines. De plus,
supposons que $\coprod U_i \to X$ et $\coprod V_i \to Y$ soient surjectifs, voir
Propriétés d'espaces,
Lemme \ref{spaces-properties-lemma-cover-by-union-affines}.
D'après le
Lemme \ref{spaces-more-morphisms-lemma-formally-smooth},
nous voyons que $f$ est formellement lisse si et seulement si chacun des morphismes
$\psi_i$ est formellement lisse. Nous nous ramenons donc au cas d'un morphisme
de schémas affines. Dans ce cas, le résultat résulte de
Algèbre, Lemme \ref{algebra-lemma-descent-formally-smooth}.
Certains détails sont omis.
\end{proof}

\begin{lemma}
\label{spaces-more-morphisms-lemma-triangle-differentials-formally-smooth}
Soit $S$ un schéma.
Soient $f : X \to Y$ et $g : Y \to Z$ des morphismes d'espaces algébriques sur $S$.
Supposons que $f$ soit formellement lisse. Alors la suite
$$
0 \to f^*\Omega_{Y/Z} \to \Omega_{X/Z} \to \Omega_{X/Y} \to 0
$$
du Lemme \ref{spaces-more-morphisms-lemma-triangle-differentials}
est exacte courte.
\end{lemma}

\begin{proof}
Cela résulte du cas des schémas, voir
Compléments sur les morphismes,
Lemme \ref{more-morphisms-lemma-triangle-differentials-formally-smooth},
par localisation étale, voir les
Lemmes \ref{spaces-more-morphisms-lemma-formally-smooth} et \ref{spaces-more-morphisms-lemma-localize-differentials}.
\end{proof}

\begin{lemma}
\label{spaces-more-morphisms-lemma-differentials-formally-unramified-formally-smooth}
Soit $S$ un schéma. Soit $B$ un espace algébrique sur $S$.
Soit $h : Z \to X$ un morphisme formellement non ramifié d'espaces algébriques
sur $B$.
Supposons que $Z$ soit formellement lisse sur $B$. Alors la
suite exacte canonique
$$
0 \to \mathcal{C}_{Z/X} \to h^*\Omega_{X/B} \to \Omega_{Z/B} \to 0
$$
du
Lemme \ref{spaces-more-morphisms-lemma-universally-unramified-differentials-sequence}
est exacte courte.
\end{lemma}

\begin{proof}
Soit $Z \to Z'$ l'épaississement universel du premier ordre de $Z$ sur $X$.
La démonstration du
Lemme \ref{spaces-more-morphisms-lemma-universally-unramified-differentials-sequence}
montre que notre suite s'identifie à la suite
$$
\mathcal{C}_{Z/Z'} \to \Omega_{Z'/B} \otimes \mathcal{O}_Z \to
\Omega_{Z/B} \to 0.
$$
Puisque $Z \to S$ est formellement lisse, nous pouvons, localement pour la topologie étale sur $Z'$, trouver
un inverse à gauche $Z' \to Z$ sur $B$ du morphisme d'inclusion $Z \to Z'$.
Ainsi, la suite est localement scindée pour la topologie étale, voir le
Lemme \ref{spaces-more-morphisms-lemma-differentials-relative-immersion-section}.
\end{proof}

\begin{lemma}
\label{spaces-more-morphisms-lemma-two-unramified-morphisms-formally-smooth}
Soit $S$ un schéma. Soit
$$
\xymatrix{
Z \ar[r]_i \ar[rd]_j & X \ar[d]^f \\
& Y
}
$$
un diagramme commutatif d'espaces algébriques sur $S$,
où $i$ et $j$ sont formellement non ramifiés et où $f$ est formellement lisse.
Alors la suite exacte canonique
$$
0 \to
\mathcal{C}_{Z/Y} \to
\mathcal{C}_{Z/X} \to
i^*\Omega_{X/Y} \to 0
$$
du
Lemme \ref{spaces-more-morphisms-lemma-two-unramified-morphisms}
est exacte et localement scindée.
\end{lemma}

\begin{proof}
Notons $Z \to Z'$ l'épaississement universel du premier ordre de $Z$ sur $X$.
Notons $Z \to Z''$ l'épaississement universel du premier ordre de $Z$ sur $Y$.
D'après le
Lemme \ref{spaces-more-morphisms-lemma-universally-unramified-differentials-sequence},
il existe un morphisme canonique $Z' \to Z''$ tel que nous ayons un diagramme
commutatif
$$
\xymatrix{
Z \ar[r]_{i'} \ar[rd]_{j'} & Z' \ar[r]_a \ar[d]^k & X \ar[d]^f \\
& Z'' \ar[r]^b & Y
}
$$
La suite ci-dessus s'identifie à la suite
$$
\mathcal{C}_{Z/Z''} \to
\mathcal{C}_{Z/Z'} \to
(i')^*\Omega_{Z'/Z''} \to 0
$$
par nos définitions des faisceaux conormaux aux morphismes formellement non ramifiés.
Soit $U'' \to Z''$ un morphisme étale avec $U''$ affine.
Notons $U \to Z$ et $U' \to Z'$ les schémas affines correspondants,
étales sur $Z$ et $Z'$.
Comme $f$ est formellement lisse, il existe un morphisme $h : U'' \to X$
qui coïncide avec $i$ sur $U$ et tel que $f \circ h$ soit égal à $b|_{U''}$.
Puisque $Z'$ est l'épaississement universel du premier ordre, nous obtenons un unique
morphisme $g : U'' \to Z'$ tel que $g = a \circ h$. La propriété
universelle de $Z''$ implique que $k \circ g$ est le morphisme d'inclusion
$U'' \to Z''$. Ainsi, $g$ est un inverse à gauche de $k$. La situation est
$$
\xymatrix{
U \ar[d] \ar[r] & Z' \ar[d]^k \\
U'' \ar[r] \ar[ru]^g & Z''
}
$$
Ainsi, $g$ induit un morphisme $\mathcal{C}_{Z/Z'}|_U \to \mathcal{C}_{Z/Z''}|_U$
qui est un inverse à gauche du morphisme
$\mathcal{C}_{Z/Z''} \to \mathcal{C}_{Z/Z'}$ sur $U$.
\end{proof}











\section{Lissité sur une base noethérienne}
\label{spaces-more-morphisms-section-smooth-Noetherian}

\noindent
Cette section est l'analogue de
Compléments sur les morphismes, section \ref{more-morphisms-section-smooth-Noetherian}.

\begin{lemma}
\label{spaces-more-morphisms-lemma-lifting-along-artinian-at-point}
Soit $S$ un schéma. Soit $f : X \to Y$ un morphisme d'espaces algébriques
sur $S$. Soit $x \in |X|$.
Supposons que $Y$ soit localement noethérien et $f$ localement de type fini.
Les assertions suivantes sont équivalentes :
\begin{enumerate}
\item $f$ est lisse en $x$,
\item pour tout diagramme commutatif de flèches pleines
$$
\xymatrix{
X \ar[d]_f & \Spec(B) \ar[d]^i \ar[l]^-\alpha \\
Y & \Spec(B') \ar[l]_-{\beta} \ar@{-->}[lu]
}
$$
où $B' \to B$ est une surjection d'anneaux locaux telle que
$\Ker(B' \to B)$ soit de carré nul, et où $\alpha$ envoie le
point fermé de $\Spec(B)$ sur $x$, il existe
une flèche pointillée qui rend le diagramme commutatif, et
\item la même condition qu'en (2), mais où $B' \to B$ parcourt les petites
extensions (voir Algèbre, Définition \ref{algebra-definition-small-extension}).
\end{enumerate}
\end{lemma}

\begin{proof}
La condition (1) signifie qu'il existe un sous-espace ouvert $X' \subset X$
tel que $X' \to Y$ soit lisse. Ainsi, (1) implique les conditions (2) et (3) d'après le
Lemme \ref{spaces-more-morphisms-lemma-smooth-formally-smooth}. La condition (2) implique
trivialement la condition (3). Supposons (3). Choisissons un diagramme commutatif
$$
\xymatrix{
X \ar[d] & U \ar[l] \ar[d] \\
Y & V \ar[l]
}
$$
où $U$ et $V$ sont affines, où les flèches horizontales sont étales, et
tel qu'il existe un point $u \in U$ envoyé sur $x$. Considérons ensuite
un diagramme
$$
\xymatrix{
X \ar[d] & U \ar[l] \ar[d] & \Spec(B) \ar[d]^i \ar[l]^-\alpha  \\
Y & V \ar[l] & \Spec(B') \ar[l]_-{\beta}
}
$$
comme en (3), mais pour $u \in U \to V$. Soit $\gamma : \Spec(B') \to X$
la flèche fournie par notre hypothèse que (3) vaut pour $X$.
Puisque $U \to X$ est étale et
donc formellement étale (Lemme \ref{spaces-more-morphisms-lemma-etale-formally-etale}),
le morphisme $\gamma$
possède un unique relèvement à $U$ compatible avec $\alpha$. Puisque
$V \to Y$ est étale, donc formellement étale, ce relèvement est compatible
avec $\beta$. Ainsi, (3) vaut pour $u \in U \to V$, et nous concluons
que $U \to V$ est lisse en $u$ d'après
Compléments sur les morphismes, Lemme
\ref{more-morphisms-lemma-lifting-along-artinian-at-point}.
Cela montre que $X \to Y$ est lisse en $x$, et
achève la démonstration.
\end{proof}

\noindent
Il est parfois utile de savoir qu'il suffit de vérifier le
critère de relèvement pour les petites extensions « centrées » en des points
de type fini (voir
Morphismes d'espaces, section \ref{spaces-morphisms-section-points-finite-type}).

\begin{lemma}
\label{spaces-more-morphisms-lemma-lifting-along-artinian}
Soit $S$ un schéma. Soit $f : X \to Y$ un morphisme d'espaces algébriques
sur $S$. Supposons que $Y$ soit localement noethérien et $f$ localement de type fini.
Les assertions suivantes sont équivalentes :
\begin{enumerate}
\item $f$ est lisse,
\item pour tout diagramme commutatif de flèches pleines
$$
\xymatrix{
X \ar[d]_f & \Spec(B) \ar[d]^i \ar[l]^-\alpha \\
Y & \Spec(B') \ar[l]_-{\beta} \ar@{-->}[lu]
}
$$
où $B' \to B$ est une petite extension d'anneaux locaux artiniens
et où $\beta$ est de type fini (!), il existe une flèche pointillée qui rend
le diagramme commutatif.
\end{enumerate}
\end{lemma}

\begin{proof}
Si $f$ est lisse, alors le critère infinitésimal de relèvement
(Lemme \ref{spaces-more-morphisms-lemma-smooth-formally-smooth}) affirme que
$f$ est formellement lisse, et (2) est vérifiée.

\medskip\noindent
Supposons que $f$ ne soit pas lisse. L'ensemble des points $x \in X$ où $f$ n'est pas lisse
forme une partie fermée $T$ de $|X|$. D'après
Morphismes d'espaces, Lemme 
\ref{spaces-morphisms-lemma-enough-finite-type-points}, il existe
un point $x \in T \subset X$ tel que $x \in X_{\text{ft-pts}}$. Choisissons un
diagramme commutatif
$$
\xymatrix{
X \ar[d] & U \ar[l] \ar[d] & u \ar@{|->}[d] \\
Y & V \ar[l] & v
}
$$
où $U$ et $V$ sont affines, où les flèches horizontales sont étales, et
tel qu'il existe un point $u \in U$ envoyé sur $x$. Alors $u$
est un point de type fini de $U$. Puisque $U \to V$ n'est pas lisse au
point $u$, d'après
Compléments sur les morphismes,
Lemme \ref{more-morphisms-lemma-lifting-along-artinian-at-point},
il existe un diagramme
$$
\xymatrix{
X \ar[d] & U \ar[l] \ar[d] & \Spec(B) \ar[d]^i \ar[l]^-\alpha  \\
Y & V \ar[l] & \Spec(B') \ar[l]_-{\beta} \ar@{-->}[lu]
}
$$
où $B' \to B$ est une petite extension d'anneaux locaux (artiniens),
telle que le corps résiduel de $B$ soit égal à $\kappa(v)$ et
que la flèche pointillée n'existe pas. Puisque $U \to V$ est
de type fini, nous voyons que $v$ est un point de type fini de $V$. D'après
Morphismes, Lemme \ref{morphisms-lemma-artinian-finite-type},
le morphisme $\beta$ est de type fini ; par conséquent, la composée
$\Spec(B) \to Y$ est elle aussi de type fini.
En raisonnant exactement comme dans la démonstration du
Lemme \ref{spaces-more-morphisms-lemma-lifting-along-artinian-at-point}
(en utilisant que $U \to X$ et $V \to Y$ sont étales, donc
formellement étales),
nous voyons qu'il ne peut exister de flèche $\Spec(B) \to X$
s'insérant dans le rectangle extérieur du dernier diagramme affiché.
Autrement dit, (2) n'est pas vérifiée, et la démonstration est terminée.
\end{proof}

\noindent
Voici une application utile.

\begin{lemma}
\label{spaces-more-morphisms-lemma-check-smoothness-on-infinitesimal-nbhds}
Soit $S$ un schéma. Soit $f : X \to Y$ un morphisme
d'espaces algébriques sur $S$. Supposons que $f$ soit
localement de type fini et $Y$ localement noethérien.
Soit $Z \subset Y$ un sous-espace fermé dont le voisinage infinitésimal d'ordre $n$ est
$Z_n \subset Y$. Posons $X_n = Z_n \times_Y X$.
\begin{enumerate}
\item Si $X_n \to Z_n$ est lisse pour tout $n$, alors $f$
est lisse en tout point de $f^{-1}(Z)$.
\item Si $X_n \to Z_n$ est étale pour tout $n$, alors $f$
est étale en tout point de $f^{-1}(Z)$.
\end{enumerate}
\end{lemma}

\begin{proof}
Supposons que $X_n \to Z_n$ soit lisse pour tout $n$.
Soit $x \in X$ un point situé au-dessus d'un point de $Z$.
Étant donnés une petite extension $B' \to B$ et des morphismes $\alpha$, $\beta$ comme dans le
Lemme \ref{spaces-more-morphisms-lemma-lifting-along-artinian-at-point}, partie (3),
l'idéal maximal de $B'$ est nilpotent (puisque $B'$ est artinien),
et le morphisme $\beta$ se factorise donc par $Z_n$, tandis que $\alpha$ se factorise
par $X_n$ pour un $n$ convenable. Ainsi, la propriété de relèvement pour
$X_n \to Z_n$ s'applique et fournit la flèche pointillée souhaitée dans le diagramme.
Cela démontre (1). La partie (2) résulte de (1) et du fait qu'un morphisme
est étale si et seulement s'il est lisse de dimension relative $0$.
\end{proof}











\section{Le complexe cotangent naïf}
\label{spaces-more-morphisms-section-netherlander}

\noindent
Cette section prolonge
Modules sur les sites, section \ref{sites-modules-section-netherlander},
qui prolonge à son tour la discussion menée dans
Algèbre, section \ref{algebra-section-netherlander}.

\begin{definition}
\label{spaces-more-morphisms-definition-netherlander}
Soit $S$ un schéma.
Soit $f : X \to Y$ un morphisme d'espaces algébriques sur $S$.
Le {\it complexe cotangent naïf de $f$}
est le complexe défini dans Modules sur les sites, Définition
\ref{sites-modules-definition-cotangent-complex-morphism-ringed-topoi}
pour le morphisme de topos annelés $f_{small}$ entre les
petits sites étales de $X$ et $Y$, voir
Propriétés d'espaces, Lemme
\ref{spaces-properties-lemma-morphism-ringed-topoi}.
Notation : $\NL_f$ ou $\NL_{X/Y}$.
\end{definition}

\noindent
Les lemmes suivants montrent que cette définition est compatible avec celle
pour les homomorphismes d'anneaux et pour les schémas, et que $\NL_{X/Y}$ est un
objet de $D_\QCoh(\mathcal{O}_X)$.

\begin{lemma}
\label{spaces-more-morphisms-lemma-NL-etale-localization}
Soit $S$ un schéma. Considérons un diagramme commutatif
$$
\xymatrix{
U \ar[d]_p \ar[r]_g & V \ar[d]^q \\
X \ar[r]^f & Y
}
$$
d'espaces algébriques sur $S$ où $p$ et $q$ sont étales.
Alors il existe une identification canonique
$\NL_{X/Y}|_{U_\etale} = \NL_{U/V}$ dans $D(\mathcal{O}_U)$.
\end{lemma}

\begin{proof}
La formation du complexe cotangent naïf commute à l'image inverse
(Modules sur les sites, Lemme \ref{sites-modules-lemma-pullback-NL})
et nous avons $p_{small}^{-1}\mathcal{O}_X = \mathcal{O}_U$ ainsi que
$g_{small}^{-1}\mathcal{O}_{V_\etale} =
p_{small}^{-1}f_{small}^{-1}\mathcal{O}_{Y_\etale}$
car $q_{small}^{-1}\mathcal{O}_{Y_\etale} =
\mathcal{O}_{V_\etale}$ d'après Propriétés d'espaces, Lemme
\ref{spaces-properties-lemma-etale-exact-pullback}.
En suivant les définitions, nous concluons que
$\NL_{X/Y}|_{U_\etale} = \NL_{U/V}$.
\end{proof}

\begin{lemma}
\label{spaces-more-morphisms-lemma-NL-compare-spaces-schemes}
Soit $S$ un schéma. Soit $f : X \to Y$ un morphisme d'espaces algébriques
sur $S$. Supposons $X$ et $Y$ représentés par des schémas $X_0$ et $Y_0$.
Alors il existe une identification canonique
$\NL_{X/Y} = \epsilon^*\NL_{X_0/Y_0}$ dans $D(\mathcal{O}_X)$
où $\epsilon$ est celui de Catégories dérivées d'espaces, section
\ref{spaces-perfect-section-derived-quasi-coherent-etale}
et $\NL_{X_0/Y_0}$ est défini dans
Compléments sur les morphismes, Définition
\ref{more-morphisms-definition-netherlander}.
\end{lemma}

\begin{proof}
Soit $f_0 : X_0 \to Y_0$ le morphisme de schémas correspondant à $f$.
Il existe un morphisme canonique
$\epsilon^{-1}f_0^{-1}\mathcal{O}_{Y_0} \to f_{small}^{-1}\mathcal{O}_Y$
compatible avec
$\epsilon^\sharp : \epsilon^{-1}\mathcal{O}_{X_0} \to \mathcal{O}_X$
car il existe un diagramme commutatif
$$
\xymatrix{
X_{0, Zar} \ar[d]_{f_0} & X_\etale \ar[l]^\epsilon \ar[d]^f \\
Y_{0, Zar} & Y_\etale \ar[l]_\epsilon
}
$$
voir Catégories dérivées d'espaces, Remarque
\ref{spaces-perfect-remark-match-total-direct-images}.
Nous obtenons ainsi un morphisme canonique
$$
\epsilon^{-1}\NL_{X_0/Y_0} =
\epsilon^{-1}\NL_{\mathcal{O}_{X_0}/f_0^{-1}\mathcal{O}_{Y_0}} =
\NL_{\epsilon^{-1}\mathcal{O}_{X_0}/\epsilon^{-1}f_0^{-1}\mathcal{O}_{Y_0}}
\to
\NL_{\mathcal{O}_X/f^{-1}_{small}\mathcal{O}_Y} = \NL_{X/Y}
$$
par fonctorialité du complexe cotangent naïf.
Pour voir que le morphisme induit $\epsilon^*\NL_{X_0/Y_0} \to \NL_{X/Y}$ est un
isomorphisme dans $D(\mathcal{O}_X)$, il suffit de le vérifier sur les fibres aux points géométriques
(Propriétés d'espaces, Théorème
\ref{spaces-properties-theorem-exactness-stalks}).
Soit $\overline{x} : \Spec(k) \to X_0$
un point géométrique situé au-dessus de $x \in X_0$, et soit
$\overline{y} = f \circ \overline{x}$ situé au-dessus de $y \in Y_0$. Alors
$$
\NL_{X/Y, \overline{x}} =
\NL_{\mathcal{O}_{X, \overline{x}}/\mathcal{O}_{Y, \overline{y}}}
$$
Cela est vrai parce que le passage à la fibre en $\overline{x}$ revient
à prendre l'image inverse par $\overline{x} : \Spec(k) \to X$
et que nous pouvons appliquer Modules sur les sites, Lemme \ref{sites-modules-lemma-pullback-NL}.
D'autre part, nous avons
$$
(\epsilon^*\NL_{X_0/Y_0})_{\overline{x}} =
\NL_{X_0/Y_0, x} \otimes_{\mathcal{O}_{X_0, x}}
\mathcal{O}_{X, \overline{x}} =
\NL_{\mathcal{O}_{X_0, x}/\mathcal{O}_{Y_0, y}}
\otimes_{\mathcal{O}_{X_0, x}} \mathcal{O}_{X, \overline{x}}
$$
Certains détails sont omis (indication : utiliser le fait que la fibre d'une image inverse
est la fibre au point image, voir
Sites, Lemme \ref{sites-lemma-point-morphism-sites},
ainsi que le résultat correspondant pour les modules, voir
Modules sur les sites, Lemme \ref{sites-modules-lemma-pullback-stalk}).
Observons que $\mathcal{O}_{X, \overline{x}}$ est l'hensélisé
strict de $\mathcal{O}_{X_0, x}$, et de même
pour $\mathcal{O}_{Y, \overline{y}}$
(Propriétés d'espaces, Lemme
\ref{spaces-properties-lemma-describe-etale-local-ring}).
Le résultat découle donc de
Compléments sur l'algèbre,
Lemme \ref{more-algebra-lemma-henselization-NL}.
\end{proof}

\begin{lemma}
\label{spaces-more-morphisms-lemma-netherlander-quasi-coherent}
Soit $S$ un schéma. Soit $f : X \to Y$ un morphisme
d'espaces algébriques sur $S$. Les faisceaux de cohomologie
du complexe $\NL_{X/Y}$ sont quasi-cohérents et nuls en dehors
des degrés $-1$, $0$ ; le faisceau $\Omega_{X/Y}$ est sa cohomologie en degré $0$.
\end{lemma}

\begin{proof}
Par construction du complexe cotangent naïf dans
Modules sur les sites, section \ref{sites-modules-section-netherlander},
nous savons que $\NL_{X/Y}$ est un complexe nul en dehors des degrés $-1$, $0$
et que sa cohomologie en degré $0$ est $\Omega_{X/Y}$ (d'après notre
construction de $\Omega_{X/Y}$ à la section \ref{spaces-more-morphisms-section-sheaf-differentials}).
Le faisceau des différentielles est quasi-cohérent (d'après le
Lemme \ref{spaces-more-morphisms-lemma-module-differentials-quasi-coherent}).
Pour achever la démonstration, il suffit de montrer que $H^{-1}(\NL_{X/Y})$
est quasi-cohérent. Cela résulte d'une vérification locale pour la topologie étale
(permise par le Lemme \ref{spaces-more-morphisms-lemma-NL-etale-localization} et
Propriétés d'espaces, Lemme
\ref{spaces-properties-lemma-characterize-quasi-coherent})
qui nous ramène au cas des schémas
(Lemme \ref{spaces-more-morphisms-lemma-NL-compare-spaces-schemes}),
puis de l'emploi du résultat dans le cas des schémas
(Compléments sur les morphismes, Lemme
\ref{more-morphisms-lemma-netherlander-quasi-coherent}).
\end{proof}

\begin{lemma}
\label{spaces-more-morphisms-lemma-netherlander-fp}
Soit $S$ un schéma.
Soit $f : X \to Y$ un morphisme d'espaces algébriques sur $S$.
Si $f$ est localement de
présentation finie, alors $\NL_{X/Y}$ est, localement pour la topologie étale sur $X$,
quasi-isomorphe à un complexe
$$
\ldots \to 0 \to \mathcal{F}^{-1} \to \mathcal{F}^0 \to 0 \to \ldots
$$
de $\mathcal{O}_X$-modules quasi-cohérents
où $\mathcal{F}^0$ est de présentation finie
et $\mathcal{F}^{-1}$ de type fini.
\end{lemma}

\begin{proof}
La formation du complexe cotangent naïf commute à la
localisation étale d'après le Lemme \ref{spaces-more-morphisms-lemma-NL-etale-localization}.
Nous nous ramenons ainsi au cas des schémas par le
Lemme \ref{spaces-more-morphisms-lemma-NL-compare-spaces-schemes}.
Dans le cas des schémas, le résultat est
Compléments sur les morphismes, Lemme
\ref{more-morphisms-lemma-netherlander-fp}.
\end{proof}

\begin{lemma}
\label{spaces-more-morphisms-lemma-NL-formally-smooth}
Soit $S$ un schéma.
Soit $f : X \to Y$ un morphisme d'espaces algébriques sur $S$.
Les assertions suivantes sont équivalentes :
\begin{enumerate}
\item $f$ est formellement lisse,
\item $H^{-1}(\NL_{X/Y}) = 0$ et le module $H^0(\NL_{X/Y}) = \Omega_{X/Y}$
est localement projectif.
\end{enumerate}
\end{lemma}

\begin{proof}
Cela résulte du
Lemme \ref{spaces-more-morphisms-lemma-formally-smooth}, du
Lemme \ref{spaces-more-morphisms-lemma-NL-etale-localization}, du
Lemme \ref{spaces-more-morphisms-lemma-NL-compare-spaces-schemes}
et du cas des schémas, traité dans
Compléments sur les morphismes, Lemme \ref{more-morphisms-lemma-NL-formally-smooth}.
\end{proof}

\begin{lemma}
\label{spaces-more-morphisms-lemma-NL-formally-etale}
Soit $f : X \to Y$ un morphisme de schémas. Les assertions suivantes sont équivalentes :
\begin{enumerate}
\item $f$ est formellement étale,
\item $H^{-1}(\NL_{X/Y}) = H^0(\NL_{X/Y}) = 0$.
\end{enumerate}
\end{lemma}

\begin{proof}
Supposons (1). Tout morphisme formellement étale est formellement lisse.
Ainsi $H^{-1}(\NL_{X/Y}) = 0$ d'après le Lemme \ref{spaces-more-morphisms-lemma-NL-formally-smooth}.
D'autre part, un morphisme formellement étale est formellement non ramifié ;
par conséquent, $\Omega_{X/Y} = 0$ d'après le
Lemme \ref{spaces-more-morphisms-lemma-characterize-formally-unramified}.
Réciproquement, si (2) est vérifiée, alors $f$ est formellement lisse d'après le
Lemme \ref{spaces-more-morphisms-lemma-NL-formally-smooth}
et formellement non ramifié d'après le
Lemme \ref{spaces-more-morphisms-lemma-characterize-formally-unramified},
et donc formellement étale d'après le
Lemme \ref{spaces-more-morphisms-lemma-formally-etale-unramified-smooth}.
\end{proof}

\begin{lemma}
\label{spaces-more-morphisms-lemma-NL-smooth}
Soit $f : X \to Y$ un morphisme de schémas. Les assertions suivantes sont équivalentes :
\begin{enumerate}
\item $f$ est lisse, et
\item $f$ est localement de présentation finie,
$H^{-1}(\NL_{X/Y}) = 0$, et le module $H^0(\NL_{X/Y}) = \Omega_{X/Y}$
est localement libre de type fini.
\end{enumerate}
\end{lemma}

\begin{proof}
Cela résulte du
Lemme \ref{spaces-more-morphisms-lemma-formally-smooth}, du
Lemme \ref{spaces-more-morphisms-lemma-NL-etale-localization}, du
Lemme \ref{spaces-more-morphisms-lemma-NL-compare-spaces-schemes}
et du cas des schémas, traité dans
Compléments sur les morphismes, Lemme \ref{more-morphisms-lemma-NL-smooth}.
\end{proof}











\section{Ouverture du lieu de platitude}
\label{spaces-more-morphisms-section-open-flat}

\noindent
Cette section est l'analogue de
Compléments sur les morphismes, section \ref{more-morphisms-section-open-flat}.
Notons que nous avons défini la notion de platitude pour les modules
quasi-cohérents sur les espaces algébriques dans
Morphismes d'espaces, section \ref{spaces-morphisms-section-flat-modules}.

\begin{theorem}
\label{spaces-more-morphisms-theorem-openness-flatness}
Soit $S$ un schéma.
Soit $f : X \to Y$ un morphisme d'espaces algébriques sur $S$.
Soit $\mathcal{F}$ un faisceau quasi-cohérent sur $X$.
Supposons que $f$ soit localement de présentation finie et que
$\mathcal{F}$ soit un $\mathcal{O}_X$-module
localement de présentation finie. Alors
$$
\{x \in |X| : \mathcal{F}\text{ est plat sur }Y\text{ en }x\}
$$
est ouvert dans $|X|$.
\end{theorem}

\begin{proof}
Choisissons un diagramme commutatif
$$
\xymatrix{
U \ar[d]_p \ar[r]_\alpha &
V \ar[d]^q \\
X \ar[r]^a & Y
}
$$
où $U$, $V$ sont des schémas et $p$, $q$ sont étales et surjectifs comme dans
Espaces, Lemme \ref{spaces-lemma-lift-morphism-presentations}.
D'après
Compléments sur les morphismes, Théorème \ref{more-morphisms-theorem-openness-flatness},
l'ensemble
$U' = \{u \in |U| : p^*\mathcal{F}\text{ est plat sur }V\text{ en }u\}$
est ouvert dans $U$. D'après
Morphismes d'espaces, Définition \ref{spaces-morphisms-definition-flat-module},
l'image de $U'$ dans $|X|$ est l'ensemble
du théorème. La démonstration est donc terminée, car l'application $|U| \to |X|$ est
ouverte, voir
Propriétés d'espaces, Lemme \ref{spaces-properties-lemma-topology-points}.
\end{proof}

\begin{lemma}
\label{spaces-more-morphisms-lemma-flat-locus-base-change}
Soit $S$ un schéma. Soit
$$
\xymatrix{
X' \ar[r]_{g'} \ar[d]_{f'} & X \ar[d]^f \\
Y' \ar[r]^g & Y
}
$$
un diagramme cartésien d'espaces algébriques sur $S$.
Soit $\mathcal{F}$ un $\mathcal{O}_X$-module quasi-cohérent.
Supposons que $g$ soit plat, que $f$ soit localement de présentation finie,
et que $\mathcal{F}$ soit localement de présentation finie.
Alors
$$
\{x' \in |X'| : (g')^*\mathcal{F}\text{ est plat sur }Y'\text{ en }x'\}
$$
est l'image inverse de la partie ouverte du
Théorème \ref{spaces-more-morphisms-theorem-openness-flatness}
par l'application continue $|g'| : |X'| \to |X|$.
\end{lemma}

\begin{proof}
Cela découle de
Morphismes d'espaces, Lemme
\ref{spaces-morphisms-lemma-base-change-module-flat}.
\end{proof}












\section{Crit\`ere de platitude par fibres}
\label{spaces-more-morphisms-section-criterion-flat-fibres}

\noindent
Soit $S$ un schéma. Considérons un diagramme commutatif d'espaces algébriques
sur $S$
$$
\xymatrix{
X \ar[rr]_f \ar[dr]_g & & Y \ar[dl]^h \\
& Z
}
$$
et un $\mathcal{O}_X$-module quasi-cohérent $\mathcal{F}$.
Étant donné un point $x \in |X|$, nous nous demandons si
$\mathcal{F}$ est plat sur $Y$ en $x$. Si $\mathcal{F}$ est plat
sur $Z$ en $x$, le théorème ci-dessous affirme que cette question est
étroitement liée à celle de savoir si la restriction de
$\mathcal{F}$ à la fibre de $X \to Z$ au-dessus de $g(x)$
est plate sur la fibre de $Y \to Z$ au-dessus de $g(x)$. Pour donner un sens
à cet énoncé, nous proposons le lemme préliminaire suivant.

\begin{lemma}
\label{spaces-more-morphisms-lemma-flat-on-fibres-at-point}
Dans la situation ci-dessus, les assertions suivantes sont équivalentes :
\begin{enumerate}
\item Choisissons un point géométrique $\overline{x}$ de $X$ situé au-dessus de $x$.
Posons $\overline{y} = f \circ \overline{x}$ et
$\overline{z} = g \circ \overline{x}$. Alors le module
$\mathcal{F}_{\overline{x}}/
\mathfrak m_{\overline{z}}\mathcal{F}_{\overline{x}}$
est plat sur
$\mathcal{O}_{Y, \overline{y}}/
\mathfrak m_{\overline{z}}\mathcal{O}_{Y, \overline{y}}$.
\item Choisissons un morphisme $x : \Spec(K) \to X$ dans la classe d'équivalence de
$x$. Posons $z = g \circ x$, $X_z = \Spec(K) \times_{z, Z} X$,
$Y_z = \Spec(K) \times_{z, Z} Y$, et soit $\mathcal{F}_z$ l'image inverse
de $\mathcal{F}$ sur $X_z$. Alors $\mathcal{F}_z$ est plat en $x$ sur
$Y_z$ (au sens de Morphismes d'espaces,
Définition \ref{spaces-morphisms-definition-flat-module}).
\item Choisissons un diagramme commutatif
$$
\xymatrix{
& & & U \ar[llld]_a \ar[rr] \ar[dr] & & V \ar[llld]_>>>>>>>b \ar[dl] \\
X \ar[rr]_f \ar[dr]_g & & Y \ar[dl]^h &  & W \ar[llld]_c \\
& Z
}
$$
où $U, V, W$ sont des schémas et $a, b, c$ sont étales,
ainsi qu'un point $u \in U$ envoyé sur $x$. Soit $w \in W$ l'image de
$u$. Soit $\mathcal{F}_w$ l'image inverse de $\mathcal{F}$ sur
la fibre $U_w$ de $U \to W$ en $w$. Alors $\mathcal{F}_w$
est plat sur $V_w$ en $u$.
\end{enumerate}
\end{lemma}

\begin{proof}
Notons qu'en (2), le morphisme $x : \Spec(K) \to X$ définit un
point $K$-rationnel de $X_z$ ; l'énoncé a donc bien un sens. En outre,
la condition de (2) est indépendante du choix de $\Spec(K) \to X$
dans la classe d'équivalence de $x$ (les détails sont omis ; cela résultera aussi
des arguments ci-dessous, puisque les autres conditions ne dépendent pas
de ce choix). Notons également que nous pouvons toujours choisir un diagramme comme en
(3) de la manière suivante : choisissons d'abord
un schéma $W$ et un morphisme étale surjectif $W \to Z$, puis
un schéma $V$ et un morphisme étale surjectif $V \to W \times_Z Y$, et
enfin un schéma $U$ et un morphisme étale surjectif
$U \to V \times_Y X$. Ces choix étant faits, nous prenons pour $U \to W$
la composée $U \to V \to W$ et pouvons choisir un point $u \in U$ envoyé
sur $x$, puisque le morphisme $U \to X$ est surjectif.

\medskip\noindent
Supposons donnés à la fois un diagramme comme en (3) et un point géométrique
$\overline{x} : \Spec(k) \to X$ comme en (1). D'après
Propriétés d'espaces, Lemme
\ref{spaces-properties-lemma-geometric-lift-to-usual}
nous pouvons choisir un point géométrique $\overline{u} : \Spec(k) \to U$
au-dessus de $u$ tel que $\overline{x} = a \circ \overline{u}$.
Notons $\overline{v} : \Spec(k) \to V$ et
$\overline{w} : \Spec(k) \to W$ les points géométriques induits de
$V$ et $W$. Dans cette situation, nous savons que
$\mathcal{O}_{X, \overline{x}} = \mathcal{O}_{U, u}^{sh}$
et de même pour $Y$ et $Z$, voir
Propriétés d'espaces,
Lemme \ref{spaces-properties-lemma-describe-etale-local-ring}.
De même, nous avons
$$
\mathcal{F}_{\overline{x}} =
(a^*\mathcal{F})_u \otimes_{\mathcal{O}_{U, u}}
\mathcal{O}_{U, u}^{sh}
$$
voir
Propriétés d'espaces, Lemme \ref{spaces-properties-lemma-stalk-quasi-coherent}.
Notons que la fibre de $\mathcal{F}_w$ en $u$ est donnée par
$$
(\mathcal{F}_w)_u = (a^*\mathcal{F})_u/\mathfrak m_w(a^*\mathcal{F})_u
$$
et que l'anneau local de $V_w$ en $v$ est donné par
$$
\mathcal{O}_{V_w, v} = \mathcal{O}_{V, v}/\mathfrak m_w\mathcal{O}_{V, v}.
$$
Puisque $\mathfrak m_{\overline{z}} =
\mathfrak m_w \mathcal{O}_{Z, \overline{z}} =
\mathfrak m_w \mathcal{O}_{W, w}^{sh}$
nous voyons que
\begin{align*}
\mathcal{F}_{\overline{x}}/
\mathfrak m_{\overline{z}}\mathcal{F}_{\overline{x}} & =
(a^*\mathcal{F})_u \otimes_{\mathcal{O}_{U, u}}
\mathcal{O}_{X, \overline{x}}/
\mathfrak m_{\overline{z}}\mathcal{O}_{X, \overline{x}} \\
& =
(\mathcal{F}_w)_u \otimes_{\mathcal{O}_{U_w, u}}
\mathcal{O}_{U, u}^{sh}/\mathfrak m_w\mathcal{O}_{U, u}^{sh} \\
& = (\mathcal{F}_w)_u \otimes_{\mathcal{O}_{U_w, u}}
\mathcal{O}_{U_w, \overline{u}}^{sh} \\
& = (\mathcal{F}_w)_{\overline{u}}
\end{align*}
l'avant-dernière égalité résultant du
Lemme \ref{algebra-lemma-quotient-strict-henselization} d'Algèbre
et la dernière du
Lemme \ref{spaces-properties-lemma-stalk-quasi-coherent} de Propriétés d'espaces.
Les mêmes arguments appliqués aux faisceaux structuraux de $V$ et $Y$
montrent que
$$
\mathcal{O}_{V_w, \overline{v}}^{sh} =
\mathcal{O}_{V, v}^{sh}/\mathfrak m_w \mathcal{O}_{V, v}^{sh} =
\mathcal{O}_{Y, \overline{y}}/
\mathfrak m_{\overline{z}}\mathcal{O}_{Y, \overline{y}}.
$$
Nous pouvons maintenant appliquer
Morphismes d'espaces, Lemme \ref{spaces-morphisms-lemma-flat-at-point},
pour voir que (1) est équivalente à (3).

\medskip\noindent
Enfin, démontrons l'équivalence de (2) et (3).
Pour cela, choisissons une extension de corps $\tilde K$ de $K$
et un morphisme $\tilde x : \Spec(\tilde K) \to U$ qui
est situé au-dessus de $u$ (c'est possible car $u \times_{X, x} \Spec(K)$
est un schéma non vide). Soit $\tilde z : \Spec(\tilde K) \to U \to W$
la composée. Nous obtenons un diagramme commutatif
$$
\xymatrix{
& & & U_w \times_w \tilde z \ar[llld]_a \ar[rr] \ar[dr] & &
V_w \times_w \tilde z \ar[llld]_>>>>>>>b \ar[dl] \\
X_z \ar[rr]_f \ar[dr]_g & & Y_z \ar[dl]^h &  & \tilde z \ar[llld]_c \\
& z
}
$$
où $z = \Spec(K)$ et $w = \Spec(\kappa(w))$. Il est alors
clair que $\mathcal{F}_w$ et $\mathcal{F}_z$ ont pour images inverses le
même module sur $U_w \times_w \tilde z$. Nous obtenons ainsi un diagramme
commutatif
$$
\xymatrix{
X_z \ar[d] & U_w \times_w \tilde z \ar[l] \ar[d] \ar[r] & U_w \ar[d] \\
Y_z & V_w \times_w \tilde z \ar[l] \ar[r] & V_w
}
$$
dont les deux carrés sont cartésiens et dont les flèches horizontales du bas
sont plates : la flèche horizontale inférieure gauche est la composée
du morphisme $Y \times_Z \tilde z \to Y \times_Z z = Y_z$ (changement de base
d'un morphisme plat), du morphisme étale
$V \times_Z \tilde z \to Y \times_Z \tilde z$, et
du morphisme étale $V \times_W \tilde z \to V \times_Z \tilde z$.
Il résulte donc de
Morphismes d'espaces,
Lemme \ref{spaces-morphisms-lemma-base-change-module-flat},
que
$$
\mathcal{F}_z\text{ plat en }x\text{ sur }Y_z
\Leftrightarrow
\mathcal{F}|_{U_w \times_w \tilde z}
\text{ plat en }\tilde x\text{ sur }V_w \times_w \tilde z
\Leftrightarrow
\mathcal{F}_w\text{ plat en }u\text{ sur }V_w
$$
ce qui conclut.
\end{proof}

\begin{definition}
\label{spaces-more-morphisms-definition-module-flat-on-fibre}
Soit $S$ un schéma. Soient $X \to Y \to Z$ des morphismes d'espaces
algébriques sur $S$. Soit $\mathcal{F}$ un $\mathcal{O}_X$-module quasi-cohérent.
Soit $x \in |X|$ un point et notons $z \in |Z|$ son image.
\begin{enumerate}
\item Nous disons que {\it la restriction de $\mathcal{F}$ à sa fibre au-dessus de $z$
est plate en $x$ sur la fibre de $Y$ au-dessus de $z$} si les conditions équivalentes du
Lemme \ref{spaces-more-morphisms-lemma-flat-on-fibres-at-point}
sont vérifiées.
\item Nous disons que {\it la fibre de $X$ au-dessus de $z$ est plate en $x$ sur la fibre de
$Y$ au-dessus de $z$} si les conditions équivalentes du
Lemme \ref{spaces-more-morphisms-lemma-flat-on-fibres-at-point}
sont vérifiées pour $\mathcal{F} = \mathcal{O}_X$.
\item Nous disons que {\it la fibre de $X$ au-dessus de $z$ est plate sur la fibre de $Y$
au-dessus de $z$} si, pour tout $x \in |X|$ situé au-dessus de $z$, la fibre de $X$ au-dessus de $z$
est plate en $x$ sur la fibre de $Y$ au-dessus de $z$.
\end{enumerate}
\end{definition}

\noindent
Cette définition nous permet d'énoncer la version suivante du critère.
La version noethérienne se trouve à la
section \ref{spaces-more-morphisms-section-flat-over-Noetherian}.

\begin{theorem}
\label{spaces-more-morphisms-theorem-criterion-flatness-fibre}
Soit $S$ un schéma.
Soient $f : X \to Y$ et $Y \to Z$ des morphismes d'espaces algébriques sur $S$.
Soit $\mathcal{F}$ un $\mathcal{O}_X$-module quasi-cohérent.
Supposons que
\begin{enumerate}
\item $X$ soit localement de présentation finie sur $Z$,
\item $\mathcal{F}$ soit un $\mathcal{O}_X$-module de présentation finie, et
\item $Y$ soit localement de type fini sur $Z$.
\end{enumerate}
Soit $x \in |X|$ et soient $y \in |Y|$ et $z \in |Z|$ les images de
$x$. Si $\mathcal{F}_{\overline{x}} \not = 0$, les assertions suivantes sont
équivalentes :
\begin{enumerate}
\item $\mathcal{F}$ est plat sur $Z$ en $x$ et
la restriction de $\mathcal{F}$ à sa fibre au-dessus de $z$
est plate en $x$ sur la fibre de $Y$ au-dessus de $z$, et
\item $Y$ est plat sur $Z$ en $y$ et $\mathcal{F}$ est
plat sur $Y$ en $x$.
\end{enumerate}
En outre, l'ensemble des points $x$ où (1) et (2) sont vérifiées est ouvert dans
$\text{Supp}(\mathcal{F})$.
\end{theorem}

\begin{proof}
Choisissons un diagramme comme dans le
Lemme \ref{spaces-more-morphisms-lemma-flat-on-fibres-at-point}, partie (3).
Il résulte des définitions que l'on se ramène au
théorème correspondant pour les morphismes de schémas
$U \to V \to W$, le faisceau quasi-cohérent $a^*\mathcal{F}$
et le point $u \in U$. Le théorème découle donc du
résultat correspondant pour les schémas, à savoir
Compléments sur les morphismes,
Théorème \ref{more-morphisms-theorem-criterion-flatness-fibre}.
\end{proof}

\begin{lemma}
\label{spaces-more-morphisms-lemma-morphism-between-flat}
Soit $S$ un schéma.
Soient $f : X \to Y$ et $Y \to Z$ des morphismes d'espaces algébriques sur $S$.
Supposons que
\begin{enumerate}
\item $X$ soit localement de présentation finie sur $Z$,
\item $X$ soit plat sur $Z$,
\item pour tout $z \in |Z|$, la fibre de $X$ au-dessus de $z$
soit plate sur la fibre de $Y$ au-dessus de $z$, et
\item $Y$ soit localement de type fini sur $Z$.
\end{enumerate}
Alors $f$ est plat. Si $f$ est en outre surjectif, alors $Y$ est plat sur $Z$.
\end{lemma}

\begin{proof}
C'est un cas particulier du
Théorème \ref{spaces-more-morphisms-theorem-criterion-flatness-fibre}.
\end{proof}

\begin{lemma}
\label{spaces-more-morphisms-lemma-base-change-criterion-flatness-fibre}
Soit $S$ un schéma. Soient $f : X \to Y$ et $Y \to Z$ des morphismes
d'espaces algébriques sur $S$. Soit $\mathcal{F}$ un
$\mathcal{O}_X$-module quasi-cohérent.
Supposons que
\begin{enumerate}
\item $X$ soit localement de présentation finie sur $Z$,
\item $\mathcal{F}$ soit un $\mathcal{O}_X$-module de présentation finie,
\item $\mathcal{F}$ soit plat sur $Z$, et
\item $Y$ soit localement de type fini sur $Z$.
\end{enumerate}
Alors l'ensemble
$$
A = \{x \in |X| : \mathcal{F} \text{ plat en }x \text{ sur }Y\}.
$$
Cet ensemble est ouvert dans $|X|$ et sa formation commute à tout changement de base :
si $Z' \to Z$ est un morphisme d'espaces algébriques et si $A'$ est l'ensemble des
points de $X' = X \times_Z Z'$ où $\mathcal{F}' = \mathcal{F} \times_Z Z'$
est plat sur $Y' = Y \times_Z Z'$, alors $A'$ est l'image inverse de
$A$ par l'application continue $|X'| \to |X|$.
\end{lemma}

\begin{proof}
Une façon de le démontrer consiste à transposer la démonstration donnée dans
Compléments sur les morphismes, Lemme \ref{more-morphisms-lemma-morphism-between-flat},
à la catégorie des espaces algébriques. Nous le démontrerons plutôt
en nous ramenant au cas des schémas. Choisissons donc un diagramme comme dans le
Lemme \ref{spaces-more-morphisms-lemma-flat-on-fibres-at-point}, partie (3),
tel que $a$, $b$ et $c$ soient surjectifs.
Il résulte des définitions que l'on se ramène au
théorème correspondant pour les morphismes de schémas
$U \to V \to W$, le faisceau quasi-cohérent $a^*\mathcal{F}$
et le point $u \in U$. Le seul point à préciser est que,
étant donné un morphisme d'espaces algébriques $Z' \to Z$, nous choisissons un schéma
$W'$ et un morphisme étale surjectif $W' \to W \times_Z Z'$.
Posons alors $U' = W' \times_W U$ et $V' = W' \times_W V$.
Notons $a', b', c'$ les morphismes de $U', V', W'$ dans
$X', Y', Z'$. Dans ce cas, $A$, resp.\ $A'$ sont les images des parties
ouvertes de $U$, resp.\ $U'$ associées à
$a^*\mathcal{F}$, resp.\ $(a')^*\mathcal{F}'$.
Cela ramène bien le lemme à
Compléments sur les morphismes, Lemme \ref{more-morphisms-lemma-morphism-between-flat}.
\end{proof}

\begin{lemma}
\label{spaces-more-morphisms-lemma-base-change-flatness-fibres}
Soit $S$ un schéma.
Soient $f : X \to Y$ et $Y \to Z$ des morphismes d'espaces algébriques sur $S$.
Supposons que
\begin{enumerate}
\item $X$ soit localement de présentation finie sur $Z$,
\item $X$ soit plat sur $Z$, et
\item $Y$ soit localement de type fini sur $Z$.
\end{enumerate}
Alors l'ensemble
$$
\{x \in |X| : X\text{ plat en }x \text{ sur }Y\}.
$$
Cet ensemble est ouvert dans $|X|$ et sa formation commute à tout changement de base
$Z' \to Z$.
\end{lemma}

\begin{proof}
C'est un cas particulier du
Lemme \ref{spaces-more-morphisms-lemma-base-change-criterion-flatness-fibre}.
\end{proof}

\begin{lemma}
\label{spaces-more-morphisms-lemma-flat-and-free-at-point-fibre}
Soit $S$ un schéma. Soit $f : X \to Y$ un morphisme d'espaces algébriques
sur $S$ qui est localement de présentation finie.
Soit $\mathcal{F}$ un $\mathcal{O}_X$-module de présentation finie.
Soit $x \in |X|$, d'image $y \in |Y|$. Si $\mathcal{F}$ est plat en $x$
sur $Y$, les assertions suivantes sont équivalentes :
\begin{enumerate}
\item $(\mathcal{F}_{\overline{y}})_{\overline{x}}$ est un module plat sur
$\mathcal{O}_{X_{\overline{y}}, \overline{x}}$,
\item $(\mathcal{F}_{\overline{y}})_{\overline{x}}$ est un module libre sur
$\mathcal{O}_{X_{\overline{y}}, \overline{x}}$,
\item $\mathcal{F}_{\overline{y}}$ est libre de type fini sur un
voisinage étale de $\overline{x}$ dans $X_{\overline{y}}$, et
\item $\mathcal{F}$ est libre de type fini sur un voisinage étale de $x$ dans $X$.
\end{enumerate}
Ici, $\overline{x}$ est un point géométrique de $X$ situé au-dessus de $x$
et $\overline{y} = f \circ \overline{x}$.
\end{lemma}

\begin{proof}
Choisissons un diagramme commutatif
$$
\xymatrix{
U \ar[d] \ar[r] & V \ar[d] \\
X \ar[r] & Y
}
$$
où $U$ et $V$ sont des schémas et les flèches verticales sont étales,
et tel qu'il existe un point $u \in U$ envoyé sur $x$. Soit $v \in V$
l'image de $u$. En appliquant le Lemme \ref{spaces-more-morphisms-lemma-flat-on-fibres-at-point}
à $\text{id} : X \to X$ sur $Y$, nous voyons que (1) devient
la condition « $\mathcal{F}|_{U_v}$ est plat sur $U_v$ en $u$ ».
Autrement dit, (1) équivaut à ce que $(\mathcal{F}|_{U_v})_u$
soit un $\mathcal{O}_{U_v, u}$-module plat.
D'après le cas des schémas (Compléments sur les morphismes, Lemme
\ref{more-morphisms-lemma-flat-and-free-at-point-fibre}),
nous voyons que cela implique que
$\mathcal{F}|_U$ est libre de type fini sur un voisinage ouvert
de $u$. Nous voyons ainsi que (1) implique (4).
Les implications (4) $\Rightarrow$ (3) et
(2) $\Rightarrow$ (1) sont immédiates.
Pour l'implication (3) $\Rightarrow$ (2), utilisons
la description des anneaux locaux et des fibres dans
Propriétés d'espaces, Lemmes
\ref{spaces-properties-lemma-describe-etale-local-ring} et
\ref{spaces-properties-lemma-stalk-quasi-coherent}.
\end{proof}

\begin{lemma}
\label{spaces-more-morphisms-lemma-finite-free-open}
Soit $S$ un schéma. Soit $f : X \to Y$ un morphisme d'espaces algébriques
sur $S$ qui est localement de présentation finie.
Soit $\mathcal{F}$ un $\mathcal{O}_X$-module de présentation finie
plat sur $Y$. Alors l'ensemble
$$
\{x \in |X| : \mathcal{F}\text{ libre dans un voisinage étale de }x\}
$$
est ouvert dans $|X|$ et sa formation commute à tout changement de base
$Y' \to Y$.
\end{lemma}

\begin{proof}
L'ouverture est immédiate. Soit $Y' \to Y$ un morphisme d'espaces algébriques,
posons $X' = Y' \times_Y X$,
et soit $x' \in |X'|$ un point situé au-dessus de $x \in |X|$.
D'après le Lemme \ref{spaces-more-morphisms-lemma-flat-and-free-at-point-fibre},
nous voyons que $x$ appartient à notre ensemble si et seulement si
$(\mathcal{F}_{\overline{y}})_{\overline{x}}$ est un module plat sur
$\mathcal{O}_{X_{\overline{y}}, \overline{x}}$.
De même, $x'$ appartient à l'analogue de notre ensemble pour l'image inverse
$\mathcal{F}'$ de $\mathcal{F}$ sur $X'$ si et seulement si
$(\mathcal{F}'_{\overline{y}'})_{\overline{x}'}$ est un module plat sur
$\mathcal{O}_{X'_{\overline{y}'}, \overline{x}'}$
(avec les notations évidentes). Ces deux assertions sont équivalentes
d'après le Lemme \ref{spaces-more-morphisms-lemma-flat-on-fibres-at-point} appliqué au
morphisme $\text{id} : X \to X$ sur $Y$.
L'assertion sur le changement de base est donc démontrée.
\end{proof}







\section{Platitude sur une base noethérienne}
\label{spaces-more-morphisms-section-flat-over-Noetherian}

\noindent
Voici le « critère de platitude par fibres » dans le cas noethérien.

\begin{theorem}
\label{spaces-more-morphisms-theorem-criterion-flatness-fibre-Noetherian}
Soit $S$ un schéma.
Soient $f : X \to Y$ et $Y \to Z$ des morphismes d'espaces algébriques sur $S$.
Soit $\mathcal{F}$ un $\mathcal{O}_X$-module quasi-cohérent.
Supposons que
\begin{enumerate}
\item $X$, $Y$, $Z$ soient localement noethériens, et
\item $\mathcal{F}$ soit un $\mathcal{O}_X$-module cohérent.
\end{enumerate}
Soit $x \in |X|$ et soient $y \in |Y|$ et $z \in |Z|$ les images de
$x$. Si $\mathcal{F}_{\overline{x}} \not = 0$, les assertions suivantes sont
équivalentes :
\begin{enumerate}
\item $\mathcal{F}$ est plat sur $Z$ en $x$ et
la restriction de $\mathcal{F}$ à sa fibre au-dessus de $z$
est plate en $x$ sur la fibre de $Y$ au-dessus de $z$, et
\item $Y$ est plat sur $Z$ en $y$ et $\mathcal{F}$ est
plat sur $Y$ en $x$.
\end{enumerate}
\end{theorem}

\begin{proof}
Choisissons un diagramme comme dans le
Lemme \ref{spaces-more-morphisms-lemma-flat-on-fibres-at-point}, partie (3).
Il résulte des définitions que l'on se ramène au
théorème correspondant pour les morphismes de schémas
$U \to V \to W$, le faisceau quasi-cohérent $a^*\mathcal{F}$
et le point $u \in U$. Le théorème découle donc du
résultat correspondant pour les schémas, à savoir
Compléments sur les morphismes,
Théorème \ref{more-morphisms-theorem-criterion-flatness-fibre-Noetherian}.
\end{proof}

\begin{lemma}
\label{spaces-more-morphisms-lemma-morphism-between-flat-Noetherian}
Soit $S$ un schéma.
Soient $f : X \to Y$ et $Y \to Z$ des morphismes d'espaces algébriques sur $S$.
Supposons que
\begin{enumerate}
\item $X$, $Y$, $Z$ soient localement noethériens,
\item $X$ soit plat sur $Z$,
\item pour tout $z \in |Z|$, la fibre de $X$ au-dessus de $z$
soit plate sur la fibre de $Y$ au-dessus de $z$.
\end{enumerate}
Alors $f$ est plat. Si $f$ est en outre surjectif, alors $Y$ est plat sur $Z$.
\end{lemma}

\begin{proof}
C'est un cas particulier du
Théorème \ref{spaces-more-morphisms-theorem-criterion-flatness-fibre-Noetherian}.
\end{proof}

\noindent
Tout comme pour vérifier la lissité, lorsque la base est noethérienne, il suffit
de vérifier la platitude sur les anneaux artiniens. Voici un énoncé type.

\begin{lemma}
\label{spaces-more-morphisms-lemma-flatness-over-Noetherian-ring}
Soit $A$ un anneau noethérien. Soit $I \subset A$ un idéal.
Soit $X$ un espace algébrique localement de présentation finie sur
$S = \Spec(A)$. Pour $n \geq 1$, posons $S_n = \Spec(A/I^n)$ et
$X_n = S_n \times_S X$. Soit $\mathcal{F}$ un $\mathcal{O}_X$-module cohérent.
Si, pour tout $n \geq 1$, l'image inverse $\mathcal{F}_n$ de $\mathcal{F}$ sur $X$
est plate sur $S_n$, alors le lieu (ouvert) où $\mathcal{F}$
est plate sur $X$ contient l'image inverse de $V(I)$ par $X \to S$.
\end{lemma}

\begin{proof}
Le lieu où $\mathcal{F}$ est plat sur $S$ est ouvert dans $|X|$ d'après le
Théorème \ref{spaces-more-morphisms-theorem-openness-flatness}.
L'énoncé est invariant lorsqu'on remplace $X$ par les membres d'un
recouvrement étale ; nous pouvons donc supposer que $X$ est un schéma affine.
Dans ce cas, le résultat découle immédiatement du
Lemme \ref{algebra-lemma-flat-module-powers} d'Algèbre.
Certains détails sont omis.
\end{proof}







\section{Retour sur la normalisation}
\label{spaces-more-morphisms-section-normalization}

\noindent
La normalisation commute au changement de base lisse.

\begin{lemma}
\label{spaces-more-morphisms-lemma-integral-closure-smooth-pullback}
Soit $S$ un schéma. Soit $f : Y \to X$ un morphisme lisse
d'espaces algébriques sur $S$. Soit $\mathcal{A}$ un faisceau quasi-cohérent
de $\mathcal{O}_X$-algèbres. La clôture intégrale
de $\mathcal{O}_Y$ dans $f^*\mathcal{A}$ est égale à $f^*\mathcal{A}'$,
où $\mathcal{A}' \subset \mathcal{A}$ est la clôture intégrale de
$\mathcal{O}_X$ dans $\mathcal{A}$.
\end{lemma}

\begin{proof}
Par notre construction de la clôture intégrale, voir
Morphismes d'espaces, Définition
\ref{spaces-morphisms-definition-integral-closure},
on se ramène immédiatement au cas où $X$ et $Y$ sont affines.
Dans ce cas, le résultat est le
Lemme \ref{algebra-lemma-integral-closure-commutes-smooth} d'Algèbre.
\end{proof}

\begin{lemma}[La normalisation commute au changement de base lisse]
\label{spaces-more-morphisms-lemma-normalization-smooth-localization}
Soit $S$ un schéma. Soit
$$
\xymatrix{
Y_2 \ar[r] \ar[d] & Y_1 \ar[d]^f \\
X_2 \ar[r]^\varphi & X_1
}
$$
un carré cartésien d'espaces algébriques sur $S$. Supposons que $f$ soit quasi-compact
et quasi-séparé, et que $\varphi$ soit lisse.
Soit $Y_i \to X_i' \to X_i$ la normalisation de $X_i$ dans $Y_i$.
Alors $X_2' \cong X_2 \times_{X_1} X_1'$.
\end{lemma}

\begin{proof}
La factorisation $Y_1 \to X_1' \to X_1$, après changement de base à $X_2$,
devient une factorisation $Y_2 \to X_2 \times_{X_1} X_1' \to X_1$ et le morphisme
$X_2 \times_{X_1} X_1' \to X_1$ est entier
(Morphismes d'espaces, Lemme \ref{spaces-morphisms-lemma-base-change-integral}).
Nous obtenons donc un morphisme
$h : X_2' \to X_2 \times_{X_1} X_1'$ par la propriété universelle de
Morphismes d'espaces, Lemme
\ref{spaces-morphisms-lemma-characterize-normalization}.
Notons que $X_2'$ est le spectre relatif de la clôture intégrale
de $\mathcal{O}_{X_2}$ dans $f_{2, *}\mathcal{O}_{Y_2}$.
Si $\mathcal{A}' \subset f_{1, *}\mathcal{O}_{Y_1}$ désigne la clôture
intégrale de $\mathcal{O}_{X_2}$, alors $X_2 \times_{X_1} X_1'$ est le
spectre relatif de $\varphi^*\mathcal{A}'$, puisque la formation du
spectre relatif commute à tout changement de base. D'après
Cohomologie des espaces, Lemme
\ref{spaces-cohomology-lemma-flat-base-change-cohomology}
nous savons que $f_{2, *}\mathcal{O}_{Y_2} = \varphi^*f_{1, *}\mathcal{O}_{Y_1}$.
Le résultat découle donc du
Lemme \ref{spaces-more-morphisms-lemma-integral-closure-smooth-pullback}.
\end{proof}









\section{Morphismes de Cohen-Macaulay}
\label{spaces-more-morphisms-section-CM}

\noindent
Ceci est l'analogue de Compléments sur les morphismes, section
\ref{more-morphisms-section-CM}.

\begin{lemma}
\label{spaces-more-morphisms-lemma-CM-local-ring-fibre}
La propriété suivante des morphismes de germes de schémas
\begin{align*}
& \mathcal{P}((X, x) \to (S, s)) = \\
& \text{l'anneau local }
\mathcal{O}_{X_s, x}
\text{ de la fibre est noethérien et Cohen-Macaulay}
\end{align*}
est de nature locale pour la topologie étale sur la source et le but (Descente, Définition
\ref{descent-definition-local-source-target-at-point}).
\end{lemma}

\begin{proof}
Étant donné un diagramme comme dans
Descente, Définition \ref{descent-definition-local-source-target-at-point},
nous obtenons un morphisme étale entre les fibres
$U'_{v'} \to U_v$ qui envoie $u'$ sur $u$ ; voir
Descente, Lemme \ref{descent-lemma-etale-on-fiber}.
Les hensélisés stricts des anneaux locaux
$\mathcal{O}_{U'_{v'}, u'}$ et $\mathcal{O}_{U_v, u}$
sont donc les mêmes. Nous concluons par le
Lemme \ref{more-algebra-lemma-henselization-CM} de Compléments sur l'algèbre.
\end{proof}

\begin{definition}
\label{spaces-more-morphisms-definition-CM}
Soit $S$ un schéma.
Soit $f : X \to Y$ un morphisme d'espaces algébriques sur $S$.
Supposons que les fibres de $f$ soient localement noethériennes
(Diviseurs sur les espaces, Définition
\ref{spaces-divisors-definition-locally-Noetherian-fibre}).
\begin{enumerate}
\item Soient $x \in |X|$ et $y = f(x)$. Nous disons que $f$ est
{\it Cohen-Macaulay en $x$} si $f$ est plat en $x$ et si
les conditions équivalentes de
Morphismes d'espaces, Lemme
\ref{spaces-morphisms-lemma-local-source-target-at-point}
sont satisfaites pour la propriété $\mathcal{P}$ décrite dans le
Lemme \ref{spaces-more-morphisms-lemma-CM-local-ring-fibre}.
\item Nous disons que $f$ est un {\it morphisme de Cohen-Macaulay} si $f$ est
Cohen-Macaulay en tout point de $X$.
\end{enumerate}
\end{definition}

\noindent
Voici une reformulation.

\begin{lemma}
\label{spaces-more-morphisms-lemma-CM}
Soit $S$ un schéma. Soit $f : X \to Y$ un morphisme d'espaces algébriques
sur $S$. Supposons que les fibres de $f$ soient localement noethériennes.
Les conditions suivantes sont équivalentes :
\begin{enumerate}
\item $f$ est Cohen-Macaulay,
\item $f$ est plat et il existe un morphisme étale surjectif $V \to Y$
où $V$ est un schéma, tel que les fibres de $X_V \to V$
soient des espaces algébriques de Cohen-Macaulay, et
\item $f$ est plat et, pour tout morphisme étale $V \to Y$
où $V$ est un schéma, les fibres de $X_V \to V$
sont des espaces algébriques de Cohen-Macaulay.
\end{enumerate}
Étant donné $x \in |X|$ d'image $y \in |Y|$, les conditions suivantes sont
équivalentes :
\begin{enumerate}
\item[(a)] $f$ est Cohen-Macaulay en $x$, et
\item[(b)] $\mathcal{O}_{Y, \overline{y}} \to \mathcal{O}_{X, \overline{x}}$
est plat et
$\mathcal{O}_{X, \overline{x}}/
\mathfrak m_{\overline{y}}\mathcal{O}_{X, \overline{x}}$ est Cohen-Macaulay.
\end{enumerate}
\end{lemma}

\begin{proof}
Étant donné un morphisme étale $V \to Y$ où $V$ est un schéma,
choisissons un schéma $U$ et un morphisme étale surjectif $U \to X \times_Y V$.
Considérons le diagramme commutatif
$$
\xymatrix{
U \ar[d] \ar[r] & V \ar[d] \\
X \ar[r] & Y
}
$$
Soit $u \in U$ d'images $x \in |X|$, $y \in |Y|$ et $v \in V$.
Alors $f$ est Cohen-Macaulay en $x$ si et seulement si $U \to V$ l'est
en $u$ (par définition). De plus, le morphisme
$U_v \to X_v = (X_V)_v$ est étale surjectif. Par conséquent, le schéma $U_v$ est
Cohen-Macaulay si et seulement si l'espace algébrique $X_v$ l'est.
L'équivalence de (1), (2) et (3) découle donc de
l'équivalence correspondante pour les morphismes de
schémas ; voir Compléments sur les morphismes, Lemme \ref{more-morphisms-lemma-CM} ;
l'argument est formel.

\medskip\noindent
Démonstration de l'équivalence de (a) et (b). L'équivalence correspondante
pour la platitude est donnée par Morphismes d'espaces, Lemme
\ref{spaces-morphisms-lemma-flat-at-point-etale-local-rings}.
Nous pouvons donc supposer $f$ plat en $x$ pour démontrer l'équivalence.
Considérons un diagramme et $x, y, u, v$ comme ci-dessus. Alors
$\mathcal{O}_{Y, \overline{y}} \to \mathcal{O}_{X, \overline{x}}$
est égal à l'homomorphisme
$\mathcal{O}_{V, v}^{sh} \to \mathcal{O}_{U, u}^{sh}$
entre les hensélisés stricts des anneaux locaux ; voir
Propriétés des espaces, Lemme
\ref{spaces-properties-lemma-describe-etale-local-ring}.
Nous avons donc
$$
\mathcal{O}_{X, \overline{x}}/
\mathfrak m_{\overline{y}}\mathcal{O}_{X, \overline{x}} =
(\mathcal{O}_{U, u}/\mathfrak m_v \mathcal{O}_{U, u})^{sh}
$$
d'après Algèbre, Lemme \ref{algebra-lemma-quotient-strict-henselization}.
Il reste donc à montrer que l'anneau local noethérien
$\mathcal{O}_{U, u}/\mathfrak m_v \mathcal{O}_{U, u}$
est Cohen-Macaulay si et seulement si son hensélisé strict l'est.
C'est Compléments sur l'algèbre, Lemme \ref{more-algebra-lemma-henselization-CM}.
\end{proof}

\begin{lemma}
\label{spaces-more-morphisms-lemma-composition-CM}
Soit $S$ un schéma.
Soient $f : X \to Y$ et $g : Y \to Z$ des morphismes d'espaces algébriques
sur $S$. Supposons que les
fibres de $f$, $g$ et $g \circ f$ soient localement noethériennes.
Soit $x \in |X|$ d'images $y \in |Y|$ et $z \in |Z|$.
\begin{enumerate}
\item Si $f$ est Cohen-Macaulay en $x$ et $g$ est Cohen-Macaulay
en $f(x)$, alors $g \circ f$ est Cohen-Macaulay en $x$.
\item Si $f$ et $g$ sont Cohen-Macaulay, alors $g \circ f$ l'est.
\item Si $g \circ f$ est Cohen-Macaulay en $x$ et $f$ est plat en $x$,
alors $f$ est Cohen-Macaulay en $x$ et $g$ l'est en $f(x)$.
\item Si $f \circ g$ est Cohen-Macaulay et $f$ est plat, alors
$f$ est Cohen-Macaulay et $g$ l'est en tout point de
l'image de $f$.
\end{enumerate}
\end{lemma}

\begin{proof}
Le résultat correspondant pour les schémas, appliqué localement pour la
topologie étale, donne l'assertion ; voir
Compléments sur les morphismes, Lemme \ref{more-morphisms-lemma-composition-CM}.
On peut aussi utiliser l'équivalence de (a) et (b) du
Lemme \ref{spaces-more-morphisms-lemma-CM}. Nous considérons alors l'homomorphisme local
d'anneaux locaux noethériens
$$
\mathcal{O}_{Y, \overline{y}}/
\mathfrak m_{\overline{z}}\mathcal{O}_{Y, \overline{y}}
\longrightarrow
\mathcal{O}_{X, \overline{x}}/
\mathfrak m_{\overline{z}}\mathcal{O}_{X, \overline{x}}
$$
dont la fibre est
$$
\mathcal{O}_{X, \overline{x}}/
\mathfrak m_{\overline{y}}\mathcal{O}_{X, \overline{x}}
$$
et nous appliquons Algèbre, Lemme \ref{algebra-lemma-CM-goes-up}.
\end{proof}

\begin{lemma}
\label{spaces-more-morphisms-lemma-flat-morphism-from-CM}
Soit $S$ un schéma.
Soit $f : X \to Y$ un morphisme plat entre espaces algébriques localement noethériens
sur $S$.
Si $X$ est Cohen-Macaulay, alors $f$ est Cohen-Macaulay et
$\mathcal{O}_{Y, f(\overline{x})}$ est Cohen-Macaulay pour tout $x \in |X|$.
\end{lemma}

\begin{proof}
En traduisant l'énoncé en algèbre à l'aide du Lemme \ref{spaces-more-morphisms-lemma-CM}
(comparer avec la démonstration du
Lemme \ref{spaces-more-morphisms-lemma-composition-CM}), le résultat découle du
Lemme \ref{algebra-lemma-CM-goes-up} d'Algèbre.
\end{proof}

\begin{lemma}
\label{spaces-more-morphisms-lemma-base-change-CM}
Soit $S$ un schéma.
Soit $f : X \to Y$ un morphisme d'espaces algébriques sur $S$.
Supposons que les fibres de $f$ soient localement noethériennes.
Supposons que $Y' \to Y$ soit localement de type fini. Soit $f' : X' \to Y'$
le changement de base de $f$.
Soit $x' \in |X'|$ un point d'image $x \in |X|$.
\begin{enumerate}
\item Si $f$ est Cohen-Macaulay en $x$, alors
$f' : X' \to Y'$ est Cohen-Macaulay en $x'$.
\item Si $f$ est plat en $x$ et $f'$ est Cohen-Macaulay en $x'$, alors $f$
est Cohen-Macaulay en $x$.
\item Si $Y' \to Y$ est plat en $f'(x')$ et $f'$ est Cohen-Macaulay en
$x'$, alors $f$ est Cohen-Macaulay en $x$.
\end{enumerate}
\end{lemma}

\begin{proof}
Notons $y \in |Y|$ et $y' \in |Y'|$ les images de $x'$.
Choisissons un morphisme étale surjectif $V \to Y$ où $V$ est un schéma.
Choisissons un morphisme étale surjectif $U \to X \times_Y V$ où
$U$ est un schéma.
Choisissons un morphisme étale surjectif $V' \to Y' \times_Y V$
où $V'$ est un schéma.
Alors $U' = U \times_V V'$ est un schéma muni
d'un morphisme étale surjectif $U' \to X'$.
Choisissons $u' \in U'$ qui s'envoie sur $x'$. Notons $u \in U$ l'image de $u'$.
Le lemme découle alors du lemme appliqué à
$U \to V$, à son changement de base $U' \to V'$ et aux
points $u'$ et $u$ (cela découle des définitions).
Il découle donc du cas des schémas ; voir
Compléments sur les morphismes, Lemme \ref{more-morphisms-lemma-base-change-CM}.
\end{proof}

\begin{lemma}
\label{spaces-more-morphisms-lemma-flat-finite-presentation-CM-open}
Soit $S$ un schéma.
Soit $f : X \to Y$ un morphisme d'espaces algébriques sur $S$
qui est plat et localement de présentation finie. Posons
$$
W = \{x \in |X| : f\text{ est Cohen-Macaulay en }x\}
$$
Alors $W$ est ouvert dans $|X|$ et la formation de $W$
commute à tout changement de base de $f$ :
pour tout morphisme $g : Y' \to Y$, considérons
le changement de base $f' : X' \to Y'$ de $f$ et la
projection $g' : X' \to X$. Alors l'ensemble correspondant
$W'$ pour le morphisme $f'$ est donné par $W' = (g')^{-1}(W)$.
\end{lemma}

\begin{proof}
Choisissons un diagramme commutatif
$$
\xymatrix{
U \ar[d] \ar[r] & V \ar[d] \\
X \ar[r] & Y
}
$$
dont les flèches verticales sont étales et où $U$ et $V$ sont des schémas.
Soit $u \in U$ d'image $x \in |X|$.
Alors $f$ est Cohen-Macaulay en $x$ si et seulement si $U \to V$ l'est
en $u$ (par définition).
Nous nous ramenons ainsi au cas du morphisme $U \to V$.
Voir Compléments sur les morphismes, Lemme
\ref{more-morphisms-lemma-flat-finite-presentation-CM-open}.
\end{proof}

\begin{lemma}
\label{spaces-more-morphisms-lemma-lfp-CM-relative-dimension}
Soient $S$ un schéma et $f : X \to Y$ un
morphisme d'espaces algébriques sur $S$. Supposons que
$f$ soit localement de présentation finie et Cohen-Macaulay.
Il existe alors des sous-schémas ouverts et fermés $X_d \subset X$
tels que $X = \coprod_{d \geq 0} X_d$ et que $f|_{X_d} : X_d \to Y$
soit de dimension relative $d$.
\end{lemma}

\begin{proof}
Choisissons un diagramme commutatif
$$
\xymatrix{
U \ar[d] \ar[r] & V \ar[d] \\
X \ar[r] & Y
}
$$
dont les flèches verticales sont étales et où $U$ et $V$ sont des schémas.
Alors $U \to V$ est localement de présentation finie et Cohen-Macaulay
(cela résulte immédiatement de nos définitions).
Nous avons donc une décomposition $U = \coprod_{d \geq 0} U_d$
en sous-schémas ouverts et fermés telle que $f|_{U_d} : U_d \to V$
soit de dimension relative $d$ ; voir Morphismes, Lemme
\ref{morphisms-lemma-flat-finite-presentation-CM-fibres-relative-dimension}.
Soit $u \in U$ d'image $x \in |X|$. Alors
$f$ est de dimension relative $d$ en $x$ si et seulement si
$U \to V$ est de dimension relative $d$ en $u$
(cela découle de nos définitions).
Nous voyons ainsi que $U_d$ est l'image inverse
d'une partie $X_d \subset |X|$ qui est nécessairement
ouverte et fermée. Notons $X_d$ le sous-espace algébrique
ouvert et fermé correspondant de $X$ ; le lemme est démontré.
\end{proof}







\section{Morphismes de Gorenstein}
\label{spaces-more-morphisms-section-gorenstein}

\noindent
Ceci est l'analogue de Dualité pour les schémas, section
\ref{duality-section-gorenstein-morphisms}.

\begin{lemma}
\label{spaces-more-morphisms-lemma-gorenstein-local-ring-fibre}
La propriété suivante des morphismes de germes de schémas
\begin{align*}
& \mathcal{P}((X, x) \to (S, s)) = \\
& \text{l'anneau local }
\mathcal{O}_{X_s, x}
\text{ de la fibre est noethérien et de Gorenstein}
\end{align*}
est de nature locale pour la topologie étale sur la source et le but (Descente, Définition
\ref{descent-definition-local-source-target-at-point}).
\end{lemma}

\begin{proof}
Étant donné un diagramme comme dans
Descente, Définition \ref{descent-definition-local-source-target-at-point},
nous obtenons un morphisme étale entre les fibres
$U'_{v'} \to U_v$ qui envoie $u'$ sur $u$ ; voir
Descente, Lemme \ref{descent-lemma-etale-on-fiber}.
Ainsi, $\mathcal{O}_{U_v, u} \to \mathcal{O}_{U'_{v'}, u'}$
est la localisation d'un homomorphisme étale d'anneaux. Par conséquent,
le premier anneau est noethérien si et seulement si le second l'est ; voir
Compléments sur l'algèbre, Lemme \ref{more-algebra-lemma-Noetherian-etale-extension}.
Puisque $\mathcal{O}_{U'_{v'}, u'}/\mathfrak m_u \mathcal{O}_{U'_{v'}, u'}
= \kappa(u')$ (Algèbre, Lemme \ref{algebra-lemma-etale-at-prime})
est un anneau de Gorenstein, nous voyons que
$\mathcal{O}_{U_v, u}$ est de Gorenstein si et seulement si
$\mathcal{O}_{U'_{v'}, u'}$ l'est, d'après
Complexes dualisants, Lemme \ref{dualizing-lemma-flat-under-gorenstein}.
\end{proof}

\begin{definition}
\label{spaces-more-morphisms-definition-gorenstein}
Soit $S$ un schéma.
Soit $f : X \to Y$ un morphisme d'espaces algébriques sur $S$.
Supposons que les fibres de $f$ soient localement noethériennes
(Diviseurs sur les espaces, Définition
\ref{spaces-divisors-definition-locally-Noetherian-fibre}).
\begin{enumerate}
\item Soient $x \in |X|$ et $y = f(x)$. Nous disons que $f$ est
{\it de Gorenstein en $x$} si $f$ est plat en $x$ et si
les conditions équivalentes de
Morphismes d'espaces, Lemme
\ref{spaces-morphisms-lemma-local-source-target-at-point}
sont satisfaites pour la propriété $\mathcal{P}$ décrite dans le
Lemme \ref{spaces-more-morphisms-lemma-gorenstein-local-ring-fibre}.
\item Nous disons que $f$ est un {\it morphisme de Gorenstein} si $f$ est
de Gorenstein en tout point de $X$.
\end{enumerate}
\end{definition}

\noindent
Voici une reformulation.

\begin{lemma}
\label{spaces-more-morphisms-lemma-gorenstein}
Soit $S$ un schéma. Soit $f : X \to Y$ un morphisme d'espaces algébriques
sur $S$. Supposons que les fibres de $f$ soient localement noethériennes.
Les conditions suivantes sont équivalentes :
\begin{enumerate}
\item $f$ est de Gorenstein,
\item $f$ est plat et il existe un morphisme étale surjectif $V \to Y$
où $V$ est un schéma, tel que les fibres de $X_V \to V$
soient des espaces algébriques de Gorenstein, et
\item $f$ est plat et, pour tout morphisme étale $V \to Y$
où $V$ est un schéma, les fibres de $X_V \to V$
sont des espaces algébriques de Gorenstein.
\end{enumerate}
Étant donné $x \in |X|$ d'image $y \in |Y|$, les conditions suivantes sont
équivalentes :
\begin{enumerate}
\item[(a)] $f$ est de Gorenstein en $x$, et
\item[(b)] $\mathcal{O}_{Y, \overline{y}} \to \mathcal{O}_{X, \overline{x}}$
est plat et
$\mathcal{O}_{X, \overline{x}}/
\mathfrak m_{\overline{y}}\mathcal{O}_{X, \overline{x}}$ est de Gorenstein.
\end{enumerate}
\end{lemma}

\begin{proof}
Étant donné un morphisme étale $V \to Y$ où $V$ est un schéma,
choisissons un schéma $U$ et un morphisme étale surjectif $U \to X \times_Y V$.
Considérons le diagramme commutatif
$$
\xymatrix{
U \ar[d] \ar[r] & V \ar[d] \\
X \ar[r] & Y
}
$$
Soit $u \in U$ d'images $x \in |X|$, $y \in |Y|$ et $v \in V$.
Alors $f$ est de Gorenstein en $x$ si et seulement si $U \to V$ l'est
en $u$ (par définition). De plus, le morphisme
$U_v \to X_v = (X_V)_v$ est étale surjectif. Par conséquent, le schéma $U_v$ est
de Gorenstein si et seulement si l'espace algébrique $X_v$ l'est.
L'équivalence de (1), (2) et (3) découle donc de
l'équivalence correspondante pour les morphismes de
schémas ; voir Dualité pour les schémas, Lemme \ref{duality-lemma-gorenstein} ;
l'argument est formel.

\medskip\noindent
Démonstration de l'équivalence de (a) et (b). L'équivalence correspondante
pour la platitude est donnée par Morphismes d'espaces, Lemme
\ref{spaces-morphisms-lemma-flat-at-point-etale-local-rings}.
Nous pouvons donc supposer $f$ plat en $x$ pour démontrer l'équivalence.
Considérons un diagramme et $x, y, u, v$ comme ci-dessus. Alors
$\mathcal{O}_{Y, \overline{y}} \to \mathcal{O}_{X, \overline{x}}$
est égal à l'homomorphisme
$\mathcal{O}_{V, v}^{sh} \to \mathcal{O}_{U, u}^{sh}$
entre les hensélisés stricts des anneaux locaux ; voir
Propriétés des espaces, Lemme
\ref{spaces-properties-lemma-describe-etale-local-ring}.
Nous avons donc
$$
\mathcal{O}_{X, \overline{x}}/
\mathfrak m_{\overline{y}}\mathcal{O}_{X, \overline{x}} =
(\mathcal{O}_{U, u}/\mathfrak m_v \mathcal{O}_{U, u})^{sh}
$$
d'après Algèbre, Lemme \ref{algebra-lemma-quotient-strict-henselization}.
Il reste donc à montrer que l'anneau local noethérien
$\mathcal{O}_{U, u}/\mathfrak m_v \mathcal{O}_{U, u}$
est de Gorenstein si et seulement si son hensélisé strict l'est.
Cela découle immédiatement de
Complexes dualisants, Lemme
\ref{dualizing-lemma-completion-henselization-dualizing}
et de la définition d'un anneau local de Gorenstein comme un
anneau local noethérien qui est un complexe dualisant sur lui-même.
\end{proof}

\begin{lemma}
\label{spaces-more-morphisms-lemma-composition-gorenstein}
Soit $S$ un schéma.
Soient $f : X \to Y$ et $g : Y \to Z$ des morphismes d'espaces algébriques
sur $S$. Supposons que les
fibres de $f$, $g$ et $g \circ f$ soient localement noethériennes.
Soit $x \in |X|$ d'images $y \in |Y|$ et $z \in |Z|$.
\begin{enumerate}
\item Si $f$ est de Gorenstein en $x$ et $g$ est de Gorenstein
en $f(x)$, alors $g \circ f$ est de Gorenstein en $x$.
\item Si $f$ et $g$ sont de Gorenstein, alors $g \circ f$ l'est.
\item Si $g \circ f$ est de Gorenstein en $x$ et $f$ est plat en $x$,
alors $f$ est de Gorenstein en $x$ et $g$ l'est en $f(x)$.
\item Si $f \circ g$ est de Gorenstein et $f$ est plat, alors
$f$ est de Gorenstein et $g$ l'est en tout point de
l'image de $f$.
\end{enumerate}
\end{lemma}

\begin{proof}
Le résultat correspondant pour les schémas, appliqué localement pour la
topologie étale, donne l'assertion ; voir Dualité pour les schémas, Lemme
\ref{duality-lemma-composition-gorenstein}.
On peut aussi utiliser l'équivalence de (a) et (b) du
Lemme \ref{spaces-more-morphisms-lemma-gorenstein}. Nous considérons alors l'homomorphisme local
d'anneaux locaux noethériens
$$
\mathcal{O}_{Y, \overline{y}}/
\mathfrak m_{\overline{z}}\mathcal{O}_{Y, \overline{y}}
\longrightarrow
\mathcal{O}_{X, \overline{x}}/
\mathfrak m_{\overline{z}}\mathcal{O}_{X, \overline{x}}
$$
dont la fibre est
$$
\mathcal{O}_{X, \overline{x}}/
\mathfrak m_{\overline{y}}\mathcal{O}_{X, \overline{x}}
$$
et nous appliquons Complexes dualisants, Lemme
\ref{dualizing-lemma-flat-under-gorenstein}.
\end{proof}

\begin{lemma}
\label{spaces-more-morphisms-lemma-flat-morphism-from-gorenstein}
Soit $S$ un schéma.
Soit $f : X \to Y$ un morphisme plat entre espaces algébriques localement noethériens
sur $S$.
Si $X$ est de Gorenstein, alors $f$ est de Gorenstein et
$\mathcal{O}_{Y, f(\overline{x})}$ est de Gorenstein pour tout $x \in |X|$.
\end{lemma}

\begin{proof}
En traduisant l'énoncé en algèbre à l'aide du Lemme \ref{spaces-more-morphisms-lemma-gorenstein}
(comparer avec la démonstration du
Lemme \ref{spaces-more-morphisms-lemma-composition-gorenstein}), le résultat découle du
Lemme \ref{dualizing-lemma-flat-under-gorenstein} de Complexes dualisants.
\end{proof}

\begin{lemma}
\label{spaces-more-morphisms-lemma-base-change-gorenstein}
Soit $S$ un schéma.
Soit $f : X \to Y$ un morphisme d'espaces algébriques sur $S$.
Supposons que les fibres de $f$ soient localement noethériennes.
Supposons que $Y' \to Y$ soit localement de type fini. Soit $f' : X' \to Y'$
le changement de base de $f$.
Soit $x' \in |X'|$ un point d'image $x \in |X|$.
\begin{enumerate}
\item Si $f$ est de Gorenstein en $x$, alors
$f' : X' \to Y'$ est de Gorenstein en $x'$.
\item Si $f$ est plat en $x$ et $f'$ est de Gorenstein en $x'$, alors $f$
est de Gorenstein en $x$.
\item Si $Y' \to Y$ est plat en $f'(x')$ et $f'$ est de Gorenstein en
$x'$, alors $f$ est de Gorenstein en $x$.
\end{enumerate}
\end{lemma}

\begin{proof}
Notons $y \in |Y|$ et $y' \in |Y'|$ les images de $x'$.
Choisissons un morphisme étale surjectif $V \to Y$ où $V$ est un schéma.
Choisissons un morphisme étale surjectif $U \to X \times_Y V$ où
$U$ est un schéma.
Choisissons un morphisme étale surjectif $V' \to Y' \times_Y V$
où $V'$ est un schéma.
Alors $U' = U \times_V V'$ est un schéma muni
d'un morphisme étale surjectif $U' \to X'$.
Choisissons $u' \in U'$ qui s'envoie sur $x'$. Notons $u \in U$ l'image de $u'$.
Le lemme découle alors du lemme appliqué à
$U \to V$, à son changement de base $U' \to V'$ et aux
points $u'$ et $u$ (cela découle des définitions).
Il découle donc du cas des schémas ; voir
Dualité pour les schémas, Lemme \ref{duality-lemma-base-change-gorenstein}.
\end{proof}

\begin{lemma}
\label{spaces-more-morphisms-lemma-flat-finite-presentation-gorenstein-open}
Soit $S$ un schéma.
Soit $f : X \to Y$ un morphisme d'espaces algébriques sur $S$
qui est plat et localement de présentation finie. Posons
$$
W = \{x \in |X| : f\text{ est de Gorenstein en }x\}
$$
Alors $W$ est ouvert dans $|X|$ et la formation de $W$
commute à tout changement de base de $f$ :
pour tout morphisme $g : Y' \to Y$, considérons
le changement de base $f' : X' \to Y'$ de $f$ et la
projection $g' : X' \to X$. Alors l'ensemble correspondant
$W'$ pour le morphisme $f'$ est donné par $W' = (g')^{-1}(W)$.
\end{lemma}

\begin{proof}
Choisissons un diagramme commutatif
$$
\xymatrix{
U \ar[d] \ar[r] & V \ar[d] \\
X \ar[r] & Y
}
$$
Soit $u \in U$ d'image $x \in |X|$.
Alors $f$ est de Gorenstein en $x$ si et seulement si $U \to V$ l'est
en $u$ (par définition).
Nous nous ramenons ainsi au cas du morphisme $U \to V$
de schémas. L'ouverture est démontrée dans
Dualité pour les schémas, Lemme
\ref{duality-lemma-flat-finite-presentation-characterize-gorenstein}
et la compatibilité au changement de base dans
Dualité pour les schémas, Lemme \ref{duality-lemma-flat-lft-base-change-gorenstein}.
\end{proof}










\section{Tranchage des morphismes de Cohen-Macaulay}
\label{spaces-more-morphisms-section-slice}

\noindent
Soit $S$ un schéma. Soit $X$ un espace algébrique sur $S$.
Soient $f_1, \ldots, f_r \in \Gamma(X, \mathcal{O}_X)$. Dans ce cas, nous
notons $V(f_1, \ldots, f_r)$ le {\it sous-espace fermé de $X$ défini par
$f_1, \ldots, f_r$}. Plus précisément, on peut définir $V(f_1, \ldots, f_r)$
comme le sous-espace fermé de $X$ correspondant au faisceau quasi-cohérent
d'idéaux engendré par $f_1, \ldots, f_r$ ; voir
Morphismes d'espaces, Lemme
\ref{spaces-morphisms-lemma-closed-immersion-ideals}.
On peut aussi choisir une présentation $X = U/R$ et considérer le
sous-schéma fermé $Z \subset U$ défini par $f_1|U, \ldots, f_r|_U$.
Il est clair que $Z$ est un sous-schéma fermé invariant par $R$ (voir
Groupoïdes, Définition \ref{groupoids-definition-invariant-open})
et l'on peut poser $V(f_1, \ldots, f_r) = Z/R_Z$.

\begin{lemma}
\label{spaces-more-morphisms-lemma-slice}
Soit $S$ un schéma. Considérons un diagramme cartésien
$$
\xymatrix{
X \ar[d] & F \ar[l]^p \ar[d] \\
Y & \Spec(k) \ar[l]
}
$$
où $X \to Y$ est un morphisme d'espaces algébriques sur $S$
qui est plat et localement de présentation finie, et où
$k$ est un corps sur $S$. Soient $f_1, \ldots, f_r \in \Gamma(X, \mathcal{O}_X)$
et $z \in |F|$ tels que les images de $f_1, \ldots, f_r$ forment une suite régulière
dans l'anneau local $\mathcal{O}_{F, \overline{z}}$.
Alors, quitte à remplacer $X$ par un sous-espace ouvert contenant $p(z)$, le morphisme
$$
V(f_1, \ldots, f_r) \longrightarrow Y
$$
est plat et localement de présentation finie.
\end{lemma}

\begin{proof}
Posons $Z = V(f_1, \ldots, f_r)$. Il est clair que $Z \to X$ est localement de
présentation finie ; la composée $Z \to Y$ est donc localement de
présentation finie ; voir
Morphismes d'espaces,
Lemme \ref{spaces-morphisms-lemma-composition-finite-presentation}.
Il suffit donc de montrer que $Z \to Y$ est plat dans un voisinage de $p(z)$.
Soit $k'/k$ une extension de corps. Alors le morphisme de
$F' = F \times_{\Spec(k)} \Spec(k')$ vers
$F$ est surjectif et plat ; on peut donc trouver un point $z' \in |F'|$ qui s'envoie sur $z$
et l'homomorphisme local
$\mathcal{O}_{F, \overline{z}} \to \mathcal{O}_{F', \overline{z}'}$ est
plat ; voir
Morphismes d'espaces,
Lemme \ref{spaces-morphisms-lemma-flat-at-point-etale-local-rings}.
Par conséquent, les images de $f_1, \ldots, f_r$ dans
$\mathcal{O}_{F', \overline{z}'}$ forment encore une suite régulière ; voir
Algèbre, Lemme \ref{algebra-lemma-flat-increases-depth}.
Ainsi, au cours de la démonstration, nous pouvons remplacer $k$ par une extension de corps.
En particulier, nous pouvons supposer que $z \in |F|$ provient d'une section
$z : \Spec(k) \to F$ du morphisme structural $F \to \Spec(k)$.

\medskip\noindent
Choisissons un schéma $V$ et un morphisme étale surjectif
$V \to Y$. Choisissons un schéma $U$ et un morphisme étale surjectif
$U \to X \times_Y V$. Après avoir éventuellement agrandi encore $k$, nous pouvons
supposer que $\Spec(k) \to F \to X$ se factorise par $U$ (puisque
$U \to X$ est surjectif). Soit
$u : \Spec(k) \to U$ une telle factorisation et notons $v \in V$
l'image de $u$. Notons que les morphismes
$$
U_v \times_{\Spec(\kappa(v))} \Spec(k) =
U \times_V \Spec(k) \to U \times_Y \Spec(k) \to F
$$
sont étales (le premier comme changement de base de $V \to V \times_Y V$ et
le second comme changement de base de $U \to X$). De plus, par construction,
le point $u : \Spec(k) \to U$ définit un point de l'espace tout à gauche
qui s'envoie sur $z$ à droite. Par conséquent, les éléments
$f_1, \ldots, f_r$ ont pour images une suite régulière dans l'anneau local
situé à droite de l'homomorphisme suivant
$$
\mathcal{O}_{U_v, u}
\longrightarrow
\mathcal{O}_{U_v \times_{\Spec(\kappa(v)} \Spec(k), \overline{u}}
=
\mathcal{O}_{U \times_V \Spec(k), \overline{u}}.
$$
Mais, puisque la flèche affichée est plate (en combinant
Compléments sur la platitude, Lemme \ref{flat-lemma-flat-up-down-henselization}
et
Morphismes d'espaces,
Lemme \ref{spaces-morphisms-lemma-flat-at-point-etale-local-rings}),
nous déduisons de
Algèbre, Lemme \ref{algebra-lemma-flat-increases-depth},
que les images de $f_1, \ldots, f_r$ forment une suite régulière dans
$\mathcal{O}_{U_v, u}$. D'après
Compléments sur les morphismes, Lemme \ref{more-morphisms-lemma-slice-given-elements},
nous concluons que le morphisme de schémas
$$
V(f_1, \ldots, f_r) \times_X U = V(f_1|_U, \ldots, f_r|_U) \to V
$$
est plat dans un voisinage ouvert $U'$ de $u$. Soit $X' \subset X$
le sous-espace ouvert correspondant à l'image de
$|U'| \to |X|$ (voir
Propriétés des espaces, Lemmes
\ref{spaces-properties-lemma-topology-points} et
\ref{spaces-properties-lemma-open-subspaces}).
Nous concluons que $V(f_1, \ldots, f_r) \cap X' \to Y$ est plat
(voir
Morphismes d'espaces, Définition \ref{spaces-morphisms-definition-flat}),
car
nous avons le diagramme commutatif
$$
\xymatrix{
V(f_1, \ldots, f_r) \times_X U' \ar[d]_a \ar[r] & V \ar[d]^b \\
V(f_1, \ldots, f_r) \cap X' \ar[r] & Y
}
$$
où $a, b$ sont étales et $a$ est surjectif.
\end{proof}






\section{Fibres réduites}
\label{spaces-more-morphisms-section-reduced}

\noindent
Cette section est l'analogue de
Compléments sur les morphismes, section \ref{more-morphisms-section-reduced}.

\begin{lemma}
\label{spaces-more-morphisms-lemma-geometrically-reduced-fibre}
Soit $S$ un schéma. Soit $f : X \to Y$ un
morphisme d'espaces algébriques sur $S$. Soit $y \in |Y|$.
Les conditions suivantes sont équivalentes :
\begin{enumerate}
\item pour un certain morphisme $\Spec(k) \to Y$ dans la classe d'équivalence
de $y$, l'espace algébrique $X_k$ est géométriquement réduit sur $k$ ;
\item pour tout morphisme $\Spec(k) \to Y$ dans la classe d'équivalence
de $y$, l'espace algébrique $X_k$ est géométriquement réduit sur $k$ ;
\item pour tout morphisme $\Spec(k) \to Y$ dans la classe d'équivalence
de $y$, l'espace algébrique $X_k$ est réduit.
\end{enumerate}
\end{lemma}

\begin{proof}
Cela résulte immédiatement du Lemme
\ref{spaces-over-fields-lemma-geometrically-reduced-upstairs} d'Espaces sur les corps
et de la définition de la relation d'équivalence qui définit $|X|$,
donnée dans
Propriétés des espaces, section \ref{spaces-properties-section-points}.
\end{proof}

\begin{definition}
\label{spaces-more-morphisms-definition-geometrically-reduced-fibre}
Soit $S$ un schéma. Soit $f : X \to Y$ un
morphisme d'espaces algébriques sur $S$. Soit $y \in |Y|$.
Nous disons que {\it la fibre de $f : X \to Y$ en $y$ est géométriquement réduite}
si les conditions équivalentes du
Lemme \ref{spaces-more-morphisms-lemma-geometrically-reduced-fibre} sont satisfaites.
\end{definition}

\noindent
Voici les lemmes obligatoires.

\begin{lemma}
\label{spaces-more-morphisms-lemma-base-change-fibres-geometrically-reduced}
Soit $S$ un schéma. Soient $f : X \to Y$ et $g : Y' \to Y$
des morphismes d'espaces algébriques sur $S$. Notons
$f' : X' \to Y'$ le changement de base de $f$ par $g$. Alors
\begin{align*}
\{y' \in |Y'| :
\text{la fibre de }f' : X' \to Y'\text{ en }y'
\text{ est géométriquement réduite}\} \\
= g^{-1}(\{y \in |Y| :
\text{la fibre de }f : X \to Y\text{ en }y
\text{ est géométriquement réduite}\}).
\end{align*}
\end{lemma}

\begin{proof}
Pour $y' \in |Y'|$, choisissons un morphisme $\Spec(k) \to Y'$
dans la classe d'équivalence de $y'$. Alors $g(y')$ est
représenté par la composée $\Spec(k) \to Y' \to Y$.
Ainsi $X' \times_{Y'} \Spec(k) = X \times_Y \Spec(k)$
et le résultat découle de la définition.
\end{proof}

\begin{lemma}
\label{spaces-more-morphisms-lemma-geometrically-reduced-constructible}
Soit $S$ un schéma. Soit $f : X \to Y$ un morphisme d'espaces algébriques
sur $S$ qui est quasi-compact et
localement de présentation finie. Alors l'ensemble
$$
E = \{y \in |Y| : \text{la fibre de }f : X \to Y\text{ en }y
\text{ est géométriquement réduite}\}
$$
est localement constructible pour la topologie étale.
\end{lemma}

\begin{proof}
Choisissons un schéma affine $V$ et un morphisme étale $V  \to Y$.
L'énoncé signifie que l'image inverse de $E$
dans $|V|$ est constructible. D'après le
Lemme \ref{spaces-more-morphisms-lemma-base-change-fibres-geometrically-reduced},
nous pouvons remplacer $Y$ par $V$, c'est-à-dire supposer que $Y$
est un schéma affine. Alors $X$ est quasi-compact. Choisissons un
schéma affine $U$ et un morphisme étale surjectif $U \to X$.
Pour un morphisme $\Spec(k) \to Y$, le morphisme entre les fibres
$U_k \to X_k$ est étale surjectif. Ainsi $U_k$ est géométriquement
réduit sur $k$ si et seulement si $X_k$ est géométriquement réduit
sur $k$ ; voir Espaces sur les corps, Lemme
\ref{spaces-over-fields-lemma-geometrically-reduced-etale-local}.
Ainsi, l'ensemble $E$ associé à $X \to Y$ est le même que l'ensemble $E$
associé à $U \to Y$.
Nous voyons de cette manière que le lemme découle du cas
des schémas ; voir Compléments sur les morphismes, Lemme
\ref{more-morphisms-lemma-geometrically-reduced-constructible}.
\end{proof}

\begin{lemma}
\label{spaces-more-morphisms-lemma-proper-flat-over-dvr-reduced-fibre}
Soit $X$ un espace algébrique sur un anneau de valuation discrète $R$
dont le morphisme structural $X \to \Spec(R)$ est propre et plat.
Si la fibre spéciale est réduite, alors
$X$ et la fibre générique $X_\eta$ sont tous deux réduits.
\end{lemma}

\begin{proof}
Choisissons un morphisme étale $U \to X$ où $U$ est un schéma affine.
Alors $U$ est de type fini sur $R$. Soit $u \in U$ un point de la fibre spéciale.
L'anneau local $A = \mathcal{O}_{U, u}$ est essentiellement de type
fini sur $R$, donc noethérien. Soit $\pi \in R$ une uniformisante.
Comme $X$ est plat sur $R$, nous voyons que $\pi \in \mathfrak m_A$
est un non diviseur de zéro dans $A$ et, comme la fibre spéciale de $X$
est réduite, l'anneau $A/\pi A$ est réduit.
Si $a \in A$, $a \not = 0$, alors il existe un $n \geq 0$ et un élément
$a' \in A$ tels que $a = \pi^n a'$ et $a' \not \in \pi A$.
Cela résulte du théorème d'intersection de Krull
(Algèbre, Lemme \ref{algebra-lemma-intersect-powers-ideal-module-zero}).
Si $a$ est nilpotent, $a'$ l'est aussi, car $\pi$ est un non diviseur de zéro.
Mais $a'$ s'envoie sur un élément non nul de l'anneau réduit $A/\pi A$,
ce qui est impossible. Ainsi $A$ est réduit. Il s'ensuit qu'il
existe un voisinage ouvert de $u$ dans $U$ qui est réduit
(petit détail omis ; utiliser le fait que $U$ est noethérien).
On peut donc trouver un morphisme étale $U \to X$ où $U$ est un schéma
réduit, tel que tout point de la fibre spéciale de $X$ appartienne
à l'image. Comme $X$ est propre sur $R$, il s'ensuit que
$U \to X$ est surjectif. Ainsi $X$ est réduit. Comme la fibre générique de
$U \to \Spec(R)$ est également réduite (sur les morceaux affines,
elle se calcule par localisation), nous concluons qu'il en est de même
de la fibre générique.
\end{proof}

\begin{lemma}
\label{spaces-more-morphisms-lemma-geometrically-reduced-open}
Soit $S$ un schéma. Soit $f : X \to Y$ un morphisme d'espaces
algébriques sur $S$. Si $f$ est plat, propre et de présentation finie,
alors l'ensemble
$$
E = \{y \in |Y| : \text{la fibre de }f : X \to Y\text{ en }y
\text{ est géométriquement réduite}\}
$$
est ouvert dans $|Y|$.
\end{lemma}

\begin{proof}
Par le Lemme \ref{spaces-more-morphisms-lemma-base-change-fibres-geometrically-reduced},
la formation de $E$ commute au changement de base. Pour vérifier qu'une partie
de $|Y|$ est ouverte, nous pouvons remplacer $Y$ par les membres d'un
recouvrement étale. Nous pouvons donc supposer $Y$ affine.
Alors $Y$ est une limite projective filtrante de schémas
affines de type fini sur $\mathbf{Z}$.
Par conséquent, nous pouvons supposer que $X \to Y$ est le
changement de base de $X_0 \to Y_0$, où $Y_0$ est le spectre d'une algèbre
de type fini sur $\mathbf{Z}$ et $X_0 \to Y_0$ est plat et propre.
Voir Limites d'espaces, Lemmes
\ref{spaces-limits-lemma-descend-finite-presentation},
\ref{spaces-limits-lemma-descend-flat}, et
\ref{spaces-limits-lemma-eventually-proper}. Puisque la formation de
$E$ commute au changement de base (voir ci-dessus),
nous pouvons supposer la base noethérienne.

\medskip\noindent
Supposons $Y$ noethérien. Cet ensemble est constructible d'après le
Lemme \ref{spaces-more-morphisms-lemma-geometrically-reduced-constructible}.
Il suffit donc de montrer que l'ensemble est stable par générisation
(Topologie, Lemme \ref{topology-lemma-characterize-closed-Noetherian}). D'après
Propriétés, Lemme \ref{properties-lemma-locally-Noetherian-specialization-dvr},
nous nous ramenons au cas où $Y = \Spec(R)$, où $R$ est un anneau de
valuation discrète et où la fibre fermée $X_y$ est géométriquement
réduite. Il reste à montrer que la fibre générique $X_\eta$ est géométriquement réduite.

\medskip\noindent
Sinon, il existe une extension finie $L$ du corps des
fractions de $R$ telle que $X_L$ ne soit pas réduit ; voir
Espaces sur les corps, Lemmes
\ref{spaces-over-fields-lemma-perfect-reduced} (caractéristique zéro) et
\ref{spaces-over-fields-lemma-geometrically-reduced-positive-characteristic}
(caractéristique positive). Il existe un anneau de valuation discrète
$R' \subset L$ de corps des fractions $L$ qui domine $R$ ; voir
Algèbre, Lemme \ref{algebra-lemma-integral-closure-Dedekind}.
Après avoir remplacé $R$ par $R'$, nous nous ramenons au
Lemme \ref{spaces-more-morphisms-lemma-proper-flat-over-dvr-reduced-fibre}.
\end{proof}







\section{Composantes connexes des fibres}
\label{spaces-more-morphisms-section-connected}

\noindent
Cette section est l'analogue de
Compléments sur les morphismes, section \ref{more-morphisms-section-connected}.

\begin{lemma}
\label{spaces-more-morphisms-lemma-base-change-fibres-nr-geometrically-connected-components}
Soit $S$ un schéma.
Soit $f : X \to Y$ un morphisme d'espaces algébriques sur $S$. Soit
$$
n_{X/Y} : |Y| \to \{0, 1, 2, 3, \ldots, \infty\}
$$
la fonction qui associe à $y \in Y$ le nombre de composantes connexes
de $X_k$, où $\Spec(k) \to Y$ appartient à la classe d'équivalence
de $y$ et où $k$ est algébriquement clos.
Cette fonction est bien définie et, si $g : Y' \to Y$ est un morphisme, alors
$$
n_{X'/Y'} = n_{X/Y} \circ g
$$
où $X' \to Y'$ est le changement de base de $f$.
\end{lemma}

\begin{proof}
Supposons que $y' \in Y'$ ait pour image $y \in Y$. Soit $\Spec(k') \to Y'$
un morphisme dans la classe d'équivalence de $y'$, avec $k'$ algébriquement clos.
On peut alors choisir un diagramme commutatif
$$
\xymatrix{
\Spec(K) \ar[r] \ar[rd] &
\Spec(k') \ar[r] & Y' \ar[d] \\
& \Spec(k) \ar[r] & Y
}
$$
où $K$ est un corps algébriquement clos.
Le résultat vient de ce que les morphismes de schémas
$$
\xymatrix{
X'_{k'} & (X'_{k'})_K = (X_k)_K \ar[l] \ar[r] & X_k
}
$$
induisent des bijections entre les composantes connexes ; voir
Espaces sur les corps, Lemme
\ref{spaces-over-fields-lemma-separably-closed-field-connected-components}.
Pour en déduire que la fonction est bien définie, prendre $Y' = Y$.
\end{proof}







\section{Dimension des fibres}
\label{spaces-more-morphisms-section-dimension}

\noindent
Cette section est l'analogue de
Compléments sur les morphismes, section \ref{more-morphisms-section-dimension}.

\begin{lemma}
\label{spaces-more-morphisms-lemma-dimension-fibre}
Soit $S$ un schéma. Soit $f : X \to Y$ un morphisme de type fini
d'espaces algébriques sur $S$. Soit $y \in |Y|$. Les grandeurs suivantes
sont égales :
\begin{enumerate}
\item $d = -\infty$ si $y$ n'appartient pas à l'image de $|f|$ et,
sinon, le plus petit entier $d$ tel que $f$ soit de dimension relative $\leq d$
en tout $x \in |X|$ s'envoyant sur $y$ ;
\item la dimension de l'espace algébrique $X_k = \Spec(k) \times_Y X$
pour tout morphisme $\Spec(k) \to Y$ dans la classe d'équivalence qui définit $y$.
\end{enumerate}
\end{lemma}

\begin{proof}
Pour interpréter cet énoncé, il faut consulter
Morphismes d'espaces, Définition
\ref{spaces-morphisms-definition-dimension-fibre},
Propriétés des espaces,
Définition \ref{spaces-properties-definition-dimension},
Propriétés des espaces,
Définition \ref{spaces-properties-definition-dimension-at-point}.
Nous allons montrer que les nombres de (1) et (2) sont égaux pour
un morphisme fixé $\Spec(k) \to Y$.
Choisissons un morphisme étale $V \to Y$, où $V$ est un schéma
affine, ainsi qu'un point $v \in V$ s'envoyant sur $y$.
Puisque $V \times_Y \Spec(k) \to \Spec(k)$ est étale surjectif
(d'après Propriétés des espaces, Lemme \ref{spaces-properties-lemma-points-cartesian}),
on peut trouver une extension séparable finie $k'/k$
(d'après Morphismes, Lemme \ref{morphisms-lemma-etale-over-field})
et un diagramme commutatif
$$
\xymatrix{
\Spec(k') \ar[r] \ar[d] & V \ar[d] \\
\Spec(k) \ar[r] & Y
}
$$
Nous pouvons remplacer $X \to Y$ par $V \times_Y X \to V$ et
$X_k$ par $X_{k'} = \Spec(k') \times_V (V \times_Y X)$,
car cela ne change pas les dimensions considérées d'après
Propriétés des espaces, Lemme
\ref{spaces-properties-lemma-dimension-decent-invariant-under-etale}
et Morphismes d'espaces, Lemme
\ref{spaces-morphisms-lemma-dimension-fibre-after-base-change}.
Ainsi, nous pouvons supposer que $Y$ est un schéma affine.
Dans ce cas, nous pouvons supposer que $k = \kappa(y)$,
car les dimensions de $X_{\kappa(y)}$ et de $X_k$
sont égales d'après le Lemme
\ref{spaces-morphisms-lemma-dimension-fibre-after-base-change} de Morphismes d'espaces
et le fait que, pour un espace algébrique $Z$ sur un corps $K$, la
dimension relative de $Z$ en un point $z \in |Z|$
est égale à $\dim_z(Z)$ par définition.
Supposons $Y$ affine et $k = \kappa(y)$. Alors,
comme $X$ est quasi-compact, nous pouvons choisir un schéma affine $U$ et
un morphisme étale surjectif $U \to X$.
Alors $\dim(X_k) = \dim(U_k) = \max \dim_u(U_k)$
est égal, par définition, au nombre donné en (1).
\end{proof}

\begin{lemma}
\label{spaces-more-morphisms-lemma-base-change-dimension-fibres}
Soit $S$ un schéma. Soit $f : X \to Y$ un morphisme de type fini
d'espaces algébriques sur $S$. Soit
$$
n_{X/Y} : |Y| \to \{-\infty, 0, 1, 2, 3, \ldots\}
$$
la fonction qui associe à $y \in |Y|$
l'entier étudié au Lemme \ref{spaces-more-morphisms-lemma-dimension-fibre}.
Si $g : Y' \to Y$ est un morphisme, alors
$$
n_{X'/Y'} = n_{X/Y} \circ |g|
$$
où $X' \to Y'$ est le changement de base de $f$.
\end{lemma}

\begin{proof}
Cela résulte immédiatement du
Lemme \ref{spaces-more-morphisms-lemma-dimension-fibre}.
\end{proof}

\begin{lemma}
\label{spaces-more-morphisms-lemma-dimension-fibres-flat}
Soit $S$ un schéma.
Soit $f : X \to Y$ un morphisme plat et de présentation finie
d'espaces algébriques sur $S$. Soit
$n_{X/Y}$ la fonction sur $Y$ donnant la dimension des fibres de $f$
introduite au Lemme \ref{spaces-more-morphisms-lemma-base-change-dimension-fibres}.
Alors $n_{X/Y}$ est semi-continue inférieurement.
\end{lemma}

\begin{proof}
Soit $V \to Y$ un morphisme étale surjectif, où $V$ est un schéma.
Si nous pouvons montrer que la composée $n_{X/Y} \circ |g|$
est semi-continue inférieurement, alors le lemme en découle, car l'application $|g|$ est ouverte.
Ainsi, nous pouvons supposer que $Y$ est un schéma.
En raisonnant localement, nous pouvons supposer que $V$ est un schéma affine.
On peut alors choisir un schéma affine $U$ et un morphisme
étale surjectif $U \to X$. Alors $n_{X/Y} = n_{U/Y}$.
Ainsi, nous pouvons supposer que $X$ et $Y$ sont tous deux des schémas.
Dans ce cas, le lemme résulte de
Compléments sur les morphismes, Lemme \ref{more-morphisms-lemma-dimension-fibres-flat}.
\end{proof}

\begin{lemma}
\label{spaces-more-morphisms-lemma-dimension-fibres-proper}
Soit $S$ un schéma.
Soit $f : X \to Y$ un morphisme propre d'espaces algébriques sur $S$. Soit
$n_{X/Y}$ la fonction sur $Y$ donnant la dimension des fibres de $f$
introduite au Lemme \ref{spaces-more-morphisms-lemma-base-change-dimension-fibres}.
Alors $n_{X/Y}$ est semi-continue supérieurement.
\end{lemma}

\begin{proof}
Soit $Z_d = \{x \in |X| :
\text{la fibre de }f\text{ en }x\text{ est de dimension }> d\}$.
Alors $Z_d$ est une partie fermée de $|X|$ d'après
Morphismes d'espaces, Lemme
\ref{spaces-morphisms-lemma-openness-bounded-dimension-fibres}.
Puisque $f$ est propre, $f(Z_d)$ est fermé dans $|Y|$.
Comme $y \in f(Z_d) \Leftrightarrow n_{X/Y}(y) > d$,
le lemme est démontré.
\end{proof}

\begin{lemma}
\label{spaces-more-morphisms-lemma-dimension-fibres-proper-flat}
Soit $S$ un schéma. Soit $f : X \to Y$ un morphisme d'espaces algébriques
sur $S$, propre, plat et de présentation finie.
Soit $n_{X/Y}$ la fonction sur $Y$ donnant la dimension des fibres de $f$
introduite au Lemme \ref{spaces-more-morphisms-lemma-base-change-dimension-fibres}.
Alors $n_{X/Y}$ est localement constante.
\end{lemma}

\begin{proof}
Conséquence immédiate des
Lemmes \ref{spaces-more-morphisms-lemma-dimension-fibres-flat} et
\ref{spaces-more-morphisms-lemma-dimension-fibres-proper}.
\end{proof}





\section{Espaces algébriques caténaires}
\label{spaces-more-morphisms-section-catenary}

\noindent
Cette section poursuit la discussion commencée dans
Espaces décents, section \ref{decent-spaces-section-catenary}.
Le lemme qui suit sera utilisé dans la démonstration du suivant.

\begin{lemma}
\label{spaces-more-morphisms-lemma-construct-glueing}
Soit $S$ un schéma. Soit $f : X \to Y$ un morphisme entier
d'espaces algébriques sur $S$.
Soit $y \in |Y|$ un point représentable par une immersion fermée
$y : \Spec(k) \to Y$. Il existe alors
une factorisation $X \to X' \to Y$ de $f$ telle que
\begin{enumerate}
\item $X' \to Y$ soit entier ;
\item $X \to X'$ soit un isomorphisme au-dessus de $X' \setminus X'_y$ ;
\item $X'_y$ possède un unique point $x'$ tel que $\kappa(x') = k$.
\end{enumerate}
De plus, si $f$ est fini et $Y$ est localement noethérien, alors
$X' \to Y$ est fini.
\end{lemma}

\begin{proof}
D'après Morphismes d'espaces, Lemme \ref{spaces-morphisms-lemma-pushforward},
les faisceaux $f_*\mathcal{O}_X$, $(X_y \to Y)_*\mathcal{O}_{X_y}$ et
$y_*\mathcal{O}_{\Spec(k)}$ sont des faisceaux quasi-cohérents
d'algèbres sur $\mathcal{O}_Y$. Considérons les applications
$$
f_*\mathcal{O}_Y \longrightarrow
(X_y \to Y)_*\mathcal{O}_{X_y} \longleftarrow
y_*\mathcal{O}_{\Spec(k)}
$$
Le produit fibré est un faisceau quasi-cohérent d'algèbres sur $\mathcal{O}_Y$,
noté $\mathcal{A}'$, et l'on peut définir $X' \to Y$ comme le spectre relatif
de $\mathcal{A}'$ sur $Y$ ; voir
Morphismes, Lemme \ref{morphisms-lemma-affine-equivalence-algebras}.
Cette construction commute à tout changement de base.
En particulier, il est clair que, sur le sous-espace ouvert
$|Y| \setminus \{y\}$, le morphisme $X \to X'$ est un isomorphisme
et que, sur $|Y| \setminus \{y\}$, le morphisme $X' \to Y$ est entier.
Il reste à démontrer les assertions dans un petit voisinage de $y$.
Choisissons un schéma affine $V = \Spec(R)$ et un morphisme
étale $\varphi : V \to Y$ tel que $y$ appartienne à l'image de
$\varphi$. Alors $V_y$ est un sous-schéma fermé de $V$,
étale sur $k$, et se compose donc d'un nombre fini de points
dont chacun a un corps résiduel séparable sur $k$
(voir Espaces décents, Remarque \ref{decent-spaces-remark-recall}).
Après avoir rétréci $V$,
nous pouvons supposer qu'il existe un unique point fermé $v = \Spec(l) \to V$
s'envoyant sur $y$, avec $l/k$ fini séparable.
On peut écrire $V \times_Y X = \Spec(C)$, où $R \to C$
est un homomorphisme d'anneaux entier. La compatibilité énoncée avec le
changement de base donne $U \times_X Y' = \Spec(C')$, où
$$
C' = C \times_{C \otimes_R l} l
$$
Comme $R \to l$ est surjectif, $C \to C \otimes_R l$
l'est aussi, et nous voyons qu'il s'agit d'un produit fibré du
genre étudié dans Compléments sur l'algèbre, Situation
\ref{more-algebra-situation-module-over-fibre-product}
(où $A, A', B, B'$ correspondent à $C \otimes_R l, C, l, C'$).
Remarquons que $C'$ est une $R$-sous-algèbre de $C$ et est donc
entière sur $R$ ; cela démontre (1).
Enfin, Compléments sur l'algèbre, Lemme
\ref{more-algebra-lemma-points-of-fibre-product}
montre que $V \times_X Y' = \Spec(C')$
possède un unique point $y''$ au-dessus de $v$, de corps résiduel $l$
(cela correspond à l'application surjective évidente $C' \to l$).
Ainsi, $X_y \times_{\Spec(k)} \Spec(l)$ possède un unique point
de corps résiduel $l$. Puisque $l/k$ est fini séparable,
cela implique que $X'_y$ possède un unique point de corps résiduel $k$, c'est-à-dire
que (3) est vérifiée.

\medskip\noindent
Pour démontrer la dernière assertion, remarquons que si $Y$ est localement noethérien,
alors $R$ est un anneau noethérien et, si $f$ est fini, $R \to C$ est
fini. Alors $C'$ est une $R$-algèbre de type fini d'après Compléments sur l'algèbre,
Lemme \ref{more-algebra-lemma-fibre-product-finite-type}.
Cela démontre que $X' \to Y$ est fini.
\end{proof}

\begin{lemma}
\label{spaces-more-morphisms-lemma-universally-catenary-dimension-function}
Soit $S$ un schéma. Soit $B$ un espace algébrique sur $S$.
Soit $\delta : |B| \to \mathbf{Z}$ une fonction.
Supposons que $B$ soit décent, localement noethérien et
universellement caténaire, et que $\delta$ soit une fonction de dimension.
Si $X$ est un espace algébrique décent sur $B$ dont le morphisme structural
$f : X \to B$ est localement de type fini, nous définissons
$\delta_X : |X| \to \mathbf{Z}$ par la règle
$$
\delta_X(x) = \delta(f(x)) + \text{degré de transcendance de }x/f(x)
$$
(Morphismes d'espaces, Définition
\ref{spaces-morphisms-definition-dimension-fibre}).
Alors $\delta_X$ est une fonction de dimension.
\end{lemma}

\begin{proof}
La question est locale sur $B$. Nous pouvons donc supposer $B$ quasi-compact.
D'après Espaces décents, Lemme
\ref{decent-spaces-lemma-locally-Noetherian-decent-quasi-separated}
nous voyons que $B$ est quasi-séparé. D'après
Limites d'espaces, Proposition
\ref{spaces-limits-proposition-there-is-a-scheme-finite-over}
nous pouvons choisir un morphisme fini surjectif $\pi : Y \to X$,
où $Y$ est un schéma.
Affirmation : $\delta_Y$ est une fonction de dimension.

\medskip\noindent
L'affirmation implique le lemme. Pour $X \to B$ comme dans le lemme,
posons $Z = Y \times_B X$, de projections $p : Z \to Y$ et $q : Z \to X$.
Nous avons alors
$$
\delta_Z(z) = \delta_Y(p(z)) + \text{degré de transcendance de }z/p(z)
$$
et $\delta_Z(z) = \delta_X(q(z))$. Cela résulte de
Morphismes d'espaces, Lemme 
\ref{spaces-morphisms-lemma-dimension-fibre-at-a-point-additive}
et du fait que ces degrés de transcendance sont nuls
pour les morphismes finis. D'après Espaces décents, Lemme
\ref{decent-spaces-lemma-scheme-with-dimension-function}
et l'affirmation, nous trouvons que $\delta_Z$ est une fonction de dimension.
Nous en déduisons que $\delta_X$ est une fonction de dimension d'après
Espaces décents, Lemme
\ref{decent-spaces-lemma-check-dimension-function-finite-cover}.

\medskip\noindent
Démonstration de l'affirmation. Considérons une spécialisation $y \leadsto y'$,
$y \not = y'$, de points du schéma noethérien $Y$.
Alors $\delta_Y(y) > \delta_Y(y')$, car il n'existe
aucune spécialisation entre deux points d'une même fibre de $Y$
(voir Espaces décents, Lemme
\ref{decent-spaces-lemma-conditions-on-fibre-and-qf}).
En appliquant ceci à une chaîne de spécialisations, nous obtenons
$$
\delta_Y(y) - \delta_Y(y') \geq
\text{codim}(\overline{\{y'\}}, \overline{\{y\}})
$$
Il nous reste à démontrer l'égalité. D'après
Propriétés, Lemme \ref{properties-lemma-locally-Noetherian-closed-point},
nous pouvons choisir une spécialisation $y' \leadsto y_0$.
Il suffit de montrer que
$\delta_Y(y) - \delta_Y(y_0) =
\text{codim}(\overline{\{y_0\}}, \overline{\{y\}})$
car cela entraînera l'égalité à la fois pour
$y \leadsto y'$ et pour $y' \leadsto y_0$.

\medskip\noindent
Choisissons une chaîne maximale
$y = y_c \leadsto y_{c - 1} \leadsto \ldots \leadsto y_0$
de spécialisations dans $Y$.
Posons $b = \pi(y)$ et $b_0 = \pi(y_0)$.
Choisissons une chaîne maximale
$b = b_e \leadsto b_{e - 1} \leadsto \ldots \leadsto b_0$
de spécialisations dans $|B|$.
Nous devons montrer que $e = c$.
Puisque $\pi$ est fermé
(Morphismes d'espaces, Lemme \ref{spaces-morphisms-lemma-finite-proper}),
nous pouvons trouver une suite de spécialisations
$y = y'_e \leadsto y'_{e - 1} \leadsto \ldots \leadsto y'_0$
s'envoyant sur
$b = b_e \leadsto b_{e - 1} \leadsto \ldots \leadsto b_0$.
Remarquons que $y'_e \leadsto y'_{e - 1} \leadsto \ldots \leadsto y'_0$
est également une chaîne maximale.
Si $y_0 = y'_0$, alors, comme $Y$ est caténaire, nous concluons
que $e = c$, comme souhaité. Dans le paragraphe suivant, nous nous ramenons à
ce cas par une astuce et concluons de la même manière.

\medskip\noindent
Puisque $\pi$ est fermé, nous voyons que $b_0$ est un point fermé de $|B|$.
D'après Espaces décents, Lemme \ref{decent-spaces-lemma-decent-space-closed-point},
nous pouvons représenter $b_0$ par une immersion fermée $b_0 : \Spec(k) \to B$.
Par le Lemme \ref{spaces-more-morphisms-lemma-construct-glueing},
nous pouvons trouver une factorisation
$$
Y \to Y' \to X
$$
où $\pi' : Y' \to X$ est fini et $Y \to Y'$ est un morphisme qui
envoie $y_0$ et $y'_0$ sur le même point et est un isomorphisme
hors de l'image inverse de $b_0$.
(Bien entendu, $Y'$ ne sera pas un schéma, mais cela n'affecte pas
l'argument qui suit.)
Il est clair que les chaînes maximales de spécialisations
$y_c \leadsto y_{c - 1} \leadsto \ldots \leadsto y_0$ et
$y'_e \leadsto y'_{e - 1} \leadsto \ldots \leadsto y'_0$
s'envoient sur des chaînes maximales de spécialisations dans $Y'$
ayant le même début et la même fin.
Comme $B$ est universellement caténaire, nous voyons que
$|Y'|$ est caténaire et concluons comme précédemment.
\end{proof}






\section{Localisation étale des morphismes}
\label{spaces-more-morphisms-section-etale-localization}

\noindent
Cette section est l'analogue de
Compléments sur les morphismes, section \ref{more-morphisms-section-etale-localization}.

\begin{lemma}
\label{spaces-more-morphisms-lemma-etale-splits-off-quasi-finite-part}
Soit $S$ un schéma.
Soit $f : X \to Y$ un morphisme d'espaces algébriques sur $S$.
Soit $y \in |Y|$. Soient $x_1, \ldots, x_n \in |X|$ des points s'envoyant sur $y$.
Supposons que
\begin{enumerate}
\item $f$ soit localement de type fini ;
\item $f$ soit séparé ;
\item $f$ soit quasi-fini en $x_1, \ldots, x_n$ ;
\item $f$ soit quasi-compact ou $Y$ soit décent.
\end{enumerate}
Il existe alors un morphisme étale $(U, u) \to (Y, y)$
d'espaces algébriques pointés et une décomposition
$$
U \times_Y X = W \amalg V
$$
en sous-espaces ouverts et fermés telle que le morphisme $V \to U$ soit fini,
que tout point de la fibre de $|V| \to |U|$ au-dessus de $u$ s'envoie sur l'un des $x_i$,
et que la fibre de $|W| \to |U|$ au-dessus de $u$ ne contienne aucun point
s'envoyant sur l'un des $x_i$.
\end{lemma}

\begin{proof}
Soit $(U, u) \to (Y, y)$ un morphisme étale d'espaces algébriques,
et considérons l'ensemble des $w \in |U \times_Y X|$ s'envoyant sur $u \in |U|$
et sur l'un des $x_i \in |X|$. D'après
Espaces décents, Lemme \ref{decent-spaces-lemma-qf-and-qc-finite-fibre}
(si $f$ est de type fini) ou
Espaces décents, Lemme \ref{decent-spaces-lemma-decent-finite-fibre}
(si $Y$ est décent), cet ensemble est fini.
Nous pouvons donc remplacer $f$ par le changement de base
$U \times_Y X \to U$ et $x_1, \ldots, x_n$ par l'ensemble de ces $w$.
En particulier, nous pouvons supposer, et nous le faisons, que $Y$ est un schéma affine ;
alors $X$ est un espace algébrique séparé.

\medskip\noindent
Choisissons un schéma affine $Z$ et un morphisme étale $Z \to X$ tels que
$x_1, \ldots, x_n$ appartiennent à l'image de $|Z| \to |X|$. Les fibres de
$|Z| \to |X|$ sont finies ; voir Propriétés des espaces, Lemme
\ref{spaces-properties-lemma-finite-fibres-presentation}
(ou la discussion plus générale dans Espaces décents, section
\ref{decent-spaces-section-reasonable-decent}).
Soit $\{z_1, \ldots, z_m\} \subset |Z|$ l'image inverse de
$\{x_1, \ldots, x_n\}$. D'après Compléments sur les morphismes, Lemme
\ref{more-morphisms-lemma-etale-splits-off-quasi-finite-part-technical}
il existe un morphisme étale $(U, u) \to (Y, y)$ tel
que $U \times_Y Z = Z_1 \amalg Z_2$, avec $Z_1 \to U$ fini et
$(Z_1)_y = \{z_1, \ldots, z_m\}$. Nous pouvons supposer que $U$ est affine,
et donc que $Z_1$ est affine lui aussi.

\medskip\noindent
Puisque $f$ est séparé, l'image $V$ de $Z_1 \to X$ est à la fois ouverte et fermée
(Morphismes d'espaces, Lemme
\ref{spaces-morphisms-lemma-universally-closed-permanence}).
Posons $W = X \setminus V$ pour obtenir une décomposition comme dans le lemme.
Pour achever la démonstration, il reste à montrer que $V \to U$ est fini.
Comme $Z_1 \to V$ est étale surjectif, $V$ est le quotient de
$Z_1$ par la relation d'équivalence étale $R = Z_1 \times_V Z_1$ ; voir
Espaces, Lemme \ref{spaces-lemma-space-presentation}.
Puisque $f$ est séparé, $V \to U$ est séparé et $R$ est fermé dans
$Z_1 \times_U Z_1$. Puisque $Z_1 \to U$ est fini,
les projections $s, t : R \to Z_1$ sont finies.
Ainsi, $V$ est un schéma affine d'après
Groupoïdes, Proposition \ref{groupoids-proposition-finite-flat-equivalence}.
D'après Morphismes, Lemme \ref{morphisms-lemma-image-proper-is-proper},
nous concluons que $V \to U$ est propre et, d'après
Morphismes, Lemme \ref{morphisms-lemma-finite-proper},
nous concluons que $V \to U$ est fini,
ce qui achève la démonstration.
\end{proof}

\begin{lemma}
\label{spaces-more-morphisms-lemma-etale-splits-off-just-one-quasi-finite-part}
Soit $S$ un schéma.
Soit $f : X \to Y$ un morphisme d'espaces algébriques sur $S$.
Soit $x \in |X|$ d'image $y \in |Y|$. Supposons que
\begin{enumerate}
\item $f$ soit localement de type fini ;
\item $f$ soit séparé ;
\item $f$ soit quasi-fini en $x$.
\end{enumerate}
Il existe alors un morphisme étale $(U, u) \to (Y, y)$
d'espaces algébriques pointés et une décomposition
$$
U \times_Y X = W \amalg V
$$
en sous-espaces ouverts et fermés telle que le morphisme $V \to U$ soit fini
et qu'il existe un point $v \in |V|$ qui s'envoie sur $x$ dans $|X|$
et sur $u$ dans $|U|$.
\end{lemma}

\begin{proof}
Choisissons un schéma $U$, un point $u \in U$ et un morphisme étale
$U \to Y$ qui envoie $u$ sur $y$. Il existe un point $x' \in |U \times_Y X|$
qui s'envoie sur $x$ dans $|X|$ et sur $u$ dans $|U|$ (Propriétés des espaces,
Lemme \ref{spaces-properties-lemma-points-cartesian}).
Pour terminer, appliquons le Lemme \ref{spaces-more-morphisms-lemma-etale-splits-off-quasi-finite-part}
au morphisme $U \times_Y X \to U$ et au point $x'$.
Il s'applique parce que $U$ est un schéma et que $u$ est donc représenté
par le monomorphisme $\Spec(\kappa(u)) \to U$.
\end{proof}










\section{Théorème principal de Zariski}
\label{spaces-more-morphisms-section-structure-quasi-finite}

\noindent
Dans cette section, nous appliquons les résultats de la section précédente pour
démontrer le théorème principal de Zariski pour les morphismes d'espaces algébriques.
Cette section est l'analogue de la section
\ref{more-morphisms-section-application-etale-neighbourhoods} de Compléments sur les morphismes.

\begin{lemma}
\label{spaces-more-morphisms-lemma-finite-type-separated}
Soit $S$ un schéma. Soit $f : X \to Y$ un morphisme d'espaces
algébriques sur $S$, de type fini et séparé.
Soit $Y'$ la normalisation de $Y$ dans $X$. Diagramme :
$$
\xymatrix{
X \ar[rd]_f \ar[rr]_{f'} & & Y' \ar[ld]^\nu \\
& Y &
}
$$
Il existe alors un sous-espace ouvert $U' \subset Y'$ tel que
\begin{enumerate}
\item $(f')^{-1}(U') \to U'$ soit un isomorphisme ;
\item $(f')^{-1}(U') \subset X$ soit l'ensemble des points où
$f$ est quasi-fini.
\end{enumerate}
\end{lemma}

\begin{proof}
D'après Morphismes d'espaces, Lemme
\ref{spaces-morphisms-lemma-locally-finite-type-quasi-finite-part}
il existe un sous-espace ouvert $U \subset X$ correspondant aux points
de $|X|$ où $f$ est quasi-fini. Nous devons démontrer que
\begin{enumerate}
\item[(a)] l'image de $|U| \to |Y'|$ est $|U'|$ pour un certain sous-espace ouvert
$U'$ de $Y'$ ;
\item[(b)] $U = f^{-1}(U')$ ;
\item[(c)] $U \to U'$ est un isomorphisme.
\end{enumerate}
Puisque la formation de $U$ commute à tout changement de base
(Morphismes d'espaces, Lemme
\ref{spaces-morphisms-lemma-locally-finite-type-quasi-finite-part}),
puisque la formation de la normalisation $Y'$ commute au changement de base
lisse (Lemme \ref{spaces-more-morphisms-lemma-normalization-smooth-localization}),
puisque les morphismes étales sont ouverts, et
puisque la propriété « être un isomorphisme » est locale pour la topologie fpqc sur la base
(Descente sur les espaces, Lemme
\ref{spaces-descent-lemma-descending-property-isomorphism}),
il suffit de démontrer (a), (b), (c) localement pour la topologie étale sur $Y$
(nous omettons certains détails). Ainsi,
nous pouvons supposer que $Y$ est un schéma affine. Cela implique que $Y'$ est
lui aussi un schéma (affine).

\medskip\noindent
Soit $x \in |U|$. Affirmation : il existe un voisinage
ouvert $f'(x) \in V \subset Y'$ tel que $(f')^{-1}V \to V$ soit un
isomorphisme. Montrons d'abord que l'affirmation implique le lemme.
En effet, alors $(f')^{-1}V \cong V$ est un schéma (comme ouvert
de $Y'$), localement de type fini sur $Y$ (comme sous-espace ouvert de $X$),
et, pour $v \in V$, l'extension de corps résiduels
$\kappa(v)/\kappa(\nu(v))$ est algébrique (car
$V \subset Y'$ et $Y'$ est entier sur $Y$). Par conséquent, les fibres
de $V \to Y$ sont discrètes (Morphismes, Lemme
\ref{morphisms-lemma-algebraic-residue-field-extension-closed-point-fibre})
et $(f')^{-1}V \to Y$ est localement quasi-fini
(Morphismes, Lemme \ref{morphisms-lemma-locally-quasi-finite-fibres}).
Cela implique $(f')^{-1}V \subset U$ et $V \subset U'$. Comme $x$ était
arbitraire, nous voyons que (a), (b) et (c) sont vraies.

\medskip\noindent
Soit $y = f(x) \in |Y|$. Soit $(T, t) \to (Y, y)$ un morphisme étale
de schémas pointés. Notons par un indice inférieur ${}_T$ le changement de base à $T$.
Soit $z \in X_T$ un point de la fibre $X_t$ situé au-dessus de $x$.
Remarquons que $U_T \subset X_T$ est l'ensemble des points où $f_T$ est
quasi-fini ; voir Morphismes d'espaces, Lemme
\ref{spaces-morphisms-lemma-locally-finite-type-quasi-finite-part}.
Remarquons que
$$
X_T \xrightarrow{f'_T} Y'_T \xrightarrow{\nu_T} T
$$
est la normalisation de $T$ dans $X_T$ ; voir
Lemme \ref{spaces-more-morphisms-lemma-normalization-smooth-localization}.
Supposons l'affirmation vraie pour $z \in U_T \subset X_T \to Y'_T \to T$, c'est-à-dire
supposons que nous puissions trouver un voisinage ouvert
$f'_T(z) \in V' \subset Y'_T$ tel que $(f'_T)^{-1}V' \to V'$ soit un
isomorphisme. Le morphisme $Y'_T \to Y'$ est étale ; par conséquent, l'image
$V \subset Y'$ de $V'$ est ouverte. Remarquons que $f'(x) \in V$ puisque $f'_T(z) \in V'$.
Remarquons que
$$
\xymatrix{
(f'_T)^{-1}V' \ar[r] \ar[d] & (f')^{-1}(V) \ar[d] \\
V' \ar[r] & V
}
$$
est un carré cartésien (car $Y'_T \times_{Y'} X = X_T$).
Puisque la flèche verticale gauche est un isomorphisme
et que $\{V' \to V\}$ est un recouvrement étale, nous concluons que la flèche verticale
droite est un isomorphisme d'après
Descente sur les espaces, Lemme
\ref{spaces-descent-lemma-descending-property-isomorphism}.
Autrement dit, l'affirmation est vraie pour $x \in U \subset X \to Y' \to Y$.

\medskip\noindent
D'après le paragraphe précédent, pour démontrer l'affirmation pour
$x \in |U|$, nous pouvons remplacer $Y$ par un voisinage étale $T$ de
$y = f(x)$ et $x$ par n'importe quel point situé au-dessus de $x$ dans $T \times_Y X$.
Ainsi, nous pouvons supposer qu'il existe une décomposition
$$
X = V \amalg W
$$
en sous-espaces ouverts et fermés, où $V \to Y$ est fini et $x \in V$ ;
voir le Lemme \ref{spaces-more-morphisms-lemma-etale-splits-off-quasi-finite-part}.
Puisque $X$ est la somme disjointe de $V$ et $W$ sur $Y$ et que
$V \to Y$ est fini, nous voyons que la
normalisation de $Y$ dans $X$ est le morphisme
$$
X = V \amalg W \longrightarrow V \amalg W' \longrightarrow S
$$
où $W'$ est la normalisation de $Y$ dans $W$ ; voir
Morphismes d'espaces, Lemmes
\ref{spaces-morphisms-lemma-normalization-in-disjoint-union},
\ref{spaces-morphisms-lemma-finite-integral}, et
\ref{spaces-morphisms-lemma-normalization-in-integral}.
L'affirmation en résulte, et la démonstration est terminée.
\end{proof}

\noindent
Le lemme suivant est un doublon de
Morphismes d'espaces, Lemme
\ref{spaces-morphisms-lemma-quasi-finite-separated-quasi-affine}.
La présence de deux exemplaires du même lemme s'explique par le fait que les
démonstrations diffèrent quelque peu. La démonstration donnée ci-dessous
repose sur le théorème principal de Zariski pour les morphismes non représentables
d'espaces algébriques sous la forme établie ci-dessus, tandis que la démonstration de
Morphismes d'espaces, Lemme
\ref{spaces-morphisms-lemma-quasi-finite-separated-quasi-affine}
repose sur Morphismes d'espaces, Proposition
\ref{spaces-morphisms-proposition-locally-quasi-finite-separated-over-scheme}
pour se ramener au cas des morphismes de schémas.

\begin{lemma}
\label{spaces-more-morphisms-lemma-quasi-finite-separated-quasi-affine}
Soit $S$ un schéma.
Soit $f : X \to Y$ un morphisme d'espaces algébriques sur $S$.
Supposons que $f$ soit quasi-fini et séparé.
Soit $Y'$ la normalisation de $Y$ dans $X$.
Diagramme :
$$
\xymatrix{
X \ar[rd]_f \ar[rr]_{f'} & & Y' \ar[ld]^\nu \\
& Y &
}
$$
Alors $f'$ est une immersion ouverte quasi-compacte et $\nu$ est entier.
En particulier, $f$ est quasi-affine.
\end{lemma}

\begin{proof}
Cela résulte du Lemme \ref{spaces-more-morphisms-lemma-finite-type-separated}. En effet, d'après
ce lemme, il existe un sous-espace ouvert $U' \subset Y'$ tel que
$(f')^{-1}(U') = X$ (!) et que $X \to U'$ soit un isomorphisme ! Autrement
dit, $f'$ est une immersion ouverte. Remarquons que $f'$ est quasi-compact puisque
$f$ est quasi-compact et que $\nu : Y' \to Y$ est séparé
(Morphismes d'espaces, Lemme
\ref{spaces-morphisms-lemma-quasi-compact-permanence}).
Ainsi, pour tout schéma affine $Z$ et tout morphisme $Z \to Y$, le
produit fibré $Z \times_Y X$ est un sous-schéma ouvert quasi-compact
du schéma affine $Z \times_Y Y'$. Par conséquent, $f$ est quasi-affine par
définition.
\end{proof}

\begin{lemma}[Théorème principal de Zariski]
\label{spaces-more-morphisms-lemma-quasi-finite-separated-pass-through-finite}
Soit $S$ un schéma.
Soit $f : X \to Y$ un morphisme d'espaces algébriques sur $S$.
Supposons que $f$ soit quasi-fini et séparé, et que
$Y$ soit quasi-compact et quasi-séparé. Il existe alors
une factorisation
$$
\xymatrix{
X \ar[rd]_f \ar[rr]_j & & T \ar[ld]^\pi \\
& Y &
}
$$
où $j$ est une immersion ouverte quasi-compacte et $\pi$ est fini.
\end{lemma}

\begin{proof}
Soit $X \to Y' \to Y$ comme dans la conclusion du
Lemme \ref{spaces-more-morphisms-lemma-quasi-finite-separated-quasi-affine}.
D'après
Limites d'espaces, Lemme
\ref{spaces-limits-lemma-integral-algebra-directed-colimit-finite}
nous pouvons écrire
$\nu_*\mathcal{O}_{Y'} = \colim_{i \in I} \mathcal{A}_i$ comme une
limite inductive filtrante d'algèbres finies quasi-cohérentes sur $\mathcal{O}_X$,
$\mathcal{A}_i \subset \nu_*\mathcal{O}_{Y'}$. Alors
$\pi_i : T_i = \underline{\Spec}_Y(\mathcal{A}_i) \to Y$
est un morphisme fini pour tout $i$.
Remarquons que les morphismes de transition $T_{i'} \to T_i$ sont affines
et que $Y' = \lim T_i$.

\medskip\noindent
D'après Limites d'espaces, Lemme \ref{spaces-limits-lemma-descend-opens},
il existe un $i$ et un ouvert quasi-compact
$U_i \subset T_i$ dont l'image inverse dans $Y'$ est
$f'(X)$. Pour $i' \geq i$, notons $U_{i'}$ l'image inverse
de $U_i$ dans $T_{i'}$. Alors $X \cong f'(X) = \lim_{i' \geq i} U_{i'}$ ; voir
Limites d'espaces, Lemme
\ref{spaces-limits-lemma-directed-inverse-system-has-limit}.
D'après
Limites d'espaces, Lemme
\ref{spaces-limits-lemma-finite-type-eventually-closed}, nous voyons que
$X \to U_{i'}$ est une immersion fermée pour un certain $i' \geq i$.
(En fait, $X \cong U_{i'}$ pour $i'$ suffisamment
grand, mais nous n'en avons pas besoin.) Ainsi, $X \to T_{i'}$ est une immersion. D'après
Morphismes d'espaces, Lemme
\ref{spaces-morphisms-lemma-factor-the-other-way}
nous pouvons factoriser ce morphisme en $X \to T \to T_{i'}$, où la première flèche
est une immersion ouverte et la seconde une immersion fermée. Cela termine la démonstration.
\end{proof}

\begin{lemma}
\label{spaces-more-morphisms-lemma-quasi-finite-separated-pass-through-finite-addendum}
Avec les notations et les hypothèses du
Lemme \ref{spaces-more-morphisms-lemma-quasi-finite-separated-pass-through-finite},
supposons de plus que $f$ soit localement de présentation finie. Nous pouvons alors
choisir la factorisation de telle sorte que $T$ soit fini et de
présentation finie sur $Y$.
\end{lemma}

\begin{proof}
D'après Limites d'espaces, Lemme
\ref{spaces-limits-lemma-finite-in-finite-and-finite-presentation}, nous pouvons écrire
$T = \lim T_i$, où tous les $T_i$ sont finis et de présentation finie
sur $Y$ et où les morphismes de transition $T_{i'} \to T_i$ sont des
immersions fermées. D'après
Limites d'espaces, Lemme \ref{spaces-limits-lemma-descend-opens},
il existe un $i$ et un sous-schéma ouvert $U_i \subset T_i$ dont l'image
inverse dans $T$ est $X$. D'après
Limites d'espaces, Lemme
\ref{spaces-limits-lemma-finite-type-eventually-closed}
nous voyons que $X \cong U_i$ pour $i$ assez grand.
Remplacer $T$ par $T_i$ achève la démonstration.
\end{proof}












\section{Applications du théorème principal de Zariski, I}
\label{spaces-more-morphisms-section-applications-zmt}

\noindent
Une première application est la caractérisation des morphismes
finis comme morphismes propres à fibres finies.

\begin{lemma}
\label{spaces-more-morphisms-lemma-characterize-finite}
Soit $S$ un schéma. Soit $f : X \to Y$ un morphisme d'espaces algébriques
sur $S$. Les propriétés suivantes sont équivalentes :
\begin{enumerate}
\item $f$ est fini ;
\item $f$ est propre et localement quasi-fini ;
\item $f$ est propre et $|X_k|$ est un espace discret pour tout morphisme
$\Spec(k) \to Y$, où $k$ est un corps ;
\item $f$ est universellement fermé, séparé, localement de type fini,
et $|X_k|$ est un espace discret pour tout morphisme $\Spec(k) \to Y$,
où $k$ est un corps.
\end{enumerate}
\end{lemma}

\begin{proof}
Nous avons (1) $\Rightarrow$ (2) d'après
Morphismes d'espaces, Lemmes \ref{spaces-morphisms-lemma-finite-proper},
\ref{spaces-morphisms-lemma-finite-quasi-finite}.
Nous avons (2) $\Rightarrow$ (3) d'après
Morphismes d'espaces, Lemme
\ref{spaces-morphisms-lemma-locally-quasi-finite}.
Par définition, (3) implique (4).

\medskip\noindent
Supposons (4). Puisque $f$ est universellement fermé, il est quasi-compact
(Morphismes d'espaces, Lemme
\ref{spaces-morphisms-lemma-universally-closed-quasi-compact}).
Choisissons un point $y$ de $|Y|$. Représentons $y$ par un
morphisme $\Spec(k) \to Y$. Remarquons que $|X_k|$ est fini et discret,
car il est quasi-compact et discret. L'application $|X_k| \to |X|$ est surjective
sur la fibre de $|X| \to |Y|$ au-dessus de $y$
(Propriétés des espaces, Lemme \ref{spaces-properties-lemma-points-cartesian}).
D'après
Morphismes d'espaces, Lemme \ref{spaces-morphisms-lemma-quasi-finite-at-point},
nous voyons que $X \to Y$ est quasi-fini en tous les points de la fibre
de $|X| \to |Y|$ au-dessus de $y$.
Choisissons un voisinage étale élémentaire $(U, u) \to (Y, y)$
et une décomposition $X_U = V \amalg W$ comme dans le
Lemme \ref{spaces-more-morphisms-lemma-etale-splits-off-quasi-finite-part},
adaptée à tous les points de $|X|$ situés au-dessus de $y$.
Remarquons que $W_u = \emptyset$, puisque nous avons utilisé tous les points
de la fibre de $|X| \to |Y|$ au-dessus de $y$.
Puisque $f$ est universellement fermé, nous voyons que
l'image de $|W|$ dans $|U|$ est un fermé qui ne contient pas $u$.
Après avoir rétréci $U$, nous pouvons supposer que $W = \emptyset$.
Autrement dit, nous voyons que $X_U = V$ est fini sur $U$.
Comme $y \in |Y|$ était arbitraire,
il existe donc une famille $\{U_i \to Y\}$
de morphismes étales dont les images recouvrent $Y$ et telle que
les changements de base $X_{U_i} \to U_i$ soient finis.
Nous concluons que $f$ est fini d'après
Morphismes d'espaces, Lemme \ref{spaces-morphisms-lemma-integral-local}.
\end{proof}

\noindent
Nous en déduisons le résultat utile suivant.

\begin{lemma}
\label{spaces-more-morphisms-lemma-proper-finite-fibre-finite-in-neighbourhood}
Soit $S$ un schéma. Soit $f : X \to Y$ un morphisme d'espaces algébriques
sur $S$. Soit $y \in |Y|$. Supposons que
\begin{enumerate}
\item $f$ soit propre ;
\item $f$ soit quasi-fini en tout $x \in |X|$ situé au-dessus de $y$
(Espaces décents, Lemme \ref{decent-spaces-lemma-conditions-on-fibre-and-qf}).
\end{enumerate}
Il existe alors un voisinage ouvert $V \subset Y$ de $y$
tel que $f|_{f^{-1}(V)} : f^{-1}(V) \to V$ soit fini.
\end{lemma}

\begin{proof}
D'après
Morphismes d'espaces, Lemme
\ref{spaces-morphisms-lemma-locally-finite-type-quasi-finite-part},
l'ensemble des points où $f$ est quasi-fini est un ouvert $U \subset X$.
Soit $Z = X \setminus U$. Alors $y \not \in f(Z)$. Puisque $f$ est propre,
l'ensemble $f(Z) \subset Y$ est fermé. Choisissons un voisinage ouvert quelconque
$V \subset Y$ de $y$ tel que $Z \cap V = \emptyset$. Alors
$f^{-1}(V) \to V$ est localement quasi-fini et propre.
Par conséquent, $f^{-1}(V) \to V$ est fini d'après le Lemme \ref{spaces-more-morphisms-lemma-characterize-finite}.
\end{proof}

\begin{lemma}
\label{spaces-more-morphisms-lemma-flat-proper-family-cannot-collapse-fibre}
\begin{slogan}
La contraction d'une fibre d'une famille propre force également la contraction des fibres voisines.
\end{slogan}
Soit $S$ un schéma. Soit
$$
\xymatrix{
X \ar[rr]_h \ar[rd]_f & & Y \ar[ld]^g \\
& B
}
$$
un diagramme commutatif de morphismes d'espaces algébriques sur $S$.
Soit $b \in B$ et soit $\Spec(k) \to B$ un morphisme appartenant à la classe
d'équivalence de $b$. Supposons que
\begin{enumerate}
\item $X \to B$ soit un morphisme propre ;
\item $Y \to B$ soit séparé et localement de type fini ;
\item l'une des propriétés suivantes soit vraie :
\begin{enumerate}
\item l'image de $|X_k| \to |Y_k|$ est finie ;
\item l'image de $|f|^{-1}(\{b\})$ dans $|Y|$ est finie
et $B$ est décent.
\end{enumerate}
\end{enumerate}
Il existe alors un
sous-espace ouvert $B' \subset B$ contenant $b$ tel que $X_{B'} \to Y_{B'}$
se factorise par un sous-espace fermé $Z \subset Y_{B'}$ fini sur $B'$.
\end{lemma}

\begin{proof}
Soit $Z \subset Y$ l'image schématique de $h$ ; voir
Morphismes d'espaces, section
\ref{spaces-morphisms-section-scheme-theoretic-image}.
D'après Morphismes d'espaces, Lemme
\ref{spaces-morphisms-lemma-scheme-theoretic-image-is-proper}
le morphisme $X \to Z$ est surjectif et $Z \to B$ est propre.
Ainsi,
$$
\{x \in |X|\text{ situé au-dessus de }b\} \to
\{z \in |Z|\text{ situé au-dessus de }b\}
$$
et $|X_k| \to |Z_k|$ sont surjectifs. Nous voyons que l'une ou l'autre
des conditions (3)(a) et (3)(b) implique que $Z \to B$ est quasi-fini
en tous les points de $|Z|$ situés au-dessus de $b$, d'après
Espaces décents, Lemme \ref{decent-spaces-lemma-conditions-on-fibre-and-qf}.
Ainsi, $Z \to B$ est fini dans un voisinage ouvert de $b$ d'après le
Lemme \ref{spaces-more-morphisms-lemma-proper-finite-fibre-finite-in-neighbourhood}.
\end{proof}







\section{Factorisation de Stein}
\label{spaces-more-morphisms-section-stein-factorization}

\noindent
La factorisation de Stein affirme qu'un morphisme propre $f : X \to S$
tel que $f_*\mathcal{O}_X = \mathcal{O}_S$ a des fibres connexes.

\begin{lemma}
\label{spaces-more-morphisms-lemma-stein-universally-closed}
Soit $S$ un schéma. Soit $f : X \to Y$ un morphisme universellement fermé et
quasi-séparé d'espaces algébriques sur $S$.
Il existe une factorisation
$$
\xymatrix{
X \ar[rr]_{f'} \ar[rd]_f & & Y' \ar[dl]^\pi \\
& Y &
}
$$
ayant les propriétés suivantes :
\begin{enumerate}
\item le morphisme $f'$ est universellement fermé, quasi-compact, quasi-séparé
et surjectif ;
\item le morphisme $\pi : Y' \to Y$ est entier ;
\item nous avons $f'_*\mathcal{O}_X = \mathcal{O}_{Y'}$ ;
\item nous avons $Y' = \underline{\Spec}_Y(f_*\mathcal{O}_X)$ ;
\item $Y'$ est la normalisation de $Y$ dans $X$ au sens de
Morphismes d'espaces, Définition
\ref{spaces-morphisms-definition-normalization-X-in-Y}.
\end{enumerate}
La formation de la factorisation $f = \pi \circ f'$ commute au changement
de base plat.
\end{lemma}

\begin{proof}
D'après Morphismes d'espaces, Lemme
\ref{spaces-morphisms-lemma-universally-closed-quasi-compact}
le morphisme $f$ est quasi-compact.
Nous définissons simplement $Y'$ comme la normalisation de $Y$ dans $X$ ; (5) et (2) sont donc
automatiquement vérifiés. D'après
Morphismes d'espaces, Lemme
\ref{spaces-morphisms-lemma-normalization-in-universally-closed}
nous voyons que (4) est vérifié. Le morphisme $f'$ est universellement fermé d'après
Morphismes d'espaces, Lemme
\ref{spaces-morphisms-lemma-universally-closed-permanence}.
Il est quasi-compact d'après
Morphismes d'espaces, Lemme
\ref{spaces-morphisms-lemma-quasi-compact-permanence}
et quasi-séparé d'après
Morphismes d'espaces, Lemme
\ref{spaces-morphisms-lemma-compose-after-separated}.

\medskip\noindent
Pour démontrer les assertions restantes, nous pouvons supposer que la base $Y$ est affine
(car la normalisation commute à la localisation étale).
Écrivons $Y = \Spec(R)$. Alors $Y' = \Spec(A)$, où
$A = \Gamma(X, \mathcal{O}_X)$ est une $R$-algèbre entière.
Il est donc clair que $f'_*\mathcal{O}_X$
est $\mathcal{O}_{Y'}$ (car $f'_*\mathcal{O}_X$ est quasi-cohérent
d'après Morphismes d'espaces, Lemme
\ref{spaces-morphisms-lemma-pushforward},
et est donc égal à $\widetilde{A}$). Cela démontre (3).

\medskip\noindent
Montrons que $f'$ est surjectif. Comme $f'$ est universellement fermé (voir ci-dessus),
l'image de $f'$ est un fermé
$V(I) \subset Y' = \Spec(A)$. Choisissons $h \in I$. Alors
$h|_X = f^\sharp(h)$ est une section globale du faisceau structural de
$X$ qui s'annule en tout point. Puisque $X$ est quasi-compact, cela signifie
que $h|_X$ est une section nilpotente, c'est-à-dire que $h^n|X = 0$ pour un certain $n > 0$.
Mais $A = \Gamma(X, \mathcal{O}_X)$, donc $h^n = 0$.
Autrement dit, $I$ est contenu dans le radical de Jacobson de $A$, et nous concluons
que $V(I) = Y'$, comme souhaité.
\end{proof}

\begin{lemma}
\label{spaces-more-morphisms-lemma-stein-universally-closed-residue-fields}
Dans le Lemme \ref{spaces-more-morphisms-lemma-stein-universally-closed}, supposons de plus que
$f$ soit localement de type fini et que $Y$ soit affine. Alors, pour $y \in Y$, la fibre
$\pi^{-1}(\{y\}) = \{y_1, \ldots, y_n\}$ est finie et les extensions de corps
$\kappa(y_i)/\kappa(y)$ sont finies.
\end{lemma}

\begin{proof}
Rappelons qu'il n'existe pas de spécialisations entre les points de $\pi^{-1}(\{y\})$ ;
voir Algèbre, Lemme \ref{algebra-lemma-integral-no-inclusion}.
Puisque $f'$ est surjectif, nous obtenons que $|X_y| \to \pi^{-1}(\{y\})$ est surjectif.
Remarquons que $X_y$ est un espace algébrique quasi-séparé de type fini
sur un corps (la quasi-compacité a été démontrée dans la preuve du
lemme cité). Ainsi, $|X_y|$ est un espace topologique noethérien
(Morphismes d'espaces, Lemme
\ref{spaces-morphisms-lemma-finite-presentation-noetherian}).
Un argument topologique (omis) montre alors que $\pi^{-1}(\{y\})$ est fini.
Pour chaque $i$, nous pouvons choisir un point de type fini $x_i \in |X_y|$
qui s'envoie sur $y_i$ (Morphismes d'espaces, Lemme
\ref{spaces-morphisms-lemma-enough-finite-type-points}).
Nous concluons que $\kappa(y_i)/\kappa(y)$ est finie :
$x_i$ peut être représenté par un morphisme $\Spec(k_i) \to X_y$
de type fini (d'après notre définition des points de type fini) ;
par conséquent, $\Spec(k_i) \to y = \Spec(\kappa(y))$ est de type fini
(comme composée de morphismes de type fini),
donc $k_i/\kappa(y)$ est finie (Morphismes, Lemme
\ref{morphisms-lemma-point-finite-type}).
\end{proof}

\noindent
Soit $f : X \to Y$ un morphisme d'espaces algébriques et soit
$\overline{y} : \Spec(k) \to Y$ un point géométrique. La
{\it fibre de $f$ au-dessus de $\overline{y}$} est alors l'espace algébrique
$X_{\overline{y}} = X \times_{Y, \overline{y}} \Spec(k)$ sur $k$.
Si $Y$ est un schéma et si $y \in Y$ est un point, nous notons comme d'habitude
$X_y = X \times_Y \Spec(\kappa(y))$ la fibre.

\begin{lemma}
\label{spaces-more-morphisms-lemma-characterize-geometrically-connected-fibres}
Soit $S$ un schéma. Soit $f : X \to Y$ un morphisme d'espaces algébriques
sur $S$. Soit $\overline{y}$ un point géométrique de $Y$. Alors
$X_{\overline{y}}$ est connexe si et seulement si, pour tout voisinage
étale $(V, \overline{v}) \to (Y, \overline{y})$ tel que $V$
soit un schéma, le changement de base $X_V \to V$ a une fibre connexe $X_v$.
\end{lemma}

\begin{proof}
Puisque la catégorie des voisinages étales de $\overline{y}$ est
cofiltrante et contient une famille cofinale de schémas
(Propriétés des espaces, Lemme \ref{spaces-properties-lemma-cofinal-etale}),
nous pouvons remplacer $Y$ par l'un de ces voisinages et supposer que
$Y$ est un schéma. Soit $y \in Y$ le point correspondant à $\overline{y}$.
Alors $X_y$ est géométriquement connexe sur $\kappa(y)$ si et seulement si
$X_{\overline{y}}$ est connexe, et si et seulement si $(X_y)_{k'}$
est connexe pour toute extension séparable finie $k'$ de $\kappa(y)$.
Voir Espaces sur les corps, section
\ref{spaces-over-fields-section-geometrically-connected}, et en particulier le
Lemme \ref{spaces-over-fields-lemma-characterize-geometrically-disconnected}.
D'après Compléments sur les morphismes, Lemme
\ref{more-morphisms-lemma-realize-prescribed-residue-field-extension-etale}
il existe un voisinage étale affine $(V, v) \to (Y, y)$ tel que
$\kappa(s) \subset \kappa(u)$ s'identifie à $\kappa(s) \subset k'$
pour toute extension séparable finie donnée. Le lemme en résulte.
\end{proof}

\begin{theorem}[Factorisation de Stein ; cas noethérien]
\label{spaces-more-morphisms-theorem-stein-factorization-Noetherian}
Soit $S$ un schéma. Soit $f : X \to Y$ un morphisme propre d'espaces
algébriques sur $S$, avec $Y$ localement noethérien.
Il existe une factorisation
$$
\xymatrix{
X \ar[rr]_{f'} \ar[rd]_f & & Y' \ar[dl]^\pi \\
& Y &
}
$$
ayant les propriétés suivantes :
\begin{enumerate}
\item le morphisme $f'$ est propre à fibres géométriques connexes ;
\item le morphisme $\pi : Y' \to Y$ est fini ;
\item nous avons $f'_*\mathcal{O}_X = \mathcal{O}_{Y'}$ ;
\item nous avons $Y' = \underline{\Spec}_Y(f_*\mathcal{O}_X)$ ;
\item $Y'$ est la normalisation de $Y$ dans $X$ ; voir
Morphismes, Définition \ref{morphisms-definition-normalization-X-in-Y}.
\end{enumerate}
\end{theorem}

\begin{proof}
Soit $f = \pi \circ f'$ la factorisation du
Lemme \ref{spaces-more-morphisms-lemma-stein-universally-closed}. Remarquons qu'outre les
conclusions du Lemme \ref{spaces-more-morphisms-lemma-stein-universally-closed}, nous
savons aussi que $f'$ est séparé
(Morphismes d'espaces, Lemme
\ref{spaces-morphisms-lemma-compose-after-separated})
et de type fini
(Morphismes d'espaces, Lemme
\ref{spaces-morphisms-lemma-permanence-finite-type}).
Ainsi, $f'$ est propre. D'après
Cohomologie des espaces, Lemme
\ref{spaces-cohomology-lemma-proper-pushforward-coherent}
nous voyons que $f_*\mathcal{O}_X$ est un $\mathcal{O}_Y$-module cohérent.
Nous voyons donc que $\pi$ est fini, c'est-à-dire que (2) est vérifié.

\medskip\noindent
Cela démontre toutes les assertions sauf la plus intéressante, à savoir que
les fibres géométriques de $f'$ sont connexes. Il ressort clairement de la
discussion précédente que nous pouvons remplacer $Y$ par $Y'$.
Alors $Y$ est localement noethérien,
$f : X \to Y$ est propre et $f_*\mathcal{O}_X = \mathcal{O}_Y$.
Soit $\overline{y}$ un point géométrique de $Y$.
Nous appliquons ici le théorème des fonctions formelles,
plus précisément Cohomologie des espaces, Lemme
\ref{spaces-cohomology-lemma-formal-functions-stalk}.
Il nous dit que
$$
\mathcal{O}^\wedge_{Y, \overline{y}} =
\lim_n H^0(X_n, \mathcal{O}_{X_n})
$$
où $X_n =
\Spec(\mathcal{O}_{Y, \overline{y}}/\mathfrak m_{\overline{y}}^n) \times_Y X$.
Remarquons que $X_1 = X_{\overline{y}} \to X_n$ est un épaississement (d'ordre fini)
et que l'espace topologique sous-jacent à $X_n$ est donc égal à celui
de $X_{\overline{y}}$. Ainsi, si $X_{\overline{y}} = T_1 \amalg T_2$
est une somme disjointe de sous-espaces ouverts et fermés non vides, alors, de même,
$X_n = T_{1, n} \amalg T_{2, n}$ pour tout $n$. Cela signifie à son tour que
$H^0(X_n, \mathcal{O}_{X_n})$ contient un idempotent $e_{1, n}$ distinct de 0 et de 1,
à savoir la fonction identiquement égale à $1$ sur $T_{1, n}$ et
identiquement égale à $0$ sur $T_{2, n}$. Il est clair que $e_{1, n + 1}$
se restreint à $e_{1, n}$ sur $X_n$. Par conséquent, $e_1 = \lim e_{1, n}$
est un idempotent de la limite distinct de 0 et de 1. Cela contredit le fait
que $\mathcal{O}^\wedge_{Y, \overline{y}}$ est un anneau local. Notre
hypothèse était donc fausse, c'est-à-dire que $X_{\overline{y}}$ est connexe,
comme souhaité.
\end{proof}

\begin{theorem}[Factorisation de Stein ; cas général]
\label{spaces-more-morphisms-theorem-stein-factorization-general}
Soit $S$ un schéma. Soit $f : X \to Y$ un morphisme propre d'espaces
algébriques sur $S$. Il existe une factorisation
$$
\xymatrix{
X \ar[rr]_{f'} \ar[rd]_f & & Y' \ar[dl]^\pi \\
& Y &
}
$$
ayant les propriétés suivantes :
\begin{enumerate}
\item le morphisme $f'$ est propre à fibres géométriques connexes ;
\item le morphisme $\pi : Y' \to Y$ est entier ;
\item nous avons $f'_*\mathcal{O}_X = \mathcal{O}_{Y'}$ ;
\item nous avons $Y' = \underline{\Spec}_Y(f_*\mathcal{O}_X)$ ;
\item $Y'$ est la normalisation de $Y$ dans $X$ (Morphismes d'espaces, Définition
\ref{spaces-morphisms-definition-normalization-X-in-Y}).
\end{enumerate}
\end{theorem}

\begin{proof}
Nous pouvons appliquer le Lemme \ref{spaces-more-morphisms-lemma-stein-universally-closed} pour obtenir le
morphisme $f' : X \to Y'$.
Remarquons qu'outre les
conclusions du Lemme \ref{spaces-more-morphisms-lemma-stein-universally-closed}, nous
savons aussi que $f'$ est séparé
(Morphismes d'espaces, Lemme
\ref{spaces-morphisms-lemma-compose-after-separated})
et de type fini
(Morphismes d'espaces, Lemme
\ref{spaces-morphisms-lemma-permanence-finite-type}).
Ainsi, $f'$ est propre. Nous avons alors démontré toutes les
assertions, sauf celle
selon laquelle $f'$ a des fibres géométriques connexes.

\medskip\noindent
Il ressort clairement de la discussion que nous pouvons remplacer $Y$ par $Y'$.
Alors $f : X \to Y$ est propre et $f_*\mathcal{O}_X = \mathcal{O}_Y$.
Remarquons que ces conditions sont préservées par changement de base plat
(Morphismes d'espaces, Lemme \ref{spaces-morphisms-lemma-base-change-proper}
et
Cohomologie des espaces, Lemme
\ref{spaces-cohomology-lemma-flat-base-change-cohomology}).
Soit $\overline{y}$ un point géométrique de $Y$. D'après le
Lemme \ref{spaces-more-morphisms-lemma-characterize-geometrically-connected-fibres}
et la remarque précédente, nous nous ramenons au cas où $Y$ est un
schéma, $y \in Y$ est un point, $f : X \to Y$ est un espace algébrique propre
sur $Y$ avec $f_*\mathcal{O}_X = \mathcal{O}_Y$,
et où nous devons montrer que la fibre $X_y$ est connexe.
En remplaçant $Y$ par un voisinage affine de $y$, nous pouvons
supposer que $Y = \Spec(R)$ est affine. Alors $f_*\mathcal{O}_X = \mathcal{O}_Y$
signifie que l'homomorphisme d'anneaux
$R \to \Gamma(X, \mathcal{O}_X)$ est bijectif.

\medskip\noindent
D'après Limites d'espaces, Lemme
\ref{spaces-limits-lemma-proper-limit-of-proper-finite-presentation-noetherian}
nous pouvons écrire $(X \to Y) = \lim (X_i \to Y_i)$, où $X_i \to Y_i$
est propre et de présentation finie et où $Y_i$ est noethérien. Pour $i$ assez
grand, $Y_i$ est affine (Limites d'espaces, Lemme
\ref{spaces-limits-lemma-limit-is-affine}).
Écrivons $Y_i = \Spec(R_i)$. Posons $R'_i = \Gamma(X_i, \mathcal{O}_{X_i})$.
Remarquons que nous avons des homomorphismes d'anneaux $R_i \to R_i' \to R$. En effet, nous avons
le premier parce que $X_i$ est un espace algébrique sur $R_i$, et le second parce que
nous avons $X \to X_i$ et $R = \Gamma(X, \mathcal{O}_X)$. Remarquons que
$R = \colim R'_i$ d'après Limites d'espaces, Lemme
\ref{spaces-limits-lemma-descend-section}.
Alors
$$
\xymatrix{
X \ar[d]  \ar[r] & X_i \ar[d] \\
Y \ar[r] & Y'_i \ar[r] & Y_i
}
$$
est commutatif, avec $Y'_i = \Spec(R'_i)$.
Soit $y'_i \in Y'_i$ l'image de $y$.
Nous avons $X_y = \lim X_{i, y'_i}$ parce que $X = \lim X_i$,
$Y = \lim Y'_i$ et $\kappa(y) = \colim \kappa(y'_i)$.
Écrivons maintenant $X_y = U \amalg V$, où $U$ et $V$ sont ouverts et fermés.
Alors $U, V$ sont les images inverses d'ouverts $U_i, V_i$ de $X_{i, y'_i}$
(Limites d'espaces, Lemme \ref{spaces-limits-lemma-descend-opens}).
D'après le Théorème \ref{spaces-more-morphisms-theorem-stein-factorization-Noetherian}, les fibres
de $X_i \to Y'_i$ sont connexes ; ainsi, $U$ ou $V$ est vide.
Cela achève la démonstration.
\end{proof}

\noindent
Voici une application.

\begin{lemma}
\label{spaces-more-morphisms-lemma-geometrically-connected-fibres-towards-normal}
Soit $S$ un schéma. Soit $f : X \to Y$ un morphisme d'espaces
algébriques sur $S$. Supposons que
\begin{enumerate}
\item $f$ soit propre ;
\item $Y$ soit intègre (Espaces sur les corps, Définition
\ref{spaces-over-fields-definition-integral-algebraic-space})
de point générique $\xi$ ;
\item $Y$ soit normal ;
\item $X$ soit réduit ;
\item tout point générique d'une composante irréductible de $|X|$ s'envoie sur $\xi$ ;
\item nous ayons $H^0(X_\xi, \mathcal{O}) = \kappa(\xi)$.
\end{enumerate}
Alors $f_*\mathcal{O}_X = \mathcal{O}_Y$ et $f$
a des fibres géométriquement connexes.
\end{lemma}

\begin{proof}
Appliquons le Théorème \ref{spaces-more-morphisms-theorem-stein-factorization-general} pour obtenir une
factorisation $X \to Y' \to Y$. Il suffit de montrer que $Y' = Y$.
Il suffit de montrer que $Y' \times_Y V \to V$ est un isomorphisme,
où $V \to Y$ est un morphisme étale et $V$ un schéma affine intègre ;
voir Espaces sur les corps, Lemme
\ref{spaces-over-fields-lemma-normal-integral-cover-by-affines}.
La formation de $Y'$ commute au changement de base étale ; voir
Morphismes d'espaces, Lemme
\ref{spaces-morphisms-lemma-properties-normalization}.
Les points génériques de $X \times_Y V$ se trouvent au-dessus des points génériques de $X$
(Espaces décents, Lemme \ref{decent-spaces-lemma-decent-generic-points})
et s'envoient donc sur le point générique de $V$ d'après l'hypothèse (5). De plus, la condition
(6) est préservée par le changement de base $V \to Y$, par exemple par
changement de base plat (Cohomologie des espaces, Lemme
\ref{spaces-cohomology-lemma-flat-base-change-cohomology}).
Il suffit donc de démontrer le lemme lorsque $Y$ est un schéma affine
normal et intègre.

\medskip\noindent
Supposons que $Y$ soit un schéma affine normal et intègre. Nous allons montrer que $Y' \to Y$
est un isomorphisme en appliquant Morphismes, Lemme
\ref{morphisms-lemma-finite-birational-over-normal}.
En effet, $Y'$ est réduit parce que $X$ est réduit (Morphismes d'espaces, Lemme
\ref{spaces-morphisms-lemma-normalization-in-reduced}).
Le morphisme $Y' \to Y$ est entier d'après le théorème cité plus haut.
Puisque $Y$ est décent et que $X \to Y$ est séparé, nous voyons que
$X$ est lui aussi décent ; pour le voir, utilisons
Espaces décents, Lemmes
\ref{decent-spaces-lemma-properties-trivial-implications} et
\ref{decent-spaces-lemma-property-over-property}. D'après l'hypothèse (5),
Morphismes d'espaces, Lemme \ref{spaces-morphisms-lemma-normalization-generic},
et Espaces décents, Lemme \ref{decent-spaces-lemma-decent-generic-points},
nous voyons que tout point générique d'une composante irréductible de $|Y'|$
s'envoie sur $\xi$. D'autre part, puisque $Y'$ est le spectre relatif
de $f_*\mathcal{O}_X$, nous voyons que la fibre schématique
$Y'_\xi$ est le spectre de $H^0(X_\xi, \mathcal{O})$, qui est
égal à $\kappa(\xi)$ par hypothèse. Ainsi, $Y'$ est un schéma
intègre dont le corps des fonctions est égal à celui de $Y$.
Cela achève la démonstration.
\end{proof}

\noindent
Voici une autre application.

\begin{lemma}
\label{spaces-more-morphisms-lemma-proper-flat-nr-geom-conn-comps-lower-semicontinuous}
Soit $S$ un schéma.
Soit $X \to Y$ un morphisme d'espaces algébriques sur $S$.
Notons $f$ ce morphisme. S'il est propre, plat et de présentation finie, alors la fonction
$n_{X/Y} : |Y| \to \mathbf{Z}$ qui compte le nombre de composantes
connexes géométriques des fibres de $f$
(Lemme \ref{spaces-more-morphisms-lemma-base-change-fibres-nr-geometrically-connected-components})
est semi-continue inférieurement.
\end{lemma}

\begin{proof}
La question est locale pour la topologie étale sur $Y$ ; nous pouvons donc supposer, et supposons,
que $Y$ est un schéma affine. Soit $y \in Y$. Posons $n = n_{X/S}(y)$.
Remarquons que $n < \infty$, car la fibre géométrique de $X \to Y$ en $y$
est un espace algébrique propre sur un corps, donc noethérien, et possède
donc un nombre fini de composantes connexes.
Nous devons trouver un voisinage ouvert $V$ de $y$ tel que $n_{X/S}|_V \geq n$.
Soit $X \to Y' \to Y$ la factorisation de Stein du
Théorème \ref{spaces-more-morphisms-theorem-stein-factorization-general}.
D'après le Lemme \ref{spaces-more-morphisms-lemma-stein-universally-closed-residue-fields},
il existe un nombre fini de points
$y'_1, \ldots, y'_m \in Y'$ situés au-dessus de $y$,
et les extensions $\kappa(y'_i)/\kappa(y)$ sont finies.
Compléments sur les morphismes, Lemme \ref{more-morphisms-lemma-etale-makes-integral-split},
nous dit qu'après avoir remplacé $Y$ par un voisinage étale
de $y$, nous pouvons supposer que $Y' = V_1 \amalg \ldots \amalg V_m$ est un schéma,
avec $y'_i \in V_i$ et $\kappa(y'_i)/\kappa(y)$ purement inséparable.
Alors les espaces algébriques $X_{y_i'}$ sont géométriquement connexes sur
$\kappa(y)$, donc $m = n$. Les espaces algébriques
$X_i = (f')^{-1}(V_i)$, $i = 1, \ldots, n$,
sont plats et de présentation finie sur $Y$. L'image de $X_i \to Y$
est donc ouverte (Morphismes d'espaces, Lemme \ref{spaces-morphisms-lemma-fppf-open}).
Ainsi, dans un voisinage de $y$, nous voyons que $n_{X/Y}$ est
au moins égal à $n$.
\end{proof}

\begin{lemma}
\label{spaces-more-morphisms-lemma-proper-flat-geom-red}
Soit $S$ un schéma.
Soit $f : X \to Y$ un morphisme d'espaces algébriques sur $S$. Supposons que
\begin{enumerate}
\item $f$ soit propre, plat et de présentation finie ;
\item les fibres géométriques de $f$ soient réduites.
\end{enumerate}
Alors la fonction $n_{X/S} : |Y| \to \mathbf{Z}$
qui compte le nombre de composantes connexes géométriques
des fibres de $f$
(Lemme \ref{spaces-more-morphisms-lemma-base-change-fibres-nr-geometrically-connected-components})
est localement constante.
\end{lemma}

\begin{proof}
D'après le Lemme \ref{spaces-more-morphisms-lemma-proper-flat-nr-geom-conn-comps-lower-semicontinuous},
la fonction $n_{X/Y}$ est semi-continue inférieurement.
Il suffit donc de montrer qu'elle est semi-continue supérieurement.
Pour cela, nous pouvons travailler localement pour la topologie étale sur $Y$ ; nous
pouvons donc supposer que $Y$ est un schéma affine.
Pour $y \in Y$, considérons la $\kappa(y)$-algèbre
$$
A = H^0(X_y, \mathcal{O}_{X_y})
$$
D'après Espaces sur les corps, Lemme
\ref{spaces-over-fields-lemma-proper-geometrically-reduced-global-sections}
et puisque $X_y$ est géométriquement réduit,
$A$ est un produit fini d'extensions finies séparables de $\kappa(y)$.
Ainsi, $A \otimes_{\kappa(y)} \kappa(\overline{y})$ est un produit
de $\beta_0(y) = \dim_{\kappa(y)} A$ exemplaires de $\kappa(\overline{y})$.
Par conséquent, $X_{\overline{y}}$ a $\beta_0(y)$ composantes connexes.
Autrement dit, nous avons $n_{X/S} = \beta_0$ comme fonctions sur $Y$.
Ainsi, $n_{X/Y}$ est semi-continue supérieurement d'après
Catégories dérivées des espaces, Lemme
\ref{spaces-perfect-lemma-jump-loci-geometric}.
Cela achève la démonstration.
\end{proof}

\begin{lemma}
\label{spaces-more-morphisms-lemma-stein-factorization-etale}
Soit $S$ un schéma.
Soit $f : X \to Y$ un morphisme propre d'espaces algébriques sur $S$.
Soit $X \to Y' \to Y$ la factorisation de Stein de $f$
(Théorème \ref{spaces-more-morphisms-theorem-stein-factorization-general}).
Si $f$ est de présentation finie, plat et à fibres
géométriquement réduites (Définition \ref{spaces-more-morphisms-definition-geometrically-reduced-fibre}),
alors $Y' \to Y$ est fini étale.
\end{lemma}

\begin{proof}
La formation de la factorisation de Stein commute au changement de base plat ;
voir Lemme \ref{spaces-more-morphisms-lemma-stein-universally-closed}.
Nous pouvons donc travailler localement pour la topologie étale sur $Y$ et supposer que $Y$
est un schéma affine. Alors $Y'$ est un schéma affine et $Y' \to Y$
est entier.

\medskip\noindent
Soit $y \in Y$. Soit $n$ le nombre de composantes connexes de
la fibre géométrique $X_{\overline{y}}$. Remarquons que $n < \infty$,
car la fibre géométrique de $X \to Y$ en $y$ est un espace
algébrique propre sur un corps, donc noethérien,
et possède donc un nombre fini de composantes connexes.
D'après le Lemme \ref{spaces-more-morphisms-lemma-stein-universally-closed-residue-fields},
il existe un nombre fini de points $y'_1, \ldots, y'_m \in Y'$ situés au-dessus de $y$ ;
pour chaque $i$, nous pouvons choisir un point de type fini $x_i \in |X_y|$
qui s'envoie sur $y'_i$, et l'extension $\kappa(y'_i)/\kappa(y)$ est finie.
Ainsi, Compléments sur les morphismes,
Lemme \ref{more-morphisms-lemma-etale-makes-integral-split},
nous dit qu'après avoir remplacé $Y$ par un voisinage étale
de $y$, nous pouvons supposer que $Y' = V_1 \amalg \ldots \amalg V_m$ est un schéma,
avec $y'_i \in V_i$ et $\kappa(y'_i)/\kappa(y)$ purement inséparable.
Dans ce cas, les espaces algébriques $X_{y_i'}$
sont géométriquement connexes sur $\kappa(y)$, donc $m = n$.
Les espaces algébriques $X_i = (f')^{-1}(V_i)$, $i = 1, \ldots, n$,
sont propres, plats, de présentation finie et à fibres géométriquement
réduites sur $Y$. Il suffit de démontrer le lemme
pour chacun des morphismes $X_i \to Y$. Nous nous ramenons ainsi au cas où
$X_{\overline{y}}$ est connexe.

\medskip\noindent
Supposons que $X_{\overline{y}}$ soit connexe. D'après le
Lemme \ref{spaces-more-morphisms-lemma-proper-flat-geom-red},
nous voyons que $X \to Y$ a des fibres géométriquement connexes
dans un voisinage de $y$. Ainsi,
nous pouvons supposer que les fibres de $X \to Y$ sont géométriquement connexes.
Alors $f_*\mathcal{O}_X = \mathcal{O}_Y$ d'après
Catégories dérivées des espaces, Lemme
\ref{spaces-perfect-lemma-proper-flat-geom-red-connected},
ce qui achève la démonstration.
\end{proof}

\noindent
La démonstration du lemme suivant utilise la factorisation de Stein pour les schémas ;
c'est pourquoi il se trouve dans cette section.

\begin{lemma}
\label{spaces-more-morphisms-lemma-split-off-proper-part-henselian}
Soit $(A, I)$ un couple hensélien. Soit $X$ un espace algébrique
séparé et de type fini sur $A$. Posons
$X_0 = X \times_{\Spec(A)} \Spec(A/I)$.
Soit $Y \subset X_0$ un sous-espace ouvert et fermé tel que
$Y \to \Spec(A/I)$ soit propre. Il existe alors un sous-espace ouvert et fermé
$W \subset X$ qui est propre sur $A$ et tel que
$W \times_{\Spec(A)} \Spec(A/I) = Y$.
\end{lemma}

\begin{proof}
Nous noterons $T \mapsto T_0$ le changement de base par
$\Spec(A/I) \to \Spec(A)$.
D'après une forme faible du lemme de Chow (sous la forme de
Cohomologie des espaces, Lemme \ref{spaces-cohomology-lemma-weak-chow}),
il existe un morphisme propre surjectif $\varphi : X' \to X$ tel
que $X'$ admette une immersion dans $\mathbf{P}^n_A$.
Posons $Y' = \varphi^{-1}(Y)$. C'est un sous-schéma ouvert et fermé
de $X'_0$. Le lemme est vrai pour $(X', Y')$ d'après
Compléments sur les morphismes, Lemme
\ref{more-morphisms-lemma-split-off-proper-part-henselian}.
Soit $W' \subset X'$ le sous-schéma ouvert et fermé, propre
sur $A$, tel que $Y' = W'_0$.
D'après Morphismes d'espaces, Lemme
\ref{spaces-morphisms-lemma-universally-closed-permanence}
$Q_1 = \varphi(|W'|) \subset |X|$ et
$Q_2 = \varphi(|X' \setminus W'|) \subset |X|$
sont des fermés et, d'après
Morphismes d'espaces, Lemme \ref{spaces-morphisms-lemma-image-proper-is-proper},
toute structure de sous-espace fermé sur $Q_1$ est propre sur $A$.
L'image de $Q_1 \cap Q_2$ dans $\Spec(A)$ est fermée.
Puisque $(A, I)$ est hensélien, si $Q_1 \cap Q_2$ est non vide, alors nous
trouvons que $Q_1 \cap Q_2$ a un point situé au-dessus de $\Spec(A/I)$.
C'est impossible puisque $W'_0 = Y' = \varphi^{-1}(Y)$.
Nous concluons que $Q_1$ est ouvert et fermé dans $|X|$.
Soit $W \subset X$ le sous-espace ouvert et fermé
correspondant. Alors $W$ est propre sur $A$ et $W_0 = Y$.
\end{proof}










\section{Prolongement de propriétés à partir d'un ouvert}
\label{spaces-more-morphisms-section-extending-properties}

\noindent
Dans cette section, nous rassemblons plusieurs résultats de la forme suivante : si $f : X \to Y$
est un morphisme plat d'espaces algébriques et si $f$ possède une certaine propriété au-dessus
d'un ouvert dense
de $Y$, alors $f$ possède la même propriété au-dessus de $Y$ tout entier.

\begin{lemma}
\label{spaces-more-morphisms-lemma-flat-finite-type-finitely-presented-over-dense-open}
Soit $S$ un schéma.
Soit $f : X \to Y$ un morphisme d'espaces algébriques sur $S$.
Soit $\mathcal{F}$ un $\mathcal{O}_X$-module quasi-cohérent.
Soit $V \subset Y$ un sous-espace ouvert. Supposons que
\begin{enumerate}
\item $f$ soit localement de présentation finie,
\item $\mathcal{F}$ soit de type fini et plat sur $Y$,
\item $V \to Y$ soit quasi-compact et schématiquement dense,
\item $\mathcal{F}|_{f^{-1}V}$ soit de présentation finie.
\end{enumerate}
Alors $\mathcal{F}$ est de présentation finie.
\end{lemma}

\begin{proof}
Il suffit de montrer que l'image inverse de $\mathcal{F}$ sur un schéma surjectif
et étale sur $X$ est de présentation finie. Nous pouvons donc supposer que $X$
est un schéma. De même, nous pouvons remplacer $Y$ par un schéma surjectif et
étale sur $Y$ (l'image inverse de $V$ dans ce schéma est schématiquement
dense ; voir
Morphismes d'espaces, section
\ref{spaces-morphisms-section-scheme-theoretic-closure}).
Nous nous ramenons ainsi au cas des schémas, qui est
Compléments sur la platitude, Lemme
\ref{flat-lemma-flat-finite-type-finitely-presented-over-dense-open}.
\end{proof}

\begin{lemma}
\label{spaces-more-morphisms-lemma-flat-finite-type-finitely-presented-over-dense-open-X}
Soit $S$ un schéma.
Soit $f : X \to Y$ un morphisme d'espaces algébriques sur $S$.
Soit $V \subset Y$ un sous-espace ouvert.
Supposons que
\begin{enumerate}
\item $f$ soit localement de type fini et plat,
\item $V \to Y$ soit quasi-compact et schématiquement dense,
\item $f|_{f^{-1}V} : f^{-1}V \to V$ soit localement de présentation finie.
\end{enumerate}
Alors $f$ est localement de présentation finie.
\end{lemma}

\begin{proof}
La démonstration est identique à celle du
Lemme \ref{spaces-more-morphisms-lemma-flat-finite-type-finitely-presented-over-dense-open},
à ceci près qu'on utilise
Compléments sur la platitude, Lemme
\ref{flat-lemma-flat-finite-type-finitely-presented-over-dense-open-X}.
\end{proof}

\begin{lemma}
\label{spaces-more-morphisms-lemma-flat-fp-dimension-over-dense-open}
Soit $S$ un schéma. Soit $f : X \to Y$ un morphisme d'espaces algébriques
sur $S$, plat et localement de type fini. Soit $V \subset Y$ un
sous-espace ouvert tel que $|V| \subset |Y|$ soit dense et que $X_V \to V$
soit de dimension relative $\leq d$. Supposons en outre que l'une des conditions suivantes soit vérifiée :
\begin{enumerate}
\item $f$ est localement de présentation finie, ou
\item $V \to Y$ est quasi-compact.
\end{enumerate}
Alors $f : X \to Y$ est de dimension relative $\leq d$.
\end{lemma}

\begin{proof}
Nous pouvons remplacer $Y$ par son réduit et donc supposer $Y$ réduit.
Alors $V$ est schématiquement dense dans $Y$ ; voir
Morphismes d'espaces, Lemme
\ref{spaces-morphisms-lemma-quasi-compact-immersion}.
Par définition, la propriété d'être de dimension relative $\leq d$ peut
se vérifier sur un recouvrement étale ; voir
Morphismes d'espaces, section \ref{spaces-morphisms-section-relative-dimension}.
Il suffit donc de montrer que $f$ est de dimension relative $\leq d$
après avoir remplacé $X$ par un schéma surjectif et étale sur $X$.
De même, nous pouvons remplacer $Y$ par un schéma surjectif et
étale sur $Y$. L'image inverse de $V$ dans ce schéma est schématiquement
dense ; voir
Morphismes d'espaces, section
\ref{spaces-morphisms-section-scheme-theoretic-closure}.
Comme un ouvert schématiquement dense d'un schéma est en particulier
dense, nous nous ramenons au cas des schémas, qui est
Compléments sur la platitude, Lemme
\ref{flat-lemma-flat-finite-presentation-dimension-over-dense-open}.
\end{proof}

\begin{lemma}
\label{spaces-more-morphisms-lemma-proper-flat-finite-over-dense-open}
Soit $S$ un schéma. Soit $f : X \to Y$ un morphisme d'espaces algébriques
sur $S$, plat et propre. Soit $V \to Y$ un sous-espace ouvert
tel que $|V| \subset |Y|$ soit dense et que $X_V \to V$ soit fini. Si, de plus,
$f$ est localement de présentation finie ou $V \to Y$ est quasi-compact,
alors $f$ est fini.
\end{lemma}

\begin{proof}
D'après le Lemme \ref{spaces-more-morphisms-lemma-flat-fp-dimension-over-dense-open},
les fibres de $f$ sont de dimension nulle.
D'après Morphismes d'espaces, Lemme
\ref{spaces-morphisms-lemma-locally-quasi-finite-rel-dimension-0}
il s'ensuit que $f$ est localement quasi-fini.
D'après Morphismes d'espaces, Lemme
\ref{spaces-morphisms-lemma-locally-quasi-finite-separated-representable}
il s'ensuit que $f$ est représentable.
La finitude de $f$ se vérifie localement pour la topologie étale sur $Y$ ;
nous pouvons donc supposer que $Y$ est un schéma. Comme $f$ est représentable,
nous nous ramenons au cas des schémas, qui est
Compléments sur la platitude, Lemme \ref{flat-lemma-proper-flat-finite-over-dense-open}.
\end{proof}

\begin{lemma}
\label{spaces-more-morphisms-lemma-zariski}
Soit $S$ un schéma.
Soit $f : X \to Y$ un morphisme d'espaces algébriques sur $S$.
Soit $V \subset Y$ un sous-espace ouvert. Si
\begin{enumerate}
\item $f$ est séparé, localement de type fini et plat,
\item $f^{-1}(V) \to V$ est un isomorphisme, et
\item $V \to Y$ est quasi-compact et schématiquement dense,
\end{enumerate}
alors $f$ est une immersion ouverte.
\end{lemma}

\begin{proof}
En appliquant le
Lemme \ref{spaces-more-morphisms-lemma-flat-finite-type-finitely-presented-over-dense-open-X},
nous voyons que $f$ est localement de présentation finie. En appliquant le
Lemme \ref{spaces-more-morphisms-lemma-flat-fp-dimension-over-dense-open},
nous voyons que $f$ est de dimension relative $\leq 0$.
D'après Morphismes d'espaces, Lemme
\ref{spaces-morphisms-lemma-locally-quasi-finite-rel-dimension-0}
il s'ensuit que $f$ est localement quasi-fini.
D'après Morphismes d'espaces, Lemme
\ref{spaces-morphisms-lemma-locally-quasi-finite-separated-representable}
il s'ensuit que $f$ est représentable.
D'après Descente sur les espaces, Lemme
\ref{spaces-descent-lemma-descending-property-open-immersion}
le fait que $f$ soit une immersion ouverte se vérifie localement pour la topologie étale sur $Y$.
Nous pouvons donc supposer que $Y$ est un schéma. Comme $f$ est représentable,
nous nous ramenons au cas des schémas, qui est
Compléments sur la platitude, Lemme \ref{flat-lemma-zariski}.
\end{proof}







\section{Éclatement et platitude}
\label{spaces-more-morphisms-section-blowup-flat}

\noindent
Au lieu de reprendre le travail de
Compléments sur la platitude, section \ref{flat-section-blowup-flat},
nous démontrons un analogue de Compléments sur la platitude, Lemme \ref{flat-lemma-push-ideal},
qui nous dit que le problème de trouver un éclatement convenable
est souvent local pour la topologie étale sur la base.

\begin{lemma}
\label{spaces-more-morphisms-lemma-push-ideal}
Soit $S$ un schéma. Soit $X$ un espace algébrique quasi-compact et quasi-séparé
sur $S$. Soit $\varphi : W \to X$ un morphisme étale, quasi-compact
et séparé. Soit $U \subset X$ un sous-espace ouvert quasi-compact.
Soit $\mathcal{I} \subset \mathcal{O}_W$ un faisceau quasi-cohérent
d'idéaux de type fini tel que
$V(\mathcal{I}) \cap \varphi^{-1}(U) = \emptyset$.
Alors il existe un faisceau quasi-cohérent d'idéaux de type fini
$\mathcal{J} \subset \mathcal{O}_X$ tel que
\begin{enumerate}
\item $V(\mathcal{J}) \cap U = \emptyset$, et
\item $\varphi^{-1}(\mathcal{J})\mathcal{O}_W = \mathcal{I} \mathcal{I}'$
pour un certain idéal quasi-cohérent de type fini
$\mathcal{I}' \subset \mathcal{O}_W$.
\end{enumerate}
\end{lemma}

\begin{proof}
Choisissons une factorisation $W \to Y \to X$ dans laquelle $j : W \to Y$ est une
immersion ouverte quasi-compacte et $\pi : Y \to X$ est un morphisme fini
de présentation finie
(Lemme \ref{spaces-more-morphisms-lemma-quasi-finite-separated-pass-through-finite-addendum}).
Posons $V = j(W) \cup \pi^{-1}(U) \subset Y$. Notons que
$\mathcal{I}$ sur $W \cong j(W)$ et $\mathcal{O}_{\pi^{-1}(U)}$
se recollent en un faisceau quasi-cohérent d'idéaux de type fini
$\mathcal{I}_1 \subset \mathcal{O}_V$. D'après
Limites d'espaces, Lemme \ref{spaces-limits-lemma-extend},
il existe un faisceau quasi-cohérent d'idéaux de type fini
$\mathcal{I}_2 \subset \mathcal{O}_Y$ tel que
$\mathcal{I}_2|_V = \mathcal{I}_1$. Autrement dit,
$\mathcal{I}_2 \subset \mathcal{O}_Y$
est un faisceau quasi-cohérent d'idéaux de type fini tel que
$V(\mathcal{I}_2)$ soit disjoint de $\pi^{-1}(U)$ et que
$j^{-1}\mathcal{I}_2 = \mathcal{I}$. Notons $i : Z \to Y$
l'immersion fermée correspondante, qui est de présentation finie
(Morphismes d'espaces, Lemme
\ref{spaces-morphisms-lemma-closed-immersion-finite-presentation}).
En particulier, la composée $\tau = \pi \circ i : Z \to X$ est finie
et de présentation finie
(Morphismes d'espaces, Lemmes
\ref{spaces-morphisms-lemma-composition-finite-presentation} et
\ref{spaces-morphisms-lemma-composition-integral}).

\medskip\noindent
Posons $\mathcal{F} = \tau_*\mathcal{O}_Z$, que nous considérons comme
un $\mathcal{O}_X$-module quasi-cohérent. D'après
Descente sur les espaces, Lemme
\ref{spaces-descent-lemma-finite-finitely-presented-module}
nous voyons que $\mathcal{F}$ est un $\mathcal{O}_X$-module de présentation finie.
Posons $\mathcal{J} = \text{Fit}_0(\mathcal{F})$. (Insérer ici une référence aux
idéaux de Fitting sur les topos annelés.) C'est un faisceau quasi-cohérent
d'idéaux de type fini sur $X$ (puisque $\mathcal{F}$ est de présentation finie, voir
Compléments d'algèbre, Lemme \ref{more-algebra-lemma-fitting-ideal-basics}).
La partie (1) du lemme vaut parce que $|\tau|(|Z|) \cap |U| = \emptyset$
par notre choix de $\mathcal{I}_2$ et parce que l'idéal de Fitting d'indice $0$
du module nul est égal au faisceau structural. Pour démontrer (2), notons que
$\varphi^{-1}(\mathcal{J})\mathcal{O}_W = \text{Fit}_0(\varphi^*\mathcal{F})$
car la formation des idéaux de Fitting commute au changement de base.
D'autre part, puisque $\varphi : W \to X$ est séparé et étale,
nous voyons que $(1, j) : W \to W \times_X Y$ est une immersion ouverte et fermée.
Ainsi $W \times_Y Z = V(\mathcal{I}) \amalg Z'$ pour un certain morphisme d'espaces
algébriques fini et de présentation finie $\tau' : Z' \to W$.
Nous voyons donc que
\begin{align*}
\text{Fit}_0(\varphi^*\mathcal{F}) & =
\text{Fit}_0((W \times_Y Z \to W)_*\mathcal{O}_{W \times_Y Z}) \\
& =
\text{Fit}_0(\mathcal{O}_W/\mathcal{I}) \cdot
\text{Fit}_0(\tau'_*\mathcal{O}_{Z'}) \\
& =
\mathcal{I} \cdot \text{Fit}_0(\tau'_*\mathcal{O}_{Z'})
\end{align*}
où la deuxième égalité résulte de
Compléments d'algèbre, Lemme \ref{more-algebra-lemma-fitting-ideal-basics},
transposé aux faisceaux sur les topos annelés.
En posant $\mathcal{I}' = \text{Fit}_0(\tau'_*\mathcal{O}_{Z'})$,
on achève la démonstration du lemme.
\end{proof}

\begin{theorem}
\label{spaces-more-morphisms-theorem-flatten-module}
Soit $S$ un schéma. Soit $B$ un espace algébrique quasi-compact et quasi-séparé
sur $S$. Soit $X$ un espace algébrique sur $B$.
Soit $\mathcal{F}$ un module quasi-cohérent sur $X$.
Soit $U \subset B$ un sous-espace ouvert quasi-compact. Supposons que
\begin{enumerate}
\item $X$ soit quasi-compact,
\item $X$ soit localement de présentation finie sur $B$,
\item $\mathcal{F}$ soit un module de type fini,
\item $\mathcal{F}_U$ soit de présentation finie, et
\item $\mathcal{F}_U$ soit plat sur $U$.
\end{enumerate}
Alors il existe un éclatement $U$-admissible $B' \to B$ tel que la
transformée stricte $\mathcal{F}'$ de $\mathcal{F}$ soit un
$\mathcal{O}_{X \times_B B'}$-module de présentation finie et
plat sur $B'$.
\end{theorem}

\begin{proof}
Choisissons un schéma affine $V$ et un morphisme étale surjectif $V \to X$.
Puisque la transformée stricte commute à la localisation étale
(Diviseurs sur les espaces, Lemme \ref{spaces-divisors-lemma-strict-transform-local}),
il suffit de démontrer le résultat en remplaçant $X$ par $V$. Nous pouvons donc
supposer que $X \to B$ est représentable (en plus des
hypothèses du lemme).

\medskip\noindent
Supposons que $X \to B$ soit représentable. Choisissons un schéma affine $W$ et un
morphisme étale surjectif $\varphi : W \to B$. Notons que $X \times_B W$
est un schéma. D'après le cas des schémas
(Compléments sur la platitude, Théorème \ref{flat-theorem-flatten-module}),
nous pouvons trouver un faisceau quasi-cohérent d'idéaux de type fini
$\mathcal{I} \subset \mathcal{O}_W$ tel que
(a) $|V(\mathcal{I})| \cap |\varphi^{-1}(U)| = \emptyset$ et (b)
la transformée stricte de $\mathcal{F}|_{X \times_B W}$ relativement
à l'éclatement $W' \to W$ en $\mathcal{I}$ devienne plate sur $W'$
et soit un module de présentation finie. Choisissons un faisceau d'idéaux de type fini
$\mathcal{J} \subset \mathcal{O}_B$ comme dans le Lemme \ref{spaces-more-morphisms-lemma-push-ideal}.
Soit $B' \to B$ l'éclatement de $\mathcal{J}$. Nous affirmons que cet
éclatement convient. En effet, il est clair que $B' \to B$ est $U$-admissible
par notre choix de l'idéal $\mathcal{J}$. De plus, le changement de base
$B' \times_B W \to W$ est l'éclatement de $W$ en
$\varphi^{-1}\mathcal{J} = \mathcal{I}\mathcal{I}'$
(compatibilité de l'éclatement avec le changement de base plat, voir
Diviseurs sur les espaces, Lemme
\ref{spaces-divisors-lemma-flat-base-change-blowing-up}).
Il existe donc une factorisation
$$
W \times_B B' \to W' \to W
$$
où le premier morphisme est lui aussi un éclatement, voir
Diviseurs sur les espaces, Lemme \ref{spaces-divisors-lemma-blowing-up-two-ideals}).
La restriction de $\mathcal{F}'$ (qui est défini sur $B' \times_B X$)
à $W \times_B B' \times_B X$ est la transformée stricte de
$\mathcal{F}|_{X \times_B W}$
(Diviseurs sur les espaces, Lemme \ref{spaces-divisors-lemma-strict-transform-local})
et est donc la transformée stricte effectuée deux fois de
$\mathcal{F}|_{X \times_B W}$ par les deux éclatements affichés ci-dessus
(Diviseurs sur les espaces, Lemme
\ref{spaces-divisors-lemma-strict-transform-composition-blowups}).
Après le premier éclatement, notre faisceau est déjà plat sur
la base et de présentation finie (par construction). Il en est donc encore ainsi 
après la deuxième transformée stricte (puisque celle-ci est une
image inverse d'après Diviseurs sur les espaces, Lemme
\ref{spaces-divisors-lemma-strict-transform-flat}).
Nous voyons ainsi que la restriction de $\mathcal{F}'$
à un recouvrement étale de $B' \times_B X$ possède les propriétés voulues,
ce qui démontre le théorème.
\end{proof}







\section{Applications}
\label{spaces-more-morphisms-section-applications-flattening-by-blowing-up}

\noindent
Dans cette section, nous appliquons le résultat d'aplatissement par éclatement.

\begin{lemma}
\label{spaces-more-morphisms-lemma-flat-after-blowing-up}
Soit $S$ un schéma. Soit $f : X \to B$ un morphisme d'espaces
algébriques sur $S$. Soit $U \subset B$ un sous-espace ouvert. Supposons que
\begin{enumerate}
\item $B$ soit quasi-compact et quasi-séparé,
\item $U$ soit quasi-compact,
\item $f : X \to B$ soit de type fini et quasi-séparé, et
\item $f^{-1}(U) \to U$ soit plat et localement de présentation finie.
\end{enumerate}
Alors il existe un éclatement $U$-admissible $B' \to B$ tel que
le transformé strict $X'$ de $X$ soit plat et de présentation finie
sur $B'$.
\end{lemma}

\begin{proof}
Soit $B' \to B$ un éclatement $U$-admissible. Notons que le transformé strict
de $X$ est quasi-compact et quasi-séparé sur $B'$, puisque $X$ est quasi-compact
et quasi-séparé sur $B$. Il nous suffit donc de trouver
un éclatement $U$-admissible tel que le transformé strict devienne plat et
localement de présentation finie. Nous ne pouvons pas appliquer directement le
Théorème \ref{spaces-more-morphisms-theorem-flatten-module}, car $X$ n'est pas localement de
présentation finie sur $B$.

\medskip\noindent
Choisissons un schéma affine $V$ et un morphisme étale surjectif $V \to X$.
(C'est possible puisque $X$ est quasi-compact, étant un espace de type fini sur
l'espace quasi-compact $B$.) Il suffit alors de démontrer le résultat pour
le morphisme $V \to B$ (puisque le transformé strict commute à la localisation
étale, voir Diviseurs sur les espaces,
Lemme \ref{spaces-divisors-lemma-strict-transform-local}).
Nous pouvons donc supposer que $X \to B$ est séparé et de type fini.
Dans ce cas, nous pouvons trouver une immersion fermée $i : X \to Y$ avec $Y \to B$
séparé et de présentation finie, voir
Limites d'espaces, Proposition
\ref{spaces-limits-proposition-separated-closed-in-finite-presentation}.

\medskip\noindent
Appliquons le Théorème \ref{spaces-more-morphisms-theorem-flatten-module} à $\mathcal{F} = i_*\mathcal{O}_X$
sur $Y/B$. Nous trouvons un éclatement $U$-admissible $B' \to B$ tel que la
transformée stricte de $\mathcal{F}$ soit plate sur $B'$ et de présentation finie.
Soit $X'$ le transformé strict de $X$ par l'éclatement $B' \to B$.
Soit $i' : X' \to Y \times_B B'$ le morphisme induit.
Puisque la formation de la transformée stricte commute à l'image directe par les morphismes
affines (Diviseurs sur les espaces, Lemme
\ref{spaces-divisors-lemma-strict-transform-affine}),
nous voyons que $i'_*\mathcal{O}_{X'}$ est plat sur $B'$ et de
présentation finie comme $\mathcal{O}_{Y \times_B B'}$-module.
Ainsi $X' \to B'$ est plat et localement de présentation finie.
Cela implique le lemme d'après nos remarques précédentes.
\end{proof}

\begin{lemma}
\label{spaces-more-morphisms-lemma-finite-after-blowing-up}
Soit $S$ un schéma. Soit $f : X \to B$ un morphisme d'espaces
algébriques sur $S$. Soit $U \subset B$ un sous-espace ouvert. Supposons que
\begin{enumerate}
\item $B$ soit quasi-compact et quasi-séparé,
\item $U$ soit quasi-compact,
\item $f : X \to B$ soit propre, et
\item $f^{-1}(U) \to U$ soit fini localement libre.
\end{enumerate}
Alors il existe un éclatement $U$-admissible $B' \to B$ tel que
le transformé strict $X'$ de $X$ soit fini localement libre sur $B'$.
\end{lemma}

\begin{proof}
D'après le Lemme \ref{spaces-more-morphisms-lemma-flat-after-blowing-up}, nous pouvons supposer que
$X \to B$ est plat et de présentation finie. Après avoir remplacé
$B$ par un éclatement $U$-admissible si nécessaire, nous pouvons supposer
que $U \subset B$ est schématiquement dense. Alors $f$ est
fini d'après le Lemme \ref{spaces-more-morphisms-lemma-proper-flat-finite-over-dense-open}.
Par conséquent, $f$ est fini localement libre d'après
Morphismes d'espaces, Lemme \ref{spaces-morphisms-lemma-finite-flat}.
\end{proof}

\begin{lemma}
\label{spaces-more-morphisms-lemma-isomorphism-after-blowing-up}
Soit $S$ un schéma. Soit $f : X \to B$ un morphisme d'espaces
algébriques sur $S$. Soit $U \subset B$ un sous-espace ouvert. Supposons que
\begin{enumerate}
\item $B$ soit quasi-compact et quasi-séparé,
\item $U$ soit quasi-compact,
\item $f : X \to B$ soit propre, et
\item $f^{-1}(U) \to U$ soit un isomorphisme.
\end{enumerate}
Alors il existe un éclatement $U$-admissible $B' \to B$ tel que
le transformé strict $X'$ de $X$ s'envoie isomorphiquement sur $B'$.
\end{lemma}

\begin{proof}
D'après le Lemme \ref{spaces-more-morphisms-lemma-flat-after-blowing-up}, nous pouvons supposer que
$X \to B$ est plat et de présentation finie. Après avoir remplacé
$B$ par un éclatement $U$-admissible si nécessaire, nous pouvons supposer
que $U \subset B$ est schématiquement dense. Alors $f$ est
fini d'après le Lemme \ref{spaces-more-morphisms-lemma-proper-flat-finite-over-dense-open}
et une immersion ouverte d'après le Lemme \ref{spaces-more-morphisms-lemma-zariski}. Ainsi $f$
est une immersion ouverte dont l'image est fermée et contient
l'ouvert dense $U$ ; par conséquent, $f$ est un isomorphisme.
\end{proof}

\begin{lemma}
\label{spaces-more-morphisms-lemma-zariski-after-blowup}
Soit $S$ un schéma. Soit $f : X \to B$ un morphisme d'espaces algébriques
sur $S$. Soit $U \subset B$ un sous-espace ouvert. Supposons que
\begin{enumerate}
\item $B$ soit quasi-compact et quasi-séparé,
\item $U$ soit quasi-compact,
\item $f$ soit de type fini,
\item $f^{-1}(U) \to U$ soit un isomorphisme.
\end{enumerate}
Alors il existe un éclatement $U$-admissible $B' \to B$ tel que $U$
soit schématiquement dense dans $B'$ et que le transformé
strict $X'$ de $X$ s'envoie isomorphiquement sur un sous-espace ouvert de $B'$.
\end{lemma}

\begin{proof}
Ce lemme généralise le
Lemme \ref{spaces-more-morphisms-lemma-isomorphism-after-blowing-up}.
Comme la composée d'éclatements $U$-admissibles est $U$-admissible
(Diviseurs sur les espaces, Lemme
\ref{spaces-divisors-lemma-composition-admissible-blowups})
nous pouvons procéder par étapes. Choisissons un faisceau quasi-cohérent
d'idéaux de type fini $\mathcal{I} \subset \mathcal{O}_B$ tel que
$|B| \setminus |U| = |V(\mathcal{I})|$. Remplaçons $B$ par l'éclatement
de $B$ en $\mathcal{I}$ et $X$ par le transformé strict de $X$.
Après ce remplacement, $B \setminus U$ est le support d'un diviseur de
Cartier effectif $D$ (Diviseurs sur les espaces, Lemme
\ref{spaces-divisors-lemma-blowing-up-gives-effective-Cartier-divisor}).
En particulier, $U$ est schématiquement dense dans $B$
(Diviseurs sur les espaces, Lemme
\ref{spaces-divisors-lemma-complement-effective-Cartier-divisor}).
Ensuite, nous effectuons un autre éclatement $U$-admissible pour nous ramener au cas où
$X \to B$ est plat et de présentation finie, voir le
Lemme \ref{spaces-more-morphisms-lemma-flat-after-blowing-up}. Notons que $U$ reste schématiquement
dense dans $B$. Par conséquent, $X \to B$ est une immersion ouverte d'après le
Lemme \ref{spaces-more-morphisms-lemma-zariski}.
\end{proof}

\noindent
Le lemme suivant affirme qu'une modification peut être dominée
par un éclatement.

\begin{lemma}
\label{spaces-more-morphisms-lemma-dominate-modification-by-blowup}
Soit $S$ un schéma. Soit $f : X \to B$ un morphisme d'espaces algébriques
sur $S$. Soit $U \subset B$ un sous-espace ouvert. Supposons que
\begin{enumerate}
\item $B$ soit quasi-compact et quasi-séparé,
\item $U$ soit quasi-compact,
\item $f : X \to B$ soit propre,
\item $f^{-1}(U) \to U$ soit un isomorphisme.
\end{enumerate}
Alors il existe un éclatement $U$-admissible $B' \to B$ qui domine $X$,
c'est-à-dire qu'il existe une factorisation $B' \to X \to B$
du morphisme d'éclatement.
\end{lemma}

\begin{proof}
D'après le Lemme \ref{spaces-more-morphisms-lemma-isomorphism-after-blowing-up}, nous pouvons trouver un éclatement
$U$-admissible $B' \to B$ tel que le transformé strict $X'$ s'envoie isomorphiquement sur
$B'$. Nous pouvons alors prendre $B' = X' \to X$ comme factorisation.
\end{proof}

\begin{lemma}
\label{spaces-more-morphisms-lemma-get-section-after-blowup}
Soit $S$ un schéma. Soient $X$, $Y$ des espaces algébriques sur $S$.
Soient $U \subset W \subset Y$ des sous-espaces ouverts.
Soient $f : X \to W$ et $s : U \to X$ des morphismes
tels que $f \circ s = \text{id}_U$. Supposons que
\begin{enumerate}
\item $f$ soit propre,
\item $Y$ soit quasi-compact et quasi-séparé, et
\item $U$ et $W$ soient quasi-compacts.
\end{enumerate}
Alors il existe un éclatement $U$-admissible $b : Y' \to Y$ et un morphisme
$s' : b^{-1}(W) \to X$ qui prolonge $s$ avec $f \circ s' = b|_{b^{-1}(W)}$.
\end{lemma}

\begin{proof}
Nous pouvons remplacer, et remplaçons, $X$ par l'image schématique de $s$.
Alors $X \to W$ est un isomorphisme au-dessus de $U$, voir
Morphismes d'espaces, Lemme
\ref{spaces-morphisms-lemma-scheme-theoretic-image-of-partial-section}.
D'après le Lemme \ref{spaces-more-morphisms-lemma-dominate-modification-by-blowup},
il existe un éclatement $U$-admissible $W' \to W$ et un
prolongement $W' \to X$ de $s$.
Nous achevons la démonstration en appliquant
Diviseurs sur les espaces, Lemme \ref{spaces-divisors-lemma-extend-admissible-blowups},
pour prolonger $W' \to W$ en un éclatement $U$-admissible de $Y$.
\end{proof}















\section{Lemme de Chow}
\label{spaces-more-morphisms-section-chow}

\noindent
Dans cette section, nous démontrons le lemme de Chow
(Lemme \ref{spaces-more-morphisms-lemma-chow-noetherian-separated}).
Nous encourageons le lecteur à consulter
Cohomologie des espaces, section \ref{spaces-cohomology-section-weak-chow},
où se trouve une version faible du lemme de Chow, facile à démontrer et suffisante
pour de nombreuses applications.

\medskip\noindent
Comme nous n'avons pas encore défini les morphismes projectifs d'espaces algébriques,
les énoncés des lemmes (voir par exemple le
Lemme \ref{spaces-more-morphisms-lemma-blowup-to-find-embedding}) feront intervenir
des morphismes propres représentables plutôt que des morphismes projectifs.

\begin{lemma}
\label{spaces-more-morphisms-lemma-find-common-blowups}
Soit $S$ un schéma. Soit $Y$ un espace algébrique quasi-compact et quasi-séparé
sur $S$. Soient $U \to X_1$ et $U \to X_2$ des immersions ouvertes
d'espaces algébriques sur $Y$, et supposons que $U$, $X_1$, $X_2$ soient de type
fini et séparés sur $Y$. Alors il existe un diagramme commutatif
$$
\xymatrix{
X_1' \ar[d] \ar[r] & X & X_2' \ar[l] \ar[d] \\
X_1 & U \ar[l] \ar[lu] \ar[u] \ar[ru] \ar[r] & X_2
}
$$
d'espaces algébriques sur $Y$, où $X_i' \to X_i$ est un éclatement
$U$-admissible, $X_i' \to X$ est une immersion ouverte, et $X$ est séparé et de
type fini sur $Y$.
\end{lemma}

\begin{proof}
Dans toute la démonstration, tous les espaces algébriques seront séparés et de
type fini sur $Y$. Cela implique en particulier que ces espaces algébriques sont
quasi-compacts et quasi-séparés et que les morphismes entre eux
seront quasi-compacts et séparés. Voir
Morphismes d'espaces, sections
\ref{spaces-morphisms-section-separation-axioms} et
\ref{spaces-morphisms-section-quasi-compact}.
Nous utiliserons que, si $U \to W$ est une immersion de tels espaces sur $Y$,
alors l'image schématique $Z$ de $U$ dans $W$ est un sous-espace fermé
de $W$ et $U \to Z$ est une immersion ouverte, $U \subset Z$ est
schématiquement dense et $|U| \subset |Z|$ est dense. Voir
Morphismes d'espaces, Lemme
\ref{spaces-morphisms-lemma-quasi-compact-immersion}.

\medskip\noindent
Soit $X_{12} \subset X_1 \times_Y X_2$ l'image schématique
de $U \to X_1 \times_Y X_2$. Les projections $p_i : X_{12} \to X_i$
induisent des isomorphismes $p_i^{-1}(U) \to U$ d'après
Morphismes d'espaces, Lemme
\ref{spaces-morphisms-lemma-scheme-theoretic-image-of-partial-section}.
Choisissons un éclatement $U$-admissible $X_i^i \to X_i$ tel que
le transformé strict $X_{12}^i$ de $X_{12}$ soit isomorphe à un
sous-espace ouvert de $X_i^i$, voir le
Lemme \ref{spaces-more-morphisms-lemma-zariski-after-blowup}.
Soit $\mathcal{I}_i \subset \mathcal{O}_{X_i}$ le faisceau quasi-cohérent
d'idéaux de type fini correspondant.
Rappelons que $X_{12}^i \to X_{12}$ est l'éclatement en
$p_i^{-1}\mathcal{I}_i \mathcal{O}_{X_{12}}$, voir
Diviseurs sur les espaces, Lemme \ref{spaces-divisors-lemma-strict-transform}.
Soit $X_{12}'$ l'éclatement de $X_{12}$ en
$p_1^{-1}\mathcal{I}_1 p_2^{-1}\mathcal{I}_2 \mathcal{O}_{X_{12}}$, voir
Diviseurs sur les espaces, Lemme \ref{spaces-divisors-lemma-blowing-up-two-ideals},
pour les conséquences de cette construction. Nous obtenons un diagramme commutatif
$$
\xymatrix{
X_{12}' \ar[d] \ar[r] & X_{12}^2 \ar[d] \\
X_{12}^1 \ar[r] & X_{12}
}
$$
où tous les morphismes sont des éclatements $U$-admissibles. Puisque
$X_{12}^i \subset X_i^i$ est un ouvert, nous pouvons choisir un éclatement $U$-admissible
$X_i' \to X_i^i$ dont la restriction est $X_{12}' \to X_{12}^i$, voir
Diviseurs sur les espaces, Lemme
\ref{spaces-divisors-lemma-extend-admissible-blowups}.
Alors $X_{12}' \subset X_i'$ est un sous-espace ouvert et le diagramme
$$
\xymatrix{
X_{12}' \ar[d] \ar[r] & X_i' \ar[d] \\
X_{12}^i \ar[r] & X_i^i
}
$$
est commutatif, ses flèches verticales étant des éclatements et ses flèches horizontales
des immersions ouvertes. Notons que $X'_{12} \to X_1' \times_Y X_2'$ est
une immersion propre (utiliser le fait que $X'_{12} \to X_{12}$ est propre,
que $X_{12} \to X_1 \times_Y X_2$ est fermé, que $X_1' \times_Y X_2' \to
X_1 \times_Y X_2$ est séparé, et appliquer Morphismes d'espaces, Lemme
\ref{spaces-morphisms-lemma-universally-closed-permanence}).
Ainsi $X'_{12} \to  X_1' \times_Y X_2'$ est une immersion fermée.
Si nous définissons $X$ en recollant $X_1'$ et $X_2'$ le long du sous-espace
ouvert commun $X_{12}'$, alors $X \to Y$ est de type fini et
séparé\footnote{En effet, nous pouvons vérifier que la diagonale
$X \to X \times_Y X$ est fermée sur les quatre parties ouvertes $X'_i \times_Y X'_j$
de $X \times_Y X$, où cela est clair.}. Comme les composées
d'éclatements $U$-admissibles sont des éclatements $U$-admissibles
(Diviseurs sur les espaces, Lemme
\ref{spaces-divisors-lemma-composition-admissible-blowups})
le lemme est démontré.
\end{proof}

\begin{lemma}
\label{spaces-more-morphisms-lemma-blowup-to-find-embedding}
Soit $S$ un schéma. Soit $f : X \to Y$ un morphisme d'espaces algébriques
sur $S$. Soit $U \subset X$ un sous-espace ouvert. Supposons que
\begin{enumerate}
\item $U$ soit quasi-compact,
\item $Y$ soit quasi-compact et quasi-séparé,
\item il existe une immersion $U \to \mathbf{P}^n_Y$ sur $Y$,
\item $f$ soit de type fini et séparé.
\end{enumerate}
Alors il existe un diagramme commutatif
$$
\xymatrix{
& U \ar[ld] \ar[d] \ar[rd] \ar[rrd] \\
X \ar[rd] & X' \ar[l] \ar[d] \ar[r] & Z' \ar[ld] \ar[r] & Z \ar[ld] \\
& Y & \mathbf{P}^n_Y \ar[l]
}
$$
où
les flèches de source $U$ sont des immersions ouvertes,
$X' \to X$ est un éclatement $U$-admissible,
$X' \to Z'$ est une immersion ouverte,
$Z' \to Y$ est un morphisme propre et représentable d'espaces algébriques.
Plus précisément, $Z' \to Z$ est un éclatement $U$-admissible
et $Z \to \mathbf{P}^n_Y$ est une immersion fermée.
\end{lemma}

\begin{proof}
Soit $Z \subset \mathbf{P}^n_Y$ l'image schématique de
l'immersion $U \to \mathbf{P}^n_Y$. Puisque $U \to \mathbf{P}^n_Y$
est quasi-compact, nous voyons que $U \subset Z$ est un sous-espace
ouvert (schématiquement) dense
(Morphismes d'espaces, Lemme
\ref{spaces-morphisms-lemma-quasi-compact-immersion}).
Appliquons le Lemme \ref{spaces-more-morphisms-lemma-find-common-blowups} pour trouver un diagramme
$$
\xymatrix{
X' \ar[d] \ar[r] & \overline{X}' & Z' \ar[l] \ar[d] \\
X & U \ar[l] \ar[lu] \ar[u] \ar[ru] \ar[r] & Z
}
$$
ayant les propriétés énumérées dans l'énoncé de ce lemme.
Comme $X' \to X$ et $Z' \to Z$ sont des éclatements $U$-admissibles,
nous trouvons que $U$ est un ouvert schématiquement dense de
$X'$ et de $Z'$ (voir Diviseurs sur les espaces, Lemmes
\ref{spaces-divisors-lemma-blowing-up-gives-effective-Cartier-divisor} et
\ref{spaces-divisors-lemma-complement-effective-Cartier-divisor}).
Puisque $Z' \to Z \to Y$ est propre, nous voyons que $Z' \subset \overline{X}'$
est un sous-espace fermé (voir Morphismes d'espaces, Lemme
\ref{spaces-morphisms-lemma-universally-closed-permanence}).
Il s'ensuit que $X' \subset Z'$ (au sens schématique) ; ainsi $X'$
est un sous-espace ouvert de $Z'$ (un petit détail est omis), ce qui démontre le lemme.
\end{proof}

\begin{lemma}
\label{spaces-more-morphisms-lemma-chow-noetherian}
Soit $S$ un schéma. Soit $f : X \to Y$ un morphisme d'espaces algébriques
sur $S$. Supposons que $f$ soit séparé et de type fini et que $Y$ soit noethérien.
Alors il existe un sous-espace ouvert dense $U \subset X$ et
un diagramme commutatif
$$
\xymatrix{
& U \ar[ld] \ar[d] \ar[rd] \ar[rrd] \\
X \ar[rd] & X' \ar[l] \ar[d] \ar[r] & Z' \ar[ld] \ar[r] & Z \ar[ld] \\
& Y & \mathbf{P}^n_Y \ar[l]
}
$$
où
les flèches de source $U$ sont des immersions ouvertes,
$X' \to X$ est un éclatement $U$-admissible,
$X' \to Z'$ est une immersion ouverte,
$Z' \to Y$ est un morphisme propre et représentable d'espaces algébriques.
Plus précisément, $Z' \to Z$ est un éclatement $U$-admissible
et $Z \to \mathbf{P}^n_Y$ est une immersion fermée.
\end{lemma}

\begin{proof}
D'après Limites d'espaces, Lemme
\ref{spaces-limits-lemma-embedding-into-affine-over-qs}
il existe un sous-espace ouvert dense $U \subset X$ et une immersion
$U \to \mathbf{A}^n_Y$ sur $Y$. En composant avec l'immersion ouverte
$\mathbf{A}^n_Y \to \mathbf{P}^n_Y$, nous obtenons la situation du
Lemme \ref{spaces-more-morphisms-lemma-blowup-to-find-embedding}, d'où le
résultat.
\end{proof}

\begin{remark}
\label{spaces-more-morphisms-remark-chow-Noetherian}
Dans les Lemmes \ref{spaces-more-morphisms-lemma-blowup-to-find-embedding} et \ref{spaces-more-morphisms-lemma-chow-noetherian},
le morphisme $g : Z' \to Y$ est une composée de morphismes projectifs.
On peut présumer (d'après l'analogue pour les espaces algébriques de
Morphismes, Lemme \ref{morphisms-lemma-ample-composition})
qu'il existe un faisceau inversible $g$-ample sur $Z'$.
Si nous en avons un jour besoin, nous l'énoncerons et le démontrerons ici.
\end{remark}

\noindent
Le résultat suivant est \cite[IV, théorème 3.1]{Kn}. Notons que l'immersion
$X' \to \mathbf{P}^n_Y$ est quasi-compacte et peut donc se factoriser en
$X' \to Z' \to \mathbf{P}^n_Y$, où le premier morphisme est une
immersion ouverte et le second une immersion fermée
(Morphismes d'espaces, Lemme
\ref{spaces-morphisms-lemma-quasi-compact-immersion}).

\begin{lemma}[Lemme de Chow]
\label{spaces-more-morphisms-lemma-chow-noetherian-separated}
\begin{reference}
\cite[IV, théorème 3.1]{Kn}
\end{reference}
Soit $S$ un schéma. Soit $f : X \to Y$ un morphisme d'espaces algébriques
sur $S$. Supposons que $f$ soit séparé et de type fini, et que $Y$ soit séparé et
noethérien. Alors il existe un diagramme commutatif
$$
\xymatrix{
X \ar[rd] & X' \ar[l] \ar[d] \ar[r] & \mathbf{P}^n_Y \ar[ld] \\
& Y
}
$$
où $X' \to X$ est un éclatement $U$-admissible pour un certain ouvert dense
$U \subset X$, et le morphisme $X' \to \mathbf{P}^n_Y$ est une immersion.
\end{lemma}

\begin{proof}
Dans ce premier paragraphe de la démonstration, nous ramenons le lemme au cas
où $Y$ est de type fini sur $\Spec(\mathbf{Z})$.
Nous pouvons remplacer, et remplaçons, le schéma de base $S$ par $\Spec(\mathbf{Z})$.
Nous pouvons écrire $Y = \lim Y_i$ comme limite projective d'espaces algébriques
séparés et de type fini sur $\Spec(\mathbf{Z})$, voir
Limites d'espaces, Proposition \ref{spaces-limits-proposition-approximate} et
Lemme \ref{spaces-limits-lemma-descend-separated}.
Pour tout $i$ assez grand, nous pouvons trouver un morphisme séparé et de type fini
$X_i \to Y_i$ tel que $X = Y \times_{Y_i} X_i$, voir
Limites d'espaces, Lemmes
\ref{spaces-limits-lemma-descend-finite-presentation} et
\ref{spaces-limits-lemma-descend-separated-morphism}.
Soient $\eta_1, \ldots, \eta_n$ les points génériques des composantes
irréductibles de $|X|$ ($X$ est noethérien en tant qu'espace algébrique séparé
et de type fini sur l'espace algébrique noethérien $Y$ ; par conséquent,
$|X|$ est un espace topologique noethérien).
D'après Limites d'espaces, Lemme \ref{spaces-limits-lemma-topology-limit},
nous trouvons que les images de $\eta_1, \ldots, \eta_n$ dans $|X_i|$
sont distinctes pour $i$ assez grand. Nous pouvons remplacer
$X_i$ par l'image schématique du morphisme (quasi-compact, et même affine)
$X \to X_i$.
Après ce remplacement, nous voyons que les images
de $\eta_1, \ldots, \eta_n$ dans $|X_i|$ sont les points génériques des
composantes irréductibles de $|X_i|$, voir
Morphismes d'espaces, Lemme
\ref{spaces-morphisms-lemma-quasi-compact-scheme-theoretic-image}.
Cela étant dit, supposons que nous puissions trouver un diagramme
$$
\xymatrix{
X_i \ar[rd] & X_i' \ar[l] \ar[d] \ar[r] & \mathbf{P}^n_{Y_i} \ar[ld] \\
& Y
}
$$
où $X_i' \to X_i$ est un éclatement $U_i$-admissible pour un certain ouvert dense
$U_i \subset X_i$, et où le morphisme $X_i' \to \mathbf{P}^n_{Y_i}$
est une immersion. Alors le transformé strict $X' \to X$ de $X$ relativement
à $X_i' \to X_i$ est un éclatement $U$-admissible, où $U \subset X$
est l'image inverse de $U_i$ dans $X$. Grâce à notre choix soigneux de
l'indice $i$, il s'ensuit que $\eta_1, \ldots, \eta_n \in |U|$ et que
$U \subset X$ est dense. De plus, $X' \to \mathbf{P}^n_Y$ est une
immersion, car $X'$ est fermé dans
$X_i' \times_{X_i} X = X_i' \times_{Y_i} Y$
qui est muni d'une immersion dans $\mathbf{P}^n_Y$. Nous sommes ainsi ramenés
à la situation du paragraphe suivant.

\medskip\noindent
Supposons que $Y$ soit séparé et de type fini sur $\Spec(\mathbf{Z})$.
Alors $X \to \Spec(\mathbf{Z})$ est lui aussi séparé et de type fini.
Nous appliquons le Lemme \ref{spaces-more-morphisms-lemma-chow-noetherian} à $X \to \Spec(\mathbf{Z})$
pour trouver un sous-espace ouvert dense $U \subset X$ et un diagramme commutatif
$$
\xymatrix{
& U \ar[ld] \ar[d] \ar[rd] \ar[rrd] \\
X \ar[rd] & X' \ar[l] \ar[d] \ar[r] & Z' \ar[ld] \ar[r] & Z \ar[ld] \\
& \Spec(\mathbf{Z}) & \mathbf{P}^n_\mathbf{Z} \ar[l]
}
$$
ayant toutes les propriétés énumérées dans le lemme.
Notons que $Z$ possède un faisceau inversible ample, à savoir
$\mathcal{O}_{\mathbf{P}^n}(1)|_Z$. Par conséquent, $Z' \to Z$
est un morphisme H-projectif d'après Morphismes, Lemme
\ref{morphisms-lemma-projective-over-quasi-projective-is-H-projective}.
Il s'ensuit que $Z' \to \Spec(\mathbf{Z})$ est H-projectif
d'après Morphismes, Lemme \ref{morphisms-lemma-H-projective-composition}.
Il existe donc une immersion fermée
$Z' \to \mathbf{P}^m_{\Spec(\mathbf{Z})}$ pour un certain $m \geq 0$.
Il s'ensuit que le morphisme diagonal
$$
X' \to Y \times \mathbf{P}^m_\mathbf{Z} = \mathbf{P}^m_Y
$$
est une immersion (car sa composée avec la projection
sur $\mathbf{P}^m_\mathbf{Z}$ est une immersion), et nous avons gagné.
\end{proof}






\section{Variantes du lemme de Chow}
\label{spaces-more-morphisms-section-chows-lemma}

\noindent
Dans cette section, nous démontrons plusieurs variantes du lemme de Chow qui portent
sur des morphismes entre espaces algébriques non noethériens.
Les versions noethériennes sont le
Lemme \ref{spaces-more-morphisms-lemma-chow-noetherian}
et le
Lemme \ref{spaces-more-morphisms-lemma-chow-noetherian-separated}.

\begin{lemma}
\label{spaces-more-morphisms-lemma-chow-finite-type}
Soit $S$ un schéma. Soit $Y$ un espace algébrique quasi-compact et quasi-séparé
sur $S$. Soit $f : X \to Y$ un morphisme séparé de
type fini. Alors il existe un diagramme commutatif
$$
\xymatrix{
X \ar[rd] & X' \ar[l] \ar[d] \ar[r] & \overline{X}' \ar[ld] \\
& Y
}
$$
où $X' \to X$ est propre surjectif,
$X' \to \overline{X}'$ est une immersion ouverte, et
$\overline{X}' \to Y$ est un morphisme propre et représentable
d'espaces algébriques.
\end{lemma}

\begin{proof}
D'après
Limites d'espaces, Proposition
\ref{spaces-limits-proposition-separated-closed-in-finite-presentation}
nous pouvons trouver une immersion fermée $X \to X_1$ où $X_1$ est séparé
et de présentation finie sur $Y$. Si nous démontrons l'assertion
pour $X_1 \to Y$, il est clair que le résultat en découle pour $X$. Nous pouvons donc supposer que
$X$ est de présentation finie sur $Y$.

\medskip\noindent
Nous pouvons remplacer, et remplaçons, le schéma de base $S$ par $\Spec(\mathbf{Z})$.
Écrivons $Y = \lim_i Y_i$ comme limite projective d'espaces algébriques
quasi-séparés et de type fini sur $\Spec(\mathbf{Z})$, voir
Limites d'espaces,
Proposition \ref{spaces-limits-proposition-approximate}.
D'après
Limites d'espaces,
Lemme \ref{spaces-limits-lemma-descend-finite-presentation},
nous pouvons
trouver un indice $i \in I$ et un schéma $X_i \to Y_i$ de présentation finie
tel que $X = Y \times_{Y_i} X_i$.
D'après
Limites d'espaces,
Lemme \ref{spaces-limits-lemma-descend-separated-morphism},
nous pouvons supposer que $X_i \to Y_i$ est séparé.
Si nous démontrons l'assertion pour
$X_i$ sur $Y_i$, il est clair que l'assertion vaut pour $X$. Le cas
$X_i \to Y_i$ est traité par le
Lemme \ref{spaces-more-morphisms-lemma-chow-noetherian}.
\end{proof}

\begin{lemma}
\label{spaces-more-morphisms-lemma-chow-finite-type-separated}
Soit $S$ un schéma. Soit $f : X \to Y$ un morphisme d'espaces algébriques
sur $S$. Supposons que $f$ soit séparé et de type fini, et que $Y$ soit séparé et
quasi-compact. Alors il existe un diagramme commutatif
$$
\xymatrix{
X \ar[rd] & X' \ar[l] \ar[d] \ar[r] & \mathbf{P}^n_Y \ar[ld] \\
& Y
}
$$
où $X' \to X$ est un morphisme propre surjectif
et le morphisme $X' \to \mathbf{P}^n_Y$ est une immersion.
\end{lemma}

\begin{proof}
D'après
Limites d'espaces, Proposition
\ref{spaces-limits-proposition-separated-closed-in-finite-presentation}
nous pouvons trouver une immersion fermée $X \to X_1$ où $X_1$ est séparé
et de présentation finie sur $Y$. Si nous démontrons l'assertion
pour $X_1 \to Y$, il est clair que le résultat en découle pour $X$. Nous pouvons donc supposer que
$X$ est de présentation finie sur $Y$.

\medskip\noindent
Nous pouvons remplacer, et remplaçons, le schéma de base $S$ par $\Spec(\mathbf{Z})$.
Écrivons $Y = \lim_i Y_i$ comme limite projective d'espaces algébriques
quasi-séparés et de type fini sur $\Spec(\mathbf{Z})$, voir
Limites d'espaces,
Proposition \ref{spaces-limits-proposition-approximate}.
D'après Limites d'espaces, Lemme \ref{spaces-limits-lemma-descend-separated},
nous pouvons supposer que $Y_i$ est séparé pour tout $i$.
D'après
Limites d'espaces,
Lemme \ref{spaces-limits-lemma-descend-finite-presentation},
nous pouvons
trouver un indice $i \in I$ et un schéma $X_i \to Y_i$ de présentation finie
tel que $X = Y \times_{Y_i} X_i$.
D'après
Limites d'espaces,
Lemme \ref{spaces-limits-lemma-descend-separated-morphism},
nous pouvons supposer que $X_i \to Y_i$ est séparé.
Si nous démontrons l'assertion pour
$X_i$ sur $Y_i$, il est clair que l'assertion vaut pour $X$. Le cas
$X_i \to Y_i$ est traité par le
Lemme \ref{spaces-more-morphisms-lemma-chow-noetherian-separated}.
\end{proof}






\section{Théorème d'existence de Grothendieck}
\label{spaces-more-morphisms-section-existence-proper}

\noindent
Dans cette section, nous étudions le théorème d'existence de Grothendieck pour les
espaces algébriques. Au lieu de développer une théorie des « espaces algébriques formels »,
nous introduisons temporairement un peu de langage qui remplace la notion de
« module cohérent sur un espace formel adique noethérien ».

\medskip\noindent
Soit $S$ un schéma. Soit $X$ un espace algébrique noethérien sur $S$.
Soit $\mathcal{I} \subset \mathcal{O}_X$ un faisceau quasi-cohérent d'idéaux.
Nous considérerons ci-dessous des systèmes projectifs $(\mathcal{F}_n)$ de
$\mathcal{O}_X$-modules cohérents tels que
\begin{enumerate}
\item $\mathcal{F}_n$ soit annulé par $\mathcal{I}^n$, et
\item les morphismes de transition induisent des isomorphismes
$\mathcal{F}_{n + 1}/\mathcal{I}^n\mathcal{F}_{n + 1} \to \mathcal{F}_n$.
\end{enumerate}
Un morphisme $\alpha : (\mathcal{F}_n) \to (\mathcal{G}_n)$
entre de tels systèmes projectifs est simplement un système compatible de morphismes
$\alpha_n : \mathcal{F}_n \to \mathcal{G}_n$.
Désignons la catégorie de ces systèmes projectifs par
$\textit{Coh}(X, \mathcal{I})$. Nous allons développer pour
ces systèmes une théorie parallèle aux
résultats correspondants dans le cas des schémas, voir
Cohomologie des schémas, sections \ref{coherent-section-existence},
\ref{coherent-section-existence-proper},
\ref{coherent-section-existence-proper-support} et
\ref{coherent-section-algebraization}.

\medskip\noindent
Fonctorialité. Soit $f : X \to Y$ un
morphisme d'espaces algébriques noethériens sur un schéma $S$, et soit
$\mathcal{J} \subset \mathcal{O}_Y$ un faisceau quasi-cohérent
d'idéaux. Posons $\mathcal{I} = f^{-1}\mathcal{J}\mathcal{O}_X$.
Dans cette situation, il existe un foncteur
$$
f^* : \textit{Coh}(Y, \mathcal{J}) \longrightarrow \textit{Coh}(X, \mathcal{I})
$$
qui envoie $(\mathcal{G}_n)$ sur $(f^*\mathcal{G}_n)$. Comparer avec
Cohomologie des schémas, Lemme \ref{coherent-lemma-inverse-systems-pullback}.
Si $f$ est étale, nous pouvons voir ce foncteur simplement comme la restriction
du système à $X$, voir Propriétés des espaces, 
Équation \ref{spaces-properties-equation-restrict-modules}.

\medskip\noindent
Descente étale. Soit $S$ un schéma. Soit $U_0 \to X$ un morphisme étale
surjectif d'espaces algébriques noethériens. Posons
$U_1 = U_0 \times_X U_0$ et $U_2 = U_0 \times_X U_0 \times_X U_0$.
Soit $\mathcal{I} \subset \mathcal{O}_{X}$ un faisceau quasi-cohérent
d'idéaux. Posons $\mathcal{I}_i = \mathcal{I}|_{U_i}$. Dans cette situation,
nous obtenons un diagramme de catégories
$$
\xymatrix{
\textit{Coh}(X, \mathcal{I}) \ar[r] &
\textit{Coh}(U_0, \mathcal{I}_0) \ar@<0.5ex>[r] \ar@<-0.5ex>[r] &
\textit{Coh}(U_1, \mathcal{I}_1) \ar@<1ex>[r] \ar[r] \ar@<-1ex>[r] &
\textit{Coh}(U_2, \mathcal{I}_2)
}
$$
et la première flèche présente $\textit{Coh}(X, \mathcal{I})$ comme la
limite homotopique de la partie droite du diagramme. Plus précisément, étant donnée
une {\it donnée de descente}, c'est-à-dire un couple $((\mathcal{G}_n), \varphi)$ où
$(\mathcal{G}_n)$ est un objet de $\textit{Coh}(U_0, \mathcal{I}_0)$ et
$\varphi : \text{pr}_0^*(\mathcal{G}_n) \to \text{pr}_1^*(\mathcal{G}_n)$
est un isomorphisme dans $\textit{Coh}(U_1, \mathcal{I}_1)$
qui satisfait la condition de cocycle dans $\textit{Coh}(U_2, \mathcal{I}_2)$,
alors il existe un unique objet $(\mathcal{F}_n)$ de
$\textit{Coh}(X, \mathcal{I})$ dont la donnée de descente canonique associée
est isomorphe à $((\mathcal{G}_n), \varphi)$. Comparer avec
Descente sur les espaces, Définition
\ref{spaces-descent-definition-descent-datum-effective-quasi-coherent}.
La démonstration de cet énoncé résulte immédiatement de l'application de
Descente sur les espaces, Proposition
\ref{spaces-descent-proposition-fpqc-descent-quasi-coherent}
aux données de descente $(\mathcal{G}_n, \varphi_n)$ lorsque $n$ varie.

\begin{lemma}
\label{spaces-more-morphisms-lemma-inverse-systems-abelian}
Soit $S$ un schéma. Soit $X$ un espace algébrique noethérien sur $S$ et soit
$\mathcal{I} \subset \mathcal{O}_X$ un faisceau quasi-cohérent d'idéaux.
\begin{enumerate}
\item La catégorie $\textit{Coh}(X, \mathcal{I})$ est abélienne.
\item L'exactitude dans $\textit{Coh}(X, \mathcal{I})$
peut se vérifier localement pour la topologie étale.
\item Pour tout morphisme plat $f : X' \to X$ d'espaces algébriques noethériens,
le foncteur $f^* : \textit{Coh}(X, \mathcal{I}) \to
\textit{Coh}(X', f^{-1}\mathcal{I}\mathcal{O}_{X'})$ est exact.
\end{enumerate}
\end{lemma}

\begin{proof}
Preuve de (1). Choisissons un schéma affine $U_0$ et un morphisme étale surjectif
$U_0 \to X$. Posons $U_1 = U_0 \times_X U_0$ et
$U_2 = U_0 \times_X U_0 \times_X U_0$ comme dans l'étude de la
descente étale ci-dessus. Les catégories $\textit{Coh}(U_i, \mathcal{I}_i)$
sont abéliennes
(Cohomologie des schémas, Lemme \ref{coherent-lemma-inverse-systems-abelian})
et les foncteurs image inverse sont les foncteurs exacts
$\textit{Coh}(U_0, \mathcal{I}_0) \to \textit{Coh}(U_1, \mathcal{I}_1)$
et
$\textit{Coh}(U_1, \mathcal{I}_1) \to \textit{Coh}(U_2, \mathcal{I}_2)$
(Cohomologie des schémas, Lemme \ref{coherent-lemma-inverse-systems-pullback}).
Le lemme résulte alors formellement de la description de
$\textit{Coh}(X, \mathcal{I})$ comme catégorie de données de descente.
Nous omettons certains détails ; comparer avec la démonstration de
Groupoïdes, Lemme \ref{groupoids-lemma-abelian}.

\medskip\noindent
La partie (2) résulte immédiatement de l'étude du paragraphe précédent.
Dans la situation de (3), choisissons un diagramme commutatif
$$
\xymatrix{
U' \ar[d] \ar[r] & U \ar[d] \\
X' \ar[r] & X
}
$$
où $U'$ et $U$ sont des schémas affines et les morphismes verticaux sont
étales surjectifs. Alors $U' \to U$ est un morphisme plat de schémas
noethériens (Morphismes d'espaces, Lemme \ref{spaces-morphisms-lemma-flat-local}),
de sorte que le foncteur image inverse
$\textit{Coh}(U, \mathcal{I}\mathcal{O}_U) \to
\textit{Coh}(U', \mathcal{I}\mathcal{O}_{U'})$
est exact d'après
Cohomologie des schémas, Lemme \ref{coherent-lemma-inverse-systems-pullback}.
Puisque l'exactitude dans $\textit{Coh}(X, \mathcal{O}_X)$ se vérifie
sur $U$, et de même pour $X', U'$, l'assertion en résulte.
\end{proof}

\begin{lemma}
\label{spaces-more-morphisms-lemma-inverse-systems-surjective}
Soit $S$ un schéma. Soit $X$ un espace algébrique noethérien sur $S$
et soit $\mathcal{I} \subset \mathcal{O}_X$ un faisceau quasi-cohérent
d'idéaux. Un morphisme $(\mathcal{F}_n) \to (\mathcal{G}_n)$ est surjectif dans
$\textit{Coh}(X, \mathcal{I})$ si et seulement si
$\mathcal{F}_1 \to \mathcal{G}_1$ est surjectif.
\end{lemma}

\begin{proof}
On peut le vérifier sur un recouvrement étale affine de $X$ d'après le
Lemme \ref{spaces-more-morphisms-lemma-inverse-systems-abelian}.
On se ramène ainsi au cas des schémas, qui est le
Lemme \ref{coherent-lemma-inverse-systems-surjective} de Cohomologie des schémas.
\end{proof}

\noindent
Soit $S$ un schéma. Soit $X$ un espace algébrique noethérien sur $S$ et
soit $\mathcal{I} \subset \mathcal{O}_X$ un faisceau quasi-cohérent d'idéaux.
Il existe un foncteur
\begin{equation}
\label{spaces-more-morphisms-equation-completion-functor}
\textit{Coh}(\mathcal{O}_X) \longrightarrow \textit{Coh}(X, \mathcal{I}), \quad
\mathcal{F} \longmapsto \mathcal{F}^\wedge
\end{equation}
qui associe au $\mathcal{O}_X$-module cohérent $\mathcal{F}$
l'objet $\mathcal{F}^\wedge = (\mathcal{F}/\mathcal{I}^n\mathcal{F})$
de $\textit{Coh}(X, \mathcal{I})$.

\begin{lemma}
\label{spaces-more-morphisms-lemma-exact}
Le foncteur (\ref{spaces-more-morphisms-equation-completion-functor}) est exact.
\end{lemma}

\begin{proof}
Il suffit de le vérifier localement pour la topologie étale sur $X$, voir le
Lemme \ref{spaces-more-morphisms-lemma-inverse-systems-abelian}.
On se ramène ainsi au cas des schémas, qui est le
Lemme \ref{coherent-lemma-exact} de Cohomologie des schémas.
\end{proof}

\begin{lemma}
\label{spaces-more-morphisms-lemma-completion-internal-hom}
Soit $S$ un schéma. Soit $X$ un espace algébrique noethérien sur $S$ et
soit $\mathcal{I} \subset \mathcal{O}_X$ un faisceau quasi-cohérent d'idéaux.
Soient $\mathcal{F}$ et $\mathcal{G}$ des $\mathcal{O}_X$-modules cohérents. Posons
$\mathcal{H} = \SheafHom_{\mathcal{O}_X}(\mathcal{F}, \mathcal{G})$.
Alors
$$
\lim H^0(X, \mathcal{H}/\mathcal{I}^n\mathcal{H}) =
\Mor_{\textit{Coh}(X, \mathcal{I})}
(\mathcal{F}^\wedge, \mathcal{G}^\wedge).
$$
\end{lemma}

\begin{proof}
Puisque $\mathcal{H}$ est un faisceau sur $X_\etale$ et que l'on dispose de la
descente étale pour les objets de $\textit{Coh}(X, \mathcal{I})$,
il suffit de le démontrer localement pour la topologie étale.
On se ramène ainsi au cas des schémas, qui est le Lemme
\ref{coherent-lemma-completion-internal-hom} de Cohomologie des schémas.
\end{proof}

\noindent
Introduisons la situation qui nous occupera dans toute la
suite de cette section.

\begin{situation}
\label{spaces-more-morphisms-situation-existence}
Ici, $A$ est un anneau noethérien complet pour la topologie $I$-adique.
De plus, $f : X \to \Spec(A)$ est un morphisme séparé de type fini
d'espaces algébriques et $\mathcal{I} = I\mathcal{O}_X$.
\end{situation}

\noindent
Dans cette situation, nous notons
$$
\textit{Coh}_{\text{support propre sur } A}(\mathcal{O}_X)
$$
la sous-catégorie pleine de $\textit{Coh}(\mathcal{O}_X)$
formée des $\mathcal{O}_X$-modules cohérents dont le
support est propre sur $\Spec(A)$ ou, de manière équivalente, dont le
support schématique est propre sur $\Spec(A)$, voir
Catégories dérivées des espaces, Lemme
\ref{spaces-perfect-lemma-module-support-proper-over-base}.
De même, nous notons
$$
\textit{Coh}_{\text{support propre sur } A}(X, \mathcal{I})
$$
la sous-catégorie pleine de $\textit{Coh}(X, \mathcal{I})$
formée des objets $(\mathcal{F}_n)$ tels que
le support de $\mathcal{F}_1$ soit propre sur $\Spec(A)$.
Puisque le support d'un module quotient est contenu dans le support
du module, il s'ensuit que (\ref{spaces-more-morphisms-equation-completion-functor})
induit un foncteur
\begin{equation}
\label{spaces-more-morphisms-equation-completion-functor-proper-over-A}
\textit{Coh}_{\text{support propre sur }A}(\mathcal{O}_X)
\longrightarrow
\textit{Coh}_{\text{support propre sur }A}(X, \mathcal{I})
\end{equation}
Notre premier résultat affirme que ce foncteur est pleinement fidèle.

\begin{lemma}
\label{spaces-more-morphisms-lemma-fully-faithful}
Dans la Situation \ref{spaces-more-morphisms-situation-existence}.
Soient $\mathcal{F}$ et $\mathcal{G}$ des $\mathcal{O}_X$-modules cohérents.
Supposons que l'intersection des supports de
$\mathcal{F}$ et $\mathcal{G}$ soit propre sur $\Spec(A)$. Alors l'application
$$
\Mor_{\textit{Coh}(\mathcal{O}_X)}(\mathcal{F}, \mathcal{G})
\longrightarrow
\Mor_{\textit{Coh}(X, \mathcal{I})}
(\mathcal{F}^\wedge, \mathcal{G}^\wedge)
$$
provenant de (\ref{spaces-more-morphisms-equation-completion-functor}) est bijective.
En particulier, (\ref{spaces-more-morphisms-equation-completion-functor-proper-over-A})
est pleinement fidèle.
\end{lemma}

\begin{proof}
Posons $\mathcal{H} = \SheafHom_{\mathcal{O}_X}(\mathcal{G}, \mathcal{F})$.
C'est un $\mathcal{O}_X$-module cohérent, car sa restriction
aux schémas étales sur $X$ est cohérente d'après
Modules, Lemme \ref{modules-lemma-internal-hom-locally-kernel-direct-sum}.
D'après le Lemme \ref{spaces-more-morphisms-lemma-completion-internal-hom}, l'application
$$
\lim_n H^0(X, \mathcal{H}/\mathcal{I}^n\mathcal{H})
\to
\Mor_{\textit{Coh}(X, \mathcal{I})}
(\mathcal{G}^\wedge, \mathcal{F}^\wedge)
$$
est bijective. Soit $i : Z \to X$ le support schématique de
$\mathcal{H}$. Il est clair que $Z$ est un sous-espace fermé tel
que $|Z|$ soit contenu dans l'intersection des supports de $\mathcal{F}$
et $\mathcal{G}$. Par hypothèse, $Z \to \Spec(A)$ est donc propre
(voir Catégories dérivées des espaces, section
\ref{spaces-perfect-section-proper-over-base}).
Écrivons $\mathcal{H} = i_*\mathcal{H}'$ pour un certain
$\mathcal{O}_Z$-module cohérent $\mathcal{H}'$. Nous avons
$i_*(\mathcal{H}'/I^n\mathcal{H}') = \mathcal{H}/I^n\mathcal{H}$.
Nous obtenons donc
\begin{align*}
\lim_n H^0(X, \mathcal{H}/\mathcal{I}^n\mathcal{H})
& =
\lim_n H^0(Z, \mathcal{H}'/\mathcal{I}^n\mathcal{H}') \\
& =
H^0(Z, \mathcal{H}') \\
& =
H^0(X, \mathcal{H}) \\
& = 
\Mor_{\textit{Coh}(\mathcal{O}_X)}(\mathcal{F}, \mathcal{G})
\end{align*}
la deuxième égalité résultant du théorème des fonctions formelles
(Cohomologie des espaces, Lemme
\ref{spaces-cohomology-lemma-spell-out-theorem-formal-functions}).
Ceci démontre le lemme.
\end{proof}

\begin{remark}
\label{spaces-more-morphisms-remark-inverse-systems-kernel-cokernel-annihilated-by}
Soit $S$ un schéma. Soit $X$ un espace algébrique noethérien sur $S$ et soient
$\mathcal{I}, \mathcal{K} \subset \mathcal{O}_X$ des faisceaux quasi-cohérents
d'idéaux. Soit $\alpha : (\mathcal{F}_n) \to (\mathcal{G}_n)$ un morphisme de
$\textit{Coh}(X, \mathcal{I})$.
Étant donnés un schéma affine $U = \Spec(A)$ et un morphisme étale surjectif
$U \to X$, notons $I, K \subset A$ les idéaux correspondant aux restrictions
$\mathcal{I}|_U, \mathcal{K}|_U$. Notons $\alpha_U : M \to N$ le morphisme de
$A^\wedge$-modules de type fini qui correspond à $\alpha|_U$ par le
Lemme \ref{coherent-lemma-inverse-systems-affine} de Cohomologie des schémas.
Nous affirmons que les conditions suivantes sont équivalentes :
\begin{enumerate}
\item il existe un entier $t \geq 1$ tel que
$\Ker(\alpha_n)$ et $\Coker(\alpha_n)$
soient annulés par $\mathcal{K}^t$ pour tout $n \geq 1$,
\item pour tout (ou pour un) ouvert affine $\Spec(A) = U \subset X$ comme ci-dessus,
les modules $\Ker(\alpha_U)$ et $\Coker(\alpha_U)$
soient annulés par $K^t$ pour un certain entier $t \geq 1$.
\end{enumerate}
Si ces conditions équivalentes sont satisfaites, nous dirons que $\alpha$ est un
{\it morphisme dont le noyau et le conoyau sont annulés par une puissance de
$\mathcal{K}$}. Pour l'équivalence, nous renvoyons à
Cohomologie des schémas, Remarque
\ref{coherent-remark-inverse-systems-kernel-cokernel-annihilated-by}.
\end{remark}

\begin{lemma}
\label{spaces-more-morphisms-lemma-existence-easy}
Soit $S$ un schéma. Soit $X$ un espace algébrique noethérien sur $S$ et soit
$\mathcal{I} \subset \mathcal{O}_X$ un faisceau quasi-cohérent d'idéaux.
Soient $\mathcal{G}$ un $\mathcal{O}_X$-module cohérent, $(\mathcal{F}_n)$
un objet de $\textit{Coh}(X, \mathcal{I})$, et
$\alpha : (\mathcal{F}_n) \to \mathcal{G}^\wedge$
un morphisme dont le noyau et le conoyau sont annulés par une puissance de $\mathcal{I}$.
Alors il existe un unique triplet (à unique isomorphisme près)
$(\mathcal{F}, a, \beta)$ où
\begin{enumerate}
\item $\mathcal{F}$ est un $\mathcal{O}_X$-module cohérent,
\item $a : \mathcal{F} \to \mathcal{G}$ est un morphisme de $\mathcal{O}_X$-modules
dont le noyau et le conoyau sont annulés par une puissance de $\mathcal{I}$,
\item $\beta : (\mathcal{F}_n) \to \mathcal{F}^\wedge$ est un isomorphisme, et
\item $\alpha = a^\wedge \circ \beta$.
\end{enumerate}
\end{lemma}

\begin{proof}
L'unicité et la descente étale pour les objets de
$\textit{Coh}(X, \mathcal{I})$ et de $\textit{Coh}(\mathcal{O}_X)$
impliquent qu'il suffit de construire $(\mathcal{F}, a, \beta)$ localement
pour la topologie étale sur $X$. On se ramène ainsi au cas des schémas, qui est le
Lemme \ref{coherent-lemma-existence-easy} de Cohomologie des schémas.
\end{proof}

\begin{lemma}
\label{spaces-more-morphisms-lemma-existence-tricky}
Dans la Situation \ref{spaces-more-morphisms-situation-existence}. Soit $\mathcal{K} \subset \mathcal{O}_X$
un faisceau quasi-cohérent d'idéaux. Soit $X_e \subset X$ le sous-espace fermé
défini par $\mathcal{K}^e$. Soit $\mathcal{I}_e = \mathcal{I}\mathcal{O}_{X_e}$.
Soit $(\mathcal{F}_n)$ un objet de
$\textit{Coh}_{\text{support propre sur } A}(X, \mathcal{I})$.
Supposons que
\begin{enumerate}
\item le foncteur
$\textit{Coh}_{\text{support propre sur } A}(\mathcal{O}_{X_e})
\to \textit{Coh}_{\text{support propre sur } A}(X_e, \mathcal{I}_e)$
soit une équivalence pour tout $e \geq 1$, et
\item il existe un objet $\mathcal{H}$ de
$\textit{Coh}_{\text{support propre sur } A}(\mathcal{O}_X)$ et un morphisme
$\alpha : (\mathcal{F}_n) \to \mathcal{H}^\wedge$ dont le
noyau et le conoyau sont annulés par une puissance de $\mathcal{K}$.
\end{enumerate}
Alors $(\mathcal{F}_n)$ appartient à l'image essentielle de
(\ref{spaces-more-morphisms-equation-completion-functor-proper-over-A}).
\end{lemma}

\begin{proof}
Dans cette démonstration, nous utiliserons sans autre mention que, pour une immersion
fermée $i : Z \to X$, le foncteur $i_*$ donne une équivalence entre la
catégorie des modules cohérents sur $Z$ et celle des modules cohérents sur $X$ annulés
par le faisceau d'idéaux de $Z$, voir
Cohomologie des espaces, Lemme \ref{spaces-cohomology-lemma-i-star-equivalence}.
En particulier, nous considérons
$$
\textit{Coh}_{\text{support propre sur } A}(\mathcal{O}_{X_e})
\subset
\textit{Coh}_{\text{support propre sur } A}(\mathcal{O}_X)
$$
comme la sous-catégorie pleine formée des modules annulés par
$\mathcal{K}^e$, et
$$
\textit{Coh}_{\text{support propre sur } A}(X_e, \mathcal{I}_e)
\subset
\textit{Coh}_{\text{support propre sur } A}(X, \mathcal{I})
$$
comme la sous-catégorie pleine des objets annulés par $\mathcal{K}^e$.
De plus, (1) nous dit que ces deux catégories sont équivalentes par le
foncteur (\ref{spaces-more-morphisms-equation-completion-functor-proper-over-A}).

\medskip\noindent
En appliquant cette équivalence, nous obtenons un $\mathcal{O}_X$-module cohérent
$\mathcal{G}_e$ annulé par $\mathcal{K}^e$ et correspondant au système
$(\mathcal{F}_n/\mathcal{K}^e\mathcal{F}_n)$ de
$\textit{Coh}_{\text{support propre sur } A}(X, \mathcal{I})$. Les morphismes
$\mathcal{F}_n/\mathcal{K}^{e + 1}\mathcal{F}_n \to
\mathcal{F}_n/\mathcal{K}^e\mathcal{F}_n$ correspondent à des morphismes canoniques
$\mathcal{G}_{e + 1} \to \mathcal{G}_e$ qui induisent des isomorphismes
$\mathcal{G}_{e + 1}/\mathcal{K}^e\mathcal{G}_{e + 1} \to \mathcal{G}_e$.
Nous obtenons un objet $(\mathcal{G}_e)$ de la catégorie
$\textit{Coh}_{\text{support propre sur } A}(X, \mathcal{K})$.
Le morphisme $\alpha$ induit un système de morphismes
$$
\mathcal{F}_n/\mathcal{K}^e\mathcal{F}_n
\longrightarrow
\mathcal{H}/(\mathcal{I}^n + \mathcal{K}^e)\mathcal{H}
$$
d'où des morphismes $\mathcal{G}_e \to \mathcal{H}/\mathcal{K}^e\mathcal{H}$
(en utilisant encore l'équivalence de catégories).
Soit $t \geq 1$ un entier, qui existe par l'hypothèse (2),
tel que $\mathcal{K}^t$ annule le noyau et le conoyau de tous les morphismes
$\mathcal{F}_n \to \mathcal{H}/\mathcal{I}^n\mathcal{H}$.
Alors $\mathcal{K}^{2t}$ annule le noyau et le conoyau des morphismes
$\mathcal{F}_n/\mathcal{K}^e\mathcal{F}_n \to
\mathcal{H}/(\mathcal{I}^n + \mathcal{K}^e)\mathcal{H}$
(détails omis ; voir Cohomologie des schémas,
Remarque \ref{coherent-remark-inverse-systems-kernel-cokernel-annihilated-by}).
Nous en concluons que $\mathcal{K}^{4t}$ annule le noyau et
le conoyau des morphismes
$$
\mathcal{G}_e
\longrightarrow
\mathcal{H}/\mathcal{K}^e\mathcal{H},
$$
(détails omis ; voir Cohomologie des schémas,
Remarque \ref{coherent-remark-inverse-systems-kernel-cokernel-annihilated-by}).
Nous appliquons le Lemme \ref{spaces-more-morphisms-lemma-existence-easy} pour obtenir un
$\mathcal{O}_X$-module cohérent $\mathcal{F}$, un morphisme
$a : \mathcal{F} \to \mathcal{H}$ et un isomorphisme
$\beta : (\mathcal{G}_e) \to (\mathcal{F}/\mathcal{K}^e\mathcal{F})$
dans $\textit{Coh}(X, \mathcal{K})$. En remontant les constructions, pour $n$ fixé,
le triplet
$(\mathcal{F}/\mathcal{I}^n\mathcal{F}, a \bmod \mathcal{I}^n, \beta
\bmod \mathcal{I}^n)$ est un triplet comme dans le lemme pour le morphisme
$\alpha_n \bmod \mathcal{K}^e :
(\mathcal{F}_n/\mathcal{K}^e\mathcal{F}_n) \to
(\mathcal{H}/(\mathcal{I}^n + \mathcal{K}^e)\mathcal{H})$
de $\textit{Coh}(X, \mathcal{K})$. L'unicité du
Lemme \ref{spaces-more-morphisms-lemma-existence-easy}
donne donc un isomorphisme canonique
$\mathcal{F}/\mathcal{I}^n\mathcal{F} \to \mathcal{F}_n$
compatible avec tous les morphismes en présence.

\medskip\noindent
Pour achever la démonstration du lemme, il reste à montrer que le
support de $\mathcal{F}$ est propre sur $A$.
Par construction, le noyau de $a : \mathcal{F} \to \mathcal{H}$
est annulé par une puissance de $\mathcal{K}$. Par conséquent, le support de
ce noyau est contenu dans le support de $\mathcal{G}_1$. Puisque
$\mathcal{G}_1$ est un objet de
$\textit{Coh}_{\text{support propre sur } A}(\mathcal{O}_{X_1})$
nous voyons que celui-ci est propre sur $A$. Joint au fait que le
support de $\mathcal{H}$ est propre sur $A$, cela montre que le
support de $\mathcal{F}$ est propre sur $A$ d'après
Catégories dérivées des espaces, Lemme
\ref{spaces-perfect-lemma-union-closed-proper-over-base}.
\end{proof}

\begin{lemma}
\label{spaces-more-morphisms-lemma-inverse-systems-push-pull}
Soit $S$ un schéma. Soit $f : X \to Y$ un morphisme propre
représentable d'espaces algébriques noethériens sur $S$. Soient
$\mathcal{J}, \mathcal{K} \subset \mathcal{O}_Y$
des faisceaux quasi-cohérents d'idéaux.
Supposons que $f$ soit un isomorphisme au-dessus de $V = Y \setminus V(\mathcal{K})$.
Posons $\mathcal{I} = f^{-1}\mathcal{J} \mathcal{O}_X$.
Soit $(\mathcal{G}_n)$ un objet de $\textit{Coh}(Y, \mathcal{J})$,
soit $\mathcal{F}$ un $\mathcal{O}_X$-module cohérent, et soit
$\beta : (f^*\mathcal{G}_n)  \to \mathcal{F}^\wedge$ un isomorphisme dans
$\textit{Coh}(X, \mathcal{I})$. Alors il existe un morphisme
$$
\alpha :
(\mathcal{G}_n)
\longrightarrow
(f_*\mathcal{F})^\wedge
$$
dans $\textit{Coh}(Y, \mathcal{J})$ dont le noyau et le conoyau
sont annulés par une puissance de $\mathcal{K}$.
\end{lemma}

\begin{proof}
Puisque $f$ est un morphisme propre, $f_*\mathcal{F}$ est un
$\mathcal{O}_Y$-module cohérent (Cohomologie des espaces, Lemme
\ref{spaces-cohomology-lemma-proper-pushforward-coherent}).
L'énoncé du lemme a donc un sens. Considérons les composés
$$
\gamma_n : \mathcal{G}_n \to
f_*f^*\mathcal{G}_n \to
f_*(\mathcal{F}/\mathcal{I}^n\mathcal{F}).
$$
Ici, le premier morphisme est le morphisme d'adjonction et le second est $f_*\beta_n$.
Nous affirmons qu'il existe un unique $\alpha$ comme dans le lemme
tel que les composés
$$
\mathcal{G}_n \xrightarrow{\alpha_n}
f_*\mathcal{F}/\mathcal{J}^nf_*\mathcal{F} \to
f_*(\mathcal{F}/\mathcal{I}^n\mathcal{F})
$$
soient égaux à $\gamma_n$ pour tout $n$. Par unicité et par descente étale
pour $\textit{Coh}(Y, \mathcal{J})$, il suffit de le démontrer
localement pour la topologie étale sur $Y$. Nous pouvons donc supposer que $Y$
est le spectre d'un anneau noethérien. Comme $f$ est représentable,
$X$ est lui aussi un schéma. On se ramène ainsi au cas des
schémas ; voir la démonstration de
Cohomologie des schémas, Lemme \ref{coherent-lemma-inverse-systems-push-pull}.
\end{proof}

\begin{theorem}[Théorème d'existence de Grothendieck]
\label{spaces-more-morphisms-theorem-grothendieck-existence}
Dans la Situation \ref{spaces-more-morphisms-situation-existence}, le foncteur
(\ref{spaces-more-morphisms-equation-completion-functor-proper-over-A})
est une équivalence.
\end{theorem}

\begin{proof}
Nous utiliserons, sans autre mention dans la démonstration du théorème,
l'équivalence de catégories de Cohomologie des espaces, Lemme
\ref{spaces-cohomology-lemma-i-star-equivalence}.
D'après le Lemme \ref{spaces-more-morphisms-lemma-fully-faithful}, le foncteur est pleinement fidèle.
Il reste donc à démontrer que le foncteur est essentiellement surjectif.

\medskip\noindent
Considérons l'ensemble $\Xi$ des faisceaux quasi-cohérents d'idéaux
$\mathcal{K} \subset \mathcal{O}_X$ tels que l'énoncé soit vrai
pour tout objet $(\mathcal{F}_n)$ de
$\textit{Coh}_{\text{support propre sur }A}(X, \mathcal{I})$
annulé par $\mathcal{K}$. Nous voulons montrer que $(0)$ appartient à $\Xi$.
Dans le cas contraire, puisque $X$ est noethérien, il existe un
faisceau quasi-cohérent d'idéaux maximal $\mathcal{K}$ qui n'appartient pas à $\Xi$, voir
Cohomologie des espaces, Lemme \ref{spaces-cohomology-lemma-acc-coherent}.
Après avoir remplacé $X$ par le sous-schéma fermé de $X$
correspondant à $\mathcal{K}$, nous pouvons supposer que tout
$\mathcal{K}$ non nul appartient à $\Xi$. Soit $(\mathcal{F}_n)$ un objet de
$\textit{Coh}_{\text{support propre sur }A}(X, \mathcal{I})$.
Nous allons montrer que cet objet appartient à l'image essentielle, ce qui
achèvera la démonstration du théorème.

\medskip\noindent
Appliquons le lemme de Chow (Lemme \ref{spaces-more-morphisms-lemma-chow-noetherian-separated})
pour trouver un morphisme propre surjectif $f : Y \to X$ qui soit un isomorphisme
au-dessus d'un ouvert dense $U \subset X$, avec $Y$ H-quasi-projectif
sur $A$. Notons que $Y$ est un schéma et que $f$ est représentable.
Choisissons une immersion ouverte $j : Y \to Y'$ avec $Y'$ projectif sur $A$, voir
Morphismes, Lemme \ref{morphisms-lemma-H-quasi-projective-open-H-projective}.
Soit $T_n$ le support schématique de $\mathcal{F}_n$.
Notons que $|T_n| = |T_1|$ ; ainsi, $T_n$ est propre sur $A$ pour tout $n$
(Morphismes d'espaces, Lemme
\ref{spaces-morphisms-lemma-image-proper-is-proper}).
Alors $f^*\mathcal{F}_n$ est supporté sur le sous-schéma fermé
$f^{-1}T_n$, qui est propre sur $A$ (d'après
Morphismes d'espaces, Lemme \ref{spaces-morphisms-lemma-composition-proper},
et la propreté de $f$).
En particulier, le composé $f^{-1}T_n \to Y \to Y'$ est fermé
(Morphismes, Lemme \ref{morphisms-lemma-image-proper-scheme-closed}).
Soit $T'_n \subset Y'$ le sous-schéma fermé correspondant ;
il est contenu dans le sous-schéma ouvert $Y$ et égal à $f^{-1}T_n$
comme sous-schéma fermé de $Y$. Soit $\mathcal{F}_n'$
le $\mathcal{O}_{Y'}$-module cohérent correspondant à
$f^*\mathcal{F}_n$, considéré comme module cohérent sur $Y'$ au moyen de
l'immersion fermée $f^{-1}T_n = T'_n \subset Y'$.
Alors $(\mathcal{F}_n')$
est un objet de $\textit{Coh}(Y', I\mathcal{O}_{Y'})$.
D'après le cas projectif du théorème d'existence de Grothendieck
(Cohomologie des schémas, Lemme \ref{coherent-lemma-existence-projective}),
il existe un $\mathcal{O}_{Y'}$-module cohérent
$\mathcal{F}'$ et un isomorphisme
$(\mathcal{F}')^\wedge \cong (\mathcal{F}'_n)$ dans
$\textit{Coh}(Y', I\mathcal{O}_{Y'})$.
Soit $Z' \subset Y'$ le support schématique de $\mathcal{F}'$.
Puisque $\mathcal{F}'/I\mathcal{F}' = \mathcal{F}'_1$, nous voyons
que $Z' \cap V(I\mathcal{O}_{Y'}) = T'_1$ ensemblistement.
Le morphisme structural $p' : Y' \to \Spec(A)$ est propre ; par conséquent,
$p'(Z' \cap (Y' \setminus Y))$ est fermé dans $\Spec(A)$.
S'il n'était pas vide, il contiendrait un point de $V(I)$,
car $I$ est contenu dans le radical de Jacobson de $A$
(Algèbre, Lemme \ref{algebra-lemma-radical-completion}).
Mais nous avons vu ci-dessus que $Z' \cap (p')^{-1}V(I) = T'_1 \subset Y$ ;
nous en concluons donc que $Z' \subset Y$. Ainsi, $\mathcal{F}'|_Y$
est supporté sur un sous-schéma fermé de $Y$ propre sur $A$.

\medskip\noindent
Soit $\mathcal{K}$ le faisceau quasi-cohérent d'idéaux qui définit
le complémentaire réduit $X \setminus U$.
D'après Cohomologie des espaces, Lemme
\ref{spaces-cohomology-lemma-proper-pushforward-coherent}
le $\mathcal{O}_X$-module $\mathcal{H} = f_*\mathcal{F}'$ est cohérent,
et, d'après le Lemme \ref{spaces-more-morphisms-lemma-inverse-systems-push-pull},
il existe un morphisme $\alpha : (\mathcal{F}_n) \to \mathcal{H}^\wedge$
dans la catégorie $\textit{Coh}_{\text{support propre sur } A}(X, \mathcal{I})$
dont le noyau et le conoyau sont annulés par une puissance de $\mathcal{K}$.
Soit $Z_0 \subset X$ le support schématique de $\mathcal{H}$.
Il est clair que $|Z_0| \subset f(|Z'|)$. Par conséquent,
$Z_0 \to \Spec(A)$ est propre
(Morphismes d'espaces, Lemme
\ref{spaces-morphisms-lemma-image-proper-is-proper}).
Ainsi, $\mathcal{H}$ est un objet de 
$\textit{Coh}_{\text{support propre sur } A}(\mathcal{O}_X)$.
Puisque chacun des faisceaux d'idéaux $\mathcal{K}^e$ est un élément de
$\Xi$, les hypothèses du Lemme \ref{spaces-more-morphisms-lemma-existence-tricky}
sont satisfaites, ce qui permet de conclure.
\end{proof}

\begin{remark}[Explicitation du théorème d'existence de Grothendieck]
\label{spaces-more-morphisms-remark-reformulate-existence-theorem}
Soit $A$ un anneau noethérien complet pour la topologie $I$-adique.
Écrivons $S = \Spec(A)$ et $S_n = \Spec(A/I^n)$.
Soit $X \to S$ un morphisme d'espaces algébriques séparé et
de type fini. Pour $n \geq 1$, posons $X_n = X \times_S S_n$.
Diagramme :
$$
\xymatrix{
X_1 \ar[r]_{i_1} \ar[d] & X_2 \ar[r]_{i_2} \ar[d] & X_3 \ar[r] \ar[d] &
\ldots & X \ar[d] \\
S_1 \ar[r] & S_2 \ar[r] & S_3 \ar[r] & \ldots & S
}
$$
Dans cette situation, nous considérons les systèmes $(\mathcal{F}_n, \varphi_n)$
où
\begin{enumerate}
\item $\mathcal{F}_n$ est un $\mathcal{O}_{X_n}$-module cohérent,
\item $\varphi_n : i_n^*\mathcal{F}_{n + 1} \to \mathcal{F}_n$
est un isomorphisme, et
\item $\text{Supp}(\mathcal{F}_1)$ est propre sur $S_1$.
\end{enumerate}
Le Théorème \ref{spaces-more-morphisms-theorem-grothendieck-existence} affirme que le
foncteur de passage au complété
$$
\begin{matrix}
\text{modules cohérents sur }\mathcal{O}_X\text{ : }\mathcal{F} \\
\text{à support propre sur }A
\end{matrix}
\quad
\longrightarrow
\quad
\begin{matrix}
\text{systèmes }(\mathcal{F}_n) \\
\text{comme ci-dessus}
\end{matrix}
$$
est une équivalence de catégories. Dans le cas particulier où $X$ est
propre sur $A$, on peut omettre les conditions sur les supports.
\end{remark}







\section{Théorème d'algébrisation de Grothendieck}
\label{spaces-more-morphisms-section-algebraization}

\noindent
Cette section est l'analogue de
Cohomologie des schémas, section \ref{coherent-section-algebraization}.
Toutefois, il y manque le résultat sur l'algébrisation des déformations
d'espaces algébriques propres munis de faisceaux inversibles amples, car un
espace algébrique propre muni d'un faisceau inversible ample est
déjà un schéma. Nous avons néanmoins un résultat d'algébrisation pour les
espaces algébriques propres de dimension relative $1$.
Notre premier résultat est une reformulation du théorème d'existence de Grothendieck
en termes de sous-schémas fermés et de morphismes finis.

\begin{lemma}
\label{spaces-more-morphisms-lemma-algebraize-formal-closed-subscheme}
Soit $A$ un anneau noethérien complet pour la topologie $I$-adique.
Écrivons $S = \Spec(A)$ et $S_n = \Spec(A/I^n)$.
Soit $X \to S$ un morphisme d'espaces algébriques séparé
et de type fini.
Pour $n \geq 1$, posons $X_n = X \times_S S_n$.
Supposons donné un diagramme commutatif
$$
\xymatrix{
Z_1 \ar[r] \ar[d] & Z_2 \ar[r] \ar[d] & Z_3 \ar[r] \ar[d] & \ldots \\
X_1 \ar[r]^{i_1} & X_2 \ar[r]^{i_2} & X_3 \ar[r] & \ldots
}
$$
d'espaces algébriques à carrés cartésiens. Supposons que
\begin{enumerate}
\item $Z_1 \to X_1$ soit une immersion fermée, et
\item $Z_1 \to S_1$ soit propre.
\end{enumerate}
Alors il existe une immersion fermée d'espaces algébriques $Z \to X$ telle que
$Z_n = Z \times_S S_n$ pour tout $n \geq 1$. De plus, $Z$ est propre sur $S$.
\end{lemma}

\begin{proof}
Notons $j_n : Z_n \to X_n$ les morphismes verticaux.
Comme les carrés de l'énoncé sont cartésiens,
le changement de base de $j_n$ à $X_1$ est $j_1$.
Ainsi, Limites d'espaces, Lemme
\ref{spaces-limits-lemma-check-closed-infinitesimally}
montre que $j_n$ est une immersion fermée.
Posons $\mathcal{F}_n = j_{n, *}\mathcal{O}_{Z_n}$, de sorte que
$j_n^\sharp$ est une surjection $\mathcal{O}_{X_n} \to \mathcal{F}_n$.
En utilisant encore le caractère cartésien des carrés, nous voyons que
l'image inverse de $\mathcal{F}_{n + 1}$ sur $X_n$ est $\mathcal{F}_n$.
Le théorème d'existence de Grothendieck, sous la forme de la
Remarque \ref{spaces-more-morphisms-remark-reformulate-existence-theorem},
nous dit donc qu'il existe un morphisme
$\mathcal{O}_X \to \mathcal{F}$
de $\mathcal{O}_X$-modules cohérents dont la restriction à
$X_n$ redonne $\mathcal{O}_{X_n} \to \mathcal{F}_n$.
De plus, le support de $\mathcal{F}$ est propre sur $S$.
Comme le foncteur de passage au complété est exact (Lemme \ref{spaces-more-morphisms-lemma-exact}),
nous voyons que $\mathcal{O}_X \to \mathcal{F}$
est surjectif. Ainsi, $\mathcal{F} = \mathcal{O}_X/\mathcal{J}$
pour un certain faisceau quasi-cohérent d'idéaux $\mathcal{J}$.
Poser $Z = V(\mathcal{J})$ achève la démonstration.
\end{proof}

\begin{lemma}
\label{spaces-more-morphisms-lemma-algebraize-formal-algebraic-space-finite-over-proper}
Soit $A$ un anneau noethérien complet pour la topologie $I$-adique.
Écrivons $S = \Spec(A)$ et $S_n = \Spec(A/I^n)$.
Soit $X \to S$ un morphisme d'espaces algébriques séparé
et de type fini.
Pour $n \geq 1$, posons $X_n = X \times_S S_n$.
Supposons donné un diagramme commutatif
$$
\xymatrix{
Y_1 \ar[r] \ar[d] & Y_2 \ar[r] \ar[d] & Y_3 \ar[r] \ar[d] & \ldots \\
X_1 \ar[r]^{i_1} & X_2 \ar[r]^{i_2} & X_3 \ar[r] & \ldots
}
$$
d'espaces algébriques à carrés cartésiens. Supposons que
\begin{enumerate}
\item $Y_1 \to X_1$ soit un morphisme fini, et
\item $Y_1 \to S_1$ soit propre.
\end{enumerate}
Alors il existe un morphisme fini d'espaces algébriques $Y \to X$ tel que
$Y_n = Y \times_S S_n$ pour tout $n \geq 1$. De plus, $Y$ est propre sur $S$.
\end{lemma}

\begin{proof}
Notons $f_n : Y_n \to X_n$ les morphismes verticaux.
Comme les carrés de l'énoncé sont cartésiens,
le changement de base de $f_n$ à $X_1$ est $f_1$.
Ainsi, le Lemme \ref{spaces-more-morphisms-lemma-thicken-property-morphisms-cartesian}
montre que $f_n$ est un morphisme fini.
Posons $\mathcal{F}_n = f_{n, *}\mathcal{O}_{Y_n}$.
En utilisant le caractère cartésien des carrés, nous voyons que
l'image inverse de $\mathcal{F}_{n + 1}$ sur $X_n$ est $\mathcal{F}_n$.
Le théorème d'existence de Grothendieck, sous la forme de la
Remarque \ref{spaces-more-morphisms-remark-reformulate-existence-theorem},
nous dit donc qu'il existe un $\mathcal{O}_X$-module cohérent $\mathcal{F}$
dont la restriction à $X_n$ redonne $\mathcal{F}_n$.
De plus, le support de $\mathcal{F}$ est propre sur $S$.
Comme le foncteur de passage au complété est pleinement fidèle
(Théorème \ref{spaces-more-morphisms-theorem-grothendieck-existence}),
nous voyons que les applications de multiplication
$\mathcal{F}_n \otimes_{\mathcal{O}_{X_n}} \mathcal{F}_n \to
\mathcal{F}_n$ se recollent pour munir $\mathcal{F}$ d'une structure d'algèbre.
Poser $Y = \underline{\Spec}_X(\mathcal{F})$ achève la démonstration.
\end{proof}

\begin{lemma}
\label{spaces-more-morphisms-lemma-algebraize-morphism}
Soit $A$ un anneau noethérien complet pour la topologie $I$-adique.
Écrivons $S = \Spec(A)$ et $S_n = \Spec(A/I^n)$. Soient $X$, $Y$ des espaces
algébriques sur $S$. Pour $n \geq 1$, posons $X_n = X \times_S S_n$ et
$Y_n = Y \times_S S_n$. Supposons donné un système compatible de
diagrammes commutatifs
$$
\xymatrix{
& & X_{n + 1} \ar[rd] \ar[rr]_{g_{n + 1}} & & Y_{n + 1} \ar[ld] \\
X_n \ar[rru] \ar[rd] \ar[rr]_{g_n} & & Y_n \ar[rru] \ar[ld] & S_{n + 1} \\
& S_n \ar[rru]
}
$$
Supposons que
\begin{enumerate}
\item $X \to S$ soit propre, et
\item $Y \to S$ soit séparé et de type fini.
\end{enumerate}
Alors il existe un unique morphisme d'espaces algébriques $g : X \to Y$
sur $S$ tel que $g_n$ soit le changement de base de $g$ à $S_n$.
\end{lemma}

\begin{proof}
Les morphismes $(1, g_n) : X_n \to X_n \times_S Y_n$ sont des immersions fermées
parce que $Y_n \to S_n$ est séparé
(Morphismes d'espaces, Lemme \ref{spaces-morphisms-lemma-section-immersion}).
Ainsi, d'après le Lemme \ref{spaces-more-morphisms-lemma-algebraize-formal-closed-subscheme},
il existe un sous-espace fermé $Z \subset X \times_S Y$
propre sur $S$ dont le changement de base à $S_n$ redonne
$X_n \subset X_n \times_S Y_n$. La première projection $p : Z \to X$
est un morphisme propre (puisque $Z$ est propre sur $S$, voir
Morphismes d'espaces, Lemme
\ref{spaces-morphisms-lemma-universally-closed-permanence})
dont le changement de base à $S_n$ est un isomorphisme pour tout $n$.
En particulier, $p : Z \to X$ est quasi-fini sur un sous-espace ouvert
de $Z$ contenant tout point de $Z_0$, par exemple d'après
Morphismes d'espaces, Lemme
\ref{spaces-morphisms-lemma-locally-finite-type-quasi-finite-part}.
Comme $Z$ est propre sur $S$, ce voisinage ouvert est $Z$ tout entier.
Nous concluons que $p : Z \to X$ est fini par le théorème principal de Zariski
(on peut par exemple appliquer le
Lemme \ref{spaces-more-morphisms-lemma-quasi-finite-separated-pass-through-finite}
et utiliser la propreté de $Z$ sur $X$ pour voir que l'immersion est
une immersion fermée). En appliquant l'équivalence du
Théorème \ref{spaces-more-morphisms-theorem-grothendieck-existence},
nous voyons que $p_*\mathcal{O}_Z = \mathcal{O}_X$, puisque c'est vrai
modulo $I^n$ pour tout $n$. Ainsi, $p$ est un isomorphisme et nous obtenons
le morphisme $g$ comme composé $X \cong Z \to Y$.
Nous omettons la démonstration de l'unicité.
\end{proof}

\begin{remark}
\label{spaces-more-morphisms-remark-weaken-separation-axioms-question}
On peut se demander si, dans le théorème d'algébrisation de Grothendieck (sous la forme
du Lemme \ref{spaces-more-morphisms-lemma-algebraize-morphism}), on peut se contenter
d'axiomes de séparation plus faibles sur le but. Soyons plus précis.
Soient $A$, $I$, $S$, $S_n$, $X$, $Y$, $X_n$, $Y_n$ et $g_n$ comme dans
l'énoncé du Lemme \ref{spaces-more-morphisms-lemma-algebraize-morphism}, et supposons que
\begin{enumerate}
\item $X \to S$ soit propre, et
\item $Y \to S$ soit localement de type fini.
\end{enumerate}
Existe-t-il un morphisme d'espaces algébriques $g : X \to Y$
sur $S$ tel que $g_n$ soit le changement de base de $g$ à $S_n$ ?
Nous ne connaissons pas la réponse en général ; si vous la connaissez, veuillez écrire à
\href{mailto:stacks.project@gmail.com}{stacks.project@gmail.com}.
Si $Y \to S$ est séparé, le résultat
découle du lemme (on se ramène immédiatement au cas où
$X$ est de type fini sur $S$, en choisissant un ouvert quasi-compact
contenant l'image de $g_1$). Si l'on suppose seulement que $Y \to S$
est quasi-séparé, le résultat est encore vrai. D'abord, comme précédemment,
nous pouvons supposer $Y$ quasi-compact aussi bien que quasi-séparé.
On peut alors utiliser soit \cite{Bhatt-Algebraize}, soit
\cite{Hall-Rydh-coherent} pour algébriser $(g_n)$. Plus précisément, pour appliquer
la première référence, nous utilisons
$$
D_{perf}(X) \to \lim D_{perf}(X_n)
\xrightarrow{\lim Lg_n^*}
\lim D_{perf}(Y_n) = D_{perf}(Y)
$$
où la dernière étape utilise un résultat d'existence de Grothendieck pour
la catégorie dérivée de l'espace algébrique propre $Y$ sur $R$
(comparer avec
Platitude sur les espaces, Remarque \ref{spaces-flat-remark-correct-generality}).
L'article cité montre que
cette flèche détermine un morphisme $Y \to X$ comme souhaité. Pour appliquer la
seconde référence, nous utilisons le même argument avec les modules cohérents :
$$
\textit{Coh}(\mathcal{O}_X) \to \lim \textit{Coh}(\mathcal{O}_{X_n})
\xrightarrow{\lim g_n^*}
\lim \textit{Coh}(\mathcal{O}_{Y_n}) = \textit{Coh}(\mathcal{O}_Y)
$$
où l'égalité finale est une conséquence du théorème d'existence de Grothendieck
(Théorème \ref{spaces-more-morphisms-theorem-grothendieck-existence}).
La seconde référence nous dit que ce foncteur correspond à un
morphisme $Y \to X$ sur $R$. Si nous avons un jour besoin de cette généralisation,
nous énoncerons précisément et démontrerons soigneusement le résultat ici.
\end{remark}

\begin{lemma}
\label{spaces-more-morphisms-lemma-formal-algebraic-space-proper-reldim-1}
Soit $(A, \mathfrak m, \kappa)$ un anneau local noethérien complet.
Posons $S = \Spec(A)$ et $S_n = \Spec(A/\mathfrak m^n)$.
Considérons un diagramme commutatif
$$
\xymatrix{
X_1 \ar[r]_{i_1} \ar[d] & X_2 \ar[r]_{i_2} \ar[d] & X_3 \ar[r] \ar[d] &
\ldots \\
S_1 \ar[r] & S_2 \ar[r] & S_3 \ar[r] & \ldots
}
$$
d'espaces algébriques à carrés cartésiens. Si $\dim(X_1) \leq 1$,
alors il existe un morphisme projectif de schémas $X \to S$
et des isomorphismes $X_n \cong X \times_S S_n$ compatibles avec $i_n$.
\end{lemma}

\begin{proof}
D'après Espaces sur les corps, Lemme
\ref{spaces-over-fields-lemma-codim-1-point-in-schematic-locus}
l'espace algébrique $X_1$ est un schéma. Ainsi, $X_1$
est un schéma propre de dimension $\leq 1$ sur $\kappa$.
D'après Variétés, Lemme \ref{varieties-lemma-dim-1-proper-projective},
nous voyons que $X_1$ est H-projectif sur $\kappa$.
Soit $\mathcal{L}_1$ un faisceau inversible ample sur $X_1$.

\medskip\noindent
Nous allons montrer que $\mathcal{L}_1$ se relève en un système compatible
$\{\mathcal{L}_n\}$ de faisceaux inversibles sur $\{X_n\}$.
Remarquons que $X_n$ est aussi un schéma d'après le Lemme \ref{spaces-more-morphisms-lemma-thickening-scheme}.
Rappelons que $X_1 \to X_n$ induit des homéomorphismes entre les espaces
topologiques sous-jacents. Dans la suite de la démonstration, nous ne distinguons pas
les faisceaux sur $X_n$ des faisceaux sur $X_1$.
Supposons donné un relèvement $\mathcal{L}_n$ sur $X_n$. Considérons
la suite exacte
$$
1 \to
(1 + \mathfrak m^n\mathcal{O}_{X_{n + 1}})^* \to
\mathcal{O}_{X_{n + 1}}^* \to \mathcal{O}_{X_n}^* \to 1
$$
de faisceaux sur $X_{n + 1}$. La classe de $\mathcal{L}_n$ dans
$H^1(X_n, \mathcal{O}_{X_n}^*)$ (voir
Cohomologie, Lemme \ref{cohomology-lemma-h1-invertible})
se relève en un élément de $H^1(X_{n + 1}, \mathcal{O}_{X_{n + 1}}^*)$
si et seulement si l'obstruction dans
$H^2(X_{n + 1}, (1 + \mathfrak m^n\mathcal{O}_{X_{n + 1}})^*)$
est nulle. Comme $X_1$ est un schéma noethérien de dimension $\leq 1$,
ce groupe de cohomologie est nul (Cohomologie, Proposition
\ref{cohomology-proposition-vanishing-Noetherian}).

\medskip\noindent
D'après le théorème d'algébrisation de Grothendieck
(Cohomologie des schémas, Théorème \ref{coherent-theorem-algebraization}),
nous obtenons un morphisme projectif de schémas $X \to S = \Spec(A)$
et un système compatible d'isomorphismes $X_n = S_n \times_S X$.
\end{proof}

\begin{lemma}
\label{spaces-more-morphisms-lemma-projective-over-complete}
Soit $(A, \mathfrak m, \kappa)$ un anneau local noethérien complet.
Soit $X$ un espace algébrique sur $\Spec(A)$.
Si $X \to \Spec(A)$ est propre et $\dim(X_\kappa) \leq 1$, alors
$X$ est un schéma projectif sur $A$.
\end{lemma}

\begin{proof}
Posons $X_n = X \times_{\Spec(A)} \Spec(A/\mathfrak m^n)$.
D'après le Lemme \ref{spaces-more-morphisms-lemma-formal-algebraic-space-proper-reldim-1},
il existe un morphisme projectif $Y \to \Spec(A)$
et des isomorphismes compatibles
$Y \times_{\Spec(A)} \Spec(A/\mathfrak m^n) \cong
X \times_{\Spec(A)} \Spec(A/\mathfrak m^n)$.
D'après le Lemme \ref{spaces-more-morphisms-lemma-algebraize-morphism},
nous voyons que $X \cong Y$, ce qui achève la démonstration.
\end{proof}










\section{Immersions régulières}
\label{spaces-more-morphisms-section-regular-immersions}

\noindent
Cette section est l'analogue de
Diviseurs, section \ref{divisors-section-regular-immersions}
pour les morphismes d'espaces algébriques. Le lecteur est invité à lire d'abord
dans cette section ce qui concerne les immersions régulières de schémas.

\medskip\noindent
Dans
Diviseurs, section \ref{divisors-section-regular-immersions},
nous avons défini quatre types d'immersions régulières pour les morphismes de schémas.
Parmi eux, seuls trois sont (à notre connaissance) locaux sur le but pour
la topologie étale ; comme d'habitude, les immersions régulières au sens usuel ne le sont pas.
C'est pourquoi, pour les morphismes d'espaces algébriques, nous ne pouvons pas à proprement parler définir
les immersions régulières. (Ce genre de désagrément a conduit Grothendieck
et son école à remplacer la notion initiale d'immersion régulière par celle
d'immersion Koszul-régulière, voir
\cite[Exposé VII, définition 1.4]{SGA6}.)
Mais nous pouvons définir les immersions Koszul-régulières, $H_1$-régulières et quasi-régulières.
Remarquons aussi que, puisque les immersions Koszul-régulières ne sont pas préservées par
un changement de base arbitraire, nous ne pouvons pas employer la méthode de
Morphismes d'espaces, section \ref{spaces-morphisms-section-representable}
pour les définir. De même, comme les immersions Koszul-régulières ne sont pas locales
sur la source pour la topologie étale, nous ne pouvons pas non plus utiliser
Morphismes d'espaces, Lemme \ref{spaces-morphisms-lemma-local-source-target}
pour les définir. Nous remplaçons plutôt ce lemme par le
suivant.

\begin{lemma}
\label{spaces-more-morphisms-lemma-representable-etale-local-target}
Soit $\mathcal{P}$ une propriété des morphismes de schémas qui est locale
sur le but pour la topologie étale. Soit $S$ un schéma.
Soit $f : X \to Y$ un morphisme représentable d'espaces algébriques sur $S$.
Considérons les diagrammes commutatifs
$$
\xymatrix{
X \times_Y V \ar[d] \ar[r] & V \ar[d] \\
X \ar[r]^f & Y
}
$$
où $V$ est un schéma et $V \to Y$ est étale.
Les conditions suivantes sont équivalentes :
\begin{enumerate}
\item pour tout diagramme comme ci-dessus, la projection $X \times_Y V \to V$
possède la propriété $\mathcal{P}$, et
\item pour un diagramme comme ci-dessus où $V \to Y$ est surjectif,
la projection $X \times_Y V \to V$ possède la propriété $\mathcal{P}$.
\end{enumerate}
Si $X$ et $Y$ sont représentables, ces conditions équivalent aussi à ce que
$f$ (considéré comme morphisme de schémas) possède la propriété $\mathcal{P}$.
\end{lemma}

\begin{proof}
Démontrons l'équivalence de (1) et (2).
L'implication (1) $\Rightarrow$ (2) est immédiate.
Supposons que
$$
\xymatrix{
X \times_Y V \ar[d] \ar[r] & V \ar[d] \\
X \ar[r]^f & Y
}
\quad\quad
\xymatrix{
X \times_Y V' \ar[d] \ar[r] & V' \ar[d] \\
X \ar[r]^f & Y
}
$$
soient deux diagrammes comme dans le lemme. Supposons que $V \to Y$ soit
surjectif et que $X \times_Y V \to V$ possède la propriété $\mathcal{P}$.
Pour montrer que (2) implique (1), nous devons prouver que
$X \times_Y V' \to V'$ possède $\mathcal{P}$. À cette
fin, considérons le diagramme
$$
\xymatrix{
X \times_Y V \ar[d] &
(X \times_Y V) \times_X (X \times_Y V') \ar[l] \ar[d] \ar[r] &
X \times_Y V' \ar[d] \\
V &
V \times_Y V' \ar[l] \ar[r] &
V'
}
$$
D'après notre hypothèse que $\mathcal{P}$ est locale sur la source pour la topologie étale,
nous voyons que $\mathcal{P}$ est préservée par changement de base étale, voir
Descente, Lemme \ref{descent-lemma-pullback-property-local-target}.
Ainsi, si la flèche verticale de gauche possède $\mathcal{P}$, il en est de même
de la flèche verticale du milieu. Puisque $U \times_X U' \to U'$ est surjectif
et étale (et définit donc un recouvrement étale de $U'$),
il en résulte (car $\mathcal{P}$ est supposée locale sur le but pour la topologie
étale) que la flèche verticale de gauche possède $\mathcal{P}$.

\medskip\noindent
Si $X$ et $Y$ sont représentables, nous pouvons prendre
$\text{id}_Y : Y \to Y$ comme recouvrement étale pour voir que la
dernière assertion du lemme est vraie.
\end{proof}

\noindent
Remarquons qu'« être une immersion Koszul-régulière (resp.\ $H_1$-régulière,
resp.\ quasi-régulière) » est une propriété des morphismes de schémas qui est locale
sur le but pour la topologie fpqc, voir
Descente, Lemme \ref{descent-lemma-descending-property-regular-immersion}.
La définition suivante a donc un sens.

\begin{definition}
\label{spaces-more-morphisms-definition-regular-immersion}
Soit $S$ un schéma. Soit $i : X \to Y$ un morphisme d'espaces
algébriques sur $S$.
\begin{enumerate}
\item Nous disons que $i$ est une {\it immersion Koszul-régulière} si $i$ est représentable
et si les conditions équivalentes du
Lemme \ref{spaces-more-morphisms-lemma-representable-etale-local-target}
sont satisfaites avec $\mathcal{P}(f) =$ « $f$ est une immersion Koszul-régulière ».
\item Nous disons que $i$ est une {\it immersion $H_1$-régulière} si $i$ est représentable
et si les conditions équivalentes du
Lemme \ref{spaces-more-morphisms-lemma-representable-etale-local-target}
sont satisfaites avec $\mathcal{P}(f) =$ « $f$ est une immersion $H_1$-régulière ».
\item Nous disons que $i$ est une {\it immersion quasi-régulière} si $i$ est représentable
et si les conditions équivalentes du
Lemme \ref{spaces-more-morphisms-lemma-representable-etale-local-target}
sont satisfaites avec $\mathcal{P}(f) =$ « $f$ est une immersion quasi-régulière ».
\end{enumerate}
\end{definition}

\begin{lemma}
\label{spaces-more-morphisms-lemma-regular-quasi-regular-immersion}
Soit $S$ un schéma.
Soit $i : Z \to X$ une immersion d'espaces algébriques sur $S$.
Nous avons les implications suivantes :
$i$ est une immersion Koszul-régulière $\Rightarrow$
$i$ est une immersion $H_1$-régulière $\Rightarrow$
$i$ est une immersion quasi-régulière.
\end{lemma}

\begin{proof}
Par la définition, ce lemme se réduit immédiatement à
Diviseurs, Lemme \ref{divisors-lemma-regular-quasi-regular-immersion}.
\end{proof}

\begin{lemma}
\label{spaces-more-morphisms-lemma-regular-immersion-noetherian}
Soit $S$ un schéma.
Soit $i : Z \to X$ une immersion d'espaces algébriques sur $S$.
Supposons que $X$ soit localement noethérien. Alors
$i$ est une immersion Koszul-régulière $\Leftrightarrow$
$i$ est une immersion $H_1$-régulière $\Leftrightarrow$
$i$ est une immersion quasi-régulière.
\end{lemma}

\begin{proof}
Par la Définition \ref{spaces-more-morphisms-definition-regular-immersion}
(et la définition d'un espace algébrique localement noethérien
dans Propriétés des espaces, section
\ref{spaces-properties-section-types-properties})
ceci se ramène immédiatement au cas des schémas, qui est
Diviseurs, Lemme \ref{divisors-lemma-regular-immersion-noetherian}.
\end{proof}

\begin{lemma}
\label{spaces-more-morphisms-lemma-flat-base-change-regular-immersion}
\begin{slogan}
Les immersions régulières sont stables par changement de base plat.
\end{slogan}
Soit $S$ un schéma. Soit $i : Z \to X$ une immersion Koszul-régulière,
$H_1$-régulière ou quasi-régulière d'espaces algébriques sur $S$.
Soit $X' \to X$ un morphisme plat d'espaces algébriques sur $S$.
Alors le changement de base $i' : Z \times_X X' \to X'$ est une immersion
Koszul-régulière, $H_1$-régulière ou quasi-régulière.
\end{lemma}

\begin{proof}
Par la Définition \ref{spaces-more-morphisms-definition-regular-immersion}
(et la définition d'un morphisme plat d'espaces algébriques
dans Morphismes d'espaces, section
\ref{spaces-morphisms-section-flat})
ce lemme se réduit au cas des schémas, voir
Diviseurs, Lemme \ref{divisors-lemma-flat-base-change-regular-immersion}.
\end{proof}

\begin{lemma}
\label{spaces-more-morphisms-lemma-quasi-regular-immersion}
Soit $S$ un schéma. Soit $i : Z \to X$ une immersion d'espaces algébriques
sur $S$. Alors $i$ est une immersion quasi-régulière si et seulement si les
conditions suivantes sont satisfaites :
\begin{enumerate}
\item $i$ est localement de présentation finie,
\item le faisceau conormal $\mathcal{C}_{Z/X}$ est localement libre de type fini, et
\item l'application (\ref{spaces-more-morphisms-equation-conormal-algebra-quotient}) est un isomorphisme.
\end{enumerate}
\end{lemma}

\begin{proof}
Cela découle du cas des schémas
(Diviseurs, Lemme \ref{divisors-lemma-quasi-regular-immersion})
par localisation étale (utiliser la Définition \ref{spaces-more-morphisms-definition-regular-immersion}
et le
Lemme \ref{spaces-more-morphisms-lemma-etale-conormal-algebra}).
\end{proof}

\begin{lemma}
\label{spaces-more-morphisms-lemma-transitivity-conormal-quasi-regular}
Soit $S$ un schéma. Soient $Z \to Y \to X$ des immersions d'espaces algébriques
sur $S$. Supposons que $Z \to Y$ soit une immersion $H_1$-régulière. Alors la suite canonique
du Lemme \ref{spaces-more-morphisms-lemma-transitivity-conormal}
$$
0 \to i^*\mathcal{C}_{Y/X} \to
\mathcal{C}_{Z/X} \to
\mathcal{C}_{Z/Y} \to 0
$$
est exacte et localement scindée (pour la topologie étale).
\end{lemma}

\begin{proof}
Comme $\mathcal{C}_{Z/Y}$ est localement libre de type fini (voir le
Lemme \ref{spaces-more-morphisms-lemma-quasi-regular-immersion}
et le
Lemme \ref{spaces-more-morphisms-lemma-regular-quasi-regular-immersion}),
il suffit de prouver que la suite est exacte.
Il suffit de montrer que la première application est injective,
puisque la suite est déjà exacte à droite en général.
Après localisation étale sur $X$, cela se réduit au cas
des schémas, voir
Diviseurs, Lemme \ref{divisors-lemma-transitivity-conormal-quasi-regular}.
\end{proof}

\noindent
Un composé d'immersions quasi-régulières peut ne pas être une immersion quasi-régulière, voir
Algèbre, Remarque \ref{algebra-remark-join-quasi-regular-sequences}.
Les autres types d'immersions régulières sont préservés par composition.

\begin{lemma}
\label{spaces-more-morphisms-lemma-composition-regular-immersion}
Soit $S$ un schéma. Soient $i : Z \to Y$ et $j : Y \to X$ des immersions
d'espaces algébriques sur $S$.
\begin{enumerate}
\item Si $i$ et $j$ sont des immersions Koszul-régulières, il en est de même de $j \circ i$.
\item Si $i$ et $j$ sont des immersions $H_1$-régulières, il en est de même de $j \circ i$.
\item Si $i$ est une immersion $H_1$-régulière et $j$ une immersion
quasi-régulière, alors $j \circ i$ est une immersion quasi-régulière.
\end{enumerate}
\end{lemma}

\begin{proof}
Cela découle immédiatement du cas des schémas, voir
Diviseurs, Lemme \ref{divisors-lemma-composition-regular-immersion}.
\end{proof}

\begin{lemma}
\label{spaces-more-morphisms-lemma-permanence-regular-immersion}
Soit $S$ un schéma. Soient $i : Z \to Y$ et $j : Y \to X$ des immersions
d'espaces algébriques sur $S$. Supposons que $j$ soit localement de présentation finie
et que la suite
$$
0 \to i^*\mathcal{C}_{Y/X} \to
\mathcal{C}_{Z/X} \to
\mathcal{C}_{Z/Y} \to 0
$$
du Lemme \ref{spaces-more-morphisms-lemma-transitivity-conormal} soit exacte et localement scindée.
\begin{enumerate}
\item Si $j \circ i$ est une immersion quasi-régulière, il en est de même de $i$.
\item Si $j \circ i$ est une immersion $H_1$-régulière, il en est de même de $i$.
\item Si $j$ et $j \circ i$ sont toutes deux des immersions Koszul-régulières, il en est de même de $i$.
\end{enumerate}
\end{lemma}

\begin{proof}
Cela découle immédiatement du cas des schémas, voir
Diviseurs, Lemme \ref{divisors-lemma-permanence-regular-immersion}.
\end{proof}

\begin{lemma}
\label{spaces-more-morphisms-lemma-extra-permanence-regular-immersion-noetherian}
Soit $S$ un schéma. Soient $i : Z \to Y$ et $j : Y \to X$ des immersions
d'espaces algébriques sur $S$. Supposons que $X$ soit localement noethérien.
Les conditions suivantes sont équivalentes :
\begin{enumerate}
\item $i$ et $j$ sont des immersions Koszul-régulières,
\item $i$ et $j \circ i$ sont des immersions Koszul-régulières,
\item $j \circ i$ est une immersion Koszul-régulière et la suite conormale
$$
0 \to i^*\mathcal{C}_{Y/X} \to
\mathcal{C}_{Z/X} \to
\mathcal{C}_{Z/Y} \to 0
$$
est exacte et localement scindée.
\end{enumerate}
\end{lemma}

\begin{proof}
Cela découle immédiatement du cas des schémas, voir Diviseurs, Lemme
\ref{divisors-lemma-extra-permanence-regular-immersion-noetherian}.
\end{proof}














\section{Pseudo-cohérence relative}
\label{spaces-more-morphisms-section-relative-pseudo-coherence}

\noindent
Cette section est l'analogue de
Compléments sur les morphismes, section
\ref{more-morphisms-section-relative-pseudo-coherence}.
Toutefois, dans le traitement de cette matière pour les
espaces algébriques, nous avons décidé de travailler exclusivement
avec les objets de la catégorie dérivée dont les faisceaux de cohomologie
sont quasi-cohérents. Il y a deux raisons à cela :
(1) cela simplifie grandement l'exposé et
(2) nous n'avons actuellement aucun usage de la notion plus générale.

\begin{remark}
\label{spaces-more-morphisms-remark-match-relative-pseudo-coherence}
Soit $S$ un schéma. Soit $f : X \to Y$ un morphisme entre des espaces algébriques
représentables sur $S$, localement de type fini. Soit
$f_0 : X_0 \to Y_0$ un morphisme de schémas représentant $f$
(notation maladroite, mais temporaire). Alors $f_0$ est localement de type fini.
Si $E$ est un objet de $D_\QCoh(\mathcal{O}_X)$, alors $E$
est l'image inverse d'un unique objet $E_0$ de $D_\QCoh(\mathcal{O}_{X_0})$, voir
Catégories dérivées des espaces, Lemme
\ref{spaces-perfect-lemma-derived-quasi-coherent-small-etale-site}.
Dans cette situation, l'expression « $E$ est $m$-pseudo-cohérent relativement à $Y$ »
signifiera « $E_0$ est $m$-pseudo-cohérent relativement à $Y_0$ »,
au sens défini dans Compléments sur les morphismes, section
\ref{more-morphisms-section-relative-pseudo-coherence}.
\end{remark}

\begin{lemma}
\label{spaces-more-morphisms-lemma-qcoh-relative-pseudo-coherence-characterize}
Soit $S$ un schéma. Soit $f : X \to Y$ un morphisme d'espaces algébriques
sur $S$, localement de type fini. Soit $m \in \mathbf{Z}$.
Soit $E \in D_\QCoh(\mathcal{O}_X)$. Avec les conventions expliquées dans la
Remarque \ref{spaces-more-morphisms-remark-match-relative-pseudo-coherence},
les conditions suivantes sont équivalentes :
\begin{enumerate}
\item pour tout diagramme commutatif
$$
\xymatrix{
U \ar[d] \ar[r] & V \ar[d] \\
X \ar[r] & Y
}
$$
où $U$, $V$ sont des schémas et les flèches verticales sont étales, le complexe
$E|_U$ est $m$-pseudo-cohérent relativement à $V$,
\item pour un diagramme commutatif comme en (1) où $U \to X$ est
surjectif, le complexe $E|_U$ est $m$-pseudo-cohérent relativement à $V$,
\item pour tout diagramme commutatif comme en (1) où $U$ et $V$ sont
affines, le complexe $R\Gamma(U, E)$ de $\mathcal{O}_X(U)$-modules
est $m$-pseudo-cohérent relativement à $\mathcal{O}_Y(V)$.
\end{enumerate}
\end{lemma}

\begin{proof}
La partie (1) implique (3) d'après Compléments sur les morphismes, Lemme
\ref{more-morphisms-lemma-qcoh-relative-pseudo-coherence-characterize}.

\medskip\noindent
Supposons (3). Choisissons un diagramme commutatif comme en (1) où $U \to X$ est surjectif.
Choisissons un recouvrement par des ouverts affines $V = \bigcup V_j$ et des recouvrements par des ouverts affines
$(U \to V)^{-1}(V_j) = \bigcup U_{ij}$. D'après (3) et Compléments sur les morphismes, Lemme
\ref{more-morphisms-lemma-qcoh-relative-pseudo-coherence-characterize}
nous voyons que $E|_U$ est $m$-pseudo-cohérent relativement à $V$.
Ainsi (3) implique (2).

\medskip\noindent
Supposons (2). Choisissons un diagramme commutatif
$$
\xymatrix{
U \ar[d] \ar[r] & V \ar[d] \\
X \ar[r] & Y
}
$$
où $U$, $V$ sont des schémas, les flèches verticales sont étales, le
morphisme $U \to X$ est surjectif et $E|_U$ est $m$-pseudo-cohérent
relativement à $V$. Supposons ensuite donné un second diagramme commutatif
$$
\xymatrix{
U' \ar[d] \ar[r] & V' \ar[d] \\
X \ar[r] & Y
}
$$
à flèches verticales étales, où $U', V'$ sont des schémas. Nous voulons montrer
que $E|_{U'}$ est $m$-pseudo-cohérent relativement à $V'$.
Le morphisme $U'' = U \times_X U' \to U'$ est surjectif et étale
et $U'' \to V'$ se factorise par $V'' = V' \times_Y V$, qui
est étale sur $V'$. Il suffit donc de montrer que $E|_{U''}$
est $m$-pseudo-cohérent relativement à $V''$, voir
Compléments sur les morphismes, Lemmes
\ref{more-morphisms-lemma-relative-pseudo-coherent-descends-fppf} et
\ref{more-morphisms-lemma-relative-pseudo-coherent-post-compose}.
En utilisant une fois encore le second lemme, il suffit de montrer que
$E|_{U''}$ est $m$-pseudo-cohérent relativement à $V$.
C'est vrai d'après Compléments sur les morphismes, Lemme
\ref{more-morphisms-lemma-pull-relative-pseudo-coherent}
et le fait qu'un morphisme étale de schémas est pseudo-cohérent d'après
Compléments sur les morphismes, Lemme
\ref{more-morphisms-lemma-flat-finite-presentation-pseudo-coherent}.
\end{proof}

\begin{definition}
\label{spaces-more-morphisms-definition-relative-pseudo-coherence}
Soit $S$ un schéma. Soit $f : X \to Y$ un morphisme d'espaces
algébriques sur $S$, localement de type fini.
Soit $E$ un objet de $D_\QCoh(\mathcal{O}_X)$. Soit $\mathcal{F}$ un
$\mathcal{O}_X$-module quasi-cohérent. Fixons $m \in \mathbf{Z}$.
\begin{enumerate}
\item Nous disons que $E$ est {\it $m$-pseudo-cohérent relativement à $Y$}
si les conditions équivalentes du
Lemme \ref{spaces-more-morphisms-lemma-qcoh-relative-pseudo-coherence-characterize} sont satisfaites.
\item Nous disons que $E$ est {\it pseudo-cohérent relativement à $Y$}
si $E$ est $m$-pseudo-cohérent relativement à $Y$ pour tout $m \in \mathbf{Z}$.
\item Nous disons que $\mathcal{F}$ est {\it $m$-pseudo-cohérent relativement à $Y$} si
$\mathcal{F}$, considéré comme objet de $D_\QCoh(\mathcal{O}_X)$, est
$m$-pseudo-cohérent relativement à $Y$.
\item Nous disons que $\mathcal{F}$ est {\it pseudo-cohérent relativement à $Y$} si
$\mathcal{F}$, considéré comme objet de $D_\QCoh(\mathcal{O}_X)$, est
pseudo-cohérent relativement à $Y$.
\end{enumerate}
\end{definition}

\noindent
La plupart des propriétés des complexes pseudo-cohérents relativement à une base
découleront immédiatement des propriétés correspondantes dans le cas
des schémas. Nous ajouterons ici les lemmes pertinents au besoin.

\begin{lemma}
\label{spaces-more-morphisms-lemma-relative-pseudo-coherent-is-moot}
Soit $S$ un schéma. Soit $f : X \to Y$ un morphisme d'espaces
algébriques sur $S$. Soit $E$ un objet de $D_\QCoh(\mathcal{O}_X)$.
Si $f$ est plat et localement de présentation finie, alors
les conditions suivantes sont équivalentes :
\begin{enumerate}
\item $E$ est pseudo-cohérent relativement à $Y$, et
\item $E$ est pseudo-cohérent sur $X$.
\end{enumerate}
\end{lemma}

\begin{proof}
Par localisation étale et d'après les définitions, nous pouvons supposer
que $X$ et $Y$ sont des schémas. Pour le cas des schémas, cela découle de
Compléments sur les morphismes, Lemme
\ref{more-morphisms-lemma-check-relative-pseudo-coherence-on-charts}.
\end{proof}

















\section{Morphismes pseudo-cohérents}
\label{spaces-more-morphisms-section-pseudo-coherent}

\noindent
Cette section est l'analogue de
Compléments sur les morphismes, section \ref{more-morphisms-section-pseudo-coherent}
pour les morphismes de schémas. Le lecteur est invité à lire d'abord
dans cette section ce qui concerne les morphismes pseudo-cohérents de schémas.

\medskip\noindent
La propriété « pseudo-cohérent » des morphismes de schémas est
de nature locale pour la topologie étale sur la source et le but. Pour le voir, utiliser
Compléments sur les morphismes,
Lemmes \ref{more-morphisms-lemma-descending-property-pseudo-coherent} et
\ref{more-morphisms-lemma-pseudo-coherent-fppf-local-source}
et
Descente, Lemme \ref{descent-lemma-etale-local-source-target}.
D'après
Morphismes d'espaces,
Lemme \ref{spaces-morphisms-lemma-local-source-target},
nous pouvons définir la notion de morphisme pseudo-cohérent d'espaces algébriques comme
suit ; elle coïncide avec la notion déjà définie dans
Compléments sur les morphismes, section \ref{more-morphisms-section-pseudo-coherent}
lorsque les espaces algébriques considérés sont représentables.

\begin{definition}
\label{spaces-more-morphisms-definition-pseudo-coherent}
Soit $S$ un schéma.
Soit $f : X \to Y$ un morphisme d'espaces algébriques sur $S$.
\begin{enumerate}
\item Nous disons que $f$ est {\it pseudo-cohérent} si les conditions équivalentes de
Morphismes d'espaces, Lemme \ref{spaces-morphisms-lemma-local-source-target}
sont satisfaites avec $\mathcal{P} =$ « pseudo-cohérent ».
\item Soit $x \in |X|$. Nous disons que $f$ est {\it pseudo-cohérent en $x$} s'il
existe un voisinage ouvert $X' \subset X$ de $x$ tel
que $f|_{X'} : X' \to Y$ soit pseudo-cohérent.
\end{enumerate}
\end{definition}

\noindent
Attention : un changement de base d'un morphisme pseudo-cohérent n'est pas
pseudo-cohérent en général.

\begin{lemma}
\label{spaces-more-morphisms-lemma-flat-base-change-pseudo-coherent}
Un changement de base plat d'un morphisme pseudo-cohérent est pseudo-cohérent.
\end{lemma}

\begin{proof}
Démonstration omise. Indication : utiliser la version de ce lemme pour les schémas, voir
Compléments sur les morphismes,
Lemme \ref{more-morphisms-lemma-flat-base-change-pseudo-coherent}.
\end{proof}

\begin{lemma}
\label{spaces-more-morphisms-lemma-composition-pseudo-coherent}
Un composé de morphismes pseudo-cohérents est pseudo-cohérent.
\end{lemma}

\begin{proof}
Démonstration omise. Indication : utiliser la version de ce lemme pour les schémas, voir
Compléments sur les morphismes,
Lemme \ref{more-morphisms-lemma-composition-pseudo-coherent}.
\end{proof}

\begin{lemma}
\label{spaces-more-morphisms-lemma-pseudo-coherent-finite-presentation}
Un morphisme pseudo-cohérent est localement de présentation finie.
\end{lemma}

\begin{proof}
Cela résulte immédiatement des définitions.
\end{proof}

\begin{lemma}
\label{spaces-more-morphisms-lemma-flat-finite-presentation-pseudo-coherent}
Un morphisme plat qui est localement de présentation finie est pseudo-cohérent.
\end{lemma}

\begin{proof}
Démonstration omise. Indication : utiliser la version de ce lemme pour les schémas, voir
Compléments sur les morphismes,
Lemme \ref{more-morphisms-lemma-flat-finite-presentation-pseudo-coherent}.
\end{proof}

\begin{lemma}
\label{spaces-more-morphisms-lemma-permanence-pseudo-coherent}
Soit $f : X \to Y$ un morphisme d'espaces algébriques pseudo-cohérent
au-dessus d'un espace algébrique de base $B$. Alors $f$ est pseudo-cohérent.
\end{lemma}

\begin{proof}
Démonstration omise. Indication : utiliser la version de ce lemme pour les schémas, voir
Compléments sur les morphismes,
Lemme \ref{more-morphisms-lemma-permanence-pseudo-coherent}.
\end{proof}

\begin{lemma}
\label{spaces-more-morphisms-lemma-Noetherian-pseudo-coherent}
Soit $S$ un schéma. Soit $f : X \to Y$ un morphisme d'espaces algébriques
sur $S$. Si $Y$ est localement noethérien, alors $f$ est pseudo-cohérent si
et seulement si $f$ est localement de type fini.
\end{lemma}

\begin{proof}
Démonstration omise. Indication : utiliser la version de ce lemme pour les schémas, voir
Compléments sur les morphismes,
Lemme \ref{more-morphisms-lemma-Noetherian-pseudo-coherent}.
\end{proof}








\section{Morphismes parfaits}
\label{spaces-more-morphisms-section-perfect}

\noindent
Cette section est l'analogue de
Compléments sur les morphismes, section \ref{more-morphisms-section-perfect}
pour les morphismes de schémas. Le lecteur est invité à lire d'abord
dans cette section ce qui concerne les morphismes parfaits de schémas.

\medskip\noindent
La propriété « parfait » des morphismes de schémas est
de nature locale pour la topologie étale sur la source et le but. Pour le voir, utiliser
Compléments sur les morphismes,
Lemmes \ref{more-morphisms-lemma-descending-property-perfect} et
\ref{more-morphisms-lemma-perfect-fppf-local-source}
et
Descente, Lemme \ref{descent-lemma-etale-local-source-target}.
D'après
Morphismes d'espaces,
Lemme \ref{spaces-morphisms-lemma-local-source-target},
nous pouvons définir la notion de morphisme parfait d'espaces algébriques comme
suit ; elle coïncide avec la notion déjà définie dans
Compléments sur les morphismes, section \ref{more-morphisms-section-perfect}
lorsque les espaces algébriques considérés sont représentables.

\begin{definition}
\label{spaces-more-morphisms-definition-perfect}
Soit $S$ un schéma.
Soit $f : X \to Y$ un morphisme d'espaces algébriques sur $S$.
\begin{enumerate}
\item Nous disons que $f$ est {\it parfait} si les conditions équivalentes de
Morphismes d'espaces, Lemme \ref{spaces-morphisms-lemma-local-source-target}
sont satisfaites avec $\mathcal{P} =$ « parfait ».
\item Soit $x \in |X|$. Nous disons que $f$ est {\it parfait en $x$} s'il
existe un voisinage ouvert $X' \subset X$ de $x$ tel
que $f|_{X'} : X' \to Y$ soit parfait.
\end{enumerate}
\end{definition}

\noindent
Remarquons qu'un morphisme parfait est pseudo-cohérent, donc localement de
présentation finie. Attention : un changement de base d'un morphisme parfait n'est pas parfait
en général.

\begin{lemma}
\label{spaces-more-morphisms-lemma-flat-base-change-perfect}
Un changement de base plat d'un morphisme parfait est parfait.
\end{lemma}

\begin{proof}
Démonstration omise. Indication : utiliser la version de ce lemme pour les schémas, voir
Compléments sur les morphismes,
Lemme \ref{more-morphisms-lemma-flat-base-change-perfect}.
\end{proof}

\begin{lemma}
\label{spaces-more-morphisms-lemma-composition-perfect}
Un composé de morphismes parfaits est parfait.
\end{lemma}

\begin{proof}
Démonstration omise. Indication : utiliser la version de ce lemme pour les schémas, voir
Compléments sur les morphismes,
Lemme \ref{more-morphisms-lemma-composition-perfect}.
\end{proof}

\begin{lemma}
\label{spaces-more-morphisms-lemma-flat-finite-presentation-perfect}
Soit $S$ un schéma. Soit $f : X \to Y$ un morphisme d'espaces algébriques sur
$S$. Les conditions suivantes sont équivalentes :
\begin{enumerate}
\item $f$ est plat et parfait, et
\item $f$ est plat et localement de présentation finie.
\end{enumerate}
\end{lemma}

\begin{proof}
Démonstration omise. Indication : utiliser la version de ce lemme pour les schémas, voir
Compléments sur les morphismes,
Lemme \ref{more-morphisms-lemma-flat-finite-presentation-perfect}.
\end{proof}

\begin{lemma}
\label{spaces-more-morphisms-lemma-perfect-proper-perfect-direct-image}
Soit $S$ un schéma. Soit $Y$ un espace algébrique noethérien sur $S$.
Soit $f : X \to Y$ un morphisme propre parfait d'espaces algébriques.
Soit $E \in D(\mathcal{O}_X)$ un objet parfait. Alors
$Rf_*E$ est un objet parfait de $D(\mathcal{O}_Y)$.
\end{lemma}

\begin{proof}
Nous affirmons que Catégories dérivées des espaces, Lemme
\ref{spaces-perfect-lemma-perfect-direct-image} s'applique.
Les conditions (1) et (2) sont immédiates. La condition (3) est locale
sur $X$. Nous pouvons donc supposer $X$ et $Y$ affines et $E$
représenté par un complexe strictement parfait de $\mathcal{O}_X$-modules.
Il suffit donc de montrer que $\mathcal{O}_X$ est de dimension de Tor
finie en tant que faisceau de $f^{-1}\mathcal{O}_Y$-modules
sur le site étale. D'après Catégories dérivées des espaces, Lemme
\ref{spaces-perfect-lemma-tor-dimension-rel}, il suffit de
le vérifier sur le site de Zariski. Pour les morphismes de schémas de type fini,
cela équivaut à la perfection du morphisme d'après Compléments sur les morphismes,
Lemme \ref{more-morphisms-lemma-check-perfect-stalks}.
\end{proof}











\section{Morphismes localement d'intersection complète}
\label{spaces-more-morphisms-section-lci}

\noindent
Cette section est l'analogue de
Compléments sur les morphismes, section \ref{more-morphisms-section-lci},
pour les morphismes de schémas. Le lecteur est invité à lire d'abord
dans cette section ce qui concerne les morphismes de schémas localement d'intersection complète.

\medskip\noindent
La propriété « être un morphisme localement d'intersection complète » des morphismes
de schémas est de nature locale pour la topologie étale sur la source et le but. Pour le voir, utiliser
Compléments sur les morphismes,
Lemmes \ref{more-morphisms-lemma-descending-property-lci} et
\ref{more-morphisms-lemma-lci-syntomic-local-source}
et
Descente, Lemme \ref{descent-lemma-etale-local-source-target}.
D'après
Morphismes d'espaces,
Lemme \ref{spaces-morphisms-lemma-local-source-target},
nous pouvons définir la notion de morphisme localement d'intersection complète
d'espaces algébriques comme suit ; elle coïncide avec la notion déjà
définie dans
Compléments sur les morphismes, section \ref{more-morphisms-section-lci},
lorsque les espaces algébriques considérés sont représentables.

\begin{definition}
\label{spaces-more-morphisms-definition-lci}
Soit $S$ un schéma.
Soit $f : X \to Y$ un morphisme d'espaces algébriques sur $S$.
\begin{enumerate}
\item Nous disons que $f$ est un {\it morphisme de Koszul}, ou que $f$ est un
{\it morphisme localement d'intersection complète}, si les conditions équivalentes de
Morphismes d'espaces, Lemme \ref{spaces-morphisms-lemma-local-source-target},
sont satisfaites avec $\mathcal{P}(f) =$ « $f$ est un morphisme localement d'intersection complète ».
\item Soit $x \in |X|$. Nous disons que $f$ est {\it Koszul en $x$} s'il
existe un voisinage ouvert $X' \subset X$ de $x$ tel
que $f|_{X'} : X' \to Y$ soit un morphisme localement d'intersection complète.
\end{enumerate}
\end{definition}

\noindent
En un certain sens, la propriété qui définit un morphisme localement
d'intersection complète est le résultat du lemme suivant.

\begin{lemma}
\label{spaces-more-morphisms-lemma-lci}
Soit $S$ un schéma.
Soit $f : X \to Y$ un morphisme localement d'intersection complète
d'espaces algébriques sur $S$.
Soit $P$ un espace algébrique lisse sur $Y$.
Soit $U \to X$ un morphisme étale d'espaces algébriques
et soit $i : U \to P$ une immersion d'espaces algébriques sur $Y$.
Diagramme :
$$
\xymatrix{
X \ar[rd] & U \ar[l] \ar[d] \ar[r]_i & P \ar[ld] \\
& Y
}
$$
Alors $i$ est une immersion Koszul-régulière d'espaces algébriques.
\end{lemma}

\begin{proof}
Choisissons un schéma $V$ et un morphisme étale surjectif $V \to Y$.
Choisissons un schéma $W$ et un morphisme étale surjectif $W \to P \times_Y V$.
Posons $U' = U \times_P W$ ; c'est un schéma étale sur $U$.
Nous devons montrer que $U' \to W$ est une immersion Koszul-régulière de
schémas, voir
Définition \ref{spaces-more-morphisms-definition-regular-immersion}.
D'après la
Définition \ref{spaces-more-morphisms-definition-lci}
ci-dessus, le morphisme de schémas $U' \to V$ est localement d'intersection
complète. Le résultat découle donc de
Compléments sur les morphismes, Lemme \ref{more-morphisms-lemma-lci}.
\end{proof}

\noindent
Il paraît opportun de rassembler ici quelques propriétés communes
à tous les morphismes de Koszul.

\begin{lemma}
\label{spaces-more-morphisms-lemma-lci-properties}
Soit $S$ un schéma. Soit $f : X \to Y$ un morphisme localement d'intersection
complète d'espaces algébriques sur $S$. Alors
\begin{enumerate}
\item $f$ est localement de présentation finie,
\item $f$ est pseudo-cohérent, et
\item $f$ est parfait.
\end{enumerate}
\end{lemma}

\begin{proof}
Démonstration omise. Indication : utiliser la version de ce lemme pour les schémas, voir
Compléments sur les morphismes,
Lemme \ref{more-morphisms-lemma-lci-properties}.
\end{proof}

\noindent
Attention : un changement de base d'un morphisme de Koszul n'est pas Koszul en général.

\begin{lemma}
\label{spaces-more-morphisms-lemma-flat-base-change-lci}
Un changement de base plat d'un morphisme localement d'intersection complète est un
morphisme localement d'intersection complète.
\end{lemma}

\begin{proof}
Démonstration omise. Indication : utiliser la version de ce lemme pour les schémas, voir
Compléments sur les morphismes,
Lemme \ref{more-morphisms-lemma-flat-base-change-lci}.
\end{proof}

\begin{lemma}
\label{spaces-more-morphisms-lemma-composition-lci}
Un composé de morphismes localement d'intersection complète est un
morphisme localement d'intersection complète.
\end{lemma}

\begin{proof}
Démonstration omise. Indication : utiliser la version de ce lemme pour les schémas, voir
Compléments sur les morphismes,
Lemme \ref{more-morphisms-lemma-composition-lci}.
\end{proof}

\begin{lemma}
\label{spaces-more-morphisms-lemma-flat-lci}
\begin{slogan}
Syntomique équivaut à plat et localement d'intersection complète (pour les espaces algébriques).
\end{slogan}
Soit $S$ un schéma.
Soit $f : X \to Y$ un morphisme d'espaces algébriques sur $S$.
Les conditions suivantes sont équivalentes :
\begin{enumerate}
\item $f$ est plat et localement d'intersection complète, et
\item $f$ est syntomique.
\end{enumerate}
\end{lemma}

\begin{proof}
Démonstration omise. Indication : utiliser la version de ce lemme pour les schémas, voir
Compléments sur les morphismes,
Lemme \ref{more-morphisms-lemma-flat-lci}.
\end{proof}

\begin{lemma}
\label{spaces-more-morphisms-lemma-regular-immersion-lci}
Soit $S$ un schéma. Une immersion Koszul-régulière d'espaces algébriques
sur $S$ est un morphisme localement d'intersection complète.
\end{lemma}

\begin{proof}
Soit $i : X \to Y$ une immersion Koszul-régulière d'espaces algébriques
sur $S$. Par définition, il existe un morphisme étale surjectif
$V \to Y$, où $V$ est un schéma, tel que $X \times_Y V$ soit un schéma
et que le changement de base $X \times_Y V \to V$ soit une immersion Koszul-régulière de
schémas. D'après Compléments sur les morphismes, Lemme
\ref{more-morphisms-lemma-regular-immersion-lci},
$X \times_Y V \to V$ est un morphisme localement d'intersection complète.
La Définition \ref{spaces-more-morphisms-definition-lci} montre alors que $i$ est un
morphisme localement d'intersection complète d'espaces algébriques.
\end{proof}

\begin{lemma}
\label{spaces-more-morphisms-lemma-lci-permanence}
Soit $S$ un schéma. Soit
$$
\xymatrix{
X \ar[rr]_f \ar[rd] & & Y \ar[ld] \\
& Z
}
$$
un diagramme commutatif de morphismes d'espaces algébriques sur $S$.
Supposons que $Y \to Z$ soit lisse et que $X \to Z$ soit un
morphisme localement d'intersection complète.
Alors $f : X \to Y$ est un morphisme localement d'intersection complète.
\end{lemma}

\begin{proof}
Choisissons un schéma $W$ et un morphisme étale surjectif $W \to Z$.
Choisissons un schéma $V$ et un morphisme étale surjectif $V \to W \times_Z Y$.
Choisissons un schéma $U$ et un morphisme étale surjectif $U \to V \times_Y X$.
Alors $U \to W$ est un morphisme de schémas localement d'intersection complète et
$V \to W$ est un morphisme de schémas lisse. Le résultat pour les schémas
(Compléments sur les morphismes, Lemme \ref{more-morphisms-lemma-lci-permanence})
montre que $U \to V$ est un morphisme localement d'intersection complète.
Par définition, cela signifie que $f$ est un morphisme localement d'intersection complète.
\end{proof}

\begin{lemma}
\label{spaces-more-morphisms-lemma-descending-property-lci}
La propriété $\mathcal{P}(f) =$ « $f$ est un morphisme localement
d'intersection complète » est de nature locale fpqc sur la base.
\end{lemma}

\begin{proof}
Soit $S$ un schéma. Soit $f : X \to Y$ un
morphisme d'espaces algébriques sur $S$.
Soit $\{Y_i \to Y\}$ un recouvrement fpqc
(Topologies sur les espaces, Définition
\ref{spaces-topologies-definition-fpqc-covering}).
Soit $f_i : X_i \to Y_i$ le changement de base de $f$ par $Y_i \to Y$.
Si $f$ est un morphisme localement d'intersection complète,
alors chaque $f_i$ est un morphisme localement d'intersection complète
d'après le Lemme \ref{spaces-more-morphisms-lemma-flat-base-change-lci}.

\medskip\noindent
Réciproquement, supposons que chaque $f_i$ soit un morphisme localement d'intersection complète.
Nous pouvons remplacer le recouvrement par un raffinement (de nouveau parce que le
changement de base plat préserve la propriété d'être un
morphisme localement d'intersection complète). Nous pouvons donc supposer que
$Y_i$ est un schéma pour tout $i$, voir
Topologies sur les espaces, Lemme \ref{spaces-topologies-lemma-refine-fpqc-schemes}.
Choisissons un schéma $V$ et un morphisme étale surjectif $V \to Y$.
Choisissons un schéma $U$ et un morphisme étale surjectif
$U \to V \times_Y X$. Nous devons montrer que $U \to V$ est un
morphisme de schémas localement d'intersection complète.
D'après Topologies sur les espaces, Lemme
\ref{spaces-topologies-lemma-recognize-fpqc-covering}
la famille $\{Y_i \times_Y V \to V\}$ est un recouvrement fpqc
de schémas. Le cas des schémas
(Compléments sur les morphismes, Lemme \ref{more-morphisms-lemma-descending-property-lci})
montre qu'il suffit de prouver que le changement de base
$$
U \times_Y Y_i = U \times_V (V \times_Y Y_i) \longrightarrow V
$$
de $U \to V$ par $V \times_Y Y_i \to V$ est un
morphisme localement d'intersection complète. Nous pouvons l'écrire comme le
composé
$$
U \times_Y Y_i \longrightarrow
(V \times_Y X) \times_Y Y_i = V \times_Y X_i \longrightarrow
V \times_Y Y_i
$$
La première flèche est un morphisme étale de schémas (comme changement de base de
$U \to V \times_Y X$) et la seconde est un
morphisme de schémas localement d'intersection complète, comme changement de base plat de $f_i$.
Le résultat s'ensuit, car la propriété d'être un morphisme localement d'intersection complète
est de nature locale pour la topologie syntomique sur la source et les morphismes étales
sont syntomiques (Compléments sur les morphismes, Lemme
\ref{more-morphisms-lemma-lci-syntomic-local-source}
et Morphismes, Lemme \ref{morphisms-lemma-etale-syntomic}).
\end{proof}

\begin{lemma}
\label{spaces-more-morphisms-lemma-lci-syntomic-local-source}
La propriété $\mathcal{P}(f) =$ « $f$ est un morphisme localement
d'intersection complète » est de nature locale pour la topologie syntomique sur la source.
\end{lemma}

\begin{proof}
Cela résulte de
Descente sur les espaces, Lemme \ref{spaces-descent-lemma-transfer-from-schemes}, et de
Compléments sur les morphismes, Lemme \ref{more-morphisms-lemma-lci-syntomic-local-source}.
\end{proof}

\begin{lemma}
\label{spaces-more-morphisms-lemma-base-change-lci-fibres}
Soit $S$ un schéma. Considérons un diagramme commutatif
$$
\xymatrix{
X \ar[rr]_f \ar[rd]_p & & Y \ar[ld]^q \\
& Z
}
$$
d'espaces algébriques sur $S$. Supposons que $p$ et $q$
soient plats et localement de présentation finie.
Alors il existe un sous-espace ouvert $U(f) \subset X$
tel que $|U(f)| \subset |X|$ soit l'ensemble des points où $f$ est Koszul.
De plus, pour tout morphisme d'espaces algébriques $Z' \to Z$, si
$f' : X' \to Y'$ est le changement de base de $f$ par $Z' \to Z$, alors
$U(f')$ est l'image inverse de $U(f)$ par la projection $X' \to X$.
\end{lemma}

\begin{proof}
Ce lemme est l'analogue de
Compléments sur les morphismes, Lemme \ref{more-morphisms-lemma-base-change-lci-fibres},
et nous allons en fait l'en déduire. D'après la
Définition \ref{spaces-more-morphisms-definition-lci},
l'ensemble $\{x \in |X| : f \text{ est Koszul en }x\}$ est
ouvert dans $|X|$ ; d'après
Propriétés des espaces, Lemme \ref{spaces-properties-lemma-open-subspaces},
il correspond donc à un sous-espace ouvert $U(f)$ de $X$. Il suffit par conséquent de
prouver la dernière assertion.

\medskip\noindent
Choisissons un schéma $W$ et un morphisme étale surjectif $W \to Z$.
Choisissons un schéma $V$ et un morphisme étale surjectif $V \to W \times_Z Y$.
Choisissons un schéma $U$ et un morphisme étale surjectif $U \to V \times_Y X$.
Enfin, choisissons un schéma $W'$ et un morphisme étale surjectif
$W' \to W \times_Z Z'$.
Posons $V' = W' \times_W V$ et $U' = W' \times_W U$ ; nous obtenons ainsi des
morphismes étales surjectifs $V' \to Y'$ et $U' \to X'$.
Nous utiliserons sans autre mention le fait qu'un morphisme étale d'espaces algébriques
induit une application ouverte entre les espaces topologiques associés (voir
Propriétés des espaces, Lemme
\ref{spaces-properties-lemma-etale-open}).
Notons que, par définition, $U(f)$ est l'image dans $|X|$ de l'ensemble $T$
des points de $U$ où le morphisme de schémas $U \to V$ est Koszul.
De même, $U(f')$ est l'image dans $|X'|$ de l'ensemble $T'$ des points de
$U'$ où le morphisme de schémas $U' \to V'$ est Koszul. Or, par construction, le
diagramme
$$
\xymatrix{
U' \ar[r] \ar[d] & U \ar[d] \\
V' \ar[r] & V
}
$$
est cartésien (dans la catégorie des schémas). Le
Lemme \ref{more-morphisms-lemma-base-change-lci-fibres} de Compléments sur les morphismes
cité ci-dessus montre donc que $T'$ est l'image inverse de $T$. Puisque
$|U'| \to |X'|$ est surjective, cela démontre le lemme.
\end{proof}

\begin{lemma}
\label{spaces-more-morphisms-lemma-unramified-lci}
Soit $S$ un schéma. Soit $f : X \to Y$ un morphisme localement d'intersection
complète d'espaces algébriques sur $S$. Alors $f$ est non ramifié si et seulement
si $f$ est formellement non ramifié et, dans ce cas, le faisceau conormal
$\mathcal{C}_{X/Y}$ est localement libre de type fini sur $X$.
\end{lemma}

\begin{proof}
Cela découle du résultat correspondant pour les morphismes de schémas, voir
Compléments sur les morphismes, Lemme \ref{more-morphisms-lemma-unramified-lci},
par localisation étale, voir
Lemme \ref{spaces-more-morphisms-lemma-universal-thickening-localize}.
(Notons que, dans la situation de ce lemme, le morphisme $V \to U$
est non ramifié et localement d'intersection complète par définition.)
\end{proof}

\begin{lemma}
\label{spaces-more-morphisms-lemma-transitivity-conormal-lci}
Soit $S$ un schéma. Soient $Z \to Y \to X$ des morphismes formellement non ramifiés
d'espaces algébriques sur $S$. Supposons que $Z \to Y$ soit un morphisme localement
d'intersection complète. La suite exacte
$$
0 \to i^*\mathcal{C}_{Y/X} \to
\mathcal{C}_{Z/X} \to
\mathcal{C}_{Z/Y} \to 0
$$
du
Lemme \ref{spaces-more-morphisms-lemma-transitivity-conormal}
est exacte courte.
\end{lemma}

\begin{proof}
Choisissons un schéma $U$ et un morphisme étale surjectif $U \to X$.
Choisissons un schéma $V$ et un morphisme étale surjectif $V \to U \times_X Y$.
Choisissons un schéma $W$ et un morphisme étale surjectif $W \to V \times_Y Z$.
D'après le
Lemme \ref{spaces-more-morphisms-lemma-universal-thickening-localize},
les morphismes $W \to V$ et $V \to U$ sont formellement non ramifiés.
De plus, la suite
$i^*\mathcal{C}_{Y/X} \to \mathcal{C}_{Z/X} \to \mathcal{C}_{Z/Y} \to 0$
se restreint à la suite correspondante
$i^*\mathcal{C}_{V/U} \to \mathcal{C}_{W/U} \to \mathcal{C}_{W/V} \to 0$
pour $W \to V \to U$. Le résultat découle donc du cas des schémas
(Compléments sur les morphismes, Lemme \ref{more-morphisms-lemma-transitivity-conormal-lci}),
car, par définition, le morphisme $W \to V$ est localement d'intersection
complète.
\end{proof}









\section{Quand un morphisme est-il un isomorphisme ?}
\label{spaces-more-morphisms-section-when-isomorphism}

\noindent
Plus généralement, on peut demander :
« Quand un morphisme possède-t-il la propriété $\mathcal{P}$ ? »
Une question plus précise est la suivante. Supposons donné un diagramme commutatif
$$
\xymatrix{
X \ar[rr]_f \ar[rd]_p & & Y \ar[ld]^q \\
& Z
}
$$
d'espaces algébriques. Existe-t-il un monomorphisme d'espaces algébriques
$W \to Z$ vérifiant les deux propriétés suivantes :
\begin{enumerate}
\item le changement de base $f_W : X_W \to Y_W$ possède la propriété $\mathcal{P}$, et
\item tout morphisme $Z' \to Z$ d'espaces algébriques se factorise par $W$ si
et seulement si le changement de base $f_{Z'} : X_{Z'} \to Y_{Z'}$ possède la propriété
$\mathcal{P}$.
\end{enumerate}
Dans de nombreux cas, si $W \to Z$ existe, c'est une immersion, une immersion ouverte
ou une immersion fermée.

\medskip\noindent
La réponse à cette question peut dépendre de propriétés auxiliaires des
morphismes $f$, $p$ et $q$. Un exemple est $\mathcal{P}(f) =$ « $f$ est plat »,
dont nous avons traité en grand détail, pour les morphismes de schémas, dans le cas $Y = S$ au
chapitre « Compléments sur la platitude », à partir de
Compléments sur la platitude, section \ref{flat-section-flattening-functors}.

\begin{lemma}
\label{spaces-more-morphisms-lemma-where-unramified}
Considérons un diagramme commutatif
$$
\xymatrix{
X \ar[rr]_f \ar[rd]_p & & Y \ar[ld]^q \\
& Z
}
$$
d'espaces algébriques. Supposons que $p$ soit localement de type fini et fermé.
Alors il existe un sous-espace ouvert $W \subset Z$
tel qu'un morphisme $Z' \to Z$ se factorise par $W$ si et seulement si le
changement de base $f_{Z'} : X_{Z'} \to Y_{Z'}$ est non ramifié.
\end{lemma}

\begin{proof}
D'après le
Lemme \ref{spaces-morphisms-lemma-where-unramified} de Morphismes d'espaces,
il existe un sous-espace ouvert $U(f) \subset X$ qui est l'ensemble des
points où $f$ est non ramifié. De plus, la formation de $U(f)$ commute
à tout changement de base. Soit $W \subset Z$ le sous-espace ouvert
(voir
Propriétés des espaces, Lemme
\ref{spaces-properties-lemma-open-subspaces})
dont l'ensemble sous-jacent est
$$
|W| = |Z| \setminus |p|\left(|X| \setminus |U(f)|\right)
$$
c'est-à-dire que $z \in |Z|$ est un point de $W$ si et seulement si $f$ est non ramifié
en tout point de $X$ au-dessus de $z$. Cet ensemble est ouvert, car nous
avons supposé que $p$ est fermé. Puisque la formation de $U(f)$
commute à tout changement de base, on voit immédiatement (en utilisant
Propriétés des espaces, Lemme
\ref{spaces-properties-lemma-factor-through-open-subspace})
que $W$ possède la propriété universelle voulue.
\end{proof}

\begin{lemma}
\label{spaces-more-morphisms-lemma-where-unramified-universally-injective}
Considérons un diagramme commutatif
$$
\xymatrix{
X \ar[rr]_f \ar[rd]_p & & Y \ar[ld]^q \\
& Z
}
$$
d'espaces algébriques. Supposons que
\begin{enumerate}
\item $p$ soit localement de type fini,
\item $p$ soit fermé, et
\item $p_2 : X \times_Y X \to Z$ soit fermé.
\end{enumerate}
Alors il existe un sous-espace ouvert $W \subset Z$
tel qu'un morphisme $Z' \to Z$ se factorise par $W$ si et seulement si le
changement de base $f_{Z'} : X_{Z'} \to Y_{Z'}$ est non ramifié et universellement
injectif.
\end{lemma}

\begin{proof}
Après avoir remplacé $Z$ par le sous-espace ouvert trouvé dans le
Lemme \ref{spaces-more-morphisms-lemma-where-unramified},
on peut supposer que $f$ est déjà non ramifié ; notons que cela ne
remet en cause ni l'hypothèse (2) ni l'hypothèse (3). D'après
Morphismes d'espaces, Lemme
\ref{spaces-morphisms-lemma-diagonal-unramified-morphism},
on voit que $\Delta_{X/Y} : X \to X \times_Y X$ est une immersion ouverte.
Cela reste vrai après tout changement de base. Par conséquent, d'après
Morphismes d'espaces, Lemme
\ref{spaces-morphisms-lemma-universally-injective},
on voit que $f_{Z'}$ est universellement injectif si et seulement si
le changement de base de la diagonale $X_{Z'} \to (X \times_Y X)_{Z'}$
est un isomorphisme. Soit $W \subset Z$ le sous-espace ouvert
(voir
Propriétés des espaces, Lemme
\ref{spaces-properties-lemma-open-subspaces})
dont l'ensemble sous-jacent est
$$
|W| = |Z| \setminus
|p_2|\left(|X \times_Y X| \setminus \Im(|\Delta_{X/Y}|)\right)
$$
c'est-à-dire que $z \in |Z|$ est un point de $W$ si et seulement si la fibre de
$|X \times_Y X| \to |Z|$ au-dessus de $z$ est contenue dans l'image de
$|X| \to |X \times_Y X|$. Il ressort alors clairement de ce qui précède
que la restriction $p^{-1}(W) \to q^{-1}(W)$ de $f$ est
non ramifiée et universellement injective.

\medskip\noindent
Réciproquement, supposons que $f_{Z'}$ soit non ramifié et universellement injectif.
Pour montrer que $Z' \to Z$ se factorise par $W$, il suffit de montrer
que l'image de $|Z'| \to |Z|$ est contenue dans $|W|$, voir
Propriétés des espaces, Lemme
\ref{spaces-properties-lemma-factor-through-open-subspace}.
Il suffit donc de démontrer le résultat lorsque $Z'$ est le spectre d'un corps.
Notons $z \in |Z|$ l'image de $|Z'| \to |Z|$. La discussion ci-dessus montre
que
$$
|X_{Z'}| \longrightarrow |(X \times_Y X)_{Z'}|
$$
est surjective. D'après
Propriétés des espaces,
Lemme \ref{spaces-properties-lemma-points-cartesian},
dans le diagramme commutatif
$$
\xymatrix{
|X_{Z'}| \ar[d] \ar[r] &
|(X \times_Y X)_{Z'}| \ar[d] \\
|p|^{-1}(\{z\}) \ar[r] &
|p_2|^{-1}(\{z\})
}
$$
les flèches verticales sont surjectives. Il s'ensuit que $z \in |W|$, comme voulu.
\end{proof}

\begin{lemma}
\label{spaces-more-morphisms-lemma-where-closed-immersion}
Considérons un diagramme commutatif
$$
\xymatrix{
X \ar[rr]_f \ar[rd]_p & & Y \ar[ld]^q \\
& Z
}
$$
d'espaces algébriques. Supposons que
\begin{enumerate}
\item $p$ soit localement de type fini,
\item $p$ soit universellement fermé, et
\item $q : Y \to Z$ soit séparé.
\end{enumerate}
Alors il existe un sous-espace ouvert $W \subset Z$
tel qu'un morphisme $Z' \to Z$ se factorise par $W$ si et seulement si le
changement de base $f_{Z'} : X_{Z'} \to Y_{Z'}$ est une immersion fermée.
\end{lemma}

\begin{proof}
Nous utiliserons la caractérisation des immersions fermées comme morphismes
universellement fermés, non ramifiés et universellement injectifs, voir
Lemme \ref{spaces-more-morphisms-lemma-characterize-closed-immersion}.
Notons d'abord que, puisque $p$ est universellement fermé et que $q$ est
séparé, $f$ est universellement fermé, voir
Morphismes d'espaces, Lemme
\ref{spaces-morphisms-lemma-universally-closed-permanence}.
Il en résulte que tout changement de base de $f$ est universellement fermé, voir
Morphismes d'espaces, Lemme
\ref{spaces-morphisms-lemma-base-change-universally-closed}.
Ainsi, pour achever la démonstration du lemme, il suffit de prouver que
les hypothèses du
Lemme \ref{spaces-more-morphisms-lemma-where-unramified-universally-injective}
sont satisfaites. La projection $\text{pr}_0 : X \times_Y X \to X$
est universellement fermée comme changement de base de $f$, voir
Morphismes d'espaces, Lemme
\ref{spaces-morphisms-lemma-base-change-universally-closed}.
Par suite, $X \times_Y X \to Z$ est universellement fermé comme
composé de morphismes universellement fermés (voir
Morphismes d'espaces, Lemme
\ref{spaces-morphisms-lemma-composition-universally-closed}).
Cela achève la démonstration du lemme.
\end{proof}

\begin{lemma}
\label{spaces-more-morphisms-lemma-where-flat}
Considérons un diagramme commutatif
$$
\xymatrix{
X \ar[rr]_f \ar[rd]_p & & Y \ar[ld]^q \\
& Z
}
$$
d'espaces algébriques. Supposons que
\begin{enumerate}
\item $p$ soit localement de présentation finie,
\item $p$ soit plat,
\item $p$ soit fermé, et
\item $q$ soit localement de type fini.
\end{enumerate}
Alors il existe un sous-espace ouvert $W \subset Z$
tel qu'un morphisme $Z' \to Z$ se factorise par $W$ si et seulement si le
changement de base $f_{Z'} : X_{Z'} \to Y_{Z'}$ est plat.
\end{lemma}

\begin{proof}
D'après le Lemme \ref{spaces-more-morphisms-lemma-base-change-flatness-fibres},
l'ensemble
$$
A = \{x \in |X| : X\text{ est plat en }x \text{ sur }Y\}.
$$
Cet ensemble est ouvert dans $|X|$ et sa formation commute à tout changement de
base. Soit $W \subset Z$ le sous-espace ouvert
(voir
Propriétés des espaces, Lemme
\ref{spaces-properties-lemma-open-subspaces})
dont l'ensemble sous-jacent est
$$
|W| = |Z| \setminus |p|\left(|X| \setminus A\right)
$$
c'est-à-dire que $z \in |Z|$ est un point de $W$ si et seulement si la fibre entière
de $|X| \to |Z|$ au-dessus de $z$ est contenue dans $A$. Cet ensemble est ouvert car
$p$ est fermé. Puisque la formation de $A$ commute à tout changement de
base, il s'ensuit que $W$ convient.
\end{proof}

\begin{lemma}
\label{spaces-more-morphisms-lemma-where-surjective-flat}
Considérons un diagramme commutatif
$$
\xymatrix{
X \ar[rr]_f \ar[rd]_p & & Y \ar[ld]^q \\
& Z
}
$$
d'espaces algébriques. Supposons que
\begin{enumerate}
\item $p$ soit localement de présentation finie,
\item $p$ soit plat,
\item $p$ soit fermé,
\item $q$ soit localement de type fini, et
\item $q$ soit fermé.
\end{enumerate}
Alors il existe un sous-espace ouvert $W \subset Z$
tel qu'un morphisme $Z' \to Z$ se factorise par $W$ si et seulement si le
changement de base $f_{Z'} : X_{Z'} \to Y_{Z'}$ est surjectif et plat.
\end{lemma}

\begin{proof}
D'après le
Lemme \ref{spaces-more-morphisms-lemma-where-flat},
on peut supposer que $f$ est plat.
Notons que $f$ est localement de présentation finie d'après
Morphismes d'espaces,
Lemme \ref{spaces-morphisms-lemma-finite-presentation-permanence}.
Par conséquent, $f$ est ouvert, voir
Morphismes d'espaces, Lemme \ref{spaces-morphisms-lemma-fppf-open}.
Soit $W \subset Z$ le sous-espace ouvert (voir
Propriétés des espaces, Lemme
\ref{spaces-properties-lemma-open-subspaces})
dont l'ensemble sous-jacent est
$$
|W| = |Z| \setminus |q|\left(|Y| \setminus |f|(|X|)\right).
$$
Autrement dit, pour $z \in |Z|$, on a $z \in |W|$ si et seulement
si la fibre entière de $|Y| \to |Z|$ au-dessus de $z$ est contenue dans l'image de
$|X| \to |Y|$. Puisque $q$ est fermé, cet ensemble est ouvert dans $|Z|$.
Le morphisme $X_W \to Y_W$ est surjectif par construction.
Enfin, supposons que $X_{Z'} \to Y_{Z'}$ soit surjectif.
Pour montrer que $Z' \to Z$ se factorise par $W$, il suffit de montrer
que l'image de $|Z'| \to |Z|$ est contenue dans $|W|$, voir
Propriétés des espaces, Lemme
\ref{spaces-properties-lemma-factor-through-open-subspace}.
Il suffit donc de démontrer le résultat lorsque $Z'$ est le spectre d'un corps.
Notons $z \in |Z|$ l'image de $|Z'| \to |Z|$. D'après
Propriétés des espaces,
Lemme \ref{spaces-properties-lemma-points-cartesian},
dans le diagramme commutatif
$$
\xymatrix{
|X_{Z'}| \ar[d] \ar[r] &
|Y_{Z'}| \ar[d] \\
|p|^{-1}(\{z\}) \ar[r] &
|q|^{-1}(\{z\})
}
$$
les flèches verticales sont surjectives. Il s'ensuit que $z \in |W|$, comme voulu.
\end{proof}

\begin{lemma}
\label{spaces-more-morphisms-lemma-where-isomorphism}
Considérons un diagramme commutatif
$$
\xymatrix{
X \ar[rr]_f \ar[rd]_p & & Y \ar[ld]^q \\
& Z
}
$$
d'espaces algébriques. Supposons que
\begin{enumerate}
\item $p$ soit localement de présentation finie,
\item $p$ soit plat,
\item $p$ soit universellement fermé,
\item $q$ soit localement de type fini,
\item $q$ soit fermé, et
\item $q$ soit séparé.
\end{enumerate}
Alors il existe un sous-espace ouvert $W \subset Z$
tel qu'un morphisme $Z' \to Z$ se factorise par $W$ si et seulement si le
changement de base $f_{Z'} : X_{Z'} \to Y_{Z'}$ est un isomorphisme.
\end{lemma}

\begin{proof}
D'après le
Lemme \ref{spaces-more-morphisms-lemma-where-surjective-flat},
il existe un sous-espace ouvert $W_1 \subset Z$ tel que
$f_{Z'}$ soit surjectif et plat si et seulement si $Z' \to Z$
se factorise par $W_1$. D'après le
Lemme \ref{spaces-more-morphisms-lemma-where-closed-immersion},
il existe un sous-espace ouvert $W_2 \subset Z$ tel que
$f_{Z'}$ soit une immersion fermée si et seulement si $Z' \to Z$
se factorise par $W_2$. Nous affirmons que $W = W_1 \cap W_2$ convient.
Si $f_{Z'}$ est un isomorphisme, alors $Z' \to Z$
se factorise par $W$. Il suffit donc de montrer que
$f_W$ est un isomorphisme. Par construction, $f_W$ est une
immersion fermée surjective et plate. En particulier, $f_W$ est
représentable. Puisqu'une immersion fermée surjective et plate de
schémas est un isomorphisme (voir
Morphismes, Lemme \ref{morphisms-lemma-characterize-flat-closed-immersions}),
cela achève la démonstration. (Notons qu'en fait $f_W$ est localement de présentation finie,
donc ouvert, de sorte que l'on peut éviter d'utiliser ce lemme si on le souhaite.)
\end{proof}

\begin{lemma}
\label{spaces-more-morphisms-lemma-where-lci}
Considérons un diagramme commutatif
$$
\xymatrix{
X \ar[rr]_f \ar[rd]_p & & Y \ar[ld]^q \\
& Z
}
$$
d'espaces algébriques. Supposons que
\begin{enumerate}
\item $p$ soit plat et localement de présentation finie,
\item $p$ soit fermé, et
\item $q$ soit plat et localement de présentation finie.
\end{enumerate}
Alors il existe un sous-espace ouvert $W \subset Z$
tel qu'un morphisme $Z' \to Z$ se factorise par $W$ si et seulement si le
changement de base $f_{Z'} : X_{Z'} \to Y_{Z'}$ est un morphisme localement d'intersection
complète.
\end{lemma}

\begin{proof}
D'après le Lemme \ref{spaces-more-morphisms-lemma-base-change-lci-fibres},
il existe un sous-espace ouvert $U(f) \subset X$ qui est l'ensemble des
points où $f$ est Koszul. De plus, la formation de $U(f)$ commute
à tout changement de base. Soit $W \subset Z$ le sous-espace ouvert
(voir
Propriétés des espaces, Lemme
\ref{spaces-properties-lemma-open-subspaces})
dont l'ensemble sous-jacent est
$$
|W| = |Z| \setminus |p|\left(|X| \setminus |U(f)|\right)
$$
c'est-à-dire que $z \in |Z|$ est un point de $W$ si et seulement si $f$ est Koszul
en tout point de $X$ au-dessus de $z$. Cet ensemble est ouvert car nous
avons supposé que $p$ est fermé. Puisque la formation de $U(f)$
commute à tout changement de base, on voit immédiatement (en utilisant
Propriétés des espaces, Lemme
\ref{spaces-properties-lemma-factor-through-open-subspace})
que $W$ possède la propriété universelle voulue.
\end{proof}








\section{Suites exactes de différentielles et de faisceaux conormaux}
\label{spaces-more-morphisms-section-exact}

\noindent
Dans cette section, nous rassemblons quelques résultats sur les suites exactes de faisceaux
conormaux et de faisceaux de différentielles. En un certain sens, ce sont toutes des
réalisations du triangle des complexes cotangents associé à des
morphismes composables d'espaces algébriques.

\medskip\noindent
Dans les suites ci-dessous, chacune des applications
est construite soit dans le
lemme \ref{spaces-more-morphisms-lemma-functoriality-differentials}, soit dans le
lemme \ref{spaces-more-morphisms-lemma-universal-thickening-functorial}.
Soit $S$ un schéma. Soient $g : Z \to Y$ et $f : Y \to X$ des morphismes
d'espaces algébriques sur $S$.
\begin{enumerate}
\item Il existe une suite exacte canonique
$$
g^*\Omega_{Y/X} \to \Omega_{Z/X} \to \Omega_{Z/Y} \to 0,
$$
voir le
lemme \ref{spaces-more-morphisms-lemma-triangle-differentials}.
Si $g : Z \to Y$ est formellement lisse, cette suite est une suite
exacte courte, voir le
lemme \ref{spaces-more-morphisms-lemma-triangle-differentials-formally-smooth}.
\item Si $g$ est formellement non ramifié, il existe une suite exacte canonique
$$
\mathcal{C}_{Z/Y} \to g^*\Omega_{Y/X} \to \Omega_{Z/X} \to 0,
$$
voir le
lemme \ref{spaces-more-morphisms-lemma-universally-unramified-differentials-sequence}.
Si $f \circ g : Z \to X$ est formellement lisse, cette suite est une suite
exacte courte, voir le
lemme \ref{spaces-more-morphisms-lemma-differentials-formally-unramified-formally-smooth}.
\item Si $g$ et $f \circ g$ sont formellement non ramifiés, il existe une
suite exacte canonique
$$
\mathcal{C}_{Z/X} \to \mathcal{C}_{Z/Y} \to g^*\Omega_{Y/X} \to 0,
$$
voir le
lemme \ref{spaces-more-morphisms-lemma-two-unramified-morphisms}.
Si $f : Y \to X$ est formellement lisse, cette suite est une suite
exacte courte, voir le
lemme \ref{spaces-more-morphisms-lemma-two-unramified-morphisms-formally-smooth}.
\item Si $g$ et $f$ sont formellement non ramifiés, il existe une suite
exacte canonique
$$
g^*\mathcal{C}_{Y/X} \to \mathcal{C}_{Z/X} \to \mathcal{C}_{Z/Y} \to 0.
$$
voir le
lemme \ref{spaces-more-morphisms-lemma-transitivity-conormal-unramified}.
Si $g : Z \to Y$ est un morphisme localement d'intersection complète,
alors cette suite est une suite exacte courte, voir le
lemme \ref{spaces-more-morphisms-lemma-transitivity-conormal-lci}.
\end{enumerate}










\section{Caractérisation des complexes pseudo-cohérents, II}
\label{spaces-more-morphisms-section-characterize-pseudo-coherent}

\noindent
Dans cette section, nous étudions une caractérisation des complexes pseudo-cohérents
en termes de cohomologie. On trouvera des résultats antérieurs sur les complexes pseudo-cohérents
sur les espaces algébriques dans
Catégories dérivées des espaces, section
\ref{spaces-perfect-section-spell-out}
et dans
Catégories dérivées des espaces, section
\ref{spaces-perfect-section-pseudo-coherent-hocolim}.
L'analogue de cette section pour les schémas est Compléments sur les morphismes, section
\ref{more-morphisms-section-characterize-pseudo-coherent}.
Un outil fondamental consistera à nous ramener au cas de l'espace projectif
au moyen d'une version dérivée du lemme de Chow, voir
Lemme \ref{spaces-more-morphisms-lemma-derived-chow}.

\begin{lemma}
\label{spaces-more-morphisms-lemma-case-of-tor-independence}
Soit $S$ un schéma. Considérons un diagramme commutatif d'espaces algébriques
$$
\xymatrix{
Z' \ar[d] \ar[r] & Y' \ar[d] \\
X' \ar[r] & B'
}
$$
sur $S$.
Soit $B \to B'$ un morphisme. Notons $X$ et $Y$ les changements
de base de $X'$ et $Y'$ à $B$.
Supposons que $Y' \to B'$ et $Z' \to X'$ soient plats.
Alors $X \times_B Y$ et $Z'$ sont Tor-indépendants sur $X' \times_{B'} Y'$.
\end{lemma}

\begin{proof}
D'après Catégories dérivées des espaces, Lemme
\ref{spaces-perfect-lemma-tor-independent},
on peut vérifier l'indépendance pour Tor localement pour la topologie \'etale sur $X \times_B Y$
et $Z'$. Cela\footnote{Voici l'argument plus en détail.
Choisissons un morphisme \'etale surjectif $W' \to B'$
où $W'$ est un schéma. Choisissons un morphisme \'etale surjectif
$W \to B \times_{B'} W'$ où $W$ est un schéma. Choisissons un
morphisme \'etale surjectif
$U' \to X' \times_{B'} W'$ où $U'$ est un schéma. Choisissons un
morphisme \'etale surjectif $V' \to Y' \times_{B'} W'$ où $V'$ est un schéma.
Remarquons que $U' \times_{W'} V' \to X' \times_{B'} Y'$ est surjectif et
\'etale. Choisissons un morphisme \'etale surjectif
$T' \to Z' \times_{X' \times_{B'} Y'} U' \times_{W'} V'$
où $T'$ est un schéma. Notons $U$ et $V$ les changements de base de $U'$ et $V'$
à $W$. Le lemme affirme alors que $X \times_B Y$ et $Z'$
sont Tor-indépendants sur $X' \times_{B'} Y'$ comme espaces algébriques
si et seulement si $U \times_W V$ et $T'$ sont Tor-indépendants sur
$U' \times_{W'} V'$ comme schémas. Ainsi,
il suffit de démontrer le lemme pour
le carré de sommets $T', U', V', W'$ et le changement de base par $W \to W'$.
La platitude de $Y' \to B'$ et de $Z' \to X'$ entraîne celle
de $V' \to W'$ et de $T' \to U'$.}
ramène le lemme au cas des schémas,
qui est Compléments sur les morphismes, Lemme
\ref{more-morphisms-lemma-case-of-tor-independence}.
\end{proof}

\begin{lemma}[Lemme de Chow dérivé]
\label{spaces-more-morphisms-lemma-derived-chow}
Soit $A$ un anneau. Soit $X$ un espace algébrique séparé
de présentation finie sur $A$. Soit $x \in |X|$. Alors il existe
un $n \geq 0$,
un sous-espace fermé $Z \subset X \times_A \mathbf{P}^n_A$,
un point $z \in |Z|$,
un ouvert $V \subset \mathbf{P}^n_A$, et
un objet $E$ de $D(\mathcal{O}_{X \times_A \mathbf{P}^n_A})$ tels que
\begin{enumerate}
\item $Z \to X \times_A \mathbf{P}^n_A$ est de présentation finie,
\item $c : Z \to \mathbf{P}^n_A$ est une immersion fermée au-dessus de $V$,
posons $W = c^{-1}(V)$,
\item la restriction de $b : Z \to X$ à $W$ est \'etale,
$z \in W$, et $b(z) = x$,
\item $E|_{X \times_A V} \cong
(b, c)_*\mathcal{O}_Z|_{X \times_A V}$,
\item $E$ est pseudo-cohérent et à support dans $Z$.
\end{enumerate}
\end{lemma}

\begin{proof}
Nous pouvons trouver une sous-$\mathbf{Z}$-algèbre de type fini $A' \subset A$
et un espace algébrique $X'$ séparé et de présentation finie sur $A'$
dont le changement de base à $A$ est $X$. Voir
Limites d'espaces, Lemmes
\ref{spaces-limits-lemma-descend-finite-presentation} et
\ref{spaces-limits-lemma-descend-separated-morphism}.
Soit $x' \in |X'|$ l'image de $x$.
Si nous pouvons démontrer le lemme pour $(X'/A', x')$, alors
le lemme en découle pour $(X/A, x)$.
En effet, si $n', Z', z', V', E'$ fournissent les données voulues
pour $(X'/A', x')$, nous pouvons poser
$n = n'$,
prendre pour $Z \subset X \times \mathbf{P}^n$ l'image inverse de $Z'$,
prendre pour $z \in Z$ l'unique point d'image $x$,
prendre pour $V \subset \mathbf{P}^n_A$ l'image inverse de $V'$, et
prendre pour $E$ l'image inverse dérivée de $E'$.
Remarquons que $E$ est pseudo-cohérent d'après
Cohomologie sur les sites, Lemme
\ref{sites-cohomology-lemma-pseudo-coherent-pullback}.
Il ne reste qu'à vérifier (5). Pour cela,
posons $W = c^{-1}(V)$ et $W' = (c')^{-1}(V')$
et considérons le carré cartésien
$$
\xymatrix{
W \ar[d]_{(b, c)} \ar[r] & W' \ar[d]^{(b', c')} \\
X \times_A V \ar[r] & X' \times_{A'} V'
}
$$
D'après le Lemme \ref{spaces-more-morphisms-lemma-case-of-tor-independence}, $X \times_A V$ et $W'$
sont Tor-indépendants sur $X' \times_{A'} V'$.
Ainsi, l'image inverse dérivée de
$(b', c')_*\mathcal{O}_{W'}$ sur $X \times_A V$
est $(b, c)_*\mathcal{O}_W$ d'après
Catégories dérivées des espaces,
Lemme \ref{spaces-perfect-lemma-compare-base-change}.
On utilise aussi ici le fait que $R(b', c')_*\mathcal{O}_{Z'} = (b', c')_*\mathcal{O}_{Z'}$
parce que $(b', c')$ est une immersion fermée, et de même pour
$(b, c)_*\mathcal{O}_Z$.
Comme $E'|_{U' \times_{A'} V'} =
(b', c')_*\mathcal{O}_{W'}$, on obtient
$E|_{U \times_A V} = (b, c)_*\mathcal{O}_W$,
et (5) est vérifié.
Nous sommes ainsi ramenés à la situation décrite au
paragraphe suivant.

\medskip\noindent
Supposons que $A$ soit de type fini sur $\mathbf{Z}$.
Choisissons un morphisme \'etale $U \to X$, où $U$ est un schéma affine,
et un point $u \in U$ d'image $x$. Alors $U$ est de type fini sur $A$.
Choisissons une immersion fermée $U \to \mathbf{A}^n_A$ et notons
$j : U \to \mathbf{P}^n_A$ l'immersion obtenue en la composant
avec l'immersion ouverte $\mathbf{A}^n_A \to \mathbf{P}^n_A$.
Soit $Z$ l'adhérence schématique de
$$
(\text{id}_U, j) : U \longrightarrow X \times_A \mathbf{P}^n_A
$$
Soit $z \in Z$ l'image de $u$.
Soit $Y \subset \mathbf{P}^n_A$ l'adhérence schématique
de $j$. Il est alors clair que $Z \subset X \times_A Y$
est l'adhérence schématique de
$(\text{id}_U, j) : U \to X \times_A Y$.
Comme $X$ est séparé, le morphisme
$X \times_A Y \to Y$ est lui aussi séparé.
On voit donc que $Z \to Y$ est un isomorphisme au-dessus
du sous-schéma ouvert $j(U) \subset Y$ d'après
Morphismes d'espaces, Lemme
\ref{spaces-morphisms-lemma-scheme-theoretic-image-of-partial-section}.
Choisissons un ouvert $V \subset \mathbf{P}^n_A$ tel que $V \cap Y = j(U)$.
On voit alors que (2) est vérifié, que $W = (\text{id}_U, j)(U)$, et donc
que (3) est vérifié. La partie (1) est vérifiée parce que $A$ est noethérien.

\medskip\noindent
Comme $A$ est noethérien, $X$ et $X \times_A \mathbf{P}^n_A$
sont des espaces algébriques noethériens. Nous pouvons donc prendre $E = (b, c)_*\mathcal{O}_Z$
dans ce cas : (4) est clair et, pour (5), voir
Catégories dérivées des espaces, Lemme
\ref{spaces-perfect-lemma-identify-pseudo-coherent-noetherian}.
Ceci achève la démonstration.
\end{proof}

\begin{lemma}
\label{spaces-more-morphisms-lemma-compute-Fourier-Mukai-for-derived-chow}
Soient $X/A$, $x \in |X|$, et
$n, Z, z, V, E$ comme dans le Lemme \ref{spaces-more-morphisms-lemma-derived-chow}.
Pour tout $K \in D_\QCoh(\mathcal{O}_X)$, on a
$$
Rq_*(Lp^*K \otimes^\mathbf{L} E)|_V = R(W \to V)_*K|_W
$$
où $p : X \times_A \mathbf{P}^n_A \to X$ et
$q : X \times_A \mathbf{P}^n_A \to \mathbf{P}^n_A$ sont
les projections et où le morphisme $W \to V$ est
l'immersion fermée de présentation finie $c|_W : W \to V$.
\end{lemma}

\begin{proof}
Comme $W = c^{-1}(V)$ et que $c$ est une immersion fermée
au-dessus de $V$, on voit que $c|_W$ est une immersion fermée.
Elle est de présentation finie parce que $W$ et $V$ sont de présentation
finie sur $A$, voir
Morphismes d'espaces, Lemme
\ref{spaces-morphisms-lemma-finite-presentation-permanence}.
Tout d'abord, on a
$$
Rq_*(Lp^*K \otimes^\mathbf{L} E)|_V =
Rq'_*\left((Lp^*K \otimes^\mathbf{L} E)|_{X \times_A V}\right)
$$
où $q' : X \times_A V \to V$ est la projection, car
la formation de l'image directe totale commute à la localisation.
Notons $i = (b, c)|_W : W \to X \times_A V$ l'immersion fermée ainsi obtenue.
Alors
$$
Rq'_*\left((Lp^*K \otimes^\mathbf{L} E)|_{X \times_A V}\right) =
Rq'_*(Lp^*K|_{X \times_A V} \otimes^\mathbf{L} i_*\mathcal{O}_W)
$$
d'après la propriété (5). Comme $i$ est une immersion fermée, on a
$i_*\mathcal{O}_W = Ri_*\mathcal{O}_W$.
En utilisant
Catégories dérivées des espaces,
Lemme \ref{spaces-perfect-lemma-cohomology-base-change},
nous pouvons réécrire cette expression sous la forme
$$
Rq'_* Ri_* Li^* Lp^*K|_{X \times_A V} =
R(q' \circ i)_* Lb^*K|_W =
R(W \to V)_* K|_W
$$
ce qui est l'égalité cherchée. (Notons que se restreindre à $W$
et prendre l'image inverse dérivée par $W \to X$ reviennent au même
puisque $W$ est \'etale sur $X$.)
\end{proof}

\begin{lemma}
\label{spaces-more-morphisms-lemma-characterize-pseudo-coherent}
Soit $A$ un anneau. Soit $X$ un espace algébrique séparé et
de présentation finie sur $A$. Soit $K \in D_\QCoh(\mathcal{O}_X)$.
Si $R\Gamma(X, E \otimes^\mathbf{L} K)$ est pseudo-cohérent
dans $D(A)$ pour tout $E$ pseudo-cohérent dans $D(\mathcal{O}_X)$,
alors $K$ est pseudo-cohérent relativement à $A$
(Définition \ref{spaces-more-morphisms-definition-relative-pseudo-coherence}).
\end{lemma}

\begin{proof}
Supposons que $K \in D_\QCoh(\mathcal{O}_X)$ et que
$R\Gamma(X, E \otimes^\mathbf{L} K)$ soit pseudo-cohérent
dans $D(A)$ pour tout $E$ pseudo-cohérent dans $D(\mathcal{O}_X)$.
Soit $x \in |X|$. Nous allons montrer que $K$ est pseudo-cohérent relativement à $A$
dans un voisinage \'etale de $x$. Cela démontrera le lemme
par notre définition de la pseudo-cohérence relative.

\medskip\noindent
Choisissons $n, Z, z, V, E$ comme dans le Lemme \ref{spaces-more-morphisms-lemma-derived-chow}.
Notons $p : X \times \mathbf{P}^n \to X$ et
$q : X \times \mathbf{P}^n \to \mathbf{P}^n_A$
les projections.
Alors, pour tout $i \in \mathbf{Z}$, on a
\begin{align*}
& R\Gamma(\mathbf{P}^n_A,
Rq_*(Lp^*K \otimes^\mathbf{L} E)
\otimes^\mathbf{L}
\mathcal{O}_{\mathbf{P}^n_A}(i)) \\
& =
R\Gamma(X \times \mathbf{P}^n,
Lp^*K \otimes^\mathbf{L} E \otimes^\mathbf{L}
Lq^*\mathcal{O}_{\mathbf{P}^n_A}(i)) \\
& =
R\Gamma(X,
K \otimes^\mathbf{L} Rp_*(E \otimes^\mathbf{L}
Lq^*\mathcal{O}_{\mathbf{P}^n_A}(i)))
\end{align*}
d'après
Catégories dérivées des espaces,
Lemme \ref{spaces-perfect-lemma-cohomology-base-change}.
D'après
Catégories dérivées des espaces, Lemme
\ref{spaces-perfect-lemma-flat-proper-pseudo-coherent-direct-image-general},
le complexe $Rp_*(E \otimes^\mathbf{L} Lq^*\mathcal{O}_{\mathbf{P}^n_A}(i))$
est pseudo-cohérent sur $X$. L'hypothèse nous dit donc que l'expression
figurant dans la formule affichée est un objet pseudo-cohérent de $D(A)$.
D'après
Catégories dérivées des schémas,
Lemme \ref{perfect-lemma-pseudo-coherent-on-projective-space},
nous en concluons que $Rq_*(Lp^*K \otimes^\mathbf{L} E)$ est
pseudo-cohérent sur $\mathbf{P}^n_A$.
D'après le Lemme \ref{spaces-more-morphisms-lemma-compute-Fourier-Mukai-for-derived-chow},
nous avons
$$
Rq_*(Lp^*K \otimes^\mathbf{L} E)|_{X \times_A V} =
R(W \to V)_*K|_W
$$
Comme $W \to V$ est une immersion fermée dans un sous-schéma ouvert de
$\mathbf{P}^n_A$, cela signifie que $K|_W$ est pseudo-cohérent relativement à $A$,
par exemple d'après
Compléments sur les morphismes,
Lemme \ref{more-morphisms-lemma-check-relative-pseudo-coherence-on-charts}.
\end{proof}

\begin{lemma}
\label{spaces-more-morphisms-lemma-characterize-pseudo-coh-improved}
Soit $A$ un anneau. Soit $X$ un espace algébrique séparé et
de présentation finie sur $A$. Soit $K \in D_\QCoh(\mathcal{O}_X)$. Si
$R \Gamma (X, E \otimes ^{\mathbf{L}} K)$ est
pseudo-cohérent dans $D(A)$ pour tout objet parfait
$E \in D(\mathcal{O}_X)$, alors $K$ est pseudo-cohérent
relativement à $A$.
\end{lemma}

\begin{proof}
Compte tenu du Lemme \ref{spaces-more-morphisms-lemma-characterize-pseudo-coherent}, il suffit
de montrer que $R \Gamma (X, E \otimes ^{\mathbf{L}} K)$ est
pseudo-cohérent dans $D(A)$ pour tout objet pseudo-cohérent
$E \in D(\mathcal{O}_X)$. D'après Catégories dérivées des espaces,
Proposition \ref{spaces-perfect-proposition-detecting-bounded-above},
il en résulte que $K \in D^-_\QCoh (\mathcal{O}_X)$. Le
résultat découle alors de Catégories dérivées des espaces, Lemme
\ref{spaces-perfect-lemma-perfect-enough}.
\end{proof}












\section{Objets relativement parfaits}
\label{spaces-more-morphisms-section-relatively-perfect}

\noindent
Dans cette section, nous introduisons une notion provenant de \cite{lieblich-complexes}.
Cette notion a été étudiée pour les morphismes de schémas dans
Catégories dérivées de schémas, section
\ref{perfect-section-relatively-perfect}.

\begin{definition}
\label{spaces-more-morphisms-definition-relatively-perfect}
Soit $S$ un schéma. Soit $f : X \to Y$ un morphisme d'espaces algébriques
sur $S$ qui soit plat et localement de présentation finie.
Un objet $E$ de $D(\mathcal{O}_X)$ est {\it parfait relativement à $Y$} ou
{\it $Y$-parfait} si $E$ est pseudo-cohérent
(Cohomologie sur les sites, définition
\ref{sites-cohomology-definition-pseudo-coherent}) et si
$E$ est localement de dimension de Tor finie en tant qu'objet de
$D(f^{-1}\mathcal{O}_Y)$
(Cohomologie sur les sites, définition
\ref{sites-cohomology-definition-tor-amplitude}).
\end{definition}

\noindent
Voir Catégories dérivées de schémas,
remarque \ref{perfect-remark-discuss-rel-perfect}, pour une discussion ;
signalons seulement ici que la pseudo-cohérence de $E$ est
équivalente à la pseudo-cohérence de $E$ relativement à $Y$ d'après le
lemme \ref{spaces-more-morphisms-lemma-relative-pseudo-coherent-is-moot}.
De plus, la pseudo-cohérence de $E$ implique $E \in D_\QCoh(\mathcal{O}_X)$ ; voir
Catégories dérivées d'espaces, lemme
\ref{spaces-perfect-lemma-pseudo-coherent}.

\begin{example}
\label{spaces-more-morphisms-example-relatively-perfect-field}
Soit $k$ un corps. Soit $X$ un espace algébrique de présentation finie
sur $k$ (en particulier, $X$ est quasi-compact). Alors un objet $E$ de
$D(\mathcal{O}_X)$ est $k$-parfait si et seulement s'il est borné et
pseudo-cohérent (par définition), c'est-à-dire si et seulement s'il appartient à
$D^b_{\textit{Coh}}(X)$ (d'après
Catégories dérivées d'espaces, lemme
\ref{spaces-perfect-lemma-identify-pseudo-coherent-noetherian}).
Ainsi, être relativement parfait ne signifie {\bf pas} ``être parfait sur les fibres''.
\end{example}

\noindent
La notion algébrique correspondante est étudiée dans
Compléments d'algèbre, section \ref{more-algebra-section-relatively-perfect}.
Nous pouvons relier la notion pour les espaces algébriques à la
notion algébrique comme suit.

\begin{lemma}
\label{spaces-more-morphisms-lemma-affine-locally-rel-perfect}
Soit $S$ un schéma. Soit $f : X \to Y$ un morphisme d'espaces algébriques
sur $S$ qui soit plat et localement de présentation finie.
Soit $E \in D_\QCoh(\mathcal{O}_X)$. Les conditions suivantes sont équivalentes :
\begin{enumerate}
\item $E$ est $Y$-parfait,
\item pour tout diagramme commutatif
$$
\xymatrix{
U \ar[d] \ar[r]_g & V \ar[d] \\
X \ar[r]^f & Y
}
$$
où $U$ et $V$ sont des schémas et où les flèches verticales sont \'etales, le complexe
$E|_U$ est $V$-parfait au sens de Catégories dérivées de schémas,
définition \ref{perfect-definition-relatively-perfect},
\item pour un diagramme commutatif comme en (2) tel que $U \to X$ soit
surjectif, le complexe $E|_U$ est $V$-parfait au sens de
Catégories dérivées de schémas,
définition \ref{perfect-definition-relatively-perfect},
\item pour tout diagramme commutatif comme en (2) avec $U$ et $V$
affines, le complexe $R\Gamma(U, E)$ est $\mathcal{O}_Y(V)$-parfait.
\end{enumerate}
\end{lemma}

\begin{proof}
Pour donner un sens aux assertions (2), (3), (4) du lemme, remarquons que
l'objet $E|_U$ de $D_\QCoh(\mathcal{O}_U)$ correspond à
un objet $E_0$ de $D_\QCoh(\mathcal{O}_{U_0})$, où $U_0$
désigne le schéma sous-jacent à $U$ ; voir Catégories dérivées d'espaces,
lemme \ref{spaces-perfect-lemma-derived-quasi-coherent-small-etale-site}.
De plus, dans ce cas, $E_0$ est pseudo-cohérent si et seulement si
$E|_U$ est pseudo-cohérent ; voir Catégories dérivées d'espaces,
lemme \ref{spaces-perfect-lemma-descend-pseudo-coherent}.
De même, $E|_U$ est localement de dimension de Tor finie sur
$f^{-1}\mathcal{O}_Y|_U = g^{-1}\mathcal{O}_V$ si et seulement si
$E_0$ est localement de dimension de Tor finie sur $g_0^{-1}\mathcal{O}_{V_0}$
d'après Catégories dérivées d'espaces, lemme
\ref{spaces-perfect-lemma-tor-dimension-rel}.
Ici, $g_0 : U_0 \to V_0$ est le morphisme de schémas représentant
$g : U \to V$ (notations de Catégories dérivées d'espaces,
remarque \ref{spaces-perfect-remark-match-total-direct-images}).
Enfin, remarquons que ``être pseudo-cohérent'' est une propriété locale pour la topologie \'etale et
que, bien entendu, ``avoir localement une dimension de Tor finie'' est une propriété locale pour la topologie \'etale.
Nous voyons ainsi qu'il suffit de vérifier localement pour la topologie \'etale que l'objet est $Y$-parfait
et, d'après la discussion ci-dessus, que (1) implique (2) et que
(3) implique (1). Comme l'assertion (4) est équivalente
à (2) et (3) d'après
Catégories dérivées de schémas, lemme
\ref{perfect-lemma-affine-locally-rel-perfect},
la démonstration est achevée.
\end{proof}

\begin{lemma}
\label{spaces-more-morphisms-lemma-triangulated}
Soit $S$ un schéma. Soit $f : X \to Y$ un morphisme
d'espaces algébriques sur $S$ qui soit plat et localement de présentation finie.
La sous-catégorie pleine de $D(\mathcal{O}_X)$ formée des objets $Y$-parfaits est
une sous-catégorie triangulée saturée\footnote{Catégories dérivées, définition
\ref{derived-definition-saturated}.}.
\end{lemma}

\begin{proof}
Cela résulte des lemmes de Cohomologie sur les sites
\ref{sites-cohomology-lemma-cone-pseudo-coherent},
\ref{sites-cohomology-lemma-summands-pseudo-coherent},
\ref{sites-cohomology-lemma-cone-tor-amplitude} et
\ref{sites-cohomology-lemma-summands-tor-amplitude}.
\end{proof}

\begin{lemma}
\label{spaces-more-morphisms-lemma-perfect-relatively-perfect}
Soit $S$ un schéma.
Soit $f : X \to Y$ un morphisme d'espaces algébriques sur $S$
qui soit plat et localement de présentation finie.
Tout objet parfait de $D(\mathcal{O}_X)$ est $Y$-parfait.
Si $K, M \in D(\mathcal{O}_X)$, alors $K \otimes_{\mathcal{O}_X}^\mathbf{L} M$
est $Y$-parfait si $K$ est parfait et si $M$ est $Y$-parfait.
\end{lemma}

\begin{proof}
Réduisons-nous au cas des schémas à l'aide du
lemme \ref{spaces-more-morphisms-lemma-affine-locally-rel-perfect},
puis appliquons
Catégories dérivées de schémas, lemme
\ref{perfect-lemma-perfect-relatively-perfect}.
\end{proof}

\begin{lemma}
\label{spaces-more-morphisms-lemma-base-change-relatively-perfect}
Soit $S$ un schéma.
Soit $f : X \to Y$ un morphisme d'espaces algébriques sur $S$
qui soit plat et localement de présentation finie.
Soit $g : Y' \to Y$ un morphisme d'espaces algébriques sur $S$.
Posons $X' = Y' \times_Y X$ et notons $g' : X' \to X$ la projection.
Si $K \in D(\mathcal{O}_X)$ est $Y$-parfait, alors $L(g')^*K$
est $Y'$-parfait.
\end{lemma}

\begin{proof}
Réduisons-nous au cas des schémas à l'aide du
lemme \ref{spaces-more-morphisms-lemma-affine-locally-rel-perfect},
puis appliquons
Catégories dérivées de schémas, lemme
\ref{perfect-lemma-base-change-relatively-perfect}.
\end{proof}

\begin{situation}
\label{spaces-more-morphisms-situation-relative-descent}
Soit $S$ un schéma.
Soit $Y = \lim_{i \in I} Y_i$ la limite projective d'un système filtrant
d'espaces algébriques sur $S$
à morphismes de transition affines $g_{i'i} : Y_{i'} \to Y_i$.
Nous supposons que $Y_i$ est quasi-compact et quasi-séparé pour tout $i \in I$.
Notons $g_i : Y \to Y_i$ la projection. Fixons un élément $0 \in I$
et un morphisme plat de présentation finie $X_0 \to Y_0$.
Posons $X_i = Y_i \times_{Y_0} X_0$ et $X = Y \times_{Y_0} X_0$
et notons les morphismes de transition $f_{i'i} : X_{i'} \to X_i$
et les projections $f_i : X \to X_i$.
\end{situation}

\begin{lemma}
\label{spaces-more-morphisms-lemma-relative-descend-homomorphisms}
Plaçons-nous dans la Situation \ref{spaces-more-morphisms-situation-relative-descent}.
Soient $K_0$ et $L_0$ des objets de $D(\mathcal{O}_{X_0})$.
Posons $K_i = Lf_{i0}^*K_0$ et $L_i = Lf_{i0}^*L_0$ pour $i \geq 0$
et posons $K = Lf_0^*K_0$ et $L = Lf_0^*L_0$. Alors l'application
$$
\colim_{i \geq 0} \Hom_{D(\mathcal{O}_{X_i})}(K_i, L_i)
\longrightarrow
\Hom_{D(\mathcal{O}_X)}(K, L)
$$
est un isomorphisme si $K_0$ est pseudo-cohérent et si
$L_0 \in D_\QCoh(\mathcal{O}_{X_0})$ est (localement) de
dimension de Tor finie en tant qu'objet de
$D((X_0 \to Y_0)^{-1}\mathcal{O}_{Y_0})$.
\end{lemma}

\begin{proof}
Pour tout objet quasi-compact et quasi-séparé $U_0$ de
$(X_0)_{spaces, \etale}$, considérons la condition $P$ selon laquelle
$$
\colim_{i \geq 0} \Hom_{D(\mathcal{O}_{U_i})}(K_i|_{U_i}, L_i|_{U_i})
\longrightarrow
\Hom_{D(\mathcal{O}_U)}(K|_U, L|_U)
$$
est un isomorphisme, où $U = X \times_{X_0} U_0$ et
$U_i = X_i \times_{X_0} U_0$. Nous allons démontrer que $P$ vaut pour tout $U_0$.

\medskip\noindent
Supposons que $(U_0 \subset W_0, V_0 \to W_0)$ soit un carré
distingué élémentaire de $(X_0)_{spaces, \etale}$
et que $P$ vaille pour les trois objets $U_0, V_0, U_0 \times_{W_0} V_0$.
Alors $P$ vaut pour $W_0$ par Mayer--Vietoris
pour les Hom dans la catégorie dérivée ; voir Catégories dérivées d'espaces,
lemme \ref{spaces-perfect-lemma-mayer-vietoris-hom}.

\medskip\noindent
Considérons d'abord $U_0 = W_0 \times_{Y_0} X_0$, où $W_0$ est un objet
quasi-compact et quasi-séparé de $(Y_0)_{spaces, \etale}$.
D'après le principe d'induction de Catégories dérivées d'espaces,
lemme \ref{spaces-perfect-lemma-induction-principle},
appliqué à ces $W_0$ et au paragraphe précédent,
nous constatons qu'il suffit de démontrer
$P$ pour $U_0 = W_0 \times_{Y_0} X_0$ avec $W_0$ affine.
Autrement dit, nous nous sommes ramenés au cas où $Y_0$ est affine.
Ensuite, appliquons de nouveau le principe d'induction, cette fois à tous les
ouverts quasi-compacts et quasi-séparés de $X_0$, pour nous ramener au
cas où $X_0$ est également affine.

\medskip\noindent
Si $X_0$ et $Y_0$ sont affines, nous sommes ramenés au cas
des schémas, démontré dans
Catégories dérivées de schémas, lemme
\ref{perfect-lemma-relative-descend-homomorphisms}.
Le lecteur pourra utiliser
Catégories dérivées d'espaces, lemmes
\ref{spaces-perfect-lemma-pseudo-coherent},
\ref{spaces-perfect-lemma-derived-quasi-coherent-small-etale-site},
\ref{spaces-perfect-lemma-descend-pseudo-coherent} et
\ref{spaces-perfect-lemma-tor-dimension-rel}
pour traduire l'énoncé en un énoncé
ne faisant intervenir que des schémas et des catégories dérivées de modules sur des schémas.
\end{proof}

\begin{lemma}
\label{spaces-more-morphisms-lemma-descend-relatively-perfect}
Dans la Situation \ref{spaces-more-morphisms-situation-relative-descent}, la catégorie des
objets $Y$-parfaits de $D(\mathcal{O}_X)$ est la limite inductive des catégories
d'objets $Y_i$-parfaits de $D(\mathcal{O}_{X_i})$.
\end{lemma}

\begin{proof}
Pour tout objet quasi-compact et quasi-séparé $U_0$ de
$(X_0)_{spaces, \etale}$, considérons la condition $P$
selon laquelle le foncteur
$$
\colim_{i \geq 0} D_{Y_i\text{-parfait}}(\mathcal{O}_{U_i})
\longrightarrow
D_{Y\text{-parfait}}(\mathcal{O}_U)
$$
est une équivalence, où $U = X \times_{X_0} U_0$ et
$U_i = X_i \times_{X_0} U_0$.
Nous savons déjà que ce foncteur est pleinement fidèle
d'après le lemme \ref{spaces-more-morphisms-lemma-relative-descend-homomorphisms}. Il suffit donc de montrer
qu'il est essentiellement surjectif.

\medskip\noindent
Supposons que $(U_0 \subset W_0, V_0 \to W_0)$ soit un carré
distingué élémentaire de $(X_0)_{spaces, \etale}$
et que $P$ vaille pour les trois objets $U_0, V_0, U_0 \times_{W_0} V_0$.
Nous affirmons que $P$ vaut pour $W_0$. Nous emploierons les notations
$U_i = X_i \times_{X_0} U_0$, $U = X \times_{X_0} U_0$,
et de même pour $V_0$ et $W_0$. Par abus de notation, nous emploierons le symbole
$f_i$ pour tous les morphismes $U \to U_i$, $V \to V_i$,
$U \times_W V \to U_i \times_{W_i} V_i$ et $W \to W_i$.
Supposons que $E$ soit un objet $Y$-parfait de $D(\mathcal{O}_W)$.
But : montrer que $E$ appartient à l'image essentielle du foncteur.
Par hypothèse,
nous pouvons trouver $i \geq 0$, un objet $Y_i$-parfait $E_{U, i}$ sur $U_i$,
un objet $Y_i$-parfait $E_{V, i}$ sur $V_i$, et
des isomorphismes $Lf_i^*E_{U, i} \to E|_U$ et $Lf_i^*E_{V, i} \to E|_V$.
Soient
$$
a : E_{U, i} \to (Rf_{i, *}E)|_{U_i}
\quad\text{et}\quad
b : E_{V, i} \to (Rf_{i, *}E)|_{V_i}
$$
les morphismes adjoints aux isomorphismes $Lf_i^*E_{U, i} \to E|_U$
et $Lf_i^*E_{V, i} \to E|_V$.
Puisque le foncteur est pleinement fidèle, quitte à augmenter $i$,
nous pouvons trouver un isomorphisme
$c : E_{U, i}|_{U_i \times_{W_i} V_i} \to E_{V, i}|_{U_i \times_{W_i} V_i}$
qui, par image inverse, donne les identifications
$$
Lf_i^*E_{U, i}|_{U \times_W V} \to E|_{U \times_W V} \to
Lf_i^*E_{V, i}|_{U \times_W V}.
$$
Appliquons Catégories dérivées d'espaces, lemme
\ref{spaces-perfect-lemma-glue},
pour obtenir un objet $E_i$ sur $W_i$ et un morphisme $d : E_i \to Rf_{i, *}E$
dont les restrictions sont les morphismes $a$ et $b$ sur $U_i$ et $V_i$.
Il est alors clair que $E_i$ est $Y_i$-parfait (car la propriété d'être
relativement parfait est locale pour la topologie \'etale) et que
$d$ est adjoint à un isomorphisme $Lf_i^*E_i \to E$.

\medskip\noindent
En reprenant exactement l'argument employé dans
la démonstration du lemme \ref{spaces-more-morphisms-lemma-relative-descend-homomorphisms},
en utilisant le principe d'induction
(Catégories dérivées d'espaces, lemme
\ref{spaces-perfect-lemma-induction-principle}),
nous nous ramenons au cas où $X_0$ et $Y_0$ sont tous deux
affines : travaillons d'abord avec les objets quasi-compacts et quasi-séparés
de $(Y_0)_{spaces, \etale}$ pour nous ramener au cas où
$Y_0$ est affine, puis avec les objets quasi-compacts et quasi-séparés
de $(X_0)_{spaces, \etale}$ pour nous ramener au cas où $X_0$ est affine.
Dans le cas affine, le résultat découle du cas des schémas, à savoir
Catégories dérivées de schémas, lemme
\ref{perfect-lemma-descend-relatively-perfect}.
Le passage au cas des schémas se fait grâce au
lemme \ref{spaces-more-morphisms-lemma-affine-locally-rel-perfect}.
\end{proof}

\begin{lemma}
\label{spaces-more-morphisms-lemma-derived-pushforward-rel-perfect}
Soit $S$ un schéma.
Soit $f : X \to Y$ un morphisme d'espaces algébriques sur $S$
qui soit plat, propre et
de présentation finie. Soit $E \in D(\mathcal{O}_X)$ un objet $Y$-parfait.
Alors $Rf_*E$ est un objet parfait de $D(\mathcal{O}_Y)$
et sa formation commute à tout changement de base.
\end{lemma}

\begin{proof}
L'assertion relative au changement de base est Catégories dérivées d'espaces,
lemme \ref{spaces-perfect-lemma-base-change-tensor} (avec
$\mathcal{G}^\bullet$ égal à $\mathcal{O}_X$ placé en degré $0$).
Il suffit donc de montrer que $Rf_*E$ est un objet parfait. Nous allons nous ramener
au cas où $Y$ est affine noethérien par un argument de passage à la limite.

\medskip\noindent
La question est locale sur $Y$ pour la topologie \'etale ; nous pouvons donc supposer $Y$ affine.
Posons $Y = \Spec(R)$. Écrivons $R = \colim R_i$ comme limite inductive filtrante
d'anneaux noethériens $R_i$. D'après Limites d'espaces, lemme
\ref{spaces-limits-lemma-descend-finite-presentation},
il existe un $i$ et un espace algébrique
$X_i$ de présentation finie sur $R_i$
dont le changement de base à $R$ est $X$. D'après
Limites d'espaces, lemmes \ref{spaces-limits-lemma-eventually-proper} et
\ref{spaces-limits-lemma-descend-flat},
on peut supposer que $X_i$ est propre et plat sur $R_i$.
D'après le lemme \ref{spaces-more-morphisms-lemma-descend-relatively-perfect},
on peut supposer qu'il existe un objet $R_i$-parfait $E_i$ de
$D(\mathcal{O}_{X_i})$ dont l'image inverse sur $X$ est $E$.
En appliquant Catégories dérivées d'espaces,
lemme \ref{spaces-perfect-lemma-perfect-direct-image},
à $X_i \to \Spec(R_i)$ et à $E_i$, et en utilisant la propriété de
changement de base déjà démontrée, on obtient le résultat.
\end{proof}

\begin{lemma}
\label{spaces-more-morphisms-lemma-compute-ext-rel-perfect}
Soit $S$ un schéma.
Soit $f : X \to Y$ un morphisme d'espaces algébriques sur $S$.
Soient $E, K \in D(\mathcal{O}_X)$.
Supposons que
\begin{enumerate}
\item $Y$ soit quasi-compact et quasi-séparé,
\item $f$ soit propre, plat et de présentation finie,
\item $E$ soit $Y$-parfait,
\item $K$ soit pseudo-cohérent.
\end{enumerate}
Alors il existe un objet pseudo-cohérent $L \in D(\mathcal{O}_Y)$ tel que
$$
Rf_*R\SheafHom(K, E) = R\SheafHom(L, \mathcal{O}_Y)
$$
et il en va de même après tout changement de base : étant donné
$$
\vcenter{
\xymatrix{
X' \ar[r]_{g'} \ar[d]_{f'} &
X \ar[d]^f \\
Y' \ar[r]^g &
Y
}
}
\quad\quad
\begin{matrix}
\text{cartésien, alors on a } \\
Rf'_*R\SheafHom(L(g')^*K, L(g')^*E) \\
= R\SheafHom(Lg^*L, \mathcal{O}_{Y'})
\end{matrix}
$$
\end{lemma}

\begin{proof}
Puisque $Y$ est quasi-compact et quasi-séparé, il en va de même de $X$.
D'après Catégories dérivées d'espaces, lemme
\ref{spaces-perfect-lemma-pseudo-coherent-hocolim}, on peut écrire
$K = \text{hocolim} K_n$, avec $K_n$ parfait et $K_n \to K$ induisant
un isomorphisme sur les troncations $\tau_{\geq -n}$. Soit $K_n^\vee$
le complexe parfait dual
(Cohomologie sur les sites, lemme \ref{sites-cohomology-lemma-dual-perfect-complex}).
On obtient un système projectif $\ldots \to K_3^\vee \to K_2^\vee \to K_1^\vee$
d'objets parfaits. D'après le lemme \ref{spaces-more-morphisms-lemma-perfect-relatively-perfect},
on voit que $K_n^\vee \otimes_{\mathcal{O}_X} E$ est $Y$-parfait.
On peut donc appliquer le lemme \ref{spaces-more-morphisms-lemma-derived-pushforward-rel-perfect}
à $K_n^\vee \otimes_{\mathcal{O}_X} E$, et on obtient un système projectif
$$
\ldots \to M_3 \to M_2 \to M_1
$$
de complexes parfaits sur $Y$, avec
$$
M_n = Rf_*(K_n^\vee \otimes_{\mathcal{O}_X}^\mathbf{L} E) =
Rf_*R\SheafHom(K_n, E)
$$
De plus, la formation de ces complexes commute à tout
changement de base, à savoir $Lg^*M_n =
Rf'_*((L(g')^*K_n)^\vee \otimes_{\mathcal{O}_{X'}}^\mathbf{L} L(g')^*E) =
Rf'_*R\SheafHom(L(g')^*K_n, L(g')^*E)$.

\medskip\noindent
Comme $K_n \to K$ induit un isomorphisme sur $\tau_{\geq -n}$, on voit que
$K_n \to K_{n + 1}$ induit un isomorphisme sur $\tau_{\geq -n}$.
Il s'ensuit que $K_{n + 1}^\vee \to K_n^\vee$
induit un isomorphisme sur $\tau_{\leq n}$, puisque
$K_n^\vee = R\SheafHom(K_n, \mathcal{O}_X)$.
Supposons que $E$ soit de Tor-amplitude dans $[a, b]$ en tant que complexe
de $f^{-1}\mathcal{O}_Y$-modules. Il en va alors de même après
tout changement de base, voir
Catégories dérivées d'espaces, lemme
\ref{spaces-perfect-lemma-tor-independence-and-tor-amplitude}.
On trouve que
$K_{n + 1}^\vee \otimes_{\mathcal{O}_X} E \to
K_n^\vee \otimes_{\mathcal{O}_X} E$
induit un isomorphisme sur $\tau_{\leq n + a}$,
et il en va de même après tout changement de base.
En appliquant le foncteur dérivé à droite $Rf_*$,
on conclut que les applications $M_{n + 1} \to M_n$
induisent des isomorphismes sur $\tau_{\leq n + a}$,
et il en va de même après tout changement de base.
Choisissons un triangle distingué
$$
M_{n + 1} \to M_n \to C_n \to M_{n + 1}[1]
$$
Prenons $y \in |Y|$. Choisissons un voisinage \'etale élémentaire
$(U, u) \to (Y, y)$ ; cela est possible d'après
Espaces décents, lemme
\ref{decent-spaces-lemma-decent-space-elementary-etale-neighbourhood}.
Prenons pour $Y'$ le spectre du corps résiduel en $u$.
En prenant l'image inverse, on voit que $C_n|_U \otimes_{\mathcal{O}_U}^\mathbf{L} \kappa(u)$
n'a de cohomologie non nulle qu'en degrés $\geq n + a$. D'après
Compléments d'algèbre, lemme
\ref{more-algebra-lemma-lift-perfect-from-residue-field},
on voit que le complexe parfait $C_n|_U$ est de Tor-amplitude dans
$[n + a, m_n]$ pour un certain entier $m_n$, quitte à rétrécir $U$.
Ainsi, $C_n$ est de Tor-amplitude dans $[n + a, m_n]$ pour un certain entier $m_n$
(car $Y$ est quasi-compact).
En particulier, le complexe parfait dual $C_n^\vee$ est de Tor-amplitude dans
$[-m_n, -n - a]$.

\medskip\noindent
Soit $L_n = M_n^\vee$ le complexe parfait dual. La
conclusion de la discussion de l'alinéa précédent est que
$L_n \to L_{n + 1}$ induit des isomorphismes sur $\tau_{\geq -n - a}$.
Ainsi, $L = \text{hocolim} L_n$ est pseudo-cohérent, voir
Catégories dérivées d'espaces, lemme
\ref{spaces-perfect-lemma-pseudo-coherent-hocolim}.
Puisque l'on a
$$
R\SheafHom(K, E) = R\SheafHom(\text{hocolim} K_n, E) =
R\lim R\SheafHom(K_n, E) = R\lim K_n^\vee \otimes_{\mathcal{O}_X} E
$$
(Cohomologie sur les sites, lemme
\ref{sites-cohomology-lemma-colim-and-lim-of-duals})
et puisque $R\lim$ commute à $Rf_*$, on trouve que
$$
Rf_*R\SheafHom(K, E) = R\lim M_n = R\lim R\SheafHom(L_n, \mathcal{O}_Y) =
R\SheafHom(L, \mathcal{O}_Y)
$$
Ceci démontre la formule sur $Y$. Puisque la construction de $M_n$ est
compatible avec le changement de base, la formule reste vraie après
tout changement de base.
\end{proof}

\begin{remark}
\label{spaces-more-morphisms-remark-compare-L}
Le lecteur aura peut-être remarqué la similitude entre le
lemme \ref{spaces-more-morphisms-lemma-compute-ext-rel-perfect} et
Catégories dérivées d'espaces, lemme \ref{spaces-perfect-lemma-compute-ext}.
En effet, le complexe pseudo-cohérent $L$ du
lemme \ref{spaces-more-morphisms-lemma-compute-ext-rel-perfect}
peut être caractérisé comme l'unique complexe pseudo-cohérent
sur $Y$ tel qu'il existe des isomorphismes fonctoriels
$$
\Ext^i_{\mathcal{O}_Y}(L, \mathcal{F}) \longrightarrow
\Ext^i_{\mathcal{O}_X}(K,
E \otimes_{\mathcal{O}_X}^\mathbf{L} Lf^*\mathcal{F})
$$
compatibles avec les homomorphismes cobord, lorsque $\mathcal{F}$ parcourt
$\QCoh(\mathcal{O}_Y)$. Si cela est un jour nécessaire, on
formulera ici un résultat précis et en donnera une démonstration détaillée.
\end{remark}

\begin{lemma}
\label{spaces-more-morphisms-lemma-bounded-on-fibres}
Soit $S$ un schéma. Soit $X$ un espace algébrique sur $S$ tel que
le morphisme structural $f : X \to S$ soit plat et
localement de présentation finie. Soit $E$ un objet pseudo-cohérent
de $D(\mathcal{O}_X)$. Les conditions suivantes sont équivalentes :
\begin{enumerate}
\item $E$ est $S$-parfait, et
\item $E$ est localement borné inférieurement et, pour tout point $s \in S$,
l'objet $L(X_s \to X)^*E$ de $D(\mathcal{O}_{X_s})$
est localement borné inférieurement.
\end{enumerate}
\end{lemma}

\begin{proof}
Comme tout est local, on se ramène immédiatement au
cas où $X$ et $S$ sont affines, voir le lemme
\ref{spaces-more-morphisms-lemma-affine-locally-rel-perfect}.
Ce cas est traité dans
Catégories dérivées de schémas, lemme \ref{perfect-lemma-bounded-on-fibres}.
\end{proof}

\begin{lemma}
\label{spaces-more-morphisms-lemma-characterize-relatively-perfect}
Soit $A$ un anneau. Soit $X$ un espace algébrique séparé, de
présentation finie et plat sur $A$. Soit $K \in D_\QCoh(\mathcal{O}_X)$.
Si $R \Gamma (X, E \otimes^\mathbf{L} K)$ est parfait dans
$D(A)$ pour tout objet parfait $E \in D(\mathcal{O}_X)$, alors $K$ est
$\Spec(A)$-parfait.
\end{lemma}

\begin{proof}
D'après le lemme \ref{spaces-more-morphisms-lemma-characterize-pseudo-coh-improved},
$K$ est pseudo-cohérent relativement à $A$. D'après le lemme
\ref{spaces-more-morphisms-lemma-relative-pseudo-coherent-is-moot},
$K$ est pseudo-cohérent dans $D(\mathcal{O}_X)$.
D'après Catégories dérivées d'espaces, proposition
\ref{spaces-perfect-proposition-detecting-bounded-below},
on voit que $K$ appartient à $D^-(\mathcal{O}_X)$.
Soit $\mathfrak{p}$ un idéal premier de $A$, et notons
$i : Y \to X$ l'inclusion de la fibre schématique
au-dessus de $\mathfrak{p}$, c'est-à-dire que $Y$ est un schéma sur $\kappa(\mathfrak p)$.
D'après le lemme \ref{spaces-more-morphisms-lemma-bounded-on-fibres},
il suffit de montrer que $Li^*(K)$ est borné inférieurement.
Soit $G \in D_{perf} (\mathcal{O}_X)$ un complexe
parfait qui engendre $D_\QCoh (\mathcal{O}_X)$,
voir Catégories dérivées d'espaces, théorème
\ref{spaces-perfect-theorem-bondal-van-den-Bergh}.
On a
\begin{align*}
R\Hom _{\mathcal{O}_Y}(Li^*(G), Li^*(K))
& =
R\Gamma(Y, Li^*(G ^\vee \otimes ^\mathbf{L} K)) \\
& =
R\Gamma(X, G^\vee \otimes ^{\mathbf{L}} K)
\otimes^\mathbf{L}_A \kappa(\mathfrak{p})
\end{align*}
La première égalité utilise le fait que $Li^*$ préserve les objets parfaits et les duaux,
ainsi que Cohomologie sur les sites, lemme
\ref{sites-cohomology-lemma-dual-perfect-complex} ; on omet
certains détails. La seconde égalité résulte de
Catégories dérivées d'espaces, lemme
\ref{spaces-perfect-lemma-compare-base-change},
puisque $X$ est plat sur $A$. Il résulte de notre hypothèse qu'il s'agit d'un
objet parfait de $D(\kappa(\mathfrak{p}))$. L'objet
$Li^*(G) \in D_{perf}(\mathcal{O}_Y)$ engendre $D_\QCoh(\mathcal{O}_Y)$ d'après
Catégories dérivées d'espaces, remarque
\ref{spaces-perfect-remark-pullback-generator}.
Par conséquent, Catégories dérivées d'espaces, proposition
\ref{spaces-perfect-proposition-detecting-bounded-below},
implique maintenant que $Li^*(K)$ est borné inférieurement, ce qui achève la démonstration.
\end{proof}














\section{Théorème du cube}
\label{spaces-more-morphisms-section-theorem-cube}

\noindent
Cette section est l'analogue de Compléments sur les morphismes, section
\ref{more-morphisms-section-theorem-cube}.
Le lemme suivant montre que la diagonale du foncteur de Picard
est représentable par des immersions localement fermées sous
les hypothèses du lemme.

\begin{lemma}
\label{spaces-more-morphisms-lemma-diagonal-picard-flat-proper}
Soit $S$ un schéma.
Soit $f : X \to Y$ un morphisme plat, propre et de présentation finie
d'espaces algébriques sur $S$.
Soit $\mathcal{E}$ un $\mathcal{O}_X$-module localement libre de type fini.
Pour un morphisme $g : Y' \to Y$, considérons le diagramme de changement de base
$$
\xymatrix{
X' \ar[d]_{f'} \ar[r]_{g'} & X \ar[d]^f \\
Y' \ar[r]^g & Y
}
$$
Supposons que $\mathcal{O}_{Y'} \to f'_*\mathcal{O}_{X'}$ soit un
isomorphisme pour tout $g : Y' \to Y$.
Alors il existe une immersion $j : Z \to Y$ de présentation finie
telle qu'un morphisme $g : Y' \to Y$ se factorise par $Z$ si et seulement s'il
existe un $\mathcal{O}_{Y'}$-module $\mathcal{N}$ localement libre de type fini
tel que $(f')^*\mathcal{N} \cong (g')^*\mathcal{L}$.
\end{lemma}

\begin{proof}
Soit $y : \Spec(k) \to Y$ un point à valeurs dans un corps.
Alors la fibre $X_y$ de $f$ en $y$ est connexe puisque, par hypothèse,
$H^0(X_y, \mathcal{O}_{X_y}) = k$. Ainsi, le rang de
$\mathcal{E}$ est constant sur les fibres. Puisque $f$ est ouvert
(Morphismes d'espaces, lemme \ref{spaces-morphisms-lemma-fppf-open})
et fermé, on conclut qu'il existe une décomposition $Y = \coprod Y_r$
de $Y$ en sous-espaces ouverts et fermés telle que $\mathcal{E}$
soit de rang constant $r$ sur l'image inverse de $Y_r$.
On peut donc supposer que $\mathcal{E}$ est de rang constant $r$.
On notera $\mathcal{E}^\vee = \SheafHom(\mathcal{E}, \mathcal{O}_X)$
le module dual de rang $r$.

\medskip\noindent
D'après le théorème de cohomologie et changement de base (plus précisément, d'après
Catégories dérivées d'espaces, lemme
\ref{spaces-perfect-lemma-flat-proper-perfect-direct-image-general}),
on voit que $E = Rf_*\mathcal{E}$ est un objet parfait de la
catégorie dérivée de $Y$ et que sa formation commute à
tout changement de base. De même pour $E' = Rf_*\mathcal{E}^\vee$.
Comme il n'y a jamais de cohomologie en degrés $< 0$, on voit que
$E$ et $E'$ sont (localement) de Tor-amplitude dans $[0, b]$ pour un certain $b$.
Observons que, pour tout $g : Y' \to Y$, on a
$f'_*((g')^*\mathcal{E}) = H^0(Lg^*E)$ et
$f'_*((g')^*\mathcal{E}^\vee) = H^0(Lg^*E')$.
Soient $j : Z \to Y$ et $j' : Z' \to Y$ les immersions localement
fermées construites dans Catégories dérivées d'espaces, lemme
\ref{spaces-perfect-lemma-locally-closed-where-H0-locally-free},
pour $E$ et $E'$, avec $a = 0$ et $r = r$ ; elles sont caractérisées
par le fait que $H^0(Lj^*E)$ et $H^0((j')^*E')$
sont des modules localement libres de rang $r$ dont la formation commute à l'image inverse.

\medskip\noindent
Soit $g : Y' \to Y$ un morphisme. S'il existe un module $\mathcal{N}$
comme dans le lemme, alors, en utilisant la formule de projection
Cohomologie sur les sites, lemme \ref{sites-cohomology-lemma-projection-formula},
on voit que les modules
$$
f'_*((g')^*\mathcal{E}) \cong
f'_*((f')^*\mathcal{N}) \cong
\mathcal{N} \otimes_{\mathcal{O}_{Y'}} f'_*\mathcal{O}_{X'} \cong
\mathcal{N}\quad\text{et de même }\quad
f'_*((g')^*\mathcal{E}^\vee) \cong \mathcal{N}^\vee
$$
sont localement libres de rang $r$ et restent localement libres de rang $r$
après tout changement de base ultérieur $Y'' \to Y'$.
Ainsi, dans ce cas, $g : Y' \to Y$ se factorise par $j$ et par $j'$.
On peut donc remplacer $Y$ par $Z \times_Y Z'$ et supposer que
$f_*\mathcal{E}$ et $f_*\mathcal{E}^\vee$ sont des
$\mathcal{O}_Y$-modules localement libres de rang $r$
dont la formation commute à tout changement de base.

\medskip\noindent
Dans cette situation, si $g : Y' \to Y$ est un morphisme et s'il existe un
module $\mathcal{N}$ comme dans le lemme, alors l'application (le cup-produit en degré $0$)
$$
f'_*((g')^*\mathcal{E})
\otimes_{\mathcal{O}_{Y'}}
f'_*((g')^*\mathcal{E}^\vee)
\longrightarrow \mathcal{O}_{Y'}
$$
est un accouplement parfait. Réciproquement, si cette application de cup-produit est un
accouplement parfait, on voit que, localement sur $Y'$, il existe
une base de sections
$\sigma_1, \ldots, \sigma_r$ dans $f'_*((g')^*\mathcal{L})$ et
$\tau_1, \ldots, \tau_r$ dans
$f'_*((g')^*\mathcal{E}^\vee)$ dont les produits vérifient
$\sigma_i \tau_j = \delta_{ij}$. En considérant $\sigma_i$
comme une section de $(g')^*\mathcal{L}$ sur $X'$
et $\tau_j$ comme une section de $(g')^*\mathcal{L}^\vee$ sur $X'$,
on conclut que
$$
\sigma_1, \ldots, \sigma_r :
\mathcal{O}_{X'}^{\oplus r}
\longrightarrow
(g')^*\mathcal{E}
$$
est un isomorphisme dont l'inverse est donné par
$$
\tau_1, \ldots, \tau_r :
(g')^*\mathcal{E}
\longrightarrow
\mathcal{O}_{X'}^{\oplus r}
$$
Autrement dit, on voit que
$(f')^*f'_*(g')^*\mathcal{E} \cong (g')^*\mathcal{E}$.
Mais la condition que le cup-produit soit non dégénéré
définit un sous-schéma ouvert rétrocompact (à savoir le lieu où un déterminant
convenable est non nul), ce qui achève la démonstration.
\end{proof}









\section{Descente des propriétés de finitude des complexes}
\label{spaces-more-morphisms-section-descent-finiteness}

\noindent
Cette section est l'analogue de Compléments sur les morphismes,
section \ref{more-morphisms-section-descent-finiteness},
et de
Catégories dérivées de schémas, section
\ref{perfect-section-descent-finiteness}.

\begin{lemma}
\label{spaces-more-morphisms-lemma-pseudo-coherent-descends-fpqc}
Soit $S$ un schéma. Soit $\{f_i : X_i \to X\}$ un recouvrement fpqc
d'espaces algébriques sur $S$. Soit $E \in D_\QCoh(\mathcal{O}_X)$.
Soit $m \in \mathbf{Z}$. Alors $E$ est $m$-pseudo-cohérent si et seulement si chaque
$Lf_i^*E$ est $m$-pseudo-cohérent.
\end{lemma}

\begin{proof}
L'image inverse préserve toujours la $m$-pseudo-cohérence, voir
Cohomologie sur les sites, lemme
\ref{sites-cohomology-lemma-pseudo-coherent-pullback}.
Il suffit donc de supposer que $Lf_i^*E$ est $m$-pseudo-cohérent
et de démontrer que $E$ est $m$-pseudo-cohérent.
On peut tout d'abord supposer que $X_i$ est un schéma pour tout $i$, voir
Topologies sur les espaces, lemme \ref{spaces-topologies-lemma-refine-fpqc-schemes}.
Ensuite, choisissons un morphisme \'etale surjectif $U \to X$, où $U$ est un schéma.
Alors $U_i = U \times_X X_i$ est un schéma et nous obtenons un recouvrement fpqc
$\{U_i \to U\}$ de schémas, voir
Topologies sur les espaces, Lemme
\ref{spaces-topologies-lemma-recognize-fpqc-covering}.
Nous savons que le résultat est vrai pour
$\{U_i \to U\}_{i \in I}$ d'après le cas des schémas, voir
Catégories dérivées des schémas, Lemme
\ref{perfect-lemma-pseudo-coherent-descends-fpqc}.
D'autre part, la restriction $E|_U$ provient
d'un objet de $D_\QCoh(\mathcal{O}_U)$ (défini au moyen de la topologie de Zariski
et du faisceau structural ``usuel'' de $U$), voir
Catégories dérivées des espaces, Lemme
\ref{spaces-perfect-lemma-derived-quasi-coherent-small-etale-site}.
Le lemme en résulte, puisque les deux notions de pseudo-cohérence
(\'etale et de Zariski) coïncident d'après
Catégories dérivées des espaces,
Lemme \ref{spaces-perfect-lemma-descend-pseudo-coherent}.
\end{proof}

\begin{lemma}
\label{spaces-more-morphisms-lemma-tor-amplitude-descends-fppf}
Soit $S$ un schéma. Soit $\{g_i : Y_i \to Y\}$ un recouvrement fpqc par des
espaces algébriques sur $S$. Soit $f : X \to Y$ un morphisme d'espaces algébriques
et posons $X_i = Y_i \times_Y X$, avec les projections $f_i : X_i \to Y_i$
et $g'_i : X_i \to X$. Soit $E \in D_\QCoh(\mathcal{O}_X)$.
Soient $a, b \in \mathbf{Z}$. Alors les conditions suivantes sont équivalentes :
\begin{enumerate}
\item $E$ est d'amplitude de Tor dans $[a, b]$ en tant qu'objet de
$D(f^{-1}\mathcal{O}_Y)$, et
\item $L(g'_i)^*E$ est d'amplitude de Tor dans $[a, b]$ en tant qu'objet de
$D(f_i^{-1}\mathcal{O}_{Y_i})$ pour tout $i$.
\end{enumerate}
L'énoncé reste vrai si ``d'amplitude de Tor dans $[a, b]$'' est remplacé par
``de dimension de Tor localement finie''.
\end{lemma}

\begin{proof}
L'image réciproque préserve la propriété ``être d'amplitude de Tor dans $[a, b]$'' d'après
Catégories dérivées des espaces, Lemme
\ref{spaces-perfect-lemma-tor-independence-and-tor-amplitude}.
Remarquons que $Y_i$ et $X$ sont Tor-indépendants au-dessus de $Y$
puisque $Y_i \to Y$ est plat. Supposons (2) et démontrons (1).
Nous pouvons calculer la dimension de Tor sur les fibres, voir
Cohomologie sur les sites, Lemme \ref{sites-cohomology-lemma-tor-amplitude-stalk}
et Propriétés des espaces, Théorème
\ref{spaces-properties-theorem-exactness-stalks}.
Soit $\overline{x}$ un point géométrique de $X$. Choisissons
un $i$ et un point géométrique $\overline{x}_i$ de $X_i$ d'image
$\overline{x}$ dans $X$. Alors
$$
(L(g_i')^*E)_{\overline{x}_i} =
E_{\overline{x}}
\otimes_{\mathcal{O}_{X, \overline{x}}}^\mathbf{L}
\mathcal{O}_{X_, \overline{x}_i}
$$
Soient $\overline{y}_i$ dans $Y_i$ et $\overline{y}$ dans $Y$
les images de $\overline{x}_i$ et $\overline{x}$.
Puisque $X$ et $Y_i$ sont Tor-indépendants au-dessus de $Y$, nous pouvons appliquer
Compléments d'algèbre, Lemme \ref{more-algebra-lemma-base-change-comparison}
pour voir que le membre de droite de la formule affichée est égal à
$E_{\overline{x}}
\otimes_{\mathcal{O}_{Y, \overline{y}}}^\mathbf{L}
\mathcal{O}_{Y_i, \overline{y}_i}$
dans $D(\mathcal{O}_{Y_i, \overline{y}_i})$.
Puisque nous avons supposé que son amplitude de Tor est contenue dans
$[a, b]$, nous concluons que l'amplitude de Tor de
$E_{\overline{x}}$ dans $D(\mathcal{O}_{Y, \overline{y}})$
est contenue dans $[a, b]$ d'après Compléments d'algèbre, Lemme
\ref{more-algebra-lemma-flat-descent-tor-amplitude}.
Ainsi (1) en résulte.

\medskip\noindent
Un peu de topologie élémentaire donne aussi le cas
``de dimension de Tor localement finie''.
\end{proof}

\noindent
Les lemmes suivants ne trouvent pas vraiment leur place dans cette section.

\begin{lemma}
\label{spaces-more-morphisms-lemma-thickening-pseudo-coherent}
Soit $S$ un schéma. Soit $i : X \to X'$ un épaississement d'ordre fini entre
espaces algébriques. Soit $K' \in D(\mathcal{O}_{X'})$ un objet tel que
$K = Li^*K'$ soit pseudo-cohérent. Alors $K'$ est pseudo-cohérent.
\end{lemma}

\begin{proof}
Nous démontrons d'abord que $K'$ a des faisceaux de cohomologie quasi-cohérents ; nous invitons
le lecteur à omettre cette partie.
Pour cela, nous pouvons nous ramener au cas d'un épaississement du premier ordre, voir
section \ref{spaces-more-morphisms-section-thickenings}. Soit $\mathcal{I} \subset \mathcal{O}_{X'}$
le faisceau quasi-cohérent d'idéaux définissant $X$.
En tensorisant la suite exacte courte
$$
0 \to \mathcal{I} \to \mathcal{O}_{X'} \to i_*\mathcal{O}_X \to 0
$$
par $K'$, nous obtenons un triangle distingué
$$
K' \otimes_{\mathcal{O}_{X'}}^\mathbf{L} \mathcal{I}
\to K' \to
K' \otimes_{\mathcal{O}_{X'}}^\mathbf{L} i_*\mathcal{O}_X
\to
(K' \otimes_{\mathcal{O}_{X'}}^\mathbf{L} \mathcal{I})[1]
$$
Puisque $i_* = Ri_*$ et que nous pouvons considérer $\mathcal{I}$
comme un $\mathcal{O}_X$-module quasi-cohérent (puisque nous avons un épaississement du
premier ordre), nous pouvons récrire ceci sous la forme
$$
i_*(K \otimes_{\mathcal{O}_X}^\mathbf{L} \mathcal{I})
\to K' \to
i_*K \to
i_*(K \otimes_{\mathcal{O}_X}^\mathbf{L} \mathcal{I})[1]
$$
Utilisons Cohomologie des espaces, Lemme
\ref{spaces-cohomology-lemma-projection-formula-finite}
pour identifier les termes. Puisque $K$ appartient à
$D_\QCoh(\mathcal{O}_X)$, nous concluons que
$K'$ appartient à $D_\QCoh(\mathcal{O}_{X'})$ ; ceci utilise
Catégories dérivées des espaces, Lemmes
\ref{spaces-perfect-lemma-pseudo-coherent},
\ref{spaces-perfect-lemma-quasi-coherence-tensor-product}, et
\ref{spaces-perfect-lemma-quasi-coherence-direct-image}.

\medskip\noindent
Supposons que $K'$ appartienne à $D_\QCoh(\mathcal{O}_{X'})$.
La question est locale pour la topologie \'etale sur $X'$ ;
nous pouvons donc supposer que $X'$ est affine.
Dans ce cas, le résultat découle du cas des schémas
(Compléments sur les morphismes, Lemme
\ref{more-morphisms-lemma-thickening-pseudo-coherent}).
Le passage au langage des schémas utilise
Catégories dérivées des espaces, Lemmes
\ref{spaces-perfect-lemma-derived-quasi-coherent-small-etale-site} et
\ref{spaces-perfect-lemma-descend-pseudo-coherent} et
Remarque \ref{spaces-perfect-remark-match-total-direct-images}.
\end{proof}

\begin{lemma}
\label{spaces-more-morphisms-lemma-thickening-relatively-perfect}
Soit $S$ un schéma. Considérons un diagramme cartésien
$$
\xymatrix{
X \ar[r]_i \ar[d]_f & X' \ar[d]^{f'} \\
Y \ar[r]^j & Y'
}
$$
d'espaces algébriques sur $S$. Supposons que $X' \to Y'$ soit plat et localement
de présentation finie, et que $Y \to Y'$ soit un épaississement d'ordre fini.
Soit $E' \in D(\mathcal{O}_{X'})$. Si $E = Li^*(E')$ est $Y$-parfait,
alors $E'$ est $Y'$-parfait.
\end{lemma}

\begin{proof}
Rappelons qu'être $Y$-parfait signifie, pour $E$, que $E$ est
pseudo-cohérent et localement de dimension de Tor finie en tant que complexe
de $f^{-1}\mathcal{O}_Y$-modules
(Définition \ref{spaces-more-morphisms-definition-relatively-perfect}).
D'après le Lemme \ref{spaces-more-morphisms-lemma-thickening-pseudo-coherent},
nous voyons que $E'$ est pseudo-cohérent.
En particulier, $E'$ appartient à $D_\QCoh(\mathcal{O}_{X'})$, voir
Catégories dérivées des espaces, Lemme
\ref{spaces-perfect-lemma-pseudo-coherent}.
Le Lemme \ref{spaces-more-morphisms-lemma-affine-locally-rel-perfect}
nous ramène au cas des schémas.
Le cas des schémas est
Compléments sur les morphismes, Lemme
\ref{more-morphisms-lemma-thickening-relatively-perfect}.
\end{proof}

\begin{lemma}
\label{spaces-more-morphisms-lemma-henselian-relatively-perfect}
Soit $(R, I)$ un couple formé d'un anneau et d'un idéal $I$
contenu dans le radical de Jacobson. Posons $S = \Spec(R)$ et $S_0 = \Spec(R/I)$.
Soit $X$ un espace algébrique sur $R$ dont le morphisme structural
$f : X \to S$ est propre, plat et de présentation finie.
Notons $X_0 = S_0 \times_S X$. Soit $E \in D(\mathcal{O}_X)$
pseudo-cohérent. Si la restriction dérivée $E_0$ de $E$
à $X_0$ est $S_0$-parfaite, alors $E$ est $S$-parfait.
\end{lemma}

\begin{proof}
Choisissons un morphisme \'etale surjectif $U \to X$ avec $U$ affine.
Choisissons une immersion fermée $U \to \mathbf{A}^d_S$.
Posons $U_0 = S_0 \times_S U$.
Le complexe $E_0|_{U_0}$ est d'amplitude de Tor
dans $[a, b]$ pour certains $a, b \in \mathbf{Z}$.
Soit $\overline{x}$ un point géométrique de $X$.
Nous allons montrer que l'amplitude de Tor de
$E_{\overline{x}}$ sur $R$ est contenue dans $[a - d, b]$.
Cela achèvera la démonstration, puisque l'amplitude de Tor peut se
lire sur les fibres d'après
Cohomologie sur les sites, Lemme \ref{sites-cohomology-lemma-tor-amplitude-stalk}
et Propriétés des espaces, Théorème
\ref{spaces-properties-theorem-exactness-stalks}.

\medskip\noindent
Soit $x \in |X|$ le point déterminé par $\overline{x}$.
Rappelons que $|X| \to |S|$ est fermé (par définition des morphismes propres).
Puisque $I$ est contenu dans le radical de Jacobson, toute partie fermée non vide
de $S$ contient un point du sous-schéma fermé $S_0$.
Nous pouvons donc trouver une spécialisation $x \leadsto x_0$ dans $|X|$
avec $x_0 \in |X_0|$. Choisissons $u_0 \in U_0$ s'envoyant sur $x_0$.
D'après Espaces décents, Lemme \ref{decent-spaces-lemma-generalizations-lift-flat}
(ou Espaces décents, Lemme \ref{decent-spaces-lemma-specialization},
qui s'applique directement aux morphismes \'etales),
nous trouvons une spécialisation $u \leadsto u_0$ dans $U$
telle que $u$ s'envoie sur $x$. Nous pouvons relever $\overline{x}$
en un point géométrique $\overline{u}$ de $U$ situé au-dessus de $u$.
Alors $E_{\overline{x}} = (E|_U)_{\overline{u}}$.

\medskip\noindent
Écrivons $U = \Spec(A)$. Alors $A$ est plate et de présentation finie
comme $R$-algèbre, et c'est un quotient d'une $R$-algèbre polynomiale en
$d$ variables. La restriction $E|_U$ correspond
(d'après Catégories dérivées des espaces, Lemmes
\ref{spaces-perfect-lemma-pseudo-coherent},
\ref{spaces-perfect-lemma-derived-quasi-coherent-small-etale-site}, et
\ref{spaces-perfect-lemma-descend-pseudo-coherent},
et
Catégories dérivées des schémas, Lemme
\ref{perfect-lemma-affine-compare-bounded} et
\ref{perfect-lemma-pseudo-coherent-affine})
à un objet pseudo-cohérent $K$ de $D(A)$.
Remarquons que $E_0$ correspond à $K \otimes_A^\mathbf{L} A/IA$.
Soient $\mathfrak q \subset \mathfrak q_0 \subset A$ les idéaux premiers
correspondant à $u \leadsto u_0$.
Alors
$$
E_{\overline{x}} =
(E|_U)_{\overline{u}} =
E_u \otimes_{\mathcal{O}_{U, u}}^\mathbf{L} \mathcal{O}_{U, \overline{u}} =
K_{\mathfrak q} \otimes_{A_\mathfrak q}^\mathbf{L} A_{\mathfrak q}^{sh}
$$
(certains détails sont omis). Puisque $A_\mathfrak q \to A_\mathfrak q^{sh}$
est plat, l'amplitude de Tor de cet objet comme $R$-module est la même que
l'amplitude de Tor de $K_\mathfrak q$ comme $R$-module
(Compléments d'algèbre, Lemme \ref{more-algebra-lemma-no-change-tor-amplitude}).
De plus, $K_{\mathfrak q}$ est un localisé de $K_{\mathfrak q_0}$.
Il suffit donc de montrer que $K_{\mathfrak q_0}$ est d'amplitude de Tor dans
$[a - d, b]$ en tant que complexe de $R$-modules.

\medskip\noindent
Soit $I \subset \mathfrak p_0 \subset R$ l'idéal premier
correspondant à $f(x_0)$. Alors nous avons
\begin{align*}
K \otimes_R^\mathbf{L} \kappa(\mathfrak p_0)
& =
(K \otimes_R^\mathbf{L} R/I) \otimes_{R/I}^\mathbf{L}
\kappa(\mathfrak p_0) \\
& =
(K \otimes_A^\mathbf{L} A/IA) \otimes_{R/I}^\mathbf{L} \kappa(\mathfrak p_0)
\end{align*}
la seconde égalité résultant de ce que $R \to A$ est plat.
Par notre choix de $a, b$, ce complexe n'a de cohomologie
que dans les degrés de l'intervalle $[a, b]$.
Nous pouvons donc finalement appliquer
Compléments d'algèbre, Lemme
\ref{more-algebra-lemma-lift-from-fibre-relatively-perfect}
à $R \to A$, $\mathfrak q_0$, $\mathfrak p_0$ et $K$
pour conclure.
\end{proof}







\section{Familles de courbes nodales}
\label{spaces-more-morphisms-section-families-nodal}

\noindent
Cette section prolonge
Courbes algébriques, section \ref{curves-section-families-nodal}.
Nous renvoyons aussi à cette section pour notre choix
de terminologie.

\medskip\noindent
La propriété ``à singularités au plus nodales de dimension relative $1$''
des morphismes de schémas est locale pour la topologie \'etale sur la source et le but, voir
Descente, Lemme \ref{descent-lemma-etale-local-source-target}
et
Courbes algébriques, Lemmes \ref{curves-lemma-nodal-family-etale-local-source},
\ref{curves-lemma-nodal-family-fpqc-local-target}, et
\ref{curves-lemma-nodal-family-postcompose-etale}.
Elle est aussi stable par changement de base et locale pour la topologie fpqc sur le but, voir
Courbes algébriques, Lemmes \ref{curves-lemma-base-change-nodal-family} et
\ref{curves-lemma-nodal-family-fpqc-local-target}. Par conséquent, d'après
Morphismes d'espaces, Lemme \ref{spaces-morphisms-lemma-local-source-target},
nous pouvons définir comme suit la notion de morphisme à singularités au plus nodales de dimension relative
$1$ pour les espaces algébriques ; elle coïncide avec la
notion déjà existante définie dans
Morphismes d'espaces, section \ref{spaces-morphisms-section-representable},
lorsque le morphisme est représentable.

\begin{definition}
\label{spaces-more-morphisms-definition-nodal-family}
Soit $S$ un schéma.
Soit $f : X \to Y$ un morphisme d'espaces algébriques sur $S$.
Nous disons que $f$ est {\it à singularités au plus nodales de dimension relative $1$}
si les conditions équivalentes de
Morphismes d'espaces, Lemme \ref{spaces-morphisms-lemma-local-source-target},
sont satisfaites avec $\mathcal{P} =$``à singularités au plus nodales de dimension relative $1$''.
\end{definition}

\begin{lemma}
\label{spaces-more-morphisms-lemma-base-change-nodal}
La propriété d'être à singularités au plus nodales de dimension relative $1$
est préservée par changement de base.
\end{lemma}

\begin{proof}
Voir
Morphismes d'espaces, Remarque \ref{spaces-morphisms-remark-base-change-P}
et
Courbes algébriques, Lemme \ref{curves-lemma-base-change-nodal-family}.
\end{proof}

\begin{lemma}
\label{spaces-more-morphisms-lemma-nodal-local}
Soit $S$ un schéma.
Soit $f : X \to Y$ un morphisme d'espaces algébriques sur $S$.
Les conditions suivantes sont équivalentes :
\begin{enumerate}
\item $f$ est à singularités au plus nodales de dimension relative $1$,
\item pour tout schéma $Z$ et tout morphisme $Z \to Y$, le morphisme
$Z \times_Y X \to Z$ est à singularités au plus nodales de dimension relative $1$,
\item pour tout schéma affine $Z$ et tout morphisme
$Z \to Y$, le morphisme $Z \times_Y X \to Z$ est
à singularités au plus nodales de dimension relative $1$,
\item il existe un schéma $V$ et un morphisme \'etale surjectif
$V \to Y$ tels que $V \times_Y X \to V$ soit
à singularités au plus nodales de dimension relative $1$,
\item il existe un schéma $U$ et un morphisme \'etale surjectif
$\varphi : U \to X$ tels que le composé $f \circ \varphi$
soit à singularités au plus nodales de dimension relative $1$,
\item pour tout diagramme commutatif
$$
\xymatrix{
U \ar[d] \ar[r] & V \ar[d] \\
X \ar[r] & Y
}
$$
où $U$, $V$ sont des schémas et les flèches verticales sont \'etales,
la flèche horizontale supérieure est à singularités au plus nodales de dimension relative $1$,
\item il existe un diagramme commutatif
$$
\xymatrix{
U \ar[d] \ar[r] & V \ar[d] \\
X \ar[r] & Y
}
$$
où $U$, $V$ sont des schémas, les flèches verticales sont \'etales et
$U \to X$ est surjectif, tel que la flèche horizontale supérieure soit
à singularités au plus nodales de dimension relative $1$, et
\item il existe des recouvrements de Zariski $Y = \bigcup_{i \in I} Y_i$
et $f^{-1}(Y_i) = \bigcup X_{ij}$ tels que
chaque morphisme $X_{ij} \to Y_i$ soit
à singularités au plus nodales de dimension relative $1$.
\end{enumerate}
\end{lemma}

\begin{proof}
La démonstration est omise.
\end{proof}

\noindent
Le lemme suivant nous dit que nous pouvons vérifier fibre par fibre si un morphisme
est à singularités au plus nodales de dimension relative $1$.

\begin{lemma}
\label{spaces-more-morphisms-lemma-locus-where-nodal}
Soit $S$ un schéma.
Soit $f : X \to Y$ un morphisme d'espaces algébriques sur $S$,
plat et localement de présentation finie. Il existe alors
un plus grand sous-espace ouvert $X' \subset X$ tel que $f|_{X'} : X' \to Y$
soit à singularités au plus nodales de dimension relative $1$. De plus, la formation
de $X'$ commute à tout changement de base.
\end{lemma}

\begin{proof}
Choisissons un diagramme commutatif
$$
\xymatrix{
U \ar[d] \ar[r]_h & V \ar[d] \\
X \ar[r]^f & Y
}
$$
où $U$, $V$ sont des schémas, les flèches verticales sont \'etales et
$U \to X$ est surjectif. D'après le lemme pour le cas des schémas
(Courbes algébriques, Lemme \ref{curves-lemma-locus-where-nodal}),
nous trouvons un plus grand sous-schéma ouvert $U' \subset U$
tel que $h|_{U'} : U' \to V$ soit à singularités au plus nodales de dimension relative $1$
et que la formation de $U'$ commute au changement de base.
Soit $X' \subset X$ le sous-espace ouvert dont les points correspondent
à l'ouvert $\Im(|U'| \to |X|)$.
Le Lemme \ref{spaces-more-morphisms-lemma-nodal-local} montre que $X' \to Y$ est
à singularités au plus nodales de dimension relative $1$ et que $X'$ est le
plus grand sous-espace ouvert possédant cette propriété (ceci implique aussi
que $U'$ est l'image réciproque de $X'$ dans $U$, mais nous n'en avons
pas besoin). Puisque la même assertion vaut après changement de base,
la démonstration est achevée.
\end{proof}




\section{La propriété de résolution}
\label{spaces-more-morphisms-section-resolution-property}

\noindent
Nous poursuivons la discussion de Catégories dérivées des espaces, section
\ref{spaces-perfect-section-resolution-property}.

\begin{situation}
\label{spaces-more-morphisms-situation-descend-resolution-property}
Soit $S$ un schéma. Soit $X$ un espace algébrique quasi-compact et quasi-séparé
sur $S$. Soit $V \to X$ un morphisme \'etale surjectif,
où $V$ est un schéma affine (un tel morphisme existe d'après
Propriétés des espaces, Lemme
\ref{spaces-properties-lemma-quasi-compact-affine-cover}).
Choisissons un diagramme commutatif
$$
\xymatrix{
V \ar[rd]_\varphi \ar[rr]_j & & Y \ar[ld]^\pi \\
& X
}
$$
où $j$ est une immersion ouverte et $\pi$ un morphisme fini
d'espaces algébriques (un tel diagramme existe d'après
le Lemme \ref{spaces-more-morphisms-lemma-quasi-finite-separated-pass-through-finite}).
Soit $\mathcal{I} \subset \mathcal{O}_Y$ un faisceau quasi-cohérent d'idéaux
de type fini sur $Y$ tel que
$V(\mathcal{I}) = Y \setminus j(V)$ (un tel faisceau d'idéaux existe
d'après Limites d'espaces, Lemme
\ref{spaces-limits-lemma-quasi-coherent-finite-type-ideals}).
\end{situation}

\begin{lemma}
\label{spaces-more-morphisms-lemma-surjection-in-Noetherian-case}
Dans la Situation \ref{spaces-more-morphisms-situation-descend-resolution-property}, supposons que
$X$ soit noethérien. Alors, pour tout $\mathcal{O}_X$-module cohérent
$\mathcal{F}$, il existe un entier $r \geq 0$,
des entiers $n_1, \ldots, n_r \geq 0$ et une surjection
$$
\bigoplus\nolimits_{i = 1, \ldots, r} \pi_*(\mathcal{I}^{n_i})
\longrightarrow \mathcal{F}
$$
de $\mathcal{O}_X$-modules.
\end{lemma}

\begin{proof}
Notons $\omega_{Y/X}$ le $\mathcal{O}_Y$-module cohérent tel qu'il existe
un isomorphisme
$$
\pi_*\omega_{Y/X} \cong
\SheafHom_{\mathcal{O}_X}(\pi_*\mathcal{O}_Y, \mathcal{O}_X)
$$
de $\pi_*\mathcal{O}_Y$-modules, voir Morphismes d'espaces, lemme
\ref{spaces-morphisms-lemma-affine-equivalence-modules} et
Descente sur les espaces, lemme \ref{spaces-descent-lemma-finite-over-finite-module}.
Le morphisme canonique $\mathcal{O}_X \to \pi_*\mathcal{O}_Y$
induit un morphisme canonique
$$
\text{Tr}_\pi : \pi_*\omega_{Y/X} \longrightarrow \mathcal{O}_X
$$
Comme $V$ est noethérien affine, nous pouvons choisir des sections
$$
s_1, \ldots, s_r \in
\Gamma(V, \pi^*\mathcal{F} \otimes_{\mathcal{O}_Y} \omega_{Y/X})
$$
qui engendrent le module cohérent
$\pi^*\mathcal{F} \otimes_{\mathcal{O}_X} \omega_{Y/X}$
sur $V$. D'après
Cohomologie des espaces, lemme \ref{spaces-cohomology-lemma-homs-over-open},
nous pouvons choisir des entiers $n_i \geq 0$ tels que $s_i$
se prolonge en un morphisme
$s_i' : \mathcal{I}^{n_i} \to
\pi^*\mathcal{F} \otimes_{\mathcal{O}_Y} \omega_{Y/X}$.
En prenant les images directes sur $X$, nous obtenons des morphismes
$$
\sigma_i :
\pi_*\mathcal{I}^{n_i}
\xrightarrow{\pi_*s'_i}
\pi_*(\pi^*\mathcal{F} \otimes_{\mathcal{O}_Y} \omega_{Y/X}) =
\mathcal{F} \otimes_{\mathcal{O}_X} \pi_*\omega_{Y/X}
\xrightarrow{\text{Tr}_\pi} \mathcal{F}
$$
où l'égalité résulte de
Cohomologie des espaces, lemme \ref{spaces-cohomology-lemma-finite-rings}.
Pour achever la démonstration, nous allons montrer que la somme de ces morphismes est surjective.

\medskip\noindent
Soit $x \in |X|$ un point de $X$. Soit $v \in |V|$ un point d'image $x$.
Nous pouvons choisir un voisinage \'etale $(U, u) \to (X, x)$ tel que
$$
U \times_X Y = W \coprod W'
$$
(somme disjointe d'espaces algébriques), où $W \to U$ est un
isomorphisme et où l'unique point $w \in W$ au-dessus de
$u$ s'envoie sur $v$ dans $V \subset Y$. Cela résulte du
lemme \ref{spaces-more-morphisms-lemma-etale-splits-off-just-one-quasi-finite-part} et de
Morphismes \'etales, lemme \ref{etale-lemma-etale-etale-local}.
Après avoir encore rétréci $U$ si nécessaire, nous pouvons supposer que la projection envoie $W$ dans
$V \subset Y$. Comme la formation de $\omega_{Y/X}$
commute à la localisation \'etale, nous voyons que
$$
\pi_*\omega_{Y/X}|_U = (\pi|_W)_*\omega_{W/U} \oplus (\pi|_{W'})_*\omega_{W'/U}
$$
Nous avons $(\pi|_W)_*\omega_{W/U} = \mathcal{O}_U$, et cet isomorphisme
est donné par l'application trace $\text{Tr}_\pi|_U$ restreinte
au premier facteur de la décomposition ci-dessus. Comme $W$ s'envoie
dans $V$, nous voyons que $\mathcal{I}^{n_i}|_W = \mathcal{O}_W$.
Par conséquent
$$
\pi_*(\mathcal{I}^{n_i})|_U = \mathcal{O}_U \oplus
(W' \to U)_*(\mathcal{I}^{n_i}|_{W'})
$$
Une chasse au diagramme montre (détails omis) que $\sigma_i|_U$
sur le facteur $\mathcal{O}_U$ est le morphisme $\mathcal{O}_U \to \mathcal{F}$
correspondant à la section
$$
s_i|_W  \in
\Gamma(W,\pi^*\mathcal{F} \otimes_{\mathcal{O}_Y} \omega_{Y/X})
=
\Gamma(W, \mathcal{F}|_W \otimes_{\mathcal{O}_W} \omega_{W/U})
=
\Gamma(U, \mathcal{F})
$$
Comme les sections $s_i$ engendrent le module
$\pi^*\mathcal{F} \otimes_{\mathcal{O}_Y} \omega_{Y/X}$
sur $V$ et comme $W$ s'envoie dans $V$, nous concluons
que la restriction de $\bigoplus \sigma_i$ à $U$ est surjective.
Comme $x$ était arbitraire, la démonstration est achevée.
\end{proof}

\begin{lemma}
\label{spaces-more-morphisms-lemma-resolution-property-in-Noetherian-case}
Dans la situation \ref{spaces-more-morphisms-situation-descend-resolution-property}, supposons
$X$ noethérien. Alors $X$ possède la propriété de résolution si et seulement
si $\pi_*\mathcal{I}$ est quotient d'un module localement libre de type fini
sur $\mathcal{O}_X$.
\end{lemma}

\begin{proof}
Le module $\pi_*\mathcal{I}$ est cohérent d'après
Cohomologie des espaces, lemme
\ref{spaces-cohomology-lemma-finite-pushforward-coherent}.
Par conséquent, si $X$ possède la propriété de résolution, alors $\pi_*\mathcal{I}$
est quotient d'un $\mathcal{O}_X$-module localement libre de type fini.
Réciproquement, supposons donnée une surjection $\mathcal{E} \to \pi_*\mathcal{I}$
pour un $\mathcal{O}_X$-module localement libre de type fini $\mathcal{E}$.
Remarquons que, pour tout $n \geq 1$, il existe une surjection
$$
\pi_*\mathcal{I} \otimes_{\mathcal{O}_X} \pi_*\mathcal{I}^n
\longrightarrow
\pi_*\mathcal{I}^{n + 1}
$$
Ainsi, $\mathcal{E}^{\otimes n}$ s'envoie surjectivement sur
$\pi_*\mathcal{I}^n$ pour tout $n \geq 1$. Nous concluons que $X$
possède la propriété de résolution en combinant ceci avec le
résultat du lemme \ref{spaces-more-morphisms-lemma-surjection-in-Noetherian-case}.
\end{proof}

\begin{lemma}
\label{spaces-more-morphisms-lemma-resolution-property-one-sheaf}
Dans la situation \ref{spaces-more-morphisms-situation-descend-resolution-property},
l'espace algébrique $X$ possède la propriété de résolution
si et seulement si $\pi_*\mathcal{I}$ est quotient d'un
$\mathcal{O}_X$-module localement libre de type fini.
\end{lemma}

\begin{proof}
L'image directe $\pi_*\mathcal{G}$ d'un module
quasi-cohérent de type fini sur $\mathcal{O}_Y$, noté $\mathcal{G}$, est un module
quasi-cohérent de type fini sur $\mathcal{O}_X$ d'après
Descente sur les espaces, lemme
\ref{spaces-descent-lemma-finite-over-finite-module}.
En particulier, si $X$ possède la propriété de résolution, alors $\pi_*\mathcal{I}$
est quotient d'un $\mathcal{O}_X$-module localement libre de type fini d'après
Catégories dérivées des espaces, définition
\ref{spaces-perfect-definition-resolution-property}.

\medskip\noindent
Supposons que nous ayons une surjection $\mathcal{E} \to \pi_*\mathcal{I}$
pour un $\mathcal{O}_X$-module localement libre de type fini $\mathcal{E}$.
Dans la suite de la démonstration, nous montrons que $X$ possède la propriété de résolution
en nous ramenant au cas noethérien traité dans le
lemme \ref{spaces-more-morphisms-lemma-resolution-property-in-Noetherian-case}.
Nous suggérons au lecteur d'omettre la suite de la démonstration.

\medskip\noindent
Une première réduction consiste à considérer $X$ comme un espace algébrique sur
$\Spec(\mathbf{Z})$, voir
Espaces, définition \ref{spaces-definition-base-change}.
(Cela n'affecte ni les conditions ni la conclusion du lemme.)

\medskip\noindent
D'après Limites d'espaces, lemme
\ref{spaces-limits-lemma-finite-in-finite-and-finite-presentation},
nous pouvons écrire $Y = \lim Y_i$, où $Y_i$ est fini et de
présentation finie sur $X$ et où les morphismes de transition sont des
immersions fermées. Considérons le sous-espace fermé $Z = V(\mathcal{I})$ de $Y$.
Comme $\mathcal{I}$ est de type fini, le morphisme $Z \to Y$ est de
présentation finie. Nous pouvons donc trouver un $i$ et un morphisme
$Z_i \to Y_i$ de présentation finie dont le changement de base à $Y$
est $Z \to Y$, voir Limites d'espaces, lemme
\ref{spaces-limits-lemma-descend-finite-presentation}.
Pour $i' \geq i$, notons $Z_{i'} = Z_i \times_{Y_i} Y_{i'}$.
Après avoir augmenté $i$, nous pouvons supposer que $Z_i \to Y_i$ est une
immersion fermée (de présentation finie), voir Limites d'espaces, lemme
\ref{spaces-limits-lemma-descend-closed-immersion}. Notons
$\mathcal{I}_i \subset \mathcal{O}_{Y_i}$ le faisceau d'idéaux
de $Z_i$ et $\pi_i : Y_i \to X$ le morphisme structural.
De même pour $i' \geq i$. Comme $Z = \lim_{i' \geq i} Z_{i'}$, nous avons
$$
\pi_*\mathcal{I} = \colim \pi_{i', *}\mathcal{I}_{i'}
$$
Les morphismes de transition du système sont tous surjectifs,
comme il résulte de la surjectivité des morphismes
$\pi_{i, *}\mathcal{O}_{Y_i} \to \pi_{i', *}\mathcal{O}_{Y_{i'}}$
et du fait que $Z_{i'} = Z_i \times_{Y_i} Y_{i'}$.
D'après Cohomologie des espaces, lemme
\ref{spaces-cohomology-lemma-finite-presentation-quasi-compact-colimit},
pour un certain $i' \geq i$, le morphisme $\mathcal{E} \to \pi_*\mathcal{I}$
se relève en un morphisme $\mathcal{E} \to \pi_{i', *}\mathcal{I}_{i'}$.
Après avoir augmenté $i'$, ce morphisme
$\mathcal{E} \to \pi_{i', *}\mathcal{I}_{i'}$
devient surjectif (sinon la limite inductive des conoyaux,
dont les morphismes de transition sont surjectifs, serait non nulle).
Cela nous ramène au cas examiné au paragraphe suivant.

\medskip\noindent
Supposons que $X$ soit un espace algébrique sur $\mathbf{Z}$ et que
$Y \to X$ soit de présentation finie. Par approximation noethérienne
absolue, nous pouvons écrire $X = \lim X_i$ comme une limite projective filtrante,
où chaque $X_i$ est un espace algébrique quasi-séparé de type fini
sur $\mathbf{Z}$ et les morphismes de transition sont affines, voir
Limites d'espaces, proposition \ref{spaces-limits-proposition-approximate}.
Comme $\pi : Y \to X$ est de présentation finie, nous pouvons trouver un $i$
et un morphisme $\pi_i : Y_i \to X_i$ de présentation finie
dont le changement de base à $X$ est $\pi$, voir
Limites d'espaces, lemme
\ref{spaces-limits-lemma-descend-finite-presentation}.
Après avoir augmenté $i$, nous pouvons supposer que $\pi_i$ est fini, voir
Limites d'espaces, lemme
\ref{spaces-limits-lemma-descend-finite}.
Ensuite, nous pouvons supposer qu'il existe un module localement
libre de type fini sur $\mathcal{O}_{X_i}$, noté $\mathcal{E}_i$,
dont l'image inverse sur $X$ est $\mathcal{E}$, voir
Limites d'espaces, lemme \ref{spaces-limits-lemma-descend-invertible-modules}.
Nous pouvons aussi supposer qu'il existe un morphisme
$\mathcal{E}_i \to \pi_{i, *}\mathcal{O}_{Y_i}$
dont l'image inverse sur $X$ est la composée
$\mathcal{E} \to \pi_*\mathcal{I} \to \pi_*\mathcal{O}_Y$, voir
Limites d'espaces, lemme
\ref{spaces-limits-lemma-descend-modules-finite-presentation}.
Le conoyau
$$
\mathcal{E}_i \to \pi_{i, *}\mathcal{O}_{Y_i} \to \mathcal{Q}_i \to 0
$$
est un $\mathcal{O}_{Y_i}$-module cohérent dont l'image inverse sur $X$
est le conoyau (de présentation finie) $\mathcal{Q}$ du morphisme
$\mathcal{E} \to \pi_*\mathcal{O}_Y$. Autrement dit, nous avons
$\mathcal{Q} = \pi_*(\mathcal{O}_Y/\mathcal{I})$. Considérons
le morphisme
$$
\mathcal{E}_i
\otimes_{\mathcal{O}_{X_i}}
\pi_{i, *} \mathcal{O}_{Y_i}
\longrightarrow
\pi_{i, *}\mathcal{O}_{Y_i}
\otimes_{\mathcal{O}_{X_i}}
\pi_{i, *} \mathcal{O}_{Y_i}
\to
\pi_{i, *} \mathcal{O}_{Y_i}
\to
\mathcal{Q}_i
$$
où la seconde flèche est donnée par la structure d'algèbre
sur $\pi_{i, *}\mathcal{O}_{Y_i}$. L'image inverse de ce morphisme
sur $Y$ est nulle, car l'image de $\mathcal{E} \to \pi_*\mathcal{O}_Y$
est l'idéal $\pi_*\mathcal{I}$. Par conséquent, d'après
Limites d'espaces, lemme
\ref{spaces-limits-lemma-descend-modules-finite-presentation},
après avoir augmenté $i$, nous pouvons supposer que la composée affichée
est nulle. Cela signifie exactement que l'image de
$\mathcal{E}_i \to \pi_{i, *}\mathcal{O}_{Y_i}$
est de la forme $\pi_{i, *}\mathcal{I}_i$
pour un faisceau d'idéaux cohérent $\mathcal{I}_i \subset \mathcal{O}_{Y_i}$.
Comme $\mathcal{E}_i \to \pi_{i, *}\mathcal{O}_{Y_i}$
s'identifie après changement de base à $\mathcal{E} \to \pi_*\mathcal{O}_Y$, nous
voyons que l'image inverse de $\mathcal{I}_i$ sur $Y$ engendre $\mathcal{I}$.
Notons $V_i \subset Y_i$ le sous-espace ouvert dont le complémentaire
est $V(\mathcal{I}_i) \subset Y_i$. Alors $V$ est l'image inverse
de $V_i$ d'après les remarques ci-dessus. Après avoir augmenté $i$, nous pouvons supposer
que $V_i$ est affine et que $\pi_i|_{V_i} : V_i \to X_i$ est 
\'etale, voir
Limites d'espaces, lemmes
\ref{spaces-limits-lemma-limit-is-affine} et
\ref{spaces-limits-lemma-descend-etale}.
Ceci étant établi, nous pouvons appliquer le
lemme \ref{spaces-more-morphisms-lemma-resolution-property-in-Noetherian-case}
pour conclure que $X_i$ possède la
propriété de résolution. Comme $X \to X_i$ est affine,
nous concluons que $X$ possède lui aussi la propriété de résolution
d'après Catégories dérivées des espaces, lemme
\ref{spaces-perfect-lemma-resolution-property-goes-up-affine}.
\end{proof}

\begin{lemma}
\label{spaces-more-morphisms-lemma-resolution-property-descends}
Soit $S$ un schéma.
Soit $X = \lim X_i$ la limite d'un système projectif d'espaces algébriques quasi-compacts
et quasi-séparés sur $S$, dont les morphismes de transition sont affines.
Alors $X$ possède la propriété de résolution si et seulement si $X_i$ possède
la propriété de résolution pour un certain $i$.
\end{lemma}

\begin{proof}
Si $X_i$ possède la propriété de résolution, alors $X$ la possède d'après
Catégories dérivées des espaces, lemme
\ref{spaces-perfect-lemma-resolution-property-goes-up-affine}.
Supposons que $X$ possède la propriété de résolution.
Choisissons $i \in I$. Nous pouvons choisir un schéma affine $V_i$ et un morphisme
\'etale surjectif $V_i \to X_i$ 
(Propriétés des espaces, lemme
\ref{spaces-properties-lemma-quasi-compact-affine-cover}).
Nous pouvons choisir une immersion $j : V_i \to Y_i$
où $Y_i$ est fini et de présentation finie sur $X_i$
(lemme \ref{spaces-more-morphisms-lemma-quasi-finite-separated-pass-through-finite-addendum}).
Nous pouvons choisir un idéal quasi-cohérent de type fini
$\mathcal{I}_i \subset \mathcal{O}_{Y_i}$
tel que $V_i = Y_i \setminus V(\mathcal{I}_i)$
(Limites d'espaces, lemme
\ref{spaces-limits-lemma-quasi-coherent-finite-type-ideals}).
Notons $V \to Y \to X$ les changements de base de $V_i \to Y_i \to X_i$
à $X$. Notons $\mathcal{I} \subset \mathcal{O}_Y$ l'image inverse
de l'idéal $\mathcal{I}_i$. D'après l'implication facile du
lemme \ref{spaces-more-morphisms-lemma-resolution-property-one-sheaf},
il existe un $\mathcal{O}_X$-module localement libre de type fini
$\mathcal{E}$ et une surjection $\mathcal{E} \to \pi_*\mathcal{I}$.
Remarquons que, puisque $\pi_i : Y_i \to X_i$ est fini et de présentation finie,
le morphisme $\pi : Y \to X$ est lui aussi fini et de présentation finie,
et les $\mathcal{O}_{X_i}$-modules
$\pi_{i, *}\mathcal{O}_{Y_i}$ et
$\pi_{i, *}(\mathcal{O}_{Y_i}/\mathcal{I}_i)$
sont de présentation finie et donnent après changement de base à $X$
$\pi_*\mathcal{O}_Y$ et
$\pi_*(\mathcal{O}_Y/\mathcal{I})$. Ainsi, d'après Limites d'espaces, lemme
\ref{spaces-limits-lemma-descend-modules-finite-presentation},
après avoir augmenté $i$, nous pouvons trouver un module localement libre
de type fini sur $\mathcal{O}_{X_i}$, noté $\mathcal{E}_i$,
et un morphisme $\mathcal{E}_i \to \pi_{i, *}\mathcal{O}_{Y_i}$
dont le changement de base à $X$ redonne la composée
$\mathcal{E} \to \pi_*\mathcal{I} \to \pi_*\mathcal{O}_Y$.
Les images inverses des $\mathcal{O}_{X_i}$-modules de présentation finie
$\Coker(\mathcal{E}_i \to \pi_{i, *}\mathcal{O}_{Y_i})$
et $\pi_{i, *}(\mathcal{O}_{Y_i}/\mathcal{I}_i)$
sur $X$ coïncident comme quotients de $\pi_*\mathcal{O}_Y$.
Ainsi, d'après Limites d'espaces, lemme
\ref{spaces-limits-lemma-descend-modules-finite-presentation},
nous pouvons supposer qu'ils coïncident, autrement dit que
l'image de $\mathcal{E}_i \to \pi_{i, *}\mathcal{O}_{X_i}$
est égale à $\pi_{i, *}\mathcal{I}_i$.
Nous concluons alors que $X_i$ possède la propriété de résolution d'après le
lemme \ref{spaces-more-morphisms-lemma-resolution-property-one-sheaf}.
\end{proof}

\begin{lemma}
\label{spaces-more-morphisms-lemma-resolution-property-affine-diagonal}
Soit $S$ un schéma.
Soit $X$ un espace algébrique quasi-compact et quasi-séparé qui possède
la propriété de résolution. Alors $X$ a une diagonale affine sur $\mathbf{Z}$
(au sens de Propriétés des espaces, définition
\ref{spaces-properties-definition-separated}).
\end{lemma}

\begin{proof}
Nous pourrions le démontrer comme dans le cas des schémas, mais nous allons plutôt
déduire le lemme du cas des schémas.
Tout d'abord, nous pouvons et voulons supposer $S = \Spec(\mathbf{Z})$.
Ensuite, choisissons un schéma $Y$ et un morphisme entier surjectif
$f : Y \to X$, voir
Espaces décents, lemme \ref{decent-spaces-lemma-there-is-a-scheme-integral-over}.
Alors $f$ est affine, donc $Y$ possède la propriété de résolution d'après
Catégories dérivées des espaces, lemme
\ref{spaces-perfect-lemma-resolution-property-goes-up-affine}.
Ainsi, d'après le cas des schémas, le schéma $Y$ a une diagonale affine, voir
Catégories dérivées des schémas, lemme
\ref{perfect-lemma-resolution-property-affine-diagonal}.
Considérons ensuite le diagramme commutatif
$$
\xymatrix{
Y \ar[d] \ar[rr]_{\Delta_Y} & & Y \times_{\mathbf{Z}} Y \ar[d] \\
X \ar[rr]^{\Delta_X} & & X \times_{\mathbf{Z}} X
}
$$
Remarquons que la flèche verticale de droite est entière, donc en particulier
affine. Soit $W \to X \times_{\mathbf{Z}} X$ un morphisme avec $W$ affine.
Nous voyons alors que
$$
Y \times_{X \times_{\mathbf{Z}} X} W =
Y \times_{\Delta_Y, Y \times_{\mathbf{Z}} Y}
(Y \times_{\mathbf{Z}} Y) \times_{X \times_{\mathbf{Z}} X} W
$$
est affine. D'autre part, $Y \to X$ est entier et surjectif,
donc
$$
Y \times_{X \times_{\mathbf{Z}} X} W
\longrightarrow
X \times_{X \times_{\mathbf{Z}} X} W
$$
est entier et surjectif comme changement de base de $Y \to X$ à $W$.
Nous concluons que le but de cette flèche est affine d'après
Limites d'espaces, proposition \ref{spaces-limits-proposition-affine}.
Il s'ensuit que $\Delta_X$ est affine, comme souhaité.
\end{proof}







\section{Éclatements et propriété de résolution}
\label{spaces-more-morphisms-section-resolution-property-by-blowing-up}

\noindent
Nous démontrons que la propriété de résolution est vérifiée après un éclatement.

\begin{lemma}
\label{spaces-more-morphisms-lemma-blow-up-resolution-property}
Soit $S$ un schéma. Soit $X$ un espace algébrique quasi-compact et quasi-séparé
sur $S$. Supposons que $|X|$ ait un nombre fini de composantes
irréductibles. Il existe un ouvert dense quasi-compact $U \subset X$
et un éclatement $U$-admissible $X' \to X$ tels que l'espace algébrique
$X'$ possède la propriété de résolution.
\end{lemma}

\begin{proof}
D'après Limites d'espaces, lemme
\ref{spaces-limits-lemma-there-is-a-scheme-finite-over-refined},
il existe un morphisme surjectif, fini et de présentation finie
$f : Y \to X$, où $Y$ est un schéma, et un ouvert dense quasi-compact
$U \subset X$ tel que $f^{-1}(U) \to U$ soit fini \'etale.
D'après Compléments sur les morphismes, lemme \ref{more-morphisms-lemma-blow-up-ample-family},
il existe un ouvert dense quasi-compact $V \subset Y$ et un éclatement $V$-admissible
$Y' \to Y$ tel que $Y'$ possède une famille ample de modules
inversibles. Après avoir rétréci $U$, nous pouvons supposer que $f^{-1}(U) \subset V$
(détails omis). Ainsi $f' : Y' \to X$ est fini \'etale au-dessus de $U$
et, en particulier, le morphisme $(f')^{-1}(U) \to U$ est fini localement libre.
D'après le lemme \ref{spaces-more-morphisms-lemma-finite-after-blowing-up}, il existe un éclatement $U$-admissible
$X' \to X$ tel que le transformé strict $Y''$ de $Y'$
soit fini localement libre sur $X'$. Considérons le diagramme
$$
\xymatrix{
Y'' \ar[d] \ar[r]_g &
Y' \ar[r] &
Y \ar[d] \\
X' \ar[rr] & & X
}
$$
Comme $g : Y'' \to Y'$ est un éclatement dont le centre est contenu
(Diviseurs sur les espaces, lemme \ref{spaces-divisors-lemma-strict-transform})
dans l'image inverse du centre de $X' \to X$, nous voyons que
$g : Y'' \to Y'$ est projectif et qu'il existe un module inversible 
$g$-ample sur $Y''$. Par conséquent, d'après
Compléments sur les morphismes, lemme \ref{more-morphisms-lemma-ample-family-ample-relative},
nous voyons que $Y''$ possède une famille ample de modules inversibles.
Ainsi $Y''$ possède la propriété de résolution, voir
Catégories dérivées des schémas, lemme
\ref{perfect-lemma-resolution-property-ample-family}.
Nous concluons que $X'$ possède la propriété de résolution
d'après
Catégories dérivées des espaces, lemme
\ref{spaces-perfect-lemma-resolution-property-goes-down-finite-flat}.
\end{proof}

\begin{lemma}
\label{spaces-more-morphisms-lemma-blow-up-resolution-property-filtered}
Soit $S$ un schéma. Soit $X$ un espace algébrique quasi-compact et quasi-séparé
sur $S$. Il existe un $t \geq 0$ et des sous-espaces
fermés
$$
X \supset Z_0 \supset Z_1 \supset \ldots \supset Z_t = \emptyset
$$
tels que $Z_i \to X$ soit de présentation finie,
$Z_0 \subset X$ soit un épaississement et que, pour chaque $i = 0, \ldots t - 1$,
il existe un éclatement $(Z_i \setminus Z_{i - 1})$-admissible
$Z'_i \to Z_i$ tel que $Z'_i$ possède la propriété de résolution.
\end{lemma}

\begin{proof}
Dans ce paragraphe, nous utilisons l'approximation noethérienne absolue pour nous ramener au
cas des espaces algébriques de présentation finie sur $\Spec(\mathbf{Z})$.
Nous pouvons considérer $X$ comme un espace algébrique sur $\Spec(\mathbf{Z})$, voir
Espaces, définition \ref{spaces-definition-base-change} et
Propriétés des espaces, définition \ref{spaces-properties-definition-separated}.
Nous pouvons donc appliquer Limites d'espaces, proposition
\ref{spaces-limits-proposition-approximate}.
Il s'ensuit que nous pouvons trouver un morphisme affine $X \to X_0$
avec $X_0$ de présentation finie sur $\mathbf{Z}$.
Si nous pouvons démontrer le lemme pour $X_0$, nous pouvons alors prendre les images inverses
de la stratification et des centres des éclatements sur $X$
et obtenir le résultat pour $X$; nous utilisons ici le fait que la propriété de résolution
monte le long des morphismes affines (Catégories dérivées des espaces,
lemme \ref{spaces-perfect-lemma-resolution-property-goes-up-affine}) et
que le transformé strict d'un morphisme affine est affine -- détails omis.
Cela nous ramène au cas examiné au paragraphe suivant.

\medskip\noindent
Supposons que $X$ soit de présentation finie sur $\mathbf{Z}$.
Alors $X$ est noethérien et $|X|$ est un espace topologique
noethérien (ayant un nombre fini de composantes irréductibles) de dimension finie.
Nous pouvons donc raisonner par récurrence sur $\dim(|X|)$.
D'après le lemme \ref{spaces-more-morphisms-lemma-blow-up-resolution-property},
il existe un ouvert dense $U \subset X$ et un éclatement $U$-admissible
$X' \to X$ tel que $X'$ possède la propriété de résolution.
Posons $Z_0 = X$ et soit $Z_1 \subset X$ le sous-espace fermé réduit
tel que $|Z_1| = |X| \setminus |U|$.
Par récurrence, nous trouvons un entier $t \geq 0$ et une filtration
$$
Z_1 \supset Z_{1, 0} \supset Z_{1, 1} \supset \ldots
\supset Z_{1, t} = \emptyset
$$
par des sous-espaces fermés, où $Z_{1, 0} \to Z_1$ est un épaississement
et il existe des éclatements $(Z_{1, i} \setminus Z_{1, i + 1})$-admissibles
$Z'_{1, i} \to Z_{1, i}$ tels que $Z'_{1, i}$
possède la propriété de résolution. Comme $Z_1$ est réduit, nous avons $Z_1 = Z_{1, 0}$.
Nous pouvons donc poser $Z_i = Z_{1, i - 1}$ et $Z'_i = Z'_{1, i - 1}$
pour $i \geq 1$, et le lemme est démontré.
\end{proof}














% D'accord, commençons ici.
%

\chapter{Platitude sur les espaces algébriques}


\phantomsection
\label{spaces-flat-section-phantom}




\section{Introduction}
\label{spaces-flat-section-introduction}

\noindent
Dans ce chapitre, nous étudions quelques résultats avancés sur les modules plats et
les morphismes plats dans le cadre des espaces algébriques. Nous recommandons vivement
au lecteur de consulter d'abord le chapitre correspondant
dans le cadre des schémas, voir
Compléments sur la platitude, section \ref{flat-section-introduction}.
L'article \cite{GruRay} de Raynaud et Gruson constitue une référence.







\section{Impuretés}
\label{spaces-flat-section-impure}

\noindent
Cette section est l'analogue de
Compléments sur la platitude, section \ref{flat-section-impure}.

\begin{situation}
\label{spaces-flat-situation-pre-pure}
Soit $S$ un schéma. Soit $f : X \to Y$ un morphisme
de type fini et décent\footnote{Les morphismes quasi-séparés sont décents, voir
Espaces décents, Lemme
\ref{decent-spaces-lemma-properties-trivial-implications}.
Pour tout morphisme $\Spec(k) \to Y$, où $k$ est un corps,
l'espace algébrique $X_k$ est de présentation finie sur $k$
car il est de type fini sur $k$ et quasi-séparé d'après
Espaces décents, Lemme
\ref{decent-spaces-lemma-locally-Noetherian-decent-quasi-separated}.}
entre espaces algébriques sur $S$. De plus, $\mathcal{F}$ est un $\mathcal{O}_X$-module
quasi-cohérent de type fini. Enfin, $y \in |Y|$ est un point de $Y$.
\end{situation}

\noindent
Dans cette situation, considérons un schéma $T$, un morphisme $g : T \to Y$,
un point $t \in T$ tel que $g(t) = y$, une spécialisation $t' \leadsto t$ dans
$T$, et un point $\xi \in |X_T|$ au-dessus de $t'$. Ici $X_T = T \times_Y X$.
On a le diagramme
\begin{equation}
\label{spaces-flat-equation-impurity}
\vcenter{
\xymatrix{
\xi \ar@{|->}[d] & \\
t' \ar@{~>}[r] & t \ar@{|->}[r] \ar[r] & y
}
}
\quad\quad
\vcenter{
\xymatrix{
X_T \ar[d]_{f_T} \ar[r] & X \ar[d]^f \\
T \ar[r]^g & Y
}
}
\end{equation}
De plus, notons $\mathcal{F}_T$ l'image réciproque de $\mathcal{F}$ sur $X_T$.

\begin{definition}
\label{spaces-flat-definition-impurity}
Dans la
Situation \ref{spaces-flat-situation-pre-pure},
nous disons qu'un diagramme (\ref{spaces-flat-equation-impurity}) définit une
{\it impureté de $\mathcal{F}$ au-dessus de $y$}
si $\xi \in \text{Ass}_{X_T/T}(\mathcal{F}_T)$ et
$t \not \in f_T(\overline{\{\xi\}})$. Nous l'indiquerons
en disant : « soit $(g : T \to Y, t' \leadsto t, \xi)$
une impureté de $\mathcal{F}$ au-dessus de $y$ ».
\end{definition}

\noindent
Autrement dit, $(g : T \to Y, t' \leadsto t, \xi)$ est
une impureté de $\mathcal{F}$ au-dessus de $y$ s'il n'existe aucune spécialisation
$\xi \leadsto \theta$ dans l'espace topologique $|X_T|$ telle que
$f_T(\theta) = t$. Les spécialisations se comportent mal dans les espaces algébriques
non décents. Si le morphisme $f$ est décent, alors $X_T$
est un espace algébrique décent pour tout morphisme $g : T \to Y$ comme ci-dessus, voir
Espaces décents, Définition \ref{decent-spaces-definition-relative-conditions}.

\begin{lemma}
\label{spaces-flat-lemma-impure-limit}
Dans la Situation \ref{spaces-flat-situation-pre-pure}.
Soit $(g : T \to S, t' \leadsto t, \xi)$ une impureté de
$\mathcal{F}$ au-dessus de $y$. Supposons que $T = \lim_{i \in I} T_i$ soit une limite dirigée
de schémas affines sur $Y$. Alors, pour un certain $i$, le triplet
$(T_i \to Y, t'_i \leadsto t_i, \xi_i)$ est une impureté de
$\mathcal{F}$ au-dessus de $y$.
\end{lemma}

\begin{proof}
Les notations de l'énoncé signifient ceci : soient $p_i : T \to T_i$
les morphismes de projection, et posons $t_i = p_i(t)$ et $t'_i = p_i(t')$.
Enfin, $\xi_i \in |X_{T_i}|$ est l'image de $\xi$. D'après
Diviseurs sur les espaces, Lemme
\ref{spaces-divisors-lemma-base-change-relative-assassin}
nous avons $\xi_i \in \text{Ass}_{X_{T_i}/T_i}(\mathcal{F}_{T_i})$.
Il reste donc seulement à montrer que
$t_i \not \in f_{T_i}(\overline{\{\xi_i\}})$ pour un certain $i$.

\medskip\noindent
Soit $Z_i \subset X_{T_i}$ la structure de schéma réduite induite
sur $\overline{\{\xi_i\}} \subset |X_{T_i}|$,
et soit $Z \subset X_T$ la structure de schéma réduite induite sur
$\overline{\{\xi\}} \subset |X_T|$.
Alors $Z = \lim Z_i$ d'après
Limites d'espaces, Lemme \ref{spaces-limits-lemma-inverse-limit-irreducibles}
(le lemme s'applique car chaque $X_{T_i}$ est décent).
Choisissons un corps $k$ et un morphisme $\Spec(k) \to T$ dont l'image est $t$.
Alors
$$
\emptyset =
Z \times_T \Spec(k) = (\lim Z_i) \times_{(\lim T_i)} \Spec(k)
= \lim Z_i \times_{T_i} \Spec(k)
$$
car les limites commutent aux produits fibrés (les limites commutent aux limites).
Chaque $Z_i \times_{T_i} \Spec(k)$ est quasi-compact, car $X_{T_i} \to T_i$
est de type fini, et donc $Z_i \to T_i$ est de type fini.
Ainsi $Z_i \times_{T_i} \Spec(k)$ est vide pour un certain $i$ d'après
Limites d'espaces, Lemme \ref{spaces-limits-lemma-limit-nonempty}.
Puisque l'image du composé $\Spec(k) \to T \to T_i$ est $t_i$,
nous obtenons le résultat voulu.
\end{proof}

\noindent
Les impuretés remontent par changement de base plat.

\begin{lemma}
\label{spaces-flat-lemma-flat-ascent-impurity}
Dans la Situation \ref{spaces-flat-situation-pre-pure}.
Soit $(Y_1, y_1) \to (Y, y)$ un morphisme d'espaces algébriques
pointés sur $S$. Supposons que $Y_1 \to Y$ soit plat en $y_1$.
Si $(T \to Y, t' \leadsto t, \xi)$ est une impureté de
$\mathcal{F}$ au-dessus de $y$, alors il existe une impureté
$(T_1 \to Y_1, t_1' \leadsto t_1, \xi_1)$ de l'image réciproque
$\mathcal{F}_1$ de $\mathcal{F}$ sur $X_1 = Y_1 \times_Y X$
au-dessus de $y_1$, telle que $T_1$ soit \'etale sur $Y_1 \times_Y T$.
\end{lemma}

\begin{proof}
Choisissons un morphisme \'etale $T_1 \to Y_1 \times_Y T$, où $T_1$
est un schéma, et un point $t_1 \in T_1$ s'envoyant sur $y_1$ et $t$.
L'existence d'un tel couple $(T_1, t_1)$ résulte de
Propriétés des espaces, Lemme \ref{spaces-properties-lemma-points-cartesian}.
Le morphisme de schémas $T_1 \to T$ est plat en $t_1$
(utiliser Morphismes d'espaces, Lemme \ref{spaces-morphisms-lemma-base-change-flat}
et la définition des morphismes plats d'espaces algébriques) ;
il existe donc une spécialisation $t'_1 \leadsto t_1$ au-dessus de
$t' \leadsto t$, voir
Morphismes, Lemme \ref{morphisms-lemma-generalizations-lift-flat}.
Choisissons un point $\xi_1 \in |X_{T_1}|$ s'envoyant sur $t'_1$
et sur $\xi$, avec $\xi_1 \in \text{Ass}_{X_{T_1}/T_1}(\mathcal{F}_{T_1})$.
Autrement dit, un point de $\Spec(\kappa(t'_1) \otimes_{\kappa(t')} \kappa(\xi))$.
C'est possible d'après
Diviseurs sur les espaces, Lemme
\ref{spaces-divisors-lemma-base-change-relative-assassin}.
Comme l'adhérence $Z_1$ de $\{\xi_1\}$ dans $|X_{T_1}|$ s'envoie dans
l'adhérence de $\{\xi\}$ dans $|X_T|$, nous en concluons que
l'image de $Z_1$ dans $|T_1|$ ne peut contenir $t_1$.
Ainsi $(T_1 \to Y_1, t'_1 \leadsto t_1, \xi_1)$
est une impureté de $\mathcal{F}_1$ au-dessus de $Y_1$.
\end{proof}

\begin{lemma}
\label{spaces-flat-lemma-pure-along-X-y}
Dans la Situation \ref{spaces-flat-situation-pre-pure}. Soit $\overline{y}$ un point
géométrique au-dessus de $y$. Soit $\mathcal{O} = \mathcal{O}_{Y, \overline{y}}$
l'anneau local étale de $Y$ en $\overline{y}$. Notons
$Y^{sh} = \Spec(\mathcal{O})$, $X^{sh} = X \times_Y Y^{sh}$, et
$\mathcal{F}^{sh}$ l'image réciproque de $\mathcal{F}$ sur $X^{sh}$.
Les conditions suivantes sont équivalentes :
\begin{enumerate}
\item il existe une impureté
$(Y^{sh} \to Y, y' \leadsto \overline{y}, \xi)$
de $\mathcal{F}$ au-dessus de $y$,
\item tout point de $\text{Ass}_{X^{sh}/Y^{sh}}(\mathcal{F}^{sh})$
se spécialise en un point de la fibre fermée $X_{\overline{y}}$,
\item il existe une impureté $(T \to Y, t' \leadsto t, \xi)$
de $\mathcal{F}$ au-dessus de $y$ telle que $(T, t) \to (Y, y)$ soit un
voisinage \'etale, et
\item il existe une impureté $(T \to Y, t' \leadsto t, \xi)$
de $\mathcal{F}$ au-dessus de $y$ telle que $T \to Y$ soit quasi-fini en $t$.
\end{enumerate}
\end{lemma}

\begin{proof}
L'équivalence de (1) et (2) résulte immédiatement de la définition.

\medskip\noindent
Rappelons que $\mathcal{O} = \mathcal{O}_{Y, \overline{y}}$
est la colimite filtrante des $\mathcal{O}(V)$ sur la catégorie
des voisinages \'etales $(V, \overline{v}) \to (Y, \overline{y})$
(Propriétés des espaces, Lemme \ref{spaces-properties-lemma-cofinal-etale}).
De plus, il suffit de considérer les voisinages \'etales affines $V$.
Par conséquent, $Y^{sh} = \Spec(\mathcal{O}) = \lim \Spec(\mathcal{O}(V)) = \lim V$.
Le Lemme \ref{spaces-flat-lemma-impure-limit} montre donc que (1) implique (3).

\medskip\noindent
Puisqu'un morphisme \'etale est localement quasi-fini
(Morphismes d'espaces, Lemme
\ref{spaces-morphisms-lemma-etale-locally-quasi-finite})
nous voyons que (3) implique (4).

\medskip\noindent
Enfin, supposons (4). Après avoir remplacé $T$ par un voisinage ouvert de $t$,
nous pouvons supposer que $T \to Y$ est localement quasi-fini.
Le Lemme \ref{spaces-flat-lemma-flat-ascent-impurity}
fournit une impureté
$(T_1 \to Y^{sh}, t_1' \leadsto t_1, \xi_1)$
pour laquelle $T_1 \to T \times_Y Y^{sh}$ est
\'etale. Puisqu'un morphisme \'etale est localement quasi-fini,
Morphismes d'espaces, Lemme
\ref{spaces-morphisms-lemma-base-change-quasi-finite}, et
Morphismes, Lemme \ref{morphisms-lemma-composition-quasi-finite},
montrent que $T_1 \to Y^{sh}$ est localement quasi-fini.
Comme $\mathcal{O}$ est strictement hensélien, nous pouvons appliquer Compléments sur les morphismes, Lemme
\ref{more-morphisms-lemma-etale-makes-quasi-finite-finite-at-point}
pour voir qu'après avoir remplacé $T_1$ par un voisinage ouvert et fermé
de $t_1$, nous pouvons supposer que $T_1 \to Y^{sh} = \Spec(\mathcal{O})$
est fini. Soit $\theta \in |X^{sh}|$ l'image de
$\xi_1$, et soit $y' \in \Spec(\mathcal{O})$ l'image
de $t_1'$. D'après Diviseurs sur les espaces, Lemme
\ref{spaces-divisors-lemma-base-change-relative-assassin}
nous voyons que $\theta \in \text{Ass}_{X^{sh}/Y^{sh}}(\mathcal{F}^{sh})$.
Comme $\pi : X_{T_1} \to X^{sh}$
est fini, il induit une application fermée $|X_{T_1}| \to |X^{sh}|$.
L'image de $\overline{\{\xi_1\}}$ est donc $\overline{\{\theta\}}$.
Il s'ensuit que $(Y^{sh} \to Y, y' \leadsto \overline{y}, \theta)$
est une impureté de $\mathcal{F}$ au-dessus de $y$, ce qui achève la démonstration.
\end{proof}








\section{Modules relativement purs}
\label{spaces-flat-section-pure}

\noindent
Cette section est l'analogue de
Compléments sur la platitude, section \ref{flat-section-pure}.

\begin{definition}
\label{spaces-flat-definition-pure}
Dans la Situation \ref{spaces-flat-situation-pre-pure}.
\begin{enumerate}
\item Nous disons que $\mathcal{F}$ est {\it pur au-dessus de $y$} si {\bf aucune} des
conditions équivalentes du Lemme \ref{spaces-flat-lemma-pure-along-X-y} n'est satisfaite.
\item Nous disons que $\mathcal{F}$ est {\it universellement pur au-dessus de $y$}
s'il n'existe aucune impureté de $\mathcal{F}$ au-dessus de $y$.
\item Nous disons que $X$ est {\it pur au-dessus de $y$} si $\mathcal{O}_X$
est pur au-dessus de $y$.
\item Nous disons que $\mathcal{F}$ est {\it universellement $Y$-pur}, ou
{\it universellement pur relativement à $Y$}, si $\mathcal{F}$ est universellement
pur au-dessus de $y$ pour tout $y \in |Y|$.
\item Nous disons que $\mathcal{F}$ est {\it $Y$-pur}, ou
{\it pur relativement à $Y$}, si $\mathcal{F}$ est pur au-dessus de $y$
pour tout $y \in |Y|$.
\item Nous disons que $X$ est {\it $Y$-pur} ou {\it pur relativement à $Y$}
si $\mathcal{O}_X$ est pur relativement à $Y$.
\end{enumerate}
\end{definition}

\noindent
Voici les lemmes obligés.

\begin{lemma}
\label{spaces-flat-lemma-base-change-universally}
Dans la Situation \ref{spaces-flat-situation-pre-pure}.
\begin{enumerate}
\item $\mathcal{F}$ est universellement pur au-dessus de $y$, et
\item pour tout morphisme $(Y', y') \to (Y, y)$ d'espaces algébriques pointés,
l'image réciproque $\mathcal{F}_{Y'}$ est pure au-dessus de $y'$.
\end{enumerate}
En particulier, $\mathcal{F}$ est universellement pur relativement à $Y$ si et
seulement si tout changement de base $\mathcal{F}_{Y'}$ de $\mathcal{F}$ est
pur relativement à $Y'$.
\end{lemma}

\begin{proof}
C'est formel.
\end{proof}

\begin{lemma}
\label{spaces-flat-lemma-quasi-finite-base-change}
Dans la Situation \ref{spaces-flat-situation-pre-pure}.
Soit $(Y', y') \to (Y, y)$ un morphisme d'espaces algébriques pointés.
Si $Y' \to Y$ est quasi-fini en $y'$ et si $\mathcal{F}$ est pur au-dessus de $y$,
alors $\mathcal{F}_{Y'}$ est pur au-dessus de $y'$.
\end{lemma}

\begin{proof}
Si $(T \to Y', t' \leadsto t, \xi)$ est une impureté de
$\mathcal{F}_{Y'}$ au-dessus de $y'$, avec $T \to Y'$ quasi-fini en $t$,
alors $(T \to Y, t' \to t, \xi)$ est une impureté de $\mathcal{F}$
au-dessus de $y$, avec $T \to Y$ quasi-fini en $t$ ; voir
Morphismes d'espaces, Lemme
\ref{spaces-morphisms-lemma-composition-quasi-finite}.
Le lemme résulte donc immédiatement de la définition de la pureté.
\end{proof}

\noindent
La pureté satisfait la descente plate.

\begin{lemma}
\label{spaces-flat-lemma-flat-descend-pure}
Dans la Situation \ref{spaces-flat-situation-pre-pure}.
Soit $(Y_1, y_1) \to (Y, y)$ un morphisme d'espaces algébriques pointés.
Supposons que $Y_1 \to Y$ soit plat en $y_1$.
\begin{enumerate}
\item Si $\mathcal{F}_{Y_1}$ est pur au-dessus de $y_1$,
alors $\mathcal{F}$ est pur au-dessus de $y$.
\item Si $\mathcal{F}_{Y_1}$ est universellement pur au-dessus de $y_1$,
alors $\mathcal{F}$ est universellement pur au-dessus de $y$.
\end{enumerate}
\end{lemma}

\begin{proof}
Cela vient de ce que les impuretés remontent par changement de base plat, voir
Lemme \ref{spaces-flat-lemma-flat-ascent-impurity}. Par exemple,
pour (1), toute impureté $(T \to Y, t' \leadsto t, \xi)$
de $\mathcal{F}$ au-dessus de $y$, avec $T \to Y$ quasi-fini en $t$,
conduit par ce lemme à une impureté
$(T_1 \to Y_1, t_1' \leadsto t_1, \xi_1)$ de l'image réciproque
$\mathcal{F}_1$ de $\mathcal{F}$ sur $X_1 = Y_1 \times_Y X$
au-dessus de $y_1$, telle que $T_1$ soit \'etale sur $Y_1 \times_Y T$.
Ainsi $T_1 \to Y_1$ est quasi-fini en $t_1$, car
les morphismes \'etales sont localement quasi-finis et les composés
de morphismes localement quasi-finis sont localement quasi-finis
(Morphismes d'espaces, Lemmes
\ref{spaces-morphisms-lemma-etale-locally-quasi-finite} et
\ref{spaces-morphisms-lemma-composition-quasi-finite}).
Le raisonnement est analogue pour (2).
\end{proof}

\begin{lemma}
\label{spaces-flat-lemma-supported-on-closed}
Dans la Situation \ref{spaces-flat-situation-pre-pure}. Soit $i : Z \to X$ une immersion fermée,
et supposons que $\mathcal{F} = i_*\mathcal{G}$ pour un faisceau
quasi-cohérent de type fini $\mathcal{G}$ sur $Z$.
Alors $\mathcal{G}$ est (universellement) pur au-dessus de $y$
si et seulement si $\mathcal{F}$ est (universellement) pur au-dessus de $y$.
\end{lemma}

\begin{proof}
Cela résulte de Diviseurs sur les espaces, Lemme
\ref{spaces-divisors-lemma-relative-weak-assassin-finite}.
\end{proof}

\begin{lemma}
\label{spaces-flat-lemma-proper-pure}
Dans la Situation \ref{spaces-flat-situation-pre-pure}.
\begin{enumerate}
\item Si le support de $\mathcal{F}$ est propre sur $Y$, alors
$\mathcal{F}$ est universellement pur relativement à $Y$.
\item Si $f$ est propre, alors
$\mathcal{F}$ est universellement pur relativement à $Y$.
\item Si $f$ est propre, alors $X$ est universellement pur relativement à $Y$.
\end{enumerate}
\end{lemma}

\begin{proof}
Réduisons d'abord (1) à (2). Soit en effet $Z \subset X$ le
support schématique de $\mathcal{F}$
(Morphismes d'espaces, Définition
\ref{spaces-morphisms-definition-scheme-theoretic-support}). Soit $i : Z \to X$
l'immersion fermée correspondante, et écrivons
$\mathcal{F} = i_*\mathcal{G}$ pour un $\mathcal{O}_Z$-module $\mathcal{G}$
quasi-cohérent de type fini.
Dans le cas (1), $Z \to Y$ est propre par hypothèse.
Le Lemme \ref{spaces-flat-lemma-supported-on-closed} réduit donc le cas (1) au cas (2).

\medskip\noindent
Supposons $f$ propre.
Soit $(g : T \to Y, t' \leadsto t, \xi)$ une impureté de $\mathcal{F}$
au-dessus de $y$. Puisque $f$ est propre, il est universellement fermé. Par conséquent,
$f_T : X_T \to T$ est fermé. Comme $f_T(\xi) = t'$, cela implique que
$t \in f(\overline{\{\xi\}})$, contradiction.
\end{proof}







\section{Modules plats de type fini}
\label{spaces-flat-section-finite-type-flat}

\noindent
Comparer avec
Compléments sur la platitude, sections
\ref{flat-section-finite-type-flat-I},
\ref{flat-section-finite-type-flat-II}, et
\ref{flat-section-finite-type-flat-III}.
La plupart de ces résultats ont des conséquences immédiates
pour les espaces algébriques par localisation \'etale.

\begin{lemma}
\label{spaces-flat-lemma-existence-complete}
Soit $S$ un schéma.
Soit $X \to Y$ un morphisme de type fini d'espaces algébriques sur $S$.
Soit $\mathcal{F}$ un $\mathcal{O}_X$-module quasi-cohérent de type fini.
Soit $y \in |Y|$ un point. Il existe un morphisme \'etale
$(Y', y') \to (Y, y)$, avec $Y'$ schéma affine, et des morphismes \'etales
$h_i : W_i \to X_{Y'}$, $i = 1, \ldots, n$, tels que, pour tout
$i$, il existe un dévissage complet de $\mathcal{F}_i/W_i/Y'$ au-dessus de $y'$,
où $\mathcal{F}_i$ est l'image réciproque de $\mathcal{F}$ sur $W_i$,
et tels que $|(X_{Y'})_{y'}| \subset \bigcup h_i(W_i)$.
\end{lemma}

\begin{proof}
La question est locale pour la topologie \'etale sur $Y$ ; nous pouvons donc supposer que $Y$
est un schéma affine. Alors $X$ est quasi-compact, de sorte que nous pouvons
choisir un schéma affine $X'$ et un morphisme \'etale surjectif
$X' \to X$. Nous pouvons alors appliquer
Compléments sur la platitude, Lemme \ref{flat-lemma-existence-complete}
à $X' \to Y$, $(X' \to Y)^*\mathcal{F}$ et $y$ pour
obtenir le résultat voulu.
\end{proof}

\begin{lemma}
\label{spaces-flat-lemma-open-in-fibre-where-flat}
Soit $S$ un schéma.
Soit $f : X \to Y$ un morphisme d'espaces algébriques sur $S$
qui est localement de type fini.
Soit $\mathcal{F}$ un $\mathcal{O}_X$-module quasi-cohérent de type fini.
Soit $y \in |Y|$ et $F = f^{-1}(\{y\}) \subset |X|$. Alors l'ensemble
$$
\{x \in F \mid \mathcal{F} \text{ est plat sur }Y\text{ en }x\}
$$
est ouvert dans $F$.
\end{lemma}

\begin{proof}
Choisissons un schéma $V$, un point $v \in V$ et un morphisme \'etale $V \to Y$
envoyant $v$ sur $y$. Choisissons un schéma $U$ et un morphisme \'etale surjectif
$U \to V \times_Y X$. Alors $|U_v| \to F$ est une application continue ouverte
d'espaces topologiques, puisque $|U| \to |X|$ est continue et ouverte.
Le résultat découle donc du cas des schémas, à savoir
Compléments sur la platitude, Lemme \ref{flat-lemma-open-in-fibre-where-flat}.
\end{proof}

\begin{lemma}
\label{spaces-flat-lemma-bourbaki-finite-type-general-base-at-point}
Soit $S$ un schéma.
Soit $f : X \to Y$ un morphisme d'espaces algébriques sur $S$ qui est
localement de type fini. Soit $x \in |X|$, d'image $y \in |Y|$.
Soit $\mathcal{F}$ un faisceau quasi-cohérent de type fini sur $X$.
Soit $\mathcal{G}$ un faisceau quasi-cohérent sur $Y$.
Si $\mathcal{F}$ est plat en $x$ sur $Y$, alors
$$
x \in \text{WeakAss}_X(\mathcal{F} \otimes_{\mathcal{O}_X} f^*\mathcal{G})
\Leftrightarrow
y \in \text{WeakAss}_Y(\mathcal{G})
\text{ et }
x \in \text{Ass}_{X/Y}(\mathcal{F}).
$$
\end{lemma}

\begin{proof}
Choisissons un diagramme commutatif
$$
\xymatrix{
U \ar[d] \ar[r]_g & V \ar[d] \\
X \ar[r]^f & Y
}
$$
où $U$ et $V$ sont des schémas et les flèches verticales sont \'etales surjectives.
Choisissons $u \in U$ s'envoyant sur $x$. Posons $\mathcal{E} = \mathcal{F}|_U$
et $\mathcal{H} = \mathcal{G}|_V$.
Soit $v \in V$ l'image de $u$. Alors
$x \in \text{WeakAss}_X(\mathcal{F} \otimes_{\mathcal{O}_X} f^*\mathcal{G})$
si et seulement si
$u \in \text{WeakAss}_X(\mathcal{E} \otimes_{\mathcal{O}_X} g^*\mathcal{H})$
d'après Diviseurs sur les espaces, Définition
\ref{spaces-divisors-definition-weakly-associated}.
De même, $y \in \text{WeakAss}_Y(\mathcal{G})$ si et seulement si
$v \in \text{WeakAss}_V(\mathcal{H})$.
Enfin, $x \in \text{Ass}_{X/Y}(\mathcal{F})$ si et seulement si
$u \in \text{Ass}_{U_v}(\mathcal{E}|_{U_v})$ d'après
Diviseurs sur les espaces, Définition
\ref{spaces-divisors-definition-relative-weak-assassin}.
Observons que la platitude de $\mathcal{F}$ en $x$ est
équivalente à celle de $\mathcal{E}$ en $u$, voir
Morphismes d'espaces, Définition \ref{spaces-morphisms-definition-flat-module}.
L'équivalence pour $g : U \to V$, $\mathcal{E}$, $\mathcal{H}$, $u$ et $v$
est Compléments sur la platitude, Lemme
\ref{flat-lemma-bourbaki-finite-type-general-base-at-point}.
\end{proof}

\begin{lemma}
\label{spaces-flat-lemma-bourbaki-finite-type-general-base}
Soit $S$ un schéma. Soit $f : X \to Y$ un morphisme d'espaces algébriques
sur $S$ qui est localement de type fini.
Soit $\mathcal{F}$ un faisceau quasi-cohérent de type fini sur $X$
qui est plat sur $Y$. Soit $\mathcal{G}$ un faisceau quasi-cohérent sur $Y$.
Alors
$$
\text{WeakAss}_X(\mathcal{F} \otimes_{\mathcal{O}_X} f^*\mathcal{G}) =
\text{Ass}_{X/Y}(\mathcal{F}) \cap
|f|^{-1}(\text{WeakAss}_Y(\mathcal{G}))
$$
\end{lemma}

\begin{proof}
C'est une conséquence immédiate du
Lemme \ref{spaces-flat-lemma-bourbaki-finite-type-general-base-at-point}.
\end{proof}

\begin{theorem}
\label{spaces-flat-theorem-finite-type-flat}
Soit $S$ un schéma.
Soit $f : X \to Y$ un morphisme d'espaces algébriques sur $S$.
Soit $\mathcal{F}$ un $\mathcal{O}_X$-module quasi-cohérent.
Supposons que
\begin{enumerate}
\item $X \to Y$ soit localement de présentation finie,
\item $\mathcal{F}$ soit un $\mathcal{O}_X$-module de type fini, et
\item l'ensemble des points faiblement associés de $Y$ soit localement fini dans $Y$.
\end{enumerate}
Alors $U = \{x \in |X| : \mathcal{F}\text{ est plat en }x\text{ sur }Y\}$
est ouvert dans $X$ et $\mathcal{F}|_U$ est un $\mathcal{O}_U$-module
de présentation finie et plat sur $Y$.
\end{theorem}

\begin{proof}
La condition (3) signifie que, si $V \to Y$ est un morphisme \'etale surjectif,
où $V$ est un schéma, alors les points faiblement associés de $V$
sont localement en nombre fini sur le schéma $V$. (Rappelons que les points faiblement
associés de $V$ sont exactement l'image réciproque des points faiblement associés
de $Y$ d'après Diviseurs sur les espaces, Définition
\ref{spaces-divisors-definition-weakly-associated}.)
Cela étant, la question
est locale pour la topologie \'etale sur $X$ et $Y$ ; nous pouvons donc supposer que $X$ et $Y$
sont des schémas. Le résultat découle alors de
Compléments sur la platitude, Théorème \ref{flat-theorem-finite-type-flat}.
\end{proof}

\begin{lemma}
\label{spaces-flat-lemma-finite-type-flat-along-fibre-free-variant}
Soit $S$ un schéma. Soit $f : X \to Y$ un morphisme d'espaces algébriques
sur $S$. Soit $\mathcal{F}$ un faisceau quasi-cohérent sur $X$.
Soit $y \in |Y|$. Posons $F = f^{-1}(\{y\}) \subset |X|$. Supposons que
\begin{enumerate}
\item $f$ soit de type fini,
\item $\mathcal{F}$ soit de type fini, et
\item $\mathcal{F}$ soit plat sur $Y$ en tout $x \in F$.
\end{enumerate}
Alors il existe un morphisme \'etale $(Y', y') \to (Y, y)$,
où $Y'$ est un schéma, et un diagramme commutatif d'espaces algébriques
$$
\xymatrix{
X \ar[d] & X' \ar[l]^g \ar[d] \\
Y & \Spec(\mathcal{O}_{Y', y'}) \ar[l]
}
$$
tels que $X' \to X \times_Y \Spec(\mathcal{O}_{Y', y'})$
soit \'etale, que $|X'_{y'}| \to F$ soit surjectif, que $X'$ soit affine,
et que $\Gamma(X', g^*\mathcal{F})$ soit un $\mathcal{O}_{Y', y'}$-module libre.
\end{lemma}

\begin{proof}
Choisissons un morphisme \'etale $(Y', y') \to (Y, y)$, où $Y'$ est un
schéma affine. Alors $X \times_Y Y'$ est quasi-compact.
Choisissons un schéma affine $X'$ et un morphisme \'etale surjectif
$X' \to X \times_Y Y'$. On a le diagramme
$$
\xymatrix{
X \ar[d] & X' \ar[l]^g \ar[d] \\
Y & Y' \ar[l]
}
$$
Alors $\mathcal{F}' = g^*\mathcal{F}$ est plat sur $Y'$ en tout
point de $X'_{y'}$, voir Morphismes d'espaces, Lemme
\ref{spaces-morphisms-lemma-base-change-module-flat}.
Nous pouvons donc appliquer le lemme dans le cas des schémas
(Compléments sur la platitude, Lemme
\ref{flat-lemma-finite-type-flat-along-fibre-free-variant})
au morphisme
$X' \to Y'$, au faisceau quasi-cohérent $g^*\mathcal{F}$ et au point $y'$.
On obtient un morphisme \'etale $(Y'', y'') \to (Y', y')$ et un
diagramme commutatif
$$
\xymatrix{
X \ar[d] & X' \ar[l]^g \ar[d] & X'' \ar[l]^{g'} \ar[d] \\
Y & Y' \ar[l] & \Spec(\mathcal{O}_{Y'', y''}) \ar[l]
}
$$
Pour obtenir le résultat voulu, on prend $(Y'', y'') \to (Y, y)$
et $g \circ g' : X'' \to X$.
\end{proof}

\begin{theorem}
\label{spaces-flat-theorem-check-flatness-at-associated-points}
Soit $S$ un schéma.
Soit $f : X \to Y$ un morphisme d'espaces algébriques sur $S$
qui est localement de type fini.
Soit $\mathcal{F}$ un $\mathcal{O}_X$-module quasi-cohérent de type fini.
Soit $x \in |X|$, d'image $y \in |Y|$.
Posons $F = f^{-1}(\{y\}) \subset |X|$.
Considérons les conditions suivantes :
\begin{enumerate}
\item $\mathcal{F}$ est plat en $x$ sur $Y$, et
\item pour tout $x' \in F \cap \text{Ass}_{X/Y}(\mathcal{F})$ qui
se spécialise en $x$, le module $\mathcal{F}$ est plat en $x'$ sur $Y$.
\end{enumerate}
On a toujours (2) $\Rightarrow$ (1). Si $X$ et $Y$ sont
décents, alors (1) $\Rightarrow$ (2).
\end{theorem}

\begin{proof}
Supposons (2).
Choisissons un schéma $V$ et un morphisme \'etale surjectif $V \to Y$.
Choisissons un schéma $U$ et un morphisme \'etale surjectif $U \to V \times_Y X$.
Choisissons un point $u \in U$ s'envoyant sur $x$. Soit $v \in V$ l'image de $u$.
Nous déduirons le résultat du résultat correspondant pour
$\mathcal{F}|_U = (U \to X)^*\mathcal{F}$ et le point $u$.
$U_v$. Cela fonctionne car $\text{Ass}_{U/V}(\mathcal{F}|_U) \cap |U_v|$
est égal à $\text{Ass}_{U_v}(\mathcal{F}|_{U_v})$, ainsi qu'à l'image
réciproque de $F \cap \text{Ass}_{X/Y}(\mathcal{F})$.
Comme l'application $|U_v| \to F$ est continue, nous voyons que
les spécialisations dans $|U_v|$ s'envoient sur des spécialisations dans $F$ ;
la condition (2) est donc héritée par $U \to V$,
$\mathcal{F}|_U$ et le point $u$.
Ainsi Compléments sur la platitude, Théorème
\ref{flat-theorem-check-flatness-at-associated-points} s'applique
et nous concluons que (1) est satisfaite.

\medskip\noindent
Si $Y$ est décent, nous pouvons représenter
$y$ par un monomorphisme quasi-compact $\Spec(k) \to Y$
(par définition des espaces décents, voir
Espaces décents, Définition \ref{decent-spaces-definition-very-reasonable}).
Alors $F = |X_k|$, voir
Espaces décents, Lemme \ref{decent-spaces-lemma-topology-fibre}.
Si, de plus, $X$ est décent (ou, plus généralement, si $f$ est décent, voir
Espaces décents, Définition \ref{decent-spaces-definition-relative-conditions}
et Espaces décents, Lemme
\ref{decent-spaces-lemma-property-for-morphism-out-of-property}),
alors $X_y$ est lui aussi un espace décent. De plus, les spécialisations dans
$F$ se relèvent en spécialisations
dans $U_v \to X_y$, voir
Espaces décents, Lemme \ref{decent-spaces-lemma-decent-specialization}.
Cela étant, il est clair que l'implication réciproque
est vraie, puisqu'elle l'est dans le cas des schémas.
\end{proof}

\begin{lemma}
\label{spaces-flat-lemma-check-along-closed-fibre}
Soit $S$ un schéma local de point fermé $s$.
Soit $f : X \to S$ un morphisme d'un espace algébrique $X$ vers $S$
qui est localement de type fini.
Soit $\mathcal{F}$ un $\mathcal{O}_X$-module quasi-cohérent de type fini.
Supposons que
\begin{enumerate}
\item tout point de $\text{Ass}_{X/S}(\mathcal{F})$ se spécialise
en un point de la fibre fermée $X_s$\footnote{C'est par exemple le cas si
$f$ est de type fini et $\mathcal{F}$ est pur le long de $X_s$, ou
si $f$ est propre.},
\item $\mathcal{F}$ est plat sur $S$ en tout point de $X_s$.
\end{enumerate}
Alors $\mathcal{F}$ est plat sur $S$.
\end{lemma}

\begin{proof}
Cela résulte immédiatement du fait qu'il suffit de vérifier la
platitude aux points de l'assassin relatif de $\mathcal{F}$
sur $S$, d'après le
Théorème \ref{spaces-flat-theorem-check-flatness-at-associated-points}.
\end{proof}




\section{Modules plats de présentation finie}
\label{spaces-flat-section-finitely-presented-flat}

\noindent
Cette section est l'analogue de Compléments sur la platitude, section
\ref{flat-section-finitely-presented-flat}.

\begin{proposition}
\label{spaces-flat-proposition-finite-presentation-flat-at-point}
Soit $S$ un schéma.
Soit $f : X \to Y$ un morphisme d'espaces algébriques sur $S$.
Soit $\mathcal{F}$ un faisceau quasi-cohérent sur $X$.
Soit $x \in |X|$, d'image $y \in |Y|$.
Supposons que
\begin{enumerate}
\item $f$ soit localement de présentation finie,
\item $\mathcal{F}$ soit de présentation finie, et
\item $\mathcal{F}$ soit plat en $x$ sur $Y$.
\end{enumerate}
Alors il existe un diagramme commutatif de schémas pointés
$$
\xymatrix{
(X, x) \ar[d] & (X', x') \ar[l]^g \ar[d] \\
(Y, y) & (Y', y') \ar[l]
}
$$
dont les flèches horizontales sont \'etales, tel que $X'$ et $Y'$
soient affines et que
$\Gamma(X', g^*\mathcal{F})$ soit un
$\Gamma(Y', \mathcal{O}_{Y'})$-module projectif.
\end{proposition}

\begin{proof}
Sous cette forme, la proposition se réduit immédiatement
au cas des schémas, qui est
Compléments sur la platitude, Proposition
\ref{flat-proposition-finite-presentation-flat-at-point}.
\end{proof}

\begin{lemma}
\label{spaces-flat-lemma-finite-presentation-flat-along-fibre}
Soit $S$ un schéma.
Soit $f : X \to Y$ un morphisme d'espaces algébriques sur $S$.
Soit $\mathcal{F}$ un faisceau quasi-cohérent sur $X$.
Soit $y \in |Y|$. Posons $F = f^{-1}(\{y\}) \subset |X|$. Supposons que
\begin{enumerate}
\item $f$ soit de présentation finie,
\item $\mathcal{F}$ soit de présentation finie, et
\item $\mathcal{F}$ soit plat sur $Y$ en tout $x \in F$.
\end{enumerate}
Alors il existe un diagramme commutatif d'espaces algébriques
$$
\xymatrix{
X \ar[d] & X' \ar[l]^g \ar[d] \\
Y & Y' \ar[l]_h
}
$$
tel que $h$ et $g$ soient \'etales, qu'il existe un point
$y' \in |Y'|$ s'envoyant sur $y$, que $F \subset g(|X'|)$,
que les espaces algébriques $X'$ et $Y'$ soient affines, et que
$\Gamma(X', g^*\mathcal{F})$ soit un
$\Gamma(Y', \mathcal{O}_{Y'})$-module projectif.
\end{lemma}

\begin{proof}
Sous cette forme, le lemme se réduit immédiatement
au cas des schémas, qui est
Compléments sur la platitude, Lemme
\ref{flat-lemma-finite-presentation-flat-along-fibre}.
\end{proof}




\section{Un critère de pureté}
\label{spaces-flat-section-criterion-purity}

\noindent
Cette section est l'analogue de Compléments sur la platitude, section
\ref{flat-section-criterion-purity}.

\begin{lemma}
\label{spaces-flat-lemma-associated-point-specializes}
Soit $S$ un schéma. Soit $X$ un espace algébrique décent
localement de type fini sur $S$.
Soit $\mathcal{F}$ un $\mathcal{O}_X$-module quasi-cohérent de type fini.
Soit $s \in S$ tel que $\mathcal{F}$ soit plat sur $S$ en tout point de $X_s$.
Soit $x' \in \text{Ass}_{X/S}(\mathcal{F})$. Si
l'adhérence de $\{x'\}$ dans $|X|$ rencontre $|X_s|$, alors elle
rencontre $\text{Ass}_{X/S}(\mathcal{F}) \cap |X_s|$.
\end{lemma}

\begin{proof}
Observons que $|X_s| \subset |X|$ est l'ensemble des points de $|X|$
au-dessus de $s \in S$, voir
Espaces décents, Lemme \ref{decent-spaces-lemma-topology-fibre}.
Soit $t \in |X_s|$ une spécialisation de $x'$ dans $|X|$.
Choisissons un schéma affine $U$, un point $u \in U$ et
un morphisme \'etale $\varphi : U \to X$ envoyant $u$ sur $t$.
D'après Espaces décents, Lemme \ref{decent-spaces-lemma-decent-specialization},
nous pouvons choisir une spécialisation $u' \leadsto u$
telle que $u'$ s'envoie sur $x'$. Posons $g = f \circ \varphi$.
Observons que $s' = g(u') = f(x')$ se spécialise en $s$.
Par notre définition de $\text{Ass}_{X/S}(\mathcal{F})$,
nous avons $u' \in \text{Ass}_{U/S}(\varphi^*\mathcal{F})$.
La version pour les schémas de ce lemme
(Compléments sur la platitude, Lemme \ref{flat-lemma-associated-point-specializes})
montre qu'il existe une spécialisation $u' \leadsto u$ avec
$u \in \text{Ass}_{U_s}(\varphi^*\mathcal{F}_s) =
\text{Ass}_{U/S}(\varphi^*\mathcal{F}) \cap U_s$.
Ainsi $x = \varphi(u) \in \text{Ass}_{X/S}(\mathcal{F})$
est au-dessus de $s$, ce qui démontre le lemme.
\end{proof}

\begin{lemma}
\label{spaces-flat-lemma-contains-relative-ass-after-base-change}
Soit $Y$ un espace algébrique sur un schéma $S$. Soit $g : X' \to X$ un
morphisme d'espaces algébriques sur $Y$, où $X$ est localement de type fini sur $Y$.
Soit $\mathcal{F}$ un $\mathcal{O}_X$-module quasi-cohérent.
Si $\text{Ass}_{X/Y}(\mathcal{F}) \subset g(|X'|)$, alors, pour tout morphisme
$Z \to Y$, on a $\text{Ass}_{X_Z/Z}(\mathcal{F}_Z) \subset g_Z(|X'_Z|)$.
\end{lemma}

\begin{proof}
D'après Propriétés des espaces, Lemme \ref{spaces-properties-lemma-points-cartesian},
l'application $|X'_Z| \to |X_Z| \times_{|X|} |X'|$ est surjective puisque
$X'_Z$ est égal à $X_Z \times_X X'$.
D'après Diviseurs sur les espaces, Lemme
\ref{spaces-divisors-lemma-base-change-relative-assassin}
l'application $|X_Z| \to |X|$ envoie $\text{Ass}_{X_Z/Z}(\mathcal{F}_Z)$
dans $\text{Ass}_{X/Y}(\mathcal{F})$. Le lemme en résulte.
\end{proof}

\begin{lemma}
\label{spaces-flat-lemma-pure-on-top}
Soit $Y$ un espace algébrique sur un schéma $S$. Soit $g : X' \to X$ un
morphisme \'etale d'espaces algébriques sur $Y$. Supposons que les morphismes
structuraux $X' \to Y$ et $X \to Y$ soient décents et de type fini.
Soit $\mathcal{F}$ un $\mathcal{O}_X$-module quasi-cohérent de type fini.
Soit $y \in |Y|$. Posons $F = f^{-1}(\{y\}) \subset |X|$.
\begin{enumerate}
\item Si $\text{Ass}_{X/Y}(\mathcal{F}) \subset g(|X'|)$
et si $g^*\mathcal{F}$ est (universellement) pur au-dessus de $y$, alors
$\mathcal{F}$ est (universellement) pur au-dessus de $y$.
\item Si $\mathcal{F}$ est pur au-dessus de $y$, si $g(|X'|)$ contient $F$, et si
$Y$ est affine local de point fermé $y$, alors
$\text{Ass}_{X/Y}(\mathcal{F}) \subset g(|X'|)$.
\item Si $\mathcal{F}$ est pur au-dessus de $y$, si $\mathcal{F}$ est plat
en tout point de $F$, si $g(|X'|)$ contient
$\text{Ass}_{X/Y}(\mathcal{F}) \cap F$, et si $Y$ est affine local
de point fermé $y$, alors
$\text{Ass}_{X/Y}(\mathcal{F}) \subset g(|X'|)$.
\item Ajouter d'autres assertions ici.
\end{enumerate}
\end{lemma}

\begin{proof}
Les hypothèses sur $X \to Y$ et $X' \to Y$ garantissent que
nous pouvons appliquer les résultats des sections \ref{spaces-flat-section-impure} et
\ref{spaces-flat-section-pure}
à ces morphismes et aux faisceaux $\mathcal{F}$ et $g^*\mathcal{F}$.
Comme $g$ est \'etale, nous voyons que
$\text{Ass}_{X'/Y}(g^*\mathcal{F})$
est l'image réciproque de $\text{Ass}_{X/Y}(\mathcal{F})$,
et cela reste vrai après tout changement de base.

\medskip\noindent
Preuve de (1). Supposons $\text{Ass}_{X/Y}(\mathcal{F}) \subset g(|X'|)$.
Supposons que $(T \to Y, t' \leadsto t, \xi)$
soit une impureté de $\mathcal{F}$ au-dessus de $y$. Comme
$\text{Ass}_{X_T/T}(\mathcal{F}_T) \subset g_T(|X'_T|)$ d'après
le Lemme \ref{spaces-flat-lemma-contains-relative-ass-after-base-change}, nous pouvons
choisir
un point $\xi' \in |X'_T|$ s'envoyant sur $\xi$. D'après ce qui précède,
$(T \to Y, t' \leadsto t, \xi')$ est une impureté de
$g^*\mathcal{F}$ au-dessus de $y'$. Cela prouve (1).

\medskip\noindent
Preuve de (2). Cela résulte de ce que $g(|X'|)$ est ouvert
dans $|X|$ et de ce que, par pureté, tout point de
$\text{Ass}_{X/Y}(\mathcal{F})$ se spécialise en un point de $F$.

\medskip\noindent
Preuve de (3). Cela résulte de ce que $g(|X'|)$ est ouvert
dans $|X|$ et de ce que la pureté, combinée au
Lemme \ref{spaces-flat-lemma-associated-point-specializes}, assure que tout point de
$\text{Ass}_{X/Y}(\mathcal{F})$ se spécialise en un point de
$\text{Ass}_{X/Y}(\mathcal{F}) \cap F$.
\end{proof}

\begin{lemma}
\label{spaces-flat-lemma-finite-type-flat-pure-along-fibre-is-universal}
Soit $S$ un schéma. Soit $f : X \to Y$ un morphisme d'espaces
algébriques sur $S$.
Soit $\mathcal{F}$ un $\mathcal{O}_X$-module quasi-cohérent.
Soit $y \in |Y|$.
Supposons que
\begin{enumerate}
\item $f$ soit décent et de type fini,
\item $\mathcal{F}$ soit de type fini,
\item $\mathcal{F}$ soit plat sur $Y$ en tout point au-dessus de $y$, et
\item $\mathcal{F}$ soit pur au-dessus de $y$.
\end{enumerate}
Alors $\mathcal{F}$ est universellement pur au-dessus de $y$.
\end{lemma}

\begin{proof}
Considérons le morphisme $\Spec(\mathcal{O}_{Y, \overline{y}}) \to Y$.
C'est un morphisme plat partant du spectre d'un anneau local
strictement hensélien et envoyant le point fermé sur $y$.
Le Lemme \ref{spaces-flat-lemma-flat-descend-pure} nous ramène au cas
décrit dans le paragraphe suivant.

\medskip\noindent
Supposons que $Y$ soit le spectre d'un anneau local strictement hensélien $R$
de point fermé $y$.
D'après le Lemme \ref{spaces-flat-lemma-finite-type-flat-along-fibre-free-variant},
il existe un morphisme \'etale $g : X' \to X$ tel que
$g(|X'|) \supset |X_y|$, avec $X'$ affine et
$\Gamma(X', g^*\mathcal{F})$ est un module libre sur $R$.
Alors $g^*\mathcal{F}$ est universellement pur relativement à $Y$, voir
Compléments sur la platitude, Lemme \ref{flat-lemma-affine-locally-projective-pure}.
Il suffit donc de démontrer que
$g(|X'|)$ contient $\text{Ass}_{X/Y}(\mathcal{F})$,
d'après le Lemme \ref{spaces-flat-lemma-pure-on-top}, partie (1).
Cela résulte à son tour du
Lemme \ref{spaces-flat-lemma-pure-on-top}, partie (2).
\end{proof}

\begin{lemma}
\label{spaces-flat-lemma-finite-type-flat-pure-is-universal}
Soit $S$ un schéma.
Soit $f : X \to Y$ un morphisme décent et de type fini d'espaces
algébriques sur $S$.
Soit $\mathcal{F}$ un $\mathcal{O}_X$-module quasi-cohérent de type fini.
Supposons $\mathcal{F}$ plat sur $Y$. Dans ce cas,
$\mathcal{F}$ est pur relativement à $Y$ si et seulement si $\mathcal{F}$
est universellement pur relativement à $Y$.
\end{lemma}

\begin{proof}
C'est une conséquence immédiate du
Lemme \ref{spaces-flat-lemma-finite-type-flat-pure-along-fibre-is-universal}
et des définitions.
\end{proof}

\begin{lemma}
\label{spaces-flat-lemma-universally-separating}
Soit $Y$ un espace algébrique sur un schéma $S$.
Soit $g : X' \to X$ un morphisme plat d'espaces algébriques sur $Y$,
où $X$ est localement de type fini sur $Y$.
Soit $\mathcal{F}$ un $\mathcal{O}_X$-module quasi-cohérent de type fini
qui est plat sur $Y$. Si $\text{Ass}_{X/Y}(\mathcal{F}) \subset g(|X'|)$,
alors l'application canonique
$$
\mathcal{F} \longrightarrow g_*g^*\mathcal{F}
$$
est injective, et le reste après tout changement de base.
\end{lemma}

\begin{proof}
La dernière assertion signifie que $\mathcal{F}_Z \to (g_Z)_*g_Z^*\mathcal{F}_Z$
est injective pour tout morphisme $Z \to Y$. Comme l'hypothèse sur
l'assassin relatif est préservée par changement de base
(Lemme \ref{spaces-flat-lemma-contains-relative-ass-after-base-change}),
il suffit de démontrer l'injectivité de la flèche affichée.

\medskip\noindent
Soit $\mathcal{K} = \Ker(\mathcal{F} \to g_*g^*\mathcal{F})$.
Notre but est de démontrer que $\mathcal{K} = 0$.
Pour cela, il suffit de démontrer que
$\text{WeakAss}_X(\mathcal{K}) = \emptyset$, voir
Diviseurs sur les espaces, Lemme \ref{spaces-divisors-lemma-weakly-ass-zero}.
Nous avons
$\text{WeakAss}_X(\mathcal{K}) \subset \text{WeakAss}_X(\mathcal{F})$, voir
Diviseurs sur les espaces, Lemme \ref{spaces-divisors-lemma-ses-weakly-ass}.
Comme $\mathcal{F}$ est plat, le
Lemme \ref{spaces-flat-lemma-bourbaki-finite-type-general-base}
donne $\text{WeakAss}_X(\mathcal{F}) \subset \text{Ass}_{X/Y}(\mathcal{F})$.
Par hypothèse, tout point $x$ de $\text{Ass}_{X/Y}(\mathcal{F})$
est l'image d'un certain $x' \in |X'|$. Puisque $g$ est plat, l'homomorphisme
d'anneaux locaux
$\mathcal{O}_{X, \overline{x}} \to \mathcal{O}_{X', \overline{x}'}$
est fidèlement plat ; par conséquent, l'application
$$
\mathcal{F}_{\overline{x}}
\longrightarrow
(g^*\mathcal{F})_{\overline{x}'} =
\mathcal{F}_{\overline{x}} \otimes_{\mathcal{O}_{X, \overline{x}}}
\mathcal{O}_{X', \overline{x}'}
$$
est injective (voir
Algèbre, Lemme \ref{algebra-lemma-faithfully-flat-universally-injective}).
Puisque la flèche affichée se factorise par
$\mathcal{F}_{\overline{x}} \to (g_*g^*\mathcal{F})_{\overline{x}}$,
nous en concluons que
$\mathcal{K}_{\overline{x}} = 0$. Ainsi $x$ ne peut être un point faiblement
associé de $\mathcal{K}$, ce qui donne le résultat.
\end{proof}









\section{Foncteurs d'aplatissement}
\label{spaces-flat-section-F-zero}

\noindent
Cette section est l'analogue de
Compléments sur la platitude, section \ref{flat-section-flattening-functors}.
Nous conseillons au lecteur de l'omettre en première lecture.

\begin{situation}
\label{spaces-flat-situation-iso}
Soit $S$ un schéma.
Soit $f : X \to B$ un morphisme d'espaces algébriques sur $S$.
Soit $u : \mathcal{F} \to \mathcal{G}$ un homomorphisme de
$\mathcal{O}_X$-modules quasi-cohérents. Pour tout schéma $T$ sur
$B$, nous noterons $u_T : \mathcal{F}_T \to \mathcal{G}_T$ le
changement de base de $u$ à $T$ ; autrement dit, $u_T$ est l'image réciproque
de $u$ par le morphisme de projection $X_T = X \times_B T \to X$.
Dans cette situation, nous pouvons considérer le foncteur
\begin{equation}
\label{spaces-flat-equation-iso}
F_{iso} : (\Sch/B)^{opp} \longrightarrow \textit{Ens}, \quad
T \longrightarrow \left\{
\begin{matrix}
\{*\} & \text{si }u_T\text{ est un isomorphisme}, \\
\emptyset & \text{sinon.}
\end{matrix}
\right.
\end{equation}
Il existe des variantes $F_{inj}$, $F_{surj}$, $F_{zero}$ pour lesquelles on exige que
$u_T$ soit injectif, surjectif ou nul.
\end{situation}

\noindent
Dans la Situation \ref{spaces-flat-situation-iso}, nous considérons parfois les foncteurs
$F_{iso}$, $F_{inj}$, $F_{surj}$ et $F_{zero}$ comme des foncteurs
$(\Sch/S)^{opp} \to \textit{Ens}$ munis d'un morphisme
$F_{iso} \to B$, $F_{inj} \to B$, $F_{surj} \to B$ et $F_{zero} \to B$.
En effet, si $T$ est un schéma sur $S$, un élément $h \in F_{iso}(T)$
est un morphisme $h : T \to B$
tel que le changement de base de $u$ par $h$ soit un isomorphisme.
En particulier, lorsque nous disons
que $F_{iso}$ est un espace algébrique, nous entendons que le foncteur
correspondant $(\Sch/S)^{opp} \to \textit{Ens}$ est un espace algébrique.

\begin{lemma}
\label{spaces-flat-lemma-iso-sheaf}
Dans la Situation \ref{spaces-flat-situation-iso}.
Chacun des foncteurs $F_{iso}$, $F_{inj}$, $F_{surj}$, $F_{zero}$
vérifie la propriété de faisceau pour la topologie fpqc.
\end{lemma}

\begin{proof}
Soit $\{T_i \to T\}_{i \in I}$ un recouvrement fpqc de schémas sur $B$.
Posons $X_i = X_{T_i} = X \times_S T_i$ et $u_i = u_{T_i}$.
Remarquons que $\{X_i \to X_T\}_{i \in I}$ est un recouvrement fpqc de $X_T$, voir
Topologies sur les espaces, Lemme \ref{spaces-topologies-lemma-fpqc}.
En particulier, pour tout $x \in |X_T|$, il existe $i \in I$
et $x_i \in |X_i|$ s'envoyant sur $x$. Comme
$\mathcal{O}_{X_T, \overline{x}} \to \mathcal{O}_{X_i, \overline{x_i}}$
est plat, donc fidèlement plat (voir
Morphismes d'espaces, section \ref{spaces-morphisms-section-flat}),
nous concluons que $(u_i)_{x_i}$ est injectif, surjectif, bijectif ou nul
si et seulement si $(u_T)_x$ est injectif, surjectif, bijectif ou nul.
Le lemme en résulte.
\end{proof}

\begin{lemma}
\label{spaces-flat-lemma-iso-go-up}
Dans la Situation \ref{spaces-flat-situation-iso}, soit $X' \to X$ un morphisme plat
d'espaces algébriques. Notons $u' : \mathcal{F}' \to \mathcal{G}'$
l'image réciproque de $u$ sur $X'$. Notons $F'_{iso}$, $F'_{inj}$, $F'_{surj}$,
$F'_{zero}$ les foncteurs sur $\Sch/B$ associés à $u'$.
\begin{enumerate}
\item Si $\mathcal{G}$ est de type fini et si l'image de $|X'| \to |X|$
contient le support de $\mathcal{G}$, alors $F_{surj} = F'_{surj}$
et $F_{zero} = F'_{zero}$.
\item Si $\mathcal{F}$ est de type fini et si l'image de $|X'| \to |X|$
contient le support de $\mathcal{F}$, alors $F_{inj} = F'_{inj}$
et $F_{zero} = F'_{zero}$.
\item Si $\mathcal{F}$ et $\mathcal{G}$ sont de type fini et si l'image de
$|X'| \to |X|$ contient les supports de $\mathcal{F}$ et de $\mathcal{G}$,
alors $F_{iso} = F'_{iso}$.
\end{enumerate}
\end{lemma}

\begin{proof}
Soit $v : \mathcal{H} \to \mathcal{E}$ une application de modules
quasi-cohérents sur un espace algébrique $Y$, et soit $\varphi : Y' \to Y$ un
morphisme plat surjectif d'espaces algébriques. Alors $v$ est
un isomorphisme, injectif, surjectif ou nul si et seulement si $\varphi^*v$ l'est.
En effet,
pour tout $y \in |Y|$, il existe $y' \in |Y'|$, et l'homomorphisme
d'anneaux locaux
$\mathcal{O}_{Y, \overline{y}} \to \mathcal{O}_{Y', \overline{y'}}$
est fidèlement plat (voir
Morphismes d'espaces, section \ref{spaces-morphisms-section-flat}).
Bien entendu, pour vérifier l'injectivité ou la nullité, il suffit de considérer
les points du support de $\mathcal{H}$ ; pour vérifier la
surjectivité, il suffit de considérer les points du support de $\mathcal{E}$.
De plus, sous les hypothèses de type fini de l'énoncé du
lemme, la formation des supports commute au changement de base, voir
Morphismes d'espaces, Lemme \ref{spaces-morphisms-lemma-support-finite-type}.
Le lemme est donc clair.
\end{proof}

\noindent
Rappelons que nous avons défini le support schématique d'un module
quasi-cohérent de type fini dans Morphismes d'espaces, Définition
\ref{spaces-morphisms-definition-scheme-theoretic-support}.

\begin{lemma}
\label{spaces-flat-lemma-iso-limits}
Dans la Situation \ref{spaces-flat-situation-iso}.
\begin{enumerate}
\item Si $\mathcal{G}$ est de type fini et si le support schématique
de $\mathcal{G}$ est quasi-compact sur $B$, alors $F_{surj}$ commute aux
limites.
\item Si $\mathcal{F}$ est de type fini et si le support schématique
de $\mathcal{F}$ est quasi-compact sur $B$, alors
$F_{zero}$ commute aux limites.
\item Si $\mathcal{F}$ est de type fini,
si $\mathcal{G}$ est de présentation finie et si les
supports schématiques de $\mathcal{F}$ et de $\mathcal{G}$ sont
quasi-compacts sur $B$, alors $F_{iso}$ commute aux limites.
\end{enumerate}
\end{lemma}

\begin{proof}
Preuve de (1). Soit $i : Z \to X$ le support schématique de
$\mathcal{G}$, et considérons $\mathcal{G}$ comme un module quasi-cohérent
de type fini sur $Z$. Nous pouvons remplacer $X$ par $Z$ et $u$ par l'application
$i^*\mathcal{F} \to \mathcal{G}$ (détails omis). Nous pouvons donc supposer que
$f$ est quasi-compact et $\mathcal{G}$ de type fini.
Soit $T = \lim_{i \in I} T_i$ une limite dirigée de $B$-schémas affines,
et supposons $u_T$ surjectif.
Posons $X_i = X_{T_i} = X \times_S T_i$ et
$u_i = u_{T_i} : \mathcal{F}_i = \mathcal{F}_{T_i}
\to \mathcal{G}_i = \mathcal{G}_{T_i}$.
Pour démontrer (1), il faut montrer que $u_i$ est surjectif pour un certain $i$.
Choisissons $0 \in I$ et remplaçons $I$ par $\{i \mid i \geq 0\}$.
Puisque $f$ est quasi-compact, $X_0$ est quasi-compact.
Nous pouvons donc choisir un morphisme \'etale surjectif $\varphi_0 : W_0 \to X_0$,
où $W_0$ est un schéma affine. Posons $W = W_0 \times_{T_0} T$
et $W_i = W_0 \times_{T_0} T_i$ pour $i \geq 0$. Ce
sont des schémas affines munis
de morphismes \'etales surjectifs $\varphi : W \to X_T$ et
$\varphi_i : W_i \to X_i$. Remarquons que $W = \lim W_i$.
Ainsi $\varphi^*u_T$ est surjectif, et il suffit de démontrer que
$\varphi_i^*u_i$ est surjectif pour un certain $i$. Nous avons donc ramené
le problème au cas affine, qui est
Algèbre, Lemme \ref{algebra-lemma-module-map-property-in-colimit}, partie (2).

\medskip\noindent
Preuve de (2). Supposons que $\mathcal{F}$ soit de type fini, de support
schématique $Z \subset B$ quasi-compact sur $B$. Soit $T = \lim_{i \in I} T_i$
une limite dirigée de $B$-schémas affines, et supposons $u_T$ nul.
Posons $X_i = T_i \times_B X$ et notons $u_i : \mathcal{F}_i \to \mathcal{G}_i$
l'image réciproque. Choisissons $0 \in I$ et remplaçons $I$ par
$\{i \mid i \geq 0\}$. Posons $Z_0 = Z \times_X X_0$. D'après
Morphismes d'espaces, Lemme \ref{spaces-morphisms-lemma-support-finite-type},
le support de $\mathcal{F}_i$ est $|Z_0|$. Puisque $|Z_0|$ est quasi-compact,
on peut trouver un schéma affine $W_0$ et un morphisme \'etale $W_0 \to X_0$
tels que $|Z_0| \subset \Im(|W_0| \to |X_0|)$.
Posons $W = W_0 \times_{T_0} T$ et $W_i = W_0 \times_{T_0} T_i$ pour $i \geq 0$.
Ce sont des schémas affines munis
de morphismes \'etales $\varphi : W \to X_T$ et
$\varphi_i : W_i \to X_i$. Remarquons que $W = \lim W_i$
et que les supports de $\mathcal{F}_T$ et de $\mathcal{F}_i$
sont contenus respectivement dans les images de $|W| \to |X_T|$ et $|W_i| \to |X_i|$.
Alors $\varphi^*u_T$ est injectif, et il suffit de démontrer que
$\varphi_i^*u_i$ est injectif pour un certain $i$.
Nous avons ainsi ramené le problème au cas affine, qui est
Algèbre, Lemme \ref{algebra-lemma-module-map-property-in-colimit}, partie (1).

\medskip\noindent
Preuve de (3). On peut la démontrer exactement comme dans les
deux paragraphes précédents, en utilisant
Algèbre, Lemme \ref{algebra-lemma-module-map-property-in-colimit}, partie (3).
On peut aussi la déduire de (1) et (2) comme suit.
Soit $T = \lim_{i \in I} T_i$ une limite dirigée de $B$-schémas affines,
et supposons que $u_T$ soit un isomorphisme. D'après (1), il existe
un indice $0 \in I$ tel que $u_{T_0}$ soit surjectif. Posons
$\mathcal{K} = \Ker(u_{T_0})$ et considérons l'application de modules quasi-cohérents
$v : \mathcal{K} \to \mathcal{F}_{T_0}$. Pour $i \geq 0$, le changement de
base $v_{T_i}$ est nul si et seulement si $u_i$ est un isomorphisme. De plus,
$v_T$ est nul. Puisque $\mathcal{G}_{T_0}$
est de présentation finie, que $\mathcal{F}_{T_0}$ est de type fini et que
$u_{T_0}$ est surjectif, nous concluons que $\mathcal{K}$ est de type fini
(Modules sur les sites, Lemme
\ref{sites-modules-lemma-kernel-surjection-finite-onto-finite-presentation}).
Il est clair que le support de $\mathcal{K}$ est contenu dans celui
de $\mathcal{F}_{T_0}$, qui est quasi-compact sur $T_0$.
On peut donc appliquer (2) pour voir que $v_{T_i}$ est nul pour un certain $i$.
\end{proof}

\begin{lemma}
\label{spaces-flat-lemma-relate-zero-iso}
Dans la Situation \ref{spaces-flat-situation-iso}, supposons donnée une suite exacte
$$
\mathcal{F} \xrightarrow{u} \mathcal{G} \xrightarrow{v} \mathcal{H} \to 0
$$
Alors $F_{v, iso} = F_{u, zero}$, avec les notations évidentes.
\end{lemma}

\begin{proof}
Puisque l'image réciproque est exacte à droite, nous voyons que
$\mathcal{F}_T \to \mathcal{G}_T \to \mathcal{H}_T \to 0$
est exacte pour tout schéma $T$ sur $B$. Ainsi $u_T$ est
surjectif si et seulement si $v_T$ est un isomorphisme.
\end{proof}

\begin{lemma}
\label{spaces-flat-lemma-relate-zero-affine}
Dans la Situation \ref{spaces-flat-situation-iso}, supposons donnés un morphisme affine
$i : Z \to X$ et un $\mathcal{O}_Z$-module quasi-cohérent $\mathcal{H}$
tels que $\mathcal{G} = i_*\mathcal{H}$. Soit
$v : i^*\mathcal{F} \to \mathcal{H}$ l'application adjointe de $u$.
Alors
\begin{enumerate}
\item $F_{v, zero} = F_{u, zero}$, et
\item si $i$ est une immersion fermée, alors $F_{v, surj} = F_{u, surj}$.
\end{enumerate}
\end{lemma}

\begin{proof}
Soit $T$ un schéma sur $B$. Notons $i_T : Z_T \to X_T$
le changement de base de $i$, et $\mathcal{H}_T$ l'image réciproque de $\mathcal{H}$
sur $Z_T$. Observons que $(i^*\mathcal{F})_T = i_T^*\mathcal{F}_T$
et $i_{T, *}\mathcal{H}_T = (i_*\mathcal{H})_T$.
La première assertion résulte de la commutativité des images réciproques,
et la seconde de Cohomologie des espaces, Lemme
\ref{spaces-cohomology-lemma-affine-base-change}.
Nous voyons donc que $u_T$ et $v_T$ sont eux aussi des applications adjointes.
Ainsi $u_T = 0$ si et seulement si $v_T = 0$. Cela démontre (1).
Dans le cas (2), nous voyons que $u_T$ est surjectif si et seulement si
$v_T$ est surjectif, car $u_T$ se factorise sous la forme
$$
\mathcal{F}_T \to
i_{T, *}i_T^*\mathcal{F}_T \xrightarrow{i_{T, *}v_T} i_{T, *}\mathcal{H}_T
$$
et parce que $i_{T, *}$ est un foncteur exact
qui plonge pleinement fidèlement la catégorie des modules quasi-cohérents sur
$Z_T$ dans la catégorie des $\mathcal{O}_{X_T}$-modules quasi-cohérents.
Voir Morphismes d'espaces, Lemme \ref{spaces-morphisms-lemma-i-star-equivalence}.
\end{proof}

\begin{lemma}
\label{spaces-flat-lemma-relate-zero-affine-push}
Dans la Situation \ref{spaces-flat-situation-iso}, supposons donné un morphisme affine
$g : X \to X'$. Posons $u' = f_*u : f_*\mathcal{F} \to f_*\mathcal{G}$.
Alors $F_{u, iso} = F_{u', iso}$, $F_{u, inj} = F_{u', inj}$,
$F_{u, surj} = F_{u', surj}$ et $F_{u, zero} = F_{u', zero}$.
\end{lemma}

\begin{proof}
D'après Cohomologie des espaces, Lemme
\ref{spaces-cohomology-lemma-affine-base-change}
nous avons $g_{T, *}u_T = u'_T$.
De plus, $g_{T, *} : \QCoh(\mathcal{O}_{X_T}) \to \QCoh(\mathcal{O}_X)$
est un foncteur fidèle et exact qui reflète les isomorphismes, les applications
injectives et les applications surjectives.
\end{proof}

\begin{situation}
\label{spaces-flat-situation-flat}
Soit $S$ un schéma.
Soit $f : X \to Y$ un morphisme d'espaces algébriques sur $S$.
Soit $\mathcal{F}$ un $\mathcal{O}_X$-module quasi-cohérent.
Pour tout schéma $T$ sur $Y$, nous noterons $\mathcal{F}_T$ le
changement de base de $\mathcal{F}$ à $T$ ; autrement dit, $\mathcal{F}_T$
est l'image réciproque de $\mathcal{F}$ par le morphisme de projection
$X_T = X \times_Y T \to X$. Puisque tout changement de base d'un module plat
est plat, nous obtenons un foncteur
\begin{equation}
\label{spaces-flat-equation-flat}
F_{flat} : (\Sch/Y)^{opp} \longrightarrow \textit{Ens}, \quad
T \longrightarrow \left\{
\begin{matrix}
\{*\} & \text{si } \mathcal{F}_T \text{ est plat sur }T, \\
\emptyset & \text{sinon.}
\end{matrix}
\right.
\end{equation}
\end{situation}

\noindent
Dans la Situation \ref{spaces-flat-situation-flat}, nous considérons parfois $F_{flat}$ comme un foncteur
$(\Sch/S)^{opp} \to \textit{Ens}$ muni d'un morphisme
$F_{flat} \to Y$. En effet, si $T$ est un schéma sur $S$,
un élément $h \in F_{flat}(T)$ est un morphisme $h : T \to Y$
tel que le changement de base de $\mathcal{F}$ par $h$ soit plat sur $T$.
En particulier, lorsque nous disons
que $F_{flat}$ est un espace algébrique, nous entendons que le foncteur
correspondant $(\Sch/S)^{opp} \to \textit{Ens}$ est un espace algébrique.

\begin{lemma}
\label{spaces-flat-lemma-flat}
Dans la Situation \ref{spaces-flat-situation-flat}.
\begin{enumerate}
\item Le foncteur $F_{flat}$ vérifie la propriété de faisceau pour la topologie fpqc.
\item Si $f$ est quasi-compact et localement de présentation finie
et si $\mathcal{F}$ est de présentation finie, alors le foncteur
$F_{flat}$ commute aux limites.
\end{enumerate}
\end{lemma}

\begin{proof}
La partie (1) résulte de l'assertion suivante : si $T' \to T$ est un morphisme
plat surjectif d'espaces algébriques sur $Y$, alors
$\mathcal{F}_{T'}$ est plat sur $T'$ si et seulement si
$\mathcal{F}_T$ est plat sur $T$, voir
Morphismes d'espaces, Lemme \ref{spaces-morphisms-lemma-base-change-module-flat}.
La partie (2) résulte de
Limites d'espaces, Lemme \ref{spaces-limits-lemma-descend-flat},
si $f$ est en outre quasi-séparé (c'est-à-dire si $f$ est de présentation finie).
Dans le cas général, on se ramène d'abord au cas où la
base est affine, puis on recouvre $X$ par un nombre fini d'ouverts affines
pour se ramener au cas quasi-séparé. Détails omis.
\end{proof}










\section{Annulation d'une application}
\label{spaces-flat-section-zero-map}

\noindent
Cette section n'a pas d'analogue dans le chapitre correspondant sur les schémas.

\begin{situation}
\label{spaces-flat-situation-somewhat-closed}
Soit $S = \Spec(R)$ un schéma affine. Soit $X$ un espace algébrique sur
$S$. Soit $u : \mathcal{F} \to \mathcal{G}$ une application de
$\mathcal{O}_X$-modules quasi-cohérents. Supposons $\mathcal{G}$ plat sur $S$.
\end{situation}

\begin{lemma}
\label{spaces-flat-lemma-F-zero-somewhat-closed}
Dans la Situation \ref{spaces-flat-situation-somewhat-closed}. Soit $T \to S$
un morphisme quasi-compact de schémas tel que le changement de base $u_T$ soit
nul. Alors il existe un sous-schéma fermé $Z \subset S$ tel que
(a) $T \to S$ se factorise par $Z$ et (b) le changement de base $u_Z$ soit nul.
Si $\mathcal{F}$ est un $\mathcal{O}_X$-module de type fini et si
le support schématique de $\mathcal{F}$ est quasi-compact,
on peut prendre $Z \to S$ de présentation finie.
\end{lemma}

\begin{proof}
Soit $U \to X$ un morphisme \'etale surjectif d'espaces algébriques,
où $U = \coprod U_i$ est une réunion disjointe de schémas affines (voir
Propriétés des espaces, Lemme
\ref{spaces-properties-lemma-cover-by-union-affines}).
D'après le Lemme \ref{spaces-flat-lemma-iso-go-up}, nous pouvons
remplacer $X$ par $U$. Autrement dit, nous pouvons supposer que $X = \coprod X_i$
est une réunion disjointe de schémas affines $X_i$. Supposons que nous sachions démontrer
le lemme pour $u_i = u|_{X_i}$. Nous obtenons alors un sous-schéma fermé
$Z_i \subset S$ tel que $T \to S$ se factorise par $Z_i$ et que
$u_{i, Z_i}$ soit nul. Si
$Z_i = \Spec(R/I_i) \subset \Spec(R) = S$, alors prendre
$Z = \Spec(R/\sum I_i)$ convient. Nous pouvons donc supposer que
$X = \Spec(A)$ est affine.

\medskip\noindent
Choisissons un recouvrement ouvert affine fini $T = T_1 \cup \ldots \cup T_m$.
Il est clair que nous pouvons remplacer $T$ par $\coprod_{j = 1, \ldots, m} T_j$.
Nous pouvons donc supposer $T$ affine. Écrivons $T = \Spec(R')$.
Soit $u : M \to N$ l'homomorphisme de $A$-modules
correspondant à $u : \mathcal{F} \to \mathcal{G}$.
Alors $N$ est un $R$-module plat puisque $\mathcal{G}$ est plat sur $S$.
L'hypothèse du lemme signifie que le composé
$$
M \otimes_R R' \to N \otimes_R R'
$$
est nul. Soit $z \in M$. D'après le théorème de Lazard
(Algèbre, Théorème \ref{algebra-theorem-lazard}) et puisque
$\otimes$ commute aux colimites, on peut trouver un $R$-module libre
$F_z$, un élément $\tilde z \in F_z$ et une application $F_z \to N$ tels que
$u(z)$ soit l'image de $\tilde z$ et que $\tilde z$ s'envoie sur zéro dans
$F_z \otimes_R R'$. Choisissons une base $\{e_{z, \alpha}\}$ de $F_z$ et écrivons
$\tilde z = \sum f_{z, \alpha} e_{z, \alpha}$ avec $f_{z, \alpha} \in R$.
Soit $I \subset R$ l'idéal engendré par les éléments $f_{z, \alpha}$,
où $z$ parcourt tous les éléments de $M$.
Par construction, $I$ s'envoie sur zéro dans $R'$ et les éléments $\tilde z$
s'envoient sur zéro dans $F_z/IF_z$, donc dans $N/IN$. Ainsi $Z = \Spec(R/I)$
résout le problème dans ce cas.

\medskip\noindent
Supposons que $\mathcal{F}$ soit de type fini, de support schématique
quasi-compact. Écrivons $Z = \Spec(R/I)$.
Écrivons $I = \bigcup I_\lambda$ comme réunion filtrante d'idéaux de type fini.
Posons $Z_\lambda = \Spec(R/I_\lambda)$, de sorte que $Z = \colim Z_\lambda$.
Puisque $u_Z$ est nul, nous voyons que $u_{Z_\lambda}$ est nul
pour un certain $\lambda$, d'après le Lemme \ref{spaces-flat-lemma-iso-limits}.
Cela achève la démonstration du lemme.
\end{proof}

\begin{lemma}
\label{spaces-flat-lemma-F-zero-module-map}
Soit $A$ un anneau. Soit $u : M \to N$ une application de $A$-modules.
Si $N$ est projectif comme $A$-module, alors il existe un idéal
$I \subset A$ tel que, pour tout homomorphisme d'anneaux $\varphi : A \to B$,
les conditions suivantes soient équivalentes :
\begin{enumerate}
\item $u \otimes 1 : M \otimes_A B \to N \otimes_A B$ est nul, et
\item $\varphi(I) = 0$.
\end{enumerate}
\end{lemma}

\begin{proof}
Puisque $N$ est projectif, on peut trouver un $A$-module projectif $C$
tel que $F = N \oplus C$ soit un $R$-module libre.
En remplaçant $u$ par $u \oplus 1 : F = M \oplus C \to N \oplus C$,
nous voyons que nous pouvons supposer $N$ libre. Dans ce cas, soit $I$
l'idéal de $A$ engendré par les coefficients de tous les éléments de
$\Im(u)$ relativement à une base (fixée) de $N$.
\end{proof}

\begin{lemma}
\label{spaces-flat-lemma-F-zero-somewhat-closed-points}
Dans la Situation \ref{spaces-flat-situation-somewhat-closed}.
Soit $T \subset S$ une partie. Soit $s \in S$ dans l'adhérence de $T$.
Pour $t \in T$, soit $u_t$ l'image réciproque de $u$ sur $X_t$,
et soit $u_s$ l'image réciproque de $u$ sur $X_s$.
Si $X$ est localement de présentation finie sur $S$,
si $\mathcal{G}$ est de présentation finie\footnote{Il suffirait
que $X$ soit localement de type fini sur $S$
et que $\mathcal{G}$ soit de présentation finie relativement à $S$,
mais cette notion n'a pas encore été définie dans le cadre
des espaces algébriques. La définition pour les schémas est
donnée dans Compléments sur les morphismes, section
\ref{more-morphisms-section-finite-type-finite-presentation}.}, et si
$u_t = 0$ pour tout $t \in T$, alors $u_s = 0$.
\end{lemma}

\begin{proof}
Vérifier si $u_s$ est nul est une question locale pour la topologie \'etale sur la fibre $X_s$.
Nous pouvons donc choisir un point $x \in |X_s| \subset |X|$ et effectuer la vérification
dans un voisinage \'etale. Choisissons
$$
\xymatrix{
(X, x) \ar[d] & (X', x') \ar[l]^g \ar[d] \\
(S, s) & (S', s') \ar[l]
}
$$
comme dans la Proposition \ref{spaces-flat-proposition-finite-presentation-flat-at-point}.
Soit $T' \subset S'$ l'image réciproque de $T$. Observons que
$s'$ appartient à l'adhérence de $T'$, car $S' \to S$ est ouvert.
Nous sommes donc ramenés au problème d'algèbre décrit dans le
paragraphe suivant.

\medskip\noindent
Nous avons une application de $R$-modules $u : M \to N$ telle que $N$ soit projectif
comme $R$-module et telle que
$u_t : M \otimes_R \kappa(t) \to N \otimes_R \kappa(t)$
soit nulle pour tout $t \in T$. Il s'agit de montrer que $u_s = 0$.
Soit $I \subset R$ l'idéal défini dans le Lemme \ref{spaces-flat-lemma-F-zero-module-map}.
Alors $I$ s'envoie sur zéro dans $\kappa(t)$ pour tout $t \in T$.
Ainsi $T \subset V(I)$. Puisque $s$ appartient à l'adhérence de $T$,
nous avons $s \in V(I)$. Par conséquent, $u_s = 0$.
\end{proof}

\noindent
Il serait intéressant de trouver une démonstration directe « simple » de l'un des
Lemmes \ref{spaces-flat-lemma-F-zero-closed-pure} ou
\ref{spaces-flat-lemma-F-zero-closed-proper},
en utilisant des arguments analogues à ceux des
Lemmes \ref{spaces-flat-lemma-F-zero-somewhat-closed} et
\ref{spaces-flat-lemma-F-zero-somewhat-closed-points}.
Une démonstration « classique » de ce lemme lorsque $f : X \to B$ est un morphisme
projectif et $B$ un schéma noethérien serait la suivante : (a) choisir un faisceau inversible
relativement ample $\mathcal{O}_X(1)$, (b) poser
$u_n : f_*\mathcal{F}(n) \to f_*\mathcal{G}(n)$,
(c) observer que $f_*\mathcal{G}(n)$ est un faisceau localement libre de rang fini
pour tout $n \gg 0$, et (d) que $F_{zero}$ est représenté par le lieu des zéros
de $u_n$ pour un certain $n \gg 0$.

\begin{lemma}
\label{spaces-flat-lemma-F-zero-closed-pure}
Dans la Situation \ref{spaces-flat-situation-iso}. Supposons que
\begin{enumerate}
\item $f$ soit de présentation finie, et
\item $\mathcal{G}$ soit de présentation finie,
plat sur $B$ et pur relativement à $B$.
\end{enumerate}
Alors $F_{zero}$ est un espace algébrique et $F_{zero} \to B$
est une immersion fermée. Si $\mathcal{F}$ est de type fini, alors
$F_{zero} \to B$ est de présentation finie.
\end{lemma}

\begin{proof}
D'après le Lemme \ref{spaces-flat-lemma-finite-type-flat-pure-is-universal},
le module $\mathcal{G}$ est universellement pur relativement à $B$.
Pour démontrer que $F_{zero}$ est un espace algébrique,
il suffit de montrer que $F_{zero} \to B$ est représentable, voir
Espaces, Lemme \ref{spaces-lemma-representable-over-space}.
Soit $B' \to B$ un morphisme, où $B'$ est un schéma, et soit
$u' : \mathcal{F}' \to \mathcal{G}'$ l'image réciproque de $u$ sur $X' = X_{B'}$.
Alors le foncteur associé $F'_{zero}$ est égal à $F_{zero} \times_B B'$.
Cela nous ramène au cas où $B$ est un schéma.

\medskip\noindent
Supposons que $B$ soit un schéma. Nous allons montrer que $F_{zero}$ est représentable
par un sous-schéma fermé de $B$. D'après le Lemme \ref{spaces-flat-lemma-iso-sheaf} et
Descente, Lemmes \ref{descent-lemma-closed-immersion} et
\ref{descent-lemma-descent-data-sheaves}
la question est locale pour la topologie \'etale sur $B$. Soit $b \in B$.
Remplaçons d'abord $B$ par un voisinage affine de $b$.
Choisissons un diagramme
$$
\xymatrix{
X \ar[d] & X' \ar[l]^g \ar[d] \\
B & B' \ar[l]
}
$$
et un point $b' \in B'$ s'envoyant sur $b \in B$,
comme dans le Lemme \ref{spaces-flat-lemma-finite-presentation-flat-along-fibre}.
Puisque nous travaillons localement pour la topologie \'etale, nous pouvons remplacer
$B$ par $B'$ et supposer que nous disposons d'un diagramme
$$
\xymatrix{
X \ar[rd] & & X' \ar[ll]^g \ar[ld] \\
& B
}
$$
où $B$ et $X'$ sont affines, tel que $\Gamma(X', g^*\mathcal{G})$
soit un $\Gamma(B, \mathcal{O}_B)$-module projectif et que
$g(|X'|) \supset |X_b|$. Soit $U \subset X$ le sous-espace ouvert
tel que $|U| = g(|X'|)$. D'après
Diviseurs sur les espaces, Lemme
\ref{spaces-divisors-lemma-relative-assassin-constructible}, l'ensemble
$$
E = \{t \in B : \text{Ass}_{X_t}(\mathcal{G}_t) \subset |U_t|\} =
\{t \in B : \text{Ass}_{X/B}(\mathcal{G}) \cap |X_t| \subset |U_t|\}
$$
est constructible dans $B$. D'après le Lemme \ref{spaces-flat-lemma-pure-on-top}, partie (2),
nous voyons que $E$ contient $\Spec(\mathcal{O}_{B, b})$. D'après
Morphismes, Lemme \ref{morphisms-lemma-constructible-containing-open},
nous voyons que $E$ contient un voisinage ouvert de $b$. Ainsi, après avoir
remplacé $B$ par un voisinage affine plus petit de $b$, nous pouvons supposer que
$\text{Ass}_{X/B}(\mathcal{G}) \subset g(|X'|)$.

\medskip\noindent
Il résulte du Lemme \ref{spaces-flat-lemma-universally-separating}
que $u : \mathcal{F} \to \mathcal{G}$ est injectif si et seulement si
$g^*u : g^*\mathcal{F} \to g^*\mathcal{G}$ est injectif, et cela reste
vrai après tout changement de base. Nous nous sommes donc ramenés au cas où,
outre les hypothèses du théorème, $X \to B$ est un morphisme de
schémas affines et $\Gamma(X, \mathcal{G})$ est un
$\Gamma(B, \mathcal{O}_B)$-module projectif. Ce cas résulte immédiatement du
Lemme \ref{spaces-flat-lemma-F-zero-module-map}.

\medskip\noindent
Il reste à montrer que $F_{zero} \to B$ est de présentation finie si
$\mathcal{F}$ est de type fini. Cela résulte du
Lemme \ref{spaces-flat-lemma-iso-limits}, combiné à
Limites d'espaces, Proposition
\ref{spaces-limits-proposition-characterize-locally-finite-presentation}.
\end{proof}

\begin{lemma}
\label{spaces-flat-lemma-F-zero-closed-proper}
Dans la Situation \ref{spaces-flat-situation-iso}. Supposons que
\begin{enumerate}
\item $f$ soit localement de présentation finie,
\item $\mathcal{G}$ soit un $\mathcal{O}_X$-module de présentation finie
plat sur $B$,
\item le support de $\mathcal{G}$ soit propre sur $B$.
\end{enumerate}
Alors le foncteur $F_{zero}$ est un espace algébrique et $F_{zero} \to B$
est une immersion fermée. Si $\mathcal{F}$ est de type fini, alors
$F_{zero} \to B$ est de présentation finie.
\end{lemma}

\begin{proof}
Si $f$ est de présentation finie, cela résulte immédiatement des
Lemmes \ref{spaces-flat-lemma-F-zero-closed-pure} et \ref{spaces-flat-lemma-proper-pure}.
C'est le seul cas intéressant, et nous conseillons au lecteur d'omettre
la suite de la démonstration, qui traite de la possibilité (autorisée
par les hypothèses du présent lemme)
que $f$ ne soit ni quasi-séparé ni quasi-compact.

\medskip\noindent
Soit $i : Z \to X$ le sous-espace fermé défini par l'idéal de Fitting
d'indice zéro de $\mathcal{G}$
(Diviseurs sur les espaces, section
\ref{spaces-divisors-section-fitting-ideals}).
Alors $Z \to B$ est propre par hypothèse (voir
Catégories dérivées des espaces, section
\ref{spaces-perfect-section-proper-over-base}).
D'autre part, $i$ est de présentation finie
(Diviseurs sur les espaces, Lemme
\ref{spaces-divisors-lemma-fitting-ideal-of-finitely-presented}, et
Morphismes d'espaces, Lemme
\ref{spaces-morphisms-lemma-closed-immersion-finite-presentation}).
Il existe un $\mathcal{O}_Z$-module quasi-cohérent
$\mathcal{H}$ de type fini tel que $i_*\mathcal{H} = \mathcal{G}$
(Diviseurs sur les espaces, Lemme
\ref{spaces-divisors-lemma-on-subscheme-cut-out-by-Fit-0}).
En fait, $\mathcal{H}$ est de présentation finie comme $\mathcal{O}_Z$-module
d'après Algèbre, Lemme \ref{algebra-lemma-finitely-presented-over-subring}
(détails omis).
Alors $F_{zero}$ est le même que le foncteur $F_{zero}$
pour l'application $i^*\mathcal{F} \to \mathcal{H}$ adjointe de $u$, voir
Lemme \ref{spaces-flat-lemma-relate-zero-affine}.
Le faisceau $\mathcal{H}$ est plat sur $B$, car
il en est de même de $\mathcal{G}$ (vérifier sur les germes ; détails omis).
De plus, remarquons que si $\mathcal{F}$ est de type fini,
alors $i^*\mathcal{F}$ est de type fini.
Nous avons donc ramené le lemme au cas
étudié dans le premier paragraphe de la démonstration.
\end{proof}








\section{Aplatissement d'une application}
\label{spaces-flat-section-flattening-map}

\noindent
Cette section est l'analogue de
Compléments sur la platitude, section \ref{flat-section-flattening-map}.
En particulier, le résultat suivant est une variante de
Compléments sur la platitude, Théorème \ref{flat-theorem-flattening-map}.

\begin{theorem}
\label{spaces-flat-theorem-flattening-map}
Dans la Situation \ref{spaces-flat-situation-iso}, supposons que
\begin{enumerate}
\item $f$ soit de présentation finie,
\item $\mathcal{F}$ soit de présentation finie, plat sur $B$ et
pur relativement à $B$, et
\item $u$ soit surjectif.
\end{enumerate}
Alors $F_{iso}$ est représentable par une immersion fermée $Z \to B$.
De plus, $Z \to S$ est de présentation finie si $\mathcal{G}$ est
de présentation finie.
\end{theorem}

\begin{proof}
Soit $\mathcal{K} = \Ker(u)$ et notons $v : \mathcal{K} \to \mathcal{F}$
l'inclusion. Le Lemme \ref{spaces-flat-lemma-relate-zero-iso} donne
$F_{u, iso} = F_{v, zero}$. En appliquant le Lemme \ref{spaces-flat-lemma-F-zero-closed-pure}
à $v$, nous voyons que $F_{u, iso} = F_{v, zero}$ est représentable
par un sous-espace fermé de $B$. Remarquons que $\mathcal{K}$ est de type fini
si $\mathcal{G}$ est de présentation finie, voir
Modules sur les sites, Lemme
\ref{sites-modules-lemma-kernel-surjection-finite-onto-finite-presentation}.
Nous obtenons donc également la dernière assertion du théorème.
\end{proof}

\begin{lemma}
\label{spaces-flat-lemma-F-iso-closed}
Dans la Situation \ref{spaces-flat-situation-iso}. Supposons que
\begin{enumerate}
\item $f$ soit localement de présentation finie,
\item $\mathcal{F}$ soit localement de présentation finie et plat sur $B$,
\item le support de $\mathcal{F}$ soit propre sur $B$, et
\item $u$ soit surjectif.
\end{enumerate}
Alors le foncteur $F_{iso}$ est un espace algébrique et $F_{iso} \to B$
est une immersion fermée. Si $\mathcal{G}$ est de présentation finie, alors
$F_{iso} \to B$ est de présentation finie.
\end{lemma}

\begin{proof}
Soit $\mathcal{K} = \Ker(u)$ et notons $v : \mathcal{K} \to \mathcal{F}$
l'inclusion. Le Lemme \ref{spaces-flat-lemma-relate-zero-iso} donne
$F_{u, iso} = F_{v, zero}$. En appliquant le Lemme \ref{spaces-flat-lemma-F-zero-closed-proper}
à $v$, nous voyons que $F_{u, iso} = F_{v, zero}$ est représentable
par un sous-espace fermé de $B$. Remarquons que $\mathcal{K}$ est de type fini
si $\mathcal{G}$ est de présentation finie, voir
Modules sur les sites, Lemme
\ref{sites-modules-lemma-kernel-surjection-finite-onto-finite-presentation}.
Nous obtenons donc également la dernière assertion du lemme.
\end{proof}

\noindent
Nous utiliserons le résultat (facile) suivant dans l'étude du foncteur Quot.

\begin{lemma}
\label{spaces-flat-lemma-F-surj-open}
Dans la Situation \ref{spaces-flat-situation-iso}. Supposons que
\begin{enumerate}
\item $f$ soit localement de présentation finie,
\item $\mathcal{G}$ soit de type fini,
\item le support de $\mathcal{G}$ soit propre sur $B$.
\end{enumerate}
Alors $F_{surj}$ est un espace algébrique et $F_{surj} \to B$
est une immersion ouverte.
\end{lemma}

\begin{proof}
Considérons $\Coker(u)$. Observons que
$\Coker(u_T) = \Coker(u)_T$ pour tout $T/B$.
Remarquons que la formation du support d'un module
quasi-cohérent de type fini commute à l'image réciproque
(Morphismes d'espaces, Lemme \ref{spaces-morphisms-lemma-support-covering}).
Ainsi $F_{surj}$ est représentable par le sous-espace ouvert de $B$
correspondant à l'ouvert
$$
|B| \setminus |f|(\text{Supp}(\Coker(u)))
$$
voir Propriétés des espaces, Lemme \ref{spaces-properties-lemma-open-subspaces}.
C'est un ouvert, car $|f|$ est fermé sur $\text{Supp}(\mathcal{G})$
et $\text{Supp}(\Coker(u))$ est une partie fermée de
$\text{Supp}(\mathcal{G})$.
\end{proof}






\section{Aplatissement dans le cas local}
\label{spaces-flat-section-flattening-local}

\noindent
Cette section est l'analogue de Compléments sur la platitude, section
\ref{flat-section-flattening-local}.

\begin{lemma}
\label{spaces-flat-lemma-freebie}
Soit $S$ le spectre d'un anneau local hensélien de point fermé $s$.
Soit $X \to S$ un morphisme d'espaces algébriques
localement de type fini.
Soit $\mathcal{F}$ un $\mathcal{O}_X$-module quasi-cohérent de type fini.
Soit $E \subset |X_s|$ une partie. Il existe un sous-schéma fermé
$Z \subset S$ ayant la propriété suivante : pour tout morphisme de schémas
pointés $(T, t) \to (S, s)$, les conditions suivantes sont équivalentes :
\begin{enumerate}
\item $\mathcal{F}_T$ est plat sur $T$ en tout point de
$|X_t|$ qui s'envoie sur un point de $E \subset |X_s|$, et
\item $\Spec(\mathcal{O}_{T, t}) \to S$ se factorise par $Z$.
\end{enumerate}
De plus, si $X \to S$ est localement de présentation finie,
si $\mathcal{F}$ est de présentation finie et si $E \subset |X_s|$ est
fermé et quasi-compact, alors $Z \to S$ est de présentation finie.
\end{lemma}

\begin{proof}
Choisissons un schéma $U$ et un morphisme \'etale $\varphi : U \to X$.
Soit $E' \subset |U_s|$ l'image réciproque de $E$. Si
$E' \to E$ est surjectif, alors la condition (1) équivaut à ceci :
$(\varphi^*\mathcal{F})_T$ est plat sur $T$ en tout point de
$|U_t|$ qui s'envoie sur un point de $E' \subset |U_t|$.
En choisissant $\varphi$ surjectif, nous sommes ramenés au cas des schémas, à savoir
Compléments sur la platitude, Lemme \ref{flat-lemma-freebie}.
Si $E$ est fermé et quasi-compact, nous pouvons choisir $U$
affine de telle sorte que $E' \to E$ soit surjectif. Alors $E'$ est fermé
et quasi-compact, et la dernière assertion résulte de la
dernière assertion de
Compléments sur la platitude, Lemme \ref{flat-lemma-freebie}.
\end{proof}









\section{Aplatissement universel}
\label{spaces-flat-section-flattening-final}

\noindent
Cette section est l'analogue de
Compléments sur la platitude, section \ref{flat-section-flattening-final}.
Notre but principal est de démontrer le
Lemme \ref{spaces-flat-lemma-when-universal-flattening}.
Cependant, nous ne voyons pas comment déduire directement ce résultat
du résultat correspondant pour les schémas.
Nous devons donc reprendre ici une partie du développement.
Commençons toutefois par une définition.

\begin{definition}
\label{spaces-flat-definition-flattening}
Soit $S$ un schéma.
Soit $X \to Y$ un morphisme d'espaces algébriques sur $S$.
Soit $\mathcal{F}$ un $\mathcal{O}_X$-module quasi-cohérent.
Nous disons que {\it l'aplatissement universel de $\mathcal{F}$ existe}
si le foncteur $F_{flat}$ défini dans la Situation \ref{spaces-flat-situation-flat}
est un espace algébrique.
Nous disons que {\it l'aplatissement universel de $X$ existe}
si l'aplatissement universel de $\mathcal{O}_X$ existe.
\end{definition}

\noindent
Cette définition est quelque peu insatisfaisante, car la définition de
l'aplatissement universel ne coïncide pas ici avec celle utilisée dans le
cas des schémas, puisque nous ignorons si tout monomorphisme d'espaces
algébriques est représentable (Compléments sur les morphismes d'espaces,
section \ref{spaces-more-morphisms-section-monomorphisms}).
Espérons qu'il n'en résultera jamais de confusion.

\begin{lemma}
\label{spaces-flat-lemma-pre-flat-dimension-n}
Soit $S$ un schéma.
Soit $f : X \to Y$ un morphisme d'espaces algébriques
localement de type fini.
Soit $\mathcal{F}$ un $\mathcal{O}_X$-module quasi-cohérent de
type fini. Soit $n \geq 0$. Les conditions suivantes sont équivalentes :
\begin{enumerate}
\item pour un certain diagramme commutatif
$$
\xymatrix{
U \ar[d]_\varphi \ar[r] & V \ar[d] \\
X \ar[r] & Y
}
$$
dont les flèches verticales sont \'etales et surjectives, où $U$ et $V$ sont
des schémas, le faisceau $\varphi^*\mathcal{F}$ est plat sur $V$
en dimensions $\geq n$ (Compléments sur la platitude, Définition
\ref{flat-definition-flat-dimension-n}),
\item pour tout diagramme commutatif
$$
\xymatrix{
U \ar[d]_\varphi \ar[r] & V \ar[d] \\
X \ar[r] & Y
}
$$
dont les flèches verticales sont \'etales, où $U$ et $V$ sont des schémas,
le faisceau $\varphi^*\mathcal{F}$ est plat sur $V$ en dimensions $\geq n$, et
\item pour $x \in |X|$ tel que $\mathcal{F}$ ne soit pas plat en $x$
sur $Y$, le degré de transcendance de $x/f(x)$ est $< n$ (Morphismes d'espaces,
Définition \ref{spaces-morphisms-definition-dimension-fibre}).
\end{enumerate}
Si ces conditions sont satisfaites, elles le restent après tout changement de base $Y' \to Y$.
\end{lemma}

\begin{proof}
Supposons disposer d'un diagramme comme en (1). L'équivalence des
conditions dans Compléments sur la platitude, Lemme \ref{flat-lemma-pre-flat-dimension-n},
montre alors que (1) et (3) sont équivalentes. Mais la condition (3) est héritée
par $\varphi^*\mathcal{F}$ pour tout $U \to V$ comme en (2).
Le résultat pour les schémas montre donc à nouveau que (3) implique (2).
Il implique également l'assertion relative au changement de base.
\end{proof}

\begin{definition}
\label{spaces-flat-definition-flat-dimension-n}
Soit $S$ un schéma. Soit $f : X \to Y$ un morphisme d'espaces algébriques
sur $S$ qui est localement de type fini.
Soit $\mathcal{F}$ un $\mathcal{O}_X$-module quasi-cohérent de type fini.
Soit $n \geq 0$.
Nous disons que {\it $\mathcal{F}$ est plat sur $Y$ en dimensions $\geq n$}
si les conditions équivalentes du Lemme \ref{spaces-flat-lemma-pre-flat-dimension-n}
sont satisfaites.
\end{definition}

\begin{situation}
\label{spaces-flat-situation-flat-dimension-n}
Soit $S$ un schéma. Soit $f : X \to Y$ un morphisme d'espaces algébriques
sur $S$ qui est localement de type fini. Soit $\mathcal{F}$ un
$\mathcal{O}_X$-module quasi-cohérent de type fini. Pour tout schéma $T$ sur $Y$, nous
noterons $\mathcal{F}_T$ le changement de base de $\mathcal{F}$ à $T$ ; autrement dit,
$\mathcal{F}_T$ est l'image réciproque de $\mathcal{F}$ par le morphisme de projection
$X_T = X \times_Y T \to X$. Remarquons que $f_T : X_T \to T$ est de type fini
et que $\mathcal{F}_T$ est un $\mathcal{O}_{X_T}$-module de type fini
(Morphismes d'espaces, Lemme
\ref{spaces-morphisms-lemma-base-change-finite-type}
et
Modules sur les sites, Lemme \ref{sites-modules-lemma-local-pullback}).
Soit $n \geq 0$. La Définition \ref{spaces-flat-definition-flat-dimension-n} et le
Lemme \ref{spaces-flat-lemma-pre-flat-dimension-n} fournissent un foncteur
\begin{equation}
\label{spaces-flat-equation-flat-dimension-n}
F_n : (\Sch/Y)^{opp} \longrightarrow \textit{Ens}, \quad
T \longrightarrow \left\{
\begin{matrix}
\{*\} & \text{si }\mathcal{F}_T\text{ est plat sur }T\text{ en }\dim \geq n, \\
\emptyset & \text{sinon.}
\end{matrix}
\right.
\end{equation}
\end{situation}

\noindent
Dans la Situation \ref{spaces-flat-situation-flat-dimension-n}, nous considérons parfois
$F_n$ comme un foncteur $(\Sch/S)^{opp} \to \textit{Ens}$ muni d'un morphisme
$F_n \to Y$. En effet, si $T$ est un schéma sur $S$,
un élément $h \in F_n(T)$ est un morphisme $h : T \to Y$
tel que le changement de base de $\mathcal{F}$ par $h$ soit plat sur $T$
en $\dim \geq n$. En particulier, lorsque nous disons
que $F_n$ est un espace algébrique, nous entendons que le foncteur
correspondant $(\Sch/S)^{opp} \to \textit{Ens}$ est un espace algébrique.

\begin{lemma}
\label{spaces-flat-lemma-flat-dimension-n}
Dans la Situation \ref{spaces-flat-situation-flat-dimension-n}.
\begin{enumerate}
\item Le foncteur $F_n$ vérifie la propriété de faisceau pour la topologie fpqc.
\item Si $f$ est quasi-compact et localement de présentation finie
et si $\mathcal{F}$ est de présentation finie, alors le foncteur $F_n$
commute aux limites.
\end{enumerate}
\end{lemma}

\begin{proof}
Preuve de (1). Supposons que $\{T_i \to T\}$ soit un recouvrement fpqc
d'un schéma $T$ sur $Y$. Nous devons montrer que, si $F_n(T_i)$ est non vide
pour tout $i$, alors $F_n(T)$ est non vide.
Choisissons un diagramme comme dans la partie (1) du Lemme \ref{spaces-flat-lemma-pre-flat-dimension-n}.
Notons $F'_n$ le foncteur correspondant pour
$\varphi^*\mathcal{F}$ et le morphisme $U \to V$.
D'après Compléments sur la platitude, Lemme \ref{flat-lemma-flat-dimension-n},
$F'_n$ vérifie la propriété de faisceau.
Il en est donc de même de $F_n$, car
pour $T \to Y$, on a $F_n(T) = F'_n(V \times_Y T)$
d'après le Lemme \ref{spaces-flat-lemma-pre-flat-dimension-n},
et parce que $\{V \times_Y T_i \to V \times_Y T\}$
est un recouvrement fpqc.

\medskip\noindent
Preuve de (2). Supposons que $T = \lim_{i \in I} T_i$ soit une limite filtrante
de schémas affines $T_i$ sur $Y$, et supposons que $F_n(T)$ soit non vide.
Nous devons montrer que $F_n(T_i)$ est non vide pour un certain $i$.
Choisissons un diagramme comme dans la
partie (1) du Lemme \ref{spaces-flat-lemma-pre-flat-dimension-n}.
Fixons $i \in I$ et choisissons un ouvert affine $W_i \subset V \times_Y T_i$
s'envoyant surjectivement sur $T_i$. Pour $i' \geq i$, soit $W_{i'}$
l'image réciproque de $W_i$ dans $V \times_Y T_{i'}$, et
soit $W \subset V \times_Y T$ l'image réciproque de $W_i$.
Alors $W = \lim_{i' \geq i} W_i$ est une limite filtrante de schémas
affines sur $V$. À nouveau, d'après le
Lemme \ref{spaces-flat-lemma-pre-flat-dimension-n},
il suffit de montrer que $F'_n(W_{i'})$ est non vide pour un certain
$i' \geq i$. Or nous savons que $F'_n(W)$ est non vide,
car, par hypothèse, $F_n(T) = F'_n(V \times_Y T)$
est non vide. Nous pouvons donc appliquer
Compléments sur la platitude, Lemme \ref{flat-lemma-flat-dimension-n},
et conclure.
\end{proof}

\begin{lemma}
\label{spaces-flat-lemma-localize-flat-dimension-n}
Dans la Situation \ref{spaces-flat-situation-flat-dimension-n}.
Soit $h : X' \to X$ un morphisme \'etale.
Posons $\mathcal{F}' = h^*\mathcal{F}$ et $f' = f \circ h$.
Soit $F_n'$ le foncteur (\ref{spaces-flat-equation-flat-dimension-n})
associé à $(f' : X' \to Y, \mathcal{F}')$.
Alors $F_n$ est un sous-foncteur de $F_n'$ et, si
$h(X') \supset \text{Ass}_{X/Y}(\mathcal{F})$, alors $F_n = F'_n$.
\end{lemma}

\begin{proof}
Choisissons $U \to X$, $V \to Y$, $U \to V$ comme dans la partie (1) du
Lemme \ref{spaces-flat-lemma-pre-flat-dimension-n}. Choisissons un morphisme \'etale
surjectif $U' \to U \times_X X'$, où $U'$ est un schéma.
Nous disposons alors du lemme pour les deux foncteurs
$F_{U, n}$ et $F_{U', n}$ déterminés par $U' \to U$ et $\mathcal{F}|_U$
sur $V$, voir
Compléments sur la platitude, Lemme \ref{flat-lemma-localize-flat-dimension-n}.
D'autre part, le Lemme \ref{spaces-flat-lemma-pre-flat-dimension-n}
nous dit que, pour tout $T \to Y$, on a
$F_n(T) = F_{U, n}(V \times_Y T)$
et
$F'_n(T) = F_{U', n}(V \times_Y T)$.
Cela démontre le lemme.
\end{proof}

\begin{theorem}
\label{spaces-flat-theorem-flat-dimension-n-representable}
Dans la Situation \ref{spaces-flat-situation-flat-dimension-n}.
Supposons en outre que $f$ soit de présentation finie, que
$\mathcal{F}$ soit un $\mathcal{O}_X$-module de présentation finie
et que $\mathcal{F}$ soit pur relativement à $Y$.
Alors $F_n$ est un espace algébrique et
$F_n \to Y$ est un monomorphisme de présentation finie.
\end{theorem}

\begin{proof}
Le foncteur $F_n$ est un faisceau pour la topologie fppf d'après le
Lemme \ref{spaces-flat-lemma-flat-dimension-n}.
Puisque $F_n \to Y$ est un monomorphisme de faisceaux sur
$(\Sch/S)_{fppf}$, nous voyons que $\Delta : F_n \to F_n \times F_n$
est l'image réciproque de la diagonale $\Delta_Y : Y \to Y \times_S Y$.
La représentabilité de $\Delta_Y$ implique donc celle
de la diagonale de $F_n$. Il suffit par conséquent de démontrer
qu'il existe un schéma $W$ sur $S$ et un morphisme \'etale surjectif
$W \to F_n$.

\medskip\noindent
Pour construire $W \to F_n$, choisissons un recouvrement \'etale $\{Y_i \to Y\}$
où $Y_i$ est un schéma. Soit $X_i = X \times_Y Y_i$ et soit
$\mathcal{F}_i$ l'image réciproque de $\mathcal{F}$ sur $X_i$.
Alors $\mathcal{F}_i$ est pur relativement à $Y_i$, soit par définition,
soit d'après le Lemme \ref{spaces-flat-lemma-quasi-finite-base-change}.
Les autres hypothèses du théorème sont également préservées.
Enfin, la restriction de $F_n$ à $Y_i$ est
le foncteur $F_n$ correspondant à $X_i \to Y_i$ et $\mathcal{F}_i$.
Il suffit donc de montrer ceci : étant donnés $\mathcal{F}$ et $f : X \to Y$
comme dans l'énoncé du théorème, avec $Y$ schéma, le
foncteur $F_n$ est représentable par un schéma $Z_n$ et
$Z_n \to Y$ est un monomorphisme de présentation finie.

\medskip\noindent
Observons qu'un monomorphisme de présentation finie est
séparé et quasi-fini (Morphismes, Lemme
\ref{morphisms-lemma-monomorphism-loc-finite-type-loc-quasi-finite}).
Ainsi, en combinant
Descente, Lemme \ref{descent-lemma-descent-data-sheaves},
Compléments sur les morphismes, Lemme
\ref{more-morphisms-lemma-separated-locally-quasi-finite-morphisms-fppf-descend}
et
Descente, Lemmes \ref{descent-lemma-descending-property-monomorphism} et
\ref{descent-lemma-descending-property-finite-presentation}
nous voyons que la question est locale pour la topologie \'etale sur $Y$.

\medskip\noindent
En particulier, la situation est locale pour la topologie de Zariski sur $Y$,
et nous pouvons supposer $Y$ affine. Dans ce cas, la dimension des
fibres de $f$ est majorée ; $F_n$ est donc représentable
pour $n$ assez grand. Nous pouvons ainsi raisonner par récurrence descendante sur $n$.
Supposons que nous sachions $F_{n + 1}$ représentable par un monomorphisme
$Z_{n + 1} \to Y$ de présentation finie. Considérons le changement de base
$X_{n + 1} = Z_{n + 1} \times_Y X$ et l'image réciproque $\mathcal{F}_{n + 1}$
de $\mathcal{F}$ sur $X_{n + 1}$. Le morphisme $Z_{n + 1} \to Y$ est
quasi-fini, car c'est un monomorphisme de présentation finie ; le
Lemme \ref{spaces-flat-lemma-quasi-finite-base-change}
implique donc que $\mathcal{F}_{n + 1}$ est pur relativement à $Z_{n + 1}$.
Puisque $F_n$ est un sous-foncteur de $F_{n + 1}$, nous concluons que, pour
démontrer le résultat pour $F_n$, il suffit de le démontrer pour le
foncteur correspondant à la situation
$\mathcal{F}_{n + 1}/X_{n + 1}/Z_{n + 1}$.
Nous sommes ainsi ramenés à démontrer le résultat pour $F_n$ dans le cas
$Y_{n + 1} = Y$, c'est-à-dire que nous pouvons supposer $\mathcal{F}$ plat
en dimensions $\geq n + 1$ sur $Y$.

\medskip\noindent
Fixons $n$ et supposons que $\mathcal{F}$ soit plat en dimensions $\geq n + 1$
sur le schéma affine $Y$.
Pour achever la démonstration, nous devons montrer que $F_n$ est représentable
par un monomorphisme $Z_n \to S$ de présentation finie.
Puisque la question est locale pour la topologie \'etale sur $Y$, il suffit de
montrer que, pour tout $y \in Y$, il existe un voisinage \'etale
$(Y', y') \to (Y, y)$ tel que le résultat soit vrai après changement de base à $Y'$.
Ainsi, d'après le
Lemme \ref{spaces-flat-lemma-existence-complete},
nous pouvons supposer qu'il existe des morphismes \'etales $h_j : W_j \to X$,
$j = 1, \ldots, m$, tels que, pour chaque $j$, il existe un dévissage
complet de $\mathcal{F}_j/W_j/Y$ au-dessus de $y$,
où $\mathcal{F}_j$ est l'image réciproque de $\mathcal{F}$ sur $W_j$,
et tels que $|X_y| \subset \bigcup h_j(W_j)$. Puisque
$h_j$ est \'etale, le
Lemme \ref{spaces-flat-lemma-pre-flat-dimension-n}
montre que les faisceaux $\mathcal{F}_j$ sont encore plats
en dimensions $\geq n + 1$ sur $Y$.
Posons $W = \bigcup h_j(W_j)$, qui est un ouvert quasi-compact de $X$.
Comme $\mathcal{F}$ est pur le long de $X_y$, nous voyons que
$$
E = \{t \in |Y| : \text{Ass}_{X_t}(\mathcal{F}_t) \subset W \}.
$$
contient toutes les générisations de $y$. D'après
Diviseurs sur les espaces,
Lemme \ref{spaces-divisors-lemma-relative-assassin-constructible},
$E$ est une partie constructible de $Y$. Nous avons vu que
$\Spec(\mathcal{O}_{Y, y}) \subset E$. D'après
Morphismes, Lemme \ref{morphisms-lemma-constructible-containing-open},
nous voyons que $E$ contient un voisinage ouvert de $y$.
Ainsi, après avoir restreint $Y$, nous pouvons supposer que $E = Y$.
Il résulte du
Lemme \ref{spaces-flat-lemma-localize-flat-dimension-n}
qu'il suffit de démontrer le théorème pour le foncteur $F_n$ associé à
$X = \coprod W_j$ et $\mathcal{F} = \coprod \mathcal{F}_j$.
Si $F_{j, n}$ désigne le foncteur associé à $W_j \to Y$ et au faisceau
$\mathcal{F}_j$, alors $F_n = \prod F_{j, n}$. Il suffit donc
de démontrer que chaque $F_{j, n}$ est représentable par un monomorphisme
$Z_{j, n} \to Y$ de présentation finie, puisque alors
$$
Z_n = Z_{1, n} \times_Y \ldots \times_Y Z_{m, n}
$$
Nous avons ainsi ramené le théorème au cas particulier traité dans
Compléments sur la platitude, Lemme \ref{flat-lemma-flat-dimension-n-representable}.
\end{proof}

\noindent
Nous obtenons enfin le résultat souhaité.

\begin{lemma}
\label{spaces-flat-lemma-when-universal-flattening}
Soit $S$ un schéma.
Soit $f : X \to Y$ un morphisme d'espaces algébriques sur $S$.
Soit $\mathcal{F}$ un $\mathcal{O}_X$-module quasi-cohérent.
\begin{enumerate}
\item Si $f$ est de présentation finie, si $\mathcal{F}$ est un
$\mathcal{O}_X$-module de présentation finie et si $\mathcal{F}$ est
pur relativement à $Y$, alors il existe un aplatissement universel
$Y' \to Y$ de $\mathcal{F}$. De plus, $Y' \to Y$ est un monomorphisme
de présentation finie.
\item Si $f$ est de présentation finie et si $X$ est pur relativement à $Y$,
alors il existe un aplatissement universel $Y' \to Y$ de $X$.
De plus, $Y' \to Y$ est un monomorphisme de présentation finie.
\item Si $f$ est propre et de présentation finie et si $\mathcal{F}$ est un
$\mathcal{O}_X$-module de présentation finie, alors il existe un
aplatissement universel $Y' \to Y$ de $\mathcal{F}$. De plus, $Y' \to Y$ est
un monomorphisme de présentation finie.
\item Si $f$ est propre et de présentation finie,
alors il existe un aplatissement universel $Y' \to Y$ de $X$.
\end{enumerate}
\end{lemma}

\begin{proof}
Ces assertions résultent immédiatement du
Théorème \ref{spaces-flat-theorem-flat-dimension-n-representable}
appliqué à $F_0 = F_{flat}$,
et du fait que, si $f$ est propre, alors $\mathcal{F}$ est automatiquement
pur sur la base, voir
Lemme \ref{spaces-flat-lemma-proper-pure}.
\end{proof}






\section{Théorème d'existence de Grothendieck}
\label{spaces-flat-section-existence}

\noindent
Cette section est l'analogue de
Compléments sur la platitude, section \ref{flat-section-existence},
et poursuit l'étude menée dans
Compléments sur les morphismes d'espaces, section
\ref{spaces-more-morphisms-section-existence-proper}.
Nous travaillerons dans la situation suivante.

\begin{situation}
\label{spaces-flat-situation-existence}
Nous avons ici un système projectif d'anneaux $(A_n)$ dont les applications
de transition sont surjectives et ont des noyaux localement nilpotents. Posons $A = \lim A_n$. Soit
$X$ un espace algébrique séparé et de présentation finie sur $A$.
Posons $X_n = X \times_{\Spec(A)} \Spec(A_n)$ et considérons-le
comme un sous-espace fermé de $X$. Supposons en outre donné un système
$(\mathcal{F}_n, \varphi_n)$, où $\mathcal{F}_n$ est un
$\mathcal{O}_{X_n}$-module de présentation finie, plat sur $A_n$, de support propre sur $A_n$,
et où
$$
\varphi_n :
\mathcal{F}_n \otimes_{\mathcal{O}_{X_n}} \mathcal{O}_{X_{n - 1}}
\longrightarrow
\mathcal{F}_{n - 1}
$$
est un isomorphisme (notation utilisant l'équivalence de
Morphismes d'espaces, Lemme \ref{spaces-morphisms-lemma-i-star-equivalence}).
\end{situation}

\noindent
Notre but est de déterminer si l'on peut trouver un faisceau quasi-cohérent $\mathcal{F}$
sur $X$ tel que
$\mathcal{F}_n = \mathcal{F} \otimes_{\mathcal{O}_X} \mathcal{O}_{X_n}$
pour tout $n$.

\begin{lemma}
\label{spaces-flat-lemma-compute-what-it-should-be}
Dans la Situation \ref{spaces-flat-situation-existence}, considérons
$$
K = R\lim_{D_\QCoh(\mathcal{O}_X)}(\mathcal{F}_n) =
DQ_X(R\lim_{D(\mathcal{O}_X)}\mathcal{F}_n)
$$
Alors $K$ appartient à $D^b_{\QCoh}(\mathcal{O}_X)$ et, en fait,
$K$ n'a de faisceaux de cohomologie non nuls qu'en degrés $\geq 0$.
\end{lemma}

\begin{proof}
C'est un cas particulier de
Catégories dérivées des espaces, Exemple
\ref{spaces-perfect-example-inverse-limit-quasi-coherent}.
\end{proof}

\begin{lemma}
\label{spaces-flat-lemma-compute-against-perfect}
Dans la Situation \ref{spaces-flat-situation-existence}, soit $K$ comme dans le
Lemme \ref{spaces-flat-lemma-compute-what-it-should-be}. Pour tout objet parfait
$E$ de $D(\mathcal{O}_X)$, on a :
\begin{enumerate}
\item $M = R\Gamma(X, K \otimes^\mathbf{L} E)$ est un objet parfait de $D(A)$,
et il existe un isomorphisme canonique
$R\Gamma(X_n, \mathcal{F}_n \otimes^\mathbf{L} E|_{X_n}) =
M \otimes_A^\mathbf{L} A_n$
dans $D(A_n)$,
\item $N = R\Hom_X(E, K)$ est un objet parfait de $D(A)$,
et il existe un isomorphisme canonique
$R\Hom_{X_n}(E|_{X_n}, \mathcal{F}_n) = N \otimes_A^\mathbf{L} A_n$
dans $D(A_n)$.
\end{enumerate}
Dans les deux assertions, $E|_{X_n}$ désigne l'image réciproque dérivée
de $E$ sur $X_n$.
\end{lemma}

\begin{proof}
Preuve de (2). Écrivons $E_n = E|_{X_n}$ et
$N_n = R\Hom_{X_n}(E_n, \mathcal{F}_n)$.
Rappelons que $R\Hom_{X_n}(-, -)$ est égal à
$R\Gamma(X_n, R\SheafHom(-, -))$, voir
Cohomologie sur les sites, section \ref{sites-cohomology-section-global-RHom}.
Ainsi, d'après Catégories dérivées des espaces, Lemme
\ref{spaces-perfect-lemma-base-change-RHom-perfect}
le complexe $N_n$ est un objet parfait de $D(A_n)$
dont la formation commute au changement de base. Ainsi, les applications
$N_n \otimes_{A_n}^\mathbf{L} A_{n - 1} \to N_{n - 1}$
provenant de $\varphi_n$ sont des isomorphismes.
D'après Compléments d'algèbre, Lemme \ref{more-algebra-lemma-Rlim-perfect-gives-perfect},
$R\lim N_n$ est parfait et
son changement de base à $A_n$ redonne $N_n$.
D'autre part, le foncteur exact
$R\Hom_X(E, -) : D_\QCoh(\mathcal{O}_X) \to D(A)$
de catégories triangulées commute aux produits,
donc aux limites dérivées ; par conséquent,
$$
R\Hom_X(E, K) =
R\lim R\Hom_X(E, \mathcal{F}_n) =
R\lim R\Hom_X(E_n, \mathcal{F}_n) =
R\lim N_n
$$
Cela démontre (2). Pour obtenir (1), on le reformule sous la forme (2)
à l'aide de Cohomologie sur les sites, Lemme
\ref{sites-cohomology-lemma-dual-perfect-complex}.
\end{proof}

\begin{lemma}
\label{spaces-flat-lemma-relative-pseudo-coherence}
Dans la Situation \ref{spaces-flat-situation-existence}, soit $K$ comme dans le
Lemme \ref{spaces-flat-lemma-compute-what-it-should-be}. Alors $K$
est pseudo-cohérent relativement à $A$.
\end{lemma}

\begin{proof}
En combinant le Lemme \ref{spaces-flat-lemma-compute-against-perfect} et
Catégories dérivées des espaces, Lemme \ref{spaces-perfect-lemma-perfect-enough},
nous voyons que $R\Gamma(X, K \otimes^\mathbf{L} E)$
est pseudo-cohérent dans $D(A)$ pour tout objet pseudo-cohérent
$E$ de $D(\mathcal{O}_X)$. Le lemme résulte donc de
Compléments sur les morphismes d'espaces, Lemme
\ref{spaces-more-morphisms-lemma-characterize-pseudo-coherent}.
\end{proof}

\begin{lemma}
\label{spaces-flat-lemma-compute-over-affine}
Dans la Situation \ref{spaces-flat-situation-existence}, soit $K$ comme dans le
Lemme \ref{spaces-flat-lemma-compute-what-it-should-be}. Pour tout
morphisme \'etale $U \to X$, avec $U$ quasi-compact et quasi-séparé, on a
$$
R\Gamma(U, K) \otimes_A^\mathbf{L} A_n =
R\Gamma(U_n, \mathcal{F}_n)
$$
dans $D(A_n)$, où $U_n = U \times_X X_n$.
\end{lemma}

\begin{proof}
Fixons $n$. D'après Catégories dérivées des espaces, Lemme
\ref{spaces-perfect-lemma-computing-sections-as-colim}
il existe un système de complexes parfaits $E_m$
sur $X$ tel que
$R\Gamma(U, K) = \text{hocolim} R\Gamma(X, K \otimes^\mathbf{L} E_m)$.
En fait, cette formule vaut non seulement pour $K$, mais pour tout objet de
$D_\QCoh(\mathcal{O}_X)$.
En l'appliquant à $\mathcal{F}_n$,
nous obtenons
\begin{align*}
R\Gamma(U_n, \mathcal{F}_n)
& =
R\Gamma(U, \mathcal{F}_n) \\
& =
\text{hocolim}_m R\Gamma(X, \mathcal{F}_n \otimes^\mathbf{L} E_m) \\
& =
\text{hocolim}_m R\Gamma(X_n, \mathcal{F}_n \otimes^\mathbf{L} E_m|_{X_n})
\end{align*}
En utilisant le Lemme \ref{spaces-flat-lemma-compute-against-perfect}
et le fait que $- \otimes_A^\mathbf{L} A_n$
commute aux colimites homotopiques, nous obtenons le résultat.
\end{proof}

\begin{lemma}
\label{spaces-flat-lemma-finitely-presented}
Dans la Situation \ref{spaces-flat-situation-existence}, soit $K$ comme dans le
Lemme \ref{spaces-flat-lemma-compute-what-it-should-be}.
Notons $X_0 \subset |X|$ la partie fermée
formée des points situés au-dessus de la partie fermée
$\Spec(A_1) = \Spec(A_2) = \ldots$ de $\Spec(A)$.
Il existe un sous-espace ouvert $W \subset X$ contenant $X_0$
tel que
\begin{enumerate}
\item $H^i(K)|_W$ soit nul sauf si $i = 0$,
\item $\mathcal{F} = H^0(K)|_W$ soit de présentation finie, et
\item $\mathcal{F}_n = \mathcal{F} \otimes_{\mathcal{O}_X} \mathcal{O}_{X_n}$.
\end{enumerate}
\end{lemma}

\begin{proof}
Fixons $n \geq 1$. Par construction, il existe une application canonique
$K \to \mathcal{F}_n$ dans $D_\QCoh(\mathcal{O}_X)$,
et donc une application canonique $H^0(K) \to \mathcal{F}_n$
de faisceaux quasi-cohérents. Cela précise le sens de la partie (3).

\medskip\noindent
Soit $x \in X_0$ un point. Nous allons trouver un voisinage ouvert $W$
de $x$ où (1), (2) et (3) sont vraies. Puisque $X_0$ est quasi-compact,
cela démontrera le lemme. Soit $U \to X$ un morphisme \'etale,
avec $U$ affine, et soit $u \in U$ un point s'envoyant sur $x$. Puisque $|U| \to |X|$
est ouvert, il suffit de trouver un voisinage ouvert de $u$ dans $U$
où (1), (2) et (3) sont vraies. Écrivons $U = \Spec(B)$.
Choisissons une surjection $P \to B$, avec $P$ lisse sur $A$.
D'après le Lemme \ref{spaces-flat-lemma-relative-pseudo-coherence}
et la définition de la pseudo-cohérence relative,
il existe un complexe borné supérieurement $F^\bullet$
de $P$-modules libres de rang fini représentant
$Ri_*K$, où $i : U \to \Spec(P)$ est l'immersion
fermée induite par la présentation.
Soit $M_n$ le $B$-module correspondant à $\mathcal{F}_n|_U$.
D'après le Lemme \ref{spaces-flat-lemma-compute-over-affine},
$$
H^i(F^\bullet \otimes_A A_n) =
\left\{
\begin{matrix}
0 & \text{si} & i \not = 0 \\
M_n & \text{si} & i = 0
\end{matrix}
\right.
$$
Soit $i$ l'indice maximal tel que $F^i$ soit non nul.
Si $i \leq 0$, alors (1), (2) et (3) sont vraies.
Sinon, $i > 0$, et nous voyons que le rang de l'application
$$
F^{i - 1} \to F^i
$$
au point $u$ est maximal. Il est donc maximal dans un voisinage ouvert
de $u$ dans $\Spec(P)$. Ainsi, après avoir remplacé
$P$ par une localisation principale, nous pouvons supposer que l'application affichée
est surjective. Puisque $F^i$ est libre de rang fini, nous pouvons choisir
une décomposition $F^{i - 1} = F' \oplus F^i$. Nous pouvons alors
remplacer $F^\bullet$ par le complexe
$$
\ldots \to F^{i - 2} \to F' \to 0 \to \ldots
$$
et conclure par récurrence sur $i$.
\end{proof}

\begin{lemma}
\label{spaces-flat-lemma-proper-support}
Dans la Situation \ref{spaces-flat-situation-existence}, soit $K$ comme dans le
Lemme \ref{spaces-flat-lemma-compute-what-it-should-be}. Soit $W \subset X$
comme dans le Lemme \ref{spaces-flat-lemma-finitely-presented}.
Posons $\mathcal{F} = H^0(K)|_W$. Alors, quitte à restreindre l'ouvert $W$,
le support de $\mathcal{F}$ est propre sur $A$.
\end{lemma}

\begin{proof}
Fixons $n \geq 1$. Soit $I_n = \Ker(A \to A_n)$.
D'après Compléments d'algèbre, Lemme \ref{more-algebra-lemma-limit-henselian},
la paire $(A, I_n)$ est hensélienne.
Soit $Z \subset W$ le support schématique de $\mathcal{F}$.
C'est un sous-espace fermé, puisque $\mathcal{F}$ est de présentation finie.
D'après la partie (3) du Lemme \ref{spaces-flat-lemma-finitely-presented},
$Z \times_{\Spec(A)} \Spec(A_n)$
est égal au support de $\mathcal{F}_n$, et est donc
propre sur $\Spec(A/I)$.
D'après Compléments sur les morphismes d'espaces, Lemme
\ref{spaces-more-morphisms-lemma-split-off-proper-part-henselian}
nous pouvons écrire $Z = Z_1 \amalg Z_2$, avec $Z_1, Z_2$ ouverts et
fermés dans $Z$, avec $Z_1$ propre
sur $A$, et avec $Z_1 \times_{\Spec(A)} \Spec(A/I_n)$
égal au support de $\mathcal{F}_n$.
Autrement dit, $|Z_2|$ ne rencontre pas $X_0$.
Ainsi, en remplaçant $W$ par $W \setminus Z_2$, nous obtenons le lemme.
\end{proof}

\begin{theorem}[Théorème d'existence de Grothendieck]
\label{spaces-flat-theorem-existence}
Dans la Situation \ref{spaces-flat-situation-existence},
il existe un $\mathcal{O}_X$-module de présentation finie
$\mathcal{F}$, plat sur $A$, de support propre sur $A$,
tel que
$\mathcal{F}_n = \mathcal{F} \otimes_{\mathcal{O}_X} \mathcal{O}_{X_n}$
pour tout $n$, de manière compatible avec les applications $\varphi_n$.
\end{theorem}

\begin{proof}
Appliquons les Lemmes \ref{spaces-flat-lemma-compute-what-it-should-be},
\ref{spaces-flat-lemma-compute-against-perfect},
\ref{spaces-flat-lemma-relative-pseudo-coherence},
\ref{spaces-flat-lemma-compute-over-affine},
\ref{spaces-flat-lemma-finitely-presented} et
\ref{spaces-flat-lemma-proper-support}
pour obtenir un sous-espace ouvert $W \subset X$ contenant tous les points
au-dessus de $\Spec(A_n)$
et un $\mathcal{O}_W$-module $\mathcal{F}$ de présentation finie,
dont le support est propre sur $A$ et tel que
$\mathcal{F}_n = \mathcal{F} \otimes_{\mathcal{O}_W} \mathcal{O}_{X_n}$
pour tout $n \geq 1$. (Cela a un sens puisque $X_n \subset W$.)
D'après le Lemme \ref{spaces-flat-lemma-proper-pure}, $\mathcal{F}$
est universellement pur relativement à $\Spec(A)$.
D'après le Théorème \ref{spaces-flat-theorem-flat-dimension-n-representable}
(pour des explications, voir le Lemme \ref{spaces-flat-lemma-when-universal-flattening}),
il existe un aplatissement universel $S' \to \Spec(A)$
de $\mathcal{F}$ et, de plus, le morphisme $S' \to \Spec(A)$
est un monomorphisme de présentation finie.
En particulier, $S'$ est un schéma (cela résulte de la démonstration
du théorème, mais aussi a posteriori de
Morphismes d'espaces, Proposition
\ref{spaces-morphisms-proposition-locally-quasi-finite-separated-over-scheme}).
Puisque le changement de base de $\mathcal{F}$ à $\Spec(A_n)$
est $\mathcal{F}_n$, nous voyons que $\Spec(A_n) \to \Spec(A)$
se factorise (uniquement) par $S'$ pour tout $n$.
D'après Compléments sur la platitude, Lemme \ref{flat-lemma-monomorphism-isomorphism},
nous avons $S' = \Spec(A)$.
Cela signifie que $\mathcal{F}$ est plat sur $A$.
Enfin, puisque le support schématique $Z$ de $\mathcal{F}$
est propre sur $\Spec(A)$, le morphisme $Z \to X$ est fermé.
Par conséquent, l'image directe $(W \to X)_*\mathcal{F}$ est supportée
sur $W$ et possède toutes les propriétés voulues.
\end{proof}







\section{Théorème d'existence de Grothendieck, bis}
\label{spaces-flat-section-existence-derived}

\noindent
Dans cette section, nous démontrons un analogue du théorème d'existence de Grothendieck
dans la catégorie dérivée, en suivant la méthode utilisée à la
section \ref{spaces-flat-section-existence} pour les modules quasi-cohérents.
Cette section est l'analogue, pour les espaces algébriques, de
Compléments sur la platitude, section \ref{flat-section-existence-derived}.
Le cas classique (pour les espaces algébriques)
est étudié dans Compléments sur les morphismes d'espaces, section
\ref{spaces-more-morphisms-section-existence-proper}.
Nous travaillerons dans la situation suivante.

\begin{situation}
\label{spaces-flat-situation-existence-derived}
Nous avons ici un système projectif d'anneaux $(A_n)$ dont les applications
de transition sont surjectives et ont des noyaux localement nilpotents. Posons $A = \lim A_n$. Soit
$X$ un espace algébrique propre, plat et de présentation finie sur $A$.
Posons $X_n = X \times_{\Spec(A)} \Spec(A_n)$ et considérons-le
comme un sous-espace fermé de $X$. Supposons en outre donné un système
$(K_n, \varphi_n)$, où $K_n$ est un objet pseudo-cohérent de
$D(\mathcal{O}_{X_n})$ et
$$
\varphi_n : K_n \longrightarrow K_{n - 1}
$$
est une application dans $D(\mathcal{O}_{X_n})$ qui induit un isomorphisme
$K_n \otimes_{\mathcal{O}_{X_n}}^\mathbf{L} \mathcal{O}_{X_{n - 1}}
\to K_{n - 1}$ dans $D(\mathcal{O}_{X_{n - 1}})$.
\end{situation}

\noindent
Plus précisément, il faudrait écrire
$\varphi_n : K_n \to Ri_{n - 1, *}K_{n - 1}$
où $i_{n - 1} : X_{n - 1} \to X_n$ est le morphisme d'inclusion ;
avec cette notation, la condition est que l'application adjointe
$Li_{n - 1}^*K_n \to K_{n - 1}$ soit un isomorphisme.
Notre but est de trouver un $K \in D(\mathcal{O}_X)$ pseudo-cohérent
tel que $K_n = K \otimes_{\mathcal{O}_X}^\mathbf{L} \mathcal{O}_{X_n}$
pour tout $n$ (avec le même abus de notation).

\begin{lemma}
\label{spaces-flat-lemma-compute-what-it-should-be-derived}
Dans la Situation \ref{spaces-flat-situation-existence-derived}, considérons
$$
K = R\lim_{D_\QCoh(\mathcal{O}_X)}(K_n) =
DQ_X(R\lim_{D(\mathcal{O}_X)} K_n)
$$
Alors $K$ appartient à $D^-_{\QCoh}(\mathcal{O}_X)$.
\end{lemma}

\begin{proof}
Le foncteur $DQ_X$ existe car $X$ est quasi-compact et
quasi-séparé, voir Catégories dérivées des espaces, Lemme
\ref{spaces-perfect-lemma-better-coherator}.
Puisque $DQ_X$ est un adjoint à droite, il commute aux produits,
et donc aux limites dérivées. On obtient ainsi l'égalité
de l'énoncé du lemme.

\medskip\noindent
D'après Catégories dérivées des espaces,
Lemme \ref{spaces-perfect-lemma-boundedness-better-coherator},
le foncteur $DQ_X$ est de dimension cohomologique bornée.
Il suffit donc de montrer que $R\lim K_n \in D^-(\mathcal{O}_X)$.
Pour cela, soit $U \to X$ un morphisme \'etale, avec $U$ affine.
Il existe alors une suite exacte canonique
$$
0 \to
R^1\lim H^{m - 1}(U, K_n) \to H^m(U, R\lim K_n) \to
\lim H^m(U, K_n) \to 0
$$
d'après Cohomologie sur les sites, Lemme
\ref{sites-cohomology-lemma-RGamma-commutes-with-Rlim}.
Puisque $U$ est affine et $K_n$ pseudo-cohérent (et possède donc des
faisceaux de cohomologie quasi-cohérents d'après
Catégories dérivées des espaces, Lemme \ref{spaces-perfect-lemma-pseudo-coherent}),
nous avons $H^m(U, K_n) = H^m(K_n)(U)$ d'après
Catégories dérivées des schémas, Lemme \ref{perfect-lemma-affine-compare-bounded}.
Il suffit donc de montrer que les $K_n$ sont
bornés supérieurement indépendamment de $n$.

\medskip\noindent
Puisque $K_n$ est pseudo-cohérent, on a $K_n \in D^-(\mathcal{O}_{X_n})$.
Supposons que $a_n$ soit maximal tel que $H^{a_n}(K_n)$ soit non nul.
Bien entendu, $a_1 \leq a_2 \leq a_3 \leq \ldots$.
Remarquons que $H^{a_n}(K_n)$ est un
$\mathcal{O}_{X_n}$-module de présentation finie
(Cohomologie sur les sites, Lemme
\ref{sites-cohomology-lemma-finite-cohomology}).
Nous avons $H^{a_n}(K_{n - 1}) =
H^{a_n}(K_n) \otimes_{\mathcal{O}_{X_n}} \mathcal{O}_{X_{n - 1}}$.
Puisque $X_{n - 1} \to X_n$ est un épaississement, le
lemme de Nakayama (Algèbre, Lemme \ref{algebra-lemma-NAK}) montre que, si
$H^{a_n}(K_n) \otimes_{\mathcal{O}_{X_n}} \mathcal{O}_{X_{n - 1}}$
est nul, alors $H^{a_n}(K_n)$ l'est aussi (on peut par exemple vérifier
sur les germes ; ce petit détail est omis).
Ainsi $a_{n - 1} = a_n$ pour tout $n$, ce qui permet de conclure.
\end{proof}

\begin{lemma}
\label{spaces-flat-lemma-compute-against-perfect-derived}
Dans la Situation \ref{spaces-flat-situation-existence-derived}, soit $K$ comme dans le
Lemme \ref{spaces-flat-lemma-compute-what-it-should-be-derived}. Pour tout objet parfait
$E$ de $D(\mathcal{O}_X)$, la cohomologie
$$
M = R\Gamma(X, K \otimes^\mathbf{L} E)
$$
est un objet pseudo-cohérent de $D(A)$ et il existe un isomorphisme canonique
$$
R\Gamma(X_n, K_n \otimes^\mathbf{L} E|_{X_n}) = M \otimes_A^\mathbf{L} A_n
$$
dans $D(A_n)$. Ici, $E|_{X_n}$ désigne l'image réciproque dérivée de $E$ sur $X_n$.
\end{lemma}

\begin{proof}
Écrivons $E_n = E|_{X_n}$ et
$M_n = R\Gamma(X_n, K_n \otimes^\mathbf{L} E|_{X_n})$.
D'après Catégories dérivées des espaces, Lemme
\ref{spaces-perfect-lemma-flat-proper-pseudo-coherent-direct-image-general}
nous voyons que $M_n$ est un objet pseudo-cohérent de $D(A_n)$
dont la formation commute au changement de base. Ainsi, les applications
$M_n \otimes_{A_n}^\mathbf{L} A_{n - 1} \to M_{n - 1}$
provenant de $\varphi_n$ sont des isomorphismes. D'après
Compléments d'algèbre, Lemme
\ref{more-algebra-lemma-Rlim-pseudo-coherent-gives-pseudo-coherent}
$R\lim M_n$ est pseudo-cohérent et
son changement de base à $A_n$ redonne $M_n$.
D'autre part, le foncteur exact
$R\Gamma(X, -) : D_\QCoh(\mathcal{O}_X) \to D(A)$
de catégories triangulées commute aux produits,
donc aux limites dérivées ; par conséquent,
$$
R\Gamma(X, E \otimes^\mathbf{L} K) =
R\lim R\Gamma(X,  E \otimes^\mathbf{L} K_n) =
R\lim R\Gamma(X_n, E_n \otimes^\mathbf{L} K_n) =
R\lim M_n
$$
comme souhaité.
\end{proof}

\begin{lemma}
\label{spaces-flat-lemma-relative-pseudo-coherence-derived}
Dans la Situation \ref{spaces-flat-situation-existence-derived}, soit $K$ comme dans le
Lemme \ref{spaces-flat-lemma-compute-what-it-should-be-derived}. Alors $K$
est pseudo-cohérent sur $X$.
\end{lemma}

\begin{proof}
En combinant le Lemme \ref{spaces-flat-lemma-compute-against-perfect-derived} et
Catégories dérivées des espaces, Lemme
\ref{spaces-perfect-lemma-perfect-enough}
nous voyons que $R\Gamma(X, K \otimes^\mathbf{L} E)$
est pseudo-cohérent dans $D(A)$ pour tout objet pseudo-cohérent
$E$ de $D(\mathcal{O}_X)$. Il résulte donc de
Compléments sur les morphismes d'espaces, Lemme
\ref{spaces-more-morphisms-lemma-characterize-pseudo-coherent}
que $K$ est pseudo-cohérent relativement à $A$.
Puisque $X$ est plat et de présentation finie
sur $A$, cela équivaut à être pseudo-cohérent sur $X$, voir
Compléments sur les morphismes d'espaces, Lemme
\ref{spaces-more-morphisms-lemma-relative-pseudo-coherent-is-moot}.
\end{proof}

\begin{lemma}
\label{spaces-flat-lemma-compute-over-affine-derived}
Dans la Situation \ref{spaces-flat-situation-existence-derived}, soit $K$ comme dans le
Lemme \ref{spaces-flat-lemma-compute-what-it-should-be-derived}. Pour tout
morphisme \'etale $U \to X$, avec $U$ quasi-compact et quasi-séparé, on a
$$
R\Gamma(U, K) \otimes_A^\mathbf{L} A_n =
R\Gamma(U_n, K_n)
$$
dans $D(A_n)$, où $U_n = U \times_X X_n$.
\end{lemma}

\begin{proof}
Fixons $n$. D'après Catégories dérivées des espaces, Lemme
\ref{spaces-perfect-lemma-computing-sections-as-colim}
il existe un système de complexes parfaits $E_m$
sur $X$ tel que
$R\Gamma(U, K) = \text{hocolim} R\Gamma(X, K \otimes^\mathbf{L} E_m)$.
En fait, cette formule vaut non seulement pour $K$, mais pour tout objet de
$D_\QCoh(\mathcal{O}_X)$.
En l'appliquant à $K_n$,
nous obtenons
\begin{align*}
R\Gamma(U_n, K_n)
& =
R\Gamma(U, K_n) \\
& =
\text{hocolim}_m R\Gamma(X, K_n \otimes^\mathbf{L} E_m) \\
& =
\text{hocolim}_m R\Gamma(X_n, K_n \otimes^\mathbf{L} E_m|_{X_n})
\end{align*}
En utilisant le Lemme \ref{spaces-flat-lemma-compute-against-perfect-derived}
et le fait que $- \otimes_A^\mathbf{L} A_n$
commute aux colimites homotopiques, nous obtenons le résultat.
\end{proof}

\begin{theorem}[Théorème d'existence de Grothendieck dérivé]
\label{spaces-flat-theorem-existence-derived}
Dans la Situation \ref{spaces-flat-situation-existence-derived},
il existe un $K$ pseudo-cohérent dans $D(\mathcal{O}_X)$
tel que $K_n = K \otimes_{\mathcal{O}_X}^\mathbf{L} \mathcal{O}_{X_n}$
pour tout $n$, de manière compatible avec les applications $\varphi_n$.
\end{theorem}

\begin{proof}
Appliquons les Lemmes \ref{spaces-flat-lemma-compute-what-it-should-be-derived},
\ref{spaces-flat-lemma-compute-against-perfect-derived},
\ref{spaces-flat-lemma-relative-pseudo-coherence-derived}
pour obtenir un objet pseudo-cohérent $K$ de $D(\mathcal{O}_X)$.
En choisissant $U$ affine dans le Lemme
\ref{spaces-flat-lemma-compute-over-affine-derived}
on voit immédiatement que $K$ se restreint en $K_n$ sur $X_n$.
\end{proof}

\begin{remark}
\label{spaces-flat-remark-correct-generality}
Le résultat de cette section peut être généralisé. Il est probablement vrai
si l'on suppose seulement que $X \to \Spec(A)$ est séparé et de présentation finie,
et que $K_n$ est pseudo-cohérent relativement à $A_n$ et supporté par une partie
fermée de $X_n$ propre sur $A_n$. On obtiendra alors un $K$
pseudo-cohérent relativement à $A$ et supporté par une partie fermée
propre sur $A$. Si cela nous est un jour nécessaire, nous
formulerons ici un énoncé précis et le démontrerons.
\end{remark}















% OK, start here.
%

\chapter{Groupoïdes dans la catégorie des espaces algébriques}



\phantomsection
\label{spaces-groupoids-section-phantom}


\section{Introduction}
\label{spaces-groupoids-section-introduction}

\noindent
Ce chapitre est consacré à des généralités sur les groupoïdes dans la catégorie des espaces algébriques. Nous recommandons la lecture du bel article \cite{K-M} de Keel et Mori.

\medskip\noindent
Une grande partie de ce qui suit reprend le chapitre consacré aux groupoïdes dans la catégorie des schémas ; voir Groupoïdes, section \ref{groupoids-section-introduction}. En revanche, la discussion des champs quotients est nouvelle.


\section{Conventions}
\label{spaces-groupoids-section-conventions}

\noindent
Nous supposons une fois pour toutes que tous les schémas sont contenus dans un grand site fppf $\Sch_{fppf}$ et que tout anneau $A$ considéré est tel que $\Spec(A)$ soit (isomorphe à) un objet de ce grand site.

\medskip\noindent
Soit $S$ un schéma et soit $X$ un espace algébrique sur $S$. Dans ce chapitre et les suivants, nous écrirons $X \times_S X$ pour le produit de $X$ avec lui-même (dans la catégorie d'espaces algébriques sur $S$), au lieu de $X \times X$.

\medskip\noindent
Nous maintenons notre convention consistant à numéroter les projections à partir de l'indice $0$ ; nous avons donc $\text{pr}_0 : X \times_S Y \to X$ et $\text{pr}_1 : X \times_S Y \to Y$.




\section{Notations}
\label{spaces-groupoids-section-notation}

\noindent
Soit $S$ un schéma ; ce sera notre schéma de base, et tous les espaces algébriques seront sur $S$. Soit $B$ un espace algébrique sur $S$ ; ce sera notre espace algébrique de base, et d'autres espaces algébriques, ainsi que des schémas, seront souvent considérés sur $B$. Dire que $X$ est un espace algébrique sur $B$ signifie que $X$ est un espace algébrique sur $S$ muni d'un morphisme structural $X \to B$. Nous nous efforçons en outre de réserver la lettre $T$ à un schéma ``test'' sur $B$, c'est-à-dire à un schéma muni d'un morphisme structural $T \to B$. Dans cette situation, nous notons $X(T)$ l'ensemble des points de $X$ à valeurs dans $T$ {\it au-dessus de} $B$. En formule :
$$
X(T) = \Mor_B(T, X).
$$
De même, étant donné un second espace algébrique $Y$ sur $B$, nous posons
$$
X(Y) = \Mor_B(Y, X).
$$
Supposons donnés des espaces algébriques $X$, $Y$ sur $B$ comme ci-dessus, ainsi qu'un morphisme $f : X \to Y$ sur $B$. Pour tout schéma $T$ sur $B$, nous obtenons une application induite entre ensembles
$$
f : X(T) \longrightarrow Y(T)
$$
qui varie fonctoriellement avec le schéma $T$ sur $B$. Puisque $f$ est un morphisme de faisceaux sur $(\Sch/S)_{fppf}$ au-dessus du faisceau $B$, il est clair que $f$ détermine cette règle et est déterminé par elle. Plus généralement, nous employons la même notation pour les applications entre produits fibrés. Par exemple, si $X$, $Y$, $Z$ sont des espaces algébriques sur $B$ et si $m : X \times_B Y \to Z \times_B Z$ est un morphisme d'espaces algébriques sur $B$, nous considérons que $m$ correspond à une famille d'applications entre points à valeurs dans $T$
$$
X(T) \times Y(T) \longrightarrow Z(T) \times Z(T).
$$
Et ainsi de suite.

\medskip\noindent
Enfin, soient deux morphismes $f, g : X \to Y$ d'espaces algébriques sur $B$. Si les applications induites $f, g : X(T) \to Y(T)$ sont égales pour tout schéma $T$ sur $B$, alors $f = g$ ; par suite, les applications $f, g : X(Z) \to Y(Z)$ sont aussi égales pour tout autre espace algébrique $Z$ sur $B$. Ainsi, pour vérifier les axiomes d'un espace algébrique en groupes $G$ sur $B$, il suffit par exemple de vérifier la commutativité des diagrammes sur les points à valeurs dans $T$, pour tout schéma $T$ sur $B$, comme nous le faisons dans la définition \ref{spaces-groupoids-definition-group-space} ci-dessous.




\section{Relations d'équivalence}
\label{spaces-groupoids-section-equivalence-relations}

\noindent
Voir Groupoïdes, section \ref{groupoids-section-equivalence-relations}, pour les notations.

\begin{definition}
\label{spaces-groupoids-definition-equivalence-relation}
Soit $B \to S$ comme dans la section \ref{spaces-groupoids-section-notation}. Soit $U$ un espace algébrique sur $B$.
\begin{enumerate}
\item Une {\it pré-relation} sur $U$ au-dessus de $B$ est tout morphisme $j : R \to U \times_B U$ d'espaces algébriques sur $B$. Dans ce cas, nous posons $t = \text{pr}_0 \circ j$ et $s = \text{pr}_1 \circ j$, de sorte que $j = (t, s)$.
\item Une {\it relation} sur $U$ au-dessus de $B$ est un monomorphisme $j : R \to U \times_B U$ d'espaces algébriques sur $B$.
\item Une {\it pré-relation d'équivalence} est une pré-relation $j : R \to U \times_B U$ telle que l'image de $j : R(T) \to U(T) \times U(T)$ soit une relation d'équivalence pour tout schéma $T$ sur $B$.
\item Nous disons qu'un morphisme $R \to U \times_B U$ d'espaces algébriques sur $B$ est une {\it relation d'équivalence sur $U$ au-dessus de $B$} si et seulement si, pour tout $T$ sur $B$, les points de $R$ à valeurs dans $T$ définissent une relation d'équivalence sur l'ensemble des points de $U$ à valeurs dans $T$.
\end{enumerate}
\end{definition}

\noindent
En d'autres termes, une relation d'équivalence est une pré-relation d'équivalence telle que $j$ soit une relation.

\begin{lemma}
\label{spaces-groupoids-lemma-restrict-relation}
Soit $B \to S$ comme dans la section \ref{spaces-groupoids-section-notation}. Soit $U$ un espace algébrique sur $B$. Soit $j : R \to U \times_B U$ une pré-relation. Soit $g : U' \to U$ un morphisme d'espaces algébriques sur $B$. Enfin, posons
$$
R' = (U' \times_B U')\times_{U \times_B U} R
\xrightarrow{j'}
U' \times_B U'
$$
$j'$ est une pré-relation sur $U'$ au-dessus de $B$. Si $j$ est une relation, alors $j'$ est une relation. Si $j$ est une pré-relation d'équivalence, alors $j'$ est une pré-relation d'équivalence. Si $j$ est une relation d'équivalence, alors $j'$ est une relation d'équivalence.
\end{lemma}

\begin{proof}
Omis.
\end{proof}

\begin{definition}
\label{spaces-groupoids-definition-restrict-relation}
Soit $B \to S$ comme dans la section \ref{spaces-groupoids-section-notation}. Soit $U$ un espace algébrique sur $B$. Soit $j : R \to U \times_B U$ une pré-relation. Soit $g : U' \to U$ un morphisme d'espaces algébriques sur $B$. La pré-relation $j' : R' \to U' \times_B U'$ du lemme \ref{spaces-groupoids-lemma-restrict-relation} est appelée la {\it restriction}, ou l'{\it image réciproque}, de la pré-relation $j$ à $U'$. Dans cette situation, nous écrivons parfois $R' = R|_{U'}$.
\end{definition}

\begin{lemma}
\label{spaces-groupoids-lemma-pre-equivalence-equivalence-relation-points}
Soit $B \to S$ comme dans la section \ref{spaces-groupoids-section-notation}. Soit $j : R \to U \times_B U$ une pré-relation d'espaces algébriques sur $B$. Considérons la relation sur $|U|$ définie par la règle
$$
x \sim y
\Leftrightarrow
\exists\ r \in |R| :
t(r) = x,
s(r) = y.
$$
Si $j$ est une pré-relation d'équivalence, alors c'est une relation d'équivalence.
\end{lemma}

\begin{proof}
Supposons que $x \sim y$ et $y \sim z$. Choisissons $r \in |R|$ tel que $t(r) = x$, $s(r) = y$, et $r' \in |R|$ tel que $t(r') = y$, $s(r') = z$. On peut choisir un corps $K$ tel que $r$ et $r'$ soient représentés par des morphismes $r, r' : \Spec(K) \to R$ vérifiant $s \circ r = t \circ r'$. Posons $x = t \circ r$, $y = s \circ r = t \circ r'$ et $z = s \circ r'$ ; ainsi, $x, y, z : \Spec(K) \to U$. Par construction, $(x, y) \in j(R(K))$ et $(y, z) \in j(R(K))$. Puisque $j$ est une pré-relation d'équivalence, $(x, z) \in j(R(K))$ également. Il en résulte clairement que $x \sim z$.

\medskip\noindent
La preuve que $\sim$ est réflexive et symétrique est omise.
\end{proof}















\section{Espaces algébriques en groupes}
\label{spaces-groupoids-section-group-spaces}

\noindent
Voir Groupoïdes, section \ref{groupoids-section-group-schemes}, pour les notations.

\begin{definition}
\label{spaces-groupoids-definition-group-space}
Soit $B \to S$ comme dans la section \ref{spaces-groupoids-section-notation}.
\begin{enumerate}
\item Un {\it espace algébrique en groupes sur $B$} est une paire $(G, m)$, où $G$ est un espace algébrique sur $B$ et $m : G \times_B G \to G$ est un morphisme d'espaces algébriques sur $B$ ayant la propriété suivante : pour tout schéma $T$ sur $B$, la paire $(G(T), m)$ est un groupe.
\item Un morphisme {\it $\psi : (G, m) \to (G', m')$ d'espaces algébriques en groupes sur $B$} est un morphisme $\psi : G \to G'$ d'espaces algébriques sur $B$ tel que, pour tout $T/B$, l'application induite $\psi : G(T) \to G'(T)$ soit un homomorphisme de groupes.
\end{enumerate}
\end{definition}

\noindent
Soit $(G, m)$ un espace algébrique en groupes sur l'espace algébrique $B$. La discussion de Groupoïdes, section \ref{groupoids-section-group-schemes}, fournit des morphismes d'espaces algébriques sur $B$, l'unité $e : B \to G$ et l'inversion $i : G \to G$, tels que, pour tout $T$, le quadruplet $(G(T), m, e, i)$ satisfasse aux axiomes d'un groupe.

\medskip\noindent
Soient $(G, m)$ et $(G', m')$ des espaces algébriques en groupes sur $B$. Soit $f : G \to G'$ un morphisme d'espaces algébriques sur $B$. Il résulte de la définition que $f$ est un morphisme d'espaces algébriques en groupes sur $B$ si et seulement si le diagramme suivant est commutatif :
$$
\xymatrix{
G \times_B G \ar[r]_-{f \times f} \ar[d]_m &
G' \times_B G' \ar[d]^m \\
G \ar[r]^f & G'
}
$$

\begin{lemma}
\label{spaces-groupoids-lemma-base-change-group-space}
Soit $B \to S$ comme dans la section \ref{spaces-groupoids-section-notation}. Soit $(G, m)$ un espace algébrique en groupes sur $B$. Soit $B' \to B$ un morphisme d'espaces algébriques. L'image réciproque $(G_{B'}, m_{B'})$ est un espace algébrique en groupes sur $B'$.
\end{lemma}

\begin{proof}
Omis.
\end{proof}








\section{Propriétés des espaces algébriques en groupes}
\label{spaces-groupoids-section-properties-group-spaces}

\noindent
Dans cette section, nous rassemblons quelques propriétés simples des espaces algébriques en groupes, valables sur une base quelconque.

\begin{lemma}
\label{spaces-groupoids-lemma-group-scheme-separated}
Soit $S$ un schéma. Soit $B$ un espace algébrique sur $S$. Soit $G$ un espace algébrique en groupes sur $B$. Alors $G \to B$ est séparé (resp.\ quasi-séparé, resp.\ localement séparé) si et seulement si le morphisme unité $e : B \to G$ est une immersion fermée (resp.\ quasi-compact, resp.\ une immersion).
\end{lemma}

\begin{proof}
Rappelons que, d'après Morphismes d'espaces, lemme \ref{spaces-morphisms-lemma-section-immersion}, $e$ est une immersion fermée (resp.\ quasi-compact, resp.\ une immersion) si $G \to B$ est séparé (resp.\ quasi-séparé, resp.\ localement séparé). Réciproquement, considérons le diagramme
$$
\xymatrix{
G \ar[r]_-{\Delta_{G/B}} \ar[d] &
G \times_B G \ar[d]^{(g, g') \mapsto m(i(g), g')} \\
B \ar[r]^e & G
}
$$
Le point de vue fonctoriel en géométrie algébrique permet de vérifier que ce diagramme est cartésien. Autrement dit, $\Delta_{G/B}$ est un changement de base de $e$. Par conséquent, si $e$ est une immersion fermée (resp.\ quasi-compact, resp.\ une immersion), il en va de même de $\Delta_{G/B}$ ; voir Espaces, lemme \ref{spaces-lemma-base-change-immersions} (resp.\ Morphismes d'espaces, lemme \ref{spaces-morphisms-lemma-base-change-quasi-compact}, resp.\ Espaces, lemme \ref{spaces-lemma-base-change-immersions}).
\end{proof}

\begin{lemma}
\label{spaces-groupoids-lemma-group-scheme-unramified-or-lqf}
Soit $S$ un schéma. Soit $B$ un espace algébrique sur $S$. Soit $G$ un espace algébrique en groupes sur $B$. Supposons $G \to B$ localement de type fini. Alors $G \to B$ est non ramifié (resp.\ localement quasi-fini) si et seulement si $G \to B$ est non ramifié (resp.\ quasi-fini) en $e(b)$ pour tout $b \in |B|$.
\end{lemma}

\begin{proof}
D'après Morphismes d'espaces, lemme \ref{spaces-morphisms-lemma-where-unramified} (resp.\ Morphismes d'espaces, lemme \ref{spaces-morphisms-lemma-base-change-quasi-finite-locus}), il existe un plus grand sous-espace ouvert $U \subset G$ tel que $U \to B$ soit non ramifié (resp.\ localement quasi-fini), et la formation de $U$ commute aux changements de base. On se ramène ainsi au cas où $B = \Spec(k)$ est le spectre d'un corps. Soit $g \in G(K)$ un point à valeurs dans une extension $K/k$. Pour vérifier si $g$ appartient à $U$, nous pouvons effectuer le changement de base à $K$. Il suffit donc de montrer
$$
G \to \Spec(k)\text{ est non ramifié en }e
\Leftrightarrow
G \to \Spec(k)\text{ est non ramifié en }g
$$
pour un point $k$-rationnel $g$ (resp.\ de même pour la quasi-finitude en $g$ et en $e$). Cela est clair, puisque la translation par $g$ est un automorphisme de $G$ sur $k$.
\end{proof}

\begin{lemma}
\label{spaces-groupoids-lemma-open-over-which-unramified-or-lqf}
Soit $S$ un schéma. Soit $B$ un espace algébrique sur $S$. Soit $G$ un espace algébrique en groupes sur $B$. Supposons que $G \to B$ est localement de type fini.
\begin{enumerate}
\item Il existe un plus grand sous-espace ouvert $U \subset B$ tel que $G_U \to U$ soit non ramifié, et la formation de $U$ commute aux changements de base.
\item Il existe un plus grand sous-espace ouvert $U \subset B$ tel que $G_U \to U$ soit localement quasi-fini, et la formation de $U$ commute aux changements de base.
\end{enumerate}
\end{lemma}

\begin{proof}
D'après Morphismes d'espaces, lemme \ref{spaces-morphisms-lemma-where-unramified} (resp.\ Morphismes d'espaces, lemme \ref{spaces-morphisms-lemma-base-change-quasi-finite-locus}), il existe un plus grand sous-espace ouvert $W \subset G$ tel que $W \to B$ soit non ramifié (resp.\ localement quasi-fini). De plus, la formation de $W$ commute aux changements de base. Le lemme \ref{spaces-groupoids-lemma-group-scheme-unramified-or-lqf} donne alors $U = e^{-1}(W)$ dans les deux cas.
\end{proof}












\section{Exemples d'espaces algébriques en groupes}
\label{spaces-groupoids-section-examples-group-spaces}

\noindent
Si $G \to S$ est un schéma en groupes sur le schéma de base $S$, alors, pour tout espace algébrique $B$ sur $S$, son changement de base $G_B$ est un espace algébrique en groupes sur $B$, d'après le lemme \ref{spaces-groupoids-lemma-base-change-group-space}. Nous utiliserons souvent ce fait dans les exemples ci-dessous.

\begin{example}[Espace algébrique en groupes multiplicatif]
\label{spaces-groupoids-example-multiplicative-group}
Soit $B \to S$ comme dans la section \ref{spaces-groupoids-section-notation}. Considérons le foncteur qui associe à tout schéma $T$ sur $B$ le groupe $\Gamma(T, \mathcal{O}_T^*)$ des éléments inversibles parmi les sections globales du faisceau structural. Ce foncteur est représenté par l'espace algébrique en groupes
$$
\mathbf{G}_{m, B} = B \times_S \mathbf{G}_{m, S}
$$
sur $B$. Ici $\mathbf{G}_{m, S}$ est le schéma en groupes multiplicatif sur $S$ ; voir Groupoïdes, exemple \ref{groupoids-example-multiplicative-group}.
\end{example}

\begin{example}[Racines de l'unité comme espace algébrique en groupes]
\label{spaces-groupoids-example-roots-of-unity}
Soit $B \to S$ comme dans la section \ref{spaces-groupoids-section-notation}. Soit $n \in \mathbf{N}$. Considérons le foncteur qui associe à tout schéma $T$ sur $B$ le sous-groupe de $\Gamma(T, \mathcal{O}_T^*)$ formé des racines $n$-ièmes de l'unité. Ce foncteur est représenté par l'espace algébrique en groupes
$$
\mu_{n, B} = B \times_S \mu_{n, S}
$$
sur $B$. Ici $\mu_{n, S}$ est le schéma en groupes des racines $n$-ièmes de l'unité sur $S$ ; voir Groupoïdes, exemple \ref{groupoids-example-roots-of-unity}.
\end{example}

\begin{example}[Espace algébrique en groupes additif]
\label{spaces-groupoids-example-additive-group}
Soit $B \to S$ comme dans la section \ref{spaces-groupoids-section-notation}. Considérons le foncteur qui associe à tout schéma $T$ sur $B$ le groupe $\Gamma(T, \mathcal{O}_T)$ de sections globales du faisceau structural. Ceci est représentable par l'espace algébrique en groupes
$$
\mathbf{G}_{a, B} = B \times_S \mathbf{G}_{a, S}
$$
sur $B$. Ici $\mathbf{G}_{a, S}$ est le schéma en groupes additif sur $S$, voir groupoïdes, exemple \ref{groupoids-example-additive-group}.
\end{example}

\begin{example}[Espace algébrique en groupes linéaire général]
\label{spaces-groupoids-example-general-linear-group}
Soit $B \to S$ comme dans la section \ref{spaces-groupoids-section-notation}. Soit $n \geq 1$. Considérons le foncteur qui associe à tout schéma $T$ sur $B$ le groupe
$$
\text{GL}_n(\Gamma(T, \mathcal{O}_T))
$$
des matrices inversibles de taille $n \times n$ à coefficients dans les sections globales du faisceau structural. Ce foncteur est représenté par l'espace algébrique en groupes
$$
\text{GL}_{n, B} = B \times_S \text{GL}_{n, S}
$$
sur $B$. Ici, $\mathbf{G}_{m, S}$ est le schéma en groupes linéaire général sur $S$ ; voir Groupoïdes, exemple \ref{groupoids-example-general-linear-group}.
\end{example}

\begin{example}
\label{spaces-groupoids-example-determinant}
Soit $B \to S$ comme dans la section \ref{spaces-groupoids-section-notation}. Soit $n \geq 1$. Le déterminant définit un morphisme d'espaces algébriques en groupes
$$
\det : \text{GL}_{n, B} \longrightarrow \mathbf{G}_{m, B}
$$
sur $B$. C'est le changement de base du morphisme déterminant sur $S$ de Groupoïdes, exemple \ref{groupoids-example-determinant}.
\end{example}

\begin{example}[Espace algébrique en groupes constant]
\label{spaces-groupoids-example-constant-group}
Soit $B \to S$ comme dans la section \ref{spaces-groupoids-section-notation}. Soit $G$ un groupe abstrait. Considérons le foncteur qui associe à tout schéma $T$ sur $B$ le groupe des applications localement constantes $T \to G$ (où $T$ est muni de la topologie de Zariski et $G$ de la topologie discrète). Ce foncteur est représenté par l'espace algébrique en groupes
$$
G_B = B \times_S G_S
$$
sur $B$. Ici, $G_S$ est le schéma en groupes constant introduit dans Groupoïdes, exemple \ref{groupoids-example-constant-group}.
\end{example}





\section{Actions d'espaces algébriques en groupes}
\label{spaces-groupoids-section-action-group-space}

\noindent
Voir Groupoïdes, section \ref{groupoids-section-action-group-scheme}, pour les notations.

\begin{definition}
\label{spaces-groupoids-definition-action-group-space}
Soit $B \to S$ comme dans la section \ref{spaces-groupoids-section-notation}. Soit $(G, m)$ un espace algébrique en groupes sur $B$. Soit $X$ un espace algébrique sur $B$.
\begin{enumerate}
\item Une {\it action de $G$ sur l'espace algébrique $X/B$} est un morphisme $a : G \times_B X \to X$ sur $B$ tel que, pour tout schéma $T$ sur $B$, l'application $a : G(T) \times X(T) \to X(T)$ définisse sur $X(T)$ une structure de $G(T)$-ensemble.
\item Supposons que $X$ et $Y$ soient des espaces algébriques sur $B$, chacun muni d'une action de $G$. Un morphisme {\it équivariant}, ou plus précisément {\it $G$-équivariant}, $\psi : X \to Y$ est un morphisme d'espaces algébriques sur $B$ tel que, pour tout $T$ sur $B$, l'application $\psi : X(T) \to Y(T)$ soit un morphisme de $G(T)$-ensembles.
\end{enumerate}
\end{definition}

\noindent
Dans la situation (1), cela signifie que les diagrammes
\begin{equation}
\label{spaces-groupoids-equation-action}
\xymatrix{
G \times_B G \times_B X \ar[r]_-{1_G \times a} \ar[d]_{m \times 1_X} &
G \times_B X \ar[d]^a \\
G \times_B X \ar[r]^a & X
}
\quad
\xymatrix{
G \times_B X \ar[r]_-a & X \\
X\ar[u]^{e \times 1_X} \ar[ru]_{1_X}
}
\end{equation}
sont commutatifs. Dans la situation (2), cela signifie simplement que le diagramme
$$
\xymatrix{
G \times_B X \ar[r]_-{\text{id} \times f} \ar[d]_a &
G \times_B Y \ar[d]^a \\
X \ar[r]^f & Y
}
$$
est commutatif.

\begin{definition}
\label{spaces-groupoids-definition-free-action}
Soient $B \to S$, $G \to B$ et $X \to B$ comme dans la définition \ref{spaces-groupoids-definition-action-group-space}. Soit $a : G \times_B X \to X$ une action de $G$ sur $X/B$. Nous disons que cette action est {\it libre} si, pour tout schéma $T$ sur $B$, l'action $a : G(T) \times X(T) \to X(T)$ est une action libre du groupe $G(T)$ sur l'ensemble $X(T)$.
\end{definition}

\begin{lemma}
\label{spaces-groupoids-lemma-free-action}
Dans la situation de la définition \ref{spaces-groupoids-definition-free-action}, l'action $a$ est libre si et seulement si
$$
G \times_B X \to X \times_B X, \quad (g, x) \mapsto (a(g, x), x)
$$
est un monomorphisme d'espaces algébriques.
\end{lemma}

\begin{proof}
Cela résulte immédiatement des définitions.
\end{proof}









\section{Espaces principaux homogènes}
\label{spaces-groupoids-section-principal-homogeneous}

\noindent
Cette section est l'analogue de Groupoïdes, section \ref{groupoids-section-principal-homogeneous}. Nous suggérons de lire d'abord cette dernière.

\begin{definition}
\label{spaces-groupoids-definition-pseudo-torsor}
Soit $S$ un schéma. Soit $B$ un espace algébrique sur $S$. Soit $(G, m)$ un espace algébrique en groupes sur $B$. Soit $X$ un espace algébrique sur $B$ et soit $a : G \times_B X \to X$ une action de $G$ sur $X$.
\begin{enumerate}
\item Nous disons que $X$ est un {\it pseudo-$G$-torseur}, ou que $X$ est {\it formellement principal homogène sous $G$}, si le morphisme induit $G \times_B X \to X \times_B X$, $(g, x) \mapsto (a(g, x), x)$, est un isomorphisme.
\item Un pseudo-$G$-torseur $X$ est dit {\it trivial} s'il existe un isomorphisme $G$-équivariant $G \to X$ sur $B$, où $G$ agit sur lui-même par multiplication à gauche.
\end{enumerate}
\end{definition}

\noindent
Il est clair que, si $B' \to B$ est un morphisme d'espaces algébriques, l'image réciproque $X_{B'}$ d'un pseudo-$G$-torseur sur $B$ est un pseudo-$G_{B'}$-torseur sur $B'$.

\begin{lemma}
\label{spaces-groupoids-lemma-characterize-trivial-pseudo-torsors}
Dans la situation de la définition \ref{spaces-groupoids-definition-pseudo-torsor}, on a les propriétés suivantes.
\begin{enumerate}
\item L'espace algébrique $X$ est un pseudo-$G$-torseur si et seulement si, pour tout schéma $T$ sur $B$, l'ensemble $X(T)$ est vide ou bien l'action du groupe $G(T)$ sur $X(T)$ est simplement transitive.
\item Un pseudo-$G$-torseur $X$ est trivial si et seulement si le morphisme $X \to B$ admet une section.
\end{enumerate}
\end{lemma}

\begin{proof}
Omis.
\end{proof}

\begin{definition}
\label{spaces-groupoids-definition-principal-homogeneous-space}
Soit $S$ un schéma. Soit $B$ un espace algébrique sur $S$. Soit $(G, m)$ un espace algébrique en groupes sur $B$. Soit $X$ un pseudo-$G$-torseur sur $B$.
\begin{enumerate}
\item Nous disons que $X$ est un {\it espace principal homogène}, ou plus précisément un {\it espace principal homogène sur $B$, de groupe structural $G$}, s'il existe un recouvrement fpqc\footnote{Dans Groupoïdes, définition \ref{groupoids-definition-principal-homogeneous-space}, le type de torseur retenu par défaut est un pseudo-torseur qui devient trivial sur un recouvrement fpqc. Puisque $G$, en tant qu'espace algébrique, peut être considéré comme un faisceau en groupes, il existe déjà une notion de $G$-torseur qui correspond au torseur fppf ; voir le lemme \ref{spaces-groupoids-lemma-torsor}. Nous employons donc « espace principal homogène » pour désigner un pseudo-torseur localement trivial pour la topologie fpqc, et nous essayons d'éviter le mot « torseur » dans cette situation.} $\{B_i \to B\}_{i \in I}$ tel que chaque $X_{B_i} \to B_i$ admette une section (c'est-à-dire soit un pseudo-$G_{B_i}$-torseur trivial).
\item Soit $\tau \in \{Zariski, \etale, lisse, syntomique, fppf\}$. Nous disons que $X$ est un {\it $G$-torseur pour la topologie $\tau$}, ou un {\it $\tau$-$G$-torseur}, ou simplement un {\it $\tau$-torseur}, s'il existe un recouvrement pour la topologie $\tau$, $\{B_i \to B\}_{i \in I}$, tel que chaque $X_{B_i} \to B_i$ admette une section.
\item Si $X$ est un espace principal homogène sur $B$, de groupe structural $G$, nous disons qu'il est {\it quasi-isotrivial} s'il est un torseur pour la topologie étale.
\item Si $X$ est un espace principal homogène sur $B$, de groupe structural $G$, nous disons qu'il est {\it localement trivial} s'il est un torseur pour la topologie de Zariski.
\end{enumerate}
\end{definition}

\noindent
Nous disons parfois « soit $X$ un espace principal homogène sur $B$, de groupe structural $G$ » pour indiquer que $X$ est un espace algébrique sur $B$, muni d'une action de $G$ qui en fait un espace principal homogène sur $B$. Nous montrons ensuite que cette terminologie coïncide, lorsque les deux s'appliquent, avec celle introduite précédemment.

\begin{lemma}
\label{spaces-groupoids-lemma-torsor}
Soit $S$ un schéma. Soit $(G, m)$ un espace algébrique en groupes sur $S$. Soit $X$ un espace algébrique sur $S$ et soit $a : G \times_S X \to X$ une action de $G$ sur $X$. Alors $X$ est un $G$-torseur pour la topologie $fppf$ au sens de la définition \ref{spaces-groupoids-definition-principal-homogeneous-space} si et seulement si $X$ est un $G$-torseur sur $(\Sch/S)_{fppf}$ au sens de Cohomologie sur les sites, définition \ref{sites-cohomology-definition-torsor}.
\end{lemma}

\begin{proof}
Omis.
\end{proof}

\begin{lemma}
\label{spaces-groupoids-lemma-pseudo-torsor-implications}
Soit $S$ un schéma. Soit $B$ un espace algébrique sur $S$. Soit $G$ un espace algébrique en groupes sur $B$. Soit $X$ un pseudo-$G$-torseur sur $B$. Supposons $G$ et $X$ localement de type fini sur $B$.
\begin{enumerate}
\item Si $G \to B$ est non ramifié, alors $X \to B$ est non ramifié.
\item Si $G \to B$ est localement quasi-fini, alors $X \to B$ est localement quasi-fini.
\end{enumerate}
\end{lemma}

\begin{proof}
Démontrons (1). Par Morphismes d'espaces, lemme \ref{spaces-morphisms-lemma-where-unramified}, nous nous ramenons au cas où $B$ est le spectre d'un corps. Si $X$ est vide, le résultat est acquis. Sinon, après une extension du corps, nous pouvons supposer que $X$ possède un point. On a alors $G \cong X$, d'où le résultat.

\medskip\noindent
La démonstration de (2) est exactement la même, en utilisant Morphismes d'espaces, lemme \ref{spaces-morphisms-lemma-base-change-quasi-finite-locus}.
\end{proof}
















\section{Faisceaux quasi-cohérents équivariants}
\label{spaces-groupoids-section-equivariant}

\noindent
Comparer avec Groupoïdes, section \ref{groupoids-section-equivariant}.

\begin{definition}
\label{spaces-groupoids-definition-equivariant-module}
Soit $B \to S$ comme dans la section \ref{spaces-groupoids-section-notation}. Soit $(G, m)$ un espace algébrique en groupes sur $B$ et soit $a : G \times_B X \to X$ une action de $G$ sur l'espace algébrique $X$ sur $B$. Un {\it $\mathcal{O}_X$-module quasi-cohérent $G$-équivariant}, ou simplement un {\it $\mathcal{O}_X$-module quasi-cohérent équivariant}, est un couple $(\mathcal{F}, \alpha)$, où $\mathcal{F}$ est un $\mathcal{O}_X$-module quasi-cohérent et $\alpha$ un morphisme de $\mathcal{O}_{G \times_B X}$-modules
$$
\alpha : a^*\mathcal{F} \longrightarrow \text{pr}_1^*\mathcal{F}
$$
où $\text{pr}_1 : G \times_B X \to X$ est la projection, tel que
\begin{enumerate}
\item le diagramme
$$
\xymatrix{
(1_G \times a)^*\text{pr}_2^*\mathcal{F} \ar[r]_-{\text{pr}_{12}^*\alpha} &
\text{pr}_2^*\mathcal{F} \\
(1_G \times a)^*a^*\mathcal{F} \ar[u]^{(1_G \times a)^*\alpha} \ar@{=}[r] &
(m \times 1_X)^*a^*\mathcal{F} \ar[u]_{(m \times 1_X)^*\alpha}
}
$$
soit commutatif dans la catégorie des $\mathcal{O}_{G \times_B G \times_B X}$-modules, et
\item l'image réciproque
$$
(e \times 1_X)^*\alpha : \mathcal{F} \longrightarrow \mathcal{F}
$$
soit l'identité.
\end{enumerate}
Pour l'interprétation, comparer avec les diagrammes correspondants de l'équation (\ref{spaces-groupoids-equation-action}).
\end{definition}

\noindent
Remarquons que la commutativité du premier diagramme garantit que $(e \times 1_X)^*\alpha$ est un endomorphisme idempotent de $\mathcal{F}$ ; la condition (2) revient donc à demander que cet endomorphisme soit un isomorphisme.

\begin{lemma}
\label{spaces-groupoids-lemma-pullback-equivariant}
Soit $B \to S$ comme dans la section \ref{spaces-groupoids-section-notation}. Soit $G$ un espace algébrique en groupes sur $B$. Soit $f : X \to Y$ un morphisme $G$-équivariant entre des espaces algébriques sur $B$ munis d'actions de $G$. Alors l'image réciproque $f^*$, donnée par $(\mathcal{F}, \alpha) \mapsto (f^*\mathcal{F}, (1_G \times f)^*\alpha)$, définit un foncteur de la catégorie des faisceaux quasi-cohérents $G$-équivariants sur $Y$ vers celle des faisceaux quasi-cohérents $G$-équivariants sur $X$.
\end{lemma}

\begin{proof}
Omis.
\end{proof}





\section{Groupoïdes en espaces algébriques}
\label{spaces-groupoids-section-groupoids}

\noindent
Voir Groupoïdes, section \ref{groupoids-section-groupoids}, pour les notations.

\begin{definition}
\label{spaces-groupoids-definition-groupoid}
Soit $B \to S$ comme dans la section \ref{spaces-groupoids-section-notation}.
\begin{enumerate}
\item Un {\it groupoïde en espaces algébriques sur $B$} est un quintuple $(U, R, s, t, c)$ où $U$ et $R$ sont des espaces algébriques sur $B$, et où $s, t : R \to U$ et $c : R \times_{s, U, t} R \to R$ sont des morphismes d'espaces algébriques sur $B$ ayant la propriété suivante : pour tout schéma $T$ sur $B$, le quintuple
$$
(U(T), R(T), s, t, c)
$$
est une catégorie qui est un groupoïde.
\item Un {\it morphisme $f : (U, R, s, t, c) \to (U', R', s', t', c')$ de groupoïdes en espaces algébriques sur $B$} est donné par des morphismes d'espaces algébriques $f : U \to U'$ et $f : R \to R'$ sur $B$ ayant la propriété suivante : pour tout schéma $T$ sur $B$, les applications $f$ définissent un foncteur de la catégorie-groupoïde $(U(T), R(T), s, t, c)$ vers la catégorie-groupoïde $(U'(T), R'(T), s', t', c')$.
\end{enumerate}
\end{definition}

\noindent
Soit $(U, R, s, t, c)$ un groupoïde en espaces algébriques sur $B$. Remarquons qu'il existe d'uniques morphismes d'espaces algébriques $e : U \to R$ et $i : R \to R$ sur $B$ tels que, pour tout schéma $T$ sur $B$, l'application induite $e : U(T) \to R(T)$ donne les identités et $i : R(T) \to R(T)$ donne les inverses de la catégorie-groupoïde. Le septuple $(U, R, s, t, c, e, i)$ satisfait aux diagrammes commutatifs correspondant à chacun des axiomes (1), (2)(a), (2)(b), (3)(a) et (3)(b) de Groupoïdes, section \ref{groupoids-section-groupoids}. Réciproquement, pour tout septuple ayant cette propriété, le quintuple $(U, R, s, t, c)$ est un groupoïde en espaces algébriques sur $B$. Remarquons que $i$ est un isomorphisme et que $e$ est une section de $s$ et de $t$. De plus, étant donné un groupoïde en espaces algébriques sur $B$, nous notons
$$
j = (t, s) : R \longrightarrow U \times_B U
$$
conformément aux conventions de la section \ref{spaces-groupoids-section-equivalence-relations} ci-dessus. Nous disons parfois « soit $(U, R, s, t, c, e, i)$ un groupoïde en espaces algébriques sur $B$ » pour souligner l'existence de l'identité et de l'inversion.

\begin{lemma}
\label{spaces-groupoids-lemma-groupoid-pre-equivalence}
Soit $B \to S$ comme dans la section \ref{spaces-groupoids-section-notation}. Étant donné un groupoïde en espaces algébriques $(U, R, s, t, c)$ sur $B$, le morphisme $j : R \to U \times_B U$ est une pré-relation d'équivalence.
\end{lemma}

\begin{proof}
Omis. C'est un joli exercice sur les définitions.
\end{proof}

\begin{lemma}
\label{spaces-groupoids-lemma-equivalence-groupoid}
Soit $B \to S$ comme dans la section \ref{spaces-groupoids-section-notation}. Étant donné une relation d'équivalence $j : R \to U \times_B U$ sur $B$, il existe une façon unique de l'étendre à un groupoïde en espaces algébriques $(U, R, s, t, c)$ sur $B$.
\end{lemma}

\begin{proof}
Omis. C'est un joli exercice sur les définitions.
\end{proof}

\begin{lemma}
\label{spaces-groupoids-lemma-diagram}
Soit $B \to S$ comme dans la section \ref{spaces-groupoids-section-notation}. Soit $(U, R, s, t, c)$ un groupoïde en espaces algébriques sur $B$. Dans le diagramme commutatif
$$
\xymatrix{
& U & \\
R \ar[d]_s \ar[ru]^t &
R \times_{s, U, t} R
\ar[l]^-{\text{pr}_0} \ar[d]^{\text{pr}_1} \ar[r]_-c &
R \ar[d]^s \ar[lu]_t \\
U & R \ar[l]_t \ar[r]^s & U
}
$$
les deux carrés inférieurs sont cartésiens. De plus, le triangle supérieur (qui est en réalité un carré) est lui aussi cartésien.
\end{lemma}

\begin{proof}
Omis. C'est un exercice sur les définitions et le point de vue fonctoriel en géométrie algébrique.
\end{proof}

\begin{lemma}
\label{spaces-groupoids-lemma-diagram-pull}
Soit $B \to S$ comme dans la section \ref{spaces-groupoids-section-notation}. Soit $(U, R, s, t, c, e, i)$ un groupoïde en espaces algébriques sur $B$. Le diagramme
\begin{equation}
\label{spaces-groupoids-equation-pull}
\xymatrix{
R \times_{t, U, t} R
\ar@<1ex>[r]^-{\text{pr}_1} \ar@<-1ex>[r]_-{\text{pr}_0}
\ar[d]_{\text{pr}_0 \times c \circ (i, 1)} &
R \ar[r]^t \ar[d]^{\text{id}_R} &
U \ar[d]^{\text{id}_U} \\
R \times_{s, U, t} R
\ar@<1ex>[r]^-c \ar@<-1ex>[r]_-{\text{pr}_0} \ar[d]_{\text{pr}_1} &
R \ar[r]^t \ar[d]^s &
U \\
R \ar@<1ex>[r]^s \ar@<-1ex>[r]_t &
U
}
\end{equation}
est commutatif. Les deux lignes supérieures sont isomorphes par les flèches verticales indiquées. Les deux carrés inférieurs de gauche sont cartésiens.
\end{lemma}

\begin{proof}
La commutativité du diagramme résulte des axiomes d'un groupoïde. Remarquons que, en termes de groupoïdes, la flèche verticale supérieure de gauche associe à une paire de morphismes $(\alpha, \beta)$ de même but la paire $(\alpha, \alpha^{-1} \circ \beta)$. Dans tout groupoïde, cela définit une bijection entre $\text{Flèches} \times_{t, \text{Ob}, t} \text{Flèches}$ et $\text{Flèches} \times_{s, \text{Ob}, t} \text{Flèches}$. Cela démontre la deuxième assertion du lemme. La dernière résulte du lemme \ref{spaces-groupoids-lemma-diagram}.
\end{proof}

\begin{lemma}
\label{spaces-groupoids-lemma-base-change-groupoid}
Soit $B \to S$ comme dans la section \ref{spaces-groupoids-section-notation}. Soit $(U, R, s, t, c)$ un groupoïde en espaces algébriques sur $B$. Soit $B' \to B$ un morphisme d'espaces algébriques. Alors les changements de base $U' = B' \times_B U$, $R' = B' \times_B R$ dotés des changements de base $s'$, $t'$, $c'$ des morphismes $s, t, c$ forment un groupoïde en espaces algébriques $(U', R', s', t', c')$ sur $B'$ et les projections déterminent un morphisme $(U', R', s', t', c') \to (U, R, s, t, c)$ de groupoïdes en espaces algébriques sur $B$.
\end{lemma}

\begin{proof}
Omis. Indice : $R' \times_{s', U', t'} R' = B' \times_B (R \times_{s, U, t} R)$.
\end{proof}





\section{Faisceaux quasi-cohérents sur les groupoïdes}
\label{spaces-groupoids-section-groupoids-quasi-coherent}

\noindent
Comparer avec Groupoïdes, section \ref{groupoids-section-groupoids-quasi-coherent}.

\begin{definition}
\label{spaces-groupoids-definition-groupoid-module}
Soit $B \to S$ comme dans la section \ref{spaces-groupoids-section-notation}. Soit $(U, R, s, t, c)$ un groupoïde en espaces algébriques sur $B$. Un {\it module quasi-cohérent sur $(U, R, s, t, c)$} est un couple $(\mathcal{F}, \alpha)$, où $\mathcal{F}$ est un $\mathcal{O}_U$-module quasi-cohérent et $\alpha$ un morphisme de $\mathcal{O}_R$-modules
$$
\alpha : t^*\mathcal{F} \longrightarrow s^*\mathcal{F}
$$
tel que
\begin{enumerate}
\item le diagramme
$$
\xymatrix{
& \text{pr}_1^*t^*\mathcal{F} \ar[r]_-{\text{pr}_1^*\alpha} &
\text{pr}_1^*s^*\mathcal{F} \ar@{=}[rd] & \\
\text{pr}_0^*s^*\mathcal{F} \ar@{=}[ru] & & & c^*s^*\mathcal{F} \\
& \text{pr}_0^*t^*\mathcal{F} \ar[lu]^{\text{pr}_0^*\alpha} \ar@{=}[r] &
c^*t^*\mathcal{F} \ar[ru]_{c^*\alpha}
}
$$
soit commutatif dans la catégorie des $\mathcal{O}_{R \times_{s, U, t} R}$-modules, et
\item l'image réciproque
$$
e^*\alpha : \mathcal{F} \longrightarrow \mathcal{F}
$$
soit l'identité.
\end{enumerate}
Comparer avec les diagrammes commutatifs du lemme \ref{spaces-groupoids-lemma-diagram}.
\end{definition}

\noindent
La commutativité du premier diagramme impose à l'endomorphisme $e^*\alpha$ d'être idempotent. La deuxième condition peut donc se reformuler en disant que $e^*\alpha$ est un isomorphisme. En fait, cette condition implique que $\alpha$ est un isomorphisme.

\begin{lemma}
\label{spaces-groupoids-lemma-isomorphism}
Soit $B \to S$ comme dans la section \ref{spaces-groupoids-section-notation}. Soit $(U, R, s, t, c)$ un groupoïde en espaces algébriques sur $B$. Si $(\mathcal{F}, \alpha)$ est un module quasi-cohérent sur $(U, R, s, t, c)$, alors $\alpha$ est un isomorphisme.
\end{lemma}

\begin{proof}
Prenons l'image réciproque du diagramme commutatif de la définition \ref{spaces-groupoids-definition-groupoid-module} par le morphisme $(i, 1) : R \to R \times_{s, U, t} R$. Nous obtenons $i^*\alpha \circ \alpha = s^*e^*\alpha$. En prenant l'image réciproque par $(1, i)$, nous obtenons la relation $\alpha \circ i^*\alpha = t^*e^*\alpha$. Par la seconde hypothèse, ces morphismes sont les identités. Ainsi, $i^*\alpha$ est un inverse de $\alpha$.
\end{proof}

\begin{lemma}
\label{spaces-groupoids-lemma-pullback}
Soit $B \to S$ comme dans la section \ref{spaces-groupoids-section-notation}. Considérons un morphisme $f : (U, R, s, t, c) \to (U', R', s', t', c')$ de groupoïdes en espaces algébriques sur $B$. L'image réciproque $f^*$, donnée par
$$
(\mathcal{F}, \alpha) \mapsto (f^*\mathcal{F}, f^*\alpha)
$$
définit un foncteur de la catégorie des faisceaux quasi-cohérents sur $(U', R', s', t', c')$ vers celle des faisceaux quasi-cohérents sur $(U, R, s, t, c)$.
\end{lemma}

\begin{proof}
Omis.
\end{proof}

\begin{lemma}
\label{spaces-groupoids-lemma-pushforward}
Soit $B \to S$ comme dans la section \ref{spaces-groupoids-section-notation}. Considérons un morphisme $f : (U, R, s, t, c) \to (U', R', s', t', c')$ de groupoïdes en espaces algébriques sur $B$. Supposons que
\begin{enumerate}
\item $f : U \to U'$ est quasi-compact et quasi-séparé,
\item le carré
$$
\xymatrix{
R \ar[d]_t \ar[r]_f & R' \ar[d]^{t'} \\
U \ar[r]^f & U'
}
$$
est cartésien, et
\item $s'$ et $t'$ sont plats.
\end{enumerate}
Alors l'image directe $f_*$, donnée par
$$
(\mathcal{F}, \alpha) \mapsto (f_*\mathcal{F}, f_*\alpha)
$$
définit un foncteur de la catégorie des faisceaux quasi-cohérents sur $(U, R, s, t, c)$ vers celle des faisceaux quasi-cohérents sur $(U', R', s', t', c')$ ; ce foncteur est adjoint à droite au foncteur image réciproque défini dans le lemme \ref{spaces-groupoids-lemma-pullback}.
\end{lemma}

\begin{proof}
Puisque $U \to U'$ est quasi-compact et quasi-séparé, $f_*$ transforme les faisceaux quasi-cohérents en faisceaux quasi-cohérents (Morphismes d'espaces, lemme \ref{spaces-morphisms-lemma-pushforward}). De plus, puisque les carrés
$$
\vcenter{
\xymatrix{
R \ar[d]_t \ar[r]_f & R' \ar[d]^{t'} \\
U \ar[r]^f & U'
}
}
\quad\text{et}\quad
\vcenter{
\xymatrix{
R \ar[d]_s \ar[r]_f & R' \ar[d]^{s'} \\
U \ar[r]^f & U'
}
}
$$
sont cartésiens, nous obtenons $(t')^*f_*\mathcal{F} = f_*t^*\mathcal{F}$ et $(s')^*f_*\mathcal{F} = f_*s^*\mathcal{F}$ ; voir Cohomologie des espaces, lemme \ref{spaces-cohomology-lemma-flat-base-change-cohomology}. On peut donc considérer $f_*\alpha$ comme un morphisme $(t')^*f_*\mathcal{F} \to (s')^*f_*\mathcal{F}$. Un argument semblable montre que $f_*\alpha$ satisfait à la condition de cocycle. Ce foncteur est adjoint au foncteur image réciproque parce que, pour les modules sur des espaces annelés, image réciproque et image directe sont adjointes. Nous omettons quelques détails.
\end{proof}


\begin{lemma}
\label{spaces-groupoids-lemma-colimits}
Soit $B \to S$ comme dans la section \ref{spaces-groupoids-section-notation}. Soit $(U, R, s, t, c)$ un groupoïde en espaces algébriques sur $B$. La catégorie des modules quasi-cohérents sur $(U, R, s, t, c)$ admet des colimites.
\end{lemma}

\begin{proof}
Soit $i \mapsto (\mathcal{F}_i, \alpha_i)$ un diagramme indexé par la catégorie $\mathcal{I}$. Nous pouvons former la colimite $\mathcal{F} = \colim \mathcal{F}_i$, qui est un faisceau quasi-cohérent sur $U$ ; voir Propriétés des espaces, lemme \ref{spaces-properties-lemma-properties-quasi-coherent}. Comme les colimites commutent aux images réciproques, on a $s^*\mathcal{F} = \colim s^*\mathcal{F}_i$ et, de même, $t^*\mathcal{F} = \colim t^*\mathcal{F}_i$. Nous pouvons donc poser $\alpha = \colim \alpha_i$. Nous omettons de vérifier que $(\mathcal{F}, \alpha)$ est la colimite du diagramme dans la catégorie des modules quasi-cohérents sur $(U, R, s, t, c)$.
\end{proof}

\begin{lemma}
\label{spaces-groupoids-lemma-abelian}
Soit $B \to S$ comme dans la section \ref{spaces-groupoids-section-notation}. Soit $(U, R, s, t, c)$ un groupoïde en espaces algébriques sur $B$. Si $s$, $t$ sont plats, alors la catégorie de modules quasi-cohérents sur $(U, R, s, t, c)$ est abélienne.
\end{lemma}

\begin{proof}
Soit $\varphi : (\mathcal{F}, \alpha) \to (\mathcal{G}, \beta)$ un homomorphisme de modules quasi-cohérents sur $(U, R, s, t, c)$. Puisque $s$ est plat, la suite
$$
0 \to s^*\Ker(\varphi)
\to s^*\mathcal{F} \to s^*\mathcal{G} \to s^*\Coker(\varphi) \to 0
$$
est exacte, et il en va de même après image réciproque par $t$. Ainsi, $\alpha$ et $\beta$ induisent des isomorphismes $\kappa : t^*\Ker(\varphi) \to s^*\Ker(\varphi)$ et $\lambda : t^*\Coker(\varphi) \to s^*\Coker(\varphi)$ qui satisfont à la condition de cocycle. Il est alors immédiat de vérifier que $(\Ker(\varphi), \kappa)$ et $(\Coker(\varphi), \lambda)$ sont respectivement un noyau et un conoyau dans la catégorie des modules quasi-cohérents sur $(U, R, s, t, c)$. De plus, l'égalité $\Coim(\varphi) = \Im(\varphi)$ résulte de ce qu'elle vaut sur $U$.
\end{proof}










\section{Colimites de modules quasi-cohérents}
\label{spaces-groupoids-section-colimits}

\noindent
Cette section est l'analogue de Groupoïdes, section \ref{groupoids-section-colimits}.

\begin{lemma}
\label{spaces-groupoids-lemma-construct-quasi-coherent}
Soit $B \to S$ comme dans la section \ref{spaces-groupoids-section-notation}. Soit $(U, R, s, t, c)$ un groupoïde en espaces algébriques sur $B$. Supposons $s$ et $t$ plats, quasi-compacts et quasi-séparés. Pour tout module quasi-cohérent $\mathcal{G}$ sur $U$, il existe un isomorphisme canonique $\alpha : t^*s_*t^*\mathcal{G} \to s^*s_*t^*\mathcal{G}$ qui munit $(s_*t^*\mathcal{G}, \alpha)$ d'une structure de module quasi-cohérent sur $(U, R, s, t, c)$. Cette construction définit un foncteur
$$
\QCoh(\mathcal{O}_U) \longrightarrow \QCoh(U, R, s, t, c)
$$
adjoint à droite au foncteur d'oubli $(\mathcal{F}, \beta) \mapsto \mathcal{F}$.
\end{lemma}

\begin{proof}
L'image directe d'un module quasi-cohérent par un morphisme quasi-compact et quasi-séparé est quasi-cohérente ; voir Morphismes d'espaces, lemme \ref{spaces-morphisms-lemma-pushforward}. Ainsi, $s_*t^*\mathcal{G}$ est quasi-cohérent. Avec les notations du lemme \ref{spaces-groupoids-lemma-diagram}, nous avons
$$
t^*s_*t^*\mathcal{G} =
\text{pr}_{1, *}\text{pr}_0^*t^*\mathcal{G} =
\text{pr}_{1, *}c^*t^*\mathcal{G} =
s^*s_*t^*\mathcal{G}
$$
L'égalité du milieu vient de ce que $t \circ c = t \circ \text{pr}_0$ comme morphismes $R \times_{s, U, t} R \to U$ ; la première et la dernière viennent de la commutation du changement de base et de l'image directe dans ces cas, par Cohomologie des espaces, lemme \ref{spaces-cohomology-lemma-flat-base-change-cohomology}.

\medskip\noindent
Pour vérifier la condition de cocycle de la définition \ref{spaces-groupoids-definition-groupoid-module} pour $\alpha$, ainsi que la propriété d'adjonction, décrivons autrement la construction $\mathcal{G} \mapsto (s_*t^*\mathcal{G}, \alpha)$. Considérons le groupoïde en schémas $(R, R \times_{t, U, t} R, \text{pr}_0, \text{pr}_1, \text{pr}_{02})$ associé à la relation d'équivalence $R \times_{t, U, t} R$ sur $R$ ; voir le lemme \ref{spaces-groupoids-lemma-equivalence-groupoid}. Il existe un morphisme
$$
f :
(R, R \times_{t, U, t} R, \text{pr}_1, \text{pr}_0, \text{pr}_{02})
\longrightarrow
(U, R, s, t, c)
$$
de groupoïdes en schémas, donné par $s : R \to U$ et par le morphisme $R \times_{t, U, t} R \to R$, $(r_0, r_1) \mapsto r_0^{-1} \circ r_1$ ; nous omettons de vérifier la commutativité des diagrammes requis. Puisque $t, s : R \to U$ sont quasi-compacts, quasi-séparés et plats, et puisque le carré
$$
\xymatrix{
R \times_{t, U, t} R \ar[d]_{\text{pr}_0}
\ar[rr]_-{(r_0, r_1) \mapsto r_0^{-1} \circ r_1} & & R \ar[d]^t \\
R \ar[rr]^s & & U
}
$$
est cartésien, le lemme \ref{spaces-groupoids-lemma-diagram-pull} montre que le lemme \ref{spaces-groupoids-lemma-pushforward} s'applique à $f$. Ainsi, les foncteurs image directe et image réciproque des modules quasi-cohérents suivant $f$ sont adjoints. Pour terminer la démonstration, identifions-les aux foncteurs décrits ci-dessus. Remarquons que
$$
t^* :
\QCoh(\mathcal{O}_U)
\longrightarrow
\QCoh(R, R \times_{t, U, t} R, \text{pr}_1, \text{pr}_0, \text{pr}_{02})
$$
est une équivalence par la théorie de la descente des faisceaux quasi-cohérents, puisque $\{t : R \to U\}$ est un recouvrement fpqc ; voir Descente sur les espaces, proposition \ref{spaces-descent-proposition-fpqc-descent-quasi-coherent}.

\medskip\noindent
L'image directe suivant $f$, précomposée avec l'équivalence $t^*$, envoie $\mathcal{G}$ sur $(s_*t^*\mathcal{G}, \alpha)$ ; nous omettons de vérifier que l'isomorphisme $\alpha$ ainsi obtenu est celui construit ci-dessus.

\medskip\noindent
L'image réciproque suivant $f$, postcomposée avec l'inverse de l'équivalence $t^*$, envoie $(\mathcal{F}, \beta)$ sur le module obtenu par descente, relativement à $\{t : R \to U\}$, à partir de $s^*\mathcal{F}$ muni de la donnée de descente $\gamma$ sur $R \times_{t, U, t} R$ qui est l'image réciproque de $\beta$ par $(r_0, r_1) \mapsto r_0^{-1} \circ r_1$. Considérons l'isomorphisme $\beta : t^*\mathcal{F} \to s^*\mathcal{F}$. Par transport au moyen de $\beta$, la donnée de descente canonique (Descente sur les espaces, définition \ref{spaces-descent-definition-descent-datum-effective-quasi-coherent}) sur $t^*\mathcal{F}$ relativement à $\{t : R \to U\}$ devient le morphisme
$$
\text{pr}_0^*s^*\mathcal{F}
\xrightarrow{\text{pr}_0^*\beta^{-1}}
\text{pr}_0^*t^*\mathcal{F}
\xrightarrow{can}
\text{pr}_1^*t^*\mathcal{F}
\xrightarrow{\text{pr}_1^*\beta}
\text{pr}_1^*s^*\mathcal{F}
$$
Puisque $\beta$ satisfait à la condition de cocycle, ce morphisme est égal à l'image réciproque de $\beta$ par $(r_0, r_1) \mapsto r_0^{-1} \circ r_1$. Pour le voir, prenons la relation de cocycle de la définition \ref{spaces-groupoids-definition-groupoid-module} et son image réciproque par le morphisme $(\text{pr}_0, c \circ (i, 1)) : R \times_{t, U, t} R \to R \times_{s, U, t} R$, qui intervient aussi dans le diagramme commutatif du lemme \ref{spaces-groupoids-lemma-diagram-pull}. Il s'ensuit que $(s^*\mathcal{F}, \gamma)$ est isomorphe à $(t^*\mathcal{F}, can)$. En définitive, l'image réciproque suivant $f$, postcomposée avec l'inverse de l'équivalence $t^*$, est isomorphe au foncteur d'oubli $(\mathcal{F}, \beta) \mapsto \mathcal{F}$.
\end{proof}

\begin{remark}
\label{spaces-groupoids-remark-adjunction-map}
Dans la situation du lemme \ref{spaces-groupoids-lemma-construct-quasi-coherent}, notons
$$
F : \QCoh(U, R, s, t, c) \to \QCoh(\mathcal{O}_U),\quad
(\mathcal{F}, \beta) \mapsto \mathcal{F}
$$
le foncteur d'oubli, et notons
$$
G : \QCoh(\mathcal{O}_U) \to \QCoh(U, R, s, t, c),\quad
\mathcal{G} \mapsto (s_*t^*\mathcal{G}, \alpha)
$$
l'adjoint à droite construit dans le lemme. Alors l'unité $\eta : \text{id} \to G \circ F$ de l'adjonction, évaluée en $(\mathcal{F}, \beta)$, est donnée par le morphisme
$$
\mathcal{F} \to s_*s^*\mathcal{F} \xrightarrow{\beta^{-1}} s_*t^*\mathcal{F}
$$
Nous omettons la vérification.
\end{remark}

\begin{lemma}
\label{spaces-groupoids-lemma-push-pull}
Soit $S$ un schéma. Soit $f : Y \to X$ un morphisme d'espaces algébriques sur $S$. Soit $\mathcal{F}$ un $\mathcal{O}_X$-module quasi-cohérent, soit $\mathcal{G}$ un $\mathcal{O}_Y$-module quasi-cohérent et soit $\varphi : \mathcal{G} \to f^*\mathcal{F}$ un morphisme de modules. Supposons que
\begin{enumerate}
\item $\varphi$ est injectif,
\item $f$ est quasi-compact, quasi-séparé, plat, et surjectif,
\item $X$ et $Y$ sont localement noethériens, et
\item $\mathcal{G}$ est un $\mathcal{O}_Y$-module cohérent.
\end{enumerate}
Alors $\mathcal{F} \cap f_*\mathcal{G}$, défini par le carré cartésien
$$
\xymatrix{
\mathcal{F} \ar[r] & f_*f^*\mathcal{F} \\
\mathcal{F} \cap f_*\mathcal{G} \ar[u] \ar[r] &
f_*\mathcal{G} \ar[u]
}
$$
est un $\mathcal{O}_X$-module cohérent.
\end{lemma}

\begin{proof}
Nous utiliserons librement la caractérisation des modules cohérents donnée dans Cohomologie des espaces, lemme \ref{spaces-cohomology-lemma-coherent-Noetherian}, ainsi que le fait que les modules cohérents forment une sous-catégorie de Serre de $\QCoh(\mathcal{O}_X)$ ; voir Cohomologie des espaces, lemme \ref{spaces-cohomology-lemma-coherent-Noetherian-quasi-coherent-sub-quotient}. Si $f$ admet une section $\sigma$, alors $\mathcal{F} \cap f_*\mathcal{G}$ est contenu dans l'image de $\sigma^*\mathcal{G} \to \sigma^*f^*\mathcal{F} = \mathcal{F}$, donc est cohérent. En général, pour montrer que $\mathcal{F} \cap f_*\mathcal{G}$ est cohérent, il suffit de montrer que $f^*(\mathcal{F} \cap f_*\mathcal{G})$ l'est (voir Descente sur les espaces, lemme \ref{spaces-descent-lemma-finite-type-descends}). Comme $f$ est plat, ce dernier est égal à $f^*\mathcal{F} \cap f^*f_*\mathcal{G}$. Puisque $f$ est plat, quasi-compact et quasi-séparé, on a $f^*f_*\mathcal{G} = p_*q^*\mathcal{G}$, où $p, q : Y \times_X Y \to Y$ sont les projections ; voir Cohomologie des espaces, lemme \ref{spaces-cohomology-lemma-flat-base-change-cohomology}. Comme $p$ admet une section, on conclut.
\end{proof}

\noindent
Soit $B \to S$ comme dans la section \ref{spaces-groupoids-section-notation}. Soit $(U, R, s, t, c)$ un groupoïde en espaces algébriques sur $B$. Supposons $U$ localement noethérien. Dans le lemme ci-dessous, nous disons qu'un faisceau quasi-cohérent $(\mathcal{F}, \alpha)$ sur $(U, R, s, t, c)$ est {\it cohérent} si $\mathcal{F}$ est un $\mathcal{O}_U$-module cohérent.

\begin{lemma}
\label{spaces-groupoids-lemma-colimit-coherent}
Soit $B \to S$ comme dans la section \ref{spaces-groupoids-section-notation}. Soit $(U, R, s, t, c)$ un groupoïde en espaces algébriques sur $B$. Supposons que
\begin{enumerate}
\item $U$ et $R$ sont noethériens,
\item $s$ et $t$ sont plats, quasi-compacts et quasi-séparés.
\end{enumerate}
Tout module quasi-cohérent $(\mathcal{F}, \alpha)$ sur $(U, R, s, t, c)$ est une colimite filtrante de modules cohérents.
\end{lemma}

\begin{proof}
Nous utiliserons sans autre mention la caractérisation des modules cohérents sur les espaces algébriques localement noethériens donnée dans Cohomologie des espaces, lemme \ref{spaces-cohomology-lemma-coherent-Noetherian}. Écrivons $\mathcal{F} = \colim \mathcal{H}_i$ comme colimite filtrante de sous-modules cohérents $\mathcal{H}_i \subset \mathcal{F}$ ; voir Cohomologie des espaces, lemme \ref{spaces-cohomology-lemma-directed-colimit-coherent}. Pour tout faisceau quasi-cohérent $\mathcal{H}$ sur $U$, notons $(s_*t^*\mathcal{H}, \alpha)$ le faisceau quasi-cohérent sur $(U, R, s, t, c)$ du lemme \ref{spaces-groupoids-lemma-construct-quasi-coherent}. Considérons le morphisme d'adjonction $(\mathcal{F}, \beta) \to (s_*t^*\mathcal{F}, \alpha)$ dans $\QCoh(U, R, s, t, c)$ ; voir la remarque \ref{spaces-groupoids-remark-adjunction-map}. Posons
$$
(\mathcal{F}_i, \beta_i) =
(\mathcal{F}, \beta)
\times_{(s_*t^*\mathcal{F}, \alpha)}
(s_*t^*\mathcal{H}_i, \alpha)
$$
dans $\QCoh(U, R, s, t, c)$. Puisque la restriction à $U$ est un foncteur exact sur $\QCoh(U, R, s, t, c)$ d'après la démonstration du lemme \ref{spaces-groupoids-lemma-abelian}, nous obtenons un carré cartésien
$$
\xymatrix{
\mathcal{F} \ar[r] & s_*t^*\mathcal{F} \\
\mathcal{F}_i \ar[r] \ar[u] & s_*t^*\mathcal{H}_i \ar[u]
}
$$
Autrement dit, $\mathcal{F}_i = \mathcal{F} \cap s_*t^*\mathcal{H}_i$. Par la description du morphisme d'adjonction dans la remarque \ref{spaces-groupoids-remark-adjunction-map}, ce diagramme est isomorphe au diagramme
$$
\xymatrix{
\mathcal{F} \ar[r] & s_*s^*\mathcal{F} \\
\mathcal{F}_i \ar[r] \ar[u] & s_*t^*\mathcal{H}_i \ar[u]
}
$$
où la flèche verticale de droite s'obtient en appliquant $s_*$ au morphisme
$$
t^*\mathcal{H}_i \to t^*\mathcal{F} \xrightarrow{\beta} s^*\mathcal{F}
$$
Ce morphisme est injectif puisque $t$ est plat. Il s'ensuit que $\mathcal{F}_i$ est cohérent, par le lemme \ref{spaces-groupoids-lemma-push-pull}. Enfin, comme $s$ est quasi-compact et quasi-séparé, $s_*$ commute aux colimites (voir Cohomologie des schémas, lemme \ref{coherent-lemma-colimit-cohomology}). Ainsi, $s_*t^*\mathcal{F} = \colim s_*t^*\mathcal{H}_i$, puis $(\mathcal{F}, \beta) = \colim (\mathcal{F}_i, \beta_i)$, comme voulu.
\end{proof}








\section{Cristaux en faisceaux quasi-cohérents}
\label{spaces-groupoids-section-crystals}

\noindent
Soit $(I, \Phi, j)$ un couple formé d'un ensemble $I$ et d'une pré-relation $j : \Phi \to I \times I$. Supposons donné, pour tout $i \in I$, un schéma $X_i$ et, pour tout $\phi \in \Phi$, un morphisme de schémas $f_\phi : X_{i'} \to X_i$, où $j(\phi) = (i, i')$. Posons $X = (\{X_i\}_{i \in I}, \{f_\phi\}_{\phi \in \Phi})$. Un {\it cristal en modules quasi-cohérents sur $X$} est, par définition, une règle qui associe à tout $i \in \Ob(\mathcal{I})$ un faisceau quasi-cohérent $\mathcal{F}_i$ sur $X_i$ et, à tout $\phi \in \Phi$ tel que $j(\phi) = (i, i')$, un isomorphisme
$$
\alpha_\phi : f_\phi^*\mathcal{F}_i \longrightarrow \mathcal{F}_{i'}
$$
de faisceaux quasi-cohérents sur $X_{i'}$. Ces cristaux en modules quasi-cohérents forment une catégorie additive $\textit{CQC}(X)$\footnote{Nous pourrions distinguer un ensemble de triplets $\phi, \phi', \phi'' \in \Phi$ tels que $j(\phi) = (i, i')$, $j(\phi') = (i', i'')$ et $j(\phi'') = (i, i'')$, avec $f_{\phi''} = f_\phi \circ f_{\phi'}$, et imposer $\alpha_{\phi'} \circ f_{\phi'}^*\alpha_\phi = \alpha_{\phi''}$ pour ces triplets. Cela définirait une sous-catégorie additive. Par exemple, les données $(I, \Phi)$ pourraient être l'ensemble des objets et des flèches d'une catégorie d'indices, et $X$ un diagramme de schémas sur cette catégorie. Le lemme \ref{spaces-groupoids-lemma-crystals-in-quasi-coherent-modules} donnerait aussitôt le résultat correspondant dans cette sous-catégorie.}. Cette catégorie admet des colimites (la démonstration est celle du lemme \ref{spaces-groupoids-lemma-colimits}). Si tous les morphismes $f_\phi$ sont plats, alors $\textit{CQC}(X)$ est abélienne (la démonstration est celle du lemme \ref{spaces-groupoids-lemma-abelian}). Soit $\kappa$ un cardinal. Nous disons qu'un cristal en modules quasi-cohérents $\mathcal{F}$ sur $X$ est {\it $\kappa$-généré} si chaque $\mathcal{F}_i$ est $\kappa$-généré (voir Propriétés, définition \ref{properties-definition-kappa-generated}).

\begin{lemma}
\label{spaces-groupoids-lemma-crystals-in-quasi-coherent-modules}
Dans la situation ci-dessus, si tous les morphismes $f_\phi$ sont plats, il existe un cardinal $\kappa$ tel que tout objet $(\{\mathcal{F}_i\}_{i \in I}, \{\alpha_\phi\}_{\phi \in \Phi})$ de $\textit{CQC}(X)$ soit la colimite filtrante de ses sous-modules $\kappa$-générés.
\end{lemma}

\begin{proof}
Dans l'énoncé et dans la démonstration, un {\it sous-module} de $(\{\mathcal{F}_i\}_{i \in I}, \{\alpha_\phi\}_{\phi \in \Phi})$ désigne la donnée, pour tout $i$, d'un sous-module quasi-cohérent $\mathcal{G}_i \subset \mathcal{F}_i$ tel que $\alpha_\phi(f_\phi^*\mathcal{G}_i) = \mathcal{G}_{i'}$ comme sous-faisceaux de $\mathcal{F}_{i'}$, pour tout $\phi \in \Phi$. Cela a un sens car, $f_\phi$ étant plat, l'image réciproque $f^*_\phi$ est exacte, c'est-à-dire qu'elle préserve les sous-faisceaux. La démonstration sera une variante de celle de Propriétés, lemme \ref{properties-lemma-colimit-kappa}. Nous invitons le lecteur à lire d'abord cette dernière.

\medskip\noindent
Nous affirmons qu'il suffit de démontrer le lemme lorsque tous les schémas $X_i$ sont affines. En effet, posons
$$
J = \coprod\nolimits_{i \in I} \{U \subset X_i\text{ ouvert affine}\}
$$
et
\begin{align*}
\Psi = & \coprod\nolimits_{\phi \in \Phi}
\{
(U, V) \mid
U \subset X_i, V \subset X_{i'}\text{ ouverts affines tels que } f_\phi(U) \subset V
\} \\
&
\amalg \coprod\nolimits_{i \in I}
\{
(U, U') \mid
U, U' \subset X_i\text{ ouverts affines tels que } U \subset U'
\}
\end{align*}
muni de l'application évidente $\Psi \to J \times J$. Notre $(\mathcal{F}, \alpha)$ induit alors un cristal en faisceaux quasi-cohérents $(\{\mathcal{H}_j\}_{j \in J}, \{\beta_\psi\}_{\psi \in \Psi})$ sur $Y = (J, \Psi)$ : on pose $\mathcal{H}_{(i, U)} = \mathcal{F}_i|_U$ pour $(i, U) \in J$ ; pour $\psi \in \Psi$, on prend pour $\beta_\psi$ la restriction de $\alpha_\phi$ à $U$ si $\psi = (\phi, U, V)$, et l'identité $\text{id} : (\mathcal{F}_i|_{U'})|_U \to \mathcal{F}_i|_U$ si $\psi = (i, U, U')$. De plus, les sous-modules de $(\{\mathcal{H}_j\}_{j \in J}, \{\beta_\psi\}_{\psi \in \Psi})$ correspondent bijectivement aux sous-modules de $(\{\mathcal{F}_i\}_{i \in I}, \{\alpha_\phi\}_{\phi \in \Phi})$. Nous omettons la démonstration (indication : utiliser Faisceaux, section \ref{sheaves-section-bases}). Enfin, il est clair que, si $\kappa$ convient pour $Y$, le même $\kappa$ convient pour $X$ (par la définition des modules $\kappa$-générés). Il suffit donc de démontrer le lemme pour les cristaux en faisceaux quasi-cohérents sur $Y$.

\medskip\noindent
Supposons désormais tous les schémas $X_i$ affines. Soit $\kappa$ un cardinal infini supérieur aux cardinaux de $I$ et de $\Phi$. Soit $(\{\mathcal{F}_i\}_{i \in I}, \{\alpha_\phi\}_{\phi \in \Phi})$ un objet de $\textit{CQC}(X)$. Pour tout $i$, écrivons $X_i = \Spec(A_i)$ et $M_i = \Gamma(X_i, \mathcal{F}_i)$. Pour tout $\phi \in \Phi$ tel que $j(\phi) = (i, i')$, le morphisme $\alpha_\phi$ se traduit par un isomorphisme de $A_{i'}$-modules
$$
\alpha_\phi : M_i \otimes_{A_i} A_{i'} \longrightarrow M_{i'}
$$
À l'aide de l'axiome du choix, choisissons une règle
$$
(\phi, m) \longmapsto S(\phi, m')
$$
dont la source est l'ensemble des couples $(\phi, m')$ tels que $\phi \in \Phi$, $j(\phi) = (i, i')$ et $m' \in M_{i'}$, et dont la valeur est une partie finie $S(\phi, m') \subset M_i$ telle que
$$
m' = \alpha_\phi\left(\sum\nolimits_{m \in S(\phi, m')} m \otimes a'_m\right)
$$
pour certains $a'_m \in A_{i'}$.

\medskip\noindent
Ces choix faits, nous affirmons que toute section de l'un quelconque des $\mathcal{F}_i$ sur $X_i$ appartient à un sous-module $\kappa$-généré. Soit en effet une famille $\mathcal{S} = \{S_i\}_{i \in I}$ de parties $S_i \subset M_i$, chacune de cardinal au plus $\kappa$. Définissons une nouvelle famille $\mathcal{S}' = \{S'_i\}_{i \in I}$ par
$$
S'_i = S_i \cup
\bigcup\nolimits_{(\phi, m'),\ j(\phi) = (i, i'),\ m' \in S_{i'}} S(\phi, m')
$$
Chaque $S'_i$ est encore de cardinal au plus $\kappa$. Posons $\mathcal{S}^{(0)} = \mathcal{S}$, $\mathcal{S}^{(1)} = \mathcal{S}'$ et, par récurrence, $\mathcal{S}^{(n + 1)} = (\mathcal{S}^{(n)})'$. Posons ensuite $S_i^{(\infty)} = \bigcup_{n \geq 0} S_i^{(n)}$ et $\mathcal{S}^{(\infty)} = \{S_i^{(\infty)}\}_{i \in I}$. Par construction, pour tout $\phi \in \Phi$ tel que $j(\phi) = (i, i')$ et tout $m' \in S^{(\infty)}_{i'}$, on peut écrire $m'$ comme combinaison linéaire finie d'éléments $\alpha_\phi(m \otimes 1)$ avec $m \in S_i^{(\infty)}$. Si $N_i$ est le sous-$A_i$-module de $M_i$ engendré par $S_i^{(\infty)}$, les sous-modules quasi-cohérents correspondants $\widetilde{N_i} \subset \mathcal{F}_i$ forment donc un sous-module $\kappa$-généré. Ceci achève la démonstration.
\end{proof}

\begin{lemma}
\label{spaces-groupoids-lemma-set-generators}
Soit $B \to S$ comme dans la section \ref{spaces-groupoids-section-notation}. Soit $(U, R, s, t, c)$ un groupoïde en espaces algébriques sur $B$. Si $s$ et $t$ sont plats, il existe un ensemble $T$ et une famille d'objets $(\mathcal{F}_t, \alpha_t)_{t \in T}$ de $\QCoh(U, R, s, t, c)$ tels que tout objet $(\mathcal{F}, \alpha)$ soit la colimite filtrante de ses sous-modules isomorphes à l'un des objets $(\mathcal{F}_t, \alpha_t)$.
\end{lemma}

\begin{proof}
Ce lemme généralise Groupoïdes, lemme \ref{groupoids-lemma-colimit-kappa}, qui traite le cas d'un groupoïde en schémas. Nous ne pouvons pas tout à fait reprendre le même argument ; nous allons donc utiliser les résultats sur les « cristaux en faisceaux quasi-cohérents » développés ci-dessus.

\medskip\noindent
Choisissons un schéma $W$ et un morphisme étale surjectif $W \to U$. Choisissons un schéma $V$ et un morphisme étale surjectif $V \to W \times_{U, s} R$. Choisissons un schéma $V'$ et un morphisme étale surjectif $V' \to R \times_{t, U} W$. Considérons la famille de schémas
$$
I = \{W, W \times_U W, V, V', V \times_R V'\}
$$
et l'ensemble des morphismes de schémas
$$
\Phi = \{\text{pr}_i : W \times_U W \to W, V \to W, V' \to W,
V \times_R V' \to V, V \times_R V' \to V'\}
$$
Posons $X = (I, \Phi)$. Rappelons que nous avons défini une catégorie $\textit{CQC}(X)$ de cristaux en faisceaux quasi-cohérents sur $X$. Il existe un foncteur
$$
\QCoh(U, R, s, t, c) \longrightarrow \textit{CQC}(X)
$$
qui associe à $(\mathcal{F}, \alpha)$ le faisceau $\mathcal{F}|_W$ sur $W$, le faisceau $\mathcal{F}|_{W \times_U W}$ sur $W \times_U W$, l'image réciproque de $\mathcal{F}$ sur $V$ par $V \to W \times_{U, s} R \to W \to U$, l'image réciproque de $\mathcal{F}$ sur $V'$ par $V' \to R \times_{t, U} W \to W \to U$, et enfin l'image réciproque de $\mathcal{F}$ sur $V \times_R V'$ par $V \times_R V' \to V \to W \times_{U, s} R \to W \to U$. Pour morphismes de comparaison $\{\alpha_\phi\}_{\phi \in \Phi}$, nous prenons les morphismes évidents, issus de l'associativité des images réciproques, sauf lorsque $\phi = \text{pr}_{V'} : V \times_R V' \to V'$ : nous prenons alors l'image réciproque sur $V \times_R V'$ de $\alpha : t^*\mathcal{F} \to s^*\mathcal{F}$. Cette définition est licite en vertu du diagramme commutatif
$$
\xymatrix{
& V \times_R V' \ar[ld] \ar[rd] \\
V \ar[rd] \ar[dd] & & V' \ar[ld] \ar[dd] \\
& R \ar@<-1ex>[dd]_s \ar@<1ex>[dd]^t \\
W \ar[rd] & & W \ar[ld] \\
& U
}
$$
Le foncteur ci-dessus n'est pas une équivalence de catégories. Toutefois, puisque $W \to U$ est étale surjectif, il est fidèle\footnote{En fait, le foncteur est pleinement fidèle, mais nous n'en aurons pas besoin.}. Comme tous les morphismes du diagramme précédent sont plats, c'est un foncteur exact entre catégories abéliennes. De plus, si $(\mathcal{F}, \alpha)$ a pour image $(\{\mathcal{F}_i\}_{i \in I}, \{\alpha_\phi\}_{\phi \in \Phi})$, nous affirmons qu'il existe une correspondance bijective entre les sous-modules quasi-cohérents de $(\mathcal{F}, \alpha)$ et ceux de $(\{\mathcal{F}_i\}_{i \in I}, \{\alpha_\phi\}_{\phi \in \Phi})$. En effet, pour un sous-module de $(\{\mathcal{F}_i\}_{i \in I}, \{\alpha_\phi\}_{\phi \in \Phi})$, la compatibilité du sous-module sur $W$ avec les projections $W \times_U W \to W$ garantit qu'il provient d'un sous-module quasi-cohérent de $\mathcal{F}$ (par Propriétés des espaces, proposition \ref{spaces-properties-proposition-quasi-coherent}) ; sa compatibilité avec $\alpha_{\text{pr}_{V'}}$ garantit ensuite que ce sous-faisceau est compatible avec $\alpha$. Nous omettons les détails.

\medskip\noindent
Choisissons un cardinal $\kappa$ comme dans le lemme \ref{spaces-groupoids-lemma-crystals-in-quasi-coherent-modules} pour le système $X = (I, \Phi)$. D'après Propriétés, lemme \ref{properties-lemma-set-of-iso-classes}, les classes d'isomorphisme des cristaux en faisceaux quasi-cohérents $\kappa$-générés sur $X$ forment un ensemble. Le résultat s'ensuit.
\end{proof}









\section{Groupoïdes et espaces algébriques en groupes}
\label{spaces-groupoids-section-groupoids-group-spaces}

\noindent
Comparer avec Groupoïdes, section \ref{groupoids-section-groupoids-group-schemes}.

\begin{lemma}
\label{spaces-groupoids-lemma-groupoid-from-action}
Soit $B \to S$ comme dans la section \ref{spaces-groupoids-section-notation}. Soit $(G, m)$ un espace algébrique en groupes sur $B$, d'unité $e_G$ et d'inversion $i_G$. Soit $X$ un espace algébrique sur $B$ et soit $a : G \times_B X \to X$ une action de $G$ sur $X$ au-dessus de $B$. Nous obtenons un groupoïde en espaces algébriques $(U, R, s, t, c, e, i)$ sur $B$ comme suit :
\begin{enumerate}
\item Nous posons $U = X$ et $R = G \times_B X$.
\item Nous définissons $s : R \to U$ par $(g, x) \mapsto x$.
\item Nous définissons $t : R \to U$ par $(g, x) \mapsto a(g, x)$.
\item Nous définissons $c : R \times_{s, U, t} R \to R$ par $((g, x), (g', x')) \mapsto (m(g, g'), x')$.
\item Nous définissons $e : U \to R$ par $x \mapsto (e_G(x), x)$.
\item Nous définissons $i : R \to R$ par $(g, x) \mapsto (i_G(g), a(g, x))$.
\end{enumerate}
\end{lemma}

\begin{proof}
Omis. Indication : il suffit de vérifier l'assertion au niveau des ensembles. Pour cela, utiliser la description précédant le lemme, où $g$ est vu comme une flèche de $v$ vers $a(g, v)$.
\end{proof}

\begin{lemma}
\label{spaces-groupoids-lemma-action-groupoid-modules}
Soit $B \to S$ comme dans la section \ref{spaces-groupoids-section-notation}. Soit $(G, m)$ un espace algébrique en groupes sur $B$. Soit $X$ un espace algébrique sur $B$ et soit $a : G \times_B X \to X$ une action de $G$ sur $X$ au-dessus de $B$. Soit $(U, R, s, t, c)$ le groupoïde en espaces algébriques construit dans le lemme \ref{spaces-groupoids-lemma-groupoid-from-action}. La règle $(\mathcal{F}, \alpha) \mapsto (\mathcal{F}, \alpha)$ définit une équivalence de catégories entre les $\mathcal{O}_X$-modules quasi-cohérents $G$-équivariants et la catégorie des modules quasi-cohérents sur $(U, R, s, t, c)$.
\end{lemma}

\begin{proof}
L'assertion a un sens parce que $t = a$ et $s = \text{pr}_1$ comme morphismes $R = G \times_B X \to X$ ; voir les définitions \ref{spaces-groupoids-definition-equivariant-module} et \ref{spaces-groupoids-definition-groupoid-module}. Au moyen de la traduction du lemme \ref{spaces-groupoids-lemma-groupoid-from-action}, les conditions de commutativité des deux définitions coïncident exactement.
\end{proof}





\section{L'espace algébrique en groupes stabilisateur}
\label{spaces-groupoids-section-stabilizer}

\noindent
Comparer avec Groupoïdes, section \ref{groupoids-section-stabilizer}. À tout groupoïde en espaces algébriques, on associe un espace algébrique en groupes de la manière suivante.

\begin{lemma}
\label{spaces-groupoids-lemma-groupoid-stabilizer}
Soit $B \to S$ comme dans la section \ref{spaces-groupoids-section-notation}. Soit $(U, R, s, t, c)$ un groupoïde en espaces algébriques sur $B$. L'espace algébrique $G$ défini par le carré cartésien
$$
\xymatrix{
G \ar[r] \ar[d] & R \ar[d]^{j = (t, s)} \\
U \ar[r]^-{\Delta} & U \times_B U
}
$$
est un espace algébrique en groupes sur $U$, dont la loi de composition $m$ est induite par la loi de composition $c$.
\end{lemma}

\begin{proof}
Cela résulte de ce que, dans une catégorie-groupoïde, l'ensemble des endomorphismes de tout objet est un groupe.
\end{proof}

\noindent
Puisque $\Delta$ est un monomorphisme, $G = j^{-1}(\Delta_{U/B})$ est un sous-faisceau de $R$. De ce point de vue, le morphisme structural $G = j^{-1}(\Delta_{U/B}) \to U$ est induit indifféremment par $s$ ou par $t$, et $m$ est induit par $c$.

\begin{definition}
\label{spaces-groupoids-definition-stabilizer-groupoid}
Soit $B \to S$ comme dans la section \ref{spaces-groupoids-section-notation}. Soit $(U, R, s, t, c)$ un groupoïde en espaces algébriques sur $B$. L'espace algébrique en groupes $j^{-1}(\Delta_{U/B}) \to U$ est appelé le {\it stabilisateur du groupoïde en espaces algébriques $(U, R, s, t, c)$}.
\end{definition}

\noindent
Dans la littérature, l'espace algébrique en groupes stabilisateur est souvent noté $S$ (sans doute parce que « stabilisateur » commence par un « s ») ; nous ne pouvons adopter cette notation, car $S$ désigne déjà le schéma de base.

\begin{lemma}
\label{spaces-groupoids-lemma-groupoid-action-stabilizer}
Soit $B \to S$ comme dans la section \ref{spaces-groupoids-section-notation}. Soit $(U, R, s, t, c)$ un groupoïde en espaces algébriques sur $B$ et soit $G/U$ son stabilisateur. Notons $R_t/U$ l'espace algébrique $R$ considéré comme espace algébrique sur $U$ par le morphisme $t : R \to U$. Il existe une action à gauche canonique
$$
a : G \times_U R_t \longrightarrow R_t
$$
induite par la loi de composition $c$.
\end{lemma}

\begin{proof}
En termes de points sur $T/B$, nous définissons $a(g, r) = c(g, r)$.
\end{proof}








\section{Restriction des groupoïdes}
\label{spaces-groupoids-section-restrict-groupoid}

\noindent
Voir Groupoïdes, section \ref{groupoids-section-restrict-groupoid}, pour les notations.

\begin{lemma}
\label{spaces-groupoids-lemma-restrict-groupoid}
Soit $B \to S$ comme dans la section \ref{spaces-groupoids-section-notation}. Soit $(U, R, s, t, c)$ un groupoïde en espaces algébriques sur $B$. Soit $g : U' \to U$ un morphisme d'espaces algébriques. Considérons le diagramme suivant
$$
\xymatrix{
R' \ar[d] \ar[r] \ar@/_3pc/[dd]_{t'} \ar@/^1pc/[rr]^{s'}&
R \times_{s, U} U' \ar[r] \ar[d] &
U' \ar[d]^g \\
U' \times_{U, t} R \ar[d] \ar[r] &
R \ar[r]^s \ar[d]_t &
U \\
U' \ar[r]^g &
U
}
$$
où tous les carrés sont cartésiens. Il existe alors une loi de composition canonique $c' : R' \times_{s', U', t'} R' \to R'$ telle que $(U', R', s', t', c')$ soit un groupoïde en espaces algébriques sur $B$ et que $U' \to U$, $R' \to R$ définissent un morphisme $(U', R', s', t', c') \to (U, R, s, t, c)$ de groupoïdes en espaces algébriques sur $B$. De plus, pour tout schéma $T$ sur $B$, le foncteur de groupoïdes
$$
(U'(T), R'(T), s', t', c') \to (U(T), R(T), s, t, c)
$$
est la restriction (voir Groupoïdes, section \ref{groupoids-section-restrict-groupoid}) de $(U(T), R(T), s, t, c)$ suivant l'application $U'(T) \to U(T)$.
\end{lemma}

\begin{proof}
Omis.
\end{proof}

\begin{definition}
\label{spaces-groupoids-definition-restrict-groupoid}
Soit $B \to S$ comme dans la section \ref{spaces-groupoids-section-notation}. Soit $(U, R, s, t, c)$ un groupoïde en espaces algébriques sur $B$. Soit $g : U' \to U$ un morphisme d'espaces algébriques sur $B$. Le morphisme de groupoïdes en espaces algébriques $(U', R', s', t', c') \to (U, R, s, t, c)$ construit dans le lemme \ref{spaces-groupoids-lemma-restrict-groupoid} est appelé la {\it restriction de $(U, R, s, t, c)$ à $U'$}. Dans cette situation, nous écrivons parfois $R' = R|_{U'}$.
\end{definition}

\begin{lemma}
\label{spaces-groupoids-lemma-restrict-groupoid-relation}
Les notions de restriction des groupoïdes et des (pré-)relations d'équivalence définies dans les définitions \ref{spaces-groupoids-definition-restrict-groupoid} et \ref{spaces-groupoids-definition-restrict-relation} coïncident par les constructions des lemmes \ref{spaces-groupoids-lemma-groupoid-pre-equivalence} et \ref{spaces-groupoids-lemma-equivalence-groupoid}.
\end{lemma}

\begin{proof}
L'assertion signifie que le $R'$ du lemme \ref{spaces-groupoids-lemma-restrict-groupoid} est aussi égal à
$$
R' = (U' \times_B U')\times_{U \times_B U} R
\longrightarrow
U' \times_B U'
$$
Cette formulation aurait d'ailleurs peut-être rendu le lemme plus clair.
\end{proof}





\section{Sous-espaces invariants}
\label{spaces-groupoids-section-invariant}

\noindent
Dans cette section, nous discutons brièvement de la notion de sous-espace invariant.

\begin{definition}
\label{spaces-groupoids-definition-invariant-open}
Soit $B \to S$ comme dans la section \ref{spaces-groupoids-section-notation}. Soit $(U, R, s, t, c)$ un groupoïde en espaces algébriques sur la base $B$.
\begin{enumerate}
\item Nous disons qu'un sous-espace ouvert $W \subset U$ est {\it $R$-invariant} si $t(s^{-1}(W)) \subset W$.
\item Un sous-espace localement fermé $Z \subset U$ est appelé {\it $R$-invariant} si $t^{-1}(Z) = s^{-1}(Z)$ comme sous-espaces localement fermés de $R$.
\item Un monomorphisme d'espaces algébriques $T \to U$ est {\it $R$-invariant} si $T \times_{U, t} R = R \times_{s, U} T$ comme espaces algébriques sur $R$.
\end{enumerate}
\end{definition}

\noindent
Pour un sous-espace ouvert $W \subset U$, l'invariance sous $R$ équivaut aussi à l'égalité $s^{-1}(W) = t^{-1}(W)$. Si $W \subset U$ est $R$-invariant, la restriction de $R$ à $W$ est simplement $R_W = s^{-1}(W) = t^{-1}(W)$. De même, si $Z \subset U$ est un sous-espace localement fermé $R$-invariant, la restriction de $R$ à $Z$ est simplement $R_Z = s^{-1}(Z) = t^{-1}(Z)$.

\begin{lemma}
\label{spaces-groupoids-lemma-constructing-invariant-opens}
Soit $B \to S$ comme dans la section \ref{spaces-groupoids-section-notation}. Soit $(U, R, s, t, c)$ un groupoïde en espaces algébriques sur $B$.
\begin{enumerate}
\item Si $s$ et $t$ sont ouverts, alors, pour tout ouvert $W \subset U$, l'ouvert $s(t^{-1}(W))$ est $R$-invariant.
\item Si $s$ et $t$ sont ouverts et quasi-compacts, alors $U$ admet un recouvrement ouvert par des sous-espaces ouverts quasi-compacts et $R$-invariants.
\end{enumerate}
\end{lemma}

\begin{proof}
Supposons $s$ et $t$ ouverts et $W \subset U$ ouvert. Puisque $s$ est ouvert, $W' = s(t^{-1}(W))$ est un sous-espace ouvert de $U$. Il est assez facile de voir, du point de vue fonctoriel, que c'est un ouvert $R$-invariant de $U$ ; nous allons néanmoins le démontrer directement au moyen de quelques diagrammes, car l'argument est instructif. Remarquons que $t^{-1}(W')$ est l'image du morphisme
$$
A := t^{-1}(W) \times_{s|_{t^{-1}(W)}, U, t} R
\xrightarrow{\text{pr}_1} R
$$
et que $s^{-1}(W')$ est l'image du morphisme
$$
B := R \times_{s, U, s|_{t^{-1}(W)}} t^{-1}(W)
\xrightarrow{\text{pr}_0} R.
$$
Les espaces algébriques $A$ et $B$, à gauche des flèches ci-dessus, sont respectivement des sous-espaces ouverts de $R \times_{s, U, t} R$ et de $R \times_{s, U, s} R$. D'après le lemme \ref{spaces-groupoids-lemma-diagram}, le diagramme
$$
\xymatrix{
R \times_{s, U, t} R \ar[rd]_{\text{pr}_1} \ar[rr]_{(\text{pr}_1, c)} & &
R \times_{s, U, s} R \ar[ld]^{\text{pr}_0} \\
& R &
}
$$
est commutatif et la flèche horizontale est un isomorphisme. De plus, il est clair que $(\text{pr}_1, c)(A) = B$. Nous en concluons que $s^{-1}(W') = t^{-1}(W')$ et que $W'$ est $R$-invariant. Cela démontre (1).

\medskip\noindent
Supposons maintenant $s$ et $t$ tous deux ouverts et quasi-compacts. Si $W \subset U$ est un ouvert quasi-compact, alors $W' = s(t^{-1}(W))$ est lui aussi ouvert quasi-compact, et il est invariant d'après ce qui précède. En faisant parcourir à $W$ les images des schémas affines étales sur $U$, on obtient (2).
\end{proof}





\section{Faisceaux quotients}
\label{spaces-groupoids-section-quotient-sheaves}

\noindent
Soit $S$ un schéma et soit $B$ un espace algébrique sur $S$. Soit $j : R \to U \times_B U$ une pré-relation sur $B$. Pour tout schéma $S'$ sur $S$, considérons la relation d'équivalence $\sim_{S'}$ engendrée par l'image de $j(S') : R(S') \to U(S') \times U(S')$. Nous obtenons ainsi un préfaisceau
\begin{equation}
\label{spaces-groupoids-equation-quotient-presheaf}
\begin{matrix}
(\Sch/S)^{opp}_{fppf} &
\longrightarrow &
\textit{Ens}, \\
S' &
\longmapsto  &
U(S')/\sim_{S'}
\end{matrix}
\end{equation}
Puisque $j$ est un morphisme d'espaces algébriques sur $B$ à valeurs dans $U \times_B U$, il existe une transformation canonique de préfaisceaux du préfaisceau (\ref{spaces-groupoids-equation-quotient-presheaf}) vers $B$.

\begin{definition}
\label{spaces-groupoids-definition-quotient-sheaf}
Soient $B \to S$ et la pré-relation $j : R \to U \times_B U$ comme ci-dessus. Dans cette situation, le {\it faisceau quotient $U/R$} associé à $j$ est le faisceau associé au préfaisceau (\ref{spaces-groupoids-equation-quotient-presheaf}) sur $(\Sch/S)_{fppf}$. Si $j : R \to U \times_B U$ provient de l'action sur $U$ d'un espace algébrique en groupes $G$ sur $B$, comme dans le lemme \ref{spaces-groupoids-lemma-groupoid-from-action}, nous notons le faisceau quotient $U/G$.
\end{definition}

\noindent
Cela signifie exactement que le diagramme
$$
\xymatrix{
R \ar@<1ex>[r] \ar@<-1ex>[r] &
U \ar[r] &
U/R
}
$$
est un diagramme coégalisateur dans la catégorie des faisceaux d'ensembles sur $(\Sch/S)_{fppf}$. Là encore, il existe un morphisme canonique de faisceaux $U/R \to B$, puisque $j$ est un morphisme d'espaces algébriques sur $B$ à valeurs dans $U \times_B U$.

\begin{remark}
\label{spaces-groupoids-remark-quotient-variant}
Une variante de la construction précédente consisterait à prendre le faisceau associé au foncteur
$$
\begin{matrix}
(\textit{Espaces}/B)^{opp}_{fppf} &
\longrightarrow &
\textit{Ens}, \\
X &
\longmapsto  &
U(X)/\sim_X
\end{matrix}
$$
où $\sim_X \subset U(X) \times U(X)$ est désormais la relation d'équivalence engendrée par l'image de $j : R(X) \to U(X) \times U(X)$. Bien entendu, $U(X) = \Mor_B(X, U)$ et $R(X) = \Mor_B(X, R)$. En fait, le résultat serait le même, au moyen des identifications de (insérer ici une référence future dans Topologies sur les espaces).
\end{remark}

\begin{definition}
\label{spaces-groupoids-definition-representable-quotient}
Dans la situation de la définition \ref{spaces-groupoids-definition-quotient-sheaf}, nous disons que la pré-relation $j$ admet un {\it quotient représentable par un espace algébrique} si le faisceau $U/R$ est un espace algébrique. Nous disons qu'elle admet un {\it quotient représentable} si le faisceau $U/R$ est représentable par un schéma. Nous dirons qu'un groupoïde en espaces algébriques $(U, R, s, t, c)$ sur $B$ admet un {\it quotient représentable} (resp.\ un {\it quotient représentable par un espace algébrique}) si le quotient $U/R$, pour $j = (t, s)$, est représentable (resp.\ est un espace algébrique).
\end{definition}

\noindent
Si le quotient $U/R$ est représenté par $M$ (schéma ou espace algébrique sur $S$), alors $M$ est muni d'un morphisme structural canonique $M \to B$, comme nous l'avons vu ci-dessus.

\medskip\noindent
Le lemme suivant caractérise les $M$ qui représentent le quotient. Il s'applique par exemple si $U \to M$ est plat, de présentation finie et surjectif, et si $R \cong U \times_M U$.

\begin{lemma}
\label{spaces-groupoids-lemma-criterion-quotient-representable}
Dans la situation de la définition \ref{spaces-groupoids-definition-quotient-sheaf}, supposons donnés un espace algébrique $M$ sur $S$ et un morphisme $U \to M$ tels que
\begin{enumerate}
\item le morphisme $U \to M$ coégalise $s$ et $t$,
\item l'application $U \to M$ est une surjection de faisceaux, et
\item l'application induite $(t, s) : R \to U \times_M U$ est une surjection de faisceaux.
\end{enumerate}
Alors $M$ représente le faisceau quotient $U/R$.
\end{lemma}

\begin{proof}
La condition (1) dit que $U \to M$ se factorise par $U/R$. La condition (2) dit que $U/R \to M$ est surjectif comme morphisme de faisceaux. La condition (3) dit qu'il est injectif. Le lemme en résulte.
\end{proof}

\noindent
Le lemme suivant est faux si l'on ne suppose pas que $j$ est une pré-relation d'équivalence, mais seulement une pré-relation.

\begin{lemma}
\label{spaces-groupoids-lemma-quotient-pre-equivalence}
Soit $S$ un schéma. Soit $B$ un espace algébrique sur $S$. Soit $j : R \to U \times_B U$ une pré-relation d'équivalence sur $B$. Pour un schéma $S'$ sur $S$ et $a, b \in U(S')$, les conditions suivantes sont équivalentes :
\begin{enumerate}
\item $a$ et $b$ ont même image dans $(U/R)(S')$, et
\item il existe un recouvrement $\{f_i : S_i \to S'\}$ fppf de $S'$ et des morphismes $r_i : S_i \to R$ tels que $a \circ f_i = s \circ r_i$ et $b \circ f_i = t \circ r_i$.
\end{enumerate}
En d'autres termes, dans ce cas, le morphisme de faisceaux
$$
R \longrightarrow U \times_{U/R} U
$$
est surjectif.
\end{lemma}

\begin{proof}
Omis. Indication : l'argument fonctionne parce que, dans ce cas, le préfaisceau (\ref{spaces-groupoids-equation-quotient-presheaf}) est réellement donné par $T \mapsto U(T)/j(R(T))$, puisque $j(R(T)) \subset U(T) \times U(T)$ est une relation d'équivalence ; voir la définition \ref{spaces-groupoids-definition-equivalence-relation}.
\end{proof}

\begin{lemma}
\label{spaces-groupoids-lemma-quotient-pre-equivalence-relation-restrict}
Soit $S$ un schéma. Soit $B$ un espace algébrique sur $S$. Soit $j : R \to U \times_B U$ une pré-relation d'équivalence sur $B$ et soit $g : U' \to U$ un morphisme d'espaces algébriques sur $B$. Soit $j' : R' \to U' \times_B U'$ la restriction de $j$ à $U'$. Le morphisme de faisceaux quotients
$$
U'/R' \longrightarrow U/R
$$
est injectif. Si $U' \to U$ est surjectif comme morphisme de faisceaux, par exemple si $\{g : U' \to U\}$ est un recouvrement fppf (voir Topologies sur les espaces, définition \ref{spaces-topologies-definition-fppf-covering}), alors $U'/R' \to U/R$ est un isomorphisme de faisceaux.
\end{lemma}

\begin{proof}
Supposons que $\xi, \xi' \in (U'/R')(S')$ aient même image dans $U/R$. Il existe alors un recouvrement fppf $\mathcal{S} = \{S_i \to S'\}$ de $S'$ tel que $\xi|_{S_i}$ et $\xi'|_{S_i}$ soient représentés par $a_i, a_i' \in U'(S_i)$. Par le lemme \ref{spaces-groupoids-lemma-quotient-pre-equivalence} et les axiomes d'un site, nous pouvons, après avoir raffiné par $\mathcal{T}$, supposer qu'il existe des morphismes $r_i : S_i \to R$ tels que $g \circ a_i = s \circ r_i$ et $g \circ a_i' = t \circ r_i$. Comme, par construction, $R' = R \times_{U \times_S U} (U' \times_S U')$, on a $(r_i, (a_i, a_i')) \in R'(S_i)$ ; ainsi, $a_i$ et $a_i'$ définissent la même section de $U'/R'$ sur $S_i$. La condition de faisceau entraîne $\xi = \xi'$.

\medskip\noindent
Si $U' \to U$ est surjectif comme morphisme de faisceaux, alors $U'/R' \to U/R$ est lui aussi surjectif. Enfin, si $\{g : U' \to U\}$ est un recouvrement fppf, le morphisme de faisceaux $U' \to U$ est surjectif ; voir Topologies sur les espaces, lemme \ref{spaces-topologies-lemma-fppf-covering-surjective}.
\end{proof}

\begin{lemma}
\label{spaces-groupoids-lemma-quotient-groupoid-restrict}
Soit $S$ un schéma. Soit $B$ un espace algébrique sur $S$. Soit $(U, R, s, t, c)$ un groupoïde en espaces algébriques sur $B$. Soit $g : U' \to U$ un morphisme d'espaces algébriques sur $B$. Soit $(U', R', s', t', c')$ la restriction de $(U, R, s, t, c)$ à $U'$. Le morphisme de faisceaux quotients
$$
U'/R' \longrightarrow U/R
$$
est injectif. Si le composé
$$
\xymatrix{
U' \times_{g, U, t} R \ar[r]_-{\text{pr}_1} \ar@/^3ex/[rr]^h
& R \ar[r]_s & U
}
$$
est une surjection de faisceaux fppf, alors ce morphisme est bijectif. C'est par exemple le cas si $\{h : U' \times_{g, U, t} R \to U\}$ est un recouvrement $fppf$, si $U' \to U$ est une surjection de faisceaux, ou si $\{g : U' \to U\}$ est un recouvrement pour la topologie fppf.
\end{lemma}

\begin{proof}
L'injectivité résulte des lemmes \ref{spaces-groupoids-lemma-groupoid-pre-equivalence} et \ref{spaces-groupoids-lemma-quotient-pre-equivalence-relation-restrict}. Pour la surjectivité (voir Sites, section \ref{sites-section-sheaves-injective}, pour une caractérisation des morphismes surjectifs de faisceaux), raisonnons comme suit. Soit $T$ un schéma et soit $\sigma \in U/R(T)$. Il existe un recouvrement $\{T_i \to T\}$ tel que $\sigma|_{T_i}$ soit l'image d'un élément $f_i \in U(T_i)$. Nous pouvons donc supposer que $\sigma$ est l'image d'un élément $f \in U(T)$. Par l'hypothèse que $h$ est une surjection de faisceaux, il existe un recouvrement fppf $\{\varphi_i : T_i \to T\}$ et des morphismes $f_i : T_i \to U' \times_{g, U, t} R$ tels que $f \circ \varphi_i = h \circ f_i$. Notons $f'_i = \text{pr}_0 \circ f_i : T_i \to U'$. Alors $f'_i \in U'(T_i)$ a pour image $g \circ f'_i \in U(T_i)$, et $g \circ f'_i \sim_{T_i} h \circ f_i = f \circ \varphi_i$, avec les notations de (\ref{spaces-groupoids-equation-quotient-presheaf}). L'élément de $R(T_i)$ qui donne cette relation est $\text{pr}_1 \circ f_i$. Ainsi, la restriction de $\sigma$ à $T_i$ appartient à l'image de $U'/R'(T_i) \to U/R(T_i)$, comme voulu.

\medskip\noindent
Si $\{h\}$ est un recouvrement fppf, il induit une surjection de faisceaux ; voir Topologies sur les espaces, lemme \ref{spaces-topologies-lemma-fppf-covering-surjective}. Si $U' \to U$ est surjectif, alors $h$ l'est aussi, puisque $s$ admet une section (à savoir l'élément neutre $e$ du groupoïde en schémas).
\end{proof}





\section{Champs quotients}
\label{spaces-groupoids-section-stacks}

\noindent
Dans cette section et les quelques suivantes, nous décrivons une sorte de généralisation de la section \ref{spaces-groupoids-section-quotient-sheaves} ci-dessus et de Groupoïdes, section \ref{groupoids-section-quotient-sheaves}. La différence est que nous allons considérer des champs quotients plutôt que des faisceaux quotients.

\medskip\noindent
Soient un schéma $S$, un espace algébrique $B$ sur $S$ et un groupoïde en espaces algébriques $(U, R, s, t, c)$ sur $B$. À partir de ces données, considérons le foncteur
\begin{equation}
\label{spaces-groupoids-equation-quotient-stack}
\begin{matrix}
(\Sch/S)_{fppf}^{opp} &
\longrightarrow &
\textit{Groupoïdes} \\
S' &
\longmapsto &
(U(S'), R(S'), s, t, c)
\end{matrix}
\end{equation}
D'après Catégories, exemple \ref{categories-example-functor-groupoids}, ce « préfaisceau en groupoïdes » correspond à une catégorie fibrée en groupoïdes sur $(\Sch/S)_{fppf}$. Dans ce chapitre, nous la noterons
$$
[U/_{\!p}R] \to (\Sch/S)_{fppf}
$$
où l'indice ${}_p$ sert à la distinguer du champ quotient.

\begin{definition}
\label{spaces-groupoids-definition-quotient-stack}
Champs quotients. Soit $B \to S$ comme ci-dessus.
\begin{enumerate}
\item Soit $(U, R, s, t, c)$ un groupoïde en espaces algébriques sur $B$. Le {\it champ quotient}
$$
p : [U/R] \longrightarrow (\Sch/S)_{fppf}
$$
de $(U, R, s, t, c)$ est le champ associé (voir Champs, lemme \ref{stacks-lemma-stackify-groupoids}) à la catégorie fibrée en groupoïdes $[U/_{\!p}R]$ sur $(\Sch/S)_{fppf}$ provenant de (\ref{spaces-groupoids-equation-quotient-stack}).
\item Soit $(G, m)$ un espace algébrique en groupes sur $B$. Soit $a : G \times_B X \to X$ une action de $G$ sur un espace algébrique sur $B$. Le {\it champ quotient}
$$
p : [X/G] \longrightarrow (\Sch/S)_{fppf}
$$
est le champ quotient associé au groupoïde en espaces algébriques $(X, G \times_B X, s, t, c)$ sur $B$ du lemme \ref{spaces-groupoids-lemma-groupoid-from-action}.
\end{enumerate}
\end{definition}

\noindent
Ainsi, $[U/R]$ et $[X/G]$ sont des champs en groupoïdes sur $(\Sch/S)_{fppf}$. Ces champs joueront un rôle très important par la suite ; il convient donc de les décrire en détail. Rappelons que, pour un espace algébrique $X$ sur $S$, nous notons $\mathcal{S}_X \to (\Sch/S)_{fppf}$ le champ en ensembles associé au faisceau $X$ ; voir Catégories, lemme \ref{categories-lemma-2-category-fibred-sets}, et Champs, lemme \ref{stacks-lemma-when-stack-in-sets}.

\begin{lemma}
\label{spaces-groupoids-lemma-quotient-stack-arrows}
Supposons $B \to S$ et $(U, R, s, t, c)$ comme dans la définition \ref{spaces-groupoids-definition-quotient-stack} (1). Il existe des $1$-morphismes canoniques $\pi : \mathcal{S}_U \to [U/R]$ et $[U/R] \to \mathcal{S}_B$ de champs en groupoïdes sur $(\Sch/S)_{fppf}$. Le composé $\mathcal{S}_U \to \mathcal{S}_B$ est le $1$-morphisme associé au morphisme structural $U \to B$.
\end{lemma}

\begin{proof}
Dans cette démonstration, notons $[U/_{\!p}R]$ la catégorie fibrée en groupoïdes associée au préfaisceau en groupoïdes (\ref{spaces-groupoids-equation-quotient-stack}). Par construction du champ associé, il existe un $1$-morphisme $[U/_{\!p}R] \to [U/R]$. Le $1$-morphisme $\mathcal{S}_U \to [U/R]$ est simplement le composé $\mathcal{S}_U \to [U/_{\!p}R] \to [U/R]$, où la première flèche associe, au schéma $S'/S$ et au morphisme $x : S' \to U$ sur $S$, l'objet $x \in U(S')$ de la catégorie fibre de $[U/_{\!p}R]$ au-dessus de $S'$.

\medskip\noindent
Pour construire le $1$-morphisme $[U/R] \to \mathcal{S}_B$, il suffit de construire le $1$-morphisme $[U/_{\!p}R] \to \mathcal{S}_B$ ; voir Champs, lemme \ref{stacks-lemma-stackify-groupoids-universal-property}. Sur les objets au-dessus de $S'/S$, nous utilisons simplement l'application
$$
U(S') \longrightarrow B(S')
$$
induite par le morphisme structural $U \to B$. Si $a \in R(S')$ est une « flèche » de source $s(a) \in U(S')$ et de but $t(a) \in U(S')$, alors, puisque $s$ et $t$ sont des morphismes {\it sur} $B$, ces deux éléments ont même image $\overline{a}$ dans $B(S')$. Nous pouvons donc envoyer la flèche $a \in R(S')$ sur le morphisme identité de $\overline{a}$ (ce qui convient, puisque la catégorie fibre $(\mathcal{S}_B)_{S'}$ ne contient que des identités). Nous omettons de vérifier que cette règle est compatible aux images réciproques dans ces catégories fibrées scindées et définit donc le $1$-morphisme voulu $[U/_{\!p}R] \to \mathcal{S}_B$.

\medskip\noindent
Nous omettons la vérification de la dernière assertion.
\end{proof}

\begin{lemma}
\label{spaces-groupoids-lemma-quotient-stack-2-arrow}
Hypothèses et notations comme dans le lemme \ref{spaces-groupoids-lemma-quotient-stack-arrows}. Il existe un $2$-morphisme canonique $\alpha : \pi \circ s \to \pi \circ t$ qui rend le diagramme
$$
\xymatrix{
\mathcal{S}_R \ar[r]_s \ar[d]_t & \mathcal{S}_U \ar[d]^\pi \\
\mathcal{S}_U \ar[r]^-\pi & [U/R]
}
$$
$2$-commutatif.
\end{lemma}

\begin{proof}
Soit $S'$ un schéma sur $S$. Soit $r : S' \to R$ un morphisme sur $S$. Alors $r \in R(S')$ est un isomorphisme entre les objets $s \circ r, t \circ r \in U(S')$. De plus, cette construction est compatible aux images réciproques. Elle donne un $2$-morphisme canonique $\alpha_p : \pi_p \circ s \to \pi_p \circ t$, où $\pi_p : \mathcal{S}_U \to [U/_{\!p}R]$ est celui de la démonstration du lemme \ref{spaces-groupoids-lemma-quotient-stack-arrows}. Ainsi, le diagramme lui-même
$$
\xymatrix{
\mathcal{S}_R \ar[r]_s \ar[d]^t & \mathcal{S}_U \ar[d]^{\pi_p} \\
\mathcal{S}_U \ar[r]^-{\pi_p} & [U/_{\!p}R]
}
$$
est $2$-commutatif. A fortiori, le diagramme du lemme est $2$-commutatif.
\end{proof}

\begin{remark}
\label{spaces-groupoids-remark-fundamental-square}
Dans les chapitres suivants, nous emploierons la notation ambiguë qui consiste à écrire simplement $X$, au lieu de $\mathcal{S}_X$, pour le champ en ensembles associé à $X$. Avec cette notation, le diagramme du lemme \ref{spaces-groupoids-lemma-quotient-stack-2-arrow} devient le diagramme familier
$$
\xymatrix{
R \ar[r]_s \ar[d]_t & U \ar[d]^\pi \\
U \ar[r]^-\pi & [U/R]
}
$$
Dans les sections suivantes, nous montrerons que ce diagramme possède de nombreuses propriétés utiles. En particulier, c'est un $2$-produit fibré (section \ref{spaces-groupoids-section-quotient-stack-2-cartesian}) et il est proche d'être un $2$-coégalisateur de $s$ et de $t$ (section \ref{spaces-groupoids-section-quotient-stacks-2-coequalize}).
\end{remark}




\section{Fonctorialité des champs quotients}
\label{spaces-groupoids-section-functoriality-quotient-stacks}

\noindent
Un morphisme de groupoïdes en espaces algébriques induit un morphisme de champs quotients.

\begin{lemma}
\label{spaces-groupoids-lemma-quotient-stack-functorial}
Soit $S$ un schéma. Soit $B$ un espace algébrique sur $S$. Soit $f : (U, R, s, t, c) \to (U', R', s', t', c')$ un morphisme de groupoïdes en espaces algébriques sur $B$. Alors $f$ induit un $1$-morphisme canonique de champs quotients
$$
[f] : [U/R] \longrightarrow [U'/R'].
$$
\end{lemma}

\begin{proof}
Notons $[U/_{\!p}R]$ et $[U'/_{\!p}R']$ les catégories fibrées en groupoïdes sur le site de base $(\Sch/S)_{fppf}$ associées aux foncteurs (\ref{spaces-groupoids-equation-quotient-stack}). Il est clair que $f$ définit un $1$-morphisme $[U/_{\!p}R] \to [U'/_{\!p}R']$ ; en le composant avec le foncteur de passage au champ associé $[U'/R']$, nous obtenons $[U/_{\!p}R] \to [U'/R']$. La propriété universelle du foncteur de passage au champ associé $[U/_{\!p}R] \to [U/R]$ donne alors $[U/R] \to [U'/R']$ ; voir Champs, lemme \ref{stacks-lemma-stackify-groupoids-universal-property}.
\end{proof}

\noindent
Soient $B \to S$ et $f : (U, R, s, t, c) \to (U', R', s', t', c')$ comme dans le lemme \ref{spaces-groupoids-lemma-quotient-stack-functorial}. Dans cette situation, définissons un troisième groupoïde en espaces algébriques sur $B$ comme suit, en termes de points à valeurs dans un schéma variable $T$ sur $B$ :
\begin{enumerate}
\item $U'' = U \times_{f, U', t'} R'$, de sorte qu'un point à valeurs dans $T$ soit un couple $(u, r')$ tel que $f(u) = t'(r')$,
\item $R'' = R \times_{f \circ s, U', t'} R'$, de sorte qu'un point à valeurs dans $T$ soit un couple $(r, r')$ tel que $f(s(r)) = t'(r')$,
\item $s'' : R'' \to U''$ est donné par $s''(r, r') = (s(r), r')$,
\item $t'' : R'' \to U''$ est donné par $t''(r, r') = (t(r), c'(f(r), r'))$,
\item $c'' : R'' \times_{s'', U'', t''} R'' \to R''$ est donné par $c''((r_1, r'_1), (r_2, r'_2)) = (c(r_1, r_2), r'_2)$.
\end{enumerate}
La formule définissant $c''$ a un sens puisque $s''(r_1, r'_1) = t''(r_2, r'_2)$. Il est clair que $c''$ est associative. L'identité $e''$ est donnée par $e''(u, r) = (e(u), r)$, et l'inverse de $(r, r')$ par $(i(r), c'(f(r), r'))$. Nous obtenons bien ainsi un groupoïde en espaces algébriques sur $B$.

\medskip\noindent
Les applications $U'' \to U$ et $R'' \to R$ définissent clairement un morphisme $g : (U'', R'', s'', t'', c'') \to (U, R, s, t, c)$ de groupoïdes en espaces algébriques sur $B$. De plus, les applications $U'' \to U'$, $(u, r') \mapsto s'(r')$, et $R'' \to U'$, $(r, r') \mapsto s'(r')$, montrent que $(U'', R'', s'', t'', c'')$ est en fait un groupoïde en espaces algébriques sur $U'$.

\begin{lemma}
\label{spaces-groupoids-lemma-cartesian-square-of-morphism}
Notations et hypothèses comme dans le lemme \ref{spaces-groupoids-lemma-quotient-stack-functorial}. Soit $(U'', R'', s'', t'', c'')$ le groupoïde en espaces algébriques sur $B$ construit ci-dessus. Il existe un carré $2$-commutatif
$$
\xymatrix{
[U''/R''] \ar[d] \ar[r]_{[g]} & [U/R] \ar[d]^{[f]} \\
\mathcal{S}_{U'} \ar[r] & [U'/R']
}
$$
qui identifie $[U''/R'']$ au $2$-produit fibré.
\end{lemma}

\begin{proof}
Les morphismes $[f]$ et $[g]$ proviennent du lemme \ref{spaces-groupoids-lemma-quotient-stack-functorial}, et les deux autres du lemme \ref{spaces-groupoids-lemma-quotient-stack-arrows} (ainsi que du fait que $(U'', R'', s'', t'', c'')$ est défini sur $U'$). Pour établir la propriété de $2$-produit fibré, il suffit de démontrer le lemme pour le diagramme
$$
\xymatrix{
[U''/_{\!p}R''] \ar[d] \ar[r]_{[g]} & [U/_{\!p}R] \ar[d]^{[f]} \\
\mathcal{S}_{U'} \ar[r] & [U'/_{\!p}R']
}
$$
de catégories fibrées en groupoïdes ; voir Champs, lemme \ref{stacks-lemma-stackification-fibre-product-categories-fibred-in-groupoids}. Autrement dit, il suffit de montrer qu'un objet du $2$-produit fibré $\mathcal{S}_U \times_{[U'/_{\!p}R']} [U/_{\!p}R]$ au-dessus de $T$ correspond à un point de $U''$ à valeurs dans $T$, et de même pour les morphismes. C'est précisément ainsi que $U''$ et $R''$ ont été construits.

\medskip\noindent
Plus précisément, un objet de $\mathcal{S}_U \times_{[U'/_{\!p}R']} [U/_{\!p}R]$ au-dessus de $T$ est un triplet $(u', u, r')$, où $u'$ est un point de $U'$ à valeurs dans $T$, $u$ un point de $U$ à valeurs dans $T$ et $r'$ un morphisme de $u'$ vers $f(u)$ dans $[U'/R']_T$ ; autrement dit, $r'$ est un point de $R$ à valeurs dans $T$ tel que $s'(r') = u'$ et $t'(r') = f(u)$. On peut oublier $u'$ sans perte d'information ; ces objets correspondent donc bijectivement aux points de $R''$ à valeurs dans $T$.

\medskip\noindent
De même pour les morphismes. Soient $(u'_1, u_1, r'_1)$ et $(u'_2, u_2, r'_2)$ deux objets du produit fibré au-dessus de $T$. Un morphisme de $(u'_2, u_2, r'_2)$ vers $(u'_1, u_1, r'_1)$ est donné par $(1, r)$, où $1 : u'_1 \to u'_2$ signifie simplement que $u'_1 = u'_2$ (car $\mathcal{S}_U$ est fibrée en ensembles), et où $r$ est un point de $R$ à valeurs dans $T$ tel que $s(r) = u_2$, $t(r) = u_1$ et $c'(f(r), r'_2) = r'_1$. Par conséquent, la flèche
$$
(1, r) : (u'_2, u_2, r'_2) \to (u'_1, u_1, r'_1)
$$
est entièrement déterminée par le couple $(r, r'_2)$. Le foncteur des flèches est donc représenté par $R''$ ; de plus, les morphismes $s''$, $t''$ et $c''$ correspondent respectivement à la source, au but et à la composition dans le $2$-produit fibré $\mathcal{S}_U \times_{[U'/_{\!p}R']} [U/_{\!p}R]$.
\end{proof}







\section{Le carré 2-cartésien d'un champ quotient}
\label{spaces-groupoids-section-quotient-stack-2-cartesian}

\noindent
Dans cette section, nous calculons les faisceaux $\mathit{Isom}$ d'un champ quotient et nous en déduisons que le diagramme qui le définit est un $2$-produit fibré.

\begin{lemma}
\label{spaces-groupoids-lemma-quotient-stack-morphisms}
Supposons que $B \to S$, $(U, R, s, t, c)$ et $\pi : \mathcal{S}_U \to [U/R]$ soient comme dans le lemme \ref{spaces-groupoids-lemma-quotient-stack-arrows}. Soit $S'$ un schéma sur $S$. Soient $x, y \in \Ob([U/R]_{S'})$ des objets du champ quotient sur $S'$. Si $x = \pi(x')$ et $y = \pi(y')$ pour certains morphismes $x', y' : S' \to U$, alors
$$
\mathit{Isom}(x, y) = S' \times_{(y', x'), U \times_S U} R
$$
comme faisceaux sur $S'$.
\end{lemma}

\begin{proof}
Soit $[U/_{\!p}R]$ la catégorie fibrée en groupoïdes associée au préfaisceau en groupoïdes (\ref{spaces-groupoids-equation-quotient-stack}), comme dans la démonstration du lemme \ref{spaces-groupoids-lemma-quotient-stack-arrows}. Par construction, le faisceau $\mathit{Isom}(x, y)$ est le faisceau associé au préfaisceau $\mathit{Isom}(x', y')$. D'autre part, par définition des morphismes dans $[U/_{\!p}R]$, nous avons
$$
\mathit{Isom}(x', y') = S' \times_{(y', x'), U \times_S U} R
$$
et le membre de droite est un espace algébrique, donc un faisceau.
\end{proof}

\begin{lemma}
\label{spaces-groupoids-lemma-quotient-stack-2-cartesian}
Supposons que $B \to S$, $(U, R, s, t, c)$ et $\pi : \mathcal{S}_U \to [U/R]$ soient comme dans le lemme \ref{spaces-groupoids-lemma-quotient-stack-arrows}. Le carré $2$-commutatif
$$
\xymatrix{
\mathcal{S}_R \ar[r]_s \ar[d]_t & \mathcal{S}_U \ar[d]^\pi \\
\mathcal{S}_U \ar[r]^-\pi & [U/R]
}
$$
du lemme \ref{spaces-groupoids-lemma-quotient-stack-2-arrow} est un $2$-produit fibré de champs en groupoïdes sur $(\Sch/S)_{fppf}$.
\end{lemma}

\begin{proof}
D'après Champs, lemme \ref{stacks-lemma-2-product-stacks-in-groupoids}, l'énoncé a un sens. Ce lemme montre aussi qu'il nous faut prouver que le foncteur
$$
\mathcal{S}_R \longrightarrow \mathcal{S}_U \times_{[U/R]} \mathcal{S}_U
$$
qui envoie $r : T \to R$ sur $(T, t(r), s(r), \alpha(r))$ est une équivalence, le membre de droite étant le $2$-produit fibré décrit dans Catégories, lemme \ref{categories-lemma-2-product-categories-over-C}. Une fois les définitions explicitées, c'est exactement le contenu du lemme \ref{spaces-groupoids-lemma-quotient-stack-morphisms}. (Autre démonstration : expliciter le lemme \ref{spaces-groupoids-lemma-cartesian-square-of-morphism} dans cette situation donne également le résultat.)
\end{proof}

\begin{lemma}
\label{spaces-groupoids-lemma-quotient-stack-isom}
Supposons que $B \to S$ et $(U, R, s, t, c)$ soient comme dans la définition \ref{spaces-groupoids-definition-quotient-stack} (1). Pour tout schéma $T$ sur $S$ et tous objets $x, y$ de $[U/R]$ sur $T$, le faisceau $\mathit{Isom}(x, y)$ sur $(\Sch/T)_{fppf}$ a la propriété suivante : il existe un recouvrement fppf $\{T_i \to T\}_{i \in I}$ tel que $\mathit{Isom}(x, y)|_{(\Sch/T_i)_{fppf}}$ soit représentable par un espace algébrique.
\end{lemma}

\begin{proof}
Cela découle immédiatement du lemme \ref{spaces-groupoids-lemma-quotient-stack-morphisms} et du fait que, localement pour la topologie fppf, $x$ et $y$ proviennent tous deux d'objets de $\mathcal{S}_U$, par construction du champ quotient.
\end{proof}














\section{La propriété de 2-coégalisateur d'un champ quotient}
\label{spaces-groupoids-section-quotient-stacks-2-coequalize}

\noindent
La composition d'un groupoïde conduit à une condition de cocycle pour le $2$-morphisme canonique du lemme précédent. Pour l'énoncer précisément, nous utiliserons les notations introduites dans Catégories, sections \ref{categories-section-formal-cat-cat} et \ref{categories-section-2-categories}.

\begin{lemma}
\label{spaces-groupoids-lemma-quotient-stack-cocycle}
Hypothèses et notations comme dans les lemmes \ref{spaces-groupoids-lemma-quotient-stack-arrows} et \ref{spaces-groupoids-lemma-quotient-stack-2-arrow}. Le composé vertical de
$$
\xymatrix@C=15pc{
\mathcal{S}_{R \times_{s, U, t} R}
\ruppertwocell^{\pi \circ s \circ \text{pr}_1 = \pi \circ s \circ c}{\ \ \ \ \ \ \alpha \star \text{id}_{\text{pr}_1}}
\ar[r]_(.3){\pi \circ t \circ \text{pr}_1 = \pi \circ s \circ \text{pr}_0}
\rlowertwocell_{\pi \circ t \circ \text{pr}_0 = \pi \circ t \circ c}{\ \ \ \ \ \ \alpha \star \text{id}_{\text{pr}_0}}
&
[U/R]
}
$$
est le $2$-morphisme $\alpha \star \text{id}_c$. Autrement dit, $\alpha \star \text{id}_c = (\alpha \star \text{id}_{\text{pr}_0}) \circ (\alpha \star \text{id}_{\text{pr}_1}) $.
\end{lemma}

\begin{proof}
Nous faisons deux remarques :
\begin{enumerate}
\item La formule $\alpha \star \text{id}_c = (\alpha \star \text{id}_{\text{pr}_0}) \circ (\alpha \star \text{id}_{\text{pr}_1})$ n'a de sens qu'en tenant compte des {\it égalités} $\pi \circ s \circ \text{pr}_1 = \pi \circ s \circ c$, $\pi \circ t \circ \text{pr}_1 = \pi \circ s \circ \text{pr}_0$ et $\pi \circ t \circ \text{pr}_0 = \pi \circ t \circ c$. La deuxième rend possible la composition verticale $\circ$, et les deux autres garantissent que les deux membres de la formule sont des $2$-morphismes de même source et de même but.
\item Le lemme vaut parce que la composition dans la catégorie fibrée en groupoïdes $[U/_{\!p}R]$ associée au préfaisceau en groupoïdes (\ref{spaces-groupoids-equation-quotient-stack}) provient de la loi de composition $c : R \times_{s, U, t} R \to R$.
\end{enumerate}
Nous omettons la démonstration du lemme.
\end{proof}

\noindent
Remarquons que, dans la situation du lemme, les égalités $s \circ \text{pr}_1 = s \circ c$, $t \circ \text{pr}_1 = s \circ \text{pr}_0$ et $t \circ \text{pr}_0 = t \circ c$ valent déjà avant composition par $\pi$. La formule du lemme suivant a donc un sens pour les mêmes raisons que celle du lemme précédent.

\begin{lemma}
\label{spaces-groupoids-lemma-quotient-stack-2-coequalizer}
Hypothèses et notations comme dans les lemmes \ref{spaces-groupoids-lemma-quotient-stack-arrows} et \ref{spaces-groupoids-lemma-quotient-stack-2-arrow}. Le diagramme $2$-commutatif du lemme \ref{spaces-groupoids-lemma-quotient-stack-2-arrow} est un $2$-coégalisateur au sens suivant : étant donnés
\begin{enumerate}
\item un champ en groupoïdes $\mathcal{X}$ sur $(\Sch/S)_{fppf}$,
\item un $1$-morphisme $f : \mathcal{S}_U \to \mathcal{X}$, et
\item une $2$-flèche $\beta : f \circ s \to f \circ t$
\end{enumerate}
tels que
$$
\beta \star \text{id}_c
=
(\beta \star \text{id}_{\text{pr}_0})
\circ
(\beta \star \text{id}_{\text{pr}_1})
$$
alors il existe un $1$-morphisme $[U/R] \to \mathcal{X}$ qui rend le diagramme
$$
\xymatrix{
\mathcal{S}_R \ar[r]_s \ar[d]^t & \mathcal{S}_U \ar[d] \ar[ddr]^f \\
\mathcal{S}_U \ar[r] \ar[rrd]_f & [U/R] \ar[rd] \\
& & \mathcal{X}
}
$$
$2$-commutatif.
\end{lemma}

\begin{proof}
Supposons $\mathcal{X}$, $f$ et $\beta$ donnés comme dans le lemme. D'après Champs, lemme \ref{stacks-lemma-stackify-groupoids-universal-property}, il suffit de construire un $1$-morphisme $g : [U/_{\!p}R] \to \mathcal{X}$. Remarquons d'abord que le $1$-morphisme $\mathcal{S}_U \to [U/_{\!p}R]$ est bijectif sur les objets. Nous pouvons donc poser $g(x) = f(x)$ sur les objets, pour $x \in \Ob(\mathcal{S}_U) = \Ob([U/_{\!p}R])$. Un morphisme $\varphi : x \to y$ de $[U/_{\!p}R]$ provient d'un diagramme commutatif
$$
\xymatrix{
S_2 \ar[dd]_h \ar[r]_x \ar[dr]_\varphi & U \\
& R \ar[u]_s \ar[d]^t \\
S_1 \ar[r]^y & U.
}
$$
Nous pouvons donc définir $g(\varphi)$ comme le composé
$$
\xymatrix{
f(x) \ar@{=}[r] \ar[rrrrrd] &
f(s \circ \varphi) \ar@{=}[r] &
(f \circ s)(\varphi) \ar[r]^\beta &
(f \circ t)(\varphi) \ar@{=}[r] &
f(t \circ \varphi) \ar@{=}[r] &
f(y \circ h) \ar[d] \\
& & & & & f(y).
}
$$
La flèche verticale s'obtient en appliquant le foncteur $f$ au morphisme canonique $y \circ h \to y$ dans $\mathcal{S}_U$ (c'est-à-dire au morphisme fortement cartésien relevant $h$ et de but $y$). Vérifions que $f$ ainsi défini est compatible à la composition, au moins dans les catégories fibres. Soit $S'$ un schéma sur $S$ et soit $a : S' \to R \times_{s, U, t} R$ un morphisme. Posons $x = s \circ \text{pr}_1 \circ a = s \circ c \circ a$, $y = t \circ \text{pr}_1 \circ a = s \circ \text{pr}_0 \circ a$ et $z = t \circ \text{pr}_0 \circ a = t \circ \text{pr}_0 \circ c$. Nous obtenons un diagramme commutatif
$$
\xymatrix{
x \ar[rr]_{c \circ a} \ar[rd]_{\text{pr}_1 \circ a} & & z \\
& y \ar[ru]_{\text{pr}_0 \circ a}
}
$$
dans la catégorie fibre $[U/_{\!p}R]_{S'}$. De plus, tout triangle commutatif de cette catégorie fibre est de cette forme. D'après les définitions précédentes, $f$ envoie ce triangle sur un diagramme commutatif si et seulement si le diagramme
$$
\xymatrix{
& (f \circ s)(c \circ a) \ar[r]_-{\beta} &
(f \circ t)(c \circ a) \ar@{=}[rd] & \\
(f \circ s)(\text{pr}_1 \circ a) \ar[rd]^\beta \ar@{=}[ru] & & &
(f \circ t)(\text{pr}_0 \circ a) \\
& (f \circ t)(\text{pr}_1 \circ a) \ar@{=}[r] &
(f \circ s)(\text{pr}_0 \circ a) \ar[ru]^\beta
}
$$
est commutatif, ce qui est exactement la condition exprimée par la formule du lemme. Nous omettons de vérifier que $f$ envoie les identités sur les identités et qu'il est compatible à la composition pour des morphismes quelconques.
\end{proof}












\section{Description explicite des champs quotients}
\label{spaces-groupoids-section-explicit-quotient-stacks}

\noindent
Pour formuler le résultat, introduisons quelques notations. Supposons que $B \to S$ et $(U, R, s, t, c)$ soient comme dans la définition \ref{spaces-groupoids-definition-quotient-stack} (1). Soit $T$ un schéma sur $S$. Soit $\mathcal{T} = \{T_i \to T\}_{i \in I}$ un recouvrement fppf. Une {\it donnée de descente pour $[U/R]$} relativement à $\mathcal{T}$ est un système $(u_i, r_{ij})$ où
\begin{enumerate}
\item pour tout $i$, $u_i : T_i \to U$ est un morphisme, et
\item pour tous $i, j$, $r_{ij} : T_i \times_T T_j \to R$ est un morphisme,
\end{enumerate}
tels que
\begin{enumerate}
\item [(a)] comme morphismes $T_i \times_T T_j \to U$, on a
$$
s \circ r_{ij} = u_i \circ \text{pr}_0
\quad\text{et}\quad
t \circ r_{ij} = u_j \circ \text{pr}_1,
$$
\item [(b)] comme morphismes $T_i \times_T T_j \times_T T_k \to R$, on a
$$
c \circ (r_{jk} \circ \text{pr}_{12}, r_{ij} \circ \text{pr}_{01})
= r_{ik} \circ \text{pr}_{02}.
$$
\end{enumerate}
Un {\it morphisme $(u_i, r_{ij}) \to (u'_i, r'_{ij})$ entre deux données de descente pour $[U/R]$} relativement au même recouvrement $\mathcal{T}$ est une famille $(r_i : T_i \to R)$ telle que
\begin{enumerate}
\item [] $(\alpha)$ comme morphismes $T_i \to U$, on ait
$$
u_i = s \circ r_i
\quad\text{et}\quad
u'_i = t \circ r_i
$$
\item [] $(\beta)$ comme morphismes $T_i \times_T T_j \to R$, on ait
$$
c \circ (r'_{ij}, r_i \circ \text{pr}_0)
=
c \circ (r_j \circ \text{pr}_1, r_{ij}).
$$
\end{enumerate}
Les morphismes de données de descente relativement à un recouvrement fixé sont munis d'une loi de composition naturelle ; on obtient ainsi une catégorie de données de descente. Cette catégorie est un groupoïde. Enfin, si $\mathcal{T}' = \{T'_j \to T\}_{j \in J}$ est un second recouvrement fppf qui raffine $\mathcal{T}$, on dispose d'une notion d'image réciproque des données de descente. Dans ce cas, elle se décrit très simplement. Si $\alpha : J \to I$ et les $\varphi_j : T'_j \to T_{\alpha(i)}$ définissent le morphisme de recouvrements, alors l'image réciproque de la donnée de descente $(u_i, r_{ii'})$ est
$$
(u_{\alpha(i)} \circ \varphi_j,
r_{\alpha(j)\alpha(j')} \circ \varphi_j \times \varphi_{j'}).
$$
L'image réciproque ainsi définie donne un foncteur de la catégorie des données de descente relativement à $\mathcal{T}$ vers celle des données de descente relativement à $\mathcal{T}'$.

\begin{lemma}
\label{spaces-groupoids-lemma-quotient-stack-objects}
Supposons que $B \to S$ et $(U, R, s, t, c)$ soient comme dans la définition \ref{spaces-groupoids-definition-quotient-stack} (1). Soit $\pi : \mathcal{S}_U \to [U/R]$ comme dans le lemme \ref{spaces-groupoids-lemma-quotient-stack-arrows}. Soit $T$ un schéma sur $S$.
\begin{enumerate}
\item pour tout objet $x$ de la catégorie fibre $[U/R]_T$, il existe un recouvrement fppf $\{f_i : T_i \to T\}_{i \in I}$ tel que $f_i^*x \cong \pi(u_i)$ pour certains $u_i \in U(T_i)$,
\item le composé des isomorphismes
$$
\pi(u_i \circ \text{pr}_0)
=
\text{pr}_0^*\pi(u_i)
\cong
\text{pr}_0^*f_i^*x
\cong
\text{pr}_1^*f_j^*x
\cong
\text{pr}_1^*\pi(u_j)
=
\pi(u_j \circ \text{pr}_1)
$$
est de la forme $\pi(r_{ij})$ pour certains morphismes $r_{ij} : T_i \times_T T_j \to R$,
\item le système $(u_i, r_{ij})$ est une donnée de descente pour $[U/R]$ au sens défini ci-dessus,
\item toute donnée de descente $(u_i, r_{ij})$ pour $[U/R]$ s'obtient de cette façon,
\item si $x$ correspond à $(u_i, r_{ij})$ comme ci-dessus et si $y \in \Ob([U/R]_T)$ correspond à $(u'_i, r'_{ij})$, alors il existe une bijection canonique
$$
\Mor_{[U/R]_T}(x, y)
\longleftrightarrow
\left\{
\begin{matrix}
\text{morphismes }(u_i, r_{ij}) \to (u'_i, r'_{ij})\\
\text{de données de descente pour }[U/R]\text{}
\end{matrix}
\right\}
$$
\item cette correspondance est compatible aux raffinements des recouvrements fppf.
\end{enumerate}
\end{lemma}

\begin{proof}
L'assertion (1) fait partie de la construction du champ associé. L'assertion (2) résulte du lemme \ref{spaces-groupoids-lemma-quotient-stack-morphisms}. Nous omettons la vérification de (3). L'assertion (4) traduit le fait que, dans un champ, toute donnée de descente est effective. Nous omettons les vérifications de (5) et (6).
\end{proof}








\section{Restriction et champs quotients}
\label{spaces-groupoids-section-quotient-stack-restrict}

\noindent
Dans cette section, nous étudions l'effet d'une restriction sur le champ quotient.

\begin{lemma}
\label{spaces-groupoids-lemma-quotient-stack-restrict}
Notations et hypothèses comme dans le lemme \ref{spaces-groupoids-lemma-quotient-stack-functorial}. Le morphisme de champs quotients
$$
[f] : [U/R] \longrightarrow [U'/R']
$$
est pleinement fidèle si et seulement si $R$ est la restriction de $R'$ suivant le morphisme $f : U \to U'$.
\end{lemma}

\begin{proof}
Soient $x, y$ des objets de $[U/R]$ sur un schéma $T/S$. Soient $x', y'$ leurs images dans la catégorie $[U'/R']_T$. Le foncteur $[f]$ est pleinement fidèle si et seulement si le morphisme de faisceaux
$$
\mathit{Isom}(x, y) \longrightarrow \mathit{Isom}(x', y')
$$
est un isomorphisme pour tous $T, x, y$. Cette propriété se vérifie localement sur $T$ pour la topologie fppf. D'après le lemme \ref{spaces-groupoids-lemma-quotient-stack-objects}, nous pouvons donc supposer que $x$ et $y$ proviennent de $a, b \in U(T)$. Alors $x'$ et $y'$ correspondent à $f \circ a$ et $f \circ b$. Par le lemme \ref{spaces-groupoids-lemma-quotient-stack-morphisms}, le morphisme de faisceaux ci-dessus devient
$$
T \times_{(a, b), U \times_B U} R
\longrightarrow
T \times_{f \circ a, f \circ b, U' \times_B U'} R'.
$$
C'est un isomorphisme si $R$ est la restriction, car alors $R = (U \times_B U) \times_{U' \times_B U'} R'$ ; voir le lemme \ref{spaces-groupoids-lemma-restrict-groupoid-relation} et sa démonstration. Réciproquement, si le dernier morphisme affiché est un isomorphisme pour tous $T, a, b$, il s'ensuit que $R = (U \times_B U) \times_{U' \times_B U'} R'$, c'est-à-dire que $R$ est la restriction de $R'$.
\end{proof}

\begin{lemma}
\label{spaces-groupoids-lemma-quotient-stack-restrict-equivalence}
Notations et hypothèses comme dans le lemme
\ref{spaces-groupoids-lemma-quotient-stack-functorial}. Le morphisme de champs quotients
$$
[f] : [U/R] \longrightarrow [U'/R']
$$
est une équivalence si et seulement si les conditions suivantes sont satisfaites :
\begin{enumerate}
\item $(U, R, s, t, c)$ est la restriction de $(U', R', s', t', c')$
suivant $f : U \to U'$, et
\item le morphisme
$$
\xymatrix{
U \times_{f, U', t'} R' \ar[r]_-{\text{pr}_1} \ar@/^3ex/[rr]^h
& R' \ar[r]_{s'} & U'
}
$$
est une surjection de faisceaux.
\end{enumerate}
La condition (2) est par exemple satisfaite si
$\{h : U \times_{f, U', t'} R' \to U'\}$ est un recouvrement fppf,
si $f : U \to U'$ est une surjection de faisceaux, ou si
$\{f : U \to U'\}$ est un recouvrement fppf.
\end{lemma}

\begin{proof}
Nous savons déjà, par le lemme \ref{spaces-groupoids-lemma-quotient-stack-restrict}, que la
condition (1) équivaut à la pleine fidélité. Nous pouvons donc supposer que
(1) est satisfaite et que $[f]$ est pleinement fidèle. Il s'agit alors de
montrer que $[f]$ est une équivalence si et seulement si (2) est satisfaite.
Nous pouvons appliquer Champs, lemme
\ref{stacks-lemma-characterize-essentially-surjective-when-ff}, qui
caractérise les équivalences.

\medskip\noindent
Supposons (2). Nous allons appliquer Champs, lemme
\ref{stacks-lemma-characterize-essentially-surjective-when-ff}, pour montrer
que $[f]$ est une équivalence. Soient $T$ un schéma et
$x' \in \Ob([U'/R']_T)$. Il existe un recouvrement
$\{g_i : T_i \to T\}$ tel que $g_i^*x'$ soit l'image d'un élément
$a'_i \in U'(T_i)$ ; voir le lemme \ref{spaces-groupoids-lemma-quotient-stack-objects}.
Nous pouvons donc supposer que $x'$ est l'image de $a' \in U'(T)$. Comme
$h$ est une surjection de faisceaux, il existe un recouvrement fppf
$\{\varphi_i : T_i \to T\}$ et des morphismes
$b_i : T_i \to U \times_{g, U', t'} R'$ tels que
$a' \circ \varphi_i = h \circ b_i$. Posons
$a_i = \text{pr}_0 \circ b_i : T_i \to U$. Alors $a_i \in U(T_i)$ a pour
image $f \circ a_i \in U'(T_i)$, et
$f \circ a_i \cong_{T_i} h \circ b_i = a' \circ \varphi_i$, où
$\cong_{T_i}$ désigne un isomorphisme dans la catégorie fibre
$[U'/R']_{T_i}$. L'élément de $R'(T_i)$ qui donne cet isomorphisme est
$\text{pr}_1 \circ b_i$. Ainsi, la restriction de $x$ à $T_i$ appartient à
l'image essentielle du foncteur
$[U/R]_{T_i} \to [U'/R']_{T_i}$, comme voulu.

\medskip\noindent
Supposons que $[f]$ soit une équivalence. Soit
$\xi' \in [U'/R']_{U'}$ l'objet correspondant au morphisme identité de
$U'$. En appliquant Champs, lemme
\ref{stacks-lemma-characterize-essentially-surjective-when-ff}, nous voyons
qu'il existe un recouvrement fppf
$\mathcal{U}' = \{g'_i : U'_i \to U'\}$ tel que
$(g'_i)^*\xi' \cong [f](\xi_i)$ pour un objet $\xi_i$ de
$[U/R]_{U'_i}$. Après avoir raffiné le recouvrement $\mathcal{U}'$ (au moyen
du lemme \ref{spaces-groupoids-lemma-quotient-stack-objects}), nous pouvons supposer que
$\xi_i$ provient d'un morphisme $a_i : U'_i \to U$. L'isomorphisme
$[f](\xi_i) \cong (g'_i)^*\xi'$ signifie qu'après avoir éventuellement
raffiné encore une fois $\mathcal{U}'$, il existe des morphismes
$r'_i : U'_i \to R'$ tels que $t' \circ r'_i = f \circ a_i$ et
$s' \circ r'_i = \text{id}_{U'} \circ g'_i$. La situation est représentée
par le diagramme
$$
\xymatrix{
U \ar[d]^f & & U'_i \ar[ll]^{a_i} \ar[ld]_{r'_i} \ar[d]^{g'_i} \\
U' & R' \ar[l]_{t'} \ar[r]^{s'} & U'
}
$$
Les $(a_i, r'_i) : U'_i \to U \times_{g, U', t'} R'$ sont donc des
morphismes tels que $h \circ (a_i, r'_i) = g'_i$. Par conséquent, le
recouvrement fppf $\mathcal{U}'$ raffine
$\{h : U \times_{g, U', t'} R' \to U'\}$, ce qui signifie que $h$ induit
une surjection de faisceaux ; voir Topologies sur les espaces, lemme
\ref{spaces-topologies-lemma-fppf-covering-surjective}.

\medskip\noindent
Si $\{h\}$ est un recouvrement fppf, alors il induit une surjection de
faisceaux ; voir Topologies sur les espaces, lemme
\ref{spaces-topologies-lemma-fppf-covering-surjective}. Si $U' \to U$ est
surjectif, alors $h$ l'est aussi, car $s$ possède une section (à savoir
l'élément neutre $e$ du groupoïde en espaces algébriques).
\end{proof}

\begin{lemma}
\label{spaces-groupoids-lemma-criterion-fibre-product}
Notations et hypothèses comme dans le lemme
\ref{spaces-groupoids-lemma-quotient-stack-functorial}. Supposons que le carré
$$
\xymatrix{
R \ar[d]_s \ar[r]_f & R' \ar[d]^{s'} \\
U \ar[r]^f & U'
}
$$
soit cartésien. Alors le carré
$$
\xymatrix{
\mathcal{S}_U \ar[d] \ar[r] & [U/R] \ar[d]^{[f]} \\
\mathcal{S}_{U'} \ar[r] & [U'/R']
}
$$
est $2$-cartésien.
\end{lemma}

\begin{proof}
En appliquant les isomorphismes d'inversion $i : R \to R$ et
$i' : R' \to R'$ au (premier) carré cartésien de l'énoncé, nous voyons que
$$
\xymatrix{
R \ar[d]_t \ar[r]_f & R' \ar[d]^{t'} \\
U \ar[r]^f & U'
}
$$
est lui aussi cartésien. Par le lemme
\ref{spaces-groupoids-lemma-cartesian-square-of-morphism}, nous disposons d'un carré de
$2$-produit fibré
$$
\xymatrix{
[U''/R''] \ar[d] \ar[r] & [U/R] \ar[d] \\
\mathcal{S}_{U'} \ar[r] & [U'/R']
}
$$
où $U'' = U \times_{f, U', t'} R'$ et
$R'' = R \times_{f \circ s, U', t'} R'$. Il résulte de ce qui précède que
$(t, f) : R \to U''$ est un isomorphisme et que
$$
R'' =
R \times_{f \circ s, U', t'} R' =
R \times_{s, U} U \times_{f, U', t'} R' =
R \times_{s, U, t} \times R.
$$
Explicitement, l'isomorphisme $R \times_{s, U, t} R \to R''$ est donné par
la règle $(r_0, r_1) \mapsto (r_0, f(r_1))$. De plus, $s'', t'', c''$
s'identifient aux applications
$$
R \times_{s, U, t} R \to R,
\quad
s''(r_0, r_1) = r_1, \quad t''(r_0, r_1) = c(r_0, r_1)
$$
et
$$
\begin{matrix}
c'' : &
(R \times_{s, U, t} R) \times_{s'', R, t''} (R \times_{s, U, t} R)
&
\longrightarrow
&
R \times_{s, U, t} R, \\
&
((r_0, r_1), (r_2, r_3)) &
\longmapsto & (c(r_0, r_2), r_3).
\end{matrix}
$$
Précomposition avec l'isomorphisme
$$
R \times_{s, U, s} R \longrightarrow R \times_{s, U, t} R,
\quad
(r_0, r_1) \longmapsto (c(r_0, i(r_1)), r_1)
$$
En précomposant, nous voyons que $t''$ et $s''$ deviennent
$\text{pr}_0$ et $\text{pr}_1$, tandis que $c''$ devient
$\text{pr}_{02} : R \times_{s, U, s} R \times_{s, U, s} R \to R \times_{s, U, s} R$.
Il existe donc un isomorphisme
$[U''/R''] \cong [R/R \times_{s, U, s} R]$, où le groupoïde en espaces
algébriques $(R, R \times_{s, U, s} R, s'', t'', c'')$ est la restriction,
suivant $s : R \to U$, du groupoïde trivial
$(U, U, \text{id}, \text{id}, \text{id})$. Puisque $s : R \to U$ est une
surjection de faisceaux fppf (il possède en effet un inverse à droite), le
morphisme
$$
[U''/R''] \cong [R/R \times_{s, U, s} R]
\longrightarrow
[U/U] = \mathcal{S}_U
$$
est une équivalence par le lemme
\ref{spaces-groupoids-lemma-quotient-stack-restrict-equivalence}. Cela démontre le lemme.
\end{proof}





\section{Inertie et champs quotients}
\label{spaces-groupoids-section-inertia}

\noindent
Le champ d'inertie (relatif) d'un champ en groupoïdes est défini dans Champs,
section \ref{stacks-section-the-inertia-stack}. Sa construction proprement
dite dans le cadre des catégories fibrées, ainsi que certaines de ses
propriétés, se trouvent dans Catégories, section
\ref{categories-section-inertia}.

\begin{lemma}
\label{spaces-groupoids-lemma-presentation-inertia}
Supposons que $B \to S$ et $(U, R, s, t, c)$ soient comme dans la
définition \ref{spaces-groupoids-definition-quotient-stack} (1). Soit $G/U$ l'espace
algébrique en groupes stabilisateur du groupoïde
$(U, R, s, t, c, e, i)$ ; voir la définition
\ref{spaces-groupoids-definition-stabilizer-groupoid}. Posons $R' = R \times_{s, U} G$ et
définissons
\begin{enumerate}
\item $s' : R' \to G$, $(r, g) \mapsto g$,
\item $t' : R' \to G$, $(r, g) \mapsto c(r, c(g, i(r)))$,
\item $c' : R' \times_{s', G, t'} R' \to R'$, $((r_1, g_1), (r_2, g_2) \mapsto (c(r_1, r_2), g_1)$.
\end{enumerate}
Alors $(G, R', s', t', c')$ est un groupoïde en espaces algébriques sur
$B$ et
$$
\mathcal{I}_{[U/R]} = [G/ R'].
$$
Autrement dit, le champ quotient associé est le champ d'inertie de $[U/R]$.
\end{lemma}

\begin{proof}
Par Champs, lemme \ref{stacks-lemma-stackification-inertia}, il suffit de
démontrer que $\mathcal{I}_{[U/_{\!p}R]} = [G/_{\!p} R']$. Soit $T$ un
schéma sur $S$. Rappelons qu'un objet de la catégorie fibrée d'inertie de
$[U/_{\!p}R]$ sur $T$ est une paire $(x, g)$, où $x$ est un objet de
$[U/_{\!\!p}R]$ sur $T$ et $g$ un automorphisme de $x$ dans sa catégorie
fibre sur $T$. Autrement dit, $x : T \to U$ et $g : T \to R$ vérifient
$x = s \circ g = t \circ g$. C'est exactement dire que $g : T \to G$.
Un morphisme, dans la catégorie fibrée d'inertie, de $(x, g)$ vers $(y, h)$
sur $T$ est donné par $r : T \to R$ tel que $s(r) = x$, $t(r) = y$ et
$c(r, g) = c(h, r)$ ; voir le diagramme commutatif de Catégories, lemme
\ref{categories-lemma-inertia-fibred-category}. Sous forme d'une formule,
$$
h = c(r, c(g, i(r))) = c(c(r, g), i(r)).
$$
La notation $s(r)$, etc., abrège $s \circ r$, etc. La composée de
$r_1 : (x_2, g_2) \to (x_1, g_1)$ et de
$r_2 : (x_1, g_1) \to (x_2, g_2)$ est
$c(r_1, r_2) : (x_1, g_1) \to (x_3, g_3)$.

\medskip\noindent
Remarquons que, dans ce qui précède, nous aurions pu noter simplement $g$,
au lieu de $(x, g)$, un objet de $\mathcal{I}_{[U/_{\!p}R]}$ sur $T$, car
$x$ est l'image de $g$ par le morphisme structural $G \to U$. Les morphismes
$g \to h$ de $\mathcal{I}_{[U/_{\!p}R]}$ sur $T$ correspondent alors
exactement aux morphismes $r' : T \to R'$ tels que $s'(r') = g$ et
$t'(r') = h$. En outre, la composition correspond à la règle décrite en
(3). Le lemme est ainsi démontré.
\end{proof}

\begin{lemma}
\label{spaces-groupoids-lemma-2-cartesian-inertia}
Supposons que $B \to S$ et $(U, R, s, t, c)$ soient comme dans la
définition \ref{spaces-groupoids-definition-quotient-stack} (1). Soit $G/U$ l'espace
algébrique en groupes stabilisateur du groupoïde
$(U, R, s, t, c, e, i)$ ; voir la définition
\ref{spaces-groupoids-definition-stabilizer-groupoid}. Il existe un diagramme canonique
$2$-cartésien
$$
\xymatrix{
\mathcal{S}_G \ar[r] \ar[d] & \mathcal{S}_U \ar[d] \\
\mathcal{I}_{[U/R]} \ar[r] & [U/R]
}
$$
de champs en groupoïdes sur $(\Sch/S)_{fppf}$.
\end{lemma}

\begin{proof}
Par le lemme \ref{spaces-groupoids-lemma-criterion-fibre-product}, il suffit de démontrer que
le morphisme $s' : R' \to G$ du lemme \ref{spaces-groupoids-lemma-presentation-inertia} est
isomorphe au changement de base de $s$ par le morphisme structural
$G \to U$. Cette propriété de changement de base ressort immédiatement de
la construction de $s'$.
\end{proof}







\section{Gerbes et champs quotients}
\label{spaces-groupoids-section-gerbes}

\noindent
Dans cette section, nous mettons les champs quotients en relation avec la
discussion de Champs, section \ref{stacks-section-gerbes}, et en particulier
avec les gerbes telles qu'elles sont définies dans Champs, définition
\ref{stacks-definition-gerbe-over-stack-in-groupoids}. Les champs en
groupoïdes qui apparaissent dans cette section ne sont en général pas des
champs algébriques !

\begin{lemma}
\label{spaces-groupoids-lemma-when-gerbe}
Notations et hypothèses comme dans le lemme
\ref{spaces-groupoids-lemma-quotient-stack-functorial}. Le morphisme de champs quotients
$$
[f] : [U/R] \longrightarrow [U'/R']
$$
fait de $[U/R]$ une gerbe sur $[U'/R']$ si $f : U \to U'$ et
$R \to R'|_U$ sont des morphismes surjectifs de faisceaux fppf. Ici,
$R'|_U$ désigne la restriction de $R'$ à $U$ suivant $f : U \to U'$.
\end{lemma}

\begin{proof}
Nous allons vérifier les propriétés (2)(a) et (2)(b) de Champs, lemme
\ref{stacks-lemma-when-gerbe}. La propriété (2)(a) est satisfaite parce que
$U \to U'$ est un morphisme surjectif de faisceaux (utiliser le lemme
\ref{spaces-groupoids-lemma-quotient-stack-objects} pour voir que les objets de $[U'/R']$
proviennent localement de $U'$). Pour démontrer (2)(b), soient $x, y$ des
objets de $[U/R]$ sur un schéma $T/S$. Soient $x', y'$ les images de $x, y$
dans la catégorie $[U'/'R]_T$. La condition (2)(b) demande de vérifier que
le morphisme de faisceaux
$$
\mathit{Isom}(x, y) \longrightarrow \mathit{Isom}(x', y')
$$
sur $(\Sch/T)_{fppf}$ est surjectif. Pour cela, nous pouvons travailler
localement pour la topologie fppf sur $T$ et supposer que $x, y$ proviennent
de $a, b \in U(T)$. Alors $x', y'$ correspondent à
$f \circ a, f \circ b$. Par le lemme
\ref{spaces-groupoids-lemma-quotient-stack-morphisms}, le morphisme de faisceaux affiché
devient dans ce cas
$$
T \times_{(a, b), U \times_B U} R
\longrightarrow
T \times_{f \circ a, f \circ b, U' \times_B U'} R' =
T \times_{(a, b), U \times_B U} R'|_U.
$$
L'hypothèse selon laquelle $R \to R'|_U$ est un morphisme surjectif de
faisceaux fppf sur $(\Sch/S)_{fppf}$ entraîne donc la surjectivité voulue.
\end{proof}

\begin{lemma}
\label{spaces-groupoids-lemma-group-quotient-gerbe}
Soit $S$ un schéma. Soit $B$ un espace algébrique sur $S$, et soit $G$ un
espace algébrique en groupes sur $B$. Munissons $B$ de l'action triviale de
$G$. Le morphisme
$$
[B/G] \longrightarrow \mathcal{S}_B
$$
(lemme \ref{spaces-groupoids-lemma-quotient-stack-arrows}) fait de $[B/G]$ une gerbe sur $B$.
\end{lemma}

\begin{proof}
Cela résulte immédiatement du lemme \ref{spaces-groupoids-lemma-when-gerbe}, car les
morphismes $B \to B$ et $B \times_B G \to B$ sont surjectifs en tant que
morphismes de faisceaux.
\end{proof}








\section{Champs quotients et changement de grand site}
\label{spaces-groupoids-section-bigger-site}

\noindent
Nous suggérons d'omettre cette section lors d'une première lecture. Les images
inverses de champs sont définies dans Champs, section
\ref{stacks-section-inverse-image}.

\begin{lemma}
\label{spaces-groupoids-lemma-quotient-stack-change-big-site}
Supposons donnés deux grands sites $\Sch_{fppf}$ et $\Sch'_{fppf}$, et que
$\Sch_{fppf}$ soit contenu dans $\Sch'_{fppf}$ ; voir Topologies, section
\ref{topologies-section-change-alpha}. Soit $S \in \Ob(\Sch_{fppf})$.
Soient $B, U, R \in \Sh((\Sch/S)_{fppf})$ des espaces algébriques, et soit
$(U, R, s, t, c)$ un groupoïde en espaces algébriques sur $B$. Soit
$f : (\Sch'/S)_{fppf} \to (\Sch/S)_{fppf}$ le morphisme de sites
correspondant au foncteur d'inclusion
$u : \Sch_{fppf} \to \Sch'_{fppf}$. Il existe alors une équivalence canonique
$$
[f^{-1}U/f^{-1}R]
\longrightarrow
f^{-1}[U/R]
$$
de champs en groupoïdes sur $(\Sch'/S)_{fppf}$.
\end{lemma}

\begin{proof}
Remarquons que $f^{-1}B, f^{-1}U, f^{-1}R \in
\Sh((\Sch'/S)_{fppf})$ sont des espaces algébriques par Espaces, lemme
\ref{spaces-lemma-change-big-site}. Par conséquent,
$(f^{-1}U, f^{-1}R, f^{-1}s, f^{-1}t, f^{-1}c)$ est un groupoïde en
espaces algébriques sur $f^{-1}B$. L'énoncé a donc bien un sens.

\medskip\noindent
La catégorie $u_p[U/_{\!p}R]$ est la localisation de la catégorie
$u_{pp}[U/_{\!p}R]$ par rapport au système multiplicatif à droite $I$ de
morphismes. Un objet de $u_{pp}[U/_{\!p}R]$ est un triplet
$$
(T', \phi : T' \to T, x)
$$
où $T' \in \Ob((\Sch'/S)_{fppf})$, $T \in \Ob((\Sch/S)_{fppf})$,
$\phi$ est un morphisme de schémas sur $S$, et $x : T \to U$ un
morphisme de faisceaux sur $(\Sch/S)_{fppf}$. Remarquons que le morphisme de
schémas $\phi : T' \to T$ revient à un morphisme
$\phi : T' \to u(T)$ et, puisque $u(T)$ représente $f^{-1}T$, à un
morphisme $T' \to f^{-1}T$. En outre, puisque $f^{-1}$ est pleinement fidèle
sur les espaces algébriques, voir Espaces, lemme
\ref{spaces-lemma-fully-faithful}, nous pouvons également considérer $x$
comme un morphisme $x : f^{-1}T \to f^{-1}U$. Nous effectuerons désormais
ces identifications sans les mentionner. Un morphisme
$$
(a, a', \alpha) :
(T'_1, \phi_1 : T'_1 \to T_1, x_1)
\longrightarrow
(T'_2, \phi_2 : T'_2 \to T_2, x_2)
$$
de $u_{pp}[U/_{\!p}R]$ est un diagramme commutatif
$$
\xymatrix{
& & U \\
T'_1 \ar[d]_{a'} \ar[r]_{\phi_1} &
T_1 \ar[d]_a \ar[ru]^{x_1} \ar[r]_\alpha &
R \ar[d]^t \ar[u]_s \\
T'_2 \ar[r]^{\phi_2} &
T_2 \ar[r]^{x_2} &
U
}
$$
Un tel morphisme appartient à $I$ si et seulement si $T'_1 = T'_2$ et
$a' = \text{id}$. Nous définissons un foncteur
$$
u_{pp}[U/_{\!p}R] \longrightarrow [f^{-1}U/_{\!p}f^{-1}R]
$$
par les règles
$$
(T', \phi : T' \to T, x) \longmapsto (x \circ \phi : T' \to f^{-1}U)
$$
sur les objets et
$$
(a, a', \alpha) \longmapsto (\alpha \circ \phi_1 : T'_1 \to f^{-1}R)
$$
sur les morphismes ci-dessus. Il est clair que les éléments de $I$ sont
envoyés sur des isomorphismes, puisque
$(f^{-1}U, f^{-1}R, f^{-1}s, f^{-1}t, f^{-1}c)$ est un groupoïde en
espaces algébriques sur $f^{-1}B$. Ce foncteur se factorise donc canoniquement
par un foncteur
$$
u_p[U/_{\!p}R] \longrightarrow [f^{-1}U/_{\!p}f^{-1}R]
$$
En passant aux champs associés, nous obtenons un foncteur entre champs
$$
f^{-1}[U/R] \longrightarrow [f^{-1}U/f^{-1}R]
$$
sur $(\Sch'/S)_{fppf}$, car, d'après Champs, lemme
\ref{stacks-lemma-technical-up}, le champ $f^{-1}[U/R]$ est le champ associé
à $u_p[U/_{\!p}R]$.

\medskip\noindent
Nous avons maintenant un morphisme de champs. Pour vérifier que c'est une
équivalence, il suffit de montrer qu'il est pleinement fidèle et que les
objets appartiennent localement à l'image essentielle ; voir Champs, lemmes
\ref{stacks-lemma-characterize-ff} et
\ref{stacks-lemma-characterize-essentially-surjective-when-ff}. L'assertion
sur les objets découle de ce que $f^{-1}R$ admet un morphisme étale surjectif
$f^{-1}W \to f^{-1}R$ pour un objet $W$ de $(\Sch/S)_{fppf}$. Pour montrer
que le foncteur est « plein », il suffit de montrer que les morphismes
appartiennent localement à son image ; cela découle de ce que $f^{-1}U$ admet
un morphisme étale surjectif $f^{-1}W \to f^{-1}U$ pour un objet $W$ de
$(\Sch/S)_{fppf}$. Nous omettons la démonstration de la fidélité.
\end{proof}










\section{Conditions de séparation}
\label{spaces-groupoids-section-separation}

\noindent
Il s'agit en fait de conditions portant sur le morphisme
$j : R \to U \times_B U$ lorsqu'un groupoïde en espaces algébriques
$(U, R, s, t, c)$ sur $B$ est donné. Comme dans la section précédente, nous
commençons par formuler le diagramme correspondant.

\begin{lemma}
\label{spaces-groupoids-lemma-diagram-diagonal}
Soit $B \to S$ comme dans la section \ref{spaces-groupoids-section-notation}. Soit
$(U, R, s, t, c)$ un groupoïde en espaces algébriques sur $B$, et soit
$G \to U$ l'espace algébrique en groupes stabilisateur. Dans le diagramme
commutatif
$$
\xymatrix{
R \ar[d]^{\Delta_{R/U \times_B U}} \ar[rrr]_{f \mapsto (f, s(f))} & & &
R \times_{s, U} U \ar[d] \ar[r] & U \ar[d] \\
R \times_{(U \times_B U)} R \ar[rrr]^{(f, g) \mapsto (f, f^{-1} \circ g)} & & &
R \times_{s, U} G \ar[r] & G
}
$$
les deux flèches horizontales de gauche sont des isomorphismes et le carré de
droite est cartésien.
\end{lemma}

\begin{proof}
Démonstration omise. C'est un exercice sur les définitions et le point de vue
fonctoriel en géométrie algébrique.
\end{proof}

\begin{lemma}
\label{spaces-groupoids-lemma-diagonal}
Soit $B \to S$ comme dans la section \ref{spaces-groupoids-section-notation}. Soit
$(U, R, s, t, c)$ un groupoïde en espaces algébriques sur $B$, et soit
$G \to U$ l'espace algébrique en groupes stabilisateur.
\begin{enumerate}
\item Les conditions suivantes sont équivalentes :
\begin{enumerate}
\item $j : R \to U \times_B U$ est séparé,
\item $G \to U$ est séparé, et
\item $e : U \to G$ est une immersion fermée.
\end{enumerate}
\item Les conditions suivantes sont équivalentes :
\begin{enumerate}
\item $j : R \to U \times_B U$ est localement séparé,
\item $G \to U$ est localement séparé, et
\item $e : U \to G$ est une immersion.
\end{enumerate}
\item Les conditions suivantes sont équivalentes :
\begin{enumerate}
\item $j : R \to U \times_B U$ est quasi-séparé,
\item $G \to U$ est quasi-séparé, et
\item $e : U \to G$ est quasi-compact.
\end{enumerate}
\end{enumerate}
\end{lemma}

\begin{proof}
L'espace algébrique en groupes $G \to U$ s'obtient à partir de
$R \to U \times_B U$ par changement de base suivant le morphisme diagonal
$U \to U \times_B U$ ; voir le lemme \ref{spaces-groupoids-lemma-groupoid-stabilizer}. Ainsi,
si $j$ est séparé (resp.\ localement séparé, resp.\ quasi-séparé), alors
$G \to U$ est séparé (resp.\ localement séparé, resp.\ quasi-séparé) ; voir
Morphismes d'espaces, lemme
\ref{spaces-morphisms-lemma-base-change-separated}. Nous obtenons donc
(a) $\Rightarrow$ (b) dans (1), (2) et (3).

\medskip\noindent
Réciproquement, si $G \to U$ est séparé (resp.\ localement séparé,
resp.\ quasi-séparé), alors le morphisme $e : U \to G$, qui est une section
du morphisme structural $G \to U$, est une immersion fermée (resp.\ une
immersion, resp.\ quasi-compacte) ; voir Morphismes d'espaces, lemme
\ref{spaces-morphisms-lemma-section-immersion}. Nous obtenons donc
(b) $\Rightarrow$ (c) dans (1), (2) et (3).

\medskip\noindent
Si $e$ est une immersion fermée (resp.\ une immersion, resp.\ quasi-compacte),
alors le lemme \ref{spaces-groupoids-lemma-diagram-diagonal} (ainsi que Espaces, lemme
\ref{spaces-lemma-base-change-immersions}, et Morphismes d'espaces, lemme
\ref{spaces-morphisms-lemma-base-change-quasi-compact}) montre que
$\Delta_{R/U \times_B U}$ est une immersion fermée (resp.\ une immersion,
resp.\ quasi-compacte). Nous obtenons donc (c) $\Rightarrow$ (a) dans (1),
(2) et (3).
\end{proof}



















% OK, start here.
%

\chapter{Compléments sur les groupoïdes en espaces algébriques}



\phantomsection
\label{spaces-more-groupoids-section-phantom}


\section{Introduction}
\label{spaces-more-groupoids-section-introduction}

\noindent
Ce chapitre est consacré à des questions avancées sur les groupoïdes
en espaces algébriques.
Bien que les résultats soient énoncés en termes de groupoïdes en
espaces algébriques, le
lecteur doit garder à l'esprit le diagramme $2$-cartésien
\begin{equation}
\label{spaces-more-groupoids-equation-quotient-stack}
\vcenter{
\xymatrix{
R \ar[r] \ar[d] & U \ar[d] \\
U \ar[r] & [U/R]
}
}
\end{equation}
où $[U/R]$ est le champ quotient; voir
Groupoïdes dans les espaces, remarque \ref{spaces-groupoids-remark-fundamental-square}.
Nombre des résultats sont motivés par l'étude de ce diagramme.
Voir par exemple le bel article \cite{K-M} de Keel et Mori.





\section{Notations}
\label{spaces-more-groupoids-section-notation}

\noindent
Nous continuons à suivre les conventions et notations introduites dans
Groupoïdes dans les espaces, section \ref{spaces-groupoids-section-notation}.





\section{Diagrammes utiles}
\label{spaces-more-groupoids-section-diagrams}

\noindent
Rappelons brièvement les résultats de
Groupoïdes dans les espaces, lemmes \ref{spaces-groupoids-lemma-diagram} et
\ref{spaces-groupoids-lemma-diagram-pull}
afin de pouvoir nous y référer aisément dans ce chapitre.
Soit $S$ un schéma. Soit $B$ un espace algébrique au-dessus de $S$.
Soit $(U, R, s, t, c)$ un groupoïde en espaces algébriques au-dessus de $B$.
Dans le diagramme commutatif
\begin{equation}
\label{spaces-more-groupoids-equation-diagram}
\vcenter{
\xymatrix{
& U & \\
R \ar[d]_s \ar[ru]^t &
R \times_{s, U, t} R
\ar[l]^-{\text{pr}_0} \ar[d]^{\text{pr}_1} \ar[r]_-c &
R \ar[d]^s \ar[lu]_t \\
U & R \ar[l]_t \ar[r]^s & U
}
}
\end{equation}
les deux carrés inférieurs sont cartésiens.
De plus, le triangle supérieur (qui est en réalité un carré)
est lui aussi cartésien.

\medskip\noindent
Le diagramme
\begin{equation}
\label{spaces-more-groupoids-equation-pull}
\vcenter{
\xymatrix{
R \times_{t, U, t} R
\ar@<1ex>[r]^-{\text{pr}_1} \ar@<-1ex>[r]_-{\text{pr}_0}
\ar[d]_{\text{pr}_0 \times c \circ (i, 1)} &
R \ar[r]^t \ar[d]^{\text{id}_R} &
U \ar[d]^{\text{id}_U} \\
R \times_{s, U, t} R
\ar@<1ex>[r]^-c \ar@<-1ex>[r]_-{\text{pr}_0} \ar[d]_{\text{pr}_1} &
R \ar[r]^t \ar[d]^s &
U \\
R \ar@<1ex>[r]^s \ar@<-1ex>[r]_t &
U
}
}
\end{equation}
est commutatif. Les deux lignes supérieures sont isomorphes par les flèches
verticales indiquées. Les deux carrés inférieurs de gauche sont cartésiens.






\section{Structure locale}
\label{spaces-more-groupoids-section-local}

\noindent
Soit $S$ un schéma.
Soit $(U, R, s, t, c, e, i)$ un groupoïde en espaces algébriques au-dessus de $S$.
Soit $\overline{u}$ un point géométrique de $U$. Dans cette section, nous
décrivons la structure obtenue sur les anneaux locaux
(Propriétés des espaces, définition
\ref{spaces-properties-definition-etale-local-rings})
$$
A = \mathcal{O}_{U, \overline{u}}
\quad\text{et}\quad
B = \mathcal{O}_{R, e(\overline{u})}
$$
Nous convenons de désigner par les lettres correspondantes les homomorphismes
d'anneaux locaux induits par les morphismes $s, t, c, e, i$.
En particulier, nous avons un diagramme commutatif
$$
\xymatrix{
A \ar[rd]_t \ar[rrd]^1 \\
& B \ar[r]^e & A \\
A \ar[ru]^s \ar[rru]_1
}
$$
d'anneaux locaux. Ainsi, si $I \subset B$ désigne le noyau de $e : B \to A$,
alors $B = s(A) \oplus I = t(A) \oplus I$. Posons
$$
C = \mathcal{O}_{R \times_{s, U, t} R, (e, e)(\overline{u})}
$$
Alors
$$
C =
(B \otimes_{s, A, t} B)_{\mathfrak m_B \otimes B + B \otimes \mathfrak m_B}^h
$$
car le localisé
$(B \otimes_{s, A, t} B)_{\mathfrak m_B \otimes B + B \otimes \mathfrak m_B}$
a un corps résiduel séparablement clos.
Soit $J \subset C$ l'idéal de $C$ engendré par $I \otimes B + B \otimes I$.
Alors $J$ est aussi le noyau de l'homomorphisme d'anneaux locaux
$$
(e, e) : C \longrightarrow A
$$
La loi de composition $c : R \times_{s, U, t} R \to R$ correspond à un
homomorphisme d'anneaux
$$
c : B \longrightarrow C
$$
qui envoie $I$ dans $J$.

\begin{lemma}
\label{spaces-more-groupoids-lemma-first-order-structure-c}
L'application $I/I^2 \to J/J^2$ induite par $c$ est la composée
$$
I/I^2 \xrightarrow{(1, 1)} I/I^2 \oplus I/I^2 \to J/J^2
$$
où la seconde flèche provient de l'égalité
$J = (I \otimes B + B \otimes I)C$.
L'application $i : B \to B$ induit l'application $-1 : I/I^2 \to I/I^2$.
\end{lemma}

\begin{proof}
Pour décrire un homomorphisme local de $C$ vers un autre
anneau local hensélien, il suffit de préciser l'image
des éléments de la forme $b_1 \otimes b_2$, par exemple d'après
Algèbre, lemme \ref{algebra-lemma-henselian-functorial}.
Compte tenu de cela, nous avons les deux applications canoniques
$$
e_2 : C \to B,\ b_1 \otimes b_2 \mapsto b_1s(e(b_2)),\quad
e_1 : C \to B,\ b_1 \otimes b_2 \mapsto t(e(b_1))b_2
$$
correspondant aux immersions
$R \to R \times_{s, U, t} R$ données par
$r \mapsto (r, e(s(r)))$ et $r \mapsto (e(t(r)), r)$.
Ces applications définissent des applications $J/J^2 \to I/I^2$ qui,
prises ensemble, fournissent un inverse de l'application
$I/I^2 \oplus I/I^2 \to J/J^2$ du lemme. Pour démontrer l'assertion, il suffit
donc de montrer que $e_1 \circ c : B \to B$ et $e_2 \circ c : B \to B$
sont les applications identiques. Cela résulte du fait que les deux
composées $R \to R \times_{s, U, t} R \to R$ sont des identités.

\medskip\noindent
L'assertion concernant $i$ résulte de celle concernant $c$ et du
fait que $c \circ (1, i) = e \circ t$. Certains détails sont omis.
\end{proof}








\section{Groupoïde des sections}
\label{spaces-more-groupoids-section-groupoid-sections}

\noindent
Supposons donné un groupoïde $(\text{Ob}, \text{Flèches}, s, t, c, e, i)$.
On peut alors construire un monoïde $\Gamma$ dont les éléments sont les
applications $\delta : \text{Ob} \to \text{Flèches}$ telles que
$s \circ \delta = \text{id}_{\text{Ob}}$, la composition étant donnée par
$$
\delta_1 \circ \delta_2 = c(\delta_1 \circ t \circ \delta_2, \delta_2)
$$
Autrement dit, un élément de $\Gamma$ est une règle $\delta$ qui
prescrit une flèche issue de chaque objet, et la composition
est celle que l'on attend naturellement. Par exemple,
$$
\vcenter{
\xymatrix{
& \bullet \ar[dl] \\
\bullet \ar@(ul, dl)[] & \bullet \ar[d] \\
& \bullet \ar[lu]
}
}
\quad\circ\quad
\vcenter{
\xymatrix{
& \bullet \ar[d] \\
\bullet \ar[ru] & \bullet \ar[d] \\
& \bullet \ar[lu]
}
}
\quad = \quad\quad
\vcenter{
\xymatrix{
& \bullet \ar@/^/[dd] \\
\bullet \ar@(ul, dl)[] & \bullet \ar[l] \\
& \bullet \ar[lu]
}
}
$$
avec des notations évidentes.

\medskip\noindent
Le même procédé s'applique à un groupoïde en espaces
algébriques $(U, R, s, t, c, e, i)$ au-dessus d'un schéma $S$.
Plus précisément, on prend pour éléments de $\Gamma$ l'ensemble
$$
\Gamma = \{\delta : U \to R \mid s \circ \delta = \text{id}_U\}
$$
et la composition $\circ : \Gamma \times \Gamma \to \Gamma$
est donnée par la règle ci-dessus :
\begin{equation}
\label{spaces-more-groupoids-equation-composition}
\delta_1 \circ \delta_2 = c(\delta_1 \circ t \circ \delta_2, \delta_2)
\end{equation}
L'élément neutre est $e \in \Gamma$. En général, le monoïde $\Gamma$ n'est pas
un groupe, car certains éléments $\delta \in \Gamma$ peuvent ne pas
avoir d'inverse. En effet, il est clair que $\delta \in \Gamma$
a un inverse si et seulement si $t \circ \delta$ est un automorphisme
de $U$; dans ce cas,
$\delta^{-1} = i \circ \delta \circ (t \circ \delta)^{-1}$.

\medskip\noindent
Pour un usage ultérieur, examinons le sous-groupe
$\Gamma_0$ de $\Gamma$ formé des sections infinitésimalement
proches de l'identité $e$. Plus précisément, supposons donné un
sous-espace fermé $R$-invariant $U_0 \subset U$ tel que $U$
soit un épaississement du premier ordre de $U_0$. Notons
$R_0 = s^{-1}(U_0) = t^{-1}(U_0)$ et soit
$(U_0, R_0, s_0, t_0, c_0, e_0, i_0)$ le groupoïde
en espaces algébriques correspondant. Posons
$$
\Gamma_0 = \{\delta \in \Gamma \mid \delta|_{U_0} = e_0\}
$$
Si $s$ et $t$ sont plats, tout élément de $\Gamma_0$
est inversible. En effet, $t \circ \delta$ est alors un
morphisme $U \to U$ qui induit l'identité sur
$\mathcal{O}_{U_0}$ et sur $\mathcal{C}_{U_0/U}$
(lemme \ref{spaces-more-groupoids-lemma-idenity-on-conormal});
on conclut grâce à la suite exacte
courte
$0 \to \mathcal{C}_{U_0/U} \to \mathcal{O}_U \to \mathcal{O}_{U_0} \to 0$.

\begin{lemma}
\label{spaces-more-groupoids-lemma-idenity-on-conormal}
Dans la situation étudiée dans cette section, soient $\delta \in \Gamma_0$
et $f = t \circ \delta : U \to U$. Si $s, t$ sont plats, alors l'application
canonique $\mathcal{C}_{U_0/U} \to \mathcal{C}_{U_0/U}$ induite par $f$
(Compléments sur les morphismes d'espaces, lemme
\ref{spaces-more-morphisms-lemma-conormal-functorial})
est l'application identique.
\end{lemma}

\begin{proof}
Pour le voir, prolongeons comme suit la partie inférieure du diagramme
(\ref{spaces-more-groupoids-equation-pull}) :
$$
\xymatrix{
Y \ar[r] \ar[d] &
R \times_{s, U, t} R
\ar@<1ex>[r]^-c \ar@<-1ex>[r]_-{\text{pr}_0} \ar[d]_{\text{pr}_1} &
R \ar[r]^t \ar[d]^s &
U \\
U \ar[r]_\delta &
R \ar@<1ex>[r]^s \ar@<-1ex>[r]_t &
U
}
$$
où le carré de gauche est cartésien; ceci constitue notre définition
de $Y$; nous n'aurons pas besoin d'en savoir davantage sur $Y$.
Un changement de base partout à $U_0$ donne un diagramme analogue
ayant les mêmes propriétés.
Nous voulons montrer que $\text{id}_U = s \circ \delta$
et $f = t \circ \delta$ induisent les mêmes applications sur les faisceaux
conormaux. Comme $s$ est plat et surjectif, il suffit de démontrer la même
chose pour les deux composées $a, b : Y \to R$ de la ligne supérieure.
Observons que $a_0 = b_0$ et que l'un de $a$ et $b$ est un isomorphisme,
puisque nous savons que $s \circ \delta$ est un isomorphisme. Par conséquent,
les deux morphismes $a, b : Y \to R$ sont des morphismes entre espaces
algébriques plats au-dessus de $U$ (par le morphisme $t : R \to U$
et le morphisme $t \circ a = t \circ b : Y \to U$).
Cela donne le résultat voulu. En effet, par compatibilité aux compositions dans
Compléments sur les morphismes d'espaces, lemme
\ref{spaces-more-morphisms-lemma-conormal-functorial-more}
on conclut que les deux applications
$a_0^*\mathcal{C}_{R_0/R} \to \mathcal{C}_{Y_0/Y}$
s'insèrent dans un diagramme commutatif
$$
\xymatrix{
a_0^*\mathcal{C}_{R_0/R} \ar[rr] & & \mathcal{C}_{Y_0/Y} \\
a_0^*t_0^*\mathcal{C}_{U_0/U} \ar[u] \ar@{=}[rr] & &
(t_0 \circ a_0)^*\mathcal{C}_{U_0/U} \ar[u]
}
$$
dont les flèches verticales sont des isomorphismes d'après
Compléments sur les morphismes d'espaces, lemme
\ref{spaces-more-morphisms-lemma-deform}.
Le lemme en résulte.
\end{proof}

\noindent
Identifions maintenant le groupe $\Gamma_0$.
Appliquons la discussion de Compléments sur les morphismes d'espaces, remarques
\ref{spaces-more-morphisms-remark-action-by-derivations} et
\ref{spaces-more-morphisms-remark-another-special-case}
au diagramme
$$
\xymatrix{
(U_0 \subset U) \ar@{..>}[rr]_{(e_0, \delta)}
\ar[rd]_{(\text{id}_{U_0}, \text{id}_U)} & &
(R_0 \subset R) \ar[ld]^{(s_0, s)} \\
& (U_0 \subset U)
}
$$
On voit que $\delta = \theta \cdot e$ pour une unique application
$\mathcal{O}_{U_0}$-linéaire
$\theta : e_0^*\Omega_{R_0/U_0} \to \mathcal{C}_{U_0/U}$.
On obtient ainsi une bijection
\begin{equation}
\label{spaces-more-groupoids-equation-isomorphism}
\Hom_{\mathcal{O}_{U_0}}(e_0^*\Omega_{R_0/U_0}, \mathcal{C}_{U_0/U})
\longrightarrow
\Gamma_0
\end{equation}
en appliquant Compléments sur les morphismes d'espaces, lemme
\ref{spaces-more-morphisms-lemma-action-sheaf}.

\begin{lemma}
\label{spaces-more-groupoids-lemma-composition-is-addition}
La bijection (\ref{spaces-more-groupoids-equation-isomorphism}) est un isomorphisme
de groupes.
\end{lemma}

\begin{proof}
Soient $\delta_1, \delta_2 \in \Gamma_0$ correspondant à $\theta_1, \theta_2$
comme ci-dessus, et la composée $\delta = \delta_1 \circ \delta_2$
dans $\Gamma_0$ correspondant à $\theta$. Il faut montrer que
$\theta = \theta_1 + \theta_2$. Rappelons
(Compléments sur les morphismes d'espaces, lemme
\ref{spaces-more-morphisms-lemma-action-by-derivations})
que $\theta_1, \theta_2, \theta$ correspondent aux dérivations
$D_1, D_2, D : e_0^{-1}\mathcal{O}_{R_0} \to \mathcal{C}_{U_0/U}$
données par $D_1 = \theta_1 \circ \text{d}_{R_0/U_0}$, et ainsi de suite.
Il suffit de vérifier que $D = D_1 + D_2$.

\medskip\noindent
On peut vérifier l'égalité sur les fibres.
Soit $\overline{u}$ un point géométrique de $U$ et utilisons les anneaux
locaux $A, B, C$ introduits dans la section \ref{spaces-more-groupoids-section-local}.
Aux morphismes $\delta_i$ correspondent des homomorphismes d'anneaux
$\delta_i : B \to A$. Soit $K \subset A$ l'idéal de carré nul tel que
$A/K = \mathcal{O}_{U_0, \overline{u}}$.
Autrement dit, $K$ est la fibre de $\mathcal{C}_{U_0/U}$ en $\overline{u}$.
La condition $\delta_i \in \Gamma_0$
signifie exactement que $\delta_i(I) \subset K$.
La dérivation $D_i$ est simplement l'application $\delta_i - e : B \to A$.
Comme $B = s(A) \oplus I$, la dérivation $D_i$ est déterminée par
sa restriction à $I$, laquelle est précisément donnée par
$\delta_i|_I$. De plus, $D_i$, et donc $\delta_i$, annule $I^2$
puisque $\delta_i(I) \subset K$ et $K^2 = 0$.

\medskip\noindent
Pour achever la démonstration, observons que $\delta$ correspond à la
composée
$$
B \to C =
(B \otimes_{s, A, t} B)^h_{\mathfrak m_B \otimes B + B \otimes \mathfrak m_B}
\to A
$$
où la première flèche est $c$ et la seconde est déterminée
par la règle
$b_1 \otimes b_2 \mapsto \delta_2(t(\delta_1(b_1))) \delta_2(b_2)$
comme il résulte de (\ref{spaces-more-groupoids-equation-composition}).
D'après le lemme \ref{spaces-more-groupoids-lemma-first-order-structure-c},
on voit qu'un élément $\zeta$ de $I$ est envoyé sur
$\zeta \otimes 1 + 1 \otimes \zeta$, à des termes d'ordre supérieur près.
On en conclut que
$$
D(\zeta) = (\delta_2 \circ t)\left(D_1(\zeta)\right) + D_2(\zeta)
$$
Or, d'après le lemme \ref{spaces-more-groupoids-lemma-idenity-on-conormal},
l'action de $\delta_2 \circ t$ sur
$K = \mathcal{C}_{U_0/U, \overline{u}}$ est l'identité, ce qui
achève la preuve.
\end{proof}






\section{Propriétés des groupoïdes}
\label{spaces-more-groupoids-section-technical-lemma}

\noindent
Cette section est l'analogue de
Compléments sur les groupoïdes, section \ref{more-groupoids-section-technical-lemma}.
Il est vivement conseillé au lecteur de commencer par lire cette section.

\medskip\noindent
Le lemme suivant est l'analogue de
Compléments sur les groupoïdes, lemme \ref{more-groupoids-lemma-property-invariant}.

\begin{lemma}
\label{spaces-more-groupoids-lemma-property-invariant}
Soit $B \to S$ comme dans la section \ref{spaces-more-groupoids-section-notation}.
Soit $(U, R, s, t, c)$ un groupoïde en espaces algébriques au-dessus de $B$.
Soit
$\tau \in \{fppf, \linebreak[0] \etale, \linebreak[0]
lisse, \linebreak[0] syntomique\}$.
Soit $\mathcal{P}$ une propriété des morphismes d'espaces algébriques
qui est $\tau$-locale sur le but
(Descente sur les espaces,
définition \ref{spaces-descent-definition-property-morphisms-local}).
Supposons que $\{s : R \to U\}$ et $\{t : R \to U\}$ soient des recouvrements
pour la topologie $\tau$. Soit $W \subset U$ le plus grand sous-espace ouvert
tel que $s^{-1}(W) \to W$ possède la propriété $\mathcal{P}$.
Alors $W$ est $R$-invariant
(Groupoïdes dans les espaces,
définition \ref{spaces-groupoids-definition-invariant-open}).
\end{lemma}

\begin{proof}
L'existence et les propriétés de l'ouvert $W \subset U$ sont décrites dans
Descente sur les espaces, lemme \ref{spaces-descent-lemma-largest-open-of-the-base}.
Dans le
diagramme (\ref{spaces-more-groupoids-equation-diagram}),
soit $W_1 \subset R$ le plus grand sous-schéma ouvert sur lequel le morphisme
$\text{pr}_1 : R \times_{s, U, t} R \to R$ possède la propriété $\mathcal{P}$.
Il résulte du lemme précité de
Descente sur les espaces, lemme \ref{spaces-descent-lemma-largest-open-of-the-base},
et de l'hypothèse que $\{s : R \to U\}$ et $\{t : R \to U\}$ sont des
recouvrements pour la topologie $\tau$, que $t^{-1}(W)=W_1=s^{-1}(W)$, comme voulu.
\end{proof}

\begin{lemma}
\label{spaces-more-groupoids-lemma-property-G-invariant}
Soit $B \to S$ comme dans la section \ref{spaces-more-groupoids-section-notation}.
Soit $(U, R, s, t, c)$ un groupoïde en espaces algébriques au-dessus de $B$.
Soit $G \to U$ son espace algébrique en groupes stabilisateur.
Soit
$\tau \in \{fppf, \linebreak[0] \etale, \linebreak[0]
lisse, \linebreak[0] syntomique\}$.
Soit $\mathcal{P}$ une propriété des morphismes d'espaces algébriques
qui est $\tau$-locale sur le but.
Supposons que $\{s : R \to U\}$ et $\{t : R \to U\}$ soient des recouvrements
pour la topologie $\tau$. Soit $W \subset U$ le plus grand sous-espace ouvert
tel que $G_W \to W$ possède la propriété $\mathcal{P}$.
Alors $W$ est $R$-invariant (voir
Groupoïdes dans les espaces,
définition \ref{spaces-groupoids-definition-invariant-open}).
\end{lemma}

\begin{proof}
L'existence et les propriétés de l'ouvert $W \subset U$ sont décrites dans
Descente sur les espaces, lemme \ref{spaces-descent-lemma-largest-open-of-the-base}.
Le morphisme
$$
G \times_{U, t} R \longrightarrow R \times_{s, U} G, \quad
(g, r) \longmapsto (r, r^{-1} \circ g \circ r)
$$
est un isomorphisme d'espaces algébriques au-dessus de $R$ (où $\circ$ désigne
la composition dans le groupoïde). Ainsi $s^{-1}(W)=t^{-1}(W)$ d'après les
propriétés de $W$ démontrées dans le lemme précité de
Descente sur les espaces, lemme \ref{spaces-descent-lemma-largest-open-of-the-base}.
\end{proof}




\section{Comparaison des fibres}
\label{spaces-more-groupoids-section-fibres}

\noindent
Cette section est l'analogue de
Compléments sur les groupoïdes, section \ref{more-groupoids-section-fibres}.
Il est vivement conseillé au lecteur de commencer par lire cette section.

\begin{lemma}
\label{spaces-more-groupoids-lemma-two-fibres}
Soit $B \to S$ comme dans la section \ref{spaces-more-groupoids-section-notation}.
Soit $(U, R, s, t, c)$ un groupoïde en espaces algébriques au-dessus de $B$.
Soit $K$ un corps et soient $r, r' : \Spec(K) \to R$
des morphismes tels que $t \circ r = t \circ r' : \Spec(K) \to U$.
Posons $u = s \circ r$, $u' = s \circ r'$ et notons
$F_u = \Spec(K) \times_{u, U, s} R$ et
$F_{u'} = \Spec(K) \times_{u', U, s} R$ les produits fibrés.
Alors $F_u \cong F_{u'}$ comme espaces algébriques au-dessus de $K$.
\end{lemma}

\begin{proof}
Nous utilisons l'existence et les propriétés du
diagramme (\ref{spaces-more-groupoids-equation-diagram}).
Il existe un morphisme $\xi : \Spec(K) \to R \times_{s, U, t} R$
tel que $\text{pr}_0 \circ \xi = r$ et $c \circ \xi = r'$.
Soit $\tilde r = \text{pr}_1 \circ \xi : \Spec(K) \to R$.
En examinant alors les deux carrés inférieurs du
diagramme (\ref{spaces-more-groupoids-equation-diagram}),
on voit que $F_u$ et $F_{u'}$ s'identifient tous deux à l'espace algébrique
$\Spec(K) \times_{\tilde r, R, \text{pr}_1} (R \times_{s, U, t} R)$.
\end{proof}

\noindent
En fait, dans la situation du lemme, les morphismes de paires
$s : (R, r) \to (U, u)$ et $s : (R, r') \to (U, u')$ sont
localement isomorphes pour la topologie $\tau$, pourvu que $\{s: R \to U\}$
soit un recouvrement $\tau$. Nous insérerons ici un énoncé précis au besoin.








\section{Restriction des groupoïdes}
\label{spaces-more-groupoids-section-restricting-groupoids}

\noindent
Dans cette section, nous rassemblons plusieurs lemmes sur les
propriétés des groupoïdes qui se transmettent aux restrictions.
La plupart se démontrent en considérant le diagramme
définissant la restriction :
\begin{equation}
\label{spaces-more-groupoids-equation-restriction}
\vcenter{
\xymatrix{
R' \ar[d] \ar[r] \ar@/_3pc/[dd]_{t'} \ar@/^1pc/[rr]^{s'}&
R \times_{s, U} U' \ar[r] \ar[d] &
U' \ar[d]^g \\
U' \times_{U, t} R \ar[d] \ar[r] &
R \ar[r]^s \ar[d]_t &
U \\
U' \ar[r]^g &
U
}
}
\end{equation}
Voir
Groupoïdes dans les espaces, lemme \ref{spaces-groupoids-lemma-restrict-groupoid}.

\begin{lemma}
\label{spaces-more-groupoids-lemma-restrict-preserves-type}
Soit $S$ un schéma. Soit $B$ un espace algébrique au-dessus de $S$.
Soit $(U, R, s, t, c)$ un groupoïde en espaces algébriques au-dessus de $B$.
Soit $g : U' \to U$ un morphisme d'espaces algébriques au-dessus de $B$.
Soit $(U', R', s', t', c')$ la restriction de
$(U, R, s, t, c)$ par $g$.
\begin{enumerate}
\item Si $s, t$ sont localement de type fini et $g$ est localement de type
fini, alors $s', t'$ sont localement de type fini.
\item Si $s, t$ sont localement de présentation finie et $g$ est localement de
présentation finie, alors $s', t'$ sont localement de présentation finie.
\item Si $s, t$ sont plats et $g$ est plat, alors $s', t'$ sont plats.
\item Ajouter d'autres cas ici.
\end{enumerate}
\end{lemma}

\begin{proof}
La propriété d'être localement de type fini est stable par composition
et par changement de base arbitraire; voir
Morphismes d'espaces,
lemmes \ref{spaces-morphisms-lemma-composition-finite-type} et
\ref{spaces-morphisms-lemma-base-change-finite-type}.
Le point (1) résulte donc du diagramme (\ref{spaces-more-groupoids-equation-restriction}).
Pour les autres cas, voir
Morphismes d'espaces,
lemmes \ref{spaces-morphisms-lemma-composition-finite-presentation},
\ref{spaces-morphisms-lemma-base-change-finite-presentation},
\ref{spaces-morphisms-lemma-composition-flat} et
\ref{spaces-morphisms-lemma-base-change-flat}.
\end{proof}









\section{Propriétés des groupes sur les corps et des groupoïdes sur les corps}
\label{spaces-more-groupoids-section-properties-groupoids-on-fields}

\noindent
Il est conseillé au lecteur de consulter d'abord les sections correspondantes
pour les schémas en groupoïdes; voir
Groupoïdes, section \ref{groupoids-section-properties-group-schemes-field}
et
Compléments sur les groupoïdes,
section \ref{more-groupoids-section-properties-groupoids-on-fields}.

\begin{situation}
\label{spaces-more-groupoids-situation-group-over-field}
Ici, $S$ est un schéma, $k$ est un corps au-dessus de $S$, et
$(G, m)$ est un espace algébrique en groupes au-dessus de $\Spec(k)$.
\end{situation}

\begin{situation}
\label{spaces-more-groupoids-situation-groupoid-on-field}
Ici, $S$ est un schéma, $B$ est un espace algébrique, et
$(U, R, s, t, c)$ est un groupoïde en espaces algébriques au-dessus de $B$
avec $U = \Spec(k)$ pour un certain corps $k$.
\end{situation}

\noindent
Remarquons que, dans la
situation \ref{spaces-more-groupoids-situation-group-over-field},
on obtient un groupoïde en espaces algébriques
\begin{equation}
\label{spaces-more-groupoids-equation-groupoid-from-group}
(\Spec(k), G, p, p, m)
\end{equation}
où $p : G \to \Spec(k)$ est le morphisme structural de $G$; voir
Groupoïdes dans les espaces, lemme \ref{spaces-groupoids-lemma-groupoid-from-action}.
On se trouve alors dans la
situation \ref{spaces-more-groupoids-situation-groupoid-on-field}.
Nous l'utiliserons sans autre mention dans la suite de cette section.

\begin{lemma}
\label{spaces-more-groupoids-lemma-groupoid-on-field-open-multiplication}
Dans la
situation \ref{spaces-more-groupoids-situation-groupoid-on-field},
le morphisme de composition $c : R \times_{s, U, t} R \to R$ est plat et
universellement ouvert.
Dans la
situation \ref{spaces-more-groupoids-situation-group-over-field},
la loi de groupe $m : G \times_k G \to G$ est plate et
universellement ouverte.
\end{lemma}

\begin{proof}
La composition est isomorphe à la projection
$\text{pr}_1 : R \times_{t, U, t} R \to R$ d'après le
diagramme (\ref{spaces-more-groupoids-equation-pull}).
La projection est plate, comme changement de base du morphisme plat $t$,
et ouverte d'après
Morphismes d'espaces,
lemme \ref{spaces-morphisms-lemma-space-over-field-universally-open}.
La seconde assertion résulte aussitôt de la première, puisque
$m$ correspond à $c$ dans (\ref{spaces-more-groupoids-equation-groupoid-from-group}).
\end{proof}

\noindent
Remarquons que le lemme suivant s'applique notamment lorsqu'on travaille
avec des espaces algébriques quasi-séparés ou localement séparés
(Espaces décents, lemme \ref{decent-spaces-lemma-locally-separated-decent}).

\begin{lemma}
\label{spaces-more-groupoids-lemma-group-scheme-over-field-separated}
Dans la situation \ref{spaces-more-groupoids-situation-groupoid-on-field},
supposons que $R$ soit un espace décent. Alors $R$ est un espace algébrique séparé.
Dans la situation \ref{spaces-more-groupoids-situation-group-over-field}, supposons que
$G$ soit un espace algébrique décent. Alors $G$ est un espace algébrique séparé.
\end{lemma}

\begin{proof}
Commençons par démontrer la seconde assertion. D'après Groupoïdes dans les espaces,
lemme \ref{spaces-groupoids-lemma-group-scheme-separated},
il faut montrer que $e : S \to G$ est une immersion fermée.
Cela résulte d'Espaces décents, lemme
\ref{decent-spaces-lemma-finite-residue-field-extension-finite}.

\medskip\noindent
Démontrons ensuite la première assertion. On peut pour cela remplacer $B$ par $S$.
D'après le paragraphe précédent, le schéma en groupes stabilisateur $G \to U$ est séparé. D'après
Groupoïdes dans les espaces, lemme \ref{spaces-groupoids-lemma-diagonal},
le morphisme $j = (t, s) : R \to U \times_S U$ est séparé.
Comme $U$ est le spectre d'un corps, le schéma
$U \times_S U$ est affine (par la construction des produits fibrés dans
Schémas, section \ref{schemes-section-fibre-products}).
Par conséquent, $R$ est séparé; voir
Morphismes d'espaces, lemme
\ref{spaces-morphisms-lemma-separated-over-separated}.
\end{proof}

\begin{lemma}
\label{spaces-more-groupoids-lemma-restrict-groupoid-on-field}
Dans la
situation \ref{spaces-more-groupoids-situation-groupoid-on-field}.
Soit $k'/k$ une extension de corps, $U' = \Spec(k')$,
et soit $(U', R', s', t', c')$ la restriction de
$(U, R, s, t, c)$ par $U' \to U$. Dans le diagramme de définition
$$
\xymatrix{
R' \ar[d] \ar[r] \ar@/_3pc/[dd]_{t'} \ar@/^1pc/[rr]^{s'} \ar@{..>}[rd] &
R \times_{s, U} U' \ar[r] \ar[d] &
U' \ar[d] \\
U' \times_{U, t} R \ar[d] \ar[r] &
R \ar[r]^s \ar[d]_t &
U \\
U' \ar[r] &
U
}
$$
tous les morphismes sont surjectifs, plats et universellement ouverts.
La flèche pointillée $R' \to R$ est de plus affine.
\end{lemma}

\begin{proof}
Le morphisme $U' \to U$ est $\Spec(k') \to \Spec(k)$;
il est donc affine, surjectif et plat. Les morphismes $s, t : R \to U$
et le morphisme $U' \to U$ sont universellement ouverts d'après
Morphismes, lemme \ref{morphisms-lemma-scheme-over-field-universally-open}.
Comme $R$ est non vide et $U$ est le spectre d'un corps, les morphismes
$s, t : R \to U$ sont surjectifs et plats. On conclut alors en utilisant
Morphismes d'espaces, lemmes
\ref{spaces-morphisms-lemma-base-change-surjective},
\ref{spaces-morphisms-lemma-composition-surjective},
\ref{spaces-morphisms-lemma-composition-open},
\ref{spaces-morphisms-lemma-base-change-affine},
\ref{spaces-morphisms-lemma-composition-affine},
\ref{spaces-morphisms-lemma-base-change-flat} et
\ref{spaces-morphisms-lemma-composition-flat}.
\end{proof}

\begin{lemma}
\label{spaces-more-groupoids-lemma-groupoid-on-field-explain-points}
Dans la
situation \ref{spaces-more-groupoids-situation-groupoid-on-field}.
Pour tout point $r \in |R|$, il existe
\begin{enumerate}
\item une extension de corps $k'/k$ avec $k'$ algébriquement clos,
\item un point $r' : \Spec(k') \to R'$, où
$(U', R', s', t', c')$ est la restriction de $(U, R, s, t, c)$
par $\Spec(k') \to \Spec(k)$,
\end{enumerate}
tels que
\begin{enumerate}
\item le point $r'$ s'envoie sur $r$ par le morphisme $R' \to R$, et
\item les applications
$s' \circ r', t' \circ r' : \Spec(k') \to \Spec(k')$
sont des automorphismes.
\end{enumerate}
\end{lemma}

\begin{proof}
Représentons $r$ par un morphisme $r : \Spec(K) \to R$ pour un certain
corps $K$. Pour démontrer le lemme, il faut trouver un corps
algébriquement clos $k'$ et un diagramme commutatif
$$
\xymatrix{
k' & k' \ar[l]^1 & \\
k' \ar[u]^\tau & K \ar[lu]^\sigma & k \ar[l]^-s \ar[lu]_i \\
& k \ar[lu]^i \ar[u]_t
}
$$
où $s, t : k \to K$ sont les homomorphismes de corps provenant de
$s \circ r$ et $t \circ r$. La démonstration de
Compléments sur les groupoïdes,
lemme \ref{more-groupoids-lemma-groupoid-on-field-explain-points},
explique comment construire un tel diagramme.
\end{proof}

\begin{lemma}
\label{spaces-more-groupoids-lemma-groupoid-on-field-move-point}
Dans la
situation \ref{spaces-more-groupoids-situation-groupoid-on-field}.
Si $r : \Spec(k) \to R$ est un morphisme tel que
$s \circ r, t \circ r$ soient des automorphismes de $\Spec(k)$, alors l'application
$$
R \longrightarrow R, \quad
x \longmapsto c(r, x)
$$
est un automorphisme $R \to R$ qui envoie $e$ sur $r$.
\end{lemma}

\begin{proof}
La démonstration est identique à celle de
Compléments sur les groupoïdes,
lemme \ref{more-groupoids-lemma-groupoid-on-field-move-point}.
\end{proof}

\begin{lemma}
\label{spaces-more-groupoids-lemma-groupoid-on-field-geometrically-irreducible}
Dans la
situation \ref{spaces-more-groupoids-situation-groupoid-on-field},
l'espace algébrique $R$ est géométriquement unibranche. Dans la
situation \ref{spaces-more-groupoids-situation-group-over-field},
l'espace algébrique $G$ est géométriquement unibranche.
\end{lemma}

\begin{proof}
Soit $r \in |R|$. Il faut montrer que $R$ est géométriquement unibranche
en $r$. En combinant le
lemme \ref{spaces-more-groupoids-lemma-restrict-groupoid-on-field}
avec
Descente sur les espaces, lemme \ref{spaces-descent-lemma-descend-unibranch},
on voit qu'il suffit de le démontrer lorsque $k$ est algébriquement clos
et que $r$ provient d'un morphisme $r : \Spec(k) \to R$ tel que
$s \circ r$ et $t \circ r$
soient des automorphismes de $\Spec(k)$. D'après le
lemme \ref{spaces-more-groupoids-lemma-groupoid-on-field-move-point},
on se ramène au cas où $r = e$ est l'identité de $R$ et où $k$ est
algébriquement clos.

\medskip\noindent
Supposons que $r = e$ et que $k$ soit algébriquement clos. Soit
$A = \mathcal{O}_{R, e}$ l'anneau local étale de
$R$ en $e$, et soit
$C = \mathcal{O}_{R \times_{s, U, t} R, (e, e)}$
l'anneau local étale de $R \times_{s, U, t} R$ en $(e, e)$.
D'après Compléments d'algèbre, lemme
\ref{more-algebra-lemma-minimal-primes-tensor-strictly-henselian}
les idéaux premiers minimaux $\mathfrak q$ de $C$ correspondent $1$ à $1$
aux couples de premiers minimaux $\mathfrak p, \mathfrak p' \subset A$.
D'autre part, la loi de composition induit un homomorphisme plat d'anneaux
$$
\xymatrix{
A \ar[r]_{c^\sharp} & C & \mathfrak q \\
& A \otimes_{s^\sharp, k, t^\sharp} A \ar[u] &
\mathfrak p \otimes A + A \otimes \mathfrak p' \ar@{|}[u]
}
$$
Remarquons que $(c^\sharp)^{-1}(\mathfrak q)$ contient à la fois $\mathfrak p$ et
$\mathfrak p'$, puisque les diagrammes
$$
\xymatrix{
A \ar[r]_{c^\sharp} & C \\
A \otimes_{s^\sharp, k} k \ar[u] &
A \otimes_{s^\sharp, k, t^\sharp} A \ar[l]_{1 \otimes e^\sharp} \ar[u]
}
\quad\quad
\xymatrix{
A \ar[r]_{c^\sharp} & C \\
k \otimes_{k, t^\sharp} A \ar[u] &
A \otimes_{s^\sharp, k, t^\sharp} A \ar[l]_{e^\sharp \otimes 1} \ar[u]
}
$$
sont commutatifs d'après (\ref{spaces-more-groupoids-equation-diagram}).
Comme $c^\sharp$ est plat (puisque $c$ est un morphisme plat d'après le
lemme \ref{spaces-more-groupoids-lemma-groupoid-on-field-open-multiplication}),
on voit que $(c^\sharp)^{-1}(\mathfrak q)$ est un idéal premier minimal
de $A$. Ainsi $\mathfrak p = (c^\sharp)^{-1}(\mathfrak q) = \mathfrak p'$.
\end{proof}

\noindent
Dans le lemme suivant, nous utilisons la dimension des espaces algébriques
(en un point), telle qu'elle est définie dans
Propriétés des espaces, section \ref{spaces-properties-section-dimension}.
Nous utilisons aussi la dimension de l'anneau local définie dans
Propriétés des espaces, section
\ref{spaces-properties-section-dimension-local-ring},
ainsi que le degré de transcendance des points; voir
Morphismes d'espaces, section \ref{spaces-morphisms-section-relative-dimension}.

\begin{lemma}
\label{spaces-more-groupoids-lemma-groupoid-on-field-locally-finite-type-dimension}
Dans la
situation \ref{spaces-more-groupoids-situation-groupoid-on-field},
supposons $s, t$ localement de type fini.
Pour tout $r \in |R|$,
\begin{enumerate}
\item $\dim(R) = \dim_r(R)$,
\item le degré de transcendance de $r$ sur $\Spec(k)$
par $s$ est égal au degré de transcendance de $r$ sur $\Spec(k)$
par $t$, et
\item si le degré de transcendance mentionné en (2) est $0$, alors
$\dim(R) = \dim(\mathcal{O}_{R, \overline{r}})$.
\end{enumerate}
\end{lemma}

\begin{proof}
Soit $r \in |R|$. Notons $\text{trdeg}(r/_{\!\! s}k)$ le degré de
transcendance de $r$ sur $\Spec(k)$ par $s$. Choisissons un morphisme étale
$\varphi : V \to R$, où $V$ est un schéma et $v \in V$ s'envoie sur $r$.
Les définitions rappelées avant le lemme donnent
$$
\dim_r(R) = \dim_v(V) =
\dim(\mathcal{O}_{V, v}) + \text{trdeg}_{s(k)}(\kappa(v)) =
\dim(\mathcal{O}_{R, \overline{r}}) + \text{trdeg}(r/_{\!\! s}k)
$$
et de même pour $t$ (la deuxième égalité provenant de
Morphismes, lemme \ref{morphisms-lemma-dimension-fibre-at-a-point}).
On voit donc que $\text{trdeg}(r/_{\!\! s}k) = \text{trdeg}(r/_{\!\! t}k)$,
c'est-à-dire que (2) est vérifié.

\medskip\noindent
Soit $k'/k$ une extension de corps. Remarquons que la restriction $R'$
de $R$ à $\Spec(k')$ (voir le
lemme \ref{spaces-more-groupoids-lemma-restrict-groupoid-on-field})
s'obtient à partir de $R$ par deux changements de base par des morphismes de corps. Ainsi,
Morphismes d'espaces,
lemme \ref{spaces-morphisms-lemma-dimension-fibre-after-base-change},
montre que la dimension de $R$ en un point reste inchangée par cette opération.
Pour démontrer (1), on peut donc supposer, d'après le
lemme \ref{spaces-more-groupoids-lemma-groupoid-on-field-explain-points},
que $r$ est représenté par un morphisme $r : \Spec(k) \to R$ tel
que $s \circ r$ et $t \circ r$ soient tous deux des automorphismes de $\Spec(k)$.
Il existe alors un automorphisme $R \to R$ qui envoie $r$ sur $e$
(lemme \ref{spaces-more-groupoids-lemma-groupoid-on-field-move-point}).
On voit donc que $\dim_r(R) = \dim_e(R)$ pour tout $r$. Par définition, cela
signifie que $\dim_r(R) = \dim(R)$.

\medskip\noindent
Le point (3) est une conséquence formelle des résultats obtenus dans la
discussion précédente.
\end{proof}

\begin{lemma}
\label{spaces-more-groupoids-lemma-group-over-field-locally-finite-type-dimension}
Dans la
situation \ref{spaces-more-groupoids-situation-group-over-field},
supposons $G$ localement de type fini.
Pour tout $g \in |G|$,
\begin{enumerate}
\item $\dim(G) = \dim_g(G)$,
\item si le degré de transcendance de $g$ sur $k$ est $0$, alors
$\dim(G) = \dim(\mathcal{O}_{G, \overline{g}})$.
\end{enumerate}
\end{lemma}

\begin{proof}
Cela résulte immédiatement du
lemme \ref{spaces-more-groupoids-lemma-groupoid-on-field-locally-finite-type-dimension}
par (\ref{spaces-more-groupoids-equation-groupoid-from-group}).
\end{proof}

\begin{lemma}
\label{spaces-more-groupoids-lemma-groupoid-on-field-dimension-equal-stabilizer}
Dans la
situation \ref{spaces-more-groupoids-situation-groupoid-on-field},
supposons $s, t$ localement de type fini.
Soit
$G = \Spec(k)
\times_{\Delta, \Spec(k) \times_B \Spec(k), t \times s} R$
l'espace algébrique en groupes stabilisateur.
Alors $\dim(R) = \dim(G)$.
\end{lemma}

\begin{proof}
Comme $G$ et $R$ sont équidimensionnels (voir les
lemmes \ref{spaces-more-groupoids-lemma-groupoid-on-field-locally-finite-type-dimension} et
\ref{spaces-more-groupoids-lemma-group-over-field-locally-finite-type-dimension})
il suffit de montrer que $\dim_e(R) = \dim_e(G)$. Soit $V$ un schéma affine,
soit $v \in V$, et soit $\varphi : V \to R$ un morphisme étale de schémas
tel que $\varphi(v) = e$. Remarquons que $V$ est noethérien, car
$s \circ \varphi$ est localement de type fini comme composée de morphismes
localement de type fini et parce que $V$ est quasi-compact (utiliser
Morphismes d'espaces, lemmes
\ref{spaces-morphisms-lemma-composition-finite-type},
\ref{spaces-morphisms-lemma-etale-locally-finite-presentation}, et
\ref{spaces-morphisms-lemma-finite-presentation-finite-type}
et
Morphismes, lemme \ref{morphisms-lemma-finite-type-noetherian}).
Ainsi, $V$ est localement connexe (voir
Propriétés, lemme \ref{properties-lemma-Noetherian-topology},
et
Topologie, lemme \ref{topology-lemma-locally-Noetherian-locally-connected}).
On peut donc remplacer $V$ par la composante connexe contenant $v$ (elle
reste affine, étant un sous-schéma ouvert et fermé de $V$).
Posons $T = V_{red}$, la réduction de $V$. Considérons les deux
morphismes $a, b : T \to \Spec(k)$ donnés par
$a = s \circ \varphi|_T$ et $b = t \circ \varphi|_T$. Remarquons que
$a, b$ induisent le même homomorphisme de corps $k \to \kappa(v)$, car $\varphi(v) = e$ !
Soit $k_a \subset \Gamma(T, \mathcal{O}_T)$ la clôture intégrale de
$a^\sharp(k) \subset \Gamma(T, \mathcal{O}_T)$. De même, soit
$k_b \subset \Gamma(T, \mathcal{O}_T)$ la clôture intégrale de
$b^\sharp(k) \subset \Gamma(T, \mathcal{O}_T)$. D'après
Variétés, proposition \ref{varieties-proposition-unique-base-field},
on a $k_a = k_b$. On obtient ainsi le diagramme commutatif suivant :
$$
\xymatrix{
k \ar[rd]^a \ar[rrrd] \\
& k_a = k_b \ar[r] & \Gamma(T, \mathcal{O}_T) \ar[r] & \kappa(v) \\
k \ar[ru]_b \ar[rrru]
}
$$
Comme on l'a vu ci-dessus, les longues flèches sont égales.
Puisque $k_a = k_b \to \kappa(v)$ est injectif, on en conclut que
les deux morphismes $a$ et $b$ sont égaux. Ainsi $T \to R$ se factorise par $G$.
Il s'ensuit que $R_{red} = G_{red}$ dans un voisinage ouvert de $e$,
ce qui implique bien que $\dim_e(R) = \dim_e(G)$.
\end{proof}






\section{Espaces algébriques en groupes sur les corps}
\label{spaces-more-groupoids-section-group-algebraic-spaces-over-fields}

\noindent
Il existe un espace algébrique en groupes non séparé sur un corps,
à savoir $\mathbf{G}_a/\mathbf{Z}$ sur un corps de caractéristique zéro; voir
Exemples, section \ref{examples-section-non-separated-group-space}.
En fait, tout schéma en groupes sur un corps est séparé
(lemme \ref{spaces-more-groupoids-lemma-group-scheme-over-field-separated});
ainsi, tout espace algébrique en groupes non séparé sur un corps
n'est pas représentable. D'autre part, un espace algébrique en
groupes sur un corps est séparé dès qu'il est décent; voir le
lemme \ref{spaces-more-groupoids-lemma-group-scheme-over-field-separated}.
Dans cette section, nous allons montrer qu'un
espace algébrique en groupes séparé sur un corps
est représentable, c'est-à-dire qu'il est un schéma.

\begin{lemma}
\label{spaces-more-groupoids-lemma-group-space-scheme-over-kbar}
Soit $k$ un corps muni d'une clôture algébrique $\overline{k}$.
Soit $G$ un espace algébrique en groupes sur $k$
qui est séparé\footnote{Il suffit de supposer $G$ décent,
par exemple localement séparé ou quasi-séparé d'après le
lemme \ref{spaces-more-groupoids-lemma-group-scheme-over-field-separated}.}.
Alors $G_{\overline{k}}$ est un schéma.
\end{lemma}

\begin{proof}
D'après Espaces sur les corps, lemme
\ref{spaces-over-fields-lemma-when-scheme-after-base-change}
il suffit de montrer que $G_K$ est un schéma pour une certaine extension de corps
$K/k$. Notons $G_K' \subset G_K$ le lieu schématique
de $G_K$, comme dans Propriétés des espaces, lemme
\ref{spaces-properties-lemma-subscheme}.
D'après Propriétés des espaces, proposition
\ref{spaces-properties-proposition-locally-quasi-separated-open-dense-scheme}
on voit que $G_K' \subset G_K$ est un ouvert dense, en particulier non vide.
Choisissons un schéma $U$ et un morphisme étale surjectif $U \to G$. D'après
Variétés, lemme \ref{varieties-lemma-make-Jacobson},
si $K$ est un corps algébriquement clos de degré de transcendance assez
grand, alors $U_K$ est un schéma de Jacobson et
tout point fermé de $U_K$ est $K$-rationnel.
Ainsi $G_K'$ possède un point $K$-rationnel et il suffit
de montrer que tout point $K$-rationnel de $G_K$ appartient à $G_K'$.
Si $g \in G_K(K)$ est un point $K$-rationnel et $g' \in G_K'(K)$
un point $K$-rationnel du lieu schématique, alors on voit que
$g$ appartient à l'image de $G_K'$ par l'automorphisme
$$
G_K \longrightarrow G_K,\quad
h \longmapsto g(g')^{-1}h
$$
de $G_K$. Comme les automorphismes de l'espace algébrique $G_K$ préservent
$G_K'$, on conclut que $g \in G_K'$, comme voulu.
\end{proof}

\begin{lemma}
\label{spaces-more-groupoids-lemma-group-space-scheme-locally-finite-type-over-k}
Soit $k$ un corps. Soit $G$ un espace algébrique en groupes sur $k$.
Si $G$ est séparé et localement de type fini sur $k$,
alors $G$ est un schéma.
\end{lemma}

\begin{proof}
Cela résulte du
lemme \ref{spaces-more-groupoids-lemma-group-space-scheme-over-kbar},
de Groupoïdes, lemme \ref{groupoids-lemma-points-in-affine}, et
d'Espaces sur les corps, lemme
\ref{spaces-over-fields-lemma-scheme-over-algebraic-closure-enough-affines}.
\end{proof}

\begin{proposition}
\label{spaces-more-groupoids-proposition-group-space-scheme-over-field}
Soit $k$ un corps. Soit $G$ un espace algébrique en groupes sur $k$.
Si $G$ est séparé, alors $G$ est un schéma.
\end{proposition}

\begin{proof}
Ce lemme généralise le
lemme \ref{spaces-more-groupoids-lemma-group-space-scheme-locally-finite-type-over-k}
(qui couvre tous les cas utiles en pratique).
La démonstration est très semblable à celle d'
Espaces sur les corps, lemme
\ref{spaces-over-fields-lemma-scheme-over-algebraic-closure-enough-affines}
utilisé dans la démonstration du
lemme \ref{spaces-more-groupoids-lemma-group-space-scheme-locally-finite-type-over-k};
nous invitons le lecteur à commencer par lire cette démonstration.

\medskip\noindent
D'après le lemme \ref{spaces-more-groupoids-lemma-group-space-scheme-over-kbar}, le changement
de base $G_{\overline{k}}$ est un schéma.
Soit $K/k$ une extension purement transcendante de très grand
degré de transcendance. D'après Espaces sur les corps, lemme
\ref{spaces-over-fields-lemma-scheme-after-purely-transcendental-base-change}
il suffit de montrer que $G_K$ est un schéma.
Soit $K^{perf}$ la clôture parfaite de $K$. D'après
Espaces sur les corps, lemme
\ref{spaces-over-fields-lemma-scheme-after-purely-inseparable-base-change}
il suffit de montrer que $G_{K^{perf}}$ est un schéma.
Soit $K \subset K^{perf} \subset \overline{K}$, cette dernière étant la clôture algébrique
de $K$. On peut choisir un plongement
$\overline{k} \to \overline{K}$ au-dessus de $k$, de sorte que $G_{\overline{K}}$
soit le changement de base du schéma $G_{\overline{k}}$ par $\overline{k} \to \overline{K}$. D'après
Variétés, lemme \ref{varieties-lemma-make-Jacobson},
on voit que $G_{\overline{K}}$ est un schéma de Jacobson dont tous les
points fermés ont pour corps résiduel $\overline{K}$.

\medskip\noindent
Comme $G_{\overline{K}} \to G_{K^{perf}}$ est surjectif, il suffit de
montrer que l'image $g \in |G_{K^{perf}}|$ d'un point fermé arbitraire de
$G_{\overline{K}}$ appartient au lieu schématique de $G_K$.
En particulier, on peut représenter $g$ par un morphisme
$g : \Spec(L) \to G_{K^{perf}}$, où $L/K^{perf}$ est algébrique séparable
(on peut par exemple prendre $L = \overline{K}$). Le schéma
\begin{align*}
T & = \Spec(L) \times_{G_{K^{perf}}} G_{\overline{K}} \\
& =
\Spec(L) \times_{\Spec(K^{perf})} \Spec(\overline{K}) \\
& =
\Spec(L \otimes_{K^{perf}} \overline{K})
\end{align*}
est donc le spectre d'une $\overline{K}$-algèbre qui est une limite inductive
filtrante d'algèbres, produits finis de copies de $\overline{K}$.
Ainsi, d'après Groupoïdes, lemme \ref{groupoids-lemma-compact-set-in-affine},
on peut trouver un ouvert affine $W \subset G_{\overline{K}}$ contenant
l'image de $g_{\overline{K}} : T \to G_{\overline{K}}$.

\medskip\noindent
Choisissons un ouvert quasi-compact $V \subset G_{K^{perf}}$ contenant
l'image de $W$. D'après Espaces sur les corps,
lemme \ref{spaces-over-fields-lemma-when-scheme-after-base-change},
on voit que $V_{K'}$ est un schéma pour une certaine extension finie $K'/K^{perf}$.
En agrandissant $K'$, on peut supposer qu'il existe un ouvert affine
$U' \subset V_{K'} \subset G_{K'}$ dont le changement de base à $\overline{K}$
redonne $W$ (utiliser que $V_{\overline{K}}$ est la limite des schémas
$V_{K''}$ pour $K' \subset K'' \subset \overline{K}$ fini, ainsi que
Limites, lemmes \ref{limits-lemma-descend-opens} et
\ref{limits-lemma-limit-affine}). On peut supposer
que $K'/K^{perf}$ est une extension galoisienne (prendre la clôture normale,
Corps, lemme \ref{fields-lemma-normal-closure}, et utiliser
que $K^{perf}$ est parfait). Posons $H = \text{Gal}(K'/K^{perf})$.
Par construction, le sous-schéma fermé $H$-invariant
$\Spec(L) \times_{G_{K^{perf}}} G_{K'}$ est contenu dans $U'$.
D'après
Espaces sur les corps, lemmes
\ref{spaces-over-fields-lemma-base-change-by-Galois} et
\ref{spaces-over-fields-lemma-when-quotient-scheme-at-point}, on conclut.
\end{proof}





\section{Absence de courbes rationnelles sur les groupes}
\label{spaces-more-groupoids-section-no-rational-curves}

\noindent
Dans cette section, nous montrons qu'il n'existe pas de morphisme non constant
de $\mathbf{P}^1$ vers un espace algébrique en groupes localement de type
fini sur un corps.

\begin{lemma}
\label{spaces-more-groupoids-lemma-factor-through-over-open}
Soit $S$ un schéma. Soit $B$ un espace algébrique au-dessus de $S$.
Soient $f : X \to Y$ et $g : X \to Z$ des morphismes d'espaces
algébriques au-dessus de $B$. Supposons que
\begin{enumerate}
\item $Y \to B$ soit séparé,
\item $g$ soit surjectif, plat et localement de présentation finie,
\item il existe un ouvert schématiquement dense $V \subset Z$
tel que $f|_{g^{-1}(V)} : g^{-1}(V) \to Y$ se factorise par $V$.
\end{enumerate}
Alors $f$ se factorise par $g$.
\end{lemma}

\begin{proof}
Posons $R = X \times_Z X$. D'après (2), on a $Z = X/R$ comme faisceaux.
De plus, (2) implique que l'image inverse de $V$ dans $R$ est
schématiquement dense dans $R$ (Morphismes d'espaces, lemme
\ref{spaces-morphisms-lemma-flat-morphism-scheme-theoretically-dense-open}).
On voit alors que les deux composées
$R \to X \to Y$ sont égales d'après Morphismes d'espaces, lemme
\ref{spaces-morphisms-lemma-equality-of-morphisms}.
Le lemme en résulte.
\end{proof}

\begin{lemma}
\label{spaces-more-groupoids-lemma-quotient-power-P1}
\begin{slogan}
Un morphisme d'un produit non vide de droites projectives sur un corps vers
un espace algébrique séparé de type fini sur un corps se factorise en un
morphisme fini après projection sur un produit de droites projectives.
\end{slogan}
Soit $k$ un corps. Soit $n \geq 1$ et soit $(\mathbf{P}^1_k)^n$
le produit de $n$ copies au-dessus de $\Spec(k)$. Soit
$f : (\mathbf{P}^1_k)^n \to Z$ un morphisme d'espaces algébriques sur $k$.
Si $Z$ est séparé et de type fini sur $k$, alors $f$ se factorise sous la forme
$$
(\mathbf{P}^1_k)^n \xrightarrow{projection}
(\mathbf{P}^1_k)^m \xrightarrow{fini} Z.
$$
\end{lemma}

\begin{proof}
On peut supposer $k$ algébriquement clos (détails omis); nous ne le faisons
que pour pouvoir raisonner avec des points rationnels, mais le lecteur peut
contourner cet artifice s'il le souhaite. Dans la preuve, les produits sont pris sur $k$.
L'espace algébrique en groupes des automorphismes de $(\mathbf{P}^1_k)^n$ contient
$G = (\text{GL}_{2, k})^n$. Si $C \subset (\mathbf{P}^1_k)^n$ est une
sous-variété fermée (donc en particulier irréductible sur $k$) envoyée
sur un point, on peut appliquer
Compléments sur les morphismes d'espaces,
lemme \ref{spaces-more-morphisms-lemma-flat-proper-family-cannot-collapse-fibre}
au morphisme
$$
G \times C \to G \times Z,\quad (g, c) \mapsto (g, f(g \cdot c))
$$
au-dessus de $G$. Ainsi, $g(C)$ est envoyé sur un point pour $g \in G(k)$
appartenant à un ouvert de Zariski $U \subset G$. Supposons que
$x = (x_1, \ldots, x_n)$, $y = (y_1, \ldots, y_n)$
soient des points à valeurs dans $k$ de $(\mathbf{P}^1_k)^n$. Soit
$I \subset \{1, \ldots, n\}$ l'ensemble des indices $i$
tels que $x_i = y_i$. Alors
$$
\{g(x) \mid g(y) = y,\ g \in U(k)\}
$$
est dense pour la topologie de Zariski dans la fibre de la projection
$\pi_I : (\mathbf{P}^1_k)^n \to \prod_{i \in I} \mathbf{P}^1_k$
(exercice). Ainsi, si $x, y \in C(k)$ sont distincts, on conclut
que $f$ envoie sur un seul point toute la fibre de $\pi_I$ contenant $x, y$.
De plus, l'orbite de $C$ sous $U(k)$ rencontre un ensemble ouvert de Zariski
de fibres de $\pi_I$. D'après le lemme \ref{spaces-more-groupoids-lemma-factor-through-over-open},
le morphisme $f$ se factorise par $\pi_I$.
Après avoir répété ce procédé un nombre fini de fois, on atteint
le stade où toutes les fibres de $f$ au-dessus des points de $k$ sont finies.
Dans ce cas, $f$ est fini d'après
Compléments sur les morphismes d'espaces, lemme
\ref{spaces-more-morphisms-lemma-proper-finite-fibre-finite-in-neighbourhood}
et le fait que les points de $k$ sont denses dans $Z$
(Espaces sur les corps, lemme
\ref{spaces-over-fields-lemma-smooth-separable-closed-points-dense}).
\end{proof}

\begin{lemma}
\label{spaces-more-groupoids-lemma-no-nonconstant-morphism-from-P1-to-group}
\begin{slogan}
Pas de courbes rationnelles complètes sur les groupes.
\end{slogan}
Soit $k$ un corps. Soit $G$ un espace algébrique en groupes séparé,
localement de type fini sur $k$. Il n'existe aucun
morphisme non constant $f : \mathbf{P}^1_k \to G$ au-dessus de $\Spec(k)$.
\end{lemma}

\begin{proof}
Supposons $f$ non constant. Considérons les morphismes
$$
\mathbf{P}^1_k \times_{\Spec(k)} \ldots \times_{\Spec(k)} \mathbf{P}^1_k
\longrightarrow G,
\quad (t_1, \ldots, t_n) \longmapsto f(t_1) \ldots f(t_n)
$$
où l'on utilise au membre de droite la multiplication du groupe.
D'après le lemme \ref{spaces-more-groupoids-lemma-quotient-power-P1} et l'hypothèse que $f$
est non constant, ce morphisme est fini sur son image.
Ainsi $\dim(G) \geq n$ pour tout $n$, ce qui est impossible d'après le
lemme \ref{spaces-more-groupoids-lemma-group-over-field-locally-finite-type-dimension}
et le fait que $G$ est localement de type fini sur $k$.
\end{proof}




\section{La partie finie d'un morphisme}
\label{spaces-more-groupoids-section-finite}

\noindent
Soit $S$ un schéma.
Soit $f : X \to Y$ un morphisme d'espaces algébriques au-dessus de $S$.
Pour un espace algébrique ou un schéma $T$ au-dessus de $S$, considérons les couples
$(a, Z)$ où
\begin{equation}
\label{spaces-more-groupoids-equation-finite-conditions}
\begin{matrix}
a : T \to Y\text{ est un morphisme au-dessus de }S, \\
Z \subset T \times_Y X\text{ est un sous-espace ouvert} \\
\text{tel que }\text{pr}_0|_Z : Z \to T\text{ soit fini.}
\end{matrix}
\end{equation}
Supposons que $h : T' \to T$ soit un morphisme d'espaces algébriques au-dessus de $S$
et que $(a, Z)$ soit un couple comme dans (\ref{spaces-more-groupoids-equation-finite-conditions}) au-dessus de $T$. Posons
$a' = a \circ h$ et $Z' = (h \times \text{id}_X)^{-1}(Z) = T' \times_T Z$.
Alors $(a', Z')$ est un couple comme dans (\ref{spaces-more-groupoids-equation-finite-conditions}) au-dessus de $T'$.
Cela résulte de ce que les morphismes finis sont préservés par changement de base; voir
Morphismes d'espaces, lemme \ref{spaces-morphisms-lemma-base-change-integral}.
On obtient ainsi un foncteur
\begin{equation}
\label{spaces-more-groupoids-equation-finite}
\begin{matrix}
(X/Y)_{fin} : &
(\Sch/S)^{opp} &
\longrightarrow &
\textit{Ensembles} \\
& T & \longmapsto &
\{(a, Z)\text{ comme ci-dessus}\}
\end{matrix}
\end{equation}
Pour les applications, nous nous intéressons surtout à ce foncteur $(X/Y)_{fin}$
lorsque $f$ est séparé et localement de type fini. Pour se faire une idée
de la construction, consulter la
remarque \ref{spaces-more-groupoids-remark-finite-quasi-finite-separated-morphism-schemes}.

\begin{lemma}
\label{spaces-more-groupoids-lemma-finite-sheaf}
Soit $S$ un schéma.
Soit $f : X \to Y$ un morphisme d'espaces algébriques au-dessus de $S$.
Alors
\begin{enumerate}
\item Le préfaisceau $(X/Y)_{fin}$ satisfait la condition de faisceau pour
la topologie fppf.
\item Si $T$ est un espace algébrique au-dessus de $S$, il existe une
bijection canonique
$$
\Mor_{\Sh((\Sch/S)_{fppf})}(T, (X/Y)_{fin})
=
\{(a, Z)\text{ satisfaisant \ref{spaces-more-groupoids-equation-finite-conditions}}\}
$$
\end{enumerate}
\end{lemma}

\begin{proof}
Soit $T$ un espace algébrique au-dessus de $S$.
Soit $\{T_i \to T\}$ un recouvrement fppf (par des espaces algébriques).
Soient $s_i = (a_i, Z_i)$ des couples au-dessus de $T_i$
satisfaisant \ref{spaces-more-groupoids-equation-finite-conditions}
tels que $s_i|_{T_i \times_T T_j} = s_j|_{T_i \times_T T_j}$.
Premièrement, cela implique notamment que $a_i$ et $a_j$ définissent le même
morphisme $T_i \times_T T_j \to Y$. D'après
Descente sur les espaces,
lemme \ref{spaces-descent-lemma-fpqc-universal-effective-epimorphisms},
on en déduit qu'il existe un unique morphisme $a : T \to Y$
tel que $a_i$ soit la composée $T_i \to T \to Y$.
Deuxièmement, cela implique que les $Z_i \subset T_i \times_Y X$ sont des sous-espaces
ouverts dont les images inverses dans $(T_i \times_T T_j) \times_Y X$ sont égales.
Comme $\{T_i \times_Y X \to T \times_Y X\}$ est un recouvrement fppf,
on en déduit qu'il existe un unique sous-espace ouvert $Z \subset T \times_Y X$
dont la restriction à $T_i$ est $Z_i$; voir
Descente sur les espaces, lemme \ref{spaces-descent-lemma-open-fpqc-covering}.
Nous affirmons que la projection $Z \to T$ est finie.
Cela résulte de ce que la finitude est locale pour la topologie fpqc; voir
Descente sur les espaces, lemme \ref{spaces-descent-lemma-descending-property-finite}.

\medskip\noindent
Remarquons que le résultat du paragraphe précédent implique notamment (1).

\medskip\noindent
Soit $T$ un espace algébrique au-dessus de $S$. Pour démontrer (2), nous allons
construire des applications réciproques entre les ensembles affichés. Dans la
suite, le mot « couple » désigne un couple satisfaisant les
conditions \ref{spaces-more-groupoids-equation-finite-conditions}.

\medskip\noindent
Soit $v : T \to (X/Y)_{fin}$ une transformation naturelle.
Choisissons un schéma $U$ et un morphisme étale surjectif $p : U \to T$.
Alors $v(p) \in (X/Y)_{fin}(U)$ correspond à un couple $(a_U, Z_U)$
au-dessus de $U$. Soit $R = U \times_T U$, de projections $t, s : R \to U$.
Comme $v$ est une transformation de foncteurs, les images inverses de
$(a_U, Z_U)$ par $s$ et $t$ sont égales. Puisque $\{U \to T\}$ est un
recouvrement fppf, on peut donc appliquer le résultat du premier paragraphe et
en déduire qu'il existe un unique couple $(a, Z)$ au-dessus de $T$.

\medskip\noindent
Réciproquement, soit $(a, Z)$ un couple au-dessus de $T$.
Soient $U \to T$, $R = U \times_T U$ et $t, s : R \to U$ comme
ci-dessus. La restriction $(a, Z)|_U$ donne alors naissance à une
transformation de foncteurs $v : h_U \to (X/Y)_{fin}$ par le
lemme de Yoneda
(Catégories, lemme \ref{categories-lemma-yoneda}).
Comme les deux images inverses $s^*(a, Z)|_U$ et $t^*(a, Z)|_U$
sont égales, on voit que $v$ coégalise les deux applications
$h_t, h_s : h_R \to h_U$. Puisque $T = U/R$ est le faisceau quotient fppf d'après
Espaces, lemme \ref{spaces-lemma-space-presentation},
et que $(X/Y)_{fin}$ est un faisceau fppf d'après (1), on conclut
que $v$ se factorise par une application $T \to (X/Y)_{fin}$.

\medskip\noindent
Nous omettons la vérification que les deux constructions ci-dessus sont
réciproques.
\end{proof}

\begin{lemma}
\label{spaces-more-groupoids-lemma-finite-open}
Soit $S$ un schéma. Considérons un diagramme commutatif
$$
\xymatrix{
X' \ar[rr]_j \ar[rd] & & X \ar[ld] \\
& Y
}
$$
d'espaces algébriques au-dessus de $S$. Si $j$ est une immersion ouverte, alors
il existe une application injective canonique de faisceaux
$j : (X'/Y)_{fin} \to (X/Y)_{fin}$.
\end{lemma}

\begin{proof}
Si $(a, Z)$ est un couple au-dessus de $T$ pour $X'/Y$, alors
$(a, j(Z))$ est un couple au-dessus de $T$ pour $X/Y$.
\end{proof}

\begin{lemma}
\label{spaces-more-groupoids-lemma-finite-lives-on-locally-quasi-finite-part}
Soit $S$ un schéma.
Soit $f : X \to Y$ un morphisme d'espaces algébriques au-dessus de $S$ qui est
localement de type fini.
Soit $X' \subset X$ le plus grand sous-espace ouvert sur lequel $f$ est
localement quasi-fini; voir
Morphismes d'espaces,
lemme \ref{spaces-morphisms-lemma-locally-finite-type-quasi-finite-part}.
Alors $(X/Y)_{fin} = (X'/Y)_{fin}$.
\end{lemma}

\begin{proof}
Le lemme \ref{spaces-more-groupoids-lemma-finite-open}
fournit une application injective $(X'/Y)_{fin} \to (X/Y)_{fin}$.
Morphismes d'espaces,
lemme \ref{spaces-morphisms-lemma-locally-finite-type-quasi-finite-part},
assure que la formation de $X'$ commute au changement de base.
Il suffit donc de montrer que, si
$Z \subset X$ est un sous-espace ouvert tel que $f|_Z : Z \to Y$ soit fini,
alors $Z \subset X'$. Cela est vrai puisqu'un morphisme fini
est localement quasi-fini; voir
Morphismes d'espaces, lemme \ref{spaces-morphisms-lemma-finite-quasi-finite}.
\end{proof}

\begin{lemma}
\label{spaces-more-groupoids-lemma-finite-separated}
Soit $S$ un schéma.
Soit $f : X \to Y$ un morphisme d'espaces algébriques au-dessus de $S$.
Soit $T$ un espace algébrique au-dessus de $S$, et soit $(a, Z)$
un couple comme dans \ref{spaces-more-groupoids-equation-finite-conditions}.
Si $f$ est séparé, alors $Z$ est fermé dans $T \times_Y X$.
\end{lemma}

\begin{proof}
Un morphisme fini d'espaces algébriques est universellement fermé d'après
Morphismes d'espaces, lemme \ref{spaces-morphisms-lemma-finite-proper}.
Comme $f$ est séparé, le morphisme $T \times_Y X \to T$ l'est aussi; voir
Morphismes d'espaces, lemme \ref{spaces-morphisms-lemma-base-change-separated}.
Le caractère fermé de $Z$ résulte donc de
Morphismes d'espaces,
lemme \ref{spaces-morphisms-lemma-universally-closed-permanence}.
\end{proof}

\begin{remark}
\label{spaces-more-groupoids-remark-finite-monoid}
Soit $f : X \to Y$ un morphisme séparé d'espaces algébriques.
Le faisceau $(X/Y)_{fin}$ est muni d'une application naturelle
$(X/Y)_{fin} \to Y$ qui envoie le couple $(a, Z) \in (X/Y)_{fin}(T)$
sur l'élément $a \in Y(T)$. On peut utiliser le
lemme \ref{spaces-more-groupoids-lemma-finite-separated}
pour définir les opérations
$$
\star_i : (X/Y)_{fin} \times_Y (X/Y)_{fin} \longrightarrow (X/Y)_{fin}
$$
par les règles
\begin{align*}
\star_1 : ((a, Z_1), (a, Z_2)) & \longmapsto (a, Z_1 \cup Z_2) \\
\star_2 : ((a, Z_1), (a, Z_2)) & \longmapsto (a, Z_1 \cap Z_2) \\
\star_3 : ((a, Z_1), (a, Z_2)) & \longmapsto (a, Z_1 \setminus Z_2) \\
\star_4 : ((a, Z_1), (a, Z_2)) & \longmapsto (a, Z_2 \setminus Z_1).
\end{align*}
Cette construction fonctionne parce que $Z_1 \cap Z_2$ est à la fois ouvert et fermé
dans $Z_1$ et $Z_2$ (ce qui implique aussi que $Z_1 \cup Z_2$ est
la réunion disjointe des trois autres morceaux).
On peut donc considérer $(X/Y)_{fin}$ comme une $\mathbf{F}_2$-algèbre
(sans unité) au-dessus de $Y$, dont la multiplication est donnée par
$ss' = \star_2(s, s')$ et l'addition par
$$
s + s' = \star_1(\star_3(s, s'), \star_4(s, s'))
$$
ce qui revient à prendre la différence symétrique.
Remarquons que, dans ce faisceau d'algèbres, $0 = (1_Y, \emptyset)$
et que l'on a bien $s + s = 0$ pour toute section locale $s$.
Si $f : X \to Y$ est fini, cette algèbre possède une unité, à savoir
$1 = (1_Y, X)$, et $\star_3(s, s') = s(1 + s')$, tandis que
$\star_4(s, s') = (1 + s)s'$.
\end{remark}

\begin{remark}
\label{spaces-more-groupoids-remark-finite-quasi-finite-separated-morphism-schemes}
Soit $f : X \to Y$ un morphisme de schémas séparé et localement
quasi-fini. Dans ce cas, le faisceau $(X/Y)_{fin}$
est étroitement lié au faisceau $f_!\mathbf{F}_2$
(insérer ici une référence future) sur $Y_\etale$.
En effet, si $V \to Y$ est étale et $s \in \Gamma(V, f_!\mathbf{F}_2)$,
alors $s \in \Gamma(V \times_Y X, \mathbf{F}_2)$ est une section
de support propre $Z = \text{Supp}(s)$ sur $V$. Comme $f$ est
aussi localement quasi-fini, la projection $Z \to V$ est en fait
finie. Comme le support d'une section d'un faisceau abélien constant est ouvert,
on voit que le couple $(V \to Y, \text{Supp}(s))$ satisfait
\ref{spaces-more-groupoids-equation-finite-conditions}.
En fait, $f_!\mathbf{F}_2 \cong (X/Y)_{fin}|_{Y_\etale}$
dans ce cas, ce qui explique aussi la structure de $\mathbf{F}_2$-algèbre
introduite dans la remarque \ref{spaces-more-groupoids-remark-finite-monoid}.
\end{remark}

\begin{lemma}
\label{spaces-more-groupoids-lemma-finite-diagonal}
Soit $S$ un schéma.
Soit $f : X \to Y$ un morphisme d'espaces algébriques au-dessus de $S$.
La diagonale de $(X/Y)_{fin} \to Y$
$$
(X/Y)_{fin} \longrightarrow (X/Y)_{fin} \times_Y (X/Y)_{fin}
$$
est représentable (par des schémas) et est une immersion ouverte; la diagonale
« absolue »
$$
(X/Y)_{fin} \longrightarrow (X/Y)_{fin} \times (X/Y)_{fin}
$$
est représentable (par des schémas).
\end{lemma}

\begin{proof}
La seconde assertion résulte de la première, car la diagonale absolue
est la composée de la diagonale relative et d'un changement de base
de la diagonale de $Y$ (qui est représentable par des schémas); voir
Espaces, section \ref{spaces-section-representable}.
Pour démontrer la première assertion, il faut établir ce qui suit.
Étant donnés un schéma $T$ et deux couples $(a, Z_1)$ et $(a, Z_2)$ au-dessus de $T$
ayant la même première composante et
satisfaisant \ref{spaces-more-groupoids-equation-finite-conditions},
il existe un sous-schéma ouvert $V \subset T$ possédant la propriété suivante :
pour tout morphisme de schémas $h : T' \to T$, on a
$$
h(T') \subset V \Leftrightarrow
\Big(T' \times_T Z_1 = T' \times_T Z_2
\text{ comme sous-espaces de }T' \times_Y X\Big)
$$
Construisons $V$. Remarquons que $Z_1 \cap Z_2$ est ouvert dans $Z_1$
et dans $Z_2$. Comme $\text{pr}_0|_{Z_i} : Z_i \to T$ est fini,
donc propre (voir
Morphismes d'espaces, lemme \ref{spaces-morphisms-lemma-finite-proper}),
on voit que
$$
E =
\text{pr}_0|_{Z_1}\left(Z_1 \setminus Z_1 \cap Z_2)\right)
\cup
\text{pr}_0|_{Z_2}\left(Z_2 \setminus Z_1 \cap Z_2)\right)
$$
est fermé dans $T$. Il est alors clair que $V = T \setminus E$ convient.
\end{proof}

\begin{lemma}
\label{spaces-more-groupoids-lemma-finite-criterion-etale}
Soit $S$ un schéma.
Soit $f : X \to Y$ un morphisme d'espaces algébriques au-dessus de $S$.
Supposons que $U$ soit un schéma, que $U \to Y$ soit un morphisme étale et que
$Z \subset U \times_Y X$ soit un sous-espace ouvert fini sur $U$.
Alors le morphisme induit $U \to (X/Y)_{fin}$ est étale.
\end{lemma}

\begin{proof}
Cela résulte formellement de la description de la diagonale dans le
lemme \ref{spaces-more-groupoids-lemma-finite-diagonal},
mais nous en donnons les détails, car il s'agit d'une étape importante dans le
développement de la théorie. Il faut vérifier que, pour tout schéma $T$ au-dessus de $S$ et tout morphisme
$T \to (X/Y)_{fin}$, la projection
$$
T \times_{(X/Y)_{fin}} U \longrightarrow T
$$
est étale. Remarquons que
$$
T \times_{(X/Y)_{fin}} U
=
(X/Y)_{fin} \times_{((X/Y)_{fin} \times_Y (X/Y)_{fin})} (T \times_Y U)
$$
En appliquant le résultat du
lemme \ref{spaces-more-groupoids-lemma-finite-diagonal},
on voit que $T \times_{(X/Y)_{fin}} U$ est représenté par un sous-schéma ouvert de
$T \times_Y U$. Comme la projection $T \times_Y U \to T$ est étale d'après
Morphismes d'espaces, lemme \ref{spaces-morphisms-lemma-base-change-etale},
on conclut.
\end{proof}

\begin{lemma}
\label{spaces-more-groupoids-lemma-finite-pullback}
Soit $S$ un schéma.
Soit
$$
\xymatrix{
X' \ar[d] \ar[r] & X \ar[d] \\
Y' \ar[r] & Y
}
$$
un carré de produits fibrés d'espaces algébriques au-dessus de $S$. Alors
$$
\xymatrix{
(X'/Y')_{fin} \ar[d] \ar[r] & (X/Y)_{fin} \ar[d] \\
Y' \ar[r] & Y
}
$$
est un carré de produits fibrés de faisceaux sur $(\Sch/S)_{fppf}$.
\end{lemma}

\begin{proof}
Il résulte immédiatement des définitions que
le faisceau $(X'/Y')_{fin}$ est égal au faisceau
$Y' \times_Y (X/Y)_{fin}$.
\end{proof}

\begin{lemma}
\label{spaces-more-groupoids-lemma-finite-surjective-etale-cover}
Soit $S$ un schéma.
Soit $f : X \to Y$ un morphisme d'espaces algébriques au-dessus de $S$.
Si $f$ est séparé et localement quasi-fini, il existe un
schéma $U$ étale sur $Y$ et un morphisme étale surjectif
$U \to (X/Y)_{fin}$ au-dessus de $Y$.
\end{lemma}

\begin{proof}
Remarquons que l'assertion a un sens d'après le résultat du
lemme \ref{spaces-more-groupoids-lemma-finite-diagonal}
sur la diagonale de $(X/Y)_{fin}$; voir
Espaces, lemme \ref{spaces-lemma-representable-diagonal}.
Soit $V$ un schéma et soit $V \to Y$ un morphisme étale surjectif. D'après le
lemme \ref{spaces-more-groupoids-lemma-finite-pullback},
le morphisme $(V \times_Y X/V)_{fin} \to (X/Y)_{fin}$ est
un changement de base de l'application $V \to Y$ et est donc surjectif et étale; voir
Espaces,
lemme \ref{spaces-lemma-base-change-representable-transformations-property}.
Il suffit donc de démontrer le lemme pour $(V \times_Y X/V)_{fin}$.
(On utilise ici implicitement que la composée de transformations de foncteurs
représentables, surjectives et étales est encore représentable, surjective et
étale; voir
Espaces, lemmes \ref{spaces-lemma-composition-representable-transformations} et
\ref{spaces-lemma-composition-representable-transformations-property}, ainsi que
Morphismes, lemmes \ref{morphisms-lemma-composition-surjective} et
\ref{morphisms-lemma-composition-etale}.)
Remarquons que les propriétés d'être séparé et localement quasi-fini
sont préservées par changement de base; voir
Morphismes d'espaces,
lemmes \ref{spaces-morphisms-lemma-base-change-separated} et
\ref{spaces-morphisms-lemma-base-change-quasi-finite}.
Ainsi $V \times_Y X \to V$ est lui aussi séparé et localement quasi-fini,
et, d'après
Morphismes d'espaces, proposition
\ref{spaces-morphisms-proposition-locally-quasi-finite-separated-over-scheme}
on voit que $V \times_Y X$ est également un schéma.
On peut donc supposer que $f : X \to Y$ est un morphisme de schémas
séparé et localement quasi-fini.

\medskip\noindent
Choisissons un point $y \in Y$. Choisissons des points $x_1, \ldots, x_n \in X$
au-dessus de $y$. Choisissons un voisinage étale $a : (U, u) \to (Y, y)$ et une
décomposition
$$
U \times_S X =
W \amalg
\ \coprod\nolimits_{i = 1, \ldots, n}
\ \coprod\nolimits_{j = 1, \ldots, m_j}
V_{i, j}
$$
comme dans
Compléments sur les morphismes, lemme
\ref{more-morphisms-lemma-etale-splits-off-quasi-finite-part-technical-variant}.
Choisissons un sous-ensemble quelconque
$$
I \subset \{(i, j) \mid 1 \leq i \leq n, \ 1 \leq j \leq m_i\}.
$$
Ces choix donnent un couple $(a, Z)$ avec
$Z = \bigcup_{(i, j) \in I} V_{i, j}$
qui satisfait les conditions \ref{spaces-more-groupoids-equation-finite-conditions}. Autrement dit,
on obtient un morphisme $U \to (X/Y)_{fin}$. La construction de ce morphisme
dépend de tous les choix effectués ci-dessus, de sorte qu'il faudrait écrire
$$
U(y, n, x_1, \ldots, x_n, a, I) \longrightarrow (X/Y)_{fin}
$$
Ce morphisme est étale d'après le lemme \ref{spaces-more-groupoids-lemma-finite-criterion-etale}.

\medskip\noindent
Affirmation : leur réunion disjointe est surjective sur $(X/Y)_{fin}$.
Il est clair que, si l'affirmation est vraie, le lemme en résulte.

\medskip\noindent
Pour établir la surjectivité, il faut montrer ce qui suit (voir
Espaces, remarque \ref{spaces-remark-warning}) : étant donnés un schéma $T$ au-dessus de
$S$, un point $t \in T$ et une application $T \to (X/Y)_{fin}$, on peut trouver une donnée
$(y, n, x_1, \ldots, x_n, a, I)$ comme ci-dessus telle que
$t$ appartienne à l'image de la projection
$$
U(y, n, x_1, \ldots, x_n, a, I) \times_{(X/Y)_{fin}} T \longrightarrow T.
$$
Pour le démontrer, on peut clairement remplacer $T$ par
$\Spec(\overline{\kappa(t)})$ et
$T \to (X/Y)_{fin}$ par la composée
$\Spec(\overline{\kappa(t)}) \to T \to (X/Y)_{fin}$.
Autrement dit, on peut supposer que $T$ est
le spectre d'un corps algébriquement clos.

\medskip\noindent
Soit $T = \Spec(k)$ le spectre d'un corps
algébriquement clos $k$. Le morphisme $T \to (X/Y)_{fin}$ est donné par un couple
$(T \to Y, Z)$ satisfaisant les conditions \ref{spaces-more-groupoids-equation-finite-conditions}.
Voici un diagramme :
$$
\xymatrix{
& Z \ar[d] \ar[r] & X \ar[d] \\
\Spec(k) \ar@{=}[r] & T \ar[r] & Y
}
$$
Soit $y \in Y$ le point image de $T \to Y$.
Comme $Z$ est fini sur $k$, il possède un nombre fini de points.
Il existe donc des points $x_1, \ldots, x_n \in X$ en nombre fini
tels que l'image de $Z$ dans $X$ soit contenue dans $\{x_1, \ldots, x_n\}$.
Choisissons $a : (U, u) \to (Y, y)$ adapté à $y$ et $x_1, \ldots, x_n$ comme
ci-dessus; on obtient le diagramme
$$
\xymatrix{
W \amalg
\ \coprod\nolimits_{i = 1, \ldots, n}
\ \coprod\nolimits_{j = 1, \ldots, m_j}
V_{i, j} \ar[d] \ar[r] &
X \ar[d] \\
U \ar[r] & Y.
}
$$
Comme $k$ est algébriquement clos et que
$\kappa(y) \subset \kappa(u)$ est une extension finie séparable,
on peut factoriser le morphisme
$T = \Spec(k) \to Y$ par le morphisme
$u = \Spec(\kappa(u)) \to \Spec(\kappa(y)) = y \subset Y$.
Avec ce choix, on obtient le diagramme commutatif :
$$
\xymatrix{
Z \ar[d] \ar[r] &
W \amalg
\ \coprod\nolimits_{i = 1, \ldots, n}
\ \coprod\nolimits_{j = 1, \ldots, m_j}
V_{i, j} \ar[d] \ar[r] &
X \ar[d] \\
\Spec(k) \ar[r] &
U \ar[r] & Y
}
$$
Nous savons que l'image de la flèche supérieure gauche est contenue dans
$\coprod V_{i, j}$. Rappelons aussi que $Z$ est un sous-schéma ouvert
de $\Spec(k) \times_Y X$ par définition de $(X/Y)_{fin}$
et que le carré de droite est un carré de produits fibrés.
On voit donc que
$$
Z \subset
\coprod\nolimits_{i = 1, \ldots, n}\ \coprod\nolimits_{j = 1, \ldots, m_j}
\Spec(k) \times_U V_{i, j}
$$
est un sous-schéma ouvert. Par construction (voir
Compléments sur les morphismes, lemme
\ref{more-morphisms-lemma-etale-splits-off-quasi-finite-part-technical-variant})
chaque $V_{i, j}$ possède un unique point $v_{i, j}$ au-dessus de $u$
dont l'extension de corps résiduels
$\kappa(v_{i, j})/\kappa(u)$ est purement inséparable. Ainsi, chaque
schéma $\Spec(k) \times_U V_{i, j}$ possède exactement un
point. On voit donc que
$$
Z = \coprod\nolimits_{(i, j) \in I} \Spec(k) \times_U V_{i, j}
$$
pour un unique sous-ensemble
$I \subset \{(i, j) \mid 1 \leq i \leq n, \ 1 \leq j \leq m_i\}$.
En explicitant les définitions, cela montre que
$$
U(y, n, x_1, \ldots, x_n, a, I) \times_{(X/Y)_{fin}} T
$$
avec le $I$ trouvé ci-dessus est non vide, comme voulu.
\end{proof}

\begin{proposition}
\label{spaces-more-groupoids-proposition-finite-algebraic-space}
Soit $S$ un schéma.
Soit $f : X \to Y$ un morphisme d'espaces algébriques au-dessus de $S$ qui
est séparé et localement de type fini. Alors $(X/Y)_{fin}$
est un espace algébrique. De plus, le morphisme
$(X/Y)_{fin} \to Y$ est étale.
\end{proposition}

\begin{proof}
D'après le
lemme \ref{spaces-more-groupoids-lemma-finite-lives-on-locally-quasi-finite-part},
on peut remplacer $X$ par le sous-schéma ouvert localement quasi-fini
sur $Y$. On peut donc supposer $f$ séparé et localement quasi-fini.
Vérifions les trois conditions de
Espaces, définition \ref{spaces-definition-algebraic-space}.
La condition (1) résulte du
lemme \ref{spaces-more-groupoids-lemma-finite-sheaf}.
La condition (2) résulte du
lemme \ref{spaces-more-groupoids-lemma-finite-diagonal}.
Enfin, la condition (3) résulte du
lemme \ref{spaces-more-groupoids-lemma-finite-surjective-etale-cover}.
Ainsi $(X/Y)_{fin}$ est un espace algébrique.
De plus, ce lemme montre qu'il existe un
diagramme commutatif
$$
\xymatrix{
U \ar[rr] \ar[rd] & & (X/Y)_{fin} \ar[ld] \\
& Y
}
$$
dont la flèche horizontale est surjective et étale et la flèche sud-est
est étale. D'après
Propriétés des espaces, lemme \ref{spaces-properties-lemma-etale-local},
cela implique que la flèche sud-ouest est elle aussi étale.
\end{proof}

\begin{remark}
\label{spaces-more-groupoids-remark-warning}
L'hypothèse que $f$ soit séparé ne peut être supprimée dans la
proposition \ref{spaces-more-groupoids-proposition-finite-algebraic-space}.
Un exemple consiste à prendre pour $X$ la droite affine à origine dédoublée; voir
Schémas, exemple \ref{schemes-example-affine-space-zero-doubled},
pour $Y = \mathbf{A}^1_k$ la droite affine, et pour $X \to Y$ l'application évidente.
Rappelons qu'au-dessus de $0 \in Y$, il y a deux points $0_1$ et $0_2$
dans $X$. Ainsi, $(X/Y)_{fin}$ possède quatre points au-dessus de $0$, à savoir
$\emptyset, \{0_1\}, \{0_2\}, \{0_1, 0_2\}$.
De ces quatre points, seuls trois peuvent se relever en un sous-schéma ouvert
de $U \times_Y X$ fini sur $U$ lorsque $U \to Y$ est étale,
à savoir $\emptyset, \{0_1\}, \{0_2\}$. Cela montre que $(X/Y)_{fin}$,
s'il est représentable par un espace algébrique, n'est pas étale sur $Y$.
Des arguments analogues montrent que $(X/Y)_{fin}$ n'est réellement pas un espace
algébrique. Les détails sont omis.
\end{remark}

\begin{remark}
\label{spaces-more-groupoids-remark-not-scheme}
Soit $Y = \mathbf{A}^1_{\mathbf{R}}$ la droite affine sur le corps des nombres
réels, et soit $X = \Spec(\mathbf{C})$ envoyé sur le point
$\mathbf{R}$-rationnel $0$ de $Y$. Dans ce cas, le morphisme
$f : X \to Y$ est fini, mais $(X/Y)_{fin}$ n'est pas
un schéma. En effet, on peut montrer que, dans ce cas, l'espace
algébrique $(X/Y)_{fin}$ est isomorphe à l'espace algébrique d'
Espaces, exemple \ref{spaces-example-non-representable-descent},
associé à l'extension $\mathbf{R} \subset \mathbf{C}$.
Il est donc réellement nécessaire de sortir de la catégorie des schémas
pour représenter le faisceau $(X/Y)_{fin}$, même lorsque $f$
est un morphisme fini.
\end{remark}

\begin{lemma}
\label{spaces-more-groupoids-lemma-finite-separated-flat-locally-finite-presentation}
Soit $S$ un schéma.
Soit $f : X \to Y$ un morphisme d'espaces algébriques au-dessus de $S$ qui
est séparé, plat et localement de présentation finie.
Dans ce cas,
\begin{enumerate}
\item $(X/Y)_{fin} \to Y$ est séparé, représentable et étale, et
\item si $Y$ est un schéma, alors $(X/Y)_{fin}$ est (représenté par) un schéma.
\end{enumerate}
\end{lemma}

\begin{proof}
Comme $f$ est en particulier séparé et localement de type fini (voir
Morphismes d'espaces,
lemme \ref{spaces-morphisms-lemma-finite-presentation-finite-type}),
on voit que $(X/Y)_{fin}$ est un espace algébrique d'après la
proposition \ref{spaces-more-groupoids-proposition-finite-algebraic-space}.
Pour montrer que $(X/Y)_{fin} \to Y$ est séparé, il faut établir
ce qui suit : étant donnés un schéma $T$ et deux couples $(a, Z_1)$ et $(a, Z_2)$
au-dessus de $T$
ayant la même première composante et satisfaisant \ref{spaces-more-groupoids-equation-finite-conditions},
il existe un sous-schéma fermé $V \subset T$ possédant la propriété suivante :
pour tout morphisme de schémas $h : T' \to T$, on a
$$
h \text{ se factorise par } V \Leftrightarrow
\Big(T' \times_T Z_1 = T' \times_T Z_2
\text{ comme sous-espaces de }T' \times_Y X\Big)
$$
Dans la démonstration du
lemme \ref{spaces-more-groupoids-lemma-finite-diagonal},
on a vu que $V = T \setminus E$ est un sous-schéma ouvert de $T$
de complémentaire fermé
$$
E =
\text{pr}_0|_{Z_1}\left(Z_1 \setminus Z_1 \cap Z_2)\right)
\cup
\text{pr}_0|_{Z_2}\left(Z_2 \setminus Z_1 \cap Z_2)\right).
$$
Il suffit donc de montrer que $E$ est aussi ouvert. D'après le
lemme \ref{spaces-more-groupoids-lemma-finite-separated},
on voit que $Z_1$ et $Z_2$ sont fermés dans $T \times_Y X$. Ainsi,
$Z_1 \setminus Z_1 \cap Z_2$ est ouvert dans $Z_1$. Comme $f$ est plat et
localement de présentation finie, $\text{pr}_0|_{Z_1}$ l'est aussi.
Cela est vrai parce que $Z_1$ est un sous-espace ouvert du changement de base
$T \times_Y X$, et d'après
Morphismes d'espaces,
lemmes \ref{spaces-morphisms-lemma-base-change-finite-presentation} et
lemme \ref{spaces-morphisms-lemma-base-change-flat}.
Ainsi $\text{pr}_0|_{Z_1}$ est ouvert; voir
Morphismes d'espaces, lemme \ref{spaces-morphisms-lemma-fppf-open}.
Par conséquent, $\text{pr}_0|_{Z_1}\left(Z_1 \setminus Z_1 \cap Z_2)\right)$ est
ouvert et il s'ensuit que $E$ est ouvert, comme voulu.

\medskip\noindent
Nous avons déjà vu que $(X/Y)_{fin} \to Y$ est étale; voir la
proposition \ref{spaces-more-groupoids-proposition-finite-algebraic-space}.
Nous savons donc qu'il est localement quasi-fini (voir
Morphismes d'espaces,
lemme \ref{spaces-morphisms-lemma-etale-locally-quasi-finite})
et séparé, donc représentable d'après
Morphismes d'espaces, lemme
\ref{spaces-morphisms-lemma-locally-quasi-finite-separated-representable}.
La dernière assertion est claire (on peut, si l'on veut, utiliser
Morphismes d'espaces, proposition
\ref{spaces-morphisms-proposition-locally-quasi-finite-separated-over-scheme}).
\end{proof}

\noindent
Variante : soit $S$ un schéma.
Soit $f : X \to Y$ un morphisme d'espaces algébriques au-dessus de $S$.
Soit $\sigma : Y \to X$ une section de $f$.
Pour un espace algébrique ou un schéma $T$ au-dessus de $S$, considérons les couples
$(a, Z)$ où
\begin{equation}
\label{spaces-more-groupoids-equation-finite-conditions-variant}
\begin{matrix}
a : T \to Y\text{ est un morphisme au-dessus de }S, \\
Z \subset T \times_Y X\text{ est un sous-espace ouvert} \\
\text{tel que }\text{pr}_0|_Z : Z \to T\text{ soit fini et} \\
(1_T, \sigma \circ a) : T \to T \times_Y X\text{ se factorise par }Z.
\end{matrix}
\end{equation}
Nous noterons $(X/Y, \sigma)_{fin}$ le sous-foncteur de $(X/Y)_{fin}$
qui paramètre ces couples.

\begin{lemma}
\label{spaces-more-groupoids-lemma-finite-plus-section}
Soit $S$ un schéma.
Soit $f : X \to Y$ un morphisme d'espaces algébriques au-dessus de $S$.
Soit $\sigma : Y \to X$ une section de $f$. Considérons la
transformation de foncteurs
$$
t : (X/Y, \sigma)_{fin} \longrightarrow (X/Y)_{fin}.
$$
définie ci-dessus. Alors
\begin{enumerate}
\item $t$ est représentable par des immersions ouvertes,
\item si $f$ est séparé, alors $t$ est représentable par des immersions
ouvertes et fermées,
\item si $(X/Y)_{fin}$ est un espace algébrique, alors
$(X/Y, \sigma)_{fin}$ est un espace algébrique et
un sous-espace ouvert de $(X/Y)_{fin}$, et
\item si $(X/Y)_{fin}$ est un schéma, alors $(X/Y, \sigma)_{fin}$ en est un
sous-schéma ouvert.
\end{enumerate}
\end{lemma}

\begin{proof}
Démonstration omise. Indication : étant donné un couple $(a, Z)$ au-dessus de $T$ comme dans
(\ref{spaces-more-groupoids-equation-finite-conditions}), l'image inverse de
$Z$ par $(1_T, \sigma \circ a) : T \to T \times_Y X$ est le sous-schéma ouvert
de $T$ recherché.
\end{proof}





\section{Collections finies de flèches}
\label{spaces-more-groupoids-section-finite-set-arrows}

\noindent
Soit $\mathcal{C}$ un groupoïde, voir
Catégories, définition \ref{categories-definition-groupoid}.
Comme expliqué dans
Groupoïdes, section \ref{groupoids-section-groupoids},
cela correspond à un septuple $(\text{Ob}, \text{Flèches}, s, t, c, e, i)$.

\medskip\noindent
À partir de ces données, nous pouvons construire un autre groupoïde
$\mathcal{C}_{fin}$ comme suit :
\begin{enumerate}
\item Un objet de $\mathcal{C}_{fin}$ consiste en une partie finie
$Z \subset \text{Flèches}$ possédant les propriétés suivantes :
\begin{enumerate}
\item $s(Z) = \{u\}$ est un singleton, et
\item $e(u) \in Z$.
\end{enumerate}
\item Un morphisme de $\mathcal{C}_{fin}$ consiste en un couple
$(Z, z)$, où $Z$ est un objet de $\mathcal{C}_{fin}$ et
$z \in Z$.
\item La source de $(Z, z)$ est $Z$.
\item Le but de $(Z, z)$ est $t(Z, z) = \{z' \circ z^{-1}; z' \in Z\}$.
\item Étant donnés $(Z_1, z_1)$ et $(Z_2, z_2)$ tels que
$s(Z_1, z_1) = t(Z_2, z_2)$, la composée $(Z_1, z_1) \circ (Z_2, z_2)$ est $(Z_2, z_1 \circ z_2)$.
\end{enumerate}
Nous omettons la vérification que ceci définit un groupoïde.
Graphiquement, un objet de $\mathcal{C}_{fin}$ peut se représenter
par un diagramme
$$
\xymatrix{
& \bullet \\
\bullet \ar@(ul, dl)[]_e \ar[ru] \ar[r] \ar[rd] & \bullet \\
& \bullet
}
$$
Pour construire un morphisme de $\mathcal{C}_{fin}$, on choisit l'une des
flèches et l'on précompose les autres flèches par son inverse. Par exemple,
si l'on choisit la flèche horizontale du milieu, le but est représenté par
$$
\xymatrix{
& \bullet \\
\bullet & \bullet \ar[l] \ar[u] \ar@(dr, ur)[]_e \ar[d] \\
& \bullet
}
$$
Notons que les cardinalités de $s(Z, z)$ et de $t(Z, z)$ sont égales.
Ainsi, $\mathcal{C}_{fin}$ est en fait une réunion disjointe dénombrable
de groupoïdes.




\section{La partie finie d'un groupoïde}
\label{spaces-more-groupoids-section-finite-part-groupoid}

\noindent
Dans cette section, nous allons utiliser l'idée expliquée dans
la section \ref{spaces-more-groupoids-section-finite-set-arrows}
pour prendre la partie finie d'un groupoïde en espaces algébriques.

\medskip\noindent
Soit $S$ un schéma.
Soit $B$ un espace algébrique au-dessus de $S$.
Soit $(U, R, s, t, c, e, i)$ un groupoïde en espaces algébriques au-dessus de $B$.
Hypothèse : les morphismes $s, t$ sont séparés et localement de type fini.
Cette notation et cette hypothèse resteront fixées dans toute la section.

\medskip\noindent
Notons $R_s$ l'espace algébrique $R$ considéré comme un
espace algébrique au-dessus de $U$ par l'intermédiaire de $s$. Posons
$U' = (R_s/U, e)_{fin}$. Puisque $s$ est séparé et localement
de type fini, la
proposition \ref{spaces-more-groupoids-proposition-finite-algebraic-space} et le
lemme \ref{spaces-more-groupoids-lemma-finite-plus-section}
montrent que $U'$ est un espace algébrique muni d'un morphisme \'etale
$g : U' \to U$. De plus, d'après le
lemme \ref{spaces-more-groupoids-lemma-finite-sheaf},
il existe un sous-espace ouvert universel
$Z_{univ} \subset R \times_{s, U, g} U'$, fini au-dessus de $U'$,
tel que $(1_{U'}, e \circ g) : U' \to R \times_{s, U, g} U'$
se factorise par $Z_{univ}$. En outre, d'après le
lemme \ref{spaces-more-groupoids-lemma-finite-separated},
le sous-espace ouvert $Z_{univ}$ est aussi fermé dans
$R \times_{s, U, g} U'$. Voici la situation obtenue :
$$
\xymatrix{
Z_{univ} \ar[d] \ar[rd] & \\
R \times_{s, U, g} U' \ar[d] \ar[r] & U' \ar[d]^g \\
R \ar[r]^s & U
}
$$
Soit $T$ un schéma au-dessus de $B$. On voit qu'un point de
$Z_{univ}$ à valeurs dans $T$ peut être considéré comme un triplet
$(u, Z, z)$ où
\begin{enumerate}
\item $u : T \to U$ est un point de $U$ à valeurs dans $T$,
\item $Z \subset R \times_{s, U, u} T$ est un sous-espace ouvert et fermé,
fini au-dessus de $T$, par lequel $(e \circ u, 1_T)$ se factorise, et
\item $z : T \to R$ est un point de $R$ à valeurs dans $T$ tel que
$s \circ z = u$ et tel que $(z, 1_T)$ se factorise par $Z$.
\end{enumerate}
Cela étant dit, la discussion de la
section \ref{spaces-more-groupoids-section-finite-set-arrows}
montre moralement que l'on peut munir $(Z_{univ}, U')$ d'une structure de
groupoïde en espaces algébriques au-dessus de $B$. Pour nous en assurer,
définissons un à un les morphismes $s', t', c', e', i'$ du point de vue
fonctoriel. (Nous invitons le lecteur à ne pas lire ce qui suit avant d'avoir
lu et compris la construction simple de la
section \ref{spaces-more-groupoids-section-finite-set-arrows}.)

\medskip\noindent
Le morphisme $s' : Z_{univ} \to U'$ correspond à la règle
$$
s' : (u, Z, z) \mapsto (u, Z).
$$
Le morphisme $t' : Z_{univ} \to U'$ est donné par la règle
$$
t' : (u, Z, z) \mapsto (t \circ z, c(Z, i \circ z)).
$$
Le terme $c(Z, i \circ z)$
a un sens, car l'application
$c(-, i \circ z) : R \times_{s, U, u} T \to R \times_{s, U, t \circ z} T$
est un isomorphisme d'inverse $c(-, z)$.
Le morphisme $e' : U' \to Z_{univ}$ est donné par la règle
$$
e' : (u, Z) \mapsto (u, Z, (e \circ u, 1_T)).
$$
Notons que ceci a un sens en vertu de la condition imposant que
$(e \circ u, 1_T)$ se factorise par $Z$.
Le morphisme $i' : Z_{univ} \to Z_{univ}$ est donné par la règle
$$
i' : (u, Z, z) \mapsto (t \circ z, c(Z, i \circ z), i \circ z).
$$
Enfin, la composition est définie par la règle
$$
c' : ((u_1, Z_1, z_1), (u_2, Z_2, z_2)) \mapsto (u_2, Z_2, z_1 \circ z_2).
$$
Nous omettons la vérification que les axiomes d'un groupoïde en espaces
algébriques sont satisfaits par $(U', Z_{univ}, s', t', c', e', i')$.

\medskip\noindent
Enfin, il existe un morphisme canonique de groupoïdes
$$
(U', Z_{univ}, s', t', c', e', i')
\longrightarrow
(U, R, s, t, c, e, i)
$$
En effet, le morphisme $U' \to U$ est le morphisme $g : U' \to U$,
qui est défini par la règle $(u, Z) \mapsto u$. Le morphisme
$Z_{univ} \to R$ est défini par la règle $(u, Z, z) \mapsto z$.
Ceci achève la construction. Résumons les résultats obtenus.

\begin{lemma}
\label{spaces-more-groupoids-lemma-finite-part-groupoid}
Soit $S$ un schéma.
Soit $B$ un espace algébrique au-dessus de $S$.
Soit $(U, R, s, t, c, e, i)$ un groupoïde en espaces algébriques au-dessus de $B$.
Supposons que les morphismes $s, t$ soient séparés et localement de type fini.
Il existe un morphisme canonique
$$
(U', Z_{univ}, s', t', c', e', i')
\longrightarrow
(U, R, s, t, c, e, i)
$$
de groupoïdes en espaces algébriques au-dessus de $B$, où
\begin{enumerate}
\item $g : U' \to U$ s'identifie à $(R_s/U, e)_{fin} \to U$, et
\item $Z_{univ} \subset R \times_{s, U, g} U'$ est le sous-espace
ouvert (et fermé) universel, fini au-dessus de $U'$, qui contient le
changement de base de l'unité $e$.
\end{enumerate}
\end{lemma}

\begin{proof}
Voir la discussion ci-dessus.
\end{proof}









\section{Localisation \'etale des schémas en groupoïdes}
\label{spaces-more-groupoids-section-etale-localize}

\noindent
Dans cette section, nous démontrons des résultats analogues à
\cite[Proposition 4.2]{K-M}. Nous cherchons à être un peu plus généraux et
à éviter l'emploi des schémas de Hilbert en utilisant plutôt la partie finie
d'un morphisme. Le but est de « scinder » un groupoïde en espaces algébriques
au-dessus d'un point après localisation \'etale. Voici la définition
(très proche de \cite[définition 4.1]{K-M}).

\begin{definition}
\label{spaces-more-groupoids-definition-split-at-point}
Soit $S$ un schéma. Soit $B$ un espace algébrique au-dessus de $S$.
Soit $(U, R, s, t, c)$ un groupoïde en espaces algébriques au-dessus de $B$.
Soit $u \in |U|$ un point.
\begin{enumerate}
\item On dit que $R$ est {\it fortement scindé au-dessus de $u$} s'il existe
un sous-espace ouvert $P \subset R$ tel que
\begin{enumerate}
\item $(U, P, s|_P, t|_P, c|_{P \times_{s, U, t} P})$ est un
groupoïde en espaces algébriques au-dessus de $B$,
\item $s|_P$, $t|_P$ sont finis, et
\item $\{r \in |R| : s(r) = u, t(r) = u\} \subset |P|$.
\end{enumerate}
Le choix d'un tel $P$ sera appelé un
{\it scindage fort de $R$ au-dessus de $u$}.
\item On dit que $R$ est {\it scindé au-dessus de $u$} s'il existe un
sous-espace ouvert $P \subset R$ tel que
\begin{enumerate}
\item $(U, P, s|_P, t|_P, c|_{P \times_{s, U, t} P})$ est un
groupoïde en espaces algébriques au-dessus de $B$,
\item $s|_P$, $t|_P$ sont finis, et
\item $\{g \in |G| : g\text{ est au-dessus de }u\} \subset |P|$, où
$G \to U$ est le stabilisateur.
\end{enumerate}
Le choix d'un tel $P$ sera appelé un
{\it scindage de $R$ au-dessus de $u$}.
\item On dit que $R$ est {\it quasi-scindé au-dessus de $u$} s'il existe
un sous-espace ouvert $P \subset R$ tel que
\begin{enumerate}
\item $(U, P, s|_P, t|_P, c|_{P \times_{s, U, t} P})$ est un
groupoïde en espaces algébriques au-dessus de $B$,
\item $s|_P$, $t|_P$ sont finis, et
\item $e(u) \in |P|$\footnote{Cette condition découle de (a).}.
\end{enumerate}
Le choix d'un tel $P$ sera appelé un {\it quasi-scindage de $R$ au-dessus de $u$}.
\end{enumerate}
\end{definition}

\noindent
Notons la similitude entre les conditions imposées à $P$ et celles imposées
aux couples dans (\ref{spaces-more-groupoids-equation-finite-conditions}). En particulier, si
$s, t$ sont séparés, alors $P$ est aussi fermé dans $R$ (voir le
lemme \ref{spaces-more-groupoids-lemma-finite-separated}).

\medskip\noindent
Supposons donné un groupoïde en espaces algébriques
$(U, R, s, t, c)$ au-dessus de $B$ et un point $u \in |U|$.
Puisque le but est de scinder le groupoïde après localisation \'etale,
on peut tout aussi bien remplacer $U$ par un schéma affine (nous voulons dire
que cela est sans conséquence pour toute application éventuelle).
En outre, les hypothèses supplémentaires que nous devrons
imposer feront de $R$ un schéma, au moins dans un voisinage
de $\{r \in |R| : s(r) = u, t(r) = u\}$ ou de $e(u)$. C'est pourquoi
nous partons d'un schéma en groupoïdes comme décrit ci-dessous.
Cependant, notre méthode de démonstration nous fait sortir de la catégorie
des schémas, raison pour laquelle nous avons formulé ci-dessus le scindage
pour les groupoïdes en espaces algébriques.
D'autre part, nous ne connaissons d'applications que
dans le cas où les morphismes $s$, $t$
sont aussi plats et de présentation finie ; dans ce cas,
nous revenons finalement dans la catégorie des schémas.

\begin{situation}[Scindage fort]
\label{spaces-more-groupoids-situation-strong-splitting}
Soit $S$ un schéma.
Soit $(U, R, s, t, c)$ un schéma en groupoïdes au-dessus de $S$.
Soit $u \in U$ un point. Supposons que
\begin{enumerate}
\item $s, t : R \to U$ soient séparés,
\item $s$, $t$ soient localement de type fini,
\item l'ensemble $\{r \in R : s(r) = u, t(r) = u\}$ soit fini, et
\item $s$ soit quasi-fini en chaque point de l'ensemble de (3).
\end{enumerate}
Notons que les hypothèses (3) et (4) découlent de l'hypothèse
que la fibre $s^{-1}(\{u\})$ est finie, voir
Morphismes, lemme \ref{morphisms-lemma-finite-fibre}.
\end{situation}

\begin{situation}[Scindage]
\label{spaces-more-groupoids-situation-splitting}
Soit $S$ un schéma.
Soit $(U, R, s, t, c)$ un schéma en groupoïdes au-dessus de $S$.
Soit $u \in U$ un point. Supposons que
\begin{enumerate}
\item $s, t : R \to U$ soient séparés,
\item $s$, $t$ soient localement de type fini,
\item l'ensemble $\{g \in G : g\text{ est au-dessus de }u\}$ soit fini, où $G \to U$
est le stabilisateur, et
\item $s$ soit quasi-fini en chaque point de l'ensemble de (3).
\end{enumerate}
\end{situation}

\begin{situation}[Quasi-scindage]
\label{spaces-more-groupoids-situation-quasi-splitting}
Soit $S$ un schéma.
Soit $(U, R, s, t, c)$ un schéma en groupoïdes au-dessus de $S$.
Soit $u \in U$ un point. Supposons que
\begin{enumerate}
\item $s, t : R \to U$ soient séparés,
\item $s$, $t$ soient localement de type fini, et
\item $s$ soit quasi-fini en $e(u)$.
\end{enumerate}
\end{situation}

\noindent
Pour notre application aux théorèmes d'existence des espaces algébriques,
le cas des quasi-scindages suffit. En outre, le cas du quasi-scindage
nous permettra de démontrer, pour les champs quasi-DM, un théorème de
structure locale pour la topologie \'etale. Le cas du scindage servira à démontrer une version
du théorème de Keel--Mori. Le cas du scindage fort fournit
un théorème de structure locale pour la topologie \'etale des champs
algébriques quasi-DM à diagonale quasi-compacte.

\begin{lemma}[Existence d'un scindage fort]
\label{spaces-more-groupoids-lemma-strong-splitting}
Dans la situation \ref{spaces-more-groupoids-situation-strong-splitting},
il existe un espace algébrique $U'$, un morphisme \'etale
$U' \to U$ et un point $u' : \Spec(\kappa(u)) \to U'$
au-dessus de $u : \Spec(\kappa(u)) \to U$,
tels que la restriction $R' = R|_{U'}$ de $R$ à $U'$
soit fortement scindée au-dessus de $u'$.
\end{lemma}

\begin{proof}
Soit $f : (U', Z_{univ}, s', t', c') \to (U, R, s, t, c)$ construit dans le
lemme \ref{spaces-more-groupoids-lemma-finite-part-groupoid}.
Rappelons que $R' = R \times_{(U \times_S U)} (U' \times_S U')$.
On obtient ainsi un morphisme $(f, t', s') : Z_{univ} \to R'$ de groupoïdes
en espaces algébriques
$$
(U', Z_{univ}, s', t', c') \to (U', R', s', t', c')
$$
(par abus de notation, nous désignons les morphismes des deux groupoïdes
par les mêmes symboles). Or $Z_{univ} \subset R \times_{s, U, g} U'$ est ouvert,
et $R' \to R \times_{s, U, g} U'$ est \'etale (comme changement de base
de $U' \to U$) ; ainsi, $Z_{univ} \to R'$ est une immersion ouverte.
Par construction, les morphismes $s', t' : Z_{univ} \to U'$ sont finis.
Il reste à trouver le point $u'$ de $U'$.

\medskip\noindent
Considérons $u$ comme un morphisme $\Spec(\kappa(u)) \to U$, comme dans
l'énoncé du lemme. Posons $F_u = R \times_{s, U} \Spec(\kappa(u))$.
L'ensemble $\{r \in R : s(r) = u, t(r) = u\}$ est fini par hypothèse
et $F_u \to \Spec(\kappa(u))$ est quasi-fini en chacun
de ses éléments par hypothèse. On peut donc trouver une décomposition en
sous-schémas ouverts et fermés
$$
F_u = Z_u \amalg Reste
$$
où $Z_u$ est un schéma fini au-dessus de $\kappa(u)$ dont le support est
$\{r \in R : s(r) = u, t(r) = u\}$. Notons que $e(u) \in Z_u$.
Par conséquent, d'après la construction de $U'$ dans la
section \ref{spaces-more-groupoids-section-finite-part-groupoid},
$(u, Z_u)$ définit un point de $U'$ à valeurs dans
$\Spec(\kappa(u))$, que nous notons $u'$.

\medskip\noindent
Il nous reste à montrer que l'ensemble
$\{r' \in |R'| : s'(r') = u', t'(r') = u'\}$
est contenu dans $|Z_{univ}|$.
Choisissons un point quelconque $r'$ de cet ensemble et représentons-le par
un morphisme $z' : \Spec(k) \to R'$. Notons $z : \Spec(k) \to R$
la composée de $z'$ et de l'application $R' \to R$.
Il est clair que $z$ définit un élément de l'ensemble
$\{r \in R : s(r) = u, t(r) = u\}$.
De plus, les composées $s \circ z, t \circ z : \Spec(k) \to U$
se factorisent par $u$ ; nous pouvons donc considérer $s \circ z, t \circ z$
comme des morphismes $\Spec(k) \to \Spec(\kappa(u))$. Alors
$z' = (z, u' \circ t \circ z, u'\circ s \circ z)$ comme morphismes vers
$R' = R \times_{(U \times_S U)} (U' \times_S U')$.
Considérons le triplet
$$
(s \circ z, Z_u \times_{\Spec(\kappa(u)), s \circ z} \Spec(k), z)
$$
où $Z_u$ est défini comme ci-dessus. Ceci définit un point de $Z_{univ}$
à valeurs dans $\Spec(k)$, dont l'image dans $U'$ par $s', t'$ est $u'$ et
dont l'image par
$Z_{univ} \to R'$ est le point $r'$, en vertu de la relation entre
$z$ et $z'$ mentionnée ci-dessus.
Ceci achève la démonstration.
\end{proof}

\begin{lemma}[Existence d'un scindage]
\label{spaces-more-groupoids-lemma-splitting}
Dans la situation \ref{spaces-more-groupoids-situation-splitting},
il existe un espace algébrique $U'$, un morphisme \'etale
$U' \to U$ et un point $u' : \Spec(\kappa(u)) \to U'$
au-dessus de $u : \Spec(\kappa(u)) \to U$,
tels que la restriction $R' = R|_{U'}$ de $R$ à $U'$
soit scindée au-dessus de $u'$.
\end{lemma}

\begin{proof}
Soit $f : (U', Z_{univ}, s', t', c') \to (U, R, s, t, c)$ construit dans le
lemme \ref{spaces-more-groupoids-lemma-finite-part-groupoid}.
Rappelons que $R' = R \times_{(U \times_S U)} (U' \times_S U')$.
On obtient ainsi un morphisme $(f, t', s') : Z_{univ} \to R'$ de groupoïdes
en espaces algébriques
$$
(U', Z_{univ}, s', t', c') \to (U', R', s', t', c')
$$
(par abus de notation, nous désignons les morphismes des deux groupoïdes
par les mêmes symboles). Or $Z_{univ} \subset R \times_{s, U, g} U'$ est ouvert,
et $R' \to R \times_{s, U, g} U'$ est \'etale (comme changement de base
de $U' \to U$) ; ainsi, $Z_{univ} \to R'$ est une immersion ouverte.
Par construction, les morphismes $s', t' : Z_{univ} \to U'$ sont finis.
Il reste à trouver le point $u'$ de $U'$.

\medskip\noindent
Considérons $u$ comme un morphisme $\Spec(\kappa(u)) \to U$, comme dans
l'énoncé du lemme. Posons $F_u = R \times_{s, U} \Spec(\kappa(u))$.
Soit $G_u \subset F_u$ la fibre schématique de $G \to U$ au-dessus de $u$.
Par hypothèse, $G_u$ est fini et $F_u \to \Spec(\kappa(u))$
est quasi-fini en chaque point de $G_u$.
On peut donc trouver une décomposition en sous-schémas ouverts et fermés
$$
F_u = Z_u \amalg Reste
$$
où $Z_u$ est un schéma fini au-dessus de $\kappa(u)$ dont le support est $G_u$.
Notons que $e(u) \in Z_u$. Par conséquent, d'après la construction de $U'$ dans la
section \ref{spaces-more-groupoids-section-finite-part-groupoid},
$(u, Z_u)$ définit un point de $U'$ à valeurs dans
$\Spec(\kappa(u))$, que nous notons $u'$.

\medskip\noindent
Il nous reste à montrer que l'ensemble $\{g' \in |G'| : g'\text{ est au-dessus de }u'\}$
est contenu dans $|Z_{univ}|$. Choisissons un point $g'$ de cet ensemble et
représentons-le par un morphisme $z' : \Spec(k) \to G'$. Notons
$z : \Spec(k) \to G$ la composée de $z'$ et de l'application $G' \to G$.
Il est clair que $z$ définit un point de $G_u$. Écrivons plus précisément
$\tilde u : \Spec(k) \to u \to U$ pour l'application correspondante vers $u$ ou $U$.
Considérons le triplet
$$
(\tilde u, Z_u \times_{u, \tilde u} \Spec(k), z)
$$
où $Z_u$ est défini comme ci-dessus. Ceci définit un point de $Z_{univ}$
à valeurs dans $\Spec(k)$, dont l'image dans $U'$ par $s', t'$ est $u'$ et
dont l'image par $Z_{univ} \to R'$ est le point $z'$
(car son image dans $R$ est $z$).
Ceci achève la démonstration.
\end{proof}

\begin{lemma}[Existence d'un quasi-scindage]
\label{spaces-more-groupoids-lemma-quasi-splitting}
Dans la situation \ref{spaces-more-groupoids-situation-quasi-splitting},
il existe un espace algébrique $U'$, un morphisme \'etale
$U' \to U$ et un point $u' : \Spec(\kappa(u)) \to U'$
au-dessus de $u : \Spec(\kappa(u)) \to U$,
tels que la restriction $R' = R|_{U'}$ de $R$ à $U'$
soit quasi-scindée au-dessus de $u'$.
\end{lemma}

\begin{proof}
Soit $f : (U', Z_{univ}, s', t', c') \to (U, R, s, t, c)$ construit dans le
lemme \ref{spaces-more-groupoids-lemma-finite-part-groupoid}.
Rappelons que $R' = R \times_{(U \times_S U)} (U' \times_S U')$.
On obtient ainsi un morphisme $(f, t', s') : Z_{univ} \to R'$ de groupoïdes
en espaces algébriques
$$
(U', Z_{univ}, s', t', c') \to (U', R', s', t', c')
$$
(par abus de notation, nous désignons les morphismes des deux groupoïdes
par les mêmes symboles). Or $Z_{univ} \subset R \times_{s, U, g} U'$ est ouvert,
et $R' \to R \times_{s, U, g} U'$ est \'etale (comme changement de base
de $U' \to U$) ; ainsi, $Z_{univ} \to R'$ est une immersion ouverte.
Par construction, les morphismes $s', t' : Z_{univ} \to U'$ sont finis.
Il reste à trouver le point $u'$ de $U'$.

\medskip\noindent
Considérons $u$ comme un morphisme $\Spec(\kappa(u)) \to U$, comme dans
l'énoncé du lemme. Posons $F_u = R \times_{s, U} \Spec(\kappa(u))$.
Le morphisme $F_u \to \Spec(\kappa(u))$ est quasi-fini en $e(u)$
par hypothèse. On peut donc trouver une décomposition en sous-schémas
ouverts et fermés
$$
F_u = Z_u \amalg Reste
$$
où $Z_u$ est un schéma fini au-dessus de $\kappa(u)$ dont le support est $e(u)$.
Par conséquent, d'après la construction de $U'$ dans la
section \ref{spaces-more-groupoids-section-finite-part-groupoid},
$(u, Z_u)$ définit un point de $U'$ à valeurs dans
$\Spec(\kappa(u))$, que nous notons $u'$. Pour achever la démonstration,
il suffit de montrer que $e'(u') \in Z_{univ}$, ce qui est clair.
\end{proof}

\noindent
Enfin, sous des hypothèses supplémentaires, nous obtenons des schémas.

\begin{lemma}
\label{spaces-more-groupoids-lemma-strong-splitting-scheme}
Dans la situation \ref{spaces-more-groupoids-situation-strong-splitting}, supposons en outre que
$s, t$ soient plats et localement de présentation finie.
Alors il existe un schéma $U'$, un morphisme \'etale séparé
$U' \to U$ et un point $u' \in U'$
au-dessus de $u$, avec $\kappa(u) = \kappa(u')$,
tels que la restriction $R' = R|_{U'}$ de $R$ à $U'$
soit fortement scindée au-dessus de $u'$.
\end{lemma}

\begin{proof}
Ceci résulte de la construction de $U'$ dans la démonstration du
lemme \ref{spaces-more-groupoids-lemma-strong-splitting},
car, dans ce cas, $U' = (R_s/U, e)_{fin}$ est un schéma séparé au-dessus
de $U$ d'après les
lemmes \ref{spaces-more-groupoids-lemma-finite-separated-flat-locally-finite-presentation} et
\ref{spaces-more-groupoids-lemma-finite-plus-section}.
\end{proof}

\begin{lemma}
\label{spaces-more-groupoids-lemma-splitting-scheme}
Dans la situation \ref{spaces-more-groupoids-situation-splitting}, supposons en outre que
$s, t$ soient plats et localement de présentation finie.
Alors il existe un schéma $U'$, un morphisme \'etale séparé
$U' \to U$ et un point $u' \in U'$
au-dessus de $u$, avec $\kappa(u) = \kappa(u')$,
tels que la restriction $R' = R|_{U'}$ de $R$ à $U'$
soit scindée au-dessus de $u'$.
\end{lemma}

\begin{proof}
Ceci résulte de la construction de $U'$ dans la démonstration du
lemme \ref{spaces-more-groupoids-lemma-splitting},
car, dans ce cas, $U' = (R_s/U, e)_{fin}$ est un schéma séparé au-dessus
de $U$ d'après les
lemmes \ref{spaces-more-groupoids-lemma-finite-separated-flat-locally-finite-presentation} et
\ref{spaces-more-groupoids-lemma-finite-plus-section}.
\end{proof}

\begin{lemma}
\label{spaces-more-groupoids-lemma-quasi-splitting-scheme}
Dans la situation \ref{spaces-more-groupoids-situation-quasi-splitting}, supposons en outre que
$s, t$ soient plats et localement de présentation finie.
Alors il existe un schéma $U'$, un morphisme \'etale séparé
$U' \to U$ et un point $u' \in U'$ au-dessus de $u$, avec
$\kappa(u) = \kappa(u')$, tels que la restriction $R' = R|_{U'}$ de
$R$ à $U'$ soit quasi-scindée au-dessus de $u'$.
\end{lemma}

\begin{proof}
Ceci résulte de la construction de $U'$ dans la démonstration du
lemme \ref{spaces-more-groupoids-lemma-quasi-splitting},
car, dans ce cas, $U' = (R_s/U, e)_{fin}$ est un schéma séparé
au-dessus de $U$ d'après les
lemmes \ref{spaces-more-groupoids-lemma-finite-separated-flat-locally-finite-presentation} et
\ref{spaces-more-groupoids-lemma-finite-plus-section}.
\end{proof}

\noindent
En fait, nous pouvons obtenir des schémas affines en appliquant un résultat
antérieur sur les groupoïdes finis localement libres.

\begin{lemma}
\label{spaces-more-groupoids-lemma-strong-splitting-affine-scheme}
Dans la situation \ref{spaces-more-groupoids-situation-strong-splitting}, supposons en outre que
$s, t$ soient plats et localement de présentation finie et que $U$ soit affine.
Alors il existe un schéma affine $U'$, un morphisme \'etale
$U' \to U$ et un point $u' \in U'$ au-dessus de $u$, avec
$\kappa(u) = \kappa(u')$, tels que la restriction $R' = R|_{U'}$ de
$R$ à $U'$ soit fortement scindée au-dessus de $u'$.
\end{lemma}

\begin{proof}
Soient $U' \to U$ le morphisme \'etale séparé de schémas et $u' \in U'$ le point
obtenus dans le lemme \ref{spaces-more-groupoids-lemma-strong-splitting-scheme}.
Soit $P \subset R'$ le scindage fort de $R'$ au-dessus de $u'$. D'après
Compléments sur les groupoïdes, lemme \ref{more-groupoids-lemma-restrict-preserves-type},
les morphismes $s'|_P, t'|_P : P \to U'$ sont plats et localement de présentation finie.
Ils sont finis par définition du scindage. Ainsi, $s'|_P, t'|_P$ sont finis
localement libres, voir
Morphismes, lemme \ref{morphisms-lemma-finite-flat}.
En particulier, $(t'|_P)((s'|_P)^{-1}(u'))$ est un ensemble fini de points
$\{u'_1, u'_2, \ldots, u'_n\}$ de $U'$. Choisissons un ouvert quasi-compact
$W \subset U'$ contenant chaque $u'_i$. Puisque $U$ est affine, le morphisme
$W \to U$ est quasi-compact (voir
Schémas, lemme \ref{schemes-lemma-quasi-compact-affine}).
Le morphisme $W \to U$ est aussi localement quasi-fini (voir
Morphismes, lemme \ref{morphisms-lemma-etale-locally-quasi-finite})
et séparé. Par conséquent, d'après
Compléments sur les morphismes,
lemme \ref{more-morphisms-lemma-quasi-finite-separated-quasi-affine}
(une version du théorème principal de Zariski),
nous concluons que $W$ est quasi-affine. D'après
Propriétés, lemme \ref{properties-lemma-ample-finite-set-in-affine},
l'ensemble $\{u'_1, \ldots, u'_n\}$ est contenu dans un ouvert affine
de $U'$. Nous pouvons donc appliquer
Groupoïdes, lemme \ref{groupoids-lemma-find-invariant-affine}
pour conclure qu'il existe un ouvert affine $U'' \subset U'$, invariant par
$P$, qui contient $u'$.

\medskip\noindent
Pour achever la démonstration, notons $R'' = R|_{U''}$ la restriction de $R$
à $U''$. C'est aussi la restriction de $R'$ à $U''$.
Puisque $P \subset R'$ est un sous-schéma ouvert et fermé, il en va de même de
$P|_{U''} \subset R''$. Par construction, le sous-schéma ouvert $U'' \subset U'$
est invariant par $P$, ce qui signifie que
$P|_{U''} = (s'|_P)^{-1}(U'') = (t'|_P)^{-1}(U'')$
(voir la discussion dans
Groupoïdes, section \ref{groupoids-section-invariant}) ;
les restrictions de $s''$ et de $t''$ à $P|_{U''}$ sont donc encore finies.
Le sous-schéma en groupoïdes $P|_{U''}$ est encore un scindage fort de
$R''$ au-dessus de $u'$ : les conditions (a) et (b) ont été vérifiées ci-dessus,
et (c) résulte immédiatement de $\{r' \in R': t'(r') = u', s'(r') = u'\} =
\{r'' \in R'': t''(r'') = u', s''(r'') = u'\}$.
Le lemme est démontré.
\end{proof}

\begin{lemma}
\label{spaces-more-groupoids-lemma-splitting-affine-scheme}
Dans la situation \ref{spaces-more-groupoids-situation-splitting}, supposons en outre que
$s, t$ soient plats et localement de présentation finie et que $U$ soit affine.
Alors il existe un schéma affine $U'$, un morphisme \'etale
$U' \to U$ et un point $u' \in U'$ au-dessus de $u$, avec
$\kappa(u) = \kappa(u')$, tels que la restriction $R' = R|_{U'}$ de
$R$ à $U'$ soit scindée au-dessus de $u'$.
\end{lemma}

\begin{proof}
La démonstration de ce lemme est littéralement la même que celle du
lemme \ref{spaces-more-groupoids-lemma-strong-splitting-affine-scheme},
si ce n'est que « scindage fort » doit être remplacé par « scindage »
(deux fois) et que le renvoi au
lemme \ref{spaces-more-groupoids-lemma-strong-splitting-scheme}
doit être remplacé par un renvoi au
lemme \ref{spaces-more-groupoids-lemma-splitting-scheme}.
\end{proof}

\begin{lemma}
\label{spaces-more-groupoids-lemma-quasi-splitting-affine-scheme}
Dans la situation \ref{spaces-more-groupoids-situation-quasi-splitting}, supposons en outre que
$s, t$ soient plats et localement de présentation finie et que $U$ soit affine.
Alors il existe un schéma affine $U'$, un morphisme \'etale
$U' \to U$ et un point $u' \in U'$ au-dessus de $u$, avec
$\kappa(u) = \kappa(u')$, tels que la restriction $R' = R|_{U'}$ de
$R$ à $U'$ soit quasi-scindée au-dessus de $u'$.
\end{lemma}

\begin{proof}
La démonstration de ce lemme est littéralement la même que celle du
lemme \ref{spaces-more-groupoids-lemma-strong-splitting-affine-scheme},
si ce n'est que « scindage fort » doit être remplacé par « quasi-scindage »
(deux fois) et que le renvoi au
lemme \ref{spaces-more-groupoids-lemma-strong-splitting-scheme}
doit être remplacé par un renvoi au
lemme \ref{spaces-more-groupoids-lemma-quasi-splitting-scheme}.
\end{proof}





\bibliography{my}
\bibliographystyle{amsalpha}
\end{document}
